% Options for packages loaded elsewhere
\PassOptionsToPackage{unicode}{hyperref}
\PassOptionsToPackage{hyphens}{url}
\documentclass[
]{book}
\usepackage{xcolor}
\usepackage{amsmath,amssymb}
\setcounter{secnumdepth}{5}
\usepackage{iftex}
\ifPDFTeX
  \usepackage[T1]{fontenc}
  \usepackage[utf8]{inputenc}
  \usepackage{textcomp} % provide euro and other symbols
\else % if luatex or xetex
  \usepackage{unicode-math} % this also loads fontspec
  \defaultfontfeatures{Scale=MatchLowercase}
  \defaultfontfeatures[\rmfamily]{Ligatures=TeX,Scale=1}
\fi
\usepackage{lmodern}
\ifPDFTeX\else
  % xetex/luatex font selection
\fi
% Use upquote if available, for straight quotes in verbatim environments
\IfFileExists{upquote.sty}{\usepackage{upquote}}{}
\IfFileExists{microtype.sty}{% use microtype if available
  \usepackage[]{microtype}
  \UseMicrotypeSet[protrusion]{basicmath} % disable protrusion for tt fonts
}{}
\makeatletter
\@ifundefined{KOMAClassName}{% if non-KOMA class
  \IfFileExists{parskip.sty}{%
    \usepackage{parskip}
  }{% else
    \setlength{\parindent}{0pt}
    \setlength{\parskip}{6pt plus 2pt minus 1pt}}
}{% if KOMA class
  \KOMAoptions{parskip=half}}
\makeatother
\usepackage{color}
\usepackage{fancyvrb}
\newcommand{\VerbBar}{|}
\newcommand{\VERB}{\Verb[commandchars=\\\{\}]}
\DefineVerbatimEnvironment{Highlighting}{Verbatim}{commandchars=\\\{\}}
% Add ',fontsize=\small' for more characters per line
\newenvironment{Shaded}{}{}
\newcommand{\AlertTok}[1]{\textcolor[rgb]{1.00,0.00,0.00}{\textbf{#1}}}
\newcommand{\AnnotationTok}[1]{\textcolor[rgb]{0.38,0.63,0.69}{\textbf{\textit{#1}}}}
\newcommand{\AttributeTok}[1]{\textcolor[rgb]{0.49,0.56,0.16}{#1}}
\newcommand{\BaseNTok}[1]{\textcolor[rgb]{0.25,0.63,0.44}{#1}}
\newcommand{\BuiltInTok}[1]{\textcolor[rgb]{0.00,0.50,0.00}{#1}}
\newcommand{\CharTok}[1]{\textcolor[rgb]{0.25,0.44,0.63}{#1}}
\newcommand{\CommentTok}[1]{\textcolor[rgb]{0.38,0.63,0.69}{\textit{#1}}}
\newcommand{\CommentVarTok}[1]{\textcolor[rgb]{0.38,0.63,0.69}{\textbf{\textit{#1}}}}
\newcommand{\ConstantTok}[1]{\textcolor[rgb]{0.53,0.00,0.00}{#1}}
\newcommand{\ControlFlowTok}[1]{\textcolor[rgb]{0.00,0.44,0.13}{\textbf{#1}}}
\newcommand{\DataTypeTok}[1]{\textcolor[rgb]{0.56,0.13,0.00}{#1}}
\newcommand{\DecValTok}[1]{\textcolor[rgb]{0.25,0.63,0.44}{#1}}
\newcommand{\DocumentationTok}[1]{\textcolor[rgb]{0.73,0.13,0.13}{\textit{#1}}}
\newcommand{\ErrorTok}[1]{\textcolor[rgb]{1.00,0.00,0.00}{\textbf{#1}}}
\newcommand{\ExtensionTok}[1]{#1}
\newcommand{\FloatTok}[1]{\textcolor[rgb]{0.25,0.63,0.44}{#1}}
\newcommand{\FunctionTok}[1]{\textcolor[rgb]{0.02,0.16,0.49}{#1}}
\newcommand{\ImportTok}[1]{\textcolor[rgb]{0.00,0.50,0.00}{\textbf{#1}}}
\newcommand{\InformationTok}[1]{\textcolor[rgb]{0.38,0.63,0.69}{\textbf{\textit{#1}}}}
\newcommand{\KeywordTok}[1]{\textcolor[rgb]{0.00,0.44,0.13}{\textbf{#1}}}
\newcommand{\NormalTok}[1]{#1}
\newcommand{\OperatorTok}[1]{\textcolor[rgb]{0.40,0.40,0.40}{#1}}
\newcommand{\OtherTok}[1]{\textcolor[rgb]{0.00,0.44,0.13}{#1}}
\newcommand{\PreprocessorTok}[1]{\textcolor[rgb]{0.74,0.48,0.00}{#1}}
\newcommand{\RegionMarkerTok}[1]{#1}
\newcommand{\SpecialCharTok}[1]{\textcolor[rgb]{0.25,0.44,0.63}{#1}}
\newcommand{\SpecialStringTok}[1]{\textcolor[rgb]{0.73,0.40,0.53}{#1}}
\newcommand{\StringTok}[1]{\textcolor[rgb]{0.25,0.44,0.63}{#1}}
\newcommand{\VariableTok}[1]{\textcolor[rgb]{0.10,0.09,0.49}{#1}}
\newcommand{\VerbatimStringTok}[1]{\textcolor[rgb]{0.25,0.44,0.63}{#1}}
\newcommand{\WarningTok}[1]{\textcolor[rgb]{0.38,0.63,0.69}{\textbf{\textit{#1}}}}
\usepackage{longtable,booktabs,array}
\newcounter{none} % for unnumbered tables
\usepackage{calc} % for calculating minipage widths
% Correct order of tables after \paragraph or \subparagraph
\usepackage{etoolbox}
\makeatletter
\patchcmd\longtable{\par}{\if@noskipsec\mbox{}\fi\par}{}{}
\makeatother
% Allow footnotes in longtable head/foot
\IfFileExists{footnotehyper.sty}{\usepackage{footnotehyper}}{\usepackage{footnote}}
\makesavenoteenv{longtable}
\usepackage{graphicx}
\makeatletter
\newsavebox\pandoc@box
\newcommand*\pandocbounded[1]{% scales image to fit in text height/width
  \sbox\pandoc@box{#1}%
  \Gscale@div\@tempa{\textheight}{\dimexpr\ht\pandoc@box+\dp\pandoc@box\relax}%
  \Gscale@div\@tempb{\linewidth}{\wd\pandoc@box}%
  \ifdim\@tempb\p@<\@tempa\p@\let\@tempa\@tempb\fi% select the smaller of both
  \ifdim\@tempa\p@<\p@\scalebox{\@tempa}{\usebox\pandoc@box}%
  \else\usebox{\pandoc@box}%
  \fi%
}
% Set default figure placement to htbp
\def\fps@figure{htbp}
\makeatother
% definitions for citeproc citations
\NewDocumentCommand\citeproctext{}{}
\NewDocumentCommand\citeproc{mm}{%
  \begingroup\def\citeproctext{#2}\cite{#1}\endgroup}
\makeatletter
 % allow citations to break across lines
 \let\@cite@ofmt\@firstofone
 % avoid brackets around text for \cite:
 \def\@biblabel#1{}
 \def\@cite#1#2{{#1\if@tempswa , #2\fi}}
\makeatother
\newlength{\cslhangindent}
\setlength{\cslhangindent}{1.5em}
\newlength{\csllabelwidth}
\setlength{\csllabelwidth}{3em}
\newenvironment{CSLReferences}[2] % #1 hanging-indent, #2 entry-spacing
 {\begin{list}{}{%
  \setlength{\itemindent}{0pt}
  \setlength{\leftmargin}{0pt}
  \setlength{\parsep}{0pt}
  % turn on hanging indent if param 1 is 1
  \ifodd #1
   \setlength{\leftmargin}{\cslhangindent}
   \setlength{\itemindent}{-1\cslhangindent}
  \fi
  % set entry spacing
  \setlength{\itemsep}{#2\baselineskip}}}
 {\end{list}}
\usepackage{calc}
\newcommand{\CSLBlock}[1]{\hfill\break\parbox[t]{\linewidth}{\strut\ignorespaces#1\strut}}
\newcommand{\CSLLeftMargin}[1]{\parbox[t]{\csllabelwidth}{\strut#1\strut}}
\newcommand{\CSLRightInline}[1]{\parbox[t]{\linewidth - \csllabelwidth}{\strut#1\strut}}
\newcommand{\CSLIndent}[1]{\hspace{\cslhangindent}#1}
\setlength{\emergencystretch}{3em} % prevent overfull lines
\providecommand{\tightlist}{%
  \setlength{\itemsep}{0pt}\setlength{\parskip}{0pt}}
\DeclareUnicodeCharacter{00B1}{\ensuremath{\pm}}
\DeclareUnicodeCharacter{0393}{\ensuremath{\Gamma}}
\DeclareUnicodeCharacter{03A9}{\ensuremath{\Omega}}
\DeclareUnicodeCharacter{03C3}{\ensuremath{\sigma}}
\DeclareUnicodeCharacter{03C4}{\ensuremath{\tau}}
\DeclareUnicodeCharacter{2192}{\ensuremath{\to}}
\DeclareUnicodeCharacter{2194}{\ensuremath{\leftrightarrow}}
\DeclareUnicodeCharacter{219D}{\ensuremath{\rightsquigarrow}}
\DeclareUnicodeCharacter{227A}{\ensuremath{\prec}}
\setlength{\emergencystretch}{3em}
\usepackage{bookmark}
\IfFileExists{xurl.sty}{\usepackage{xurl}}{} % add URL line breaks if available
\urlstyle{same}
\hypersetup{
  pdftitle={A Mathematical Atlas of Adaptive Intelligence},
  pdfauthor={Grand Challenge Labs},
  hidelinks,
  pdfcreator={LaTeX via pandoc}}

\ifXeTeX
\special{pdf:trailerid [ <a8420a94033497f38185a3cb7df471e5> <a8420a94033497f38185a3cb7df471e5> ]}
\fi
\ifPDFTeX
\pdftrailerid{}
\pdftrailer{/ID [ <a8420a94033497f38185a3cb7df471e5> <a8420a94033497f38185a3cb7df471e5> ]}
\fi
\ifLuaTeX
\pdfvariable trailerid {[ <a8420a94033497f38185a3cb7df471e5> <a8420a94033497f38185a3cb7df471e5> ]}
\fi

\title{A Mathematical Atlas of Adaptive Intelligence}
\author{Grand Challenge Labs}
\date{2026-10-07}

\begin{document}
\frontmatter
\maketitle

{
\setcounter{tocdepth}{2}
\tableofcontents
}
\mainmatter
\chapter{The Adaptive-System Thesis}\label{the-adaptive-system-thesis}

\textbf{Epistemic status:} Atlas Synthesis\\
\textbf{Specification:}
manuscript/specifications/ATLAS-CH-THESIS-001.md\\
\textbf{Documentary packet:}
mathematics/derivations/ATLAS-CH-THESIS-001-DERIVATIONS.md\\
\textbf{Source lock:} sources/source-locks/ATLAS-CH-THESIS-001.yaml

\section{What this book is trying to
explain}\label{what-this-book-is-trying-to-explain}

Machine learning is often introduced as a catalogue of model classes.

Linear models lead to multilayer networks. Convolutional networks lead
to attention. Attention leads to Transformers. Transformers lead to
mixture-of-experts systems, retrieval, tools, agents, and increasingly
elaborate training pipelines.

That chronology is useful, but it hides a structural fact.

The objects that matter are no longer contained inside one function
approximator.

A capable system may include:

\begin{itemize}
\tightlist
\item
  learned states;
\item
  transformations acting on those states;
\item
  persistent or external memory;
\item
  numerical solvers;
\item
  routing and scheduling;
\item
  tools and other models;
\item
  interaction protocols;
\item
  evidence stores;
\item
  human or machine review;
\item
  mechanisms governing how the system may change itself.
\end{itemize}

The Atlas begins from that observation.

Its governing thesis is:

\begin{quote}
Adaptive intelligence is more usefully studied here as organized
computation over states, operators, dynamics, memory, interfaces,
evidence, and governance than as a catalogue of model classes.
\end{quote}

This is the \textbf{Adaptive-System Thesis}.

It is the thesis of this monograph.

It is not introduced as an externally established theorem.

\section{An atlas, not a catalogue}\label{an-atlas-not-a-catalogue}

A catalogue tells us what objects exist.

An atlas tells us how a territory is charted.

It contains:

\begin{itemize}
\tightlist
\item
  coordinate systems;
\item
  local maps;
\item
  overlaps;
\item
  landmarks;
\item
  routes;
\item
  scales;
\item
  boundaries;
\item
  regions that remain poorly mapped.
\end{itemize}

This book is organized in the second way.

A chapter on geometry does not exist merely because manifolds are
mathematically important.

It exists because later chapters need a language for legal motion in
constrained state spaces.

A chapter on non-normality does not exist merely because pseudospectra
are elegant.

It exists because asymptotic eigenvalue stability can miss
finite-horizon amplification.

A chapter on replayable evidence does not exist merely because
provenance is administratively desirable.

It exists because an adaptive research system must know which claim is
supported by which reconstructable path.

The atlas allegory has a limit.

A research programme is not static geography.

The territory can change when new algorithms, new evidence, or new forms
of computation appear.

The map therefore has to be versioned.

\section{The recurring shifts}\label{the-recurring-shifts}

The source inventory that seeded this monograph records a recurring
sequence of conceptual shifts.

They are not presented as discoveries of this chapter.

They are the programme-level pattern the Atlas is being built to
examine.

\subsection{Vectors to operators}\label{vectors-to-operators}

A vector tells us where we are in a representation space.

An operator tells us how states are transformed.

Many modern mechanisms become clearer when the transformation is treated
as the central object rather than the stored vector alone.

Attention is one example.

Relative position may also be better understood as an operator acting on
representation space rather than as an appended coordinate.

\subsection{Layers to flows}\label{layers-to-flows}

A stack of layers is a discrete sequence.

The same stack can sometimes be interpreted as an approximation to an
evolution law.

That viewpoint introduces questions of:

\begin{itemize}
\tightlist
\item
  stability;
\item
  step size;
\item
  reversibility;
\item
  stiffness;
\item
  splitting;
\item
  error control.
\end{itemize}

Depth can then be studied partly as computational time.

This does not mean every network is literally a discretized differential
equation.

It means numerical dynamics can reveal structure that layer counting
alone hides.

\subsection{Parameters to geometry}\label{parameters-to-geometry}

Parameter values live in spaces with constraints, symmetries, scales,
and equivalences.

If a state or parameter is constrained to a sphere, Stiefel manifold,
simplex, quotient space, or other structured domain, the geometry
changes what constitutes a legal or meaningful update.

The Geometry keystone develops this point formally.

\subsection{Routing to decision
making}\label{routing-to-decision-making}

A router does more than compute a score.

It allocates scarce computation under uncertainty and capacity
constraints.

That opens the door to online decision theory, regret, optionality, and
coordination costs.

\subsection{Context to compiled memory
access}\label{context-to-compiled-memory-access}

A context window is one memory surface.

It is not the only one.

An adaptive system can combine:

\begin{itemize}
\tightlist
\item
  parametric memory;
\item
  external retrieval;
\item
  associative stores;
\item
  symbolic records;
\item
  shared institutional memory;
\item
  compiled context.
\end{itemize}

The Atlas therefore treats memory as infrastructure as well as model
anatomy.

\subsection{Agents to distributed
systems}\label{agents-to-distributed-systems}

Once multiple models, tools, memories, humans, and validators interact,
the relevant object is no longer a lone agent.

It is a coordinated system.

Questions of concurrency, communication, provenance, shared state,
failure isolation, and authority become mathematically and operationally
relevant.

\subsection{Training to dynamical
identification}\label{training-to-dynamical-identification}

An optimizer has state.

A model has state.

The data stream and scheduler can have state.

The combined training process can therefore be studied as a dynamical
system rather than only as repeated minimization of a scalar loss.

\subsection{Interpretability to
intervention}\label{interpretability-to-intervention}

Description is weaker than controlled substitution.

If a proposed internal object is said to matter, the stronger question
is whether changing, replacing, suppressing, or preserving it changes
behavior in the predicted way.

The Atlas therefore treats intervention and substitution as important
diagnostic standards.

\subsection{Architecture to composition of
contracts}\label{architecture-to-composition-of-contracts}

A system can fail even when every tensor shape matches.

Composition requires assumptions about meaning, geometry, perturbations,
and numerical behavior.

The Boundary Contracts keystone develops this as an interface problem.

\subsection{Result to evidence
object}\label{result-to-evidence-object}

A research conclusion is not only a sentence.

It is supported by sources, methods, environments, artifacts,
observations, interpretations, and review state.

The Replayable Evidence Objects keystone turns that support path into an
explicit object.

\section{A provisional system
decomposition}\label{a-provisional-system-decomposition}

For orientation, the Atlas uses the explanatory decomposition

\[
\mathcal A
=
(\mathcal S,\mathcal O,\mathcal D,\mathcal M,\mathcal I,\mathcal E,\mathcal G).
\]

Here:

\[
\mathcal S
=
\text{states},
\]

\[
\mathcal O
=
\text{operators},
\]

\[
\mathcal D
=
\text{dynamics},
\]

\[
\mathcal M
=
\text{memory},
\]

\[
\mathcal I
=
\text{interfaces},
\]

\[
\mathcal E
=
\text{evidence},
\]

and

\[
\mathcal G
=
\text{governance}.
\]

This tuple is not a universal ontology.

It is a map legend.

Later chapters will refine each term and, where necessary, expose the
places where the decomposition becomes inadequate.

\section{Why the model boundary is too
small}\label{why-the-model-boundary-is-too-small}

Consider a retrieval-augmented system.

Its answer can depend on:

\begin{itemize}
\tightlist
\item
  the model weights;
\item
  the tokenizer;
\item
  the query transformation;
\item
  the retrieval index;
\item
  the current contents of the index;
\item
  ranking;
\item
  context compilation;
\item
  tool outputs;
\item
  runtime policies.
\end{itemize}

If we ask only what is stored in the model parameters, we have chosen a
system boundary that excludes part of the causal machinery.

The same issue appears in agent systems.

A planner that calls a solver and writes to shared memory cannot be
understood entirely by inspecting the planner weights.

Capability is partly distributed over the arrangement.

This motivates a broader unit of analysis:

\[
\boxed{
\text{model}
\subseteq
\text{adaptive computational system}.
}
\]

The inclusion can be strict.

\section{A useful thesis must be
vulnerable}\label{a-useful-thesis-must-be-vulnerable}

A thesis that explains everything after the fact explains very little.

The Adaptive-System Thesis therefore has pressure points.

It should become less attractive if:

\begin{itemize}
\tightlist
\item
  the proposed object distinctions repeatedly obscure rather than
  clarify important mechanisms;
\item
  memory, interfaces, or governance turn out to be accidental
  implementation details rather than load-bearing computational
  structure;
\item
  the operator/dynamics viewpoint adds no predictive or design value;
\item
  composition can be handled adequately by shape-level software
  interfaces alone;
\item
  evidence provenance proves irrelevant to the reliability of adaptive
  research systems;
\item
  later chapters require entirely different primitives that cannot be
  expressed naturally through the Atlas roles.
\end{itemize}

The book is not arranged to prevent these outcomes.

It is arranged to make them visible.

\section{Description versus
programme}\label{description-versus-programme}

Some claims in the Atlas are descriptive.

For example:

\begin{itemize}
\tightlist
\item
  attention uses state-dependent mixing operators;
\item
  momentum introduces optimizer state;
\item
  a source hash identifies bytes;
\item
  a Jacobian describes local differential sensitivity.
\end{itemize}

Other claims are programme-level.

For example:

\begin{itemize}
\tightlist
\item
  intelligent systems should externalize more world knowledge into
  governed shared memory;
\item
  adaptive depth should be treated as numerical error control;
\item
  boundary contracts should become first-class architectural objects;
\item
  research systems should preserve replayable evidence objects.
\end{itemize}

The second class is not smuggled into the first.

A research programme can be useful before it becomes established theory.

It must still be labeled correctly.

\section{Dependencies in the argument}\label{the-dependency-graph-is-part-of-the-argument}

The dependency graph is more than a way to navigate chapters: it records which
mathematical or documentary objects a later argument is entitled to assume.
An editorial reading order need not coincide with that dependency order.
When a chapter borrows a definition, theorem, example, or evidence object,
the relevant prerequisite must remain inspectable. That principle is part of
the Atlas's method, not evidence that its broader thesis is true.

\section{Tests the thesis must pass}\label{what-later-chapters-must-earn}

The decomposition is useful only where it changes the questions we can ask
and the errors we can detect. Later chapters must put four propositions at risk:
(1) geometry and operators expose behavior hidden by coordinate descriptions;
(2) dynamics and optimization explain finite-time effects missed by static
summaries; (3) explicit interfaces, memory and coordination expose system-level
failure modes; and (4) evidence records delimit which conclusions can be
reconstructed and reused.

These are research obligations, not results proved by announcing a taxonomy.
The counterexamples and limits in the chapters ahead are part of the test.

\section{The map ahead}\label{the-map-ahead}

The immediate next chapter introduces four recurring roles:

\begin{itemize}
\tightlist
\item
  states;
\item
  operators;
\item
  flows;
\item
  interfaces.
\end{itemize}

After that, the mathematical substrate develops:

\begin{itemize}
\tightlist
\item
  linear maps;
\item
  probability and information;
\item
  geometry;
\item
  dynamics;
\item
  numerical analysis;
\item
  local-to-global reasoning.
\end{itemize}

Only then does the Atlas move deeply into representation, architecture,
optimization, memory, routing, agents, and governed adaptation.

The recurring conceptual spine is:

\[
\boxed{
\text{geometry}
\to
\text{operators}
\to
\text{dynamics}
\to
\text{optimization}
\to
\text{composition}
\to
\text{memory}
\to
\text{coordination}
\to
\text{diagnostics}
\to
\text{governed adaptation}.
}
\]

This is a route through the atlas.

It is not the only possible route.

\section{Closing view}\label{closing-view}

The central wager of this book is that the next useful abstraction
boundary for machine intelligence is larger than the model.

It includes the transformations the model performs, the geometry in
which those transformations occur, the dynamics through which they
accumulate, the memory they consult, the interfaces through which
components compose, the evidence by which claims are supported, and the
governance by which adaptation is constrained.

Whether that wager survives is the work of the chapters ahead.

The Atlas begins by naming the territory.

It will spend the rest of the book trying to deserve the map.

\section{References used in this
chapter}\label{references-used-in-this-chapter}

This chapter is Atlas synthesis. Its load-bearing documentary sources
are the source inventory, Atlas Map, Editorial Profile, and Chapter
Composition Protocol bound in
sources/source-locks/ATLAS-CH-THESIS-001.yaml. External mathematical
claims are deferred to the downstream chapters that source them
directly.

\chapter{How to Read a Mathematical
Atlas}\label{how-to-read-a-mathematical-atlas}

\textbf{Epistemic status:} Atlas Documentary Synthesis\\
\textbf{Specification:} manuscript/specifications/ATLAS-CH-MAP-001.md\\
\textbf{Documentary packet:}
mathematics/derivations/ATLAS-CH-MAP-001-DOCUMENTARY.md\\
\textbf{Source lock:} sources/source-locks/ATLAS-CH-MAP-001.yaml

\section{An atlas is not a road}\label{an-atlas-is-not-a-road}

A road tells you where to go next.

An atlas gives you enough structure to choose.

That distinction matters for this book.

\emph{A Mathematical Atlas of Adaptive Intelligence} has an editorial
order, but it is not designed to have only one valid reading order.

Different readers arrive with different objectives:

\begin{itemize}
\tightlist
\item
  learn the foundations;
\item
  understand one technical result;
\item
  reconstruct a computational witness;
\item
  evaluate a research frontier;
\item
  follow one concept across several domains;
\item
  enter through a figure and descend into the mathematics.
\end{itemize}

The Atlas must support all of these without allowing convenience to
erase prerequisites.

The governing rule is simple:

\begin{quote}
reading order is flexible; dependency and claim boundaries are not.
\end{quote}

\section*{Start here: choose a reading route}

Readers seeking a connected introduction can follow the printed order.
For a technical target, follow its hard prerequisites backward before
reading optional cross-links. For a conceptual question, trace the relevant
theme across parts. For an open research question, inspect the chapter's
claim boundary, exact sources, witnesses and failure modes. For formal or
computational work, begin with the derivation and replay record. For visual
learning, start with a figure and its declared literal and nonliteral meanings.
The fuller route descriptions and a chapter-entry checklist appear below.

\section{Map and territory}\label{map-and-territory}

The word \textbf{atlas} is more than decoration.

A geographic atlas contains:

\begin{itemize}
\tightlist
\item
  maps at several scales;
\item
  legends;
\item
  coordinates;
\item
  routes;
\item
  landmarks;
\item
  survey records;
\item
  boundaries.
\end{itemize}

The mathematical correspondence is:

\begin{itemize}
\tightlist
\item
  territory ↔ the mathematical and scientific subject;
\item
  map ↔ an organized representation of that subject;
\item
  scale ↔ level of resolution;
\item
  legend ↔ notation and epistemic labels;
\item
  coordinates ↔ stable chapter and object IDs;
\item
  route ↔ a reader-selected path;
\item
  landmark ↔ a chapter of unusual structural importance;
\item
  survey record ↔ source locks, witnesses, and audits.
\end{itemize}

The limit is equally important.

The territory is not exhausted by the map.

A route is not a unique curriculum.

A landmark is not evidence for every claim around it.

A legend tells us how to read a mark; it does not make the mark true.

\section{Read at several
resolutions}\label{read-at-several-resolutions}

The Atlas is designed for movement among scales.

At the largest scale:

\[
\text{Atlas}
\]

shows the whole conceptual landscape.

Then come:

\[
\text{Part}
\to
\text{Chapter}
\to
\text{claim}
\to
\text{support object}.
\]

A reader may begin broadly and zoom inward.

Or begin with one exact derivation and zoom outward.

Changing scale changes the amount of visible detail.

It must not change the epistemic status of the underlying claim.

A conjecture does not become established because it is drawn on a clean
overview plate.

A theorem does not become weaker because we first encounter it through
intuition.

\section{Stable identity}\label{stable-identity}

The Atlas is expected to evolve.

Chapters may move.

Parts may be reorganized.

A chapter may gain or lose neighbors.

For that reason, identity is not carried primarily by final numbering.

It is carried by stable IDs such as

\[
\texttt{ATLAS-CH-DYN-001}.
\]

The Atlas Map explicitly treats numbering and physical order as
revisable while stable IDs survive reordering.

This gives the book something like a coordinate grid.

A citation to a stable object should remain meaningful even when the
local page neighborhood changes.

\section{The map is not the dependency
graph}\label{the-map-is-not-the-dependency-graph}

Two structures coexist.

The \textbf{Atlas Map} answers:

\begin{quote}
What exists, and what role does it play in the book?
\end{quote}

The \textbf{Dependency Graph} answers:

\begin{quote}
What may this chapter assume without rebuilding it?
\end{quote}

They overlap.

They are not identical.

A chapter can appear nearby in the narrative without being a hard
prerequisite.

A concept can illuminate another chapter without authorizing it to
assume the concept has already been established.

\section{Hard dependency}\label{hard-dependency}

Suppose chapter (A) is a hard dependency of chapter (B).

Then (B) may consume the declared content of (A) without reconstructing
it from first principles.

Write this schematically as

\[
A \longrightarrow B.
\]

The arrow is permission.

It says:

\begin{quote}
the downstream chapter may stand on this upstream object.
\end{quote}

That permission creates an authorial obligation.

The upstream chapter must actually contain what the downstream chapter
is entitled to assume.

This is why dependency audits matter.

\section{Soft cross-link}\label{soft-cross-link}

A soft cross-link means something weaker.

Two chapters may illuminate one another:

\[
A \leftrightarrow B.
\]

But neither is automatically permitted to assume the other has been
read.

A soft link is a conceptual bridge.

It is not prerequisite authority.

This distinction lets the Atlas remain richly cross-connected without
turning every conceptual resonance into a cycle of mandatory reading.

\section{Narrative is not mathematical
authority}\label{narrative-is-not-mathematical-authority}

The dependency graph states this explicitly:

\begin{quote}
the graph is a dependency structure, not the book's narrative.
\end{quote}

The converse is also true.

Narrative order is not dependency authority.

This gives the author freedom to place an accessible orientation chapter
early while still exposing the technical mathematics that later chapters
truly require.

It gives the reader freedom to choose a route while still knowing when a
detour is mathematically necessary.

\section{A compact epistemic legend}\label{the-legend-epistemic-labels}

The Atlas distinguishes definitions, established external results, Atlas
derivations, bounded computational witnesses, observations, interpretations,
public project evidence, programme context, conjectures, open problems, and
institutional statuses. These are distinct roles, not steps on a single scale.
A green audit, a reproducible computation, or polished typography does not
by itself prove a theorem or upgrade the scope of a claim.

An exact figure represents declared mathematics; a data-derived figure
summarizes observations within a stated scope; a schematic figure makes a
structural relationship legible without pretending to be measured geometry.
The figure manifest determines what is literal and what is not.

A source lock identifies the consumed source and what that source can actually
support. Exact identity is not proof of adequate authority. A replay confirms
reconstruction under stated conditions, not truth, independent replication
or certification. For definitions, examples, promotion failures and the
claim-support discipline, continue to the next chapter, \emph{Claims,
Evidence, and Computational Witnesses}.

% Historical subsection anchors retained for links from the released text.
\phantomsection\label{definition}
\phantomsection\label{established-result}
\phantomsection\label{atlas-derivation}
\phantomsection\label{computational-witness}
\phantomsection\label{observation-and-interpretation}
\phantomsection\label{gcl-project-evidence-and-programme-context}
\phantomsection\label{conjecture-and-open-problem}
\phantomsection\label{institutional-status}
\phantomsection\label{presentation-creates-no-epistemic-status}
\phantomsection\label{proof-derivation-and-witness}
\phantomsection\label{replay-is-not-truth}
\phantomsection\label{figures-have-representation-classes}
\phantomsection\label{exact-figures}
\phantomsection\label{data-derived-figures}
\phantomsection\label{schematic-figures}
\phantomsection\label{the-source-lock}
\phantomsection\label{documentary-provenance}
\phantomsection\label{audits}

\section{Six ways through the
Atlas}\label{six-ways-through-the-atlas}

There is no single privileged route.

The best route depends on what you are trying to do.

\subsection{Route A --- linear
foundations}\label{route-a-linear-foundations}

Follow the editorial sequence.

Best for:

\begin{itemize}
\tightlist
\item
  broad learning;
\item
  conceptual continuity;
\item
  readers who want the argument to unfold gradually.
\end{itemize}

Risk:

you may read more prerequisite material than a targeted question
requires.

\subsection{Route B ---
dependency-first}\label{route-b-dependency-first}

Choose a target chapter.

Follow hard dependencies backward until the path closes.

Then read forward.

Best for:

\begin{itemize}
\tightlist
\item
  technical work;
\item
  minimal prerequisite reconstruction;
\item
  verifying that a later result is actually supported.
\end{itemize}

Risk:

the route can be mathematically efficient but narratively austere.

\subsection{Route C --- concept/theme}\label{route-c-concepttheme}

Choose a concept such as:

\begin{itemize}
\tightlist
\item
  geometry;
\item
  operators;
\item
  memory;
\item
  optimization;
\item
  coordination;
\item
  governance.
\end{itemize}

Follow soft cross-links and recurring objects across Parts.

Best for:

\begin{itemize}
\tightlist
\item
  synthesis;
\item
  comparative reading;
\item
  seeing one idea change form.
\end{itemize}

Risk:

soft links do not replace hard prerequisites.

\subsection{Route D ---
research-frontier}\label{route-d-research-frontier}

Start at an open problem, conjecture, or GCL programme chapter.

Then descend.

Read:

\begin{enumerate}
\def\labelenumi{\arabic{enumi}.}
\tightlist
\item
  the claim boundary;
\item
  the source lock;
\item
  the immediate mathematical prerequisites;
\item
  the computational witnesses;
\item
  the failure modes;
\item
  the relevant audit records.
\end{enumerate}

Best for:

\begin{itemize}
\tightlist
\item
  evaluating a live research idea;
\item
  separating established substrate from programme ambition.
\end{itemize}

Risk:

frontier prose can feel more mature than the evidence if the epistemic
labels are ignored.

\subsection{Route E --- proof/replay}\label{route-e-proofreplay}

Begin with the result you want to trust.

Trace:

\[
\text{claim}
\to
\text{derivation/proof}
\to
\text{witness}
\to
\text{source lock}
\to
\text{audit/replay record}.
\]

Best for:

\begin{itemize}
\tightlist
\item
  reconstruction;
\item
  adjudication;
\item
  technical reuse.
\end{itemize}

Risk:

replayability can be mistaken for correctness unless proof and authority
are read separately.

\subsection{Route F ---
visual/Atlas-plate}\label{route-f-visualatlas-plate}

Begin with a figure or conceptual plate.

Read its:

\begin{enumerate}
\def\labelenumi{\arabic{enumi}.}
\tightlist
\item
  representation class;
\item
  literal semantics;
\item
  nonliteral semantics;
\item
  supporting derivation or witness;
\item
  claim boundary.
\end{enumerate}

Best for:

\begin{itemize}
\tightlist
\item
  geometric intuition;
\item
  operator intuition;
\item
  rapid orientation.
\end{itemize}

Risk:

visual fluency can create false certainty if schematic elements are read
literally.

\section{Choosing a route}\label{choosing-a-route}

A practical routing table is:

{\def\LTcaptype{none} % do not increment counter
\begin{longtable}[]{@{}
  >{\raggedright\arraybackslash}p{(\linewidth - 4\tabcolsep) * \real{0.3333}}
  >{\raggedright\arraybackslash}p{(\linewidth - 4\tabcolsep) * \real{0.3333}}
  >{\raggedright\arraybackslash}p{(\linewidth - 4\tabcolsep) * \real{0.3333}}@{}}
\toprule\noalign{}
\begin{minipage}[b]{\linewidth}\raggedright
Goal
\end{minipage} & \begin{minipage}[b]{\linewidth}\raggedright
Start with
\end{minipage} & \begin{minipage}[b]{\linewidth}\raggedright
Then follow
\end{minipage} \\
\midrule\noalign{}
\endhead
\bottomrule\noalign{}
\endlastfoot
Learn the field systematically & editorial order & hard dependencies as
they appear \\
Understand one technical chapter & target chapter & hard dependencies
backward \\
Compare one concept across fields & theme & soft cross-links plus local
prerequisites \\
Evaluate a frontier proposal & open-problem/programme chapter & source
lock, prerequisites, witnesses, failures \\
Reproduce a result & claim/support object & derivation, witness,
provenance, audit \\
Build intuition first & figure/plate & manifest semantics, mathematics,
source boundary \\
\end{longtable}
}

The table chooses a route.

It does not choose what is true.

\section{How to read a keystone}\label{how-to-read-a-keystone}

A keystone chapter has unusually large downstream reach.

That makes it important.

It does not make it infallible.

When reading a keystone:

\begin{enumerate}
\def\labelenumi{\arabic{enumi}.}
\tightlist
\item
  identify the formal object it stabilizes;
\item
  inspect its exact derivations;
\item
  read its failure boundaries;
\item
  note which downstream chapters consume it;
\item
  inspect its audit disposition.
\end{enumerate}

A defect in a high-leverage chapter can propagate widely.

That is why keystones deserve disproportionate scrutiny.

\section{How to read a frontier
chapter}\label{how-to-read-a-frontier-chapter}

A frontier chapter should be read in the opposite direction from a
theorem.

Do not begin by asking:

\begin{quote}
Is this exciting?
\end{quote}

Begin by asking:

\begin{quote}
Which parts are already established?
\end{quote}

\begin{quote}
Which parts are Atlas derivation?
\end{quote}

\begin{quote}
Which parts are computational evidence?
\end{quote}

\begin{quote}
Which parts are interpretation?
\end{quote}

\begin{quote}
Which sentence is the actual conjecture or open problem?
\end{quote}

The Atlas is intentionally allowed to be imaginative.

Its epistemic labels keep imagination from borrowing the authority of
proof.

\section{Allegory as a bounded
instrument}\label{allegory-as-a-bounded-instrument}

This book uses allegory deliberately.

Examples include:

\begin{itemize}
\tightlist
\item
  maps and territories;
\item
  flows and film frames;
\item
  interfaces and engineering contracts;
\item
  memory and civic archives;
\item
  distributed intelligence and scientific institutions.
\end{itemize}

The rule is always:

\[
\text{allegory}
\to
\text{structural correspondence}
\to
\text{formal object}
\to
\text{limit}.
\]

An allegory is successful when it makes the formal structure easier to
see.

It has failed when the reader can no longer tell which part is metaphor.

\section{What changed in our
picture?}\label{what-changed-in-our-picture}

The Atlas thesis says adaptive intelligence is better understood as a
system of geometry, operators, dynamics, memory, composition,
coordination, diagnostics, and governed adaptation.

This chapter adds a second-order lesson.

The book itself has to be adaptive without becoming unstable.

It needs:

\begin{itemize}
\tightlist
\item
  stable identity;
\item
  revisable architecture;
\item
  explicit dependencies;
\item
  bounded evidence;
\item
  recoverable provenance.
\end{itemize}

That is why the Atlas is versioned like a technical system rather than
treated as a frozen sequence of pages.

\section{Five reading errors to
avoid}\label{five-reading-errors-to-avoid}

\subsection{Chapter order treated as dependency
order}\label{chapter-order-treated-as-dependency-order}

Narrative placement does not authorize hidden assumptions.

\subsection{Soft cross-link treated as
prerequisite}\label{soft-cross-link-treated-as-prerequisite}

Conceptual resonance is not a dependency edge.

\subsection{Exact computation treated as proof
automatically}\label{exact-computation-treated-as-proof-automatically}

A witness establishes only its bounded claim.

\subsection{Replay treated as
certification}\label{replay-treated-as-certification}

Reconstructability and authority are different dimensions.

\subsection{Schematic figure treated as measured
geometry}\label{schematic-figure-treated-as-measured-geometry}

Layout is not data unless the manifest says it is.

\section{A compact reader protocol}\label{a-compact-reader-protocol}

When entering any unfamiliar Atlas chapter, ask:

\begin{enumerate}
\def\labelenumi{\arabic{enumi}.}
\tightlist
\item
  What problem is this chapter solving?
\item
  What are its hard prerequisites?
\item
  What formal or documentary object does it introduce?
\item
  What is established externally?
\item
  What is derived here?
\item
  What is only witnessed computationally?
\item
  What is interpretation or programme context?
\item
  What would falsify or limit the claim?
\item
  Which figures are literal, data-derived, or schematic?
\item
  What downstream chapter is allowed to assume this work?
\end{enumerate}

If those ten answers are visible, the chapter is doing its job.

\section{Closing view}\label{closing-view-1}

A mathematical atlas should make exploration possible without making
rigor optional.

The reader is free to choose a route.

The author is not free to hide a dependency.

The reader is free to begin from a figure.

The figure is not free to conceal which marks are schematic.

The book is free to propose conjectures and programmes.

They are not free to masquerade as established results.

That is the compact bargain of the Atlas:

\begin{quote}
many routes, stable coordinates, explicit legends, reconstructable
evidence.
\end{quote}

The territory will remain larger than the book.

The task of the Atlas is to help us move through it without confusing
navigation for proof.

\section{References used in this
chapter}\label{references-used-in-this-chapter-1}

This chapter is a documentary synthesis of project-internal canonical
objects.

Primary locked sources:

\begin{itemize}
\tightlist
\item
  \texttt{governance/ATLAS\_MAP.md};
\item
  \texttt{governance/DEPENDENCY\_GRAPH.md};
\item
  \texttt{governance/CHAPTER\_COMPOSITION\_PROTOCOL.md};
\item
  \texttt{governance/EPISTEMIC\_STATUS.yaml};
\item
  \texttt{governance/ATLAS\_EDITORIAL\_PROFILE.md};
\item
  \texttt{governance/FIGURE\_REGISTER.yaml};
\item
  \texttt{governance/SOURCE\_REGISTER.yaml}.
\end{itemize}

See \texttt{sources/source-locks/ATLAS-CH-MAP-001.yaml} for exact Git
blob identities and declared authority.

\chapter{States, Operators, Flows, and
Interfaces}\label{states-operators-flows-and-interfaces}

\textbf{Epistemic status:} Atlas Synthesis\\
\textbf{Specification:}
manuscript/specifications/ATLAS-CH-OBJECTS-001.md\\
\textbf{Derivation packet:}
mathematics/derivations/ATLAS-CH-OBJECTS-001-DERIVATIONS.md\\
\textbf{Source lock:} sources/source-locks/ATLAS-CH-OBJECTS-001.yaml

\section{The same tensor can be four different
things}\label{the-same-tensor-can-be-four-different-things}

Suppose a program contains a square array

\[
A\in\mathbb R^{n\times n}.
\]

What is it?

It might be:

\begin{itemize}
\tightlist
\item
  a covariance state;
\item
  a learned parameter;
\item
  a linear operator;
\item
  an adjacency matrix;
\item
  a metric tensor in coordinates;
\item
  an attention score matrix;
\item
  an evidence artifact.
\end{itemize}

The shape does not answer the question.

The Atlas therefore begins with a discipline:

\begin{quote}
identify an object by the role it plays in the argument, not only by the
container used to store it.
\end{quote}

Four roles recur throughout this book:

\[
\boxed{
\text{state},
\quad
\text{operator},
\quad
\text{flow},
\quad
\text{interface}.
}
\]

These are explanatory categories.

They are not a universal ontology.

\section{Noun, verb, motion,
handshake}\label{noun-verb-motion-handshake}

A useful first analogy is grammatical.

A \textbf{state} behaves like a noun: something has a value.

An \textbf{operator} behaves like a verb: something is done to a state.

A \textbf{flow} behaves like unfolding motion: transformations
accumulate through an evolution parameter.

An \textbf{interface} behaves like a handshake: two objects agree on
what may cross a boundary.

The analogy is useful because it catches category errors.

It also fails quickly.

A mathematical state can itself encode an operator.

An operator can be state-dependent.

A family of operators may fail to satisfy the exact algebraic law
required of a true flow.

An interface can be learned or adaptive.

Natural-language grammar is therefore only an entry point.

The mathematics determines the object.

\section{State}\label{state}

A state is an information-bearing value

\[
x\in\mathcal X.
\]

The state space \(\mathcal X\) may be:

\begin{itemize}
\tightlist
\item
  Euclidean;
\item
  discrete;
\item
  probabilistic;
\item
  constrained;
\item
  manifold-valued;
\item
  a product space;
\item
  a space of functions;
\item
  a structured record.
\end{itemize}

A state answers:

\begin{quote}
What information is currently instantiated?
\end{quote}

Examples include:

\begin{itemize}
\tightlist
\item
  a hidden activation;
\item
  an optimizer momentum buffer;
\item
  a memory record;
\item
  a probability distribution;
\item
  a graph;
\item
  a model checkpoint;
\item
  a tuple-space configuration.
\end{itemize}

The word ``state'' is deliberately broader than ``vector.''

\section{Operator}\label{operator}

An operator is a transformation

\[
F:\mathcal X\to\mathcal Y.
\]

It answers:

\begin{quote}
What transformation is being applied?
\end{quote}

A linear operator has the special form

\[
F(x)=Ax.
\]

But many operators in the Atlas are nonlinear or state-dependent.

Examples include:

\begin{itemize}
\tightlist
\item
  attention;
\item
  normalization;
\item
  a retrieval rule;
\item
  an optimizer update;
\item
  a solver step;
\item
  a routing map;
\item
  a projection or retraction.
\end{itemize}

A matrix is not automatically an operator.

It becomes an operator when the model assigns it a transformation role.

\subsection{First category error: array equals
role}\label{first-category-error-array-equals-role}

Let

\[
A=
\begin{pmatrix}
1&0\\
0&4
\end{pmatrix}.
\]

If \(A\) is a covariance matrix, it is a state describing second-order
structure.

If \(A\) acts by

\[
x\mapsto Ax,
\]

it is a linear operator.

The stored entries are identical.

The semantics are different.

This distinction becomes important later when the Atlas moves from
vectors to operators.

Without role discipline, the phrase ``the representation matrix'' can
hide whether we mean data, a transformation, or a coordinate encoding of
something basis-independent.

\section{Flow and dynamics}\label{flow-and-dynamics}

A global autonomous flow is a family of maps
\(\Phi_t:\mathcal X\to\mathcal X\), indexed by \(t\in\mathbb R\), such that
\(\Phi_0=\mathrm{id}_{\mathcal X}\) and

\[
\Phi_{t+s}=\Phi_t\circ\Phi_s
\qquad (t,s\in\mathbb R).
\]

If only \(t,s\geq 0\) are admitted, with the same identity and
composition law, the corresponding object is a \emph{semiflow}.
A \emph{local flow} instead permits \(\Phi_t(x)\) only on a
state-dependent time domain; the composition law is required only when
both sides are defined. Existence for all real times must not be inferred
from local existence.

If a differential equation

\[
\dot x=f(x)
\]

has a well-defined evolution, \(\Phi_t(x_0)\) denotes the state after
time \(t\).

This gives a new kind of object.

A single operator says:

\begin{quote}
transform this state once.
\end{quote}

A flow says:

\begin{quote}
here is how transformations accumulate across an evolution parameter.
\end{quote}

This distinction becomes central when network depth is treated as
computational time.

\subsection{Second category error: one step equals
dynamics}\label{second-category-error-one-step-equals-dynamics}

Consider a residual update

\[
x_{k+1}
=
x_k+h f(x_k).
\]

One update map is

\[
F_h(x)=x+h f(x).
\]

The depth-indexed computation is the sequence

\[
x_0,
x_1,
x_2,
\ldots
\]

generated by repeated application.

These are not the same object.

The single-step operator can be locally well behaved while repeated
application has different long-horizon behavior.

Later chapters on numerical intelligence and non-normality make this
distinction quantitative.

\subsection{Flow-like is weaker than
flow}\label{flow-like-is-weaker-than-flow}

Neural networks often invite continuous-time language.

That language can be useful.

It must be scoped.

If layer \(k\) uses a different transformation from layer \(k+1\), the
depth-indexed family need not satisfy

\[
\Phi_{t+s}
=
\Phi_t\circ\Phi_s
\]

for one autonomous flow.

The safer statement may be:

\begin{quote}
the architecture admits a dynamical or integrator interpretation.
\end{quote}

That is weaker than claiming an exact continuous-time flow.

The Atlas will use the stronger word only when the stronger structure is
present.

\section{Interface}\label{interface}

Suppose

\[
F:\mathcal X\to\mathcal Y
\]

feeds

\[
G:\mathcal Y\to\mathcal Z.
\]

At minimum, the codomain of \(F\) must be admissible input for \(G\).

Software type systems often check part of this requirement.

But a meaningful interface can also depend on:

\begin{itemize}
\tightlist
\item
  units;
\item
  scale;
\item
  normalization;
\item
  geometric constraints;
\item
  semantic interpretation;
\item
  uncertainty;
\item
  precision;
\item
  differential sensitivity.
\end{itemize}

An interface therefore answers:

\begin{quote}
What is another component allowed to assume about what crosses this
boundary?
\end{quote}

The Boundary Contracts chapter later turns this question into a formal
programme.

\subsection{Third category error: matching shape equals
composability}\label{third-category-error-matching-shape-equals-composability}

Let

\[
F(x)
=
\frac{x}{\|x\|_2}
\]

return direction only.

Let

\[
G(y)
=
\|y\|_2
\]

interpret magnitude as meaningful amplitude.

The shapes match.

The semantics do not.

For

\[
x=(3,4),
\]

the original magnitude is \(5\), but

\[
G(F(x))=1.
\]

The interface failed before any matrix multiplication failed.

\subsection{Interfaces are not
interiors}\label{interfaces-are-not-interiors}

A contract can describe:

\begin{itemize}
\tightlist
\item
  accepted inputs;
\item
  guarantees;
\item
  error bounds;
\item
  invariants;
\item
  provenance.
\end{itemize}

It does not need to expose the full internal mechanism.

This separation is important for modular intelligence.

If every component must know every other component's interior,
composition does not scale.

But abstraction has a price.

A contract can omit something the downstream system actually needs.

The problem is therefore not:

\begin{quote}
hide as much as possible.
\end{quote}

It is:

\begin{quote}
expose enough for the property the composition is supposed to preserve.
\end{quote}

\section{State-dependent operators}\label{state-dependent-operators}

The four roles are not rigid boxes.

Attention provides a useful example.

For fixed queries and keys, the normalized attention matrix \(A\) acts
linearly on values:

\[
V\mapsto AV.
\]

But in the full layer,

\[
A=A(X)
\]

depends on the current state \(X\).

The operator itself is state-dependent.

This is not a contradiction.

It is exactly why role language is useful.

We can distinguish:

\begin{itemize}
\tightlist
\item
  the current state;
\item
  the operator instantiated from that state;
\item
  the transformation applied to another part of the state.
\end{itemize}

\subsection{Operator-valued state}\label{operator-valued-state}

Sometimes a state contains an object that later acts as an operator.

An optimizer state may contain a preconditioner.

A memory system may store a learned routing matrix.

A hypernetwork may emit weights that define another map.

The same mathematical object can therefore change role across a
pipeline.

The rule is not:

\begin{quote}
every object has one eternal category.
\end{quote}

The rule is:

\begin{quote}
state the role being used in the present argument.
\end{quote}

\subsection{Evidence is another object
class}\label{evidence-is-another-object-class}

Later in the Atlas, evidence records become explicit.

An evidence object may describe an operator:

\begin{itemize}
\tightlist
\item
  which source defined it;
\item
  which code implemented it;
\item
  which run exercised it;
\item
  which output was observed.
\end{itemize}

The evidence record is not the operator.

Confusing those categories produces a different kind of error.

A green replay tells us that an evidence path reconstructed under its
declared conditions.

It does not transform the evidence record into the mathematical object
it describes.

\section{Why this vocabulary
matters}\label{why-this-vocabulary-matters}

The rest of the Atlas repeatedly changes levels of description.

Linear algebra moves between coordinates and operators.

Geometry distinguishes ambient coordinates from admissible states.

Attention distinguishes scores from mixing operators.

Optimization adds optimizer state to model state.

Numerical analysis distinguishes one update from accumulated dynamics.

Memory distinguishes stored content from retrieval operators.

Composition distinguishes component interiors from boundary obligations.

Scientific method distinguishes claims from their evidence paths.

Without a small amount of object discipline, these chapters would reuse
the same words for incompatible roles.

\subsection{The four questions}\label{the-four-questions}

When encountering a new object, ask four questions.

\subsection{State}\label{state-1}

What information is instantiated?

\subsection{Operator}\label{operator-1}

What transformation acts?

\subsection{Flow}\label{flow}

How do transformations accumulate?

\subsection{Interface}\label{interface-1}

What may another component assume at the boundary?

Not every object requires all four answers.

But asking the questions often exposes hidden assumptions.

\section{Handoff to linear algebra}\label{handoff-to-linear-algebra}

The next mathematical chapter specializes the operator role.

For finite-dimensional linear maps, bases turn operators into matrices.

That creates a powerful computational language:

\begin{itemize}
\tightlist
\item
  projections;
\item
  eigendecompositions;
\item
  singular values;
\item
  pseudoinverses;
\item
  norms;
\item
  conditioning;
\item
  low-rank structure.
\end{itemize}

The important point is already in place:

\[
\boxed{
\text{matrix coordinates}
\neq
\text{operator role by definition}.
}
\]

Linear algebra will now make that distinction precise.

\subsection{Handoff to information}\label{handoff-to-information}

Probability introduces another kind of state.

A distribution is an information-bearing mathematical object.

Entropy, KL divergence, and mutual information are functionals of such
states.

Again, storage format is secondary.

What matters is the declared probability model and the operation being
performed on it.

\section{Closing view}\label{closing-view-2}

The first conceptual discipline of the Atlas is simple:

Do not ask only what shape an object has.

Ask what role it plays.

A state is what is instantiated.

An operator transforms.

A flow organizes transformation through evolution.

An interface states what composition may assume.

These distinctions are not the final theory.

They are the grammar that lets the rest of the book say what kind of
thing it is talking about.

\section{References used in this
chapter}\label{references-used-in-this-chapter-2}

This chapter is Atlas synthesis. Its project-local vocabulary is bound
through sources/source-locks/ATLAS-CH-OBJECTS-001.yaml to the source
inventory, Mathematical Lexicon, and Atlas Map. Standard mathematical
instances are developed and externally sourced in downstream chapters.

\chapter{Claims, Evidence, and Computational
Witnesses}\label{claims-evidence-and-computational-witnesses}

\textbf{Epistemic status:} Atlas documentary synthesis with one exact
finite computational witness.\\
\textbf{Specification:}
manuscript/specifications/ATLAS-CH-EVIDENCE-001.md\\
\textbf{Documentary packet:}
mathematics/derivations/ATLAS-CH-EVIDENCE-001-DOCUMENTARY.md\\
\textbf{Computational witness:}
mathematics/computational-witnesses/ATLAS-CW-EVIDENCE-001.md\\
\textbf{Source lock:} sources/source-locks/ATLAS-CH-EVIDENCE-001.yaml

A mathematical atlas does not merely contain statements. It tells the
reader why some statements may be used as foundations, why others are
observations, why still others are proposals, and where the boundary
between them lies.

That distinction matters because polished exposition can make unlike
things look deceptively similar. A theorem and a numerical plot may
occupy adjacent pages. A conjecture may be written in notation as crisp
as a proved result. A project record may contain exact hashes and still
say nothing about whether the mathematics is true. A computational
witness may be exact down to the last bit and still establish only a
finite claim.

The first discipline of evidence is therefore separation.

A claim is one object. Its support is another. The interpretation placed
on that support is another. The institutional state attached to the work
is another.

The Atlas will often connect these objects. It must not fuse them.

\section{Start with the exact
claim}\label{start-with-the-exact-claim}

Before asking whether evidence is strong, ask what the evidence is
supposed to support.

Compare:

\begin{itemize}
\tightlist
\item
  this matrix has spectral norm greater than one;
\item
  this optimizer can amplify perturbations in this toy system;
\item
  this training run exhibited a transient spike;
\item
  non-normality caused the spike;
\item
  non-normal transient amplification is prevalent in frontier training.
\end{itemize}

These statements are not interchangeable. They differ in object,
quantifier, domain, and ambition. A calculation sufficient for the first
may merely illustrate the second. An experiment relevant to the third
may be suggestive but insufficient for the fourth. None automatically
establishes the fifth.

Evidence becomes intelligible only after the claim is made precise
enough to fail.

For that reason the Atlas treats claim identity as part of the
scientific object, not as decorative prose wrapped around a result.

\section{The claim-support packet}\label{the-claim-support-packet}

The chapter will use a compact documentary object:

\[
K=(q,\tau,S,\Omega,N,D).
\]

Its fields are:

q --- the exact claim.

τ --- the declared epistemic class.

S --- the support objects or support routes.

Ω --- the scope: hypotheses, domain, dataset, run, environment, version,
numerical convention, or other boundary that controls where the support
applies.

N --- explicit non-entailments: stronger conclusions the support does
not license.

D --- downstream permission: what a later chapter may consume without
silently strengthening the claim.

We write

\[
S \mathrel{\rightsquigarrow}_{\Omega,\tau} q.
\]

This records a scoped evidentiary support judgment, not a logical entailment.
The support relation must be assessed against the particular claim and scope.

to mean that S supports q under the declared scope and epistemic class.

The arrow is intentionally not a symbol for logical entailment.

Sometimes the support route is a proof, and then a genuine logical
relation can be stated. Sometimes it is a finite computation, an
empirical observation, a source-locked public artifact, or a review
record. These objects can support claims without all functioning as
proofs.

The notation exists to prevent a common category error: treating every
reason for belief as though it were the same kind of mathematical
implication.

\section{Eleven labels, not eleven
rungs}\label{eleven-labels-not-eleven-rungs}

The Atlas uses eleven canonical reader-facing epistemic classes.

{\def\LTcaptype{none} % do not increment counter
\begin{longtable}[]{@{}
  >{\raggedright\arraybackslash}p{(\linewidth - 4\tabcolsep) * \real{0.3333}}
  >{\raggedright\arraybackslash}p{(\linewidth - 4\tabcolsep) * \real{0.3333}}
  >{\raggedright\arraybackslash}p{(\linewidth - 4\tabcolsep) * \real{0.3333}}@{}}
\toprule\noalign{}
\begin{minipage}[b]{\linewidth}\raggedright
Class
\end{minipage} & \begin{minipage}[b]{\linewidth}\raggedright
What it records
\end{minipage} & \begin{minipage}[b]{\linewidth}\raggedright
What it does not create by itself
\end{minipage} \\
\midrule\noalign{}
\endhead
\bottomrule\noalign{}
\endlastfoot
Definition & A stipulated mathematical or documentary object & External
existence or empirical truth \\
Established Result & A result supported by appropriate external
scholarly or formal authority & Permission to drop the source
hypotheses \\
Atlas Derivation & A derivation carried out in this book from stated
assumptions & External attribution or independent validation \\
Computational Witness & A reproducible symbolic, numerical,
finite-search, simulation, graphical, or replay object & Proof or
certification merely by being exact or reproducible \\
Observation & What was empirically or computationally observed under
stated conditions & Universal mechanism or prevalence \\
Interpretation & A reasoned reading of results, observations, or
witnesses & Replacement for the underlying evidence \\
GCL Public Project Evidence & A claim grounded in an exact public GCL
artifact & Authority beyond what that artifact actually contains \\
GCL Programme & A research direction, hypothesis, remembered programme
term, or planned construction & Public implementation or completed
result \\
Conjecture & A substantive proposed claim not yet sufficiently proved &
Established result \\
Open Problem & A precise unresolved question or target & Existence of a
solution \\
Institutional Status & A governed state such as reviewed, accepted,
superseded, adjudicated, or certified & Mathematical truth by status
alone \\
\end{longtable}
}

It is tempting to place these classes on a ladder from weak to strong.
That would be a mistake.

A definition is not weak evidence; it is not evidence of that kind at
all. An open problem is not a low-confidence theorem. Institutional
status records a relation to an authority process, not a probability
that a mathematical sentence is true. Interpretation sits on top of
evidence without replacing it. Programme context can be valuable for
choosing research directions while remaining deliberately barred from
public implementation claims.

The vocabulary is therefore closer to a set of typed roles than to a
single confidence score. The classes do not form a scalar ladder.

Within a particular role, support can certainly be better or worse. A
proof can be correct or defective. An experiment can be well or poorly
controlled. A source can be primary or remote. But those judgments do
not turn the role taxonomy itself into one scalar order.

\section{Definition is not
discovery}\label{definition-is-not-discovery}

Definitions are unusually powerful because they let a subject become
precise.

Suppose the Atlas defines a boundary contract, a replayable evidence
object, or a Residual. The definition can be exact and useful
immediately. It can support derivations about any object satisfying the
definition.

What it cannot do is establish that the defined object occurs in nature,
in a deployed system, or in a frontier model.

This sounds obvious when stated explicitly. It is easy to forget when a
definition is elegant.

A name gives a phenomenon a handle. It does not give the phenomenon
empirical prevalence.

\section{Established results and Atlas
derivations}\label{established-results-and-atlas-derivations}

An established result enters the Atlas from an appropriate external
scholarly or formal source. The source matters, but so do its
hypotheses.

If a theorem holds for compact operators, smooth manifolds, convex
objectives, independent samples, finite state spaces, or asymptotic
limits, those conditions are part of the object being imported. Citation
does not dissolve them.

An Atlas derivation is different. It is mathematics carried out in the
book from declared assumptions or earlier results.

The distinction is not an insult to either class. It is provenance.

A correct Atlas derivation may be entirely rigorous. Its label tells the
reader where the derivation was performed and where responsibility for
the argument lies. If the same statement is later located in external
literature, the source record can be enriched. The original derivation
does not retroactively become someone else's result.

\section{Computational witnesses}\label{computational-witnesses}

A computational witness is one of the Atlas's most useful evidence forms
because many structures become visible only when symbolic, numerical,
finite-search, graphical, simulation, or replay machinery is allowed to
participate.

The rule is simple:

A witness must be as exact about its boundary as it is about its output.

Consider the elementary parity example recorded in
ATLAS-CW-EVIDENCE-001.

We enumerate every integer from -10 through 10 and compute n(n-1) modulo
2 using exact integer arithmetic. Every residue is zero.

That computation establishes a complete finite statement:

for each of those 21 integers, n(n-1) is even.

It does not prove the universal statement:

for every integer n, n(n-1) is even.

The universal statement has a separate proof. Two consecutive integers
have opposite parity, so one is even, and their product is even.

The witness and the proof agree on the tested values. They do different
logical work.

This is the central discipline of computational evidence. Exactness
protects the computation from one class of doubt. Reproducibility
protects the reconstruction route from another. Neither property
silently enlarges the domain of the claim.

\subsection{What a witness can be excellent
at}\label{what-a-witness-can-be-excellent-at}

A computational witness can:

\begin{itemize}
\tightlist
\item
  verify a finite family exhaustively;
\item
  expose a counterexample;
\item
  check an algebraic identity at a declared set of exact inputs;
\item
  evaluate a numerical quantity under a stated convention;
\item
  reconstruct an artifact from pinned sources;
\item
  make geometry visible;
\item
  replay a deterministic procedure;
\item
  falsify a proposed universal statement with one valid counterexample.
\end{itemize}

\subsection{What a witness cannot acquire by
presentation}\label{what-a-witness-cannot-acquire-by-presentation}

A witness does not become a theorem because the plot is attractive.

It does not become independent replication because another process reran
the same internal pipeline.

It does not become certification because CI is green.

It does not become causal explanation because the observed pattern
matches an intuition.

The support route remains what it is.

\section{Observation is not
interpretation}\label{observation-is-not-interpretation}

An observation records what happened under stated conditions.

For example:

\begin{itemize}
\tightlist
\item
  a loss curve rose sharply at a particular step;
\item
  one intervention reduced a measured quantity;
\item
  a router changed assignments under a perturbation;
\item
  an ablation changed test accuracy on a declared benchmark;
\item
  a diagnostic statistic concentrated in a particular layer.
\end{itemize}

Interpretation asks what the observation means.

Perhaps the spike reflects transient amplification. Perhaps the
intervention removed a mechanism. Perhaps the router statistic indicates
specialization. These can be good interpretations. They are still
interpretations.

The distinction is especially important when several mechanisms can
produce similar observables.

A useful discipline is to ask whether another mechanism could produce
the same observation. If yes, the observation may constrain the
explanation without identifying it.

This is why later diagnostic chapters will need interventions,
substitutions, controls, counterfactuals, and competing explanations.
The evidence grammar comes first so those chapters can state exactly
which inferential step each experiment supplies.

\section{Public project evidence and programme
context}\label{public-project-evidence-and-programme-context}

Research programmes develop language faster than every idea becomes a
stable public artifact.

The Atlas therefore separates two project-local classes.

GCL Public Project Evidence refers to an exact public object: a
repository file, commit, artifact, result, or other inspectable object
whose identity can be locked.

GCL Programme refers to a research direction, remembered term, planned
construction, hypothesis, or context that is not being asserted as exact
public evidence.

The distinction protects both.

Public evidence should not be weakened by mixing it with recollection.
Programme context should not be forced into false certainty merely to
make the prose sound complete.

A programme idea can motivate a chapter, suggest an experiment, or
identify a useful open problem. It simply cannot be laundered into a
statement that an implementation or result exists publicly when no exact
public object has been bound.

\section{Conjecture and open
problem}\label{conjecture-and-open-problem-1}

A conjecture and an open problem are related but different.

A conjecture says: here is a substantive statement we suspect may be
true.

An open problem says: here is an unresolved question or target.

A good research programme often contains both.

For example, one may conjecture that a certain invariant controls
transfer, while the corresponding open problem asks for necessary and
sufficient conditions under which reconstruction succeeds.

Precision does not demote either object. A conjecture can be
mathematically sharp. An open problem can have a detailed attack plan.
What remains forbidden is pretending that a plausible route is already a
solution.

\section{Institutional status is a different
axis}\label{institutional-status-is-a-different-axis}

The Atlas also records institutional states.

A result may be reviewed, audited, accepted, superseded, adjudicated,
certified, or merged.

These states matter. They tell the reader what process the object has
passed through and which authority made the decision.

But process state and mathematical truth are not synonyms.

Green CI may show that declared automated checks passed.

A documentary audit may show that source identities, scope statements,
and manuscript claims agree.

A formal certificate may establish something much stronger if the
certification authority and formal statement are themselves specified.

The error is not in having institutional status. The error is in
allowing the status word to float free of the authority and checks that
give it meaning.

The same discipline applies in reverse. A true mathematical argument can
exist before any institution has reviewed it. Lack of institutional
status is not a proof of falsehood.

\section{Source identity is not source
authority}\label{source-identity-is-not-source-authority}

A source lock answers two different questions:

\begin{enumerate}
\def\labelenumi{\arabic{enumi}.}
\tightlist
\item
  Which exact object did the chapter consume?
\item
  What is that object authoritative for?
\end{enumerate}

These fields must remain separate.

A commit hash can identify a README perfectly. That does not make the
README a primary authority for a historical theorem.

A DOI can identify a paper perfectly. That does not make every sentence
in later secondary commentary a theorem of that paper.

A source can therefore be perfectly pinned and still be insufficient for
the claim being made.

This is why the Atlas rejects citation laundering. Precision of identity
is necessary for reconstructability. Adequacy of authority is a separate
judgment.

\section{Promotion requires new
support}\label{promotion-requires-new-support}

Suppose a chapter begins with an Observation and later acquires a proof
of a general theorem.

That is a legitimate promotion because the support changed.

Suppose instead that the observation is rewritten in more confident
prose, placed in a polished figure, cited repeatedly, and merged after
green CI.

Nothing epistemic has changed.

Presentation can reveal status. It cannot create status.

In the claim-support packet, a genuine promotion changes τ only when S
and usually Ω change in a way that licenses the stronger q.

This yields a useful editorial test:

If the claim sounds stronger after revision, point to the new support
object that justifies the stronger wording.

If no such object exists, the prose has outrun the evidence.

\section{Five common promotion
failures}\label{five-common-promotion-failures}

\subsection{Finite check to universal
theorem}\label{finite-check-to-universal-theorem}

A finite exhaustive search proves a finite exhaustive claim.

It proves an infinite claim only when there is an additional argument
reducing the infinite domain to the finite search.

Timeout is not a nonexistence proof. Large coverage is not an infinite
quantifier.

\subsection{Repeatable observation to universal
mechanism}\label{repeatable-observation-to-universal-mechanism}

Repeated observation within one family improves confidence that the
observation is stable in that family.

It does not alone prove that one mechanism caused it, that the mechanism
is unique, or that the phenomenon is prevalent outside the tested
regime.

\subsection{Narrow proof to broad
prose}\label{narrow-proof-to-broad-prose}

A proof can be flawless and still fail to support the sentence written
above it.

Typical causes include a wider domain, missing hypothesis, stronger
quantifier, causal language, or an implicit approximation-to-exactness
jump.

The cure is not to distrust proof. It is to align q with the theorem
actually proved.

\subsection{Precise citation to excessive
authority}\label{precise-citation-to-excessive-authority}

A source can be exact and still be the wrong kind of authority.

Source identity and source authority must both survive inspection.

\subsection{Procedure to truth}\label{procedure-to-truth}

CI, replay, review, audit, merge, adjudication, acceptance, and
certification are procedural or institutional facts.

Each can support important claims about what was checked.

None may be treated as a generic truth operator.

\section{What later chapters may now
assume}\label{what-later-chapters-may-now-assume}

Four architectural consumers depend directly on this chapter.

ATLAS-CH-DATA-001 --- Data Quality, Mixtures, and Contamination may
assume that empirical claims must retain dataset, mixture, benchmark,
contamination, and run scope.

ATLAS-CH-MECHDIAG-001 --- Mechanistic Diagnostics may assume the
distinction between Observation and Interpretation and may require
stronger interventions before promoting correlation to mechanism.

ATLAS-CH-EVIDEX-001 --- Evidence Extraction and Claim Discipline may
assume the claim-support packet and can refine it into operational
extraction and adjudication procedures.

ATLAS-CH-EXPERIMENT-001 --- Experiments as Arguments may assume that
controls, ablations, counterfactuals, effect sizes, seeds, and
reproducibility are support routes whose force depends on the exact
claim and scope.

The handoff is deliberately modest.

Later chapters inherit the grammar. They do not inherit domain evidence
they have not yet produced.

\section{The durable habit}\label{the-durable-habit}

The strongest practical lesson of this chapter is a question sequence.

When encountering a technical statement, ask:

What exactly is the claim?

What class is it in?

What supports it?

Under what scope?

What stronger inference is not licensed?

What may a later chapter safely reuse?

If those questions can be answered, evidence becomes inspectable rather
than atmospheric.

That is enough structure to let a large monograph combine theorem,
derivation, computation, experiment, interpretation, programme, and
governance without pretending they are the same thing.

\section{References used in this
chapter}\label{references-used-in-this-chapter-3}

This chapter is a project-local documentary synthesis. Its load-bearing
sources are exact Atlas governance and prerequisite objects pinned in:

sources/source-locks/ATLAS-CH-EVIDENCE-001.yaml

The chapter makes no claim that the Atlas taxonomy is a universal or
uniquely correct theory of scientific epistemology.

\chapter{Linear Maps and
Decompositions}\label{linear-maps-and-decompositions}

\textbf{Epistemic status:} Established Theory + Atlas Derivation\\
\textbf{Specification:}
manuscript/specifications/ATLAS-CH-LINALG-001.md\\
\textbf{Derivation packet:}
mathematics/derivations/ATLAS-CH-LINALG-001-DERIVATIONS.md\\
\textbf{Computational witness:}
mathematics/computational-witnesses/ATLAS-CW-LINALG-001.md\\
\textbf{Source lock:} sources/source-locks/ATLAS-CH-LINALG-001.yaml

\section{Linear algebra as transformation
language}\label{linear-algebra-as-transformation-language}

A matrix is often introduced as a rectangular table of numbers.

That description is correct and insufficient.

For the Atlas, the important object is the linear map.

Given finite-dimensional vector spaces \(V\) and \(W\), a linear
transformation is a map

\[
T:V\to W
\]

satisfying

\[
T(\alpha x+\beta y)
=
\alpha T(x)+\beta T(y).
\]

Once bases are chosen, \(T\) is represented by a matrix \(A\).

The matrix is a coordinate description of the map.

Changing basis changes the coordinates used to describe the same
transformation.

Changing the operator changes the transformation itself.

This distinction is the first reason linear algebra matters here.

\section{Changing lenses versus changing the
machine}\label{changing-lenses-versus-changing-the-machine}

Imagine observing one machine through different camera angles.

The pictures change.

The machine does not.

A change of basis plays a similar role.

The coordinates of a vector and the matrix representing an operator
change together.

The underlying linear relation can remain the same.

The analogy has a limit.

A basis change is an exact algebraic transformation, not merely a visual
viewpoint.

And many learned transformations are nonlinear, state-dependent, or not
globally representable by one fixed matrix.

The analogy is useful only for separating coordinate description from
operator identity.

\section{Bases and coordinate
representations}\label{bases-and-coordinate-representations}

Let

\[
B=\{b_1,\ldots,b_n\}
\]

be a basis of \(V\).

Any vector \(x\in V\) has a unique coordinate vector

\[
[x]_B.
\]

If \(T:V\to W\) and \(C\) is a basis of \(W\), then there is a matrix

\[
[T]_{C\leftarrow B}
\]

such that

\[
[T(x)]_C
=
[T]_{C\leftarrow B}[x]_B.
\]

This equation is the bridge between abstract linear maps and matrix
computation.

The Atlas will often work directly with matrices.

But the operator viewpoint remains useful because it prevents
basis-dependent coordinates from being mistaken for invariant structure.

\section{Subspaces}\label{subspaces}

A subspace is a subset closed under linear combination.

Three recurring subspaces are:

\[
\operatorname{range}(A),
\]

\[
\operatorname{null}(A),
\]

and, for a matrix acting on an inner-product space, orthogonal
complements of these spaces.

Subspaces matter because many later problems are really questions about
which directions:

\begin{itemize}
\tightlist
\item
  survive a transformation;
\item
  are erased;
\item
  are amplified;
\item
  are identifiable;
\item
  are reachable.
\end{itemize}

Rank is the dimension of the image:

\[
\operatorname{rank}(A)
=
\dim \operatorname{range}(A).
\]

A low-rank operator therefore communicates through a lower-dimensional
effective channel.

That can be useful compression.

It can also be information loss.

\section{Orthogonal projection}\label{orthogonal-projection}

Suppose \(Q\in\mathbb R^{n\times k}\) has orthonormal columns:

\[
Q^\top Q=I.
\]

Then

\[
P=QQ^\top
\]

is the orthogonal projector onto the column space of \(Q\).

Indeed,

\[
P^2
=
QQ^\top QQ^\top
=
QIQ^\top
=
P,
\]

and

\[
P^\top=P.
\]

Projection is one of the simplest examples of an operator whose geometry
is visible.

It keeps components inside one subspace and removes orthogonal
components.

Later chapters use this language for tangent projections, subspace
methods, low-rank structure, and constrained updates.

\section{Eigendecomposition}\label{eigendecomposition}

For a square matrix \(A\), an eigenpair satisfies

\[
Av=\lambda v.
\]

An eigenvector is a direction mapped back into its own span.

If \(A\) is diagonalizable,

\[
A=V\Lambda V^{-1}.
\]

Then powers satisfy

\[
A^k
=
V\Lambda^k V^{-1}.
\]

This is powerful.

It is also easy to overuse.

Eigenvalues directly describe asymptotic modal behavior only under
assumptions that later chapters will make explicit.

If \(V\) is badly conditioned, eigenvalues can be a poor guide to
finite-time amplification.

The Non-normality keystone exists largely because that distinction
matters.

\section{Normality}\label{normality}

A real matrix is normal when

\[
A^\top A
=
AA^\top.
\]

In complex spaces the transpose is replaced by the adjoint.

Normal matrices admit orthonormal eigenvector bases.

This property greatly simplifies spectral reasoning.

A non-normal matrix need not.

The Atlas will therefore use the word ``normal'' carefully.

It does not mean statistically normal.

It describes an operator relation.

\section{Singular-value
decomposition}\label{singular-value-decomposition}

For any real matrix

\[
A\in\mathbb R^{m\times n},
\]

the singular-value decomposition is

\[
\boxed{
A=U\Sigma V^\top
}
\]

with orthonormal columns in \(U\) and \(V\), and nonnegative singular
values in \(\Sigma\).

The right singular vectors identify input directions.

The left singular vectors identify corresponding output directions.

The singular values quantify amplification:

\[
Av_i
=
\sigma_i u_i.
\]

This makes the SVD one of the Atlas's central linear tools.

It works for rectangular matrices.

It does not require diagonalizability.

And it directly answers Euclidean gain questions.

Standard treatments appear in Trefethen and Bau (Trefethen and David Bau
1997), Horn and Johnson (Horn and Johnson 2012), and Golub and Van Loan
(Golub and Loan 2013).

\section{Operator norm}\label{operator-norm}

The induced Euclidean norm is

\[
\|A\|_2
=
\sup_{\|x\|_2=1}\|Ax\|_2.
\]

The SVD gives

\[
\boxed{
\|A\|_2
=
\sigma_{\max}(A).
}
\]

This is not merely notation.

It says the largest singular value is the greatest one-step
amplification of a unit vector under the Euclidean norm.

Later chapters on boundary sensitivity, optimization, and transient
growth use exactly this interpretation.

\section{Eigenvalues and singular values answer different
questions}\label{eigenvalues-and-singular-values-answer-different-questions}

Consider

\[
A=
\begin{pmatrix}
1&1\\
0&1
\end{pmatrix}.
\]

Both eigenvalues equal \(1\).

But the singular values are

\[
\sigma_1
=
\frac{1+\sqrt5}{2},
\]

and

\[
\sigma_2
=
\frac{\sqrt5-1}{2}.
\]

Therefore

\[
\|A\|_2
=
\frac{1+\sqrt5}{2}
>
1.
\]

A unit vector can be amplified in one step even though every eigenvalue
has modulus \(1\).

This is the smallest important lesson:

\[
\boxed{
\text{spectral radius}
\neq
\text{operator norm}
}
\]

in general.

The Non-normality chapter develops the consequences.

\section{Pseudoinverse and least
squares}\label{pseudoinverse-and-least-squares}

If a linear system

\[
Ax=b
\]

has no exact solution, one may seek

\[
x_\star
=
\arg\min_x \|Ax-b\|_2.
\]

Using

\[
A=U\Sigma V^\top,
\]

define the Moore--Penrose pseudoinverse

\[
A^+
=
V\Sigma^+U^\top,
\]

where each nonzero singular value is inverted.

Then

\[
x_\star=A^+b
\]

gives the minimum-norm least-squares solution under the standard
finite-dimensional assumptions.

This object recurs in inverse problems, linear probes, regression, and
local approximations.

\section{Low-rank approximation}\label{low-rank-approximation}

Suppose

\[
A=U\Sigma V^\top
\]

with singular values

\[
\sigma_1\ge\sigma_2\ge\cdots.
\]

Truncating after \(k\) singular components gives

\[
A_k
=
U_k\Sigma_kV_k^\top.
\]

The Eckart--Young theorem identifies \(A_k\) as a best rank-\(k\)
approximation in the spectral and Frobenius norms.

In the spectral norm,

\[
\|A-A_k\|_2
=
\sigma_{k+1}.
\]

For the witness matrix above, the best rank-one error is

\[
\frac{\sqrt5-1}{2}.
\]

Low rank is therefore a precise approximation concept.

It is not automatically semantic compression.

A small matrix approximation error can still erase a downstream feature
that matters.

\section{Conditioning}\label{conditioning}

Some problems are intrinsically sensitive.

Suppose

\[
D=
\begin{pmatrix}
1&0\\
0&1/100
\end{pmatrix}.
\]

Then

\[
\sigma_{\max}(D)=1,
\qquad
\sigma_{\min}(D)=1/100.
\]

The \(2\)-norm condition number is

\[
\boxed{
\kappa_2(D)=100.
}
\]

The weak singular direction is much more sensitive to inversion.

This is a property of the problem.

It must be distinguished from the stability of a particular algorithm
used to solve it.

That distinction becomes central in Numerical Intelligence.

\section{Condition number is not optimization
difficulty}\label{condition-number-is-not-optimization-difficulty}

A large matrix condition number can matter greatly.

But it does not by itself prove that a nonlinear training problem will
be difficult.

Training dynamics also depend on:

\begin{itemize}
\tightlist
\item
  parameterization;
\item
  optimizer state;
\item
  curvature variation;
\item
  stochastic gradients;
\item
  non-normality;
\item
  constraints;
\item
  numerical precision;
\item
  data ordering.
\end{itemize}

The Atlas therefore treats conditioning as one diagnostic quantity, not
a universal explanation.

\section{Block structure}\label{block-structure}

Large transformations often have meaningful block structure.

For example,

\[
A=
\begin{pmatrix}
A_{11}&A_{12}\\
A_{21}&A_{22}
\end{pmatrix}
\]

can represent coupling among subsystems.

Block Jacobians later make optimizer-state dynamics explicit.

Boundary composition can also be analyzed through block sensitivity.

The point is structural:

a matrix need not be an undifferentiated table.

Its partition can encode a model of interactions.

\section{Kronecker structure}\label{kronecker-structure}

For matrices \(A\) and \(B\), the Kronecker product

\[
A\otimes B
\]

constructs a larger structured linear map.

Kronecker products appear in:

\begin{itemize}
\tightlist
\item
  separable transformations;
\item
  tensor-product bases;
\item
  covariance approximations;
\item
  structured preconditioners.
\end{itemize}

This chapter does not develop the full theory.

It establishes the notation and the idea that large operators may be
described compositionally.

\section{What linearization does and does not buy
us}\label{what-linearization-does-and-does-not-buy-us}

Even nonlinear systems have local linear structure.

If

\[
F:\mathbb R^n\to\mathbb R^m
\]

is differentiable, then near \(x\),

\[
F(x+\delta)
\approx
F(x)+J_F(x)\delta.
\]

The Jacobian is a linear operator.

This lets the Atlas reuse:

\begin{itemize}
\tightlist
\item
  singular values;
\item
  operator norms;
\item
  projections;
\item
  condition estimates.
\end{itemize}

But local linearization is local.

A Jacobian at one point does not automatically describe a global
nonlinear system.

The Boundary Contracts and optimizer-dynamics chapters both rely on this
limitation.

\section{Computational witness}\label{computational-witness-1}

\begin{figure}
\centering
\pandocbounded{\includegraphics[keepaspectratio,alt={The unit circle mapped by a non-normal matrix into an ellipse beside an exact conditioning example with condition number 100.}]{figures/masters/ATLAS-FIG-LINALG-001.png}}
\caption{The unit circle mapped by a non-normal matrix into an ellipse
beside an exact conditioning example with condition number 100.}
\end{figure}

The companion witness replays three exact facts.

For

\[
A=
\begin{pmatrix}
1&1\\
0&1
\end{pmatrix},
\]

Wolfram returns eigenvalues

\[
1,1
\]

and singular values

\[
\frac{1+\sqrt5}{2},
\qquad
\frac{\sqrt5-1}{2}.
\]

For the projection

\[
P=
\begin{pmatrix}
1&0\\
0&0
\end{pmatrix},
\]

it verifies

\[
P^2=P.
\]

For

\[
D=
\begin{pmatrix}
1&0\\
0&1/100
\end{pmatrix},
\]

it verifies

\[
\kappa_2(D)=100.
\]

The witness is exact finite-dimensional mathematics.

It does not imply a neural-network mechanism.

\section{Five failure modes}\label{five-failure-modes}

\subsection{Coordinates mistaken for
invariants}\label{coordinates-mistaken-for-invariants}

A matrix entry can change under basis change while the operator remains
the same.

\subsection{Eigenvalues mistaken for all
amplification}\label{eigenvalues-mistaken-for-all-amplification}

Non-normal operators can have large singular gain despite benign-looking
eigenvalues.

\subsection{Low rank mistaken for semantic
sufficiency}\label{low-rank-mistaken-for-semantic-sufficiency}

A good norm approximation can still lose task-critical information.

\subsection{Conditioning mistaken for algorithm
failure}\label{conditioning-mistaken-for-algorithm-failure}

An ill-conditioned problem and an unstable algorithm are different
objects.

\subsection{Local linearization mistaken for global
behavior}\label{local-linearization-mistaken-for-global-behavior}

A Jacobian is a local differential object.

\section{Atlas connections}\label{atlas-connections}

\textbf{Geometry.}\\
Tangent spaces are linear spaces attached to nonlinear manifolds.

\textbf{Attention.}\\
Attention matrices act as state-dependent linear mixing operators on
value fields when the mixing weights are fixed.

\textbf{Optimization.}\\
Hessians, preconditioners, and update operators are linear-algebraic
objects embedded in nonlinear dynamics.

\textbf{Numerical intelligence.}\\
Conditioning, projections, factorizations, and Krylov methods are core
tools.

\textbf{Memory and representation.}\\
Low-rank and subspace structure can describe compressed state, but
semantic adequacy is a separate question.

\section{Closing view}\label{closing-view-3}

Linear algebra is not here because neural networks contain matrices.

It is here because adaptive systems repeatedly ask linear questions
inside larger nonlinear ones.

Which directions survive?

Which directions are amplified?

Which subspace matters?

Which information is lost?

Which coordinate change is merely descriptive?

Which operator is ill-conditioned?

The SVD, projections, pseudoinverses, norms, and decompositions give
precise answers to those questions.

They become useful when we remember what kind of object they describe.

\section{References used in this
chapter}\label{references-used-in-this-chapter-4}

\begin{itemize}
\tightlist
\item
  (Trefethen and David Bau 1997)
\item
  (Horn and Johnson 2012)
\item
  (Golub and Loan 2013)
\end{itemize}

See sources/source-locks/ATLAS-CH-LINALG-001.yaml for exact source roles
and claim scope.

\chapter{Normality, Pseudospectra, and Transient
Growth}\label{normality-pseudospectra-and-transient-growth}

\textbf{Epistemic status:} mathematical exposition based on standard
matrix analysis and pseudospectral theory, with Atlas-owned
finite-dimensional derivations.\\
\textbf{Primary figure:} `ATLAS-FIG-PSPECTRUM-001`\\
\textbf{Derivation packet:}
`mathematics/derivations/ATLAS-CH-NONNORMAL-001-DERIVATIONS.md`

\section{Stable eigenvalues can tell an incomplete
story}\label{stable-eigenvalues-can-tell-an-incomplete-story}

Suppose every eigenvalue of a matrix lies strictly inside the unit disk.

For a normal matrix, this strongly controls the powers of the matrix.
For a non-normal matrix, it may not control what happens over the first
few steps.

The key distinction is between:

\[
\text{asymptotic spectral behavior}
\]

and

\[
\text{finite-horizon amplification}.
\]

A matrix can satisfy

\[
\rho(A)<1
\]

and still have

\[
\|A^k\|_2>1
\]

for some finite \(k\).

This chapter explains exactly how.

The useful allegory is \textbf{aligned currents in a harbor}.

Each current may eventually decay, but if their directions are badly
aligned, motion can transfer among them in a way that produces a
temporary surge.

The correspondence is:

\begin{itemize}
\tightlist
\item
  modal decay ↔ eigenvalues inside the unit disk;
\item
  non-orthogonal modal geometry ↔ non-normality;
\item
  temporary surge ↔ transient norm growth.
\end{itemize}

The limit is strict. A matrix need not represent a fluid. The analogy is
only about geometric interaction among directions.

\section{Normality}\label{normality-1}

A real matrix \(A\) is normal when

\[
A^\top A=AA^\top.
\]

For complex matrices, transpose is replaced by conjugate transpose.

Normal matrices admit an orthogonal/unitary eigenbasis. In the real
symmetric or complex normal setting,

\[
A=U\Lambda U^*,
\]

with \(U\) unitary.

Then

\[
A^k=U\Lambda^kU^*,
\]

so

\[
\|A^k\|_2
=
\|\Lambda^k\|_2
=
\max_i|\lambda_i|^k
=
\rho(A)^k.
\]

Thus if

\[
\rho(A)<1,
\]

the operator norm decays monotonically with \(k\).

For normal matrices, eigenvalue geometry and singular-value growth align
unusually well (Horn and Johnson 2012).

\section{Non-normality breaks that
alignment}\label{non-normality-breaks-that-alignment}

If \(A\) is diagonalizable but non-normal,

\[
A=V\Lambda V^{-1},
\]

then

\[
A^k
=
V\Lambda^kV^{-1},
\]

so

\[
\|A^k\|_2
\le
\kappa_2(V)\rho(A)^k,
\]

where

\[
\kappa_2(V)=\|V\|_2\|V^{-1}\|_2.
\]

A badly conditioned eigenvector matrix can therefore permit substantial
finite-horizon amplification even when every eigenmode decays.

The deeper contrast is not simply

\[
\text{diagonalizable}
\quad\text{versus}\quad
\text{non-diagonalizable}.
\]

It is closer to

\[
\text{orthogonal modal geometry}
\quad\text{versus}\quad
\text{non-orthogonal modal geometry}.
\]

Normality removes the amplification associated with geometric
misalignment of eigen-directions.

\section{The smallest useful
counterexample}\label{the-smallest-useful-counterexample}

Consider

\[
A=
\begin{pmatrix}
a&K\\
0&a
\end{pmatrix},
\qquad
|a|<1.
\]

Write

\[
A=aI+KN,
\qquad
N=
\begin{pmatrix}
0&1\\
0&0
\end{pmatrix},
\qquad
N^2=0.
\]

For each integer \(n\ge 2\), the nilpotent binomial expansion truncates:

\[
A^n
=
a^nI
+
n a^{n-1}KN.
\]

For these \(n\), therefore

\[
\boxed{
A^n
=
\begin{pmatrix}
a^n&nKa^{n-1}\\
0&a^n
\end{pmatrix}.
}
\]

The low powers are \(A^0=I\) and \(A^1=A\). Treating them separately
avoids the \(a^0\) convention at \(a=0\).

The eigenvalue is always \(a\), with algebraic multiplicity two, so

\[
\rho(A)=|a|.
\]

If \(|a|<1\), then \(A^n\to0\).

But the off-diagonal term

\[
nKa^{n-1}
\]

can first grow before eventual decay.

The matrix is asymptotically stable and can still amplify some
perturbations strongly over a finite horizon.

\section{Singular values expose the
amplification}\label{singular-values-expose-the-amplification}

For

\[
B=
\begin{pmatrix}
p&q\\
0&p
\end{pmatrix},
\]

we have

\[
B^\top B
=
\begin{pmatrix}
p^2&pq\\
pq&p^2+q^2
\end{pmatrix}.
\]

Its eigenvalues are

\[
\lambda_{\pm}
=
\frac{
2p^2+q^2
\pm
|q|\sqrt{4p^2+q^2}
}{2}.
\]

The spectral norm is therefore

\[
\boxed{
\|B\|_2
=
\sqrt{
\frac{
2p^2+q^2
+
|q|\sqrt{4p^2+q^2}
}{2}
}.
}
\]

For \(B=A^n\) with \(n\ge 2\),

\[
p=a^n,
\qquad
q=nKa^{n-1}.
\]

This is an exact finite-horizon gain formula.

Nothing in the eigenvalue set alone displays the \(q\)-term.

\section{A matched normal
comparison}\label{a-matched-normal-comparison}

The Atlas witness fixes

\[
a=\frac45,
\qquad
K=4.
\]

Then

\[
A=
\begin{pmatrix}
4/5&4\\
0&4/5
\end{pmatrix}.
\]

Compare it with

\[
N_0=\frac45I.
\]

Both matrices have exactly the same eigenvalue set:

\[
\left\{\frac45,\frac45\right\}.
\]

For the normal matrix,

\[
\|N_0^n\|_2
=
\left(\frac45\right)^n.
\]

For the non-normal matrix, Wolfram evaluation gives

\[
\max_{0\le n\le40}\|A^n\|_2
=
8.2124290544\ldots
\]

at

\[
n=4.
\]

The two matrices have the same eigenvalue story and radically different
finite-time stories.

\section{The primary figure}\label{the-primary-figure}

\begin{figure}
\centering
\pandocbounded{\includegraphics[keepaspectratio,alt={Matched-eigenvalue normal and non-normal matrices compared through transient 2-norm gain and pseudospectral contours.}]{figures/derivatives/ATLAS-FIG-PSPECTRUM-001-v0.1.1.png}}
\caption{Matched-eigenvalue normal and non-normal matrices compared
through transient 2-norm gain and pseudospectral contours.}
\end{figure}

The left panel compares

\[
\|A^n\|_2
\]

against

\[
\|N_0^n\|_2.
\]

The right panel compares their \(\varepsilon\)-pseudospectral
boundaries.

The plate is not merely numerical decoration. For this \(2\times2\)
example, the pseudospectral radii are available in closed form.

\section{The pseudospectrum}\label{the-pseudospectrum}

We use the closed \(2\)-norm convention

\[
\boxed{
\Lambda_\varepsilon(A)
=
\left\{
z\in\mathbb C:
\sigma_{\min}(zI-A)\le\varepsilon
\right\}.
}
\]

Outside the spectrum,

\[
\|(zI-A)^{-1}\|_2
=
\frac1{\sigma_{\min}(zI-A)},
\]

so equivalently,

\[
\Lambda_\varepsilon(A)
=
\sigma(A)
\cup
\left\{
z:
\|(zI-A)^{-1}\|_2\ge\varepsilon^{-1}
\right\}.
\]

Another equivalent finite-dimensional characterization is:

\[
z\in\Lambda_\varepsilon(A)
\]

if and only if \(z\) is an eigenvalue of \(A+E\) for some perturbation
satisfying

\[
\|E\|_2\le\varepsilon.
\]

These equivalent views connect pseudospectra to both resolvent growth
and eigenvalue sensitivity (Trefethen and Embree 2005).

\section{\texorpdfstring{Exact pseudospectrum of the \(2\times2\)
example}{Exact pseudospectrum of the 2\textbackslash times2 example}}\label{exact-pseudospectrum-of-the-2times2-example}

For

\[
A=
\begin{pmatrix}
a&K\\
0&a
\end{pmatrix},
\qquad
K>0,
\]

write

\[
z-a=\delta,
\qquad
r=|\delta|.
\]

Then

\[
zI-A
=
\begin{pmatrix}
\delta&-K\\
0&\delta
\end{pmatrix}.
\]

Its squared singular values are

\[
\sigma_{\pm}^2
=
\frac{
2r^2+K^2
\pm
K\sqrt{K^2+4r^2}
}{2}.
\]

On the \(\varepsilon\)-pseudospectral boundary,

\[
\sigma_{\min}=\varepsilon.
\]

Because the product of the singular values is

\[
|\det(zI-A)|=r^2,
\]

we have

\[
\sigma_{\max}
=
\frac{r^2}{\varepsilon}.
\]

The sum of squared singular values is

\[
\sigma_{\max}^2+\sigma_{\min}^2
=
2r^2+K^2.
\]

Substituting gives

\[
\frac{r^4}{\varepsilon^2}
+
\varepsilon^2
=
2r^2+K^2.
\]

Hence

\[
(r^2-\varepsilon^2)^2
=
K^2\varepsilon^2.
\]

The admissible boundary branch is

\[
\boxed{
r^2
=
\varepsilon(\varepsilon+K).
}
\]

Therefore

\[
\boxed{
\Lambda_\varepsilon(A)
=
\left\{
z:
|z-a|
\le
\sqrt{\varepsilon(\varepsilon+K)}
\right\}.
}
\]

For the normal comparison \(N_0=aI\),

\[
\boxed{
\Lambda_\varepsilon(N_0)
=
\{z:|z-a|\le\varepsilon\}.
}
\]

The eigenvalue plot is identical.

The perturbation neighborhood is not.

For small \(\varepsilon\),

\[
\sqrt{\varepsilon(\varepsilon+K)}
\sim
\sqrt{K\varepsilon},
\]

which can be much larger than \(\varepsilon\).

\section{Why pseudospectra matter}\label{why-pseudospectra-matter}

Eigenvalues answer:

\begin{quote}
Where are the exact modal growth factors?
\end{quote}

Pseudospectra ask:

\begin{quote}
How sensitive is that spectral picture, and how large can the resolvent
become nearby?
\end{quote}

For normal matrices, \(\varepsilon\)-pseudospectra are simply
\(\varepsilon\)-neighborhoods of the spectrum in the \(2\)-norm.

For non-normal matrices they can expand far beyond those neighborhoods.

This enlargement signals that small perturbations can move eigenvalues
much farther than the ordinary eigenvalue plot suggests, and that
resolvent response can be large.

Pseudospectra are therefore a way to visualize \textbf{spectral
fragility}.

They are not a causal diagnosis by themselves.

\section{Transient growth in broader non-normal
systems}\label{transient-growth-in-broader-non-normal-systems}

The importance of non-normal transient growth is not peculiar to the
Atlas toy matrix. It has a substantial literature in hydrodynamic
stability, where decaying eigenmodes can combine to produce substantial
transient energy growth (Reddy et al. 1993; Trefethen et al. 1993).

Those examples motivate the general lesson without making a direct
identification with neural optimization.

The Atlas will use the mathematics, not import the application
wholesale.

Later, when an optimizer-state Jacobian or router linearization is
non-normal, the correct question is not:

\begin{quote}
Is this secretly a fluid?
\end{quote}

It is:

\begin{quote}
Can non-orthogonal dynamical directions create finite-horizon
amplification that an eigenvalue-only analysis misses?
\end{quote}

\section{What the Wolfram witness
establishes}\label{what-the-wolfram-witness-establishes}

For the exact matrix pair

\[
A=
\begin{pmatrix}
4/5&4\\
0&4/5
\end{pmatrix},
\qquad
N_0=\frac45I,
\]

the witness establishes:

\begin{enumerate}
\def\labelenumi{\arabic{enumi}.}
\tightlist
\item
  identical eigenvalues;
\item
  exact powers of \(A\);
\item
  large finite-horizon \(2\)-norm gain for \(A\);
\item
  monotone decay for \(N_0\);
\item
  exact non-normal pseudospectral radius \[
  \sqrt{\varepsilon(\varepsilon+4)};
  \]
\item
  normal pseudospectral radius \[
  \varepsilon.
  \]
\end{enumerate}

The witness does \textbf{not} establish:

\begin{itemize}
\tightlist
\item
  that every non-normal system exhibits large transient growth;
\item
  that transient growth implies asymptotic instability;
\item
  that broad pseudospectra identify a causal mechanism in a trained
  model;
\item
  that one norm is universally privileged.
\end{itemize}

\section{Five mistakes to avoid}\label{five-mistakes-to-avoid}

\subsection{Mistake 1: ``Non-normal means
unstable.''}\label{mistake-1-non-normal-means-unstable.}

No.~The worked matrix satisfies \(\rho(A)<1\) and \(A^n\to0\).

\subsection{Mistake 2: ``Stable eigenvalues imply monotone
decay.''}\label{mistake-2-stable-eigenvalues-imply-monotone-decay.}

No.~The worked matrix amplifies strongly before decaying.

\subsection{Mistake 3: ``A transient spike proves
divergence.''}\label{mistake-3-a-transient-spike-proves-divergence.}

No.~Finite-horizon gain and asymptotic divergence are different.

\subsection{Mistake 4: ``Pseudospectra replace
eigenvalues.''}\label{mistake-4-pseudospectra-replace-eigenvalues.}

No.~They supplement the spectral picture by adding perturbation and
resolvent information.

\subsection{Mistake 5: ``Broad pseudospectra explain a neural
failure.''}\label{mistake-5-broad-pseudospectra-explain-a-neural-failure.}

No.~They expose a possible amplification mechanism. Causal attribution
requires additional evidence.

\section{Atlas connections}\label{atlas-connections-1}

\textbf{Spectral shaping.}\\
An optimizer that changes singular structure may alter transient
behavior even when eigenvalue summaries appear similar.

\textbf{Optimizer-State Dynamics.}\\
The augmented optimizer Jacobian can inherit exactly the
stable-but-amplifying structure developed here.

\textbf{Router dynamics.}\\
Routing systems can possess non-normal local update operators and
finite-horizon amplification.

\textbf{Spectral diagnostics.}\\
Eigenvalues, singular values, pseudospectra, and finite-horizon
propagators become complementary diagnostic objects.

The recurring Atlas shift is:

\[
\boxed{
\text{Where are the eigenvalues?}
\longrightarrow
\text{What can this operator do to perturbations over the horizon that matters?}
}
\]

\section{Closing view}\label{closing-view-4}

Eigenvalues are not wrong.

They are sometimes incomplete.

Normal matrices make eigenvalue geometry unusually transparent.
Non-normal matrices can hide large finite-time behavior in the geometry
of their directions.

Pseudospectra reveal part of that hidden sensitivity.

The aligned-currents allegory is only an intuition pump.

The operator norm, resolvent, singular values, and pseudospectrum are
the objects.

\section{References used in this
chapter}\label{references-used-in-this-chapter-5}

\begin{itemize}
\tightlist
\item
  (Horn and Johnson 2012)
\item
  (Trefethen and Embree 2005)
\item
  (Reddy et al. 1993)
\item
  (Trefethen et al. 1993)
\end{itemize}

See `sources/source-locks/ATLAS-CH-NONNORMAL-001.yaml` for exact source
identities and claim scope.

\chapter{Krylov Subspaces and Iterative
Solves}\label{krylov-subspaces-and-iterative-solves}

\textbf{Epistemic status:} established numerical linear algebra +
audited Atlas substrate + Atlas synthesis.\\
\textbf{Specification:}
manuscript/specifications/ATLAS-CH-KRYLOV-001.md\\
\textbf{Derivation packet:}
mathematics/derivations/ATLAS-CH-KRYLOV-001-DERIVATIONS.md\\
\textbf{Computational witness:}
mathematics/computational-witnesses/ATLAS-CW-KRYLOV-001.md\\
\textbf{Source lock:} sources/source-locks/ATLAS-CH-KRYLOV-001.yaml

Large linear problems are often attacked without solving the whole space
at once.

The key move is to build a sequence of informative directions from
repeated application of the same operator to a starting residual. The
resulting trial space is not arbitrary. It is generated by the operator
itself.

For a square operator A and starting residual r0,

K\_m(A,r0) = span\{r0, A r0, \ldots, A\^{}(m-1) r0\}.

For Ax=b with starting point x0 and r0=b-Ax0, an iterative subspace
method seeks

xm in x0 + K\_m(A,r0).

This chapter develops the classical structure behind that idea.

\section{Why operator-generated directions
matter}\label{why-operator-generated-directions-matter}

A direct method tries to expose enough of the full matrix structure to
solve the system outright.

An iterative subspace method instead asks:

\begin{itemize}
\tightlist
\item
  which directions become visible after repeated operator action;
\item
  what reduced problem appears in those directions;
\item
  whether the reduced correction meaningfully reduces the target error
  or residual;
\item
  how numerical conditioning changes that interpretation.
\end{itemize}

This distinction matters when A is large, sparse, structured, or
available only through a routine that computes A v.

\section{Arnoldi: general matrices}\label{arnoldi-general-matrices}

Starting from a normalized residual v1, Arnoldi repeatedly applies A and
orthogonalizes the result against the basis already built.

In exact arithmetic it produces orthonormal columns V\_m and an
upper-Hessenberg reduced relation

A V\_m = V\_(m+1) Hbar\_m.

Three objects must remain distinct:

\begin{itemize}
\tightlist
\item
  the generated trial space;
\item
  the orthonormal basis used to represent it;
\item
  the reduced projected problem solved in that basis.
\end{itemize}

The Arnoldi relation is general for square matrices. It does not require
symmetry.

\section{Lanczos: the symmetric
specialization}\label{lanczos-the-symmetric-specialization}

When A is symmetric, or Hermitian in the complex case, the Arnoldi
projection becomes symmetric as well as upper Hessenberg.

A symmetric Hessenberg matrix is tridiagonal.

The exact-arithmetic recurrence therefore collapses to three terms:

beta\_(j-1) v\_(j-1) + alpha\_j v\_j + beta\_j v\_(j+1) = A v\_j.

This is the Lanczos structure.

The short recurrence is a structural consequence of symmetry in exact
arithmetic. It is not a promise that floating-point basis vectors remain
exactly orthogonal forever.

\section{Galerkin projection}\label{galerkin-projection}

Write

xm=x0+V\_m y.

Then

rm=r0-A V\_m y.

A Galerkin method chooses y so that the new residual is orthogonal to
the trial space:

V\_m\^{}T rm=0.

This gives the reduced system

(V\_m\^{}T A V\_m)y = V\_m\^{}T r0.

The method therefore solves a smaller problem whose dimension is m
rather than the ambient dimension.

\section{Minimum residual is a different projection
rule}\label{minimum-residual-is-a-different-projection-rule}

A minimum-residual method instead chooses y to minimize

\textbar\textbar r0-A V\_m y\textbar\textbar\_2.

With an Arnoldi basis, this becomes a small least-squares problem
involving Hbar\_m.

The distinction is important:

\begin{itemize}
\tightlist
\item
  Galerkin imposes an orthogonality condition;
\item
  minimum residual explicitly minimizes residual norm;
\item
  neither statement is automatically the same as minimizing solution
  error in Euclidean norm.
\end{itemize}

\section{Residual and error are related, not
identical}\label{residual-and-error-are-related-not-identical}

If x* is the exact solution and

em=x*-xm,

then

A em=rm.

For nonsingular A,

em=A\^{}(-1) rm.

Therefore

\textbar\textbar em\textbar\textbar\_2 \textless=
\textbar\textbar A\^{}(-1)\textbar\textbar\_2
\textbar\textbar rm\textbar\textbar\_2.

A small residual can support a useful error bound when the inverse
operator is well behaved.

A badly conditioned operator can magnify residual-to-error conversion.

This is why convergence statements must say which norm and which
quantity are being controlled.

\section{Polynomial viewpoint}\label{polynomial-viewpoint}

Every correction in the m-step generated subspace can be written as

q\_(m-1)(A) r0

for a polynomial of degree at most m-1.

The residual has the form

rm=p\_m(A) r0

with p\_m(0)=1.

This viewpoint connects iterative linear algebra to polynomial
approximation.

It also explains why spectral distribution, clustering, non-normality,
and the initial residual can strongly affect convergence.

\section{Exact two-dimensional
witness}\label{exact-two-dimensional-witness}

Take

A=diag(1,2,4),

b=(1,1,1)\^{}T,

x0=0.

Then

r0=b,

Ab=(1,2,4)\^{}T.

Use the two-column basis

V={[}b,Ab{]} = {[}{[}1,1{]},{[}1,2{]},{[}1,4{]}{]}.

The reduced Galerkin system is

(V\^{}T A V)c=V\^{}T b.

Direct arithmetic gives

V\^{}T A V={[}{[}7,21{]},{[}21,73{]}{]},

V\^{}T b={[}3,7{]}\^{}T,

and

c={[}36/35,-1/5{]}\^{}T.

Hence

x2={[}29/35,22/35,8/35{]}\^{}T.

The residual is

r2={[}6,-9,3{]}\^{}T/35.

Now check the defining orthogonality:

b\^{}T r2=0,

(Ab)\^{}T r2=0.

Thus

V\^{}T r2=0.

The squared residual norm is

\textbar\textbar r2\textbar\textbar\_2\^{}2=18/175.

The exact solution is

x*={[}1,1/2,1/4{]}\^{}T.

The resulting error satisfies

A e2=r2,

as it must.

\section{The same witness in an orthonormal symmetric
basis}\label{the-same-witness-in-an-orthonormal-symmetric-basis}

Normalize

v1=(1,1,1)\^{}T/sqrt(3).

Then

alpha1=7/3,

beta1=sqrt(14)/3,

and

v2=(-4,-1,5)\^{}T/sqrt(42).

The next projected diagonal coefficient is

alpha2=59/21.

So the two-dimensional projected matrix is

T2={[}{[}7/3,sqrt(14)/3{]},{[}sqrt(14)/3,59/21{]}{]}.

This is the symmetric tridiagonal reduction in explicit form.

\section{A conditioning control}\label{a-conditioning-control}

Now take

Ac=diag(1,100,10000),

with the same b and x0.

The one-dimensional SPD Galerkin/conjugate-gradient step has

alpha=1/3367.

Its Euclidean residual is

r1={[}3366,3267,-6633{]}\^{}T/3367,

with

\textbar\textbar r1\textbar\textbar\_2\^{}2=19602/3367.

That is larger than the initial residual norm squared:

\textbar\textbar r0\textbar\textbar\_2\^{}2=3.

So a valid conjugate-gradient step can increase Euclidean residual norm.

There is no contradiction.

For SPD systems, the relevant optimality statement for conjugate
gradient concerns the A-norm of the error over the affine trial space,
not monotone decrease of Euclidean residual norm.

The example therefore separates:

\begin{itemize}
\tightlist
\item
  residual norm;
\item
  error norm;
\item
  energy norm;
\item
  method-specific optimality.
\end{itemize}

\section{Matrix-free access}\label{matrix-free-access}

Nothing in the definition requires an explicitly materialized dense
matrix.

If the implementation can compute

v -\textgreater{} A v,

it can generate operator-derived directions.

This matters for large sparse systems and for implicit operators.

The real cost then includes:

\begin{itemize}
\tightlist
\item
  operator-vector products;
\item
  orthogonalization;
\item
  projected solves;
\item
  storage;
\item
  preconditioner application;
\item
  restarts.
\end{itemize}

\section{Preconditioning changes the problem seen by the
iteration}\label{preconditioning-changes-the-problem-seen-by-the-iteration}

A left-preconditioned system uses

M\^{}(-1) A x = M\^{}(-1)b.

A right-preconditioned system uses

A M\^{}(-1)y=b,

then recovers

x=M\^{}(-1)y.

The generated space is therefore built from a different effective
operator.

Preconditioning is not equivalent to merely taking more iterations.

A good preconditioner changes the geometry or spectrum of the problem in
a way that makes the reduced iteration more effective.

\section{Finite precision}\label{finite-precision}

Exact Arnoldi basis vectors are orthogonal.

Finite-precision implementations can lose orthogonality.

For symmetric recurrences, this can produce repeated or distorted
spectral information if not controlled.

Possible responses include reorthogonalization, selective
reorthogonalization, restarting, or other stabilization strategies.

Each changes cost and information retention.

\section{Restarting is a tradeoff}\label{restarting-is-a-tradeoff}

Unrestarted Arnoldi stores more basis vectors and pays increasing
orthogonalization cost.

Restarting limits memory and work per cycle.

But restarting can also discard useful directions accumulated earlier.

The right tradeoff depends on the operator, spectrum, target accuracy,
and memory budget.

\section{What this chapter does not
claim}\label{what-this-chapter-does-not-claim}

This chapter does not claim:

\begin{itemize}
\tightlist
\item
  that low-dimensional generated spaces always converge rapidly;
\item
  that small residual always means small solution error;
\item
  that eigenvalues alone determine convergence for non-normal operators;
\item
  that exact-arithmetic orthogonality survives floating point
  automatically;
\item
  that more iterations are equivalent to better preconditioning;
\item
  that a learned nonlinear architecture inherits classical convergence
  results merely because it uses iterative transport or
  operator-generated directions.
\end{itemize}

\section{Downstream handoff}\label{downstream-handoff}

Approximate and Structured Attention may inherit the principle that
repeated operator action and projection can expose useful
low-dimensional structure without full dense computation.

Neural iterative transport may inherit the classical trial-space,
residual, projection, and preconditioning vocabulary.

It must independently justify any learned, nonlinear, adaptive, or
representation-level extension.

\section{References used in this
chapter}\label{references-used-in-this-chapter-6}

\begin{itemize}
\tightlist
\item
  (Arnoldi 1951)
\item
  (Lanczos 1950)
\item
  (Saad 2003)
\item
  (Trefethen and David Bau 1997)
\item
  (Golub and Loan 2013)
\end{itemize}

See sources/source-locks/ATLAS-CH-KRYLOV-001.yaml for exact source
identities and claim scope.

\chapter{Probability, Information, and Statistical
Structure}\label{probability-information-and-statistical-structure}

\textbf{Epistemic status:} Established Theory + Atlas Derivation\\
\textbf{Specification:} manuscript/specifications/ATLAS-CH-INFO-001.md\\
\textbf{Derivation packet:}
mathematics/derivations/ATLAS-CH-INFO-001-DERIVATIONS.md\\
\textbf{Computational witness:}
mathematics/computational-witnesses/ATLAS-CW-INFO-001.md\\
\textbf{Source lock:} sources/source-locks/ATLAS-CH-INFO-001.yaml

\section{Information requires a probability
model}\label{information-requires-a-probability-model}

Machine learning uses the word ``information'' constantly.

The word can mean:

\begin{itemize}
\tightlist
\item
  a random variable reduces uncertainty about another;
\item
  a code uses fewer bits;
\item
  a hidden state preserves a task-relevant distinction;
\item
  a message is semantically meaningful;
\item
  a variable causally influences another;
\item
  a dataset contains predictive signal.
\end{itemize}

These meanings are not interchangeable.

This chapter begins with the probabilistic meaning.

A probability model states how uncertainty is represented.

Information-theoretic quantities are then functions of that model.

The chapter's first discipline is:

\begin{quote}
never ask what entropy or mutual information means before asking what
probability model it is computed under.
\end{quote}

\section{Uncertainty budget, not semantic
content}\label{uncertainty-budget-not-semantic-content}

A useful allegory is an uncertainty budget.

Entropy says how uncertain a random variable is under a declared
distribution.

Conditional entropy says how much uncertainty remains after another
variable is known.

Mutual information says how much uncertainty is reduced, symmetrically,
by knowing the other variable.

The allegory fails if ``information'' is taken to mean semantic
significance.

A random identifier can have high entropy and almost no meaning to a
user.

A one-bit alarm can carry enormous operational importance.

Information theory measures structure in a probability model.

It does not directly measure meaning.

\section{Random variables and
distributions}\label{random-variables-and-distributions}

A random variable

\[
X:\Omega\to\mathcal X
\]

maps outcomes in a sample space to values.

In the discrete case, a mass function satisfies

\[
p(x)\ge0,
\]

and

\[
\sum_x p(x)=1.
\]

For continuous variables, probability is described relative to a measure
through a density where one exists.

This distinction matters.

Expressions such as entropy and KL divergence behave differently in
discrete and continuous settings.

The Atlas will not silently move between them.

\section{Conditional probability}\label{conditional-probability}

For events \(A\) and \(B\) with \(P(B)>0\),

\[
P(A\mid B)
=
\frac{P(A\cap B)}{P(B)}.
\]

For random variables, conditional distributions express how the law of
one variable changes when another is known.

This is the foundation for:

\begin{itemize}
\tightlist
\item
  Bayesian inference;
\item
  predictive modeling;
\item
  filtering;
\item
  probabilistic memory;
\item
  uncertainty-aware decisions.
\end{itemize}

Conditional probability is not a causal operator by itself.

It is a relation inside a probability model.

\section{Expectation}\label{expectation}

For discrete \(X\),

\[
\mathbb E[X]
=
\sum_x x\,p(x).
\]

Expectation is linear:

\[
\mathbb E[aX+bY]
=
a\mathbb E[X]+b\mathbb E[Y].
\]

That algebraic simplicity is one reason expectations organize so much
statistical learning.

Loss functions, gradients, risk, calibration, and uncertainty estimates
are often expectations under one distribution or another.

The distribution must be named.

\section{Entropy}\label{entropy}

For a finite random variable,

\[
\boxed{
H(X)
=
-\sum_x p(x)\log_2 p(x).
}
\]

Entropy is measured in bits when the logarithm has base \(2\).

A fair bit has

\[
H(X)=1.
\]

A deterministic bit has

\[
H(X)=0.
\]

These statements are exact.

They do not say that a fair bit is ``more meaningful.''

They say its value is more uncertain before observation.

Standard information-theory definitions and identities are developed in
Cover and Thomas (Cover and Thomas 2006).

\section{Joint and conditional
entropy}\label{joint-and-conditional-entropy}

For jointly distributed \(X,Y\),

\[
H(X,Y)
=
-\sum_{x,y}p(x,y)\log_2p(x,y).
\]

Conditional entropy satisfies

\[
H(X\mid Y)
=
H(X,Y)-H(Y).
\]

The entropy chain rule is

\[
\boxed{
H(X,Y)
=
H(X)+H(Y\mid X).
}
\]

This is one of the first examples of information decomposing according
to conditional structure.

\section{Cross-entropy}\label{cross-entropy}

If data are distributed as \(P\) but a model assigns distribution \(Q\),
the cross-entropy is

\[
H(P,Q)
=
-\mathbb E_{x\sim P}\log q(x).
\]

In the discrete case,

\[
H(P,Q)
=
-\sum_x p(x)\log q(x).
\]

Cross-entropy is therefore a model-relative quantity.

It depends on both the data distribution and the proposed predictive
distribution.

\section{KL divergence}\label{kl-divergence}

The Kullback--Leibler divergence is

\[
\boxed{
D_{\rm KL}(P\|Q)
=
\sum_x p(x)\log\frac{p(x)}{q(x)}
}
\]

when the support conditions make the expression well defined.

KL divergence is nonnegative.

It is not symmetric:

\[
D_{\rm KL}(P\|Q)
\neq
D_{\rm KL}(Q\|P)
\]

in general.

Therefore KL is not a metric.

It also can be infinite.

If \(P\) places positive mass where \(Q\) assigns zero mass, then

\[
D_{\rm KL}(P\|Q)=\infty.
\]

The direction matters.

\section{Cross-entropy
decomposition}\label{cross-entropy-decomposition}

For discrete distributions,

\[
H(P,Q)
=
H(P)
+
D_{\rm KL}(P\|Q).
\]

If \(P\) is fixed, minimizing cross-entropy with respect to \(Q\) is
equivalent to minimizing

\[
D_{\rm KL}(P\|Q).
\]

This identity is simple and widely used.

Its assumptions should remain visible.

The expression depends on the declared target distribution \(P\) and
model distribution \(Q\).

\section{Mutual information}\label{mutual-information}

Mutual information is

\[
\boxed{
I(X;Y)
=
D_{\rm KL}(P_{XY}\|P_XP_Y).
}
\]

Equivalently,

\[
I(X;Y)
=
H(X)-H(X\mid Y),
\]

and symmetrically,

\[
I(X;Y)
=
H(Y)-H(Y\mid X).
\]

If \(X\) and \(Y\) are independent,

\[
P_{XY}=P_XP_Y,
\]

so

\[
I(X;Y)=0.
\]

Mutual information therefore measures statistical dependence.

It does not specify a causal direction.

\section{Exact correlated-bit
witness}\label{exact-correlated-bit-witness}

\begin{figure}
\centering
\pandocbounded{\includegraphics[keepaspectratio,alt={A 2 by 2 exact joint probability table with uniform marginals and mutual information about 0.1887 bits, explicitly labeled as dependence rather than causal direction.}]{figures/masters/ATLAS-FIG-INFO-001.png}}
\caption{A 2 by 2 exact joint probability table with uniform marginals
and mutual information about 0.1887 bits, explicitly labeled as
dependence rather than causal direction.}
\end{figure}

Consider

\[
P_{XY}
=
\begin{pmatrix}
3/8&1/8\\
1/8&3/8
\end{pmatrix}.
\]

Both marginals are fair:

\[
P_X=P_Y=(1/2,1/2).
\]

The companion Wolfram witness evaluates

\[
I(X;Y)
=
\frac{\log(27/16)}{\log16}
\approx
0.1887218755408671
\]

bits.

For deterministic equality \(Y=X\) with a fair bit,

\[
I(X;Y)=1
\]

bit.

The statistic is symmetric in \(X\) and \(Y\).

The result therefore cannot by itself tell us whether \(X\) causes
\(Y\), \(Y\) causes \(X\), or both share another cause.

\section{Same entropy, different
structure}\label{same-entropy-different-structure}

Two distributions can have the same entropy while assigning probability
to different events.

Entropy therefore does not uniquely identify a distribution.

Likewise, two learned representations can have similar aggregate
information measures while differing dramatically in:

\begin{itemize}
\tightlist
\item
  geometry;
\item
  localization;
\item
  robustness;
\item
  causal role;
\item
  downstream accessibility.
\end{itemize}

Information-theoretic summaries are powerful.

They are not complete descriptions.

\section{Concentration}\label{concentration}

Expected behavior is not enough.

We often need to know how far a random quantity can deviate from its
mean.

Concentration inequalities answer questions of the form:

\[
P(|X-\mathbb E X|\ge t)
\le
\text{tail bound}.
\]

The assumptions matter.

For example, standard Hoeffding inequalities rely on bounded independent
variables.

High-dimensional probability develops many concentration tools with
different assumptions and regimes (Vershynin 2018).

The Atlas will carry those assumptions forward rather than quoting a
tail formula detached from its conditions.

\section{Why concentration matters for machine
learning}\label{why-concentration-matters-for-machine-learning}

Concentration appears in:

\begin{itemize}
\tightlist
\item
  generalization arguments;
\item
  randomized embeddings;
\item
  sketching;
\item
  stochastic gradients;
\item
  random-feature approximations;
\item
  sampling-based estimators;
\item
  high-dimensional geometry.
\end{itemize}

A concentration bound is an uncertainty statement about a random object.

It is not a guarantee that a learned system behaves well under
distribution shift.

That requires a different argument.

\section{Sufficient statistics}\label{sufficient-statistics}

A statistic

\[
T(X)
\]

compresses observed data.

It is sufficient for a parameter \(\theta\) in a statistical family
when, informally, it retains all information in the sample relevant to
\(\theta\).

Sufficiency is therefore relative to a model.

A statistic can be sufficient for one family and inadequate for another.

This is a useful early example of a theme that later reappears in
boundary compression:

\begin{quote}
sufficient for what?
\end{quote}

\section{Exponential families}\label{exponential-families}

An exponential family has form

\[
p_\eta(x)
=
h(x)
\exp\{
\eta^\top T(x)-A(\eta)
\}.
\]

Here:

\begin{itemize}
\tightlist
\item
  \(\eta\) is a natural parameter;
\item
  \(T(x)\) is a sufficient-statistic vector in the standard regular
  setting;
\item
  \(A(\eta)\) is the log-partition function.
\end{itemize}

Exponential-family structure connects:

\begin{itemize}
\tightlist
\item
  statistics;
\item
  convexity;
\item
  entropy;
\item
  duality.
\end{itemize}

Wainwright and Jordan develop this framework systematically (Wainwright
and Jordan 2008).

\section{Bernoulli as an exponential
family}\label{bernoulli-as-an-exponential-family}

For \(x\in\{0,1\}\),

\[
p(x)
=
p^x(1-p)^{1-x}.
\]

Define

\[
\eta
=
\log\frac{p}{1-p}.
\]

Then

\[
p(x)
=
\exp\{
x\eta-A(\eta)
\},
\]

where

\[
A(\eta)
=
\log(1+e^\eta).
\]

Thus

\[
T(x)=x
\]

is the natural sufficient statistic in this representation.

This small example is useful because it connects probability, sufficient
statistics, and convex log-partition structure in one line.

\section{Information geometry waits for
later}\label{information-geometry-waits-for-later}

Exponential families also support a geometric viewpoint.

The Hessian of the log-partition function is related to covariance and
Fisher information under standard regularity conditions.

That leads toward information geometry.

This chapter does not develop that theory.

It establishes the probabilistic and statistical objects the later
chapter will require.

\section{Three overinterpretations to
avoid}\label{three-overinterpretations-to-avoid}

\subsection{Entropy equals meaning}\label{entropy-equals-meaning}

False.

Entropy quantifies uncertainty under a probability model.

\subsection{Mutual information equals
causality}\label{mutual-information-equals-causality}

False.

Mutual information is symmetric dependence information.

\subsection{KL equals distance}\label{kl-equals-distance}

False.

KL is asymmetric and can be infinite.

These are not pedantic distinctions.

They determine which downstream arguments are legal.

\section{Atlas connections}\label{atlas-connections-2}

\textbf{Representation.}\\
Information measures can characterize statistical dependence or
compression, but representation geometry and identifiability remain
separate questions.

\textbf{Attention.}\\
Softmax creates normalized weights, but probabilistic-looking
normalization does not automatically make every attention object a
probability model with the semantics assumed here.

\textbf{Decision theory.}\\
Expected utility and uncertainty require explicit distributions or
uncertainty sets.

\textbf{Memory.}\\
Retrieval and compression can be studied in terms of sufficient
statistics and information preservation, but semantic adequacy is
task-relative.

\textbf{Optimization.}\\
KL and related divergences can induce optimization geometry.

\section{Closing view}\label{closing-view-5}

Probability gives the Atlas a language for uncertainty.

Information theory gives it a language for statistical dependence and
coding structure.

Exponential families give it a language for structured statistical
models.

Concentration gives it a language for deviations.

These tools are powerful precisely because their meanings are narrow.

The Atlas will use them most effectively by refusing to ask them to mean
more than they do.

\section{References used in this
chapter}\label{references-used-in-this-chapter-7}

\begin{itemize}
\tightlist
\item
  (Cover and Thomas 2006)
\item
  (Wainwright and Jordan 2008)
\item
  (Vershynin 2018)
\end{itemize}

\chapter{Geometry of Constrained State
Spaces}\label{geometry-of-constrained-state-spaces}

\textbf{Epistemic status:} mathematical exposition based on standard
Riemannian and matrix-manifold geometry, with Atlas-owned elementary
derivations.\\
\textbf{Primary figure:} `ATLAS-FIG-MANIFOLD-001`\\
\textbf{Derivation packet:}
`mathematics/derivations/ATLAS-CH-GEOM-001-DERIVATIONS.md`

\section{The space you write in is not always the space you may move
in}\label{the-space-you-write-in-is-not-always-the-space-you-may-move-in}

Machine learning often presents objects in ambient Euclidean
coordinates. A unit vector lives in \(\mathbb R^d\). An orthonormal
frame is stored as an array in \(\mathbb R^{n\times p}\). A subspace may
be represented by many different matrices.

Coordinates are convenient. They can also hide constraints.

If a state must satisfy

\[
\|x\|=1,
\]

then an arbitrary Euclidean step \(x\mapsto x-\eta g\) will generally
leave the admissible set. If a matrix must satisfy

\[
X^\top X=I,
\]

then an arbitrary matrix update will generally destroy orthogonality.

The geometric question is therefore:

\begin{quote}
What are the legal local directions, and how do we turn them into legal
finite moves?
\end{quote}

The useful allegory is a \textbf{cartographer walking on a mountain}.

The map coordinates are the ambient coordinates. The mountain surface is
the admissible state space. A tangent plane gives a locally linear set
of legal directions. A geodesic stays intrinsically on the surface. A
retraction takes a locally legal direction and returns a nearby ambient
step to the constraint.

The limit of the allegory is important. A manifold need not be a visible
surface in three-dimensional space, and a retraction need not be an
orthogonal projection. The image teaches constrained motion, not all of
differential geometry.

\section{Ambient space and admissible
space}\label{ambient-space-and-admissible-space}

A smooth embedded manifold \(M\subseteq\mathbb R^N\) is locally well
approximated by a linear space even though \(M\) need not itself be
linear (Lee 2018).

At a point \(x\in M\), the tangent space

\[
T_xM
\]

contains the velocities of differentiable curves through \(x\).

If

\[
\gamma:(-\varepsilon,\varepsilon)\to M,
\qquad
\gamma(0)=x,
\]

then

\[
\dot\gamma(0)\in T_xM.
\]

This gives a recurring Atlas pattern:

\[
\boxed{
\text{constrained object}
\longrightarrow
\text{tangent linearization}
\longrightarrow
\text{local update}
\longrightarrow
\text{return to the constraint}.
}
\]

The tangent space is linear. The manifold is not.

That distinction will later let us separate ``what direction should the
optimizer choose?'' from ``how should the chosen direction be
represented as a legal state update?''

\section{The unit sphere}\label{the-unit-sphere}

Consider

\[
S^{d-1}
=
\{x\in\mathbb R^d:x^\top x=1\}.
\]

Let \(\gamma(t)\in S^{d-1}\) with

\[
\gamma(0)=x,
\qquad
\dot\gamma(0)=v.
\]

Because the curve remains on the sphere,

\[
\gamma(t)^\top\gamma(t)=1.
\]

Differentiate at \(t=0\):

\[
2x^\top v=0.
\]

Hence

\[
\boxed{
T_xS^{d-1}
=
\{v\in\mathbb R^d:x^\top v=0\}.
}
\]

The tangent space is the hyperplane orthogonal to \(x\).

If \(g\in\mathbb R^d\) is an arbitrary ambient direction, its tangent
component is

\[
g_{\rm tan}
=
g-(x^\top g)x.
\]

The geometry has already changed the update rule. The radial component
is not a legal first-order direction.

\section{Retractions: step locally, return
legally}\label{retractions-step-locally-return-legally}

A retraction \(R_x:T_xM\to M\) is a smooth local map satisfying, at
minimum,

\[
R_x(0)=x
\]

and

\[
DR_x(0)=\operatorname{id}_{T_xM}.
\]

For the sphere, a simple retraction is normalization:

\[
\boxed{
R_x(v)
=
\frac{x+v}{\|x+v\|},
\qquad
v\in T_xS^{d-1}.
}
\]

Why is this a valid first-order retraction?

Set

\[
r(t)
=
\frac{x+tv}{\|x+tv\|}.
\]

Because \(x^\top v=0\) and \(\|x\|=1\),

\[
r'(0)=v.
\]

So the map agrees with the tangent direction to first order while
returning every sufficiently small step to the sphere.

This is a common numerical idea: do the difficult reasoning in a linear
tangent space, then use a controlled map back to the nonlinear
constraint.

Matrix-manifold optimization develops this pattern systematically (Absil
et al. 2008).

\section{Retraction is not exponential
map}\label{retraction-is-not-exponential-map}

The Riemannian exponential map moves along a geodesic.

On the unit sphere, for nonzero \(v\in T_xS^{d-1}\),

\[
\operatorname{Exp}_x(v)
=
\cos(\|v\|)x
+
\sin(\|v\|)
\frac{v}{\|v\|}.
\]

The normalized retraction is

\[
R_x(v)
=
\frac{x+v}{\sqrt{1+\|v\|^2}},
\]

because \(x^\top v=0\).

Both are legal maps back to the sphere. They are not the same map.

For small \(v\),

\[
\operatorname{Exp}_x(v)
=
x+v-\frac{\|v\|^2}{2}x+O(\|v\|^3),
\]

and

\[
R_x(v)
=
x+v-\frac{\|v\|^2}{2}x+O(\|v\|^3).
\]

They agree through the first-order retraction requirement and share the
same quadratic radial correction, but differ at higher order.

This distinction matters later. ``Geometry-aware'' should not become a
vague synonym for ``normalize after every step.''

\section{The primary geometric
plate}\label{the-primary-geometric-plate}

The Atlas figure fixes

\[
x=(0,0,1),
\qquad
v=(0.8,0,0).
\]

Since

\[
x^\top v=0,
\]

\(v\) is tangent to the sphere at \(x\).

The figure shows:

\begin{itemize}
\tightlist
\item
  the unit sphere;
\item
  the tangent plane at \(x\);
\item
  the tangent vector \(v\);
\item
  the spherical exponential-map endpoint;
\item
  the normalized-retraction endpoint;
\item
  the great-circle direction.
\end{itemize}

\begin{figure}
\centering
\pandocbounded{\includegraphics[keepaspectratio,alt={Unit sphere with the tangent plane at x, a tangent vector v, and separate exponential-map and normalized-retraction endpoints.}]{figures/derivatives/ATLAS-FIG-MANIFOLD-001-v0.1.1.png}}
\caption{Unit sphere with the tangent plane at x, a tangent vector v,
and separate exponential-map and normalized-retraction endpoints.}
\end{figure}

The plate is partly schematic. Perspective, tangent-plane extent, and
label placement are pedagogical. The plotted points themselves are
generated from the stated equations.

Its purpose is to make one fact visible:

\begin{quote}
a tangent vector belongs to a linear space, but a finite update must
still be interpreted through the manifold.
\end{quote}

\section{Geodesics and spherical
distance}\label{geodesics-and-spherical-distance}

On a Riemannian manifold, a geodesic generalizes the idea of a locally
straight path.

For unit vectors \(x,y\in S^{d-1}\), define their angle

\[
\theta
=
\arccos(x^\top y),
\qquad
0\le\theta\le\pi.
\]

For non-antipodal endpoints, the shorter great-circle distance is

\[
d_{S}(x,y)=\theta.
\]

The ambient chord distance is

\[
\|x-y\|
=
2\sin\frac{\theta}{2}.
\]

For small \(\theta\),

\[
2\sin\frac{\theta}{2}
=
\theta+O(\theta^3),
\]

so local Euclidean distance approximates geodesic distance.

Globally they are different.

This is a simple example of a recurring principle: local linearization
can be excellent without making the global geometry flat.

\section{SLERP: interpolate on the
sphere}\label{slerp-interpolate-on-the-sphere}

Let

\[
x^\top y=\cos\theta,
\qquad
0<\theta<\pi.
\]

Spherical linear interpolation is

\[
\boxed{
\operatorname{SLERP}(x,y;t)
=
\frac{\sin((1-t)\theta)}{\sin\theta}x
+
\frac{\sin(t\theta)}{\sin\theta}y,
\qquad
0\le t\le1.
}
\]

Shoemake introduced SLERP in quaternion interpolation for computer
animation (Shoemake 1985). The underlying formula is simply great-circle
interpolation on a sphere.

Writing

\[
\alpha
=
\frac{\sin((1-t)\theta)}{\sin\theta},
\qquad
\beta
=
\frac{\sin(t\theta)}{\sin\theta},
\]

we obtain

\[
\|\alpha x+\beta y\|^2
=
\alpha^2+\beta^2+2\alpha\beta\cos\theta
=
1.
\]

So the path remains on the unit sphere exactly.

There are two important boundaries:

\begin{itemize}
\tightlist
\item
  as \(\theta\to0\), the apparent singularity is removable;
\item
  at \(\theta=\pi\), the shortest geodesic is not unique, so an
  additional directional choice is required.
\end{itemize}

The antipodal singularity is geometric information, not merely a
numerical nuisance.

\section{Stiefel geometry: orthonormal
frames}\label{stiefel-geometry-orthonormal-frames}

For \(X\in\mathbb R^{n\times p}\), the Stiefel manifold is

\[
\operatorname{St}(n,p)
=
\{X:X^\top X=I_p\}.
\]

Its points are ordered orthonormal \(p\)-frames.

Let \(X(t)\in\operatorname{St}(n,p)\) and set

\[
X(0)=X,
\qquad
\dot X(0)=Z.
\]

Differentiate

\[
X(t)^\top X(t)=I_p.
\]

At \(t=0\),

\[
Z^\top X+X^\top Z=0.
\]

Therefore

\[
\boxed{
T_X\operatorname{St}(n,p)
=
\{Z:X^\top Z+Z^\top X=0\}.
}
\]

This is the matrix analogue of the sphere condition \(x^\top v=0\).

The Stiefel manifold is especially important in optimization with
orthogonality constraints; Edelman, Arias, and Smith provide a canonical
geometric treatment (Edelman et al. 1998).

\section{Grassmann geometry: the frame is not the
subspace}\label{grassmann-geometry-the-frame-is-not-the-subspace}

A Stiefel matrix \(X\) stores an ordered orthonormal basis.

But if the object of interest is only the \(p\)-dimensional subspace
spanned by its columns, then

\[
X
\]

and

\[
XQ,
\qquad
Q\in O(p),
\]

represent the same subspace.

The Grassmann manifold identifies all such frames:

\[
\boxed{
\operatorname{Gr}(n,p)
\cong
\operatorname{St}(n,p)/O(p).
}
\]

This is the first important quotient-geometry example in the Atlas.

A parameterization can contain distinctions that the represented object
does not.

Later, when we discuss semantic equivalence and benign nonconvexity,
this pattern returns in a broader form:

\begin{quote}
if two parameter states represent the same functional object, should the
optimization landscape distinguish them?
\end{quote}

The Grassmannian gives a clean finite-dimensional case where the answer
is no.

\section{Geometry changes what an optimizer is allowed to
do}\label{geometry-changes-what-an-optimizer-is-allowed-to-do}

Suppose \(f:M\to\mathbb R\).

In ambient coordinates, automatic differentiation may give an ordinary
gradient \(\nabla f\). A manifold-aware method must convert that ambient
information into a tangent direction and then map a finite step back to
\(M\).

Schematically,

\[
\nabla f
\longrightarrow
\operatorname{Proj}_{T_xM}(\nabla f)
\longrightarrow
v
\longrightarrow
R_x(v).
\]

The exact Riemannian gradient depends on the chosen metric, and later
chapters will treat that carefully.

The present chapter establishes the more basic point:

\begin{quote}
constraints are part of the state geometry, so they should appear in the
update mathematics rather than only as after-the-fact repairs.
\end{quote}

\section{Four mistakes to avoid}\label{four-mistakes-to-avoid}

\subsection{Mistake 1: ``The tangent space is the
manifold.''}\label{mistake-1-the-tangent-space-is-the-manifold.}

No.~It is a local linearization at one point.

\subsection{Mistake 2: ``Normalization is the exponential
map.''}\label{mistake-2-normalization-is-the-exponential-map.}

No.~On the sphere it is a useful retraction, not the same finite map as
geodesic motion.

\subsection{Mistake 3: ``A Stiefel point and a Grassmann point are the
same
object.''}\label{mistake-3-a-stiefel-point-and-a-grassmann-point-are-the-same-object.}

No.~The first is an ordered frame; the second is a subspace modulo basis
change.

\subsection{Mistake 4: ``Local Euclidean behavior means the space is
globally
Euclidean.''}\label{mistake-4-local-euclidean-behavior-means-the-space-is-globally-euclidean.}

No.~Local tangent approximations coexist with global curvature and
topological structure.

\section{Atlas connections}\label{atlas-connections-3}

\textbf{Normalized representations.}\\
Unit-norm states move naturally on spheres rather than unconstrained
Euclidean space.

\textbf{Optimization on manifolds.}\\
Tangent gradients, retractions, and matrix manifolds turn the present
geometry into algorithms.

\textbf{Position geometry.}\\
Rotations and spherical constructions recur in positional encodings.

\textbf{Representation as transport.}\\
A representation update can be treated as motion through a constrained
state space.

\textbf{Quotient geometry.}\\
Grassmann geometry supplies the simplest durable example of removing
representational redundancy.

The recurring Atlas shift is:

\[
\boxed{
\text{parameters as coordinates}
\longrightarrow
\text{parameters as points in a structured state space}.
}
\]

\section{Closing view}\label{closing-view-6}

The mountain allegory is useful only until the mathematics becomes
visible.

A constrained representation has an admissible state space. At each
point, tangent geometry tells us which first-order motions are legal.
Geodesics describe intrinsic motion. Retractions provide practical
finite updates. Quotient geometry tells us when different coordinates
describe the same underlying object.

The map is not the territory.

The ambient coordinates are not the admissible state space.

\section{References used in this
chapter}\label{references-used-in-this-chapter-8}

\begin{itemize}
\tightlist
\item
  (Lee 2018)
\item
  (Absil et al. 2008)
\item
  (Edelman et al. 1998)
\item
  (Shoemake 1985)
\end{itemize}

See `sources/source-locks/ATLAS-CH-GEOM-001.yaml` for exact source
identities and claim scope.

\chapter{Flows, Stability, and
Bifurcation}\label{flows-stability-and-bifurcation}

\textbf{Epistemic status:} Established Theory + Atlas Synthesis + Exact
Computational Witness\\
\textbf{Specification:} manuscript/specifications/ATLAS-CH-DYN-001.md\\
\textbf{Derivation packet:}
mathematics/derivations/ATLAS-CH-DYN-001-DERIVATIONS.md\\
\textbf{Computational witness:}
mathematics/computational-witnesses/ATLAS-CW-DYN-001.md\\
\textbf{Source lock:} sources/source-locks/ATLAS-CH-DYN-001.yaml

\section{A state is not yet a
dynamics}\label{a-state-is-not-yet-a-dynamics}

A vector tells us where a system is.

A vector field tells us how the state wants to move.

A trajectory tells us one realized history.

A flow tells us exact transport through time.

These are different mathematical objects.

The distinction is simple enough to overlook and important enough to
organize much of the Atlas.

Later chapters will discuss:

\begin{itemize}
\tightlist
\item
  residual depth;
\item
  optimizer state;
\item
  adaptive computation;
\item
  recurrent systems;
\item
  reinforcement learning;
\item
  numerical solvers.
\end{itemize}

All of them can be described badly if ``state,'' ``update,'' and
``dynamics'' are treated as synonyms.

\section{Law of motion versus film
frames}\label{law-of-motion-versus-film-frames}

A useful allegory is a film.

The vector field is the local law telling each point how to move.

A trajectory is one actor's path through the scene.

A collection of sampled states is a sequence of film frames.

The flow map is the exact rule carrying every admissible initial state
forward by a chosen amount of time.

A numerical integrator is different again.

It manufactures approximate frames according to its own discrete update
rule.

The allegory fails if taken too literally.

Many neural layers do not arise from one autonomous continuous-time law.

The point is to distinguish the objects, not to force every architecture
into an ODE.

\section{Autonomous ordinary differential
equations}\label{autonomous-ordinary-differential-equations}

Consider

\[
\dot x
=
f(x),
\]

with initial condition

\[
x(0)=x_0.
\]

The function

\[
f:\mathcal X\to T\mathcal X
\]

is the vector field, written here in Euclidean notation.

At a state \(x\), the vector

\[
f(x)
\]

specifies the instantaneous direction of motion.

The differential equation asks for a curve

\[
t\mapsto x(t)
\]

whose derivative agrees with that field.

\section{Trajectory}\label{trajectory}

A trajectory is one solution beginning at one initial state:

\[
x(t;x_0).
\]

Changing \(x_0\) can change the entire future path.

A trajectory is therefore not the vector field.

It is one realization generated by the field.

This matters in machine learning because a training run is a trajectory.

The optimizer update law is part of the dynamics that generates it.

One run does not identify the whole dynamical system.

\section{Flow map}\label{flow-map}

Where existence and uniqueness hold, define

\[
\Phi_t(x_0)
=
x(t;x_0).
\]

The map

\[
\Phi_t
\]

acts on initial states.

For an autonomous true flow,

\[
\boxed{
\Phi_{t+s}
=
\Phi_t\circ\Phi_s.
}
\]

The identity holds only for states and times on which the relevant
compositions are defined. A forward-only system has a semiflow, with
\(s,t\ge 0\), rather than a globally defined two-sided flow.

This composition law is stronger than merely having a depth-indexed
family of transformations.

The word \textbf{flow} should therefore not be used automatically for
any sequence of neural layers.

\section{Exact flow of a linear
system}\label{exact-flow-of-a-linear-system}

For a constant linear system

\[
\dot x=Ax,
\]

the exact flow is

\[
\Phi_t(x_0)
=
e^{tA}x_0.
\]

The matrix exponential therefore converts a generator \(A\) into
finite-time transport.

This is one reason linear algebra and dynamics are tightly coupled.

Eigenvalues describe local modal growth for the linear system.

But finite-time behavior can still depend on non-normal structure, as
the Non-normality chapter emphasized.

\section{Equilibrium}\label{equilibrium}

An equilibrium \(x_\star\) satisfies

\[
f(x_\star)=0.
\]

If the system starts there, it remains there.

Equilibria are fixed points of the flow:

\[
\Phi_t(x_\star)
=
x_\star.
\]

The important next question is not only whether an equilibrium exists.

It is what nearby trajectories do.

\section{Stability is a neighborhood
question}\label{stability-is-a-neighborhood-question}

An equilibrium is Lyapunov stable if sufficiently small initial
perturbations remain small for all future time.

Informally:

\begin{quote}
start close enough, remain close.
\end{quote}

Asymptotic stability adds convergence:

\begin{quote}
start close enough, remain close, and eventually return.
\end{quote}

Exponential stability strengthens the convergence rate:

\[
\|x(t)-x_\star\|
\le
C e^{-\alpha t}
\|x(0)-x_\star\|
\]

for suitable positive constants \(C,\alpha\) in a declared neighborhood.

These are not interchangeable labels.

\section{Linearization}\label{linearization}

Suppose \(f\) is differentiable and

\[
f(x_\star)=0.
\]

Write

\[
x=x_\star+\delta.
\]

Then

\[
f(x_\star+\delta)
=
J_f(x_\star)\delta
+
o(\|\delta\|).
\]

So the local perturbation dynamics are

\[
\dot\delta
=
J_f(x_\star)\delta
+
o(\|\delta\|).
\]

The little-\(o\) remainder is the general differentiability guarantee. A
quadratic \(O(\|\delta\|^2)\) bound needs additional regularity, such as
a locally Lipschitz Jacobian.

The Jacobian provides the first local approximation to the nonlinear
system.

\section{Hyperbolic local
stability}\label{hyperbolic-local-stability}

For a (C\^{}1) vector field near a hyperbolic equilibrium, the Jacobian
spectrum gives strong local information.

If every eigenvalue of

\[
J_f(x_\star)
\]

has strictly negative real part, the equilibrium is locally
exponentially stable.

If at least one eigenvalue has strictly positive real part, the
equilibrium is unstable.

The difficult boundary is the imaginary axis.

If an eigenvalue has zero real part, linearization alone may be
inconclusive.

Nonlinear terms can decide the outcome.

\section{Zero eigenvalue is not ``neutral
stability''}\label{zero-eigenvalue-is-not-neutral-stability}

This mistake is common.

Suppose

\[
\dot x
=
-x^3.
\]

At the origin,

\[
f'(0)=0.
\]

The linearized system says nothing.

Yet the nonlinear system attracts nearby states to zero.

Now consider

\[
\dot x
=
x^3.
\]

Again,

\[
f'(0)=0.
\]

But now the origin is unstable.

The same zero eigenvalue supports opposite nonlinear behavior.

This is why nonhyperbolic points require more than first-order spectral
reasoning.

\section{Lyapunov functions}\label{lyapunov-functions}

A Lyapunov function is a scalar witness for stability.

For an equilibrium at the origin, choose

\[
V(x)
\]

such that

\[
V(0)=0
\]

and

\[
V(x)>0
\]

for nearby nonzero \(x\).

Its derivative along trajectories is

\[
\dot V(x)
=
\nabla V(x)^\top f(x).
\]

If \(V\) decreases along trajectories, the system cannot wander
arbitrarily outward while remaining in the certified neighborhood.

Khalil gives the standard nonlinear stability framework used here
(Khalil 2002).

\section{Negative semidefinite versus negative
definite}\label{negative-semidefinite-versus-negative-definite}

The sign of

\[
\dot V
\]

matters.

If

\[
\dot V\le0,
\]

a standard Lyapunov theorem can establish stability under appropriate
local hypotheses.

If

\[
\dot V<0
\]

for every nearby nonzero state, the same framework can establish
asymptotic stability.

Negative semidefinite derivative alone does not automatically prove
convergence to the equilibrium.

Additional arguments, such as invariant-set reasoning, may be required.

The Atlas will keep that distinction visible.

\section{Exact scalar stable flow}\label{exact-scalar-stable-flow}

Consider

\[
\dot x=-2x,
\qquad
x(0)=3.
\]

The exact solution is

\[
\boxed{
x(t)=3e^{-2t}.
}
\]

Choose

\[
V(x)
=
\frac12x^2.
\]

Then

\[
\dot V
=
x\dot x
=
-2x^2.
\]

Therefore

\[
\dot V<0
\]

for every nonzero state.

The origin is exponentially stable.

\begin{figure}
\centering
\pandocbounded{\includegraphics[keepaspectratio,alt={Three exact dynamics panels: exponential attraction under x'=-2x, stable and unstable saddle-node branches x*=±sqrt(mu), and a circular harmonic-oscillator energy orbit.}]{figures/masters/ATLAS-FIG-DYN-001.png}}
\caption{Three exact dynamics panels: exponential attraction under
x'=-2x, stable and unstable saddle-node branches x*=±sqrt(mu), and a
circular harmonic-oscillator energy orbit.}
\end{figure}

\section{Phase portrait}\label{phase-portrait}

A phase portrait shows trajectories or direction fields in state space.

It is a compressed picture of qualitative behavior.

It can reveal:

\begin{itemize}
\tightlist
\item
  attraction;
\item
  repulsion;
\item
  separatrices;
\item
  periodic orbits;
\item
  invariant regions.
\end{itemize}

But a phase portrait is still a representation.

Its appearance depends on:

\begin{itemize}
\tightlist
\item
  coordinates;
\item
  projection;
\item
  dimensional reduction;
\item
  parameter choice.
\end{itemize}

A two-dimensional picture of a high-dimensional system can hide
essential structure.

\section{Parameters can change qualitative
behavior}\label{parameters-can-change-qualitative-behavior}

Consider a family

\[
\dot x
=
f(x;\mu).
\]

Changing the parameter \(\mu\) may deform trajectories continuously.

Sometimes the qualitative structure changes.

Equilibria can:

\begin{itemize}
\tightlist
\item
  appear;
\item
  disappear;
\item
  change stability;
\item
  split;
\item
  give rise to cycles.
\end{itemize}

A bifurcation is a qualitative change of this kind.

Kuznetsov provides the formal bifurcation framework behind the normal
forms used here (Kuznetsov 2004).

\section{Saddle-node bifurcation}\label{saddle-node-bifurcation}

Take

\[
\dot x
=
\mu-x^2.
\]

Equilibria satisfy

\[
x^2=\mu.
\]

For

\[
\mu>0,
\]

there are two equilibria:

\[
x_\pm
=
\pm\sqrt{\mu}.
\]

For

\[
\mu<0,
\]

there are no real equilibria.

At

\[
\mu=0,
\]

the two branches meet.

This is the canonical saddle-node picture.

\section{Saddle-node stability}\label{saddle-node-stability}

Differentiate:

\[
f'(x)
=
-2x.
\]

At the positive branch,

\[
f'(+\sqrt{\mu})
=
-2\sqrt{\mu}<0.
\]

So the positive branch is locally stable.

At the negative branch,

\[
f'(-\sqrt{\mu})
=
2\sqrt{\mu}>0.
\]

So the negative branch is unstable.

At the collision point,

\[
f'(0)=0.
\]

The equilibrium is nonhyperbolic exactly where the qualitative structure
changes.

\section{Pitchfork bifurcation}\label{pitchfork-bifurcation}

Now consider

\[
\dot x
=
\mu x-x^3.
\]

Equilibria satisfy

\[
x(\mu-x^2)=0.
\]

The central branch

\[
x=0
\]

exists for all \(\mu\).

For

\[
\mu>0,
\]

two additional branches appear:

\[
x_\pm
=
\pm\sqrt{\mu}.
\]

This is the standard supercritical pitchfork normal form.

\section{Pitchfork stability
exchange}\label{pitchfork-stability-exchange}

The derivative is

\[
f'(x)
=
\mu-3x^2.
\]

At the origin:

\[
f'(0)=\mu.
\]

Thus the origin is locally stable for

\[
\mu<0
\]

and unstable for

\[
\mu>0.
\]

At the nonzero branches:

\[
f'(\pm\sqrt{\mu})
=
-2\mu<0.
\]

Thus both new branches are locally stable.

A qualitative reorganization has occurred.

\section{Why bifurcation language matters for
learning}\label{why-bifurcation-language-matters-for-learning}

Training systems contain parameters that change over time:

\begin{itemize}
\tightlist
\item
  learning rate;
\item
  momentum;
\item
  weight decay;
\item
  normalization scales;
\item
  loss weights;
\item
  routing temperature;
\item
  data mixture;
\item
  model state itself.
\end{itemize}

Qualitative changes in observed dynamics can resemble bifurcation
phenomena.

The analogy can be useful.

It must not be overclaimed.

Observing a phase transition in training metrics does not by itself
establish a mathematically identified bifurcation normal form.

That requires an actual reduced dynamical model.

\section{Hopf bifurcation: the minimal
picture}\label{hopf-bifurcation-the-minimal-picture}

The supercritical Hopf normal form can be written in polar coordinates
as

\[
\dot r
=
\mu r-r^3,
\]

\[
\dot\theta
=
\omega.
\]

For

\[
\mu<0,
\]

the origin is radially stable.

For

\[
\mu>0,
\]

the origin becomes unstable and a stable nonzero radius appears:

\[
r_\star
=
\sqrt{\mu}.
\]

The angular equation turns that radius into a periodic orbit.

This is the qualitative picture.

The chapter does not attempt the full Hopf theorem.

\section{Periodic orbit is not
convergence}\label{periodic-orbit-is-not-convergence}

A stable periodic orbit attracts nearby trajectories to a cycle.

It does not make the state converge to a fixed point.

This distinction becomes useful later when optimization dynamics
oscillate.

A training system can settle into:

\begin{itemize}
\tightlist
\item
  a fixed point;
\item
  a limit cycle;
\item
  a slowly drifting attractor;
\item
  more complicated recurrent behavior.
\end{itemize}

``Stable'' therefore needs an object:

stable \textbf{what}?

\section{Conservative dynamics}\label{conservative-dynamics}

Not every useful system is dissipative.

Some dynamics preserve quantities.

Hamiltonian mechanics provides the canonical example.

Let

\[
z=
\begin{pmatrix}
q\\
p
\end{pmatrix}.
\]

Define

\[
J=
\begin{pmatrix}
0&I\\
-I&0
\end{pmatrix}.
\]

For Hamiltonian \(H(z)\),

\[
\boxed{
\dot z
=
J\nabla H(z).
}
\]

The matrix \(J\) is skew-symmetric:

\[
J^\top=-J.
\]

\section{Hamiltonian conservation}\label{hamiltonian-conservation}

Along the exact autonomous Hamiltonian flow,

\[
\frac{dH}{dt}
=
\nabla H^\top\dot z.
\]

Substitute

\[
\dot z=J\nabla H.
\]

Then

\[
\frac{dH}{dt}
=
\nabla H^\top J\nabla H.
\]

For any real vector \(v\),

\[
v^\top Jv=0
\]

when \(J\) is skew-symmetric.

Therefore

\[
\boxed{
\frac{dH}{dt}=0.
}
\]

The Hamiltonian is conserved along the exact flow.

\section{Harmonic oscillator}\label{harmonic-oscillator}

Take

\[
H(q,p)
=
\frac12(q^2+p^2).
\]

Then

\[
\dot q=p,
\]

and

\[
\dot p=-q.
\]

In matrix form:

\[
\dot z
=
Az,
\]

with

\[
A=
\begin{pmatrix}
0&1\\
-1&0
\end{pmatrix}.
\]

Since

\[
A^2=-I,
\]

the exact flow is

\[
\boxed{
M(t)
=
e^{tA}
=
\begin{pmatrix}
\cos t&\sin t\\
-\sin t&\cos t
\end{pmatrix}.
}
\]

The state rotates.

It does not decay.

\section{Symplectic structure}\label{symplectic-structure}

A map \(M\) is symplectic in canonical coordinates when

\[
M^\top J M=J.
\]

For the exact harmonic-oscillator flow,

\[
\boxed{
M(t)^\top J M(t)=J.
}
\]

The computational witness verifies this identity symbolically.

Symplecticity means the canonical two-form is preserved.

In two dimensions, it also implies oriented phase-area preservation.

\section{Energy conservation and symplecticity are not
synonyms}\label{energy-conservation-and-symplecticity-are-not-synonyms}

The harmonic oscillator exact flow preserves:

\begin{itemize}
\tightlist
\item
  the symplectic form;
\item
  the Hamiltonian energy.
\end{itemize}

But these are different properties.

A numerical method can be symplectic without preserving the original
Hamiltonian exactly at every step.

Instead, symplectic integrators often exhibit favorable long-time
behavior associated with a nearby modified Hamiltonian.

That belongs in the Numerics chapter.

The distinction matters because ``structure-preserving'' has more than
one meaning.

\section{Dissipative versus conservative
pictures}\label{dissipative-versus-conservative-pictures}

The stable scalar flow satisfies

\[
x(t)\to0.
\]

It loses the scalar Lyapunov energy

\[
V(x)=x^2/2.
\]

The harmonic oscillator does not converge.

It moves around a constant-energy orbit.

These are qualitatively different kinds of dynamics.

Later optimization chapters will encounter mixtures:

\begin{itemize}
\tightlist
\item
  damping;
\item
  momentum;
\item
  oscillation;
\item
  dissipation;
\item
  near-conservative transport.
\end{itemize}

One scalar notion of ``stability'' cannot describe all of them.

\section{Exact flow versus numerical
step}\label{exact-flow-versus-numerical-step}

Suppose the exact flow over time \(h\) is

\[
\Phi_h.
\]

A numerical method defines another map:

\[
\Psi_h.
\]

A method of order \(p\) may satisfy locally

\[
\Psi_h(x)
=
\Phi_h(x)
+
O(h^{p+1}).
\]

That does not make the maps identical.

The numerical method is itself a discrete dynamical system.

It can possess behavior absent from the continuous system.

\section{Discretization can change
stability}\label{discretization-can-change-stability}

Consider

\[
\dot x=-\lambda x,
\qquad
\lambda>0.
\]

The exact flow is stable for every positive \(\lambda\).

Explicit Euler gives

\[
x_{n+1}
=
(1-h\lambda)x_n.
\]

This discrete update is stable only when

\[
|1-h\lambda|<1.
\]

So the continuous system can be stable while the discretization
explodes.

This is the bridge to numerical stability.

The next Numerics chapter will make it systematic.

\section{Discretization can change
geometry}\label{discretization-can-change-geometry}

A generic numerical step may fail to preserve:

\begin{itemize}
\tightlist
\item
  energy;
\item
  symplectic form;
\item
  norm;
\item
  invariant sets;
\item
  positivity;
\item
  constraints.
\end{itemize}

Sometimes that failure is acceptable.

Sometimes it destroys the property the model was built to exploit.

This motivates structure-preserving methods.

Hairer, Lubich, and Wanner provide the standard geometric-numerical
framework used in the later chapter (Hairer et al. 2006).

\section{Depth as computational
time}\label{depth-as-computational-time}

Residual architectures often take the form

\[
x_{k+1}
=
x_k
+
h f_k(x_k).
\]

If

\[
f_k=f
\]

and the interpretation is consistent, this resembles an explicit Euler
step for

\[
\dot x=f(x).
\]

This makes depth interpretable as a form of computational time.

But the caveat is load-bearing.

Generic neural layers can vary with depth.

They can contain:

\begin{itemize}
\tightlist
\item
  normalization;
\item
  attention;
\item
  routing;
\item
  stochasticity;
\item
  discontinuous choices;
\item
  data-dependent operators.
\end{itemize}

So:

\[
\boxed{
\text{residual update resemblance}
\neq
\text{proof of one underlying autonomous ODE}.
}
\]

\section{Flow-like architectures}\label{flow-like-architectures}

The weaker claim is often the useful one.

An architecture can be studied with dynamical tools if:

\begin{itemize}
\tightlist
\item
  repeated updates accumulate state;
\item
  local Jacobians matter;
\item
  stability matters;
\item
  long-horizon transport matters;
\item
  step size or depth controls computation.
\end{itemize}

This does not require literal identification with classical mechanics.

It requires only that the dynamical viewpoint makes testable predictions
or design constraints.

\section{Optimizer dynamics}\label{optimizer-dynamics}

An optimizer can have state of its own.

Momentum methods introduce variables such as

\[
v_t.
\]

Adaptive optimizers introduce running statistics.

The combined training state can therefore be written schematically as

\[
s_t
=
(\theta_t,m_t,v_t,\ldots).
\]

The update acts on this larger state.

The Optimizer-State Dynamics keystone already uses that principle.

The present chapter supplies the continuous-time language needed to
interpret it.

\section{Reinforcement-learning
dynamics}\label{reinforcement-learning-dynamics}

Reinforcement learning adds further state:

\begin{itemize}
\tightlist
\item
  environment state;
\item
  belief state;
\item
  policy state;
\item
  value estimates;
\item
  replay buffers;
\item
  target networks.
\end{itemize}

Learning itself becomes an evolving coupled system.

The RL foundations chapter will build on the same distinctions:

\begin{itemize}
\tightlist
\item
  state;
\item
  flow/update;
\item
  stability;
\item
  fixed point.
\end{itemize}

\section{Bifurcation as diagnostic
hypothesis}\label{bifurcation-as-diagnostic-hypothesis}

Suppose a training process suddenly changes regime when a control
parameter crosses a threshold.

Possible explanations include:

\begin{itemize}
\tightlist
\item
  true local bifurcation;
\item
  numerical instability;
\item
  data-mixture shift;
\item
  optimizer-state transition;
\item
  representation collapse;
\item
  finite-sample noise.
\end{itemize}

Bifurcation language generates hypotheses.

It does not adjudicate them automatically.

The correct next step is to identify a reduced state and test the
implied local dynamics.

\section{Lyapunov functions as
certificates}\label{lyapunov-functions-as-certificates}

The same logic applies to Lyapunov ideas.

A training loss that decreases is not automatically a Lyapunov function
for the full training state.

If the optimizer has hidden state, the scalar loss may rise temporarily
while the full system remains stable.

A genuine Lyapunov certificate must be defined on the actual state whose
stability is being claimed.

This becomes important in governed adaptation.

\section{Four failure modes}\label{four-failure-modes}

\subsection{Local linearization treated as global
truth}\label{local-linearization-treated-as-global-truth}

The Jacobian describes first-order local behavior.

\subsection{Stable equilibrium confused with stable
orbit}\label{stable-equilibrium-confused-with-stable-orbit}

Fixed points and cycles are different invariant objects.

\subsection{Hamiltonian structure asserted from oscillation
alone}\label{hamiltonian-structure-asserted-from-oscillation-alone}

Oscillation does not imply a canonical symplectic/Hamiltonian model.

\subsection{Numerical method identified with continuous
flow}\label{numerical-method-identified-with-continuous-flow}

An integrator is a separate discrete system.

These errors all arise from collapsing distinct mathematical objects.

\section{Atlas connections}\label{atlas-connections-4}

\textbf{Numerics.}\\
The next chapter studies how discrete methods approximate or distort
exact flows.

\textbf{Architecture history.}\\
Residual, recurrent, and continuous-depth models can be organized by how
they accumulate transformations.

\textbf{Latent computational time.}\\
Depth can become an adaptive integration horizon when the analogy is
justified.

\textbf{Optimization.}\\
Momentum, damping, oscillation, and optimizer state are dynamical
phenomena.

\textbf{Reinforcement learning.}\\
Bellman and policy updates define iterative dynamics on value/policy
objects.

\textbf{SPINDLE.}\\
Operator splitting and adaptive depth require a clean distinction
between flow generator, split update, and integrator.

\section{What would falsify a dynamical
interpretation?}\label{what-would-falsify-a-dynamical-interpretation}

A proposed dynamical model should make predictions.

Examples:

\begin{itemize}
\tightlist
\item
  perturbation decay rates;
\item
  stability boundaries;
\item
  conserved or slowly varying quantities;
\item
  phase relationships;
\item
  bifurcation thresholds;
\item
  step-size sensitivity.
\end{itemize}

If those predictions fail systematically, the proposed state or flow
model is wrong or incomplete.

The dynamical viewpoint earns its keep by constraining what should
happen.

\section{Closing view}\label{closing-view-7}

The mathematical transition in this chapter is from static
transformation to organized evolution.

We began with:

\[
\dot x=f(x).
\]

From that one object emerged:

\begin{itemize}
\tightlist
\item
  trajectories;
\item
  exact flows;
\item
  equilibria;
\item
  stability;
\item
  Lyapunov certificates;
\item
  bifurcations;
\item
  periodic orbits;
\item
  Hamiltonian conservation;
\item
  symplectic structure.
\end{itemize}

The most important lesson is not any one theorem.

It is object discipline.

The vector field is the law.

The trajectory is one history.

The flow is exact transport.

The integrator is a different discrete map.

Once those are separated, the later Atlas can ask a sharper question:

\begin{quote}
When machine computation unfolds through repeated updates, which
properties belong to the intended dynamics, and which were created by
the way we chose to discretize it?
\end{quote}

That is where Numerical Intelligence begins.

\section{References used in this
chapter}\label{references-used-in-this-chapter-9}

\begin{itemize}
\tightlist
\item
  (Khalil 2002)
\item
  (Kuznetsov 2004)
\item
  (Hairer et al. 2006)
\end{itemize}

See `sources/source-locks/ATLAS-CH-DYN-001.yaml` for exact source roles
and claim boundaries.

\chapter{Discretization, Stability, and
Splitting}\label{discretization-stability-and-splitting}

\textbf{Epistemic status:} Established Theory + Atlas Synthesis + Exact
Computational Witness\\
\textbf{Specification:}
manuscript/specifications/ATLAS-CH-NUMERICS-001.md\\
\textbf{Derivation packet:}
mathematics/derivations/ATLAS-CH-NUMERICS-001-DERIVATIONS.md\\
\textbf{Computational witness:}
mathematics/computational-witnesses/ATLAS-CW-NUMERICS-001.md\\
\textbf{Source lock:} sources/source-locks/ATLAS-CH-NUMERICS-001.yaml

\section{The continuous system is not the
algorithm}\label{the-continuous-system-is-not-the-algorithm}

The previous chapter separated the exact flow

\[
\Phi_h
\]

from a numerical update

\[
\Psi_h.
\]

That distinction is the starting point of numerical analysis.

A differential equation can be perfectly stable while a poor
discretization explodes.

A numerical method can be stable while remaining inaccurate.

A high-order method can fail if used outside the regime for which its
assumptions hold.

The numerical question is therefore not:

\begin{quote}
Does this update look like the differential equation?
\end{quote}

It is:

\begin{quote}
In what precise sense does the discrete map approximate the exact flow,
and what happens to errors as steps accumulate?
\end{quote}

\section{Map versus terrain}\label{map-versus-terrain}

A useful allegory is navigation.

The exact flow is the true route through terrain.

A numerical method is a stepping rule based on a map.

\textbf{Consistency} says one small step points approximately where the
true route points.

\textbf{Stability} says small errors in the stepping process are not
catastrophically amplified.

\textbf{Convergence} says the full discrete route approaches the true
route as step size shrinks.

The analogy has a limit.

Numerical error is not geographic error.

The actual theory requires norms, regularity assumptions, and
method-specific bounds.

\section{One-step methods}\label{one-step-methods}

Consider

\[
\dot x
=
f(t,x).
\]

A one-step method has the form

\[
\boxed{
x_{n+1}
=
\Psi_h(t_n,x_n).
}
\]

The step size is

\[
h=t_{n+1}-t_n.
\]

The exact flow over one step is

\[
\Phi_h(t_n,x).
\]

The method is an approximation to this exact transport.

\section{Local defect}\label{local-defect}

Start one numerical step from the exact state.

Define

\[
\delta_{n+1}
=
\Phi_h(t_n,x(t_n))
-
\Psi_h(t_n,x(t_n)).
\]

This is the \textbf{local defect}.

It asks:

\begin{quote}
If previous error were magically removed, how wrong would this one step
be?
\end{quote}

For an order-\(p\) method,

\[
\delta_{n+1}
=
O(h^{p+1})
\]

under the required smoothness assumptions.

\section{Global error}\label{global-error}

The actual numerical state is not restarted from truth.

It starts from the previous approximation.

Define

\[
e_n
=
x(t_n)-x_n.
\]

This is the \textbf{global error}.

Global error combines:

\begin{itemize}
\tightlist
\item
  new local defect;
\item
  propagation of old error.
\end{itemize}

That second term is where stability enters.

\section{Consistency is not enough}\label{consistency-is-not-enough}

A method can have a small one-step defect and still behave badly when
errors are repeatedly amplified.

So:

\[
\boxed{
\text{consistency}
\not\Rightarrow
\text{convergence}
}
\]

without the relevant stability and regularity assumptions.

For a well-behaved one-step method over a finite time interval, a local
defect

\[
O(h^{p+1})
\]

combined with suitable Lipschitz error propagation yields global error

\[
O(h^p).
\]

Iserles develops this numerical ODE viewpoint systematically (Iserles
2008).

\section{Explicit Euler}\label{explicit-euler}

Taylor-expand the exact solution:

\[
x(t_n+h)
=
x(t_n)
+
h x'(t_n)
+
\frac{h^2}{2}x''(t_n)
+
O(h^3).
\]

Using

\[
x'(t_n)=f(t_n,x(t_n)),
\]

drop the higher-order terms:

\[
\boxed{
x_{n+1}
=
x_n
+
h f(t_n,x_n).
}
\]

This is explicit Euler.

\section{Explicit Euler local
defect}\label{explicit-euler-local-defect}

If the step begins from the exact state, then

\[
\delta_{n+1}^{\rm EE}
=
\frac{h^2}{2}x''(t_n)
+
O(h^3).
\]

So the local defect is second order.

The accumulated global error is first order:

\[
e_n
=
O(h)
\]

under standard finite-time regularity and stability conditions.

The distinction between local \(h^2\) and global \(h\) is fundamental.

\section{Implicit Euler}\label{implicit-euler}

Implicit Euler evaluates the vector field at the future state:

\[
\boxed{
x_{n+1}
=
x_n
+
h f(t_{n+1},x_{n+1}).
}
\]

The future state appears on both sides.

So a step generally requires solving an equation.

That extra computational cost buys a different stability geometry.

\section{Explicit versus implicit is not fast versus
slow}\label{explicit-versus-implicit-is-not-fast-versus-slow}

``Explicit'' means the new state can be computed directly from known
quantities.

``Implicit'' means the new state is defined by an equation that must be
solved.

The distinction is algebraic.

It does not automatically determine:

\begin{itemize}
\tightlist
\item
  wall-clock cost;
\item
  accuracy;
\item
  parallelism;
\item
  implementation quality.
\end{itemize}

A cheap implicit solve can outperform a tiny explicit step.

A difficult nonlinear implicit solve can dominate the computation.

\section{The scalar test equation}\label{the-scalar-test-equation}

To study numerical stability, use

\[
y'
=
\lambda y.
\]

Its exact one-step factor is

\[
e^{h\lambda}.
\]

Define

\[
z=h\lambda.
\]

A one-step numerical method gives

\[
y_{n+1}
=
R(z)y_n,
\]

where \(R\) is the \textbf{stability function}.

The method's behavior on this scalar problem reveals important
properties of its discrete dynamics.

\section{Explicit Euler stability
function}\label{explicit-euler-stability-function}

Explicit Euler gives

\[
y_{n+1}
=
y_n+h\lambda y_n.
\]

Therefore

\[
\boxed{
R_{\rm EE}(z)=1+z.
}
\]

Use the standard non-growth absolute-stability set

\[
\mathcal S
=
\{z:|R(z)|\le1\}.
\]

For explicit Euler,

\[
|1+z|\le1
\]

is the absolute-stability disk.

For a decaying exact mode with

\[
\operatorname{Re}(\lambda)<0,
\]

strict asymptotic decay of the discrete mode requires the interior
condition

\[
|1+z|<1.
\]

\section{The explicit stability
disk}\label{the-explicit-stability-disk}

The non-growth condition

\[
|1+z|\le1
\]

defines the closed disk of radius \(1\) centered at

\[
-1
\]

in the complex plane.

Its open interior

\[
|1+z|<1
\]

gives strict asymptotic decay.

A stable continuous eigenmode can therefore become unstable numerically
if

\[
h\lambda
\]

falls outside this disk.

The differential equation did not become unstable.

The discretization did.

\begin{figure}
\centering
\pandocbounded{\includegraphics[keepaspectratio,alt={Explicit and implicit Euler absolute-stability regions beside log-log Lie-Trotter and Strang local-defect scaling for the exact split witness.}]{figures/masters/ATLAS-FIG-NUMERICS-001.png}}
\caption{Explicit and implicit Euler absolute-stability regions beside
log-log Lie-Trotter and Strang local-defect scaling for the exact split
witness.}
\end{figure}

\section{Implicit Euler stability
function}\label{implicit-euler-stability-function}

Implicit Euler gives

\[
y_{n+1}
=
y_n+h\lambda y_{n+1}.
\]

Rearrange:

\[
(1-z)y_{n+1}
=
y_n.
\]

Thus

\[
\boxed{
R_{\rm IE}(z)
=
\frac{1}{1-z}.
}
\]

The standard non-growth absolute-stability condition is

\[
\left|
\frac{1}{1-z}
\right|
\le1.
\]

Equivalently,

\[
|1-z|\ge1.
\]

Strict decay uses the corresponding strict inequalities.

\section{A-stability}\label{a-stability}

Every point in the closed left half-plane satisfies

\[
|1-z|\ge1.
\]

Therefore implicit Euler is A-stable in the standard non-growth sense.

For every point in the open left half-plane,

\[
|1-z|>1,
\]

so an exponentially decaying scalar mode remains strictly decaying for
any positive step size.

That is a strong stability property.

It is not a claim of arbitrary accuracy.

\section{L-stability}\label{l-stability}

For implicit Euler,

\[
R_{\rm IE}(z)
=
\frac{1}{1-z}.
\]

As

\[
|z|\to\infty
\]

within the left half-plane,

\[
R_{\rm IE}(z)\to0.
\]

Thus strongly damped continuous modes are also strongly damped
numerically.

Implicit Euler is L-stable.

This matters in stiff problems, where very fast decaying modes can
otherwise pollute a computation.

\section{Stability is not accuracy}\label{stability-is-not-accuracy}

Take

\[
y'=-100y.
\]

Use

\[
h=0.05.
\]

Then

\[
z=-5.
\]

The exact one-step factor is

\[
e^{-5}
\approx
0.00673795.
\]

Explicit Euler gives

\[
R_{\rm EE}(-5)
=
-4.
\]

The numerical mode grows in magnitude.

Implicit Euler gives

\[
R_{\rm IE}(-5)
=
\frac16
\approx
0.1667.
\]

It decays.

But it is not close to the exact factor.

So:

\[
\boxed{
\text{stable}
\neq
\text{accurate}.
}
\]

\section{The stiff scalar witness}\label{the-stiff-scalar-witness}

The same continuous mode can therefore produce three qualitatively
different one-step behaviors:

\[
\text{exact: }
0.0067,
\]

\[
\text{explicit Euler: }
-4,
\]

\[
\text{implicit Euler: }
0.1667.
\]

Explicit Euler is unstable.

Implicit Euler is stable.

Neither numerical factor equals the exact one.

This is why numerical analysis needs both stability and error analysis.

\section{What stiffness means}\label{what-stiffness-means}

Stiffness is not simply ``large derivative.''

A useful operational picture is:

\begin{quote}
the system contains rapidly decaying modes that force a chosen numerical
method to use a much smaller stable step than the slower behavior of
interest would otherwise require.
\end{quote}

The word is therefore partly problem-relative and partly
method-relative.

Hairer and Wanner develop the stiff ODE framework in detail (Hairer and
Wanner 1996).

\section{Multiple timescales}\label{multiple-timescales}

Suppose

\[
\dot x
=
\begin{pmatrix}
-1&0\\
0&-1000
\end{pmatrix}
x.
\]

The first mode evolves on timescale roughly \(1\).

The second decays on timescale roughly \(10^{-3}\).

After the fast mode has essentially vanished, an explicit method may
still be forced to use tiny steps because of its stability restriction.

The numerical method is spending effort respecting a mode that no longer
matters much to the observable slow trajectory.

This is the classic stiffness tension.

\section{Implicit methods move the stability
boundary}\label{implicit-methods-move-the-stability-boundary}

An implicit method can tolerate much larger stable steps for decaying
modes.

But the price can include:

\begin{itemize}
\tightlist
\item
  linear solves;
\item
  nonlinear solves;
\item
  Jacobians;
\item
  preconditioning;
\item
  memory;
\item
  implementation complexity.
\end{itemize}

Numerical method choice is therefore a systems tradeoff.

This becomes relevant when neural architectures contain implicit or
equilibrium-style blocks.

\section{Absolute stability is not nonlinear
stability}\label{absolute-stability-is-not-nonlinear-stability}

The region

\[
|R(z)|<1
\]

is defined for the scalar linear test equation.

It provides powerful local/mode-level information.

It is not a universal theorem about a nonlinear problem.

A nonlinear method can leave the test region satisfied while
encountering:

\begin{itemize}
\tightlist
\item
  changing Jacobians;
\item
  constraints;
\item
  bifurcations;
\item
  non-normal amplification;
\item
  state-dependent stiffness.
\end{itemize}

The stability region is a diagnostic object, not a complete nonlinear
certificate.

\section{From one vector field to
two}\label{from-one-vector-field-to-two}

Suppose the dynamics decompose:

\[
\dot x
=
(A+B)x.
\]

Perhaps:

\begin{itemize}
\tightlist
\item
  \(A\) is easy to integrate;
\item
  \(B\) is easy to integrate;
\item
  \(A+B\) is expensive;
\item
  each subflow preserves a useful structure.
\end{itemize}

This motivates splitting.

Instead of approximating the full vector field in one step, we compose
simpler exact or approximate subflows.

\section{Lie--Trotter splitting}\label{lietrotter-splitting}

For constant matrices, define

\[
\boxed{
S_{\rm LT}(h)
=
e^{hA}e^{hB}.
}
\]

This means the \(B\) subflow acts first on a state vector, followed by
the \(A\) subflow, under the usual right-to-left composition convention.

The result approximates

\[
e^{h(A+B)}.
\]

How accurate is it?

The answer depends on commutation.

\section{Noncommutativity is the
obstruction}\label{noncommutativity-is-the-obstruction}

Expand both subflows:

\[
e^{hA}
=
I+hA+\frac{h^2}{2}A^2+O(h^3),
\]

\[
e^{hB}
=
I+hB+\frac{h^2}{2}B^2+O(h^3).
\]

Their product has second-order term

\[
\frac12A^2+AB+\frac12B^2.
\]

The exact flow has second-order term

\[
\frac12
(A^2+AB+BA+B^2).
\]

Subtracting gives

\[
\boxed{
S_{\rm LT}(h)
-
e^{h(A+B)}
=
\frac{h^2}{2}[A,B]
+
O(h^3),
}
\]

where

\[
[A,B]
=
AB-BA.
\]

\section{Commuting subflows}\label{commuting-subflows}

If

\[
[A,B]=0,
\]

then for constant matrices

\[
e^{hA}e^{hB}
=
e^{h(A+B)}
\]

exactly.

The split introduces no error.

This is an important lesson for architecture.

Splitting quality is not determined only by the size of the components.

It depends on their algebraic interaction.

\section{Exact noncommuting
witness}\label{exact-noncommuting-witness}

Choose

\[
A=
\begin{pmatrix}
0&1\\
0&0
\end{pmatrix},
\]

and

\[
B=
\begin{pmatrix}
0&0\\
1&0
\end{pmatrix}.
\]

Then

\[
[A,B]
=
\begin{pmatrix}
1&0\\
0&-1
\end{pmatrix}.
\]

The operators do not commute.

The exact symbolic series therefore exposes a nonzero second-order Lie
defect.

\section{Lie local defect}\label{lie-local-defect}

For the declared witness,

\[
S_{\rm LT}(h)
-
e^{h(A+B)}
\]

has leading term

\[
\frac{h^2}{2}
\begin{pmatrix}
1&0\\
0&-1
\end{pmatrix}.
\]

The off-diagonal terms begin at order

\[
h^3.
\]

The local defect is therefore

\[
O(h^2).
\]

Under the usual finite-time stability assumptions, Lie--Trotter is
globally first order.

\section{Strang splitting}\label{strang-splitting}

A symmetric composition is

\[
\boxed{
S_{\rm S}(h)
=
e^{hA/2}
e^{hB}
e^{hA/2}.
}
\]

The half-step symmetry cancels the second-order local defect.

For sufficiently regular bounded operators,

\[
S_{\rm S}(h)
-
e^{h(A+B)}
=
O(h^3).
\]

This yields second-order global accuracy under the usual stability
assumptions.

\section{Exact Strang witness}\label{exact-strang-witness}

For the same \(A,B\), the symbolic defect begins as

\[
\begin{pmatrix}
O(h^4)&h^3/12\\
-h^3/6&O(h^4)
\end{pmatrix}.
\]

The first nonzero terms are cubic.

The figure's log--log plot evaluates exact matrix exponentials across a
declared \(h\)-grid.

The asymptotic slopes reveal the expected local orders:

\begin{itemize}
\tightlist
\item
  Lie--Trotter: \(2\);
\item
  Strang: \(3\).
\end{itemize}

\section{Symmetry is doing mathematical
work}\label{symmetry-is-doing-mathematical-work}

The advantage of Strang is not merely that it performs more substeps.

The composition is time-symmetric in its outer \(A\) factors.

That symmetry removes an entire order of local defect.

Architectural ordering can therefore change approximation order even
when the same primitive operations are used.

This is the core mathematical motivation for later split-operator neural
architectures.

\section{Splitting is not arbitrary
modularity}\label{splitting-is-not-arbitrary-modularity}

Suppose a model has two modules.

Calling their composition a ``split method'' does not make numerical
splitting theory applicable.

A genuine splitting interpretation requires a declared underlying
evolution such as

\[
\dot x=(A+B)x
\]

or its nonlinear/operator analogue.

The substeps should correspond to meaningful component flows or
approximations.

Otherwise ``splitting'' is metaphor, not derivation.

\section{Splitting can preserve
structure}\label{splitting-can-preserve-structure}

Suppose each subflow is symplectic.

The composition of symplectic maps is symplectic.

Thus a split integrator can preserve the symplectic form exactly even
when it only approximates the full exact flow.

This is one way discrete algorithms can preserve qualitative geometry.

McLachlan and Quispel survey this splitting viewpoint (McLachlan and
Quispel 2002).

\section{Structure-preserving means name the
structure}\label{structure-preserving-means-name-the-structure}

A method may preserve:

\begin{itemize}
\tightlist
\item
  symplectic form;
\item
  norm;
\item
  positivity;
\item
  mass;
\item
  reversibility;
\item
  constraint manifold;
\item
  volume.
\end{itemize}

These are different promises.

A method can preserve one while violating another.

So the phrase

\begin{quote}
structure-preserving
\end{quote}

is incomplete unless the structure is named.

Hairer, Lubich, and Wanner develop this geometric integration
perspective (Hairer et al. 2006).

\section{Backward error viewpoint}\label{backward-error-viewpoint}

A powerful later idea is to ask whether a discrete method exactly
follows a nearby modified differential equation.

This is backward error analysis.

Instead of saying

\begin{quote}
the numerical path is slightly wrong,
\end{quote}

we ask

\begin{quote}
for which nearby dynamics is this path more nearly exact?
\end{quote}

This can explain excellent long-time behavior of geometric methods.

The present chapter only introduces the viewpoint.

A later Numerical Intelligence chapter can develop it more deeply.

\section{Step size is a modeling
parameter}\label{step-size-is-a-modeling-parameter}

The step size

\[
h
\]

controls more than numerical precision.

It can alter:

\begin{itemize}
\tightlist
\item
  stability;
\item
  damping;
\item
  phase error;
\item
  effective depth;
\item
  computational cost.
\end{itemize}

In learned systems, scale parameters that resemble step sizes can
therefore influence qualitative computation.

But an analogy becomes useful only when the update admits a coherent
numerical interpretation.

\section{Depth as number of steps}\label{depth-as-number-of-steps}

If a residual block approximates

\[
x_{n+1}
=
x_n
+
h f(x_n),
\]

then increasing depth while reducing \(h\) can resemble time refinement.

The natural questions become:

\begin{itemize}
\tightlist
\item
  Does the output converge as depth is refined?
\item
  Is the method stable?
\item
  Which invariant changes with depth?
\item
  Is the architecture approximating one flow or changing the vector
  field itself?
\end{itemize}

The later Adaptive Depth chapter will inherit this language.

\section{Adaptive depth as error
control}\label{adaptive-depth-as-error-control}

Classical adaptive integration allocates more steps where error
estimates demand them.

A learned analogue might use:

\begin{itemize}
\tightlist
\item
  residual magnitude;
\item
  local defect proxy;
\item
  uncertainty;
\item
  contract violation;
\item
  commutator magnitude.
\end{itemize}

to decide whether more computation is needed.

That is an architectural research programme.

This chapter supplies the numerical grammar.

It does not validate a specific controller.

\section{Commutator as split
diagnostic}\label{commutator-as-split-diagnostic}

For Lie--Trotter,

\[
[A,B]
\]

appears directly in the leading local defect.

This suggests a diagnostic principle:

\begin{quote}
if two learned operators are intended to be split, measure how strongly
their order matters.
\end{quote}

In nonlinear systems, Lie brackets or lifted commutators can play
analogous roles.

This is one bridge to SPLICE-style diagnostics.

Again, the exact numerical theorem applies only after the operator model
has been declared.

\section{Stiff learned dynamics}\label{stiff-learned-dynamics}

A learned update can exhibit widely separated effective timescales.

Examples might include:

\begin{itemize}
\tightlist
\item
  fast normalization variables;
\item
  slow feature evolution;
\item
  rapidly decaying auxiliary states;
\item
  momentum transients.
\end{itemize}

Calling such a system ``stiff'' requires an actual dynamical and
numerical analysis.

The useful test is whether stability forces a step/depth scale much
smaller than the behavior of interest requires.

\section{Implicit computation}\label{implicit-computation}

Implicit Euler solves for a future state satisfying an equation.

This creates a conceptual link to:

\begin{itemize}
\tightlist
\item
  equilibrium networks;
\item
  implicit layers;
\item
  fixed-point solvers;
\item
  proximal updates.
\end{itemize}

But the relationship is not identity.

An implicit neural block can solve a different equation with different
stability and convergence properties.

Numerical language should clarify the contract, not erase differences.

\section{Four numerical failure
modes}\label{four-numerical-failure-modes}

\subsection{Small local defect, unstable
propagation}\label{small-local-defect-unstable-propagation}

Consistency without stability.

\subsection{Stable but inaccurate coarse
steps}\label{stable-but-inaccurate-coarse-steps}

The stiff implicit-Euler witness.

\subsection{High-order formula outside
assumptions}\label{high-order-formula-outside-assumptions}

Formal order does not protect against nonsmoothness, singularity, or
unbounded-operator domain failures.

\subsection{Structure-preserving label without
structure}\label{structure-preserving-label-without-structure}

A claim is meaningless until the preserved invariant or geometry is
stated.

These are distinct failures.

\section{What would falsify a numerical
interpretation?}\label{what-would-falsify-a-numerical-interpretation}

Suppose a neural architecture is claimed to approximate a particular
flow.

Useful falsification tests include:

\begin{itemize}
\tightlist
\item
  depth refinement fails to converge;
\item
  halving the apparent step size does not reduce error at the predicted
  order;
\item
  stability boundaries disagree with the proposed test equation or
  Jacobian model;
\item
  split ordering has no relation to predicted commutator effects;
\item
  claimed invariants drift despite a supposed structure-preserving
  construction.
\end{itemize}

A numerical interpretation should generate quantitative predictions.

\section{Atlas connections}\label{atlas-connections-5}

\textbf{Dynamics.}\\
The exact flow is the object being approximated.

\textbf{Adaptive Depth.}\\
Depth can become an error-controlled step budget when the numerical
model is justified.

\textbf{Split-operator architectures.}\\
Commutators and composition order become design variables.

\textbf{Variational optimization.}\\
Discrete updates can be derived from geometric or variational
principles.

\textbf{Numerical Intelligence.}\\
Solvers, residuals, error estimates, and backward error become
computational mechanisms rather than background implementation details.

\textbf{SPINDLE.}\\
The controlled-operator-splitting programme sits directly on the
distinctions developed here.

\section{Closing view}\label{closing-view-8}

Numerical analysis asks a disciplined question:

\begin{quote}
What discrete computation faithfully represents the behavior we intend?
\end{quote}

The answer is not contained in one error order.

A method must negotiate:

\begin{itemize}
\tightlist
\item
  local approximation;
\item
  error propagation;
\item
  stability;
\item
  stiffness;
\item
  computational cost;
\item
  algebraic interaction;
\item
  geometric structure.
\end{itemize}

The scalar test equation showed why a stable continuous system can
become numerically unstable.

The stiff witness showed why stable does not mean accurate.

The split witness showed why operator order matters.

And Strang splitting showed that composition symmetry can improve
accuracy without changing the primitive operators.

These are not implementation details.

They are mathematical properties of computation.

The next Atlas chapters can now ask a sharper architectural question:

\begin{quote}
If depth, modules, and residual updates are treated as numerical
machinery, which integration contract are they actually satisfying?
\end{quote}

\section{References used in this
chapter}\label{references-used-in-this-chapter-10}

\begin{itemize}
\tightlist
\item
  (Iserles 2008)
\item
  (Hairer and Wanner 1996)
\item
  (McLachlan and Quispel 2002)
\item
  (Hairer et al. 2006)
\end{itemize}

See `sources/source-locks/ATLAS-CH-NUMERICS-001.yaml` for exact source
roles and claim boundaries.

\chapter{Local-to-Global Mathematics}\label{local-to-global-mathematics}

\textbf{Epistemic status:} standard sheaf/local-to-global mathematics +
audited Objects prerequisite + Atlas synthesis + exact finite witness.\\
\textbf{Specification:}
manuscript/specifications/ATLAS-CH-LOCALGLOBAL-001.md\\
\textbf{Formal packet:}
mathematics/derivations/ATLAS-CH-LOCALGLOBAL-001-DERIVATIONS.md\\
\textbf{Computational witness:}
mathematics/computational-witnesses/ATLAS-CW-LOCALGLOBAL-001.md\\
\textbf{Source lock:} sources/source-locks/ATLAS-CH-LOCALGLOBAL-001.yaml

Many systems are easier to understand locally than globally.

A sensor sees one region.

A module exposes one interface.

A worker holds one piece of state.

A chart describes one patch.

A subsystem enforces one constraint.

The difficult question is not whether each piece makes sense by itself.

It is:

\begin{quote}
do the pieces agree where they meet, and if they do, is there one
coherent global object that realizes them?
\end{quote}

That is the local-to-global problem.

This chapter introduces only the amount of graph and sheaf language
needed to make that question precise.

\section{Local validity is not global
coherence}\label{local-validity-is-not-global-coherence}

Suppose two teams each provide a locally valid configuration.

One says the shared boundary value is 2.

The other says it is 3.

Nothing is wrong inside either local object.

The failure appears only when the objects are compared on their common
interface.

This is the central structural lesson:

\begin{quote}
consistency lives on overlaps.
\end{quote}

\section{Start with incidence}\label{start-with-incidence}

Before assigning data, we need to know which pieces touch.

A graph:

\texttt{G=(V,E)}

can record pairwise incidence.

A cover:

\texttt{U=union\_i\ U\_i}

can record which local regions together describe a larger region.

These structures answer:

\begin{itemize}
\tightlist
\item
  what is local?
\item
  which pieces overlap?
\item
  which interfaces exist?
\end{itemize}

They do not yet say what data live there.

\section{A graph is not a sheaf}\label{a-graph-is-not-a-sheaf}

A graph can tell us that vertex A is adjacent to vertex B.

It cannot, by itself, tell us:

\begin{itemize}
\tightlist
\item
  what values A and B carry;
\item
  what part of those values is visible at the edge;
\item
  what equality or transformation the edge requires.
\end{itemize}

A sheaf-like object adds data spaces and maps over the incidence
structure.

This is why graph topology and carried algebra must remain distinct.

\section{Local sections}\label{local-sections}

For a region \texttt{U}, let:

\texttt{F(U)}

denote the allowable data on that region.

An element:

\texttt{s\ in\ F(U)}

is called a local section.

The word \textbf{section} is useful because it emphasizes that the data
belong to a particular region.

The same numerical vector can be a section on one region and not even be
well-typed on another.

\section{Restriction maps}\label{restriction-maps}

If:

\texttt{V\ subseteq\ U},

a restriction map:

\texttt{rho\^{}U\_V:F(U)-\textgreater{}F(V)}

asks:

\begin{quote}
what does the data on U say when viewed only on V?
\end{quote}

Restriction must be coherent.

Restricting to the same region does nothing:

\texttt{rho\^{}U\_U=id}.

Restricting in stages must agree with restricting directly:

\texttt{rho\^{}V\_W\ o\ rho\^{}U\_V\ =\ rho\^{}U\_W}

for:

\texttt{W\ subseteq\ V\ subseteq\ U}.

\section{Why restriction is an interface
map}\label{why-restriction-is-an-interface-map}

OBJECTS-001 introduced interfaces as the obligations composition may
assume across a boundary.

Restriction makes that idea mathematical.

Two larger local objects do not need to be identical.

They need to induce compatible data on what they share.

The shared region is the comparison interface.

\section{Matching families}\label{matching-families}

Let:

\texttt{\{U\_i\}}

cover U.

Choose one local section:

\texttt{s\_i\ in\ F(U\_i)}

for each region.

The family is matching when every pair agrees after restriction to its
overlap:

\texttt{rho\^{}\{U\_i\}\_\{U\_i\ intersect\ U\_j\}(s\_i)\ =\ rho\^{}\{U\_j\}\_\{U\_i\ intersect\ U\_j\}(s\_j)}.

This is exact compatibility.

\section{Compatibility is typed}\label{compatibility-is-typed}

The two local sections may live in different spaces:

\texttt{F(U\_i)}

and:

\texttt{F(U\_j)}.

They become comparable only after both are mapped into:

\texttt{F(U\_i\ intersect\ U\_j)}.

This prevents a common category error:

\begin{quote}
compare whole local states when only boundary data are supposed to
agree.
\end{quote}

\section{A presheaf gives the restriction
machinery}\label{a-presheaf-gives-the-restriction-machinery}

A presheaf supplies:

\begin{itemize}
\tightlist
\item
  local data spaces;
\item
  restriction maps;
\item
  identity and composition laws.
\end{itemize}

That is already useful.

It gives a disciplined way to talk about local views.

But it does not yet guarantee that compatible local views come from one
global object.

\section{A sheaf adds unique
gluing}\label{a-sheaf-adds-unique-gluing}

A sheaf adds the local-to-global property.

For every matching family over a cover, there must exist a unique:

\texttt{s\ in\ F(U)}

whose restriction to each \texttt{U\_i} is \texttt{s\_i}.

The word \textbf{unique} matters.

The condition has two parts:

\begin{enumerate}
\def\labelenumi{\arabic{enumi}.}
\tightlist
\item
  existence;
\item
  uniqueness.
\end{enumerate}

\section{Existence and uniqueness answer different
questions}\label{existence-and-uniqueness-answer-different-questions}

Existence asks:

\begin{quote}
can these local pieces be assembled at all?
\end{quote}

Uniqueness asks:

\begin{quote}
if they can, do they determine one global object or several?
\end{quote}

A system can fail either way.

Keeping these clauses separate is useful far beyond sheaf theory.

\section{The equalizer picture}\label{the-equalizer-picture}

For a finite cover, gather all local restrictions into a map:

\texttt{R:F(U)-\textgreater{}prod\_i\ F(U\_i)}.

Gather all overlap discrepancies into a map:

\texttt{Delta:prod\_i\ F(U\_i)-\textgreater{}prod\_\{i\textless{}j\}F(U\_i\ intersect\ U\_j)}.

Matching families satisfy:

\texttt{Delta=0}.

The sheaf condition says, informally:

\begin{quote}
the local families that actually come from global sections are exactly
the matching families, and the global source is unique.
\end{quote}

For vector-space-valued examples:

\texttt{im(R)=ker(Delta)}

with R injective.

\section{Why this is a useful computational
form}\label{why-this-is-a-useful-computational-form}

The equalizer picture turns local-to-global reasoning into explicit
operations:

\begin{itemize}
\tightlist
\item
  restrict;
\item
  subtract on overlaps;
\item
  test a kernel condition;
\item
  reconstruct.
\end{itemize}

That is exactly the amount of sheaf theory needed for many finite
engineering examples.

The deeper theory remains available later if needed.

\section{Exact witness: three
points}\label{exact-witness-three-points}

Let:

\texttt{X=\{a,b,c\}}.

Use the cover:

\texttt{U=\{a,b\}}

and:

\texttt{V=\{b,c\}}.

Their overlap is:

\texttt{W=\{b\}}.

This is the smallest example that has:

\begin{itemize}
\tightlist
\item
  two local regions;
\item
  one genuine overlap;
\item
  one global union.
\end{itemize}

\section{Sections are functions}\label{sections-are-functions}

For every subset \texttt{S} of X, define:

\texttt{F(S)=R\^{}S}.

A section is simply a real-valued function on S.

Restriction means forgetting coordinates outside the smaller subset.

Nothing abstract is hidden.

\section{The global-to-local map}\label{the-global-to-local-map}

Write a global section as:

\texttt{(x,y,z)}.

Its restrictions are:

\texttt{(x,y)}

on U and:

\texttt{(y,z)}

on V.

So:

\texttt{R(x,y,z)=(x,y,y,z)}.

\section{Uniqueness is injectivity}\label{uniqueness-is-injectivity}

Suppose two global sections have the same restrictions to U and V.

They agree at:

\begin{itemize}
\tightlist
\item
  a because U contains a;
\item
  b because both regions contain b;
\item
  c because V contains c.
\end{itemize}

Therefore the global sections are equal.

In matrix language, R has rank 3.

The local restrictions determine at most one global function.

\section{The overlap-discrepancy
map}\label{the-overlap-discrepancy-map}

A general pair of local sections is:

\texttt{((u\_a,u\_b),(v\_b,v\_c))}.

Define:

\texttt{Delta(u\_a,u\_b,v\_b,v\_c)\ =\ u\_b-v\_b}.

The pair is compatible exactly when:

\texttt{Delta=0}.

The discrepancy lives precisely on the overlap coordinate b.

\section{Image equals kernel}\label{image-equals-kernel}

The image of R consists of vectors:

\texttt{(x,y,y,z)}.

The kernel of Delta consists of vectors satisfying:

\texttt{u\_b=v\_b}.

These are the same set.

Therefore:

\texttt{im(R)=ker(Delta)}.

This is the exact gluing statement in the finite witness.

\section{A compatible family}\label{a-compatible-family}

Choose:

\texttt{s\_U=(1,2)}

and:

\texttt{s\_V=(2,4)}.

The overlap values both equal 2.

So:

\texttt{Delta=0}.

The unique global glue is:

\texttt{s=(1,2,4)}.

\section{An incompatible family}\label{an-incompatible-family}

Now choose:

\texttt{q\_U=(1,2)}

and:

\texttt{q\_V=(3,4)}.

The overlap discrepancy is:

\texttt{Delta=2-3=-1}.

It is nonzero.

Therefore the pair is not in the image of R.

No global function can restrict to both local sections.

\section{The failure is localized}\label{the-failure-is-localized}

The value at a is not a problem.

The value at c is not a problem.

The incompatibility is entirely at b.

This is one practical virtue of explicit overlap maps:

\begin{quote}
they tell us where coherence fails.
\end{quote}

\section{Gluing is stronger than
aggregation}\label{gluing-is-stronger-than-aggregation}

Suppose we simply concatenate the two local records:

\texttt{(1,2,3,4)}.

We have stored all the data.

We have not resolved the fact that the two copies of b disagree.

Aggregation collects.

Gluing enforces compatibility.

\section{The graph view of the
witness}\label{the-graph-view-of-the-witness}

Represent U and V as two vertices.

Connect them by an edge because their overlap is nonempty.

That graph records interaction.

But the numerical obstruction \texttt{-1} does not come from the graph
alone.

It comes from:

\begin{itemize}
\tightlist
\item
  data spaces;
\item
  restrictions;
\item
  the two local sections.
\end{itemize}

The carried algebra matters.

\section{Cellular sheaves make graph-local algebra
explicit}\label{cellular-sheaves-make-graph-local-algebra-explicit}

Cellular sheaves assign data spaces to cells and maps along incidences
under a declared convention.

This creates a finite algebraic object over a graph or cell complex.

Curry develops cellular sheaves as a computable combinatorial sheaf
framework, and Hansen--Ghrist develop further graph/cell-complex
structure.

The Atlas needs only the elementary idea here.

\section{We do not need cohomology
yet}\label{we-do-not-need-cohomology-yet}

Sheaf cohomology is a powerful obstruction and structural tool.

It is not required to understand the basic gluing mechanism.

For this chapter, a nonzero overlap discrepancy already provides an
exact finite obstruction.

The more advanced machinery would distract from the foundational idea.

\section{Pairwise compatibility needs the sheaf
context}\label{pairwise-compatibility-needs-the-sheaf-context}

A common slogan says:

\begin{quote}
pairwise agreement does not imply global agreement.
\end{quote}

That slogan is too crude here.

For an actual sheaf over an open cover, agreement on every pairwise
overlap is precisely the matching-family condition, and the sheaf axiom
supplies a unique global glue.

The correct warning is different:

\begin{quote}
pairwise plausibility in an arbitrary constraint system is not enough
unless the restriction structure and gluing property are specified.
\end{quote}

\section{Local-to-global claims are
model-relative}\label{local-to-global-claims-are-model-relative}

A global section is global relative to the chosen sheaf.

Change:

\begin{itemize}
\tightlist
\item
  the regions;
\item
  the data spaces;
\item
  the restriction maps;
\item
  the overlap semantics;
\end{itemize}

and the set of global sections can change.

Therefore a successful glue is evidence about one declared
local-to-global model.

It is not universal coherence.

\section{Exact consistency is brittle by
design}\label{exact-consistency-is-brittle-by-design}

If one local sensor reports 2.000 and another reports 2.001, exact
equality fails.

That is not a defect in the sheaf condition.

It means the mathematical claim being tested is exact.

If the application tolerates small discrepancies, it needs a different
object:

\begin{itemize}
\tightlist
\item
  a metric;
\item
  a tolerance;
\item
  a likelihood;
\item
  a loss;
\item
  an optimization problem.
\end{itemize}

\section{Approximate consistency is a second
layer}\label{approximate-consistency-is-a-second-layer}

Robinson's sensor-integration work makes this distinction operational.

Local measurements can be compared through a sheaf model even when they
are not exactly consistent.

One can then define an approximate-section or fusion problem.

The Atlas keeps the layers separate:

\begin{enumerate}
\def\labelenumi{\arabic{enumi}.}
\tightlist
\item
  exact structural compatibility;
\item
  approximate repair/fusion.
\end{enumerate}

\section{An approximate glue is not a sheaf-theoretic
equality}\label{an-approximate-glue-is-not-a-sheaf-theoretic-equality}

Suppose the overlap values are:

\texttt{2}

and:

\texttt{2.1}.

A least-squares procedure might replace both with:

\texttt{2.05}.

That can be sensible.

But it changes the data.

The original family still did not satisfy exact compatibility.

The optimization result and the sheaf-gluing result are different
claims.

\section{Local truth is not
required}\label{local-truth-is-not-required}

A section can be mathematically valid while factually wrong.

Suppose two sensors both report 100 when the true value is 0.

Their data can glue perfectly.

The sheaf has checked consistency, not truth.

This distinction will matter whenever sheaf language is used for
evidence fusion.

\section{Global consistency is not semantic
adequacy}\label{global-consistency-is-not-semantic-adequacy}

A software system can satisfy all declared interface equalities and
still implement the wrong specification.

A scientific model can integrate all measurements consistently while
omitting a latent variable.

A local-to-global formalism checks the contracts it was given.

It does not prove the contracts are complete.

\section{Local-to-global structure can expose hidden
assumptions}\label{local-to-global-structure-can-expose-hidden-assumptions}

The advantage is not only successful reconstruction.

Writing explicit overlaps forces questions such as:

\begin{itemize}
\tightlist
\item
  what must two modules agree on?
\item
  which values are shared?
\item
  which transformations compare them?
\item
  where can inconsistency appear?
\end{itemize}

That makes assumptions inspectable.

\section{Graphs provide the
skeleton}\label{graphs-provide-the-skeleton}

A graph or hypergraph can organize:

\begin{itemize}
\tightlist
\item
  subsystems;
\item
  agents;
\item
  sensors;
\item
  patches;
\item
  data sources;
\item
  interfaces.
\end{itemize}

It gives the combinatorial skeleton.

A sheaf-like structure then says what algebra or information moves
across that skeleton.

Keeping those layers distinct makes models easier to audit.

\section{Covers provide another
skeleton}\label{covers-provide-another-skeleton}

Sometimes the natural structure is spatial or semantic rather than
graph-first.

A collection:

\texttt{\{U\_i\}}

covers a larger domain.

Overlaps:

\texttt{U\_i\ intersect\ U\_j}

become the places where local descriptions must agree.

This view is natural for:

\begin{itemize}
\tightlist
\item
  coordinate charts;
\item
  distributed observations;
\item
  overlapping datasets;
\item
  decomposed physical domains.
\end{itemize}

\section{Interface data can be
lower-dimensional}\label{interface-data-can-be-lower-dimensional}

Two large subsystems may share only a small boundary.

They need not expose all internal state.

Restriction maps formalize this compression:

\texttt{F(U\_i)-\textgreater{}F(U\_i\ intersect\ U\_j)}.

The local-to-global problem can therefore be much smaller than comparing
whole subsystem states.

This anticipates later boundary-contract and separator ideas without
importing them as prerequisites.

\section{Compatibility can be transformed rather than literal
identity}\label{compatibility-can-be-transformed-rather-than-literal-identity}

In the function witness, both sides restrict to the same scalar and
equality is literal.

More general sheaves can use maps that translate local coordinates into
a common comparison space.

Thus the real principle is not:

\begin{quote}
local arrays must be numerically equal.
\end{quote}

It is:

\begin{quote}
their induced boundary data must agree after the declared restriction
maps.
\end{quote}

\section{Coordinate changes fit
naturally}\label{coordinate-changes-fit-naturally}

Two charts can describe the same underlying object with different
coordinates.

The overlap maps translate each local representation before comparison.

This is one reason local-to-global mathematics appears throughout
geometry.

The Atlas will use the same structural idea later in more computational
settings.

\section{Composition requires more than
adjacency}\label{composition-requires-more-than-adjacency}

Connecting modules with arrows is not enough.

A meaningful composition needs:

\begin{itemize}
\tightlist
\item
  compatible types;
\item
  compatible boundary semantics;
\item
  agreement on shared quantities;
\item
  a rule for assembling the composite.
\end{itemize}

Sheaf gluing provides one mathematically clean archetype for such
assembly.

It is not the only possible compositional formalism.

\section{Obstruction certificates are valuable
outputs}\label{obstruction-certificates-are-valuable-outputs}

When gluing fails, the answer need not be only:

\begin{quote}
no.
\end{quote}

A useful system should expose:

\begin{itemize}
\tightlist
\item
  which overlap failed;
\item
  by how much;
\item
  under which restriction maps;
\item
  which local sections were involved.
\end{itemize}

In the witness, the certificate is:

\texttt{Delta=-1}

at b.

That is already actionable.

\section{A larger graph produces a discrepancy
vector}\label{a-larger-graph-produces-a-discrepancy-vector}

For many overlaps, collect all boundary disagreements into:

\texttt{Delta(s)}.

Then:

\texttt{Delta(s)=0}

means exact compatibility.

Nonzero coordinates localize violated interfaces.

This is a finite linear-algebra pattern that later chapters can reuse.

\section{Zero discrepancy can have several
meanings}\label{zero-discrepancy-can-have-several-meanings}

A zero discrepancy might arise because:

\begin{itemize}
\tightlist
\item
  the data are genuinely consistent;
\item
  the restrictions are too weak to detect a conflict;
\item
  the comparison space has discarded important information;
\item
  two wrong local models happen to agree.
\end{itemize}

Therefore the quality of the restriction maps matters.

An easy compatibility test can be weak evidence.

\section{Richer interfaces detect more and cost
more}\label{richer-interfaces-detect-more-and-cost-more}

A boundary representation that exposes more information can detect more
incompatibilities.

It can also increase:

\begin{itemize}
\tightlist
\item
  communication;
\item
  storage;
\item
  computation;
\item
  coupling.
\end{itemize}

This creates a general design tradeoff:

\begin{quote}
expose enough boundary structure to protect composition, but not so much
that modularity disappears.
\end{quote}

The chapter does not solve this tradeoff.

It only makes it visible.

\section{Sheaf language should earn its
keep}\label{sheaf-language-should-earn-its-keep}

The word \textbf{sheaf} is useful when the problem genuinely has:

\begin{itemize}
\tightlist
\item
  local data;
\item
  overlaps/incidences;
\item
  restriction maps;
\item
  a local-to-global question.
\end{itemize}

If those objects are absent, sheaf language can become decorative.

The Atlas therefore applies a discipline:

\begin{quote}
name the section spaces and restriction maps before claiming a sheaf
viewpoint.
\end{quote}

\section{A practical local-to-global
ledger}\label{a-practical-local-to-global-ledger}

Before accepting a gluing claim, record:

{\def\LTcaptype{none} % do not increment counter
\begin{longtable}[]{@{}
  >{\raggedright\arraybackslash}p{(\linewidth - 2\tabcolsep) * \real{0.5000}}
  >{\raggedright\arraybackslash}p{(\linewidth - 2\tabcolsep) * \real{0.5000}}@{}}
\toprule\noalign{}
\begin{minipage}[b]{\linewidth}\raggedright
Field
\end{minipage} & \begin{minipage}[b]{\linewidth}\raggedright
Question
\end{minipage} \\
\midrule\noalign{}
\endhead
\bottomrule\noalign{}
\endlastfoot
domain & What is the global object/space? \\
local regions & Which pieces cover or compose it? \\
incidence & Which pieces overlap or meet? \\
local data & What is \texttt{F(U\_i)}? \\
overlap data & What is \texttt{F(U\_i\ intersect\ U\_j)}? \\
restrictions & How is local data mapped to the overlap? \\
compatibility & Which equalities must hold? \\
existence & Does every matching family considered have a glue? \\
uniqueness & Can two global objects have the same local restrictions? \\
obstruction & What records a failure to match? \\
approximation & If disagreement is tolerated, which metric/loss? \\
interpretation & What does a global section mean in the application? \\
boundary & What does consistency not prove? \\
\end{longtable}
}

This ledger prevents local-to-global language from becoming metaphor.

\section{Failure modes}\label{failure-modes}

\subsection{Graph-equals-sheaf}\label{graph-equals-sheaf}

Incidence structure is confused with the data and maps carried over it.

\subsection{Local-equals-global}\label{local-equals-global}

Locally valid pieces are assumed to assemble without overlap checks.

\subsection{Pairwise-plausible-equals-matching}\label{pairwise-plausible-equals-matching}

Informal compatibility substitutes for exact restriction equations.

\subsection{Matching-equals-truth}\label{matching-equals-truth}

A global section is treated as factual correctness.

\subsection{Approximate-equals-exact}\label{approximate-equals-exact}

A low-loss fusion is described as exact gluing.

\subsection{Zero-obstruction-equals-complete-model}\label{zero-obstruction-equals-complete-model}

The declared restrictions agree, so omitted constraints are assumed not
to exist.

\subsection{Sheaf-as-decoration}\label{sheaf-as-decoration}

The terminology is used without explicit section spaces or restriction
maps.

\section{What the exact witness
establishes}\label{what-the-exact-witness-establishes}

The companion witness proves:

\begin{itemize}
\tightlist
\item
  the global-to-local map \texttt{R:R\^{}3-\textgreater{}R\^{}4} is
  injective;
\item
  the overlap discrepancy is
  \texttt{Delta(u\_a,u\_b,v\_b,v\_c)=u\_b-v\_b};
\item
  \texttt{im(R)=ker(Delta)};
\item
  \texttt{(1,2)} on U and \texttt{(2,4)} on V glue uniquely to
  \texttt{(1,2,4)};
\item
  \texttt{(1,2)} on U and \texttt{(3,4)} on V have discrepancy
  \texttt{-1} and cannot glue;
\item
  the obstruction is localized to the shared coordinate b.
\end{itemize}

Nothing about noisy fusion, semantic truth, or global system stability
is inferred.

\section{Downstream handoff}\label{downstream-handoff-1}

Later chapters may now assume the following local-to-global vocabulary:

\begin{itemize}
\tightlist
\item
  graph/cover incidence;
\item
  local section;
\item
  restriction map;
\item
  matching family;
\item
  gluing existence;
\item
  gluing uniqueness;
\item
  discrepancy/obstruction;
\item
  exact versus approximate compatibility.
\end{itemize}

They must independently define their domain-specific section spaces,
restriction maps, and interpretation.

The local-to-global lesson is simple:

\begin{quote}
make the interfaces explicit, compare what each local piece says there,
and do not call the whole coherent until the gluing condition is
actually satisfied.
\end{quote}

\section{References used in this
chapter}\label{references-used-in-this-chapter-11}

\begin{itemize}
\tightlist
\item
  Justin Michael Curry, \emph{Sheaves, Cosheaves and Applications},
  University of Pennsylvania doctoral thesis, 2014, arXiv:1303.3255.
\item
  Michael Robinson, \emph{Sheaves are the canonical data structure for
  sensor integration}, Information Fusion 36 (2017), 208--224, DOI
  10.1016/j.inffus.2016.12.002.
\item
  Jakob Hansen and Robert Ghrist, \emph{Toward a spectral theory of
  cellular sheaves}, Journal of Applied and Computational Topology 3
  (2019), 315--358, DOI 10.1007/s41468-019-00038-7.
\end{itemize}

Exact source identities and authority boundaries are recorded in:

sources/source-locks/ATLAS-CH-LOCALGLOBAL-001.yaml

\chapter{Representations and
Invariants}\label{representations-and-invariants}

\textbf{Epistemic status:} Established Theory + Atlas Synthesis +
Computational Witness\\
\textbf{Specification:} manuscript/specifications/ATLAS-CH-REP-001.md\\
\textbf{Derivation packet:}
mathematics/derivations/ATLAS-CH-REP-001-DERIVATIONS.md\\
\textbf{Computational witness:}
mathematics/computational-witnesses/ATLAS-CW-REP-001.md\\
\textbf{Source lock:} sources/source-locks/ATLAS-CH-REP-001.yaml

\section{A representation is not merely a
vector}\label{a-representation-is-not-merely-a-vector}

A learned representation is often described as a vector

\[
z\in\mathbb R^d.
\]

That description tells us where the coordinates live.

It does not tell us:

\begin{itemize}
\tightlist
\item
  what transformations preserve meaning;
\item
  which coordinates are arbitrary;
\item
  which symmetries are respected;
\item
  what downstream tasks can recover;
\item
  which latent factors are identifiable;
\item
  whether sparse features exist;
\item
  whether multiple features share dimensions.
\end{itemize}

The Atlas therefore uses a stronger working definition.

A representation is an information-bearing encoding whose geometry,
invariances, equivalences, and downstream transformation laws matter.

This framing is consistent with the broad representation-learning
literature (Bengio et al. 2013), but the exact organization used here is
Atlas synthesis.

\section{Languages that preserve the same
message}\label{languages-that-preserve-the-same-message}

Consider two languages.

Their words look different.

Their grammar may differ.

Yet a translation can preserve the distinctions needed by a reader.

This is a useful analogy for representation equivalence.

Two encodings may differ coordinate-by-coordinate while preserving the
structure needed by downstream computation.

The analogy fails if taken literally.

Natural-language translation is approximate, context-dependent, and not
generally invertible.

Representation equivalence must be defined mathematically for the
property we care about.

\section{Representation map}\label{representation-map}

Let

\[
r:\mathcal X\to\mathcal Z
\]

map inputs into a representation space.

A representation is useful only relative to some downstream purpose.

That purpose might involve:

\begin{itemize}
\tightlist
\item
  classification;
\item
  control;
\item
  reconstruction;
\item
  retrieval;
\item
  prediction;
\item
  planning;
\item
  communication.
\end{itemize}

There is therefore no universal scalar called ``representation
quality.''

Different tasks can prefer different structures.

\section{Coordinates versus represented
structure}\label{coordinates-versus-represented-structure}

Suppose

\[
z=r(x).
\]

For a finite-dimensional real representation space
\(\mathcal Z=\mathbb R^d\), apply an invertible linear transformation

\[
T:\mathbb R^d\to\mathbb R^d
\]

and define

\[
\tilde z
=
Tz.
\]

The coordinates changed.

Did the representation change?

The answer depends on the downstream computation.

If a linear readout is

\[
y=w^\top z,
\]

then choosing

\[
\tilde w
=
T^{-\top}w
\]

gives

\[
\tilde w^\top\tilde z
=
w^\top z.
\]

For this readout class, the recoding is behavior-preserving.

This motivates an important distinction:

\[
\boxed{
\text{coordinate identity}
\neq
\text{functional equivalence}.
}
\]

\section{Exact recoding witness}\label{exact-recoding-witness}

Take

\[
T=
\begin{pmatrix}
1&1\\
0&1
\end{pmatrix},
\]

\[
z=(2,3),
\]

and

\[
w=(4,-1).
\]

The original readout is

\[
w^\top z=5.
\]

The transformed representation is

\[
\tilde z=Tz=(5,3).
\]

The transformed readout is

\[
\tilde w=T^{-\top}w=(4,-5).
\]

Then

\[
\tilde w^\top\tilde z=5.
\]

Nothing in this example says the two coordinate systems are semantically
identical for every possible downstream task.

It establishes one exact equivalence relation for a declared readout
class.

\section{Invariance}\label{invariance}

Suppose a group \(G\) acts on inputs.

A representation is invariant when

\[
\boxed{
r(g\cdot x)=r(x)
}
\]

for the transformations under consideration.

Invariance intentionally removes distinctions.

That can be useful.

If image classification should not change under a small translation, a
translation-invariant feature can reduce unnecessary variability.

But invariance has a cost.

If the downstream task requires the transformed quantity, the
information has been erased.

A position-invariant representation is poor for a task that must report
position.

\section{Equivariance}\label{equivariance}

Equivariance preserves transformation structure rather than erasing it.

Let

\[
\rho(g)
\]

be the action of group element \(g\) in representation space.

Then equivariance means

\[
\boxed{
r(g\cdot x)
=
\rho(g)r(x).
}
\]

The input transforms.

The representation transforms predictably.

Group-equivariant neural architectures provide concrete examples of this
principle (Cohen and Welling 2016).

\section{Reflection witness}\label{reflection-witness}

\begin{figure}
\centering
\pandocbounded{\includegraphics[keepaspectratio,alt={A reflection example showing an equivariant vector and invariant norm beside an invertible coordinate recoding that preserves a linear readout.}]{figures/masters/ATLAS-FIG-REP-001.png}}
\caption{A reflection example showing an equivariant vector and
invariant norm beside an invertible coordinate recoding that preserves a
linear readout.}
\end{figure}

Let

\[
G=
\begin{pmatrix}
-1&0\\
0&1
\end{pmatrix}
\]

reflect the first coordinate.

For

\[
x=(2,3),
\]

\[
Gx=(-2,3).
\]

The identity representation

\[
r(x)=x
\]

is equivariant:

\[
r(Gx)=Gr(x).
\]

The scalar representation

\[
r_{\rm inv}(x)
=
\|x\|_2^2
\]

is invariant:

\[
r_{\rm inv}(Gx)
=
r_{\rm inv}(x)
=
13.
\]

These are different design choices.

One preserves the transformation.

The other removes it.

\section{Invariance and equivariance are not
synonyms}\label{invariance-and-equivariance-are-not-synonyms}

The difference matters.

Invariance says:

\begin{quote}
the representation should not change.
\end{quote}

Equivariance says:

\begin{quote}
the representation should change in a predictable corresponding way.
\end{quote}

A system that confuses these can destroy useful structure.

The right choice depends on the task.

\section{Geometry enters
representation}\label{geometry-enters-representation}

The Geometry chapter established that admissible states can live on
constrained spaces.

Representation learning inherits that fact.

Examples include:

\begin{itemize}
\tightlist
\item
  unit-normalized embeddings;
\item
  orthogonal frames;
\item
  quotient representations;
\item
  spherical positions;
\item
  subspace-valued features.
\end{itemize}

A representation is therefore not always best modeled as an
unconstrained point in \(\mathbb R^d\).

Its geometry can carry semantics.

This is the direct handoff from the Geometry keystone.

\section{Information enters
representation}\label{information-enters-representation}

The previous chapter supplied probability and information structure.

Representation learning often asks whether an encoding preserves or
discards statistically relevant information.

But an information measure does not completely determine representation
quality.

Two encodings can have the same mutual information with a target and
still differ in:

\begin{itemize}
\tightlist
\item
  geometry;
\item
  robustness;
\item
  linear accessibility;
\item
  sparsity;
\item
  causal role;
\item
  compositional compatibility.
\end{itemize}

Information is one axis.

Representation structure is richer.

\section{Identifiability}\label{identifiability}

Suppose latent variables \(z\) generate observations through a decoder

\[
x=f(z).
\]

Now choose an invertible transformation \(T\) and define

\[
\tilde z=Tz.
\]

The decoder can be rewritten as

\[
\tilde f(\tilde z)
=
f(T^{-1}\tilde z).
\]

The same observations can therefore be explained through different
latent coordinates.

Without additional assumptions, there may be no unique latent
representation.

This is an identifiability problem.

\section{Disentanglement needs
assumptions}\label{disentanglement-needs-assumptions}

A common hope is that unsupervised learning will recover independent or
semantically meaningful generative factors automatically.

That hope requires care.

Locatello et al.~show that unsupervised disentanglement is not
identifiable without inductive biases under the setting they analyze
(Locatello et al. 2019).

The Atlas uses that result as a boundary.

It does not conclude that disentanglement is useless.

It concludes that claims of uniquely recovered semantic factors require
assumptions, supervision, or other structure.

\section{Sparse representations}\label{sparse-representations}

A representation is sparse when relatively few coefficients are active.

One classical model writes

\[
x\approx D\alpha,
\]

where \(D\) is a dictionary and \(\alpha\) is sparse.

Olshausen and Field gave a historically influential example in which
sparse coding of natural images produced localized oriented
receptive-field-like structures (Olshausen and Field 1996).

This is evidence for one representation regime.

It is not proof that neural representations are universally sparse.

\section{Dictionaries and features}\label{dictionaries-and-features}

A dictionary representation introduces another distinction.

The coordinates \(\alpha_i\) are meaningful only relative to dictionary
elements \(d_i\).

Changing the dictionary changes the coordinate interpretation.

The Atlas will later use this observation when discussing:

\begin{itemize}
\tightlist
\item
  feature dictionaries;
\item
  sparse autoencoders;
\item
  superposition;
\item
  reconstruction;
\item
  transferable residual structure.
\end{itemize}

Again, coordinates and represented structure are not identical objects.

\section{Superposition}\label{superposition}

If the number of potentially useful features exceeds the number of
representational dimensions, a system may encode multiple features
non-orthogonally.

Toy Models of Superposition studies this phenomenon in deliberately
simplified systems and provides geometric and empirical toy-model
evidence (Elhage et al. 2022).

The source is useful.

Its scope is narrow.

The Atlas does not infer:

\begin{quote}
all frontier neural networks store features exactly according to the toy
model.
\end{quote}

The chapter uses superposition as a research concept, not a universal
established law.

\section{Representation
equivalence}\label{representation-equivalence}

What does it mean for two representations to be ``the same''?

There is no single answer.

Possible equivalence relations include:

\begin{itemize}
\tightlist
\item
  exact coordinate equality;
\item
  invertible linear recoding;
\item
  invertible nonlinear recoding;
\item
  equality up to group action;
\item
  equality of pairwise distances;
\item
  equality of downstream predictions;
\item
  equality of recoverable sufficient statistics.
\end{itemize}

Different scientific questions require different equivalence relations.

The Atlas therefore treats ``representation equivalence'' as a
declaration, not a default.

\section{Quotient thinking}\label{quotient-thinking}

If multiple coordinates represent the same object, the mathematically
natural object may be an equivalence class.

That leads toward quotient spaces.

For example, the Geometry chapter distinguished orthonormal frames from
the subspaces they span.

Different frames can represent the same subspace.

The later Quotient Representations chapter generalizes this lesson.

\section{Representation collapse}\label{representation-collapse}

A representation can discard too much.

If

\[
r(x)=c
\]

for every input, the representation is perfectly invariant to
everything.

It is also useless for tasks that need to distinguish inputs.

This is the simplest counterexample to the idea that more invariance is
always better.

The useful question is:

\begin{quote}
invariant to what, while preserving what?
\end{quote}

\section{Identical task performance does not imply identical
mechanism}\label{identical-task-performance-does-not-imply-identical-mechanism}

Two representations can support the same benchmark score.

That does not mean they encode the same internal structure.

One may use:

\begin{itemize}
\tightlist
\item
  sparse features;
\item
  distributed features;
\item
  different symmetries;
\item
  different causal pathways;
\item
  different robustness margins.
\end{itemize}

Behavioral equivalence on one task is therefore a weak equivalence
relation.

Mechanistic claims require stronger evidence.

\section{The Residual question}\label{the-residual-question}

The Atlas later asks a more ambitious question:

\begin{quote}
What is the minimal transferable structure that survives relevant
representational changes?
\end{quote}

That programme uses the provisional term \textbf{Residual}.

This chapter does not answer the question.

It prepares the mathematics required to ask it correctly.

If representations can differ under recoding while preserving
capability, then the transferable object may not be any one coordinate
system.

\section{Five failure modes}\label{five-failure-modes-1}

\subsection{Coordinates mistaken for
meaning}\label{coordinates-mistaken-for-meaning}

A basis component is treated as a semantic primitive without an
identifiability argument.

\subsection{Invariance mistaken for universal
virtue}\label{invariance-mistaken-for-universal-virtue}

Task-relevant variation is erased.

\subsection{Equivariance claimed without group
action}\label{equivariance-claimed-without-group-action}

The transformation law is not specified.

\subsection{Disentanglement claimed without
assumptions}\label{disentanglement-claimed-without-assumptions}

Non-identifiability is ignored.

\subsection{Toy superposition promoted to universal
mechanism}\label{toy-superposition-promoted-to-universal-mechanism}

A bounded research model is treated as established empirical prevalence.

\section{Atlas connections}\label{atlas-connections-6}

\textbf{Normalized representations.}\\
Once geometry matters, unit-norm and hyperspherical representations
become natural special cases.

\textbf{Quotient representations.}\\
Equivalence under transformations suggests quotient structure.

\textbf{Attention and position.}\\
Representations are acted on by state-dependent and position-dependent
operators.

\textbf{Memory.}\\
External stores may preserve semantic objects under different internal
encodings.

\textbf{Diagnostics.}\\
Interpretability must distinguish coordinate features from invariant or
causally relevant structure.

\textbf{Residual programme.}\\
Transfer under representation change becomes a first-class question.

\section{Closing view}\label{closing-view-9}

A representation is not just where a vector lives.

It is a contract between an input, an encoding, a transformation law,
and a downstream use.

Its coordinates can change while capability survives.

Its invariances can help or destroy.

Its latent factors may not be identifiable.

Its features may be sparse, distributed, or superposed.

The Atlas therefore asks a stronger question than:

\begin{quote}
What is the representation?
\end{quote}

It asks:

\begin{quote}
Which properties of the representation survive the transformations that
should not matter, and which distinctions remain available to the
computations that do?
\end{quote}

That question opens the path to normalized geometry, quotient spaces,
and the search for transferable residual structure.

\section{References used in this
chapter}\label{references-used-in-this-chapter-12}

\begin{itemize}
\tightlist
\item
  (Bengio et al. 2013)
\item
  (Cohen and Welling 2016)
\item
  (Locatello et al. 2019)
\item
  (Olshausen and Field 1996)
\item
  (Elhage et al. 2022)
\end{itemize}

\chapter{Normalized and Hyperspherical
Representations}\label{normalized-and-hyperspherical-representations}

\textbf{Epistemic status:} Established Theory + Atlas Derivation +
Source-Scoped Research Evidence\\
\textbf{Specification:}
manuscript/specifications/ATLAS-CH-NORMREP-001.md\\
\textbf{Derivation packet:}
mathematics/derivations/ATLAS-CH-NORMREP-001-DERIVATIONS.md\\
\textbf{Computational witness:}
mathematics/computational-witnesses/ATLAS-CW-NORMREP-001.md\\
\textbf{Source lock:} sources/source-locks/ATLAS-CH-NORMREP-001.yaml

\section{Normalization is a modeling
decision}\label{normalization-is-a-modeling-decision}

A vector

\[
x\in\mathbb R^d
\]

contains at least two kinds of Euclidean information:

\begin{itemize}
\tightlist
\item
  direction;
\item
  norm.
\end{itemize}

The normalization map

\[
N(x)
=
\frac{x}{\|x\|_2}
\]

keeps the first and removes the second.

This is often useful.

It is never neutral.

If two vectors lie on the same positive ray,

\[
x_2=cx_1,
\qquad
c>0,
\]

then

\[
N(x_2)=N(x_1).
\]

After normalization, positive radial scale has become an invariance.

The central question of this chapter is therefore not:

\begin{quote}
Why is normalization good?
\end{quote}

It is:

\begin{quote}
What mathematical structure do we create when we deliberately remove
radial freedom?
\end{quote}

\section{Direction without volume}\label{direction-without-volume}

The useful allegory is an arrow whose length has been erased.

Before normalization, the arrow has:

\begin{itemize}
\tightlist
\item
  a direction;
\item
  a magnitude.
\end{itemize}

After normalization, only direction remains.

The structural correspondence is:

\begin{itemize}
\tightlist
\item
  direction ↔ angular information;
\item
  arrow length ↔ radial information;
\item
  tangent step ↔ change of direction;
\item
  renormalization ↔ return to the sphere.
\end{itemize}

The limit is essential.

Length is not noise by definition.

A task may depend on:

\begin{itemize}
\tightlist
\item
  magnitude;
\item
  confidence encoded in norm;
\item
  energy;
\item
  count;
\item
  scale;
\item
  uncertainty.
\end{itemize}

Normalization should therefore be understood as a declared invariance,
not an automatic improvement.

\section{The unit hypersphere}\label{the-unit-hypersphere}

For unit vectors,

\[
u\in S^{d-1}
=
\{u\in\mathbb R^d:\|u\|_2=1\}.
\]

The Geometry chapter established the tangent space

\[
T_uS^{d-1}
=
\{\xi:u^\top\xi=0\}.
\]

Once representations are constrained to the sphere, legal first-order
motion is tangential.

This changes the meaning of an update.

An unconstrained Euclidean step can point partly:

\begin{itemize}
\tightlist
\item
  along the sphere;
\item
  away from the sphere.
\end{itemize}

Only the first component changes direction to first order.

\section{The normalization map}\label{the-normalization-map}

For

\[
x\neq0,
\]

define

\[
u=N(x)
=
\frac{x}{\|x\|_2}.
\]

Let

\[
r=\|x\|_2.
\]

Then

\[
N(x)=r^{-1}x.
\]

Differentiating gives

\[
\boxed{
J_N(x)
=
\frac{1}{\|x\|_2}
\left(
I-uu^\top
\right).
}
\]

The matrix

\[
I-uu^\top
\]

is the Euclidean projector onto the tangent space at \(u\).

Normalization therefore has a clean local interpretation:

\begin{quote}
project away the radial differential component, then scale the tangent
component by \(1/\|x\|_2\).
\end{quote}

\section{Radial and tangential
decomposition}\label{radial-and-tangential-decomposition}

Any perturbation

\[
\delta\in\mathbb R^d
\]

can be decomposed as

\[
\delta
=
(uu^\top)\delta
+
(I-uu^\top)\delta.
\]

The first term is radial.

The second is tangent.

Applying the normalization Jacobian gives

\[
J_N(x)\delta
=
\frac{1}{\|x\|_2}
(I-uu^\top)\delta.
\]

The radial term vanishes.

This is not metaphorical.

It is the differential of the normalization map.

\section{\texorpdfstring{Exact witness at
\((3,4)\)}{Exact witness at (3,4)}}\label{exact-witness-at-34}

Take

\[
x=(3,4).
\]

Then

\[
\|x\|_2=5,
\]

and

\[
u=(3/5,4/5).
\]

The normalization Jacobian is

\[
J_N(x)
=
\begin{pmatrix}
16/125&-12/125\\
-12/125&9/125
\end{pmatrix}.
\]

The exact computational witness verifies

\[
J_N(x)x=0.
\]

The radial direction disappears to first order.

Now take

\[
t=(-4,3).
\]

Since

\[
x^\top t=0,
\]

this is tangent to the radius through \(x\).

The witness verifies

\[
J_N(x)t
=
(-4/5,3/5).
\]

The direction survives.

Its scale changes by \(1/5\).

\begin{figure}
\centering
\pandocbounded{\includegraphics[keepaspectratio,alt={A unit-circle figure showing the normalized direction (3/5,4/5), a radial direction removed by normalization, an orthogonal tangent direction, and a second panel comparing the chord, great-circle arc, and SLERP midpoint for two unit vectors separated by pi over three.}]{figures/masters/ATLAS-FIG-NORMREP-001.png}}
\caption{A unit-circle figure showing the normalized direction
(3/5,4/5), a radial direction removed by normalization, an orthogonal
tangent direction, and a second panel comparing the chord, great-circle
arc, and SLERP midpoint for two unit vectors separated by pi over
three.}
\end{figure}

\section{A singularity at zero}\label{a-singularity-at-zero}

Normalization is undefined at

\[
x=0.
\]

This is not a minor implementation detail.

The map

\[
x\mapsto\frac{x}{\|x\|_2}
\]

has no direction to preserve at the origin.

Practical systems may:

\begin{itemize}
\tightlist
\item
  add an epsilon;
\item
  avoid zero states by construction;
\item
  use guarded normalization;
\item
  switch to another chart or convention.
\end{itemize}

Each choice changes the exact mathematical object.

The Atlas will therefore distinguish idealized normalization from
epsilon-regularized implementations where the difference matters.

\section{Norm erasure is information
erasure}\label{norm-erasure-is-information-erasure}

Take

\[
x_1=(1,0),
\qquad
x_2=(2,0).
\]

Then

\[
N(x_1)=N(x_2)=(1,0).
\]

Suppose the target is

\[
y=\|x\|_2.
\]

Before normalization,

\[
y_1=1,
\qquad
y_2=2.
\]

After normalization, the two states are identical.

No downstream function of the normalized representation alone can
recover which norm was present.

This is the smallest counterexample to the idea that norm is always
redundant.

\section{Angular similarity}\label{angular-similarity}

For unit vectors \(u,v\),

\[
u^\top v
=
\cos\theta,
\]

where \(\theta\) is the angle between them.

Thus cosine similarity reduces to the dot product.

This is one reason normalized representations are attractive:

scale no longer contaminates angular comparison.

But the statement should be read literally.

Normalization makes the comparison scale-invariant because scale has
been removed.

It does not prove that scale was irrelevant to the original task.

\section{Chord distance and angle}\label{chord-distance-and-angle}

For unit vectors,

\[
\begin{aligned}
\|u-v\|_2^2
&=
(u-v)^\top(u-v)\\
&=
\|u\|_2^2+\|v\|_2^2-2u^\top v\\
&=
2-2\cos\theta.
\end{aligned}
\]

Therefore

\[
\boxed{
\|u-v\|_2^2
=
2-2\cos\theta.
}
\]

Chord distance and angular distance are monotonically related for

\[
0\le\theta\le\pi.
\]

They remain different metrics.

The Geometry chapter's geodesic distance on the unit sphere is

\[
d_{\rm geo}(u,v)=\theta.
\]

The ambient chord distance is

\[
d_{\rm chord}(u,v)
=
\sqrt{2-2\cos\theta}.
\]

\section{Sixty-degree witness}\label{sixty-degree-witness}

Take

\[
a=(1,0),
\]

and

\[
b=
\left(
\frac12,
\frac{\sqrt3}{2}
\right).
\]

Then

\[
a^\top b=\frac12,
\]

so

\[
\theta=\frac{\pi}{3}.
\]

The chord length is

\[
\|a-b\|_2=1.
\]

The geodesic distance is

\[
\frac{\pi}{3}.
\]

One pair of unit vectors therefore carries at least two natural
distances.

Which one matters depends on the problem.

\section{SLERP}\label{slerp}

Let \(a,b\) be unit vectors with angle

\[
\theta=\arccos(a^\top b).
\]

For \(0<\theta<\pi\), spherical linear interpolation is

\[
\operatorname{SLERP}(a,b;t)
=
\frac{\sin((1-t)\theta)}{\sin\theta}a
+
\frac{\sin(t\theta)}{\sin\theta}b.
\]

For coincident endpoints \(a=b\), the displayed quotient is \(0/0\), so
define \(\operatorname{SLERP}(a,a;t)=a\) by continuous extension. At
antipodal endpoints \(a=-b\), the shortest interpolation path is not
unique without an additional directional choice.

For the sixty-degree witness at

\[
t=\frac12,
\]

the exact midpoint is

\[
\left(
\frac{\sqrt3}{2},
\frac12
\right).
\]

It remains unit norm.

It also lies halfway along the great-circle arc.

Straight interpolation in ambient coordinates does not generally do
that.

\section{Retraction}\label{retraction}

The Geometry chapter introduced the normalized sphere retraction

\[
R_u(\xi)
=
\frac{u+\xi}{\|u+\xi\|_2},
\]

for tangent

\[
u^\top\xi=0.
\]

This provides a simple update pattern:

\begin{enumerate}
\def\labelenumi{\arabic{enumi}.}
\tightlist
\item
  compute a tangent direction;
\item
  take an ambient step;
\item
  renormalize.
\end{enumerate}

The result stays on the sphere.

A retraction approximates the manifold exponential map locally.

It is not generally identical to the exponential map.

\section{Angular learning rate}\label{angular-learning-rate}

Once a state lives on a sphere, raw Euclidean step norm is no longer the
only natural update magnitude.

A more intrinsic quantity is angular displacement.

If unit states \(u\) and \(u'\) satisfy

\[
u^\top u'=\cos\theta,
\]

then

\[
\theta
\]

directly measures rotational change.

This makes ``learning rate'' interpretable in angular units when the
update rule is explicitly designed that way.

The public GCL MODULUS Hyperball module implements this kind of
machinery:

\begin{itemize}
\tightlist
\item
  optional tangent projection;
\item
  target angular/update scaling;
\item
  normalization-based retraction;
\item
  norm control.
\end{itemize}

Its source explicitly describes the module as designed for ``nGPT /
MODULUS-style training.''

That is a real public implementation connection.

It is not evidence of a separate public GCL nGPT repository.

\section{Hyperspherical contrastive
representations}\label{hyperspherical-contrastive-representations}

Normalized features also appear outside Transformer architectures.

Wang and Isola analyze contrastive representations in terms of:

\begin{itemize}
\tightlist
\item
  alignment of positive pairs;
\item
  uniformity of normalized features on the hypersphere (Wang and Isola
  2020).
\end{itemize}

This provides a peer-reviewed example in which hyperspherical geometry
is not decorative.

The distribution of normalized features on the sphere becomes part of
the analysis.

The result is still setting-specific.

It does not establish that every useful representation should be
uniformly distributed on a hypersphere.

\section{nGPT}\label{ngpt}

The nGPT paper proposes a normalized Transformer with representation
learning on the hypersphere (Loshchilov et al. 2024).

The paper's central architectural move is to constrain major
representation and parameter vectors to unit-norm geometry and interpret
layerwise changes as movement on the hypersphere.

A later preprint develops a practical training recipe and reports
scaling experiments on larger hybrid architectures (Loshchilov and
Ginsburg 2026).

These are important examples because they make the geometry operational.

The Atlas keeps three levels separate:

\begin{enumerate}
\def\labelenumi{\arabic{enumi}.}
\tightlist
\item
  sphere geometry is established mathematics;
\item
  nGPT is a specific research architecture and empirical programme;
\item
  broader claims that normalized Transformers are universally superior
  remain unsupported.
\end{enumerate}

\section{Geometry does not choose the task for
us}\label{geometry-does-not-choose-the-task-for-us}

Suppose two hidden states differ only in norm.

On the unit sphere they are identified.

This can be desirable when the system should care only about direction.

It can be harmful when norm encodes:

\begin{itemize}
\tightlist
\item
  evidence strength;
\item
  uncertainty;
\item
  frequency;
\item
  value magnitude;
\item
  resource level.
\end{itemize}

The right representation geometry therefore depends on what distinctions
the downstream task should preserve.

\section{Normalization can relocate
information}\label{normalization-can-relocate-information}

Removing norm from one representation does not necessarily remove
magnitude from the whole system.

A model may reintroduce scale through:

\begin{itemize}
\tightlist
\item
  separate scalar channels;
\item
  learned temperatures;
\item
  biases;
\item
  logits;
\item
  gating coefficients;
\item
  residual amplitudes;
\item
  external metadata.
\end{itemize}

This leads to a more useful question:

\begin{quote}
Where should radial information live?
\end{quote}

A normalized architecture may be best understood not as destroying all
scale, but as assigning scale to explicit channels rather than letting
it remain entangled with direction.

That is an architectural hypothesis.

It must be tested system by system.

\section{Collapse on the sphere}\label{collapse-on-the-sphere}

Unit norm does not prevent representational collapse.

Every input can map to the same unit vector.

Then

\[
\|r(x)\|_2=1
\]

for all \(x\), yet the representation preserves no input distinctions.

Normalization controls radius.

It does not guarantee:

\begin{itemize}
\tightlist
\item
  diversity;
\item
  separability;
\item
  identifiability;
\item
  useful semantics.
\end{itemize}

This is why hyperspherical representation objectives often care about
angular distribution as well as norm.

\section{Antipodal ambiguity}\label{antipodal-ambiguity}

SLERP has a special case when

\[
a=-b.
\]

Then

\[
\theta=\pi,
\]

and the shortest interpolation path is not unique. On \(S^1\) there are
two shortest semicircular arcs; on \(S^{d-1}\) for \(d\ge3\) there are
infinitely many great-circle semicircles joining the endpoints.

This is another reminder that normalization simplifies one degree of
freedom without making all geometry trivial.

\section{Five failure modes}\label{five-failure-modes-2}

\subsection{Radial information assumed
irrelevant}\label{radial-information-assumed-irrelevant}

Normalization removes a quantity the task needed.

\subsection{Unit norm mistaken for
non-collapse}\label{unit-norm-mistaken-for-non-collapse}

All states can collapse to one point on the sphere.

\subsection{Cosine similarity mistaken for complete
similarity}\label{cosine-similarity-mistaken-for-complete-similarity}

Angular proximity can ignore task-relevant scale or other structure.

\subsection{Retraction mistaken for exponential
map}\label{retraction-mistaken-for-exponential-map}

The normalized update is a local manifold-compatible device, not exact
geodesic flow in general.

\subsection{nGPT evidence
overgeneralized}\label{ngpt-evidence-overgeneralized}

Paper-scoped empirical results are promoted into a universal
architectural law.

\section{Atlas connections}\label{atlas-connections-7}

\textbf{Representation.}\\
Normalization declares positive radial scaling an invariance.

\textbf{Geometry.}\\
The sphere supplies tangent spaces, geodesics, retractions, and SLERP.

\textbf{Optimization.}\\
Updates can be measured by angular displacement and projected
tangentially.

\textbf{Attention.}\\
Normalized queries, keys, values, or parameter vectors alter the
scale/angular decomposition of attention computations.

\textbf{Position.}\\
Spherical geometry provides a natural substrate for directional or
harmonic positional constructions.

\textbf{Quotient thinking.}\\
Removing radial scale can be understood as identifying states along
positive rays, though the precise quotient and sphere cross-section
should not be conflated automatically.

\section{Closing view}\label{closing-view-10}

Normalization changes the question a representation can answer.

Before normalization, a vector can say:

\begin{quote}
where am I pointing, and how large am I?
\end{quote}

After normalization, it can say only:

\begin{quote}
where am I pointing?
\end{quote}

That trade is sometimes exactly what we want.

It exposes clean angular geometry.

It makes tangent motion and spherical interpolation natural.

It can stabilize or clarify specific architectures.

But it is still a trade.

The correct mathematical stance is not to celebrate norm removal.

It is to account for what has been removed, what structure has been
gained, and where the missing information must go if the system still
needs it.

\section{References used in this
chapter}\label{references-used-in-this-chapter-13}

\begin{itemize}
\tightlist
\item
  (Lee 2018)
\item
  (Absil et al. 2008)
\item
  (Wang and Isola 2020)
\item
  (Loshchilov et al. 2024)
\item
  (Loshchilov and Ginsburg 2026)
\end{itemize}

See `sources/source-locks/ATLAS-CH-NORMREP-001.yaml` for exact source
roles, the public GCL MODULUS implementation-context object, and claim
scope.

\chapter{Equivalence and Quotient
Geometry}\label{equivalence-and-quotient-geometry}

\textbf{Epistemic status:} Established Theory + Atlas Derivation + Open
Research Question\\
\textbf{Specification:}
manuscript/specifications/ATLAS-CH-QUOTIENT-001.md\\
\textbf{Derivation packet:}
mathematics/derivations/ATLAS-CH-QUOTIENT-001-DERIVATIONS.md\\
\textbf{Computational witness:}
mathematics/computational-witnesses/ATLAS-CW-QUOTIENT-001.md\\
\textbf{Source lock:} sources/source-locks/ATLAS-CH-QUOTIENT-001.yaml

\section{Different parameters can mean the same
thing}\label{different-parameters-can-mean-the-same-thing}

Suppose two neural networks have different parameter arrays.

Are they different models?

At the level of stored bytes, yes.

At the level of the represented function, perhaps not.

A hidden-unit permutation can reorder internal coordinates without
changing the input-output map.

A positive rescaling in a ReLU network can change incoming and outgoing
weights while preserving the function exactly.

A basis change can alter a frame without changing the subspace it spans.

These are not numerical coincidences.

They are redundancies in representation.

Quotient geometry asks us to stop treating redundant representatives as
distinct objects when the scientific question concerns the underlying
equivalence class.

\section{Many addresses, one place}\label{many-addresses-one-place}

Imagine one physical location with several valid addresses.

Different naming conventions can point to the same place.

The structural correspondence is:

\begin{itemize}
\tightlist
\item
  address ↔ parameter representative;
\item
  place ↔ represented function or equivalence class;
\item
  address conversion ↔ symmetry action;
\item
  quotient map ↔ forgetting redundant labels.
\end{itemize}

The analogy is useful because it separates identity of coordinates from
identity of the object.

It also has a strict limit.

Neural parameter spaces can contain:

\begin{itemize}
\tightlist
\item
  fixed points;
\item
  singular strata;
\item
  hidden symmetries;
\item
  changing stabilizers;
\item
  regions where a smooth quotient description fails.
\end{itemize}

The quotient is not automatically one smooth map of a tidy city.

\section{Equivalence relations}\label{equivalence-relations}

A relation

\[
\sim
\]

on a set \(X\) is an equivalence relation if it is:

\begin{itemize}
\tightlist
\item
  reflexive;
\item
  symmetric;
\item
  transitive.
\end{itemize}

The equivalence class of \(x\) is

\[
[x]
=
\{y\in X:y\sim x\}.
\]

The quotient set is

\[
\boxed{
X/{\sim}
=
\{[x]:x\in X\}.
}
\]

The quotient keeps the distinctions that survive the chosen equivalence
relation and removes the distinctions declared redundant.

Everything depends on that declaration.

Choose an equivalence relation that is too coarse and functionally
different objects can be collapsed together.

\section{Group actions and orbits}\label{group-actions-and-orbits}

Let a group \(G\) act on \(X\).

The orbit of \(x\) is

\[
G\cdot x
=
\{g\cdot x:g\in G\}.
\]

A natural equivalence relation is

\[
x\sim y
\iff
y=g\cdot x
\text{ for some }g\in G.
\]

Then equivalence classes are group orbits.

This construction is central to symmetry reduction.

It also immediately raises a geometric question:

\begin{quote}
Does the orbit space inherit a well-behaved smooth structure?
\end{quote}

Sometimes yes.

Sometimes only on a regular subset.

Sometimes singularities remain.

\section{Grassmann: the canonical
example}\label{grassmann-the-canonical-example}

The Geometry chapter distinguished:

\[
\operatorname{St}(n,p),
\]

the space of ordered orthonormal \(p\)-frames, from

\[
\operatorname{Gr}(n,p),
\]

the space of \(p\)-dimensional subspaces.

If

\[
X\in\operatorname{St}(n,p),
\]

and

\[
Q\in O(p),
\]

then

\[
XQ
\]

is another orthonormal frame spanning the same subspace.

Therefore

\[
\boxed{
\operatorname{Gr}(n,p)
\cong
\operatorname{St}(n,p)/O(p).
}
\]

The quotient removes basis choice.

The subspace remains.

This is not merely philosophical.

It changes which directions are treated as genuine motion.

\section{Parameter equality versus functional
equality}\label{parameter-equality-versus-functional-equality}

Let

\[
f_\theta
\]

denote the function represented by parameters \(\theta\).

Parameter equality is

\[
\theta=\theta'.
\]

Functional equality is

\[
f_\theta(x)=f_{\theta'}(x)
\]

for every admissible input \(x\).

The second is weaker.

A parameterization can therefore be many-to-one:

\[
\theta
\mapsto
f_\theta.
\]

When that happens, Euclidean geometry in parameter space contains
directions that may change coordinates without changing the represented
function.

This is the opening for quotient analysis.

\section{Hidden-unit permutation
symmetry}\label{hidden-unit-permutation-symmetry}

Consider

\[
f(x)
=
W_2\sigma(W_1x),
\]

with elementwise activation \(\sigma\).

Let \(P\) permute hidden units.

Define

\[
W_1'=PW_1,
\]

and

\[
W_2'=W_2P^{-1}.
\]

Because elementwise activation commutes with permutation,

\[
\sigma(Pz)=P\sigma(z).
\]

Therefore

\[
\begin{aligned}
f'(x)
&=
W_2P^{-1}\sigma(PW_1x)\\
&=
W_2P^{-1}P\sigma(W_1x)\\
&=
f(x).
\end{aligned}
\]

The hidden units were relabeled.

The function did not change.

\section{Exact two-unit example}\label{exact-two-unit-example}

Take scalar input and

\[
W_1=
\begin{pmatrix}
1\\
2
\end{pmatrix},
\qquad
W_2=
\begin{pmatrix}
3&4
\end{pmatrix}.
\]

With ReLU activation,

\[
f(x)
=
3\operatorname{ReLU}(x)
+
4\operatorname{ReLU}(2x).
\]

For positive \(x\),

\[
f(x)=11x.
\]

For nonpositive \(x\),

\[
f(x)=0.
\]

Swap the hidden units:

\[
W_1'=
\begin{pmatrix}
2\\
1
\end{pmatrix},
\qquad
W_2'=
\begin{pmatrix}
4&3
\end{pmatrix}.
\]

The parameter arrays change.

The function remains

\[
f(x)=11\max(0,x).
\]

\section{Positive-rescaling
symmetry}\label{positive-rescaling-symmetry}

ReLU is positively homogeneous:

\[
\operatorname{ReLU}(cz)
=
c\operatorname{ReLU}(z),
\qquad
c>0.
\]

For positive diagonal \(D\),

\[
\operatorname{ReLU}(Dz)
=
D\operatorname{ReLU}(z).
\]

Define

\[
W_1''=DW_1,
\]

and

\[
W_2''=W_2D^{-1}.
\]

Then

\[
f''(x)=f(x).
\]

For

\[
D=\operatorname{diag}(2,1/3),
\]

the exact witness becomes

\[
W_1''=
\begin{pmatrix}
2\\
2/3
\end{pmatrix},
\]

\[
W_2''=
\begin{pmatrix}
3/2&12
\end{pmatrix}.
\]

These parameters look very different.

They still compute the same function.

\begin{figure}
\centering
\pandocbounded{\includegraphics[keepaspectratio,alt={Three distinct parameter representatives, original, hidden-unit permuted, and positively rescaled, all mapping to the exact same ReLU function f(x)=11 max(0,x).}]{figures/masters/ATLAS-FIG-QUOTIENT-001.png}}
\caption{Three distinct parameter representatives, original, hidden-unit
permuted, and positively rescaled, all mapping to the exact same ReLU
function f(x)=11 max(0,x).}
\end{figure}

\section{Exact replay}\label{exact-replay}

The computational witness evaluates all three parameterizations on

\[
\{-2,-3/2,-1,-1/2,0,1/2,1,3/2,2\}.
\]

Every representative gives

\[
\{0,0,0,0,0,11/2,11,33/2,22\}.
\]

That finite replay is a check.

The algebra above is the reason equality holds for every real input.

This distinction mirrors the Atlas evidence doctrine:

computation and proof can support one another without becoming the same
object.

\section{Why quotienting is
tempting}\label{why-quotienting-is-tempting}

Suppose an optimizer moves through parameter space.

If many parameter directions correspond only to symmetry motion, then
ambient Euclidean distance can overstate meaningful change.

Two far-separated parameter vectors may compute the same function.

This suggests a cleaner object:

\[
[\theta]
\in
\Theta/G.
\]

Instead of asking how far parameters moved, ask how far the equivalence
class moved.

This can remove redundancy from the description.

It does not automatically remove difficulty from the optimization
problem.

\section{Vertical and horizontal
directions}\label{vertical-and-horizontal-directions}

On a regular quotient manifold, tangent directions can be separated
conceptually into:

\begin{itemize}
\tightlist
\item
  \textbf{vertical directions}, tangent to the symmetry orbit;
\item
  \textbf{horizontal directions}, transverse directions that represent
  quotient motion.
\end{itemize}

The terminology is standard in quotient geometry (Absil et al. 2008).

A vertical move changes the representative.

A horizontal move changes the equivalence class.

This distinction is especially useful when the scientific object is the
represented function rather than the raw parameter vector.

\section{Local null directions}\label{local-null-directions}

Suppose

\[
F(\theta)
\]

is constant along a smooth symmetry orbit

\[
\theta(s).
\]

Then

\[
F(\theta(s))
=
F(\theta(0)).
\]

Differentiating at zero gives

\[
J_F(\theta)
\dot\theta(0)
=
0.
\]

The orbit tangent lies in the local null space of the function map.

If a loss depends only on \(F(\theta)\), its directional derivative also
vanishes along the orbit.

An actual Hessian null vector requires a stronger statement than ``one
path has constant loss.''

Assume a smooth one-parameter group action \(g_s\) is an exact symmetry
of the loss in a neighborhood:

\[
L(g_s\theta)=L(\theta).
\]

If \(\theta_\star\) is a smooth critical point, symmetry maps it to
other critical points.

For the orbit tangent

\[
v
=
\frac{d}{ds}g_s\theta_\star
\Big|_{s=0},
\]

differentiating

\[
\nabla L(g_s\theta_\star)=0
\]

gives

\[
\boxed{
H_L(\theta_\star)v=0.
}
\]

A minimal exact example is

\[
L(a,b)=\frac12(ab-y)^2.
\]

Positive rescaling

\[
(a,b)\mapsto(a/c,cb)
\]

preserves the model product.

On the critical manifold \(y=ab\),

\[
H
=
\begin{pmatrix}
b^2&ab\\
ab&a^2
\end{pmatrix},
\]

and the scaling-orbit tangent

\[
v=(-a,b)
\]

satisfies

\[
Hv=0.
\]

This is a precise source of Hessian degeneracy.

It does not mean every flat direction is a symmetry direction.

\section{Hidden symmetries}\label{hidden-symmetries}

Permutation and positive scaling are not necessarily the whole story.

Grigsby, Lindsey, and Rolnick study hidden symmetries in ReLU networks
and show that the redundancy structure can be richer and
parameter-dependent (Grigsby et al. 2023).

This matters for quotient reasoning.

If the symmetry group changes across parameter space, a single globally
regular manifold model may fail.

The orbit space can have different local structure in different regions.

\section{Symmetry can affect learning
dynamics}\label{symmetry-can-affect-learning-dynamics}

Symmetry is not only a bookkeeping problem.

It can constrain optimization and interact with regularization and
noise.

Ziyin studies consequences of loss symmetries for learned parameter
structure under explicit assumptions (Ziyin 2024).

The Atlas uses this literature to support a limited point:

\begin{quote}
symmetry can influence learning dynamics.
\end{quote}

It does not infer that one universal quotient geometry determines all
neural training.

\section{Quotienting removes declared
redundancy}\label{quotienting-removes-declared-redundancy}

Suppose the group \(G\) captures an exact functional symmetry.

Then moving within one orbit does not change the function.

Passing to

\[
\Theta/G
\]

removes those orbit directions from the represented object.

This is a genuine simplification of description.

It can reduce:

\begin{itemize}
\tightlist
\item
  redundant coordinates;
\item
  gauge-like degrees of freedom;
\item
  ambiguity in comparing models.
\end{itemize}

That is already useful.

A stronger claim would be:

\begin{quote}
the quotient objective is easier to optimize.
\end{quote}

That does not follow automatically.

\section{Redundancy removed does not mean
convexity}\label{redundancy-removed-does-not-mean-convexity}

A quotient can still contain:

\begin{itemize}
\tightlist
\item
  multiple functionally distinct minima;
\item
  saddles;
\item
  barriers;
\item
  poor conditioning;
\item
  bifurcations;
\item
  singular regions.
\end{itemize}

Therefore

\[
\boxed{
\text{symmetry removed}
\not\Rightarrow
\text{convex landscape}.
}
\]

The quotient can be the more intrinsic object while remaining
mathematically difficult.

\section{Singular quotients}\label{singular-quotients}

A group action is easiest to quotient when it acts freely and regularly.

Neural networks can violate that simplicity.

Examples include:

\begin{itemize}
\tightlist
\item
  a hidden unit with zero incoming or outgoing weights;
\item
  two hidden units that become identical;
\item
  parameter points with extra stabilizer symmetries;
\item
  architecture-dependent hidden symmetries.
\end{itemize}

At such points, orbit dimension can change.

The quotient can develop singular structure.

This is why the chapter refuses to speak of ``the neural-network
quotient manifold'' without qualifications.

\section{Function space is an even stronger
quotient}\label{function-space-is-an-even-stronger-quotient}

One possible equivalence relation is

\[
\theta\sim\theta'
\iff
f_\theta=f_{\theta'}.
\]

This identifies every parameterization of the same function.

Conceptually, that is attractive.

Practically, it may be difficult to characterize.

The exact function-equivalence relation can include more than the
obvious architecture symmetries.

The quotient can also lose useful information about optimization
trajectories, regularization, or implementation cost.

The most intrinsic representation is not automatically the most useful
operational representation.

\section{Quotienting by the wrong
relation}\label{quotienting-by-the-wrong-relation}

Suppose we declare two models equivalent because they agree on a finite
training set.

That relation can merge functions that differ elsewhere.

If the intended object is global predictor behavior, the equivalence is
too coarse.

A quotient is only as meaningful as its equivalence relation.

This is the quotient analogue of the boundary-contract question:

\begin{quote}
sufficient for what?
\end{quote}

\section{Gauge language}\label{gauge-language}

Physics often uses the word \textbf{gauge} for representational freedom
that does not change observable quantities.

That language can be useful in machine learning when parameter
transformations leave the represented function invariant.

But it should be used carefully.

Not every neural symmetry has the full mathematical structure of a gauge
theory.

The safe statement is:

\begin{quote}
some neural parameter redundancies are gauge-like in the limited sense
that multiple representatives encode the same functional object.
\end{quote}

The exact group action should be stated.

\section{Representation equivalence
revisited}\label{representation-equivalence-revisited}

The Representation chapter introduced invertible recoding:

\[
\tilde z=Tz,
\]

with compatible downstream transformation.

Quotient thinking generalizes the lesson.

A representation is often meaningful only up to some family of
transformations.

Possible equivalences include:

\begin{itemize}
\tightlist
\item
  orthogonal basis change;
\item
  invertible linear recoding;
\item
  permutation;
\item
  positive scaling;
\item
  group action;
\item
  exact functional equivalence.
\end{itemize}

Different questions require different quotients.

There is no universal ``correct quotient'' for every analysis.

\section{The benign-nonconvexity
question}\label{the-benign-nonconvexity-question}

The Atlas now reaches a frontier question:

\begin{quote}
What portion of neural nonconvexity is merely redundancy induced by
representation or parameter symmetries, and what portion remains
intrinsic after those redundancies are removed?
\end{quote}

This is a central motivation for the constructive benign-nonconvexity
programme.

The hope is not that quotienting magically makes every problem convex.

The more modest possibility is:

\begin{itemize}
\tightlist
\item
  identify symmetry-induced degeneracy;
\item
  remove it explicitly;
\item
  measure the remaining landscape;
\item
  separate representational difficulty from intrinsic optimization
  difficulty.
\end{itemize}

That programme remains open.

\section{A reconstruction test}\label{a-reconstruction-test}

Suppose two parameter vectors are declared equivalent.

A useful test is:

\begin{quote}
Can every downstream property we intend to preserve be reconstructed
from the equivalence class?
\end{quote}

If not, the quotient discarded too much.

This reconstruction test connects quotient geometry to the Atlas's
broader Residual programme:

find the smallest structure that can change representation while
preserving the capability of interest.

\section{Six failure modes}\label{six-failure-modes}

\subsection{Parameter difference mistaken for functional
difference}\label{parameter-difference-mistaken-for-functional-difference}

Symmetry-related representatives are counted as distinct models.

\subsection{Quotient relation too
coarse}\label{quotient-relation-too-coarse}

Functionally relevant distinctions are erased.

\subsection{Quotient relation too
narrow}\label{quotient-relation-too-narrow}

Redundant coordinates remain and contaminate geometry.

\subsection{Regular quotient assumed
globally}\label{regular-quotient-assumed-globally}

Singular orbit structure is ignored.

\subsection{Flat direction assumed to be
symmetry}\label{flat-direction-assumed-to-be-symmetry}

Other sources of degeneracy are overlooked.

\subsection{Quotienting assumed to solve
optimization}\label{quotienting-assumed-to-solve-optimization}

Redundancy removal is promoted into an unsupported global simplicity
theorem.

\section{Atlas connections}\label{atlas-connections-8}

\textbf{Geometry.}\\
Grassmann is the canonical quotient example.

\textbf{Representation.}\\
Functional or task-relative equivalence can survive coordinate change.

\textbf{Normalized representations.}\\
Normalization removes one radial degree of freedom, but a sphere
cross-section and a quotient by positive scaling should be distinguished
mathematically.

\textbf{Optimization.}\\
Symmetry directions can create redundant ambient motion and local
degeneracy.

\textbf{Diagnostics.}\\
Comparing models may require symmetry alignment before interpreting
parameter distance.

\textbf{Residual.}\\
The next chapter asks what transferable structure survives
representation and parameter changes.

\textbf{Benign nonconvexity.}\\
Quotient structure provides a constructive route for separating
redundant from intrinsic landscape complexity.

\section{Closing view}\label{closing-view-11}

A parameter vector is often an address, not the place itself.

Two different addresses can point to one function.

Quotient geometry gives us a disciplined way to say which distinctions
are representational and which remain intrinsic.

That discipline is already valuable.

It prevents parameter distance from being mistaken for functional
distance.

It exposes exact symmetry directions.

It clarifies why some flatness is coordinate-dependent.

But the deeper question remains open.

After we remove the redundancy we understand, what difficulty is left?

That remainder is where the Atlas turns next.

\section{References used in this
chapter}\label{references-used-in-this-chapter-14}

\begin{itemize}
\tightlist
\item
  (Lee 2018)
\item
  (Absil et al. 2008)
\item
  (Edelman et al. 1998)
\item
  (Grigsby et al. 2023)
\item
  (Ziyin 2024)
\end{itemize}

See `sources/source-locks/ATLAS-CH-QUOTIENT-001.yaml` for exact source
roles and claim scope.

\chapter{The Residual}\label{the-residual}

\textbf{Epistemic status:} Atlas/GCL Research Synthesis + Established
Supporting Theory + Exact Toy Witness\\
\textbf{Specification:}
manuscript/specifications/ATLAS-CH-RESIDUAL-001.md\\
\textbf{Derivation packet:}
mathematics/derivations/ATLAS-CH-RESIDUAL-001-DERIVATIONS.md\\
\textbf{Computational witness:}
mathematics/computational-witnesses/ATLAS-CW-RESIDUAL-001.md\\
\textbf{Source lock:} sources/source-locks/ATLAS-CH-RESIDUAL-001.yaml

\section{What survives when representation
changes?}\label{what-survives-when-representation-changes}

Suppose two systems implement the same useful capability through very
different internal coordinates.

One uses dense vectors.

Another uses sparse features.

A third uses an operator basis.

A fourth has been permuted, rescaled, normalized, compressed, or
otherwise reparameterized.

What, if anything, must survive?

The Atlas uses the provisional name \textbf{Residual} for the answer we
are trying to find:

\begin{quote}
the smallest transferable structure sufficient to reconstruct a declared
capability after admissible changes of representation.
\end{quote}

The word is intentionally singular.

The theory is not.

There may be no unique universal Residual.

There may be no finite-dimensional one.

There may be no efficiently computable one.

The object can only be discussed after the capability, transformations,
and admissible descriptor class have been declared.

\section{The thing that cannot be transformed
away}\label{the-thing-that-cannot-be-transformed-away}

A useful GCL phrase is:

\begin{quote}
The Residual is the thing that cannot be transformed away.
\end{quote}

Taken literally, that phrase is too strong.

Many objects can survive a transformation while being irrelevant.

A constant survives almost everything.

What matters is not survival alone.

The proposed object must satisfy at least three demands:

\begin{enumerate}
\def\labelenumi{\arabic{enumi}.}
\tightlist
\item
  it survives the transformations we have declared irrelevant;
\item
  it still contains enough structure to reconstruct the capability we
  care about;
\item
  among admissible such descriptions, it retains no unnecessary
  distinctions.
\end{enumerate}

So the stronger working phrase is:

\begin{quote}
The Residual is the least admissible structure that survives the
declared representational changes while remaining sufficient to
reconstruct the declared capability.
\end{quote}

That is the object this chapter formalizes.

\section{The message that survives
translation}\label{the-message-that-survives-translation}

The pedagogical allegory is a message translated between languages.

The surface symbols change.

Word order may change.

The alphabet may disappear entirely.

Yet a receiver can still recover what is needed to act.

The structural correspondence is:

\begin{itemize}
\tightlist
\item
  language ↔ representation;
\item
  translation ↔ admissible representation change;
\item
  task-relevant meaning ↔ declared capability;
\item
  compact adequate message ↔ Residual.
\end{itemize}

The limit is important.

Human meaning is not a formal mathematical map.

Translation is rarely invertible.

``Shortest message'' depends on the coding language and admissible
operations.

The allegory teaches one point only:

\begin{quote}
what survives is meaningful here only because it remains
reconstructively sufficient for something declared in advance.
\end{quote}

\section{From representation equivalence to
reconstruction}\label{from-representation-equivalence-to-reconstruction}

The Representation chapter established that coordinate identity is
stronger than functional equivalence.

The Quotient chapter established that multiple parameter representatives
can encode one function.

Those results remove a common mistake:

\[
\text{different coordinates}
\not\Rightarrow
\text{different capability}.
\]

The Residual question goes further.

It asks:

\begin{quote}
Once representation-specific redundancy is removed, what portable
structure is still enough to rebuild the capability?
\end{quote}

This is not merely a quotient question.

A quotient tells us which distinctions we are willing to ignore.

A Residual must also pass a reconstruction test.

\section{Declaring the problem}\label{declaring-the-problem}

Let

\[
X
\]

be a representation or system-state space.

Let

\[
G
\]

be a declared family of admissible transformations acting on \(X\).

Let

\[
\mathcal Q
\]

be a declared family of capability probes, contexts, or interventions.

Let

\[
B:X\times\mathcal Q\to\mathcal Y
\]

be the declared behavior profile.

For each state \(x\), write

\[
B_x(q)=B(x,q).
\]

Finally, let

\[
\mathfrak D
\]

be a declared class of admissible portable descriptors.

This last object matters more than it first appears.

Without it, the problem can collapse into a tautology.

\section{Transformation invariance}\label{transformation-invariance}

A descriptor

\[
R:X\to\mathcal Z,
\qquad
R\in\mathfrak D,
\]

is invariant to the declared representation changes when

\[
\boxed{
R(g\cdot x)=R(x)
}
\]

for every admissible \(g\in G\).

Invariance is necessary for the intended notion of transfer.

It is not sufficient.

The constant descriptor

\[
R_0(x)=0
\]

is invariant under every transformation.

It usually preserves no useful capability.

\section{Capability sufficiency}\label{capability-sufficiency}

A descriptor \(R\) is sufficient for the declared behavior profile when
there exists a reconstructor

\[
D:\mathcal Z\times\mathcal Q\to\mathcal Y
\]

such that

\[
\boxed{
B(x,q)=D(R(x),q)
}
\]

for every declared state/probe pair.

This is capability reconstruction.

It is weaker than full-state reconstruction.

The Residual is allowed to forget anything the declared capability does
not need.

That forgetting is not a bug.

It is the point.

\section{The factorization
preorder}\label{the-factorization-preorder}

Suppose we have two descriptors,

\[
R_1:X\to Z_1,
\]

and

\[
R_2:X\to Z_2.
\]

Write

\[
R_1\preceq R_2
\]

when there exists an admissible post-processing map \(\phi\) from a
declared class \(\mathfrak M\), with

\[
\phi:Z_2\to Z_1
\]

such that

\[
\boxed{
R_1=\phi\circ R_2.
}
\]

Interpretation:

\(R_1\) contains no more distinctions than \(R_2\), because everything
in \(R_1\) can be recovered from \(R_2\).

This is a preorder when the declared post-processing class contains
identities and is closed under composition.

If \(\mathfrak M\) is the class of all set-theoretic maps, the preorder
is purely semantic.

If \(\mathfrak M\) is restricted to efficiently computable,
bounded-cost, differentiable, or otherwise admissible maps, the preorder
becomes operationally stronger.

The leastness argument below uses unrestricted set-theoretic
post-processing. If \(\mathfrak M\) is restricted, an abstract decoder
witnessing sufficiency need not belong to \(\mathfrak M\); it cannot be
used to certify operational leastness without that additional check.

Two different codings can factor through one another.

That is exactly what we want: coordinate choice should not automatically
create a new intrinsic object.

\section{Least, not merely small}\label{least-not-merely-small}

Let

\[
\mathcal S_{B,G,\mathfrak D}
\]

be the admissible descriptors that are both:

\begin{itemize}
\tightlist
\item
  invariant under \(G\);
\item
  sufficient for \(B\).
\end{itemize}

A provisional Residual is a \textbf{least element} of this class under
\(\preceq\).

Thus \(R\) must satisfy

\[
R\preceq S
\]

for every

\[
S\in\mathcal S_{B,G,\mathfrak D}.
\]

This is stronger than saying \(R\) is ``small.''

It is also stronger than saying no obviously redundant field remains.

Every other admissible invariant sufficient descriptor must be able to
produce \(R\).

\section{Why leastness is useful}\label{why-leastness-is-useful}

Suppose \(R\) and \(R'\) are both least.

Then

\[
R=\phi\circ R',
\]

and

\[
R'=\psi\circ R
\]

for some maps \(\phi,\psi\).

On the realized images,

\[
\psi\circ\phi
=
\operatorname{id},
\]

and

\[
\phi\circ\psi
=
\operatorname{id}.
\]

So two least Residuals are equivalent up to invertible recoding on what
the system actually realizes.

This is the kind of non-uniqueness the Atlas can tolerate.

Different coordinates need not mean different Residual structure.

\section{The behavior-table trap}\label{the-behavior-table-trap}

There is an immediate problem.

If \(\mathfrak D\) allows arbitrary descriptors, define

\[
R_B(x)=B_x.
\]

That is, store the complete behavior profile itself.

Then a decoder can simply evaluate

\[
R_B(x)(q).
\]

If \(B_x\) is invariant under the declared representation changes,
\(R_B\) is an invariant sufficient descriptor.

And because every sufficient descriptor must be able to reconstruct
\(B_x\), the full behavior profile is already least under factorization.

Mathematically, this is fine.

Scientifically, it can be useless.

We wanted transferable mechanism.

We got a behavior table.

\section{Descriptor admissibility is part of the theorem
statement}\label{descriptor-admissibility-is-part-of-the-theorem-statement}

A nontrivial Residual programme therefore has to declare what counts as
an admissible portable descriptor.

Possible restrictions include:

\begin{itemize}
\tightlist
\item
  bounded description length;
\item
  finite-dimensionality;
\item
  efficient computability;
\item
  architectural independence;
\item
  intervention stability;
\item
  bounded reconstruction cost;
\item
  portability across model families;
\item
  composability with a fixed interpreter.
\end{itemize}

The choice is not housekeeping.

It changes the object.

The Residual is therefore relative to

\[
(B,G,\mathfrak D,\mathcal Q).
\]

Suppressing those arguments may be convenient later.

It is dangerous at the beginning.

\section{Classical sufficiency as a
model}\label{classical-sufficiency-as-a-model}

Statistics gives a mature example of this kind of reasoning.

A sufficient statistic preserves all sample information needed for
inference about a declared parameter under a declared statistical model.

Minimal sufficiency then asks for a statistic that is no more
informative than any other sufficient statistic, in the appropriate
factorization sense (Lehmann and Casella 1998).

The Residual borrows this pattern:

\begin{itemize}
\tightlist
\item
  declare what must be preserved;
\item
  identify nuisance variation;
\item
  compress without losing what matters;
\item
  compare candidate descriptions by factorization.
\end{itemize}

The Residual is broader.

Therefore classical sufficiency theorems do not automatically apply.

\section{Minimal sufficient
representations}\label{minimal-sufficient-representations}

Achille and Soatto study task-relative representations in terms of:

\begin{itemize}
\tightlist
\item
  sufficiency;
\item
  minimality;
\item
  invariance to nuisance factors (Achille and Soatto 2018).
\end{itemize}

This is close in spirit to the Residual programme.

It supports a key lesson:

\begin{quote}
minimality is not meaningful without a declared task and nuisance
structure.
\end{quote}

The Atlas extends the question beyond one probabilistic
representation-learning setup toward transferable computational
structure.

That extension is ours.

It is not a theorem imported from the paper.

\section{Information bottleneck}\label{information-bottleneck}

The information bottleneck programme asks for compressed encodings that
preserve information relevant to a declared target (Tishby et al. 2000).

Again, the structural analogy is strong:

\[
\text{compress}
\quad\text{while preserving what matters}.
\]

But the Residual is not defined here as an information-bottleneck
optimum.

A Residual may be:

\begin{itemize}
\tightlist
\item
  symbolic;
\item
  algorithmic;
\item
  operator-valued;
\item
  graph-structured;
\item
  programmatic;
\item
  hybrid.
\end{itemize}

Mutual information is one possible formalization, not the definition.

\section{A calibration toy model}\label{a-calibration-toy-model}

Let

\[
X=\mathbb R^2,
\]

with state

\[
x=(s,n).
\]

Interpret:

\begin{itemize}
\tightlist
\item
  \(s\): task-relevant signal;
\item
  \(n\): nuisance.
\end{itemize}

Let admissible transformations translate only the nuisance coordinate:

\[
g_a(s,n)
=
(s,n+a).
\]

Use a singleton probe family for this calibration example and define
capability

\[
B(s,n)=s.
\]

Now define

\[
R(s,n)=s.
\]

This is deliberately simple.

It is a calibration case for the definitions, not evidence for a
frontier-model Residual.

\section{Invariance of the toy
Residual}\label{invariance-of-the-toy-residual}

For any \(a\),

\[
R(g_a(s,n))
=
R(s,n+a)
=
s.
\]

Therefore

\[
R(g_a x)=R(x).
\]

The Residual is constant along vertical nuisance orbits.

\section{Sufficiency of the toy
Residual}\label{sufficiency-of-the-toy-residual}

Choose decoder

\[
D(r)=r.
\]

Then

\[
D(R(s,n))
=
s
=
B(s,n).
\]

The Residual reconstructs the declared capability exactly.

It does not reconstruct \(n\).

It does not need to.

\section{Leastness in the toy
model}\label{leastness-in-the-toy-model}

Let \(S\) be any sufficient descriptor.

Then by sufficiency there exists a decoder \(D_S\) such that

\[
B=D_S\circ S.
\]

But in this toy model,

\[
R=B.
\]

Therefore

\[
R
=
D_S\circ S.
\]

Hence, for unrestricted semantic post-processing,

\[
R\preceq S.
\]

So \(R=s\) is least among the declared sufficient descriptors in the
semantic preorder. For an operational preorder, this argument
additionally requires \(D_S\in\mathfrak M\) for every compared
sufficient descriptor \(S\).

The proof is almost embarrassingly simple.

That is useful.

It shows exactly what the formal definition does before we ask it to
carry frontier complexity.

\section{Invariance without
sufficiency}\label{invariance-without-sufficiency}

Now define

\[
S_0(s,n)=0.
\]

This descriptor is invariant.

Take

\[
x_1=(2,0),
\qquad
x_2=(3,0).
\]

Then

\[
S_0(x_1)
=
S_0(x_2)
=
0.
\]

But

\[
B(x_1)=2,
\]

and

\[
B(x_2)=3.
\]

No deterministic decoder from the constant descriptor can produce both
capability values.

This is the simplest failure mode:

\[
\boxed{
\text{invariant}
\not\Rightarrow
\text{sufficient}.
}
\]

\section{Sufficiency without
minimality}\label{sufficiency-without-minimality}

Consider the full-state descriptor

\[
S_{\rm full}(s,n)
=
(s,n).
\]

It is sufficient for \(B=s\).

A decoder can discard \(n\).

But it stores more than the declared capability needs.

Indeed,

\[
R
=
\pi_1\circ S_{\rm full}.
\]

Full-state reconstruction is stronger than capability reconstruction.

The Residual programme explicitly does not demand the stronger object
unless the capability itself requires it.

\section{Representation change}\label{representation-change}

Now change coordinates.

Let

\[
T=
\begin{pmatrix}
1&1\\
1&-1
\end{pmatrix}.
\]

Define

\[
z=Tx.
\]

If

\[
x=(s,n),
\]

then

\[
z_1=s+n,
\]

and

\[
z_2=s-n.
\]

The original signal is no longer stored in one coordinate.

But it remains reconstructable:

\[
\boxed{
s
=
\frac{z_1+z_2}{2}.
}
\]

This is the first genuinely Residual-like lesson of the toy model.

The transferable object is not a coordinate slot.

It is reconstructable structure.

\section{Exact recoding witness}\label{exact-recoding-witness-1}

Take

\[
x=(2,5).
\]

Then

\[
z
=
T x
=
(7,-3).
\]

Reconstruct:

\[
\frac{7+(-3)}{2}=2.
\]

The exact computational witness also moves

\[
(2,5)
\]

along its nuisance orbit to

\[
(2,2),
\]

while preserving the Residual value \(2\).

\begin{figure}
\centering
\pandocbounded{\includegraphics[keepaspectratio,alt={Vertical nuisance orbits in signal-nuisance space collapse to their signal coordinate, while an invertible recoding sends x=(2,5) to z=(7,-3) and the Residual is reconstructed as (z1+z2)/2=2.}]{figures/masters/ATLAS-FIG-RESIDUAL-001.png}}
\caption{Vertical nuisance orbits in signal-nuisance space collapse to
their signal coordinate, while an invertible recoding sends x=(2,5) to
z=(7,-3) and the Residual is reconstructed as (z1+z2)/2=2.}
\end{figure}

The surface representation changed.

The declared capability did not.

\section{Task dependence}\label{task-dependence}

Now change the capability.

Suppose

\[
B'(s,n)=(s,n).
\]

Then

\[
R(s,n)=s
\]

is no longer sufficient.

The nuisance has stopped being nuisance.

This is why the Residual cannot be declared before the capability is
declared.

A structure can be residual for one task and inadequate for another.

\section{Transformation dependence}\label{transformation-dependence}

The same issue applies to \(G\).

Suppose we enlarge the admissible transformation family to permit
changing \(s\).

Then \(R=s\) is no longer invariant.

If the transformation family identifies capability-distinct states, the
problem itself becomes inconsistent:

no descriptor can be both invariant to that transformation and exactly
sufficient for the capability.

This gives a useful compatibility condition:

\begin{quote}
admissible transformations must preserve the declared capability if an
exact invariant sufficient Residual is to exist.
\end{quote}

\section{Residual as equivalence
class}\label{residual-as-equivalence-class}

Sometimes the natural Residual may be an equivalence class rather than a
chosen coordinate.

If

\[
x\sim x'
\]

whenever every declared capability probe gives the same result, then one
can consider the class

\[
[x].
\]

That quotient is canonical relative to the behavior relation.

But it can be enormous.

It may be impossible to compute.

It may still be a behavior-table object in disguise.

The research challenge is to find a smaller structural representative
that supports reconstruction.

\section{Residual as operator
structure}\label{residual-as-operator-structure}

The transferable object need not be a vector.

Suppose capability depends on a transformation rule rather than on a
stored state.

Then what survives may be:

\begin{itemize}
\tightlist
\item
  a linear operator;
\item
  an algebra of operators;
\item
  a transition kernel;
\item
  a control law;
\item
  a program;
\item
  a graph of conditional dependencies.
\end{itemize}

This is why the Atlas calls the target \textbf{computational structure}
rather than ``the minimal vector.''

\section{Residual as program}\label{residual-as-program}

Consider two neural systems that implement the same algorithm with
different hidden bases.

A basis-independent program or operator composition may transfer more
naturally than any activation vector.

This possibility connects the Residual to the Atlas shift:

\[
\text{vectors}
\to
\text{operators}.
\]

It also raises a harder question:

\begin{quote}
What is the right language in which minimality should be measured?
\end{quote}

Description length depends on representation language.

There is no free universal metric.

\section{Exact versus approximate
Residuals}\label{exact-versus-approximate-residuals}

Frontier systems rarely invite exact equivalence.

A practical Residual may need an error tolerance.

One can imagine requiring

\[
d_{\mathcal Y}
\big(
B(x,q),
D(R(x),q)
\big)
\le
\varepsilon
\]

for all probes in a declared test family, or in expectation under a
declared distribution.

Then minimality becomes approximate.

Tradeoffs appear among:

\begin{itemize}
\tightlist
\item
  descriptor size;
\item
  reconstruction error;
\item
  computation cost;
\item
  robustness;
\item
  transfer range.
\end{itemize}

The exact theory in this chapter is the zero-error baseline.

The approximate theory is a later research problem.

\section{Reconstruction is stronger than
correlation}\label{reconstruction-is-stronger-than-correlation}

Suppose a candidate feature correlates strongly with capability.

That is useful evidence.

It is not yet a Residual.

The candidate must support reconstruction under the declared changes of
representation.

A feature that predicts capability in one model but disappears under a
basis change fails transferability.

A feature that survives transformations but loses task-critical
distinctions fails sufficiency.

The Residual standard is deliberately stricter.

\section{Intervention strengthens the
test}\label{intervention-strengthens-the-test}

A compelling Residual candidate should survive more than passive
measurement.

Where possible:

\begin{enumerate}
\def\labelenumi{\arabic{enumi}.}
\tightlist
\item
  isolate the candidate structure;
\item
  alter non-Residual structure;
\item
  reconstruct or transplant the candidate;
\item
  test whether capability returns or remains;
\item
  repeat across admissible representations.
\end{enumerate}

This turns the question from:

\begin{quote}
Can we read the capability from this feature?
\end{quote}

into:

\begin{quote}
Is this structure enough to rebuild the capability under the
transformations we claim should not matter?
\end{quote}

That is closer to the intended programme.

\section{Finite replay can fool us}\label{finite-replay-can-fool-us}

Suppose a descriptor reproduces outputs on a finite benchmark.

That does not prove it reconstructs the capability.

It may simply encode the benchmark.

The declared probe family \(\mathcal Q\) therefore matters.

A serious reconstruction test must include enough variation to
distinguish:

\begin{itemize}
\tightlist
\item
  memorized outputs;
\item
  interpolated behavior;
\item
  genuine transferable mechanism.
\end{itemize}

The Replayable Evidence chapter tells us how to preserve the test.

It does not tell us that the test family was sufficient.

\section{Minimal curriculum}\label{minimal-curriculum}

The Residual idea connects directly to the minimal-curriculum programme.

If later capability can be efficiently reconstructed from a small
early-learned computational basis, then that basis is a candidate
Residual across training trajectories.

The question becomes:

\begin{quote}
Which early structures survive retraining, recoding, or architectural
change while remaining sufficient to rebuild later capability?
\end{quote}

That is a reconstruction problem.

\section{Minimal reasoning basis}\label{minimal-reasoning-basis}

The same pattern appears in mathematical reasoning.

Perhaps many surface reasoning traces can be generated from a smaller
set of transferable operations.

Then the target is not a canonical chain of thought.

It is a basis of operations sufficient to reconstruct broad reasoning
capability.

Again:

\begin{itemize}
\tightlist
\item
  invariance to surface trace;
\item
  sufficiency for capability;
\item
  leastness within an admissible operation language.
\end{itemize}

The Residual formalism gives the question a shape.

It does not answer it.

\section{Model interoperability}\label{model-interoperability}

Suppose two models use incompatible hidden coordinates.

Directly exchanging activations may fail.

A transferable Residual could act as a lingua franca if both systems
can:

\begin{itemize}
\tightlist
\item
  compile their internal state into the Residual;
\item
  reconstruct the required capability from it.
\end{itemize}

This is one possible route to model interoperability that does not
require one universal latent coordinate system.

\section{Memory}\label{memory}

The Residual also changes how we think about memory.

A memory store need not preserve every internal activation that produced
a result.

It may preserve only the structure required to reconstruct the
capability later.

This suggests a compression principle:

\begin{quote}
store what future reconstruction needs, not necessarily the historical
representation that happened to carry it.
\end{quote}

Whether such a sufficient Residual can be identified is the hard part.

\section{Benign nonconvexity}\label{benign-nonconvexity}

The Quotient chapter asked which nonconvexity disappears after symmetry
reduction.

The Residual asks an adjacent question:

\begin{quote}
Which computational distinctions survive after all
representation-specific redundancy is removed?
\end{quote}

If a capability can be reconstructed from a small quotient-stable
object, then large regions of parameter variation may be irrelevant to
that capability.

This could help separate:

\begin{itemize}
\tightlist
\item
  representational complexity;
\item
  intrinsic computational complexity.
\end{itemize}

That remains a research programme.

\section{Governed adaptation}\label{governed-adaptation}

A self-changing system needs to know what it must preserve.

The Residual offers one possible object for that preservation contract.

If a system can modify:

\begin{itemize}
\tightlist
\item
  architecture;
\item
  parameterization;
\item
  memory layout;
\item
  routing;
\item
  tool topology;
\end{itemize}

while retaining a certified Residual sufficient for declared
capabilities, then adaptation may become safer to reason about.

This is an aspiration.

No such universal certificate is established here.

\section{What would falsify a proposed
Residual?}\label{what-would-falsify-a-proposed-residual}

A candidate fails if any of the following occurs.

\subsection{Invariance failure}\label{invariance-failure}

An admissible representation change alters the descriptor.

\subsection{Sufficiency failure}\label{sufficiency-failure}

The declared capability cannot be reconstructed.

\subsection{Leastness failure}\label{leastness-failure}

A strictly coarser admissible invariant descriptor remains sufficient.

\subsection{Transfer failure}\label{transfer-failure}

The descriptor works only in one coordinate system or architecture
despite the claimed transfer range.

\subsection{Intervention failure}\label{intervention-failure}

Capability disappears when supposedly nonessential structure changes.

\subsection{Probe failure}\label{probe-failure}

The descriptor passes a narrow benchmark but fails on the declared
capability family.

These are operational failure modes.

They make the programme falsifiable.

\section{What the Residual is not}\label{what-the-residual-is-not}

The Residual is not automatically:

\begin{itemize}
\tightlist
\item
  a hidden neuron;
\item
  a feature direction;
\item
  a latent vector;
\item
  a quotient coordinate;
\item
  a sufficient statistic in the classical sense;
\item
  an information-bottleneck optimum;
\item
  a compressed checkpoint;
\item
  a behavior table;
\item
  a universal invariant of intelligence.
\end{itemize}

Any of these could instantiate the role in a particular setting.

None is the definition.

\section{Atlas connections}\label{atlas-connections-9}

\textbf{Representation.}\\
The Residual is defined only after representation equivalence and
task-relative structure are explicit.

\textbf{Quotient geometry.}\\
Symmetry reduction removes representational redundancy but does not by
itself identify sufficient reconstructive structure.

\textbf{Information.}\\
Sufficiency and bottleneck ideas provide formal analogies and possible
special cases.

\textbf{Memory.}\\
A Residual could be what a system stores when historical coordinates are
disposable.

\textbf{Interoperability.}\\
A shared Residual could support translation without a shared latent
basis.

\textbf{Minimal curriculum / reasoning basis.}\\
The target becomes a minimal transferable basis from which capability
can be reconstructed.

\textbf{Governed adaptation.}\\
The Residual may eventually become a preservation contract across
self-modification.

\section{Closing view}\label{closing-view-12}

The Residual programme begins with a refusal.

We refuse to identify capability with the coordinates that happen to
implement it today.

We also refuse the opposite mistake: calling anything invariant ``the
essence.''

The target has to survive.

It has to reconstruct.

And, within a declared admissible language, it has to be no richer than
necessary.

In the toy model, the answer is trivial:

\[
R(s,n)=s.
\]

In frontier systems, we do not know the answer.

That ignorance is the point of naming the problem.

The Residual is not yet a discovered object.

It is a reconstruction test for finding what, after all admissible
transformations have done their work, still has to remain.

\section{References used in this
chapter}\label{references-used-in-this-chapter-15}

\begin{itemize}
\tightlist
\item
  (Lehmann and Casella 1998)
\item
  (Achille and Soatto 2018)
\item
  (Tishby et al. 2000)
\end{itemize}

See `sources/source-locks/ATLAS-CH-RESIDUAL-001.yaml` for exact source
roles, audited Atlas dependencies, and claim boundaries.

\chapter{Reconstruction and Transfer}\label{reconstruction-and-transfer}

\textbf{Epistemic status:} Established Transfer Framework + Atlas
Synthesis + Exact Toy Witness\\
\textbf{Specification:}
manuscript/specifications/ATLAS-CH-TRANSFER-001.md\\
\textbf{Derivation packet:}
mathematics/derivations/ATLAS-CH-TRANSFER-001-DERIVATIONS.md\\
\textbf{Computational witness:}
mathematics/computational-witnesses/ATLAS-CW-TRANSFER-001.md\\
\textbf{Source lock:} sources/source-locks/ATLAS-CH-TRANSFER-001.yaml

\section{Transfer is not reuse}\label{transfer-is-not-reuse}

A model begins on one task.

Some of what it learned is reused somewhere else.

That story is often called transfer learning.

But several different phenomena can hide inside the word
\textbf{transfer}:

\begin{itemize}
\tightlist
\item
  reusing parameters as initialization;
\item
  freezing a feature extractor;
\item
  fine-tuning a whole model;
\item
  adapting to a new domain;
\item
  adapting to a new task;
\item
  reusing architecture but not weights;
\item
  compiling one representation into another;
\item
  reconstructing a capability through a shared latent object.
\end{itemize}

These are not the same operation.

The Atlas therefore asks a narrower question before discussing methods:

\begin{quote}
What structure must remain recoverable for a declared target capability
to survive a change of representation, task, or domain?
\end{quote}

This chapter connects that question to the Residual.

\section{Moving a machine through a narrow
doorway}\label{moving-a-machine-through-a-narrow-doorway}

A useful allegory is shipping a machine through a doorway.

One option is to push the whole assembled machine through.

Another is to:

\begin{enumerate}
\def\labelenumi{\arabic{enumi}.}
\tightlist
\item
  disassemble it;
\item
  keep a portable core;
\item
  carry that core across;
\item
  rebuild what is needed on the other side.
\end{enumerate}

The structural correspondence is:

\begin{itemize}
\tightlist
\item
  source representation ↔ original assembly;
\item
  compiler ↔ disassembly map;
\item
  Residual ↔ portable core;
\item
  target decoder ↔ reconstruction procedure;
\item
  bottleneck ↔ doorway width.
\end{itemize}

The limit matters.

Real transfer learning also involves:

\begin{itemize}
\tightlist
\item
  finite data;
\item
  optimization;
\item
  distribution shift;
\item
  task mismatch;
\item
  decoder restrictions;
\item
  stochastic training.
\end{itemize}

Exact deterministic reconstruction is only one clean special case.

\section{Standard source and target
framing}\label{standard-source-and-target-framing}

A common transfer-learning formulation distinguishes a \textbf{domain}
from a \textbf{task} (Pan and Yang 2010).

A domain can be written as

\[
\mathcal D
=
(\mathcal X,P(X)),
\]

where \(\mathcal X\) is an input space and \(P(X)\) a distribution.

A task can be written as

\[
\mathcal T
=
(\mathcal Y,f),
\]

where \(\mathcal Y\) is an output space and \(f\) a predictive object.

A transfer setting then relates:

\[
(\mathcal D_s,\mathcal T_s)
\]

to

\[
(\mathcal D_t,\mathcal T_t).
\]

The source and target may differ in:

\begin{itemize}
\tightlist
\item
  input distribution;
\item
  feature space;
\item
  output space;
\item
  predictive task;
\item
  available labels;
\item
  data volume.
\end{itemize}

The important point is relational:

transfer is always transfer \textbf{from something to something}.

\section{Transfer needs a target
criterion}\label{transfer-needs-a-target-criterion}

Suppose we reuse a source model on a target task.

Did transfer succeed?

That question is meaningless without a target criterion.

We need at least:

\begin{itemize}
\tightlist
\item
  a target metric;
\item
  a target data budget;
\item
  an adaptation procedure;
\item
  a baseline.
\end{itemize}

For a higher-is-better score \(J\), transfer improves the target when

\[
J(M_{\rm transfer})
>
J(M_{\rm target-only}),
\]

under a matched comparison.

If

\[
J(M_{\rm transfer})
<
J(M_{\rm target-only}),
\]

the result is negative transfer.

The sign depends on the metric convention.

The concept depends on the baseline.

\section{Negative transfer is not
paradoxical}\label{negative-transfer-is-not-paradoxical}

Why can previously learned structure hurt?

Because reusable structure can be mismatched.

Possible causes include:

\begin{itemize}
\tightlist
\item
  source-specific features;
\item
  wrong invariances;
\item
  optimization bias;
\item
  insufficient adaptation;
\item
  domain shift;
\item
  label mismatch;
\item
  representational bottlenecks;
\item
  co-adaptation among features.
\end{itemize}

Pan and Yang discuss negative transfer as a central issue in transfer
learning (Pan and Yang 2010).

The Atlas will use the term operationally, not metaphysically.

A source model does not carry a universal transferable essence merely
because it learned something useful once.

\section{Source accuracy is not
transferability}\label{source-accuracy-is-not-transferability}

A source model can perform extremely well on its original task while
encoding features that transfer poorly.

Yosinski and colleagues experimentally studied how feature
transferability changes across network depth and task distance (Yosinski
et al. 2014).

Their results show at least two distinct limitations:

\begin{itemize}
\tightlist
\item
  higher-layer specialization can reduce transferability;
\item
  splitting co-adapted representations can create optimization
  difficulties.
\end{itemize}

Kornblith, Shlens, and Le later showed that source benchmark quality and
transferable-feature quality are not interchangeable, and that training
choices can substantially affect transfer behavior (Kornblith et al.
2019).

The bounded conclusion is:

\[
\boxed{
\text{source performance}
\not\Rightarrow
\text{universal transfer quality}.
}
\]

\section{Parameter reuse is weaker than capability
transfer}\label{parameter-reuse-is-weaker-than-capability-transfer}

Suppose we copy source weights into a target model.

That establishes one fact:

the target optimization began from source parameters.

It does not establish:

\begin{itemize}
\tightlist
\item
  which source capabilities survived;
\item
  which features were reused;
\item
  whether the target behavior improved;
\item
  whether the target could have learned the same behavior from scratch;
\item
  whether the transferred parameters were actually necessary.
\end{itemize}

Parameter reuse is a mechanism.

Transfer is an outcome relative to a declared target criterion.

\section{Representation transfer}\label{representation-transfer}

A common strategy is to reuse an intermediate representation

\[
z=E_s(x)
\]

and learn a new target readout.

This works only if the representation retains distinctions the target
task needs.

A feature extractor can be highly sufficient for the source task and
still erase information needed by the target.

This is the direct bridge to the Residual formalism.

The target asks:

\begin{quote}
Is the transferred representation still sufficient for the capability I
now want?
\end{quote}

\section{The common-Residual reconstruction
certificate}\label{the-common-residual-reconstruction-certificate}

Let

\[
E_s:X\to Z_s
\]

be a source representation.

Let

\[
E_t:X\to Z_t
\]

be a target representation.

Let

\[
R:X\to\mathcal R
\]

be a declared admissible Residual.

Suppose there exist compilers

\[
C_s:Z_s\to\mathcal R,
\]

and

\[
C_t:Z_t\to\mathcal R,
\]

such that

\[
\boxed{
C_s(E_s(x))
=
C_t(E_t(x))
=
R(x)
}
\]

for every declared state \(x\).

Now let the target capability family be indexed by \(q\in\mathcal Q_t\).

Suppose each target behavior factors through \(R\):

\[
\boxed{
B_t(x,q)
=
D_q(R(x)).
}
\]

Then the target capability can be reconstructed from either
representation:

\[
B_t(x,q)
=
D_q(C_s(E_s(x))),
\]

and

\[
B_t(x,q)
=
D_q(C_t(E_t(x))).
\]

This is the \textbf{common-Residual reconstruction certificate}.

\section{What the certificate says}\label{what-the-certificate-says}

The certificate says:

\begin{quote}
source and target coordinates may differ, but if both compile to the
same capability-sufficient structure, the declared capability survives.
\end{quote}

This is a sufficient condition.

It is not a characterization of all transfer.

A direct adapter

\[
A:Z_s\to Z_t
\]

might exist even when we have not identified a shared Residual.

The common Residual is useful because it separates:

\begin{itemize}
\tightlist
\item
  representation-specific coordinates;
\item
  transferable capability structure.
\end{itemize}

\section{Exact linear witness}\label{exact-linear-witness}

Let the transferable structure be

\[
r=
\begin{pmatrix}
a\\
b
\end{pmatrix}.
\]

The source representation is

\[
z_s
=
A_s r,
\]

with

\[
A_s
=
\begin{pmatrix}
1&1\\
1&-1
\end{pmatrix}.
\]

The target representation is

\[
z_t
=
A_t r,
\]

with

\[
A_t
=
\begin{pmatrix}
2&0\\
0&1/2
\end{pmatrix}.
\]

Both matrices are invertible.

Therefore define

\[
C_s=A_s^{-1},
\]

and

\[
C_t=A_t^{-1}.
\]

Both recover the same \(r\).

\section{Exact numeric
reconstruction}\label{exact-numeric-reconstruction}

Take

\[
r=(3,1).
\]

Then

\[
z_s
=
A_s r
=
(4,2).
\]

And

\[
z_t
=
A_t r
=
(6,1/2).
\]

The exact witness verifies:

\[
C_s z_s=(3,1),
\]

and

\[
C_t z_t=(3,1).
\]

Different coordinates.

Same reconstructive core.

\section{Two target tasks}\label{two-target-tasks}

Now define two target tasks:

\[
D_1(a,b)=a+b,
\]

and

\[
D_2(a,b)=2a-b.
\]

For

\[
r=(3,1),
\]

we obtain

\[
D_1=4,
\]

and

\[
D_2=5.
\]

Both source and target representations reconstruct both outputs exactly
through the same \(r\).

\begin{figure}
\centering
\pandocbounded{\includegraphics[keepaspectratio,alt={Two different representations compile to the same exact Residual r=(3,1) and reconstruct target outputs 4 and 5; a second panel shows a lossy bottleneck P(r)=1 merging two states that require different D2 outputs.}]{figures/masters/ATLAS-FIG-TRANSFER-001.png}}
\caption{Two different representations compile to the same exact
Residual r=(3,1) and reconstruct target outputs 4 and 5; a second panel
shows a lossy bottleneck P(r)=1 merging two states that require
different D2 outputs.}
\end{figure}

\section{Exact recoding is the easy
case}\label{exact-recoding-is-the-easy-case}

In the witness above, both representations are invertible transforms of
the Residual.

No information is lost.

This is the easiest kind of transfer.

It isolates one concept:

\[
\text{coordinate change}
\neq
\text{capability loss}.
\]

But exact invertibility is not required in general.

A target representation can discard information as long as it preserves
everything needed for the declared target task family.

That brings us to bottlenecks.

\section{A lossy bottleneck}\label{a-lossy-bottleneck}

Define

\[
P(a,b)=a.
\]

The second coordinate is discarded.

This can still be sufficient for tasks that depend only on \(a\).

For example,

\[
D_0(a,b)=3a
\]

can be reconstructed from \(P\).

But the same bottleneck can fail for another task.

\section{Exact bottleneck failure}\label{exact-bottleneck-failure}

Take

\[
r_A=(1,0),
\]

and

\[
r_B=(1,1).
\]

Then

\[
P(r_A)=1,
\]

and

\[
P(r_B)=1.
\]

The bottleneck makes the states identical.

Now evaluate

\[
D_2(a,b)=2a-b.
\]

We get

\[
D_2(r_A)=2,
\]

and

\[
D_2(r_B)=1.
\]

One bottleneck value would have to decode to two different outputs.

That is impossible for a deterministic decoder.

\section{The fiber criterion}\label{the-fiber-criterion}

The bottleneck example is a special case of a general exact criterion.

Let

\[
Z:X\to\mathcal Z
\]

be a representation.

Let

\[
B:X\to\mathcal Y
\]

be a deterministic capability map.

A deterministic decoder

\[
D:\mathcal Z\to\mathcal Y
\]

with

\[
B=D\circ Z
\]

exists if and only if \(B\) is constant on every fiber of \(Z\).

That is:

\[
\boxed{
Z(x_1)=Z(x_2)
\Rightarrow
B(x_1)=B(x_2).
}
\]

This is the exact reconstruction test.

\section{Why the fiber criterion is
necessary}\label{why-the-fiber-criterion-is-necessary}

Suppose

\[
B=D\circ Z.
\]

If

\[
Z(x_1)=Z(x_2),
\]

then

\[
B(x_1)
=
D(Z(x_1))
=
D(Z(x_2))
=
B(x_2).
\]

So any representation that merges capability-distinct states cannot
support exact deterministic reconstruction.

This is information loss in the most concrete possible sense.

\section{Why the fiber criterion is
sufficient}\label{why-the-fiber-criterion-is-sufficient}

Now suppose \(B\) is constant on each realized fiber of \(Z\).

For each realized representation value \(z\), choose any state \(x\)
with

\[
Z(x)=z.
\]

Define

\[
D(z)=B(x).
\]

This is well-defined because every state in the same fiber has the same
\(B\)-value.

Therefore

\[
B=D\circ Z.
\]

This is a semantic existence result.

It does not guarantee that the decoder is simple or efficiently
learnable.

\section{One-task sufficiency is not task-family
sufficiency}\label{one-task-sufficiency-is-not-task-family-sufficiency}

A representation may preserve exactly the information needed for one
task and destroy information needed by another.

The bottleneck

\[
P(a,b)=a
\]

is sufficient for any target depending only on \(a\).

It is insufficient for

\[
D_2(a,b)=2a-b.
\]

Therefore transferability should often be indexed by a task family:

\[
\mathcal Q_t.
\]

A representation that transfers to task \(q_1\) need not transfer to
\(q_2\).

There is no contradiction.

The target changed.

\section{Universal representation quality is
underspecified}\label{universal-representation-quality-is-underspecified}

People often ask whether a representation is ``good.''

Good for what?

Possible criteria include:

\begin{itemize}
\tightlist
\item
  linear separability;
\item
  reconstruction;
\item
  invariance;
\item
  robustness;
\item
  compression;
\item
  adaptation speed;
\item
  transfer across tasks;
\item
  transfer across domains.
\end{itemize}

These objectives can disagree.

The Transfer chapter therefore refuses to assign one universal scalar
quality to a representation without a declared target family and
adaptation setting.

\section{Information bottleneck
connection}\label{information-bottleneck-connection}

The Information Bottleneck programme asks for compression while
preserving information relevant to a declared target (Tishby et al.
2000).

That is structurally close to transfer through a bottleneck.

A representation can discard source variation while retaining
target-relevant structure.

But mutual information is not the definition of transfer.

Transfer also depends on:

\begin{itemize}
\tightlist
\item
  decoder class;
\item
  optimization;
\item
  sample size;
\item
  shift;
\item
  task geometry;
\item
  adaptation cost.
\end{itemize}

Information preservation is one piece.

\section{Minimal sufficient representation
connection}\label{minimal-sufficient-representation-connection}

Achille and Soatto frame representations through task sufficiency,
minimality, and nuisance invariance (Achille and Soatto 2018).

The Transfer chapter inherits a compatible lesson:

\begin{quote}
compressing away nuisance can help only if the target does not later
require what was removed.
\end{quote}

The target task determines which variation is nuisance.

A source-task nuisance can become a target-task signal.

\section{Transfer across domains}\label{transfer-across-domains}

Suppose source and target use the same task but different input
distributions:

\[
P_s(X)\neq P_t(X).
\]

A source representation can preserve all distinctions needed on the
source support while behaving poorly on target regions.

The exact fiber criterion must therefore be evaluated over the declared
target domain.

A certificate over the wrong support is not a reconstruction
certificate.

This is one reason domain shift matters.

\section{Transfer across tasks}\label{transfer-across-tasks}

Suppose source and target share the same domain but differ in target
function.

A source representation can be sufficient for

\[
f_s
\]

but insufficient for

\[
f_t.
\]

This is exactly what the bottleneck witness demonstrates in miniature.

Task change can turn discarded variation into required information.

\section{Transfer across representation
systems}\label{transfer-across-representation-systems}

The Residual framing is most distinctive when the source and target do
not share coordinates.

One model may encode a capability in:

\begin{itemize}
\tightlist
\item
  a dense hidden state;
\item
  a sparse dictionary;
\item
  an operator basis;
\item
  a normalized manifold state;
\item
  external memory.
\end{itemize}

If both can compile to the same admissible Residual, transfer does not
require one universal latent basis.

This is the bridge to model interoperability.

\section{Exact information preservation does not imply easy
learning}\label{exact-information-preservation-does-not-imply-easy-learning}

Suppose source and target representations are related by an invertible
matrix.

No information is lost.

But the adapter may still be hard to learn because of:

\begin{itemize}
\tightlist
\item
  poor conditioning;
\item
  limited samples;
\item
  nonlinear parameterization;
\item
  restricted decoder class;
\item
  optimization failure.
\end{itemize}

Therefore:

\[
\boxed{
\text{semantic transfer path}
\neq
\text{guaranteed learnable adapter}.
}
\]

The Atlas keeps semantic sufficiency and computational attainability
separate.

\section{Transfer can fail at the
interface}\label{transfer-can-fail-at-the-interface}

Suppose a source representation contains the required information.

Suppose the target decoder class could in principle use it.

Transfer can still fail if the interface contract is wrong.

Examples include:

\begin{itemize}
\tightlist
\item
  unit mismatch;
\item
  scaling mismatch;
\item
  permutation mismatch;
\item
  unsupported uncertainty semantics;
\item
  precision loss;
\item
  stale memory schema.
\end{itemize}

Transfer is therefore also a composition problem.

The Boundary Contracts chapter will later make this more explicit.

\section{Architecture transfer versus feature
transfer}\label{architecture-transfer-versus-feature-transfer}

Kornblith and colleagues provide a useful empirical warning (Kornblith
et al. 2019).

An architecture can transfer well even when a particular set of
source-trained features is less transferable than source performance
might suggest.

This separates two objects:

\begin{itemize}
\tightlist
\item
  reusable architecture;
\item
  reusable representation state.
\end{itemize}

The Atlas treats them differently.

Architecture is a transformation/composition template.

Features are instantiated states.

Transfer can occur through either.

\section{Fine-tuning changes the
object}\label{fine-tuning-changes-the-object}

When a transferred representation is fine-tuned, it does not simply
``remain transferable.''

It becomes a new target-adapted state.

The causal story now includes:

\begin{itemize}
\tightlist
\item
  initialization;
\item
  optimization path;
\item
  target data;
\item
  regularization;
\item
  adaptation schedule.
\end{itemize}

A post-fine-tuning representation cannot be treated as unchanged source
structure.

Transfer attribution requires care.

\section{A reconstruction certificate is not a training
algorithm}\label{a-reconstruction-certificate-is-not-a-training-algorithm}

The common-Residual certificate proves that exact target reconstruction
is possible under declared maps.

It does not tell us how to learn:

\[
C_s,
\qquad
C_t,
\qquad
D_q.
\]

Those may be known analytically in a toy setting and unknown in
practice.

A future Atlas programme can distinguish:

\begin{itemize}
\tightlist
\item
  \textbf{existence of transfer path};
\item
  \textbf{identifiability of transfer path};
\item
  \textbf{learnability of transfer path};
\item
  \textbf{efficiency of transfer path}.
\end{itemize}

These are different questions.

\section{Reconstruction tests}\label{reconstruction-tests}

A practical transfer claim can be stress-tested with reconstruction.

Given source and target representations:

\begin{enumerate}
\def\labelenumi{\arabic{enumi}.}
\tightlist
\item
  declare the target capability family;
\item
  declare admissible compilers/adapters;
\item
  attempt to reconstruct capability from each representation;
\item
  search for target states collapsed by the representation;
\item
  compare transfer against a matched target-only baseline;
\item
  intervene on supposedly irrelevant source structure;
\item
  repeat under domain and representation changes.
\end{enumerate}

Failure at different stages means different things.

\section{What would falsify a transfer
claim?}\label{what-would-falsify-a-transfer-claim}

\subsection{Sufficiency failure}\label{sufficiency-failure-1}

Two target states share the same transferred representation but require
different outputs.

\subsection{Domain failure}\label{domain-failure}

A certificate holds on source support but not target support.

\subsection{Adapter failure}\label{adapter-failure}

An admissible adapter class cannot recover the required object.

\subsection{Baseline failure}\label{baseline-failure}

Transfer underperforms a matched target-only baseline.

\subsection{Robustness failure}\label{robustness-failure}

The transferred capability disappears under allowed representation
changes.

\subsection{Attribution failure}\label{attribution-failure}

Target performance comes from relearning rather than transferred
structure.

These failure modes should not be collapsed into one number.

\section{Transferable Residual versus universal
Residual}\label{transferable-residual-versus-universal-residual}

The Residual chapter was task-relative.

Transfer does not remove that relativity.

A shared Residual sufficient for one target family may fail for another.

The useful object is therefore:

\[
R_{\mathcal Q}
\]

for a declared family of capability probes \(\mathcal Q\), not an
assumed universal essence of the model.

The search for broader Residuals is an empirical and mathematical
programme.

It is not a free theorem.

\section{Minimal curricula}\label{minimal-curricula}

Transfer connects directly to curriculum design.

If an early mechanism can reconstruct many later capabilities with
limited adaptation, it has high transfer value.

A minimal curriculum therefore asks:

\begin{quote}
Which learned structures create the broadest reconstructive reuse across
later tasks?
\end{quote}

This is stronger than asking which examples reduce next-step training
loss fastest.

The later Minimal Curricula chapter will combine this transfer criterion
with learning-progress search.

\section{Minimal reasoning basis}\label{minimal-reasoning-basis-1}

A reasoning operation is transferable if it can be reused across problem
families without carrying problem-specific surface structure.

The Residual/Transfer pair suggests a test:

\begin{itemize}
\tightlist
\item
  identify a candidate reasoning operation;
\item
  recode the problem;
\item
  change surface notation;
\item
  change task family;
\item
  test whether the operation still supports reconstruction of successful
  reasoning behavior.
\end{itemize}

This is how ``canonical reasoning transformations'' can become
falsifiable rather than rhetorical.

\section{Model interoperability}\label{model-interoperability-1}

Two models need not share hidden coordinates if both can compile to a
common transferable object.

This suggests an interoperability architecture:

\[
\text{Model A state}
\to
R
\to
\text{Model B capability}.
\]

The hard questions are:

\begin{itemize}
\tightlist
\item
  what is \(R\)?
\item
  which tasks does it preserve?
\item
  how costly are the compilers?
\item
  what semantics must the interface expose?
\item
  how do we certify reconstruction?
\end{itemize}

Transfer becomes a systems problem.

\section{Memory and transfer}\label{memory-and-transfer}

Persistent memory is another transfer surface.

A model may write an object today and another model may consume it
later.

If memory stores representation-specific coordinates, the second model
may not understand them.

If memory stores a transferable Residual or sufficient evidence object,
reconstruction can survive model change.

This is one reason external memory should preserve semantics and
provenance rather than raw activations alone.

\section{Compression as a transfer
probe}\label{compression-as-a-transfer-probe}

A bottleneck tests what can be discarded.

If a compressed object still transfers across many target tasks, it may
have isolated reusable structure.

If transfer collapses after one additional compression step, that step
removed capability-relevant information.

Compression can therefore become an empirical probe for transfer
structure.

This connects directly to the later Compression as Discovery chapter.

\section{Closing view}\label{closing-view-13}

Transfer learning is often described as reuse.

The Atlas asks a stricter question.

What, exactly, survived?

A copied parameter is not yet evidence.

A reused architecture is not the same thing as a reused feature.

A high source score is not a transfer guarantee.

A common Residual gives one clean sufficient answer:

\[
\text{different representation}
\to
\text{same reconstructive core}
\to
\text{target capability}.
\]

And the fiber criterion gives the corresponding impossibility test:

if a representation merges states the target must distinguish, exact
reconstruction is gone.

Transfer therefore sits at the intersection of:

\begin{itemize}
\tightlist
\item
  representation;
\item
  information;
\item
  reconstruction;
\item
  optimization;
\item
  composition.
\end{itemize}

The practical question is not whether knowledge was reused in some vague
sense.

It is whether the structure that crossed the boundary was sufficient for
what the target still needed to do.

\section{References used in this
chapter}\label{references-used-in-this-chapter-16}

\begin{itemize}
\tightlist
\item
  (Pan and Yang 2010)
\item
  (Yosinski et al. 2014)
\item
  (Kornblith et al. 2019)
\item
  (Achille and Soatto 2018)
\item
  (Tishby et al. 2000)
\end{itemize}

See `sources/source-locks/ATLAS-CH-TRANSFER-001.yaml` for exact source
roles, Atlas dependencies, and claim boundaries.

\chapter{From Layered Networks to Residual
Systems}\label{from-layered-networks-to-residual-systems}

\textbf{Epistemic status:} Primary Architecture History + Established
Structural Mathematics + Atlas Synthesis\\
\textbf{Specification:}
manuscript/specifications/ATLAS-CH-ARCHHIST-001.md\\
\textbf{Derivation packet:}
mathematics/derivations/ATLAS-CH-ARCHHIST-001-DERIVATIONS.md\\
\textbf{Computational witness:}
mathematics/computational-witnesses/ATLAS-CW-ARCHHIST-001.md\\
\textbf{Source lock:} sources/source-locks/ATLAS-CH-ARCHHIST-001.yaml

\section{Architecture is a choice of mathematical
object}\label{architecture-is-a-choice-of-mathematical-object}

A neural architecture is not merely a list of layers.

It declares:

\begin{itemize}
\tightlist
\item
  what counts as state;
\item
  which transformations are shared;
\item
  what information persists;
\item
  where bottlenecks live;
\item
  how depth transports information;
\item
  which symmetries are built in.
\end{itemize}

The useful historical question is therefore not:

\begin{quote}
Which architecture came next?
\end{quote}

It is:

\begin{quote}
What new mathematical object became worth analyzing?
\end{quote}

This chapter follows a selective lineage:

\[
\text{composition}
\to
\text{structured operators}
\to
\text{persistent state}
\to
\text{interfaces}
\to
\text{gated carry}
\to
\text{residual transport}
\to
\text{explicit continuous depth}.
\]

That sequence is pedagogical.

It is not a claim that the actual history of neural networks was linear.

\section{From assembly line to evolving
state}\label{from-assembly-line-to-evolving-state}

A useful allegory is an assembly line.

A feed-forward network moves an object through a sequence of stations.

A convolutional network reuses the same local tool at many spatial
locations.

A recurrent network gives the object a notebook that persists through
computational time.

An encoder--decoder system forces two subsystems to communicate through
an interface.

A highway layer adds learned gates deciding how much to transform and
how much to carry.

A residual block carries the current state forward and adds a
correction.

A continuous-depth model replaces the discrete sequence with an
explicitly parameterized evolution law.

The limit of the allegory is important.

Neural networks are not factories.

The value of the metaphor is only that it makes the change in
mathematical object visible.

\section{Layered composition}\label{layered-composition}

A feed-forward network can be written as

\[
x_{k+1}
=
F_k(x_k).
\]

After \(L\) layers,

\[
\boxed{
x_L
=
F_{L-1}\circ\cdots\circ F_0(x_0).
}
\]

The basic object is therefore a composition of learned maps.

This viewpoint is already enough to introduce:

\begin{itemize}
\tightlist
\item
  hidden representations;
\item
  depth;
\item
  nonlinear composition;
\item
  back-propagated derivatives.
\end{itemize}

Rumelhart, Hinton, and Williams gave one influential formulation of
efficient error back-propagation through multilayer networks (Rumelhart
et al. 1986).

The historical claim here is deliberately narrow.

The chapter does not assign unique invention of multilayer learning to
one paper.

\section{Composition creates a product of
Jacobians}\label{composition-creates-a-product-of-jacobians}

If each \(F_k\) is differentiable, the chain rule gives

\[
J_{\rm total}
=
J_{F_{L-1}}(x_{L-1})
\cdots
J_{F_0}(x_0).
\]

This simple product has enormous consequences.

Signals and gradients can be:

\begin{itemize}
\tightlist
\item
  amplified;
\item
  attenuated;
\item
  rotated;
\item
  collapsed;
\item
  made ill-conditioned.
\end{itemize}

The architecture therefore determines not only which functions can be
represented.

It also shapes how derivatives are transported through depth.

\section{The early object: a hierarchy of
representations}\label{the-early-object-a-hierarchy-of-representations}

In a layered network, each hidden state

\[
x_k
\]

is both:

\begin{itemize}
\tightlist
\item
  the output of earlier computation;
\item
  the input to later computation.
\end{itemize}

This creates a hierarchy of internal representations.

The Representation chapter taught us not to assume that a hidden
coordinate is intrinsically meaningful.

Architecture adds a further question:

\begin{quote}
Which transformations are imposed before that representation even has a
chance to form?
\end{quote}

Convolution provides the first major structural answer in this lineage.

\section{Convolution replaces arbitrary connectivity with a
structured
operator}\label{convolution-replaces-arbitrary-connectivity-with-a-structured-operator}

A dense linear map can connect every input coordinate to every output
coordinate.

A convolutional map ties weights across locations.

For a periodic one-dimensional signal, define

\[
(C_kx)_j
=
\sum_r k_r x_{j-r}.
\]

The same kernel coefficients

\[
k_r
\]

are reused across positions.

LeCun and colleagues used local receptive fields and weight sharing in a
highly influential convolutional architecture for document recognition
(LeCun et al. 1998).

The structural lesson is broader than the application:

\begin{quote}
architecture can encode a symmetry assumption through operator
structure.
\end{quote}

\section{Exact translation equivariance in the declared
model}\label{exact-translation-equivariance-in-the-declared-model}

Let cyclic shift \(S_m\) act by

\[
(S_mx)_j
=
x_{j-m}.
\]

Then

\[
\begin{aligned}
(C_kS_mx)_j
&=
\sum_r k_r(S_mx)_{j-r}\\
&=
\sum_r k_r x_{j-r-m}\\
&=
(S_mC_kx)_j.
\end{aligned}
\]

Therefore

\[
\boxed{
C_kS_m
=
S_mC_k.
}
\]

This is exact translation equivariance for the declared
circular-convolution model.

The architecture is not learning this symmetry from scratch.

The operator class already contains it.

\section{Exact convolution witness}\label{exact-convolution-witness}

Use

\[
S=
\begin{pmatrix}
0&0&0&1\\
1&0&0&0\\
0&1&0&0\\
0&0&1&0
\end{pmatrix},
\]

and

\[
C=I+2S.
\]

Because \(C\) is a polynomial in \(S\),

\[
CS=SC.
\]

For

\[
x=(1,2,3,4),
\]

the exact computational witness gives

\[
CSx
=
SCx
=
(10,9,4,7).
\]

The symmetry is algebraic in this toy model.

It is not an empirical approximation.

\section{Practical CNNs can break exact
equivariance}\label{practical-cnns-can-break-exact-equivariance}

The exact proof above is intentionally narrow.

Real pipelines may include:

\begin{itemize}
\tightlist
\item
  zero padding;
\item
  asymmetric boundaries;
\item
  striding;
\item
  pooling;
\item
  resizing;
\item
  nonlinear preprocessing;
\item
  position-dependent operations.
\end{itemize}

Any of these can alter exact equivariance.

So the safe statement is:

\begin{quote}
convolution supplies a translation-structured operator, while
full-pipeline equivariance must be checked operation by operation.
\end{quote}

This distinction matters later when positional information is added
deliberately.

\section{Recurrent systems change the object
again}\label{recurrent-systems-change-the-object-again}

A feed-forward layer consumes a state and passes a new state onward.

A recurrent system makes state persistence explicit:

\[
\boxed{
h_{t+1}
=
F(h_t,x_t;\theta).
}
\]

The same parameter set \(\theta\) can be reused across computational
time.

Now the central object is no longer only a composition of distinct maps.

It is a state-transition system.

\section{Parameter sharing across
time}\label{parameter-sharing-across-time}

Unrolling a recurrent network produces a deep computation.

But the same transition law can reappear at every time step.

This creates a different kind of depth:

\[
\text{depth by repeated state evolution}.
\]

The distinction becomes important later for:

\begin{itemize}
\tightlist
\item
  recurrent depth;
\item
  weight tying;
\item
  latent clocks;
\item
  iterative reasoning.
\end{itemize}

A recurrent architecture is naturally read as a discrete dynamical
system.

It is not automatically a continuous flow.

\section{Exact linear recurrence}\label{exact-linear-recurrence}

Consider

\[
h_{t+1}
=
a h_t+b x_t.
\]

Repeated substitution gives

\[
\boxed{
h_T
=
a^T h_0
+
b
\sum_{j=0}^{T-1}
a^{T-1-j}x_j.
}
\]

Earlier inputs are transported by powers of \(a\).

The recurrence therefore exposes a memory law directly.

If

\[
|a|<1,
\]

old contributions decay geometrically in this scalar model.

If

\[
|a|>1,
\]

they can grow.

The architecture creates a temporal transport problem.

\section{Exact recurrence witness}\label{exact-recurrence-witness}

Take

\[
a=\frac12,
\qquad
b=2,
\qquad
h_0=1,
\]

and

\[
(x_0,x_1,x_2)
=
(3,-1,4).
\]

The direct recurrence gives

\[
(h_0,h_1,h_2,h_3)
=
\left(
1,
\frac{13}{2},
\frac54,
\frac{69}{8}
\right).
\]

The closed form independently gives

\[
h_3
=
\frac{69}{8}.
\]

This is the simplest exact demonstration that recurrent state carries
weighted history.

\section{Long-term memory is not guaranteed by
recurrence}\label{long-term-memory-is-not-guaranteed-by-recurrence}

A recurrent loop can preserve state in principle.

That does not mean it preserves useful information over long intervals.

Problems can arise from:

\begin{itemize}
\tightlist
\item
  repeated Jacobian products;
\item
  contraction;
\item
  expansion;
\item
  saturation;
\item
  interference;
\item
  optimization difficulty.
\end{itemize}

This is one reason gated recurrent systems were developed.

\section{LSTM: persistent state plus
gates}\label{lstm-persistent-state-plus-gates}

The Long Short-Term Memory architecture uses gated recurrent machinery
designed to address long-lag learning difficulties (Hochreiter and
Schmidhuber 1997).

The structural shift is important:

\begin{quote}
persistence becomes an explicit controlled pathway.
\end{quote}

The state is no longer merely whatever survives repeated application of
one transition.

Gates regulate:

\begin{itemize}
\tightlist
\item
  what enters;
\item
  what remains;
\item
  what becomes externally visible.
\end{itemize}

This is an architectural move toward managed state transport.

\section{Gates are operators on information
flow}\label{gates-are-operators-on-information-flow}

A gate can be read as a data-dependent multiplicative operator.

Schematically,

\[
g_t\odot v_t.
\]

The gate value decides how strongly a quantity is passed.

This is distinct from a plain additive update.

The later Highway Networks construction will reuse the same broad design
idea across depth rather than only across recurrent time.

\section{Encoder--decoder architecture introduces an explicit
interface}\label{encoderdecoder-architecture-introduces-an-explicit-interface}

Now separate computation into two subsystems:

\[
z=E(x),
\]

\[
\hat y=D(z).
\]

The encoder produces an interface object \(z\).

The decoder is constrained to work from that object.

Sequence-to-sequence systems made this separation especially visible for
variable-length input/output mappings (Sutskever et al. 2014).

The important Atlas object is the interface itself.

\section{The interface is a
contract}\label{the-interface-is-a-contract}

Once the decoder receives only \(z\), every downstream capability must
be recoverable from \(z\).

This makes encoder--decoder architecture a concrete instance of a
broader principle:

\begin{quote}
an interface is sufficient only for the distinctions downstream
computation still needs.
\end{quote}

The Representation, Residual, and Transfer chapters formalized this idea
more generally.

Here it appears as architecture.

\section{Encoder information-loss
boundary}\label{encoder-information-loss-boundary}

Suppose

\[
E(x_1)
=
E(x_2)
=
z,
\]

but the required outputs satisfy

\[
y_1\ne y_2.
\]

A deterministic decoder receives the same \(z\) in both cases.

Therefore

\[
D(E(x_1))
=
D(E(x_2)).
\]

It cannot equal both distinct targets.

So a noninjective encoder may discard information safely only relative
to the downstream capability being preserved.

This is not an argument that useful encoders must be injective.

Compression is often the point.

\section{Bottleneck is a design
decision}\label{bottleneck-is-a-design-decision}

The interface can be:

\begin{itemize}
\tightlist
\item
  a fixed-dimensional vector;
\item
  a sequence;
\item
  a memory;
\item
  a graph;
\item
  an operator state;
\item
  a set of retrieved evidence.
\end{itemize}

What matters is what distinctions it preserves.

Later attention architectures loosen the fixed-vector bottleneck
dramatically.

But the encoder--decoder lesson remains:

\begin{quote}
communication between subsystems creates an explicit sufficiency
boundary.
\end{quote}

\section{Highway Networks: transform or
carry}\label{highway-networks-transform-or-carry}

Highway Networks use learned transform and carry gates across depth
(Srivastava et al. 2015).

A structural form is

\[
\boxed{
y
=
T(x)\odot H(x)
+
C(x)\odot x.
}
\]

A common tied-gate choice is

\[
C(x)=1-T(x).
\]

Then

\[
y
=
T(x)\odot H(x)
+
(1-T(x))\odot x.
\]

The layer can learn how much to transform and how much to carry.

\section{Highway identity limit}\label{highway-identity-limit}

In the tied-gate form, if

\[
T(x)=0,
\]

then

\[
C(x)=1,
\]

so

\[
\boxed{
y=x.
}
\]

The architecture therefore contains an exact carry path.

At the opposite limit,

\[
T(x)=1,
\]

we obtain

\[
y=H(x).
\]

With the tied gate constrained coordinatewise to ({[}0,1{]}), the layer
interpolates coordinatewise between carried and transformed state.

\section{Exact highway witness}\label{exact-highway-witness}

Take

\[
x=2,
\qquad
H(x)=5,
\]

and

\[
T=\frac14,
\qquad
C=\frac34.
\]

Then

\[
y
=
\frac14\cdot5
+
\frac34\cdot2
=
\frac{11}{4}.
\]

The exact carry limit returns

\[
2.
\]

The exact transform limit returns

\[
5.
\]

This small witness shows the gate algebra without making any claim about
training dynamics.

\section{Gated carry is not residual
addition}\label{gated-carry-is-not-residual-addition}

Highway and residual systems both create a direct path for state
transport.

But they do so differently.

A highway layer uses multiplicative gates:

\[
T\odot H
+
C\odot x.
\]

A canonical residual block uses additive identity transport:

\[
x+F(x).
\]

These are related design ideas.

They are not the same algebra.

The distinction matters when analyzing Jacobians and optimization.

\section{Residual learning}\label{residual-learning}

Residual networks made the identity-plus-increment form central:

\[
\boxed{
x_{k+1}
=
x_k
+
F_k(x_k).
}
\]

He and colleagues used residual learning to train very deep
image-recognition systems and reported substantial empirical gains in
that setting (He et al. 2016).

The Atlas uses the architecture for a narrower structural reason:

\begin{quote}
depth now looks like repeated state transport plus correction.
\end{quote}

\section{Residual Jacobian}\label{residual-jacobian}

Differentiate:

\[
x_{k+1}
=
x_k
+
F_k(x_k).
\]

Then

\[
\boxed{
J_k
=
I
+
J_{F_k}(x_k).
}
\]

Across \(L\) residual blocks,

\[
J_{\rm total}
=
\left(I+J_{F_{L-1}}\right)
\cdots
\left(I+J_{F_0}\right).
\]

The identity contribution is exact.

The product can still be badly conditioned.

\section{Identity transport}\label{identity-transport}

If the residual branch vanishes,

\[
F_k(x)=0,
\]

then

\[
x_{k+1}=x_k.
\]

The block becomes exactly the identity map.

That is the precise identity-transport claim.

It does not prove:

\begin{itemize}
\tightlist
\item
  all residual networks optimize easily;
\item
  all gradients remain well conditioned;
\item
  arbitrary depth is harmless.
\end{itemize}

The architecture creates an identity route.

Training still determines what happens around it.

\section{Exact residual witness}\label{exact-residual-witness}

For

\[
F(x)=Ax,
\qquad
A=
\begin{pmatrix}
1&2\\
-1&3
\end{pmatrix},
\]

the residual block Jacobian is

\[
I+A
=
\begin{pmatrix}
2&2\\
-1&4
\end{pmatrix}.
\]

If

\[
A=0,
\]

the block Jacobian is exactly

\[
I.
\]

The computational witness checks both identities.

\section{Structural lineage plate}\label{structural-lineage-plate}

\begin{figure}
\centering
\pandocbounded{\includegraphics[keepaspectratio,alt={Six schematic panels showing layered composition, convolutional weight sharing, recurrent state, encoder-decoder interface, highway transform/carry gating, and residual identity bypass.}]{figures/masters/ATLAS-FIG-ARCHHIST-001.png}}
\caption{Six schematic panels showing layered composition, convolutional
weight sharing, recurrent state, encoder-decoder interface, highway
transform/carry gating, and residual identity bypass.}
\end{figure}

The plate is deliberately schematic.

It does not encode:

\begin{itemize}
\tightlist
\item
  historical priority;
\item
  benchmark quality;
\item
  parameter count;
\item
  compute;
\item
  architectural superiority.
\end{itemize}

Its only job is to show how the mathematical object changes.

\section{Residual systems invite a dynamical
lens}\label{residual-systems-invite-a-dynamical-lens}

Rewrite the residual update as

\[
x_{k+1}-x_k
=
F_k(x_k).
\]

If

\[
F_k(x)
=
h f_k(x),
\]

then

\[
x_{k+1}
=
x_k
+
h f_k(x_k).
\]

This resembles an explicit Euler step.

That resemblance is useful.

It is not an identity theorem.

\section{The ODE analogy has
conditions}\label{the-ode-analogy-has-conditions}

To interpret a residual system as a discretization of one autonomous
ODE,

\[
\dot x=f(x),
\]

we would need more than the residual form.

Questions include:

\begin{itemize}
\tightlist
\item
  Is there one shared \(f\)?
\item
  Is there a meaningful small step \(h\)?
\item
  Does the discrete family converge as \(h\to0\)?
\item
  Are the layer operations compatible with a smooth vector field?
\item
  Does the numerical method preserve the relevant structure?
\end{itemize}

Generic ResNets do not answer these automatically.

So the safe statement is:

\begin{quote}
residual systems admit an integrator lens.
\end{quote}

\section{Neural ODEs make continuous depth
explicit}\label{neural-odes-make-continuous-depth-explicit}

Neural Ordinary Differential Equations specify the evolution law
directly:

\[
\frac{dz}{dt}
=
f(z,t;\theta),
\]

and obtain outputs through numerical integration (Chen et al. 2018).

Here continuous depth is part of the model definition.

The numerical solver is therefore an architectural component.

This is categorically different from retrospectively declaring every
residual network to be an ODE.

\section{Continuous depth changes where approximation
lives}\label{continuous-depth-changes-where-approximation-lives}

In a discrete network, learned layer maps are explicit model components.

In a Neural ODE, the learned object is the vector field, while
finite-time transport is produced by a numerical solver.

Approximation therefore enters in two places:

\begin{itemize}
\tightlist
\item
  function approximation in \(f\);
\item
  numerical approximation of the flow.
\end{itemize}

This creates the bridge to the Numerical Intelligence chapter.

\section{Architecture as transport}\label{architecture-as-transport}

The lineage can now be restated.

A layered network says:

\[
\text{apply another map}.
\]

A recurrent system says:

\[
\text{update persistent state}.
\]

A highway system says:

\[
\text{learn how much state to carry}.
\]

A residual system says:

\[
\text{carry state and add an increment}.
\]

A continuous-depth system says:

\[
\text{specify a law of motion and integrate it}.
\]

The object has shifted from static composition toward transport.

\section{Architecture can encode
invariance}\label{architecture-can-encode-invariance}

Convolution encodes translation structure.

Recurrence encodes temporal parameter sharing.

Encoder--decoder systems encode a communication bottleneck.

Highway layers encode gated carry.

Residual blocks encode identity transport.

These are architectural priors.

They reduce the set of possible computations before training begins.

Architecture is therefore a form of inductive bias expressed through
operators and state organization.

\section{Architecture also creates failure
modes}\label{architecture-also-creates-failure-modes}

Every structural prior excludes something.

Convolution can mishandle boundaries or nonstationary spatial structure.

Recurrence can forget.

Bottlenecks can erase task-relevant distinctions.

Gates can saturate or route poorly.

Residual products can still become ill-conditioned.

Continuous-depth models inherit numerical-solver choices.

There is no architecture without a failure surface.

\section{History is not a ranking}\label{history-is-not-a-ranking}

The sequence in this chapter should not be read as:

\[
\text{MLP}<\text{CNN}<\text{RNN}<\text{ResNet}<\text{ODE}.
\]

Different architectures solve different structural problems.

A simple MLP may be exactly right for one task.

A recurrent state may be unnecessary.

A continuous-depth solver may add cost without benefit.

The Atlas uses history only to expose reusable mathematical ideas.

\section{What changes in the object of
analysis?}\label{what-changes-in-the-object-of-analysis}

The structural transitions can be summarized as follows.

\subsection{Layered networks}\label{layered-networks}

Analyze:

\[
\text{composition}.
\]

\subsection{CNNs}\label{cnns}

Analyze:

\[
\text{structured shared operators}.
\]

\subsection{RNNs and LSTMs}\label{rnns-and-lstms}

Analyze:

\[
\text{persistent state evolution}.
\]

\subsection{Encoder--decoders}\label{encoderdecoders}

Analyze:

\[
\text{interfaces and sufficiency}.
\]

\subsection{Highway systems}\label{highway-systems}

Analyze:

\[
\text{gated transform/carry paths}.
\]

\subsection{Residual systems}\label{residual-systems}

Analyze:

\[
\text{identity transport plus increment}.
\]

\subsection{Neural ODEs}\label{neural-odes}

Analyze:

\[
\text{explicit continuous evolution plus numerical solution}.
\]

This is the architecture lineage the downstream Atlas needs.

\section{Bridge to Transformers}\label{bridge-to-transformers}

Transformers will introduce another decisive shift.

Instead of propagating only through fixed local connectivity or
recurrent state, a token representation can dynamically aggregate
information from other positions through attention.

The next architecture chapter will therefore add:

\[
\text{content-dependent global interaction}.
\]

But it inherits everything here:

\begin{itemize}
\tightlist
\item
  composition;
\item
  state;
\item
  interfaces;
\item
  residual paths;
\item
  normalization;
\item
  depth.
\end{itemize}

Transformers are not outside this lineage.

They reorganize it.

\section{Bridge to adaptive depth}\label{bridge-to-adaptive-depth}

Once depth is read as accumulated state transformation, a new question
becomes natural:

\begin{quote}
Why must every input receive the same number of updates?
\end{quote}

Adaptive-depth systems make computational horizon input-dependent.

That chapter depends on the distinction established here between:

\begin{itemize}
\tightlist
\item
  layer count;
\item
  state evolution;
\item
  computational time.
\end{itemize}

\section{Bridge to network
numerics}\label{bridge-to-network-numerics}

Residual and continuous-depth systems invite numerical questions:

\begin{itemize}
\tightlist
\item
  What is the effective step size?
\item
  Is the update stable?
\item
  Which invariants drift?
\item
  How does splitting affect the trajectory?
\item
  Can a better integrator improve computation?
\end{itemize}

Those are not metaphorical questions once the architecture is explicitly
organized as state evolution.

The Network Numerics chapter will take them literally.

\section{What would falsify an architectural
analogy?}\label{what-would-falsify-an-architectural-analogy}

A structural analogy should make testable predictions.

If we claim convolutional equivariance, test the commutator with
translation.

If we claim recurrence stores long-term information, perturb old inputs
and measure influence.

If we claim a bottleneck is sufficient, search for target-distinct
states it merges.

If we claim residual transport improves conditioning, inspect the
Jacobian product rather than the identity term alone.

If we claim a network approximates a continuous flow, test refinement
and solver dependence.

Architecture should constrain observation.

\section{Closing view}\label{closing-view-14}

The important history is not a procession of brand names.

It is a progression in what computation is allowed to remember, share,
preserve, and transport.

We began with

\[
x_{k+1}=F_k(x_k).
\]

Then architecture learned to express:

\begin{itemize}
\tightlist
\item
  spatial symmetry;
\item
  temporal persistence;
\item
  bounded interfaces;
\item
  gated carrying;
\item
  identity transport;
\item
  explicit continuous evolution.
\end{itemize}

The Atlas thesis becomes clearer here.

A modern neural model is not merely a large static function.

It is increasingly useful to treat it as an organized computational
system whose state is transformed under architectural rules.

The next chapters will keep asking what those rules buy us, what they
destroy, and which mathematical tools can certify the difference.

\section{References used in this
chapter}\label{references-used-in-this-chapter-17}

\begin{itemize}
\tightlist
\item
  (Rumelhart et al. 1986)
\item
  (LeCun et al. 1998)
\item
  (Hochreiter and Schmidhuber 1997)
\item
  (Sutskever et al. 2014)
\item
  (Srivastava et al. 2015)
\item
  (He et al. 2016)
\item
  (Chen et al. 2018)
\end{itemize}

See `sources/source-locks/ATLAS-CH-ARCHHIST-001.yaml` for exact source
roles, internal Atlas pins, and claim boundaries.

\chapter{The Transformer as Baseline
Object}\label{the-transformer-as-baseline-object}

\textbf{Epistemic status:} Established Architecture + Atlas Synthesis +
Exact Computational Witness\\
\textbf{Specification:}
manuscript/specifications/ATLAS-CH-TRANSFORMER-001.md\\
\textbf{Derivation packet:}
mathematics/derivations/ATLAS-CH-TRANSFORMER-001-DERIVATIONS.md\\
\textbf{Computational witness:}
mathematics/computational-witnesses/ATLAS-CW-TRANSFORMER-001.md\\
\textbf{Source lock:} sources/source-locks/ATLAS-CH-TRANSFORMER-001.yaml

\section{We need one reference
object}\label{we-need-one-reference-object}

The Atlas will later discuss:

\begin{itemize}
\tightlist
\item
  normalized Transformers;
\item
  sparse Transformers;
\item
  mixture-of-experts;
\item
  split-operator residual transport;
\item
  relative-position operators;
\item
  adaptive depth;
\item
  mechanistic diagnostics;
\item
  network numerics.
\end{itemize}

Those chapters can only say what they change if the baseline object is
explicit.

This chapter therefore asks a deliberately conservative question:

\begin{quote}
What is the conventional Transformer object we are modifying?
\end{quote}

The answer is not ``attention.''

A Transformer is an organized state-update system built from:

\begin{itemize}
\tightlist
\item
  token-state representations;
\item
  positional information;
\item
  attention sublayers;
\item
  position-wise feed-forward sublayers;
\item
  residual paths;
\item
  normalization;
\item
  masking;
\item
  topology;
\item
  readout.
\end{itemize}

\section{Attention is a component, not the
architecture}\label{attention-is-a-component-not-the-architecture}

The Attention-as-an-Operator chapter studies attention deeply.

Here we need less mathematics and more assembly.

One attention head computes a score-derived mixing operator and applies
it to values.

A Transformer block then wraps that operation inside:

\begin{itemize}
\tightlist
\item
  projections;
\item
  multi-head factorization;
\item
  output projection;
\item
  residual transport;
\item
  normalization;
\item
  a feed-forward sublayer.
\end{itemize}

So:

\[
\boxed{
\text{attention}
\neq
\text{Transformer block}
\neq
\text{Transformer system}.
}
\]

This distinction will matter repeatedly downstream.

\section{The shared blackboard
allegory}\label{the-shared-blackboard-allegory}

Imagine a shared blackboard.

Each token position owns one row.

Attention allows one row to read selected information from other rows.

A causal mask closes some sight lines before the read weights are
normalized.

The feed-forward network then updates each row locally.

Residual paths preserve the evolving shared state while sublayers add
corrections.

The analogy is useful because it separates:

\begin{itemize}
\tightlist
\item
  shared state;
\item
  cross-row communication;
\item
  local transformation;
\item
  access restrictions.
\end{itemize}

Its limit is equally important.

The residual stream is not literally a blackboard service, and an
attention weight is not automatically a semantic explanation.

\section{Residual-stream state}\label{residual-stream-state}

Let sequence length be \(n\) and model width be \(d\).

Write the Transformer state as

\[
\boxed{
H\in\mathbb R^{n\times d}.
}
\]

Each row represents one sequence position.

The initial state is typically formed from token embeddings plus
positional information:

\[
H_0
=
E(\text{tokens})+P.
\]

The positional mechanism \(P\) is a separate design axis.

The chapter does not assume that sinusoidal, learned, rotary, relative,
or operator-valued position is universally correct.

\section{One attention head}\label{one-attention-head}

For one head,

\[
Q=HW_Q,
\qquad
K=HW_K,
\qquad
V=HW_V.
\]

The score matrix is

\[
S
=
\frac{QK^\top}{\sqrt{d_k}}+M.
\]

The row-normalized mixing matrix is

\[
A
=
\operatorname{softmax}_{\rm row}(S).
\]

The head output is

\[
O=AV.
\]

This creates a useful object chain:

\[
H
\to
(Q,K,V)
\to
S
\to
A
\to
O.
\]

Each object answers a different question.

\section{Scores are not weights}\label{scores-are-not-weights}

The score matrix \(S\) is unnormalized.

The mixing matrix \(A\) is normalized row-wise.

The values \(V\) are the content being mixed.

The output \(O\) is the result of applying the mixing operator.

Therefore:

\[
\text{score}
\neq
\text{attention weight}
\neq
\text{value}
\neq
\text{output}.
\]

This may look elementary.

Many later diagnostic errors begin by collapsing these distinctions.

\section{The original Transformer}\label{the-original-transformer}

Vaswani and colleagues introduced the Transformer architecture built
around self-attention, multi-head attention, position-wise feed-forward
layers, residual connections, LayerNorm, positional encoding, and an
encoder-decoder topology (Vaswani et al. 2017).

The historical claim here is source-scoped.

The Atlas uses that paper as the canonical source for the original
architecture.

It does not claim every modern Transformer follows the original block
ordering.

\section{Multi-head attention}\label{multi-head-attention}

Suppose there are \(m\) heads.

Each head produces

\[
O_h.
\]

The outputs are concatenated:

\[
C
=
\operatorname{Concat}
(O_1,\ldots,O_m).
\]

The multi-head sublayer then applies an output projection:

\[
\boxed{
\operatorname{MHA}(H)
=
CW_O.
}
\]

The heads therefore do not simply remain separate.

Their projected outputs are recombined.

\section{Exact two-head witness}\label{exact-two-head-witness}

Take one token and two head outputs:

\[
O_1=(1,2),
\]

\[
O_2=(3,4).
\]

Concatenate:

\[
C=(1,2,3,4).
\]

Let

\[
W_O
=
\begin{pmatrix}
1&0\\
0&1\\
1&1\\
2&-1
\end{pmatrix}.
\]

Then

\[
\boxed{
CW_O=(12,1).
}
\]

This exact witness checks head concatenation and output projection only.

It says nothing about whether the heads learned distinct semantic roles.

\section{Multi-head does not mean one monolithic attention
matrix}\label{multi-head-does-not-mean-one-monolithic-attention-matrix}

Each head can have its own:

\begin{itemize}
\tightlist
\item
  query projection;
\item
  key projection;
\item
  value projection;
\item
  score matrix;
\item
  normalized mixing matrix.
\end{itemize}

The resulting head outputs are then concatenated and projected.

One can sometimes derive equivalent block representations.

But the baseline architecture should not be casually described as one
giant attention matrix unless that equivalence is stated precisely.

The factorization is part of the architecture.

\section{Causal masking is architectural access
control}\label{causal-masking-is-architectural-access-control}

For left-to-right autoregressive self-attention, define the ideal causal
mask:

\[
M_{ij}
=
\begin{cases}
0,&j\le i,\\
-\infty,&j>i.
\end{cases}
\]

After adding the mask to scores, future positions receive ideal score

\[
-\infty.
\]

Therefore their exponentiated contribution is zero.

After softmax,

\[
\boxed{
A_{ij}=0
\qquad
(j>i).
}
\]

The mask changes which dependencies are permitted before learned weights
are normalized.

It is not a learned preference.

\section{Mathematical mask versus implementation
mask}\label{mathematical-mask-versus-implementation-mask}

The exact mathematical notation uses

\[
-\infty.
\]

Real systems usually implement this with:

\begin{itemize}
\tightlist
\item
  masked-softmax logic;
\item
  a representable very negative number;
\item
  fused attention kernels with explicit mask semantics.
\end{itemize}

That distinction matters.

The theorem is about zero support on forbidden positions.

The implementation is about realizing that support reliably in finite
arithmetic.

\section{Exact three-token causal
witness}\label{exact-three-token-causal-witness}

Take zero unmasked scores.

Then each row distributes weight uniformly over its allowed prefix:

\[
A
=
\begin{pmatrix}
1&0&0\\
1/2&1/2&0\\
1/3&1/3&1/3
\end{pmatrix}.
\]

Let

\[
V=
\begin{pmatrix}
2\\
5\\
11
\end{pmatrix}.
\]

Then

\[
AV
=
\begin{pmatrix}
2\\
7/2\\
6
\end{pmatrix}.
\]

Now replace the third value only:

\[
V'
=
\begin{pmatrix}
2\\
5\\
101
\end{pmatrix}.
\]

Then

\[
AV'
=
\begin{pmatrix}
2\\
7/2\\
36
\end{pmatrix}.
\]

The first two outputs remain exactly unchanged.

The future value cannot affect earlier rows through this masked
attention operation.

\section{Causal mask does not imply causal
understanding}\label{causal-mask-does-not-imply-causal-understanding}

The dependency graph is causal in the narrow left-to-right architectural
sense:

future token states are unavailable.

This does not imply that:

\begin{itemize}
\tightlist
\item
  the model has discovered causal structure in the scientific sense;
\item
  attention weights identify causal explanations;
\item
  generated text is causally grounded.
\end{itemize}

``Causal mask'' describes permitted sequence dependence.

Nothing more follows automatically.

\section{Position-wise feed-forward
network}\label{position-wise-feed-forward-network}

After attention, the Transformer applies a feed-forward network
independently at each position.

For token state \(h\),

\[
\boxed{
\operatorname{FFN}(h)
=
\phi(hW_1+b_1)W_2+b_2.
}
\]

The same FFN parameters are reused across sequence positions within the
layer.

This is an important contrast:

\begin{itemize}
\tightlist
\item
  attention mixes across positions;
\item
  FFN transforms each position locally in state space.
\end{itemize}

\section{Different kinds of
sharing}\label{different-kinds-of-sharing}

The architecture now contains several kinds of repeated structure.

Convolution shares a local kernel across spatial positions.

Recurrence shares one transition across time.

Transformer FFNs share one position-wise map across sequence positions.

Attention shares projection matrices across positions while the mixing
pattern is content-dependent.

Parameter sharing is therefore not one phenomenon.

Its geometry depends on the operator.

\section{Residual transport}\label{residual-transport}

A Transformer sublayer usually sits inside a residual connection.

Schematically:

\[
R(H)
=
H+S(H).
\]

If

\[
S(H)=0,
\]

then

\[
\boxed{
R(H)=H.
}
\]

The exact identity path from the Architecture History chapter survives
here.

The sublayer adds a correction to the evolving residual-stream state.

\section{Exact residual identity
witness}\label{exact-residual-identity-witness}

Take

\[
H=(2,-1,3).
\]

Let

\[
S(H)=(0,0,0).
\]

Then

\[
H+S(H)
=
(2,-1,3).
\]

This is the precise identity-path claim.

It does not establish that the total Transformer Jacobian remains well
conditioned.

\section{Layer Normalization}\label{layer-normalization}

Layer Normalization normalizes features within a state vector (Ba et al.
2016).

For

\[
x\in\mathbb R^d,
\]

define

\[
\mu(x)
=
\frac1d
\sum_i x_i,
\]

and

\[
\sigma^2(x)
=
\frac1d
\sum_i
(x_i-\mu)^2.
\]

Then

\[
\boxed{
\operatorname{LN}(x)
=
\gamma\odot
\frac{x-\mu}
{\sqrt{\sigma^2+\varepsilon}}
+
\beta.
}
\]

The stabilizing constant

\[
\varepsilon>0
\]

belongs in the definition.

\section{LayerNorm is not
bookkeeping}\label{layernorm-is-not-bookkeeping}

LayerNorm changes the representation.

Before the learned affine parameters are applied, it centers the
declared feature vector by its sample mean and rescales it by the
regularized sample standard deviation:

\[
x
\mapsto
\frac{x-\mu(x)}
{\sqrt{\sigma^2(x)+\varepsilon}}.
\]

It then applies learned coordinatewise scale and shift through
\(\gamma\) and \(\beta\).

This operation changes:

\begin{itemize}
\tightlist
\item
  geometry;
\item
  scale;
\item
  derivative structure;
\item
  how residual and sublayer signals interact.
\end{itemize}

It should not be summarized as an unconditional exact invariance to
arbitrary affine rescaling. In particular, the explicit
\(\varepsilon>0\) term breaks exact scale invariance in general, and
learned \(\gamma,\beta\) restore trainable coordinatewise scale and
offset after normalization.

Later normalized-geometry chapters will ask whether some normalization
machinery can be replaced by more intrinsic constraints.

For now, LayerNorm is part of the baseline object.

\section{Post-LN ordering}\label{post-ln-ordering}

The original Transformer applies residual addition and then LayerNorm
around a sublayer.

Write attention sublayer \(\mathcal A\).

Then

\[
Y
=
\operatorname{LN}
\left(
H+\mathcal A(H)
\right).
\]

For feed-forward sublayer \(\mathcal F\),

\[
H'
=
\operatorname{LN}
\left(
Y+\mathcal F(Y)
\right).
\]

This is a post-LN arrangement.

\section{Pre-LN ordering}\label{pre-ln-ordering}

A common later design moves normalization before the residual branch:

\[
Y
=
H
+
\mathcal A
\left(
\operatorname{LN}(H)
\right),
\]

\[
H'
=
Y
+
\mathcal F
\left(
\operatorname{LN}(Y)
\right).
\]

Xiong and colleagues analyze pre-LN versus post-LN Transformer behavior
and optimization differences (Xiong et al. 2020).

The Atlas uses that work only to support the architectural distinction
and its training relevance.

It does not declare one ordering universally superior.

\section{Pre-LN and post-LN are not notation
variants}\label{pre-ln-and-post-ln-are-not-notation-variants}

Compare:

\[
\operatorname{LN}
\left(
H+\mathcal A(H)
\right)
\]

with

\[
H+
\mathcal A
\left(
\operatorname{LN}(H)
\right).
\]

Normalization and nonlinear sublayer evaluation occur in different
places.

The residual-stream state therefore evolves differently.

A chapter that says only ``Transformer with LayerNorm'' is
underspecified.

\section{Baseline pre-LN block}\label{baseline-pre-ln-block}

For later Atlas work, it is useful to display one common modern
baseline:

\[
Y
=
H+\operatorname{MHA}(\operatorname{LN}(H)),
\]

\[
H'
=
Y+\operatorname{FFN}(\operatorname{LN}(Y)).
\]

This chapter uses that form for its main visual plate.

The choice is descriptive.

The original Transformer remains post-LN.

\begin{figure}
\centering
\pandocbounded{\includegraphics[keepaspectratio,alt={A three-panel baseline Transformer plate showing a pre-LN residual-stream block, the exact 3x3 causal mixing matrix, and encoder-only/decoder-only/encoder-decoder topology families.}]{figures/masters/ATLAS-FIG-TRANSFORMER-001.png}}
\caption{A three-panel baseline Transformer plate showing a pre-LN
residual-stream block, the exact 3x3 causal mixing matrix, and
encoder-only/decoder-only/encoder-decoder topology families.}
\end{figure}

\section{Positional information is a separate
axis}\label{positional-information-is-a-separate-axis}

Attention without positional information is permutation-sensitive only
through content, not sequence order in the intended way.

The original Transformer adds positional encodings (Vaswani et al.
2017).

Modern systems may instead use:

\begin{itemize}
\tightlist
\item
  learned absolute position;
\item
  rotary position;
\item
  relative bias;
\item
  multidimensional schemes;
\item
  operator-valued position.
\end{itemize}

The baseline object only requires that the positional mechanism be
stated explicitly.

The Geometry of Position chapter develops this axis separately.

\section{Encoder-only topology}\label{encoder-only-topology}

An encoder-only Transformer applies self-attention and FFN blocks to an
input sequence.

The attention pattern is not autoregressively causal by architectural
necessity.

BERT is a representative encoder-only family (Devlin et al. 2019).

Its pretraining objectives and masking procedure are training/task
choices layered on top of the architecture.

Encoder-only therefore does not mean:

\begin{quote}
every token always sees every other token under every possible masking
policy.
\end{quote}

Topology and task masking are separate.

\section{Decoder-only topology}\label{decoder-only-topology}

A decoder-only autoregressive Transformer uses causally masked
self-attention.

Each position can depend on:

\begin{itemize}
\tightlist
\item
  itself;
\item
  previous positions;
\item
  not future positions under the ideal mask.
\end{itemize}

GPT-2 is a representative decoder-only autoregressive family (Radford et
al. 2019).

Again, the source is used only to document topology.

Decoder-only is not a synonym for every language model.

\section{Encoder--decoder topology}\label{encoderdecoder-topology}

The original Transformer uses two stacks.

The encoder produces memory:

\[
E
=
\operatorname{Encoder}(X).
\]

The decoder first uses masked self-attention over its generated prefix.

It then cross-attends to encoder memory.

For cross-attention:

\[
Q
=
H_{\rm dec}W_Q,
\]

while

\[
K
=
EW_K,
\]

and

\[
V
=
EW_V.
\]

Decoder queries therefore read from encoder-derived keys and values.

\section{Cross-attention is an
interface}\label{cross-attention-is-an-interface}

Encoder-decoder architecture creates an explicit computational boundary:

\[
\text{encoder state}
\to
\text{decoder access}.
\]

Cross-attention is the read mechanism.

The interface can be interpreted using the same sufficiency questions
developed earlier:

\begin{quote}
Did the encoder memory preserve what the decoder still needs?
\end{quote}

This becomes especially important in multimodal and modular systems.

\section{Three topologies, three dependency
graphs}\label{three-topologies-three-dependency-graphs}

The topology families differ in which state can directly interact.

\subsection{Encoder-only}\label{encoder-only}

Self-attention within one encoded sequence.

\subsection{Decoder-only}\label{decoder-only}

Causally restricted self-attention over a prefix.

\subsection{Encoder--decoder}\label{encoderdecoder}

Encoder self-attention, decoder causal self-attention, and
decoder-to-encoder cross-attention.

These are graph differences.

They should not be collapsed into one generic phrase such as ``a
Transformer.''

\section{Architecture versus training
objective}\label{architecture-versus-training-objective}

A decoder-only Transformer can be trained with different objectives.

An encoder-only Transformer can be trained with different
corruption/masking schemes.

An encoder--decoder Transformer can be used for different conditional
tasks.

Therefore:

\[
\boxed{
\text{architecture}
\neq
\text{training objective}.
}
\]

The distinction is necessary for interpreting experimental claims.

\section{Architecture versus
optimizer}\label{architecture-versus-optimizer}

The block equations do not determine:

\begin{itemize}
\tightlist
\item
  SGD versus Adam;
\item
  learning rate;
\item
  warmup;
\item
  clipping;
\item
  weight decay;
\item
  batch size;
\item
  precision;
\item
  initialization.
\end{itemize}

These choices alter training behavior without changing the nominal
architecture.

The First-Order Optimization chapter treats that machinery separately.

\section{Architecture versus scale}\label{architecture-versus-scale}

Two systems can share architecture and differ dramatically in:

\begin{itemize}
\tightlist
\item
  depth;
\item
  width;
\item
  head count;
\item
  context length;
\item
  vocabulary;
\item
  parameter count;
\item
  training compute.
\end{itemize}

Calling both ``Transformers'' can be correct and still conceal most of
the engineering difference.

The Atlas therefore uses architecture labels carefully.

\section{Architecture versus
tokenizer}\label{architecture-versus-tokenizer}

A Transformer consumes token states.

How text becomes tokens is external to the Transformer block.

Different tokenizers can change:

\begin{itemize}
\tightlist
\item
  sequence length;
\item
  morphology;
\item
  multilingual fertility;
\item
  information boundaries;
\item
  effective context usage.
\end{itemize}

The Tokenization chapter will treat that interface separately.

\section{The residual stream as the persistent
state}\label{the-residual-stream-as-the-persistent-state}

The Architecture History chapter moved from static composition toward
state transport.

The Transformer makes that transport explicit through repeated additive
residual updates.

One useful lens is:

\[
H_0
\to
H_1
\to
\cdots
\to
H_L.
\]

Attention and FFN sublayers write corrections into this evolving state.

This residual-stream viewpoint will matter in:

\begin{itemize}
\tightlist
\item
  mechanistic analysis;
\item
  split operators;
\item
  routing;
\item
  normalization geometry;
\item
  adaptive depth.
\end{itemize}

\section{But ``residual stream'' is not an ontological
primitive}\label{but-residual-stream-is-not-an-ontological-primitive}

The residual stream is an architectural coordinate system.

Another architecture may store equivalent information through:

\begin{itemize}
\tightlist
\item
  recurrent state;
\item
  external memory;
\item
  operator state;
\item
  sparse experts;
\item
  multiple coupled streams.
\end{itemize}

The Atlas uses the residual stream as a baseline object, not as the
unique natural representation of computation.

\section{Attention as communication, FFN as local
transformation}\label{attention-as-communication-ffn-as-local-transformation}

A useful first approximation is:

\begin{itemize}
\tightlist
\item
  attention: position-to-position communication;
\item
  FFN: position-local nonlinear transformation.
\end{itemize}

This is not a complete mechanistic theory.

Attention changes state content.

FFNs can implement complex feature transformations.

Residual mixing across layers entangles the two.

Still, the decomposition is operationally useful.

\section{Multi-head factorization creates multiple projected
interaction
spaces}\label{multi-head-factorization-creates-multiple-projected-interaction-spaces}

Each head receives its own projected Q/K/V spaces.

Thus head \(h\) operates on

\[
Q_h,
K_h,
V_h.
\]

The outputs are later recombined.

This creates a structured factorization of the communication mechanism.

It does not guarantee that every head learns a cleanly separable
function.

Later mechanistic chapters will test specialization rather than assume
it.

\section{Masking changes the operator
support}\label{masking-changes-the-operator-support}

A causal mask sets certain attention entries to zero after
normalization.

A sparse mask can remove other edges.

A local-window mask restricts interaction distance.

Therefore masking changes the support of the mixing operator.

This will connect directly to sparse computation and structured
attention.

\section{Normalization placement changes the update
law}\label{normalization-placement-changes-the-update-law}

In post-LN:

\[
H
\mapsto
\operatorname{LN}(H+S(H)).
\]

In pre-LN:

\[
H
\mapsto
H+S(\operatorname{LN}(H)).
\]

The residual state is therefore normalized at different points in the
computation.

This becomes crucial when asking whether normalized geometry should be
enforced intrinsically rather than repeatedly imposed.

\section{Residual identity path is necessary context for
depth}\label{residual-identity-path-is-necessary-context-for-depth}

Because each block contains an additive identity path, a deep
Transformer is not simply a composition of unrelated maps.

It is a sequence of corrections around a transported state.

This is one reason depth admits a dynamical interpretation.

But the earlier warning remains:

\[
\text{identity path}
\not\Rightarrow
\text{automatic stability}.
\]

The product of block Jacobians still matters.

\section{The baseline object in one
equation}\label{the-baseline-object-in-one-equation}

A simplified pre-LN decoder-only block can be written:

\[
Y_l
=
H_l
+
\operatorname{MHA}_l
\left(
\operatorname{LN}_l^{(1)}(H_l);
M_{\rm causal}
\right),
\]

\[
H_{l+1}
=
Y_l
+
\operatorname{FFN}_l
\left(
\operatorname{LN}_l^{(2)}(Y_l)
\right).
\]

This is the reference object many later Atlas chapters will modify.

They may change:

\begin{itemize}
\tightlist
\item
  normalization;
\item
  position;
\item
  attention approximation;
\item
  sparsity;
\item
  expert routing;
\item
  depth;
\item
  residual transport;
\item
  memory.
\end{itemize}

\section{What the baseline does not
contain}\label{what-the-baseline-does-not-contain}

The equation above does not specify:

\begin{itemize}
\tightlist
\item
  tokenizer;
\item
  embedding tying;
\item
  bias conventions;
\item
  dropout;
\item
  activation function;
\item
  attention kernel implementation;
\item
  KV cache;
\item
  sequence packing;
\item
  distributed parallelism;
\item
  optimizer;
\item
  objective;
\item
  data mixture.
\end{itemize}

Those are real system choices.

They are simply outside this baseline architectural abstraction.

\section{Failure modes of vague Transformer
language}\label{failure-modes-of-vague-transformer-language}

\subsection{``Attention model''}\label{attention-model}

Too coarse.

It hides residual and FFN structure.

\subsection{``Causal Transformer''}\label{causal-transformer}

Still underspecified.

Normalization placement and topology remain open.

\subsection{``Standard Transformer''}\label{standard-transformer}

Often ambiguous between the 2017 post-LN encoder-decoder and later
pre-LN decoder-only systems.

\subsection{``Transformer layer''}\label{transformer-layer}

Can refer to different arrangements of attention, cross-attention, FFN,
normalization, and residuals.

The cure is equations.

\section{Atlas connections}\label{atlas-connections-10}

\textbf{Attention as an Operator.}\\
Supplies the deeper score/operator/value analysis.

\textbf{Geometry of Position.}\\
Treats position as a mathematical object rather than an embedding
afterthought.

\textbf{Normalized representations.}\\
Questions whether LayerNorm-like machinery should be replaced by
intrinsic constraints.

\textbf{Split operators.}\\
Treats attention and FFN/residual updates as composable evolution
operators.

\textbf{Mixture-of-experts.}\\
Replaces or augments the FFN sublayer with routed conditional
computation.

\textbf{Mechanistic diagnostics.}\\
Uses the residual stream, heads, keys, queries, values, and MLP channels
as intervention sites.

\textbf{Adaptive depth.}\\
Questions whether every token/input requires the same number of block
updates.

\section{What would falsify a baseline
claim?}\label{what-would-falsify-a-baseline-claim}

A baseline statement should be architectural and checkable.

If we claim a mask is causal, inspect support.

If we claim a block is pre-LN, inspect equations.

If we claim an architecture is decoder-only, inspect whether encoder
memory/cross-attention exists.

If we claim identity transport, zero the sublayer and verify state
preservation.

If we claim position-wise FFN sharing, inspect parameter reuse across
positions.

If the implementation violates the equations, the label is wrong.

\section{Closing view}\label{closing-view-15}

The Transformer is not one idea.

It is an assembly of:

\begin{itemize}
\tightlist
\item
  projected communication;
\item
  normalized mixing;
\item
  position-local nonlinear transformation;
\item
  residual transport;
\item
  normalization;
\item
  masking;
\item
  topology.
\end{itemize}

That assembly became a baseline because it is modular enough to vary and
stable enough to serve as a reference point.

The Atlas now has a precise object to depart from.

From here, later chapters can ask sharper questions:

\begin{itemize}
\tightlist
\item
  What if normalization becomes geometry?
\item
  What if position becomes an operator?
\item
  What if depth becomes adaptive time?
\item
  What if FFNs become routed experts?
\item
  What if attention is approximated or split?
\item
  What if persistent memory moves outside the weights?
\end{itemize}

Those questions only become meaningful once ``the Transformer'' stops
being a vague noun and becomes an explicit computational object.

\section{References used in this
chapter}\label{references-used-in-this-chapter-18}

\begin{itemize}
\tightlist
\item
  (Vaswani et al. 2017)
\item
  (Ba et al. 2016)
\item
  (Xiong et al. 2020)
\item
  (Devlin et al. 2019)
\item
  (Radford et al. 2019)
\end{itemize}

See `sources/source-locks/ATLAS-CH-TRANSFORMER-001.yaml` for exact
source roles, internal Atlas provenance, and claim boundaries.

\chapter{Depth as Computational
Time}\label{depth-as-computational-time-1}

\textbf{Epistemic status:} established adaptive-depth mechanisms plus
Atlas synthesis and exact finite derivation.\\
\textbf{Specification:}
manuscript/specifications/ATLAS-CH-DEPTH-001.md\\
\textbf{Formal packet:}
mathematics/derivations/ATLAS-CH-DEPTH-001-DERIVATIONS.md\\
\textbf{Computational witness:}
mathematics/computational-witnesses/ATLAS-CW-DEPTH-001.md\\
\textbf{Source lock:} sources/source-locks/ATLAS-CH-DEPTH-001.yaml

A network layer can be read as a transformation.

A stack of layers can also be read as a history.

That second reading changes what depth means.

Instead of asking only how many transformations the architecture
contains, we can ask how much computation a representation has undergone
before the system decides it is finished.

This is the computational-time view of depth.

It is useful precisely because it is weaker than a claim about physical
time.

Depth need not be seconds.

It need not be the time variable of one hidden differential equation.

It is an ordered coordinate of computation.

Once that coordinate is explicit, fixed-depth networks, recurrent-depth
networks, adaptive halting, equilibrium models, and conditional
execution can be compared without pretending they are the same
architecture.

\section{Five notions of time}\label{five-notions-of-time}

Several indices are routinely called ``time'' in machine learning.

They should not be merged.

Physical time measures duration in the external world.

Sequence time indexes positions or events in data.

Optimization time indexes parameter updates.

Network depth indexes transformations applied within one forward
computation.

Computational time counts or weighs internal state transitions used to
produce one result.

A recurrent model may have sequence time t and computational-depth index
k at the same moment.

A training run adds optimization step n.

A deployed system adds wall-clock duration.

These coordinates interact, but they answer different questions.

The purpose of this chapter is to make the depth coordinate explicit
enough that later chapters can reason about allocating it.

\section{Fixed depth}\label{fixed-depth}

Write a depth-L computation as

\[
x_{k+1}=F_k(x_k),\qquad k=0,\ldots,L-1.
\]

Then

\[
x_L=(F_{L-1}\circ\cdots\circ F_0)(x_0).
\]

The architecture commits in advance to L sequential transformations.

Different inputs may produce different activations, but the nominal
transformation count is fixed.

This is the familiar layer-stack view.

The computational-time view does not reject it.

It notices that L is also a budget decision.

Every input is given the same number of depth transitions, whether it
needs all of them or not.

That uniformity can be desirable.

It simplifies batching, latency prediction, and implementation.

It is also only one possible allocation rule.

\section{Depth is not automatically physical
time}\label{depth-is-not-automatically-physical-time}

The Numerics chapter gave us a disciplined way to interpret repeated
updates.

A residual update can resemble a numerical step.

A tied recurrence can resemble repeated evolution under one operator.

Those analogies are useful.

But the index k does not become physical time merely because the
equations resemble an integrator.

A network may change its operator with k.

Its step size may have no declared physical units.

Training may learn transformations with no claim of approximating a
continuous process.

The Atlas therefore uses a careful phrase:

depth can serve as computational time.

That is an architectural interpretation.

A stronger claim that a particular network discretizes a particular
differential equation requires additional evidence.

\section{Recurrent depth}\label{recurrent-depth}

Now tie the transformation:

\[
x_{k+1}=F(x_k).
\]

One operator is reused across computational depth.

This separates parameter count from execution count.

Running F ten times does not require ten independently parameterized
layers.

It also introduces a new question:

what happens if we keep going?

Nothing in weight tying answers that.

The sequence may converge.

It may oscillate.

It may diverge.

It may enter a cycle.

It may converge in some regions of state space and fail in others.

Recurrent depth therefore turns ``more layers'' into an
iterative-dynamics question.

Universal Transformer is a representative example in which a
transformation is applied recurrently across depth rather than assigning
a wholly independent block to every depth position.

The important mechanism is repeated computation, not the product name.

\section{A computation can stop
itself}\label{a-computation-can-stop-itself}

Fixed depth chooses L before inference.

Adaptive depth lets the execution decide when to stop.

Let h\_k be a declared halting signal.

Let epsilon be a threshold.

Let K\_max be the largest permitted depth under the current compute
budget.

Define the bounded stopping depth

\[
\tau=\min\left(\{k\in\{0,\ldots,K_{\max}\}:h_k\le\varepsilon\}\cup\{K_{\max}\}\right).
\]

The realized output is \(x_\tau\). The termination status must also be
retained:

\begin{itemize}
\tightlist
\item
  \texttt{criterion\_met} if \(h_\tau\le\varepsilon\);
\item
  \texttt{budget\_exhausted} otherwise.
\end{itemize}

This equation separates two ideas.

The recurrence determines how state changes.

The stopping rule determines how long the recurrence is allowed to run.

That separation is essential.

A poor halting rule can stop a good recurrence too early.

A good halting rule cannot rescue a recurrence that moves in the wrong
direction.

And a halting score is not automatically an error estimate merely
because training learns it.

Graves's Adaptive Computation Time is a concrete mechanism for learning
variable numbers of recurrent steps.

Universal Transformer adds a dynamic per-position halting mechanism to
recurrent depth.

These sources show that learned variable computation is an implementable
mechanism.

They do not show that learned halting always tracks task difficulty
correctly.

\section{An exact depth-as-error
witness}\label{an-exact-depth-as-error-witness}

Consider the scalar recurrence

\[
x_{k+1}=\frac{x_k+2}{2},\qquad x_0=0.
\]

The fixed point is 2.

Subtract it:

\[
x_{k+1}-2=\frac12(x_k-2).
\]

Hence

\[
x_k=2(1-2^{-k}),\qquad |x_k-2|=2^{1-k}.
\]

Depth now has a literal approximation meaning for this one system.

At depth 1, the error is 1.

At depth 2, it is 1/2.

At depth 3, it is 1/4.

At depth 5, it is 1/16.

For tolerance epsilon in (0,2), the minimum depth satisfying the
tolerance is

\[
\tau(\varepsilon)=\left\lceil\log_2\!\left(\frac{2}{\varepsilon}\right)\right\rceil.
\]

The transformation did not change.

The requested accuracy changed how long we computed.

This is the simplest useful model of adaptive computational time.

It is also deliberately narrow.

Most learned networks do not hand us an exact error formula.

The later Adaptive Depth as Error Control chapter exists precisely
because learned stopping becomes much more interesting when h\_k must
stand in for unknown task error.

\section{Compute is more than step
count}\label{compute-is-more-than-step-count}

Suppose two systems both take four depth transitions.

They need not cost the same amount.

One may use a small operator at every step.

Another may invoke attention over a long sequence.

One may reuse cached state.

Another may move large tensors through memory.

One may parallelize internal work.

Another may be sequential.

So define realized computational cost more carefully.

For one declared resource unit,

\[
C(x_0)=\sum_{k=0}^{\tau-1}c_k(x_k),
\]

where c\_k records the cost of the executed transition in that unit.

If several heterogeneous resources are tracked, use a vector such as

\[
C_{\mathrm{vec}}=(C_{\mathrm{FLOPs}},C_{\mathrm{memory}},C_{\mathrm{energy}},C_{\mathrm{latency}}).
\]

rather than adding unlike units without a declared scalarization.

The integer tau tells us sequential depth.

It does not by itself tell us wall-clock latency.

This becomes especially important for conditional computation.

Skipping half the mathematical blocks is not useful as a speed claim
unless the implementation actually avoids the relevant work and the
hardware can exploit the irregular execution.

\section{Equilibrium depth}\label{equilibrium-depth}

Recurrent depth suggests taking more and more iterations.

Equilibrium models make a different move.

Instead of declaring an explicit number of layers, define the
representation by a fixed-point equation:

\[
x_\star=F(x_\star).
\]

The computational problem becomes:

find x*.

Deep Equilibrium Models are a representative neural architecture built
around this idea.

The phrase ``infinite depth'' is evocative but must be handled
carefully.

The model does not need to store an actually infinite sequence of layer
activations.

It solves a fixed-point/root-finding problem.

And the equilibrium equation does not prescribe one solver.

Naive fixed-point iteration

\[
x_{k+1}=F(x_k)
\]

is one possibility.

Other root-finding methods can reach the same equation through different
computational trajectories.

This distinction matters because an equilibrium can exist while naive
iteration fails to converge.

Definition of a fixed point and convergence of a solver are different
claims.

\section{Convergence is an architectural question only after
assumptions}\label{convergence-is-an-architectural-question-only-after-assumptions}

For the exact witness, F is a contraction.

Every step halves the distance to the fixed point.

That gives a clean convergence story.

A learned recurrent transformation need not have that property.

The local Jacobian may expand in some directions.

The operator may have multiple fixed points.

A solver may depend on initialization.

A stopping test may declare convergence while task-relevant quantities
remain wrong.

Equilibrium depth therefore does not abolish numerical analysis.

It makes numerical analysis part of the architecture.

The solver, tolerance, maximum iterations, residual test,
initialization, and failure handling are all computational choices.

\section{Conditional depth}\label{conditional-depth}

Adaptive halting asks when to stop a sequential computation.

Conditional execution asks which blocks to run at all.

For a residual-style block, write

\[
x_{k+1}=x_k+g_k(x_k)\Delta_k(x_k).
\]

With a hard gate

\[
g_k\in\{0,1\},
\]

the block is either applied or skipped.

If g\_k=0 and the implementation respects that decision before
evaluating Delta\_k, the block's expensive transform can be avoided.

SkipNet is a representative architecture of this kind: an
input-dependent routing policy skips residual blocks.

This creates input-dependent effective depth.

Two examples entering the same nominal network can traverse different
numbers of transformations.

\section{Soft gates are not free
skips}\label{soft-gates-are-not-free-skips}

During training, it is often convenient to relax a binary gate:

\[
g_k\in[0,1].
\]

Then the update can interpolate continuously between carry and
transform.

But the computational claim changes.

If the system must compute Delta\_k(x\_k) before multiplying it by a
small g\_k, then the expensive operation was not skipped.

A soft coefficient near zero is mathematical sparsity of influence.

It is not necessarily computational sparsity of execution.

This distinction is easy to lose when a paper reports sparse weights,
sparse gates, or small activation coefficients.

The hardware question is concrete:

which operations were never executed?

That boundary is handed directly to the later Conditional Computation
chapter.

\section{Effective depth is input-dependent execution
history}\label{effective-depth-is-input-dependent-execution-history}

For hard gating, define the active block set

\[
\mathcal A(x_0)=\{k:g_k(x_k)=1\}.
\]

A simple realized block cost is

\[
C_{\mathrm{block}}(x_0)=\sum_{k\in\mathcal A(x_0)}c_k,
\]

plus the overhead required to make the routing decisions.

The path is a function of the evolving state.

This means conditional depth can create system effects absent from a
static layer count.

Different examples may have different latency.

A batch may wait for its slowest member.

Branch divergence may reduce hardware utilization.

Routing decisions themselves cost compute.

The mean number of active blocks is therefore not the same quantity as
end-to-end throughput.

\section{Parameter depth and execution
depth}\label{parameter-depth-and-execution-depth}

A weight-tied recurrent system may have few unique parameters and many
execution steps.

A large feed-forward system may have many unique parameters and one pass
through each.

An equilibrium model may have a compact parameterization but a variable
number of solver iterations.

A conditional network may contain many blocks but execute only some of
them for one input.

This suggests at least three separate depth coordinates:

architectural depth: how many transformation slots are represented in
the nominal graph;

parameter depth: how many distinct parameterized transformations occur
along the path;

execution depth: how many sequential transformations are actually
evaluated for this input.

They coincide in a plain feed-forward stack.

They separate in recurrent, equilibrium, and conditional systems.

\section{Halting is a decision under
uncertainty}\label{halting-is-a-decision-under-uncertainty}

An adaptive system rarely knows the true remaining task error.

It sees proxies.

A confidence score.

A residual norm.

A change in representation.

A learned halting logit.

A budget counter.

A task-specific certificate.

The stop decision therefore has two possible errors.

Stop too early and computation is insufficient.

Continue too long and resources are wasted.

This is the computational analogue of a sequential decision problem.

But this chapter does not yet equate adaptive halting with numerical
error control.

A learned halting logit may correlate with difficulty without estimating
a rigorous truncation error.

The next numerical chapter will ask what additional structure is needed
before a stopping signal can be interpreted as an error-control
mechanism.

\section{More depth can make things
worse}\label{more-depth-can-make-things-worse}

The computational-time picture can tempt us into a monotone story:

more time, better answer.

That is false in general.

An unstable recurrence can amplify error.

Repeated application can oversmooth or collapse distinctions.

A learned iterative process can move away from a useful intermediate
state.

A conditional controller can route difficult inputs into pathological
loops.

An equilibrium solver can fail its tolerance or reach an unintended
root.

More computation is only useful when the dynamics and stopping semantics
support it.

Depth is a resource coordinate.

It is not a quality score.

\section{Budget exhaustion is a first-class
outcome}\label{budget-exhaustion-is-a-first-class-outcome}

Suppose execution terminates at the finite budget cap K\_max before the
declared criterion h\_k \textless= epsilon is met.

The termination status is then budget\_exhausted rather than
criterion\_met.

If the budget limit is reached first, the system has not necessarily
converged.

It has terminated because resources ended.

That outcome should be distinguishable from successful halting.

A production system may still return its best current state.

A scientific report should not label that return ``converged'' unless
the declared criterion was met.

This sounds obvious.

It becomes less obvious when adaptive computation is buried inside a
large model and the maximum-step cap is treated as an implementation
detail.

The cap is part of the semantics.

\section{Training and inference can
disagree}\label{training-and-inference-can-disagree}

Dynamic-depth mechanisms often use relaxations, penalties, or stochastic
estimators during training.

Deployment may use hard thresholds.

A training objective may reward expected compute.

Hardware may care about tail latency.

A controller may learn under one maximum depth and be deployed under
another.

Thus there are two distinct questions:

did the learning procedure optimize the intended depth policy?

does the deployed execution realize the intended compute behavior?

The first is an optimization question.

The second is a systems question.

A complete conditional-compute story requires both.

\section{What recurrent, adaptive, equilibrium, and conditional
depth
share}\label{what-recurrent-adaptive-equilibrium-and-conditional-depth-share}

These mechanisms look different.

Their common object is a computational path whose length or structure is
no longer just the static count of named layers.

Recurrent depth asks how many times to reuse an operator.

Adaptive depth asks when to stop.

Equilibrium depth asks what state satisfies a terminal fixed-point
condition and how to solve for it.

Conditional depth asks which transformations to execute.

All four make computation allocation explicit.

That is why they belong in one chapter.

\section{What the next chapters may
assume}\label{what-the-next-chapters-may-assume}

ATLAS-CH-ADAPTDEPTH-001 --- Adaptive Depth as Error Control may assume:

\begin{itemize}
\tightlist
\item
  depth as computational time;
\item
  fixed versus adaptive stopping;
\item
  the stopping variable tau;
\item
  exact distinction between a stopping signal and true error;
\item
  fixed-point/equilibrium semantics;
\item
  resource-aware compute accounting.
\end{itemize}

Its task is to add numerical error estimation and adaptive-step logic.

ATLAS-CH-SPARSE-001 --- Conditional Computation may assume:

\begin{itemize}
\tightlist
\item
  hard and relaxed gates;
\item
  active execution sets;
\item
  realized path cost;
\item
  input-dependent effective depth;
\item
  the distinction between small soft weights and actually skipped
  operations.
\end{itemize}

Its task is to extend this into sparse activation, token selection,
routing, and hardware-aware conditional execution.

\section{The boundary to remember}\label{the-boundary-to-remember}

Depth is often drawn vertically.

The Atlas needs a second picture.

Imagine depth horizontally as a computational trajectory.

A representation enters at k=0.

Each transition changes it.

A fixed-depth system stops because the architecture says L.

An adaptive system stops because a rule says enough.

An equilibrium system stops because a solver says the fixed-point
criterion is met.

A conditional system changes which transitions exist on the realized
path.

This picture does not prove that neural computation is continuous time.

It does something more practical.

It turns depth from a static architectural count into a resource that
can be allocated, reused, skipped, or terminated under explicit rules.

Once depth is treated that way, adaptive computation becomes a question
of dynamics, numerics, and systems design rather than a synonym for
``deeper network.''

\section{References used in this
chapter}\label{references-used-in-this-chapter-19}

The exact source identities and authority scopes are pinned in:

sources/source-locks/ATLAS-CH-DEPTH-001.yaml

The external mechanism basis includes Graves's Adaptive Computation
Time, Universal Transformers, Deep Equilibrium Models, and SkipNet. The
numerical and architectural interpretation is grounded in the audited
Atlas Numerics and Architecture History prerequisites.

\chapter{Split-Operator Networks}\label{split-operator-networks}

\textbf{Epistemic status:} established operator-splitting theory + Atlas
architectural synthesis.

A neural block can contain the same named ingredients and still compute
a different map when their order changes.

That sounds obvious when stated in ordinary language. It becomes more
consequential when the block is described as if attention, an MLP,
normalization, gating, transport, or another update were merely
ingredients in a bag. They are not. They are transformations applied to
state. A later transformation sees the state produced by the earlier
one.

The mathematics of operator splitting gives us a disciplined way to talk
about this fact.

It also gives us a trap.

Classical splitting theory studies compositions of flows or operators
under explicit assumptions. A Transformer sublayer is not automatically
one of those exact flows. If we borrow the notation without preserving
the assumptions, an architectural metaphor quietly turns into a false
numerical theorem.

This chapter keeps the useful part and blocks the overclaim.

The central lesson is:

\begin{quote}
\textbf{When subcomputations do not commute, ordering is part of the
architecture. Classical splitting theory can explain that structure
exactly in declared reference systems, but learned neural maps inherit
its convergence orders only when the required flow and regularity
assumptions are actually established.}
\end{quote}

The numerical substrate comes from the audited Numerics chapter and
standard splitting references (McLachlan and Quispel 2002; Hairer et al.
2006). The baseline neural anatomy comes from the audited Transformer
chapter and the original Transformer architecture (Vaswani et al. 2017).

\section{The smallest place order can
matter}\label{the-smallest-place-order-can-matter}

Suppose a state \(x\) is acted on by two transformations.

Call them \(A\) and \(B\).

There are at least three different things one might mean by ``use
both.''

One is an additive update:

\[
x^+
=
x+F_A(x)+F_B(x).
\]

Another is sequential composition:

\[
x_1=\Psi_A(x),
\qquad
x^+=\Psi_B(x_1).
\]

A third is the reversed sequence:

\[
\tilde x_1=\Psi_B(x),
\qquad
\tilde x^+=\Psi_A(\tilde x_1).
\]

These are different constructions.

Even when the two transformations are small residual increments, the
sequential map contains cross terms that the additive map does not.

Take the linear residual increments

\[
F_A(x)=hAx,
\qquad
F_B(x)=hBx.
\]

The additive update is

\[
x^+
=
\left(I+h(A+B)\right)x.
\]

If \(A\) acts first and \(B\) acts on the resulting state,

\[
x^+
=
(I+hB)(I+hA)x,
\]

so

\[
x^+
=
\left(I+h(A+B)+h^2BA\right)x.
\]

The \(h^2BA\) term is not decoration. It records that \(B\) saw a state
already changed by \(A\).

Reverse the order and the cross term becomes \(h^2AB\).

This is the first split-operator fact worth carrying into neural
architecture design:

\begin{quote}
\textbf{Sequential composition creates interaction terms that depend on
order.}
\end{quote}

\section{A telescope with two
lenses}\label{a-telescope-with-two-lenses}

A useful picture is a telescope with two different lenses.

Passing light through lens A and then lens B need not produce the same
transformation as B then A. If the lenses implement transformations that
commute, the order does not matter. If they do not commute, the ordering
is part of the instrument.

The analogy maps cleanly onto operator composition:

\begin{itemize}
\tightlist
\item
  the incoming light corresponds to the current state;
\item
  each lens corresponds to one state transformation;
\item
  lens order corresponds to composition order;
\item
  order independence corresponds to commutation.
\end{itemize}

The analogy stops there.

A neural attention sublayer is not literally an optical lens, and this
picture does not prove numerical convergence, reversibility, or
stability. Its job is only to make one structural point visible before
we formalize it.

\section{The exact reference
problem}\label{the-exact-reference-problem}

Start with the declared linear system

\[
\dot x=(A+B)x,
\]

where \(A\) and \(B\) are constant matrices.

The exact combined flow over a step \(h\) is

\[
e^{h(A+B)}.
\]

If we can solve the two parts separately, we can form the Lie-Trotter
compositions

\[
S_{AB}(h)=e^{hA}e^{hB},
\]

and

\[
S_{BA}(h)=e^{hB}e^{hA}.
\]

The first means one exact \(B\)-subflow and one exact \(A\)-subflow in
the composition order encoded by the matrix product. The second reverses
them.

Expanding both exponentials gives

\[
e^{hA}e^{hB}
=
e^{h(A+B)}
+
\frac{h^2}{2}[A,B]
+
O(h^3),
\]

where

\[
[A,B]=AB-BA.
\]

The reversed order gives

\[
e^{hB}e^{hA}
=
e^{h(A+B)}
-
\frac{h^2}{2}[A,B]
+
O(h^3).
\]

So the first commutator is the leading object that records
noncommutation.

If

\[
[A,B]=0,
\]

the two constant matrix exponentials commute and

\[
e^{hA}e^{hB}
=
e^{hB}e^{hA}
=
e^{h(A+B)}
\]

exactly.

This is not a vague statement that ``order sometimes matters.'' It tells
us what object measures the leading order difference in the declared
smooth finite-dimensional setting.

\section{The exact two-by-two
witness}\label{the-exact-two-by-two-witness}

Use

\[
A=
\begin{pmatrix}
0&1\\
0&0
\end{pmatrix},
\qquad
B=
\begin{pmatrix}
0&0\\
1&0
\end{pmatrix}.
\]

Both are nilpotent:

\[
A^2=B^2=0.
\]

Therefore

\[
e^{hA}=I+hA,
\qquad
e^{hB}=I+hB.
\]

Their products are

\[
AB=
\begin{pmatrix}
1&0\\
0&0
\end{pmatrix},
\qquad
BA=
\begin{pmatrix}
0&0\\
0&1
\end{pmatrix},
\]

so

\[
[A,B]
=
\begin{pmatrix}
1&0\\
0&-1
\end{pmatrix}.
\]

The two Lie steps are exactly

\[
S_{AB}(h)
=
\begin{pmatrix}
1+h^2&h\\
h&1
\end{pmatrix},
\]

and

\[
S_{BA}(h)
=
\begin{pmatrix}
1&h\\
h&1+h^2
\end{pmatrix}.
\]

Subtract:

\[
S_{AB}(h)-S_{BA}(h)
=
h^2[A,B].
\]

At \(h=1/2\),

\[
S_{AB}
=
\begin{pmatrix}
5/4&1/2\\
1/2&1
\end{pmatrix},
\]

while

\[
S_{BA}
=
\begin{pmatrix}
1&1/2\\
1/2&5/4
\end{pmatrix}.
\]

The two computations contain the same two exact subflows. They differ
only in order. Yet the result is different.

That is the architectural fact we want.

\section{The combined flow is not either Lie
step}\label{the-combined-flow-is-not-either-lie-step}

For this witness,

\[
A+B
=
\begin{pmatrix}
0&1\\
1&0
\end{pmatrix},
\]

and

\[
(A+B)^2=I.
\]

Hence

\[
e^{h(A+B)}
=
\begin{pmatrix}
\cosh h&\sinh h\\
\sinh h&\cosh h
\end{pmatrix}.
\]

At \(h=1/2\), neither Lie ordering equals this exact combined flow.

Their Frobenius errors are equal in this symmetric witness,
approximately

\[
0.1793148494.
\]

The equality of those two error norms is incidental to this example. The
important structural result is the signed leading commutator term.

This distinction matters when translating the picture into neural
networks. The fact that an ordered block is a useful decomposition does
not imply that it is the exact flow of some sum operator.

\section{Symmetry cancels the quadratic
defect}\label{symmetry-cancels-the-quadratic-defect}

A standard symmetric composition is Strang splitting:

\[
S_{ABA}(h)
=
e^{hA/2}e^{hB}e^{hA/2}.
\]

For sufficiently regular problems, Strang splitting has a local defect
of order \(O(h^3)\), corresponding to second-order global accuracy under
the usual numerical-analysis assumptions (McLachlan and Quispel 2002;
Hairer et al. 2006).

Our witness shows the cancellation directly.

Because \(A^2=B^2=0\),

\[
S_{ABA}(h)
=
\begin{pmatrix}
1+h^2/2&h+h^3/4\\
h&1+h^2/2
\end{pmatrix}.
\]

The exact combined-flow series is

\[
e^{h(A+B)}
=
\begin{pmatrix}
1+h^2/2+h^4/24+O(h^6)&
h+h^3/6+O(h^5)\\
h+h^3/6+O(h^5)&
1+h^2/2+h^4/24+O(h^6)
\end{pmatrix}.
\]

The quadratic mismatch has disappeared.

At \(h=1/2\),

\[
S_{ABA}
=
\begin{pmatrix}
9/8&17/32\\
1/2&9/8
\end{pmatrix},
\]

whose Frobenius error against the exact flow is approximately

\[
0.02370487547.
\]

This one number is not the reason Strang is second order. The series
derivation is the reason. The numeric checkpoint only witnesses the
declared finite case.

\section{A commuting control}\label{a-commuting-control}

Now choose

\[
A_c=\operatorname{diag}(1,2),
\qquad
B_c=\operatorname{diag}(3,4).
\]

Then

\[
[A_c,B_c]=0.
\]

Because the matrices commute,

\[
e^{hA_c}e^{hB_c}
=
e^{hB_c}e^{hA_c}
=
e^{h(A_c+B_c)}
\]

for every \(h\).

This control is important.

Without it, the chapter could leave the impression that every
composition order must differ. The correct statement is narrower:

\begin{quote}
\textbf{Noncommutation creates order dependence. Commutation removes it
in the exact setting.}
\end{quote}

The question for a learned architecture is therefore not ``does order
matter in principle?'' but ``what are the actual submaps, and how
strongly do they fail to commute on the states that matter?''

That later empirical question is not answered by this chapter.

\section{From exact flows to neural
sublayers}\label{from-exact-flows-to-neural-sublayers}

Return to the Transformer residual stream.

Let \(H\) be the current residual-stream state. A simplified
pre-normalized pair of learned submaps might be written as

\[
\Psi_A(H)=H+F_A(N_A(H)),
\]

\[
\Psi_B(H)=H+F_B(N_B(H)).
\]

Here \(F_A\) might stand for attention-like computation and \(F_B\) for
an FFN-like computation.

The ordered block is then

\[
H'
=
(\Psi_B\circ\Psi_A)(H).
\]

The reversed block is

\[
\tilde H'
=
(\Psi_A\circ\Psi_B)(H).
\]

In general,

\[
H'\neq\tilde H'.
\]

Why?

Because the second submap receives a different input state.

This remains true even before we invoke any continuous-time
interpretation.

The split-operator lens is therefore useful at two levels.

At the exact mathematical level, it gives the established theory of
flows, commutators, Lie products, symmetric compositions, and splitting
error.

At the architectural level, it gives a disciplined language for asking
how named subcomputations are isolated and ordered.

Those two levels must not be collapsed.

\section{Normalization belongs inside the
submap}\label{normalization-belongs-inside-the-submap}

A common source of false simplification is to write ``attention operator
\(A\)'' and ``MLP operator \(B\)'' while quietly omitting normalization.

In a pre-LN block, the actual maps look more like

\[
H_1
=
H+
\operatorname{Attn}(\operatorname{LN}(H)),
\]

\[
H_2
=
H_1+
\operatorname{FFN}(\operatorname{LN}(H_1)).
\]

If we call these two stages \(\Psi_A\) and \(\Psi_B\), then LayerNorm is
part of the maps.

It affects the state dependence.

The same rule applies to:

\begin{itemize}
\tightlist
\item
  masking;
\item
  gating;
\item
  clipping;
\item
  routing;
\item
  residual scaling;
\item
  stochastic dropout when active;
\item
  state-dependent normalization;
\item
  data-dependent control flow.
\end{itemize}

A commutator or order argument that drops these operations is analyzing
a different system unless an explicit approximation justifies the
omission.

\section{Symmetry is not exact
reversibility}\label{symmetry-is-not-exact-reversibility}

The word ``symmetric'' is dangerous because it has several meanings.

Strang composition is symmetric in the step sequence:

\[
A/2
\to
B
\to
A/2.
\]

For exact flows, this time-symmetric composition has important numerical
consequences.

To even write a neural analogue of that sequence, one must first define
a step-parameterized family of learned maps \[
\Psi_A(h),\qquad \Psi_B(h),
\] with a meaningful half-stage \(\Psi_A(h/2)\). Only then can one form
\[
\Psi_A(h/2)\circ\Psi_B(h)\circ\Psi_A(h/2).
\]

For a generic learned block \(\Psi_A\) with no declared step parameter,
the symbol ``half of A'' is not defined by the architecture. Halving a
residual coefficient or duplicating a block is a new design choice, not
automatically the classical half-flow.

Even when a valid step-parameterized family is declared, the resulting
neural palindromic block is not automatically invertible.

If \(\Psi_A\) or \(\Psi_B\) loses information, clips state, normalizes
non-injectively, uses non-invertible activation, drops tokens, or
otherwise fails to be bijective on the relevant domain, the palindromic
ordering does not repair that.

So we must distinguish:

\begin{itemize}
\tightlist
\item
  symmetric ordering;
\item
  self-adjoint numerical composition under a declared flow model;
\item
  exact invertibility;
\item
  practical reconstructability;
\item
  reversible-network implementation.
\end{itemize}

These are related ideas, not synonyms.

\section{Shared operators and changing
operators}\label{shared-operators-and-changing-operators}

Classical notation often suggests one fixed pair \(A,B\).

Neural networks often do not behave that way.

Layer \(k\) may use

\[
A_k,
\qquad
B_k,
\]

with different learned parameters at every depth.

Then the architecture is better viewed as stage-dependent or
nonautonomous.

A commutator such as

\[
[A_k,B_k]
\]

may still be meaningful in a local linearized model, but it does not
become one global autonomous commutator for the entire network.

This distinction separates two regimes:

\subsection{Shared operators}\label{shared-operators}

The same submaps recur across depth.

This is structurally closer to repeated stepping of one declared
dynamical system.

\subsection{Layer-varying operators}\label{layer-varying-operators}

Each stage has different submaps.

This is closer to a time-dependent system or a general composition
chain.

Both can be studied with operator language. Their mathematical contracts
differ.

\section{Splitting error is only one
error}\label{splitting-error-is-only-one-error}

Suppose a network is motivated by splitting a reference evolution.

Even then, ``the error'' is not one scalar object.

At minimum distinguish:

\begin{enumerate}
\def\labelenumi{\arabic{enumi}.}
\tightlist
\item
  \textbf{splitting/discretization error}: difference between a
  numerical composition and the declared reference flow;
\item
  \textbf{approximation error}: learned submaps may not represent the
  intended component dynamics exactly;
\item
  \textbf{estimation error}: finite data limits what can be learned;
\item
  \textbf{optimization error}: training may not reach the intended
  parameter solution;
\item
  \textbf{stochastic variation}: minibatching, initialization, routing,
  or noise may change outcomes;
\item
  \textbf{implementation error}: finite precision, kernels, masking
  details, and system behavior may differ from the mathematical model.
\end{enumerate}

Classical Lie or Strang order controls the first category under its
assumptions.

It does not automatically control the others.

This is the same epistemic discipline used throughout the Atlas: one
support route does not acquire authority over neighboring claims merely
because the notation is convenient.

\section{Ordering as an architectural degree of
freedom}\label{ordering-as-an-architectural-degree-of-freedom}

With those boundaries in place, the design lesson becomes stronger, not
weaker.

A network designer can ask:

\begin{itemize}
\tightlist
\item
  Which transformations should be isolated as distinct submaps?
\item
  Which state does each submap receive?
\item
  Should they be applied sequentially, additively, or in a symmetric
  composition?
\item
  Are parameters shared across depth?
\item
  Is there a declared continuous reference problem?
\item
  Is the architecture trying to preserve a structure?
\item
  Where can noncommutation be measured?
\item
  Does reversing order alter useful behavior?
\item
  Is a small local commutator associated with robustness, stability, or
  transfer?
\item
  When does symmetric composition help, and when is it merely extra
  compute?
\end{itemize}

These are empirical and mathematical questions.

The splitting framework organizes them without answering them in
advance.

\section{A Transformer block through the split
lens}\label{a-transformer-block-through-the-split-lens}

A baseline Transformer already presents an ordered pair of major
subcomputations: attention and position-wise feed-forward transformation
(Vaswani et al. 2017).

In a pre-LN form,

\[
H_1
=
H+
\operatorname{Attn}(\operatorname{LN}(H)),
\]

\[
H_2
=
H_1+
\operatorname{FFN}(\operatorname{LN}(H_1)).
\]

A split-operator reading does not claim that these are exponentials of
two known generators.

It asks instead:

\begin{quote}
What changes if we treat these as two named state transformations whose
order and composition law are explicit design choices?
\end{quote}

That question exposes several architecture variants immediately:

\begin{itemize}
\tightlist
\item
  attention then FFN;
\item
  FFN then attention;
\item
  additive parallel updates from the same input state;
\item
  explicitly step-parameterized palindromic attention-like / FFN-like
  compositions, when a meaningful half-step family is actually defined;
\item
  shared repeated submaps;
\item
  depth-varying submaps;
\item
  submaps with measured or constrained local interaction.
\end{itemize}

Whether any of these is better is an experimental question.

The Atlas does not promote an architectural rearrangement from
mathematical elegance alone.

\section{What a neural commutator could
mean}\label{what-a-neural-commutator-could-mean}

For linear operators, the commutator is unambiguous:

\[
[A,B]=AB-BA.
\]

For nonlinear maps, several related objects are possible.

One can compare compositions directly:

\[
\Psi_B(\Psi_A(x))
-
\Psi_A(\Psi_B(x)).
\]

One can also differentiate the nonlinear composition defect. If \[
C_\Psi(x)
=
\Psi_B(\Psi_A(x))
-
\Psi_A(\Psi_B(x)),
\] then, when the maps are differentiable, \[
DC_\Psi(x)
=
J_B(\Psi_A(x))J_A(x)
-
J_A(\Psi_B(x))J_B(x).
\]

A same-state algebraic Jacobian commutator \[
J_B(x)J_A(x)-J_A(x)J_B(x)
\] can still be used as a local diagnostic, but it is generally not the
derivative of \(C_\Psi\). Intermediate-state evaluation matters.

One can work with vector-field Lie brackets when actual differentiable
vector fields have been declared.

These are not interchangeable.

A future diagnostic chapter may define a specific operational
noncommutation probe. SPLIT-001 only establishes why such a probe could
be informative and what mathematical distinctions it must preserve.

\section{Failure modes}\label{failure-modes-1}

\subsection{Calling every residual block an
integrator}\label{calling-every-residual-block-an-integrator}

A residual form

\[
x_{k+1}=x_k+F_k(x_k)
\]

resembles an Euler step.

Resemblance is not identity.

Without a declared reference vector field, step-size semantics, and
refinement family, ``integrator'' is an interpretation rather than a
convergence theorem.

\subsection{Treating Strang order as architectural
magic}\label{treating-strang-order-as-architectural-magic}

A palindromic block can be aesthetically attractive.

The \(O(h^3)\) local defect belongs to the classical splitting setting
under the required assumptions. It does not automatically apply to
arbitrary learned nonlinear sublayers. A neural half-step must itself be
defined by a declared step-parameterized family before Strang language
has mathematical content.

\subsection{Ignoring normalization}\label{ignoring-normalization}

If normalization is present in the executed block, dropping it from the
operator model changes the object.

\subsection{Confusing symmetry and
invertibility}\label{confusing-symmetry-and-invertibility}

A symmetric sequence of non-invertible maps is still non-invertible.

\subsection{Measuring one local commutator and claiming global
control}\label{measuring-one-local-commutator-and-claiming-global-control}

A small commutator at one state, one layer, or one linearization does
not establish global near-commutation.

\subsection{Folding training error into splitting
error}\label{folding-training-error-into-splitting-error}

Optimization failure is not a discretization defect.

\subsection{Assuming shared-operator formulas for layer-varying
networks}\label{assuming-shared-operator-formulas-for-layer-varying-networks}

A different \(A_k,B_k\) at every layer is a different mathematical
object from repeated application of fixed \(A,B\).

\section{What later chapters may now
assume}\label{what-later-chapters-may-now-assume-1}

The downstream Composition chapter may assume:

\begin{itemize}
\tightlist
\item
  additive and sequential composition are distinct;
\item
  operator ordering can be architecturally load-bearing;
\item
  the commutator is the leading finite-dimensional Lie-order object in
  the declared exact setting;
\item
  symmetric composition has stronger cancellation properties under
  classical assumptions;
\item
  neural submaps must include their actual normalization, routing,
  masking, and residual semantics.
\end{itemize}

The downstream Transport chapter may assume:

\begin{itemize}
\tightlist
\item
  split computation can be interpreted as staged transport only when the
  stage maps are explicitly declared;
\item
  an exact-flow interpretation is stronger than a generic
  residual-composition interpretation;
\item
  layer-varying stages require nonautonomous semantics.
\end{itemize}

Neither chapter may inherit a claim that Transformer sublayers are exact
flows or that Strang-style neural blocks are automatically second order.

\section{Epistemic status}\label{epistemic-status}

Established theory:

\begin{itemize}
\tightlist
\item
  Lie-Trotter and Strang splitting;
\item
  commutator-controlled order effects in the appropriate operator
  setting;
\item
  composition and structure-preserving numerical methodology (McLachlan
  and Quispel 2002; Hairer et al. 2006);
\item
  baseline Transformer sublayer anatomy (Vaswani et al. 2017).
\end{itemize}

Atlas-owned derivation:

\begin{itemize}
\tightlist
\item
  the exact nilpotent two-matrix witness;
\item
  the additive-versus-sequential residual comparison;
\item
  the explicit claim firewall connecting operator-splitting theory to
  neural architecture.
\end{itemize}

Computational witness:

\begin{itemize}
\tightlist
\item
  exact matrix products;
\item
  exact commutator;
\item
  exact Lie and Strang formulas;
\item
  bounded \(h=1/2\) error values;
\item
  commuting control.
\end{itemize}

Not established here:

\begin{itemize}
\tightlist
\item
  empirical superiority of split-operator neural architectures;
\item
  small commutators in trained Transformers;
\item
  exact flow semantics for attention or FFN sublayers;
\item
  global reversibility from symmetric ordering;
\item
  general convergence guarantees for learned split networks.
\end{itemize}

That boundary is deliberate.

The value of the split-operator lens is not that it turns a neural
network into a numerical theorem.

It is that it forces us to say what the transformations are, what order
they act in, what state each sees, and which mathematical claims survive
the translation.

\section{References used in this
chapter}\label{references-used-in-this-chapter-20}

\begin{itemize}
\tightlist
\item
  (McLachlan and Quispel 2002) for Lie-Trotter/Strang splitting,
  commutators, composition, and order conditions.
\item
  (Hairer et al. 2006) for geometric/composition-method boundaries and
  structure-preserving semantics.
\item
  (Vaswani et al. 2017) for baseline Transformer attention/FFN
  residual-block anatomy.
\end{itemize}

Exact provenance and authority boundaries are locked in:

\texttt{sources/source-locks/ATLAS-CH-SPLIT-001.yaml}.

\chapter{Representation as Transport}\label{representation-as-transport}

\textbf{Epistemic status:} established geometry and dynamics + audited
prerequisite inheritance + Atlas synthesis.\\
\textbf{Specification:}
\texttt{manuscript/specifications/ATLAS-CH-TRANSPORT-001.md}\\
\textbf{Derivation packet:}
\texttt{mathematics/derivations/ATLAS-CH-TRANSPORT-001-DERIVATIONS.md}\\
\textbf{Computational witness:}
\texttt{mathematics/computational-witnesses/ATLAS-CW-TRANSPORT-001.md}\\
\textbf{Source lock:}
\texttt{sources/source-locks/ATLAS-CH-TRANSPORT-001.yaml}

A representation does not merely exist. In a deep adaptive system it is
repeatedly changed.

The same hidden state may be:

\begin{itemize}
\tightlist
\item
  mixed by attention;
\item
  transformed by a feed-forward map;
\item
  normalized;
\item
  gated;
\item
  routed;
\item
  projected;
\item
  constrained;
\item
  carried through a residual path;
\item
  or updated by a continuous-depth evolution law.
\end{itemize}

The useful question is therefore not only:

\begin{quote}
What is the representation?
\end{quote}

It is also:

\begin{quote}
How is the representation moved?
\end{quote}

This chapter develops a disciplined answer.

The word \textbf{transport} will mean the evolution of a declared
representation state through declared maps.

That definition is deliberately broad enough to cover finite neural
compositions and constrained updates, but deliberately narrow enough to
block several false identifications.

Representation transport is not automatically:

\begin{itemize}
\tightlist
\item
  optimal transport of probability mass;
\item
  parallel transport under a connection;
\item
  vector transport of tangent data;
\item
  geodesic motion;
\item
  an exact ODE flow.
\end{itemize}

Those objects may enter a model when their assumptions are actually
present. They are not consequences of the word ``transport.''

\section{The primary object is a sequence of
maps}\label{the-primary-object-is-a-sequence-of-maps}

Let the representation state at stage \(k\) be

\[
z_k\in\mathcal Z_k.
\]

A stage map is

\[
\Phi_k:\mathcal Z_k\to\mathcal Z_{k+1}.
\]

The update is

\[
z_{k+1}=\Phi_k(z_k).
\]

After \(K\) stages,

\[
z_K
=
(\Phi_{K-1}\circ\cdots\circ\Phi_0)(z_0).
\]

This composition is exact as written.

It does not require:

\begin{itemize}
\tightlist
\item
  an ODE;
\item
  an infinitesimal step;
\item
  a manifold;
\item
  invertibility;
\item
  a metric;
\item
  a probability distribution.
\end{itemize}

That is the first transport object of the Atlas:

\begin{quote}
\textbf{a finite representation trajectory is the ordered composition of
the maps that actually execute.}
\end{quote}

This sounds elementary. It is also where many conceptual mistakes begin.

If the executed map contains normalization, masking, routing, clipping,
stochasticity, or projection, those operations belong to \(\Phi_k\).
Omitting them means analyzing another system.

\section{Residual transport}\label{residual-transport-1}

A residual stage has the familiar form

\[
\Phi_k(z)
=
z+F_k(z).
\]

The identity term gives a direct state path. The learned residual
modifies it.

Residual networks made this identity-plus-increment structure an
explicit architecture (He et al. 2016).

At the discrete level, the statement is exact:

\[
\text{next state}
=
\text{current state}
+
\text{learned increment}.
\]

The dynamical interpretation is useful, but it must come second.

If we instead write

\[
\Phi_{k,h}(z)
=
z+hF_k(z),
\]

the algebra resembles an explicit Euler step.

That resemblance becomes a numerical statement only after a reference
evolution law has been declared.

The equation

\[
z_{k+1}=z_k+hF_k(z_k)
\]

does not by itself tell us that one unique continuous vector field
generated the network.

This is inherited doctrine from the audited Split-Operator and numerical
chapters.

\section{Continuous depth is a different
declaration}\label{continuous-depth-is-a-different-declaration}

Neural ODEs make a stronger object explicit:

\[
\frac{dz}{dt}
=
f(z,t;\theta).
\]

The architecture defines a continuous-depth evolution law and obtains
finite-time state changes through the associated ODE solution and
numerical integration (Chen et al. 2018).

This is a genuine continuous-time transport object.

A residual stack and a Neural ODE can be related.

They are not identical by definition.

The safe implication is:

\[
\text{explicit ODE architecture}
\Rightarrow
\text{declared continuous state evolution}.
\]

The unsafe implication is:

\[
\text{residual stack}
\Rightarrow
\text{unique exact ODE}.
\]

Haber and Ruthotto use dynamical-systems and discrete-ODE stability
ideas to motivate deep architecture design (Haber and Ruthotto 2018).
That bridge is valuable precisely because it makes assumptions visible.
It should not be turned into an identity between all deep networks and
one continuous flow.

\section{Stage index is not automatically physical
time}\label{stage-index-is-not-automatically-physical-time}

The index

\[
k=0,1,\ldots,K
\]

orders computation.

It need not represent physical time.

Even when one speaks of ``depth as time,'' several meanings are
possible:

\begin{itemize}
\tightlist
\item
  literal continuous time in an ODE;
\item
  numerical time in a discretization;
\item
  recurrent iteration count;
\item
  transformer block index;
\item
  optimization iteration;
\item
  adaptive computational depth.
\end{itemize}

Transport language becomes useful when these are distinguished rather
than blended.

A chapter that says ``the representation moves through time'' must say
what the time variable is.

\section{The constrained case}\label{the-constrained-case}

Suppose representation states lie on a manifold

\[
M\subseteq\mathbb R^n.
\]

Now the update has two separate questions:

\begin{enumerate}
\def\labelenumi{\arabic{enumi}.}
\tightlist
\item
  what direction is proposed?
\item
  how is a legal endpoint on \(M\) produced?
\end{enumerate}

At state \(z_k\), a tangent proposal satisfies

\[
\xi_k\in T_{z_k}M.
\]

A retraction maps tangent data back to the manifold:

\[
z_{k+1}
=
R_{z_k}(\xi_k).
\]

For the unit sphere,

\[
S^{n-1}
=
\{u:\|u\|_2=1\},
\]

the tangent condition is

\[
u^\top\xi=0.
\]

A standard normalized retraction is

\[
R_u(\xi)
=
\frac{u+\xi}{\|u+\xi\|_2}.
\]

This guarantees

\[
\|R_u(\xi)\|_2=1.
\]

It guarantees feasibility.

Nothing in that equation says the update is:

\begin{itemize}
\tightlist
\item
  geodesic;
\item
  optimal;
\item
  invertible;
\item
  information preserving;
\item
  dynamically stable;
\item
  task improving.
\end{itemize}

The NORMREP prerequisite established exactly this geometry and the
radial-information boundary (Lee 2018; Absil et al. 2008).

\section{Feasible motion is not geodesic
motion}\label{feasible-motion-is-not-geodesic-motion}

Take the unit circle.

Let

\[
u=(1,0),
\qquad
\xi=(0,1).
\]

The normalized retraction gives

\[
R_u(\xi)
=
\frac{(1,1)}{\sqrt2}.
\]

This point is on the circle.

Its angular displacement from \(u\) is

\[
\frac{\pi}{4}.
\]

Now use the sphere exponential map with the same unit tangent vector at
unit time.

The endpoint is

\[
\operatorname{Exp}_u(\xi)
=
(\cos1,\sin1).
\]

Its angular displacement is

\[
1.
\]

Since

\[
\frac{\pi}{4}\neq1,
\]

the endpoints differ.

This is the smallest exact warning we need:

\begin{quote}
\textbf{returning to the manifold is not the same thing as following the
manifold geodesic.}
\end{quote}

A retraction is often computationally convenient because it approximates
the exponential map locally (Absil et al. 2008).

Approximation is not identity.

\section{Moving the point and moving an arrow are different
operations}\label{moving-the-point-and-moving-an-arrow-are-different-operations}

The phrase ``transport'' is overloaded in geometry.

Suppose

\[
v_0\in T_{u_0}M.
\]

After the base point moves from \(u_0\) to \(u_1\), the same ambient
coordinate vector need not belong to

\[
T_{u_1}M.
\]

This matters whenever a system carries local information such as:

\begin{itemize}
\tightlist
\item
  momentum;
\item
  a tangent update;
\item
  a local search direction;
\item
  a basis;
\item
  a differential perturbation;
\item
  optimizer state defined intrinsically on a manifold.
\end{itemize}

The state update answers:

\begin{quote}
Where did the point move?
\end{quote}

A vector-transport or parallel-transport rule answers another question:

\begin{quote}
How should tangent information be represented at the new point?
\end{quote}

Those are not the same map.

\section{Exact tangent-ownership
witness}\label{exact-tangent-ownership-witness}

Take

\[
u_0=(1,0),
\qquad
v_0=(0,1).
\]

Initially,

\[
u_0^\top v_0=0,
\]

so

\[
v_0\in T_{u_0}S^1.
\]

Move the state by normalized retraction:

\[
u_1
=
\frac{(1,1)}{\sqrt2}.
\]

Now test the old tangent vector against the new base point:

\[
u_1^\top v_0
=
\frac1{\sqrt2}.
\]

Therefore

\[
v_0\notin T_{u_1}S^1.
\]

The same ambient coordinates are no longer legal tangent data.

The orthogonal projection onto the new tangent space is

\[
P_{u_1}v_0
=
v_0-(u_1^\top v_0)u_1
=
\left(-\frac12,\frac12\right).
\]

This projected vector is tangent because

\[
u_1^\top P_{u_1}v_0=0.
\]

Projection is one possible local operation.

It is not automatically parallel transport.

It is not automatically the best vector transport for a learning
algorithm.

The durable rule is simply:

\begin{quote}
\textbf{tangent data belongs to a base point.}
\end{quote}

\section{The local-compass
allegory}\label{the-local-compass-allegory}

Imagine moving across the surface of a globe while carrying a small
arrow that lies flat against the surface.

Two things happen:

\begin{itemize}
\tightlist
\item
  your location changes;
\item
  the meaning of ``this arrow lies tangent here'' changes with location.
\end{itemize}

Moving the location does not automatically tell you how to carry the
arrow.

That is the geometric distinction between:

\begin{itemize}
\tightlist
\item
  state transport;
\item
  tangent-data transport.
\end{itemize}

The analogy has a strict limit.

A neural representation space may not be a smooth manifold at all. Even
when we constrain vectors to a sphere, the complete learned computation
can include discrete routing, masking, non-smooth nonlinearities, and
dimension-changing interfaces.

The allegory is useful only where the declared geometry supports it.

\section{Path order can matter even when every endpoint is
feasible}\label{path-order-can-matter-even-when-every-endpoint-is-feasible}

The Split-Operator prerequisite established that ordered maps need not
commute.

Constrained transport adds another useful witness.

Start on

\[
S^2
\]

at

\[
u_0=(1,0,0).
\]

Choose

\[
a=(0,1,0),
\qquad
b=(0,0,1).
\]

Both are tangent at \(u_0\).

Use normalized retraction

\[
R_u(\xi)
=
\frac{u+\xi}{\|u+\xi\|}.
\]

Apply \(a\) first:

\[
u_a
=
\frac{(1,1,0)}{\sqrt2}.
\]

The vector \(b\) remains tangent there, because

\[
u_a^\top b=0.
\]

Apply the second stage:

\[
u_{ab}
=
R_{u_a}(b)
=
\left(
\frac12,
\frac12,
\frac1{\sqrt2}
\right).
\]

Reverse the order:

\[
u_b
=
\frac{(1,0,1)}{\sqrt2},
\]

then

\[
u_{ba}
=
R_{u_b}(a)
=
\left(
\frac12,
\frac1{\sqrt2},
\frac12
\right).
\]

Both endpoints are feasible:

\[
\|u_{ab}\|_2
=
\|u_{ba}\|_2
=
1.
\]

But

\[
u_{ab}\neq u_{ba}.
\]

The exact inner product is

\[
u_{ab}^\top u_{ba}
=
\frac14+\frac1{\sqrt2}.
\]

Constraint preservation therefore does not erase path ordering.

This is not a universal theorem that every pair of constrained neural
updates must fail to commute.

It is a finite counterexample to the opposite claim.

\section{The executed maps determine the
path}\label{the-executed-maps-determine-the-path}

If the network executes

\[
z_{k+1}
=
R_{z_k}\bigl(\xi_k(z_k)\bigr),
\]

then both objects matter:

\begin{itemize}
\tightlist
\item
  the field or proposal \(\xi_k\);
\item
  the endpoint map \(R\).
\end{itemize}

Replacing the retraction changes the transport.

Replacing the order changes the transport.

Dropping normalization changes the transport.

Changing the base point before evaluating the next stage changes the
transport.

This is the correct level at which architectural comparisons should
begin.

\section{State-dependent transport}\label{state-dependent-transport}

Many neural maps are state-dependent.

Attention is the clearest example: the operator applied at one stage
depends on the current state.

So a map such as

\[
z_{k+1}
=
\Phi_k(z_k)
\]

is generally not a fixed matrix multiplying every possible input.

Its differential at one point is

\[
J_{\Phi_k}(z_k).
\]

That Jacobian can describe local perturbation propagation.

It is another object again.

We must separate:

\begin{itemize}
\tightlist
\item
  the nonlinear state map;
\item
  its local Jacobian;
\item
  a tangent-space transport rule;
\item
  a global flow, if one exists.
\end{itemize}

These objects can interact without being interchangeable.

\section{Differential pushforward is not automatically vector
transport}\label{differential-pushforward-is-not-automatically-vector-transport}

For a differentiable map,

\[
\Phi:M\to N,
\]

the differential maps tangent vectors by

\[
D\Phi_z:
T_zM
\to
T_{\Phi(z)}N.
\]

This is a precise geometric construction when the required smooth
manifold setting exists.

In ordinary neural calculations one often works instead with an ambient
Jacobian

\[
J_\Phi(z).
\]

If constraints, quotient structure, projections, or non-smooth
operations are present, the relation between the ambient Jacobian and an
intrinsic tangent-space map must be stated.

Calling every Jacobian multiplication ``vector transport'' would erase
these distinctions.

The Atlas will not do that.

\section{Transport can lose
information}\label{transport-can-lose-information}

A constrained path can remain perfectly feasible while destroying
distinctions.

The simplest example is normalization:

\[
N(x)
=
\frac{x}{\|x\|}.
\]

For every

\[
c>0,
\]

\[
N(cx)=N(x).
\]

The positive radial degree of freedom disappears.

If the original representation stored useful information in norm, the
transport has erased it unless another channel carries that quantity.

Therefore:

\[
\text{feasible constrained transport}
\not\Rightarrow
\text{information-preserving transport}.
\]

This is inherited directly from the NORMREP chapter.

\section{Transport can also relocate
information}\label{transport-can-also-relocate-information}

Information removed from one coordinate channel can reappear elsewhere.

A normalized architecture may retain magnitude through:

\begin{itemize}
\tightlist
\item
  explicit scale parameters;
\item
  residual amplitudes;
\item
  gates;
\item
  temperatures;
\item
  separate scalar channels;
\item
  logits;
\item
  routing variables.
\end{itemize}

The right question is not whether normalization ``destroys scale'' in
some absolute sense.

It is:

\begin{quote}
Which distinctions survive in the complete state carried forward?
\end{quote}

Transport is therefore an interface-level concept.

A local geometric constraint does not describe the full system state
unless the system state has actually been defined that way.

\section{Shared and stage-dependent
transport}\label{shared-and-stage-dependent-transport}

Suppose every layer uses the same map:

\[
\Phi_k=\Phi.
\]

Then

\[
z_K=\Phi^K(z_0).
\]

This repeated-map structure supports one kind of dynamical
interpretation.

Now suppose the parameters change at every layer:

\[
\Phi_0,\Phi_1,\ldots,\Phi_{K-1}.
\]

Then the system is stage-dependent.

A continuous analogy, if one is introduced, is naturally nonautonomous:

\[
\dot z=f(z,t).
\]

The chapter does not collapse these two cases.

A single autonomous generator is a stronger statement than an ordered
family of learned maps.

\section{Stability is a separate
question}\label{stability-is-a-separate-question}

Transport language does not answer whether perturbations grow.

For two nearby states,

\[
z_k,\qquad z_k+\delta z_k,
\]

a local linearization gives

\[
\delta z_{k+1}
\approx
J_{\Phi_k}(z_k)\delta z_k.
\]

Across many stages,

\[
\delta z_K
\approx
J_{\Phi_{K-1}}(z_{K-1})
\cdots
J_{\Phi_0}(z_0)
\delta z_0.
\]

The product may:

\begin{itemize}
\tightlist
\item
  contract;
\item
  preserve;
\item
  amplify;
\item
  rotate;
\item
  shear;
\item
  become highly non-normal.
\end{itemize}

Those are stability and sensitivity questions.

They do not follow from the fact that the base representation path
remains feasible.

Haber--Ruthotto's dynamical-systems perspective is relevant here (Haber
and Ruthotto 2018), but no universal stability claim is inherited.

\section{Invertibility is also
separate}\label{invertibility-is-also-separate}

A map can transport a state forward while discarding information.

Normalization is one example.

ReLU is another familiar non-injective operation on part of its domain.

Pooling, clipping, quantization, token dropping, and many routing
decisions can also be non-invertible.

Therefore

\[
\text{state evolution}
\not\Rightarrow
\text{reversible state evolution}.
\]

A reversible architecture requires its own construction and proof
obligations.

\section{This is not optimal
transport}\label{this-is-not-optimal-transport}

Optimal transport studies the movement of mass or probability measures
under a specified cost and admissible coupling/map structure.

Nothing in

\[
z_{k+1}=\Phi_k(z_k)
\]

by itself defines:

\begin{itemize}
\tightlist
\item
  a probability measure;
\item
  a source and target distribution;
\item
  a coupling;
\item
  a transport cost;
\item
  an optimality problem.
\end{itemize}

Representation transport may eventually be connected to
optimal-transport ideas in a separately specified model.

The current chapter does not make that identification.

\section{This is not parallel
transport}\label{this-is-not-parallel-transport}

Parallel transport requires geometric structure specifying how tangent
vectors are carried along curves, usually through a connection.

The state path

\[
z_0\to z_1\to\cdots\to z_K
\]

does not automatically define such a rule.

Even on a sphere, ``project the old tangent vector onto the new tangent
plane'' and ``parallel transport the tangent vector along the connecting
geodesic'' are different constructions.

The Atlas uses the terminology literally.

\section{This is not vector transport
either}\label{this-is-not-vector-transport-either}

Optimization on manifolds often introduces a \textbf{vector transport}
to move tangent vectors from one tangent space to another in a
computationally convenient way (Absil et al. 2008).

That is closer to the machine-learning setting than abstract parallel
transport, but it remains distinct from moving the representation point
itself.

The notation should therefore preserve two maps when both are present:

\[
z_{k+1}
=
R_{z_k}(\xi_k),
\]

and

\[
v_{k+1}
=
\mathcal T_{\xi_k}(v_k).
\]

The first moves state.

The second moves tangent information.

A system that carries optimizer momentum on a manifold may need both.

\section{Transport and split
operators}\label{transport-and-split-operators}

Suppose one representation block contains two declared state maps,

\[
\Phi_A,
\qquad
\Phi_B.
\]

Then chronological A-then-B transport is

\[
\Phi_B\circ\Phi_A.
\]

Chronological B-then-A transport is

\[
\Phi_A\circ\Phi_B.
\]

The Split-Operator chapter already established why these can differ and
why execution order must be explicit.

TRANSPORT adds the constraint-aware reading:

\begin{itemize}
\tightlist
\item
  the first stage changes the base state seen by the second;
\item
  tangent spaces may also change;
\item
  a constraint-restoring map can introduce additional path dependence;
\item
  exact-flow commutator theory must not be imported unless its
  assumptions hold.
\end{itemize}

This is a synthesis, not a new numerical theorem.

\section{Transport and residual identity
paths}\label{transport-and-residual-identity-paths}

Residual architectures are especially natural to describe as transport
because the identity path makes the old state explicitly available to
the next stage.

If

\[
F_k(z)=0,
\]

then

\[
\Phi_k(z)=z.
\]

That is exact identity transport at that stage.

It does not prove:

\begin{itemize}
\tightlist
\item
  that information remains unchanged when \(F_k\neq0\);
\item
  that gradients propagate perfectly;
\item
  that the full model is invertible;
\item
  that the representation follows a geodesic;
\item
  that the architecture is stable.
\end{itemize}

Identity transport is one structural ingredient.

\section{Transport as an interface
contract}\label{transport-as-an-interface-contract}

The most useful role of this chapter is to define what a downstream
method is allowed to assume.

A downstream algorithm should be able to ask:

\begin{itemize}
\tightlist
\item
  What is the state space?
\item
  What is the current base state?
\item
  What map advances it?
\item
  Is the map linear or nonlinear?
\item
  Is it shared or stage-dependent?
\item
  Is there a constraint?
\item
  Is the update tangent?
\item
  What returns the state to the constraint set?
\item
  Is tangent history being carried?
\item
  If so, by what rule?
\item
  Is there an exact continuous flow or only a finite composition?
\item
  Which information is intentionally invariant or discarded?
\end{itemize}

Those questions turn ``representation transport'' from a metaphor into
an interface.

\section{Handoff to Neural Krylov
Transport}\label{handoff-to-neural-krylov-transport}

The next downstream chapter is

\texttt{ATLAS-CH-NEURALKRYLOV-001}.

That chapter will explore short-horizon preconditioned solves as
representation computation.

TRANSPORT supplies only the representation-side interface.

For example, if a differentiable state map is

\[
z_{k+1}=\Phi_k(z_k),
\]

one may study its local linearization

\[
J_{\Phi_k}(z_k).
\]

A Krylov construction could then be posed for a declared operator built
from that local object.

But classical Krylov theory requires its own assumptions:

\begin{itemize}
\tightlist
\item
  a specified linear operator;
\item
  a generated Krylov subspace;
\item
  a projection or residual criterion;
\item
  conditioning/preconditioning semantics;
\item
  finite-precision limits;
\item
  convergence hypotheses.
\end{itemize}

None of those arrive for free from the transport language.

The hard boundary is:

\[
\text{representation transport}
\not\Rightarrow
\text{Krylov convergence guarantee}.
\]

\section{Five distinctions to keep
visible}\label{five-distinctions-to-keep-visible}

\subsection{State transport versus tangent
transport}\label{state-transport-versus-tangent-transport}

Moving \(z_k\) is not the same problem as moving \(v_k\in T_{z_k}M\).

\subsection{Feasibility versus geodesic
exactness}\label{feasibility-versus-geodesic-exactness}

A retraction can keep a point on the manifold without following the
exponential map.

\subsection{Discrete composition versus continuous
flow}\label{discrete-composition-versus-continuous-flow}

A finite neural stack is exact as a composition of maps, not
automatically as an ODE solution.

\subsection{Constraint preservation versus information
preservation}\label{constraint-preservation-versus-information-preservation}

A normalized path can remain feasible while erasing radial distinctions.

\subsection{Order versus ingredients}\label{order-versus-ingredients}

The same named stages applied in another order can produce another
endpoint.

\section{Failure modes}\label{failure-modes-2}

The transport language fails when it is used to smuggle in stronger
mathematics than the model contains.

Common failures are:

\begin{itemize}
\tightlist
\item
  calling any sequence of activations an ODE trajectory;
\item
  calling any normalized update geodesic motion;
\item
  calling state evolution vector transport;
\item
  calling tangent projection parallel transport;
\item
  calling representation movement optimal transport without measures and
  costs;
\item
  treating unit norm as information preservation;
\item
  treating feasibility as stability;
\item
  treating repeated residual form as one autonomous generator;
\item
  treating a local Jacobian as the global nonlinear map;
\item
  transferring Krylov convergence claims to learned nonlinear transport
  by analogy.
\end{itemize}

Each mistake replaces a declared object with a stronger undeclared one.

\section{Closing view}\label{closing-view-16}

A deep model can be read as a machine for moving representations.

That reading is useful when it remains literal.

The representation at one stage is a state.

The next stage applies a map.

Constraints may restrict where the state is allowed to live.

Retractions may restore feasibility.

Tangent data may require its own transport rule.

Order may change the path.

Continuous-depth models may provide an actual flow.

Finite residual networks may provide only a composition.

The durable Atlas synthesis is therefore:

\[
\boxed{
\text{representation transport}
=
\text{declared state}
+
\text{declared stage maps}
+
\text{declared constraints}
+
\text{declared path semantics}.
}
\]

Everything stronger must be earned separately.

\section{References used in this
chapter}\label{references-used-in-this-chapter-21}

\begin{itemize}
\tightlist
\item
  (Lee 2018)
\item
  (Absil et al. 2008)
\item
  (He et al. 2016)
\item
  (Chen et al. 2018)
\item
  (Haber and Ruthotto 2018)
\end{itemize}

Exact provenance and claim boundaries are locked in:

\texttt{sources/source-locks/ATLAS-CH-TRANSPORT-001.yaml}.

\chapter{Attention as an Operator}\label{attention-as-an-operator}

\textbf{Epistemic status:} mathematical exposition built from the
standard Transformer equations, source-locked kernel interpretations,
and Atlas-owned operator decompositions.\\
\textbf{Primary figure:} `ATLAS-FIG-ATTNOP-001`\\
\textbf{Derivation packet:}
`mathematics/derivations/ATLAS-CH-ATTNOP-001-DERIVATIONS.md`

\section{From an attention picture to an attention
action}\label{from-an-attention-picture-to-an-attention-action}

Attention is often introduced as a matrix of weights. The matrix is
useful, but it is easy to confuse a picture of coefficients with the
computation those coefficients perform.

For one attention head there are at least four distinct objects:

\[
\text{scores}
\longrightarrow
\text{mixing operator}
\longrightarrow
\text{value field}
\longrightarrow
\text{output}.
\]

The score matrix records pairwise compatibility. Softmax turns each
score row into normalized mixing coefficients. Those coefficients define
an operator acting on value vectors. The full attention transformation
is nonlinear because both the operator and the values depend on the
input state.

The useful allegory is a \textbf{dynamic switching board}.

Incoming channels carry values. The present state determines how
strongly each channel contributes to each output. The board is not
fixed; queries and keys reconfigure it.

The correspondence is:

\begin{itemize}
\tightlist
\item
  compatibility signal ↔ query-key score;
\item
  switch configuration ↔ normalized mixing operator;
\item
  incoming channels ↔ value vectors;
\item
  outgoing channels ↔ mixed outputs.
\end{itemize}

The limit matters. Softmax attention is not discrete switching. It is
differentiable weighted mixing, often dense, and the weights themselves
are computed from the state. The board is a mnemonic for state-dependent
transformation, not a literal circuit.

\section{Standard scaled dot-product
attention}\label{standard-scaled-dot-product-attention}

Let

\[
X\in\mathbb R^{n\times d}
\]

contain \(n\) token representations.

For one head,

\[
Q=XW_Q,\qquad
K=XW_K,\qquad
V=XW_V.
\]

The score matrix is

\[
S(X)
=
\frac{QK^\top}{\sqrt{d_k}}.
\]

Row-wise softmax produces

\[
A(X)
=
\operatorname{softmax}_{\rm row}(S(X)),
\]

and the head output is

\[
Y=A(X)V.
\]

This is the standard scaled dot-product attention construction (Vaswani
et al. 2017).

The Atlas reserves different names for different roles:

\[
\boxed{
S=\text{score matrix}
}
\]

\[
\boxed{
A=\text{normalized mixing operator}
}
\]

\[
\boxed{
V=\text{value field}
}
\]

\[
\boxed{
F(X)=A(X)V(X)=\text{full attention transformation}.
}
\]

This distinction prevents several category errors. The score matrix is
not the normalized operator. The operator is not the value field. And
the fact that \(Y=AV\) is linear in \(V\) for fixed \(A\) does not make
\(F\) linear in \(X\).

\section{Why ``operator'' is the right
word}\label{why-operator-is-the-right-word}

Fix \(Q\) and \(K\). Then \(A\) is fixed.

For two value fields \(V_1,V_2\) and scalars \(\alpha,\beta\),

\[
A(\alpha V_1+\beta V_2)
=
\alpha AV_1+\beta AV_2.
\]

So

\[
V\mapsto AV
\]

is an ordinary linear operator on the value field.

The operator language changes the question.

A heat-map reading asks:

\begin{quote}
Which coefficients are large?
\end{quote}

An operator reading asks:

\begin{quote}
What transformation is being performed on the information carried by the
values?
\end{quote}

The second question naturally leads to rank, singular values,
invariants, sparsity, support constraints, approximation, perturbation,
and composition.

Those are properties of transformations rather than images.

\section{Softmax gives a row-stochastic mixing
stage}\label{softmax-gives-a-row-stochastic-mixing-stage}

For finite unmasked scores,

\[
A_{ij}
=
\frac{e^{S_{ij}}}
{\sum_{\ell=1}^n e^{S_{i\ell}}}.
\]

Therefore

\[
A_{ij}>0
\]

and

\[
\sum_{j=1}^n A_{ij}=1.
\]

Equivalently,

\[
\boxed{A\mathbf 1=\mathbf 1.}
\]

Each output row is a convex combination of value rows before any later
output projection.

There is a useful but limited Markov analogy here. The rows lie on
probability simplices. But an attention head is not generally a fixed
Markov chain: the operator changes with the input, values are learned
feature vectors rather than probabilities, and later transformations
change the state again.

The algebraic fact is row-stochasticity. The Markov story is only an
intuition aid.

\section{The full self-attention map is
nonlinear}\label{the-full-self-attention-map-is-nonlinear}

Conditional linearity in \(V\) can produce the opposite mistake: calling
the whole self-attention layer linear.

It is not.

Take two scalar tokens,

\[
X=
\begin{pmatrix}
1\\
0
\end{pmatrix},
\qquad
W_Q=W_K=W_V=1,
\qquad
d_k=1.
\]

Then

\[
Q=K=V=X,
\]

and

\[
S=XX^\top.
\]

The first output component is

\[
F(X)_1=\frac{e}{1+e}.
\]

Now double the input. The scores scale quadratically, so

\[
F(2X)_1
=
\frac{2e^4}{1+e^4}.
\]

But

\[
2F(X)_1
=
\frac{2e}{1+e}.
\]

Therefore

\[
F(2X)\neq 2F(X).
\]

Numerically the first-component discrepancy has magnitude about
\(0.5019104228\).

The exact statement is:

\begin{quote}
for fixed queries and keys, attention is linear in the value field;
ordinary self-attention is nonlinear in the hidden state because the
operator itself depends on that state.
\end{quote}

\section{A small operator we can inspect
completely}\label{a-small-operator-we-can-inspect-completely}

The primary figure uses three tokens:

\[
Q=K=
\begin{pmatrix}
1&0\\
0&1\\
1&1
\end{pmatrix},
\qquad
V=
\begin{pmatrix}
1&0\\
0&1\\
1&-1
\end{pmatrix},
\qquad
d_k=2.
\]

Then

\[
S
=
\frac1{\sqrt2}
\begin{pmatrix}
1&0&1\\
0&1&1\\
1&1&2
\end{pmatrix}.
\]

Wolfram evaluation gives

\[
A\approx
\begin{pmatrix}
0.401112&0.197776&0.401112\\
0.197776&0.401112&0.401112\\
0.248255&0.248255&0.503490
\end{pmatrix},
\]

and

\[
Y=AV\approx
\begin{pmatrix}
0.802224&-0.203336\\
0.598888&0\\
0.751745&-0.255235
\end{pmatrix}.
\]

\begin{figure}
\centering
\pandocbounded{\includegraphics[keepaspectratio,alt={Four numeric matrix panels showing the score matrix S, row-softmax mixing operator A, value field V, and output Y equals A V.}]{figures/masters/ATLAS-FIG-ATTNOP-001.png}}
\caption{Four numeric matrix panels showing the score matrix S,
row-softmax mixing operator A, value field V, and output Y equals A V.}
\end{figure}

The gray intensity is deliberately redundant. Every cell carries its
numerical value. The plate is intended to be read as a transformation
rather than interpreted as a color pattern.

The computational witness verifies two bounded claims:

\[
A\mathbf 1=\mathbf 1
\]

to numerical precision, and, with \(A\) fixed,

\[
A(2V)=2AV.
\]

\section{The operator changes when the query
changes}\label{the-operator-changes-when-the-query-changes}

Keep \(K\) and \(V\) fixed, but perturb the first query:

\[
q_1=(1,0)
\quad\longrightarrow\quad
q_1'=(1,\tfrac12).
\]

Only the first score row changes. Consequently, only the first row of
the mixing operator changes.

The replay gives

\[
\Delta A_{1,:}
\approx
(-0.08124593,\;0.02683053,\;0.05441540).
\]

This exposes a useful local chain:

\[
\text{state perturbation}
\longrightarrow
\text{operator perturbation}
\longrightarrow
\text{changed information mixing}.
\]

Later chapters will ask whether a perturbation is causally important.
Here the purpose is narrower: to make state dependence explicit.

\section{Permutation equivariance before position
enters}\label{permutation-equivariance-before-position-enters}

Suppose no positional signal or asymmetric mask is present.

Let \(P\) permute token order:

\[
X'=PX.
\]

Then

\[
Q'=PQ,\qquad
K'=PK,\qquad
V'=PV.
\]

Thus

\[
S'
=
\frac{Q'K'^\top}{\sqrt{d_k}}
=
PSP^\top.
\]

Row-wise softmax respects the simultaneous row/column permutation:

\[
A'=PAP^\top.
\]

Therefore

\[
Y'
=
A'V'
=
PAP^\top PV
=
PY.
\]

Hence

\[
\boxed{F(PX)=PF(X).}
\]

This is permutation \textbf{equivariance}, not invariance. Reordering
inputs reorders outputs in the same way.

The result is useful precisely because sequence models need
order-sensitive structure somewhere. Positional encodings and causal
masks modify the admissible operator family and break arbitrary
permutation symmetry.

\section{A mask is a constraint on the operator
family}\label{a-mask-is-a-constraint-on-the-operator-family}

For causal attention, define

\[
M_{ij}
=
\begin{cases}
0,&j\le i,\\
-\infty,&j>i.
\end{cases}
\]

Then

\[
A
=
\operatorname{softmax}_{\rm row}(S+M)
\]

satisfies

\[
A_{ij}=0
\qquad\text{for }j>i.
\]

The usual phrase says that a token ``cannot attend to the future.''

The operator view says something slightly sharper:

\begin{quote}
causal masking restricts the admissible support pattern of the mixing
operator.
\end{quote}

This places causal masking in the same mathematical family as local
windows, graph adjacency, block sparsity, and learned sparse support.

\section{Multihead attention is a family of state-dependent
operators}\label{multihead-attention-is-a-family-of-state-dependent-operators}

For head \(h\),

\[
Q_h=XW_Q^{(h)},\qquad
K_h=XW_K^{(h)},\qquad
V_h=XW_V^{(h)},
\]

and

\[
A_h(X)
=
\operatorname{softmax}_{\rm row}
\left(
\frac{Q_hK_h^\top}{\sqrt{d_k}}
\right).
\]

Each head produces

\[
Y_h=A_hV_h.
\]

The layer then concatenates and projects the head outputs.

A useful abstraction is therefore

\[
X
\longmapsto
\{A_1(X),\ldots,A_H(X)\}
\longmapsto
\{A_1V_1,\ldots,A_HV_H\}
\longmapsto
Y.
\]

This is more informative than saying that a layer contains several
attention maps. It contains a family of separately parameterized,
state-dependent mixing operators whose outputs are recombined.

That viewpoint prepares later questions about head specialization,
redundancy, substitution, spectral structure, and relative-position
modes.

\section{The kernel lens}\label{the-kernel-lens}

Attention can also be written as normalized similarity-weighted
smoothing. Tsai et al.~develop an explicit kernel-smoother
interpretation of Transformer attention (Tsai et al. 2019).

In generic form,

\[
y_i
=
\frac{\sum_j k(q_i,k_j)v_j}
{\sum_j k(q_i,k_j)}.
\]

For

\[
k(q,k)
=
\exp\left(
\frac{q^\top k}{\sqrt{d_k}}
\right),
\]

this recovers ordinary softmax attention.

The lens is useful because it isolates the similarity function. But it
does not erase the rest of the architecture. Queries and keys are
learned and state-dependent; masks and position alter interactions;
multihead attention constructs several such operators and recombines
them.

The disciplined conclusion is:

\begin{quote}
kernel language is one productive mathematical lens on attention, not a
universal identification of every attention mechanism with one fixed
kernel machine.
\end{quote}

\section{Feature maps and linear
attention}\label{feature-maps-and-linear-attention}

Suppose the similarity has a feature representation

\[
k(q,k)=\phi(q)^\top\phi(k).
\]

Then a numerator

\[
\sum_j
\phi(q_i)^\top\phi(k_j)v_j^\top
\]

can be reassociated:

\[
\phi(q_i)^\top
\left(
\sum_j \phi(k_j)v_j^\top
\right).
\]

This is the algebraic move exploited by linear-attention constructions
such as Katharopoulos et al. (Katharopoulos et al. 2020).

From the Atlas viewpoint, the interesting fact is not merely lower
asymptotic cost. The representation of the operator action has changed.

That prepares a later question:

\begin{quote}
when we change the factorization or approximation of attention, which
properties of the original operator survive?
\end{quote}

\section{Attention sensitivity has a concrete
Jacobian}\label{attention-sensitivity-has-a-concrete-jacobian}

For one softmax row

\[
a=\operatorname{softmax}(s),
\]

the Jacobian is

\[
J_{\rm softmax}(a)
=
\operatorname{diag}(a)-aa^\top.
\]

A query perturbation \(\delta q_i\) changes the score row by

\[
\delta s_i
=
\frac{\delta q_iK^\top}{\sqrt{d_k}}.
\]

Therefore, to first order,

\[
\delta a_i
\approx
\left(
\operatorname{diag}(a_i)-a_i a_i^\top
\right)
\frac{\delta q_iK^\top}{\sqrt{d_k}}.
\]

This is already an interface-style sensitivity calculation. A
representation perturbation induces an operator perturbation with a
measurable local Jacobian.

Nothing about this calculation says that a large coefficient is
semantically important. Sensitivity is not explanation.

\section{Three distinctions that must remain
visible}\label{three-distinctions-that-must-remain-visible}

\subsection{Score matrix versus mixing
operator}\label{score-matrix-versus-mixing-operator}

The score matrix can contain arbitrary real values.

The softmax operator is positive and row-normalized. Adding a constant
to one score row changes nothing:

\[
\operatorname{softmax}(s+c\mathbf 1)
=
\operatorname{softmax}(s).
\]

So even the score representation has a redundancy.

\subsection{Mixing operator versus full
head}\label{mixing-operator-versus-full-head}

The operator \(A\) determines how value rows mix. The value projection
decides what information is available to be mixed. A visually striking
\(A\) can have little downstream effect if the value/output pathway
contributes little or is canceled.

\subsection{Operator versus
explanation}\label{operator-versus-explanation}

An observed \(A\) answers:

\begin{quote}
what mixing coefficients were produced?
\end{quote}

It does not automatically answer:

\begin{quote}
which component caused the model's output?
\end{quote}

Causal claims require interventions, substitutions, ablations, or
comparable evidence.

\section{Why this viewpoint matters for
position}\label{why-this-viewpoint-matters-for-position}

Position is often introduced as extra information attached to vectors.

The operator view suggests a broader possibility: relative displacement
may change how representations are transformed.

RoPE already points in that direction because displacement appears
through compositions of rotations. Later the Atlas will ask whether
position can be isolated into a low-dimensional family

\[
R(\Delta)
\]

acting on representation space.

That question is natural once attention is already being described
through operator action and composition.

\section{What the computational witness
establishes}\label{what-the-computational-witness-establishes}

The figure and Wolfram replay establish:

\begin{enumerate}
\def\labelenumi{\arabic{enumi}.}
\tightlist
\item
  a concrete score matrix \(S\);
\item
  a row-softmax operator \(A\);
\item
  a concrete value field \(V\);
\item
  the operator output \(Y=AV\);
\item
  row-stochasticity of \(A\);
\item
  linearity in \(V\) when \(A\) is fixed;
\item
  change in \(A\) under a query perturbation.
\end{enumerate}

They do \textbf{not} establish:

\begin{itemize}
\tightlist
\item
  that an attention map explains a model decision;
\item
  that a high weight is a causal feature importance;
\item
  that the toy operator resembles a trained head;
\item
  that kernel language is the only useful interpretation.
\end{itemize}

The witness makes the decomposition visible. It does not settle
interpretation.

\section{Atlas connections}\label{atlas-connections-11}

\textbf{Approximate and structured attention.}\\
Approximation can now be evaluated in terms of operator error, output
error, complexity, and preserved structure.

\textbf{Position geometry.}\\
Position can modify structured transformations rather than merely
decorate vectors.

\textbf{Relative-position operators.}\\
The operator vocabulary becomes explicit: displacement may be encoded by
a low-dimensional transformation family.

\textbf{Retrieval.}\\
Attention and retrieval both perform state-dependent selection and
mixing, but differ in persistence, indexing, and externality.

\textbf{Mechanistic diagnosis.}\\
Operator substitution and counterfactual replacement are stronger tests
than coefficient inspection.

The recurring Atlas shift is now concrete:

\[
\boxed{
\text{vectors and pairwise scores}
\longrightarrow
\text{state-dependent operators acting on representations}.
}
\]

\section{Closing view}\label{closing-view-17}

The attention matrix is worth looking at.

Its deeper importance is that it \textbf{acts}.

A query-key state constructs a transformation. That transformation mixes
information carried by values. Masks restrict its support. Position
changes its structure. Heads create a family of such transformations.
Approximation changes how the action is represented.

The dynamic switching board is only an allegory.

The operator is the object.

\section{References used in this
chapter}\label{references-used-in-this-chapter-22}

\begin{itemize}
\tightlist
\item
  (Vaswani et al. 2017)
\item
  (Tsai et al. 2019)
\item
  (Katharopoulos et al. 2020)
\end{itemize}

See `sources/source-locks/ATLAS-CH-ATTNOP-001.yaml` for exact source
identities and claim scope.

\chapter{Approximate and Structured
Attention}\label{approximate-and-structured-attention}

\textbf{Epistemic status:} audited attention/Krylov substrate + primary
Performer, Linformer, Nyströmformer, and alpha-entmax sources + Atlas
synthesis.\\
\textbf{Specification:}
manuscript/specifications/ATLAS-CH-ATTNAPPROX-001.md\\
\textbf{Derivation packet:}
mathematics/derivations/ATLAS-CH-ATTNAPPROX-001-DERIVATIONS.md\\
\textbf{Computational witness:}
mathematics/computational-witnesses/ATLAS-CW-ATTNAPPROX-001.md\\
\textbf{Source lock:} sources/source-locks/ATLAS-CH-ATTNAPPROX-001.yaml

Approximate attention is not one problem.

A method may approximate the exponential kernel before normalization.

Another may compress the normalized attention matrix.

Another may avoid constructing that matrix by projecting keys and values
into a lower-dimensional sequence space.

Another may intentionally replace softmax with a sparse probability map.

All four can be called ``efficient attention.''

Mathematically, they are different operations.

This chapter imposes one rule before any speed or accuracy claim:

\begin{quote}
name the object that changed.
\end{quote}

The rule sounds administrative.

It is actually the shortest route to clear mathematics.

\section{The attention stack}\label{the-attention-stack}

For one head, let

\[
Q=XW_Q,
\qquad
K_{\rm key}=XW_K,
\qquad
V=XW_V.
\]

To avoid using K for both ``keys'' and ``kernel,'' write the score
matrix as

\[
S
=
\frac{QK_{\rm key}^\top}{\sqrt{d_k}}.
\]

Now define the positive exponential kernel matrix

\[
G_{ij}
=
e^{S_{ij}}.
\]

Its row sums are

\[
z_i
=
\sum_j G_{ij}.
\]

Let

\[
D_G
=
\operatorname{diag}(z_1,\ldots,z_n).
\]

The normalized attention operator is

\[
A
=
D_G^{-1}G.
\]

Finally,

\[
Y=AV.
\]

The audited Attention-as-an-Operator chapter already separated:

\[
S
\longrightarrow
A
\longrightarrow
V
\longrightarrow
Y.
\]

This chapter inserts one more object:

\[
S
\longrightarrow
G
\longrightarrow
A
\longrightarrow
V
\longrightarrow
Y.
\]

That extra step matters because several efficient-attention methods act
on G, while others act on A, K/V sequence structure, or the
normalization itself.

\section{Five different approximation
statements}\label{five-different-approximation-statements}

Suppose a method produces hatted objects.

The following statements are not interchangeable.

\subsection{Score approximation}\label{score-approximation}

\[
\widehat S\approx S.
\]

This concerns pairwise compatibility before exponentiation.

Small score error can be magnified by exponentiation, especially when
norms are large.

\subsection{Kernel approximation}\label{kernel-approximation}

\[
\widehat G\approx G.
\]

This is the natural target for random-feature approximations of the
exponential dot-product kernel.

\subsection{Operator approximation}\label{operator-approximation}

\[
\widehat A\approx A.
\]

This is the natural target if we care about the actual mixing
coefficients.

\subsection{Output approximation}\label{output-approximation}

\[
\widehat Y
=
\widehat A V
\approx
AV.
\]

This is weaker than operator approximation when tested on only one value
field.

\subsection{Alternative operator}\label{alternative-operator}

\[
\widetilde A
=
T(S),
\]

where T is deliberately not softmax.

Alpha-entmax belongs here.

The method can be valuable without being an estimator of softmax.

The distinction is the central organizing principle of this chapter.

\section{Why raw kernel error can
mislead}\label{why-raw-kernel-error-can-mislead}

Softmax is invariant to adding a constant to a score row.

For any row i,

\[
\operatorname{softmax}(s_i+c_i\mathbf 1)
=
\operatorname{softmax}(s_i).
\]

After exponentiation, this is positive row scaling.

Let C be any positive diagonal matrix.

Then

\[
D(CG)^{-1}CG
=
D(G)^{-1}G.
\]

So replacing G by CG leaves A unchanged.

Take the simplest case:

\[
\widehat G=cG.
\]

Then

\[
\widehat A=A,
\]

but

\[
\|\widehat G-G\|_F
=
|c-1|\|G\|_F.
\]

By taking c large, the raw kernel error becomes arbitrarily large while
the attention operator does not move at all.

This is not a pathology.

It is the algebraic expression of a symmetry already present in softmax.

The practical lesson is:

\begin{quote}
kernel error is meaningful only together with the normalization
convention and the route by which kernel error propagates to normalized
attention.
\end{quote}

\section{Kernel error needs denominator
control}\label{kernel-error-needs-denominator-control}

Take one positive kernel row g.

Let

\[
z=\mathbf 1^\top g.
\]

Suppose

\[
\widehat g=g+e,
\qquad
\widehat z=z+\delta,
\qquad
\delta=\mathbf 1^\top e.
\]

The normalized rows (assuming \(\widehat z>0\)) are

\[
a=\frac g z,
\qquad
\widehat a=\frac{g+e}{\widehat z}.
\]

This algebra does not by itself require every entry of
\(\widehat g=g+e\) to be nonnegative. If any approximate entry is
negative, \(\widehat a\) still sums to one but is not a probability row;
interpreting it as an attention distribution requires separate entrywise
nonnegativity. Positive random-feature constructions supply that
nonnegativity by design.

Direct subtraction gives

\[
\widehat a-a
=
\frac e{\widehat z}
-
\frac{g\delta}{z\widehat z}.
\]

Hence

\[
\|\widehat a-a\|_1
\le
\frac{2\|e\|_1}{\widehat z}.
\]

If

\[
\|e\|_1\le \rho z
\]

with \(0\le\rho<1\), then

\[
\widehat z\ge(1-\rho)z
\]

and therefore

\[
\|\widehat a-a\|_1
\le
\frac{2\rho}{1-\rho}.
\]

The bound is intentionally elementary.

Its purpose is not to replace Performer's analysis.

It isolates the missing ingredient in vague statements such as ``the
kernel is close, so attention is close.''

Normalization must be controlled too.

\section{Positive random features}\label{positive-random-features}

The exponential dot-product kernel has a useful probabilistic
representation.

Let

\[
\omega\sim N(0,I)
\]

and define

\[
\phi_\omega(x)
=
\exp
\left(
\omega^\top x
-
\frac12\|x\|_2^2
\right).
\]

Then

\[
\mathbb E_\omega[
\phi_\omega(q)\phi_\omega(k)
]
=
e^{q^\top k}.
\]

The derivation is one line of Gaussian moment-generating algebra after
expanding

\[
\|q+k\|_2^2.
\]

For scaled dot-product attention, choose

\[
q'
=
\frac{q}{d_k^{1/4}},
\qquad
k'
=
\frac{k}{d_k^{1/4}},
\]

so

\[
q'^\top k'
=
\frac{q^\top k}{\sqrt{d_k}}.
\]

Therefore

\[
e^{q^\top k/\sqrt{d_k}}
\]

inherits the same positive random-feature identity.

This is the clean mathematical reason random features are relevant to
softmax attention.

\section{From an expectation identity to
Performer}\label{from-an-expectation-identity-to-performer}

A finite feature approximation replaces the expectation by an average:

\[
\widehat G_{ij}
=
\frac1m
\sum_{r=1}^m
\phi_{\omega_r}(q_i')
\phi_{\omega_r}(k_j').
\]

With independent Gaussian features, the kernel estimator is unbiased
entrywise.

But ``unbiased'' does not mean ``small error in one run.''

The finite estimator has variance.

The normalized operator is a ratio involving the approximate row sums.

Finite precision can matter because exponentials span large ranges.

And the cost depends explicitly on the feature count m.

Performer addresses this setting with FAVOR+, Fast Attention Via
positive Orthogonal Random features (Choromanski et al. 2021).

The primary source develops positive orthogonal random features and
source-specific accuracy/variance guarantees for softmax-kernel
attention.

The Atlas inherits that result at its stated scope.

It does not convert ``random features'' in general into a guarantee.

\section{Why feature attention can avoid the dense
matrix}\label{why-feature-attention-can-avoid-the-dense-matrix}

Write feature matrices

\[
\Phi_Q\in\mathbb R^{n\times m},
\qquad
\Phi_K\in\mathbb R^{n\times m}.
\]

Then

\[
\widehat G
=
\Phi_Q\Phi_K^\top.
\]

The numerator can be associated as

\[
\widehat G V
=
\Phi_Q
\left(
\Phi_K^\top V
\right).
\]

The denominator is

\[
\widehat G\mathbf 1
=
\Phi_Q
\left(
\Phi_K^\top\mathbf 1
\right).
\]

So the normalized output can be computed without constructing the full
\(n\times n\) matrix.

This is the computational move.

The cost is linear in n only after exposing what is held fixed,
especially m and the feature/value widths.

The word ``linear'' is therefore shorthand for a scaling regime, not a
property independent of approximation budget.

\section{Low-rank attention is another
idea}\label{low-rank-attention-is-another-idea}

Linformer starts from a different structural premise (Wang et al. 2020).

Its motivating claim is that self-attention can often be represented or
approximated through a low-dimensional sequence structure.

Rather than Monte Carlo features for the exponential kernel, Linformer
projects key/value sequence dimensions into a smaller width.

The distinction matters.

A random-feature method asks:

\begin{quote}
can I approximate the kernel action through features?
\end{quote}

A low-rank projection method asks:

\begin{quote}
can I compress the sequence interaction into a lower-dimensional
subspace?
\end{quote}

Both reduce cost.

They do not have the same approximation variable.

For Performer-like methods the explicit budget is feature count m.

For Linformer-like methods it is a projection/rank width k.

A theorem about one budget does not automatically say anything about the
other.

\section{Low rank is a hypothesis, not a universal
property}\label{low-rank-is-a-hypothesis-not-a-universal-property}

Any matrix has a rank.

Saying that an attention operator has a useful low-rank approximation is
stronger.

It says that most of the relevant action can be captured with fewer
directions than the ambient sequence length.

That may occur.

It need not occur uniformly across:

\begin{itemize}
\tightlist
\item
  layers;
\item
  heads;
\item
  inputs;
\item
  training stages;
\item
  positional regimes;
\item
  masks;
\item
  distribution shifts.
\end{itemize}

The correct experimental question is not:

\begin{quote}
is attention low rank?
\end{quote}

It is:

\begin{quote}
for this declared target, norm, data regime, and rank budget, how much
error remains?
\end{quote}

This is the same discipline that classical low-dimensional numerical
methods demand.

The audited Krylov chapter supplies vocabulary for reduced action.

It does not supply a convergence theorem for learned attention.

\section{Landmark reconstruction:
Nyströmformer}\label{landmark-reconstruction-nystruxf6mformer}

Nyströmformer uses a landmark-based route (Xiong et al. 2021).

The broad Nyström idea is to reconstruct a large matrix from
interactions involving a smaller landmark set.

Schematically, one encounters an approximation of the form

\[
\widehat A
\approx
C W^\dagger R,
\]

where W is a reduced landmark block and C/R connect the full matrix to
the landmarks.

Nyströmformer adapts this idea to softmax self-attention.

This introduces a third approximation budget:

\[
m
=
\text{number of landmarks}.
\]

Again, the reduced representation is not the same object as random
features.

A landmark is not a Gaussian feature.

A pseudoinverse-like reduced solve is not a Monte Carlo average.

The shared idea is economy through structure, not mathematical identity.

\section{Conditioning belongs in the landmark
story}\label{conditioning-belongs-in-the-landmark-story}

Reduced systems can be small and still be numerically awkward.

Whenever reconstruction involves an inverse or pseudoinverse-like
operation, conditioning matters.

A reduced matrix that is nearly singular can amplify perturbations.

This is where the Krylov chapter contributes a useful habit:

\begin{quote}
never identify a computable residual or reduced dimension with the true
target error without the operator relation that connects them.
\end{quote}

For attention, the analogous warning is:

\begin{quote}
a small landmark system does not by itself imply a small reconstruction
error.
\end{quote}

Landmark quality, matrix structure, numerical conditioning, and the
chosen norm remain part of the claim.

\section{Sparse attention can mean a different
normalization}\label{sparse-attention-can-mean-a-different-normalization}

Softmax has a rigid property.

For finite scores,

\[
\operatorname{softmax}(s)_i>0
\]

for every coordinate.

The resulting row is dense.

The alpha-entmax family changes this geometry (Peters et al. 2019).

For alpha greater than one, entmax can produce exact zeros.

At alpha equal to two, the family gives sparsemax.

This is not merely ``softmax with a cheap approximation.''

It is a different score-to-simplex map.

That distinction is useful because exact sparsity may be the desired
inductive bias.

It is also necessary because a sparse alternative should not be
penalized as a bad softmax estimator when reproducing softmax was never
its goal.

\section{An exact sparse-support
witness}\label{an-exact-sparse-support-witness}

Take

\[
s=(2,0,-1).
\]

Softmax gives

\[
\frac1{e^2+1+e^{-1}}
(e^2,1,e^{-1}).
\]

All entries are positive.

For sparsemax, use the simplex-projection threshold form

\[
p_i=\max(s_i-\tau,0).
\]

With

\[
\tau=1,
\]

we obtain

\[
p=(1,0,0).
\]

So

\[
\operatorname{entmax}_2(2,0,-1)
=
(1,0,0).
\]

One map mixes all three positions.

The other selects one exactly.

The difference is structural, not a rounding artifact.

\section{Sparsity and low rank are
independent}\label{sparsity-and-low-rank-are-independent}

A sparse matrix can have full rank.

The identity matrix is the simplest example.

A low-rank matrix can be completely dense.

For example,

\[
\frac1n\mathbf 1\mathbf 1^\top
\]

has rank one and no zero entries.

Therefore

\[
\text{sparse}
\not\Rightarrow
\text{low rank},
\]

and

\[
\text{low rank}
\not\Rightarrow
\text{sparse}.
\]

This prevents another common collapse of vocabulary.

``Structured attention'' is a family name.

The structures inside the family remain different.

\section{Exact sparsity is not yet a complexity
theorem}\label{exact-sparsity-is-not-yet-a-complexity-theorem}

Suppose we first form

\[
S=QK_{\rm key}^\top.
\]

That step already considers all token pairs.

If a later normalization turns most coefficients into zero, the dense
pairwise score cost has already been paid.

So

\[
\text{many zeros in }A
\]

does not by itself imply

\[
\text{cheap attention}.
\]

To obtain computational savings, the implementation must exploit
structure earlier or later in a way that changes actual work.

Examples include:

\begin{itemize}
\tightlist
\item
  structured neighborhoods;
\item
  block support;
\item
  hashing/routing;
\item
  feature factorization;
\item
  low-rank projection;
\item
  landmark reconstruction;
\item
  hardware-efficient sparse kernels.
\end{itemize}

The exact mechanism matters.

\section{Three error surfaces}\label{three-error-surfaces}

The Atlas uses three primary numerical surfaces.

\subsection{Kernel error}\label{kernel-error}

\[
E_G
=
\|\widehat G-G\|.
\]

This is natural for random-feature methods.

But it is sensitive to row scale.

\subsection{Operator error}\label{operator-error}

\[
E_A
=
\|\widehat A-A\|.
\]

This measures the mixing transformation itself.

It is the central comparison surface when the target is ordinary softmax
attention.

\subsection{Output error}\label{output-error}

\[
E_Y
=
\|\widehat A V-AV\|.
\]

This measures one realized action on one value field.

For fixed V,

\[
E_Y
\le
\|\widehat A-A\|_2\|V\|_F.
\]

The reverse implication fails from one V.

A large operator discrepancy can lie in directions not excited by the
chosen value field.

\section{A small exact operator
approximation}\label{a-small-exact-operator-approximation}

Consider

\[
A=
\begin{pmatrix}
1/2&1/3&1/6\\
1/2&1/3&1/6\\
1/6&1/3&1/2
\end{pmatrix}.
\]

It is positive and row-stochastic.

Its rank is two.

Now replace all rows by their average:

\[
\widehat A=
\begin{pmatrix}
7/18&1/3&5/18\\
7/18&1/3&5/18\\
7/18&1/3&5/18
\end{pmatrix}.
\]

This approximation is rank one.

The exact error is

\[
A-\widehat A
=
\frac19
\begin{pmatrix}
1&0&-1\\
1&0&-1\\
-2&0&2
\end{pmatrix}.
\]

Therefore

\[
\|A-\widehat A\|_F^2
=
\frac4{27}
\]

and

\[
\|A-\widehat A\|_2
=
\frac{2\sqrt3}{9}.
\]

The witness is not intended to imitate a trained head.

It exists so every quantity can be inspected exactly.

\section{The same operator error can look smaller after acting on
values}\label{the-same-operator-error-can-look-smaller-after-acting-on-values}

Take

\[
V=
\begin{pmatrix}
1&0\\
0&1\\
1&-1
\end{pmatrix}.
\]

Then

\[
AV=
\begin{pmatrix}
2/3&1/6\\
2/3&1/6\\
2/3&-1/6
\end{pmatrix}
\]

and

\[
\widehat A V=
\begin{pmatrix}
2/3&1/18\\
2/3&1/18\\
2/3&1/18
\end{pmatrix}.
\]

The squared output error is

\[
\frac2{27}.
\]

This is smaller than the squared Frobenius operator error,

\[
\frac4{27}.
\]

Nothing paradoxical has happened.

V does not probe every operator direction equally.

This is why output fidelity and operator fidelity must be reported
separately.

\section{Task performance is a fourth
question}\label{task-performance-is-a-fourth-question}

Suppose a model is retrained with an approximate mechanism and reaches
the same benchmark score as a standard Transformer.

That is important empirical evidence.

It still does not prove

\[
\widehat A\approx A
\]

for the corresponding softmax operator.

The trained model may compensate elsewhere.

Its representations may reorganize.

Its heads may specialize differently.

The approximate mechanism may be better for the task while being farther
from softmax.

Thus we should ask four questions separately:

\begin{enumerate}
\def\labelenumi{\arabic{enumi}.}
\tightlist
\item
  Is the kernel close?
\item
  Is the normalized operator close?
\item
  Is the output on the declared V close?
\item
  Does the trained system perform well?
\end{enumerate}

The fourth does not subsume the first three.

\section{Complexity must expose the approximation
budget}\label{complexity-must-expose-the-approximation-budget}

Standard dense attention pays for pairwise interactions across n tokens.

Feature methods replace n-by-n interaction with a feature width m.

Low-rank methods replace sequence dimension with a projection width k.

Landmark methods replace full interaction with m landmarks.

Sparse methods introduce a support size only if the implementation
actually exploits it.

Therefore every scaling statement should look like:

\begin{quote}
cost as a function of n, with m/k/support budget declared.
\end{quote}

The phrase ``linear attention'' is only meaningful after that
declaration.

If maintaining a fixed error tolerance requires m to grow with n, the
effective complexity can change.

\section{Accuracy-cost curves are more informative than single
points}\label{accuracy-cost-curves-are-more-informative-than-single-points}

A fair comparison should sweep the approximation budget.

For random features:

\[
m=16,32,64,\ldots
\]

For low-rank projection:

\[
k=16,32,64,\ldots
\]

For landmarks:

\[
m=16,32,64,\ldots
\]

For sparse alternatives, vary alpha or the mechanism's support-control
parameter when appropriate.

At each budget, measure:

\begin{itemize}
\tightlist
\item
  wall-clock or kernel time;
\item
  memory;
\item
  kernel/operator/output error where defined;
\item
  support/rank diagnostics;
\item
  downstream performance if training is part of the experiment.
\end{itemize}

One point can hide the shape of the tradeoff.

The curve is the object.

\section{Masks must remain in the target
definition}\label{masks-must-remain-in-the-target-definition}

Causal attention is not ordinary unmasked attention followed by a
cosmetic display rule.

The mask constrains support.

If

\[
M_{ij}=-\infty
\]

for forbidden interactions, the exact target operator has zeros there.

An approximation must state whether it:

\begin{itemize}
\tightlist
\item
  preserves those zeros exactly;
\item
  approximates only the allowed region;
\item
  uses a causal recurrence;
\item
  leaks mass into forbidden positions.
\end{itemize}

A method that approximates unmasked attention accurately can still be
invalid for a masked target if support is wrong.

\section{Position changes what is being
approximated}\label{position-changes-what-is-being-approximated}

Suppose scores include a relative-position term:

\[
S_{ij}
=
\frac{q_i^\top k_j}{\sqrt{d_k}}
+
b_{ij}.
\]

Then b is part of the target.

A feature map derived only for the content exponential kernel does not
automatically approximate the combined positional operator.

The same warning applies to rotary structure, head-specific biases, and
other positional mechanisms.

The later positional chapters can inherit the approximation vocabulary
developed here.

They must restate which structure is inside the target.

\section{Finite precision is not an
afterthought}\label{finite-precision-is-not-an-afterthought}

Approximation changes numerical pathways.

Random-feature exponentials can overflow or underflow.

Reduced landmark systems can be ill-conditioned.

Low-rank factorizations can amplify errors through poorly scaled
factors.

Sparse transformations can change gradient behavior near support
transitions.

These are not reasons to reject the methods.

They are reasons to include finite precision in the contract.

The correct question is:

\begin{quote}
what numerical operation is actually performed, and which quantity
remains stable under its implementation?
\end{quote}

\section{A useful comparison table in
words}\label{a-useful-comparison-table-in-words}

The four representative routes can now be summarized without collapsing
them.

\textbf{Performer/FAVOR+.}\\
Target: exponential softmax kernel and normalized attention action.\\
Structure: positive orthogonal random features.\\
Budget: feature count.\\
Characteristic risk: finite-feature variance and normalization error.

\textbf{Linformer.}\\
Target: compressed attention action through sequence projection.\\
Structure: low-dimensional learned projection.\\
Budget: projection/rank width.\\
Characteristic risk: low-rank premise failing on relevant states.

\textbf{Nyströmformer.}\\
Target: normalized self-attention matrix/action.\\
Structure: landmarks and reduced reconstruction.\\
Budget: landmark count.\\
Characteristic risk: landmark adequacy and reduced-system conditioning.

\textbf{Alpha-entmax.}\\
Target: a different score-to-simplex operator.\\
Structure: exact sparse support for alpha greater than one.\\
Budget: not naturally an approximation rank; alpha controls
geometry/sparsity.\\
Characteristic risk: calling a changed operator a softmax approximation
when it is not intended to be one.

\section{The Krylov analogy, used
carefully}\label{the-krylov-analogy-used-carefully}

The Krylov chapter taught a productive lesson:

large operator problems can sometimes be attacked through informative
low-dimensional directions.

That lesson transfers.

The convergence theorems do not.

A Krylov subspace is generated by

\[
r_0,Ar_0,A^2r_0,\ldots
\]

for a declared operator.

Random features are sampled feature directions.

Linformer projections are learned sequence reductions.

Nyström landmarks are selected representatives.

These are all reduced representations.

They are not the same reduced representation.

The safe inheritance is methodological:

\begin{itemize}
\tightlist
\item
  define the target;
\item
  define the reduced object;
\item
  define the residual or error;
\item
  expose conditioning;
\item
  expose the budget;
\item
  test convergence or approximation rather than assuming it.
\end{itemize}

\section{What not to infer}\label{what-not-to-infer}

This chapter rejects several tempting shortcuts.

\subsection{Small kernel error does not automatically give small
operator
error}\label{small-kernel-error-does-not-automatically-give-small-operator-error}

Normalization can amplify errors if denominator control is weak.

\subsection{Small output error on one value field does not give small
operator
error}\label{small-output-error-on-one-value-field-does-not-give-small-operator-error}

The value field may fail to excite discrepant directions.

\subsection{Low rank does not imply
sparsity}\label{low-rank-does-not-imply-sparsity}

Dense rank-one matrices exist.

\subsection{Sparsity does not imply low
rank}\label{sparsity-does-not-imply-low-rank}

The identity matrix is sparse and full rank.

\subsection{Exact zeros do not imply fast
execution}\label{exact-zeros-do-not-imply-fast-execution}

Dense score formation may already have occurred.

\subsection{Few random features do not imply uniform
accuracy}\label{few-random-features-do-not-imply-uniform-accuracy}

Finite approximation quality is a statistical question.

\subsection{Good benchmark performance does not imply softmax
fidelity}\label{good-benchmark-performance-does-not-imply-softmax-fidelity}

Retraining can change the system.

These are not edge cases.

They are the basic claim boundaries.

\section{A practical evaluation
protocol}\label{a-practical-evaluation-protocol}

For a fixed trained head or controlled synthetic case:

\begin{enumerate}
\def\labelenumi{\arabic{enumi}.}
\tightlist
\item
  freeze Q, keys, V, mask, and position;
\item
  compute exact dense softmax A and Y when feasible;
\item
  choose one approximation family;
\item
  state its budget m or k;
\item
  compute its own target objects;
\item
  measure operator error to A if softmax fidelity is intended;
\item
  measure output error on V;
\item
  record row sums, positivity, support violations, and rank/sparsity;
\item
  record memory and runtime;
\item
  sweep the budget;
\item
  repeat over representative inputs;
\item
  separate approximation fidelity from retrained task performance.
\end{enumerate}

This turns ``efficient attention'' from a label into an experiment.

\section{Research bridges}\label{research-bridges}

Several later Atlas questions become sharper after this chapter.

\subsection{Spectral diagnostics}\label{spectral-diagnostics}

If A is approximated by \(\widehat A\), how do singular values,
non-normality, mixing rates, or transient amplification change?

\subsection{Context compression}\label{context-compression}

Can a low-rank or landmark structure be used deliberately as a
context-compression operator rather than merely as a cheaper attention
surrogate?

\subsection{Relative-position
operators}\label{relative-position-operators}

Can positional structure be factored exactly or approximately without
destroying support or equivariance?

\subsection{Systems}\label{systems}

When does theoretical sparsity or low rank translate into actual
accelerator speed?

\subsection{Adaptive approximation}\label{adaptive-approximation}

Can feature count, landmark count, or rank be selected dynamically using
an error indicator rather than fixed globally?

That last question connects directly to the Atlas treatment of adaptive
error control.

The error indicator would need an explicit relation to the operator or
output error being controlled.

\section{What the computational witness
establishes}\label{what-the-computational-witness-establishes-1}

The finite witness establishes exactly:

\begin{enumerate}
\def\labelenumi{\arabic{enumi}.}
\tightlist
\item
  a positive row-stochastic rank-2 operator;
\item
  a positive row-stochastic rank-1 approximation;
\item
  exact Frobenius operator error \(4/27\) in squared norm;
\item
  exact spectral operator error \(2\sqrt3/9\);
\item
  exact output error \(2/27\) in squared Frobenius norm for a declared
  V;
\item
  the standard operator-to-output inequality;
\item
  a score row for which softmax is dense;
\item
  the alpha=2 entmax/sparsemax result \((1,0,0)\);
\item
  row-scaling invariance showing raw kernel error and operator error can
  completely decouple.
\end{enumerate}

It does not establish trained-model superiority or a finite-feature
Performer guarantee.

\section{Closing view}\label{closing-view-18}

Approximate attention is best understood as a map of design choices.

One route approximates a kernel.

One route compresses an operator.

One route reconstructs from landmarks.

One route changes the normalization geometry and creates exact sparsity.

The useful unifying picture is not:

\[
\text{all efficient attention is the same}.
\]

It is:

\[
\boxed{
\text{target}
+
\text{structure}
+
\text{budget}
+
\text{error notion}
+
\text{implementation}
}.
\]

That tuple tells us what was preserved and what was changed.

Without it, ``linear,'' ``sparse,'' ``low-rank,'' and ``approximate''
are only adjectives.

With it, they become mathematical claims.

\section{References used in this
chapter}\label{references-used-in-this-chapter-23}

\begin{itemize}
\tightlist
\item
  (Choromanski et al. 2021)
\item
  (Wang et al. 2020)
\item
  (Xiong et al. 2021)
\item
  (Peters et al. 2019)
\end{itemize}

Exact source identities and claim boundaries are recorded in
sources/source-locks/ATLAS-CH-ATTNAPPROX-001.yaml.

\chapter{The Geometry of Position}\label{the-geometry-of-position}

\textbf{Epistemic status:} mathematical exposition built from audited
Attention and Geometry prerequisites, primary positional-encoding
sources, and Atlas-owned exact witnesses.\\
\textbf{Derivation packet:}
`mathematics/derivations/ATLAS-CH-POSGEOM-001-DERIVATIONS.md`\\
\textbf{Computational witness:}
`mathematics/computational-witnesses/ATLAS-CW-POSGEOM-001.md`

\section{Position is not content}\label{position-is-not-content}

A Transformer sees token representations.

Without an asymmetric mask or positional signal, ordinary self-attention
is permutation equivariant: permuting the token order permutes the
output order.

Sequence order therefore has to enter somewhere.

The first conceptual separation is:

\[
\boxed{
\text{content geometry}
\neq
\text{position geometry}.
}
\]

Content determines what a token represents.

Position determines where that token sits in a declared coordinate
system.

The two interact inside attention, but they are not the same object.

This chapter develops position as an explicit transformation acting on
attention representations.

\section{Additive sinusoidal
position}\label{additive-sinusoidal-position}

The original Transformer adds a fixed sinusoidal vector to the token
embedding (Vaswani et al. 2017).

For one angular frequency \(\omega\), consider the two-dimensional
component

\[
p_m(\omega)
=
\begin{pmatrix}
\sin(m\omega)\\
\cos(m\omega)
\end{pmatrix}.
\]

A fixed offset \(\delta\) acts linearly:

\[
p_{m+\delta}
=
A_\delta p_m,
\]

where \(A_\delta\) is a rotation matrix determined by \(\delta\omega\).

That fact explains the original motivation: relative displacement can be
represented through a position-independent linear transform of the
sinusoidal features.

But the mechanism is still additive.

The positional vector is added to the representation before later layers
act.

\section{Rotary position moves the geometry into query and key
space}\label{rotary-position-moves-the-geometry-into-query-and-key-space}

Rotary Position Embedding takes a different route (Su et al. 2021).

Instead of adding a positional vector, it rotates query and key
coordinates as a function of position.

For one 2D block define

\[
R(\phi)
=
\begin{pmatrix}
\cos\phi&-\sin\phi\\
\sin\phi&\cos\phi
\end{pmatrix}.
\]

At position \(m\), use

\[
R_m(\omega)=R(m\omega).
\]

Apply that rotation to the corresponding query and key block.

The crucial algebra is

\[
R(\phi)^\top
=
R(-\phi),
\]

so

\[
R_m(\omega)^\top R_n(\omega)
=
R((n-m)\omega).
\]

Therefore

\[
(R_m q)^\top(R_n k)
=
q^\top R((n-m)\omega)k.
\]

The absolute positions \(m\) and \(n\) disappear from this pairwise
expression except through their difference.

This is the core relative-position identity.

\section{What the identity actually
says}\label{what-the-identity-actually-says}

The identity is exact.

It says that the contribution of one rotary block to the query-key inner
product is determined by:

\begin{itemize}
\tightlist
\item
  the unrotated query block;
\item
  the unrotated key block;
\item
  the frequency;
\item
  the relative displacement \(n-m\).
\end{itemize}

It does not say that the whole model depends only on relative position.

Other components can break that symmetry:

\begin{itemize}
\tightlist
\item
  causal masks;
\item
  finite sequence boundaries;
\item
  content distributions;
\item
  layerwise nonlinearities;
\item
  caches;
\item
  architecture-specific operations;
\item
  any additional position-dependent mechanism.
\end{itemize}

The exact statement belongs to the transformed query-key score
contribution.

\section{An exact same-offset
witness}\label{an-exact-same-offset-witness}

Choose

\[
\theta=\pi/6,
\qquad
q=k=
\begin{pmatrix}
1\\
0
\end{pmatrix}.
\]

Compare positions \((0,2)\).

The relative offset is \(2\), so

\[
(R_0q)^\top(R_2k)
=
\cos(2\theta)
=
\frac12.
\]

Now compare positions \((3,5)\).

The relative offset is still \(2\), and

\[
(R_3q)^\top(R_5k)
=
\frac12.
\]

The absolute coordinates moved.

The rotary score contribution did not.

This is a finite replay of the relative-position identity.

\section{Why the rotation structure
matters}\label{why-the-rotation-structure-matters}

The result is not true for arbitrary position-dependent transforms.

Define instead

\[
S_m
=
\operatorname{diag}(2^m,1).
\]

Using the same \(q=k=(1,0)^\top\),

\[
(S_m q)^\top(S_n k)
=
2^{m+n}.
\]

For the same-offset pair \((0,2)\),

\[
\text{score}=4.
\]

For the same-offset pair \((3,5)\),

\[
\text{score}=256.
\]

The relative displacement is identical.

The score is not.

What changed?

The family \(S_m\) does not satisfy the rotary group relation

\[
T_m^\top T_n=T_{n-m}.
\]

The exact relative-position identity depends on the structure of the
transform family.

\section{Multi-frequency position}\label{multi-frequency-position}

Real rotary encodings use many 2D blocks.

For even dimension \(d=2h\),

\[
R_m
=
\operatorname{diag}
\left(
R(m\omega_1),
\ldots,
R(m\omega_h)
\right).
\]

Then

\[
R_m^\top R_n
=
\operatorname{diag}
\left(
R((n-m)\omega_1),
\ldots,
R((n-m)\omega_h)
\right).
\]

Each block sees the same displacement through a different phase rate.

The frequency schedule therefore determines how relative displacement is
represented across coordinates.

Low frequencies rotate slowly.

High frequencies rotate quickly.

The resulting positional representation is a bank of phase measurements
at multiple scales.

\section{Position as phase}\label{position-as-phase}

The phase viewpoint is more useful than thinking of RoPE as ``just
another embedding.''

For one frequency,

\[
m
\longmapsto
m\omega
\longmapsto
R(m\omega).
\]

Position becomes a group parameter.

The representation at position \(m\) is related to position \(n\) by the
phase difference

\[
(n-m)\omega.
\]

This makes several properties transparent:

\begin{itemize}
\tightlist
\item
  common translation cancels in pairwise phase difference;
\item
  different frequencies resolve displacement differently;
\item
  rescaling position changes phase;
\item
  rescaling frequency also changes phase;
\item
  those two interventions can look algebraically related while being
  operationally distinct.
\end{itemize}

\section{Common translation}\label{common-translation}

Shift every position by the same amount \(c\).

Then

\[
(n+c)-(m+c)=n-m.
\]

So

\[
R_{m+c}^\top R_{n+c}
=
R_{n-m}.
\]

The rotary query-key contribution is invariant under a common shift of
both position indices.

Again, this is not a theorem that the full sequence model is translation
invariant.

It is an exact property of the rotary score geometry.

\section{Multidimensional position}\label{multidimensional-position}

A sequence uses one coordinate.

Images, grids, fields, and other structured domains may need several.

Let a 2D position be

\[
p=(u,v).
\]

Use two independent rotary blocks:

\[
R_p
=
\operatorname{diag}
\left(
R(u\omega_x),
R(v\omega_y)
\right).
\]

For another position

\[
q=(u',v'),
\]

we obtain

\[
R_p^\top R_q
=
\operatorname{diag}
\left(
R((u'-u)\omega_x),
R((v'-v)\omega_y)
\right).
\]

The pairwise operator depends on the coordinate-wise displacement.

This is the direct-product version of the one-dimensional construction.

The important point is not that every 2D positional method must use this
formula.

It is that multidimensional position needs an explicit coordinate
action.

Duplicating a scalar position label does not automatically create
geometry.

\section{Position and attention support are
separate}\label{position-and-attention-support-are-separate}

A causal mask and positional encoding do different jobs.

The mask decides whether an attention edge is allowed.

The positional transformation affects the score attached to an allowed
query-key pair.

A pair may have a well-defined relative phase and still be masked out.

This distinction keeps support geometry separate from score geometry.

\section{Exact algebra does not guarantee long-context
generalization}\label{exact-algebra-does-not-guarantee-long-context-generalization}

The identity

\[
R_m^\top R_n=R_{n-m}
\]

continues to make algebraic sense far beyond any finite training
context.

That does not mean the trained model knows how to use those phases.

A model is trained under a finite distribution of:

\begin{itemize}
\tightlist
\item
  absolute positions;
\item
  relative displacements;
\item
  phase combinations;
\item
  sequence lengths;
\item
  attention patterns;
\item
  optimization conditions.
\end{itemize}

Outside that regime, exact continuation of the positional formula is
only one ingredient.

The model-level question is empirical.

\section{Direct extrapolation}\label{direct-extrapolation}

The simplest long-context strategy is to evaluate the same positional
rule at larger indices.

For rotary position, that means evaluating phases

\[
m\omega
\]

for \(m\) beyond the range used in training.

The formula is valid.

The behavior need not be.

This is the distinction between:

\[
\boxed{
\text{formula extrapolates}
}
\]

and

\[
\boxed{
\text{trained model generalizes}.
}
\]

They are not equivalent.

\section{Position Interpolation}\label{position-interpolation}

Position Interpolation changes the coordinate presented to the
positional mechanism (Chen et al. 2023).

Suppose the original context scale is \(L\) and the desired scale is
\(L'>L\).

A simple interpolation map is

\[
m
\longmapsto
\tilde m
=
m\frac{L}{L'}.
\]

The phase becomes

\[
\tilde m\omega
\]

instead of \(m\omega\).

The paper reports that this can extend RoPE-based pretrained models with
relatively limited fine-tuning and gives a theoretical comparison
between interpolation and direct extrapolation under its stated
analysis.

The Atlas uses that result at its actual scope.

It is not a theorem that every model, every frequency schedule, or every
target length will behave well under the same rescaling.

\section{YaRN}\label{yarn}

YaRN belongs to the same broader problem family but is not merely
``Position Interpolation again'' (Peng et al. 2023).

It combines positional rescaling ideas with additional treatment
intended to improve efficient long-context extension.

The reported results are empirical method evidence.

The relevant Atlas lesson is structural:

\begin{quote}
Long-context extension is an intervention on the positional phase regime
and often also on training.
\end{quote}

The exact rotary identity remains true before and after such
modifications.

What changes is which phases the model encounters and how well its
learned computation handles them.

\section{Frequency scaling and index scaling are different
interventions}\label{frequency-scaling-and-index-scaling-are-different-interventions}

Consider a phase

\[
m\omega.
\]

One can change it by:

\begin{itemize}
\tightlist
\item
  replacing \(m\) with \(\alpha m\);
\item
  replacing \(\omega\) with \(\alpha\omega\);
\item
  changing both nonuniformly across frequencies;
\item
  fine-tuning after the change.
\end{itemize}

These may coincide algebraically in a single isolated block under
special choices.

They are not operationally interchangeable in a full model.

A positional extension method should therefore state exactly what it
changes.

\section{The geometry of
extrapolation}\label{the-geometry-of-extrapolation}

Long-context behavior is often described as if context length were one
scalar knob.

Geometrically, several things can move at once:

\begin{itemize}
\tightlist
\item
  the range of phase angles;
\item
  the density of sampled offsets;
\item
  aliasing relationships across frequencies;
\item
  the distribution of query-key content seen at those phases;
\item
  mask/boundary effects;
\item
  the learned response to phase combinations.
\end{itemize}

This is why the positional mechanism can remain mathematically
well-defined while the trained system degrades.

The geometry survives.

The learned use of the geometry may not.

\section{Failure modes}\label{failure-modes-3}

\subsection{Additive equals rotary}\label{additive-equals-rotary}

It does not.

Both use sinusoidal structure, but one adds position features while the
other rotates query/key coordinates.

\subsection{Relative identity equals long-context
guarantee}\label{relative-identity-equals-long-context-guarantee}

It does not.

The identity is algebraic.

Generalization is model-level.

\subsection{Any position transform
works}\label{any-position-transform-works}

It does not.

The non-orthogonal control shows that same-offset dependence can fail.

\subsection{Frequency schedule equals context
policy}\label{frequency-schedule-equals-context-policy}

It does not.

The schedule defines phase rates.

A context-extension policy says how positions/frequencies/training are
changed when the target range changes.

\subsection{2D means two copies of
1D}\label{d-means-two-copies-of-1d}

Not automatically.

The coordinate action has to be declared.

\subsection{Interpolation equals
extrapolation}\label{interpolation-equals-extrapolation}

They are opposite geometric choices.

Interpolation maps a larger target range back into a smaller phase
regime.

Direct extrapolation evaluates outside that regime.

\subsection{Position explains all long-context
behavior}\label{position-explains-all-long-context-behavior}

It does not.

Attention content, optimization, retrieval, memory, architecture, and
training distribution all matter.

\section{Downstream handoff}\label{downstream-handoff-2}

Relative-Position Operators may assume:

\begin{itemize}
\tightlist
\item
  the exact block-rotation identity;
\item
  the multi-frequency decomposition;
\item
  common-translation invariance of the rotary score contribution;
\item
  multidimensional direct-product position;
\item
  the distinction between exact positional algebra and empirical context
  extension.
\end{itemize}

It may not assume that vector RoPE already provides:

\begin{itemize}
\tightlist
\item
  a low-dimensional operator decomposition;
\item
  a DC-mode theory;
\item
  frequency-mode specialization;
\item
  head-specific relative-position operators.
\end{itemize}

Those are the next chapter's obligations.

\section{References used in this
chapter}\label{references-used-in-this-chapter-24}

\begin{itemize}
\tightlist
\item
  (Vaswani et al. 2017) --- fixed additive sinusoidal positional
  encoding.
\item
  (Su et al. 2021) --- rotary position embedding and relative-position
  score construction.
\item
  (Chen et al. 2023) --- Position Interpolation for RoPE-based context
  extension.
\item
  (Peng et al. 2023) --- YaRN context-extension method.
\end{itemize}

Exact provenance and claim boundaries are locked in:

`sources/source-locks/ATLAS-CH-POSGEOM-001.yaml`

\chapter{Relative-Position
Operators}\label{relative-position-operators-1}

\textbf{Epistemic status:} audited Positional Geometry/Linear Algebra
prerequisites + primary relative-position sources + Atlas synthesis +
exact finite witness\\
\textbf{Specification:} manuscript/specifications/ATLAS-CH-RPO-001.md\\
\textbf{Derivation packet:}
mathematics/derivations/ATLAS-CH-RPO-001-DERIVATIONS.md\\
\textbf{Computational witness:}
mathematics/computational-witnesses/ATLAS-CW-RPO-001.md

Position can enter a Transformer as a vector.

It can also enter as an operator.

That shift matters because an operator can be decomposed into modes,
projected, truncated, compared across heads, and composed with other
structure. The chapter develops that viewpoint without pretending that
every positional mechanism is the same object.

\section{Two spaces, two frequency
notions}\label{two-spaces-two-frequency-notions}

The audited Positional Geometry chapter already gives a precise RoPE
identity.

For a relative offset \(\Delta\), one can write a feature-space operator

\[
\rho(\Delta)
=
\bigoplus_a R(\omega_a\Delta),
\]

where each \(R(\omega_a\Delta)\) is a 2D rotation block.

Those frequencies live in feature coordinates.

This chapter also studies translation-invariant operators over sequence
index.

Those operators have Fourier modes over positions.

The two frequency notions are not automatically the same.

That distinction is the first discipline of RPO.

\section{Relative position as an operator over sequence
index}\label{relative-position-as-an-operator-over-sequence-index}

Suppose positions lie on the cyclic set

\[
\mathbb Z_L.
\]

For one attention head, define a relative-position kernel

\[
\kappa_h(\Delta).
\]

Then let

\[
(K_hx)_i
=
\sum_{\Delta=0}^{L-1}
\kappa_h(\Delta)
x_{i+\Delta\;\mathrm{mod}\;L}.
\]

This is a circulant linear operator.

The periodic boundary is deliberate.

It makes the algebra exact.

It is not a claim that real language sequences wrap around from the end
to the beginning.

For ordinary finite nonperiodic sequences, a relative-position matrix is
generally Toeplitz-like rather than circulant.

\section{Why the Fourier basis
appears}\label{why-the-fourier-basis-appears}

A circulant operator commutes with cyclic shifts.

The natural basis is therefore the Fourier basis

\[
\phi_k(n)
=
L^{-1/2}
e^{2\pi i kn/L}.
\]

Each mode is an eigenvector:

\[
K_h\phi_k
=
\lambda_{h,k}\phi_k,
\]

with

\[
\lambda_{h,k}
=
\sum_{\Delta}
\kappa_h(\Delta)
e^{2\pi i k\Delta/L}.
\]

So the operator becomes

\[
K_h
=
\sum_k
\lambda_{h,k}
\phi_k\phi_k^\ast.
\]

This is the core low-dimensional operator view.

A relative-position rule becomes a collection of labeled mode gains.

\section{The DC mode}\label{the-dc-mode}

The zero-frequency mode is

\[
\phi_0
=
L^{-1/2}\mathbf1.
\]

Its eigenvalue is

\[
\lambda_{h,0}
=
\sum_\Delta \kappa_h(\Delta).
\]

The corresponding operator is

\[
K_h^{DC}
=
\frac{\lambda_{h,0}}{L}
\mathbf1\mathbf1^\top.
\]

It acts on the constant sequence-index mode.

This is what DC means here.

It is not an absolute-position embedding.

It contains no distinguished position label.

\section{RoPE can have a feature-space DC block
too}\label{rope-can-have-a-feature-space-dc-block-too}

There is a parallel feature-space idea.

If a RoPE-style block uses

\[
\omega=0,
\]

then

\[
R(\omega\Delta)=I
\]

for every offset.

That block is offset-invariant.

It is natural to call it a zero-frequency feature-space block.

But it is still not the same object as the sequence-index Fourier DC
mode.

One lives in feature coordinates.

The other lives over positions.

\section{Low-dimensional operator
reduction}\label{low-dimensional-operator-reduction}

Suppose only a subset of Fourier modes is retained:

\[
S\subseteq\{0,\ldots,L-1\}.
\]

Define

\[
K_{h,S}
=
\sum_{k\in S}
\lambda_{h,k}
\phi_k\phi_k^\ast.
\]

For a circulant operator,

\[
\|K_h-K_{h,S}\|_F^2
=
\sum_{k\notin S}
|\lambda_{h,k}|^2,
\]

and, when at least one mode is omitted,

\[
\|K_h-K_{h,S}\|_2
=
\max_{k\notin S}
|\lambda_{h,k}|.
\]

If \(S=\{0,\ldots,L-1\}\), no mode is omitted and \(K_{h,S}=K_h\); both
errors are exactly zero. This avoids interpreting a maximum over the
empty set as a numerical value.

These are exact approximation statements.

They do not imply that the truncated operator preserves downstream
language-model quality.

An operator norm is not a task metric.

\section{Relative position in actual attention
architectures}\label{relative-position-in-actual-attention-architectures}

Shaw, Uszkoreit, and Vaswani introduced relative-position
representations directly into self-attention rather than relying only on
absolute input-position representations (Shaw et al. 2018).

T5 uses a simpler relative-position mechanism: a learned scalar
positional bias is added to attention logits, and different heads within
a layer use different learned relative-position embeddings (Raffel et
al. 2020).

These architectures motivate an operator question:

\begin{quote}
what positional operator does each head implement before
content-dependent attention is applied?
\end{quote}

The Atlas answers that question mathematically without claiming the
papers themselves performed the Fourier decomposition developed here.

\section{Relative bias is not full
attention}\label{relative-bias-is-not-full-attention}

Suppose a head has relative-position logit bias

\[
B_h(i,j)
=
b_h(j-i).
\]

That positional component can be studied as an operator.

But the full attention logits are

\[
L_h
=
Q_hK_h^\top/\sqrt d+B_h.
\]

The attention weights are

\[
A_h
=
\operatorname{softmax}(L_h).
\]

The content term changes from input to input.

Softmax is nonlinear.

Therefore diagonalizing the positional bias does not diagonalize the
full attention mechanism in general.

This chapter diagnoses one component of attention, not the entire head.

\section{Head-specific positional mode
profiles}\label{head-specific-positional-mode-profiles}

For a head \(h\) whose declared relative-position operator is nonzero,
so that \(\sum_j|\lambda_{h,j}|^2>0\), define

\[
q_{h,k}
=
\frac{|\lambda_{h,k}|^2}
{\sum_j|\lambda_{h,j}|^2}.
\]

This is a probability distribution over labeled sequence-index modes.
When \(K_h=0\), the denominator vanishes and the normalized profile is
undefined; report a zero-energy operator instead of assigning it a
fictitious frequency preference.

It answers a narrow question:

\begin{quote}
where does this head's relative-position operator place its spectral
energy?
\end{quote}

That is a useful notion of positional specialization.

It is not yet semantic specialization.

It does not tell us whether the head is causally necessary for a task.

\section{Exact two-head witness}\label{exact-two-head-witness-1}

Take four cyclic positions.

Head A has kernel

\[
\kappa_A
=
\left(
1,\frac12,0,\frac12
\right).
\]

Its mode eigenvalues are

\[
\lambda_A=(2,1,0,1).
\]

Head B has

\[
\kappa_B
=
\left(
1,-\frac12,0,-\frac12
\right),
\]

with

\[
\lambda_B=(0,1,2,1).
\]

Both heads have the same unordered singular values:

\[
(2,1,1,0).
\]

If we looked only at the singular-value multiset, they would appear
identical.

But the mode labels reveal a sharp difference.

\section{One head is DC-dominant}\label{one-head-is-dc-dominant}

Head A has energy profile

\[
q_A
=
\left(
\frac23,\frac16,0,\frac16
\right).
\]

Two-thirds of its spectral energy lies in the DC mode.

Its DC component is

\[
K_A^{DC}
=
\frac12
\mathbf1\mathbf1^\top.
\]

A DC-only approximation has Frobenius error

\[
\sqrt2.
\]

\section{The other head is
Nyquist-dominant}\label{the-other-head-is-nyquist-dominant}

Head B has

\[
q_B
=
\left(
0,\frac16,\frac23,\frac16
\right).
\]

Its DC eigenvalue is zero.

Two-thirds of its spectral energy lies in the \(k=2\) mode, the Nyquist
mode on four cyclic positions.

A DC-only approximation is therefore the zero operator and has Frobenius
error

\[
\sqrt6.
\]

The same approximation strategy can be reasonable for one head and
useless for another.

\section{Singular values lose mode
identity}\label{singular-values-lose-mode-identity}

The witness gives the chapter's main finite warning:

\[
\boxed{
\text{same singular values}
\not\Rightarrow
\text{same positional-frequency specialization}
}
\]

Singular values retain magnitude.

They do not retain which Fourier mode carries that magnitude.

For positional operators, the label can be the interesting part.

\section{Head specialization needs a
qualifier}\label{head-specialization-needs-a-qualifier}

T5 gives an architectural precedent for different relative-position
parameters across heads.

That does not mean every head develops a meaningful positional role.

A head-specific profile can differ because of:

\begin{itemize}
\tightlist
\item
  initialization;
\item
  optimization;
\item
  architecture;
\item
  redundancy;
\item
  incidental parameterization;
\item
  genuinely useful positional structure.
\end{itemize}

Therefore this chapter uses the phrase \textbf{positional-operator
specialization}.

Semantic or causal specialization requires additional evidence.

\section{The DC component is not always
desirable}\label{the-dc-component-is-not-always-desirable}

A large DC component means strong response to the constant
sequence-index mode.

That may be useful.

It may also be an uninformative offset-like component.

The meaning depends on how the positional operator enters the model.

RPO therefore reports DC energy before interpreting it.

\section{Boundary conditions
matter}\label{boundary-conditions-matter}

The Fourier diagonalization above is exact because the toy operator is
circulant.

Real relative-position schemes may use:

\begin{itemize}
\tightlist
\item
  finite clipping;
\item
  logarithmic buckets;
\item
  causal masks;
\item
  nonperiodic boundaries;
\item
  asymmetric left/right treatment;
\item
  context-length-dependent rules.
\end{itemize}

Those operations break or modify simple circulant structure.

The right generalization may involve Toeplitz analysis, block structure,
finite-section operators, or numerical spectral diagnostics rather than
a single exact DFT.

\section{Bucketing is an operator approximation
too}\label{bucketing-is-an-operator-approximation-too}

T5 maps ranges of relative offsets into shared learned buckets.

That can be viewed as reducing the degrees of freedom of the positional
operator.

But bucket reduction is not the same as Fourier truncation.

One constrains the kernel in offset space.

The other constrains it in mode space.

The two approximations preserve different structures.

\section{Frequency truncation can be
head-specific}\label{frequency-truncation-can-be-head-specific}

Suppose a system wants a low-dimensional positional operator.

A global rule such as ``keep the DC mode first'' can fail badly when
different heads emphasize different modes.

The exact witness makes this concrete:

\begin{itemize}
\tightlist
\item
  head A: DC-only error \(\sqrt2\);
\item
  head B: DC-only error \(\sqrt6\).
\end{itemize}

A head-specific reduction can therefore be mathematically better even
before any task-specific evaluation is considered.

\section{Multi-frequency RoPE and sequence modes should not be
collapsed}\label{multi-frequency-rope-and-sequence-modes-should-not-be-collapsed}

RoPE already uses multiple frequencies.

RPO also uses multiple sequence-index Fourier modes.

The notation can look temptingly similar.

But the two decompositions diagonalize different actions.

RoPE blocks encode relative phase in feature coordinates.

Sequence Fourier modes diagonalize a translation-invariant operator
acting across positions.

A model can contain both structures at once.

\section{What a positional-operator diagnostic should
record}\label{what-a-positional-operator-diagnostic-should-record}

At minimum:

\begin{itemize}
\tightlist
\item
  the space on which the operator acts;
\item
  the sequence boundary convention;
\item
  the relative-position kernel or bias rule;
\item
  clipping or bucketing;
\item
  whether the operator is before or after content interaction;
\item
  the head/layer identity;
\item
  the basis used for decomposition;
\item
  retained mode labels, not only sorted magnitudes;
\item
  the norm used for approximation error;
\item
  any task metric used to interpret the approximation.
\end{itemize}

Without these fields, ``low-frequency positional behavior'' can refer to
several different mathematical objects.

\section{Handoff to Latent Clocks}\label{handoff-to-latent-clocks}

Observed sequence position and latent time are not the same thing.

But relative-position operators supply useful structure for asking
whether a model treats some offsets as equivalent, periodic, slowly
varying, or high-frequency.

ATLAS-CH-LATENTTIME-001 may inherit:

\begin{itemize}
\tightlist
\item
  feature-space relative-offset operators;
\item
  sequence-index relative-position kernels;
\item
  Fourier-mode decomposition;
\item
  DC components;
\item
  exact low-dimensional truncation errors;
\item
  head-specific positional-mode profiles;
\item
  the finite example separating singular values from labeled
  frequencies.
\end{itemize}

LATENTTIME must add the inference problem:

\begin{quote}
when do observed sequence offsets correspond to an underlying latent
temporal coordinate, and how should asynchronous or warped sequences be
aligned?
\end{quote}

\section{Durable lesson}\label{durable-lesson}

A positional mechanism is easier to reason about once we ask:

\[
\text{what operator acts on what space?}
\]

Then frequency, DC structure, truncation, and head-specific variation
become typed mathematical questions instead of metaphors.

The operator view does not replace empirical evaluation.

It makes the positional object explicit enough to evaluate.

\section{References used in this
chapter}\label{references-used-in-this-chapter-25}

\begin{itemize}
\tightlist
\item
  (Shaw et al. 2018)
\item
  (Raffel et al. 2020)
\end{itemize}

Exact source authority, prerequisite identities, and claim boundaries
are locked in:

sources/source-locks/ATLAS-CH-RPO-001.yaml

\chapter{Latent Clocks}\label{latent-clocks}

\textbf{Epistemic status:} audited Relative-Position Operators and
Dynamics substrates + primary time-warping sources + Atlas-owned exact
finite alignment witness.\\
\textbf{Specification:}
manuscript/specifications/ATLAS-CH-LATENTTIME-001.md\\
\textbf{Derivation packet:}
mathematics/derivations/ATLAS-CH-LATENTTIME-001-DERIVATIONS.md\\
\textbf{Computational witness:}
mathematics/computational-witnesses/ATLAS-CW-LATENTTIME-001.md\\
\textbf{Source lock:} sources/source-locks/ATLAS-CH-LATENTTIME-001.yaml

Sequence position is observed.

Time can be inferred.

Those statements are not the same.

A token, frame, sensor sample, or event has an index in an observed
stream. But two streams can progress through the same underlying phase
at different rates, with pauses, duplicates, missing samples, or
modality-specific distortions.

This chapter asks:

\begin{quote}
when should observed offsets be replaced by an inferred alignment
coordinate?
\end{quote}

The governing distinction is:

\[
\boxed{
\text{observed sequence index}
\neq
\text{latent temporal alignment}.
}
\]

And the governing caution is:

\[
\boxed{
\text{optimal alignment}
\neq
\text{proof of a uniquely true physical clock}.
}
\]

\section{The handoff from positional
operators}\label{the-handoff-from-positional-operators}

Relative-Position Operators already separates:

\begin{itemize}
\tightlist
\item
  feature-space RoPE rotations;
\item
  sequence-index relative-position kernels;
\item
  Fourier modes over positional operators;
\item
  DC structure;
\item
  low-dimensional positional approximations.
\end{itemize}

Those are positional objects.

They can tell us that a model treats certain offsets as:

\begin{itemize}
\tightlist
\item
  similar;
\item
  periodic;
\item
  slowly varying;
\item
  high-frequency.
\end{itemize}

They do not by themselves tell us that two asynchronous streams are at
the same latent phase.

LATENTTIME adds that inference problem.

\section{The handoff from Dynamics}\label{the-handoff-from-dynamics}

Dynamics already separates:

\begin{itemize}
\tightlist
\item
  state;
\item
  trajectory;
\item
  vector field;
\item
  flow;
\item
  discrete update.
\end{itemize}

That distinction matters because time-warping produces a discrete
correspondence path.

A correspondence path is not automatically a trajectory of a continuous
dynamical system.

The chapter therefore keeps:

\[
\text{alignment}
\]

separate from:

\[
\text{law of motion}.
\]

\section{Two observed streams}\label{two-observed-streams}

Let:

\[
X=(x_1,\dots,x_n),
\]

\[
Y=(y_1,\dots,y_m).
\]

The indices:

\[
i=1,\dots,n,
\qquad
j=1,\dots,m
\]

refer to positions in different observed sequences.

There is no reason in general to assume:

\[
i=j
\]

means ``same latent time.''

\section{Dynamic time warping}\label{dynamic-time-warping}

Dynamic time warping searches over monotone correspondences between two
ordered streams.

For the finite Atlas formalism, an admissible path is:

\[
P=((i_1,j_1),\dots,(i_L,j_L))
\]

with:

\[
(i_1,j_1)=(1,1),
\]

\[
(i_L,j_L)=(n,m),
\]

and local increments:

\[
(1,0),
\quad
(0,1),
\quad
(1,1).
\]

These steps allow one observed sample in one stream to align with more
than one consecutive sample in the other.

\section{The path is monotone}\label{the-path-is-monotone}

Indices never move backward.

This encodes an ordering assumption:

\begin{quote}
later aligned observations cannot be matched to earlier positions after
the path has passed them.
\end{quote}

That is already a modeling choice.

A system with genuine temporal reversals or reordering would need a
different alignment object.

\section{A local discrepancy}\label{a-local-discrepancy}

For scalar witness sequences use:

\[
c(i,j)
=
(x_i-y_j)^2.
\]

The cost of a path is:

\[
C(P)
=
\sum_{(i,j)\in P}
c(i,j).
\]

The optimal path solves:

\[
\boxed{
\operatorname{DTW}(X,Y)
=
\min_{P\in\mathcal P_{n,m}}
C(P).
}
\]

The optimum is relative to the declared cost and path class.

\section{Dynamic programming}\label{dynamic-programming}

Let:

\[
D(i,j)
\]

be the minimum cost of an admissible path ending at \((i,j)\).

Then:

\[
D(i,j)
=
c(i,j)
+
\min
\{
D(i-1,j),
D(i,j-1),
D(i-1,j-1)
\}.
\]

This is the recursive structure behind the finite alignment problem.

The terminal value:

\[
D(n,m)
\]

is the exact optimum.

\section{The exact warped witness}\label{the-exact-warped-witness}

Use:

\[
X=(0,1,2),
\]

\[
Y=(0,0,1,2).
\]

The first stream has three observations.

The second has four.

The second stream repeats the initial value.

That duplicate can represent many things:

\begin{itemize}
\tightlist
\item
  slower local progress;
\item
  repeated sampling;
\item
  a plateau;
\item
  asynchronous acquisition.
\end{itemize}

The witness does not choose among those explanations.

\section{The local-cost matrix}\label{the-local-cost-matrix}

The exact squared-distance matrix is:

\[
\begin{pmatrix}
0&0&1&4\\
1&1&0&1\\
4&4&1&0
\end{pmatrix}.
\]

The zero-cost cells are exactly:

\[
(1,1),
(1,2),
(2,3),
(3,4).
\]

\section{There are 25 admissible
paths}\label{there-are-25-admissible-paths}

Under the declared endpoint and step constraints, exhaustive enumeration
gives:

\[
\boxed{25}
\]

admissible paths from:

\[
(1,1)
\]

to:

\[
(3,4).
\]

The path-count recurrence is:

\[
N(i,j)
=
N(i-1,j)
+
N(i,j-1)
+
N(i-1,j-1).
\]

The terminal count is:

\[
N(3,4)=25.
\]

\section{The unique zero-cost
alignment}\label{the-unique-zero-cost-alignment}

Because every local cost is nonnegative, a zero-cost path can use only
zero-cost cells.

Only one admissible chain connects all required endpoints through those
cells:

\[
\boxed{
P^\star
=
((1,1),(1,2),(2,3),(3,4)).
}
\]

Thus:

\[
\boxed{
\operatorname{DTW}(X,Y)=0.
}
\]

And the optimum is unique.

\section{Observed index has separated from alignment
phase}\label{observed-index-has-separated-from-alignment-phase}

The first two observations of \(Y\):

\[
y_1=0,
\qquad
y_2=0
\]

both align to:

\[
x_1=0.
\]

Then:

\[
y_3=1
\]

aligns to:

\[
x_2=1,
\]

and:

\[
y_4=2
\]

aligns to:

\[
x_3=2.
\]

Therefore:

\[
j=2
\]

in one stream does not correspond to:

\[
i=2
\]

in the other.

This is the smallest exact latent-clock lesson.

\section{A common alignment
coordinate}\label{a-common-alignment-coordinate}

Index the path itself by:

\[
\ell=1,\dots,L.
\]

Each latent alignment position carries a pair:

\[
(i_\ell,j_\ell).
\]

This gives a common ordered phase coordinate.

It is useful to call this a latent alignment coordinate.

It is not yet justified to call it physical time.

\section{Why ``clock'' is a dangerous
word}\label{why-clock-is-a-dangerous-word}

A clock suggests something stronger than alignment.

It can suggest:

\begin{itemize}
\tightlist
\item
  elapsed physical duration;
\item
  causal progression;
\item
  an intrinsic phase variable;
\item
  a unique temporal parameterization.
\end{itemize}

DTW alone gives none of those automatically.

It gives an optimizer of an alignment objective.

The chapter therefore uses ``latent clock'' as an inference object whose
epistemic status must remain explicit.

\section{Identity/no-warp control}\label{identityno-warp-control}

Now use:

\[
X_0=(0,1,2),
\qquad
Y_0=(0,1,2).
\]

The two streams have equal length and already progress in the same
observed order.

There are:

\[
\boxed{13}
\]

admissible paths under the same rule.

\section{The diagonal is uniquely
optimal}\label{the-diagonal-is-uniquely-optimal}

The local-cost matrix is:

\[
\begin{pmatrix}
0&1&4\\
1&0&1\\
4&1&0
\end{pmatrix}.
\]

The only zero-cost cells are diagonal.

Therefore the unique zero-cost path is:

\[
\boxed{
((1,1),(2,2),(3,3)).
}
\]

So:

\[
\operatorname{DTW}(X_0,Y_0)=0
\]

without any warp.

\section{The control prevents a false
lesson}\label{the-control-prevents-a-false-lesson}

The warped witness alone could tempt us to think that latent-time
machinery always invents a nontrivial warp.

The control shows otherwise.

When observed indices already match the phase structure, the exact
optimum is the identity alignment.

Thus:

\[
\boxed{
\text{warping is inferred when supported by the declared data and cost, not imposed by default}.
}
\]

\section{Constraints are part of the
model}\label{constraints-are-part-of-the-model}

Changing the allowed local steps changes the feasible path set.

Adding a slope band or window can also change:

\begin{itemize}
\tightlist
\item
  the optimum;
\item
  its cost;
\item
  uniqueness.
\end{itemize}

Therefore path constraints are not implementation details.

They are assumptions.

Sakoe and Chiba's classical formulation makes time-normalization and
slope constraints part of the alignment problem (Sakoe and Chiba 1978).

\section{A unique optimum is still
model-relative}\label{a-unique-optimum-is-still-model-relative}

The finite witness has a unique optimum.

But that means only:

\begin{quote}
there is one best path under this observation space, squared cost,
endpoint rule, and step set.
\end{quote}

Change any of those, and the optimum may change.

Therefore:

\[
\boxed{
\text{unique optimum}
\not\Rightarrow
\text{unique true time}.
}
\]

\section{Multimodal alignment changes the comparison
problem}\label{multimodal-alignment-changes-the-comparison-problem}

Suppose one stream is:

\begin{itemize}
\tightlist
\item
  video;
\end{itemize}

and another is:

\begin{itemize}
\tightlist
\item
  motion capture;
\item
  audio;
\item
  physiology;
\item
  text events.
\end{itemize}

The raw coordinates are not generally comparable.

A reasonable framework needs representation maps:

\[
\phi_X:\mathcal X\to\mathcal Z,
\]

\[
\phi_Y:\mathcal Y\to\mathcal Z.
\]

Then local cost can be computed in the comparison space:

\[
d_{\mathcal Z}
(
\phi_X(x_i),
\phi_Y(y_j)
).
\]

\section{Generalized time warping}\label{generalized-time-warping}

Zhou and De la Torre explicitly treat multimodal temporal alignment and
allow richer monotonic warping structure (Zhou and De la Torre 2012).

That is useful precedent for LATENTTIME because the inferred clock need
not be tied to one raw sensor space.

But the Atlas does not inherit a universal multimodal representation.

The representation and discrepancy remain part of the declared model.

\section{Asynchrony is not one
phenomenon}\label{asynchrony-is-not-one-phenomenon}

Two streams can become misaligned because of:

\begin{itemize}
\tightlist
\item
  different sampling rates;
\item
  missing observations;
\item
  duplicated observations;
\item
  variable execution speed;
\item
  buffering;
\item
  latency;
\item
  sensor-specific preprocessing;
\item
  genuinely different event timing.
\end{itemize}

A low-cost warp cannot by itself identify which explanation is correct.

\section{Positional frequency is not a
clock}\label{positional-frequency-is-not-a-clock}

RPO can expose a strong low-frequency positional mode.

That can mean the operator varies slowly with sequence offset.

It does not imply that the same mode tracks:

\begin{itemize}
\tightlist
\item
  elapsed time;
\item
  behavioral phase;
\item
  event age;
\item
  causal stage.
\end{itemize}

Thus:

\[
\boxed{
\text{positional Fourier structure}
\not\Rightarrow
\text{latent temporal coordinate}.
}
\]

\section{RoPE frequency is not elapsed time
either}\label{rope-frequency-is-not-elapsed-time-either}

A RoPE frequency describes rotation in feature coordinates as position
changes.

Its phase is an algebraic function of positional offset.

That does not automatically make it a physical or semantic clock.

The same word ``phase'' can describe different mathematical objects.

LATENTTIME keeps them typed.

\section{Discrete alignment is not a continuous
trajectory}\label{discrete-alignment-is-not-a-continuous-trajectory}

The DTW path is:

\[
((i_1,j_1),\dots,(i_L,j_L)).
\]

This is a finite combinatorial object.

A continuous trajectory is something like:

\[
t\mapsto x(t)
\]

generated under an explicit dynamical model.

Those objects can be related only after adding:

\begin{itemize}
\tightlist
\item
  interpolation;
\item
  parameterization;
\item
  regularity assumptions;
\item
  possibly a law of motion.
\end{itemize}

Therefore:

\[
\boxed{
\text{DTW path}
\not\Rightarrow
\text{continuous-time flow}.
}
\]

\section{Continuous-time latent clocks require more
structure}\label{continuous-time-latent-clocks-require-more-structure}

If a model claims a scalar latent time:

\[
\tau
\]

with continuous evolution, it should specify how:

\[
i\mapsto\tau_i
\]

or:

\[
j\mapsto\tau_j
\]

is obtained.

It should also specify whether \(\tau\) is:

\begin{itemize}
\tightlist
\item
  merely monotone;
\item
  metrically meaningful;
\item
  differentiable;
\item
  dynamically generated.
\end{itemize}

A monotone ordering alone does not give a metric clock.

\section{Alignment is not
causality}\label{alignment-is-not-causality}

Suppose audio and video align tightly under a learned warp.

That supports temporal correspondence under the model.

It does not prove:

\begin{itemize}
\tightlist
\item
  audio causes video;
\item
  video causes audio;
\item
  one latent system drives both.
\end{itemize}

Thus:

\[
\boxed{
\text{alignment}
\not\Rightarrow
\text{causal direction}.
}
\]

\section{Low cost is not semantic
identity}\label{low-cost-is-not-semantic-identity}

Two sequences can be numerically easy to align while representing
different semantics.

Conversely, semantically corresponding events can look very different in
raw modalities.

So:

\[
\boxed{
\text{low alignment cost}
\not\Rightarrow
\text{semantic identity}.
}
\]

\section{Non-unique alignments should stay
non-unique}\label{non-unique-alignments-should-stay-non-unique}

Some data admit multiple optimal paths.

A responsible system should not silently pretend one arbitrary tie-break
is the discovered clock.

It should report:

\begin{itemize}
\tightlist
\item
  multiplicity;
\item
  tie-breaking;
\item
  alignment uncertainty;
\item
  downstream sensitivity where relevant.
\end{itemize}

The exact Atlas witness avoids this only because the optimum happens to
be unique.

\section{Partial and missing-data
alignment}\label{partial-and-missing-data-alignment}

The core witness uses full endpoint-to-endpoint matching.

Real streams may require:

\begin{itemize}
\tightlist
\item
  subsequence matching;
\item
  open-begin or open-end paths;
\item
  partial matching;
\item
  missing-event models.
\end{itemize}

Those are different problem definitions.

They should be declared rather than smuggled into ``DTW.''

\section{A useful latent-clock
protocol}\label{a-useful-latent-clock-protocol}

For two asynchronous streams:

\begin{enumerate}
\def\labelenumi{\arabic{enumi}.}
\tightlist
\item
  define the observed streams;
\item
  define any representation maps;
\item
  define the local discrepancy;
\item
  define endpoints and path constraints;
\item
  compute the optimal path;
\item
  test an identity/no-warp control where possible;
\item
  inspect uniqueness or multiplicity;
\item
  report the induced alignment coordinate;
\item
  keep physical-clock claims separate;
\item
  test downstream usefulness independently.
\end{enumerate}

\section{When the alignment coordinate is
useful}\label{when-the-alignment-coordinate-is-useful}

An inferred alignment coordinate can support:

\begin{itemize}
\tightlist
\item
  cross-modal fusion;
\item
  event matching;
\item
  phase-normalized comparison;
\item
  memory alignment;
\item
  asynchronous agent communication;
\item
  trajectory comparison.
\end{itemize}

Its usefulness can be empirical even when it is not a uniquely true
physical time variable.

\section{Atlas non-implications}\label{atlas-non-implications}

The chapter rejects:

\[
\text{observed position}
\not\Rightarrow
\text{latent time},
\]

\[
\text{unique optimal alignment}
\not\Rightarrow
\text{unique true clock},
\]

\[
\text{periodic positional mode}
\not\Rightarrow
\text{temporal clock},
\]

\[
\text{alignment}
\not\Rightarrow
\text{causality},
\]

\[
\text{discrete warping path}
\not\Rightarrow
\text{continuous-time flow},
\]

and:

\[
\text{low alignment cost}
\not\Rightarrow
\text{semantic identity}.
\]

\section{Atlas connections}\label{atlas-connections-12}

\textbf{Relative-Position Operators.}\\
RPO supplies operator/frequency structure over observed position;
LATENTTIME adds alignment inference across asynchronous indices.

\textbf{Dynamics.}\\
Dynamics supplies the distinction between sampled states, trajectories,
and flows; LATENTTIME keeps discrete warping separate from
continuous-time laws.

\textbf{Multimodal systems.}\\
A shared alignment coordinate can organize observations from
heterogeneous sensors or representations.

\textbf{Agent systems.}\\
Different agents can observe related events on different local clocks;
alignment can provide a common phase coordinate without assuming
synchronized wall time.

\textbf{Memory.}\\
Episodes with different sampling density can be compared after an
explicit alignment transformation.

\section{Closing view}\label{closing-view-19}

Position is cheap to observe.

Time is often not.

A sequence index tells us where a sample sits in a record.

It does not necessarily tell us where that sample sits in a shared
process.

The exact witness shows the smallest nontrivial case:

\[
X=(0,1,2),
\qquad
Y=(0,0,1,2).
\]

The second stream lingers on one phase.

The optimal correspondence therefore repeats one alignment state:

\[
(1,1)\to(1,2)\to(2,3)\to(3,4).
\]

The identity control shows the same machinery leaves already aligned
streams unwarped.

The right conclusion is:

\[
\boxed{
\text{latent time is an inferred correspondence structure whose meaning depends on the declared alignment model.}
}
\]

It can be useful without being a metaphysical clock.

\section{References used in this
chapter}\label{references-used-in-this-chapter-26}

\begin{itemize}
\tightlist
\item
  (Sakoe and Chiba 1978)
\item
  (Zhou and De la Torre 2012)
\end{itemize}

Exact source identities and claim boundaries are recorded in:

sources/source-locks/ATLAS-CH-LATENTTIME-001.yaml

\chapter{First-Order Optimization}\label{first-order-optimization}

\textbf{Epistemic status:} Established Theory + Bounded Deep-Learning
Practice + Exact Computational Witness\\
\textbf{Specification:}
manuscript/specifications/ATLAS-CH-OPTBASE-001.md\\
\textbf{Derivation packet:}
mathematics/derivations/ATLAS-CH-OPTBASE-001-DERIVATIONS.md\\
\textbf{Computational witness:}
mathematics/computational-witnesses/ATLAS-CW-OPTBASE-001.md\\
\textbf{Source lock:} sources/source-locks/ATLAS-CH-OPTBASE-001.yaml

\section{The gradient is not the
optimizer}\label{the-gradient-is-not-the-optimizer}

Under the Euclidean inner product on parameter coordinates, the gradient
answers a local question:

\begin{quote}
Among unit Euclidean directions, which infinitesimal direction increases
the objective fastest?
\end{quote}

For nonzero gradient, that direction is

\[
\frac{\nabla F}{\|\nabla F\|_2}.
\]

Its negative is the Euclidean steepest-descent direction.

Later geometry chapters will change the metric and therefore change what
``steepest'' means.

An optimizer answers a different question:

\begin{quote}
Given that signal, how should the parameters actually move?
\end{quote}

That second question introduces choices:

\begin{itemize}
\tightlist
\item
  step size;
\item
  momentum;
\item
  stochastic sampling;
\item
  coordinatewise scaling;
\item
  clipping;
\item
  weight decay;
\item
  schedules;
\item
  initialization;
\item
  optimizer state.
\end{itemize}

The Atlas therefore treats

\[
g_k
\]

and

\[
\Delta\theta_k
\]

as different objects.

The gradient is a signal.

The optimizer converts that signal into motion.

\section{Steering versus changing the
road}\label{steering-versus-changing-the-road}

A useful allegory is steering a vehicle down a sloped landscape.

The gradient indicates local downhill direction.

The learning rate controls how far we move.

Momentum carries history.

A preconditioner changes the coordinatewise steering response.

Clipping acts like a speed limiter.

Weight decay adds shrinkage.

A schedule changes the control policy over time.

The structural correspondence is useful.

The literal-mechanics analogy is not.

Optimization state is algorithmic state, not physical momentum unless a
specific mathematical correspondence is established.

\section{Gradient descent}\label{gradient-descent}

For a differentiable objective

\[
F(\theta),
\]

gradient descent uses

\[
\boxed{
\theta_{k+1}
=
\theta_k-\eta\nabla F(\theta_k).
}
\]

This is already a discrete dynamical system:

\[
\theta_{k+1}
=
\Psi(\theta_k).
\]

The learning rate \(\eta\) is not simply a ``speed'' parameter.

It changes the discrete map itself.

A step size can therefore alter stability.

\section{Exact quadratic
calibration}\label{exact-quadratic-calibration}

Take

\[
F(x)
=
\frac{\lambda}{2}x^2,
\qquad
\lambda>0.
\]

Then

\[
\nabla F(x)
=
\lambda x,
\]

so

\[
x_{k+1}
=
(1-\eta\lambda)x_k.
\]

Thus

\[
x_k
=
(1-\eta\lambda)^k x_0.
\]

The dynamics are completely determined by the scalar amplification
factor

\[
r
=
1-\eta\lambda.
\]

\section{Step-size stability}\label{step-size-stability}

For asymptotic convergence to zero,

\[
|r|<1.
\]

Therefore

\[
|1-\eta\lambda|<1,
\]

which gives

\[
\boxed{
0<\eta<\frac{2}{\lambda}.
}
\]

At

\[
\eta=0,
\]

the iterate does not move.

At

\[
\eta=\frac2\lambda,
\]

the factor is \(-1\), so the iterate alternates sign without decaying.

Beyond that boundary, the discrete method diverges.

\begin{figure}
\centering
\pandocbounded{\includegraphics[keepaspectratio,alt={Three exact panels: scalar quadratic gradient-descent amplification and its convergence interval; norm clipping of g=(3,4) to threshold 2; and the exact coupled-versus-decoupled adaptive decay next iterates 183/100 and 181/100.}]{figures/derivatives/ATLAS-FIG-OPTBASE-001-v0.1.1.png}}
\caption{Three exact panels: scalar quadratic gradient-descent
amplification and its convergence interval; norm clipping of g=(3,4) to
threshold 2; and the exact coupled-versus-decoupled adaptive decay next
iterates 183/100 and 181/100.}
\end{figure}

\emph{Corrected-edition figure candidate:} This plate is the separately
declared Matplotlib 3.10.3 editorial derivative in
\texttt{figures/derivatives/ATLAS-FIG-OPTBASE-001-v0.1.1.yaml}, not a
Wolfram replay. The original Wolfram source and master remain unchanged.
Print-scale and accessibility acceptance of the derivative is still
pending under \#324.

\section{Curvature sets the useful
scale}\label{curvature-sets-the-useful-scale}

The bound

\[
0<\eta<\frac2\lambda
\]

belongs to the declared quadratic curvature \(\lambda\).

Change the curvature and the stability range changes.

For a multidimensional quadratic with Hessian eigenvalues

\[
\lambda_i,
\]

one expects the most restrictive positive curvature scale to limit a
simple constant step.

This is the first hint that optimization should be read spectrally, not
only coordinatewise.

Later chapters will develop that point.

\section{Full gradient versus stochastic
gradient}\label{full-gradient-versus-stochastic-gradient}

Suppose

\[
F(\theta)
=
\mathbb E_\xi[\ell(\theta;\xi)].
\]

A stochastic-gradient method uses a random sample or minibatch gradient

\[
g(\theta;\xi)
\]

in place of the exact expectation.

The ideal unbiased relation is

\[
\mathbb E_\xi[g(\theta;\xi)]
=
\nabla F(\theta).
\]

Unbiasedness is an expectation statement.

It says nothing by itself about variance.

\section{Exact unbiased-SGD
witness}\label{exact-unbiased-sgd-witness}

Take two sample losses:

\[
\ell_1(x)
=
\frac12(x-1)^2,
\]

\[
\ell_2(x)
=
\frac12(x+1)^2.
\]

Sample uniformly.

The finite-average objective is

\[
F(x)
=
\frac12(\ell_1+\ell_2)
=
\frac12x^2+\frac12.
\]

Therefore

\[
\nabla F(x)=x.
\]

The sample gradients are

\[
g_1=x-1,
\qquad
g_2=x+1.
\]

Their expectation is

\[
\frac12(g_1+g_2)=x.
\]

So the estimator is unbiased.

Its variance is

\[
1.
\]

Thus:

\[
\boxed{
\text{unbiased}
\not\Rightarrow
\text{low variance}.
}
\]

Robbins and Monro provide the foundational stochastic-approximation
setting behind this distinction (Robbins and Monro 1951).

\section{Noise can be useful and
harmful}\label{noise-can-be-useful-and-harmful}

Stochasticity can:

\begin{itemize}
\tightlist
\item
  reduce per-step cost;
\item
  perturb iterates out of narrow regions;
\item
  act as implicit regularization;
\item
  increase estimator variance;
\item
  destabilize large steps;
\item
  complicate convergence diagnostics.
\end{itemize}

There is no single sign.

A gradient-noise model must be tied to:

\begin{itemize}
\tightlist
\item
  sampling scheme;
\item
  batch size;
\item
  data distribution;
\item
  optimizer;
\item
  current parameter state.
\end{itemize}

The Atlas will not treat ``SGD noise'' as one universal distribution.

\section{Momentum adds state}\label{momentum-adds-state}

One heavy-ball convention is

\[
v_{k+1}
=
\beta v_k+g_k,
\]

\[
\theta_{k+1}
=
\theta_k-\eta v_{k+1}.
\]

The optimizer state is now

\[
s_k=(\theta_k,v_k).
\]

This is a conceptual change.

The update no longer depends only on the current gradient.

It depends on history carried in \(v_k\).

Polyak's work supplies the classical acceleration context (Polyak 1964).

\section{Momentum conventions
differ}\label{momentum-conventions-differ}

Some implementations define velocity with the opposite sign.

Some absorb \(\eta\) into the momentum variable.

Some write updates directly as parameter differences.

Therefore two formulas both called ``momentum'' need not look identical.

The Atlas uses equations, not labels, as the canonical object.

This discipline becomes essential when we later study optimizer-state
Jacobians.

\section{Nesterov-style lookahead}\label{nesterov-style-lookahead}

A common lookahead form evaluates the gradient at a predicted state:

\[
g_k
=
\nabla F(\theta_k-\eta\beta v_k),
\]

followed by

\[
v_{k+1}
=
\beta v_k+g_k,
\]

\[
\theta_{k+1}
=
\theta_k-\eta v_{k+1}.
\]

The important structural change is:

\begin{quote}
the gradient is evaluated at a lookahead point rather than the current
parameter point.
\end{quote}

Nesterov's convex-optimization theory gives the classical accelerated
context (Nesterov 2004).

The Atlas does not transfer convex rate guarantees automatically to deep
nonconvex networks.

\section{Deep-learning momentum}\label{deep-learning-momentum}

Sutskever and colleagues provide an influential deep-learning study of
momentum and initialization (Sutskever et al. 2013).

That work is useful here as a bridge between classical first-order
machinery and practical neural-network training.

It is not used to rank momentum universally above all alternatives.

\section{AdaGrad}\label{adagrad}

AdaGrad accumulates squared gradients coordinatewise:

\[
G_{k,i}
=
G_{k-1,i}
+
g_{k,i}^2.
\]

A standard coordinate update is

\[
\theta_{k+1,i}
=
\theta_{k,i}
-
\eta
\frac{g_{k,i}}
{\sqrt{G_{k,i}}+\varepsilon}.
\]

Coordinates with large historical squared gradients receive smaller
effective step scales (Duchi et al. 2011).

This is more than a changing learning rate.

It is coordinate-dependent scaling.

\section{Adaptive scaling as a diagonal
geometry}\label{adaptive-scaling-as-a-diagonal-geometry}

Write a generic first-order update as

\[
\theta_{k+1}
=
\theta_k
-
\eta P_k g_k,
\]

where \(P_k\) is diagonal and positive.

When

\[
P_k=I,
\]

we recover ordinary gradient descent.

When \(P_k\) depends on accumulated gradient statistics, the algorithm
changes the relative step size across coordinates.

This is the beginning of a geometric viewpoint.

Later Atlas chapters will replace diagonal scaling with richer matrix
and manifold structure.

\section{RMSProp}\label{rmsprop}

RMSProp uses an exponential moving average of squared gradients:

\[
v_k
=
\rho v_{k-1}
+
(1-\rho)g_k^2,
\]

coordinatewise.

Then

\[
\theta_{k+1}
=
\theta_k
-
\eta
\frac{g_k}
{\sqrt{v_k}+\varepsilon}.
\]

The original RMSProp description is course lecture material rather than
a peer-reviewed paper (Tieleman and Hinton 2012).

The Atlas marks that source status explicitly.

\section{Adam}\label{adam}

Adam combines two moving averages (Kingma and Ba 2015).

First moment:

\[
m_k
=
\beta_1m_{k-1}
+
(1-\beta_1)g_k.
\]

Second raw moment:

\[
v_k
=
\beta_2v_{k-1}
+
(1-\beta_2)g_k^2.
\]

With zero initialization, both estimates are biased toward zero early in
training.

Adam corrects that initialization bias.

\section{Bias correction}\label{bias-correction}

The corrected moments are

\[
\hat m_k
=
\frac{m_k}{1-\beta_1^k},
\]

and

\[
\hat v_k
=
\frac{v_k}{1-\beta_2^k}.
\]

A standard update is

\[
\theta_{k+1}
=
\theta_k
-
\eta
\frac{\hat m_k}
{\sqrt{\hat v_k}+\varepsilon}.
\]

Bias correction is an algebraic adjustment.

It is not a proof that the full optimizer is unbiased or universally
convergent.

\section{Exact Adam first step}\label{exact-adam-first-step}

Take

\[
g_1=3,
\]

\[
\beta_1=\frac9{10},
\]

\[
\beta_2=\frac{99}{100},
\]

with

\[
m_0=v_0=0.
\]

Then

\[
m_1=\frac3{10},
\]

and

\[
v_1=\frac9{100}.
\]

After bias correction:

\[
\hat m_1=3,
\]

\[
\hat v_1=9.
\]

For the calibration

\[
\varepsilon=0,
\]

the normalized direction is exactly

\[
1.
\]

Practical implementations use nonzero epsilon.

The zero-epsilon witness exists only to isolate the correction algebra.

\section{Adaptive methods are not one
algorithm}\label{adaptive-methods-are-not-one-algorithm}

AdaGrad, RMSProp, and Adam all use gradient-history information.

But they differ in:

\begin{itemize}
\tightlist
\item
  accumulation versus exponential forgetting;
\item
  first-moment state;
\item
  bias correction;
\item
  denominator construction;
\item
  epsilon placement;
\item
  update ordering.
\end{itemize}

Treating them as one generic ``adaptive optimizer'' can hide
mathematically important differences.

\section{Gradient clipping}\label{gradient-clipping}

Global norm clipping at threshold \(\tau\) uses

\[
g_{\rm clip}
=
g
\min\left(
1,
\frac{\tau}{\|g\|_2}
\right).
\]

If the norm is below threshold, nothing changes.

If the norm is above threshold, the vector is rescaled to norm \(\tau\).

Pascanu, Mikolov, and Bengio discuss clipping in the context of
exploding gradients in recurrent networks (Pascanu et al. 2013).

\section{Exact clipping geometry}\label{exact-clipping-geometry}

Take

\[
g=(3,4).
\]

Then

\[
\|g\|_2=5.
\]

With threshold

\[
\tau=2,
\]

the scale factor is

\[
2/5.
\]

Therefore

\[
\boxed{
g_{\rm clip}
=
(6/5,8/5).
}
\]

Its norm is exactly

\[
2.
\]

The direction is unchanged in this global-norm scheme.

The magnitude is not.

\section{Clipping changes the
estimator}\label{clipping-changes-the-estimator}

Clipping is nonlinear.

Therefore in general

\[
\mathbb E[\operatorname{clip}(g)]
\neq
\operatorname{clip}(\mathbb E[g]).
\]

Even if the original stochastic gradient is unbiased,

\[
\mathbb E[g]
=
\nabla F,
\]

the clipped gradient need not remain unbiased.

This is not necessarily a defect.

It is simply a different estimator/update rule.

The distinction should be explicit.

\section{Weight decay and L2
regularization}\label{weight-decay-and-l2-regularization}

For plain gradient descent, adding an L2 penalty

\[
\frac{\lambda}{2}\|\theta\|^2
\]

produces gradient contribution

\[
\lambda\theta.
\]

The update becomes

\[
\theta^+
=
\theta-\eta(g+\lambda\theta).
\]

Rearranging,

\[
\theta^+
=
(1-\eta\lambda)\theta-\eta g.
\]

In this simple setting, L2 regularization and multiplicative weight
shrinkage coincide algebraically.

That equivalence does not survive arbitrary preconditioning.

\section{Coupled regularization under adaptive
scaling}\label{coupled-regularization-under-adaptive-scaling}

Suppose a scalar preconditioner \(p\) scales the gradient.

If L2 regularization is coupled into the gradient,

\[
\theta_{\rm c}^+
=
\theta-\eta p(g+\lambda\theta).
\]

The preconditioner acts on the shrinkage gradient too.

This changes the effective decay.

\section{Decoupled weight decay}\label{decoupled-weight-decay}

AdamW separates shrinkage from the adaptive gradient step (Loshchilov
and Hutter 2019).

In the scalar calibration:

\[
\theta_{\rm d}^+
=
(1-\eta\lambda)\theta
-
\eta p g.
\]

Subtracting coupled from decoupled gives

\[
\boxed{
\theta_{\rm c}^+
-
\theta_{\rm d}^+
=
\eta\lambda\theta(1-p).
}
\]

They coincide only under special conditions.

\section{Exact decay
counterexample}\label{exact-decay-counterexample}

Take

\[
\theta=2,
\quad
g=3,
\quad
p=\frac12,
\quad
\eta=\frac1{10},
\quad
\lambda=\frac15.
\]

Coupled regularization gives

\[
\theta_{\rm c}^+
=
\frac{183}{100}.
\]

Decoupled decay gives

\[
\theta_{\rm d}^+
=
\frac{181}{100}.
\]

The exact difference is

\[
\boxed{
\frac1{50}.
}
\]

So under adaptive scaling:

\[
\boxed{
\text{coupled L2}
\neq
\text{decoupled weight decay}
}
\]

in general.

\section{Learning-rate schedules}\label{learning-rate-schedules}

A constant learning rate uses

\[
\eta_k=\eta.
\]

A schedule makes the step scale time-dependent:

\[
\eta_k.
\]

Common patterns include:

\begin{itemize}
\tightlist
\item
  piecewise decay;
\item
  exponential or inverse-power decay;
\item
  cosine schedules;
\item
  warmup plus decay.
\end{itemize}

A schedule changes the optimization dynamics over training time.

It is therefore part of the algorithm, not cosmetic metadata.

\section{Warmup}\label{warmup}

Warmup deliberately begins with smaller step scales.

One influential example is the Transformer schedule of Vaswani et al.,
which combines warmup with inverse-square-root scaling (Vaswani et al.
2017).

The Atlas treats warmup as a training recipe with many possible
rationales:

\begin{itemize}
\tightlist
\item
  unstable early statistics;
\item
  large transient gradients;
\item
  poorly calibrated optimizer moments;
\item
  scale changes during early representation formation.
\end{itemize}

No universal theorem says that warmup is always necessary or optimal.

\section{Initialization is part of
optimization}\label{initialization-is-part-of-optimization}

Optimization starts before the first update.

The initial parameter distribution determines:

\begin{itemize}
\tightlist
\item
  activation scale;
\item
  gradient scale;
\item
  symmetry breaking;
\item
  saturation risk;
\item
  early Jacobian statistics.
\end{itemize}

Initialization therefore belongs in the optimization baseline.

It defines the starting state of the dynamical system.

\section{Fan-in second moment}\label{fan-in-second-moment}

Consider

\[
z
=
\sum_{i=1}^{n}w_i x_i.
\]

Assume:

\begin{itemize}
\tightlist
\item
  independent zero-mean weights;
\item
  independent zero-mean inputs;
\item
  weight/input independence;
\item
  common input second moment \(q\);
\item
  common weight variance \(\sigma_w^2\).
\end{itemize}

Cross terms vanish.

Therefore

\[
\boxed{
\mathbb E[z^2]
=
n\sigma_w^2q.
}
\]

This is the mathematical core of fan-in scaling.

\section{Linear second-moment
preservation}\label{linear-second-moment-preservation}

To keep

\[
\mathbb E[z^2]
=
q,
\]

choose

\[
\sigma_w^2
=
\frac1n.
\]

Glorot and Bengio analyze related variance/saturation issues in deep
feedforward networks (Glorot and Bengio 2010).

The derivation depends on its assumptions.

Real network activations are not guaranteed to remain independent or
identically distributed across depth.

\section{ReLU second moment}\label{relu-second-moment}

Suppose the preactivation distribution is symmetric around zero.

Then

\[
\operatorname{ReLU}(z)^2
=
z^2\mathbf 1_{\{z>0\}}.
\]

By symmetry,

\[
\mathbb E[\operatorname{ReLU}(z)^2]
=
\frac12\mathbb E[z^2].
\]

Therefore

\[
\mathbb E[\operatorname{ReLU}(z)^2]
=
\frac12 n\sigma_w^2q.
\]

Choose

\[
\boxed{
\sigma_w^2
=
\frac2n
}
\]

to preserve the second moment.

He and colleagues develop rectifier-aware initialization in this
direction (He et al. 2015).

\section{Second moment is not variance after
ReLU}\label{second-moment-is-not-variance-after-relu}

This wording matters.

ReLU output generally has positive mean.

Therefore

\[
\operatorname{Var}[\operatorname{ReLU}(z)]
\neq
\mathbb E[\operatorname{ReLU}(z)^2]
\]

unless the mean term is handled.

The simple \(2/n\) derivation above is an exact second-moment statement
under symmetry assumptions.

It should not be inflated into an exact theorem about full deep-network
activation distributions.

\section{Optimization is a state
machine}\label{optimization-is-a-state-machine}

The simplest gradient method has state

\[
\theta_k.
\]

Momentum has

\[
(\theta_k,v_k).
\]

Adam has

\[
(\theta_k,m_k,v_k).
\]

Schedules introduce time or step counters.

Mixed precision may introduce loss-scaling state.

Gradient accumulation introduces buffering state.

So even ``first-order optimization'' quickly becomes a structured state
machine.

This prepares the later Optimizer-State Dynamics chapter.

\section{Coordinate scaling already hints at
geometry}\label{coordinate-scaling-already-hints-at-geometry}

Ordinary gradient descent treats parameter coordinates through the
Euclidean metric implicit in

\[
\nabla F.
\]

Adaptive methods introduce coordinate-dependent scales.

That raises a question the baseline does not answer:

\begin{quote}
Which geometry should determine steepest descent?
\end{quote}

Natural gradient, mirror descent, Shampoo, Muon, manifold methods, and
MODULUS all answer that question differently.

They belong downstream.

\section{Spectral effects are still
hidden}\label{spectral-effects-are-still-hidden}

The scalar quadratic example had one curvature scale \(\lambda\).

Real parameter blocks can have many singular/eigen directions.

A diagonal adaptive method sees only coordinatewise statistics.

It does not directly diagonalize the true curvature or account for
non-normal transient behavior.

That gap motivates:

\begin{itemize}
\tightlist
\item
  second-order methods;
\item
  matrix preconditioning;
\item
  spectral shaping;
\item
  optimizer-state diagnostics.
\end{itemize}

\section{Five failure modes}\label{five-failure-modes-3}

\subsection{Large step mistaken for fast
learning}\label{large-step-mistaken-for-fast-learning}

Beyond the stability boundary, a larger step can destroy convergence.

\subsection{Unbiased mistaken for low
variance}\label{unbiased-mistaken-for-low-variance}

The exact SGD witness has variance \(1\) despite being unbiased.

\subsection{Clipping mistaken for harmless
rescaling}\label{clipping-mistaken-for-harmless-rescaling}

Clipping changes the estimator whenever the threshold is active.

\subsection{Coupled L2 mistaken for AdamW
decay}\label{coupled-l2-mistaken-for-adamw-decay}

Adaptive scaling breaks the simple equivalence.

\subsection{Initialization heuristic mistaken for distribution
theorem}\label{initialization-heuristic-mistaken-for-distribution-theorem}

Fan-in calculations rely on assumptions that deep networks can violate.

\section{What would falsify an optimizer
story?}\label{what-would-falsify-an-optimizer-story}

A useful optimizer explanation should predict something measurable.

Examples:

\begin{itemize}
\tightlist
\item
  stability boundary;
\item
  update norm;
\item
  coordinate scale;
\item
  momentum phase lag;
\item
  decay behavior;
\item
  clipping activation rate;
\item
  variance propagation;
\item
  sensitivity to schedule.
\end{itemize}

If the predicted mechanism is absent, the explanation is incomplete.

The Atlas will use this standard later when testing optimizer geometry
and state dynamics.

\section{Atlas connections}\label{atlas-connections-13}

\textbf{Dynamics.}\\
An optimizer defines a discrete dynamical system, often with hidden
state.

\textbf{Second-order optimization.}\\
Curvature information can replace or augment first-order scaling.

\textbf{Manifold optimization.}\\
The notion of allowable motion changes when parameters live on
constrained spaces.

\textbf{Spectral shaping.}\\
Matrix updates can be designed around singular/eigen structure rather
than coordinates.

\textbf{Muon.}\\
Orthogonalized or spectrally shaped updates depart sharply from diagonal
adaptive scaling.

\textbf{CPS.}\\
Optimizer-state Jacobians expose coupling and transient instability
invisible in static update formulas.

\textbf{Curriculum.}\\
The optimization state interacts with which examples or tasks are
presented.

\section{Closing view}\label{closing-view-20}

First-order optimization begins with a simple equation:

\[
\theta_{k+1}
=
\theta_k-\eta g_k.
\]

But practical training rapidly adds structure:

\[
\text{gradient}
\to
\text{stochastic estimator}
\to
\text{state}
\to
\text{preconditioner}
\to
\text{clip}
\to
\text{decay}
\to
\text{schedule}.
\]

Each addition changes the actual dynamical system.

That is the central lesson.

The baseline is not ``gradient descent plus tricks.''

It is a family of stateful rules that translate local first-order
information into parameter motion.

Later Atlas chapters will ask whether those rules should instead be
derived from:

\begin{itemize}
\tightlist
\item
  curvature;
\item
  geometry;
\item
  spectra;
\item
  manifolds;
\item
  variational principles;
\item
  compositional contracts.
\end{itemize}

To ask that question rigorously, we first needed to say exactly what the
conventional baseline is.

\section{References used in this
chapter}\label{references-used-in-this-chapter-27}

\begin{itemize}
\tightlist
\item
  (Robbins and Monro 1951)
\item
  (Polyak 1964)
\item
  (Nesterov 2004)
\item
  (Duchi et al. 2011)
\item
  (Tieleman and Hinton 2012)
\item
  (Kingma and Ba 2015)
\item
  (Loshchilov and Hutter 2019)
\item
  (Pascanu et al. 2013)
\item
  (Glorot and Bengio 2010)
\item
  (He et al. 2015)
\item
  (Sutskever et al. 2013)
\item
  (Vaswani et al. 2017)
\end{itemize}

See `sources/source-locks/ATLAS-CH-OPTBASE-001.yaml` for exact source
roles and claim boundaries.

\chapter{Curvature and Second-Order
Structure}\label{curvature-and-second-order-structure}

\textbf{Epistemic status:} mathematical exposition built from audited
First-Order Optimization and Geometry prerequisites, classical
second-order optimization sources, and Atlas-owned exact witnesses.\\
\textbf{Derivation packet:}
`mathematics/derivations/ATLAS-CH-SECOND-001-DERIVATIONS.md`\\
\textbf{Computational witness:}
`mathematics/computational-witnesses/ATLAS-CW-SECOND-001.md`\\
\textbf{Source lock:} `sources/source-locks/ATLAS-CH-SECOND-001.yaml`

\section{The gradient tells you a direction; curvature tells you how
that direction
changes}\label{the-gradient-tells-you-a-direction-curvature-tells-you-how-that-direction-changes}

First-order optimization begins with a local slope.

For a differentiable objective \(f\), the Euclidean gradient

\[
g=\nabla f(x)
\]

is the first-order signal.

It answers:

\begin{quote}
What infinitesimal direction changes the objective most rapidly under
the declared Euclidean metric?
\end{quote}

Second-order structure asks another question:

\begin{quote}
How does that first-order signal change as we move?
\end{quote}

For twice-differentiable \(f\), the Hessian

\[
H=\nabla^2 f(x)
\]

is the local linear map that differentiates the gradient.

Gradient magnitude and curvature are therefore different objects.

A small gradient can occur in a flat basin, at a sharp local extremum,
or near a saddle.

Curvature helps distinguish those cases locally.

\section{Local quadratic model}\label{local-quadratic-model}

Near \(x\), a twice-differentiable objective has the local model

\[
m_x(p)
=
f(x)
+
g^\top p
+
\frac12 p^\top H p.
\]

The first term is the current objective.

The linear term is the first-order prediction.

The quadratic term describes how the prediction bends.

If \(H\) is positive definite, the quadratic term bends upward in every
nonzero direction.

If \(H\) is indefinite, some directions bend upward and others downward.

That sign structure is central.

\section{Newton's equation}\label{newtons-equation}

A stationary point of the local quadratic model satisfies

\[
\nabla_p m_x(p)
=
g+Hp
=
0.
\]

If \(H\) is nonsingular,

\[
\boxed{
p_N=-H^{-1}g.
}
\]

This is the Newton step.

It is tempting to read this as ``divide the gradient by curvature.''

That phrase is useful only if we remember that curvature is a matrix,
not a scalar.

The inverse Hessian rotates and rescales the gradient according to local
curvature directions.

\section{Positive-definite curvature gives a descent
direction}\label{positive-definite-curvature-gives-a-descent-direction}

Suppose

\[
H\succ0
\]

and

\[
g\neq0.
\]

Then

\[
g^\top p_N
=
-g^\top H^{-1}g
<
0.
\]

So the Newton direction is locally descending.

This fact depends on positive definiteness.

It is not true for arbitrary invertible Hessians.

That boundary will matter immediately.

\section{Exact positive-definite
witness}\label{exact-positive-definite-witness}

Take

\[
H=
\begin{pmatrix}
1&0\\
0&4
\end{pmatrix},
\qquad
b=
\begin{pmatrix}
1\\
1
\end{pmatrix},
\]

and

\[
f(x)
=
\frac12 x^\top Hx-b^\top x.
\]

At

\[
x_0=0,
\]

the gradient is

\[
g_0=-b
=
\begin{pmatrix}
-1\\
-1
\end{pmatrix}.
\]

The Newton equation is

\[
Hp=-g_0=b.
\]

Thus

\[
p_N
=
H^{-1}b
=
\begin{pmatrix}
1\\
1/4
\end{pmatrix}.
\]

Because the objective itself is quadratic, the local model is exact.

The step reaches the unique minimizer in one move.

The objective changes from

\[
f(0)=0
\]

to

\[
f(p_N)=-5/8.
\]

The directional derivative is

\[
g_0^\top p_N=-5/4<0.
\]

This is the clean case.

\section{Indefinite curvature breaks the simple Newton
story}\label{indefinite-curvature-breaks-the-simple-newton-story}

Now consider

\[
f(x,y)
=
\frac12(x^2-y^2).
\]

At

\[
z_0=(0,1)^\top,
\]

we have

\[
g=
\begin{pmatrix}
0\\
-1
\end{pmatrix},
\qquad
H=
\begin{pmatrix}
1&0\\
0&-1
\end{pmatrix}.
\]

The Newton equation gives

\[
p_N=
\begin{pmatrix}
0\\
-1
\end{pmatrix}.
\]

But

\[
g^\top p_N=1>0.
\]

So the raw Newton direction is an ascent direction.

The full step moves to

\[
(0,0),
\]

where

\[
f(0,0)=0
\]

instead of the starting value

\[
f(0,1)=-1/2.
\]

The objective increases.

The Newton equation solved the stationary condition of the local
quadratic model.

That stationary point was a saddle.

This is why second-order methods need more than ``invert the Hessian.''

\section{Curvature signatures}\label{curvature-signatures}

For a symmetric Hessian:

\begin{itemize}
\tightlist
\item
  all positive eigenvalues indicate locally positive curvature;
\item
  all negative eigenvalues indicate locally negative curvature;
\item
  mixed signs indicate an indefinite saddle-like quadratic model;
\item
  zero eigenvalues indicate locally flat or degenerate directions.
\end{itemize}

This is local information.

A positive-definite Hessian at one point does not make the entire
objective globally convex.

An indefinite Hessian at one point does not completely describe the
global landscape.

\section{Damped and regularized Newton
systems}\label{damped-and-regularized-newton-systems}

One response to difficult curvature is to modify the local system.

A common form is

\[
(H+\lambda I)p=-g.
\]

For sufficiently large positive \(\lambda\), the shifted matrix may
become positive definite.

This changes the local model.

It can improve stability.

But it is no longer the exact unmodified Newton step.

The Atlas keeps that distinction explicit:

\[
\boxed{
\text{exact Hessian}
\neq
\text{regularized Hessian}.
}
\]

\section{Trust regions: constrain where the model is
trusted}\label{trust-regions-constrain-where-the-model-is-trusted}

A different globalization strategy is to keep the quadratic model but
limit how far we are willing to follow it.

A trust-region subproblem has the form

\[
\min_p
\left(
g^\top p
+
\frac12 p^\top Bp
\right)
\]

subject to

\[
\|p\|\le\Delta.
\]

The radius \(\Delta\) is not merely a clipping threshold.

It expresses a model-validity region.

The algorithm compares predicted and actual decrease and can adjust the
radius accordingly (Nocedal and Wright 2006).

\section{Exact trust-region control on the indefinite
witness}\label{exact-trust-region-control-on-the-indefinite-witness}

Return to

\[
f(x,y)=\frac12(x^2-y^2)
\]

at

\[
(0,1).
\]

Use trust-region radius

\[
\Delta=1.
\]

The model increment is

\[
q(p)
=
-p_2
+
\frac12 p_1^2
-
\frac12 p_2^2.
\]

Over the unit disk, the exact minimizer is

\[
p_{\rm TR}
=
\begin{pmatrix}
0\\
1
\end{pmatrix}.
\]

The model change is

\[
q(p_{\rm TR})=-3/2.
\]

The new point is

\[
(0,2),
\]

where

\[
f(0,2)=-2.
\]

The raw Newton step climbed toward the saddle.

The trust-region subproblem moved along negative curvature and decreased
the objective.

This one example does not prove a universal trust-region advantage.

It shows why indefinite curvature needs different logic from
positive-definite Newton.

\section{Exact Hessians can be too large to
store}\label{exact-hessians-can-be-too-large-to-store}

For \(d\) parameters, a dense Hessian has \(d^2\) entries.

That quickly becomes impractical.

But many second-order algorithms do not need the full matrix.

They need the action

\[
Hv
\]

on selected vectors.

Pearlmutter showed how such products can be computed exactly in the
cited neural-computation setting without explicitly forming the dense
Hessian (Pearlmutter 1994).

The fundamental identity is

\[
Hv
=
\frac{d}{d\varepsilon}
\nabla f(x+\varepsilon v)
\Big|_{\varepsilon=0}.
\]

This is the second-order analogue of matrix-free operator access.

\section{Matrix-free second-order
structure}\label{matrix-free-second-order-structure}

Once Hessian-vector products are available, several operations become
possible without storing \(H\):

\begin{itemize}
\tightlist
\item
  curvature probing;
\item
  iterative linear solves for Newton-like systems;
\item
  extremal-eigenvalue estimation;
\item
  conjugate-gradient methods in suitable positive-definite systems;
\item
  trust-region subproblem methods.
\end{itemize}

This does not make second-order computation free.

The relevant costs move into repeated products, iterative solves,
damping, preconditioning, and numerical stability.

\section{Fisher information is a different
matrix}\label{fisher-information-is-a-different-matrix}

Optimization discussions sometimes slide from ``curvature'' to ``Fisher
information'' as if they were synonyms.

They are not.

For a probabilistic model

\[
p_\theta(y),
\]

the Fisher information matrix is built from score covariances, for
example

\[
F(\theta)
=
\mathbb E
\left[
\nabla_\theta \log p_\theta(y)
\nabla_\theta \log p_\theta(y)^\top
\right]
\]

under the declared distribution.

The Hessian is a second derivative of a declared scalar objective.

These matrices can be related in important settings.

They are not generically identical.

\section{Natural gradient}\label{natural-gradient}

The ordinary gradient depends on the metric used to define steepest
descent.

In the Euclidean metric, the local squared norm is

\[
\|\delta\|_2^2
=
\delta^\top\delta.
\]

Under a positive-definite metric \(G(\theta)\),

\[
\|\delta\|_G^2
=
\delta^\top G(\theta)\delta.
\]

The steepest-descent direction changes to a direction proportional to

\[
-G^{-1}g.
\]

When \(G\) is the Fisher information metric, this gives the natural
gradient (Amari 1998).

This is a geometric metric statement.

It is not automatically a Newton step.

\section{Exact metric witness}\label{exact-metric-witness}

Let

\[
G=
\begin{pmatrix}
1&0\\
0&4
\end{pmatrix},
\qquad
g=
\begin{pmatrix}
1\\
1
\end{pmatrix}.
\]

Euclidean steepest descent uses

\[
-g=
\begin{pmatrix}
-1\\
-1
\end{pmatrix}.
\]

Metric steepest descent uses

\[
-G^{-1}g
=
\begin{pmatrix}
-1\\
-1/4
\end{pmatrix}.
\]

The directions differ.

Changing the metric changes what ``steepest'' means.

The numerical coincidence that this matrix resembles the
positive-definite Hessian in our earlier witness does not identify
natural gradient with Newton.

The interpretation of the matrix is different.

\section{Hessian versus Fisher}\label{hessian-versus-fisher}

A useful discipline is to ask three separate questions:

\begin{enumerate}
\def\labelenumi{\arabic{enumi}.}
\tightlist
\item
  What scalar objective is being differentiated?
\item
  What probability distribution defines the Fisher expectation?
\item
  What metric is being used to define steepest descent?
\end{enumerate}

If those objects are not explicitly aligned, the words ``curvature,''
``Fisher,'' and ``natural gradient'' should not be substituted for one
another.

This prevents a common category error.

\section{Quasi-Newton methods}\label{quasi-newton-methods}

Exact Hessians can be expensive.

Quasi-Newton methods instead update a matrix approximation using
observed changes in gradients.

Let

\[
s_k=x_{k+1}-x_k
\]

and

\[
y_k=
\nabla f(x_{k+1})
-
\nabla f(x_k).
\]

A Hessian approximation \(B_{k+1}\) may be required to satisfy the
secant equation

\[
B_{k+1}s_k=y_k.
\]

This forces the approximation to match the observed curvature action
along the most recent displacement.

It does not make \(B_{k+1}\) the exact Hessian in every direction.

\section{BFGS and curvature
conditions}\label{bfgs-and-curvature-conditions}

BFGS is one of the canonical quasi-Newton updates (Nocedal and Wright
2006).

For the standard positive-definite Hessian-approximation form, the
condition

\[
y_k^\top s_k>0
\]

is load-bearing for preserving positive definiteness when the previous
approximation is positive definite.

This matters because the sign of the approximate curvature influences
whether the resulting search direction is descending.

The update formula is not magic.

Its guarantees have hypotheses.

\section{Secant replay}\label{secant-replay}

For our positive-definite quadratic

\[
H=
\operatorname{diag}(1,4),
\]

take

\[
s=
\begin{pmatrix}
1\\
2
\end{pmatrix}.
\]

Then the exact gradient change is

\[
y=Hs
=
\begin{pmatrix}
1\\
8
\end{pmatrix}.
\]

Any quasi-Newton matrix satisfying

\[
Bs=y
\]

matches the true curvature along that one displacement.

That single equation still leaves many degrees of freedom in \(B\).

A secant condition is partial curvature information.

\section{Proximal structure is not ``more
curvature''}\label{proximal-structure-is-not-more-curvature}

Second-order optimization is often discussed alongside proximal methods
because both go beyond a bare gradient step.

But proximal structure addresses a different problem.

Suppose

\[
F(x)=f(x)+h(x),
\]

where \(f\) is smooth and \(h\) may be nonsmooth.

The proximal operator is

\[
\operatorname{prox}_{\alpha h}(v)
=
\arg\min_x
\left[
h(x)
+
\frac{1}{2\alpha}\|x-v\|_2^2
\right].
\]

It is an implicit local optimization primitive (Parikh and Boyd 2014).

It is not a Hessian approximation.

\section{Soft thresholding witness}\label{soft-thresholding-witness}

Take

\[
h(x)=\lambda |x|
\]

with

\[
\alpha\lambda=1.
\]

Then

\[
\operatorname{prox}_{\alpha h}(v)
=
\operatorname{sign}(v)
\max(|v|-1,0).
\]

Thus

\[
v=3
\longmapsto
2,
\]

while

\[
v=1/2
\longmapsto
0.
\]

The zero output comes from solving the nonsmooth proximal subproblem.

It is not ordinary gradient clipping.

\section{Proximal gradient}\label{proximal-gradient}

For composite objectives

\[
F(x)=f(x)+h(x),
\]

a proximal-gradient step takes the form

\[
x_{k+1}
=
\operatorname{prox}_{\alpha h}
\left(
x_k-\alpha\nabla f(x_k)
\right).
\]

The step explicitly separates:

\begin{itemize}
\tightlist
\item
  smooth first-order information from \(f\);
\item
  implicit nonsmooth structure from \(h\).
\end{itemize}

This is a different axis from exact second derivatives.

\section{What second-order methods actually
buy}\label{what-second-order-methods-actually-buy}

Second-order structure can provide:

\begin{itemize}
\tightlist
\item
  anisotropic rescaling;
\item
  local curvature-aware directions;
\item
  faster local convergence in favorable regimes;
\item
  curvature diagnostics;
\item
  principled handling of local quadratic models.
\end{itemize}

But the price can include:

\begin{itemize}
\tightlist
\item
  expensive curvature access;
\item
  difficult indefinite systems;
\item
  noisy estimates;
\item
  storage or iterative-solve costs;
\item
  damping/globalization logic;
\item
  metric or model misspecification.
\end{itemize}

There is no free ``use curvature'' switch.

\section{Local models and global
optimization}\label{local-models-and-global-optimization}

A second-order method is still local.

The Hessian tells us how the objective bends near the current point.

A trust region limits how far we believe that local model.

A line search tests a step along a direction.

A quasi-Newton matrix infers curvature from recent changes.

A natural gradient changes the metric.

A proximal operator solves a declared local implicit subproblem.

These mechanisms solve different problems.

They should not be merged into one vague notion of sophistication.

\section{Failure modes}\label{failure-modes-4}

\subsection{Hessian equals Fisher}\label{hessian-equals-fisher}

Not in general.

They arise from different constructions.

\subsection{Natural gradient equals
Newton}\label{natural-gradient-equals-newton}

Not in general.

Natural gradient uses a metric; Newton uses the objective Hessian.

\subsection{Invertible Hessian means safe Newton
step}\label{invertible-hessian-means-safe-newton-step}

False.

The indefinite witness gives an ascent direction.

\subsection{Trust region equals
clipping}\label{trust-region-equals-clipping}

False.

A trust region solves a constrained model subproblem and reasons about
model validity.

\subsection{Quasi-Newton equals exact
curvature}\label{quasi-newton-equals-exact-curvature}

False.

The approximation satisfies limited secant information.

\subsection{HVP means Hessian is
cheap}\label{hvp-means-hessian-is-cheap}

False.

It avoids explicit formation but repeated products and solves still cost
computation.

\subsection{Proximal method is
second-order}\label{proximal-method-is-second-order}

Not by definition.

It is an implicit optimization construction, especially useful for
nonsmooth/composite structure.

\subsection{Positive local curvature means global
convexity}\label{positive-local-curvature-means-global-convexity}

False.

Local Hessian information is local.

\subsection{Exact local model means global
optimum}\label{exact-local-model-means-global-optimum}

False.

Newton solves the local quadratic stationary condition.

Global behavior requires additional assumptions and globalization logic.

\section{Downstream handoff}\label{downstream-handoff-3}

Matrix-Aware Optimization may inherit:

\begin{itemize}
\tightlist
\item
  Hessian and curvature notation;
\item
  positive/indefinite curvature distinctions;
\item
  matrix-free operator action;
\item
  natural-gradient metric language;
\item
  quasi-Newton exact-versus-approximate discipline;
\item
  trust-region and damping boundaries.
\end{itemize}

It may not assume that these classical objects already justify:

\begin{itemize}
\tightlist
\item
  Shampoo's specific preconditioner;
\item
  polar-factor updates;
\item
  orthogonalized matrix steps;
\item
  Muon-like spectral shaping;
\item
  square-to-rectangular generalization.
\end{itemize}

Those are downstream obligations.

\section{References used in this
chapter}\label{references-used-in-this-chapter-28}

\begin{itemize}
\tightlist
\item
  (Nocedal and Wright 2006) --- Newton, quasi-Newton, trust-region, and
  large-scale smooth optimization.
\item
  (Amari 1998) --- natural gradient and information geometry.
\item
  (Pearlmutter 1994) --- exact Hessian-vector products.
\item
  (Parikh and Boyd 2014) --- proximal operators and algorithms.
\end{itemize}

Exact provenance and claim boundaries are locked in:

`sources/source-locks/ATLAS-CH-SECOND-001.yaml`

\chapter{Matrix-Aware Optimization}\label{matrix-aware-optimization}

\textbf{Epistemic status:} audited Linear Algebra and
Curvature/Second-Order prerequisites + primary Shampoo,
polar-decomposition, norm-steepest-descent, and Muon design sources +
Atlas derivation + exact rectangular witness.\\
\textbf{Specification:}
manuscript/specifications/ATLAS-CH-MATRIXOPT-001.md\\
\textbf{Derivation packet:}
mathematics/derivations/ATLAS-CH-MATRIXOPT-001-DERIVATIONS.md\\
\textbf{Computational witness:}
mathematics/computational-witnesses/ATLAS-CW-MATRIXOPT-001.md\\
\textbf{Source lock:} sources/source-locks/ATLAS-CH-MATRIXOPT-001.yaml

\section{A matrix parameter is not merely a bag of
scalars}\label{a-matrix-parameter-is-not-merely-a-bag-of-scalars}

A linear layer may have a weight

\[
W\in\mathbb R^{m\times n},
\]

with gradient

\[
G=\nabla_W f(W).
\]

An optimizer can legitimately treat every entry of \(G\) as one scalar
coordinate.

But the same matrix also has row directions, column directions, singular
vectors, singular values, rank, and aspect ratio.

Matrix-aware optimization asks whether that structure should shape the
update.

The claim is not that matrix-aware methods are automatically better.

The claim is that they define a different local geometry.

\section{Three ways to change a matrix
gradient}\label{three-ways-to-change-a-matrix-gradient}

An elementwise rule may rescale entries separately.

A row/column-aware rule may transform

\[
G\mapsto AGB.
\]

An orthogonalizing rule may preserve singular vectors while transforming
singular values.

These operations can produce different updates from the same \(G\).

They should not share one name merely because each is called
``preconditioning'' informally.

\section{Left and right
multiplication}\label{left-and-right-multiplication}

For compatible \(A,G,B\),

\[
\operatorname{vec}(AGB)
=
(B^\top\otimes A)\operatorname{vec}(G).
\]

So a left-right matrix transform is a structured operator on the
vectorized gradient.

It can couple coordinates across rows and columns.

That is richer than independent scalar scaling.

It is still not automatically a full Hessian or Fisher inverse. A
structured preconditioner and exact curvature are distinct unless a
separate result identifies them.

\section{Shampoo: accumulated matrix
geometry}\label{shampoo-accumulated-matrix-geometry}

Shampoo was introduced as a preconditioned stochastic tensor
optimization method (Gupta et al. 2018).

For a matrix gradient sequence, its structure can be seen through
accumulated left and right statistics:

\[
L_t=\sum_{s\le t}G_sG_s^\top,
\]

\[
R_t=\sum_{s\le t}G_s^\top G_s.
\]

A representative damped matrix update is

\[
P_t
=
(L_t+\varepsilon I)^{-1/4}
G_t
(R_t+\varepsilon I)^{-1/4}.
\]

This remembers row-space and column-space gradient correlations and
applies matrix inverse roots rather than only coordinatewise division.

But it deliberately uses structured factors. It is not, by definition,
an arbitrary dense full-curvature inverse over every scalar parameter.

The convergence and empirical statements in the original paper remain
scoped to that paper's assumptions and experiments.

\section{Rectangular matrices change the
algebra}\label{rectangular-matrices-change-the-algebra}

Take a tall matrix

\[
G\in\mathbb R^{m\times n},
\qquad m>n.
\]

Even if \(G\) has full column rank,

\[
\operatorname{rank}(G)=n,
\]

the left Gram matrix obeys

\[
\operatorname{rank}(GG^\top)=n<m.
\]

Thus \(GG^\top\) is singular.

Nothing is defective. The gradient's column space simply occupies only
an \(n\)-dimensional subspace of \(\mathbb R^m\).

This is why square inverse formulas cannot be transferred mechanically
to rectangular blocks.

Damping, accumulated history, pseudoinverse conventions, or
support-restricted inverse powers change the actual method.

\section{The polar factor}\label{the-polar-factor}

Let \(r=\operatorname{rank}(G)\), and take a compact SVD retaining only
the \(r\) strictly positive singular values:

\[
G=U_r\Sigma_rV_r^\top.
\]

The support-restricted polar factor is

\[
Q=U_rV_r^\top.
\]

For \(G=0\), set \(Q=0\). For full-column-rank tall \(G\), this agrees
with the usual \(UV^\top\).

For a tall full-column-rank matrix,

\[
Q^\top Q=I_n.
\]

But generally

\[
QQ^\top\ne I_m.
\]

In the full-column-rank case, the columns are orthonormal; the rows
cannot all be orthonormal because there are more rows than columns. If
\(G\) has rank \(r<n\), instead \(Q^\top Q=V_rV_r^\top\), a rank-\(r\)
projector.

For a wide full-row-rank matrix, the roles reverse.

This is semi-orthogonality.

Higham's classical treatment supplies the numerical-analysis setting for
polar decomposition and its computation (Higham 1986).

\section{Polar orthogonalization is singular-value
flattening}\label{polar-orthogonalization-is-singular-value-flattening}

Using the compact rank-\(r\) SVD, the polar factor is

\[
Q=U_rI_rV_r^\top.
\]

Every nonzero singular value is replaced by \(1\).

The nonzero singular subspaces remain; zero singular directions remain
zero. For example, \(G=\operatorname{diag}(1,0)\) has compact polar
factor \(Q=\operatorname{diag}(1,0)\), not \(I_2\).

That is exact singular-value flattening on the nonzero support.

It is a strong and specific spectral transformation.

It does not by itself prove improved conditioning, training stability,
convergence, or generalization for a learning problem.

\section{``Steepest'' depends on the
norm}\label{steepest-depends-on-the-norm}

Locally,

\[
f(W+\Delta)
\approx
f(W)+\langle G,\Delta\rangle_F.
\]

For \(G\ne0\) and \(\rho>0\), if equally large moves are defined by

\[
\|\Delta\|_F\le\rho,
\]

then a steepest direction is

\[
\Delta_F^\star
=
-\rho\frac{G}{\|G\|_F}.
\]

If equally large moves are instead defined by

\[
\|\Delta\|_2\le\rho,
\]

then one steepest direction is

\[
\Delta_2^\star
=
-\rho U_rV_r^\top.
\]

If \(G=0\), the linearized objective is constant on both norm balls;
every feasible displacement minimizes it, and choosing \(\Delta=0\) is
valid. In particular, do not divide by \(\|G\|_F=0\).

The polar factor therefore appears naturally as a spectral-norm steepest
direction for nonzero \(G\). Bernstein and Newhouse develop this
norm-dependent optimizer viewpoint (Bernstein and Newhouse 2024).

The two directions do not contradict one another. They minimize the same
linear functional over different feasible sets.

\section{Why their local decreases are not optimizer
scores}\label{why-their-local-decreases-are-not-optimizer-scores}

For the Frobenius ball,

\[
\min_{\|\Delta\|_F\le\rho}
\langle G,\Delta\rangle_F
=
-\rho\|G\|_F.
\]

For the spectral ball,

\[
\min_{\|\Delta\|_2\le\rho}
\langle G,\Delta\rangle_F
=
-\rho\|G\|_*.
\]

A larger magnitude in the second expression does not prove a better
optimizer. The constraint defining an equally large step has changed.

The norm is part of the optimization problem.

\section{A one-step inverse-root
identity}\label{a-one-step-inverse-root-identity}

On the nonzero singular support,

\[
(GG^\top)^{-1/4}
G
(G^\top G)^{-1/4}
=
UV^\top.
\]

This explains an algebraic meeting point between a two-sided
inverse-root transform and the polar factor.

But two warnings are essential.

First, a rectangular Gram matrix on the larger side is singular, so a
full-space inverse needs damping or another declared convention.

Second, Shampoo is stateful. Its \(L_t\) and \(R_t\) accumulate earlier
gradients.

The identity does not say that stateful Shampoo is ``take the polar
factor of the current gradient.''

\section{Muon adds state and numerical
approximation}\label{muon-adds-state-and-numerical-approximation}

Muon is a concrete optimizer design for matrix-shaped hidden-layer
parameters (Jordan et al. 2024).

At a high level, it combines momentum with a finite Newton-Schulz-like
matrix iteration intended to produce an approximately orthogonalized
update.

The momentum state means the transformed matrix is not simply the raw
current gradient.

And a finite polynomial iteration is not the exact SVD or polar
decomposition.

A finite iteration implements a finite singular-value map.

The exact numerical realization therefore matters.

\section{Exact versus approximate matrix
functions}\label{exact-versus-approximate-matrix-functions}

Exact expressions such as

\[
(G^\top G)^{-1/2}
\]

define mathematical targets.

Large-scale optimizers may realize related transforms through
eigendecomposition, iterative matrix roots, Newton-Schulz iterations,
low-precision polynomials, blocking, distribution, or delayed refresh.

The realized update can differ from the exact target.

An optimizer is therefore not specified by the ideal matrix function
alone. Its state, damping, numerical method, precision, and refresh
schedule are also part of the algorithm.

\section{\texorpdfstring{The exact \(3\times2\)
witness}{The exact 3\textbackslash times2 witness}}\label{the-exact-3times2-witness}

Take

\[
G=
\begin{pmatrix}
2&0\\
0&1\\
1&0
\end{pmatrix}.
\]

Then

\[
G^\top G
=
\begin{pmatrix}
5&0\\
0&1
\end{pmatrix},
\]

so the singular values are

\[
\sqrt5,\ 1.
\]

The left Gram matrix is

\[
GG^\top
=
\begin{pmatrix}
4&0&2\\
0&1&0\\
2&0&1
\end{pmatrix}.
\]

It has rank \(2\), not \(3\).

The right one-step statistic is invertible while the left one is
singular.

The aspect ratio is already visible in the algebra.

\section{Polarizing the witness}\label{polarizing-the-witness}

The exact polar factor is

\[
Q=
\begin{pmatrix}
2/\sqrt5&0\\
0&1\\
1/\sqrt5&0
\end{pmatrix}.
\]

Then

\[
Q^\top Q=I_2.
\]

But

\[
QQ^\top
=
\begin{pmatrix}
4/5&0&2/5\\
0&1&0\\
2/5&0&1/5
\end{pmatrix},
\]

which is a rank-two projector rather than \(I_3\).

The witness is deliberately rectangular so ``orthogonalized'' cannot
silently mean ``square orthogonal.''

\section{A diagonal control}\label{a-diagonal-control}

Apply

\[
D=\operatorname{diag}(1/2,1,1/2).
\]

Then

\[
DG=
\begin{pmatrix}
1&0\\
0&1\\
1/2&0
\end{pmatrix}.
\]

Its singular values are

\[
\sqrt5/2,\ 1.
\]

The original singular-value ratio was \(\sqrt5\). The diagonal control
changes it to \(\sqrt5/2\). The polar factor makes it \(1\).

Both operations reshape the spectrum, but they do so differently.

This control is generic. It is not presented as the exact rule used by
Adam or another named elementwise optimizer.

\section{Flattening is not the whole spectral
story}\label{flattening-is-not-the-whole-spectral-story}

The polar factor gives the simplest possible nonzero singular spectrum:

\[
1,1,\ldots,1.
\]

But an optimizer might instead want to suppress only large modes, temper
near-zero modes, preserve relative scale over a band, or impose another
target spectrum.

Those questions belong to ATLAS-CH-SPECTRALSHAPE-001.

This chapter supplies the boundary:

\[
\boxed{
\text{singular-value flattening}
\neq
\text{arbitrary intentional spectral shaping}.
}
\]

\section{Matrix-aware does not mean
manifold-constrained}\label{matrix-aware-does-not-mean-manifold-constrained}

An optimizer can choose a semi-orthogonal update \(Q_t\) while leaving
the parameter \(W_t\) unconstrained:

\[
W_{t+1}=W_t-\eta Q_t.
\]

Even if

\[
Q_t^\top Q_t=I,
\]

there is no general implication that

\[
W_{t+1}^\top W_{t+1}=I.
\]

Manifold optimization instead declares a parameter constraint and uses
tangent/retraction machinery to maintain it.

Orthogonalized updates and orthogonality-constrained parameters are
different objects.

\section{Matrix-aware does not mean full second
order}\label{matrix-aware-does-not-mean-full-second-order}

A matrix preconditioner can encode more structure than an elementwise
rule without becoming an exact Newton method.

A full Hessian for a matrix block acts on the vectorized parameter space
and may couple every scalar coordinate.

Left-right factors impose a structured form.

That structure may be computationally attractive. It is still only the
structure actually specified.

The Atlas therefore keeps separate:

\begin{itemize}
\tightlist
\item
  matrix-aware preconditioning;
\item
  second-order curvature;
\item
  hard manifold constraints.
\end{itemize}

\section{Optimizer state is part of the
trajectory}\label{optimizer-state-is-part-of-the-trajectory}

An instantaneous map

\[
G_t\mapsto T(G_t)
\]

does not describe a stateful optimizer.

A stateful optimizer has a law such as

\[
S_t=\Phi(S_{t-1},G_t),
\]

\[
\Delta_t=\Psi(S_t,G_t).
\]

Shampoo's accumulated factors and Muon's momentum make this explicit.

Two methods can share a similar instantaneous transform and still
produce different trajectories because their states differ.

\section{What this chapter
establishes}\label{what-this-chapter-establishes}

The chapter establishes a controlled vocabulary and exact finite algebra
for matrix-aware updates.

It supports these claims:

\begin{itemize}
\tightlist
\item
  rectangular geometry matters;
\item
  left and right Gram matrices encode different sides of a matrix block;
\item
  a rectangular polar factor is semi-orthogonal;
\item
  polar orthogonalization flattens nonzero singular values;
\item
  norm choice changes the steepest direction;
\item
  a support-restricted one-step inverse-root transform meets the polar
  factor algebraically;
\item
  state, damping, and numerical realization prevent that identity from
  collapsing named optimizers into one another.
\end{itemize}

It does not establish that one update geometry is universally best.

\section{Failure boundaries}\label{failure-boundaries}

Keep these separations intact:

\[
\text{elementwise scaling}
\neq
\text{matrix preconditioning},
\]

\[
\text{matrix preconditioning}
\neq
\text{full Hessian inversion},
\]

\[
\text{one-step inverse-root identity}
\neq
\text{stateful Shampoo equivalence},
\]

\[
\text{exact polar factor}
\neq
\text{finite Newton-Schulz realization},
\]

\[
\text{orthogonalized update}
\neq
\text{orthogonality-constrained parameter},
\]

\[
\text{singular-value flattening}
\neq
\text{all spectral shaping},
\]

and

\[
\text{instantaneous transform}
\neq
\text{stateful optimizer dynamics}.
\]

\section{References used in this
chapter}\label{references-used-in-this-chapter-29}

\begin{itemize}
\tightlist
\item
  Shampoo: Gupta, Koren, and Singer (Gupta et al. 2018).
\item
  Polar decomposition and numerical computation: Higham (Higham 1986).
\item
  Norm-dependent steepest-descent interpretation: Bernstein and Newhouse
  (Bernstein and Newhouse 2024).
\item
  Muon design and finite Newton-Schulz update construction: Jordan et
  al. (Jordan et al. 2024).
\item
  Audited Atlas prerequisites: ATLAS-CH-LINALG-001 and
  ATLAS-CH-SECOND-001.
\end{itemize}

Exact bibliographic authority, prerequisite blob identities, and claim
boundaries are recorded in:

sources/source-locks/ATLAS-CH-MATRIXOPT-001.yaml

\chapter{Optimization on Manifolds}\label{optimization-on-manifolds}

\textbf{Epistemic status:} audited Geometry + audited First-Order
Optimization prerequisites + standard matrix-manifold optimization
sources + Atlas derivation + exact finite witness.\\
\textbf{Specification:}
manuscript/specifications/ATLAS-CH-MANOPT-001.md\\
\textbf{Derivation packet:}
mathematics/derivations/ATLAS-CH-MANOPT-001-DERIVATIONS.md\\
\textbf{Computational witness:}
mathematics/computational-witnesses/ATLAS-CW-MANOPT-001.md\\
\textbf{Source lock:} sources/source-locks/ATLAS-CH-MANOPT-001.yaml

\section{Optimization changes when legal states form a curved
space}\label{optimization-changes-when-legal-states-form-a-curved-space}

Ordinary first-order optimization usually assumes parameters may move
freely in Euclidean coordinates.

Many useful objects do not have that freedom.

Examples include:

\begin{itemize}
\tightlist
\item
  unit vectors;
\item
  orthonormal frames;
\item
  rotation matrices;
\item
  fixed-rank factorizations;
\item
  positive-definite matrices;
\item
  subspaces.
\end{itemize}

If the parameter must satisfy

\[
\|x\|=1
\]

or

\[
X^\top X=I,
\]

then an arbitrary Euclidean update generally leaves the admissible set.

Optimization on manifolds changes the question from:

\begin{quote}
Which ambient direction decreases the objective?
\end{quote}

to:

\begin{quote}
Which locally legal direction decreases the objective, and how do we
turn it into a legal finite move?
\end{quote}

That separation is the chapter's organizing idea.

\section{Tangent directions are local legal
velocities}\label{tangent-directions-are-local-legal-velocities}

For a smooth manifold \(M\), the tangent space

\[
T_xM
\]

is a linear approximation to the legal directions at \(x\).

A tangent vector is not usually another point on the manifold.

It is a velocity.

This matters because the update

\[
x+\xi
\]

can leave \(M\) even when

\[
\xi\in T_xM.
\]

A constrained optimizer therefore needs two distinct operations:

\begin{enumerate}
\def\labelenumi{\arabic{enumi}.}
\tightlist
\item
  choose a tangent direction;
\item
  map the finite displacement back to the manifold.
\end{enumerate}

The first is local differential geometry.

The second is finite constrained motion.

\section{The gradient depends on the
metric}\label{the-gradient-depends-on-the-metric}

The First-Order Optimization chapter treated the Euclidean gradient as
the steepest direction under the Euclidean metric.

A manifold may have a different metric.

So the Riemannian gradient is defined by

\[
Df(x)[\xi]
=
\langle \operatorname{grad}f(x),\xi\rangle_x
\]

for every tangent vector

\[
\xi\in T_xM.
\]

The metric appears inside the definition.

Therefore:

\[
\boxed{
\text{gradient}
=
\text{objective derivative + metric}.
}
\]

Changing the metric can change the steepest direction even when the
objective function is unchanged.

\section{Embedded manifolds and tangent
projection}\label{embedded-manifolds-and-tangent-projection}

For an embedded manifold with the induced Euclidean metric, the
Riemannian gradient is obtained by projecting the ambient Euclidean
gradient onto the tangent space (Absil et al. 2008).

Write

\[
\Pi_x
\]

for the orthogonal tangent projection.

Then

\[
\operatorname{grad}f(x)
=
\Pi_x(\nabla F(x)),
\]

where \(F\) is an ambient extension of \(f\).

This formula is powerful and easy to misuse.

It is valid for the induced metric setting being described.

It should not be promoted into a universal formula for every Riemannian
metric.

\section{Sphere geometry}\label{sphere-geometry}

For the unit sphere

\[
S^{n-1}
=
\{x\in\mathbb R^n:\|x\|=1\},
\]

the tangent space is

\[
T_xS^{n-1}
=
\{\xi:x^\top\xi=0\}.
\]

The induced-metric tangent projection is

\[
\Pi_x(g)
=
g-(x^\top g)x.
\]

The normal component lies along \(x\).

The tangent component is what first-order constrained motion may use.

\section{Tangent projection is not the
update}\label{tangent-projection-is-not-the-update}

Suppose

\[
g=\nabla F(x).
\]

A projected direction

\[
-\Pi_x(g)
\]

is tangent.

But the raw finite step

\[
x-\eta\Pi_x(g)
\]

need not satisfy the original nonlinear constraint.

The tangent plane touches the manifold.

It is not the manifold.

This gives a basic boundary:

\[
\boxed{
\text{legal infinitesimal direction}
\neq
\text{legal finite endpoint}.
}
\]

\section{Retractions}\label{retractions}

A retraction is a practical map from tangent data back to the manifold
(Absil et al. 2008).

Locally it satisfies:

\[
R_x(0_x)=x
\]

and

\[
DR_x(0_x)
=
\operatorname{id}_{T_xM}.
\]

So the retraction agrees with tangent motion to first order.

It does not have to trace an exact geodesic.

This matters computationally because exact exponential maps may be
expensive while useful retractions can be cheaper.

\section{Sphere normalization as a
retraction}\label{sphere-normalization-as-a-retraction}

For tangent \(\xi\) on the unit sphere, define:

\[
R_x(\xi)
=
\frac{x+\xi}{\|x+\xi\|}.
\]

Because

\[
x^\top\xi=0,
\]

we have

\[
\|x+\xi\|^2
=
1+\|\xi\|^2.
\]

Therefore

\[
\|R_x(\xi)\|=1.
\]

The normalized endpoint is exactly feasible.

The move is not generally the sphere exponential map.

\section{Exact sphere witness}\label{exact-sphere-witness}

Take

\[
x=(1,0)^\top,
\qquad
a=(1,2)^\top,
\qquad
f(z)=a^\top z.
\]

The ambient gradient is

\[
a=(1,2)^\top.
\]

Projecting onto the tangent space at \(x\) gives

\[
\operatorname{grad}f(x)
=
(0,2)^\top.
\]

Choose

\[
\eta=\frac12.
\]

The tangent descent vector is

\[
\xi=(0,-1)^\top.
\]

The unconstrained Euclidean update is

\[
(1/2,-1)^\top,
\]

whose squared norm is

\[
5/4.
\]

It leaves the sphere.

The raw tangent displacement

\[
x+\xi=(1,-1)^\top
\]

also leaves the sphere because its squared norm is \(2\).

After normalization:

\[
R_x(\xi)
=
\frac1{\sqrt2}(1,-1)^\top,
\]

and the norm is exactly \(1\).

The witness separates three things:

\begin{itemize}
\tightlist
\item
  ambient Euclidean descent;
\item
  tangent descent;
\item
  feasible finite motion.
\end{itemize}

\section{Retraction versus exponential
map}\label{retraction-versus-exponential-map}

For the same sphere witness,

\[
\|\xi\|=1.
\]

The exact exponential endpoint is

\[
\operatorname{Exp}_x(\xi)
=
(\cos1,-\sin1)^\top.
\]

The normalized retraction endpoint is

\[
(1/\sqrt2,-1/\sqrt2)^\top.
\]

Both are legal.

They differ.

Therefore:

\[
\boxed{
\text{retraction}
\neq
\text{exponential map}.
}
\]

A retraction earns its role through local first-order agreement and
feasibility, not by being an exact geodesic solver.

\section{The Stiefel manifold}\label{the-stiefel-manifold}

For integers \(1\le p\le n\), the nonempty Stiefel manifold is

\[
\operatorname{St}(n,p)
=
\{X\in\mathbb R^{n\times p}:X^\top X=I_p\}.
\]

If \(p>n\), an \(n\times p\) real matrix cannot have \(p\) orthonormal
columns; the tangent and retraction constructions below assume
\(1\le p\le n\).

Its columns are orthonormal.

This constraint appears throughout numerical linear algebra and
optimization with orthogonality constraints (Edelman et al. 1998).

A tangent matrix \(Z\) satisfies

\[
X^\top Z+Z^\top X=0.
\]

This is obtained by differentiating

\[
X^\top X=I.
\]

The condition is linear in \(Z\), as a tangent-space constraint should
be.

\section{Stiefel tangent
projection}\label{stiefel-tangent-projection}

Under the induced Euclidean metric, an ambient matrix \(G\) can be
projected by

\[
\Pi_X(G)
=
G-X\operatorname{sym}(X^\top G),
\]

where

\[
\operatorname{sym}(A)
=
\frac12(A+A^\top).
\]

The projected matrix satisfies

\[
X^\top\Pi_X(G)
+
\Pi_X(G)^\top X
=
0.
\]

Again, this produces a legal local direction.

It does not by itself produce a legal finite orthonormal matrix after
ambient addition.

\section{Polar retraction on the Stiefel
manifold}\label{polar-retraction-on-the-stiefel-manifold}

For tangent \(\Xi\), one retraction is

\[
R_X(\Xi)
=
(X+\Xi)
(I+\Xi^\top\Xi)^{-1/2}.
\]

The tangent condition implies

\[
(X+\Xi)^\top(X+\Xi)
=
I+\Xi^\top\Xi.
\]

Therefore:

\[
R_X(\Xi)^\top R_X(\Xi)=I.
\]

The retraction restores the orthogonality constraint exactly.

\section{\texorpdfstring{Exact \(2\times2\)
witness}{Exact 2\textbackslash times2 witness}}\label{exact-2times2-witness}

Let

\[
X=I_2
\]

and

\[
\Xi=
\begin{pmatrix}
0&-1\\
1&0
\end{pmatrix}.
\]

Because \(\Xi\) is skew-symmetric,

\[
X^\top\Xi+\Xi^\top X=0.
\]

So it is tangent.

The raw step is

\[
X+\Xi
=
\begin{pmatrix}
1&-1\\
1&1
\end{pmatrix}.
\]

Its Gram matrix is

\[
2I.
\]

Thus the raw tangent displacement is not orthogonal.

The polar retraction is

\[
\frac1{\sqrt2}
\begin{pmatrix}
1&-1\\
1&1
\end{pmatrix},
\]

whose Gram matrix is exactly \(I\).

Constraint restoration is exact.

\section{Orthogonal does not mean geodesically
exact}\label{orthogonal-does-not-mean-geodesically-exact}

For the same skew generator,

\[
e^\Xi
=
\begin{pmatrix}
\cos1&-\sin1\\
\sin1&\cos1
\end{pmatrix}.
\]

This is rotation by \(1\) radian.

The polar-retracted endpoint is rotation by

\[
\pi/4.
\]

Both are orthogonal.

They are different.

This is the matrix-manifold version of the same distinction:

\[
\text{constraint preservation}
\not\Rightarrow
\text{exponential/geodesic motion}.
\]

\section{A generic first-order manifold
step}\label{a-generic-first-order-manifold-step}

A simple constrained first-order method can be written:

\[
\xi_k
=
-\eta_k\operatorname{grad}f(x_k),
\]

followed by

\[
x_{k+1}
=
R_{x_k}(\xi_k).
\]

This compact notation hides important design choices:

\begin{itemize}
\tightlist
\item
  metric;
\item
  gradient estimator;
\item
  tangent projection;
\item
  step size;
\item
  retraction;
\item
  momentum;
\item
  vector transport;
\item
  stochastic sampling.
\end{itemize}

Optimization on a manifold is not one algorithm.

It is a family of algorithms respecting a declared geometric structure.

\section{Momentum creates a tangent-space
problem}\label{momentum-creates-a-tangent-space-problem}

Suppose an optimizer stores momentum

\[
m_k\in T_{x_k}M.
\]

After moving to

\[
x_{k+1},
\]

the new tangent space is

\[
T_{x_{k+1}}M.
\]

These two vector spaces are generally different subspaces of the ambient
space.

Blindly reusing \(m_k\) as if nothing changed can ignore geometry.

A vector transport supplies a declared way to move tangent information
between tangent spaces (Absil et al. 2008).

This is distinct from retraction:

\begin{itemize}
\tightlist
\item
  retraction moves a point;
\item
  vector transport moves tangent information.
\end{itemize}

\section{Intrinsic and extrinsic
viewpoints}\label{intrinsic-and-extrinsic-viewpoints}

A manifold can be described intrinsically through its metric and curves.

It can also be embedded in a larger Euclidean space and manipulated
through ambient coordinates.

Neither viewpoint is automatically superior.

The extrinsic viewpoint can make projections and constraints concrete.

The intrinsic viewpoint helps identify coordinate-independent structure.

The Atlas uses whichever representation makes the relevant claim
precise.

\section{Constraint preservation is not optimization
quality}\label{constraint-preservation-is-not-optimization-quality}

A method can preserve the constraint perfectly and still:

\begin{itemize}
\tightlist
\item
  increase the objective;
\item
  take poor step sizes;
\item
  converge slowly;
\item
  stagnate;
\item
  approach a bad stationary point.
\end{itemize}

Therefore:

\[
\boxed{
\text{constraint preservation}
\neq
\text{optimizer quality}.
}
\]

Geometry defines legal motion.

Optimization still requires a useful search rule.

\section{Stationarity on a
manifold}\label{stationarity-on-a-manifold}

A first-order stationary point satisfies

\[
\operatorname{grad}f(x)=0.
\]

For an embedded manifold under induced metric, this means the ambient
gradient has no tangent component.

The ambient gradient can still have a normal component.

Thus constrained stationarity differs from unconstrained Euclidean
stationarity.

And first-order stationarity still does not imply global optimality.

\section{Sphere versus Stiefel}\label{sphere-versus-stiefel}

The sphere constrains one vector:

\[
x^\top x=1.
\]

The Stiefel manifold constrains multiple columns simultaneously:

\[
X^\top X=I.
\]

The sphere tangent condition is scalar:

\[
x^\top\xi=0.
\]

The Stiefel tangent condition is matrix-valued:

\[
X^\top Z+Z^\top X=0.
\]

The two geometries are related but not interchangeable.

\section{Stiefel versus Grassmann}\label{stiefel-versus-grassmann}

A Stiefel point is an orthonormal frame.

A Grassmann point represents a subspace independent of which orthonormal
basis is chosen for it.

The Geometry chapter already established this distinction.

It matters for optimization.

If the objective depends only on the spanned subspace, then treating
different bases as distinct states may introduce redundant directions.

This is where quotient geometry enters.

MANOPT uses the distinction but does not redevelop quotient optimization
in full.

\section{Why parameterization
matters}\label{why-parameterization-matters}

One can sometimes avoid an explicit constrained optimizer by
parameterizing the constraint indirectly.

Examples include:

\begin{itemize}
\tightlist
\item
  normalized vectors;
\item
  matrix exponentials;
\item
  factorizations;
\item
  Householder products;
\item
  Cayley-like transforms.
\end{itemize}

Such choices change:

\begin{itemize}
\tightlist
\item
  coordinates;
\item
  conditioning;
\item
  computational cost;
\item
  reachable sets;
\item
  singularities;
\item
  optimization dynamics.
\end{itemize}

A constraint-preserving parameterization is not automatically equivalent
to direct manifold optimization.

The effective geometry can differ.

\section{Geometry can change steepest
descent}\label{geometry-can-change-steepest-descent}

In Euclidean space, steepest descent depends on the Euclidean inner
product.

On a manifold, the Riemannian metric changes the local notion of unit
direction.

Therefore changing geometry can change the gradient even before any
adaptive optimizer state is introduced.

This is one reason later Atlas work on geometry-aware and matrix-aware
optimization must keep separate:

\begin{itemize}
\tightlist
\item
  the objective;
\item
  the metric;
\item
  the optimizer;
\item
  the parameterization.
\end{itemize}

\section{Retraction choice can
matter}\label{retraction-choice-can-matter}

Two valid retractions agree to first order at the tangent origin.

They can differ at finite step sizes.

So a large-step optimization trajectory can depend on the retraction
even when the tangent gradient is unchanged.

This is not a defect.

It is part of the algorithm.

A convergence theorem must state what retraction properties it assumes.

An experiment must state which retraction it used.

\section{Step size still matters}\label{step-size-still-matters}

Geometry does not remove ordinary optimization concerns.

A constrained method still needs to choose:

\begin{itemize}
\tightlist
\item
  step size;
\item
  stochastic batch;
\item
  line search or schedule;
\item
  momentum state;
\item
  clipping or trust restrictions;
\item
  stopping rule.
\end{itemize}

A legal step can still be too large.

A legal trajectory can still be poor.

Manifold structure constrains the search space.

It does not solve the search problem by itself.

\section{Relation to normalized
models}\label{relation-to-normalized-models}

Unit-norm states and orthogonality constraints appear naturally in
normalized neural architectures.

The mathematical substrate in this chapter is directly relevant to such
systems.

But the chapter stops at generic manifold optimization.

It does not establish:

\begin{itemize}
\tightlist
\item
  a theorem for nGPT;
\item
  a theorem for Muon;
\item
  a theorem for rectangular spectral updates;
\item
  a theorem for MODULUS;
\item
  a theorem for any specific GCL optimizer.
\end{itemize}

Those require their own objectives, metrics, updates, and evidence.

\section{Relation to variational
optimization}\label{relation-to-variational-optimization}

The direct downstream consumer is:

\textbf{Variational and Divergence-Derived Optimization ---
ATLAS-CH-VARIOPT-001.}

VARIOPT may inherit:

\begin{itemize}
\tightlist
\item
  tangent gradients;
\item
  constrained finite motion;
\item
  retractions;
\item
  metric dependence;
\item
  vector-transport distinctions;
\item
  sphere/Stiefel examples.
\end{itemize}

It must independently justify any claim that an update follows from:

\begin{itemize}
\tightlist
\item
  a divergence;
\item
  a discrete Lagrangian;
\item
  a symplectic principle;
\item
  a duality relation;
\item
  another variational construction.
\end{itemize}

Manifold optimization says how to respect a constrained geometry.

Variational optimization asks where the motion rule itself comes from.

\section{Durable boundaries}\label{durable-boundaries}

The chapter leaves ten distinctions:

\begin{enumerate}
\def\labelenumi{\arabic{enumi}.}
\tightlist
\item
  Euclidean gradient is not automatically the Riemannian gradient.
\item
  Tangent projection is not a complete optimizer.
\item
  Tangent direction is not a finite feasible endpoint.
\item
  Retraction is not the exponential map.
\item
  First-order agreement is not geodesic exactness.
\item
  Sphere and Stiefel constraints are different.
\item
  Retraction and vector transport have different roles.
\item
  Constraint preservation is not convergence.
\item
  Manifold stationarity is not global optimality.
\item
  Generic manifold optimization is not a specific GCL optimizer theorem.
\end{enumerate}

These boundaries are the reusable substrate.

\section{References used in this
chapter}\label{references-used-in-this-chapter-30}

\begin{itemize}
\tightlist
\item
  (Absil et al. 2008)
\item
  (Edelman et al. 1998)
\end{itemize}

Exact source identities and authority scopes are pinned in:

sources/source-locks/ATLAS-CH-MANOPT-001.yaml

\chapter{Spectral Shaping}\label{spectral-shaping}

\textbf{Epistemic status:} audited Matrix-Aware Optimization and
Non-normality substrates + Atlas-owned exact finite spectral-shaping
witnesses.\\
\textbf{Specification:}
manuscript/specifications/ATLAS-CH-SPECTRALSHAPE-001.md\\
\textbf{Derivation packet:}
mathematics/derivations/ATLAS-CH-SPECTRALSHAPE-001-DERIVATIONS.md\\
\textbf{Computational witness:}
mathematics/computational-witnesses/ATLAS-CW-SPECTRALSHAPE-001.md\\
\textbf{Source lock:}
sources/source-locks/ATLAS-CH-SPECTRALSHAPE-001.yaml

Spectral manipulation is often discussed as though it were one
operation.

It is not.

A matrix update can be:

\begin{itemize}
\tightlist
\item
  normalized;
\item
  clipped;
\item
  conditioned;
\item
  flattened;
\item
  or deliberately reshaped toward a non-flat target.
\end{itemize}

Those operations can have different mathematical meanings even when they
all alter singular values.

The governing rule is:

\[
\boxed{
\text{declare the target spectral map before interpreting the result}.
}
\]

\section{The handoff from Matrix-Aware
Optimization}\label{the-handoff-from-matrix-aware-optimization}

MATRIXOPT already gives the Atlas:

\begin{itemize}
\tightlist
\item
  SVD language;
\item
  polar factors;
\item
  rectangular semi-orthogonality;
\item
  matrix preconditioning;
\item
  exact-versus-approximate singular-value transformations;
\item
  the key boundary that polar flattening is not all spectral shaping.
\end{itemize}

SPECTRALSHAPE starts where that chapter deliberately stopped.

\section{The handoff from
Non-normality}\label{the-handoff-from-non-normality}

NONNORMAL adds a second constraint.

Even if eigenvalues look benign, finite-horizon behavior can be very
different when the operator is non-normal.

Therefore an instantaneous update spectrum cannot be promoted directly
into a dynamical guarantee.

\section{A general shaping map}\label{a-general-shaping-map}

Let

\[
G=U\Sigma V^\top,
\]

with singular values

\[
\sigma_1,\dots,\sigma_r.
\]

A singular-value shaping rule has the form

\[
T_f(G)
=
U f(\Sigma)V^\top,
\]

where \(f\) acts on the singular values. When \(G\) is rank-deficient,
the map must specify its behavior on zero singular directions. Taking
\(f(0)=0\) leaves those directions zero and yields an SVD-independent
singular-value transform. If \(f(0)\ne0\), a rule also needs a declared
pairing of the left and right nullspaces; the zero-singular-vector
choices in an SVD of \(G\) are not unique.

As a counterexample, let \(G=\operatorname{diag}(1,0)\) and
\(f(1)=f(0)=1\). Two valid SVDs are \(U=V=I_2\) and
\(U=\operatorname{diag}(1,-1),V=I_2\), both with
\(\Sigma=\operatorname{diag}(1,0)\). They give different transformed
matrices, \(I_2\) and \(\operatorname{diag}(1,-1)\). Hence a
rank-increasing rule requires extra structure. The full-rank witnesses
below are unaffected.

Different well-defined \(f\) give different update geometries.

\section{Exact witness matrix}\label{exact-witness-matrix}

Use

\[
G=
\begin{pmatrix}
4&0\\
0&1
\end{pmatrix}.
\]

Its singular values are:

\[
(4,1).
\]

Hence:

\[
\boxed{
\kappa_2(G)=4.
}
\]

This one matrix is enough to distinguish four common operations.

\section{Scalar normalization changes
scale}\label{scalar-normalization-changes-scale}

Normalize by the operator norm:

\[
G_{\rm scale}
=
\frac{G}{\|G\|_2}
=
\begin{pmatrix}
1&0\\
0&1/4
\end{pmatrix}.
\]

The condition number remains:

\[
\boxed{
\kappa_2(G_{\rm scale})=4.
}
\]

So:

\[
\boxed{
\text{normalization}
\not\Rightarrow
\text{better conditioning}.
}
\]

A global scalar multiplies every singular value by the same amount.

Their ratio is unchanged.

\section{Clipping can change
conditioning}\label{clipping-can-change-conditioning}

Now cap singular values at \(2\):

\[
f_{\rm clip}(\sigma)
=
\min(\sigma,2).
\]

Then:

\[
G_{\rm clip}
=
\begin{pmatrix}
2&0\\
0&1
\end{pmatrix}.
\]

Thus:

\[
\boxed{
\kappa_2(G_{\rm clip})=2.
}
\]

The spectrum is still non-flat.

Only the large mode was clipped.

\section{Polar flattening is
stronger}\label{polar-flattening-is-stronger}

For a full-rank square positive diagonal matrix, the polar factor
replaces every nonzero singular value by one.

Thus:

\[
G_{\rm polar}=I_2.
\]

Hence:

\[
\boxed{
\kappa_2(G_{\rm polar})=1.
}
\]

This is complete nonzero singular-value flattening.

\section{A non-flat target is another
operation}\label{a-non-flat-target-is-another-operation}

Suppose we deliberately choose target singular values:

\[
(3,2).
\]

Then:

\[
G_\tau
=
\begin{pmatrix}
3&0\\
0&2
\end{pmatrix}.
\]

Its condition number is:

\[
\boxed{
\kappa_2(G_\tau)=\frac32.
}
\]

This is neither normalization, clipping, nor flattening.

It is explicit spectral design.

\section{The four maps are not
synonyms}\label{the-four-maps-are-not-synonyms}

For the same input \(G\), the shaped spectra are:

\[
(1,1/4),
\]

\[
(2,1),
\]

\[
(1,1),
\]

and

\[
(3,2).
\]

Therefore:

\[
\boxed{
\text{scale control}
\neq
\text{clipping}
\neq
\text{flattening}
\neq
\text{target shaping}.
}
\]

\section{Better condition number is not automatically better
optimization}\label{better-condition-number-is-not-automatically-better-optimization}

On this witness, the condition numbers are:

\[
4,\quad4,\quad2,\quad1,\quad3/2.
\]

That is an exact matrix statement.

It is not a universal optimizer ranking.

A nonlinear optimizer also depends on:

\begin{itemize}
\tightlist
\item
  objective geometry;
\item
  step size;
\item
  optimizer state;
\item
  stochasticity;
\item
  parameterization;
\item
  update interaction across steps.
\end{itemize}

\section{Flat is not automatically
best}\label{flat-is-not-automatically-best}

The polar spectrum:

\[
(1,1)
\]

is maximally flat.

But there may be reasons to preserve anisotropy.

A target:

\[
(3,2)
\]

retains directional preference.

The Atlas therefore treats the target spectrum as a design object rather
than an article of faith.

\section{Conditioning is only one
objective}\label{conditioning-is-only-one-objective}

Possible spectral goals include:

\begin{itemize}
\tightlist
\item
  cap large modes;
\item
  raise small modes;
\item
  reduce condition number;
\item
  preserve relative scale in a band;
\item
  remove scale entirely;
\item
  impose a desired anisotropy.
\end{itemize}

These objectives need not agree.

\section{Rank matters}\label{rank-matters}

If a singular value is zero, the standard full-space condition number is
infinite.

Changing a zero singular value to a positive one changes rank.

That is no longer mere rescaling.

It is a structural intervention and must be declared as such.

\section{Singular vectors matter
too}\label{singular-vectors-matter-too}

A map

\[
U f(\Sigma)V^\top
\]

changes singular values while retaining the singular-vector frames.

A more general transform can also rotate \(U\) or \(V\).

Two matrices with the same singular values can therefore still behave
differently relative to a particular interface.

Spectral shape is not the whole operator.

\section{Non-normality supplies the dynamical
warning}\label{non-normality-supplies-the-dynamical-warning}

Consider:

\[
A=
\begin{pmatrix}
1/2&2\\
0&1/2
\end{pmatrix},
\qquad
N=
\frac12I.
\]

Both have eigenvalues:

\[
\{1/2,1/2\}.
\]

If eigenvalues were the whole story, their short-time behavior would
look alike.

It does not.

\section{Same eigenvalues, different one-step
response}\label{same-eigenvalues-different-one-step-response}

Probe both matrices with:

\[
e_2=(0,1)^\top.
\]

Then:

\[
Ae_2=(2,1/2)^\top,
\]

so:

\[
\boxed{
\|Ae_2\|_2^2=\frac{17}{4}.
}
\]

But:

\[
Ne_2=(0,1/2)^\top,
\]

so:

\[
\boxed{
\|Ne_2\|_2^2=\frac14.
}
\]

Equal eigenvalues do not imply equal interface response.

\section{Pseudospectra separate them
further}\label{pseudospectra-separate-them-further}

For the upper-triangular non-normal family inherited from NONNORMAL, the
exact closed 2-norm \(\varepsilon\)-pseudospectral radius around the
repeated eigenvalue is:

\[
r_A
=
\sqrt{\varepsilon(\varepsilon+K)}.
\]

Set:

\[
K=2,
\qquad
\varepsilon=\frac14.
\]

Then:

\[
\boxed{
r_A=\frac34.
}
\]

For the normal scalar matrix:

\[
\boxed{
r_N=\frac14.
}
\]

Matched eigenvalues coexist with a threefold difference in this
pseudospectral radius.

\section{Why that matters for optimizer
updates}\label{why-that-matters-for-optimizer-updates}

Suppose an optimizer shapes the singular spectrum of an instantaneous
update matrix:

\[
\Delta_t.
\]

That says something about \(\Delta_t\).

It does not automatically say anything about the spectrum or
pseudospectrum of the optimizer's joint state transition.

Those are different operators.

\section{Stateful optimizers enlarge the
object}\label{stateful-optimizers-enlarge-the-object}

A stateful optimizer has:

\[
S_{t+1}
=
\Phi(S_t,G_t),
\]

and perhaps:

\[
W_{t+1}
=
W_t+\Psi(S_{t+1},G_t).
\]

The relevant local dynamical operator is then a Jacobian on the coupled
state:

\[
(W,S).
\]

The singular values of one instantaneous parameter update are not that
Jacobian.

\section{Spectral shaping is not optimizer-state
dynamics}\label{spectral-shaping-is-not-optimizer-state-dynamics}

Therefore:

\[
\boxed{
\text{instantaneous update spectrum}
\not\Rightarrow
\text{optimizer-state dynamics}.
}
\]

This boundary is mandatory.

It prevents a useful matrix diagnostic from becoming an unsupported
dynamical theorem.

\section{Better conditioning is not
convergence}\label{better-conditioning-is-not-convergence}

Even if a spectral map reduces:

\[
\kappa_2(G),
\]

the nonlinear training trajectory may improve, worsen, or remain
unchanged.

A convergence theorem needs the actual update law and assumptions.

So:

\[
\boxed{
\text{better update condition number}
\not\Rightarrow
\text{faster nonlinear convergence}.
}
\]

\section{Task quality is another
layer}\label{task-quality-is-another-layer}

A shaped update can have:

\begin{itemize}
\tightlist
\item
  a lower condition number;
\item
  a flatter spectrum;
\item
  smaller operator norm;
\end{itemize}

and still harm the downstream objective.

Conversely, task quality can improve without a flatter update spectrum.

Thus:

\[
\boxed{
\text{spectral-shape quality}
\not\Rightarrow
\text{task quality}.
}
\]

\section{A practical reporting
protocol}\label{a-practical-reporting-protocol}

For a spectral-shaping method, report separately:

\begin{enumerate}
\def\labelenumi{\arabic{enumi}.}
\tightlist
\item
  the matrix/operator being shaped;
\item
  the SVD or spectral object being used;
\item
  the shaping function \(f\);
\item
  the target spectrum or clipping rule;
\item
  condition number before/after;
\item
  rank changes, if any;
\item
  singular-vector changes, if any;
\item
  optimizer-state semantics;
\item
  dynamical diagnostics;
\item
  downstream task metrics.
\end{enumerate}

This prevents one spectral statistic from impersonating the entire
optimization process.

\section{What a flat spectrum actually
certifies}\label{what-a-flat-spectrum-actually-certifies}

For a full-rank square matrix, flattening the nonzero singular values to
one certifies:

\[
\kappa_2=1.
\]

It does not certify:

\begin{itemize}
\tightlist
\item
  minimum training loss;
\item
  good generalization;
\item
  stable coupled dynamics;
\item
  absence of non-normality elsewhere;
\item
  correct task behavior.
\end{itemize}

The certificate is exactly the matrix property it proves.

\section{What a non-flat target actually
certifies}\label{what-a-non-flat-target-actually-certifies}

A target such as:

\[
(3,2)
\]

certifies only that the chosen singular-value profile has been imposed.

Whether this geometry is useful is an empirical or theoretical question
about the downstream problem.

\section{Atlas non-implications}\label{atlas-non-implications-1}

SPECTRALSHAPE rejects:

\[
\text{normalization}
\not\Rightarrow
\text{better conditioning},
\]

\[
\text{conditioning}
\not\Rightarrow
\text{flattening},
\]

\[
\text{flat spectrum}
\not\Rightarrow
\text{universal optimizer quality},
\]

\[
\text{better update conditioning}
\not\Rightarrow
\text{faster nonlinear convergence},
\]

\[
\text{instantaneous update spectrum}
\not\Rightarrow
\text{optimizer-state dynamics},
\]

\[
\text{eigenvalue shape}
\not\Rightarrow
\text{non-normal transient control},
\]

and:

\[
\text{spectral-shape quality}
\not\Rightarrow
\text{task quality}.
\]

\section{Atlas connections}\label{atlas-connections-14}

\textbf{Matrix-Aware Optimization.}\\
Supplies SVD/polar update geometry and the exact flattening boundary.

\textbf{Normality, Pseudospectra, and Transient Growth.}\\
Supplies the warning that spectral eigenvalue information may be
insufficient for finite-horizon dynamics.

\textbf{Optimizer-State Dynamics.}\\
Owns the coupled model--optimizer state and its Jacobian.

\textbf{Spectral and Operator Diagnostics.}\\
Can diagnose operator behavior but does not convert a spectral signature
into mechanistic evidence.

\section{Closing view}\label{closing-view-21}

Spectral shaping is best understood as an explicit design map:

\[
\boxed{
G
\mapsto
U f(\Sigma)V^\top.
}
\]

The important question is not:

\begin{quote}
did the spectrum change?
\end{quote}

It is:

\begin{quote}
which spectrum, by what map, toward what target, for what objective, and
with what separate evidence about dynamics and task behavior?
\end{quote}

The exact witness gives four different answers from one input matrix.

That is the point.

\section{References used in this
chapter}\label{references-used-in-this-chapter-31}

No new external academic authority is added.

Matrix-update and SVD/polar authority is inherited through audited
ATLAS-CH-MATRIXOPT-001.

Non-normal transient-growth and pseudospectral authority is inherited
through audited ATLAS-CH-NONNORMAL-001.

Exact prerequisite identities and claim boundaries are recorded in:

sources/source-locks/ATLAS-CH-SPECTRALSHAPE-001.yaml

\chapter{Optimizer-State Dynamics}\label{optimizer-state-dynamics}

\textbf{Epistemic status:} mathematical exposition built from canonical
optimizer formulations, control-theoretic optimization analysis, the
preceding Non-normality chapter, and Atlas-owned augmented-state
derivations.\\
\textbf{Primary figure:} `ATLAS-FIG-OPTDYN-001`\\
\textbf{Derivation packet:}
`mathematics/derivations/ATLAS-CH-OPTDYN-001-DERIVATIONS.md`

\section{The optimizer is part of the
state}\label{the-optimizer-is-part-of-the-state}

A common picture of training writes

\[
\theta_{t+1}
=
\theta_t-\eta\nabla L(\theta_t).
\]

That picture is exact for plain gradient descent.

For momentum, Adam, and many modern optimizers it is incomplete.

The update at time \(t+1\) depends not only on the current parameters
but also on internal state accumulated from previous gradients. Momentum
carries velocity-like state. Adam carries first- and second-moment
estimates (Kingma and Ba 2015). Other methods maintain matrices,
preconditioners, or moving statistics.

The mathematical object is therefore not only \(\theta_t\).

It is an augmented state,

\[
z_t=(\theta_t,s_t),
\]

with update

\[
z_{t+1}
=
F_t(z_t,\xi_t),
\]

where \(s_t\) is optimizer state and \(\xi_t\) represents data or
stochastic forcing.

This chapter asks:

\begin{quote}
what dynamical behavior becomes visible only when the optimizer and
model are analyzed as one coupled system?
\end{quote}

The useful allegory is a \textbf{vehicle with a flywheel}.

Position corresponds to model parameters. The flywheel corresponds to
stored optimizer state. A steering command changes position, but the
stored motion changes how the same command behaves one step later.

The limit of the allegory is strict. Optimizer state is not physical
momentum, adaptive methods have coordinatewise and state-dependent
statistics, and no conservation law is implied.

The point is only that memory in the update rule changes the state
space.

\section{Heavy-ball momentum as a state-space
map}\label{heavy-ball-momentum-as-a-state-space-map}

Polyak's multistep acceleration is a canonical source for momentum-like
iteration (Polyak 1964), and momentum later became central to
deep-network optimization (Sutskever et al. 2013).

Consider the scalar quadratic

\[
L(\theta)
=
\frac h2\theta^2,
\qquad
h>0.
\]

Then

\[
g(\theta)=h\theta.
\]

Use the recurrence

\[
v_{t+1}
=
\beta v_t+h\theta_t,
\]

followed by

\[
\theta_{t+1}
=
\theta_t-\eta v_{t+1}.
\]

Substitution gives

\[
\theta_{t+1}
=
(1-\eta h)\theta_t-\eta\beta v_t.
\]

Define

\[
z_t=
\begin{pmatrix}
\theta_t\\
v_t
\end{pmatrix}.
\]

Then

\[
z_{t+1}
=
Jz_t,
\]

with

\[
\boxed{
J=
\begin{pmatrix}
1-\eta h&-\eta\beta\\
h&\beta
\end{pmatrix}.
}
\]

Nothing has been approximated. The optimizer-model pair is literally a
two-dimensional linear dynamical system for this quadratic example.

\section{The characteristic equation belongs to the coupled
system}\label{the-characteristic-equation-belongs-to-the-coupled-system}

The trace is

\[
\operatorname{tr}J
=
1-\eta h+\beta,
\]

and

\[
\det J=\beta.
\]

Therefore

\[
\boxed{
p(\lambda)
=
\lambda^2-(1-\eta h+\beta)\lambda+\beta.
}
\]

The stability boundary depends jointly on

\[
\eta,\qquad
\beta,\qquad
h.
\]

This already changes the conceptual story.

Curvature \(h\) matters, but curvature alone does not determine the
update. Learning rate matters, but learning rate alone is not the
dynamical system. Momentum matters, but it is not a scalar correction
applied outside the dynamics.

The boundary is a property of the coupled recurrence.

\section{A stable example with transient
amplification}\label{a-stable-example-with-transient-amplification}

Choose

\[
h=1,\qquad
\eta=\frac1{10},\qquad
\beta=\frac9{10}.
\]

Then

\[
J=
\begin{pmatrix}
0.9&-0.09\\
1&0.9
\end{pmatrix}.
\]

Its eigenvalues are

\[
\lambda_\pm
=
0.9\pm0.3i.
\]

Thus

\[
|\lambda_\pm|
=
\sqrt{0.9}
\approx0.9486832981<1.
\]

The fixed point is asymptotically stable.

If eigenvalues were the whole story, one might expect every perturbation
simply to decay.

They are not the whole story.

The matrix is non-normal.

Indeed,

\[
J^\top J
=
\begin{pmatrix}
1.81&0.819\\
0.819&0.8181
\end{pmatrix},
\]

while

\[
JJ^\top
=
\begin{pmatrix}
0.8181&0.819\\
0.819&1.81
\end{pmatrix}.
\]

Therefore

\[
J^\top J\neq JJ^\top.
\]

The preceding Non-normality chapter prepared exactly this situation:
asymptotic spectral stability can coexist with finite-horizon
amplification.

\section{\texorpdfstring{Four stable steps can amplify by more than
\(2.6\times\)}{Four stable steps can amplify by more than 2.6\textbackslash times}}\label{four-stable-steps-can-amplify-by-more-than-2.6times}

For the same matrix,

\[
J^4
=
\begin{pmatrix}
\frac{567}{2500}&-\frac{729}{3125}\\
\frac{324}{125}&\frac{567}{2500}
\end{pmatrix}.
\]

The largest singular value is

\[
\|J^4\|_2
\approx
2.610090585959495.
\]

So there exists a unit perturbation \(\delta z_0\) such that

\[
\|J^4\delta z_0\|_2
\approx
2.61.
\]

The same system eventually decays because its eigenvalues are inside the
unit disk.

This is the mechanism to remember:

\[
\boxed{
\text{asymptotically stable}
\not\Rightarrow
\text{monotone finite-horizon decay}.
}
\]

The claim is exact for this linear example.

It is not yet a claim about a neural training run.

\section{The Wolfram plate}\label{the-wolfram-plate}

The primary figure puts three views of the same toy system beside one
another.

\begin{figure}
\centering
\pandocbounded{\includegraphics[keepaspectratio,alt={Phase trajectories, eigenvalues relative to the unit circle, and finite-horizon spectral norm gain for the exact momentum-state matrix.}]{figures/masters/ATLAS-FIG-OPTDYN-001.png}}
\caption{Phase trajectories, eigenvalues relative to the unit circle,
and finite-horizon spectral norm gain for the exact momentum-state
matrix.}
\end{figure}

The left panel shows trajectories in the augmented \((\theta,v)\) state
plane.

The middle panel shows both eigenvalues strictly inside the unit circle.

The right panel shows

\[
\|J^n\|_2
\]

for \(n=0,\ldots,20\), with a peak of approximately

\[
2.61009
\]

at

\[
n=4.
\]

The three panels are not three different stories. They are three
descriptions of one matrix.

The figure is therefore a compact warning against using one spectral
summary as if it exhausted the dynamics.

\section{Why the augmented state
matters}\label{why-the-augmented-state-matters}

If one looks only at \(\theta_t\), the optimizer's memory appears
indirectly.

In augmented state, it is explicit.

This matters because the local Jacobian of the full update contains
couplings between parameter and optimizer coordinates.

For a general optimizer state \(s_t\),

\[
z_t=
\begin{pmatrix}
\theta_t\\
s_t
\end{pmatrix},
\]

and a local Jacobian has block structure

\[
J_t
=
\begin{pmatrix}
\partial_\theta F_\theta&
\partial_s F_\theta\\
\partial_\theta F_s&
\partial_s F_s
\end{pmatrix}.
\]

The off-diagonal blocks are not decoration. They encode how optimizer
memory changes parameter motion and how parameter-space gradients update
optimizer memory.

A parameter-only Hessian does not contain all of this information.

\section{Curvature is not the same as optimizer
dynamics}\label{curvature-is-not-the-same-as-optimizer-dynamics}

On the quadratic example, curvature appears through \(h\).

But the state Jacobian is

\[
J(\eta,\beta,h)
=
\begin{pmatrix}
1-\eta h&-\eta\beta\\
h&\beta
\end{pmatrix}.
\]

The same \(h\) with different \(\eta\) or \(\beta\) produces a different
dynamical system.

Likewise, a large Hessian eigenvalue can matter without being a complete
explanation of instability.

The disciplined statement is:

\begin{quote}
the loss landscape contributes to the update dynamics, but the update
dynamics belong to the combined model-optimizer recurrence.
\end{quote}

This is the level at which learning-rate boundaries should be analyzed.

\section{From an autonomous toy system to time-varying
training}\label{from-an-autonomous-toy-system-to-time-varying-training}

The quadratic example has one fixed matrix \(J\).

Training in a neural network is generally nonautonomous.

Write

\[
z_{t+1}
=
F_t(z_t,\xi_t).
\]

Linearizing along a trajectory gives

\[
\delta z_{t+1}
\approx
J_t\delta z_t,
\qquad
J_t
=
\frac{\partial F_t}{\partial z}(z_t,\xi_t).
\]

After \(k\) steps,

\[
\delta z_{t+k}
\approx
J_{t+k-1}\cdots J_t\,\delta z_t.
\]

A single-step eigenvalue calculation cannot in general summarize this
product.

The relevant object is the finite-horizon propagator.

This is one reason non-normal and transient-growth thinking is useful:
it asks what perturbations can do over a finite horizon, not only what
one frozen matrix does asymptotically.

\section{Adam makes the state
richer}\label{adam-makes-the-state-richer}

Adam maintains first- and second-moment estimates (Kingma and Ba 2015).

Schematically,

\[
m_{t+1}
=
\beta_1m_t+(1-\beta_1)g_t,
\]

\[
v_{t+1}
=
\beta_2v_t+(1-\beta_2)g_t^2,
\]

followed by a parameter update using bias-corrected moment estimates.

The exact augmented state is therefore larger:

\[
z_t=(\theta_t,m_t,v_t,\ldots).
\]

The present chapter does not derive a universal Adam stability theory.
That would obscure the mechanism we are trying to expose.

Momentum is used as the exact microscope because its state is small
enough that every coupling can be seen.

Adam is included to show that the augmented-state viewpoint becomes
more, not less, relevant as optimizers accumulate richer memory.

\section{Optimization as a dynamical and control
system}\label{optimization-as-a-dynamical-and-control-system}

There is established precedent for analyzing iterative optimization as a
feedback/dynamical system. Lessard, Recht, and Packard use integral
quadratic constraints from robust control to analyze methods including
gradient descent, heavy-ball momentum, and accelerated schemes (Lessard
et al. 2016).

The Atlas does not reproduce that full framework here.

The connection is conceptual and mathematical:

\begin{itemize}
\tightlist
\item
  an optimizer is an iterative dynamical system;
\item
  the objective/gradient map enters the loop;
\item
  stability and gain are properties of the interconnection;
\item
  finite-horizon behavior can matter even when asymptotic convergence
  conditions hold.
\end{itemize}

The Atlas adds a particular emphasis: in learned systems, optimizer
state should be treated as part of the object being diagnosed.

\section{What a local Jacobian can and cannot tell
us}\label{what-a-local-jacobian-can-and-cannot-tell-us}

A local augmented-state Jacobian can reveal:

\begin{itemize}
\tightlist
\item
  instantaneous coupling;
\item
  singular amplification directions;
\item
  local spectral structure;
\item
  departure from normality;
\item
  sensitivity to hyperparameters;
\item
  finite-horizon products along a recorded trajectory.
\end{itemize}

It cannot, by itself, establish:

\begin{itemize}
\tightlist
\item
  the global nonlinear behavior of training;
\item
  the cause of an observed loss spike;
\item
  that one local mode persists over many steps;
\item
  that a toy quadratic captures a frontier model;
\item
  that non-normality is the dominant mechanism in every unstable run.
\end{itemize}

The distinction between mechanism and prevalence is essential.

A toy example can prove that a mechanism exists.

It cannot prove how often that mechanism matters in practice.

\section{The flywheel allegory, and where it
fails}\label{the-flywheel-allegory-and-where-it-fails}

The flywheel picture is useful because it blocks one mistake: pretending
that only current position matters.

The correspondence is:

\begin{itemize}
\tightlist
\item
  position ↔ model parameters;
\item
  stored rotational state ↔ optimizer memory;
\item
  control gain ↔ learning rate;
\item
  coupled overshoot ↔ transient amplification.
\end{itemize}

But the optimizer is not a physical rigid body.

Adam's state has no literal angular momentum. Stochastic gradients are
not deterministic forces. Coordinatewise normalization has no simple
mechanical analogue.

The allegory has done its job once the reader accepts that optimizer
memory enlarges the state space.

The mathematics should then take over.

\section{A handoff to Coupling-Phase
Spectroscopy}\label{a-handoff-to-coupling-phase-spectroscopy}

Suppose a training phase transition is genuinely a property of the
augmented state.

Then a parameter-only diagnostic may miss part of the mechanism.

This motivates a later instrumentation question:

\begin{quote}
can we probe the Jacobian of the coupled model-optimizer state strongly
enough to detect changes in dynamical regime?
\end{quote}

The Atlas names the downstream GCL programme \textbf{Coupling-Phase
Spectroscopy (CPS)}.

At present, this chapter makes no project-specific empirical claim
because no exact public source lock for CPS results is bound here. The
source-lock manifest records that boundary explicitly.

The chapter therefore contributes the mathematical object and the
diagnostic question, not an unverified result.

\section{Five mistakes to avoid}\label{five-mistakes-to-avoid-1}

\subsection{Mistake 1: ``The Hessian is the optimizer
dynamics.''}\label{mistake-1-the-hessian-is-the-optimizer-dynamics.}

No.~Curvature enters the recurrence, but optimizer memory changes the
state and the update Jacobian.

\subsection{Mistake 2: ``Eigenvalues inside the unit disk mean every
perturbation shrinks every
step.''}\label{mistake-2-eigenvalues-inside-the-unit-disk-mean-every-perturbation-shrinks-every-step.}

No.~The worked non-normal example has spectral radius below one and peak
finite-horizon gain above \(2.6\).

\subsection{Mistake 3: ``A transient spike proves
divergence.''}\label{mistake-3-a-transient-spike-proves-divergence.-1}

No.~Transient amplification and asymptotic divergence are different.

\subsection{Mistake 4: ``A local linearization predicts the whole
training
run.''}\label{mistake-4-a-local-linearization-predicts-the-whole-training-run.}

No.~It is local, and realistic training is time-varying and nonlinear.

\subsection{Mistake 5: ``Optimizer state is implementation
detail.''}\label{mistake-5-optimizer-state-is-implementation-detail.}

No.~If the next update depends on it, it is part of the dynamical state.

\section{Atlas connections}\label{atlas-connections-15}

\textbf{Non-normality.}\\
This chapter turns transient-growth mathematics into an optimizer-state
mechanism.

\textbf{Coupling-Phase Spectroscopy.}\\
The augmented-state Jacobian becomes a candidate probe for phase
changes.

\textbf{Router dynamics.}\\
Routing systems with memory can be analyzed through similar
coupled-state ideas.

\textbf{Spectral diagnostics.}\\
Eigenvalues, singular values, pseudospectra, and finite-horizon gains
become complementary diagnostic objects.

\textbf{Variational optimization.}\\
Later chapters will ask whether better update rules can be derived from
geometry and discrete variational structure rather than tuned only
through heuristics.

The recurring Atlas shift is:

\[
\boxed{
\text{optimizer as update rule}
\longrightarrow
\text{optimizer-model pair as dynamical system}.
}
\]

\section{Closing view}\label{closing-view-22}

A modern optimizer remembers.

Once it remembers, its memory becomes part of the state.

Once the optimizer state is part of the state, the natural object is no
longer only a loss surface or a gradient. It is a coupled recurrence
with its own Jacobian, stability region, transient behavior, and
finite-horizon propagator.

The simple quadratic example is small enough to solve exactly and rich
enough to expose the point.

Stable eigenvalues do not forbid transient amplification.

The flywheel is only an allegory.

The augmented state is the object.

\section{References used in this
chapter}\label{references-used-in-this-chapter-32}

\begin{itemize}
\tightlist
\item
  (Polyak 1964)
\item
  (Sutskever et al. 2013)
\item
  (Kingma and Ba 2015)
\item
  (Lessard et al. 2016)
\end{itemize}

See `sources/source-locks/ATLAS-CH-OPTDYN-001.yaml` for exact source
identities and claim scope.

\chapter{Coupling-Phase Spectroscopy}\label{coupling-phase-spectroscopy}

\textbf{Epistemic status:} audited Optimizer-State Dynamics + current
public GSD-001 transition/artifact evidence + Atlas-owned exact
spectroscopy and prediction protocol.\\
\textbf{Specification:} manuscript/specifications/ATLAS-CH-CPS-001.md\\
\textbf{Derivation packet:}
mathematics/derivations/ATLAS-CH-CPS-001-DERIVATIONS.md\\
\textbf{Computational witness:}
mathematics/computational-witnesses/ATLAS-CW-CPS-001.md\\
\textbf{Source lock:} sources/source-locks/ATLAS-CH-CPS-001.yaml

A training run can change behavior before we know what changed inside
the optimization dynamics.

A benchmark state can flip.

A representation can drift.

An optimizer can carry momentum, moments, scheduler state, or update
geometry that is invisible in a parameter-only snapshot.

Coupling-Phase Spectroscopy asks a narrow question:

\begin{quote}
can local structure in the augmented model-optimizer state diagnose or
predict a behavioral transition before the transition is observed?
\end{quote}

The chapter begins with a strict warning.

\[
\boxed{
\text{diagnosis}
\neq
\text{prediction}
\neq
\text{causal mechanism}.
}
\]

All three can be scientifically useful.

They require different evidence.

\section{The object is the augmented training
state}\label{the-object-is-the-augmented-training-state}

Optimizer-State Dynamics established that optimizer memory belongs in
the dynamical state whenever the next update depends on it.

Write:

\[
z_{t+1}=F_t(z_t,\xi_t).
\]

The local Jacobian is:

\[
J_t
=
\frac{\partial F_t}{\partial z}(z_t,\xi_t).
\]

A short-horizon perturbation is governed locally by:

\[
J_{t+k-1}\cdots J_t.
\]

CPS treats these local maps and products as instrumentation surfaces.

It does not claim that one frozen Jacobian globally describes training.

\section{What is spectroscopy here?}\label{what-is-spectroscopy-here}

The word spectroscopy is used by analogy with measuring several
signatures of a state rather than collapsing everything into one scalar.

A CPS diagnostic vector can contain:

\[
\mathcal S_t
=
\Bigl(
\operatorname{tr}J_t,
\det J_t,
\rho(J_t),
\widehat\sigma_{\max}(J_t),
\mathcal N(J_t),
\widehat G_{t,k},
\mathcal A_t
\Bigr).
\]

Possible coordinates include:

\begin{itemize}
\tightlist
\item
  trace;
\item
  determinant;
\item
  spectral radius;
\item
  dominant singular gain;
\item
  non-normality proxy;
\item
  finite-horizon gain;
\item
  directional alignment.
\end{itemize}

No coordinate is promoted to a universal health score.

\section{The behavioral transition is defined
elsewhere}\label{the-behavioral-transition-is-defined-elsewhere}

CPS does not get to create its own target after looking at optimizer
diagnostics.

The behavioral transition must be frozen independently.

For a checkpoint \(t\), define:

\[
Y_t
=
\mathbf 1
\{
\text{declared behavioral transition occurs in }(t,t+\Delta]
\}.
\]

CPS receives \(Y_t\) as a target variable.

It does not redefine \(Y_t\).

\section{Current GCL state changed since the first Atlas source
lock}\label{current-gcl-state-changed-since-the-first-atlas-source-lock}

The original Optimizer-State Dynamics source lock was written on October
2, 2026.

At that time it recorded no separate public GCL source for a CPS result.

That statement is no longer a complete description of current public
project state.

A fresh organization search on October 6 finds a public GSD programme
with:

\begin{itemize}
\tightlist
\item
  GSD-WP04 --- CPS transition-local dynamics;
\item
  one confirmed behavioral transition;
\item
  an explicit optimizer-artifact availability receipt.
\end{itemize}

The Atlas therefore replaces the old current-state gap with exact newer
project evidence.

\section{What GSD has actually
established}\label{what-gsd-has-actually-established}

The bound GSD-WP03R receipt establishes one robust OLMo2-1B transition
for:

truthy\_answer/surprising\_truth.

The frozen GSD protocol reports:

\begin{itemize}
\tightlist
\item
  step 2000: GENERALIZING;
\item
  step 3000: PATTERN\_MATCHING;
\item
  step 4000: PATTERN\_MATCHING;
\item
  transition localized between steps 2000 and 3000;
\item
  stable controls passed;
\item
  independent replay preserved the transition;
\item
  prompt variants preserved the confirmed direction.
\end{itemize}

This is a behavioral result.

The receipt explicitly does not establish:

\begin{itemize}
\tightlist
\item
  mechanism;
\item
  optimizer cause;
\item
  persistence state;
\item
  capacity-allocation explanation;
\item
  architecture-level conclusion.
\end{itemize}

\section{What GSD-WP04 asks}\label{what-gsd-wp04-asks}

The public GSD research programme defines GSD-WP04 as:

CPS transition-local dynamics.

Its question is:

\begin{quote}
does optimizer-state structure predict or explain confirmed transitions?
\end{quote}

Its declared outputs include:

\begin{itemize}
\tightlist
\item
  optimizer-state extraction;
\item
  augmented-state probes;
\item
  transition-local traces;
\item
  heldout predictive testing.
\end{itemize}

Its scientific exit grammar is:

\begin{itemize}
\tightlist
\item
  PREDICTIVE\_SIGNAL;
\item
  DESCRIPTIVE\_ONLY;
\item
  NO\_SIGNAL.
\end{itemize}

That is a useful experimental contract.

\section{Why there is no WP04 signal result
yet}\label{why-there-is-no-wp04-signal-result-yet}

The exact public transition-local revisions inspected are:

\begin{itemize}
\tightlist
\item
  stage1-step2000-tokens5B;
\item
  stage1-step3000-tokens7B;
\item
  stage1-step4000-tokens9B.
\end{itemize}

The public revisions expose:

\begin{itemize}
\tightlist
\item
  configuration;
\item
  tokenizer;
\item
  model weight shards;
\item
  weight index.
\end{itemize}

They do not expose:

\begin{itemize}
\tightlist
\item
  optimizer state;
\item
  trainer state;
\item
  scheduler state;
\item
  gradient artifacts;
\item
  update-direction artifacts.
\end{itemize}

Therefore the exact GSD-WP04 receipt says:

\[
\boxed{
\mathrm{BLOCKED\_EXTERNAL\_ARTIFACT\_ABSENT}.
}
\]

\section{Blocked is not negative}\label{blocked-is-not-negative}

This distinction matters.

A negative CPS result would require the selected CPS probe to be
computed and then fail.

That did not happen.

The needed artifacts are missing.

Therefore:

\[
\boxed{
\text{missing artifacts}
\not\Rightarrow
\text{no signal}.
}
\]

The scientific question remains open.

\section{The Atlas can still formalize the
protocol}\label{the-atlas-can-still-formalize-the-protocol}

The empirical lane is blocked.

The mathematical lane is not.

The Atlas can still define:

\begin{itemize}
\tightlist
\item
  what a CPS probe is;
\item
  how discovery differs from heldout testing;
\item
  what counts as prediction;
\item
  what counts only as description;
\item
  what evidence would be needed for causal promotion.
\end{itemize}

That is the purpose of the exact finite witness.

\section{Return to the exact momentum
Jacobian}\label{return-to-the-exact-momentum-jacobian}

For scalar quadratic curvature \(h\), fix:

\[
\eta=\frac1{10},
\qquad
\beta=\frac9{10}.
\]

The augmented momentum Jacobian is:

\[
J(h)
=
\begin{pmatrix}
1-\frac h{10}&-\frac9{100}\\
h&\frac9{10}
\end{pmatrix}.
\]

This is inherited directly from the audited OPTDYN derivation.

\section{One invariant is constant}\label{one-invariant-is-constant}

For every \(h\):

\[
\det J(h)=\frac9{10}.
\]

The trace is:

\[
\operatorname{tr}J(h)
=
\frac{19}{10}-\frac h{10}.
\]

Define a toy matrix-derived coordinate:

\[
\kappa(J)
=
1+\det J-\operatorname{tr}J.
\]

Then:

\[
\boxed{
\kappa(J(h))=\frac h{10}.
}
\]

This coordinate is chosen for the witness.

It is not proposed as the CPS statistic for real training.

\section{Four synthetic windows}\label{four-synthetic-windows}

Create two discovery windows:

\[
d_0:\quad h=2,\quad \kappa=\frac15,\quad Y=0,
\]

\[
d_1:\quad h=8,\quad \kappa=\frac45,\quad Y=1.
\]

Create two heldout windows:

\[
h_0:\quad h=3,\quad \kappa=\frac3{10},\quad Y=0,
\]

\[
h_1:\quad h=7,\quad \kappa=\frac7{10},\quad Y=1.
\]

The transition labels are synthetic.

They exist only to demonstrate protocol discipline.

\section{Discovery comes before heldout
evaluation}\label{discovery-comes-before-heldout-evaluation}

Using the discovery windows only:

\[
\frac15<\frac12<\frac45.
\]

Freeze threshold:

\[
\tau=\frac12.
\]

Prediction rule:

\[
\widehat Y
=
\mathbf 1\{\kappa>\tau\}.
\]

Now the rule is frozen.

The heldout labels cannot be used to change:

\begin{itemize}
\tightlist
\item
  statistic;
\item
  threshold;
\item
  state coordinates;
\item
  horizon;
\item
  normalization.
\end{itemize}

\section{Exact heldout replay}\label{exact-heldout-replay}

For \(h_0\):

\[
\frac3{10}<\frac12,
\]

so:

\[
\widehat Y=0.
\]

For \(h_1\):

\[
\frac7{10}>\frac12,
\]

so:

\[
\widehat Y=1.
\]

Therefore:

\[
\boxed{
(\widehat Y_{h_0},\widehat Y_{h_1})=(0,1).
}
\]

The heldout synthetic labels are recovered exactly.

This proves only that the protocol mechanics are coherent.

\section{Why one spectral scalar is not
enough}\label{why-one-spectral-scalar-is-not-enough}

All four toy Jacobians have:

\[
\det J=\frac9{10}.
\]

Their characteristic discriminants are all negative.

Therefore each has a complex-conjugate eigenvalue pair.

For each pair:

\[
|\lambda|^2
=
\det J
=
\frac9{10}.
\]

So:

\[
\boxed{
\rho(J)
=
\sqrt{\frac9{10}}
}
\]

for all four toy windows.

The toy labels differ.

The spectral radius does not.

\section{Same spectral radius, different local
maps}\label{same-spectral-radius-different-local-maps}

For \(h=2\):

\[
J(2)
=
\begin{pmatrix}
\frac45&-\frac9{100}\\
2&\frac9{10}
\end{pmatrix}.
\]

For \(h=8\):

\[
J(8)
=
\begin{pmatrix}
\frac15&-\frac9{100}\\
8&\frac9{10}
\end{pmatrix}.
\]

The matrices are not the same.

Their trace differs.

Their coupling coordinate differs.

Their action on perturbation directions differs.

Thus:

\[
\boxed{
\text{same spectral radius}
\not\Rightarrow
\text{same local augmented-state dynamics}.
}
\]

\section{This is why CPS is
multi-coordinate}\label{this-is-why-cps-is-multi-coordinate}

The point is not that \(\kappa\) is better than spectral radius.

The point is that different matrix summaries answer different questions.

A real CPS probe may need:

\begin{itemize}
\tightlist
\item
  singular gain;
\item
  finite-horizon gain;
\item
  non-normality;
\item
  update-direction alignment;
\item
  layer-local block structure;
\item
  time-varying products.
\end{itemize}

The exact toy merely proves that one scalar can be blind.

\section{Local Jacobians remain
local}\label{local-jacobians-remain-local}

Suppose a Jacobian changes sharply near a transition.

That supports a local dynamical observation.

It does not prove:

\[
\text{global phase change}.
\]

It does not prove:

\[
\text{optimizer state caused behavioral change}.
\]

It does not prove the same signature persists across checkpoints, tasks,
or scales.

The local/global firewall from OPTDYN remains active.

\section{Retrospective description is the weakest positive CPS
result}\label{retrospective-description-is-the-weakest-positive-cps-result}

Imagine scanning all layers, horizons, and statistics after the
transition is known.

Suppose one statistic moves dramatically at the transition.

That is interesting.

It is also selected after seeing the answer.

The correct label is:

\[
\mathrm{DESCRIPTIVE\_ONLY}.
\]

It can motivate a new preregistered test.

It is not prediction.

\section{Prediction requires temporal and selection
discipline}\label{prediction-requires-temporal-and-selection-discipline}

A predictive test needs a probe frozen before heldout evaluation.

For checkpoint \(t\), compute:

\[
X_t
=
\mathrm{Probe}_{\Pi_{\rm CPS}}
(\text{artifacts available at or before }t).
\]

Then predict:

\[
Y_t
\]

for the future window:

\[
(t,t+\Delta].
\]

Information from \(t+\Delta\) cannot enter \(X_t\).

Otherwise the future has leaked into the predictor.

\section{Discovery must be separated from
confirmation}\label{discovery-must-be-separated-from-confirmation}

The discovery set can choose:

\begin{itemize}
\tightlist
\item
  layer;
\item
  state block;
\item
  horizon;
\item
  statistic;
\item
  threshold;
\item
  normalization;
\item
  estimator.
\end{itemize}

Then those choices must freeze.

The heldout set tests the frozen object.

This is the difference between:

\begin{quote}
I found a curve that moves near transitions.
\end{quote}

and:

\begin{quote}
I can identify transition risk before the heldout transition occurs.
\end{quote}

\section{A real CPS heldout
protocol}\label{a-real-cps-heldout-protocol}

For future artifact-complete transition windows:

\begin{enumerate}
\def\labelenumi{\arabic{enumi}.}
\tightlist
\item
  freeze the behavioral transition definition;
\item
  partition discovery and heldout windows;
\item
  choose CPS features on discovery only;
\item
  freeze checkpoint lead \(\Delta\);
\item
  compute pre-transition features only;
\item
  compare against non-transition windows;
\item
  include shuffled-label controls;
\item
  include simple training-time baselines;
\item
  test the frozen predictor on heldout transitions;
\item
  retain negative results.
\end{enumerate}

This structure is already compatible with GSD-WP04's intended exit
grammar.

\section{What counts as
PREDICTIVE\_SIGNAL?}\label{what-counts-as-predictive_signal}

Use:

\[
\mathrm{PREDICTIVE\_SIGNAL}
\]

only when:

\begin{itemize}
\tightlist
\item
  exact required artifacts exist;
\item
  the probe is frozen before heldout testing;
\item
  the heldout metric clears its preregistered criterion;
\item
  relevant baselines/controls do not explain the result.
\end{itemize}

This is predictive association.

It is not causality.

\section{What counts as
DESCRIPTIVE\_ONLY?}\label{what-counts-as-descriptive_only}

Use:

\[
\mathrm{DESCRIPTIVE\_ONLY}
\]

when:

\begin{itemize}
\tightlist
\item
  a transition-local signature exists;
\item
  but it was selected retrospectively;
\item
  or it fails heldout prediction;
\item
  or heldout confirmation was never preregistered.
\end{itemize}

Descriptive signals are legitimate research outputs.

They should not be relabeled as forecasts.

\section{What counts as NO\_SIGNAL?}\label{what-counts-as-no_signal}

Use:

\[
\mathrm{NO\_SIGNAL}
\]

when:

\begin{itemize}
\tightlist
\item
  the required artifacts exist;
\item
  the declared probe family is executable;
\item
  the preregistered test is actually run;
\item
  the signal fails the declared criteria.
\end{itemize}

A negative result can save large downstream compute.

\section{What counts as an artifact
block?}\label{what-counts-as-an-artifact-block}

Use:

\[
\mathrm{BLOCKED\_EXTERNAL\_ARTIFACT\_ABSENT}
\]

when the selected test cannot be executed because required exact
artifacts are unavailable.

That is the current public GSD-WP04 state.

It is an evidence-availability disposition.

It is not a scientific result about CPS.

\section{Prediction is not cause}\label{prediction-is-not-cause}

Suppose a CPS feature predicts a transition.

The feature can still be:

\begin{itemize}
\tightlist
\item
  a correlate;
\item
  a downstream symptom;
\item
  a common-cause indicator;
\item
  a proxy for training time;
\item
  a proxy for curvature or data composition.
\end{itemize}

Therefore:

\[
\boxed{
\text{heldout prediction}
\not\Rightarrow
\text{causal optimizer mechanism}.
}
\]

\section{Causal promotion needs
intervention}\label{causal-promotion-needs-intervention}

A stronger test would manipulate the optimizer-side state or update
while matching relevant controls.

Examples can include:

\begin{itemize}
\tightlist
\item
  controlled optimizer-state reset;
\item
  controlled moment rescaling;
\item
  update-direction substitution;
\item
  matched optimizer-rule intervention;
\item
  state patching where semantics are well-defined.
\end{itemize}

The intervention must be specified before the outcome is interpreted.

Only then can the evidence move toward a causal claim.

\section{The causal hierarchy}\label{the-causal-hierarchy}

Use five levels.

\subsection{C0 --- behavioral
transition}\label{c0-behavioral-transition}

A transition exists independently of CPS.

\subsection{C1 --- local diagnostic
movement}\label{c1-local-diagnostic-movement}

A CPS feature changes near the transition.

\subsection{C2 --- heldout prediction}\label{c2-heldout-prediction}

A frozen pre-transition feature predicts heldout transition windows.

\subsection{C3 --- intervention
sensitivity}\label{c3-intervention-sensitivity}

Changing optimizer-side state/update changes transition behavior under
matched controls.

\subsection{C4 --- mechanism-specific
support}\label{c4-mechanism-specific-support}

Intervention, reconstruction, and alternative-hypothesis attacks support
a declared mechanism class.

Each level needs new evidence.

\section{Current GSD is at C0 for
CPS}\label{current-gsd-is-at-c0-for-cps}

The public source establishes one confirmed behavioral transition.

For CPS specifically, the optimizer artifacts are unavailable.

Therefore current public GSD supports:

\[
C0
\]

as a target transition.

It does not support:

\[
C1,
C2,
C3,
C4
\]

from optimizer-state evidence.

This is not a criticism.

It is the exact evidence boundary.

\section{One transition cannot validate a general
predictor}\label{one-transition-cannot-validate-a-general-predictor}

The current GSD source has one confirmed transition.

A general predictor needs:

\begin{itemize}
\tightlist
\item
  discovery examples;
\item
  heldout examples.
\end{itemize}

One event can support:

\begin{itemize}
\tightlist
\item
  artifact-readiness work;
\item
  a descriptive case study if artifacts exist.
\end{itemize}

It cannot establish broad heldout predictive value by itself.

\section{Small model does not silently become large
model}\label{small-model-does-not-silently-become-large-model}

The confirmed public transition is OLMo2 1B.

Even a successful CPS result there would not automatically transfer to:

\begin{itemize}
\tightlist
\item
  OLMo3 7B;
\item
  OLMo3 32B;
\item
  other model families;
\item
  other optimizers;
\item
  post-training settings.
\end{itemize}

Scale transfer requires its own heldout evidence.

\section{Artifact completeness is part of the
experiment}\label{artifact-completeness-is-part-of-the-experiment}

A CPS packet should identify:

\begin{itemize}
\tightlist
\item
  model revision;
\item
  optimizer state;
\item
  trainer/global step;
\item
  scheduler;
\item
  gradient or update direction if required;
\item
  evaluator/tokenizer;
\item
  code;
\item
  hardware;
\item
  transition label provenance.
\end{itemize}

If one required artifact is absent, the test is not silently weakened.

The block should be explicit.

\section{Why exact identity
matters}\label{why-exact-identity-matters}

A transition-local probe can be invalidated by stale identity.

Optimizer state from one checkpoint cannot be paired casually with model
parameters from another.

Scheduler state from the wrong run is not equivalent.

A CPS result must bind to:

\[
(\text{model},\text{optimizer},\text{step},\text{run},\text{probe code}).
\]

This is scientific provenance, not bookkeeping.

\section{Time-varying products
matter}\label{time-varying-products-matter}

A real training trajectory is not generally autonomous.

The local perturbation product is:

\[
J_{t+k-1}\cdots J_t.
\]

Therefore finite-horizon CPS may need to estimate a product or its
action.

One-step eigenvalues can miss transient amplification.

OPTDYN already proved that lesson in a fixed toy.

CPS turns it into an instrumentation question.

\section{JVP/VJP sketches can make the probe
tractable}\label{jvpvjp-sketches-can-make-the-probe-tractable}

A full frontier-scale Jacobian is too large to materialize naively.

The GSD programme therefore contemplates randomized JVP/VJP estimates.

These can target:

\begin{itemize}
\tightlist
\item
  dominant singular directions;
\item
  local gains;
\item
  selected subspaces;
\item
  layer/block couplings.
\end{itemize}

The result remains an estimate of a declared local object.

It does not become a global training theorem because the computation is
scalable.

\section{Training time is an important
baseline}\label{training-time-is-an-important-baseline}

A CPS statistic may drift smoothly with training step.

Behavioral transitions also occur at particular steps.

A naive association can therefore be explained by time itself.

A heldout protocol should compare against:

\begin{itemize}
\tightlist
\item
  checkpoint index;
\item
  tokens seen;
\item
  learning rate;
\item
  update norm;
\item
  other simple progress variables.
\end{itemize}

A sophisticated probe must outperform the relevant simple baseline to
justify its complexity.

\section{Parameter-only baselines
matter}\label{parameter-only-baselines-matter}

If a parameter-only diagnostic predicts the transition just as well as
the optimizer-state probe, then the added optimizer state may not
provide incremental value for that prediction task.

This does not make optimizer state irrelevant to training.

It narrows the claim.

CPS should report incremental predictive value when that is the
question.

\section{Stable capability controls matter
too}\label{stable-capability-controls-matter-too}

A transition-specific predictor can accidentally be a generic
degradation detector.

GSD already protects against this with stable conventional-capability
controls.

CPS should inherit that discipline.

A broad collapse is different from a selective generalization-state
transition.

\section{Negative CPS evidence is
useful}\label{negative-cps-evidence-is-useful}

Suppose optimizer artifacts become available and a carefully
preregistered probe returns:

\[
\mathrm{NO\_SIGNAL}.
\]

That is informative.

It can redirect effort toward:

\begin{itemize}
\tightlist
\item
  representation/operator diagnostics;
\item
  accessibility/control mechanisms;
\item
  data-window attribution;
\item
  other causal routes.
\end{itemize}

The Atlas should preserve such negative results.

\section{What the current chapter
establishes}\label{what-the-current-chapter-establishes}

Mathematically, the chapter establishes:

\[
\det J(h)=\frac9{10},
\]

\[
\kappa(J(h))=\frac h{10},
\]

the exact synthetic discovery/heldout threshold result, and identical
spectral radius across all four toy windows.

Scientifically, it establishes the current source boundary:

\begin{itemize}
\tightlist
\item
  one public confirmed GSD behavioral transition exists;
\item
  the public CPS work package exists;
\item
  the exact transition-local public revisions lack optimizer artifacts;
\item
  therefore no public GSD CPS signal disposition is currently
  authorized.
\end{itemize}

\section{Non-implications}\label{non-implications}

The chapter rejects:

\[
\text{behavioral transition}
\not\Rightarrow
\text{optimizer-state transition},
\]

\[
\text{local Jacobian change}
\not\Rightarrow
\text{global training phase change},
\]

\[
\text{retrospective alignment}
\not\Rightarrow
\text{prediction},
\]

\[
\text{heldout prediction}
\not\Rightarrow
\text{causal mechanism},
\]

\[
\text{same spectral radius}
\not\Rightarrow
\text{same augmented-state dynamics},
\]

\[
\text{missing optimizer artifacts}
\not\Rightarrow
\text{no CPS signal},
\]

and:

\[
\text{one 1B transition}
\not\Rightarrow
\text{general predictor across scale}.
\]

\section{Atlas connections}\label{atlas-connections-16}

\textbf{Optimizer-State Dynamics.}\\
CPS operationalizes the augmented-state Jacobian as an empirical
diagnostic object.

\textbf{Generalization-State Dynamics.}\\
GSD supplies behavioral transition targets and the heldout predictive
question.

\textbf{Spectral Diagnostics.}\\
CPS supplies a transition-local application of spectral, singular-gain,
and non-normality probes.

\textbf{Experiments as Arguments.}\\
Discovery/freeze/heldout separation is part of the claim structure, not
merely experimental hygiene.

\textbf{Governed research systems.}\\
Artifact identity and explicit block states prevent missing evidence
from being silently converted into negative findings.

\section{Closing view}\label{closing-view-23}

Coupling-Phase Spectroscopy is not a claim that optimizer state explains
every transition.

It is a way to ask the question correctly.

First establish the behavioral transition.

Then obtain the exact optimizer-state artifacts.

Then define the local augmented-state probe.

Then separate exploration from heldout prediction.

Then separate prediction from causality.

The current public GSD programme has reached the first step and
specified the second.

The required optimizer artifacts are not publicly present at the exact
transition checkpoints.

So the present result is neither success nor failure of CPS.

It is a precise boundary:

\[
\boxed{
\text{transition confirmed; optimizer-state spectroscopy currently artifact-blocked}.
}
\]

That is the right place to continue from when the missing evidence
becomes available.

\section{References used in this
chapter}\label{references-used-in-this-chapter-33}

External mathematical authority is inherited through audited
ATLAS-CH-OPTDYN-001.

Current GCL project evidence is bound to exact public grandchallenge/GSD
files at commit:

ebd4681e1815a3d7f0285cc4ce2bc090b38fae8c

See sources/source-locks/ATLAS-CH-CPS-001.yaml for exact identities and
claim boundaries.

\chapter{Variational and Divergence-Derived
Optimization}\label{variational-and-divergence-derived-optimization}

\textbf{Epistemic status:} audited Numerics and Manifold Optimization
prerequisites + primary Bregman/variational sources + exact GCL MODULUS
programme evidence + Atlas finite derivations.\\
\textbf{Specification:}
manuscript/specifications/ATLAS-CH-VARIOPT-001.md\\
\textbf{Derivation packet:}
mathematics/derivations/ATLAS-CH-VARIOPT-001-DERIVATIONS.md\\
\textbf{Computational witness:}
mathematics/computational-witnesses/ATLAS-CW-VARIOPT-001.md\\
\textbf{Source lock:} sources/source-locks/ATLAS-CH-VARIOPT-001.yaml

Optimization rules can be written down directly.

They can also be derived from a declared geometry, a divergence, an
action, or a constraint.

Those derivations are valuable because they expose which structure the
update is intended to respect.

But the derivation does not grant every desirable property at once.

The governing rule is:

\[
\boxed{
\text{derivation structure}
\neq
\text{performance guarantee}.
}
\]

This chapter develops that distinction through Bregman geometry,
variational mechanics, symplectic updates, and an exact GCL MODULUS
programme bridge.

\section{Two inherited firewalls}\label{two-inherited-firewalls}

NUMERICS-001 established that:

\[
\text{exact flow}
\neq
\text{numerical update},
\]

and that:

\begin{itemize}
\tightlist
\item
  stability is not accuracy;
\item
  structure preservation is not exact energy conservation;
\item
  order claims require assumptions;
\item
  a computational analogy is not automatically a theorem.
\end{itemize}

MANOPT-001 established that:

\begin{itemize}
\tightlist
\item
  gradients depend on metric;
\item
  tangent directions are local legal velocities;
\item
  retractions return tangent data to the constraint set;
\item
  retractions need not equal exponential maps;
\item
  constraint preservation does not imply descent;
\item
  stationarity does not imply global optimality.
\end{itemize}

VARIOPT inherits both.

\section{Why divergences enter
optimization}\label{why-divergences-enter-optimization}

Euclidean distance is only one way to measure a local change.

A convex generator can define a directed discrepancy that adapts to the
geometry of the domain.

Let:

\[
\phi:\Omega\to\mathbb R
\]

be differentiable and strictly convex.

The Bregman divergence is:

\[
D_\phi(y,x)
=
\phi(y)-\phi(x)-\langle\nabla\phi(x),y-x\rangle.
\]

Geometrically, it measures how far the graph of \(\phi\) at \(y\) sits
above the tangent hyperplane to \(\phi\) at \(x\).

\section{Divergence is not
distance}\label{divergence-is-not-distance}

The ordering matters.

Take:

\[
\phi(x)=e^x.
\]

Then:

\[
D_\phi(1,0)=e-2,
\]

but:

\[
D_\phi(0,1)=1.
\]

Thus:

\[
D_\phi(1,0)\ne D_\phi(0,1).
\]

A Bregman divergence is therefore not generally a metric.

That remains true even though it can generate useful local geometry.

\section{Local geometry from a
Hessian}\label{local-geometry-from-a-hessian}

Assume \(\phi\) has a locally Lipschitz Hessian near \(x\), for example
bounded third derivatives. Then for a small displacement \(\delta\):

\[
D_\phi(x+\delta,x)
=
\frac12
\delta^\top\nabla^2\phi(x)\delta
+
O(\|\delta\|^3).
\]

The Hessian acts as a local quadratic form. Under merely continuous
second differentiability, the same expansion holds with the weaker
remainder \(o(\|\delta\|^2)\); the displayed cubic-order remainder
requires the stated stronger regularity.

For:

\[
\phi(x)=e^x,
\]

at \(x=0\):

\[
D_\phi(\delta,0)
=
e^\delta-1-\delta
=
\frac12\delta^2
+
\frac16\delta^3
+\cdots.
\]

The local quadratic geometry can be symmetric even when the full
divergence is not.

\section{Geometry can define an
update}\label{geometry-can-define-an-update}

Suppose the objective supplies a local covector \(g_k\).

Instead of taking a Euclidean step, choose the next state by minimizing:

\[
\eta\langle g_k,x\rangle
+
D_\phi(x,x_k).
\]

The first term asks for progress under the local linearized objective.

The second penalizes movement according to the chosen geometry.

At an interior optimum:

\[
\nabla\phi(x_{k+1})
=
\nabla\phi(x_k)-\eta g_k.
\]

The step is simple in the dual coordinate \(\nabla\phi\).

\section{Euclidean descent is one special
case}\label{euclidean-descent-is-one-special-case}

For:

\[
\phi(x)=\frac12\|x\|^2,
\]

we have:

\[
\nabla\phi(x)=x.
\]

So:

\[
x_{k+1}=x_k-\eta g_k.
\]

Ordinary gradient descent therefore sits inside the broader
divergence-derived construction.

The broader construction should not be reduced back to Euclidean
distance by terminology.

\section{An entropic coordinate
system}\label{an-entropic-coordinate-system}

For positive scalar \(x\), choose:

\[
\phi(x)=x\log x-x.
\]

Then:

\[
\nabla\phi(x)=\log x.
\]

The update becomes:

\[
\log x_{k+1}
=
\log x_k-\eta g_k.
\]

Hence:

\[
x_{k+1}
=
x_k e^{-\eta g_k}.
\]

For:

\[
x_k=1,
\qquad
\eta=1,
\qquad
g_k=\log2,
\]

we obtain exactly:

\[
x_{k+1}=\frac12.
\]

A change that is additive in dual coordinates becomes multiplicative in
primal coordinates.

\section{Divergence-derived does not mean
descending}\label{divergence-derived-does-not-mean-descending}

Consider:

\[
f(x)=\frac12(x-2)^2.
\]

At \(x_0=0\):

\[
g_0=-2.
\]

With the Euclidean Bregman generator and step size:

\[
\eta=3,
\]

the update is:

\[
x_1=6.
\]

Yet:

\[
f(x_0)=2,
\qquad
f(x_1)=8.
\]

The objective increased.

Thus:

\[
\boxed{
\text{derived from a convex divergence}
\not\Rightarrow
\text{automatic descent}.
}
\]

Step-size, smoothness, convexity, and algorithm-specific conditions
still matter.

\section{From geometry to dynamics}\label{from-geometry-to-dynamics}

A first-order update uses a local objective model plus a movement
penalty.

A variational formulation can instead define a whole trajectory through
an action.

For path \(q(t)\):

\[
\mathcal A[q]
=
\int L(q,\dot q,t)\,dt.
\]

Stationarity gives the Euler-Lagrange equation:

\[
\frac{d}{dt}
\frac{\partial L}{\partial\dot q}
-
\frac{\partial L}{\partial q}
=
0.
\]

The object being made stationary is the action.

It need not be the optimization objective itself.

\section{Bregman Lagrangian}\label{bregman-lagrangian}

Wibisono, Wilson, and Jordan introduced a Bregman Lagrangian:

\[
\mathcal L(X,V,t)
=
e^{\alpha_t+\gamma_t}
\left[
D_h(X+e^{-\alpha_t}V,X)
-
e^{\beta_t}f(X)
\right].
\]

Under their declared ideal-scaling conditions and regularity
assumptions, its Euler-Lagrange dynamics organize broad families of
accelerated continuous-time optimization flows.

This is a powerful bridge:

\[
\text{divergence geometry}
+
\text{variational dynamics}
\to
\text{optimization flow}.
\]

But the source does not license every possible discretization of that
flow.

\section{Continuous flow still is not the
algorithm}\label{continuous-flow-still-is-not-the-algorithm}

The audited Numerics boundary remains active.

Even when a continuous optimization flow has a proved rate:

\[
\text{continuous theorem}
\not\Rightarrow
\text{arbitrary discrete theorem}.
\]

A discrete method must be analyzed as its own map.

That map can have different:

\begin{itemize}
\tightlist
\item
  stability;
\item
  error;
\item
  invariants;
\item
  energy behavior;
\item
  descent behavior.
\end{itemize}

\section{Discrete action}\label{discrete-action}

Variational mechanics offers a different route.

Instead of discretizing the differential equation directly, discretize
the action.

Let:

\[
L_d(q_k,q_{k+1};h)
\]

represent one discrete action contribution.

Then:

\[
\mathcal A_d
=
\sum_k
L_d(q_k,q_{k+1};h).
\]

Stationarity with respect to an interior point \(q_k\) gives:

\[
D_2L_d(q_{k-1},q_k)
+
D_1L_d(q_k,q_{k+1})
=
0.
\]

This is the discrete Euler-Lagrange equation.

\section{Why derive the discrete rule this
way?}\label{why-derive-the-discrete-rule-this-way}

A discrete variational principle can preserve geometric structure
inherited from the action formulation.

Marsden and West develop this framework systematically and show how
regular discrete variational mechanics generates symplectic schemes.

The important word is \textbf{structure}.

The theorem is not:

\begin{quote}
variational integrator = exact trajectory.
\end{quote}

Nor is it:

\begin{quote}
symplectic integrator = monotonically improving objective.
\end{quote}

\section{Exact oscillator
construction}\label{exact-oscillator-construction}

Take the harmonic-oscillator Lagrangian:

\[
L(q,\dot q)
=
\frac12\dot q^2-\frac12q^2.
\]

Use the left-endpoint discrete Lagrangian:

\[
L_d(q_k,q_{k+1};h)
=
\frac{(q_{k+1}-q_k)^2}{2h}
-
\frac h2q_k^2.
\]

This is now a finite algebraic object.

\section{Discrete Euler-Lagrange
equation}\label{discrete-euler-lagrange-equation}

The two derivatives are:

\[
D_2L_d(q_{k-1},q_k)
=
\frac{q_k-q_{k-1}}h,
\]

and:

\[
D_1L_d(q_k,q_{k+1})
=
-\frac{q_{k+1}-q_k}h
-hq_k.
\]

Set their sum to zero:

\[
\frac{q_k-q_{k-1}}h
-
\frac{q_{k+1}-q_k}h
-hq_k
=
0.
\]

Rearrange:

\[
q_{k+1}
=
(2-h^2)q_k-q_{k-1}.
\]

The update has been derived from a discrete stationarity condition
rather than guessed from the ODE.

\section{Momentum lift}\label{momentum-lift}

Define:

\[
p_k=
\frac{q_k-q_{k-1}}h.
\]

Then the recurrence becomes:

\[
p_{k+1}=p_k-hq_k,
\]

followed by:

\[
q_{k+1}=q_k+h p_{k+1}.
\]

This is one ordering of symplectic Euler.

\section{Exact map}\label{exact-map}

Write:

\[
z_k=
\begin{pmatrix}
q_k\\
p_k
\end{pmatrix}.
\]

Then:

\[
z_{k+1}=M_hz_k,
\]

with:

\[
M_h
=
\begin{pmatrix}
1-h^2 & h\\
-h & 1
\end{pmatrix}.
\]

Its determinant is exactly:

\[
1.
\]

More strongly:

\[
M_h^\top J M_h=J
\]

for the canonical symplectic matrix:

\[
J=
\begin{pmatrix}
0&1\\
-1&0
\end{pmatrix}.
\]

\section{Symplectic is a precise
property}\label{symplectic-is-a-precise-property}

Symplecticity means the map preserves the canonical symplectic form.

It is not a synonym for:

\begin{itemize}
\tightlist
\item
  stable;
\item
  accurate;
\item
  energy conserving;
\item
  reversible;
\item
  objective descending;
\item
  globally optimal.
\end{itemize}

These properties require their own evidence.

\section{Exact energy
counterexample}\label{exact-energy-counterexample}

Use:

\[
h=\frac12,
\qquad
(q_0,p_0)=(0,1).
\]

The update gives:

\[
p_1=1,
\]

\[
q_1=\frac12.
\]

For:

\[
H(q,p)=\frac12(q^2+p^2),
\]

the initial energy is:

\[
H_0=\frac12.
\]

After one exactly symplectic step:

\[
H_1=\frac58.
\]

So energy increased by:

\[
\frac18.
\]

The symplectic identity remains exact.

\section{Exact objective
counterexample}\label{exact-objective-counterexample}

Define a toy objective:

\[
F(q)=\frac12q^2.
\]

At the same states:

\[
F(q_0)=0,
\]

\[
F(q_1)=\frac18.
\]

Therefore the symplectic map moved uphill under this objective.

This is not a defect in the integrator.

The map was derived to preserve variational/symplectic structure, not to
minimize \(F\) at every step.

\section{Optimization and mechanics use different
objectives}\label{optimization-and-mechanics-use-different-objectives}

This distinction is easy to blur.

In mechanics:

\begin{itemize}
\tightlist
\item
  the action determines equations of motion;
\item
  the Hamiltonian can represent energy;
\item
  a symplectic map preserves phase-space geometry.
\end{itemize}

In optimization:

\begin{itemize}
\tightlist
\item
  an objective is to be minimized;
\item
  a step may be designed for descent, convergence, acceleration,
  constraint preservation, or another criterion.
\end{itemize}

A variational optimizer can connect these worlds.

It must still state which mathematical object is doing which job.

\section{Manifold geometry
re-enters}\label{manifold-geometry-re-enters}

MANOPT-001 established another route to geometry-derived updates.

Given a constrained state:

\begin{enumerate}
\def\labelenumi{\arabic{enumi}.}
\tightlist
\item
  compute an ambient or Riemannian gradient;
\item
  project or represent it in the tangent space;
\item
  choose a tangent step;
\item
  retract back to the manifold.
\end{enumerate}

This pipeline resembles variational construction in one broad sense:

the update is constrained by declared geometry.

But a retraction-based update need not arise from an action principle.

\section{MODULUS as programme
evidence}\label{modulus-as-programme-evidence}

At protected commit:

7ca4ffcdace32d5ff79c27ad557bfb690fac0af3

the GCL MODULUS project describes itself as an optimization research
toolkit that treats geometry as a first-class optimizer control surface.

Its Hyperball implementation wraps a base optimizer.

The base optimizer proposes an update.

Hyperball can then:

\begin{itemize}
\tightlist
\item
  project it into a tangent space;
\item
  control its norm;
\item
  apply an approximate angular-step policy;
\item
  retract the result to a sphere;
\item
  or apply ball-style norm control.
\end{itemize}

This is a concrete programme example of deriving an update from explicit
geometric constraints.

\section{The exact MODULUS
pipeline}\label{the-exact-modulus-pipeline}

For parameter group \(w\) and base proposal \(u\), the source
implements:

\[
u_\parallel
=
\frac{\langle u,w\rangle}{\|w\|^2}w,
\]

\[
u_\perp
=
u-u_\parallel.
\]

A target-angle branch sets desired tangent-update norm approximately to:

\[
\alpha\|w\|.
\]

The finite point is then retracted.

This separates:

\begin{itemize}
\tightlist
\item
  base optimizer policy;
\item
  tangent geometry;
\item
  step-control policy;
\item
  finite constraint restoration.
\end{itemize}

\section{Exact programme witness}\label{exact-programme-witness}

Take:

\[
w=(1,0),
\qquad
u=(1,1).
\]

Then:

\[
u_\parallel=(1,0),
\]

and:

\[
u_\perp=(0,1).
\]

With programme parameter:

\[
\alpha=1
\]

and unit radius, the pre-retraction tangent update has norm \(1\).

The retracted point is:

\[
w_+
=
\frac{(1,1)}{\sqrt2}.
\]

\section{The target-angle subtlety}\label{the-target-angle-subtlety}

The exact angle from \(w\) to \(w_+\) is:

\[
\frac\pi4.
\]

That is not \(1\) radian.

For unit tangent direction \(v\), the general finite update:

\[
w_+
=
\frac{w+\alpha v}{\sqrt{1+\alpha^2}}
\]

has:

\[
\theta=\arctan\alpha.
\]

For small \(\alpha\):

\[
\theta
=
\alpha-\frac{\alpha^3}{3}
+
O(\alpha^5).
\]

So the implementation's target-angle parameter acts as a tangent-norm
angular proxy.

It agrees to first order.

It is not an exact exponential-map/geodesic-angle control at finite
step.

\section{Why this distinction
matters}\label{why-this-distinction-matters}

If a parameter is called an angle, users may mentally identify it with a
geodesic angle.

The exact construction shows the actual semantics.

This is the same discipline applied throughout the Atlas:

\begin{itemize}
\tightlist
\item
  name the representation;
\item
  name the update;
\item
  derive what it actually preserves;
\item
  do not import stronger semantics from a convenient label.
\end{itemize}

\section{MODULUS is not being promoted to general
theory}\label{modulus-is-not-being-promoted-to-general-theory}

The GCL repository profile for MODULUS is explicit.

Its authority covers:

\begin{itemize}
\tightlist
\item
  versioned implementation;
\item
  benchmarks;
\item
  empirical evidence.
\end{itemize}

It has no claim-promotion or certification authority.

Accordingly, VARIOPT uses MODULUS to show what one governed GCL
geometry-first update pipeline actually does.

Established mathematical claims continue to come from established
sources or Atlas proofs.

\section{Geometry-derived is broader than
variational}\label{geometry-derived-is-broader-than-variational}

There are now at least three distinct mechanisms in view.

\subsection{Divergence-derived}\label{divergence-derived}

Choose a convex generator and derive a local proximal/mirror step.

\subsection{Variational}\label{variational}

Choose an action or Lagrangian and derive stationarity equations.

\subsection{Constraint-geometric}\label{constraint-geometric}

Choose a legal state geometry and transform an update through
projection/retraction.

All three can produce structured updates.

They are not interchangeable derivations.

\section{A variational derivation is not a convergence
certificate}\label{a-variational-derivation-is-not-a-convergence-certificate}

Suppose a discrete rule comes from a stationary discrete action.

That tells us where the rule came from.

Convergence of an optimization objective can still require:

\begin{itemize}
\tightlist
\item
  convexity;
\item
  smoothness;
\item
  coercivity;
\item
  step restrictions;
\item
  damping;
\item
  Lyapunov arguments;
\item
  stochastic assumptions;
\item
  constraint regularity.
\end{itemize}

The derivation and the convergence proof are separate artifacts.

\section{A divergence is not automatically the right
geometry}\label{a-divergence-is-not-automatically-the-right-geometry}

Bregman divergences are generated by a chosen convex function.

Different generators define different local Hessians and different dual
coordinates.

The right generator can depend on:

\begin{itemize}
\tightlist
\item
  domain;
\item
  constraints;
\item
  scale;
\item
  sparsity;
\item
  probability structure;
\item
  desired invariances.
\end{itemize}

No source in this chapter proves one universal generator.

\section{A symplectic optimizer would still need an optimization
story}\label{a-symplectic-optimizer-would-still-need-an-optimization-story}

If an optimization method is designed to preserve a symplectic form in
an enlarged state containing momentum, that can be a meaningful
structural property.

But one must still answer:

\begin{itemize}
\tightlist
\item
  what objective is minimized?
\item
  what damping is present?
\item
  what convergence notion is used?
\item
  how is the discrete step chosen?
\item
  what happens to energy-like quantities?
\item
  what does the structure buy empirically?
\end{itemize}

Symplecticity alone does not answer them.

\section{Bregman-Lagrangian acceleration remains
source-scoped}\label{bregman-lagrangian-acceleration-remains-source-scoped}

The Bregman-Lagrangian framework demonstrates that acceleration can be
organized through continuous-time variational structure under specific
scaling conditions.

That is stronger than an analogy.

But it remains a theorem about a declared framework.

The Atlas does not infer:

\[
\text{has a divergence}
\Rightarrow
\text{accelerated},
\]

nor:

\[
\text{has a Lagrangian}
\Rightarrow
\text{optimal convergence}.
\]

\section{The final architecture
lesson}\label{the-final-architecture-lesson}

The useful abstraction is not ``physics-inspired optimizer.''

It is more precise:

\[
\boxed{
\text{choose structure}
\to
\text{derive legal dynamics}
\to
\text{discretize explicitly}
\to
\text{prove or measure the property actually needed}.
}
\]

Each arrow can fail if left implicit.

\section{Evidence classes}\label{evidence-classes}

A VARIOPT claim should identify whether it is:

\begin{itemize}
\tightlist
\item
  a standard theorem from the mathematical literature;
\item
  an Atlas exact finite derivation;
\item
  a numerical model;
\item
  a GCL programme implementation fact;
\item
  an empirical benchmark result;
\item
  an architectural proposal.
\end{itemize}

This chapter deliberately contains all but the empirical-superiority
class.

They remain separate.

\section{Durable non-implications}\label{durable-non-implications}

The chapter freezes:

\[
\text{divergence}
\not\Rightarrow
\text{metric},
\]

\[
\text{divergence-derived update}
\not\Rightarrow
\text{descent},
\]

\[
\text{variational derivation}
\not\Rightarrow
\text{accuracy or stability},
\]

\[
\text{symplecticity}
\not\Rightarrow
\text{energy conservation},
\]

\[
\text{symplecticity}
\not\Rightarrow
\text{objective descent},
\]

\[
\text{constraint preservation}
\not\Rightarrow
\text{optimizer superiority},
\]

and:

\[
\text{GCL programme mechanism}
\not\Rightarrow
\text{established general theory}.
\]

\section{Closing the architecture
layer}\label{closing-the-architecture-layer}

VARIOPT is the last chapter still in architecture state at this
transaction baseline.

Its role is therefore not to add another optimizer catalog.

It closes an architectural question that runs through the Atlas:

\begin{quote}
how should adaptive systems derive updates when the state space,
numerical structure, or optimization geometry matters?
\end{quote}

The answer is not one universal rule.

It is a disciplined methodology:

\begin{enumerate}
\def\labelenumi{\arabic{enumi}.}
\tightlist
\item
  declare the geometry or action;
\item
  derive the update;
\item
  identify the preserved structure;
\item
  distinguish the discrete algorithm from its continuous source;
\item
  test descent, stability, accuracy, and empirical utility separately;
\item
  bind programme-specific mechanisms to their actual authority.
\end{enumerate}

\section{References used in this
chapter}\label{references-used-in-this-chapter-34}

\begin{itemize}
\tightlist
\item
  L. M. Bregman (1967), \emph{The relaxation method of finding the
  common point of convex sets and its application to the solution of
  problems in convex programming}, DOI 10.1016/0041-5553(67)90040-7.
\item
  J. E. Marsden and M. West (2001), \emph{Discrete mechanics and
  variational integrators}, Acta Numerica 10:357--514, DOI
  10.1017/S096249290100006X.
\item
  Andre Wibisono, Ashia C. Wilson, and Michael I. Jordan (2016), \emph{A
  variational perspective on accelerated methods in optimization}, PNAS
  113(47):E7351--E7358, DOI 10.1073/pnas.1614734113.
\item
  protected GCL MODULUS implementation/programme evidence at commit
  7ca4ffcdace32d5ff79c27ad557bfb690fac0af3.
\end{itemize}

Exact authority and claim boundaries are recorded in:

sources/source-locks/ATLAS-CH-VARIOPT-001.yaml

\chapter{Networks as Numerical
Schemes}\label{networks-as-numerical-schemes}

\textbf{Epistemic status:} established numerical analysis + primary
residual/continuous-depth architecture sources + Atlas synthesis + exact
scalar witness.\\
\textbf{Specification:}
manuscript/specifications/ATLAS-CH-NETNUM-001.md\\
\textbf{Formal packet:}
mathematics/derivations/ATLAS-CH-NETNUM-001-DERIVATIONS.md\\
\textbf{Computational witness:}
mathematics/computational-witnesses/ATLAS-CW-NETNUM-001.md\\
\textbf{Source lock:} sources/source-locks/ATLAS-CH-NETNUM-001.yaml

A residual block looks like Euler's method.

That resemblance is useful.

It is also one of the easiest analogies in modern deep learning to
overstate.

Write a residual layer as

\[x_{k+1}=x_k+F_k(x_k).\]

Write one explicit-Euler step for

\[\frac{dx}{dt}=f(t,x).\]

as

\[
x_{k+1}=x_k+h_k f(t_k,x_k).
\]

The shapes match.

But numerical analysis does not begin with a shape match.

It begins with a reference problem.

To speak meaningfully about consistency, local truncation error, global
error, stability, step size, stiffness, or reversibility, we must know
what continuous evolution the discrete map is supposed to approximate,
how one layer corresponds to one step, what changes when the mesh is
refined, and which norm or state metric is being used.

The governing rule of this chapter is therefore:

\begin{quote}
residual form opens the numerical lens; it does not, by itself, complete
the numerical interpretation.
\end{quote}

That distinction lets us use numerical analysis rigorously rather than
decoratively.

\section{From architecture to update
map}\label{from-architecture-to-update-map}

The Architecture History chapter established the residual form

\[x_{k+1}=x_k+F_k(x_k).\]

as a structural shift.

The identity path transports the current state forward.

The residual branch computes a change.

This makes depth look less like repeated replacement and more like
repeated state update.

The Numerical Methods chapter supplied the corresponding mathematical
object:

\[x_{n+1}=\Psi_h(t_n,x_n).\]

A one-step method takes the current numerical state and advances it by a
discrete rule.

These objects align naturally.

The question is not whether they look similar.

The question is what additional assumptions make the alignment
mathematically substantive.

\section{The bridge must be
declared}\label{the-bridge-must-be-declared}

Suppose the intended continuous system is

\[\frac{dx}{dt}=f(t,x).\]

Explicit Euler gives

\[x_{k+1}=x_k+h_kf(t_k,x_k).\]

A residual layer becomes Euler-compatible after the declaration

\[F_k(x)=h_k f(t_k,x).\]

This line performs the bridge.

It says:

\begin{itemize}
\tightlist
\item
  which continuous field is being sampled;
\item
  which layer corresponds to which time;
\item
  which scale is the step size.
\end{itemize}

Without that bridge, \texttt{F\_k} is simply a learned residual map.

There is nothing wrong with that.

But numerical-analysis terminology must remain attached to the object it
actually describes.

\section{Residual architecture is not a unique
ODE}\label{residual-architecture-is-not-a-unique-ode}

Could any residual stack be rewritten as some continuous system?

In a loose interpolation sense, often yes.

But that is not the same as saying the network has a unique or canonical
underlying ODE.

A finite sequence of maps

\texttt{Psi\_0,Psi\_1,...,Psi\_\{N-1\}}

can be embedded or interpolated in many ways.

Different continuous fields may agree on the sampled layer states.

Different time parameterizations may produce the same discrete sequence.

Different modified equations may explain the same discrete update to a
given asymptotic order.

Therefore the ODE is not generally recovered uniquely from residual
syntax.

The Architecture History audit already fixed the safe statement:

\begin{quote}
residual systems admit an integrator lens.
\end{quote}

This chapter sharpens the conditions under which that lens supports
numerical claims.

\section{Autonomous versus nonautonomous
depth}\label{autonomous-versus-nonautonomous-depth}

Weight sharing is sometimes treated as the dividing line between
dynamical and nondynamical networks.

That is too crude.

If

\[F_k(x)=h f(x).\]

for a common \texttt{f}, then the blocks have the direct form of
fixed-step Euler applied to an autonomous ODE.

But a continuous system can be explicitly time-dependent:

\[\frac{dx}{dt}=f(t,x).\]

Then layer-varying residuals can correspond to

\[F_k(x)=h_k f(t_k,x).\]

So untied weights do not rule out an ODE interpretation.

They may correspond to a nonautonomous field.

The real issue is whether there is a declared continuous object and a
meaningful sampling/refinement relation.

\section{A layer is not yet a time
step}\label{a-layer-is-not-yet-a-time-step}

Suppose we introduce a residual scale:

\[x_{k+1}=x_k+\alpha_kG_k(x_k).\]

It is tempting to call \(\alpha_k\) the step size.

Sometimes that is exactly right.

If \(G_k\) samples a declared vector field, then \(\alpha_k\) can play
the role of \(h_k\).

But if changing \(\alpha_k\) is accompanied by retraining \(G_k\),
changing normalization, changing feature geometry, or otherwise changing
the update field, then the comparison is not simply ``same ODE, smaller
step.''

A step size has meaning only relative to an underlying evolution.

A residual scale always changes the network map.

It does not automatically refine a fixed continuous problem.

\section{Local truncation error requires an exact
flow}\label{local-truncation-error-requires-an-exact-flow}

The Numerics chapter defined the exact flow over one step as

\(\Phi_h\)

and the numerical map as \(\Psi_h\).

The local defect is computed by starting the numerical step from the
exact state:

\[
\delta_{k+1}
=
\Phi_{h_k}(t_k,x(t_k))
-
\Psi_k(x(t_k)).
\]

For explicit Euler, under the usual smoothness assumptions,

\[
\delta_{k+1}=O(h_k^2).
\]

That statement contains more structure than it first appears.

It assumes:

\begin{itemize}
\tightlist
\item
  a continuous trajectory exists;
\item
  an exact flow is defined;
\item
  the layer is tied to a numerical step;
\item
  \(h_k\) is a meaningful small parameter.
\end{itemize}

If a residual block is merely a learned map with no declared continuous
reference, then saying it has ``small local truncation error'' is not
merely unproven.

The quantity has not been defined.

\section{Small residual does not mean small truncation
error}\label{small-residual-does-not-mean-small-truncation-error}

Suppose

\(\|F_k(x)\|\)

is small.

Does that imply the corresponding Euler defect is small?

No.

A small update norm can arise because:

\begin{itemize}
\tightlist
\item
  the vector field itself is small;
\item
  the step is small;
\item
  the representation is scaled;
\item
  normalization suppresses magnitude;
\item
  the network has learned a near-identity map;
\item
  the local coordinates make the update appear small.
\end{itemize}

Local truncation error compares a numerical step to the exact flow of a
specified continuous problem.

The norm of the residual branch alone does not perform that comparison.

This is a recurring category error:

\begin{quote}
small update is not the same statement as accurate discretization.
\end{quote}

\section{Global error is not layerwise error added
up}\label{global-error-is-not-layerwise-error-added-up}

Define the global numerical error as

\[e_k=x(t_k)-x_k.\]

This error is not simply the sum of local defects.

Each old error is transported through later steps.

For a one-step method with the declared Lipschitz error-propagation
bound, the error recursion has the schematic form

\[
\|e_{k+1}\|
\le
(1+C h_k)\|e_k\|
+
\|\delta_{k+1}\|.
\]

Here \(C\) and the norm depend on the fixed reference problem and
interval; this is not an unconditional guarantee for a learned residual
network.

The factor multiplying \texttt{e\_k} matters.

Repeated amplification can dominate small local defects.

That is why numerical analysis separates:

\begin{itemize}
\tightlist
\item
  consistency;
\item
  stability;
\item
  convergence.
\end{itemize}

A network analogy that mentions only one-step approximation misses the
accumulated-dynamics problem.

\section{More layers do not automatically mean finer
discretization}\label{more-layers-do-not-automatically-mean-finer-discretization}

Take a 12-layer residual network and a 24-layer residual network.

Can we say the second uses half the step size?

Only if the two belong to a declared refinement family.

For a numerical refinement, we normally hold the continuous problem
fixed while changing the mesh.

A deeper neural network may instead change:

\begin{itemize}
\tightlist
\item
  all learned parameters;
\item
  residual scaling;
\item
  width;
\item
  normalization;
\item
  activation structure;
\item
  feature coordinates;
\item
  training objective;
\item
  optimizer.
\end{itemize}

That is a different model, not necessarily a finer mesh.

So the statement

\begin{quote}
more layers give smaller discretization error
\end{quote}

needs an actual refinement construction.

Depth alone does not supply it.

\section{The scalar test equation enters neural
depth}\label{the-scalar-test-equation-enters-neural-depth}

The Numerical Methods chapter used the scalar test equation

\(\dot x=\lambda x\).

This small system exposes stability cleanly.

For explicit Euler,

\(\,x_{k+1}=(1+h\lambda)x_k\,\).

Define

\(z=h\lambda\).

The amplification factor is

\(R(z)=1+z\).

The standard non-growth stability condition is

\(|1+z|\le1\).

This exact calculation transfers directly to a linear residual block
when the block has been declared to represent the Euler step.

Now numerical stability has architectural meaning:

small perturbations in that scalar mode are multiplied by the same
amplification factor across depth.

\section{A stable continuous mode can become an unstable network
step}\label{a-stable-continuous-mode-can-become-an-unstable-network-step}

Set

\(\lambda=-1\).

The continuous solution decays:

\(x(t)=e^{-t}x(0)\).

Euler gives

\(x_{k+1}=(1-h)x_k\).

The non-growth condition is

\(|1-h|\le1\).

For real nonnegative \(h\), the non-growth set is

\(0\le h\le2\).

Strict asymptotic decay requires

\(0<h<2\).

At \texttt{h=2}, the factor is \texttt{-1}: bounded but non-decaying.

Now choose

\texttt{h=3}.

The exact continuous factor over one step is

\(e^{-3}\).

The Euler factor is

\(-2\).

The continuous mode decays strongly.

The discrete mode doubles in magnitude and flips sign.

Nothing happened to the ODE.

The discretization changed the dynamics.

This is one of the most useful lessons numerical analysis contributes to
neural architecture:

\begin{quote}
the update rule has its own dynamics.
\end{quote}

\section{Numerical stability is not training
stability}\label{numerical-stability-is-not-training-stability}

The previous example concerns forward propagation of state
perturbations.

That is not the same thing as optimization stability.

Training can be unstable because of:

\begin{itemize}
\tightlist
\item
  exploding gradients;
\item
  vanishing gradients;
\item
  learning-rate choice;
\item
  optimizer state;
\item
  minibatch noise;
\item
  curvature;
\item
  non-normal transients;
\item
  parameterization.
\end{itemize}

Forward numerical stability and training stability can interact.

They can share Jacobian structure.

But they are not interchangeable.

Haber and Ruthotto explicitly use dynamical-systems and discrete-ODE
stability ideas to motivate stable deep architectures (Haber and
Ruthotto 2018).

That bridge is useful precisely because the notions are related.

It should not be read as an identity between all numerical and
optimization stability concepts.

\section{The residual Jacobian is a perturbation
propagator}\label{the-residual-jacobian-is-a-perturbation-propagator}

For

\texttt{Psi\_k(x)=x+F\_k(x)},

the Jacobian is

\texttt{J\_\{Psi\_k\}(x)=I+J\_\{F\_k\}(x)}.

For a small perturbation \texttt{delta\ x\_k},

\texttt{delta\ x\_\{k+1\}\ approximately\ J\_\{Psi\_k\}(x\_k)\ delta\ x\_k}.

Across many layers,

\texttt{delta\ x\_N\ approximately\ J\_\{Psi\_\{N-1\}\}\ ...\ J\_\{Psi\_0\}\ delta\ x\_0}.

This is a discrete tangent evolution.

It is one reason residual systems are naturally analyzed dynamically.

But again the claim must remain scoped.

The Jacobian product describes local first-order perturbation
propagation along the realized discrete trajectory.

It is not, by itself, a theorem about generalization or optimization
success.

\section{Exact finite witness: same flow, different discretization
behavior}\label{exact-finite-witness-same-flow-different-discretization-behavior}

Consider

\(x'=-x\), \(x(0)=1\).

The exact solution at time \texttt{1} is

\(e^{-1}\).

Use explicit Euler with \(N\) equal steps:

\(h=1/N\).

Then

\(x_N=(1-1/N)^N\).

For \texttt{N=2}:

\(x_2=1/4\).

For \texttt{N=4}:

\(x_4=81/256\).

For \texttt{N=8}:

\(x_8=5764801/16777216\).

These values move toward

\(e^{-1}\)

over the declared finite sequence.

The exact witness records both the rational values and their error
against the exact flow.

But the finite table is not the convergence theorem.

The theorem that Euler is globally first order under standard
assumptions comes from numerical analysis (Iserles 2008).

The witness only makes the mechanism visible.

\section{Refinement is a family, not a single
network}\label{refinement-is-a-family-not-a-single-network}

The exact witness contains something ordinary neural-network comparisons
often lack.

It contains a controlled family:

\begin{itemize}
\tightlist
\item
  same ODE;
\item
  same initial state;
\item
  same final time;
\item
  same numerical method;
\item
  only the mesh changes.
\end{itemize}

That is what gives the phrase \textbf{refinement} mathematical force.

If a neural architecture study wants to make an analogous convergence
claim, it needs a corresponding family.

For example, it might declare:

\begin{itemize}
\tightlist
\item
  a continuous target field;
\item
  a sequence of grids;
\item
  parameter interpolation across grids;
\item
  a residual scaling rule;
\item
  a common output time;
\item
  a norm for comparing trajectories.
\end{itemize}

Without such a construction, comparing depths is architecture
comparison, not numerical convergence analysis.

\section{A modified-equation
viewpoint}\label{a-modified-equation-viewpoint}

Numerical analysts often ask a different question:

\begin{quote}
which nearby continuous equation is the discrete method solving more
accurately?
\end{quote}

This leads to modified-equation reasoning.

Lu, Zhong, Li, and Dong explicitly connect deep architectures with
numerical differential-equation discretizations and use numerical ideas
to motivate architecture design (Lu et al. 2018).

This perspective can be more realistic than insisting on one exact
underlying ODE.

A trained network may be profitably studied through a nearby effective
flow.

But the same boundary remains:

the modified equation is an analytical construction under declared
assumptions.

It is not automatically the network's unique hidden physical law.

\section{Residual networks and Neural ODEs are related but
different}\label{residual-networks-and-neural-odes-are-related-but-different}

The Architecture History chapter source-locks both residual networks and
Neural ODEs (He et al. 2016; Chen et al. 2018).

A residual network directly defines finitely many maps.

A Neural ODE directly defines a continuous vector field and asks a
numerical solver to construct the trajectory.

Those directions of definition differ.

For a ResNet:

\texttt{discrete\ architecture\ -\textgreater{}\ possible\ continuous\ interpretation}.

For a Neural ODE:

\texttt{continuous\ dynamics\ -\textgreater{}\ chosen\ numerical\ discretization\ at\ execution}.

The two meet through numerical analysis.

They should not be collapsed.

\section{The discretization can become part of the
model}\label{the-discretization-can-become-part-of-the-model}

In a classical numerical simulation, changing the solver while refining
the mesh is often intended to preserve the same underlying physical
model.

In learned systems, the situation can be more entangled.

The architecture, solver, parameterization, and training process may
adapt to one another.

A change of discretization may change:

\begin{itemize}
\tightlist
\item
  optimization geometry;
\item
  parameter sharing;
\item
  expressive bias;
\item
  compute cost;
\item
  numerical error;
\item
  learned representation.
\end{itemize}

So the phrase \textbf{discretization error} should be reserved for error
relative to a declared continuous reference.

Performance change after changing architecture is a broader empirical
phenomenon.

\section{Reversibility has three
meanings}\label{reversibility-has-three-meanings}

The word \textbf{reversible} causes another common collision.

We must distinguish:

\subsection{Map invertibility}\label{map-invertibility}

A layer map \texttt{Psi} is one-to-one and has an inverse.

\subsection{Computational
reversibility}\label{computational-reversibility}

Earlier states can be reconstructed from later states, perhaps to reduce
activation storage.

\subsection{Time reversibility of a numerical
method}\label{time-reversibility-of-a-numerical-method}

A one-step method is symmetric if

\texttt{Psi\_\{-h\}=Psi\_h\^{}\{-1\}}.

These properties can overlap.

They are not identical.

A layer can be invertible without being a symmetric integrator.

\section{Explicit Euler is invertible and not time
reversible}\label{explicit-euler-is-invertible-and-not-time-reversible}

Return to

\(x'=-x\).

Euler gives

\texttt{Psi\_h(x)=(1-h)x}.

For

\texttt{h!=1},

the map is invertible.

Its inverse is

\texttt{Psi\_h\^{}\{-1\}(x)=x/(1-h)}.

Now apply Euler with negative step:

\texttt{Psi\_\{-h\}(x)=(1+h)x}.

Compose:

\texttt{Psi\_\{-h\}(Psi\_h(x))\ =\ (1-h\^{}2)x}.

Except at \texttt{h=0}, this is not the identity.

So the same update formula run with \texttt{-h} does not undo the
forward step.

The method is not time reversible.

\section{Exact reversibility
witness}\label{exact-reversibility-witness}

Choose

\texttt{h=1/2}.

The forward step multiplies by

\texttt{1/2}.

Its true inverse multiplies by

\texttt{2}.

The negative-step Euler map multiplies by

\texttt{3/2}.

Forward followed by negative-step Euler gives

\texttt{3/4}.

That one line separates the concepts:

\begin{itemize}
\tightlist
\item
  invertible map: yes;
\item
  same method with negative step is inverse: no;
\item
  time-symmetric method: no.
\end{itemize}

The exact flow behaves differently:

\texttt{e\^{}\{h\}e\^{}\{-h\}=1}.

\section{Why this matters for reversible neural
networks}\label{why-this-matters-for-reversible-neural-networks}

A reversible neural architecture may guarantee that activations can be
reconstructed.

That is valuable for memory.

But it does not follow that the architecture is a reversible numerical
integrator for an underlying differential equation.

Conversely, a symmetric numerical integrator has a time-reversal
property tied to its step map.

Its software implementation may still use memory in an ordinary
irreversible way.

The adjective is the same.

The mathematical objects are different.

\section{Stability also has several
meanings}\label{stability-also-has-several-meanings}

The same linguistic problem appears with \textbf{stability}.

In this chapter we need at least:

\begin{itemize}
\tightlist
\item
  absolute stability of a numerical method;
\item
  perturbation stability of a discrete map;
\item
  Lyapunov stability of a dynamical system;
\item
  conditioning/sensitivity;
\item
  optimization stability;
\item
  robustness under data or distribution perturbation.
\end{itemize}

The Numerics chapter already warned that the scalar stability region is
not a universal nonlinear certificate.

NETNUM adds a second warning:

\begin{quote}
the word ``stable'' should not move between numerical analysis and
training dynamics without an explicit bridge.
\end{quote}

This is especially important because deep-learning papers often use one
term for several phenomena.

\section{Stiffness in networks needs a declared dynamical
problem}\label{stiffness-in-networks-needs-a-declared-dynamical-problem}

Can a neural network be stiff?

Sometimes the numerical analogy is useful.

If a declared continuous model has widely separated decay scales, an
explicit discretization may face severe step restrictions.

But ``the network has large Jacobian eigenvalues'' is not, by itself, a
complete stiffness statement.

Stiffness is partly problem-relative and method-relative (Hairer and
Wanner 1996).

To use the term rigorously we should identify:

\begin{itemize}
\tightlist
\item
  the continuous or effective dynamical problem;
\item
  the numerical method;
\item
  the scale separation;
\item
  the step restriction or solver difficulty it induces.
\end{itemize}

Otherwise \textbf{stiff} risks becoming a synonym for ``hard to train.''

\section{The numerical lens suggests architecture
questions}\label{the-numerical-lens-suggests-architecture-questions}

Once the bridge is explicit, numerical analysis supplies useful design
questions.

For a residual architecture, ask:

\begin{itemize}
\tightlist
\item
  Is the forward map in a stable regime?
\item
  Does the update approximate a declared flow?
\item
  What changes under depth refinement?
\item
  Is the residual scale acting like a step size?
\item
  Would an implicit update enlarge the stable region?
\item
  Would a symmetric or structure-preserving composition protect an
  invariant?
\item
  Does operator splitting isolate meaningful subdynamics?
\item
  Does adaptive depth correspond to local error control?
\end{itemize}

These questions do not answer themselves.

They create disciplined hypotheses.

That is more valuable than merely calling a network an ODE.

\section{What Haber--Ruthotto adds}\label{what-haberruthotto-adds}

Haber and Ruthotto develop deep architectures from the perspective of
nonlinear dynamical systems and discrete ODE stability (Haber and
Ruthotto 2018).

The important Atlas lesson is methodological.

Numerical analysis can be used not just to interpret an existing network
after the fact, but to constrain architecture design.

If a discrete forward model has undesirable stability properties, modify
the architecture.

This turns the numerical lens into engineering.

Still, source scope matters.

The paper motivates particular stable constructions.

It does not establish that every network instability is a numerical ODE
instability.

\section{What Lu et al.~adds}\label{what-lu-et-al.-adds}

Lu et al.~connect several deep architectures to numerical
differential-equation schemes and propose architectures inspired by
numerical methods (Lu et al. 2018).

This widens the design vocabulary.

A network block need not be compared only with forward Euler.

One can ask about analogues of:

\begin{itemize}
\tightlist
\item
  multistep methods;
\item
  higher-order schemes;
\item
  implicit methods;
\item
  splitting methods;
\item
  modified equations.
\end{itemize}

The Atlas uses this as a bridge, not as a license to label any
multi-branch network with a numerical-method name.

Structural correspondence must be shown.

\section{Architecture analogy versus numerical
theorem}\label{architecture-analogy-versus-numerical-theorem}

This distinction deserves a compact table.

{\def\LTcaptype{none} % do not increment counter
\begin{longtable}[]{@{}
  >{\raggedright\arraybackslash}p{(\linewidth - 2\tabcolsep) * \real{0.5000}}
  >{\raggedright\arraybackslash}p{(\linewidth - 2\tabcolsep) * \real{0.5000}}@{}}
\toprule\noalign{}
\begin{minipage}[b]{\linewidth}\raggedright
Statement
\end{minipage} & \begin{minipage}[b]{\linewidth}\raggedright
Status
\end{minipage} \\
\midrule\noalign{}
\endhead
\bottomrule\noalign{}
\endlastfoot
\(x_{k+1}=x_k+F_k(x_k)\) resembles Euler form & algebraic observation \\
\texttt{F\_k=h\_k\ f(t\_k,.)} for declared \texttt{f,h\_k} & modeling
declaration \\
local defect is \texttt{O(h\^{}2)} & numerical theorem under smoothness
assumptions \\
global error is \texttt{O(h)} & numerical theorem under
stability/regularity assumptions \\
deeper network has smaller error & not implied without refinement
semantics \\
residual Jacobian is \texttt{I+J\_F} & exact differential identity \\
stable scalar Euler factor & exact for declared linear mode \\
training is stable & separate empirical/mathematical claim \\
\end{longtable}
}

The point is not to weaken the analogy.

It is to give each statement the correct support route.

\section{A practical numerical-network
contract}\label{a-practical-numerical-network-contract}

Before using numerical language for a network, record:

{\def\LTcaptype{none} % do not increment counter
\begin{longtable}[]{@{}
  >{\raggedright\arraybackslash}p{(\linewidth - 2\tabcolsep) * \real{0.5000}}
  >{\raggedright\arraybackslash}p{(\linewidth - 2\tabcolsep) * \real{0.5000}}@{}}
\toprule\noalign{}
\begin{minipage}[b]{\linewidth}\raggedright
Field
\end{minipage} & \begin{minipage}[b]{\linewidth}\raggedright
Question
\end{minipage} \\
\midrule\noalign{}
\endhead
\bottomrule\noalign{}
\endlastfoot
reference flow & What continuous evolution is being approximated? \\
step map & What network block corresponds to one numerical step? \\
mesh & What are \texttt{t\_k} and \(h_k\)? \\
refinement & How is depth increased while preserving the same reference
problem? \\
error & What norm and exact reference define local/global error? \\
stability & Which perturbation and which stability notion are
intended? \\
reversibility & Invertibility, reconstruction, or numerical symmetry? \\
scope & Which claims are theorem, interpretation, or empirical
observation? \\
\end{longtable}
}

This contract prevents the numerical lens from silently changing
categories.

\section{Failure modes}\label{failure-modes-5}

\subsection{Residual-equals-Euler}\label{residual-equals-euler}

A skip connection is treated as proof that the network is Euler
discretization.

Failure: no continuous reference problem is specified.

\subsection{Depth-equals-time}\label{depth-equals-time}

Layer index is silently relabeled time.

Failure: no mesh or refinement relation exists.

\subsection{Small-residual-equals-small-error}\label{small-residual-equals-small-error}

A small residual update is treated as a small truncation defect.

Failure: no exact flow comparison is made.

\subsection{Stable-forward-equals-stable-training}\label{stable-forward-equals-stable-training}

A forward amplification argument is promoted into an optimizer theorem.

Failure: the objects differ.

\subsection{Invertible-equals-reversible-integrator}\label{invertible-equals-reversible-integrator}

An invertible layer is called time reversible.

Failure: \texttt{Psi\_\{-h\}=Psi\_h\^{}\{-1\}} was never checked.

\subsection{More-layers-equals-convergence}\label{more-layers-equals-convergence}

A deeper model is assumed closer to a continuous limit.

Failure: the target model and discretization family changed together.

\section{Downstream handoff: Adaptive
Depth}\label{downstream-handoff-adaptive-depth}

\textbf{Adaptive Depth as Error Control --- ATLAS-CH-ADAPTDEPTH-001} may
now assume:

\begin{itemize}
\tightlist
\item
  exact-flow versus step-map distinction;
\item
  local versus global error;
\item
  refinement-family discipline;
\item
  step-size versus residual-scale distinction;
\item
  scalar stability-function reasoning.
\end{itemize}

That chapter can ask whether learned stopping or variable compute can be
interpreted as adaptive time stepping or local error control.

It must establish that bridge itself.

A stopping score is not automatically a local truncation-error
estimator.

\section{Downstream handoff: Boundary
Probes}\label{downstream-handoff-boundary-probes}

\textbf{Boundary Probes --- ATLAS-CH-BOUNDARYPROBE-001} may now assume:

\begin{itemize}
\tightlist
\item
  one-step map \texttt{Psi\_k};
\item
  local Jacobian \texttt{J\_\{Psi\_k\}};
\item
  first-order perturbation propagation;
\item
  accumulated Jacobian products;
\item
  the distinction between local sensitivity and global dynamics.
\end{itemize}

This prepares the language needed for JVPs, VJPs, power iteration, and
interface sensitivity.

\section{What this chapter
establishes}\label{what-this-chapter-establishes-1}

The numerical interpretation of neural depth is strongest when it is
conditional and explicit.

A residual architecture can be read as a numerical scheme when the
continuous problem, step semantics, and refinement family are supplied.

Then numerical analysis contributes exact distinctions:

\begin{itemize}
\tightlist
\item
  local defect versus global error;
\item
  consistency versus stability versus convergence;
\item
  residual scale versus step size;
\item
  forward stability versus optimization stability;
\item
  invertibility versus time reversibility.
\end{itemize}

Without those declarations, the same vocabulary remains metaphor.

The goal is not to ban the metaphor.

It is to know when the metaphor has become mathematics.

\section{References used in this
chapter}\label{references-used-in-this-chapter-35}

\begin{itemize}
\tightlist
\item
  Iserles, \emph{A First Course in the Numerical Analysis of
  Differential Equations} (Iserles 2008).
\item
  Hairer and Wanner, \emph{Solving Ordinary Differential Equations II}
  (Hairer and Wanner 1996).
\item
  He et al., \emph{Deep Residual Learning for Image Recognition} (He et
  al. 2016).
\item
  Chen et al., \emph{Neural Ordinary Differential Equations} (Chen et
  al. 2018).
\item
  Haber and Ruthotto, \emph{Stable Architectures for Deep Neural
  Networks} (Haber and Ruthotto 2018).
\item
  Lu et al., \emph{Beyond Finite Layer Neural Networks: Bridging Deep
  Architectures and Numerical Differential Equations} (Lu et al. 2018).
\end{itemize}

Exact prerequisite identities and source-authority boundaries are
recorded in:

sources/source-locks/ATLAS-CH-NETNUM-001.yaml

\chapter{Adaptive Depth as Error
Control}\label{adaptive-depth-as-error-control-1}

\textbf{Epistemic status:} audited numerical/adaptive-depth substrate +
primary adaptive-computation sources + Atlas synthesis.\\
\textbf{Specification:}
manuscript/specifications/ATLAS-CH-ADAPTDEPTH-001.md\\
\textbf{Derivation packet:}
mathematics/derivations/ATLAS-CH-ADAPTDEPTH-001-DERIVATIONS.md\\
\textbf{Computational witness:}
mathematics/computational-witnesses/ATLAS-CW-ADAPTDEPTH-001.md\\
\textbf{Source lock:} sources/source-locks/ATLAS-CH-ADAPTDEPTH-001.yaml

Adaptive computation sounds like one idea.

It is not.

A numerical solver may adapt its step size because a local error
estimate is too large.

A recurrent network may run more steps because a learned halting unit
has not fired.

A residual model may spend more layers on one spatial position than
another.

A system may terminate because it exhausted its budget even though its
quality criterion was not met.

All of these change computation.

They do not carry the same semantics.

This chapter asks when adaptive depth deserves the stronger name
\textbf{error control}.

The answer is deliberately strict:

\begin{quote}
adaptive computation becomes error control only when the system names a
target error, an estimator or certified relation to that error, a
tolerance, and a controller whose stopping or step-size decision is
justified by that relation.
\end{quote}

Without those objects, the system may still be useful.

It is simply doing adaptive computation rather than certified error
control.

\section{The four control surfaces}\label{the-four-control-surfaces}

The audited Depth chapter separated computational depth from physical
time.

The audited Networks-as-Numerical-Schemes chapter separated residual
updates from literal numerical discretizations.

This chapter combines those boundaries.

Four objects recur.

\subsection{Execution depth}\label{execution-depth}

How many transformations are actually executed?

\[
z_0\to z_1\to\cdots\to z_\tau.
\]

The stopping depth \(\tau\) may depend on the input.

\subsection{Numerical step size}\label{numerical-step-size}

For a declared reference evolution,

\[
x_{n+1}
=
\Psi_{h_n}(x_n),
\]

the parameter \(h_n\) controls the numerical mesh.

\subsection{Local error estimate}\label{local-error-estimate}

A numerical controller may construct

\[
\widehat e_n
\]

to estimate a declared local truncation/error quantity.

\subsection{Learned halting signal}\label{learned-halting-signal}

A learned system may produce

\[
q_k
\]

and stop when \(q_k\) crosses a threshold.

The Atlas will not silently identify any pair of these objects.

\section{What adaptive depth already gives
us}\label{what-adaptive-depth-already-gives-us}

Graves's Adaptive Computation Time provides a differentiable mechanism
for learning how many recurrent computational steps to allocate before
producing an output (Graves 2016).

Figurnov et al.~extend adaptive computation into residual vision systems
where different spatial positions can receive different amounts of
computation (Figurnov et al. 2017).

These are important mechanisms.

They show that a network can learn to allocate compute nonuniformly.

They do not, by themselves, show that the learned halting signal
estimates:

\begin{itemize}
\tightlist
\item
  truncation error;
\item
  approximation error;
\item
  task loss;
\item
  probability of correctness;
\item
  distance to a fixed point;
\item
  distance to a target representation.
\end{itemize}

A learned stop signal can be useful without being an error estimate.

\section{Numerical error control is more
specific}\label{numerical-error-control-is-more-specific}

Now consider a declared numerical problem.

Let

\[
x_{n+1}
=
\Psi_{h_n}(x_n).
\]

A numerical error controller needs more than a variable \(h_n\).

It needs a target quantity.

For example:

\[
e_n
=
\text{local one-step error in a declared norm}.
\]

Because the exact local error is usually unavailable, the solver
constructs an estimator

\[
\widehat e_n.
\]

Then the controller compares it with a tolerance

\[
\tau_n.
\]

A basic accept/reject rule is

\[
\widehat e_n\le\tau_n.
\]

This is already much stronger than ``run until a score is small.''

The score has declared semantics.

\section{Embedded formulas}\label{embedded-formulas}

A common numerical pattern is to obtain two approximations from related
stage evaluations.

Write

\[
x_{n+1}^{[p]}
\]

and

\[
x_{n+1}^{[p+1]}.
\]

Then a local estimator may use

\[
\widehat e_n
=
\left\|
x_{n+1}^{[p+1]}
-
x_{n+1}^{[p]}
\right\|.
\]

Dormand and Prince develop a family of embedded Runge-Kutta formulas
with different orders (Dormand and Prince 1980).

The important structural point for this chapter is not the exact
Dormand-Prince coefficient table.

It is the control pattern:

\begin{itemize}
\tightlist
\item
  two approximations;
\item
  one step;
\item
  one difference;
\item
  one local error estimate;
\item
  one decision about accepting or shrinking the step.
\end{itemize}

This pattern gives the phrase ``error-controlled adaptation'' real
mathematical content.

\section{Estimator is not oracle}\label{estimator-is-not-oracle}

Even in numerical analysis, an embedded difference is not magic.

Its interpretation depends on:

\begin{itemize}
\tightlist
\item
  the declared differential equation;
\item
  the method;
\item
  smoothness;
\item
  order conditions;
\item
  the norm;
\item
  the local asymptotic regime.
\end{itemize}

The safe statement is:

\[
\text{embedded difference}
\to
\text{method-specific local error estimator}.
\]

The unsafe statement is:

\[
\text{embedded difference}
=
\text{exact total error}
\]

for every problem.

The finite witness below is chosen precisely because the equality
happens to be exact there.

\section{A one-step exact witness}\label{a-one-step-exact-witness}

Consider

\[
y'(t)=t,
\qquad
y(0)=0.
\]

The exact solution is

\[
y(t)=\frac{t^2}{2}.
\]

Take one step from \(t=0\) to \(t=h\).

Euler gives

\[
y_E=0.
\]

Why?

Because

\[
f(0,0)=0.
\]

Now use the explicit trapezoid / Heun update.

Its first slope is

\[
k_1=0.
\]

The Euler predictor is still zero.

The second slope is

\[
k_2=f(h,0)=h.
\]

Therefore

\[
y_H
=
\frac h2(k_1+k_2)
=
\frac{h^2}{2}.
\]

But the exact solution at \(h\) is also

\[
y(h)=\frac{h^2}{2}.
\]

Thus, for this witness,

\[
y_H=y(h).
\]

The pair difference is

\[
|y_H-y_E|
=
\frac{h^2}{2}.
\]

And that equals the exact Euler one-step error.

This is unusually clean.

It lets us examine the controller without hiding behind asymptotic
notation.

\section{Rejecting a step}\label{rejecting-a-step}

Set the tolerance to

\[
\tau=\frac18.
\]

Try

\[
h=1.
\]

The estimated error is

\[
\widehat e
=
\frac12.
\]

Therefore

\[
\widehat e>\tau.
\]

The step is rejected.

Notice what was declared:

\begin{itemize}
\tightlist
\item
  target: Euler one-step error in absolute value;
\item
  estimator: Euler/Heun difference;
\item
  tolerance: \(1/8\);
\item
  decision: reject if estimator exceeds tolerance.
\end{itemize}

That is an error-control interface.

\section{Shrinking the step}\label{shrinking-the-step}

For a low-order method of order \(p\), local truncation error scales
like

\[
C h^{p+1}
\]

in the asymptotic regime.

An idealized controller can use

\[
h_{\mathrm{new}}
=
h
\left(
\frac{\tau}{\widehat e}
\right)^{1/(p+1)}.
\]

For Euler,

\[
p=1.
\]

Using

\[
h=1,
\qquad
\tau=\frac18,
\qquad
\widehat e=\frac12,
\]

we obtain

\[
h_{\mathrm{new}}
=
\sqrt{
\frac{1/8}{1/2}
}
=
\frac12.
\]

Retry at \(h=1/2\).

Then

\[
\widehat e
=
\frac{(1/2)^2}{2}
=
\frac18.
\]

The retry lands exactly on tolerance.

Practical controllers use safety factors and clamps.

The witness omits them because the point is the semantics, not
production solver tuning.

\section{Error control is a feedback
loop}\label{error-control-is-a-feedback-loop}

The controller has the form

\[
\text{estimate}
\to
\text{compare}
\to
\text{accept/reject}
\to
\text{change future computation}.
\]

This is why adaptive numerical integration is conceptually useful for
adaptive neural computation.

Both are feedback systems.

But analogy is not identity.

The crucial difference is what the feedback signal means.

\section{A learned halting signal}\label{a-learned-halting-signal}

A learned computation may evolve

\[
z_{k+1}
=
F_k(z_k)
\]

and emit a score

\[
q_k.
\]

Given threshold \(\delta\) and budget cap \(K_{\max}\),

\[
\tau
=
\min
\left(
\{k\le K_{\max}:q_k\le\delta\}
\cup
\{K_{\max}\}
\right).
\]

This is a perfectly valid stopping rule.

It says nothing yet about whether

\[
q_k
\]

is:

\begin{itemize}
\tightlist
\item
  a probability;
\item
  a task-loss estimate;
\item
  a residual;
\item
  a confidence score;
\item
  a distance to convergence;
\item
  a numerical local-error estimate.
\end{itemize}

The semantics must come from somewhere else.

\section{Learned stopping is not numerical
stopping}\label{learned-stopping-is-not-numerical-stopping}

Graves ACT uses learned halting to allocate recurrent computation
(Graves 2016).

Its mechanism is not ``estimate a Runge-Kutta truncation error and
reduce the step.''

Figurnov et al.~allocate computation spatially in residual networks
(Figurnov et al. 2017).

Their mechanism is not ``refine a finite-element mesh where a PDE
residual is large.''

These systems can inspire numerical analogies.

The Atlas preserves the direction of implication:

\[
\text{adaptive computation}
\Rightarrow
\text{variable compute allocation}.
\]

It does not infer:

\[
\text{adaptive computation}
\Rightarrow
\text{numerical error control}.
\]

\section{Perfect ranking is still not
calibration}\label{perfect-ranking-is-still-not-calibration}

A common temptation is weaker.

Suppose the learned halting score tracks difficulty very well.

Perhaps smaller score always means smaller true error.

Is that enough?

No.

Consider the audited contraction from the Depth chapter:

\[
x_{k+1}
=
\frac{x_k+2}{2},
\qquad
x_0=0.
\]

The fixed point is 2.

Its exact error is

\[
E_k
=
|x_k-2|
=
2^{1-k}.
\]

Now define

\[
q_k=4^{-k}.
\]

Then

\[
q_k=\frac{E_k^2}{4}.
\]

The score is a strictly monotone function of the true error.

It ranks every depth perfectly.

But it is not on the same scale.

\section{An exact premature-halting
witness}\label{an-exact-premature-halting-witness}

At depth

\[
k=2,
\]

the score is

\[
q_2=\frac1{16}.
\]

The true error is

\[
E_2=\frac12.
\]

Suppose we want true error at most

\[
\frac1{16}.
\]

If we mistakenly use the same numerical threshold on the learned score,

\[
q_k\le\frac1{16},
\]

the system stops at \(k=2\).

But

\[
E_2=\frac12
\]

is eight times the requested tolerance.

This failure is not due to poor ranking.

The ranking is perfect.

The failure is magnitude calibration.

\section{The missing map}\label{the-missing-map}

In the witness, the relation is exactly

\[
E_k
=
2\sqrt{q_k}.
\]

Once that map is known, the halting score can support correct
thresholding.

For target error

\[
\varepsilon,
\]

we require

\[
2\sqrt{q_k}
\le
\varepsilon.
\]

So the correct score threshold is

\[
q_k
\le
\left(
\frac{\varepsilon}{2}
\right)^2.
\]

For

\[
\varepsilon=\frac1{16},
\]

this becomes

\[
q_k\le\frac1{1024}.
\]

The first matching depth is

\[
k=5,
\]

where

\[
E_5=\frac1{16}.
\]

The score has become an error-control signal only because the relation
to target error is declared.

\section{Rank quality versus magnitude
semantics}\label{rank-quality-versus-magnitude-semantics}

This distinction appears throughout machine learning.

A score can be excellent at:

\begin{itemize}
\tightlist
\item
  ranking examples by difficulty;
\item
  ranking candidate states by quality;
\item
  separating easy from hard inputs;
\item
  predicting which examples deserve more compute.
\end{itemize}

Those are useful capabilities.

They do not automatically answer:

\begin{quote}
What threshold guarantees error below epsilon?
\end{quote}

That question requires magnitude semantics.

The Atlas therefore distinguishes:

\[
\text{ranking}
\]

from

\[
\text{calibration}
\]

from

\[
\text{certified upper bound}.
\]

\section{Expected calibration may still be too
weak}\label{expected-calibration-may-still-be-too-weak}

Suppose a learned halting score satisfies

\[
\mathbb E[E_k\mid q_k=s]
=
g(s).
\]

That is a meaningful probabilistic calibration statement.

It does not imply

\[
E_k\le g(s)
\]

for every individual example.

If the controller requires a hard error guarantee, an expectation is not
enough.

The appropriate object might instead be:

\begin{itemize}
\tightlist
\item
  a high-probability bound;
\item
  a conformal upper bound;
\item
  a deterministic residual bound;
\item
  a certified interval;
\item
  a worst-case estimate.
\end{itemize}

The right guarantee depends on the application.

\section{Local error is not global
error}\label{local-error-is-not-global-error}

Numerical error control has its own version of this warning.

Suppose every accepted step satisfies

\[
\widehat e_n\le\tau_n.
\]

Does that immediately prove the final trajectory is within a desired
global error?

No.

Global error depends on:

\begin{itemize}
\tightlist
\item
  accumulation;
\item
  stability;
\item
  amplification;
\item
  the mesh sequence;
\item
  the problem dynamics;
\item
  method order;
\item
  regularity.
\end{itemize}

The NETNUM prerequisite already separated local and global error.

ADAPTDEPTH inherits that boundary intact.

\section{More depth is not always
refinement}\label{more-depth-is-not-always-refinement}

Suppose a network runs twice as many layers.

That does not automatically mean the numerical step size was halved.

For a numerical refinement statement, we need a fixed reference problem.

For example, if a horizon

\[
T
\]

is represented by steps

\[
h_0,\ldots,h_{N-1},
\]

then a fixed-horizon mesh satisfies

\[
\sum_{n=0}^{N-1}h_n=T.
\]

Increasing \(N\) can represent a finer mesh only if the \(h_n\) values
are changed accordingly.

If the network simply appends more transformations, it may be extending
computational time rather than refining the same interval.

This is exactly why Depth and NETNUM are both hard prerequisites.

\section{Adaptive depth versus adaptive step
size}\label{adaptive-depth-versus-adaptive-step-size}

The distinction can be written compactly.

\subsection{Adaptive execution depth}\label{adaptive-execution-depth}

\[
z_{k+1}=F_k(z_k),
\qquad
k<\tau(x).
\]

The input chooses how many stages execute.

\subsection{Adaptive time stepping}\label{adaptive-time-stepping}

\[
x_{n+1}=\Psi_{h_n}(x_n),
\]

where an error controller changes

\[
h_n
\]

for a declared reference evolution.

The first changes stage count.

The second changes a numerical mesh.

A system can combine them.

It must declare that combination explicitly.

\section{A learned error estimator}\label{a-learned-error-estimator}

There is a legitimate bridge between the two worlds.

Suppose a learned module predicts

\[
\widehat e_\theta(z_k)
\]

and training data provides a target numerical or task error.

Now the learned score is intended to estimate an error.

But intention is still not guarantee.

We should ask:

\begin{itemize}
\tightlist
\item
  what is the training target?
\item
  what norm?
\item
  what distribution?
\item
  what calibration?
\item
  what tail behavior?
\item
  what happens out of distribution?
\item
  is the estimate biased?
\item
  is an upper bound needed?
\item
  what does rejection do?
\end{itemize}

The bridge becomes mathematically meaningful only after these questions
are answered.

\section{Error control as an interface
contract}\label{error-control-as-an-interface-contract}

A complete adaptive error-control interface should declare:

\subsection{State}\label{state-2}

What state is being advanced?

\[
x_n
\quad\text{or}\quad
z_k.
\]

\subsection{Target error}\label{target-error}

What quantity should be small?

\[
E_n.
\]

\subsection{Estimator}\label{estimator}

What is actually observed?

\[
\widehat e_n.
\]

\subsection{Relation}\label{relation}

Why does the estimator say anything about the target error?

Examples:

\[
E_n
\le
B(\widehat e_n),
\]

or

\[
P(E_n\le B(\widehat e_n))
\ge
1-\alpha.
\]

\subsection{Tolerance}\label{tolerance}

What target is requested?

\[
\varepsilon.
\]

\subsection{Controller}\label{controller}

What action follows?

\begin{itemize}
\tightlist
\item
  stop;
\item
  continue;
\item
  shrink step;
\item
  grow step;
\item
  route elsewhere;
\item
  spend more compute.
\end{itemize}

\subsection{Budget}\label{budget}

What happens if the criterion never passes?

This interface is the central durable object of the chapter.

\section{Budget exhaustion remains a separate
state}\label{budget-exhaustion-remains-a-separate-state}

Suppose a system has maximum depth

\[
K_{\max}.
\]

If the error or halting criterion is not met before that cap, the output
may still be returned.

But the termination reason is

\[
\text{budget\_exhausted}.
\]

It is not

\[
\text{criterion\_met}.
\]

This distinction was repaired explicitly in AUDIT-018.

ADAPTDEPTH preserves it.

The same rule applies to numerical solvers.

A solver that runs out of permitted work has not thereby satisfied
tolerance.

\section{Error control is not compute
optimality}\label{error-control-is-not-compute-optimality}

A step-size controller may meet an accuracy target while using more
function evaluations than another method.

A halting policy may reduce FLOPs while increasing error.

A routing scheme may lower average cost while worsening tail latency.

Therefore two objectives remain separate:

\[
\text{error control}
\]

and

\[
\text{resource optimization}.
\]

They can be combined through a joint objective.

They should not be conflated.

\section{Spatially adaptive
computation}\label{spatially-adaptive-computation}

Figurnov et al.~show a system in which different image regions receive
different computation (Figurnov et al. 2017).

This is a useful reminder that adaptive depth need not be scalar.

We might write a local depth field

\[
\tau(i,j).
\]

That looks superficially like adaptive mesh refinement.

The resemblance is interesting.

It is not yet a PDE statement.

Adaptive mesh refinement requires:

\begin{itemize}
\tightlist
\item
  a spatial discretization;
\item
  a reference PDE or operator;
\item
  an error indicator/estimator;
\item
  refinement/coarsening semantics;
\item
  consistency across interfaces.
\end{itemize}

A spatial halting mask supplies none of these automatically.

\section{Neural error indicators}\label{neural-error-indicators}

A useful research direction is to learn an error indicator rather than
merely a halting probability.

For example, a model might predict:

\[
\widehat e_k
=
g_\theta(z_k,z_{k+1},r_k).
\]

The controller could then continue while

\[
\widehat e_k>\varepsilon.
\]

But the Atlas would still ask what supports the implication:

\[
\widehat e_k\le\varepsilon
\Rightarrow
E_k\le\varepsilon.
\]

Possible answers include:

\begin{itemize}
\tightlist
\item
  exact residual theory;
\item
  calibration on a declared distribution;
\item
  conformal upper bounds;
\item
  interval propagation;
\item
  theorem-backed a posteriori estimates.
\end{itemize}

The learned component does not erase the proof obligation.

\section{Residuals are promising but not
universal}\label{residuals-are-promising-but-not-universal}

Numerical methods often use residuals or embedded differences because
they have a known mathematical relation to error.

Neural systems also expose residual-like objects:

\begin{itemize}
\tightlist
\item
  state increments;
\item
  reconstruction residuals;
\item
  fixed-point residuals;
\item
  gradient norms;
\item
  disagreement scores.
\end{itemize}

These can be useful stopping signals.

But the relation

\[
\text{small residual}
\Rightarrow
\text{small target error}
\]

is problem dependent.

For an ill-conditioned or non-normal system, a small residual may
coexist with a larger solution error.

The error-control interface therefore requires the relation, not just
the residual.

\section{Fixed-point stopping}\label{fixed-point-stopping}

Suppose an equilibrium computation seeks

\[
z^*=F(z^*).
\]

A natural residual is

\[
r_k
=
\|F(z_k)-z_k\|.
\]

Stopping when

\[
r_k\le\varepsilon
\]

is a residual criterion.

To convert it into a state-error guarantee

\[
\|z_k-z^*\|\le C\varepsilon,
\]

we need additional structure.

For a contraction with factor \(L<1\), such relations can be derived.

Without them, the residual is not the state error.

This is another instance of the same principle.

\section{Halting as decision
making}\label{halting-as-decision-making}

A learned halting rule can also be viewed as a decision problem.

At depth \(k\), the system chooses between:

\begin{itemize}
\tightlist
\item
  stop now;
\item
  spend another unit of compute.
\end{itemize}

The decision trades expected improvement against compute cost.

That is a valid formulation.

It is not numerical error control unless the objective or constraint
explicitly references an error quantity with justified semantics.

The Atlas therefore allows adaptive depth to be interpreted in several
ways without collapsing them.

\section{A hierarchy of stopping
semantics}\label{a-hierarchy-of-stopping-semantics}

From weakest to strongest:

\subsection{Heuristic stop score}\label{heuristic-stop-score}

A learned or handcrafted number triggers stopping.

\subsection{Empirically predictive
score}\label{empirically-predictive-score}

The score correlates with future improvement or error on held-out data.

\subsection{Calibrated expected error}\label{calibrated-expected-error}

The score predicts average error under a declared distribution.

\subsection{High-probability error
control}\label{high-probability-error-control}

The score yields a probabilistic upper bound with stated coverage.

\subsection{Deterministic certified error
bound}\label{deterministic-certified-error-bound}

The score or residual implies a mathematical upper bound under declared
assumptions.

These levels should not be described with the same language.

\section{Why embedded pairs matter
conceptually}\label{why-embedded-pairs-matter-conceptually}

Embedded numerical formulas are useful beyond ODE software.

They illustrate a design pattern for adaptive intelligence:

\begin{quote}
create two computations whose discrepancy carries information about
whether more computation is warranted.
\end{quote}

The pattern might reappear as:

\begin{itemize}
\tightlist
\item
  shallow versus deep prediction disagreement;
\item
  coarse versus refined representation;
\item
  low-rank versus expanded solve;
\item
  cheap versus expensive model;
\item
  one-step versus multi-step estimate.
\end{itemize}

But the discrepancy becomes error control only when its relation to
target error is justified.

The numerical example teaches the grammar.

It does not certify every neural analogue.

\section{Error estimators can fail outside their
regime}\label{error-estimators-can-fail-outside-their-regime}

Numerical estimators themselves have domains of validity.

A step controller can struggle when:

\begin{itemize}
\tightlist
\item
  the problem is stiff;
\item
  smoothness assumptions fail;
\item
  the step is outside the asymptotic regime;
\item
  the estimator suffers cancellation;
\item
  the chosen norm misses a relevant component.
\end{itemize}

This matters for neural analogies.

Even a mathematically motivated error signal has to be checked against
the regime where its guarantee was derived.

\section{Adaptive computation and distribution
shift}\label{adaptive-computation-and-distribution-shift}

A learned halting model may be calibrated on one input distribution.

Deployment may shift.

Then the relation between score and error can change.

This parallels the Uncertainty chapter's warning:

\[
P\text{-calibration}
\not\Rightarrow
Q\text{-calibration}.
\]

ADAPTDEPTH therefore treats learned error control as
distribution-relative unless a stronger robust guarantee is established.

\section{Acceptance and rejection should be
explicit}\label{acceptance-and-rejection-should-be-explicit}

A controller is incomplete if it only says ``adapt.''

For each proposed computation, one should know:

\begin{itemize}
\tightlist
\item
  what was measured;
\item
  whether the criterion passed;
\item
  whether the state was accepted;
\item
  what changed after rejection;
\item
  whether computation was retried;
\item
  whether the budget was consumed;
\item
  whether the final output met tolerance.
\end{itemize}

This transaction-like view aligns numerical adaptation with the Atlas's
broader emphasis on explicit state and evidence.

\section{A minimal controller
grammar}\label{a-minimal-controller-grammar}

A generic adaptive controller can be written:

\begin{enumerate}
\def\labelenumi{\arabic{enumi}.}
\tightlist
\item
  propose an update;
\item
  compute an indicator;
\item
  map indicator to target error semantics;
\item
  compare against tolerance;
\item
  accept or reject;
\item
  change compute allocation;
\item
  stop on criterion or budget;
\item
  record the termination reason.
\end{enumerate}

This grammar covers:

\begin{itemize}
\tightlist
\item
  numerical step-size control;
\item
  recurrent halting;
\item
  iterative solver stopping;
\item
  conditional depth;
\item
  learned refinement.
\end{itemize}

The semantics differ at step 3.

That is where the important mathematics lives.

\section{Failure modes}\label{failure-modes-6}

The recurring category errors are now concrete.

\subsection{Halting score as error}\label{halting-score-as-error}

A learned number is thresholded and then described as tolerance
satisfaction without calibration.

\subsection{Rank as magnitude}\label{rank-as-magnitude}

A score orders easy and hard cases correctly, so its numerical value is
treated as an error bound.

\subsection{Local as global}\label{local-as-global}

A local estimator is small, so final accumulated error is claimed small
without stability analysis.

\subsection{Depth as mesh}\label{depth-as-mesh}

More layers are described as finer time discretization without a fixed
horizon and step-scale relation.

\subsection{Budget as convergence}\label{budget-as-convergence}

The system stops at \(K_{\max}\) and the output is labeled converged.

\subsection{Spatial depth as adaptive mesh
refinement}\label{spatial-depth-as-adaptive-mesh-refinement}

A learned spatial computation mask is described with PDE refinement
language without a declared discretized PDE.

\subsection{Error control as optimal
compute}\label{error-control-as-optimal-compute}

Meeting tolerance is confused with using the minimum possible resources.

\section{Research programme}\label{research-programme}

The synthesis suggests several precise research targets.

\subsection{Learned a posteriori
estimators}\label{learned-a-posteriori-estimators}

Train a model to estimate a mathematically defined error and then test
calibration and bounds.

\subsection{Two-resolution neural
computation}\label{two-resolution-neural-computation}

Construct cheap/high-resolution pairs whose discrepancy predicts the
value of more computation.

\subsection{Conformal stopping}\label{conformal-stopping}

Use distribution-free calibration machinery to turn empirical error
predictors into high-probability stopping rules under exchangeability
assumptions.

\subsection{Robust halting under
shift}\label{robust-halting-under-shift}

Test whether error/halting calibration survives changed data laws.

\subsection{Resource-aware error
control}\label{resource-aware-error-control}

Optimize compute subject to a declared error constraint instead of
mixing accuracy and cost in one opaque score.

Each target requires its own proof or evidence layer.

\section{Closing view}\label{closing-view-24}

Adaptive depth is most mathematically useful when it is treated as a
control problem rather than a slogan.

The system evolves a state.

It observes an indicator.

The indicator has declared semantics.

A controller compares that indicator with a target.

The system spends more or less computation accordingly.

The durable synthesis is:

\[
\boxed{
\text{adaptive error control}
=
\text{adaptive computation}
+
\text{target error}
+
\text{estimator/error relation}
+
\text{tolerance}
+
\text{controller}
+
\text{budget state}.
}
\]

Remove the estimator/error relation and we still have adaptive
computation.

We no longer have justified error control.

\section{References used in this
chapter}\label{references-used-in-this-chapter-36}

\begin{itemize}
\tightlist
\item
  Graves, Adaptive Computation Time for Recurrent Neural Networks
  (Graves 2016).
\item
  Figurnov et al., Spatially Adaptive Computation Time for Residual
  Networks (Figurnov et al. 2017).
\item
  Dormand and Prince, A family of embedded Runge-Kutta formulae (Dormand
  and Prince 1980).
\end{itemize}

Exact provenance and claim boundaries are locked in:

sources/source-locks/ATLAS-CH-ADAPTDEPTH-001.yaml

\chapter{Neural Krylov Transport}\label{neural-krylov-transport}

\textbf{Epistemic status:} audited classical
Krylov/numerical-linear-algebra substrate + audited
representation-transport substrate + Atlas-owned exact finite
local-solve witnesses.\\
\textbf{Specification:}
manuscript/specifications/ATLAS-CH-NEURALKRYLOV-001.md\\
\textbf{Derivation packet:}
mathematics/derivations/ATLAS-CH-NEURALKRYLOV-001-DERIVATIONS.md\\
\textbf{Computational witness:}
mathematics/computational-witnesses/ATLAS-CW-NEURALKRYLOV-001.md\\
\textbf{Source lock:}
sources/source-locks/ATLAS-CH-NEURALKRYLOV-001.yaml

A deep model already transports representations through a sequence of
maps.

Sometimes a useful local computation can be posed as a linear correction
problem:

\[
A_z\delta_z=b_z.
\]

That raises a natural question:

\begin{quote}
can a model use a few operator-generated directions to compute a better
representation correction without forming or solving a large dense
system?
\end{quote}

Classical Krylov methods provide one mathematical language for that
question.

But the analogy is useful only if its boundaries remain visible.

The governing rule is:

\[
\boxed{
\text{declare the linear operator first; inherit no nonlinear convergence theorem by analogy.}
}
\]

\section{What comes from classical Krylov
theory}\label{what-comes-from-classical-krylov-theory}

Krylov Subspaces and Iterative Solves already supplies:

\begin{itemize}
\tightlist
\item
  operator-generated trial spaces;
\item
  Arnoldi/Lanczos vocabulary;
\item
  projection and residual criteria;
\item
  conditioning boundaries;
\item
  left/right preconditioning semantics;
\item
  matrix-free operator access;
\item
  restart and finite-precision cautions.
\end{itemize}

Those objects remain classical linear algebra.

NEURALKRYLOV does not redefine them.

\section{What comes from representation
transport}\label{what-comes-from-representation-transport}

Representation as Transport already gives a separate interface:

\[
z_{k+1}=\Phi_k(z_k).
\]

A representation is a state.

A stage is a declared map.

A local Jacobian can describe differential behavior near a base point.

But the Jacobian is not the global nonlinear map.

NEURALKRYLOV joins these two languages only through a declared local
linear problem.

\section{The local-solve interface}\label{the-local-solve-interface}

At state \(z\), suppose we define:

\[
A_z\delta_z=b_z.
\]

The operator \(A_z\) might be:

\begin{itemize}
\tightlist
\item
  a local Jacobian-derived map;
\item
  a normal-equation operator;
\item
  a Hessian-like approximation;
\item
  a learned linear surrogate;
\item
  another explicitly declared linear operator.
\end{itemize}

Nothing in the phrase ``Neural Krylov'' identifies \(A_z\)
automatically.

The operator must be named.

\section{The representation update is
separate}\label{the-representation-update-is-separate}

Once an approximate correction:

\[
\delta_z^{(m)}
\]

has been computed, a representation block might apply:

\[
z^+=\Phi(z,\delta_z^{(m)}).
\]

The local solve and the representation update are different maps.

A good local solve does not automatically imply:

\begin{itemize}
\tightlist
\item
  a good representation;
\item
  a stable global trajectory;
\item
  a better task score.
\end{itemize}

Those are separate tests.

\section{A fixed left-preconditioned exact
witness}\label{a-fixed-left-preconditioned-exact-witness}

Use the local system:

\[
A=
\begin{pmatrix}
1&0\\
0&4
\end{pmatrix},
\qquad
b=
\begin{pmatrix}
1\\1
\end{pmatrix}.
\]

The exact solution is:

\[
x^\star=
\begin{pmatrix}
1\\1/4
\end{pmatrix}.
\]

Introduce fixed left preconditioner:

\[
M=
\begin{pmatrix}
1&0\\
0&2
\end{pmatrix}.
\]

Then:

\[
M^{-1}A=
\begin{pmatrix}
1&0\\
0&2
\end{pmatrix},
\]

and:

\[
M^{-1}b=
\begin{pmatrix}
1\\1/2
\end{pmatrix}.
\]

Write:

\[
B=M^{-1}A,
\qquad
c=M^{-1}b.
\]

\section{Why the preconditioner
matters}\label{why-the-preconditioner-matters}

The Krylov process now sees:

\[
Bx=c.
\]

That is not merely ``the same system with faster convergence.''

It is a transformed system with a transformed residual.

The original residual remains:

\[
r=b-Ax.
\]

The transformed residual is:

\[
\widehat r=c-Bx=M^{-1}r.
\]

These are different numerical objects.

\section{One operator-generated
direction}\label{one-operator-generated-direction}

With:

\[
x_0=0,
\]

the first Krylov space is:

\[
\mathcal K_1(B,c)=\operatorname{span}\{c\}.
\]

So every candidate has form:

\[
x=\alpha c.
\]

The best transformed-residual coefficient is:

\[
\boxed{
\alpha=\frac34.
}
\]

Hence:

\[
x_1=
\begin{pmatrix}
3/4\\
3/8
\end{pmatrix}.
\]

\section{The transformed residual is
nonzero}\label{the-transformed-residual-is-nonzero}

For \(x_1\):

\[
\widehat r_1
=
\begin{pmatrix}
1/4\\
-1/4
\end{pmatrix}.
\]

Thus:

\[
\boxed{
\|\widehat r_1\|_2^2=\frac18.
}
\]

The best one-dimensional trial-space point is still not the exact solve.

\section{The original residual is
different}\label{the-original-residual-is-different}

The original-system residual is:

\[
r_1
=
\begin{pmatrix}
1/4\\
-1/2
\end{pmatrix}.
\]

Thus:

\[
\boxed{
\|r_1\|_2^2=\frac5{16}.
}
\]

The same iterate has two different residual norms depending on which
system we evaluate.

This is not a contradiction.

It is preconditioning semantics.

\section{The true error is another
object}\label{the-true-error-is-another-object}

The true error is:

\[
e_1=x^\star-x_1
=
\begin{pmatrix}
1/4\\
-1/8
\end{pmatrix}.
\]

Hence:

\[
\boxed{
\|e_1\|_2^2=\frac5{64}.
}
\]

Now we already have three different quantities:

\[
\frac18,
\qquad
\frac5{16},
\qquad
\frac5{64}.
\]

They answer different questions.

\section{One more direction changes everything in this
toy}\label{one-more-direction-changes-everything-in-this-toy}

Compute:

\[
Bc=
\begin{pmatrix}
1\\1
\end{pmatrix}.
\]

The two generated directions are:

\[
c=
\begin{pmatrix}
1\\1/2
\end{pmatrix},
\qquad
Bc=
\begin{pmatrix}
1\\1
\end{pmatrix}.
\]

Their determinant is:

\[
\boxed{\frac12}.
\]

Therefore:

\[
\boxed{
\mathcal K_2(B,c)=\mathbb R^2.
}
\]

\section{The exact solution is in the second Krylov
space}\label{the-exact-solution-is-in-the-second-krylov-space}

Indeed:

\[
\boxed{
x^\star
=
\frac32c-\frac12Bc.
}
\]

So an exact residual-minimizing solve in \(\mathcal K_2\) can recover:

\[
x^\star.
\]

Then:

\[
\widehat r_2=0,
\qquad
r_2=0,
\qquad
e_2=0.
\]

\section{What the positive witness actually
shows}\label{what-the-positive-witness-actually-shows}

It shows one finite fact:

\begin{quote}
for this declared operator, right-hand side, and preconditioner, one
generated direction is insufficient while two are sufficient.
\end{quote}

That is a useful architecture motif.

It is not:

\begin{itemize}
\tightlist
\item
  a rate theorem;
\item
  a dimension-independent guarantee;
\item
  evidence that every neural representation has a two-step Krylov
  structure.
\end{itemize}

\section{Short horizon can be computationally
meaningful}\label{short-horizon-can-be-computationally-meaningful}

The representation-computation idea is not necessarily to solve a large
system to machine precision.

It may instead be:

\begin{itemize}
\tightlist
\item
  generate a few directions;
\item
  project a local problem;
\item
  obtain a useful correction;
\item
  stop because compute is bounded.
\end{itemize}

This makes horizon \(m\) an architectural resource.

But ``small \(m\)'' is not itself a quality certificate.

\section{A low-dimensional failure
control}\label{a-low-dimensional-failure-control}

Consider:

\[
B_{\rm bad}
=
\begin{pmatrix}
1&0\\
0&10
\end{pmatrix},
\qquad
c_{\rm bad}
=
\begin{pmatrix}
1\\1
\end{pmatrix}.
\]

Again:

\[
\mathcal K_1
=
\operatorname{span}\{c_{\rm bad}\}.
\]

The trial space is one-dimensional.

\section{The best one-dimensional residual can still be
poor}\label{the-best-one-dimensional-residual-can-still-be-poor}

The exact minimum-residual scalar is:

\[
\alpha_{\rm bad}
=
\frac{11}{101}.
\]

The residual becomes:

\[
r_{\rm bad}
=
\begin{pmatrix}
90/101\\
-9/101
\end{pmatrix}.
\]

Therefore:

\[
\boxed{
\|r_{\rm bad}\|_2^2
=
\frac{81}{101}.
}
\]

That is large.

\section{Dimension is not
convergence}\label{dimension-is-not-convergence}

Both positive and bad controls use a one-dimensional first trial space.

One gives a moderate approximation.

The other leaves a large residual.

Thus:

\[
\boxed{
\text{small trial-space dimension}
\not\Rightarrow
\text{small residual}.
}
\]

The operator geometry matters.

\section{Conditioning and spectral structure still
matter}\label{conditioning-and-spectral-structure-still-matter}

KRYLOV already established that convergence can depend on:

\begin{itemize}
\tightlist
\item
  spectrum;
\item
  field of values;
\item
  non-normality;
\item
  right-hand side;
\item
  polynomial approximation;
\item
  preconditioning;
\item
  finite precision.
\end{itemize}

A neural implementation does not make those dependencies disappear.

It may change the operator.

Then the analysis must change with it.

\section{Task quality can disagree with solve
quality}\label{task-quality-can-disagree-with-solve-quality}

Now define downstream readout:

\[
w=
\begin{pmatrix}
1\\2
\end{pmatrix}.
\]

For the exact solution:

\[
w^\top x^\star=\frac32.
\]

For the non-exact one-step iterate:

\[
w^\top x_1=\frac32.
\]

So:

\[
\boxed{
w^\top x_1=w^\top x^\star.
}
\]

\section{Yet the one-step solve is still
wrong}\label{yet-the-one-step-solve-is-still-wrong}

At the same time:

\[
r_1\neq0,
\]

and:

\[
e_1\neq0.
\]

So the downstream scalar is exact even though the local solve is not.

Therefore:

\[
\boxed{
\text{exact task readout}
\not\Rightarrow
\text{exact solve}.
}
\]

This is a minimal but important metric firewall.

\section{The converse also needs
proof}\label{the-converse-also-needs-proof}

A smaller residual does not automatically imply a better task metric.

A downstream task may:

\begin{itemize}
\tightlist
\item
  ignore the corrected direction;
\item
  amplify a tiny error direction;
\item
  saturate;
\item
  threshold;
\item
  use a nonlinear readout.
\end{itemize}

Therefore solve quality and task quality must be evaluated separately.

\section{Representation quality is another
layer}\label{representation-quality-is-another-layer}

Suppose \(x\) becomes a representation correction.

One can measure:

\begin{itemize}
\tightlist
\item
  Euclidean correction error;
\item
  angular error;
\item
  constraint violation;
\item
  retained information;
\item
  probe separability;
\item
  downstream task score.
\end{itemize}

None equals:

\[
\|b-Ax\|
\]

by definition.

A chapter that reports only residual has not yet measured representation
quality.

\section{The local Jacobian bridge}\label{the-local-jacobian-bridge}

Suppose a representation map is:

\[
z^+=\Phi(z).
\]

At base point \(z_0\), a local differential object is:

\[
J_\Phi(z_0).
\]

A local correction operator might be declared as:

\[
A_{z_0}=I-J_\Phi(z_0).
\]

Then a Krylov method can be built for \(A_{z_0}\).

But that is a Krylov problem for a local linear operator.

It is not a Krylov method for the entire nonlinear \(\Phi\).

\section{Local Jacobian is not global
transport}\label{local-jacobian-is-not-global-transport}

If the representation moves to \(z_1\), then:

\[
J_\Phi(z_1)
\]

can differ from:

\[
J_\Phi(z_0).
\]

If the operator changes at each stage, a sequence of generated
directions is not automatically:

\[
r,Ar,A^2r,\ldots
\]

for one fixed \(A\).

The classical object has changed.

\section{Learned preconditioners make the distinction
sharper}\label{learned-preconditioners-make-the-distinction-sharper}

Suppose:

\[
M=M_\theta(z).
\]

Then the preconditioner may depend on:

\begin{itemize}
\tightlist
\item
  state;
\item
  layer;
\item
  input;
\item
  training parameters.
\end{itemize}

A learned preconditioner can be a useful architecture.

But it is not automatically the fixed linear \(M\) of classical theory.

Any guarantee must match the actual semantics.

\section{Fixed left versus right versus
flexible}\label{fixed-left-versus-right-versus-flexible}

The exact witness is fixed left preconditioning:

\[
M^{-1}Ax=M^{-1}b.
\]

Right preconditioning would instead change variables.

Flexible preconditioning can use different preconditioners across
iterations.

Nonlinear preconditioning changes the mathematical problem further.

The word ``preconditioned'' is insufficient without the type.

\section{Matrix-free access fits the
architecture}\label{matrix-free-access-fits-the-architecture}

A neural system may never materialize:

\[
A_z.
\]

It may expose only:

\[
v\mapsto A_zv.
\]

For Jacobian-derived operators, this can be implemented by JVP/VJP-style
primitives.

That is compatible with classical operator-action semantics.

But the operator is still conceptually defined.

\section{Matrix-free does not mean
operator-free}\label{matrix-free-does-not-mean-operator-free}

The fact that a matrix is not explicitly stored does not license an
undefined operator.

For every Krylov claim, one should be able to answer:

\begin{quote}
what map is being applied to the vector?
\end{quote}

Without that answer, the subspace claim is underspecified.

\section{Residual and error need
conditioning}\label{residual-and-error-need-conditioning}

For an invertible linear problem:

\[
Ae=r.
\]

Therefore:

\[
e=A^{-1}r.
\]

Bounds such as:

\[
\|e\|
\le
\|A^{-1}\|\|r\|
\]

depend on the operator.

Residual is observable from \(A,b,x\).

True error requires the exact solution or a certificate.

They are not interchangeable.

\section{Preconditioned residual adds another
layer}\label{preconditioned-residual-adds-another-layer}

Under left preconditioning:

\[
\widehat r=M^{-1}r.
\]

A solver can minimize:

\[
\|\widehat r\|
\]

while the application ultimately cares about:

\[
\|r\|.
\]

Sometimes this is exactly the intended geometry.

Sometimes it is not.

The distinction should be explicit.

\section{Finite precision still exists inside neural
systems}\label{finite-precision-still-exists-inside-neural-systems}

If a block explicitly constructs orthogonal or nearly orthogonal
operator-generated directions, finite precision matters.

Possible failures include:

\begin{itemize}
\tightlist
\item
  loss of orthogonality;
\item
  direction collapse;
\item
  residual-estimate drift;
\item
  numerical dependence.
\end{itemize}

A learned implementation is still numerical computation.

\section{Restarting and truncation are architectural
choices}\label{restarting-and-truncation-are-architectural-choices}

A bounded-compute Neural Krylov block may stop after \(m\) directions.

It may also restart.

This controls:

\begin{itemize}
\tightlist
\item
  memory;
\item
  latency;
\item
  operator calls.
\end{itemize}

But it can throw away useful subspace information.

The exact positive witness already shows:

\[
m=1
\]

misses the exact solution, while:

\[
m=2
\]

captures it.

\section{A Neural Krylov block as an
interface}\label{a-neural-krylov-block-as-an-interface}

A bounded block can be specified as:

\begin{enumerate}
\def\labelenumi{\arabic{enumi}.}
\tightlist
\item
  receive \(z\);
\item
  define \(A_z\);
\item
  define \(b_z\);
\item
  define preconditioner \(M_z\), if used;
\item
  generate at most \(m\) directions;
\item
  solve or project in the trial space;
\item
  return correction \(\delta_z^{(m)}\);
\item
  update representation;
\item
  expose diagnostics.
\end{enumerate}

This is a precise architecture description.

\section{Diagnostics should remain
plural}\label{diagnostics-should-remain-plural}

A serious implementation should report, where available:

\begin{itemize}
\tightlist
\item
  transformed residual;
\item
  original residual;
\item
  estimated/known solve error;
\item
  representation change;
\item
  constraint violation;
\item
  downstream task metric;
\item
  operator/preconditioner condition diagnostics.
\end{itemize}

One scalar should not impersonate all of them.

\section{The right empirical
question}\label{the-right-empirical-question}

The most useful empirical question is not:

\begin{quote}
does this architecture ``look like Krylov''?
\end{quote}

It is:

\begin{quote}
for the exact operator and preconditioner declared by the architecture,
what does the short generated subspace buy us relative to a matched
compute baseline?
\end{quote}

That question can be falsified.

\section{The right theoretical
question}\label{the-right-theoretical-question}

For a learned architecture, the theoretical burden is:

\begin{itemize}
\tightlist
\item
  identify the operator family;
\item
  identify whether it is fixed or state-dependent;
\item
  identify the projection criterion;
\item
  identify the preconditioner class;
\item
  control nonlinearity/state drift;
\item
  state finite-precision assumptions;
\item
  connect solve quality to representation/task quality only with an
  explicit theorem.
\end{itemize}

Without these, classical convergence is only an analogy.

\section{Non-implications}\label{non-implications-1}

NEURALKRYLOV rejects:

\[
\text{low subspace dimension}
\not\Rightarrow
\text{small residual},
\]

\[
\text{small residual}
\not\Rightarrow
\text{small true error without conditioning information},
\]

\[
\text{small solve error}
\not\Rightarrow
\text{good representation/task quality},
\]

\[
\text{exact task readout}
\not\Rightarrow
\text{exact solve},
\]

\[
\text{local Jacobian Krylov step}
\not\Rightarrow
\text{global nonlinear convergence},
\]

\[
\text{representation transport}
\not\Rightarrow
\text{Krylov convergence theorem},
\]

and:

\[
\text{learned preconditioner}
\not\Rightarrow
\text{fixed-linear preconditioner guarantees}.
\]

\section{Atlas connections}\label{atlas-connections-17}

\textbf{Krylov Subspaces and Iterative Solves.}\\
Supplies the exact linear-algebra vocabulary and limits.

\textbf{Representation as Transport.}\\
Supplies the state-map/local-Jacobian interface while blocking
ODE/geodesic/Krylov overclaims.

\textbf{Adaptive depth.}\\
A future controller can treat horizon \(m\) as adaptive compute, but
must justify its stopping signal.

\textbf{Boundary probes.}\\
JVP/VJP and sensitivity machinery can expose matrix-free local operator
actions.

\textbf{Systems.}\\
Short-horizon iterative representation computation creates explicit
operator-call, memory, and latency tradeoffs.

\section{Closing view}\label{closing-view-25}

Neural Krylov Transport is not the claim that neural networks secretly
run GMRES.

It is a disciplined architecture pattern:

\[
\boxed{
\text{declared local operator}
+
\text{few operator-generated directions}
+
\text{declared projection}
+
\text{declared representation update}.
}
\]

The finite witness shows why this can matter.

One direction leaves nonzero residual and error.

Two directions recover the exact local solution.

A separate one-dimensional control remains poor.

And an exact downstream readout can coexist with a non-exact solve.

So the durable lesson is not ``Krylov always works.''

It is:

\[
\boxed{
\text{short-horizon operator-generated computation can be useful, but every guarantee is operator-, metric-, preconditioner-, and interface-relative.}
}
\]

\section{References used in this
chapter}\label{references-used-in-this-chapter-37}

No new external academic authority is added.

Classical numerical-linear-algebra authority is inherited through
audited ATLAS-CH-KRYLOV-001.

Representation-transport authority is inherited through audited
ATLAS-CH-TRANSPORT-001.

Exact prerequisite identities and claim boundaries are recorded in:

sources/source-locks/ATLAS-CH-NEURALKRYLOV-001.yaml

\chapter{Boundary Probes}\label{boundary-probes}

\textbf{Epistemic status:} established linear algebra + established
algorithmic differentiation + audited numerical-network sensitivity
substrate + Atlas synthesis + exact finite witness.\\
\textbf{Specification:}
manuscript/specifications/ATLAS-CH-BOUNDARYPROBE-001.md\\
\textbf{Formal packet:}
mathematics/derivations/ATLAS-CH-BOUNDARYPROBE-001-DERIVATIONS.md\\
\textbf{Computational witness:}
mathematics/computational-witnesses/ATLAS-CW-BOUNDARYPROBE-001.md\\
\textbf{Source lock:}
sources/source-locks/ATLAS-CH-BOUNDARYPROBE-001.yaml

A learned component can be shape-compatible with its neighbor and still
compose poorly.

One component may amplify a direction the next component cannot
tolerate. A downstream objective may be sensitive to one interface
coordinate. A finite probe may look benign because it did not explore
the dominant singular direction. A local derivative can look small while
the nonlinear map changes behavior farther away.

The governing rule is:

\begin{quote}
a boundary probe is evidence about a declared local interface question,
not a global certificate of the component or system.
\end{quote}

\section{Boundary object}\label{boundary-object}

Let a component be

\texttt{F:X-\textgreater{}Y}.

A later component may consume it:

\texttt{G:Y-\textgreater{}Z}.

At the interface, several obligations can coexist:

\begin{itemize}
\tightlist
\item
  type and shape compatibility;
\item
  numerical sensitivity;
\item
  semantic compatibility;
\item
  units and conventions;
\item
  provenance;
\item
  downstream task adequacy.
\end{itemize}

This chapter develops the numerical/local-sensitivity part only.

\section{Local linearization}\label{local-linearization}

Let

\texttt{F:R\^{}n-\textgreater{}R\^{}m}

be differentiable at operating point \texttt{x}.

Write:

\texttt{J=J\_F(x)}.

Then:

\texttt{F(x+delta)\ =\ F(x)+J\ delta+r(delta)}

with

\texttt{\textbar{}\textbar{}r(delta)\textbar{}\textbar{}/\textbar{}\textbar{}delta\textbar{}\textbar{}\ -\textgreater{}\ 0}

as \texttt{delta\ -\textgreater{}\ 0}.

LINALG-001 already established the key boundary: a Jacobian is local.

\section{Why derivative products
matter}\label{why-derivative-products-matter}

The full Jacobian can be too large to materialize.

Many interface questions require only products:

\begin{itemize}
\tightlist
\item
  response to one input perturbation;
\item
  pullback of one output objective;
\item
  dominant local Euclidean gain;
\item
  sensitivity of a selected interface mode.
\end{itemize}

This motivates JVPs and VJPs.

\section{JVP}\label{jvp}

For tangent direction

\texttt{v\ in\ R\^{}n},

define:

\texttt{JVP\_F(x;v)=J\_F(x)v}.

Equivalently:

\texttt{J\_F(x)v\ =\ d/d\ epsilon\ F(x+epsilon\ v)\textbar{}\_(epsilon=0)}.

A JVP pushes a local input perturbation forward.

Forward-mode algorithmic differentiation computes this product through
the computational graph without requiring a dense Jacobian (Griewank and
Walther 2008; Baydin et al. 2018).

\section{JVP versus finite
differences}\label{jvp-versus-finite-differences}

The finite difference

\texttt{{[}F(x+epsilon\ v)-F(x){]}/epsilon}

can approximate the JVP.

It is not the same computational object.

Finite differences depend on a chosen step and carry truncation and
roundoff effects. Algorithmic differentiation propagates derivative
rules through the represented computation, subject to ordinary numerical
arithmetic error.

\section{VJP}\label{vjp}

Under the declared Euclidean coordinate convention, represent the output
covector by

\texttt{w\ in\ R\^{}m}.

Use:

\texttt{VJP\_F(x;w)\ =\ J\_F(x)\^{}T\ w}.

If:

\texttt{ell(x)=w\^{}T\ F(x)},

then:

\texttt{grad\_x\ ell(x)=J\_F(x)\^{}T\ w}.

Reverse-mode algorithmic differentiation computes this pullback
(Griewank and Walther 2008; Baydin et al. 2018).

\section{Pairing identity}\label{pairing-identity}

For compatible \texttt{v} and \texttt{w}:

\texttt{w\^{}T(Jv)\ =\ (J\^{}T\ w)\^{}T\ v}.

This identity gives a compact consistency check between independently
implemented JVP and VJP routines.

Agreement is useful evidence.

It is not a proof of complete derivative correctness.

\section{Chain rule at an
interface}\label{chain-rule-at-an-interface}

For:

\texttt{x\ -\/-F-\/-\textgreater{}\ y\ -\/-G-\/-\textgreater{}\ z},

the local Jacobian is:

\texttt{J\_(G\ o\ F)(x)\ =\ J\_G(F(x))J\_F(x)}.

A JVP propagates forward through this product.

A VJP propagates backward through the transposed factors.

The products follow the chain rule without requiring all intermediate
Jacobians to be stored explicitly.

\section{Local Euclidean gain}\label{local-euclidean-gain}

LINALG-001 established:

\texttt{\textbar{}\textbar{}J\textbar{}\textbar{}\_2=sigma\_max(J)}.

Therefore:

\texttt{\textbar{}\textbar{}J\ delta\textbar{}\textbar{}\_2\ \textless{}=\ sigma\_max(J)\textbar{}\textbar{}delta\textbar{}\textbar{}\_2}.

The largest singular value is the largest first-order Euclidean
amplification factor at the operating point.

\section{Local does not mean
global}\label{local-does-not-mean-global}

The nonlinear map obeys:

\texttt{F(x+delta)-F(x)\ =\ J\ delta+r(delta)}.

So a known local operator norm controls the linear term.

A global Lipschitz claim needs additional evidence, such as a derivative
bound over a region or another nonlinear argument.

\section{Power iteration as a
probe}\label{power-iteration-as-a-probe}

Define:

\texttt{A=J\^{}T\ J}.

Its eigenvalues are:

\texttt{sigma\_i(J)\^{}2}.

Power iteration applies:

\[
u_{k+1}=Az_k,
\qquad
z_{k+1}=\frac{u_{k+1}}{\|u_{k+1}\|_2},
\]

\textbf{only when} \(\|Az_k\|_2>0\). If \(Az_k=0\), the normalized step
is undefined and a probe must report that breakdown rather than silently
divide by zero. For example, \(A=\operatorname{diag}(9,0)\) and
\(z_0=(0,1)^\top\) give \(Az_0=0\), even though the true dominant
singular value of the corresponding \(J=\operatorname{diag}(3,0)\) is
\(3\). A restart with a different direction may be needed.

For a nonzero probe direction, the Rayleigh quotient

\[
\rho_k=\frac{z_k^\top A z_k}{z_k^\top z_k}
\]

can approach the dominant eigenvalue under the stated spectral and
initialization conditions.

Then:

\texttt{sqrt(rho\_k)}

approaches the dominant singular value.

\section{JVP and VJP implement the power
operator}\label{jvp-and-vjp-implement-the-power-operator}

One application of \texttt{A} is:

\texttt{A\ z\ =\ J\^{}T(Jz)}.

Thus it can be evaluated as:

\begin{enumerate}
\def\labelenumi{\arabic{enumi}.}
\tightlist
\item
  JVP: \texttt{u=Jz};
\item
  VJP: \texttt{J\^{}T\ u}.
\end{enumerate}

This supports local spectral probing without explicit Jacobian
materialization.

\section{Estimator and quantity are different
objects}\label{estimator-and-quantity-are-different-objects}

\texttt{sigma\_max(J)} is a mathematical quantity.

Power iteration is an estimator procedure.

A finite run can fail to recover the dominant singular direction.

Therefore the procedure that generated an estimate is part of the
evidence.

\section{Initialization matters}\label{initialization-matters}

Let eigenpairs of \texttt{A} be:

\texttt{lambda\_1\ \textgreater{}\ lambda\_2\ \textgreater{}=\ ...}

with eigenvectors \texttt{q\_i}.

Write:

\texttt{z\_0=sum\_i\ c\_i\ q\_i}.

Then:

\texttt{A\^{}k\ z\_0\ =\ sum\_i\ c\_i\ lambda\_i\^{}k\ q\_i}.

If \texttt{c\_1} is nonzero, the dominant component can eventually
control the normalized iterate.

If \texttt{c\_1=0} exactly, ordinary power iteration cannot create the
missing component.

\section{Finite Rayleigh estimates can be
low}\label{finite-rayleigh-estimates-can-be-low}

For positive semidefinite \texttt{A}:

\texttt{rho(z)\ =\ (z\^{}T\ A\ z)/(z\^{}T\ z)\ \textless{}=\ lambda\_max(A)}.

Hence:

\texttt{sqrt(rho(z))\ \textless{}=\ sigma\_max(J)}.

A finite Rayleigh estimate is not automatically a certified upper bound
on the true local operator norm.

\section{Exact witness map}\label{exact-witness-map}

Use:

\texttt{F(x\_1,x\_2)\ =\ (x\_1\^{}2+x\_2,\ \ x\_1+2x\_2)}.

At:

\texttt{x\_0=(1,1)},

\texttt{F(x\_0)=(2,3)}.

The Jacobian is:

\texttt{J=\ {[}{[}2,1{]},\ \ {[}1,2{]}{]}}.

\section{Exact singular structure}\label{exact-singular-structure}

The eigenvectors of \texttt{J} are:

\texttt{q\_1=(1,1)/sqrt(2)}

and

\texttt{q\_2=(1,-1)/sqrt(2)}

with eigenvalues:

\texttt{3} and \texttt{1}.

The singular values are therefore:

\texttt{3,1}.

Thus:

\texttt{\textbar{}\textbar{}J\textbar{}\textbar{}\_2=3}.

\section{Exact JVP and nonlinear
remainder}\label{exact-jvp-and-nonlinear-remainder}

Choose:

\texttt{v=(1,1)}.

Then:

\texttt{Jv=(3,3)}

and the Euclidean gain is exactly \texttt{3}.

For scalar \texttt{epsilon}:

\texttt{F(x\_0+epsilon\ v)-F(x\_0)\ =\ epsilon(3,3)+(epsilon\^{}2,0)}.

The JVP is the exact first-order term. The quadratic remainder makes the
local/global distinction explicit.

\section{Exact VJP}\label{exact-vjp}

Choose:

\texttt{w=(1,2)}.

Then:

\texttt{J\^{}T\ w=(4,5)}.

For:

\texttt{ell(x)=w\^{}T\ F(x)},

the gradient at \texttt{x\_0} is exactly:

\texttt{(4,5)}.

The pairing check gives:

\texttt{w\^{}T(Jv)=9=(J\^{}T\ w)\^{}T\ v}.

\section{Exact power matrix}\label{exact-power-matrix}

For the witness:

\texttt{A=J\^{}T\ J\ =\ {[}{[}5,4{]},\ \ {[}4,5{]}{]}}.

Its eigenvalues are:

\texttt{9,1}.

The exact dominant singular target is therefore:

\texttt{sqrt(9)=3}.

\section{A mixed start converges}\label{a-mixed-start-converges}

Start with:

\texttt{z\_0=(1,0)}.

One multiplication gives:

\texttt{u\_1=(5,4)}.

The next gives:

\texttt{u\_2=(41,40)}.

The corresponding Rayleigh quotients are:

\texttt{rho\_1=365/41}

and

\texttt{rho\_2=29525/3281}.

Both approach the exact dominant eigenvalue \texttt{9} from below.

Their square roots approach the true singular norm \texttt{3}.

\section{The convergence has a closed
form}\label{the-convergence-has-a-closed-form}

Because:

\texttt{z\_0=(q\_1+q\_2)/sqrt(2)},

the unnormalized kth iterate is:

\texttt{A\^{}k\ z\_0\ =\ ((9\^{}k+1)/2,\ \ (9\^{}k-1)/2)}.

Its Rayleigh quotient is:

\texttt{rho\_k\ =\ (9*81\^{}k+1)/(81\^{}k+1)}.

Hence:

\texttt{rho\_k\ -\textgreater{}\ 9}.

This exact sequence makes the estimator dynamics visible without
numerical ambiguity.

\section{A different start misses the dominant
direction}\label{a-different-start-misses-the-dominant-direction}

Now choose:

\texttt{z\_bad=(1,-1)}.

Then:

\texttt{A\ z\_bad=z\_bad}.

Every iterate remains in the weak eigenspace.

The Rayleigh quotient is always:

\texttt{1}.

The inferred singular value is always:

\texttt{1}.

The true operator norm remains:

\texttt{3}.

Nothing about the component changed.

Only the probe initialization changed.

\section{One run is not a
certificate}\label{one-run-is-not-a-certificate}

Suppose a monitoring rule records whether the reported local gain is
below \texttt{2}.

The weak-eigenspace run reports \texttt{1}, while the true local norm is
\texttt{3}.

Therefore that finite estimate alone cannot establish the stronger
condition
\texttt{\textbar{}\textbar{}J\textbar{}\textbar{}\_2\ \textless{}=\ 2}.

The lesson is not that power iteration is unusable.

The lesson is that estimator metadata and limitations belong to the
result.

\section{Restarts and subspace
methods}\label{restarts-and-subspace-methods}

Random starts reduce the risk of exact orthogonality under ordinary
continuous sampling assumptions.

Multiple restarts improve directional coverage.

Subspace iteration can track several modes.

These are useful engineering choices.

They do not turn a local finite estimate into a global nonlinear
theorem.

\section{Probe metadata}\label{probe-metadata}

A reproducible spectral or sensitivity probe should record at least:

\begin{itemize}
\tightlist
\item
  exact component identity;
\item
  operating point or sampled region;
\item
  input/output norm convention;
\item
  tangent or cotangent initialization;
\item
  random seed;
\item
  iteration count;
\item
  restart count;
\item
  normalization rule;
\item
  stopping criterion;
\item
  Rayleigh or residual trace;
\item
  numerical precision;
\item
  runtime/software environment.
\end{itemize}

A scalar estimate without this context is weak evidence.

\section{Boundary-probe contract}\label{boundary-probe-contract}

Use:

\texttt{B=(F,X,Y,O,U,N\_X,N\_Y,P,E,tau)}.

Here:

\begin{itemize}
\tightlist
\item
  \texttt{F} is the component map;
\item
  \texttt{X,Y} are typed input/output spaces;
\item
  \texttt{O} is the operating point or region;
\item
  \texttt{U} is the perturbation or tangent class;
\item
  \texttt{N\_X,N\_Y} are norm or inner-product conventions;
\item
  \texttt{P} is the probe procedure;
\item
  \texttt{E} is estimator evidence and metadata;
\item
  \texttt{tau} is an acceptance threshold or policy.
\end{itemize}

This is the numerical-probe portion of a later full boundary contract.

\section{The perturbation class
matters}\label{the-perturbation-class-matters}

A full Euclidean operator norm ranges over every Euclidean direction.

A real interface may allow only structured perturbations.

Examples include:

\begin{itemize}
\tightlist
\item
  tangent directions on a constrained manifold;
\item
  physically admissible perturbations;
\item
  data-supported directions;
\item
  parameter-induced variations;
\item
  a full norm ball.
\end{itemize}

A sensitivity statement should name which class is being probed.

\section{Norm choice matters}\label{norm-choice-matters}

A gain measured in Euclidean norm need not have the same value in
another norm.

Interfaces may care about:

\begin{itemize}
\tightlist
\item
  Euclidean energy;
\item
  maximum coordinate error;
\item
  weighted physical units;
\item
  probability geometry;
\item
  task-specific seminorms.
\end{itemize}

The geometry belongs to the contract.

\section{Units precede numerical
norms}\label{units-precede-numerical-norms}

If different coordinates represent incompatible physical units, a raw
Euclidean norm can be meaningless.

Unit conversion or scaling conventions must be established before the
numerical gain is interpreted.

This is not something the Jacobian probe can infer.

It is a semantic interface obligation.

\section{Relation to NETNUM
sensitivity}\label{relation-to-netnum-sensitivity}

NETNUM-001 established that residual-step Jacobians propagate
first-order perturbations locally.

For:

\texttt{Psi(x)=x+F(x)},

the exact local Jacobian is:

\texttt{J\_Psi=I+J\_F}.

Boundary probes turn that local derivative into an inspectable interface
diagnostic.

The NETNUM boundary remains active:

\begin{quote}
local perturbation propagation is not a complete global stability
theorem.
\end{quote}

\section{Composite interfaces}\label{composite-interfaces}

For two local Jacobians \texttt{J\_1} and \texttt{J\_2}:

\texttt{J\_composite=J\_2\ J\_1}.

The submultiplicative norm bound gives:

\texttt{\textbar{}\textbar{}J\_2\ J\_1\textbar{}\textbar{}\_2\ \textless{}=\ \textbar{}\textbar{}J\_2\textbar{}\textbar{}\_2\ \textbar{}\textbar{}J\_1\textbar{}\textbar{}\_2}.

This can expose possible amplification chains.

The bound can be loose.

The product of individual worst-case gains need not equal the actual
composite gain.

\section{Small sensitivity is not semantic
adequacy}\label{small-sensitivity-is-not-semantic-adequacy}

Consider a component that maps every input to zero.

Its Jacobian norm is zero.

Numerically it is insensitive.

Semantically it may destroy all information needed downstream.

Therefore:

\begin{quote}
low sensitivity is not automatically good composition.
\end{quote}

Numerical compatibility and semantic adequacy are separate axes.

\section{Large sensitivity is not automatically
failure}\label{large-sensitivity-is-not-automatically-failure}

A component may intentionally amplify a weak signal.

If the downstream consumer expects that scale, a large singular value
can be acceptable.

The singular value is a diagnostic quantity.

Whether it violates a contract depends on:

\begin{itemize}
\tightlist
\item
  perturbation class;
\item
  units;
\item
  downstream tolerance;
\item
  application policy.
\end{itemize}

The threshold is not produced by the Jacobian itself.

\section{VJP magnitude is not causal
attribution}\label{vjp-magnitude-is-not-causal-attribution}

A large coordinate in:

\texttt{J\^{}T\ w}

means the chosen scalarized output is locally sensitive to that input
coordinate under the derivative model.

It does not by itself establish causal influence under interventions.

Sensitivity and causal explanation are different objects.

\section{Probe ensembles}\label{probe-ensembles}

A richer boundary profile may contain:

\begin{itemize}
\tightlist
\item
  several JVP directions;
\item
  several VJP objectives;
\item
  one or more singular-value estimates;
\item
  directional gains;
\item
  local invariant directions;
\item
  estimator residuals.
\end{itemize}

This can be more informative than a single scalar.

Each statistic still needs an explicit interpretation.

\section{Probe freshness}\label{probe-freshness}

Learned components change.

Fine-tuning, quantization, pruning, compilation, or kernel changes can
alter boundary behavior.

A probe result should therefore bind to the exact component artifact and
execution configuration.

API compatibility does not make an old probe current evidence for a
changed component.

\section{Operating-region coverage}\label{operating-region-coverage}

A derivative can be exactly computed at every sampled point while the
sampled region fails to represent later inputs.

This separates two questions:

\begin{itemize}
\tightlist
\item
  derivative/probe estimation;
\item
  operating-region coverage.
\end{itemize}

Power iteration addresses only the first, and only locally.

\section{Acceptance thresholds are
policy}\label{acceptance-thresholds-are-policy}

Suppose the measured gain is \texttt{1.8} and the threshold is
\texttt{2}.

The mathematics can establish the measurement and its procedure.

It cannot determine whether \texttt{2} is acceptable for the
application.

Threshold selection depends on downstream tolerance, uncertainty
margins, model role, and governance.

\section{Numerical versus semantic
conditions}\label{numerical-versus-semantic-conditions}

A numerical boundary probe can answer questions such as:

\begin{itemize}
\tightlist
\item
  which tangent directions amplify?
\item
  which input directions affect a chosen output objective?
\item
  what local singular gain is observed?
\item
  how stable is the spectral estimate across restarts?
\end{itemize}

It cannot by itself answer:

\begin{itemize}
\tightlist
\item
  do physical units agree?
\item
  does the output mean what the consumer assumes?
\item
  is provenance acceptable?
\item
  are business, legal, or scientific invariants preserved?
\item
  is the full nonlinear composition globally safe?
\end{itemize}

Those are separate contract dimensions.

\section{Failure modes}\label{failure-modes-7}

\subsection{Local-to-global promotion}\label{local-to-global-promotion}

A Jacobian norm at one point is presented as global stability.

\subsection{Estimate-to-certificate
promotion}\label{estimate-to-certificate-promotion}

A finite power estimate is presented as a proved upper bound.

\subsection{Initialization erasure}\label{initialization-erasure}

A spectral estimate is reported without its start or restart procedure.

\subsection{Norm erasure}\label{norm-erasure}

A gain is reported without the geometry in which it was measured.

\subsection{Semantic laundering}\label{semantic-laundering}

Low numerical sensitivity is treated as proof of interface meaning.

\subsection{VJP-as-causality}\label{vjp-as-causality}

A derivative pullback is presented as causal explanation.

\subsection{Stale evidence}\label{stale-evidence}

A probe from an earlier component version is reused after the component
changes.

\section{Practical probe ledger}\label{practical-probe-ledger}

{\def\LTcaptype{none} % do not increment counter
\begin{longtable}[]{@{}ll@{}}
\toprule\noalign{}
Field & Question \\
\midrule\noalign{}
\endhead
\bottomrule\noalign{}
\endlastfoot
component & Which exact artifact/version? \\
boundary & Which input and output spaces? \\
operating region & At which point(s) or data region? \\
perturbations & Which directions are admissible? \\
norms & Which input/output geometry? \\
JVP & Which tangent directions were pushed forward? \\
VJP & Which output covectors were pulled back? \\
spectral probe & Which power/subspace procedure? \\
initialization & Which seeds/vectors/restarts? \\
iterations & How many and with what stop rule? \\
estimate & Which Rayleigh/singular values were observed? \\
limitation & What estimator or coverage uncertainty remains? \\
threshold & Which policy determines acceptance? \\
semantic contract & Which non-numerical obligations remain separate? \\
\end{longtable}
}

This ledger converts \textbf{stable interface} into an inspectable
claim.

\section{What the witness
establishes}\label{what-the-witness-establishes}

The companion witness proves:

\begin{itemize}
\tightlist
\item
  exact Jacobian \texttt{{[}{[}2,1{]},{[}1,2{]}{]}};
\item
  exact singular values \texttt{3,1};
\item
  JVP \texttt{(3,3)} for \texttt{v=(1,1)};
\item
  VJP \texttt{(4,5)} for \texttt{w=(1,2)};
\item
  exact JVP/VJP pairing value \texttt{9};
\item
  nonlinear remainder \texttt{(epsilon\^{}2,0)} along the chosen
  tangent;
\item
  mixed-start Rayleigh values \texttt{365/41} and \texttt{29525/3281}
  approaching \texttt{9};
\item
  weak-eigenspace initialization returning Rayleigh value \texttt{1}
  forever despite true singular norm \texttt{3}.
\end{itemize}

No global nonlinear theorem is inferred from these finite facts.

\section{Downstream handoff}\label{downstream-handoff-4}

\textbf{Boundary Contracts --- ATLAS-CH-BCONTRACT-001} may now assume:

\begin{itemize}
\tightlist
\item
  JVP semantics;
\item
  VJP semantics;
\item
  local Jacobian composition;
\item
  Euclidean singular-gain interpretation;
\item
  power-iteration estimator behavior and initialization failure;
\item
  probe metadata requirements;
\item
  numerical-versus-semantic separation.
\end{itemize}

The downstream chapter must independently define:

\begin{itemize}
\tightlist
\item
  semantic interface conditions;
\item
  low-order separator variables;
\item
  admissibility obligations;
\item
  cross-component contract composition;
\item
  governance and evidence rules.
\end{itemize}

A boundary probe measures.

A boundary contract says what must hold.

\section{References used in this
chapter}\label{references-used-in-this-chapter-38}

\begin{itemize}
\tightlist
\item
  Griewank and Walther, \emph{Evaluating Derivatives: Principles and
  Techniques of Algorithmic Differentiation} (Griewank and Walther
  2008).
\item
  Baydin et al., \emph{Automatic Differentiation in Machine Learning: a
  Survey} (Baydin et al. 2018).
\item
  Trefethen and Bau, \emph{Numerical Linear Algebra} (Trefethen and
  David Bau 1997).
\item
  Golub and Van Loan, \emph{Matrix Computations} (Golub and Loan 2013).
\end{itemize}

Exact source identities and authority boundaries are recorded in:

sources/source-locks/ATLAS-CH-BOUNDARYPROBE-001.yaml

\chapter{Boundary Contracts}\label{boundary-contracts}

\textbf{Epistemic status:} mathematical exposition built from standard
sensitivity/conditioning and automatic-differentiation machinery,
local-to-global compatibility ideas, exact toy derivations, and a
source-bounded GCL project connection.\\
\textbf{Primary figure:} `ATLAS-FIG-BCONTRACT-001`\\
\textbf{Derivation packet:}
`mathematics/derivations/ATLAS-CH-BCONTRACT-001-DERIVATIONS.md`\\
\textbf{Computational witness:}
`mathematics/computational-witnesses/ATLAS-CW-BCONTRACT-001.md`

\section{Composition fails at
boundaries}\label{composition-fails-at-boundaries}

A complex learned system is rarely one undivided map.

It is assembled from components:

\[
x
\xrightarrow{f}
y
\xrightarrow{g}
z.
\]

The dimensions may match. The program may run. The result can still be
wrong.

The reason is simple: a boundary carries more than shape.

It carries meaning.

It carries admissible geometry.

It carries perturbations and gradients.

It carries numerical error.

It may also carry hidden assumptions about scale, normalization, units,
support, precision, or invariants.

The central question of this chapter is therefore:

\begin{quote}
What must one component expose at its boundary so that another component
can use it without reopening the entire interior?
\end{quote}

The Atlas answer is a \textbf{boundary contract}.

A boundary contract is not merely a software API and not automatically a
certificate. It is a structured statement of obligations that a
composition depends on.

\section{The engineered flange}\label{the-engineered-flange}

The useful allegory is an \textbf{engineered flange}.

Two physical machines may have matching bolt patterns and still be
unsafe to connect. One side may expect pressure the other cannot
tolerate. A load may arrive along the wrong axis. Two fluids may be
chemically incompatible. The dimensions fit while the interface
semantics fail.

The correspondence is:

\begin{itemize}
\tightlist
\item
  bolt geometry ↔ shape/type compatibility;
\item
  working medium ↔ semantic interpretation;
\item
  load direction ↔ perturbation direction;
\item
  rated load ↔ sensitivity bound;
\item
  tolerance ↔ numerical/error budget;
\item
  alignment constraint ↔ geometric invariant.
\end{itemize}

The limit is important.

A learned component is not a steel flange. Its behavior can depend on
state, data, adaptation, and training history. It need not obey fixed
constitutive laws.

The flange is useful only for one lesson:

\begin{quote}
compatibility is a bundle of obligations, not one matching dimension.
\end{quote}

Once that point is clear, the mathematics should take over.

\section{A provisional contract
object}\label{a-provisional-contract-object}

Let

\[
f:\mathcal X\to\mathcal Y.
\]

The Atlas uses the provisional explanatory object

\[
\boxed{
C_f=
(\mathcal X,\mathcal Y,\Sigma,\mathcal I,\mathcal S,\mathcal E).
}
\]

Its fields are:

\[
\mathcal X
=
\text{admissible input domain},
\]

\[
\mathcal Y
=
\text{admissible output domain},
\]

\[
\Sigma
=
\text{semantic declaration},
\]

\[
\mathcal I
=
\text{invariants and geometric obligations},
\]

\[
\mathcal S
=
\text{sensitivity and perturbation behavior},
\]

\[
\mathcal E
=
\text{numerical error and precision obligations}.
\]

This tuple is not proposed as a universal standard.

It is a disciplined way to ask what must cross a boundary besides bytes
and dimensions.

\section{Interface signature versus boundary
contract}\label{interface-signature-versus-boundary-contract}

An interface signature might say:

\[
f:\mathbb R^{768}\to\mathbb R^{768}.
\]

That is useful.

It tells us what shapes can be connected.

It does not tell us whether the output is:

\begin{itemize}
\tightlist
\item
  normalized;
\item
  centered;
\item
  a probability vector;
\item
  a logit vector;
\item
  an embedding whose norm carries confidence;
\item
  an embedding whose norm is intentionally discarded;
\item
  an element of a quotient class;
\item
  valid only in a bounded operating region.
\end{itemize}

A boundary contract adds these obligations explicitly.

The distinction can be written schematically as

\[
\boxed{
\text{signature}
\subset
\text{contract information}.
}
\]

A signature can be part of a contract.

It is not the whole contract.

\section{Shape-compatible, semantically
incompatible}\label{shape-compatible-semantically-incompatible}

Consider

\[
f_{\rm dir}(x)
=
\frac{x}{\|x\|_2},
\qquad
x\neq0.
\]

The output is a unit vector.

Suppose its intended semantics are:

\begin{quote}
direction only; amplitude has been intentionally discarded.
\end{quote}

Now define

\[
g_{\rm mag}(y)
=
\|y\|_2,
\]

and suppose the downstream component interprets magnitude as meaningful
amplitude.

The tensor shapes compose perfectly:

\[
\mathbb R^2
\to
\mathbb R^2
\to
\mathbb R.
\]

Take

\[
x=(3,4).
\]

Its amplitude is

\[
\|x\|_2=5.
\]

But

\[
g_{\rm mag}(f_{\rm dir}(x))
=
1.
\]

The software connection is legal.

The semantic connection is not.

This is the simplest reason a boundary contract must carry meaning, not
only type.

\section{Four obligation classes}\label{four-obligation-classes}

The Atlas separates four classes because they fail in different ways.

\subsection{Semantic obligations}\label{semantic-obligations}

A semantic obligation answers:

\begin{quote}
What does this state mean?
\end{quote}

Examples include:

\begin{itemize}
\tightlist
\item
  probability versus logit;
\item
  displacement versus absolute position;
\item
  direction versus magnitude-bearing vector;
\item
  uncertainty estimate versus raw score;
\item
  local coordinate versus globally identified state.
\end{itemize}

Semantic mismatch can survive every shape check.

\subsection{Geometric obligations}\label{geometric-obligations}

A geometric obligation answers:

\begin{quote}
What states are admissible, and what equivalences matter?
\end{quote}

Examples include:

\[
\|x\|=1,
\]

\[
X^\top X=I,
\]

simplex constraints,

quotient identifications,

positivity,

sparsity,

or a known gauge freedom.

The Geometry chapter showed that ambient coordinates and admissible
states need not be the same object.

A boundary that forgets this distinction can hand a downstream component
an invalid state even when every tensor dimension is correct.

\subsection{Differential
obligations}\label{differential-obligations}

A differential obligation answers:

\begin{quote}
How do perturbations and gradients cross this boundary?
\end{quote}

For a differentiable map

\[
y=f(x),
\]

the local derivative is

\[
J_f(x).
\]

Rather than materializing the full Jacobian, we often probe it through

\[
J_f(x)v
\]

and

\[
J_f(x)^\top u.
\]

These are Jacobian-vector products and vector-Jacobian products.
Automatic differentiation provides efficient machinery for computing
them in many practical systems (Baydin et al. 2018).

Differential obligations can include:

\begin{itemize}
\tightlist
\item
  directional gain;
\item
  dominant singular value estimates;
\item
  gradient transfer conventions;
\item
  null directions;
\item
  sensitivity to selected subspaces.
\end{itemize}

\subsection{Numerical obligations}\label{numerical-obligations}

A numerical obligation answers:

\begin{quote}
What error and conditioning assumptions are being passed across the
boundary?
\end{quote}

It may include:

\begin{itemize}
\tightlist
\item
  dtype or precision;
\item
  scaling convention;
\item
  absolute or relative error budget;
\item
  conditioning estimate;
\item
  tolerance;
\item
  solver residual;
\item
  stopping criterion;
\item
  stability region.
\end{itemize}

Numerical analysis has long distinguished problem conditioning from
algorithmic stability (Higham 2002).

A composition can fail because the mathematics is badly conditioned,
because an implementation is unstable, or because an interface discards
the information needed to tell which happened.

\section{The chain rule is the first composition
law}\label{the-chain-rule-is-the-first-composition-law}

Let

\[
y=f(x),
\]

and

\[
z=g(y).
\]

Then

\[
z=(g\circ f)(x).
\]

Locally,

\[
\boxed{
J_{g\circ f}(x)
=
J_g(f(x))J_f(x).
}
\]

This elementary identity is one of the most useful facts in the chapter.

It tells us that local perturbation behavior composes multiplicatively.

If

\[
\delta x
\]

is a small input perturbation, then

\[
\delta z
\approx
J_gJ_f\,\delta x.
\]

Using the induced \(2\)-norm,

\[
\|J_gJ_f\|_2
\le
\|J_g\|_2\|J_f\|_2.
\]

This is a local bound.

It is not automatically a global theorem.

\section{An exact two-module
witness}\label{an-exact-two-module-witness}

Take

\[
A=
\begin{pmatrix}
2&0\\
0&1/2
\end{pmatrix},
\qquad
B=
\begin{pmatrix}
1&1\\
0&2
\end{pmatrix}.
\]

Let

\[
f(x)=Ax,
\]

and

\[
g(y)=By.
\]

Then

\[
C=BA
=
\begin{pmatrix}
2&1/2\\
0&1
\end{pmatrix}.
\]

The composed local sensitivity is exactly the matrix \(C\).

For

\[
v=
\begin{pmatrix}
1\\
-1
\end{pmatrix},
\]

the forward directional sensitivity is

\[
Cv
=
\begin{pmatrix}
3/2\\
-1
\end{pmatrix}.
\]

For

\[
u=
\begin{pmatrix}
1\\
1
\end{pmatrix},
\]

the reverse sensitivity is

\[
C^\top u
=
\begin{pmatrix}
2\\
3/2
\end{pmatrix}.
\]

The same boundary can therefore be probed in the forward direction
through a JVP and in the reverse direction through a VJP.

\section{Exact local amplification}\label{exact-local-amplification}

For this \(C\),

\[
C^\top C
=
\begin{pmatrix}
4&1\\
1&5/4
\end{pmatrix}.
\]

Its eigenvalues are

\[
\lambda_\pm
=
\frac{21\pm\sqrt{185}}8.
\]

Therefore

\[
\boxed{
\|C\|_2
=
\sqrt{
\frac{21+\sqrt{185}}8
}
\approx
2.079707626949502.
}
\]

This number has a precise meaning.

Among unit perturbations, there is a direction whose local output
perturbation is amplified by about \(2.08\).

It does \textbf{not} mean:

\begin{itemize}
\tightlist
\item
  every perturbation grows by \(2.08\);
\item
  every state has the same Jacobian;
\item
  the nonlinear system is globally \(2.08\)-Lipschitz;
\item
  the composition is unsafe.
\end{itemize}

A sensitivity number is useful only when its scope remains attached to
it.

\section{The boundary figure}\label{the-boundary-figure}

\begin{figure}
\centering
\pandocbounded{\includegraphics[keepaspectratio,alt={A schematic upstream-to-downstream boundary contract listing semantic, geometric, differential, and numerical obligations beside the exact image of the unit perturbation circle under C equals {[}{[}2,1/2{]},{[}0,1{]}{]}.}]{figures/masters/ATLAS-FIG-BCONTRACT-001.png}}
\caption{A schematic upstream-to-downstream boundary contract listing
semantic, geometric, differential, and numerical obligations beside the
exact image of the unit perturbation circle under C equals
{[}{[}2,1/2{]},{[}0,1{]}{]}.}
\end{figure}

The left panel is schematic.

It names the obligation classes and shows forward JVP and reverse VJP
channels.

The right panel is literal for the toy system.

The dashed circle is the set of unit perturbations.

The solid ellipse is its image under

\[
C.
\]

The longest radius is the dominant singular value,

\[
2.0797\ldots
\]

The figure deliberately keeps the schematic contract language beside an
exact local sensitivity object.

One should not be mistaken for the other.

\section{Power iteration as a boundary
probe}\label{power-iteration-as-a-boundary-probe}

If the full Jacobian is too large to form, one may estimate a dominant
singular value using repeated applications of

\[
J
\]

and

\[
J^\top.
\]

For the toy map,

\[
M=C^\top C.
\]

Power iteration uses

\[
v_{k+1}
=
\frac{Mv_k}{\|Mv_k\|_2}.
\]

The Rayleigh quotient

\[
\mu_k
=
\frac{v_k^\top Mv_k}{v_k^\top v_k}
\]

approaches the dominant eigenvalue under standard convergence
conditions.

Then

\[
\sqrt{\mu_k}
\]

approaches the dominant singular value of \(C\).

Starting from the normalized direction proportional to \((1,1)\), the
witness gives

\[
1.9039433,
\]

\[
2.0701070,
\]

\[
2.0792647,
\]

\[
2.0796874,
\]

\[
2.0797067,
\]

and then agreement with the exact value to the displayed precision.

This is a good use of a probe because the answer is known independently.

\section{Estimation is not
certification}\label{estimation-is-not-certification}

Power iteration has assumptions.

Its convergence depends on the spectrum and the starting vector.

A finite iteration gives an estimate, not an oracle.

In a nonlinear component, the Jacobian may also change with the
operating point.

Therefore a contract should not silently transform

\begin{quote}
we measured a large singular direction near this state
\end{quote}

into

\begin{quote}
the component has a proven global gain bound.
\end{quote}

The first can be strong local evidence.

The second requires a much stronger argument.

This is where numerical conditioning discipline matters (Higham 2002).

\section{Separator variables}\label{separator-variables}

Large systems often cannot expose their entire internal state at every
interface.

They summarize.

Let

\[
s(y)
\]

be a low-dimensional boundary summary.

This can be valuable.

It can also erase exactly the information the next component needs.

Suppose

\[
s(y)=y_1.
\]

Take

\[
y=(0,1),
\qquad
y'=(0,2).
\]

Then

\[
s(y)=s(y')=0.
\]

But if downstream behavior depends on

\[
h(y)=y_2,
\]

then

\[
h(y)\neq h(y').
\]

The separator has identified two states that are operationally
different.

A low-order interface is therefore justified only relative to a
property:

\begin{quote}
sufficient for what?
\end{quote}

That question cannot be skipped.

\section{Local-to-global reasoning}\label{local-to-global-reasoning}

Boundary contracts are local objects.

They describe what neighboring components may assume about each other.

This naturally recalls local-to-global mathematics.

Sheaf theory provides a formal language for compatible local data and
gluing (Curry 2013). That language is useful because it teaches the
right kind of question:

\begin{quote}
if local pieces agree on overlaps, under what conditions do they
assemble into a coherent global object?
\end{quote}

The Atlas does not claim that learned-system contracts are automatically
sheaves.

Nor does local compatibility automatically yield:

\begin{itemize}
\tightlist
\item
  global stability;
\item
  global safety;
\item
  global optimality;
\item
  absence of feedback loops;
\item
  preservation of a system-level invariant.
\end{itemize}

The sheaf connection is a lens on compatibility.

It is not a free global theorem.

\section{Individually acceptable pieces can fail
together}\label{individually-acceptable-pieces-can-fail-together}

Suppose two components each have a finite local gain bound.

That does not settle the behavior of a large feedback system containing
them.

The problem can arise from:

\begin{itemize}
\tightlist
\item
  gain multiplication;
\item
  incompatible semantics;
\item
  hidden state;
\item
  delay;
\item
  approximation error;
\item
  non-normal coupling;
\item
  feedback;
\item
  state-dependent regime changes.
\end{itemize}

The chapter's local contract object is deliberately not named a
certificate.

It is an interface statement.

Whether the interface obligations are sufficient for a stronger global
property must be proved or tested separately.

\section{Public GCL project state}\label{public-gcl-project-state}

The Atlas design connects this chapter to the GCL MODULUS research
programme.

That connection must remain source-bounded.

The current public repository

`grandchallenge/MODULUS`

was inspected at commit

`9fc42eb5f29d5fff396f13e1a6c972af8fe64b35`.

It contains typed contract machinery, including the public file

`modulus/online/contracts.py`,

which defines a `RegretContract` and associated geometry, feedback,
comparator, loss, guarantee, constraint, and telemetry fields.

That is real public project state.

However, the remembered research names

\begin{itemize}
\tightlist
\item
  `BoundaryContract`;
\item
  `SeparatorCompiler`;
\item
  `Modula`
\end{itemize}

were not found in the inspected public tree or searchable commit
history.

The chapter therefore does \textbf{not} claim that those named
extensions are implemented in current public MODULUS.

This is an example of the source discipline the Atlas will later
formalize: project memory can motivate a question without being promoted
into documentary evidence.

\section{The GCL working slogan, properly
bounded}\label{the-gcl-working-slogan-properly-bounded}

The Atlas specification preserves the working slogan:

\begin{quote}
``Modula tells the optimizer how to move. Boundary contracts tell the
system how to compose.''
\end{quote}

It is useful rhetorically because it distinguishes two scales:

\begin{itemize}
\tightlist
\item
  an update rule governs motion;
\item
  an interface contract governs composition.
\end{itemize}

But the slogan is not a theorem.

The present chapter establishes only the bounded mathematical pieces
developed above:

\begin{itemize}
\tightlist
\item
  semantic obligations;
\item
  geometric obligations;
\item
  JVP/VJP sensitivity;
\item
  numerical conditioning;
\item
  separator sufficiency questions;
\item
  local composition bounds.
\end{itemize}

The stronger programme remains research.

\section{Six failure modes}\label{six-failure-modes-1}

\subsection{Shape without meaning}\label{shape-without-meaning}

Dimensions match while semantic interpretation differs.

\subsection{Geometry without
preservation}\label{geometry-without-preservation}

The downstream component receives a state outside its admissible set.

\subsection{Local sensitivity without global
scope}\label{local-sensitivity-without-global-scope}

A Jacobian norm measured at one state is treated as a global bound.

\subsection{Compressed boundary without
sufficiency}\label{compressed-boundary-without-sufficiency}

A separator erases state that changes downstream behavior.

\subsection{Numerical tolerance without
conditioning}\label{numerical-tolerance-without-conditioning}

An absolute tolerance is quoted without regard to problem scaling or
conditioning.

\subsection{Contract without
support}\label{contract-without-support}

An obligation is written down but not measured, derived, proved, or
otherwise supported.

A contract is only as strong as the obligations it records and the
evidence behind them.

\section{What a mature contract might
carry}\label{what-a-mature-contract-might-carry}

A future machine-readable boundary contract could record, for example:

\begin{itemize}
\tightlist
\item
  object identity;
\item
  input/output semantic types;
\item
  shape and dtype;
\item
  manifold or invariant constraints;
\item
  units or scale conventions;
\item
  local sensitivity probes;
\item
  norm conventions;
\item
  uncertainty or error budget;
\item
  approximation method;
\item
  separator variables;
\item
  operating region;
\item
  failure conditions;
\item
  evidence identity;
\item
  version identity.
\end{itemize}

The list is intentionally not canonical.

Different components need different obligations.

The purpose is to make hidden assumptions inspectable.

\section{Atlas connections}\label{atlas-connections-18}

\textbf{Geometry.}\\
A boundary can carry admissible state geometry, not merely coordinates.

\textbf{Numerical intelligence.}\\
Conditioning, residuals, and error budgets become interface obligations.

\textbf{Automatic differentiation.}\\
JVPs and VJPs provide practical directional probes (Baydin et al. 2018).

\textbf{Composition Without Catastrophe.}\\
The next step is to ask when local contracts are sufficient for
system-level reasoning.

\textbf{Governed adaptation.}\\
A self-changing component should not silently invalidate the contract
under which other components depend on it.

The recurring Atlas shift is:

\[
\boxed{
\text{architecture as connected modules}
\longrightarrow
\text{architecture as composition of explicit contracts}.
}
\]

\section{Closing view}\label{closing-view-26}

The most dangerous interface assumptions are often the ones nobody wrote
down.

A boundary contract makes them objects.

What does the state mean?

What geometry is valid?

What perturbations are amplified?

What numerical error is being inherited?

What summary has been retained, and what has been discarded?

The flange is only an allegory.

The contract is the object.

And a local contract is not the end of the argument.

It is the beginning of a compositional one.

\section{References used in this
chapter}\label{references-used-in-this-chapter-39}

\begin{itemize}
\tightlist
\item
  (Baydin et al. 2018)
\item
  (Higham 2002)
\item
  (Curry 2013)
\end{itemize}

See `sources/source-locks/ATLAS-CH-BCONTRACT-001.yaml` for exact source
identities, public GCL project-state boundaries, and claim scope.

\chapter{Composition Without
Catastrophe}\label{composition-without-catastrophe}

\textbf{Epistemic status:} audited Boundary Contracts + audited
Split-Operator Networks + Atlas-owned exact composition synthesis.\\
\textbf{Specification:}
manuscript/specifications/ATLAS-CH-COMPOSE-001.md\\
\textbf{Derivation packet:}
mathematics/derivations/ATLAS-CH-COMPOSE-001-DERIVATIONS.md\\
\textbf{Computational witness:}
mathematics/computational-witnesses/ATLAS-CW-COMPOSE-001.md\\
\textbf{Source lock:} sources/source-locks/ATLAS-CH-COMPOSE-001.yaml

A system can be built from components that are individually reasonable
and still fail as a whole.

That is not paradoxical.

The reason is that composition introduces new questions.

A component can satisfy its own input/output contract while a downstream
component assumes something different.

Two maps can each have modest local sensitivity while their product
violates a system budget.

The same two transformations can be safe in one order and unsafe in
another.

A local certificate can be correct and still be irrelevant outside its
operating region.

The governing rule of this chapter is:

\[
\boxed{
\text{component validity}
+
\text{component validity}
\not\Rightarrow
\text{system validity}.
}
\]

The missing object is the composition argument.

\section{Components do not compose by
optimism}\label{components-do-not-compose-by-optimism}

Suppose

\[
x\xrightarrow{f}y\xrightarrow{g}z.
\]

A naive implementation view asks:

\begin{quote}
do the tensor shapes match?
\end{quote}

A stronger contract view asks:

\begin{itemize}
\tightlist
\item
  does \(y\) mean what \(g\) expects?
\item
  is \(y\) inside the operating region required by \(g\)?
\item
  are invariants preserved?
\item
  are perturbation and error budgets compatible?
\item
  does ordering matter?
\item
  what happens to upstream error after downstream amplification?
\end{itemize}

These are different questions.

Passing one does not answer the others.

The audited Boundary Contracts chapter supplied the interface language.

The audited Split-Operator Networks chapter supplied the
order/noncommutativity language.

COMPOSE puts them together.

\section{A composition has obligations of its
own}\label{a-composition-has-obligations-of-its-own}

Write a provisional boundary contract for \(f\) as

\[
C_f
=
(\mathcal X,\mathcal Y,\Sigma_f,\mathcal I_f,\mathcal S_f,\mathcal E_f).
\]

Write one for \(g\) as

\[
C_g
=
(\mathcal Y,\mathcal Z,\Sigma_g,\mathcal I_g,\mathcal S_g,\mathcal E_g).
\]

The shared symbol \(\mathcal Y\) is not enough.

The connecting boundary also has to satisfy:

\begin{itemize}
\tightlist
\item
  semantic compatibility;
\item
  domain compatibility;
\item
  invariant compatibility;
\item
  sensitivity assumptions;
\item
  numerical/error accounting.
\end{itemize}

This chapter uses the explanatory notation

\[
C_f\triangleright C_g
\]

to mean that the exported guarantees of \(f\) satisfy the imported
assumptions of \(g\) for the declared composition.

This is Atlas notation.

It is not proposed as a universal contract standard.

\section{The first composition law is the chain
rule}\label{the-first-composition-law-is-the-chain-rule}

At a differentiable operating point,

\[
J_{g\circ f}(x)
=
J_g(f(x))J_f(x).
\]

Therefore

\[
\|J_{g\circ f}(x)\|_2
\le
\|J_g(f(x))\|_2
\|J_f(x)\|_2.
\]

This statement is elementary.

Its implications are not.

The local sensitivities multiply in an upper bound.

So a downstream stage can amplify upstream perturbation.

But the product bound can also be very loose.

The actual alignment of the singular directions matters.

This is already enough to show why component-local sensitivity numbers
are not a system certificate.

\section{The exact two-component
witness}\label{the-exact-two-component-witness}

Take

\[
A=
\begin{pmatrix}
3/2&0\\
0&2/3
\end{pmatrix},
\qquad
B=
\begin{pmatrix}
0&3/2\\
2/3&0
\end{pmatrix}.
\]

Both components have the same operator norm:

\[
\boxed{
\|A\|_2=\|B\|_2=\frac32.
}
\]

Suppose each component has a valid local contract allowing gain up to

\[
\frac32.
\]

Now impose a system-level gain budget

\[
\tau=2.
\]

Each component passes its local test.

The composition still has work to do.

\section{Chronological order
matters}\label{chronological-order-matters}

For column vectors, chronological \(A\) then \(B\) is represented by

\[
BA.
\]

Compute:

\[
BA
=
\begin{pmatrix}
0&1\\
1&0
\end{pmatrix}.
\]

This is an orthogonal swap.

Therefore

\[
\boxed{
\|BA\|_2=1.
}
\]

The system gain budget passes:

\[
1<2.
\]

Now reverse the chronology.

Chronological \(B\) then \(A\) is

\[
AB.
\]

Compute:

\[
AB
=
\begin{pmatrix}
0&9/4\\
4/9&0
\end{pmatrix}.
\]

Its singular values are

\[
\frac94,
\qquad
\frac49.
\]

Therefore

\[
\boxed{
\|AB\|_2=\frac94.
}
\]

Now the system gain budget fails:

\[
\frac94>2.
\]

Same components.

Same component-local contracts.

Different order.

Different system disposition.

\section{Product order is not chronological
prose}\label{product-order-is-not-chronological-prose}

The notation matters.

For column-vector states, matrix products act right-to-left.

So:

\begin{itemize}
\tightlist
\item
  chronological \(A\) then \(B\) corresponds to \(BA\);
\item
  chronological \(B\) then \(A\) corresponds to \(AB\).
\end{itemize}

This distinction was already repaired and frozen in the audited
Split-Operator chapter.

COMPOSE inherits it without ambiguity.

\section{The pair is genuinely
noncommuting}\label{the-pair-is-genuinely-noncommuting}

The difference is

\[
[A,B]
=
AB-BA.
\]

Exactly,

\[
[A,B]
=
\begin{pmatrix}
0&5/4\\
-5/9&0
\end{pmatrix}
\ne0.
\]

So the order effect is not a naming artifact.

The transformations genuinely do not commute.

This is the finite-dimensional version of a general warning:

\begin{quote}
when operations do not commute, ordering is part of the system.
\end{quote}

\section{The product norm bound can be
loose}\label{the-product-norm-bound-can-be-loose}

The component norms give the generic upper bound

\[
\|BA\|_2
\le
\|B\|_2\|A\|_2
=
\frac94.
\]

But the exact composition is

\[
\|BA\|_2=1.
\]

The upper bound overestimates the actual gain.

For the reverse product,

\[
\|AB\|_2=\frac94,
\]

so the same generic bound is attained exactly.

This gives a useful pair of facts:

\[
\boxed{
\text{component norm ceilings provide a bound, not the exact composition.}
}
\]

and

\[
\boxed{
\text{relative alignment can make the bound either loose or sharp.}
}
\]

\section{A commuting control}\label{a-commuting-control-1}

Now use

\[
A_c=
\operatorname{diag}(3/2,2/3),
\]

\[
B_c=
\operatorname{diag}(2/3,3/2).
\]

Each still has operator norm

\[
\frac32.
\]

But now

\[
[A_c,B_c]=0,
\]

and

\[
A_cB_c
=
B_cA_c
=
I.
\]

Therefore either order has gain

\[
\boxed{1.}
\]

This is the matched control.

The local component norm ceilings are unchanged.

What changed is the relationship between the transformations.

\section{The lesson of the control}\label{the-lesson-of-the-control}

The witness and control have the same individual gain ceilings.

They do not have the same composition.

Therefore a component-local statement such as

\[
\|J_f\|\le L_f
\]

does not contain enough information to reconstruct the exact composed
gain.

The interaction matters.

That interaction can involve:

\begin{itemize}
\tightlist
\item
  singular-vector alignment;
\item
  cancellation;
\item
  amplification;
\item
  noncommutativity;
\item
  invariant mismatch;
\item
  domain escape.
\end{itemize}

The contract needs to say which of these matter for the downstream
obligation.

\section{Error budgets also
compose}\label{error-budgets-also-compose}

Gain is only one failure mode.

Approximation error is another.

Suppose \(\widehat f\) approximates \(f\) and \(\widehat g\)
approximates \(g\).

Assume:

\[
\|\widehat f(x)-f(x)\|
\le
\varepsilon_f,
\]

and, on the connecting region,

\[
\|\widehat g(y)-g(y)\|
\le
\varepsilon_g.
\]

Suppose the true downstream map has Lipschitz bound

\[
\|g(y_1)-g(y_2)\|
\le
L_g\|y_1-y_2\|.
\]

Then:

\[
\begin{aligned}
\|\widehat g(\widehat f(x))-g(f(x))\|
&\le
\|\widehat g(\widehat f(x))-g(\widehat f(x))\|\\
&\quad+
\|g(\widehat f(x))-g(f(x))\|\\
&\le
\varepsilon_g+L_g\varepsilon_f.
\end{aligned}
\]

So

\[
\boxed{
E_{g\circ f}
\le
\varepsilon_g+L_g\varepsilon_f.
}
\]

The downstream stage amplifies the upstream error.

\section{The connecting domain is
load-bearing}\label{the-connecting-domain-is-load-bearing}

The previous derivation quietly used an important assumption:

\[
\widehat f(x)
\]

must remain inside the region where the downstream bounds apply.

If upstream approximation error pushes the state outside that domain,
then:

\begin{itemize}
\tightlist
\item
  the downstream Lipschitz bound may no longer hold;
\item
  the downstream implementation error bound may no longer hold;
\item
  the semantic interpretation may no longer be valid;
\item
  an invariant may be violated.
\end{itemize}

This is exactly why a boundary contract needs more than a scalar error
budget.

\section{A finite error-budget
failure}\label{a-finite-error-budget-failure}

Take the scalar maps

\[
f(x)=x,
\]

\[
g(y)=\frac32y.
\]

Suppose the local implementation errors satisfy

\[
|\widehat f(x)-f(x)|
\le
\frac1{10},
\]

and

\[
|\widehat g(y)-g(y)|
\le
\frac1{10}.
\]

The downstream gain is

\[
L_g=\frac32.
\]

The derived composition guarantee is

\[
E
\le
\frac1{10}
+
\frac32\frac1{10}
=
\frac14.
\]

Suppose the system requirement is

\[
E\le\frac15.
\]

Each component-local error ceiling is only

\[
\frac1{10}.
\]

Yet the derived system guarantee is

\[
\frac14>\frac15.
\]

The local budgets pass individually. \textbf{The upper bound alone would
show only that the system budget cannot yet be certified.} It would not
prove that a particular implementation violates the requirement.

For this witness, we can exhibit an implementation attaining that bound.
Choose

\[
\widehat f(x)=x+\frac1{10},
\qquad
\widehat g(y)=\frac32y+\frac1{10}.
\]

Each component has actual error exactly \(1/10\), on all real inputs, so
the local contracts are satisfied. But

\[
\begin{aligned}
\widehat g(\widehat f(x))-g(f(x))
&=\frac32\left(x+\frac1{10}\right)+\frac1{10}-\frac32x\\
&=\frac14
>\frac15.
\end{aligned}
\]

Thus \textbf{this explicit composition actually fails} the system error
requirement. The separate statement that the generic upper bound is
insufficient for certification remains true when no particular
implementations are supplied.

\section{This is not the same as
noncommutativity}\label{this-is-not-the-same-as-noncommutativity}

The scalar error-budget witness has no interesting order effect.

Its failure comes from amplification and accumulation.

The matrix witness has a strong order effect.

Its failure comes from noncommutativity and alignment.

These are different failure mechanisms.

They should not be collapsed into one generic word such as instability.

\section{Longer chains}\label{longer-chains}

Suppose a chain has stages

\[
f_1,f_2,\ldots,f_n.
\]

Let local approximation error at stage \(k\) be bounded by

\[
\varepsilon_k,
\]

and let the downstream true map at stage \(k\) have gain bound

\[
L_k
\]

on the connecting region.

If

\[
E_k
\le
L_kE_{k-1}+\varepsilon_k,
\]

with

\[
E_0=0,
\]

then induction gives

\[
\boxed{
E_n
\le
\sum_{j=1}^{n}
\varepsilon_j
\prod_{k=j+1}^{n}L_k.
}
\]

The formula makes one point visible:

\begin{quote}
upstream error is weighted by every downstream gain that follows it.
\end{quote}

\section{Uniform local budgets can still grow
badly}\label{uniform-local-budgets-can-still-grow-badly}

If every stage has

\[
L_k=L
\]

and

\[
\varepsilon_k=\varepsilon,
\]

then

\[
E_n
\le
\varepsilon
\sum_{r=0}^{n-1}L^r.
\]

If

\[
L=1,
\]

this becomes

\[
E_n\le n\varepsilon.
\]

If

\[
L>1,
\]

the worst-case bound grows geometrically:

\[
E_n
\le
\varepsilon
\frac{L^n-1}{L-1}.
\]

This does not mean the worst case is attained.

It means a local budget cannot be interpreted independently of chain
depth and downstream gains.

\section{Four levels of
certificate}\label{four-levels-of-certificate}

It helps to separate four evidence objects.

\subsection{Component certificate}\label{component-certificate}

A statement about one component.

For example:

\[
\|J_f(x_0)\|_2\le L_f.
\]

\subsection{Interface compatibility
certificate}\label{interface-compatibility-certificate}

A statement that the upstream output satisfies the downstream
assumptions.

For example:

\[
f(x_0)\in D_g.
\]

\subsection{Composition certificate}\label{composition-certificate}

A derived statement about the combined map.

For example:

\[
\|J_g(f(x_0))J_f(x_0)\|
\le
L_gL_f.
\]

\subsection{System-level guarantee}\label{system-level-guarantee}

A statement about the full reachable behavior of the assembled system.

For example:

\begin{quote}
every reachable state satisfies the declared safety property.
\end{quote}

These are not the same evidence level.

\section{Local does not mean
global}\label{local-does-not-mean-global-1}

A local Jacobian bound at one point,

\[
\|J_f(x_0)\|\le L,
\]

does not imply that \(f\) is globally \(L\)-Lipschitz.

The derivative can be larger elsewhere.

The operating region can change.

The system can leave the region.

A global conclusion requires a bound over the relevant region and
assumptions that connect those local bounds into a global result.

The audited Boundary Contracts chapter already enforced this firewall.

COMPOSE keeps it.

\section{Interface compatibility is not system
safety}\label{interface-compatibility-is-not-system-safety}

Suppose every adjacent pair of modules is interface-compatible.

That can still be insufficient for a full system guarantee.

Why?

Because a global property can depend on:

\begin{itemize}
\tightlist
\item
  long-horizon accumulation;
\item
  loops and feedback;
\item
  reachable-state geometry;
\item
  interaction between distant modules;
\item
  stochasticity;
\item
  adaptation;
\item
  distribution shift;
\item
  stateful memory.
\end{itemize}

Local compatibility is useful.

It is not magical.

\section{Shape compatibility is the weakest
case}\label{shape-compatibility-is-the-weakest-case}

The simplest failure is inherited from Boundary Contracts.

An upstream module can output a normalized direction.

A downstream module can interpret vector magnitude as meaningful
amplitude.

The shapes match.

The semantics do not.

The software graph composes.

The intended mathematical object does not.

This is the most basic reason the contract must carry meaning.

\section{Invariants can fail at the
boundary}\label{invariants-can-fail-at-the-boundary}

Suppose a downstream component assumes

\[
\|y\|=1.
\]

If the upstream component guarantees this exactly, the invariant
composes.

If it guarantees only

\[
|\|y\|-1|\le\delta,
\]

the downstream theorem may or may not survive.

That depends on the downstream assumptions.

The phrase approximately normalized is not a complete contract.

It needs a tolerance and a theorem that knows what to do with that
tolerance.

\section{Semantic mismatch and numerical error can
interact}\label{semantic-mismatch-and-numerical-error-can-interact}

A system can have both:

\begin{itemize}
\tightlist
\item
  a semantic mismatch;
\item
  a numerical error budget.
\end{itemize}

A small numerical error does not repair a wrong semantic interpretation.

Likewise, semantically aligned modules can still violate a numerical
budget.

The obligation classes remain separate because the remedies are
different.

\section{Order can change what the next module
sees}\label{order-can-change-what-the-next-module-sees}

In the exact matrix witness, \(B\) acts on a state already transformed
by \(A\), or vice versa.

That is the point.

Sequential composition is not a bag of independent ingredients.

The second stage sees the result of the first.

This is why order is architectural information.

The Split-Operator chapter made the same point using exact flows and
commutators (McLachlan and Quispel 2002).

COMPOSE turns it into a contract-level system-budget question.

\section{Classical splitting results stay
bounded}\label{classical-splitting-results-stay-bounded}

The existence of an exact matrix witness does not make an arbitrary
learned network an exact split flow.

A learned block can include:

\begin{itemize}
\tightlist
\item
  normalization;
\item
  masking;
\item
  stochasticity;
\item
  routing;
\item
  parameter variation;
\item
  data-dependent control flow.
\end{itemize}

Classical Lie or Strang order only transfers when the required
mathematical assumptions are actually established (McLachlan and Quispel
2002; Hairer et al. 2006).

The chapter uses the mathematics as a reference system, not as a free
theorem generator.

\section{JVP/VJP probes remain local
evidence}\label{jvpvjp-probes-remain-local-evidence}

Boundary Contracts inherited JVP and VJP machinery (Baydin et al. 2018).

These probes are useful for testing:

\begin{itemize}
\tightlist
\item
  directional gain;
\item
  gradient transfer;
\item
  dominant local sensitivity;
\item
  interface alignment.
\end{itemize}

They still report local differential behavior.

A successful local probe does not certify a whole state space.

This is an evidence boundary, not a criticism of the probe.

\section{Conditioning remains a separate
question}\label{conditioning-remains-a-separate-question}

A large observed error can come from:

\begin{itemize}
\tightlist
\item
  an ill-conditioned problem;
\item
  an unstable implementation;
\item
  upstream approximation error;
\item
  downstream amplification;
\item
  interface mismatch.
\end{itemize}

Numerical analysis keeps conditioning and algorithmic stability distinct
(Higham 2002).

COMPOSE should do the same.

The goal is not merely to detect failure.

It is to identify which obligation failed.

\section{Feedback changes the
problem}\label{feedback-changes-the-problem}

So far the chapter has mostly discussed feed-forward composition:

\[
x\to f(x)\to g(f(x)).
\]

Now suppose the composed output is fed back:

\[
x_{t+1}
=
(g\circ f)(x_t).
\]

A one-pass composition certificate does not automatically answer:

\begin{itemize}
\tightlist
\item
  whether trajectories remain bounded;
\item
  whether they converge;
\item
  whether transient growth occurs;
\item
  whether invariant sets survive;
\item
  whether noise accumulates.
\end{itemize}

These are dynamical questions.

COMPOSE marks the boundary and stops.

A full closed-loop theorem would require additional dependencies and
source authority.

\section{Repeated local validity can still miss
reachability}\label{repeated-local-validity-can-still-miss-reachability}

A component may be certified on an operating region

\[
D.
\]

If the composed system leaves \(D\), the certificate has expired.

This creates a subtle failure mode:

\begin{enumerate}
\def\labelenumi{\arabic{enumi}.}
\tightlist
\item
  every local theorem is correct;
\item
  every theorem is scoped;
\item
  the assembled trajectory exits the scope;
\item
  the system becomes uncertified.
\end{enumerate}

Nothing contradictory happened.

The composition argument failed to prove reachability containment.

\section{Budget accounting should be
explicit}\label{budget-accounting-should-be-explicit}

A system contract should say which budget is being composed.

Examples include:

\begin{itemize}
\tightlist
\item
  perturbation gain;
\item
  absolute numerical error;
\item
  relative error;
\item
  probability mass leakage;
\item
  constraint violation;
\item
  latency;
\item
  memory;
\item
  energy;
\item
  uncertainty;
\item
  privacy loss.
\end{itemize}

Different budgets obey different composition laws.

COMPOSE derives only the laws it states.

It does not assume every budget adds or multiplies in the same way.

\section{Conservative bounds are not
failures}\label{conservative-bounds-are-not-failures}

The product norm bound

\[
\|J_gJ_f\|
\le
\|J_g\|\|J_f\|
\]

can be loose.

That does not make it wrong.

A conservative bound may still be useful for screening.

But if the conservative bound violates a threshold, one should not
immediately conclude that the real composition violates it.

The exact \(BA\) witness demonstrates this:

\[
\|B\|\|A\|=\frac94,
\]

while

\[
\|BA\|=1.
\]

The right conclusion is:

\begin{quote}
the coarse certificate is inconclusive for that threshold.
\end{quote}

\section{Tightness can be
order-dependent}\label{tightness-can-be-order-dependent}

The reverse product \(AB\) reaches the full product bound:

\[
\|AB\|
=
\|A\|\|B\|
=
\frac94.
\]

So a conservative certificate can be loose in one ordering and exact in
another.

That is a useful diagnostic fact.

It means a system designer should not treat component-local norm
ceilings as a complete description of compositional geometry.

\section{A practical composition
protocol}\label{a-practical-composition-protocol}

For every boundary \(f\to g\), record:

\begin{enumerate}
\def\labelenumi{\arabic{enumi}.}
\tightlist
\item
  operating domain;
\item
  semantic meaning of the connecting state;
\item
  invariants;
\item
  local sensitivity evidence;
\item
  numerical/error budget;
\item
  precision/unit convention;
\item
  downstream assumptions.
\end{enumerate}

Then for the composition:

\begin{enumerate}
\def\labelenumi{\arabic{enumi}.}
\tightlist
\item
  verify interface compatibility;
\item
  derive the composition rule for the relevant budget;
\item
  test order if operations may not commute;
\item
  compare the derived bound with the system budget;
\item
  check whether the composed state remains in the region where the local
  contracts apply;
\item
  use an exact or higher-fidelity computation when a conservative bound
  is inconclusive.
\end{enumerate}

\section{A practical order test}\label{a-practical-order-test}

When two modules may be reordered, compare:

\[
g\circ f
\]

with

\[
f\circ g.
\]

Possible diagnostics include:

\begin{itemize}
\tightlist
\item
  exact finite witness;
\item
  output difference;
\item
  commutator in a declared linear reference model;
\item
  Jacobian commutator as a local diagnostic;
\item
  task-level difference.
\end{itemize}

Do not silently identify a same-state Jacobian commutator with the
derivative of the full nonlinear composition defect.

That distinction is inherited from the audited Split chapter.

\section{A practical error test}\label{a-practical-error-test}

For a two-stage approximation, report:

\[
\varepsilon_f,
\qquad
\varepsilon_g,
\qquad
L_g,
\]

and the derived composition budget:

\[
\varepsilon_g+L_g\varepsilon_f.
\]

Also report the connecting domain on which those quantities are valid.

Without the domain, the number is incomplete.

\section{A practical chain test}\label{a-practical-chain-test}

For a long chain, report the recurrence:

\[
E_k\le L_kE_{k-1}+\varepsilon_k.
\]

Do not only report the individual \(\varepsilon_k\).

The downstream gains determine how strongly early errors matter.

This is the composition analogue of keeping provenance attached to
evidence.

\section{What counts as catastrophe
here?}\label{what-counts-as-catastrophe-here}

The title is broader than any one metric.

A catastrophe can mean:

\begin{itemize}
\tightlist
\item
  a violated numerical tolerance;
\item
  a broken invariant;
\item
  a semantic mismatch;
\item
  a system gain over budget;
\item
  a state leaving the certified region;
\item
  order-dependent behavior not represented in the design;
\item
  accumulated error beyond specification.
\end{itemize}

The chapter does not define a universal catastrophe metric.

It requires the system-level obligation to be declared.

\section{What this chapter proves}\label{what-this-chapter-proves}

The exact finite witness proves:

\[
\|A\|=\|B\|=\frac32,
\]

\[
\|BA\|=1,
\]

\[
\|AB\|=\frac94.
\]

It proves that the same components pass a gain budget of \(2\) in one
order and fail it in the reverse order.

The commuting control proves:

\[
A_cB_c=B_cA_c=I.
\]

The error derivation proves, under its assumptions:

\[
E_{g\circ f}
\le
\varepsilon_g+L_g\varepsilon_f.
\]

The chain derivation proves, under its assumptions:

\[
E_n
\le
\sum_{j=1}^{n}
\varepsilon_j
\prod_{k=j+1}^{n}L_k.
\]

These are the load-bearing results.

\section{What this chapter does not
prove}\label{what-this-chapter-does-not-prove}

It does not prove:

\[
\text{all local contracts pass}
\Rightarrow
\text{global safety}.
\]

It does not prove that learned neural modules are exact flows.

It does not prove that a local Jacobian bound is a global Lipschitz
theorem.

It does not prove that commuting linearizations imply globally commuting
nonlinear maps.

It does not prove long-horizon stability.

It does not prove a universal contract calculus.

\section{Non-implications}\label{non-implications-2}

The chapter therefore rejects:

\[
\text{component contracts pass}
\not\Rightarrow
\text{system contract passes},
\]

\[
\|J_f(x_0)\|\le L
\not\Rightarrow
\text{global }L\text{-Lipschitz},
\]

\[
\|J_gJ_f\|\le\|J_g\|\|J_f\|
\not\Rightarrow
\text{tight estimate},
\]

\[
[A,B]=0\text{ in a reference model}
\not\Rightarrow
\text{learned modules commute globally},
\]

\[
\text{local error budget}
\not\Rightarrow
\text{long-horizon stability},
\]

and

\[
\text{interface compatibility}
\not\Rightarrow
\text{closed-loop safety}.
\]

\section{Atlas connections}\label{atlas-connections-19}

\textbf{Boundary Contracts.}\\
COMPOSE consumes the obligation language and asks whether the guarantees
actually survive assembly.

\textbf{Split-Operator Networks.}\\
COMPOSE consumes order/noncommutativity discipline and turns it into
explicit system-budget tests.

\textbf{Adaptive depth.}\\
A variable-depth system changes the number of times local gain/error
budgets compose.

\textbf{Routing and sparse systems.}\\
Different routes induce different composition chains and therefore
different contract obligations.

\textbf{Agents and tools.}\\
A tool interface can be type-correct while semantically or numerically
incompatible with the agent's assumptions.

\textbf{Governed adaptation.}\\
A self-changing system must know which local contract changes invalidate
downstream composition arguments.

\section{Closing view}\label{closing-view-27}

A component can be correct.

Another component can be correct.

Their composition can still be wrong.

Sometimes the problem is semantic.

Sometimes it is numerical.

Sometimes an error budget accumulates.

Sometimes a noncommuting order turns a harmless composition into an
amplifying one.

Sometimes every local statement remains true while the system simply
leaves the region where those statements apply.

The correct response is not to distrust modularity.

It is to make composition explicit.

\[
\boxed{
\text{safe composition requires a proved bridge between local guarantees and the declared system obligation.}
}
\]

When that bridge is absent, local validity remains local.

\section{References used in this
chapter}\label{references-used-in-this-chapter-40}

Inherited through audited prerequisites:

\begin{itemize}
\tightlist
\item
  (Higham 2002)
\item
  (Baydin et al. 2018)
\item
  (McLachlan and Quispel 2002)
\item
  (Hairer et al. 2006)
\end{itemize}

Exact source identities and claim boundaries are recorded in
sources/source-locks/ATLAS-CH-COMPOSE-001.yaml.

\chapter{Conditional Computation}\label{conditional-computation}

\textbf{Epistemic status:} established sparse/conditional-computation
mechanisms + audited Depth prerequisite + Atlas synthesis + exact finite
resource witness.\\
\textbf{Specification:}
manuscript/specifications/ATLAS-CH-SPARSE-001.md\\
\textbf{Formal packet:}
mathematics/derivations/ATLAS-CH-SPARSE-001-DERIVATIONS.md\\
\textbf{Computational witness:}
mathematics/computational-witnesses/ATLAS-CW-SPARSE-001.md\\
\textbf{Source lock:} sources/source-locks/ATLAS-CH-SPARSE-001.yaml

A sparse model is not one thing.

A matrix can contain many zeros.

A network can skip whole blocks.

A transformer can drop tokens.

A router can select only part of a module family.

An adaptive-depth system can stop early.

These mechanisms can all reduce some notion of work.

They do not reduce the same resource, and they do not imply the same
hardware behavior.

The governing rule of this chapter is:

\begin{quote}
conditional computation is about which computation actually executes for
a particular input or state; sparsity alone does not determine realized
cost.
\end{quote}

\section{From computational depth to conditional
paths}\label{from-computational-depth-to-conditional-paths}

The Depth chapter established that execution depth need not equal
architectural depth.

A network can:

\begin{itemize}
\tightlist
\item
  recur;
\item
  halt early;
\item
  solve toward an equilibrium;
\item
  skip blocks;
\item
  allocate different numbers of updates to different inputs.
\end{itemize}

Conditional computation broadens that idea.

Instead of asking only:

\begin{quote}
how many steps execute?
\end{quote}

ask:

\begin{quote}
which parts of the available computation execute, for which input, under
which routing rule, and at what resource cost?
\end{quote}

That question covers depth, width, token count, modules, and sparse
parameter structure.

\section{A typed sparsity object}\label{a-typed-sparsity-object}

Use

\texttt{S=(P,A,T,B,D,R)}.

Here:

\begin{itemize}
\tightlist
\item
  \texttt{P} is parameter sparsity;
\item
  \texttt{A} is activation sparsity;
\item
  \texttt{T} is token or item sparsity;
\item
  \texttt{B} is block, module, or expert sparsity;
\item
  \texttt{D} is conditional execution depth;
\item
  \texttt{R} is the routing or selection rule.
\end{itemize}

These coordinates are not interchangeable.

A network can be highly sparse in parameters while executing the same
sparse graph for every input.

Another network can have dense parameters but execute only two of eight
blocks for one input and all eight for another.

The first is sparse.

The second is conditional.

A system can be both.

\section{Static parameter sparsity}\label{static-parameter-sparsity}

A simple static pruning model is

\texttt{W\textquotesingle{}=M\ odot\ W}

with fixed binary mask \texttt{M}.

The zeros remain in the same positions across evaluated inputs.

Han, Mao, and Dally's Deep Compression is a representative
pruning/compression lineage (Han et al. 2016b).

Static parameter sparsity can reduce:

\begin{itemize}
\tightlist
\item
  storage;
\item
  nonzero arithmetic;
\item
  memory movement;
\end{itemize}

if the implementation and hardware exploit the structure.

It does not require a per-input router.

\section{Sparsity ratio is not execution
semantics}\label{sparsity-ratio-is-not-execution-semantics}

Suppose 90 percent of matrix entries are zero.

A dense kernel may still multiply through the zeros.

In that implementation:

\begin{itemize}
\tightlist
\item
  the representation is sparse;
\item
  the mathematical operator has many zeros;
\item
  the hardware execution may remain dense.
\end{itemize}

Therefore:

\begin{quote}
stored zero structure is not enough to establish skipped physical work.
\end{quote}

The execution path must be inspected.

\section{Conditional computation}\label{conditional-computation-1}

Conditional computation makes the active graph depend on the current
input or state.

Write:

\texttt{p(x)=Route(x)}

and

\texttt{y=Execute(p(x),x)}.

The path \texttt{p(x)} may specify:

\begin{itemize}
\tightlist
\item
  which blocks execute;
\item
  which tokens survive;
\item
  which modules activate;
\item
  how long computation continues.
\end{itemize}

The router and executor are separate objects.

That separation matters because routing costs work.

\section{Hard block skipping}\label{hard-block-skipping}

For residual-style branch \texttt{F\_k}, write

\texttt{x\_\{k+1\}=x\_k+g\_k(x\_k)F\_k(x\_k)}.

If

\texttt{g\_k(x\_k)\ in\ \{0,1\}}

and the implementation avoids evaluating \texttt{F\_k} when the gate is
zero, the branch is genuinely skipped.

SkipNet gives a concrete learned example of input-dependent
residual-block skipping (Wang et al. 2018).

The important mechanism is not merely that a coefficient becomes zero.

The block is bypassed.

\section{Soft gates are not automatically sparse
execution}\label{soft-gates-are-not-automatically-sparse-execution}

Suppose instead

\texttt{g\_k\ in\ {[}0,1{]}}.

If the system computes

\texttt{F\_k(x\_k)}

and only afterward multiplies by a small or zero coefficient, the
expensive branch has already run.

So:

\texttt{g\_k=0}

does not by itself prove:

\texttt{cost(F\_k)=0}.

A differentiable training relaxation can therefore describe a sparse
intent without producing sparse execution.

This boundary was already present in DEPTH-001 and remains load-bearing
here.

\section{Spatially adaptive
computation}\label{spatially-adaptive-computation-1}

Conditional computation need not apply uniformly across one input.

Figurnov et al.~allocate different computation to different spatial
positions in residual networks (Figurnov et al. 2017).

This introduces a finer distinction:

\begin{itemize}
\tightlist
\item
  input-level conditional computation;
\item
  region- or token-level conditional computation.
\end{itemize}

Two examples can execute different paths.

Two positions inside one example can also receive different computation.

The execution graph can vary within the input.

\section{Token sparsity}\label{token-sparsity}

A transformer-like system processes a set or sequence of tokens.

Let active tokens before stage \texttt{k} be \texttt{X\_k}.

A selector produces

\texttt{A\_k=Select\_k(X\_k)}.

Later computation can run only on \texttt{A\_k}.

DynamicViT is a concrete example of progressive, input-dependent token
sparsification in vision transformers (Rao et al. 2021).

The crucial implementation question is:

\begin{quote}
at what point does the dropped token stop participating in expensive
computation?
\end{quote}

\section{Masking after dense work is not the same as removing before
it}\label{masking-after-dense-work-is-not-the-same-as-removing-before-it}

Suppose attention over \texttt{n} tokens constructs an \texttt{n\ x\ n}
score matrix.

If all \texttt{n\^{}2} pair interactions are computed and some outputs
are masked afterward, the score-matrix work was not saved.

If selection reduces the active token count to \texttt{m\textless{}n}
before the later attention stage, the later dense pair count becomes
\texttt{m\^{}2}.

That arithmetic distinction is exact.

Whether it becomes a wall-clock speedup depends on the execution system.

\section{Structural sparsity and dynamic
sparsity}\label{structural-sparsity-and-dynamic-sparsity}

A useful taxonomy is:

\subsection{Structural/static}\label{structuralstatic}

The sparse pattern is known before seeing the particular input.

Examples:

\begin{itemize}
\tightlist
\item
  fixed pruning mask;
\item
  fixed block-sparse matrix;
\item
  fixed local attention window.
\end{itemize}

\subsection{Dynamic/input-dependent}\label{dynamicinput-dependent}

The sparse pattern depends on the current input or state.

Examples:

\begin{itemize}
\tightlist
\item
  block routing;
\item
  token dropping;
\item
  early exit;
\item
  adaptive spatial computation.
\end{itemize}

The same mathematical zero pattern can have different engineering
implications depending on when it becomes known.

\section{Why when matters}\label{why-when-matters}

Static structure can often be compiled, fused, packed, or scheduled
ahead of time.

Dynamic structure may require:

\begin{itemize}
\tightlist
\item
  a router;
\item
  synchronization;
\item
  variable shapes;
\item
  irregular memory access;
\item
  dynamic batching.
\end{itemize}

So dynamic sparsity can reduce arithmetic while increasing control
overhead.

Conditional computation is not free selection.

\section{Parameter count and active parameter
count}\label{parameter-count-and-active-parameter-count}

A conditionally executed model may contain many parameters while using
only a subset per input.

Therefore distinguish:

\begin{itemize}
\tightlist
\item
  total parameter capacity;
\item
  active parameters per input;
\item
  arithmetic operations per input;
\item
  memory residency;
\item
  parameter communication.
\end{itemize}

A model can have large total capacity without evaluating all parameters
on every example.

The downstream Mixture-of-Experts chapter will develop this in the
expert-routing setting.

\section{Routing is a computation}\label{routing-is-a-computation}

A router can be:

\begin{itemize}
\tightlist
\item
  a small neural network;
\item
  a confidence threshold;
\item
  a heuristic;
\item
  a learned policy;
\item
  a top-k selector.
\end{itemize}

Whatever its form, its cost belongs in the resource ledger.

Write, in one resource unit,

\texttt{C\_total(x)\ =\ C\_route(x)\ +\ C\_exec(p(x),x)}.

A system should not report only the work after routing and pretend
selection was free.

\section{Resource vectors}\label{resource-vectors}

One scalar rarely captures conditional-computation cost.

Use a resource vector such as

\texttt{C(x)=(F(x),K(x),M(x),E(x))},

with:

\begin{itemize}
\tightlist
\item
  arithmetic/FLOPs \texttt{F};
\item
  launches or dispatches \texttt{K};
\item
  memory traffic \texttt{M};
\item
  energy \texttt{E}.
\end{itemize}

Other deployments may add:

\begin{itemize}
\tightlist
\item
  communication volume;
\item
  synchronization;
\item
  sequential depth;
\item
  accelerator occupancy.
\end{itemize}

Do not add unlike units without a declared scalarization.

\section{Average cost}\label{average-cost}

Conditional computation is often motivated by easy and hard inputs.

For input distribution \texttt{X}, average cost is

\texttt{E{[}C(X){]}}.

On a finite benchmark it is the empirical mean.

This quantity answers:

\begin{quote}
what does a typical input cost under the evaluation distribution?
\end{quote}

It does not answer:

\begin{quote}
how much capacity is required for the worst input?
\end{quote}

\section{Peak and tail cost}\label{peak-and-tail-cost}

A service may care about:

\begin{itemize}
\tightlist
\item
  maximum active blocks;
\item
  p95 latency;
\item
  p99 communication;
\item
  worst-case memory;
\item
  deadline misses.
\end{itemize}

These are tail or peak questions.

A low mean can coexist with a dense worst-case path.

So reports should separate:

\begin{itemize}
\tightlist
\item
  average compute;
\item
  peak compute;
\item
  tail latency.
\end{itemize}

\section{Exact witness: fixed
baseline}\label{exact-witness-fixed-baseline}

Consider four sequential blocks.

Each contributes ten arithmetic units.

Assume the fixed implementation fuses them into one launch.

Its resource pair is

\texttt{C\_fixed=(40,1)}.

The coordinates are:

\begin{itemize}
\tightlist
\item
  arithmetic units;
\item
  launches.
\end{itemize}

Every input receives the same computation.

\section{Exact witness: conditional
paths}\label{exact-witness-conditional-paths}

Now use four input-dependent paths:

\begin{itemize}
\tightlist
\item
  \texttt{x\_1:\{1,4\}};
\item
  \texttt{x\_2:\{1,2,4\}};
\item
  \texttt{x\_3:\{1,4\}};
\item
  \texttt{x\_4:\{1,2,3,4\}}.
\end{itemize}

Assume:

\begin{itemize}
\tightlist
\item
  one router launch;
\item
  one launch for each active block.
\end{itemize}

The resource pairs are:

\begin{itemize}
\tightlist
\item
  \texttt{x\_1:(20,3)};
\item
  \texttt{x\_2:(30,4)};
\item
  \texttt{x\_3:(20,3)};
\item
  \texttt{x\_4:(40,5)}.
\end{itemize}

\section{Average arithmetic falls}\label{average-arithmetic-falls}

Average arithmetic is

\texttt{(20+30+20+40)/4\ =\ 55/2\ =\ 27.5}.

Compared with fixed arithmetic 40, the reduction is

\texttt{1-(55/2)/40\ =\ 5/16\ =\ 31.25\%}.

The conditional system therefore performs less arithmetic on average
under this finite workload.

That claim is exact.

\section{Peak arithmetic does not
fall}\label{peak-arithmetic-does-not-fall}

The fourth input executes all four blocks.

So:

\texttt{F\_peak\_cond=40}.

The fixed baseline also uses 40.

Thus:

\begin{itemize}
\tightlist
\item
  average arithmetic drops;
\item
  peak arithmetic is unchanged.
\end{itemize}

If hardware must provision for the worst path, the mean alone is
insufficient.

\section{Launch count gets worse}\label{launch-count-gets-worse}

Average conditional launches are

\texttt{(3+4+3+5)/4\ =\ 15/4\ =\ 3.75}.

The fixed path uses one.

So the conditional system is better in one coordinate and worse in
another:

\texttt{(40,1)}

versus

\texttt{(55/2,15/4)}.

Neither dominates componentwise.

\section{Hardware assumptions determine the
ordering}\label{hardware-assumptions-determine-the-ordering}

Define a toy scalar cost

\texttt{T\_(alpha,beta)(F,K)\ =\ alpha\ F+beta\ K}.

This is not claimed to be a realistic universal latency model.

It is a proof device.

If

\texttt{alpha=1,beta=0},

then only arithmetic matters:

\begin{itemize}
\tightlist
\item
  fixed = 40;
\item
  conditional = 27.5.
\end{itemize}

Conditional wins.

If

\texttt{alpha=1/10,beta=10},

then launches are expensive:

\begin{itemize}
\tightlist
\item
  fixed = 14;
\item
  conditional = 40.25.
\end{itemize}

Fixed wins.

The same execution traces support opposite latency conclusions under
different hardware cost models.

\section{FLOPs do not determine
latency}\label{flops-do-not-determine-latency}

The witness proves a structural point:

\begin{quote}
fewer arithmetic operations do not uniquely determine lower latency.
\end{quote}

Real hardware adds further effects:

\begin{itemize}
\tightlist
\item
  memory bandwidth;
\item
  cache behavior;
\item
  kernel launch overhead;
\item
  vectorization;
\item
  occupancy;
\item
  synchronization;
\item
  batching;
\item
  communication.
\end{itemize}

The Roofline model makes one part of this explicit by relating
attainable performance to arithmetic intensity, memory bandwidth, and
compute capability (Williams et al. 2009).

FLOPs are one coordinate.

They are not the clock.

\section{Dynamic irregularity can cost
efficiency}\label{dynamic-irregularity-can-cost-efficiency}

Hardware often prefers regular work.

A dense matrix operation can be highly optimized.

A dynamically sparse workload may introduce:

\begin{itemize}
\tightlist
\item
  tiny kernels;
\item
  irregular memory reads;
\item
  divergent branches;
\item
  underfilled accelerators;
\item
  fragmented batches.
\end{itemize}

Therefore a method can report strong nominal arithmetic reduction and
weak realized speedup.

That is not necessarily an implementation failure.

It may be the consequence of the hardware/software cost structure.

\section{Static sparsity can also fail to
accelerate}\label{static-sparsity-can-also-fail-to-accelerate}

Parameter pruning creates zeros.

But unstructured zeros only help execution if the kernels and hardware
exploit them.

Gale, Elsen, and Hooker's large-scale sparsity study reinforces the need
to evaluate sparsification methods at scale rather than treating
sparsity ratio itself as a complete performance metric (Gale et al.
2019).

The Atlas conclusion is broader:

\begin{quote}
sparsity percentage is not a hardware benchmark.
\end{quote}

\section{Conditional depth}\label{conditional-depth-1}

One form of conditional computation is to vary execution depth.

Easy input:

\texttt{tau(x)=2}.

Hard input:

\texttt{tau(x)=8}.

DEPTH-001 already supplied the stopping semantics.

SPARSE-001 adds the resource question:

\begin{itemize}
\tightlist
\item
  what work disappears when execution stops?
\item
  what routing/halting overhead remains?
\item
  what are average and tail depths?
\item
  what hardware can exploit the variable path?
\end{itemize}

\section{Early exit versus block
skipping}\label{early-exit-versus-block-skipping}

Early exit stops the remaining suffix of a network.

Block skipping can omit internal blocks while later blocks still
execute.

These are different path families.

For an eight-block network:

early exit might produce

\texttt{\{1,2,3\}}.

Skipping might produce

\texttt{\{1,3,4,7,8\}}.

Both are conditional.

Their scheduling and representation behavior can differ.

\section{Token dropping versus conditional
depth}\label{token-dropping-versus-conditional-depth}

Token dropping changes active width.

Conditional depth changes the number or identity of sequential
transformations.

A system can do both.

For example:

\begin{itemize}
\tightlist
\item
  prune tokens at stage 3;
\item
  skip block 5;
\item
  stop at block 7.
\end{itemize}

One scalar called ``sparsity'' hides these distinctions.

The typed sparsity object keeps them visible.

\section{Training-time execution can be
denser}\label{training-time-execution-can-be-denser}

Discrete routing creates optimization difficulties.

Training may therefore use:

\begin{itemize}
\tightlist
\item
  continuous gates;
\item
  straight-through estimators;
\item
  stochastic sampling;
\item
  auxiliary losses;
\item
  policy gradients;
\item
  dense masking.
\end{itemize}

Inference may use hard decisions.

The training graph and inference graph can therefore have different
cost.

Do not infer inference sparsity directly from the training relaxation.

\section{Compute and quality form a joint
claim}\label{compute-and-quality-form-a-joint-claim}

Skipping computation changes the function executed.

So an efficiency report should pair compute with quality.

At minimum report:

\texttt{(Q,C)}

for task quality \texttt{Q} and cost \texttt{C}.

A method that cuts computation in half and destroys task performance is
not a successful efficiency result under most objectives.

The acceptable frontier is application-dependent.

\section{Dynamic computation can specialize
effort}\label{dynamic-computation-can-specialize-effort}

The attraction of conditional computation is simple:

\begin{quote}
not every input may need the same amount or kind of processing.
\end{quote}

Easy examples may need less work.

Hard examples may need more.

Some tokens may be redundant.

Some spatial locations may be uninformative.

This creates an allocation problem.

The system spends compute where the router expects it to matter.

\section{The router can be wrong}\label{the-router-can-be-wrong}

A conditional system introduces a new failure surface.

The router can:

\begin{itemize}
\tightlist
\item
  drop useful tokens;
\item
  skip necessary blocks;
\item
  stop too early;
\item
  over-compute easy inputs;
\item
  under-compute hard inputs;
\item
  become unstable under distribution shift.
\end{itemize}

A dense baseline does not have exactly the same failure mode.

Efficiency comes with a control problem.

\section{Routing confidence is not
correctness}\label{routing-confidence-is-not-correctness}

A router may emit a high confidence score.

That does not prove the skipped computation was unnecessary.

The same epistemic boundary from DEPTH-001 applies:

\begin{quote}
a learned halting or routing signal is not automatically a calibrated
error estimator.
\end{quote}

If one wants guaranteed error-controlled adaptation, a separate theorem
or calibration layer is required.

\section{Batching changes the
economics}\label{batching-changes-the-economics}

Suppose active-block counts in a batch are

\texttt{2,3,2,4}.

The mean is

\texttt{11/4}.

But a synchronous implementation may effectively wait for the four-block
path.

Then wall-clock behavior tracks something closer to the maximum.

Dynamic batching, compaction, or regrouping can recover efficiency.

Those mechanisms add overhead of their own.

Average path length is not enough to predict batch latency.

\section{Communication can dominate
arithmetic}\label{communication-can-dominate-arithmetic}

In distributed conditional systems, selecting a remote module can
require communication.

Then total cost can contain:

\begin{itemize}
\tightlist
\item
  local compute;
\item
  routing;
\item
  all-to-all exchange;
\item
  synchronization;
\item
  imbalance.
\end{itemize}

A method that reduces local arithmetic can still slow down if
communication grows.

This becomes central in Mixture-of-Experts systems.

SPARSE-001 establishes the accounting discipline before that chapter.

\section{Capacity and active work can
decouple}\label{capacity-and-active-work-can-decouple}

Conditional execution allows a model to possess more total components
than one input uses.

This can increase representational capacity without proportional
per-input arithmetic.

But total capacity still affects:

\begin{itemize}
\tightlist
\item
  storage;
\item
  loading;
\item
  communication;
\item
  optimizer state during training;
\item
  deployment footprint.
\end{itemize}

``Only k modules active'' is not the same as ``only k modules exist.''

\section{Static compression and conditional capacity are different
strategies}\label{static-compression-and-conditional-capacity-are-different-strategies}

Static pruning asks:

\begin{quote}
which structure can be removed for all evaluated inputs?
\end{quote}

Conditional computation asks:

\begin{quote}
which structure can be omitted for this input?
\end{quote}

The former reduces a model globally.

The latter selects among available computation dynamically.

Combining them is possible.

Conflating them obscures the design space.

\section{Token sparsity can alter later complexity
superlinearly}\label{token-sparsity-can-alter-later-complexity-superlinearly}

If a later dense attention stage has pairwise score cost proportional to
\texttt{n\^{}2}, reducing token count from \texttt{n} to \texttt{m} can
reduce that score-matrix arithmetic roughly from \texttt{n\^{}2} to
\texttt{m\^{}2}.

But end-to-end speedup also depends on:

\begin{itemize}
\tightlist
\item
  selector cost;
\item
  projection cost;
\item
  memory movement;
\item
  fixed layers;
\item
  implementation shape efficiency.
\end{itemize}

The local arithmetic saving is not the whole pipeline.

\section{Measured speedup is contextual
evidence}\label{measured-speedup-is-contextual-evidence}

DynamicViT reports both FLOP reduction and measured throughput
improvement on its evaluated systems (Rao et al. 2021).

That is stronger evidence than FLOPs alone for those experiments.

It remains contextual:

\begin{itemize}
\tightlist
\item
  device;
\item
  software;
\item
  batch;
\item
  model;
\item
  implementation.
\end{itemize}

A measured speedup should not be promoted to a hardware-independent
theorem.

\section{Conditional computation is an allocation
policy}\label{conditional-computation-is-an-allocation-policy}

Viewed abstractly, the router allocates a compute budget.

For input \texttt{x}, it chooses path \texttt{p(x)}.

The execution then spends resources according to that path.

This suggests three separate questions:

\begin{enumerate}
\def\labelenumi{\arabic{enumi}.}
\tightlist
\item
  \textbf{selection:} which path is chosen?
\item
  \textbf{execution:} what work does that path perform?
\item
  \textbf{valuation:} was the saved work worth the quality change?
\end{enumerate}

These questions should be evaluated separately.

\section{Failure modes}\label{failure-modes-8}

\subsection{Sparsity/FLOP conflation}\label{sparsityflop-conflation}

Treat sparse weights as proof of skipped arithmetic.

\subsection{FLOP/latency conflation}\label{floplatency-conflation}

Treat fewer operations as proof of faster execution.

\subsection{Soft/hard conflation}\label{softhard-conflation}

Treat a relaxed gate as actual skipped hardware work.

\subsection{Mean/peak conflation}\label{meanpeak-conflation}

Report low average compute while hiding dense worst-case paths.

\subsection{Router-free accounting}\label{router-free-accounting}

Count only selected computation and omit selection overhead.

\subsection{Training/inference
conflation}\label{traininginference-conflation}

Measure soft training sparsity and claim hard inference savings.

\subsection{Capacity/active-work
conflation}\label{capacityactive-work-conflation}

Treat inactive modules as if they do not consume storage or training
resources.

\subsection{Quality omission}\label{quality-omission}

Report compute savings without the associated task-performance change.

\section{A practical conditional-compute
ledger}\label{a-practical-conditional-compute-ledger}

Before accepting an efficiency claim, record:

{\def\LTcaptype{none} % do not increment counter
\begin{longtable}[]{@{}
  >{\raggedright\arraybackslash}p{(\linewidth - 2\tabcolsep) * \real{0.5000}}
  >{\raggedright\arraybackslash}p{(\linewidth - 2\tabcolsep) * \real{0.5000}}@{}}
\toprule\noalign{}
\begin{minipage}[b]{\linewidth}\raggedright
Field
\end{minipage} & \begin{minipage}[b]{\linewidth}\raggedright
Question
\end{minipage} \\
\midrule\noalign{}
\endhead
\bottomrule\noalign{}
\endlastfoot
sparse object & Parameters, activations, tokens, blocks, experts,
depth? \\
static/dynamic & Fixed structure or input-dependent path? \\
selector & What chooses the active work? \\
hard/soft & Is work physically skipped or merely weighted? \\
training graph & What executes during training? \\
inference graph & What executes during deployment? \\
average work & Mean over which workload? \\
peak/tail work & Worst case or declared quantile? \\
arithmetic & How many operations are avoided? \\
memory & What traffic/storage changes? \\
launches & Does execution become more fragmented? \\
communication & Is data routed across devices? \\
hardware & Which machine/software/batch regime? \\
quality & What performance change accompanies the saving? \\
\end{longtable}
}

This ledger prevents ``sparse'' from becoming an untyped efficiency
claim.

\section{What the exact witness
establishes}\label{what-the-exact-witness-establishes-1}

The companion witness establishes:

\begin{itemize}
\tightlist
\item
  fixed baseline resource pair \texttt{(40,1)};
\item
  conditional average pair \texttt{(55/2,15/4)};
\item
  exact average arithmetic reduction \texttt{5/16=31.25\%};
\item
  unchanged peak arithmetic \texttt{40};
\item
  compute-dominated toy cost favors conditional execution;
\item
  launch-dominated toy cost favors fixed execution.
\end{itemize}

The witness does not claim either toy cost model is a real accelerator.

It proves only that arithmetic count does not uniquely identify latency.

\section{Downstream handoff}\label{downstream-handoff-5}

\textbf{Mixture-of-Experts Systems --- ATLAS-CH-MOE-001} may now assume:

\begin{itemize}
\tightlist
\item
  static versus conditional sparsity;
\item
  hard versus soft routing;
\item
  routing versus execution cost;
\item
  parameter capacity versus active work;
\item
  average versus peak/tail compute;
\item
  resource-vector accounting;
\item
  arithmetic-versus-latency boundary;
\item
  hardware-context dependence of measured speedup.
\end{itemize}

The downstream chapter must independently develop:

\begin{itemize}
\tightlist
\item
  expert routing;
\item
  capacity factors;
\item
  load balancing;
\item
  specialization;
\item
  collapse;
\item
  expert parallelism;
\item
  communication structure.
\end{itemize}

Conditional computation supplies the grammar.

Mixture-of-Experts will supply the system.

\section{References used in this
chapter}\label{references-used-in-this-chapter-41}

\begin{itemize}
\tightlist
\item
  Han, Mao, and Dally, \emph{Deep Compression: Compressing Deep Neural
  Networks with Pruning, Trained Quantization and Huffman Coding} (Han
  et al. 2016b).
\item
  Figurnov et al., \emph{Spatially Adaptive Computation Time for
  Residual Networks} (Figurnov et al. 2017).
\item
  Wang et al., \emph{SkipNet: Learning Dynamic Routing in Convolutional
  Networks} (Wang et al. 2018).
\item
  Rao et al., \emph{DynamicViT: Efficient Vision Transformers with
  Dynamic Token Sparsification} (Rao et al. 2021).
\item
  Gale, Elsen, and Hooker, \emph{The State of Sparsity in Deep Neural
  Networks} (Gale et al. 2019).
\item
  Williams, Waterman, and Patterson, \emph{Roofline: An Insightful
  Visual Performance Model for Multicore Architectures} (Williams et al.
  2009).
\end{itemize}

Exact source identities and claim boundaries are recorded in:

sources/source-locks/ATLAS-CH-SPARSE-001.yaml

\chapter{Mixture-of-Experts Systems}\label{mixture-of-experts-systems}

\textbf{Epistemic status:} established sparse-expert mechanisms +
audited Conditional Computation and Transformer prerequisites + Atlas
synthesis + exact finite witness.\\
\textbf{Specification:} manuscript/specifications/ATLAS-CH-MOE-001.md\\
\textbf{Formal packet:}
mathematics/derivations/ATLAS-CH-MOE-001-DERIVATIONS.md\\
\textbf{Computational witness:}
mathematics/computational-witnesses/ATLAS-CW-MOE-001.md\\
\textbf{Source lock:} sources/source-locks/ATLAS-CH-MOE-001.yaml

A mixture-of-experts model can contain far more parameters than it
executes for one token.

That is its central architectural attraction.

It is also the source of several confusions.

A router emits probabilities.

A capacity rule accepts or rejects assignments.

A distributed runtime moves tokens between devices.

Experts receive different training signals.

An auxiliary loss pushes traffic toward a target distribution.

These are related operations.

They are not the same object.

The governing rule of this chapter is:

\begin{quote}
an MoE system is defined by router scores, discrete dispatch, capacity
semantics, expert computation, and distributed execution together---not
by the router probabilities alone.
\end{quote}

\section{From conditional computation to
experts}\label{from-conditional-computation-to-experts}

SPARSE-001 established the broad conditional-computation picture:

\begin{itemize}
\tightlist
\item
  a model can expose more capacity than it executes for one input;
\item
  hard and soft gating must be distinguished;
\item
  routing overhead counts;
\item
  arithmetic reduction does not automatically imply latency reduction.
\end{itemize}

Mixture-of-experts systems instantiate that idea with a family of
alternative expert modules.

For each token, only a subset is selected.

\section{Where experts enter a
Transformer}\label{where-experts-enter-a-transformer}

TRANSFORMER-001 established the standard Transformer block:

\begin{itemize}
\tightlist
\item
  residual-stream state;
\item
  attention;
\item
  position-wise feed-forward transformation;
\item
  residual paths;
\item
  normalization.
\end{itemize}

A common sparse-expert Transformer replaces selected dense feed-forward
sublayers with expert families.

Attention still performs token mixing.

The expert layer performs conditionally selected token-wise
transformation.

Other designs are possible, but this is the baseline used here.

\section{Expert family}\label{expert-family}

Let a token state be:

\texttt{h\_i\ in\ R\^{}d}.

Let expert e be:

\texttt{E\_e:R\^{}d-\textgreater{}R\^{}d}.

There are:

\texttt{e=1,...,M}

experts.

Total expert parameter capacity can grow with M even if each token
visits only one or a few experts.

That decoupling between total capacity and active work is a defining
sparse-expert property.

\section{Router logits}\label{router-logits}

A router maps token state to scores:

\texttt{z\_i\ in\ R\^{}M}.

The scores can be arbitrary real values.

By themselves they do not specify which experts execute.

They are preferences or routing evidence.

\section{Router probabilities}\label{router-probabilities}

A common transformation is:

\texttt{p\_\{i,e\}\ =\ exp(z\_\{i,e\})\ /\ sum\_j\ exp(z\_\{i,j\})}.

Then:

\texttt{p\_\{i,e\}\textgreater{}=0}

and:

\texttt{sum\_e\ p\_\{i,e\}=1}.

These probabilities are useful for ranking and weighting.

They are still not the executed route.

\section{Preferred top-k route}\label{preferred-top-k-route}

Let:

\texttt{P\_i=TopK(p\_i,k)}.

For top-1:

\texttt{\textbar{}P\_i\textbar{}=1}

before capacity handling.

For top-2 or larger k, several experts are preferred.

The word \textbf{preferred} matters.

Capacity can prevent a preferred route from executing.

\section{Accepted dispatch}\label{accepted-dispatch}

Define:

\texttt{a\_\{i,e\}\ in\ \{0,1\}}.

Here:

\texttt{a\_\{i,e\}=1}

means token i is actually dispatched to expert e after capacity and
overflow policy.

This is the first object that directly describes expert execution.

A router probability can be nonzero while:

\texttt{a\_\{i,e\}=0}.

\section{Gate weights}\label{gate-weights}

Suppose token i is accepted by expert set:

\texttt{A\_i=\{e:a\_\{i,e\}=1\}}.

A generic sparse mixture writes:

\texttt{MoE(h\_i)\ =\ sum\_(e\ in\ A\_i)\ alpha\_\{i,e\}\ E\_e(h\_i)}.

The weights \texttt{alpha} might be:

\begin{itemize}
\tightlist
\item
  original router probabilities;
\item
  probabilities renormalized over accepted experts;
\item
  hard unit weights;
\item
  another declared convention.
\end{itemize}

Dispatch and aggregation weighting are separate design choices.

\section{Top-1 is simpler, not
trivial}\label{top-1-is-simpler-not-trivial}

Switch Transformer popularized a particularly simple routing case: one
expert per token.

This reduces some communication and combination complexity.

But even top-1 requires decisions about:

\begin{itemize}
\tightlist
\item
  capacity;
\item
  overflow;
\item
  balancing;
\item
  dropped tokens;
\item
  expert placement;
\item
  batch composition.
\end{itemize}

One selected expert does not mean one simple system.

\section{Capacity is not a probability
threshold}\label{capacity-is-not-a-probability-threshold}

Suppose an expert can process at most C routed tokens in the current
routing group.

Capacity constrains:

\texttt{a\_\{i,e\}}.

It does not change the fact that the router may assign high probability
to more than C tokens.

Thus:

\begin{quote}
capacity acts on accepted assignments, not on the probability simplex.
\end{quote}

This distinction is load-bearing.

\section{Nominal capacity}\label{nominal-capacity}

A common nominal form is:

\texttt{C\ =\ ceil(c\_f\ k\ N/M)}.

Here:

\begin{itemize}
\tightlist
\item
  N is token count in the routing group;
\item
  M is expert count;
\item
  k is nominal experts per token;
\item
  \texttt{c\_f} is a capacity factor.
\end{itemize}

The exact formula and routing group are implementation choices.

The Atlas does not promote one capacity formula into a universal
definition.

\section{Demand versus accepted
load}\label{demand-versus-accepted-load}

Preferred demand for expert e is:

\texttt{d\_e\ =\ sum\_i\ 1\{e\ in\ P\_i\}}.

Accepted token load is:

\texttt{n\_e\ =\ sum\_i\ a\_\{i,e\}}.

When capacity or overflow changes routing:

\texttt{d\_e\ !=\ n\_e}

can occur.

This is the first reason a router histogram can differ from executed
expert load.

\section{Overflow needs policy}\label{overflow-needs-policy}

If:

\texttt{d\_e\textgreater{}C},

something must happen.

Possible policies include:

\begin{itemize}
\tightlist
\item
  drop overflow tokens;
\item
  reroute them;
\item
  send them through a fallback dense path;
\item
  expand effective capacity;
\item
  buffer work.
\end{itemize}

The router probabilities do not answer this question.

The overflow policy is part of the model's operational semantics.

\section{Dropping changes the
computation}\label{dropping-changes-the-computation}

If a token is dropped from the expert layer, then the executed
computation differs from the nominal top-k design.

The token may:

\begin{itemize}
\tightlist
\item
  skip the expert transform;
\item
  continue through a residual path;
\item
  use some fallback;
\item
  be handled differently by implementation.
\end{itemize}

Therefore drop rate belongs in both model-quality and systems
accounting.

\section{Rerouting changes the meaning of
top-1}\label{rerouting-changes-the-meaning-of-top-1}

Suppose a token's preferred expert is full.

A reroute policy may send it to its second choice.

Then the executed expert is not the argmax expert.

This is not a contradiction.

It means:

\begin{quote}
preferred route and accepted route are different objects.
\end{quote}

Any analysis of router behavior should state which one it studies.

\section{Expert token load}\label{expert-token-load}

Define normalized accepted load:

\texttt{f\_e\ =\ n\_e\ /\ sum\_j\ n\_j}, \textbf{provided at least one
expert assignment was accepted}, so \(\sum_jn_j>0\). If every token is
dropped (or there are no dispatched tokens), the normalized
accepted-load vector is undefined; record an empty-dispatch state and
the absolute loads rather than divide by zero or claim balanced traffic.

If all top-1 tokens are accepted:

\texttt{sum\_e\ n\_e=N}.

If some are dropped:

\texttt{sum\_e\ n\_e\textless{}N}.

A load-balancing plot should therefore state its denominator.

\section{Router probability mass}\label{router-probability-mass}

Define:

\texttt{m\_e=sum\_i\ p\_\{i,e\}}

and normalized mass:

\texttt{q\_e=m\_e/N}.

Because each probability row sums to one:

\texttt{sum\_e\ q\_e=1}.

But there is no general identity:

\texttt{q\_e=f\_e}.

One measures soft router preference.

The other measures accepted discrete traffic.

\section{Balance has multiple
meanings}\label{balance-has-multiple-meanings}

``Balanced experts'' can mean several different things:

\begin{itemize}
\tightlist
\item
  equal token counts;
\item
  equal router probability mass;
\item
  equal arithmetic work;
\item
  equal bytes communicated;
\item
  equal device time;
\item
  equal training signal;
\item
  equal utility.
\end{itemize}

These are not interchangeable.

The word \textbf{balance} is incomplete without the quantity being
balanced.

\section{Auxiliary balancing
objectives}\label{auxiliary-balancing-objectives}

Sparse expert systems often add an auxiliary objective:

\texttt{L\_total\ =\ L\_task\ +\ lambda\ L\_balance}.

Different systems use different proxies and stabilizers.

The important structural point is:

\texttt{L\_balance\ !=\ L\_task}.

A router can become more balanced under its proxy while the main task
becomes worse.

Or the reverse.

\section{Balance is a systems
pressure}\label{balance-is-a-systems-pressure}

Why balance at all?

If one expert receives nearly every token while others sit idle:

\begin{itemize}
\tightlist
\item
  capacity is wasted;
\item
  one device can become a straggler;
\item
  communication patterns become uneven;
\item
  underused experts may receive little training signal.
\end{itemize}

So load balancing often serves utilization and optimization stability.

It is not evidence that every expert should receive exactly the same
semantic workload.

\section{Specialization is a different
claim}\label{specialization-is-a-different-claim}

An expert is specialized when its behavior or performance differs
systematically on some class of inputs.

Evidence might include:

\begin{itemize}
\tightlist
\item
  domain enrichment;
\item
  conditional accuracy;
\item
  feature selectivity;
\item
  stable functional differences;
\item
  systematic output differences.
\end{itemize}

Traffic imbalance alone does not prove any of these.

A popular expert may simply be a routing attractor.

\section{Equal traffic does not prove equal
usefulness}\label{equal-traffic-does-not-prove-equal-usefulness}

Two experts can each receive 10,000 tokens.

One can be critical.

The other can be nearly redundant.

Traffic counts measure use frequency.

They do not directly measure value.

This distinction will appear exactly in the finite witness.

\section{Preferred-router collapse versus accepted-load
concentration}\label{preferred-router-collapse-versus-accepted-load-concentration}

Call \textbf{preferred-router collapse} a regime where router
probability mass or preferred routes concentrate on a small expert
subset under a declared statistic and horizon.

Accepted dispatch is a different object. Capacity, rerouting, or
dropping can make accepted loads look more balanced than the underlying
preferences.

Useful diagnostics therefore distinguish:

\begin{itemize}
\tightlist
\item
  probability-mass concentration;
\item
  preferred-route concentration;
\item
  accepted-load concentration;
\item
  persistent overflow at particular experts.
\end{itemize}

``Collapse'' without naming the object and statistic is too vague.

\section{Expert underuse}\label{expert-underuse}

An expert can receive so few tokens that it gets little training signal.

That is underuse.

It may be caused by:

\begin{itemize}
\tightlist
\item
  router bias;
\item
  initialization;
\item
  data distribution;
\item
  capacity interactions;
\item
  optimization dynamics.
\end{itemize}

Underuse is not the same thing as two experts learning the same
function.

\section{Expert redundancy}\label{expert-redundancy}

Two experts can both be heavily used and still implement similar
functions.

A functional comparison might evaluate:

\texttt{E{[}\textbar{}\textbar{}E\_e(h)-E\_j(h)\textbar{}\textbar{}\^{}2{]}}

under a declared input distribution.

Low functional distance is one possible redundancy signal.

It says something traffic counts cannot.

\section{Capacity overload is different
again}\label{capacity-overload-is-different-again}

A routing system can repeatedly demand more slots than an expert is
allowed to serve.

That is a capacity/overflow problem.

It may co-occur with router concentration.

But the concepts differ:

\begin{itemize}
\tightlist
\item
  preferred-router collapse concerns where router preferences
  concentrate;
\item
  accepted-load concentration concerns executed traffic after capacity
  handling;
\item
  capacity overload concerns preferred demand relative to an execution
  limit.
\end{itemize}

Keeping them separate helps identify the correct repair.

\section{Expert parallelism}\label{expert-parallelism}

Large expert families are often partitioned across devices.

Let:

\texttt{dev(e)}

denote the device hosting expert e.

A routed token may originate elsewhere.

Then conditional computation becomes a distributed-systems problem.

The expert computation may be sparse while token movement is expensive.

\section{Dispatch communication}\label{dispatch-communication}

If token i starts on:

\texttt{src(i)}

and is accepted by expert e, inter-device movement occurs when:

\texttt{src(i)\ !=\ dev(e)}.

A simple communication-count proxy is:

\texttt{K\_comm\ =\ sum\_(i,e)\ a\_\{i,e\}\ 1\{src(i)!=dev(e)\}}.

This is only a count.

Actual cost depends on:

\begin{itemize}
\tightlist
\item
  representation size;
\item
  collective implementation;
\item
  batching;
\item
  topology;
\item
  return traffic;
\item
  synchronization.
\end{itemize}

\section{All-to-all is not free}\label{all-to-all-is-not-free}

Expert-parallel systems often group tokens by destination expert,
exchange them, compute expert outputs, and return results.

This can require collective communication.

Therefore:

\begin{quote}
low active arithmetic can coexist with high communication cost.
\end{quote}

GShard and Switch make this systems issue central to large sparse-expert
execution.

\section{Stragglers matter}\label{stragglers-matter}

Suppose expert computation is parallel.

Step completion can depend on the slowest device/expert path.

Even if average load is good, one overloaded expert can dominate
latency.

This is another reason average token count is not enough.

Peak and tail behavior matter.

\section{Equal token counts need not mean equal
compute}\label{equal-token-counts-need-not-mean-equal-compute}

If every expert has the same architecture and input shape, equal
accepted counts can equalize one arithmetic proxy.

But experts may differ in:

\begin{itemize}
\tightlist
\item
  size;
\item
  precision;
\item
  hardware placement;
\item
  kernel efficiency;
\item
  input shape;
\item
  cache locality.
\end{itemize}

Then:

\texttt{n\_e=n\_j}

does not imply:

\texttt{C\_e=C\_j}.

The load unit must match the cost question.

\section{Training and serving need not use the same routing
semantics}\label{training-and-serving-need-not-use-the-same-routing-semantics}

During training, the system may tolerate:

\begin{itemize}
\tightlist
\item
  router noise;
\item
  auxiliary balance pressure;
\item
  higher capacity factors;
\item
  token dropping;
\item
  larger batches.
\end{itemize}

Serving may prioritize:

\begin{itemize}
\tightlist
\item
  determinism;
\item
  tail latency;
\item
  replica placement;
\item
  availability;
\item
  bounded communication.
\end{itemize}

A routing analysis should therefore state the execution phase.

\section{The exact witness: six
tokens}\label{the-exact-witness-six-tokens}

Use six tokens and three experts.

The router probabilities are:

{\def\LTcaptype{none} % do not increment counter
\begin{longtable}[]{@{}lrrr@{}}
\toprule\noalign{}
token & E1 & E2 & E3 \\
\midrule\noalign{}
\endhead
\bottomrule\noalign{}
\endlastfoot
t1 & 9/10 & 1/20 & 1/20 \\
t2 & 4/5 & 3/20 & 1/20 \\
t3 & 1/2 & 1/10 & 2/5 \\
t4 & 1/20 & 9/10 & 1/20 \\
t5 & 1/20 & 3/4 & 1/5 \\
t6 & 1/20 & 1/20 & 9/10 \\
\end{longtable}
}

Every row sums to one.

\section{Naive top-1 preferences}\label{naive-top-1-preferences}

The preferred experts are:

\begin{itemize}
\tightlist
\item
  t1 -\textgreater{} E1;
\item
  t2 -\textgreater{} E1;
\item
  t3 -\textgreater{} E1;
\item
  t4 -\textgreater{} E2;
\item
  t5 -\textgreater{} E2;
\item
  t6 -\textgreater{} E3.
\end{itemize}

So preferred loads are:

\texttt{(3,2,1)}.

Let capacity be:

\texttt{C=2}.

E1 is overloaded.

\section{Keep the strongest E1
assignments}\label{keep-the-strongest-e1-assignments}

For E1:

\begin{itemize}
\tightlist
\item
  t1 has probability 9/10;
\item
  t2 has probability 4/5;
\item
  t3 has probability 1/2.
\end{itemize}

The declared overflow policy retains the two highest E1 preferences:

\texttt{t1,t2}.

Token t3 becomes overflow.

\section{Reroute the overflow
token}\label{reroute-the-overflow-token}

For t3:

\begin{itemize}
\tightlist
\item
  E2 probability is 1/10;
\item
  E3 probability is 2/5.
\end{itemize}

E3 has one spare slot.

Therefore:

\texttt{t3:E1-\textgreater{}E3}.

Final accepted assignments are:

\begin{itemize}
\tightlist
\item
  E1: t1,t2;
\item
  E2: t4,t5;
\item
  E3: t3,t6.
\end{itemize}

Accepted loads are:

\texttt{(2,2,2)}.

\section{Accepted load is perfectly
balanced}\label{accepted-load-is-perfectly-balanced}

Accepted load fractions are:

\texttt{f=(1/3,1/3,1/3)}.

For the squared count-imbalance statistic:

\texttt{B\_count\ =\ sum\_e(f\_e-1/3)\^{}2},

we get:

\texttt{B\_count=0}.

By this metric, routing is perfectly balanced.

\section{Router probability mass is not
balanced}\label{router-probability-mass-is-not-balanced}

Summing probabilities over all six tokens gives:

\texttt{m=(47/20,2,33/20)}.

Normalize by six:

\texttt{q=(47/120,1/3,11/40)}.

This is not uniform.

Therefore:

\texttt{f\ !=\ q}.

Balanced accepted counts have not made the router's soft preference
distribution uniform.

\section{Exact probability-imbalance
value}\label{exact-probability-imbalance-value}

The offsets from uniform are:

\texttt{47/120-1/3=7/120};

\texttt{1/3-1/3=0};

\texttt{11/40-1/3=-7/120}.

Therefore:

\texttt{B\_prob\ =\ 2(7/120)\^{}2\ =\ 49/7200}.

So the exact witness has:

\texttt{B\_count=0}

but:

\texttt{B\_prob\textgreater{}0}.

The two balance notions disagree.

\section{Equal traffic still does not imply equal
usefulness}\label{equal-traffic-still-does-not-imply-equal-usefulness}

Now attach a toy post-routing utility diagnostic.

Let final assignment utilities be:

\begin{itemize}
\tightlist
\item
  E1 on t1: 4;
\item
  E1 on t2: 4;
\item
  E2 on t4: 2;
\item
  E2 on t5: 2;
\item
  E3 on t3: 1;
\item
  E3 on t6: 1.
\end{itemize}

Expert totals are:

\texttt{(8,4,2)}.

Traffic remains:

\texttt{(2,2,2)}.

Equal counts do not imply equal evaluated utility.

\section{The utility numbers are not router
scores}\label{the-utility-numbers-are-not-router-scores}

This distinction matters.

The router probability says:

\begin{quote}
how strongly did the router prefer this expert?
\end{quote}

The toy utility says:

\begin{quote}
under a separate evaluation, how much value did this accepted expert
assignment contribute?
\end{quote}

These need not match.

Indeed, one research question in MoE systems is whether the router
reliably sends inputs to experts that are actually useful for them.

\section{Drop instead of reroute}\label{drop-instead-of-reroute}

Now keep the same probability table and same capacity.

Change only the overflow policy.

Drop t3 instead of rerouting it.

Then accepted loads become:

\texttt{(2,2,1)}.

Drop rate is:

\texttt{1/6}.

The router has not changed.

The executed MoE has.

\section{Overflow policy changes
arithmetic}\label{overflow-policy-changes-arithmetic}

Assume identical expert cost c per accepted token.

Under rerouting:

\texttt{C\_expert=6c}.

Under dropping:

\texttt{C\_expert=5c}.

Thus capacity semantics alter realized expert arithmetic even with
identical router probabilities.

This is a concrete version of the SPARSE-001 distinction between routing
logic and executed work.

\section{Balanced counts do not settle
communication}\label{balanced-counts-do-not-settle-communication}

Suppose E1, E2, and E3 live on different devices.

Even with:

\texttt{(2,2,2)}

accepted token counts, the communication pattern depends on where the
tokens originated.

One routing assignment can be local.

Another can require extensive cross-device exchange.

Count balance is therefore not a communication theorem.

\section{Shazeer et al.: sparse gating and
balancing}\label{shazeer-et-al.-sparse-gating-and-balancing}

Shazeer et al.~introduced a sparsely gated MoE layer at large scale,
with a trainable gating network selecting a sparse combination of
experts.

Their work also makes balancing pressure explicit because highly uneven
expert use creates practical problems.

The Atlas uses this as a primary historical mechanism.

It does not treat the original routing loss or expert layout as
universally canonical.

\section{GShard: capacity meets distributed
execution}\label{gshard-capacity-meets-distributed-execution}

GShard places sparse MoE computation inside a large Transformer and
combines it with automatic sharding.

That makes two facts impossible to ignore:

\begin{itemize}
\tightlist
\item
  routing creates per-expert capacity constraints;
\item
  experts distributed over accelerators create communication and
  placement problems.
\end{itemize}

The Atlas uses GShard to connect algorithmic routing to expert-parallel
execution.

\section{Switch: top-1 simplifies one
axis}\label{switch-top-1-simplifies-one-axis}

Switch Transformer reduces the number of selected experts per token to
one.

This can simplify routing and communication relative to multi-expert
mixtures.

It does not remove:

\begin{itemize}
\tightlist
\item
  capacity;
\item
  balancing;
\item
  training instability;
\item
  distributed communication.
\end{itemize}

The important lesson is that simplification of one routing axis leaves a
multi-layer system.

\section{ST-MoE: stability and transfer remain system
properties}\label{st-moe-stability-and-transfer-remain-system-properties}

ST-MoE studies sparse expert design with emphasis on training stability
and downstream transfer.

That matters because an MoE system cannot be judged from parameter count
or routing sparsity alone.

Useful questions include:

\begin{itemize}
\tightlist
\item
  does training remain stable?
\item
  do experts receive enough signal?
\item
  does specialization persist?
\item
  does sparse pretraining transfer?
\end{itemize}

These questions belong downstream of the basic routing semantics
established here.

\section{A practical MoE ledger}\label{a-practical-moe-ledger}

Before accepting an MoE result, record:

{\def\LTcaptype{none} % do not increment counter
\begin{longtable}[]{@{}ll@{}}
\toprule\noalign{}
Field & Question \\
\midrule\noalign{}
\endhead
\bottomrule\noalign{}
\endlastfoot
tokens & What is the routing group and token count? \\
experts & How many, and what does each expert compute? \\
router & Which logits/probabilities are produced? \\
top-k & How many experts are preferred per token? \\
gate weights & How are accepted expert outputs weighted? \\
capacity & Which unit and formula set expert capacity? \\
overflow & Drop, reroute, fallback, buffer, or other? \\
accepted load & How many assignments actually execute per expert? \\
probability mass & What soft router mass goes to each expert? \\
auxiliary loss & Which balance/stability proxy is optimized? \\
drops & Which tokens receive no expert execution? \\
specialization & What functional evidence supports the claim? \\
redundancy & How are expert functions compared? \\
placement & Which device hosts each expert? \\
communication & Which token/activation bytes move? \\
phase & Training or inference/serving? \\
quality boundary & What does balance not prove? \\
\end{longtable}
}

This ledger turns ``MoE routing'' into an auditable system description.

\section{Failure modes}\label{failure-modes-9}

\subsection{Probability-equals-dispatch}\label{probability-equals-dispatch}

Soft router mass is treated as executed expert traffic.

\subsection{Top-k-equals-execution}\label{top-k-equals-execution}

Preferred experts are reported without capacity/overflow semantics.

\subsection{Balanced-counts-equals-specialization}\label{balanced-counts-equals-specialization}

Equal traffic is used as evidence that experts learned distinct useful
functions.

\subsection{Balanced-counts-equals-cheap}\label{balanced-counts-equals-cheap}

Equal accepted loads are treated as proof of equal latency or
communication.

\subsection{Collapse-as-one-word}\label{collapse-as-one-word}

Preferred-router concentration, accepted-load concentration, underuse,
redundancy, and capacity overload are conflated.

\subsection{Auxiliary-loss-equals-task-objective}\label{auxiliary-loss-equals-task-objective}

A balancing proxy is treated as if it directly measures task quality.

\subsection{Sparse-parameters-equals-sparse-system-cost}\label{sparse-parameters-equals-sparse-system-cost}

Active expert arithmetic is counted while router, communication,
padding, and synchronization are ignored.

\section{What the exact witness
establishes}\label{what-the-exact-witness-establishes-2}

The companion witness proves:

\begin{itemize}
\tightlist
\item
  naive top-1 demand is \texttt{(3,2,1)};
\item
  capacity is 2;
\item
  t3 is the lowest-confidence E1 token;
\item
  its best available alternative is E3;
\item
  rerouting t3 yields accepted loads \texttt{(2,2,2)};
\item
  probability mass is \texttt{(47/20,2,33/20)};
\item
  normalized mass is \texttt{(47/120,1/3,11/40)};
\item
  count imbalance is 0;
\item
  probability-mass imbalance is \texttt{49/7200};
\item
  toy utility totals are \texttt{(8,4,2)};
\item
  drop-on-overflow instead yields \texttt{(2,2,1)} and drop rate
  \texttt{1/6};
\item
  under equal expert token cost, reroute and drop policies use expert
  arithmetic \texttt{6c} and \texttt{5c}.
\end{itemize}

No universal routing optimum is inferred.

\section{Downstream handoff: Router
Dynamics}\label{downstream-handoff-router-dynamics}

\textbf{ATLAS-CH-ROUTERDYN-001} may now assume:

\begin{itemize}
\tightlist
\item
  logits/probabilities versus accepted dispatch;
\item
  capacity and overflow semantics;
\item
  count-load versus probability-mass diagnostics;
\item
  preferred-router concentration versus accepted-load concentration
  versus underuse versus redundancy;
\item
  explicit evidence requirements for specialization.
\end{itemize}

It must independently develop temporal churn, instability, commutators,
spectral diagnostics, and route evolution.

\section{Downstream handoff:
Systems}\label{downstream-handoff-systems}

\textbf{ATLAS-CH-SYSTEMS-001} may now assume:

\begin{itemize}
\tightlist
\item
  expert placement;
\item
  dispatch communication;
\item
  capacity;
\item
  accepted per-expert work;
\item
  active versus total expert capacity;
\item
  straggler boundaries.
\end{itemize}

It must independently develop hardware/system cost models, communication
topology, serving constraints, and measured efficiency.

The MoE lesson is:

\begin{quote}
sparse experts create conditional capacity, not free capacity. The
router proposes; the capacity policy dispatches; the experts compute;
the network moves data; and each layer needs its own evidence.
\end{quote}

\section{References used in this
chapter}\label{references-used-in-this-chapter-42}

\begin{itemize}
\tightlist
\item
  Noam Shazeer et al., \emph{Outrageously Large Neural Networks: The
  Sparsely-Gated Mixture-of-Experts Layer}, 2017, arXiv:1701.06538.
\item
  Dmitry Lepikhin et al., \emph{GShard: Scaling Giant Models with
  Conditional Computation and Automatic Sharding}, 2020,
  arXiv:2006.16668.
\item
  William Fedus, Barret Zoph, and Noam Shazeer, \emph{Switch
  Transformers: Scaling to Trillion Parameter Models with Simple and
  Efficient Sparsity}, Journal of Machine Learning Research 23(120),
  1--39, 2022.
\item
  Barret Zoph et al., \emph{ST-MoE: Designing Stable and Transferable
  Sparse Expert Models}, 2022, arXiv:2202.08906.
\end{itemize}

Exact source identities and authority boundaries are recorded in:

sources/source-locks/ATLAS-CH-MOE-001.yaml

\chapter{Router Dynamics and
Diagnostics}\label{router-dynamics-and-diagnostics}

\textbf{Epistemic status:} audited MoE/Optimizer-Dynamics prerequisites
+ primary router-stability/routing sources + Atlas synthesis + exact
finite witnesses\\
\textbf{Specification:}
manuscript/specifications/ATLAS-CH-ROUTERDYN-001.md\\
\textbf{Derivation packet:}
mathematics/derivations/ATLAS-CH-ROUTERDYN-001-DERIVATIONS.md\\
\textbf{Computational witness:}
mathematics/computational-witnesses/ATLAS-CW-ROUTERDYN-001.md

A routing system can look balanced while changing its mind about every
token.

That is the central diagnostic problem.

Mixture-of-Experts systems expose several routing layers at once:
probabilities, preferred routes, capacity-constrained dispatch, realized
expert load, and expert behavior. Looking at only one can hide motion in
the others.

\section{Four moving objects}\label{four-moving-objects}

For N tracked tokens and E experts, let

\[
P_t
\]

be the router probability matrix at checkpoint t.

Let

\[
r_t(i)
\]

be the declared preferred top-1 expert for token i.

Let

\[
A_t
\]

be the accepted dispatch matrix after capacity, overflow, and execution
policy.

Finally, let

\[
\ell_t=A_t^\top\mathbf 1
\]

be accepted expert load.

The audited MoE chapter already established that these are not
interchangeable.

ROUTERDYN adds time.

\section{Probability drift}\label{probability-drift}

Define average rowwise probability drift by

\[
D_P(t)
=
\frac{1}{2N}
\sum_i
\|P_t(i,:)-P_{t-1}(i,:)\|_1.
\]

The factor one-half turns rowwise L1 distance between probability
vectors into total variation.

This diagnostic can move even when the preferred expert does not.

That matters because routing margins can become sharper or flatter
before a top-1 assignment flips.

\section{Preferred-route churn}\label{preferred-route-churn}

For top-1 routing, define

\[
\chi_R(t)
=
\frac1N
\sum_i
\mathbf1\{r_t(i)\ne r_{t-1}(i)\}.
\]

This is token-identity sensitive.

A value of zero says every tracked token kept the same preferred expert.

A value of one says every tracked token changed preferred expert.

It says nothing by itself about whether those changes helped the task.

\section{Accepted-dispatch churn}\label{accepted-dispatch-churn}

Capacity and overflow can change what is actually executed.

For one accepted expert per token,

\[
\chi_A(t)
=
\frac{1}{2N}
\|A_t-A_{t-1}\|_{1,\mathrm{entry}}.
\]

Preferred-route churn and accepted-dispatch churn can differ when
capacity handling intervenes.

That is why a router-dynamics dashboard should not report one as a
synonym for the other.

\section{Load drift}\label{load-drift}

Accepted load is much coarser:

\[
D_\ell(t)
=
\frac{1}{2N}
\|\ell_t-\ell_{t-1}\|_1.
\]

Load answers a systems question: how much accepted work moved among
experts?

It does not preserve token identity.

That information loss creates the first exact counterexample.

\section{Exact counterexample: perfect balance, total
churn}\label{exact-counterexample-perfect-balance-total-churn}

Take four tracked tokens.

At one checkpoint,

\[
r_0=(1,1,2,2).
\]

At the next,

\[
r_1=(2,2,1,1).
\]

The load vector is unchanged:

\[
\ell_0=\ell_1=(2,2).
\]

So

\[
D_\ell=0.
\]

But every token changed expert:

\[
\chi_R=\chi_A=1.
\]

Thus:

\[
\boxed{
\text{stable load}
\not\Rightarrow
\text{stable routing}
}
\]

in the exact finite witness.

This is why load balancing is not a temporal-stability metric.

\section{Transition operators on expert
identity}\label{transition-operators-on-expert-identity}

For a fixed tracked token panel, define the empirical expert-transition
matrix

\[
T_t(a,b)
=
\frac{
\#\{i:r_{t-1}(i)=a,\ r_t(i)=b\}
}{
\#\{i:r_{t-1}(i)=a\}
}.
\]

Every row with a positive denominator is stochastic. An expert with
\textbf{no previously assigned tracked token} has no empirical
conditional transition row: the quotient is undefined there, and any
zero/uniform fill-in would be an explicit smoothing convention rather
than an observation. The following witness has both prior expert classes
occupied, so both rows are defined.

The swap witness produces

\[
T_{\rm swap}
=
\begin{pmatrix}
0&1\\
1&0
\end{pmatrix}.
\]

A stable routing panel produces

\[
T_{\rm stable}=I.
\]

The matrices describe very different temporal behavior.

Yet both have singular values

\[
(1,1).
\]

So:

\[
\boxed{
\text{same singular spectrum}
\not\Rightarrow
\text{same routing dynamics}
}
\]

for this pair.

Spectral diagnostics are useful only when the operator being diagnosed
and the information discarded by the summary are explicit.

\section{One-step transition is not automatically a Markov
model}\label{one-step-transition-is-not-automatically-a-markov-model}

The matrix \(T_t\) is an empirical transition object.

Calling it a stationary Markov chain requires additional assumptions.

The data distribution may change.

The router parameters may change.

The experts themselves may change.

Capacity policy may change.

The token panel may not be representative.

ROUTERDYN therefore uses Markov-style algebra only as a diagnostic lens
unless stationarity and sampling assumptions are separately justified.

\section{Route churn can be zero while probabilities
move}\label{route-churn-can-be-zero-while-probabilities-move}

Now take two tokens and two experts:

\[
P_0=
\begin{pmatrix}
0.6&0.4\\
0.4&0.6
\end{pmatrix},
\qquad
P_1=
\begin{pmatrix}
0.9&0.1\\
0.1&0.9
\end{pmatrix}.
\]

The preferred experts do not change.

Therefore

\[
\chi_R=0.
\]

But the average probability drift is

\[
D_P=0.3.
\]

A routing system can therefore become much more confident without
changing its top-1 token assignment.

The converse is also possible: a small probability change near a
decision boundary can cause route churn.

\section{Routing stability and training stability are different
claims}\label{routing-stability-and-training-stability-are-different-claims}

ST-MoE studies training instabilities in large sparse models and
introduces router z-loss as one stabilization technique in that setting
(Zoph et al. 2022).

That is valuable evidence about sparse-model training.

But ``training became more stable'' is not identical to ``token
assignments became temporally stable.''

A loss spike, router-logit magnitude, assignment churn, load imbalance,
and expert specialization are different observables.

The chapter keeps them separate.

\section{Routing rules alter what can
specialize}\label{routing-rules-alter-what-can-specialize}

Expert Choice reverses a common routing perspective: instead of each
token choosing a fixed number of experts, experts choose tokens subject
to capacity (Zhou et al. 2022).

This reinforces a broader point.

Specialization is not produced by experts in isolation.

It depends on the routing rule that determines what experience reaches
them.

A temporal specialization study must therefore record the routing policy
alongside the expert profile.

\section{Specialization needs a declared
taxonomy}\label{specialization-needs-a-declared-taxonomy}

Suppose each tracked token has a declared category

\[
c_i\in\mathcal C.
\]

For expert e with nonzero accepted load, define

\[
S_t(e,c)
=
\frac{
\sum_i A_t(i,e)\mathbf1\{c_i=c\}
}{
\ell_t(e)
}.
\]

This gives the expert's accepted-token category profile.

A simple specialization score can compare that profile with the
tracked-panel category distribution.

But the category system is part of the measurement.

If the categories are language, syntax, topic, token identity, task
family, or another feature, the meaning of ``specialized'' changes.

Therefore:

\[
\boxed{
\text{load concentration}
\ne
\text{semantic specialization}
}
\]

and

\[
\boxed{
\text{specialization score}
\ne
\text{causal importance}.
}
\]

\section{Temporal specialization
drift}\label{temporal-specialization-drift}

An expert can keep the same load while the composition of that load
changes.

A total-variation profile drift can be written

\[
D_S(t,e)
=
\frac12
\sum_c
|S_t(e,c)-S_{t-1}(e,c)|.
\]

This lets the Atlas ask whether an expert's observed role is stable even
when its utilization is stable.

Again, the answer is relative to the declared taxonomy.

\section{Router dynamics include optimizer
memory}\label{router-dynamics-include-optimizer-memory}

Routing parameters are trained parameters.

If their next update depends on momentum, adaptive moments, schedules,
or other optimizer state, then the dynamical state is larger than the
router weights alone.

Let

\[
\xi_t
\]

collect the declared router-plus-optimizer state.

For a differentiable router-plus-optimizer update map \(H_t\), set
\(J_t=DH_t(\xi_t)\). Its \textbf{tangent (first-variation) dynamics} are

\[
\delta\xi_{t+1}=J_t\delta\xi_t.
\]

Here \(\delta\xi\) denotes a tangent perturbation along a fixed
reference trajectory, not an arbitrary finite difference between two
nonlinear runs. For finite input displacement \(v\), differentiability
gives \(H_t(\xi_t+v)-H_t(\xi_t)=J_tv+o(\|v\|)\) instead.

Along a changing trajectory,

\[
\delta\xi_{t+h}
=
J_{t+h-1}\cdots J_t\delta\xi_t.
\]

This product is exact for the tangent recursion with its declared
trajectory-dependent Jacobians, not a global finite-amplitude nonlinear
evolution law. This is inherited from the audited Optimizer-State
Dynamics viewpoint.

\section{Commutators detect order
sensitivity}\label{commutators-detect-order-sensitivity}

Successive local maps need not commute.

Define

\[
\mathcal C_t
=
J_tJ_{t-1}-J_{t-1}J_t.
\]

If

\[
\mathcal C_t\ne0,
\]

the order of the two local linear maps matters.

This is a diagnostic fact.

It is not a causal explanation of why the maps changed.

\section{Exact commutator witness}\label{exact-commutator-witness}

Take

\[
J_0=
\begin{pmatrix}
1&1\\
0&1
\end{pmatrix},
\qquad
J_1=
\begin{pmatrix}
1&0\\
1&1
\end{pmatrix}.
\]

Each has eigenvalues

\[
(1,1).
\]

But

\[
J_1J_0
\ne
J_0J_1.
\]

The commutator is

\[
\begin{pmatrix}
-1&0\\
0&1
\end{pmatrix},
\]

with Frobenius norm

\[
\sqrt2.
\]

Identical one-step eigenvalue multisets therefore do not determine the
two-step order-sensitive behavior.

\section{Spectral router diagnostics are
typed}\label{spectral-router-diagnostics-are-typed}

Several spectra may be interesting:

\begin{itemize}
\tightlist
\item
  the spectrum of an empirical expert-transition matrix;
\item
  singular values of that transition matrix;
\item
  eigenvalues of a local augmented router-state Jacobian;
\item
  singular values of that Jacobian;
\item
  finite-horizon gains of Jacobian products.
\end{itemize}

These are spectra of different operators.

They cannot be compared as though ``the router spectrum'' were one
universal object.

\section{Churn is not automatically
bad}\label{churn-is-not-automatically-bad}

High churn can indicate instability.

It can also indicate:

\begin{itemize}
\tightlist
\item
  adaptation to a changing data distribution;
\item
  a router still discovering useful specialization;
\item
  exchange among functionally similar experts;
\item
  movement caused by changing capacity semantics;
\item
  benign margin crossing near equivalent experts.
\end{itemize}

Low churn can be bad too.

A router can be stably wrong.

ROUTERDYN therefore reports churn before interpreting churn.

\section{Capacity can hide preference
dynamics}\label{capacity-can-hide-preference-dynamics}

The MoE audit established an important warning: accepted load can look
balanced even when preferred routing is concentrated because capacity
handling modifies execution.

The temporal version is stronger.

Accepted-load stability can arise from a capacity mechanism even while
preferred token assignments or router probabilities move substantially.

A longitudinal study should therefore preserve both pre-capacity and
post-capacity observables.

\section{A practical diagnostic
record}\label{a-practical-diagnostic-record}

A router-dynamics measurement should declare at least:

\begin{itemize}
\tightlist
\item
  checkpoint/time index;
\item
  tracked token panel or sampling distribution;
\item
  router logits/probabilities;
\item
  preferred-route rule and tie-breaking;
\item
  capacity and overflow policy;
\item
  accepted dispatch;
\item
  load vector;
\item
  route-churn metric;
\item
  probability-drift metric;
\item
  any token taxonomy used for specialization;
\item
  any transition operator;
\item
  any augmented-state Jacobian approximation;
\item
  spectral statistic and operator identity.
\end{itemize}

Without these fields, two ``router stability'' claims may be measuring
different phenomena.

\section{Handoff to Routing as Online Decision
Making}\label{handoff-to-routing-as-online-decision-making}

ATLAS-CH-REGRETROUTE-001 may now assume:

\begin{itemize}
\tightlist
\item
  probability, preferred-route, accepted-dispatch, and load temporal
  separation;
\item
  exact route-churn and probability-drift definitions;
\item
  empirical expert-transition operators;
\item
  taxonomy-relative specialization profiles;
\item
  local augmented-state commutator diagnostics;
\item
  the finite examples showing that loads and singular spectra can hide
  token-level motion.
\end{itemize}

REGRETROUTE must add the online-decision layer.

It must define:

\begin{itemize}
\tightlist
\item
  the routing action;
\item
  comparison policy/class;
\item
  reward or loss timing;
\item
  regret;
\item
  optionality;
\item
  correction capacity;
\item
  nonstationarity assumptions.
\end{itemize}

Temporal diagnostics are not regret theory by themselves.

\section{Durable lesson}\label{durable-lesson-1}

A router does not have one dynamics.

It has several coupled temporal surfaces:

\[
\text{probability}
\to
\text{preference}
\to
\text{dispatch}
\to
\text{load},
\]

plus expert behavior and optimizer memory.

The chapter's main discipline is to diagnose each surface before
interpreting the system.

\section{References used in this
chapter}\label{references-used-in-this-chapter-43}

\begin{itemize}
\tightlist
\item
  (Zoph et al. 2022)
\item
  (Zhou et al. 2022)
\end{itemize}

Exact source authority, prerequisite identities, and claim boundaries
are locked in:

sources/source-locks/ATLAS-CH-ROUTERDYN-001.yaml

\chapter{Routing as Online Decision
Making}\label{routing-as-online-decision-making}

\textbf{Epistemic status:} audited Router Dynamics substrate + primary
adversarial-bandit/shifting-expert authority + Atlas-owned exact finite
routing witnesses.\\
\textbf{Specification:}
manuscript/specifications/ATLAS-CH-REGRETROUTE-001.md\\
\textbf{Derivation packet:}
mathematics/derivations/ATLAS-CH-REGRETROUTE-001-DERIVATIONS.md\\
\textbf{Computational witness:}
mathematics/computational-witnesses/ATLAS-CW-REGRETROUTE-001.md\\
\textbf{Source lock:} sources/source-locks/ATLAS-CH-REGRETROUTE-001.yaml

A sparse router changes over time.

That does not yet make it an online decision problem.

To speak about regret, one must say:

\begin{itemize}
\tightlist
\item
  what action was actually taken;
\item
  what actions were feasible;
\item
  when loss was incurred;
\item
  what feedback arrived;
\item
  what comparator is allowed;
\item
  what kind of nonstationarity is being measured.
\end{itemize}

The governing rule is:

\[
\boxed{
\text{diagnose routing first; define regret only after the decision problem is explicit.}
}
\]

\section{Router dynamics are not regret
theory}\label{router-dynamics-are-not-regret-theory}

Router Dynamics already distinguished:

\[
\text{probability}
\to
\text{preference}
\to
\text{dispatch}
\to
\text{load}.
\]

It also showed that:

\begin{itemize}
\tightlist
\item
  stable load need not mean stable token assignment;
\item
  zero preferred-route churn need not mean zero probability drift;
\item
  equal singular spectra need not mean equal dynamics.
\end{itemize}

Those are temporal diagnostics.

Regret asks a different question:

\begin{quote}
how much cumulative loss did the routing decisions incur relative to a
declared comparator class?
\end{quote}

\section{Preferred route is not always the executed
action}\label{preferred-route-is-not-always-the-executed-action}

Let:

\[
p_t
\]

be the preferred route proposed by the router.

Let:

\[
F_t
\]

be the set of dispatch actions that remain feasible after capacity and
availability constraints.

Let:

\[
a_t\in F_t
\]

be the accepted/executed dispatch.

When capacity binds:

\[
a_t\neq p_t
\]

can occur.

That distinction matters because the executed action is the one that
actually consumes capacity and incurs the declared decision loss.

\section{Capacity is part of the decision
semantics}\label{capacity-is-part-of-the-decision-semantics}

Write:

\[
a_t=G_t(p_t,\kappa_t),
\]

where \(\kappa_t\) contains the relevant capacity/overflow state.

The map \(G_t\) can:

\begin{itemize}
\tightlist
\item
  accept the proposal;
\item
  reroute it;
\item
  reject it;
\item
  defer it;
\item
  send it to overflow capacity.
\end{itemize}

The online problem is underspecified until this behavior is declared.

\section{Loss happens after the accepted
action}\label{loss-happens-after-the-accepted-action}

Use the ordering:

\begin{enumerate}
\def\labelenumi{\arabic{enumi}.}
\tightlist
\item
  observe the pre-action context;
\item
  propose a route;
\item
  enforce feasibility/capacity;
\item
  execute accepted dispatch;
\item
  incur loss;
\item
  reveal feedback.
\end{enumerate}

Denote the accepted-action loss by:

\[
\ell_t(a_t).
\]

If future losses are used to choose \(a_t\), the object is an oracle or
hindsight construction, not an ordinary online policy.

\section{Feedback must be named}\label{feedback-must-be-named}

Under full information, after round \(t\) the learner receives:

\[
\{\ell_t(a):a\in F_t\}.
\]

Under bandit feedback, it receives only:

\[
\ell_t(a_t).
\]

Auer, Cesa-Bianchi, Freund, and Schapire treat adversarial/nonstochastic
multi-armed bandits under partial feedback (Auer, Cesa-Bianchi, Freund,
et al. 2002).

That distinction changes what an online router can infer and guarantee.

\section{Comparator comes before
regret}\label{comparator-comes-before-regret}

Let:

\[
\Pi
\]

be a declared class of comparison policies or action sequences.

Then define:

\[
R_T^{\Pi}(a)
=
\sum_{t=1}^T\ell_t(a_t)
-
\min_{u\in\Pi}
\sum_{t=1}^T\ell_t(u_t).
\]

Without \(\Pi\), ``regret'' is incomplete.

\section{Static regret asks one
question}\label{static-regret-asks-one-question}

A static comparator uses one fixed route for every round, restricted
here to routes in \(\bigcap_{t=1}^T F_t\). If the intersection is empty,
a feasible static comparator does not exist and a different comparator
or an expressly labeled counterfactual benchmark is needed.

It asks:

\begin{quote}
how much worse was the router than the best single fixed expert in
hindsight?
\end{quote}

That is useful when one expert should remain globally best.

It is not the only plausible sparse-routing benchmark.

\section{Shifting regret asks another
question}\label{shifting-regret-asks-another-question}

Suppose which expert is best changes over time.

Let:

\[
\Pi_S(F_{1:T})
=
\left\{
 u_{1:T}:
 u_t\in F_t\ \text{for every }t,\quad
 \sum_{t=2}^T\mathbf 1[u_t\neq u_{t-1}]
 \le S
\right\}.
\]

Now the comparator may switch at most \(S\) times \textbf{while
respecting the same declared feasible accepted-dispatch sets}. This
version assumes the comparator class is nonempty; otherwise the minimum
in the regret definition is not defined without an additional
convention. A comparator allowed to use capacity-infeasible actions is a
different, explicitly counterfactual benchmark.

Herbster and Warmuth study tracking a best expert that can change over
segments (Herbster and Warmuth 1998).

This is the comparator style used in the exact witness.

\section{One common finite routing
problem}\label{one-common-finite-routing-problem}

Take two accepted actions:

\[
A,\ B.
\]

Use four rounds.

Both routes remain feasible every round:

\[
F_t=\{A,B\}.
\]

Use losses:

\[
\begin{array}{c|cc}
 t&A&B\\
\hline
1&0&1\\
2&1&0\\
3&0&1\\
4&1&0
\end{array}
\]

The loss-minimizing route alternates every round.

\section{The shifting comparator}\label{the-shifting-comparator}

Allow up to three switches:

\[
\Pi_3.
\]

The unique zero-loss comparator is:

\[
(A,B,A,B).
\]

So the hindsight comparator loss is:

\[
0.
\]

\section{A perfectly stable route can regret its
stability}\label{a-perfectly-stable-route-can-regret-its-stability}

Consider:

\[
a^{\mathrm{stay}}=(A,A,A,A).
\]

Its accepted-dispatch churn is:

\[
0.
\]

Its cumulative loss is:

\[
0+1+0+1=2.
\]

Therefore its shifting regret is:

\[
\boxed{2}.
\]

The route is maximally stable.

It is not regret-optimal for the shifting comparator.

\section{A maximally churning route can be
regret-optimal}\label{a-maximally-churning-route-can-be-regret-optimal}

Now use:

\[
a^{\mathrm{track}}=(A,B,A,B).
\]

Its switch count is:

\[
3.
\]

Its cumulative loss is:

\[
0.
\]

Therefore its shifting regret is:

\[
\boxed{0}.
\]

The route churns every possible time.

It is nevertheless comparator-optimal.

\section{Churn is not a regret
surrogate}\label{churn-is-not-a-regret-surrogate}

The two policies use the same:

\begin{itemize}
\tightlist
\item
  action space;
\item
  feasible sets;
\item
  horizon;
\item
  losses;
\item
  feedback;
\item
  comparator class.
\end{itemize}

Yet:

\[
0<3
\]

for churn, while:

\[
2>0
\]

for regret.

So:

\[
\boxed{
\text{low churn}
\not\Rightarrow
\text{low regret}
}
\]

and:

\[
\boxed{
\text{high churn}
\not\Rightarrow
\text{high regret}.
}
\]

\section{Static regret tells a different
story}\label{static-regret-tells-a-different-story}

The fixed action \(A\) loses two.

The fixed action \(B\) also loses two.

So best static comparator loss is two.

Against that comparator:

\begin{itemize}
\tightlist
\item
  stay-on-\(A\) has regret zero;
\item
  the alternating tracker has regret \(-2\).
\end{itemize}

The numbers changed because the question changed.

Comparator choice is part of the meaning of regret.

\section{Negative static regret is not an
error}\label{negative-static-regret-is-not-an-error}

The alternating route beats every fixed comparator.

Therefore static regret can be negative for a realized sequence under
this definition.

That simply means the policy outperformed the restricted comparator
class.

It does not invalidate the calculation.

\section{Optionality is feasible
choice}\label{optionality-is-feasible-choice}

Define:

\[
O_t=|F_t|-1.
\]

This counts feasible accepted alternatives other than the executed
action.

In the exact witness:

\[
O_t=1
\]

for every round.

Both routes have the same optionality.

Their regret differs.

So optionality is not regret either.

\section{Correction capacity is ability to change
course}\label{correction-capacity-is-ability-to-change-course}

One operational form is remaining switch budget.

Let allowed switch count be \(S\).

Before round \(t\), let:

\[
N_t
=
\sum_{s=2}^{t-1}
\mathbf1[a_s\neq a_{s-1}].
\]

Define:

\[
K_t=\max(0,S-N_t).
\]

This says how many further route changes the controller may still make
under the declared switch budget.

\section{Optionality and correction capacity are
distinct}\label{optionality-and-correction-capacity-are-distinct}

A router can have many feasible alternatives but no remaining switching
budget.

Or it can have switching budget but only one feasible action because
capacity is saturated.

Thus:

\[
\boxed{
\text{available alternatives}
\neq
\text{ability to correct}.
}
\]

\section{The stay policy preserves correction
capacity}\label{the-stay-policy-preserves-correction-capacity}

In the exact witness, \(S=3\).

The stay route uses no switches.

Its remaining switch budget remains high.

Yet it accumulates regret two.

Preserving correction capacity is not automatically the same as using it
well.

\section{The tracking policy spends correction
capacity}\label{the-tracking-policy-spends-correction-capacity}

The alternating route consumes switch budget.

Its remaining correction capacity falls over time.

But it obtains zero shifting regret.

Again, correction capacity is a resource, not a performance score.

\section{Capacity should constrain the comparator
too}\label{capacity-should-constrain-the-comparator-too}

Suppose expert \(B\) is full at round \(t\).

Then \(B\notin F_t\).

A feasible comparator should normally obey the same restriction.

Otherwise the regret gap mixes online routing quality with access to an
impossible action.

An unconstrained oracle comparator can still be informative, but it must
be labeled as counterfactual.

\section{Preferred-route regret is a separate
counterfactual}\label{preferred-route-regret-is-a-separate-counterfactual}

Suppose the router prefers \(B\), but capacity sends the token to \(A\).

One can measure:

\begin{itemize}
\tightlist
\item
  accepted-dispatch regret;
\item
  preferred-route counterfactual regret.
\end{itemize}

They answer different questions.

The default REGRETROUTE object is accepted-dispatch regret.

\section{Capacity can hide or create
churn}\label{capacity-can-hide-or-create-churn}

Changing capacity can make accepted dispatch move even while preference
remains fixed.

It can also make accepted dispatch remain fixed while preference
changes.

Therefore ROUTERDYN's typed observables are required even after regret
enters the analysis.

\section{Load balance is not online
optimality}\label{load-balance-is-not-online-optimality}

A capacity mechanism can enforce balanced load.

Balanced load can be operationally valuable.

But it does not imply that accepted decisions minimize cumulative
routing loss.

So:

\[
\boxed{
\text{balanced load}
\not\Rightarrow
\text{low regret}.
}
\]

\section{Specialization is not
regret}\label{specialization-is-not-regret}

A specialization statistic asks which token classes or features tend to
use an expert.

Regret asks how losses compare to a benchmark.

An expert can be highly specialized and still be wrong for the current
token.

A broad expert can be the regret-minimizing action under another
objective.

\section{Probability drift is not
regret}\label{probability-drift-is-not-regret}

Router probabilities can drift substantially while the accepted route
remains unchanged.

Then accepted-dispatch regret can remain identical.

Conversely, tiny probability movement near a routing threshold can
change the accepted action and cumulative regret.

\section{If switching is costly, put it in the
loss}\label{if-switching-is-costly-put-it-in-the-loss}

Suppose switching has penalty \(\lambda\).

Use:

\[
\widetilde\ell_1=\ell_1(a_1),
\qquad
\widetilde\ell_t
=
\ell_t(a_t)
+
\lambda\mathbf1[a_t\neq a_{t-1}]
\quad(t=2,\ldots,T).
\]

This convention charges only switches \textbf{between} the \(T\)
executed actions; there is no undeclared \(a_0\) or first-round
switching penalty. A deployment that charges initialization or migration
from a pre-existing route must declare that initial route and charge
separately.

Now churn enters the decision objective directly.

For the alternating route, three switches contribute:

\[
3\lambda.
\]

A large enough switching cost can make stability preferable.

That is not a contradiction.

It is a different objective.

\section{Nonstationarity must be
typed}\label{nonstationarity-must-be-typed}

The exact witness has changing roundwise losses.

But empirical router nonstationarity can come from:

\begin{itemize}
\tightlist
\item
  changing contexts;
\item
  changing expert parameters;
\item
  changing capacity;
\item
  optimizer-state drift;
\item
  changing token mix;
\item
  reward-model drift.
\end{itemize}

Those mechanisms can produce different regret problems.

\section{Bandit feedback changes the online
problem}\label{bandit-feedback-changes-the-online-problem}

Under bandit feedback, the router sees only the loss of the accepted
action.

It does not immediately know whether another expert would have been
better.

This is why exploration matters in adversarial bandits (Auer,
Cesa-Bianchi, Freund, et al. 2002).

The exact witness computes realized regret after the loss table is
fixed; it does not claim the online policy knows the hindsight-optimal
route.

\section{Delayed feedback weakens
correction}\label{delayed-feedback-weakens-correction}

If routing loss arrives late, the controller can spend several rounds
without seeing a mistake.

That changes the operational value of correction capacity.

The exact witness uses immediate post-round feedback.

Delayed-feedback guarantees require separate analysis.

\section{Hindsight comparator is not an executable
oracle}\label{hindsight-comparator-is-not-an-executable-oracle}

The comparator may use the full realized sequence to define the
benchmark.

The online router cannot.

Regret measures the gap to hindsight performance.

It does not prove the comparator was available as a real-time policy.

\section{Downstream task loss is still
separate}\label{downstream-task-loss-is-still-separate}

Routing loss can be defined from:

\begin{itemize}
\tightlist
\item
  latency;
\item
  communication;
\item
  expert quality;
\item
  capacity violation;
\item
  token-level prediction loss;
\item
  a weighted combination.
\end{itemize}

A low routing regret under one declared loss does not automatically
imply low end-to-end task loss under another metric.

An interface theorem or empirical relation is needed.

\section{A practical routing-decision
record}\label{a-practical-routing-decision-record}

A regret-bearing routing experiment should record at least:

\begin{itemize}
\tightlist
\item
  context/token identity or sampling rule;
\item
  router probability/logit state;
\item
  preferred route;
\item
  feasible action set;
\item
  capacity/overflow state;
\item
  accepted dispatch;
\item
  incurred loss/reward;
\item
  feedback revealed;
\item
  comparator class;
\item
  switch/churn metric;
\item
  optionality;
\item
  correction-capacity state;
\item
  regret notion;
\item
  downstream task metric separately.
\end{itemize}

Without these fields, ``router regret'' is underspecified.

\section{Durable non-implications}\label{durable-non-implications-1}

REGRETROUTE rejects:

\[
\text{low churn}
\not\Rightarrow
\text{low regret},
\]

\[
\text{high churn}
\not\Rightarrow
\text{high regret},
\]

\[
\text{balanced load}
\not\Rightarrow
\text{low regret},
\]

\[
\text{stable probabilities}
\not\Rightarrow
\text{low regret},
\]

\[
\text{high optionality}
\not\Rightarrow
\text{low regret},
\]

and:

\[
\text{low routing regret}
\not\Rightarrow
\text{low downstream task loss without an explicit bridge}.
\]

\section{Closing view}\label{closing-view-28}

A router has dynamics before it has regret.

Once routing is treated as online decision making, the burden becomes
explicit:

\begin{itemize}
\tightlist
\item
  name the executed action;
\item
  constrain the feasible set;
\item
  state when loss arrives;
\item
  state what feedback is revealed;
\item
  define the comparator;
\item
  define optionality and correction capacity;
\item
  state what kind of nonstationarity is allowed.
\end{itemize}

The four-round witness then gives the smallest useful lesson:

\[
(A,A,A,A)
\]

has zero churn but regret two, while:

\[
(A,B,A,B)
\]

has three switches but regret zero under the same common frame.

So the right conclusion is not ``stable routing is good'' or ``adaptive
routing is good.''

It is:

\[
\boxed{
\text{routing quality is comparator-, loss-, feedback-, feasibility-, and capacity-relative.}
}
\]

\section{References used in this
chapter}\label{references-used-in-this-chapter-44}

\begin{itemize}
\tightlist
\item
  (Auer, Cesa-Bianchi, Freund, et al. 2002)
\item
  (Herbster and Warmuth 1998)
\end{itemize}

Exact source identities and claim boundaries are recorded in:

sources/source-locks/ATLAS-CH-REGRETROUTE-001.yaml

\chapter{A Taxonomy of Machine
Memory}\label{a-taxonomy-of-machine-memory}

\textbf{Epistemic status:} established memory mechanisms plus Atlas
synthesis and exact finite witness.\\
\textbf{Specification:}
manuscript/specifications/ATLAS-CH-MEMTAX-001.md\\
\textbf{Formal packet:}
mathematics/derivations/ATLAS-CH-MEMTAX-001-DERIVATIONS.md\\
\textbf{Computational witness:}
mathematics/computational-witnesses/ATLAS-CW-MEMTAX-001.md\\
\textbf{Source lock:} sources/source-locks/ATLAS-CH-MEMTAX-001.yaml

A machine does not have one thing called memory.

It has states that live in different places, survive for different
lengths of time, are written by different mechanisms, and can be reached
through different kinds of address.

A model's parameters can carry information across deployments.

A context can carry information across one inference.

A cache can survive seconds.

A replay buffer can preserve individual experience.

A database can survive model replacement.

A vector index can recover a record because its content resembles a
query even when the caller does not know its exact key.

Calling all of these ``memory'' is useful only if we preserve the
distinctions.

This chapter builds that vocabulary.

\section{Memory is an interface, not merely a
medium}\label{memory-is-an-interface-not-merely-a-medium}

Suppose two systems store exactly the same records.

One lets us retrieve them only by exact identifier.

The other lets us retrieve them by similarity.

The storage is the same. The memory behavior is different.

Now reverse the example.

One system stores a fact in model parameters.

Another stores the same fact in a database.

Both may answer the same query, yet their write paths, update costs,
provenance, mutability, and failure modes are radically different.

So the first question should not be merely:

Where is the memory?

It should be:

What memory contract does this state satisfy?

For the Atlas, write a memory system as

\[
M = (L, W, R, T, A, U, P, S).
\]

Here:

\begin{itemize}
\tightlist
\item
  \(L\) is storage locus;
\item
  \(W\) is the write or update path;
\item
  \(R\) is the read or retrieval path;
\item
  \(T\) is lifetime and retention policy;
\item
  \(A\) is addressability;
\item
  \(U\) is mutability and update cadence;
\item
  \(P\) is provenance;
\item
  \(S\) is sharing and synchronization scope.
\end{itemize}

The familiar memory labels are patterns across these coordinates.

They are not six boxes into which every implementation must fit exactly
once.

\section{Six roles, not six rungs}\label{six-roles-not-six-rungs}

The chapter uses six recurring labels:

parametric memory;

working memory;

episodic memory;

semantic memory;

associative memory;

external memory.

They do not form a hierarchy.

They are not even all the same kind of category.

Parametric and external primarily describe where state lives and how it
is updated.

Working primarily describes active availability and lifetime.

Episodic and semantic primarily describe how stored information relates
to experience.

Associative primarily describes how memory is addressed.

A single store can therefore occupy several roles at once.

An external vector store of timestamped experiences can be external,
episodic, and associative simultaneously.

That overlap is not a defect in the taxonomy.

It is the point.

\section{Parametric memory}\label{parametric-memory}

A trained model carries information in its parameters.

Let \(\theta\) denote the parameter state.

A training or editing procedure changes it:

\[
\theta' = W_{\mathrm{param}}(\theta, \text{data}, \text{objective}, \text{update\_rule}).
\]

A forward pass reads from that state implicitly:

\[
y = f_\theta(x).
\]

There may be no explicit record saying that one fact is stored at one
address.

Information can be distributed across many parameters and activated only
by particular inputs.

This is what the chapter means by parametric memory.

Empirical work on language models has shown that factual associations
can sometimes be elicited without supplying an external knowledge source
at inference time.

That is evidence that parameters can function as a memory substrate.

It is not evidence that parameters behave like a clean database.

A stored association can be hard to locate.

A correction can interfere with other behavior.

Provenance may be diffuse.

Freshness depends on retraining, editing, or replacement.

And the model may generate a plausible relation that was never stored as
a stable factual record.

Parametric memory is powerful partly because storage and computation are
fused.

That same fusion makes explicit memory management difficult.

\section{Working memory}\label{working-memory}

Some state exists only because a computation is currently in progress.

A recurrent hidden state.

An active context.

A scratch buffer.

A temporary plan.

A cache of intermediate keys and values.

These are examples of machine working memory when they remain directly
available to the current computation but are not intended to survive
indefinitely.

The term has a cognitive-science lineage. Baddeley and Hitch used
working memory for limited active storage and processing in human
cognition.

The Atlas borrows the engineering intuition, not the biological
architecture.

For machines, the operational question is:

what state is actively available now, and when will it disappear?

Working memory is often cheap to read because it is already in the
active computation.

It is also easy to lose.

A context can be truncated.

A process can restart.

A cache can be evicted.

A session can end.

High immediate availability and long-term persistence are different
properties.

\section{Context is not durable
memory}\label{context-is-not-durable-memory}

Modern systems often call context memory.

That can be reasonable.

But context residence should not be confused with persistence.

If a fact appears in the current prompt, the model can use it.

If the next session begins without that prompt, the fact may be gone.

The same applies to hidden state and caches.

A useful test is:

if the process is destroyed and reconstructed, does the state still
exist somewhere from which it can be recovered?

If no, the state may be working memory without being durable memory.

A system that repeatedly compiles durable state into an active context
therefore has at least two memory layers, not one.

\section{Episodic memory}\label{episodic-memory}

An episode remembers an event as an event.

A machine episodic record might contain:

event identity;

time;

task;

observation;

action;

outcome;

local context;

provenance.

Write one record as

\[
e_i = (\mathrm{id}_i, \mathrm{time}_i, \mathrm{context}_i, \mathrm{payload}_i, \mathrm{provenance}_i).
\]

The key feature is event binding.

If we later ask,

what happened on attempt 17?

we are asking an episodic question.

If we ask,

what general rule did the last thousand attempts support?

we are asking for something closer to semantic consolidation.

Neural Episodic Control is a concrete machine-learning example of
rapidly written experience records used in decision making.

The Atlas uses it as a mechanism example.

It does not infer that every replay buffer or event log has the same
semantics.

\section{Semantic memory}\label{semantic-memory}

Semantic memory stores generalized content that can be used without
reconstructing one originating event.

Suppose many episodes support a relation:

objects of type X usually require operation Y.

A consolidation procedure can produce a reusable knowledge object:

\[
K = C(\{e_i\}_{i \in I}).
\]

The new object may be more compact and easier to reuse than the raw
episodes.

But consolidation can destroy information.

Which event first supported the claim?

Which examples were exceptions?

How uncertain was the conclusion?

Which source supplied the observation?

A semantic representation may preserve the relation and discard the path
that produced it.

This is why semantic memory and provenance must remain separate
coordinates.

The term semantic memory comes from the episodic/semantic distinction
associated with Tulving.

Again, the Atlas imports a useful distinction in information
organization.

It does not claim that a knowledge graph, parameter tensor, or vector
store reproduces a human semantic-memory system.

\section{Episodic and semantic are not storage
media}\label{episodic-and-semantic-are-not-storage-media}

The same physical store can contain both.

A database can hold raw event records and consolidated facts.

A parameterized model can acquire abstractions from many episodes.

An external memory can store summaries whose provenance links back to
source events.

Thus episodic and semantic are not answers to:

where are the bits?

They are answers to:

what relationship does the stored object preserve to the experience that
produced it?

Representation role and storage substrate should not be collapsed.

\section{Associative memory}\label{associative-memory}

Exact addressing asks for a known location or key.

Associative addressing asks for a memory because the query resembles its
content or pattern.

Hopfield's classic content-addressable model makes this distinction
concrete.

A partial or noisy cue can drive the system toward a stored pattern.

In modern systems, association may be implemented through
nearest-neighbor search, attention, learned key matching, hashing,
pattern predicates, energy minimization, or hybrid symbolic/vector
indices.

The Atlas therefore treats associative memory primarily as an access
contract.

A memory can be both episodic and associative.

It can be both semantic and associative.

It can even be parametric and exhibit associative behavior.

The access path does not uniquely identify the storage locus.

\section{One store, three memory
behaviors}\label{one-store-three-memory-behaviors}

The computational witness makes this exact.

Store three immutable records.

A has identifier A, vector (1,0), time 1, and value red.

B has identifier B, vector (0,1), time 2, and value blue.

C has identifier C, vector (1,1), time 3, and value green.

Nothing about the stored records changes.

Now change only the read operator.

Ask for exact identifier B.

The answer is blue.

Ask for the record nearest to

q=(4/5,1/5).

The squared distances to A, B, and C are respectively

2/25, 32/25, 17/25.

So the associative answer is A, red.

Now ask for episodes with time at least 2.

The result is B and C.

The latest is C, green.

Same store.

Three access contracts.

This is why a memory taxonomy that looks only at storage media is
incomplete.

\section{External memory}\label{external-memory}

External memory places state outside the current model parameters and
transient active state. Externality is a locus property; persistence is
a separate lifetime choice.

That can include files, databases, journals, knowledge graphs, vector
stores, tuple spaces, object stores, or differentiable memory matrices.

``External'' is relative to the model/process boundary.

It does not necessarily mean remote over a network.

Memory Networks, Neural Turing Machines, and Differentiable Neural
Computers all make external memory an explicit part of the learned
computational architecture.

Their mechanisms differ.

The shared idea is that stored state can be read and written through an
interface distinct from ordinary parameter access.

That separation opens possibilities unavailable to purely parametric
memory:

rapid writes;

explicit record identity;

large persistent stores;

different retention policies;

independent indexing;

provenance;

sharing across models.

It also creates new systems problems.

\section{External does not mean
trustworthy}\label{external-does-not-mean-trustworthy}

A database can contain false data.

A vector store can return the wrong neighbor.

A file can be stale.

A journal can contain a bad event with perfect provenance.

A knowledge graph can encode an obsolete relation.

An external memory can be unavailable during a partition.

The important gain is controllability of the memory interface.

Truth remains a separate question.

External memory is sometimes described as if retrieval itself solved
hallucination or knowledge reliability.

It does not.

It can make information more inspectable, replaceable, and independently
updateable.

Whether the information is correct still depends on evidence,
provenance, validation, and retrieval quality.

\section{Memory Networks and explicit long-term
state}\label{memory-networks-and-explicit-long-term-state}

Memory Networks were designed around a long-term memory component used
jointly with learned inference.

That is an important architectural separation.

Instead of forcing every relevant fact into parameters or current input,
the system can maintain a memory object that is read and written as part
of the task.

Inference asks what to do with available memory.

Memory management asks what to store and how to expose it.

These can be learned jointly while remaining conceptually distinct.

That separation becomes essential in later chapters, where retrieval
quality and continual update must be analyzed independently.

\section{Differentiable external
memory}\label{differentiable-external-memory}

Neural Turing Machines and Differentiable Neural Computers expose
explicit read/write memory to a neural controller and make addressing
differentiable enough to learn from data.

This demonstrates that memory management can itself become part of
learned computation.

But differentiability does not erase memory semantics.

A read head still has an address rule.

A write still changes stored state.

A usage policy still determines overwrite.

A link structure still changes what sequences can be recovered.

The memory is learnable infrastructure, not magic persistence.

\section{Lifetime is a separate
axis}\label{lifetime-is-a-separate-axis}

Consider four pieces of information:

a token in a current context;

a record in an episodic buffer;

a fact encoded in model weights;

a row in a durable database.

Their lifetimes can differ by many orders of magnitude.

But longer is not automatically better.

A stale fact persisting forever can be worse than a fresh fact retained
briefly.

A replay buffer that never evicts may become computationally unusable.

A context that keeps every previous token may exceed its budget.

Retention is therefore a policy.

The taxonomy records lifetime T separately from semantic role.

\section{Writes have different
costs}\label{writes-have-different-costs}

Parametric writes can require training, editing, or optimization.

Working-memory writes can happen every step.

Episodic writes may be append-like.

Semantic writes may require consolidation.

External stores may support transactional updates.

These differences determine how quickly a system can incorporate new
information.

A system that knows something only after retraining has a very different
adaptation loop from one that can append a validated record immediately.

This is one reason Continual Learning depends on this taxonomy.

Continual learning is not only about retaining old information.

It is about how different write mechanisms interact over time.

\section{Provenance can be preserved or
destroyed}\label{provenance-can-be-preserved-or-destroyed}

Suppose a memory contains:

Paris is the capital of France.

The payload alone does not tell us:

who asserted it;

when;

from which source;

under which transformation;

with what evidence;

whether the source has since been superseded.

An episodic record can retain those fields.

A semantic consolidation step can drop them.

A parameterized model can make them difficult to recover.

An external governed store can preserve them explicitly.

This makes provenance a first-class memory coordinate rather than
decorative metadata.

A memory system that preserves provenance supports different downstream
claims from one that stores only payloads.

\section{Shared memory introduces
coordination}\label{shared-memory-introduces-coordination}

A memory can be private to one process.

It can be shared across agents.

It can be replicated.

It can be transactional.

It can be eventually synchronized.

These choices belong to the sharing coordinate S.

Once several actors write to one memory, the Coordination chapter's
concerns return:

ordering;

atomicity;

duplicate effects;

stale views;

partial failure.

Memory architecture and coordination architecture therefore meet at
shared state.

But they remain distinct.

A shared store tells us where state resides and who can access it.

Coordination semantics tell us how concurrent access behaves.

\section{Failure by memory role}\label{failure-by-memory-role}

Different roles fail differently.

Parametric memory can become stale, entangled, or difficult to edit
without interference.

Working memory can overflow, truncate, or vanish.

Episodic memory can duplicate events, lose identity, or grow without
bound.

Semantic memory can overgeneralize or lose provenance during
consolidation.

Associative memory can return false neighbors or suffer representation
drift.

External memory can become unavailable, inconsistent, stale, or
disconnected from the process that needs it.

These are not one failure called forgetting.

A useful memory diagnosis begins by identifying which coordinate failed.

\section{Memory is not knowledge}\label{memory-is-not-knowledge}

A stored object may be a fact, a hypothesis, a failed attempt, a
prediction, a raw observation, a proof, or an obsolete record.

The memory system does not decide which is true merely by retaining it.

The Evidence chapter taught us to separate claim from support.

The same discipline applies here.

Memory is infrastructure for persistence and access.

Knowledge is an epistemic status that requires more.

This is particularly important for agent systems that retrieve old
outputs from themselves.

Self-retrieval is not self-verification.

\section{What Retrieval may assume}\label{what-retrieval-may-assume}

ATLAS-CH-RETRIEVAL-001 can now start from a stable distinction between
storage and access.

It may assume:

M=(L,W,R,T,A,U,P,S);

exact-key versus associative addressing;

parametric versus external locus;

working versus durable lifetime;

episodic versus semantic organization;

provenance as an independent coordinate.

Its job is to develop retrieval itself:

vector search;

symbolic lookup;

hybrid retrieval;

multi-index access;

key-value memory;

ranking and selection.

It need not rebuild the taxonomy.

\section{What Continual Learning may
assume}\label{what-continual-learning-may-assume}

ATLAS-CH-CONTINUAL-001 may assume:

parametric write semantics;

episodic records;

semantic consolidation;

retention and mutability coordinates;

the distinction between preserving records and preserving behavior.

Its job is to study what happens when learning continues:

interference;

catastrophic forgetting;

replay;

consolidation;

regularization;

parameter isolation.

Again, the taxonomy becomes useful because ``remembering'' can now mean
several different things.

\section{The boundary to remember}\label{the-boundary-to-remember-1}

The most useful question is not:

Does the system have memory?

Almost every nontrivial adaptive system does.

The useful questions are:

Where does the state live?

How is it written?

How is it read?

How long does it survive?

How is it addressed?

How can it change?

What provenance survives with it?

Who else can see or modify it?

Once those coordinates are explicit, the familiar labels become precise.

Parametric memory tells us that information is carried in learned
parameters.

Working memory tells us that state is actively available for current
computation.

Episodic memory tells us that event identity matters.

Semantic memory tells us that generalized content can survive without
the full originating episode.

Associative memory tells us that access can be content-driven.

External memory tells us that state can live outside the current model
and active context; whether it is durable is a separate retention
decision.

None of these labels is the whole memory system.

Together they give us enough structure to ask the next question:

given that something has been stored, how should the system find the
right thing again?

That is where retrieval begins.

\section{References used in this
chapter}\label{references-used-in-this-chapter-45}

Exact source identities and authority scopes are pinned in:

sources/source-locks/ATLAS-CH-MEMTAX-001.yaml

The external basis includes Baddeley and Hitch for working-memory
terminology lineage, Tulving for episodic/semantic terminology lineage,
Hopfield for content-addressable memory, Memory Networks, Neural Turing
Machines, Differentiable Neural Computers, Neural Episodic Control, and
empirical work on factual associations stored in language-model
parameters.

\chapter{Retrieval and Associative
Access}\label{retrieval-and-associative-access}

\textbf{Epistemic status:} established
information-retrieval/database/vector-search foundations + audited
Memory Taxonomy and Attention prerequisites + Atlas synthesis + exact
finite witness.\\
\textbf{Specification:}
manuscript/specifications/ATLAS-CH-RETRIEVAL-001.md\\
\textbf{Formal packet:}
mathematics/derivations/ATLAS-CH-RETRIEVAL-001-DERIVATIONS.md\\
\textbf{Computational witness:}
mathematics/computational-witnesses/ATLAS-CW-RETRIEVAL-001.md\\
\textbf{Source lock:} sources/source-locks/ATLAS-CH-RETRIEVAL-001.yaml

A memory is useful only if a system can find the right part of it.

That sounds obvious.

The difficulty begins when we ask what \textbf{right} means.

A record can be right because:

\begin{itemize}
\tightlist
\item
  its exact key matches;
\item
  its metadata satisfies a predicate;
\item
  its vector is nearest to the query;
\item
  its lexical score is highest;
\item
  it survives a symbolic filter;
\item
  it wins a fused rank;
\item
  it is recent;
\item
  it comes from an admissible source.
\end{itemize}

These are different retrieval contracts.

The governing rule of this chapter is:

\begin{quote}
retrieval is a declared read operation over records, not a synonym for
relevance, truth, provenance, or usefulness.
\end{quote}

The same immutable records can support several retrieval behaviours.

That is not inconsistency.

It is the architecture.

\section{Retrieval starts where the memory taxonomy
stops}\label{retrieval-starts-where-the-memory-taxonomy-stops}

The Memory Taxonomy chapter wrote a memory object as

\texttt{M=(L,W,R,T,A,U,P,S)}

for storage locus, write path, read path, lifetime, addressability,
mutability, provenance, and sharing.

Retrieval is chiefly about the read side:

\texttt{R}

and the addressability coordinate:

\texttt{A}.

But the other coordinates still matter.

A result may be easy to retrieve but stale.

A result may be relevant but lack provenance.

An index may be external while the record is ephemeral.

A query may be associative while the store itself is exact-key
addressable.

So retrieval must not collapse the rest of the memory contract.

\section{A retrieval contract}\label{a-retrieval-contract}

Let \(\mathcal{D}\) be the underlying record set.

Write a retrieval contract as

\[
\operatorname{Retr}=(\mathcal{D}, \mathcal{Q}, F, s, \pi, k, \mathcal{O}).
\]

Here:

\begin{itemize}
\tightlist
\item
  \(\mathcal{Q}\) is the query space;
\item
  \(F(q, \mathcal{D})\) is an optional admissibility/filter operator;
\item
  \(s(q, d)\) is a score, distance, or match relation;
\item
  \(\pi\) is a ranking/selection rule;
\item
  \(k\) is a truncation or budget;
\item
  \(\mathcal{O}\) is the returned projection or set of fields.
\end{itemize}

This representation makes one fact explicit:

\begin{quote}
a retrieval result is defined by more than a score.
\end{quote}

The candidate set matters.

The filter matters.

The tie rule matters.

The top-k budget matters.

The fields returned matter.

Two systems can use the same embeddings and still implement different
retrieval semantics.

\section{Exact-key lookup}\label{exact-key-lookup}

The simplest access contract is exact addressing.

Let every record carry a unique key:

\texttt{key(d)\ in\ K}.

Then exact-key retrieval asks for the unique record satisfying

\texttt{key(d)=u}.

This is not approximate.

It is not semantic.

It does not ask whether the key's record resembles the query.

A key is an address.

Dynamo is a large-scale systems example of key-value access in which a
caller supplies a key and the system manages distributed storage and
availability around that interface (DeCandia et al. 2007).

This chapter uses only the access distinction.

Replication, versioning, failure recovery, and consistency policies are
separate systems concerns.

\section{Key-value retrieval is not semantic
retrieval}\label{key-value-retrieval-is-not-semantic-retrieval}

Suppose key

\texttt{invoice-2026-1042}

maps to a record.

The lookup can be exact even if:

\begin{itemize}
\tightlist
\item
  the record is obsolete;
\item
  the payload is wrong;
\item
  another record would be semantically more useful;
\item
  the record has poor provenance.
\end{itemize}

Exactness refers to the address relation.

It does not certify the payload.

This distinction will recur throughout the chapter.

\section{Symbolic retrieval}\label{symbolic-retrieval}

A symbolic query specifies conditions over structured fields.

Let

\texttt{P\_q(d)}

be a predicate compiled from query \texttt{q}.

Then symbolic retrieval returns

\texttt{\{d\ in\ D\ :\ P\_q(d)\}}.

Examples include:

\begin{itemize}
\tightlist
\item
  \texttt{year\textgreater{}=2024};
\item
  \texttt{type="paper"};
\item
  \texttt{source\_id\ in\ approved\_sources};
\item
  \texttt{status="audited"};
\item
  conjunctions and disjunctions of structured conditions.
\end{itemize}

Codd's relational model provides the foundational idea that data can be
addressed through declarative relations and operations rather than by
exposing physical storage layout (Codd 1970).

The Atlas takes a narrow lesson:

\begin{quote}
symbolic admissibility can be specified independently of physical
storage.
\end{quote}

\section{A predicate is not a
ranking}\label{a-predicate-is-not-a-ranking}

The result of a predicate can be a set.

If a caller needs an ordered list, another rule is required.

For example:

\begin{itemize}
\tightlist
\item
  sort by date;
\item
  sort by score;
\item
  sort by source priority;
\item
  rank by a learned model;
\item
  apply a deterministic record-ID tie break.
\end{itemize}

So a system that says

\begin{quote}
retrieve all approved papers
\end{quote}

has specified admissibility.

It has not yet specified which approved paper should come first.

Selection and ranking are distinct operators.

\section{Vector retrieval}\label{vector-retrieval}

Vector retrieval maps queries and records into a shared coordinate
space.

Let

\[
\phi_Q(q) \in \mathbb{R}^m
\]

and

\[
\phi_{\mathcal{D}}(d) \in \mathbb{R}^m.
\]

Then a distance-based retrieval rule can be written as

\[
d^* \in \arg\min_{d \in \mathcal{D}} \delta(\phi_Q(q), \phi_{\mathcal{D}}(d)).
\]

Or a similarity rule can maximize a score.

Salton, Wong, and Yang's vector-space model is a foundational
information-retrieval construction in which stored objects and queries
are compared through vector representations (Salton et al. 1975).

Modern learned embeddings change how the vectors are obtained.

They do not change the core fact that ranking is relative to a
representation and a similarity geometry.

\section{Similarity is
representation-relative}\label{similarity-is-representation-relative}

Suppose two embedding models map the same records to different
coordinates.

The nearest record can change.

Suppose the embeddings remain fixed but cosine similarity is replaced by
Euclidean distance.

The ranking can change again.

Suppose vectors are normalized before scoring.

The ranking may change again.

Therefore the statement

\begin{quote}
record A is nearest
\end{quote}

is incomplete without:

\begin{itemize}
\tightlist
\item
  query embedding;
\item
  record embedding;
\item
  normalization rule;
\item
  distance or similarity;
\item
  candidate set.
\end{itemize}

Nearest is never representation-free.

\section{A similarity score is not automatically relevance
probability}\label{a-similarity-score-is-not-automatically-relevance-probability}

A dense retriever can emit a score:

\texttt{s(q,d)}.

That score can be useful for ordering candidates.

But unless a separate calibration argument is established, it does not
mean

\texttt{P(relevant\ \textbar{}\ q,d)}.

Scores may be:

\begin{itemize}
\tightlist
\item
  inner products;
\item
  cosine similarities;
\item
  negative distances;
\item
  logits;
\item
  learned ranking outputs.
\end{itemize}

Their numerical scale can be arbitrary.

A score of 0.9 need not mean 90 percent relevant.

Ranking and probability calibration are different tasks.

\section{Top-k is a budgeted truncation
operator}\label{top-k-is-a-budgeted-truncation-operator}

Suppose candidate order is

\texttt{d\_(1),...,d\_(n)}.

Top-k returns the ordered list

\texttt{(d\_(1),...,d\_(k))}.

This is a computational and interface decision.

It does not prove:

\begin{itemize}
\tightlist
\item
  the omitted records are irrelevant;
\item
  the returned records are all relevant;
\item
  the k-th score marks a natural semantic boundary.
\end{itemize}

Top-k converts an ordering into a bounded ordered working list; its
underlying selected set is a separate view when set semantics are
needed.

That is useful precisely because downstream systems cannot consume
everything.

\section{Tie rules matter for exact
replay}\label{tie-rules-matter-for-exact-replay}

Two records can have the same score.

If the system needs reproducible ordered output, it must declare a tie
rule.

Examples include:

\begin{itemize}
\tightlist
\item
  stable record ID;
\item
  insertion order;
\item
  timestamp;
\item
  secondary score.
\end{itemize}

Without such a rule, two conforming implementations can return different
top-k lists while agreeing on all primary scores.

This is not always a problem.

But it matters when exact replay is part of the evidence contract.

\section{Exact nearest neighbour}\label{exact-nearest-neighbour}

Exact nearest-neighbor search asks for the global optimum under the
declared metric and candidate set.

For query \texttt{q}:

\texttt{NN(q,D)\ =\ argmin\_\{d\ in\ D\}\ delta(q,d)}.

If the store is small, exhaustive comparison may be acceptable.

At large scale, exhaustive search can become too expensive.

That motivates indexes and approximation.

\section{Approximate nearest
neighbour}\label{approximate-nearest-neighbour}

Approximate nearest-neighbor search changes the search procedure.

The target remains similarity search.

The algorithm accepts some risk of missing the exact global nearest
candidate in exchange for improvements such as:

\begin{itemize}
\tightlist
\item
  lower latency;
\item
  lower query compute;
\item
  better scalability;
\item
  bounded memory structures;
\item
  practical billion-scale search.
\end{itemize}

Malkov and Yashunin's HNSW algorithm is one graph-based approximate
nearest-neighbor construction (Malkov and Yashunin 2020).

Its layered navigable graph is a search structure.

The crucial Atlas boundary is:

\begin{quote}
an ANN result is not silently identical to the exact nearest-neighbor
result.
\end{quote}

\section{Search approximation and representation error are
different}\label{search-approximation-and-representation-error-are-different}

Imagine the truly task-relevant record is semantically important but
badly embedded.

Exact nearest-neighbor search can retrieve the wrong semantic record
perfectly.

That is not an index failure.

The index exactly solved the wrong geometric proxy.

Now reverse the case.

Suppose the embedding geometry puts the relevant record closest to the
query, but an approximate index misses it.

That is a search/index approximation failure.

So separate:

\begin{enumerate}
\def\labelenumi{\arabic{enumi}.}
\tightlist
\item
  representation error;
\item
  search/index error;
\item
  downstream relevance error.
\end{enumerate}

Without this separation, retrieval diagnostics become confused.

\section{Sparse and lexical
retrieval}\label{sparse-and-lexical-retrieval}

Not every retrieval system needs dense embeddings.

Classical text retrieval can score overlap and term evidence through
sparse representations.

The probabilistic relevance framework and BM25 family are an important
lineage for such ranking (Robertson and Zaragoza 2009).

Sparse lexical systems can be attractive because their matching evidence
is often easier to inspect:

\begin{itemize}
\tightlist
\item
  term occurrence;
\item
  term frequency;
\item
  document frequency;
\item
  field structure.
\end{itemize}

Dense systems can recover conceptual proximity that exact lexical
overlap misses.

Neither representation dominates universally.

\section{Symbolic, sparse, and dense are different
axes}\label{symbolic-sparse-and-dense-are-different-axes}

It is tempting to place all retrieval into two boxes:

\begin{itemize}
\tightlist
\item
  lexical;
\item
  vector.
\end{itemize}

That misses structured predicates.

A symbolic filter such as

\texttt{source="primary"\ AND\ year\textgreater{}=2024}

does not need to behave like BM25 or embedding similarity.

It declares admissibility.

A complete retrieval architecture may therefore combine:

\begin{itemize}
\tightlist
\item
  symbolic filtering;
\item
  sparse lexical scoring;
\item
  dense similarity;
\item
  recency weighting;
\item
  source constraints.
\end{itemize}

The composition rule is the system.

\section{One store can have many
indices}\label{one-store-can-have-many-indices}

A record can be reachable through:

\begin{itemize}
\tightlist
\item
  primary key;
\item
  metadata index;
\item
  inverted lexical index;
\item
  vector index;
\item
  temporal index;
\item
  provenance index.
\end{itemize}

These do not necessarily duplicate the logical record.

They are access structures.

This gives a useful invariant:

\texttt{record\_id\ !=\ index\_entry\_id}.

One logical record may have several index entries.

Changing or deleting the record can therefore require coordinated index
maintenance.

\section{Multi-index retrieval}\label{multi-index-retrieval}

The phrase \textbf{multi-index} can mean more than one thing.

At the broad architecture level, the Atlas uses it for multiple access
structures over one record set.

At the algorithm level, Babenko and Lempitsky's inverted multi-index is
a specific similarity-search construction based on product quantization
(Babenko and Lempitsky 2012).

These meanings should not be conflated.

The paper establishes one concrete structured vector-search index.

The Atlas generalization is an architectural synthesis:

\begin{quote}
one record identity can participate in several independent retrieval
access paths.
\end{quote}

\section{Hybrid retrieval requires composition
semantics}\label{hybrid-retrieval-requires-composition-semantics}

The word \textbf{hybrid} is often too vague.

A hybrid retriever might mean:

\begin{itemize}
\tightlist
\item
  symbolic filter, then vector rank;
\item
  lexical candidate generation, then dense reranking;
\item
  dense and sparse score fusion;
\item
  vector candidates intersected with policy constraints;
\item
  several retrievers feeding a learned reranker.
\end{itemize}

These pipelines are not equivalent.

A serious description should state the operator order.

\section{Exact witness store}\label{exact-witness-store}

Use three immutable records.

A:

\begin{itemize}
\tightlist
\item
  id = A;
\item
  kind = paper;
\item
  year = 2024;
\item
  vector = \texttt{(1,0)};
\item
  value = alpha.
\end{itemize}

B:

\begin{itemize}
\tightlist
\item
  id = B;
\item
  kind = paper;
\item
  year = 2022;
\item
  vector = \texttt{(4/5,3/5)};
\item
  value = beta.
\end{itemize}

C:

\begin{itemize}
\tightlist
\item
  id = C;
\item
  kind = note;
\item
  year = 2025;
\item
  vector = \texttt{(24/25,1/25)};
\item
  value = gamma.
\end{itemize}

Query vector:

\texttt{q=(19/20,1/20)}.

Use squared Euclidean distance.

Nothing about the records changes across the following reads.

\section{Exact vector distances}\label{exact-vector-distances}

To A:

\texttt{d(q,A)\^{}2=1/200}.

To B:

\texttt{d(q,B)\^{}2=13/40}.

To C:

\texttt{d(q,C)\^{}2=1/5000}.

Therefore:

\texttt{C,A,B}

is the exact vector ranking.

Unfiltered top-1 is C.

\section{Exact-key read of the same
store}\label{exact-key-read-of-the-same-store}

Ask for:

\texttt{id=B}.

The answer is B.

Vector proximity is irrelevant.

The store did not change.

The access contract did.

This is the same structural lesson established abstractly by MEMTAX-001,
now made specific to retrieval.

\section{Symbolic read of the same
store}\label{symbolic-read-of-the-same-store}

Ask for:

\texttt{year\textgreater{}=2024}.

A qualifies.

B does not.

C qualifies.

So the result set is:

\texttt{\{A,C\}}.

This query defines admissibility but not an order.

A secondary ranking clause would be required to choose which record
comes first.

\section{Hybrid retrieval does not generally
commute}\label{hybrid-retrieval-does-not-generally-commute}

Now require:

\texttt

and also rank by vector distance.

There are at least two plausible pipelines.

\subsection{Filter then rank}\label{filter-then-rank}

First keep papers:

\texttt{\{A,B\}}.

Then rank by distance.

A is closer.

Top-1 result:

\texttt{A}.

\subsection{Rank then filter}\label{rank-then-filter}

First rank the full store.

Top-1 is C.

Then apply the paper predicate.

C fails.

The result is empty.

Thus:

\texttt{FilterThenRank\ !=\ RankThenFilter}.

The difference arises only from operator order and top-k truncation.

\section{Why filter-then-rank often has cleaner
semantics}\label{why-filter-then-rank-often-has-cleaner-semantics}

If a symbolic predicate means

\begin{quote}
only these records are admissible,
\end{quote}

then applying it before ranking ensures the ranking operates inside the
permitted set.

This is particularly important for constraints such as:

\begin{itemize}
\tightlist
\item
  approved sources;
\item
  tenant boundaries;
\item
  date windows;
\item
  document status;
\item
  legal access permissions.
\end{itemize}

But the Atlas does not declare filter-first universally superior.

Other pipelines can be legitimate.

The requirement is to state what the operators mean.

\section{Score fusion is another hybrid
family}\label{score-fusion-is-another-hybrid-family}

Suppose a document has sparse score

\texttt{s\_sparse(q,d)}

and dense score

\texttt{s\_dense(q,d)}.

A fused ranker may use:

\texttt{s\_hybrid\ =\ lambda\ s\_sparse\ +\ (1-lambda)s\_dense}

after suitable normalization.

This immediately creates new questions:

\begin{itemize}
\tightlist
\item
  are the score scales commensurate?
\item
  how is \texttt{lambda} chosen?
\item
  what happens when one retriever has no score?
\item
  are ranks fused instead of raw scores?
\item
  is the rule query-dependent?
\end{itemize}

Calling the system hybrid does not answer them.

\section{Retrieval can be staged}\label{retrieval-can-be-staged}

A large corpus may use a cheap first stage to generate candidates and a
more expensive second stage to rerank them.

Write:

\texttt{D\ -\textgreater{}\ C\_k(q)\ -\textgreater{}\ Rank\_expensive(C\_k(q))}.

This reduces cost.

It can also introduce a hard ceiling:

if the relevant record never enters the candidate set, the reranker
cannot recover it.

Candidate recall and reranker quality are therefore separate quantities.

\section{Attention and retrieval overlap at the operator
level}\label{attention-and-retrieval-overlap-at-the-operator-level}

ATTNOP-001 showed attention as a state-dependent operator built from
query-key scores and values.

Retrieval can share the pattern:

\texttt{query\ -\textgreater{}\ candidate\ scores\ -\textgreater{}\ selection\ -\textgreater{}\ returned\ values}.

This resemblance is real.

It does not make attention and retrieval architecturally identical.

\section{Retrieval can exceed the active
field}\label{retrieval-can-exceed-the-active-field}

Attention normally operates over the keys and values supplied to the
active computation.

A retriever can search:

\begin{itemize}
\tightlist
\item
  a database;
\item
  a document collection;
\item
  a vector index;
\item
  a file store;
\item
  a shared organizational memory.
\end{itemize}

The persistent corpus can be much larger than the active context.

Retrieval therefore often acts as the bridge between long-lived memory
and bounded working context.

\section{Retrieval can use hard
identity}\label{retrieval-can-use-hard-identity}

Attention weights usually mix supplied values.

Retrieval can instead return explicit record identities.

That creates capabilities such as:

\begin{itemize}
\tightlist
\item
  source citation;
\item
  stable replacement;
\item
  deletion;
\item
  access control;
\item
  provenance lookup;
\item
  deduplication.
\end{itemize}

The record can remain independently addressable after retrieval.

This becomes central in the External-Memory chapter.

\section{Retrieval is not causal
explanation}\label{retrieval-is-not-causal-explanation}

A high retrieval score says:

\begin{quote}
this candidate ranked highly under this retrieval contract.
\end{quote}

It does not say:

\begin{quote}
this record caused the final output.
\end{quote}

A downstream model may:

\begin{itemize}
\tightlist
\item
  ignore it;
\item
  contradict it;
\item
  combine it with other records;
\item
  distort it;
\item
  fail for unrelated reasons.
\end{itemize}

ATTNOP-001 already enforced the analogous boundary for attention
coefficients.

The same discipline applies here.

\section{Retrieval and provenance}\label{retrieval-and-provenance}

Suppose record A ranks first.

The score does not tell us:

\begin{itemize}
\tightlist
\item
  where A came from;
\item
  who authored it;
\item
  whether it has been superseded;
\item
  which transformation produced its embedding;
\item
  whether it was audited.
\end{itemize}

Those belong to the record/provenance system.

A governed retriever should preserve a path from result to stable record
identity and source metadata when downstream claims depend on
provenance.

\section{Rank does not create
authority}\label{rank-does-not-create-authority}

A record can rank first because its vector is close.

That does not promote it to:

\begin{itemize}
\tightlist
\item
  verified;
\item
  approved;
\item
  authoritative;
\item
  current.
\end{itemize}

A retrieval index is not an adjudicator.

The Evidence and Replay chapters remain relevant whenever retrieved
material is used to support claims.

\section{Retrieval correctness, relevance, and
usefulness}\label{retrieval-correctness-relevance-and-usefulness}

Three distinct questions are:

\begin{enumerate}
\def\labelenumi{\arabic{enumi}.}
\tightlist
\item
  \textbf{contract correctness:} did the retriever return what its
  algorithm specifies?
\item
  \textbf{task relevance:} does the record bear on the user's task?
\item
  \textbf{downstream usefulness:} does giving the record to the consumer
  improve the objective?
\end{enumerate}

A result can be correct under the retrieval contract and irrelevant to
the task.

A relevant record can be too long for the available context.

A useful summary can rank below a less useful full document under raw
similarity.

Evaluation should therefore name which layer it measures.

\section{Recall and precision require a relevance
judgment}\label{recall-and-precision-require-a-relevance-judgment}

Information retrieval metrics such as precision and recall need a
declared notion of relevance.

For retrieved set \(S\) and relevant ground-truth set \(G\):

\[
\operatorname{precision} = \frac{|S \cap G|}{|S|}
\]

and

\[
\operatorname{recall} = \frac{|S \cap G|}{|G|}
\]

when denominators are nonzero.

The formulas are simple.

The hard question is how \texttt{G} is defined.

Human labels, task outcomes, source policies, and benchmark annotations
can produce different relevance notions.

Metrics inherit that choice.

\section{ANN recall is not semantic
recall}\label{ann-recall-is-not-semantic-recall}

Approximate-nearest-neighbor evaluation often asks whether the ANN
method recovered exact nearest neighbors under a metric.

That is an index-quality question.

Semantic retrieval recall asks whether task-relevant records were found.

These should not be conflated.

An ANN system can have perfect recall relative to exact vector nearest
neighbors while the vector model itself has poor semantic recall.

Again:

representation and search are separate layers.

\section{Retrieval-augmented generation is downstream
composition}\label{retrieval-augmented-generation-is-downstream-composition}

Lewis et al.~combine a parametric generator with a dense vector index
used as non-parametric memory (Lewis et al. 2020).

This is a useful architecture example because it makes retrieval a
distinct stage in a larger model.

The retrieval stage can be inspected separately from generation.

But retrieval augmentation does not create a theorem that generated
answers are true.

The pipeline can fail at:

\begin{itemize}
\tightlist
\item
  query representation;
\item
  index search;
\item
  relevance;
\item
  context assembly;
\item
  generation;
\item
  interpretation.
\end{itemize}

Each layer deserves its own evidence.

\section{Retrieval is not memory
writing}\label{retrieval-is-not-memory-writing}

A read can be excellent while the memory is stale.

A memory can be correct while the retriever is weak.

A retriever can repeatedly surface one record without changing the
store.

This preserves the MEMTAX distinction:

\begin{itemize}
\tightlist
\item
  write path;
\item
  read path;
\item
  lifetime;
\item
  mutability.
\end{itemize}

Retrieval is one side of the memory interface.

\section{Retrieval under updates}\label{retrieval-under-updates}

When records change, indexes may need maintenance.

A system should be able to state whether updates are:

\begin{itemize}
\tightlist
\item
  synchronous;
\item
  asynchronous;
\item
  transactional;
\item
  eventually indexed;
\item
  versioned.
\end{itemize}

A record can exist in the source store before it becomes visible to one
index.

That creates a freshness gap.

The chapter records the issue but does not develop distributed
consistency.

That belongs to systems/coordination layers.

\section{Deletion is also a retrieval
property}\label{deletion-is-also-a-retrieval-property}

If a record is deleted from the canonical store but remains in an index,
retrieval can surface a ghost.

If an index deletes the vector but the canonical record remains, one
access path disappears while others survive.

So deletion semantics require:

\begin{itemize}
\tightlist
\item
  canonical record identity;
\item
  index maintenance;
\item
  visibility rules.
\end{itemize}

This is another reason record and index identity must remain separate.

\section{Multi-index consistency}\label{multi-index-consistency}

Suppose one record has:

\begin{itemize}
\tightlist
\item
  primary-key entry;
\item
  lexical postings;
\item
  vector node;
\item
  time index.
\end{itemize}

An update to the record can leave these access structures temporarily
inconsistent.

The retrieval layer should therefore distinguish:

\begin{itemize}
\tightlist
\item
  logical record version;
\item
  index version;
\item
  query-time snapshot.
\end{itemize}

The chapter does not prescribe a transaction protocol.

It makes the state distinction visible.

\section{Retrieval budgets shape
cognition}\label{retrieval-budgets-shape-cognition}

A downstream model rarely receives the full store.

It receives a working set.

The retrieval budget \texttt{k} therefore shapes what information
becomes available for reasoning.

A small \texttt{k} can exclude complementary evidence.

A large \texttt{k} can consume context and introduce distractors.

So retrieval is not only a search problem.

It is a compilation boundary between persistent memory and active
computation.

The later Context Compilation chapter will build on this.

\section{The nearest record need not be the best context
item}\label{the-nearest-record-need-not-be-the-best-context-item}

Suppose one highly similar document repeats the query without adding
evidence.

A slightly less similar record may contain the decisive fact.

Vector similarity can therefore be a useful proxy without being the
downstream utility function.

Possible downstream reranking signals include:

\begin{itemize}
\tightlist
\item
  novelty;
\item
  coverage;
\item
  provenance quality;
\item
  contradiction;
\item
  diversity;
\item
  recency;
\item
  task-specific expected utility.
\end{itemize}

The Atlas keeps these as separate design layers.

\section{Diversification changes the
objective}\label{diversification-changes-the-objective}

Top-k nearest-neighbor retrieval can return near-duplicates.

A diversified retriever may intentionally sacrifice some score to cover
different aspects of a query.

Then its objective is no longer:

\begin{quote}
return the k highest individual scores.
\end{quote}

It becomes a set-selection problem.

The distinction matters when evaluating retrieval quality.

Individual rank and set utility are different objects.

\section{Retrieval over governed
records}\label{retrieval-over-governed-records}

A governed store can expose symbolic fields such as:

\begin{itemize}
\tightlist
\item
  epistemic class;
\item
  review state;
\item
  source identity;
\item
  supersession state;
\item
  authority;
\item
  validity interval.
\end{itemize}

These fields can participate in filters.

That makes retrieval policy-aware.

But the filter must not invent those statuses.

It can only consume statuses established elsewhere.

Retrieval applies governance metadata.

It does not create governance authority.

\section{Failure modes}\label{failure-modes-10}

\subsection{Address conflation}\label{address-conflation}

Treat exact key match as semantic relevance.

\subsection{Similarity laundering}\label{similarity-laundering}

Treat nearest vector as correct answer.

\subsection{Score calibration fiction}\label{score-calibration-fiction}

Interpret arbitrary ranking scores as probabilities.

\subsection{Top-k completeness
fiction}\label{top-k-completeness-fiction}

Treat omitted records as irrelevant.

\subsection{ANN/exact conflation}\label{annexact-conflation}

Describe approximate search output as the exact global nearest neighbour
without verification.

\subsection{Hybrid ambiguity}\label{hybrid-ambiguity}

Say ``hybrid retrieval'' without specifying operator order or fusion.

\subsection{Index/record conflation}\label{indexrecord-conflation}

Treat multiple index entries as multiple independent records.

\subsection{Provenance laundering}\label{provenance-laundering}

Infer source authority from retrieval rank.

\subsection{Retrieval/generation
conflation}\label{retrievalgeneration-conflation}

Treat a retrieval-augmented answer as correct because retrieval
occurred.

\subsection{Read/write conflation}\label{readwrite-conflation}

Treat a retrieval improvement as a memory freshness improvement.

\section{A practical retrieval
ledger}\label{a-practical-retrieval-ledger}

Before evaluating a retriever, record:

{\def\LTcaptype{none} % do not increment counter
\begin{longtable}[]{@{}ll@{}}
\toprule\noalign{}
Field & Question \\
\midrule\noalign{}
\endhead
\bottomrule\noalign{}
\endlastfoot
corpus & What records are eligible in principle? \\
record identity & What is the stable logical record ID? \\
query & How is the query represented? \\
filter & What makes a record admissible? \\
index & Which access structure is searched? \\
representation & Sparse, dense, symbolic, mixed? \\
score & What quantity orders candidates? \\
approximation & Exact or approximate search? \\
budget & What is k or the candidate limit? \\
tie rule & How are equal scores ordered? \\
projection & Which fields are returned? \\
provenance & How is source identity preserved? \\
relevance target & What labels or task define relevance? \\
downstream target & What consumer/usefulness metric matters? \\
\end{longtable}
}

This is the minimum grammar needed to interpret a retrieval result.

\section{What the exact witness
establishes}\label{what-the-exact-witness-establishes-3}

The companion witness proves four facts about one immutable finite
store.

First:

\texttt{id=B}

returns B by exact key.

Second:

\texttt{year\textgreater{}=2024}

returns the symbolic set

\texttt{\{A,C\}}.

Third, vector ranking under the declared squared Euclidean geometry is:

\texttt{C,A,B}.

Fourth, with the paper predicate:

\begin{itemize}
\tightlist
\item
  filter-then-rank top-1 returns A;
\item
  rank-top-1-then-filter returns nothing.
\end{itemize}

Therefore retrieval semantics depend on the read contract and operator
order even before learned embeddings, ANN approximations, or large-scale
infrastructure are introduced.

\section{Downstream handoff}\label{downstream-handoff-6}

\textbf{The External-Memory Thesis --- ATLAS-CH-EXTMEM-001} may now
assume:

\begin{itemize}
\tightlist
\item
  exact-key versus associative access;
\item
  symbolic predicate retrieval;
\item
  vector similarity retrieval;
\item
  top-k semantics;
\item
  ANN versus exact-NN distinction;
\item
  multiple indices over one record identity;
\item
  explicit hybrid-composition rules;
\item
  provenance preservation;
\item
  retrieval correctness versus relevance versus downstream usefulness.
\end{itemize}

The downstream chapter must independently argue:

\begin{itemize}
\tightlist
\item
  which information should leave model parameters;
\item
  which memory loci should persist;
\item
  how shared memory changes system architecture;
\item
  what update/provenance/governance contract makes external memory
  preferable.
\end{itemize}

Retrieval provides access.

It does not by itself establish the case for externalization.

\section{References used in this
chapter}\label{references-used-in-this-chapter-46}

\begin{itemize}
\tightlist
\item
  Salton, Wong, and Yang, \emph{A Vector Space Model for Automatic
  Indexing} (Salton et al. 1975).
\item
  Codd, \emph{A Relational Model of Data for Large Shared Data Banks}
  (Codd 1970).
\item
  Robertson and Zaragoza, \emph{The Probabilistic Relevance Framework:
  BM25 and Beyond} (Robertson and Zaragoza 2009).
\item
  Malkov and Yashunin, \emph{Efficient and Robust Approximate Nearest
  Neighbor Search Using Hierarchical Navigable Small World Graphs}
  (Malkov and Yashunin 2020).
\item
  Babenko and Lempitsky, \emph{The Inverted Multi-Index} (Babenko and
  Lempitsky 2012).
\item
  DeCandia et al., \emph{Dynamo: Amazon's Highly Available Key-value
  Store} (DeCandia et al. 2007).
\item
  Lewis et al., \emph{Retrieval-Augmented Generation for
  Knowledge-Intensive NLP Tasks} (Lewis et al. 2020).
\end{itemize}

Exact source identities and authority boundaries are recorded in:

sources/source-locks/ATLAS-CH-RETRIEVAL-001.yaml

\chapter{Continual Learning and
Forgetting}\label{continual-learning-and-forgetting}

\textbf{Epistemic status:} established continual-learning mechanisms +
audited memory/optimization prerequisites + Atlas synthesis + exact
scalar witness.\\
\textbf{Specification:}
manuscript/specifications/ATLAS-CH-CONTINUAL-001.md\\
\textbf{Formal packet:}
mathematics/derivations/ATLAS-CH-CONTINUAL-001-DERIVATIONS.md\\
\textbf{Computational witness:}
mathematics/computational-witnesses/ATLAS-CW-CONTINUAL-001.md\\
\textbf{Source lock:} sources/source-locks/ATLAS-CH-CONTINUAL-001.yaml

A learner can become better at what it sees now and worse at what it
knew before.

That sentence contains the central difficulty of continual learning.

Ordinary optimization asks how to improve an objective.

Continual learning asks what happens when the objective, data
distribution, label space, task context, or available evidence changes
over time while the same learning system must preserve useful earlier
capability.

The problem is not simply that parameters move.

Parameters are supposed to move.

The problem is that later updates can alter earlier behavior in ways the
learner no longer has enough evidence, capacity, or structural
protection to repair.

This chapter develops continual learning as a problem of
\textbf{controlled interference under sequential updates}.

Its central rule is:

\begin{quote}
forgetting is a behavioral statement under an evaluation protocol, not a
synonym for parameter change.
\end{quote}

That rule keeps the discussion grounded.

\section{Sequential learning creates a
history}\label{sequential-learning-creates-a-history}

Suppose contexts arrive in order

\texttt{1,2,...,T}.

After training through context \texttt{i}, evaluate the current model on
every context seen so far.

Write

\texttt{R\_\{i,j\}}

for the score on context \texttt{j} after training through context
\texttt{i}.

Then the basic object is not one number.

For retention and forgetting, the mandatory record is a lower-triangular
performance-through-time matrix over contexts already encountered.

For three contexts it has the form

\texttt{R\_\{1,1\}}

\texttt{R\_\{2,1\},\ R\_\{2,2\}}

\texttt{R\_\{3,1\},\ R\_\{3,2\},\ R\_\{3,3\}}.

If a protocol also evaluates not-yet-trained contexts, those
\texttt{j\textgreater{}i} entries can complete a full matrix for
forward-transfer analysis. Forward transfer additionally needs a
declared untrained or reference baseline for the future context; it
cannot be inferred from the lower triangle above.

The performance record can reveal things a final average hides.

A later task may damage one earlier task and improve another.

A model may preserve average accuracy while destroying a rare context.

A method may learn new tasks slowly but retain old ones almost
perfectly.

Another may adapt rapidly and forget sharply.

All of those behaviors can collapse to similar final averages.

So continual learning should begin from the trajectory of performance,
not only its endpoint.

\section{An Atlas forgetting
summary}\label{an-atlas-forgetting-summary}

For a sequence with \(T \ge 2\), and for an earlier context \(j < T\),
define its best prior score

\[
B_j = \max_{k=j,\dots,T-1} R_{k,j}.
\]

Define endpoint forgetting

\[
F_j = \max(0, B_j - R_{T,j}).
\]

Then define

\[
F_{\mathrm{avg}} = \frac{1}{T-1} \sum_{j < T} F_j
\]

and

\[
F_{\mathrm{max}} = \max_{j < T} F_j.
\]

These are Atlas summaries.

They are intentionally modest.

Cross-context aggregation is meaningful only when the context scores
share a commensurate scale or a declared normalization. If different
contexts use incomparable metrics or ranges, retain the per-context
values instead of averaging or taking a cross-context maximum.

They do not claim to be the one canonical continual-learning metric.

Lopez-Paz and Ranzato emphasize that continual-learning evaluation
should include transfer behavior, not only final task accuracy
(Lopez-Paz and Ranzato 2017).

That matters because the nonnegative \texttt{F\_j} definition
deliberately clips improvement.

If later learning improves an earlier task, the signed change should
remain available elsewhere in the evaluation record.

\section{Parameter movement is not
forgetting}\label{parameter-movement-is-not-forgetting}

Suppose a model moves from parameters

\texttt{theta\_old}

to

\texttt{theta\_new}.

The norm

\texttt{\textbar{}\textbar{}theta\_new-theta\_old\textbar{}\textbar{}}

does not tell us how much was forgotten.

A large movement may occur along directions that barely matter for an
earlier task.

A small movement may cross a sensitive decision boundary and cause a
large behavioral change.

Forgetting is therefore evaluated through retained performance, loss,
calibration, reward, or another declared behavioral quantity.

Parameter-space movement becomes meaningful only after it is connected
to that behavior.

This distinction is crucial for understanding consolidation methods.

They act in parameter space.

Their purpose is behavioral retention.

The two levels are related but not identical.

\section{Why the word catastrophic
exists}\label{why-the-word-catastrophic-exists}

McCloskey and Cohen studied what they called catastrophic interference
in connectionist networks under sequential learning (McCloskey and Cohen
1989).

The important phenomenon is not merely that an old score decreases.

It is that new learning can produce severe loss of earlier capability
when the same distributed parameters are reused.

Modern deep-learning literature often uses \textbf{catastrophic
forgetting} for the same broad family of failures.

The Atlas uses the word carefully.

A small noisy decline is not automatically catastrophic.

A report should state:

\begin{itemize}
\tightlist
\item
  what was learned first;
\item
  what was learned later;
\item
  how old performance changed;
\item
  over what training interval;
\item
  under what evaluation protocol.
\end{itemize}

Severity belongs in the evidence.

It should not be smuggled in by the adjective.

\section{Continual learning is not one benchmark
regime}\label{continual-learning-is-not-one-benchmark-regime}

A major source of confusion is that different continual-learning
experiments provide different information to the learner at evaluation.

Van de Ven, Tuytelaars, and Tolias distinguish task-incremental,
domain-incremental, and class-incremental scenarios (Ven et al. 2022).

\subsection{Task-incremental learning}\label{task-incremental-learning}

The system knows, or can unambiguously determine, which task is being
evaluated.

Task-specific heads, masks, or modules can therefore be selected.

\subsection{Domain-incremental
learning}\label{domain-incremental-learning}

The input distribution or context changes, but the output semantics
remain shared.

Explicit task identity need not be supplied at test time.

\subsection{Class-incremental
learning}\label{class-incremental-learning}

The output problem grows.

The model must discriminate across a global set of classes accumulated
over time.

These regimes differ materially.

A method that relies on a task selector has an easier inference contract
than one that must infer among all accumulated classes.

So a continual-learning result should always say what context
information is available.

\section{Memory enters continual learning in different
roles}\label{memory-enters-continual-learning-in-different-roles}

The Memory Taxonomy chapter established that memory should be described
by locus, write path, read path, lifetime, addressability, mutability,
provenance, and sharing.

That discipline prevents an immediate mistake.

\textbf{Replay is not a memory type.}

Replay is a use of remembered information during learning.

The remembered material might live in:

\begin{itemize}
\tightlist
\item
  a local buffer;
\item
  an external store;
\item
  a generated model;
\item
  a compressed exemplar set;
\item
  another durable or ephemeral locus.
\end{itemize}

The continual-learning mechanism is the reintroduction of earlier
evidence.

The memory architecture is a separate design decision.

\section{Replay: put old evidence back into the
objective}\label{replay-put-old-evidence-back-into-the-objective}

Let

\texttt{L\_cur(theta)}

be the current-context loss.

Let a memory \texttt{M} support an earlier-data loss

\texttt{L\_M(theta)}.

A simple replay objective is

\texttt{L\_replay(theta)\ =\ L\_cur(theta)\ +\ alpha\ L\_M(theta)}.

The logic is direct.

If later optimization sees only the new context, it may move into a
region that is excellent for the new objective and poor for the old one.

Replay makes earlier evidence visible again.

Rebuffi et al.'s iCaRL uses a bounded exemplar set in a
class-incremental learning system (Rebuffi et al. 2017).

Gradient Episodic Memory also retains episodic examples, although it
uses them differently (Lopez-Paz and Ranzato 2017).

Replay is powerful because it attacks one root cause of forgetting:

the optimizer otherwise cannot optimize against evidence it no longer
sees.

\section{Replay does not guarantee
retention}\label{replay-does-not-guarantee-retention}

Reintroducing old evidence is not the same as freezing old capability.

A replay buffer may be incomplete.

Its sampling distribution may be distorted.

Its weight in the loss may be too small.

Optimization may stop early.

The model may lack enough capacity to satisfy old and new objectives
simultaneously.

The remembered examples may not represent the full earlier distribution.

Therefore replay changes the optimization problem in a
retention-friendly direction.

It does not create an automatic theorem of zero forgetting.

\section{Consolidation: preserve important parameter
directions}\label{consolidation-preserve-important-parameter-directions}

A different strategy is to remember not the old examples themselves but
information about which parameters were important for the old task.

Elastic Weight Consolidation is a canonical example.

After learning task A, let

\texttt{theta\_A\^{}*}

be the selected parameter point.

When learning B, EWC uses

\texttt{L\_EWC(theta)\ =\ L\_B(theta)\ +\ (lambda/2)\ sum\_i\ F\_i\ (theta\_i-theta\_\{A,i\}\^{}*)\^{}2}.

Kirkpatrick et al.~motivate the weights \texttt{F\_i} through a diagonal
Fisher-information approximation to parameter importance (Kirkpatrick et
al. 2017).

The geometry is intuitive.

Moving a parameter with large \texttt{F\_i} is expensive.

Moving one with small \texttt{F\_i} is cheaper.

The old solution behaves like an anisotropic spring anchor.

\section{EWC is an approximation, not an
oracle}\label{ewc-is-an-approximation-not-an-oracle}

The EWC penalty is useful precisely because storing and replaying all
old data may be impossible or undesirable.

But compression has a price.

The old task is represented through:

\begin{itemize}
\tightlist
\item
  an anchor point;
\item
  a local quadratic approximation;
\item
  diagonal importance estimates.
\end{itemize}

That is far less information than the full old data distribution and
full old loss surface.

The original EWC paper explicitly frames the diagonal Fisher
construction as a tractable approximation (Kirkpatrick et al. 2017).

So the correct claim is not:

\begin{quote}
important parameters can never move.
\end{quote}

It is:

\begin{quote}
estimated important directions are penalized more strongly.
\end{quote}

That can reduce interference.

It does not guarantee exact behavioral preservation.

\section{The stability-plasticity
tension}\label{the-stability-plasticity-tension}

If every old parameter were frozen forever, old behavior could be
protected locally.

But new learning would eventually run out of freedom.

If every parameter remained fully plastic, adaptation would be easy but
interference could be severe.

Continual learning therefore repeatedly faces a stability-plasticity
tension.

How much should the system preserve?

How much should it change?

The answer depends on:

\begin{itemize}
\tightlist
\item
  capacity;
\item
  task compatibility;
\item
  available memory;
\item
  context information;
\item
  horizon;
\item
  evaluation priorities.
\end{itemize}

The exact witness makes this tension visible with no stochastic noise.

\section{Exact witness: two tasks, one
scalar}\label{exact-witness-two-tasks-one-scalar}

Use one parameter

\texttt{w}.

Task A has loss

\texttt{L\_A(w)=(1/2)(w+1)\^{}2}.

Task B has loss

\texttt{L\_B(w)=(1/2)(w-1)\^{}2}.

Task A wants

\texttt{w=-1}.

Task B wants

\texttt{w=1}.

There is no single scalar that makes both losses zero.

This is not an optimization failure.

It is a representational conflict.

The model has one degree of freedom and two incompatible optima.

\section{Unprotected sequential
learning}\label{unprotected-sequential-learning}

Train Task A first.

The optimum is

\texttt{w=-1}.

Then

\texttt{L\_A=0}.

Now train Task B to its own optimum.

The parameter moves to

\texttt{w=1}.

Then

\texttt{L\_B=0}

but

\texttt{L\_A=2}.

The old-task loss increased exactly by

\texttt{2}.

That is forgetting in the declared loss metric.

No randomness is involved.

No optimizer pathology is required.

Sequential interference follows from incompatible objectives sharing one
parameter.

\section{Exact quadratic
consolidation}\label{exact-quadratic-consolidation}

Set the scalar importance estimate to \(F=1\).

Use the retention objective

\[
J_\lambda(w) = L_B(w) + \frac{\lambda}{2}(w+1)^2.
\]

The exact minimizer is

\[
w_\lambda = \frac{1-\lambda}{1+\lambda}.
\]

At that point,

\[
L_A(w_\lambda) = \frac{2}{(1+\lambda)^2}
\]

and

\[
L_B(w_\lambda) = \frac{2\lambda^2}{(1+\lambda)^2}.
\]

For

\texttt{lambda=0},

we recover full adaptation to B:

\texttt{w=1}, \texttt{L\_A=2}, \texttt{L\_B=0}.

For

\texttt{lambda=1},

the compromise is

\texttt{w=0}, \texttt{L\_A=1/2}, \texttt{L\_B=1/2}.

For very large \texttt{lambda},

\texttt{w}

approaches the old optimum and Task B becomes poorly fit.

The penalty therefore does exactly what a consolidation penalty should
do in this toy:

it trades plasticity for retention.

\section{Replay and EWC coincide here for a special
reason}\label{replay-and-ewc-coincide-here-for-a-special-reason}

Now imagine exact replay with equal weight on both complete task
objectives:

\texttt{L\_joint=L\_A+L\_B}.

In this witness,

\texttt{L\_joint(w)=w\^{}2+1}.

The minimizer is again

\texttt{w=0}.

So equal replay and the \texttt{lambda=1,F=1} quadratic retention
penalty produce the same solution.

This is not a general equivalence.

It happens because Task A's loss is itself exactly the quadratic used as
the retention penalty.

In a real network:

\begin{itemize}
\tightlist
\item
  replay reevaluates remembered data under the current model;
\item
  EWC uses a local parameter-space approximation around an earlier
  solution.
\end{itemize}

The exact coincidence is useful precisely because it shows both the
connection and the boundary.

\section{Gradient constraints: protect old losses at update
time}\label{gradient-constraints-protect-old-losses-at-update-time}

Gradient Episodic Memory takes another route (Lopez-Paz and Ranzato
2017).

Let

\texttt{g}

be the current gradient.

Let

\texttt{g\_j}

be a gradient computed from episodic memory for old context \texttt{j}.

Suppose the update uses

\texttt{theta\textquotesingle{}\ =\ theta\ -\ eta\ g\_tilde}.

To first order,

\texttt{L\_j(theta\textquotesingle{})-L\_j(theta)\ approximately\ -eta\ g\_j\^{}T\ g\_tilde}.

If

\texttt{g\_j\^{}T\ g\_tilde\ \textgreater{}=\ 0},

the linearized old-context loss does not increase.

GEM therefore turns remembered examples into constraints on the update
direction.

This is neither ordinary replay nor EWC.

Replay changes the objective.

EWC adds a parameter penalty.

GEM changes the feasible gradient direction.

\section{First-order protection has a
scope}\label{first-order-protection-has-a-scope}

A gradient constraint is local.

It reasons about the immediate linearized change around the current
parameters.

Finite learning rates introduce higher-order effects.

A bounded episodic memory may not represent the full old distribution.

So even a constrained update must retain its claim boundary.

This is a recurring Atlas theme:

\begin{quote}
a local protective condition is not automatically a global retention
guarantee.
\end{quote}

The correct object and scale matter.

\section{Parameter isolation: stop sharing the contested
parameters}\label{parameter-isolation-stop-sharing-the-contested-parameters}

There is a more structural response to interference.

Do not let later tasks overwrite the same parameters.

PackNet uses iterative pruning and parameter allocation to add tasks
while preserving previously assigned weights (Mallya and Lazebnik 2018).

In the scalar witness, parameter isolation means replacing the one
shared variable with two task-selected variables:

\texttt{w\_A=-1}

and

\texttt{w\_B=1}.

Task A uses \texttt{w\_A}.

Task B uses \texttt{w\_B}.

Both losses are zero.

Interference disappears because the conflict has been removed from the
shared parameter.

\section{Isolation pays with capacity and
routing}\label{isolation-pays-with-capacity-and-routing}

The previous result sounds perfect until we count resources.

The original model had one scalar.

The isolated model has two.

It also needs a selector telling it which scalar to use.

This distinction matters in the continual-learning scenarios.

If task identity is supplied at test time, task-specific masks or heads
are natural.

If the model must classify across all accumulated classes without task
identity, routing is part of the problem.

Parameter isolation therefore trades one form of difficulty for another:

\begin{itemize}
\tightlist
\item
  less overwrite;
\item
  more capacity accounting;
\item
  more routing structure.
\end{itemize}

\section{Replay, consolidation, constraints, and isolation are
different
levers}\label{replay-consolidation-constraints-and-isolation-are-different-levers}

The four mechanism families act at different places.

\subsection{Replay}\label{replay}

Change what data or evidence is present during optimization.

\subsection{Consolidation}\label{consolidation}

Change the cost of moving through parameter space.

\subsection{Gradient constraints}\label{gradient-constraints}

Change which local update directions are allowed.

\subsection{Parameter isolation}\label{parameter-isolation}

Change which parameters can be written by later contexts.

These can be combined.

They can also fail for different reasons.

A replay method may fail because the memory is unrepresentative.

A consolidation method may fail because its importance approximation is
poor.

A gradient constraint may fail outside its local approximation or memory
support.

An isolation method may exhaust capacity.

Calling all four ``memory methods'' hides the mechanism.

\section{Capacity must be reported}\label{capacity-must-be-reported}

Continual learning is easy if we permit unlimited duplication.

For every new task, create a fresh model.

Old performance remains untouched.

But the system grows without bound and may require explicit task
routing.

That is why memory and capacity accounting are part of the scientific
claim.

Relevant resources can include:

\begin{itemize}
\tightlist
\item
  stored examples;
\item
  stored activations or transitions;
\item
  generated replay models;
\item
  Fisher or importance statistics;
\item
  task-specific heads;
\item
  masks;
\item
  reserved parameter blocks;
\item
  additional networks.
\end{itemize}

A method that retains better by spending more memory may still be the
right engineering choice.

The cost must simply remain visible.

\section{Average forgetting can hide a damaged
context}\label{average-forgetting-can-hide-a-damaged-context}

Suppose two earlier contexts have

\texttt{F\_1=0}

and

\texttt{F\_2=1}.

Then

\texttt{F\_avg=1/2}.

But

\texttt{F\_max=1}.

Another learner might have

\texttt{F\_1=1/2}

and

\texttt{F\_2=1/2}.

The average is identical.

The risk profile is not.

This is why the Atlas keeps both average and worst-context summaries.

A safety-critical application may care about the worst retained
capability.

A broad benchmark may care more about average retained performance.

Neither scalar subsumes the other.

\section{Positive transfer should not be clipped out of
existence}\label{positive-transfer-should-not-be-clipped-out-of-existence}

Forgetting metrics often focus on decline.

But later learning can improve earlier tasks.

If context B teaches a reusable representation, then

\texttt{R\_\{T,A\}\textgreater{}R\_\{A,A\}}

may occur.

A nonnegative forgetting metric reports zero in that case.

That is correct for the narrow question ``how much was lost?''

It is incomplete for the broader question ``what effect did later
learning have?''

So continual-learning reports should preserve signed transfer
information where it matters.

Lopez-Paz and Ranzato make transfer an explicit part of
continual-learning evaluation (Lopez-Paz and Ranzato 2017).

\section{Final accuracy alone is also
insufficient}\label{final-accuracy-alone-is-also-insufficient}

Suppose one learner forgets A badly while learning B, then relearns A
later.

Its final score may look good.

Another learner preserves A continuously.

If the application requires always-available capability, those histories
are different.

Similarly, a learner may retain old tasks perfectly but fail to learn
new ones.

That is not good continual learning merely because forgetting is zero.

The performance-through-time matrix makes both failures visible:

\begin{itemize}
\tightlist
\item
  instability of old capability;
\item
  failure of new acquisition.
\end{itemize}

\section{Forgetting and transfer can be
asymmetric}\label{forgetting-and-transfer-can-be-asymmetric}

Task A may help Task B.

Task B may harm Task A.

Task C may improve both.

Sequential order therefore matters.

A continual-learning benchmark is not fully described by its set of
tasks.

It also has an ordering, data schedule, context-transition structure,
and evaluation schedule.

Changing the order can change the interference pattern.

That makes continual learning a dynamical training process, not merely
multi-task learning with a shuffled dataset.

\section{Multi-task learning is the useful
counterfactual}\label{multi-task-learning-is-the-useful-counterfactual}

If all tasks and all data are jointly available from the start, we can
optimize a joint objective.

That is not the continual-learning information pattern.

It is still a valuable reference.

The gap between:

\begin{itemize}
\tightlist
\item
  joint training;
\item
  sequential training;
\item
  sequential training with replay;
\item
  sequential training with consolidation;
\end{itemize}

helps separate capacity conflict from information loss.

If joint training cannot fit all tasks, the model has a compatibility or
capacity problem.

If joint training succeeds but sequential training forgets, the schedule
and information constraints are implicated more directly.

\section{Replay versus joint
training}\label{replay-versus-joint-training}

Replay approximates some aspects of joint access.

But unless the replay store contains all old data with the same
weighting and sampling, it is not identical to joint training.

A bounded exemplar buffer changes the old-data distribution.

Generated replay adds modeling error.

Reservoir sampling, class balancing, and prioritization change which old
evidence remains visible.

So ``replay'' should always be accompanied by a memory contract:

what is stored, how much, how selected, and how sampled.

\section{EWC versus ordinary L2
regularization}\label{ewc-versus-ordinary-l2-regularization}

An ordinary quadratic penalty might constrain

\texttt{\textbar{}\textbar{}theta\textbar{}\textbar{}\^{}2}

or distance from a generic reference.

EWC instead anchors parameters to an earlier task solution with
coordinate-specific importance weights.

The source emphasizes this selective plasticity (Kirkpatrick et al.
2017).

The distinction matters.

Uniform resistance can protect unimportant coordinates too strongly and
important coordinates too weakly.

EWC attempts to allocate plasticity where the previous task is estimated
to tolerate it.

Again, the quality of that estimate is part of the method's empirical
burden.

\section{Consolidation and semantic memory are not
synonyms}\label{consolidation-and-semantic-memory-are-not-synonyms}

The Memory Taxonomy chapter uses \textbf{semantic} as a role describing
generalized knowledge rather than event-specific records.

EWC is called a consolidation method because it stabilizes parameters.

That does not mean the resulting parameters form a clean semantic store.

Likewise, episodic replay examples can contribute to generalized
semantic behavior during training.

Memory role and anti-forgetting mechanism answer different questions.

This distinction becomes important in the next chapter family.

\section{Catastrophic forgetting is partly an access
problem}\label{catastrophic-forgetting-is-partly-an-access-problem}

Why does replay help?

Because the learner regains access to evidence that the current stream
no longer supplies.

Why does external memory become relevant downstream?

Because parameters are a lossy, interference-prone place to preserve
some kinds of information.

But this chapter stops short of the stronger claim.

It does not conclude that knowledge should generally leave the weights.

It establishes the ingredients needed to ask that question:

\begin{itemize}
\tightlist
\item
  what was retained;
\item
  what was forgotten;
\item
  what evidence had to be stored;
\item
  what capacity was consumed;
\item
  what mechanism prevented interference.
\end{itemize}

The External-Memory chapter will make the architectural argument.

\section{Failure modes}\label{failure-modes-11}

\subsection{Forgetting-equals-parameter-drift}\label{forgetting-equals-parameter-drift}

Measure only parameter distance.

Failure: behavior is not evaluated.

\subsection{Zero-forgetting-by-zero-learning}\label{zero-forgetting-by-zero-learning}

Freeze everything.

Failure: old capability is retained because new capability is never
acquired.

\subsection{Replay-equals-memory-architecture}\label{replay-equals-memory-architecture}

Call every replay buffer an external persistent memory system.

Failure: training use is confused with storage locus and lifetime.

\subsection{EWC-equals-guarantee}\label{ewc-equals-guarantee}

Treat a Fisher-weighted local quadratic penalty as exact preservation.

Failure: approximation becomes certification.

\subsection{Isolation-equals-free
retention}\label{isolation-equals-free-retention}

Allocate fresh parameters indefinitely.

Failure: capacity and routing costs disappear from the comparison.

\subsection{Average-only reporting}\label{average-only-reporting}

Report one mean retention score.

Failure: severe failure on one context can be hidden.

\subsection{Regime mixing}\label{regime-mixing}

Compare task-incremental and class-incremental results as if task
identity were equally available.

Failure: the inference contracts differ.

\section{A practical continual-learning
ledger}\label{a-practical-continual-learning-ledger}

Before interpreting a continual-learning result, record:

{\def\LTcaptype{none} % do not increment counter
\begin{longtable}[]{@{}
  >{\raggedright\arraybackslash}p{(\linewidth - 2\tabcolsep) * \real{0.5000}}
  >{\raggedright\arraybackslash}p{(\linewidth - 2\tabcolsep) * \real{0.5000}}@{}}
\toprule\noalign{}
\begin{minipage}[b]{\linewidth}\raggedright
Field
\end{minipage} & \begin{minipage}[b]{\linewidth}\raggedright
Question
\end{minipage} \\
\midrule\noalign{}
\endhead
\bottomrule\noalign{}
\endlastfoot
stream & What changes over time? \\
context & Is task/context identity available during training and
evaluation? \\
write target & Which parameters or memory objects are updated? \\
old evidence & What earlier information remains accessible? \\
protection & Replay, consolidation, gradient constraint, isolation, or
combination? \\
capacity & What grows with the number of contexts? \\
evaluation & What is the full performance-through-time matrix? \\
forgetting & Average, worst-context, or another declared summary? \\
transfer & Are signed forward/backward effects retained? \\
baseline & Joint training, separate models, naive sequential training,
or another comparator? \\
\end{longtable}
}

This ledger prevents an anti-forgetting mechanism from being evaluated
without its information and capacity contract.

\section{What the exact witness
establishes}\label{what-the-exact-witness-establishes-4}

The scalar witness establishes a narrow but useful result.

Two tasks with incompatible optima share one parameter.

Sequentially optimizing the second task exactly raises the first task's
loss from

\texttt{0}

to

\texttt{2}.

A unit quadratic retention penalty changes the exact compromise to

\texttt{w=0}

with losses

\texttt{1/2}

and

\texttt{1/2}.

Equal-weight replay happens to produce the same point because the old
task loss is exactly the same quadratic as the chosen penalty.

Two isolated task-selected parameters achieve zero loss on both tasks by
spending extra capacity and routing.

Nothing stronger is implied.

The witness is not a benchmark result.

It is a controlled model of interference.

\section{Atlas handoff}\label{atlas-handoff}

\textbf{The External-Memory Thesis --- ATLAS-CH-EXTMEM-001.}

The downstream chapter may now assume:

\begin{itemize}
\tightlist
\item
  forgetting is behavioral, not parameter-distance alone;
\item
  replay restores earlier evidence to training;
\item
  consolidation protects selected parameter directions without replaying
  the full old distribution;
\item
  gradient constraints and parameter isolation are distinct mechanisms;
\item
  memory and capacity costs belong in the comparison;
\item
  task identity materially changes the difficulty of retention.
\end{itemize}

That chapter must independently answer the architectural question:

\begin{quote}
which knowledge should remain in parameters, and which knowledge should
move into persistent shared memory?
\end{quote}

The dependency direction matters.

Continual-learning difficulty motivates that question.

It does not answer it in advance.

\section{References used in this
chapter}\label{references-used-in-this-chapter-47}

\begin{itemize}
\tightlist
\item
  McCloskey and Cohen, \emph{Catastrophic Interference in Connectionist
  Networks: The Sequential Learning Problem} (McCloskey and Cohen 1989).
\item
  Kirkpatrick et al., \emph{Overcoming catastrophic forgetting in neural
  networks} (Kirkpatrick et al. 2017).
\item
  Rebuffi et al., \emph{iCaRL: Incremental Classifier and Representation
  Learning} (Rebuffi et al. 2017).
\item
  Lopez-Paz and Ranzato, \emph{Gradient Episodic Memory for Continual
  Learning} (Lopez-Paz and Ranzato 2017).
\item
  Mallya and Lazebnik, \emph{PackNet: Adding Multiple Tasks to a Single
  Network by Iterative Pruning} (Mallya and Lazebnik 2018).
\item
  van de Ven, Tuytelaars, and Tolias, \emph{Three types of incremental
  learning} (Ven et al. 2022).
\end{itemize}

Exact prerequisite identities and source-authority boundaries are
recorded in:

sources/source-locks/ATLAS-CH-CONTINUAL-001.yaml

\chapter{The External-Memory Thesis}\label{the-external-memory-thesis}

\textbf{Epistemic status:} audited retrieval + audited
continual-learning mechanisms + primary
retrieval-augmented/model-editing examples + Atlas synthesis + exact
finite witness.\\
\textbf{Specification:}
manuscript/specifications/ATLAS-CH-EXTMEM-001.md\\
\textbf{Formal packet:}
mathematics/derivations/ATLAS-CH-EXTMEM-001-DERIVATIONS.md\\
\textbf{Computational witness:}
mathematics/computational-witnesses/ATLAS-CW-EXTMEM-001.md\\
\textbf{Source lock:} sources/source-locks/ATLAS-CH-EXTMEM-001.yaml

A model can know something in at least two radically different ways.

It can carry the information implicitly in its parameters.

Or it can retrieve an explicit record when the information is needed.

Those two forms of knowing have different update semantics, provenance
semantics, sharing semantics, latency costs, and failure modes.

The governing thesis of this chapter is:

\begin{quote}
an intelligent system should not force every useful fact into its
weights when some information is better represented as explicit,
versioned, retrievable state.
\end{quote}

The thesis has an equally important converse:

\begin{quote}
external memory is not a universal substitute for learning.
\end{quote}

Stable structure, reusable procedures, compressed regularities, and
learned representations may belong in parameters.

The architecture is usually hybrid.

\section{The wrong question is parameters versus
memory}\label{the-wrong-question-is-parameters-versus-memory}

A simple binary question asks:

\begin{quote}
should knowledge live in parameters or external memory?
\end{quote}

That framing is too coarse.

The same system can use:

\begin{itemize}
\tightlist
\item
  parameters for syntax, features, procedures, and compressed
  regularities;
\item
  an external store for mutable facts, records, evidence, and shared
  organizational state;
\item
  a working context for the subset needed now.
\end{itemize}

Placement is an architectural decision about which representation best
serves which kind of information.

\section{What is parametric
knowledge?}\label{what-is-parametric-knowledge}

Call knowledge \textbf{parametric} when its use is mediated by learned
model state.

For parameters \texttt{theta}, behavior is:

\texttt{f\_theta(q)}.

The knowledge may not correspond to a single readable parameter.

It can be distributed across many coordinates and interactions.

That distribution is often the point.

It enables compression and generalization.

\section{What is external memory?}\label{what-is-external-memory}

External memory means explicit state outside the model parameters.

A record can carry:

\begin{itemize}
\tightlist
\item
  key;
\item
  value/content;
\item
  source;
\item
  timestamp;
\item
  version;
\item
  type;
\item
  access policy;
\item
  status;
\item
  supersession relation.
\end{itemize}

The important property is not merely that bytes are stored elsewhere.

It is that the system can address and update an explicit memory object
under a declared read/write contract.

\section{Retrieval is the read path, not the memory
itself}\label{retrieval-is-the-read-path-not-the-memory-itself}

RETRIEVAL-001 established:

\begin{quote}
retrieval is a declared read operation over records, not a synonym for
relevance, truth, provenance, or usefulness.
\end{quote}

That distinction remains active here.

A perfectly maintained memory can still fail if retrieval:

\begin{itemize}
\tightlist
\item
  misses the record;
\item
  applies the wrong filter;
\item
  ranks badly;
\item
  uses a stale index;
\item
  returns an unauthorized version.
\end{itemize}

External storage and successful memory use are separate layers.

\section{Continual learning is the parameter-update
side}\label{continual-learning-is-the-parameter-update-side}

CONTINUAL-001 established that sequential parameter learning can alter
earlier behavior.

It also separated:

\begin{itemize}
\tightlist
\item
  replay;
\item
  consolidation/EWC;
\item
  episodic gradient constraints;
\item
  parameter isolation.
\end{itemize}

Those mechanisms matter because moving knowledge into parameters creates
update obligations.

The external-memory thesis asks when some of those obligations can be
avoided by leaving selected knowledge explicit.

\section{The placement descriptor}\label{the-placement-descriptor}

For knowledge item or class \(k\), use:

\[
\operatorname{Place}(k)=(V,P,S,D,R,L,H,A,G).
\]

The coordinates are:

\begin{itemize}
\tightlist
\item
  volatility \(V\);
\item
  provenance/audit need \(P\);
\item
  sharing scope \(S\);
\item
  deletion/supersession need \(D\);
\item
  retrievability/addressability \(R\);
\item
  latency constraint \(L\);
\item
  availability/failure-tolerance requirement \(H\);
\item
  access-control/privacy requirement \(A\);
\item
  value of parametric generalization/compression \(G\).
\end{itemize}

The chapter defines no universal scalar score over these dimensions.
External locus and persistence also remain distinct: this thesis focuses
on persistent external records when that lifetime is deliberately
chosen; it does not redefine every external memory object as persistent.

\section{Why no scalar placement
score?}\label{why-no-scalar-placement-score}

Suppose one fact is highly volatile and provenance-critical.

That points toward external memory.

But suppose the same fact is required on every token with a hard
microsecond latency budget.

That points toward parameterization or aggressive caching.

Neither fact cancels the other.

A placement decision must expose the tradeoff rather than hide it in an
unexplained number.

\section{Volatility favors explicit
state}\label{volatility-favors-explicit-state}

Some information changes often:

\begin{itemize}
\tightlist
\item
  current prices;
\item
  active policies;
\item
  account permissions;
\item
  schedules;
\item
  inventory;
\item
  organizational decisions;
\item
  scientific records under revision.
\end{itemize}

Repeatedly changing model weights for every update may be operationally
awkward or behaviorally risky.

A versioned external record has a direct update operation.

That makes volatility one pressure toward externalization.

\section{Provenance favors explicit
state}\label{provenance-favors-explicit-state}

Some claims must answer:

\begin{itemize}
\tightlist
\item
  where did this come from?
\item
  which version?
\item
  who authorized it?
\item
  what replaced it?
\item
  when was it observed?
\end{itemize}

A record can carry those fields directly.

A bare parameter vector does not, by itself, expose a per-fact source
record.

This is a representation distinction, not a claim that parametric
systems can never maintain provenance.

They can maintain auxiliary provenance externally.

\section{Deletion and supersession favor explicit
records}\label{deletion-and-supersession-favor-explicit-records}

An explicit record can be marked:

\begin{itemize}
\tightlist
\item
  current;
\item
  superseded;
\item
  revoked;
\item
  quarantined;
\item
  deleted.
\end{itemize}

That does not make erasure trivial in replicated infrastructure.

But it gives the semantic operation a named target.

By contrast, removing one learned association from model behavior is a
model-editing problem.

The two deletion semantics are not equivalent.

\section{Sharing can favor external
memory}\label{sharing-can-favor-external-memory}

Suppose ten agents need the same mutable operational fact.

One authoritative record can, in principle, be updated once.

Every agent can then retrieve the same current version.

But this advantage exists only if synchronization, access control, and
freshness are well-defined.

A shared store with divergent caches can create multiple effective
memories.

\section{Addressability matters}\label{addressability-matters}

Some knowledge naturally has a key:

\begin{itemize}
\tightlist
\item
  customer ID;
\item
  document ID;
\item
  theorem ID;
\item
  experiment run;
\item
  date;
\item
  entity-relation pair;
\item
  configuration name.
\end{itemize}

Such information is easy to imagine as an explicit memory record.

Other knowledge is diffuse.

There may be no sensible lookup key for:

\begin{itemize}
\tightlist
\item
  fluency;
\item
  visual invariance;
\item
  a learned optimization heuristic;
\item
  a broad concept manifold.
\end{itemize}

Addressability is therefore a real placement coordinate.

\section{Parametric compression has
value}\label{parametric-compression-has-value}

Parameters can compress many examples into reusable behavior.

The system does not need to retrieve every training example whenever it
applies a learned rule.

This is a major advantage.

A database of examples is not the same thing as learning a
representation or procedure from them.

External memory should not become an excuse to avoid learning stable
structure.

\section{Low-latency ubiquitous structure can favor
parameters}\label{low-latency-ubiquitous-structure-can-favor-parameters}

Some knowledge is required constantly.

An explicit retrieval on every use may add:

\begin{itemize}
\tightlist
\item
  index lookup;
\item
  memory traffic;
\item
  network latency;
\item
  synchronization;
\item
  failure modes.
\end{itemize}

Parametric behavior is already present in the model's forward
computation.

This can favor parameterization for stable, ubiquitous structure.

No universal latency ordering is claimed.

\section{RAG makes the hybrid idea
concrete}\label{rag-makes-the-hybrid-idea-concrete}

Lewis et al.~combine a parametric sequence model with a dense
non-parametric Wikipedia index (Lewis et al. 2020).

Their formulation is useful here because it makes the two loci explicit:

\begin{itemize}
\tightlist
\item
  a learned parametric model;
\item
  a retrievable external corpus.
\end{itemize}

The source also explicitly motivates retrieval in part through
provenance and world-knowledge update limitations of parameter-only
systems.

The Atlas uses this as one primary example, not as proof that RAG is the
universal architecture.

\section{kNN-LM shows datastore
substitution}\label{knn-lm-shows-datastore-substitution}

Khandelwal et al.~augment a pretrained language model with a
nearest-neighbor datastore (Khandelwal et al. 2020).

In their reported experiments, changing the datastore supports domain
adaptation without further model training.

The bounded architectural lesson is:

\begin{quote}
some behavior can be changed by changing explicit memory while freezing
the model.
\end{quote}

That is exactly the update-separation property the external-memory
thesis cares about.

\section{RETRO shows scale, not universal
superiority}\label{retro-shows-scale-not-universal-superiority}

Borgeaud et al.~condition language modeling on chunks retrieved from a
very large text database (Borgeaud et al. 2022).

Their result demonstrates that explicit retrieved memory can complement
model parameters at large scale.

It does not prove that every knowledge item belongs in retrieval.

Scale is evidence that the architecture is viable, not a universal
placement rule.

\section{ROME gives the necessary
counterexample}\label{rome-gives-the-necessary-counterexample}

Meng et al.~show that specific factual associations can be edited
directly in model parameters in their reported setting (Meng et al.
2022).

That matters because it blocks an overstrong thesis.

It would be false to say:

\begin{quote}
mutable facts must always be external.
\end{quote}

Parametric editing is possible.

The real question is which update semantics and system guarantees are
desirable.

\section{The exact parametric
witness}\label{the-exact-parametric-witness}

Use two queries:

\(A\) and \(B\).

Let:

\[
f_\theta(A)=\theta_1+\theta_2
\]

and:

\[
f_\theta(B)=\theta_1-\theta_2.
\]

Initial target facts are:

\[
f_\theta(A)=2, \qquad f_\theta(B)=0.
\]

The unique parameter vector is:

\[
\theta=(1,1).
\]

\section{Updating one fact requires a coordinated edit
here}\label{updating-one-fact-requires-a-coordinated-edit-here}

Now change only the target for A:

\[
A: 2 \to 4
\]

while preserving:

\[
f_\theta(B)=0.
\]

The new unique parameter vector is:

\[
\theta'=(2,2).
\]

Therefore:

\[
\Delta \theta=(1,1).
\]

Both coordinates change.

This is a property of the chosen representation.

It is not a theorem that factual parameter edits always require many
changed weights.

\section{A naive coordinate-local edit
interferes}\label{a-naive-coordinate-local-edit-interferes}

Change only:

\texttt{theta\_1:1\ -\textgreater{}\ 2}

and leave:

\texttt{theta\_2=1}.

Then:

\texttt{f(A)=3}

and:

\texttt{f(B)=1}.

The edit fails twice:

\begin{itemize}
\tightlist
\item
  A does not reach 4;
\item
  B no longer remains 0.
\end{itemize}

This is the smallest exact example of representation-induced edit
coupling.

\section{The external-store
version}\label{the-external-store-version}

Represent the same two facts as records.

Initial state:

\texttt{A\ -\textgreater{}\ (2,s\_A1,v1)}

\texttt{B\ -\textgreater{}\ (0,s\_B1,v1)}.

Update A by making the version transition explicit:

\texttt{A,v1\ -\textgreater{}\ (2,s\_A1,v1,status=superseded)}

\texttt{A,v2\ -\textgreater{}\ (4,s\_A2,v2,status=current,supersedes=v1)}.

B is unchanged and current.

A latest-version exact-key read now returns:

\texttt{A=4}

\texttt{B=0}.

Only logical key A changed.

\section{Record-locality is not compute
optimality}\label{record-locality-is-not-compute-optimality}

The external witness changes one logical key.

That is a statement about state mutation locality.

It does not prove that the external update uses fewer machine operations
than a model edit.

The database may replicate records, rebuild indexes, invalidate caches,
or synchronize across regions.

Representation-locality and systems cost are different claims.

\section{Provenance is explicit in the
store}\label{provenance-is-explicit-in-the-store}

The external records retain:

\texttt{s\_A1}

for version 1 and:

\texttt{s\_A2}

for version 2.

The supersession edge preserves the update history.

The bare parameter vector:

\texttt{theta=(1,1)}

does not itself expose which source supported fact A.

Again, the correct conclusion is not:

\begin{quote}
parameters cannot have provenance.
\end{quote}

The correct conclusion is:

\begin{quote}
provenance is not encoded as an explicit per-fact record in the bare
parameter representation.
\end{quote}

\section{The stale-read
counterexample}\label{the-stale-read-counterexample}

External memory creates a new failure mode.

Suppose one consumer still reads snapshot \texttt{M\_1} after the
authoritative store advances to \texttt{M\_2}.

The authoritative latest value is:

\texttt{A=4}.

The stale consumer still observes:

\texttt{A=2}.

Externalization has made updating the authoritative record easy.

It has not made every read fresh.

\section{Freshness needs a
contract}\label{freshness-needs-a-contract}

Let:

\texttt{v\_star(k)}

be the authoritative current version for key \texttt{k}.

Let:

\texttt{v\_C(k)}

be the version observed by consumer \texttt{C}.

A strict freshness rule might require:

\texttt{v\_C(k)=v\_star(k)}.

A bounded-staleness policy might permit some lag.

The point is architectural:

\begin{quote}
freshness must be specified.
\end{quote}

It is not implied by the word \textbf{memory}.

\section{Caching creates local memory
again}\label{caching-creates-local-memory-again}

A system may externalize knowledge into a shared store and then cache
retrieved results near each model instance.

That is often sensible for latency.

It also means the system now has several memory loci:

\begin{itemize}
\tightlist
\item
  authoritative store;
\item
  retrieval index;
\item
  cache;
\item
  working context;
\item
  parameters.
\end{itemize}

The memory taxonomy has returned.

Externalization does not make the architecture simple.

It makes some state explicit.

\section{Write correctness and read correctness
differ}\label{write-correctness-and-read-correctness-differ}

A memory update can be correct while a later answer is wrong.

Possible causes include:

\begin{itemize}
\tightlist
\item
  stale cache;
\item
  stale index;
\item
  wrong query;
\item
  wrong filter;
\item
  wrong rank;
\item
  context omission;
\item
  downstream reasoning error.
\end{itemize}

This separation matters operationally.

It lets the system ask:

\begin{quote}
did we store the right fact?
\end{quote}

before asking:

\begin{quote}
did the model use it correctly?
\end{quote}

\section{Retrieval failure is not memory-write
failure}\label{retrieval-failure-is-not-memory-write-failure}

Suppose version 2 of A is present and current.

A vector retriever fails to return it.

That is a read-path failure.

Rewriting the memory record will not necessarily fix it.

The correct repair may be in:

\begin{itemize}
\tightlist
\item
  indexing;
\item
  query formulation;
\item
  ranking;
\item
  filters;
\item
  context compilation.
\end{itemize}

Typed failure boundaries prevent useless updates.

\section{Poisoning is a memory-layer
risk}\label{poisoning-is-a-memory-layer-risk}

Explicit memory can be modified.

That is an advantage for legitimate updates.

It is also an attack surface.

A malicious or low-quality record can be:

\begin{itemize}
\tightlist
\item
  inserted;
\item
  ranked highly;
\item
  given false metadata;
\item
  propagated to many consumers.
\end{itemize}

External memory therefore needs evidence and authority controls, not
only retrieval quality.

\section{Access control becomes
first-class}\label{access-control-becomes-first-class}

A shared store may contain facts that not every agent may read.

Let:

\texttt{allow(agent,key,operation)}

be an authorization predicate.

A successful lookup should not bypass this predicate merely because the
record is relevant.

This creates another separation:

\begin{quote}
retrievable does not mean authorized.
\end{quote}

The downstream Polity chapter will need this distinction.

\section{Sharing is not global
visibility}\label{sharing-is-not-global-visibility}

A memory can be shared among:

\begin{itemize}
\tightlist
\item
  one team;
\item
  one model family;
\item
  one tenant;
\item
  one organization;
\item
  one scientific programme.
\end{itemize}

Shared does not mean public.

Scope is part of the memory contract.

This is especially important when memory contains private, licensed,
embargoed, or role-restricted information.

\section{Deletion is not trivial even
externally}\label{deletion-is-not-trivial-even-externally}

An explicit record can be marked deleted or revoked.

But copies may remain in:

\begin{itemize}
\tightlist
\item
  caches;
\item
  replicas;
\item
  logs;
\item
  backups;
\item
  derived indexes;
\item
  compiled contexts.
\end{itemize}

Therefore external memory gives deletion a clear target but not
necessarily immediate total erasure.

A serious system must define deletion propagation.

\section{Parametric erasure is a different
problem}\label{parametric-erasure-is-a-different-problem}

If information has been learned into parameters, deleting the original
training record does not imply the behavior disappears.

The two operations target different objects.

This distinction is especially important for systems with both:

\begin{itemize}
\tightlist
\item
  learned parametric state;
\item
  retained external records.
\end{itemize}

Deletion policy should say which loci are in scope.

\section{Consolidation moves in the other
direction}\label{consolidation-moves-in-the-other-direction}

The external-memory thesis is not a one-way migration out of parameters.

Sometimes repeated external evidence becomes stable enough that learning
it into parameters is useful.

Call this \textbf{consolidation}.

Examples might include:

\begin{itemize}
\tightlist
\item
  repeated tool-use patterns;
\item
  stable schema structure;
\item
  durable domain regularities;
\item
  procedural skills.
\end{itemize}

Consolidation is a parameter-learning operation.

It is distinct from memory write.

\section{Externalization can precede
consolidation}\label{externalization-can-precede-consolidation}

A system can first store new evidence externally because it is:

\begin{itemize}
\tightlist
\item
  recent;
\item
  uncertain;
\item
  volatile;
\item
  provenance-sensitive.
\end{itemize}

Later, after enough evidence accumulates, some stable regularity can be
learned parametrically.

This creates a useful architecture:

\begin{quote}
explicit first, consolidate later.
\end{quote}

The chapter presents this as a design pattern, not a theorem of optimal
learning.

\section{Parameter isolation and external memory solve different
problems}\label{parameter-isolation-and-external-memory-solve-different-problems}

CONTINUAL-001 showed that parameter isolation can avoid direct overwrite
by adding capacity or routing.

External memory avoids some parameter updates by moving selected state
elsewhere.

Those are not the same strategy.

Parameter isolation still places knowledge in learned capacity.

Externalization places it in explicit records.

The costs differ.

\section{Replay uses external evidence without making it the final
answer
source}\label{replay-uses-external-evidence-without-making-it-the-final-answer-source}

Replay stores examples or exemplars externally.

During learning, those records return to the optimizer.

This is different from runtime retrieval of a fact for direct use.

So even within one system, external records can have multiple roles:

\begin{itemize}
\tightlist
\item
  training evidence;
\item
  episodic replay;
\item
  semantic knowledge;
\item
  runtime retrieval;
\item
  audit/provenance archive.
\end{itemize}

Role should not be inferred from locus alone.

\section{A record is not automatically
knowledge}\label{a-record-is-not-automatically-knowledge}

An external store can contain:

\begin{itemize}
\tightlist
\item
  false statements;
\item
  conflicting statements;
\item
  stale statements;
\item
  unsupported statements;
\item
  low-authority sources.
\end{itemize}

The system still needs:

\begin{itemize}
\tightlist
\item
  provenance;
\item
  ranking;
\item
  authority;
\item
  reconciliation;
\item
  uncertainty handling.
\end{itemize}

External memory makes records inspectable.

It does not make them true.

\section{Multiple records can legitimately
conflict}\label{multiple-records-can-legitimately-conflict}

Suppose two sources disagree.

Parameter-only consolidation may blur the conflict into one behavior.

An external store can preserve both records explicitly.

That can be useful when the disagreement itself matters.

But then the reader must decide:

\begin{itemize}
\tightlist
\item
  which source is authoritative;
\item
  whether both are shown;
\item
  whether uncertainty is retained.
\end{itemize}

Preserving conflict is not the same as resolving it.

\section{Update frequency is only one
axis}\label{update-frequency-is-only-one-axis}

A highly volatile fact often favors external memory.

But a rarely changing fact can also be external if:

\begin{itemize}
\tightlist
\item
  provenance is critical;
\item
  deletion must be possible;
\item
  sharing is required;
\item
  access policy is complex.
\end{itemize}

Conversely, some moderately changing information may still be
parameterized if retrieval is too costly.

Placement is multi-dimensional.

\section{The ``world in weights'' failure
mode}\label{the-world-in-weights-failure-mode}

A system that tries to carry every changing fact in parameters inherits
several burdens:

\begin{itemize}
\tightlist
\item
  weight updates for factual changes;
\item
  unclear per-fact provenance;
\item
  difficult selective deletion;
\item
  duplicated updates across model copies;
\item
  interference risk;
\item
  stale deployed checkpoints.
\end{itemize}

This is the core motivation for the external-memory thesis.

But the remedy must not become a slogan.

\section{The ``database as intelligence'' failure
mode}\label{the-database-as-intelligence-failure-mode}

The opposite mistake is to externalize everything.

A datastore does not replace:

\begin{itemize}
\tightlist
\item
  abstraction;
\item
  reasoning;
\item
  representation learning;
\item
  compression;
\item
  procedure learning;
\item
  transfer.
\end{itemize}

If every answer requires finding a memorized record, the system may fail
whenever exact evidence is absent.

External memory should complement learned capability.

\section{The hybrid architecture}\label{the-hybrid-architecture}

A mature system may use four layers:

\begin{enumerate}
\def\labelenumi{\arabic{enumi}.}
\tightlist
\item
  \textbf{parameters} for compressed reusable structure;
\item
  \textbf{persistent external memory} for explicit versioned records;
\item
  \textbf{retrieval} for selecting candidate evidence;
\item
  \textbf{working context} for the task-specific subset currently
  active.
\end{enumerate}

The next chapter, Context Compilation, will develop layer 4.

This chapter stops before that step.

\section{A practical placement
ledger}\label{a-practical-placement-ledger}

Before moving a knowledge class into or out of parameters, record:

{\def\LTcaptype{none} % do not increment counter
\begin{longtable}[]{@{}ll@{}}
\toprule\noalign{}
Field & Question \\
\midrule\noalign{}
\endhead
\bottomrule\noalign{}
\endlastfoot
item/class & What knowledge is being placed? \\
volatility & How often does it change? \\
provenance & Must source/version be visible per item? \\
sharing & Which agents/models need one common state? \\
supersession & Must updates, revocation, or deletion be explicit? \\
retrievability & Is there a meaningful key/query/index? \\
latency & Can retrieval cost be tolerated? \\
availability & What happens if the store is unavailable? \\
access control & Who may read/write it? \\
privacy & Does explicit storage create new exposure? \\
generalization & Is compression into reusable behavior valuable? \\
consolidation & Should stable evidence later move into parameters? \\
freshness & Which version guarantees are required? \\
fallback & What should happen on retrieval failure? \\
\end{longtable}
}

This ledger turns ``put it in memory'' into an auditable design
decision.

\section{Failure modes}\label{failure-modes-12}

\subsection{External-equals-fresh}\label{external-equals-fresh}

The record is updated but consumers use stale snapshots.

\subsection{External-equals-correct}\label{external-equals-correct}

A stored record is treated as true merely because it is explicit.

\subsection{Retrieved-equals-authorized}\label{retrieved-equals-authorized}

A relevant record bypasses access policy.

\subsection{Parameter-change-equals-forgetting}\label{parameter-change-equals-forgetting}

Any weight movement is called forgetting without behavioral evaluation.

\subsection{External-equals-cheap}\label{external-equals-cheap}

Record-local updates are assumed to imply lower wall-clock cost.

\subsection{Parametric-equals-uneditable}\label{parametric-equals-uneditable}

Direct model editing is ignored despite evidence that some factual
associations can be changed parametrically.

\subsection{Shared-equals-consistent}\label{shared-equals-consistent}

Multiple agents point to one store, so consistency is assumed without
version/synchronization semantics.

\subsection{Externalize-everything}\label{externalize-everything}

Explicit memory is treated as a replacement for learning compressed
structure.

\section{What the exact witness
establishes}\label{what-the-exact-witness-establishes-5}

The companion witness proves only that, for the chosen two-fact
representations:

\begin{itemize}
\tightlist
\item
  initial parameter state is \texttt{(1,1)};
\item
  updated parameter state is \texttt{(2,2)};
\item
  both parameter coordinates change under the exact preserving update;
\item
  the naive one-coordinate edit produces \texttt{A=3,B=1};
\item
  the external latest-version update changes only logical key A;
\item
  external provenance retains old/new source/version information;
\item
  a stale snapshot still returns A=2 after authoritative A becomes 4.
\end{itemize}

It does not prove universal systems superiority.

\section{Downstream handoff: Context
Compilation}\label{downstream-handoff-context-compilation}

\textbf{ATLAS-CH-CONTEXTCOMP-001} may now assume:

\begin{itemize}
\tightlist
\item
  hybrid parametric/external placement;
\item
  versioned external records;
\item
  freshness as an explicit policy;
\item
  retrieval success as distinct from storage correctness;
\item
  provenance-preserving record semantics.
\end{itemize}

It must independently define how retrieved records become the bounded
working context of a model invocation.

\section{Downstream handoff:
Polity}\label{downstream-handoff-polity}

\textbf{ATLAS-CH-POLITY-001} may now assume:

\begin{itemize}
\tightlist
\item
  shared memory is an explicit architectural locus;
\item
  shared visibility requires synchronization/version semantics;
\item
  access control is distinct from relevance;
\item
  common memory can coexist with private parametric and working state.
\end{itemize}

It must independently define multi-agent governance, coordination,
authority, and shared-state policy.

\section{References used in this
chapter}\label{references-used-in-this-chapter-48}

\begin{itemize}
\tightlist
\item
  Lewis et al., \emph{Retrieval-Augmented Generation for
  Knowledge-Intensive NLP Tasks} (Lewis et al. 2020).
\item
  Khandelwal et al., \emph{Generalization through Memorization: Nearest
  Neighbor Language Models} (Khandelwal et al. 2020).
\item
  Borgeaud et al., \emph{Improving language models by retrieving from
  trillions of tokens} (Borgeaud et al. 2022).
\item
  Meng et al., \emph{Locating and Editing Factual Associations in GPT}
  (Meng et al. 2022).
\end{itemize}

Exact source identities and authority boundaries are recorded in:

sources/source-locks/ATLAS-CH-EXTMEM-001.yaml

\chapter{Context Compilation}\label{context-compilation}

\textbf{Epistemic status:} audited External-Memory Thesis + Atlas-owned
finite context-compilation formalism and exact witness.\\
\textbf{Specification:}
manuscript/specifications/ATLAS-CH-CONTEXTCOMP-001.md\\
\textbf{Derivation packet:}
mathematics/derivations/ATLAS-CH-CONTEXTCOMP-001-DERIVATIONS.md\\
\textbf{Computational witness:}
mathematics/computational-witnesses/ATLAS-CW-CONTEXTCOMP-001.md\\
\textbf{Source lock:} sources/source-locks/ATLAS-CH-CONTEXTCOMP-001.yaml

A memory system can contain the right information and still hand the
model the wrong working context.

The record can exist.

Retrieval can find it.

The ranker can score it highly.

The final prompt can still be stale, unauthorized, over budget, missing
a mandatory record, stripped of provenance, or badly assembled.

This chapter treats that assembly step as a first-class transformation.

The governing claim is:

\[
\boxed{
\text{a model context is a compiled working set, not merely a prefix of retrieved text.}
}
\]

The word compiled matters because context construction has inputs,
constraints, policy, failure modes, and a replayable output.

\section{The external-memory
handoff}\label{the-external-memory-handoff}

The External-Memory Thesis established a hybrid architecture:

\begin{itemize}
\tightlist
\item
  some structure remains parametric;
\item
  some information remains explicit and retrievable;
\item
  external records can carry source, version, status, and supersession
  metadata;
\item
  freshness is a policy question;
\item
  correct storage is distinct from successful retrieval and use.
\end{itemize}

CONTEXTCOMP starts after retrieval has already produced candidates.

Its question is narrower:

\begin{quote}
which candidate records should actually become the bounded working
context of this invocation, and in what serialized form?
\end{quote}

\section{The pipeline}\label{the-pipeline}

Write:

\[
M
\xrightarrow{\mathrm{retrieve}}
R_q
\xrightarrow{\mathrm{compile}_{\Pi,B}}
K_q
\xrightarrow{\mathrm{invoke}}
y.
\]

Here:

\begin{itemize}
\tightlist
\item
  \(M\) is external memory;
\item
  \(R_q\) is the retrieved candidate set for query/invocation \(q\);
\item
  \(\Pi\) is compilation policy;
\item
  \(B\) is context budget;
\item
  \(K_q\) is the compiled ordered working context;
\item
  \(y\) is downstream model behavior.
\end{itemize}

Every arrow has different obligations.

\section{Storage is not retrieval}\label{storage-is-not-retrieval}

A correct current record can be present in memory and never retrieved.

That is a retrieval failure.

It is not a compilation failure.

The compiler cannot select a candidate it never receives.

This distinction is inherited from EXTMEM-001 and its Retrieval
prerequisite.

\section{Retrieval is not
compilation}\label{retrieval-is-not-compilation}

A retriever may return:

\begin{itemize}
\tightlist
\item
  stale versions;
\item
  duplicate keys;
\item
  unauthorized records;
\item
  many more records than the model budget permits;
\item
  records whose metadata makes them mandatory;
\item
  records useful for ranking but invalid under the current invocation
  policy.
\end{itemize}

The candidate set is therefore an input to compilation.

It is not yet the context.

\section{Ranking is one signal, not the whole
contract}\label{ranking-is-one-signal-not-the-whole-contract}

A retrieval rank can express one notion of relevance or similarity.

It does not automatically encode:

\begin{itemize}
\tightlist
\item
  current version;
\item
  authorization;
\item
  mandatory policy;
\item
  provenance requirements;
\item
  atomic serialization cost;
\item
  conflict state.
\end{itemize}

This chapter's exact witness makes that distinction literal.

The top-ranked policy record is stale.

The current mandatory policy ranks third.

A rank-only prefix fails.

\section{Candidate records}\label{candidate-records}

Represent a candidate record as

\[
r_i=
(k_i,v_i,\sigma_i,p_i,a_i,m_i,\rho_i,c_i,u_i,t_i).
\]

The fields are:

\begin{itemize}
\tightlist
\item
  logical key;
\item
  version;
\item
  status;
\item
  provenance/source;
\item
  authorization;
\item
  mandatory flag;
\item
  retrieval rank;
\item
  compiled cost;
\item
  declared compiler utility;
\item
  type.
\end{itemize}

This is a chapter-local formal object.

It is not proposed as a universal context format.

\section{Compiled cost includes
metadata}\label{compiled-cost-includes-metadata}

Suppose the payload itself costs

\[
d_i
\]

units.

Suppose required metadata costs

\[
h_i.
\]

Then the actual compiled cost is

\[
c_i=d_i+h_i.
\]

If the system budgets only the payload while still requiring provenance
in the prompt, its budget calculation is wrong.

The resource accounting must describe the object actually serialized.

\section{Provenance is not
decoration}\label{provenance-is-not-decoration}

The External-Memory Thesis made source/version identity explicit.

Context compilation should not throw that information away casually.

For the finite formalism, serialize each included record with an
envelope:

\[
E(r_i)
=
[\mathrm{id},\mathrm{key},\mathrm{version},\mathrm{source},\mathrm{status},\mathrm{payload}].
\]

The envelope lets the working context say not only what was included,
but which version and source were included.

That does not prove the source is true.

It preserves identity.

\section{Authorization comes before
optimization}\label{authorization-comes-before-optimization}

Suppose a retrieved record is highly ranked and appears highly useful.

If the invocation is not authorized to use it, the compiler must exclude
it.

Authorization is therefore modeled as a hard admissibility filter:

\[
a_i=1.
\]

It is not a utility bonus.

The compiler cannot trade privacy or authority against a few extra
utility points.

\section{Freshness comes before optional
utility}\label{freshness-comes-before-optional-utility}

Suppose two records share logical key policy:

\begin{itemize}
\tightlist
\item
  version 1 is superseded;
\item
  version 2 is current.
\end{itemize}

If the declared policy is latest-current, version 1 is removed before
optional selection.

This remains true even if version 1 has:

\begin{itemize}
\tightlist
\item
  better retrieval rank;
\item
  higher declared utility;
\item
  lower compiled cost.
\end{itemize}

A validity constraint is not a soft preference.

\section{Unresolved conflict is a failure
state}\label{unresolved-conflict-is-a-failure-state}

Now suppose two authorized records for the same logical key are both
marked current.

If the version policy does not say which wins, the compiler should not
silently guess.

It should return:

\[
\mathrm{VERSION\_CONFLICT}.
\]

Choosing the higher retrieval rank would be a new policy.

If that policy is desired, it must be declared.

\section{Mandatory records}\label{mandatory-records}

Some records are not optional under a particular invocation contract.

Examples can include:

\begin{itemize}
\tightlist
\item
  current governing policy;
\item
  task instruction;
\item
  safety constraint;
\item
  required schema;
\item
  experiment identity.
\end{itemize}

Let

\[
M
\]

be the mandatory set.

Its total cost is

\[
C_M=
\sum_{r_i\in M}c_i.
\]

This cost is paid before optional context is selected.

\section{Mandatory overflow}\label{mandatory-overflow}

If

\[
C_M>B,
\]

no atomic context can both:

\begin{itemize}
\tightlist
\item
  contain every mandatory record;
\item
  fit the budget.
\end{itemize}

The correct output is explicit failure:

\[
\boxed{
\mathrm{MANDATORY\_OVERFLOW}.
}
\]

The compiler should not make a mandatory rule half-visible merely to
keep the invocation running.

\section{Optional selection}\label{optional-selection}

After authorization, freshness, conflict resolution, and mandatory
inclusion, let \(O\) be the optional records.

Residual budget is:

\[
B'=B-C_M.
\]

This chapter uses the finite objective:

\[
\max_{S\subseteq O}
\sum_{r_i\in S}u_i
\]

subject to:

\[
\sum_{r_i\in S}c_i\le B'.
\]

This is a declared compiler objective.

It is not a definition of true context usefulness.

\section{Utility is policy input}\label{utility-is-policy-input}

The score

\[
u_i
\]

can come from:

\begin{itemize}
\tightlist
\item
  a heuristic;
\item
  a learned scorer;
\item
  a human rule;
\item
  a task-specific policy.
\end{itemize}

The finite chapter does not care where it came from.

It only proves which subset is optimal relative to those declared
scores.

Therefore:

\[
\boxed{
\text{utility-optimal under }\Pi
\not\Rightarrow
\text{universally optimal context}.
}
\]

\section{Determinism}\label{determinism}

If several feasible subsets have equal utility, the compiler needs a
tie-break if exact replay matters.

This chapter uses:

\begin{enumerate}
\def\labelenumi{\arabic{enumi}.}
\tightlist
\item
  lower total cost;
\item
  lexicographically smaller stable record-ID list.
\end{enumerate}

The choice is arbitrary but explicit.

Determinism turns the same inputs into the same compiled result.

\section{Serialization is part of
compilation}\label{serialization-is-part-of-compilation}

A selected set is unordered.

A model context is ordered.

Therefore compilation is not finished when subset selection ends.

The finite policy uses:

\[
\text{policy}
\prec
\text{evidence}
\prec
\text{detail}
\prec
\text{background}.
\]

Within a type, stable record IDs break ties.

No empirical claim is made that this order is universally best for
language models.

It is simply part of the declared context object.

\section{The exact budget-8
witness}\label{the-exact-budget-8-witness}

Use context budget:

\[
B=8.
\]

The retrieved candidates are:

{\def\LTcaptype{none} % do not increment counter
\begin{longtable}[]{@{}lllrrrl@{}}
\toprule\noalign{}
ID & key & status/version & rank & cost & utility & mandatory \\
\midrule\noalign{}
\endhead
\bottomrule\noalign{}
\endlastfoot
\(r_{p1}\) & policy & superseded v1 & 1 & 3 & 9 & no \\
\(r_f\) & task\_fact & current v1 & 2 & 4 & 7 & no \\
\(r_{p2}\) & policy & current v2 & 3 & 4 & 8 & yes \\
\(r_b\) & background & current v1 & 4 & 2 & 3 & no \\
\(r_d\) & detail & current v1 & 5 & 3 & 5 & no \\
\end{longtable}
}

All records are authorized.

The current policy v2 supersedes v1.

\section{Freshness removes the rank-1
policy}\label{freshness-removes-the-rank-1-policy}

The latest-current policy excludes:

\[
r_{p1}.
\]

That record had retrieval rank \(1\).

It also had declared utility \(9\).

Neither fact can override its superseded status under this policy.

Remaining records:

\[
r_f,r_{p2},r_b,r_d.
\]

\section{Mandatory policy consumes half the
budget}\label{mandatory-policy-consumes-half-the-budget}

The current policy v2 is mandatory.

Its cost is:

\[
4.
\]

Residual budget:

\[
8-4=4.
\]

The optional candidates are:

\[
r_f:(4,7),
\]

\[
r_b:(2,3),
\]

\[
r_d:(3,5).
\]

Each pair denotes:

\[
(\text{cost},\text{utility}).
\]

\section{Exhaustive finite
selection}\label{exhaustive-finite-selection}

Under residual budget \(4\), feasible optional subsets are:

\[
\varnothing,
\]

\[
\{r_f\},
\]

\[
\{r_b\},
\]

\[
\{r_d\}.
\]

Their utilities are:

\[
0,7,3,5.
\]

Every two-record optional subset costs more than \(4\).

Therefore the unique optimum is:

\[
\boxed{
S^\star=\{r_f\}.
}
\]

\section{The compiled set}\label{the-compiled-set}

Add the mandatory current policy:

\[
\boxed{
S_{\rm final}
=
\{r_{p2},r_f\}.
}
\]

Total cost:

\[
4+4=8.
\]

Total declared utility:

\[
8+7=15.
\]

The serialized context places:

\begin{enumerate}
\def\labelenumi{\arabic{enumi}.}
\tightlist
\item
  current policy v2;
\item
  task evidence.
\end{enumerate}

Both retain provenance metadata.

\section{What the witness actually
proves}\label{what-the-witness-actually-proves}

It proves an exact finite statement about a declared compiler.

It proves:

\begin{itemize}
\tightlist
\item
  which candidate is stale;
\item
  which record is mandatory;
\item
  which optional subsets fit;
\item
  which feasible set maximizes declared utility;
\item
  which records are serialized;
\item
  exact cost;
\item
  exact declared utility.
\end{itemize}

It does not prove that a language model will answer correctly.

\section{The naive rank-prefix
control}\label{the-naive-rank-prefix-control}

Now ignore freshness and mandatory semantics.

Take candidates by retrieval rank until the next record does not fit.

First:

\[
r_{p1}
\]

costs \(3\).

Then:

\[
r_f
\]

costs \(4\).

Cumulative cost:

\[
7.
\]

The next candidate is current mandatory policy v2 with cost \(4\).

It cannot fit.

The rank-prefix result is:

\[
\boxed{
\{r_{p1},r_f\}.
}
\]

\section{Why the rank prefix is
invalid}\label{why-the-rank-prefix-is-invalid}

It contains:

\[
r_{p1},
\]

which is superseded.

It omits:

\[
r_{p2},
\]

which is current and mandatory.

Therefore the result fails the declared compiler contract twice.

This is an exact counterexample to:

\[
\boxed{
\text{high retrieval rank}
\Rightarrow
\text{valid working context}.
}
\]

\section{Truncating the policy is not a
repair}\label{truncating-the-policy-is-not-a-repair}

After the rank prefix, one budget unit remains.

The current policy record costs \(4\).

If the compiler cuts the record into a one-unit fragment, it no longer
has the declared atomic object.

The fragment can lose:

\begin{itemize}
\tightlist
\item
  record ID;
\item
  version;
\item
  source;
\item
  status;
\item
  payload.
\end{itemize}

A valid compiler must revise selection or fail.

It must not pretend a record fragment is the full provenance-aware
record.

\section{Atomicity is
policy-relative}\label{atomicity-is-policy-relative}

This chapter treats records as atomic.

That is a modeling choice.

A different system could define safe sub-record segmentation.

If it does, the segments need their own:

\begin{itemize}
\tightlist
\item
  identity;
\item
  provenance;
\item
  ordering;
\item
  truncation semantics;
\item
  validity rules.
\end{itemize}

The chapter's non-implication is therefore bounded:

\begin{quote}
arbitrary token-level truncation is not authorized by this atomic-record
formalism.
\end{quote}

It does not prove that no principled segmentation scheme can exist.

\section{Compiler validity}\label{compiler-validity}

Define:

\[
\mathrm{Valid}_{\Pi}(K;R,B).
\]

The predicate requires:

\begin{enumerate}
\def\labelenumi{\arabic{enumi}.}
\tightlist
\item
  authorization;
\item
  freshness/version validity;
\item
  mandatory inclusion;
\item
  budget compliance;
\item
  atomicity;
\item
  required provenance;
\item
  declared serialization order.
\end{enumerate}

The witness context satisfies all seven.

The rank-prefix control does not.

\section{Valid context is not correct
answer}\label{valid-context-is-not-correct-answer}

A model can receive a compiler-valid context and still fail.

It can:

\begin{itemize}
\tightlist
\item
  ignore the evidence;
\item
  misread a version;
\item
  invent unsupported content;
\item
  follow the wrong record;
\item
  fail for reasons unrelated to memory.
\end{itemize}

Thus:

\[
\boxed{
\mathrm{Valid}_{\Pi}(K)
\not\Rightarrow
\mathrm{CorrectAnswer}.
}
\]

The compiler certifies its transformation.

It does not certify the model.

\section{Stored is not retrieved}\label{stored-is-not-retrieved}

This chapter keeps several non-implications visible:

\[
\text{stored}
\not\Rightarrow
\text{retrieved}.
\]

A record can exist but never enter the candidate set.

\section{Retrieved is not included}\label{retrieved-is-not-included}

Likewise:

\[
\text{retrieved}
\not\Rightarrow
\text{included}.
\]

A record can be excluded because it is:

\begin{itemize}
\tightlist
\item
  stale;
\item
  unauthorized;
\item
  optional but over budget;
\item
  lower utility under the declared policy;
\item
  in conflict.
\end{itemize}

\section{Included is not used}\label{included-is-not-used}

And:

\[
\text{included}
\not\Rightarrow
\text{used by model}.
\]

The context can contain the right evidence while the model ignores it.

This stage separation prevents one success metric from being misreported
as another.

\section{Budget is an abstract
resource}\label{budget-is-an-abstract-resource}

This chapter writes:

\[
B
\]

as a compiled-context budget.

It can represent:

\begin{itemize}
\tightlist
\item
  tokens;
\item
  bytes;
\item
  slots;
\item
  another declared resource.
\end{itemize}

The accounting must be explicit.

The chapter does not assume all model APIs count context identically.

\section{Prompt construction becomes
replayable}\label{prompt-construction-becomes-replayable}

A compiled context should have a receipt.

Record:

\begin{itemize}
\tightlist
\item
  query/invocation ID;
\item
  compiler policy version;
\item
  context budget;
\item
  candidates;
\item
  versions/statuses;
\item
  authorization decisions;
\item
  freshness exclusions;
\item
  mandatory records;
\item
  selected optional records;
\item
  total cost;
\item
  serialization order;
\item
  source identities;
\item
  explicit failure code.
\end{itemize}

Now context construction can be replayed.

That is stronger than saying:

\begin{quote}
we put the most relevant things in the prompt.
\end{quote}

\section{Replay does not validate retrieval
quality}\label{replay-does-not-validate-retrieval-quality}

Suppose the compiler receipt replays perfectly.

The retriever may still have omitted the most useful record.

Replayability tells us what the compiler did with its inputs.

It does not prove the inputs were ideal.

That requires separate retrieval evidence.

\section{Freshness and ranking answer different
questions}\label{freshness-and-ranking-answer-different-questions}

Freshness asks:

\begin{quote}
is this the version allowed by the invocation policy?
\end{quote}

Ranking asks:

\begin{quote}
how does the retrieval system order candidates under its scoring rule?
\end{quote}

Those questions can disagree.

The exact witness deliberately makes them disagree.

Rank 1 is stale.

Rank 3 is current and mandatory.

The compiler resolves the conflict through policy, not wishful
interpretation of the rank.

\section{Provenance and relevance answer different
questions}\label{provenance-and-relevance-answer-different-questions}

A record can be highly relevant and poorly sourced.

A record can be authoritative and irrelevant to the current task.

Compilation may need both dimensions.

This chapter preserves provenance metadata.

It does not define a universal method for scoring source quality.

\section{Shared memory needs invocation-local
compilation}\label{shared-memory-needs-invocation-local-compilation}

External memory can be shared across agents or model instances.

Working context is typically invocation-local.

Therefore the same shared store can compile differently for:

\begin{itemize}
\tightlist
\item
  different tasks;
\item
  different authorization scopes;
\item
  different budgets;
\item
  different freshness policies;
\item
  different mandatory constraints.
\end{itemize}

Context is a view over memory, not memory itself.

\section{Context compilation is not memory
deletion}\label{context-compilation-is-not-memory-deletion}

Excluding a record from one invocation does not remove it from external
memory.

The compiler chooses a working set.

It does not mutate the source store unless a separate operation
explicitly does so.

This is another reason memory and context need separate objects.

\section{Context compilation is not model
editing}\label{context-compilation-is-not-model-editing}

A compiled prompt can change behavior without changing parameters.

That does not make prompt construction equivalent to parameter editing.

The state loci and update semantics differ.

EXTMEM-001 supplied that distinction.

CONTEXTCOMP keeps it visible.

\section{Practical compiler
protocol}\label{practical-compiler-protocol}

For an invocation:

\begin{enumerate}
\def\labelenumi{\arabic{enumi}.}
\tightlist
\item
  retrieve candidates;
\item
  record candidate identities and ranks;
\item
  apply authorization;
\item
  apply freshness/version rules;
\item
  detect unresolved conflicts;
\item
  identify mandatory records;
\item
  check mandatory feasibility;
\item
  compute complete compiled costs;
\item
  select optional records;
\item
  serialize deterministically;
\item
  emit context plus compilation receipt;
\item
  invoke the model.
\end{enumerate}

If any hard constraint fails, emit an explicit compiler failure.

\section{What this chapter does not
claim}\label{what-this-chapter-does-not-claim-1}

It does not claim that:

\[
\text{retrieval rank}
=
\text{truth}.
\]

It does not claim that:

\[
\text{declared utility}
=
\text{true model usefulness}.
\]

It does not claim that:

\[
\text{more context}
=
\text{better answer}.
\]

It does not claim that:

\[
\text{compiler-valid}
=
\text{answer-correct}.
\]

It does not claim that one universal record ordering is optimal.

\section{Non-implications}\label{non-implications-3}

The chapter rejects:

\[
\text{stored}
\not\Rightarrow
\text{retrieved},
\]

\[
\text{retrieved}
\not\Rightarrow
\text{included},
\]

\[
\text{high retrieval rank}
\not\Rightarrow
\text{fresh/current},
\]

\[
\text{compiled}
\not\Rightarrow
\text{used by model},
\]

\[
\text{compiler-valid}
\not\Rightarrow
\text{answer-correct},
\]

\[
\text{utility-optimal under }\Pi
\not\Rightarrow
\text{universally optimal context},
\]

and

\[
\text{fits budget}
\not\Rightarrow
\text{provenance or authorization valid}.
\]

\section{Atlas connections}\label{atlas-connections-20}

\textbf{Retrieval.}\\
Retrieval supplies candidates; compilation determines the invocation
working set.

\textbf{External Memory.}\\
Context compilation consumes versioned provenance-aware records without
collapsing memory into prompt text.

\textbf{Agents.}\\
An agent can compile different working sets for different subgoals while
sharing one persistent store.

\textbf{Evidence systems.}\\
Compilation receipts can preserve which evidence objects were actually
placed before a model.

\textbf{Context and token compression.}\\
Compression can become one optional selection mechanism, but it remains
constrained by freshness, authority, provenance, and mandatory content.

\textbf{Governed adaptation.}\\
Changing compiler policy changes which external state enters model
computation and therefore deserves versioned governance.

\section{Closing view}\label{closing-view-29}

The model context is small.

The memory can be large.

The job between them is not merely retrieval.

It is compilation.

A serious compiler knows:

\begin{itemize}
\tightlist
\item
  which records are current;
\item
  which records are authorized;
\item
  which records are mandatory;
\item
  what provenance must survive;
\item
  what the budget actually counts;
\item
  which optional records fit;
\item
  how ties are resolved;
\item
  how the final sequence is serialized.
\end{itemize}

The exact witness is small because the principle is structural.

The highest-ranked record can be stale.

The mandatory record can rank lower.

A naive prefix can fit the budget and still be invalid.

The correct conclusion is:

\[
\boxed{
\text{context construction is a governed bounded transformation from candidate memory records to a replayable working set.}
}
\]

That transformation deserves its own mathematics, provenance, and
failure semantics.

\section{References used in this
chapter}\label{references-used-in-this-chapter-49}

External authority is inherited through the audited External-Memory
Thesis, ATLAS-CH-EXTMEM-001 / AUDIT-031.

Exact source identities and authority boundaries are recorded in:

sources/source-locks/ATLAS-CH-CONTEXTCOMP-001.yaml

\chapter{Tokenization and Representation
Boundaries}\label{tokenization-and-representation-boundaries}

\textbf{Epistemic status:} audited Information/Representation
prerequisites + primary tokenizer sources + Atlas synthesis + exact
finite witnesses\\
\textbf{Specification:}
manuscript/specifications/ATLAS-CH-TOKEN-001.md\\
\textbf{Derivation packet:}
mathematics/derivations/ATLAS-CH-TOKEN-001-DERIVATIONS.md\\
\textbf{Computational witness:}
mathematics/computational-witnesses/ATLAS-CW-TOKEN-001.md

A model never receives ``text'' in the abstract.

It receives a representation of text produced by a boundary convention.

That convention decides where symbols begin and end, which units receive
IDs, how long the model's input sequence becomes, and which distinctions
are preserved before the first learned layer acts.

Tokenization is therefore part of the representation interface.

\section{A ladder of boundaries}\label{a-ladder-of-boundaries}

Several units are routinely called characters or tokens even though they
are mathematically different:

\begin{enumerate}
\def\labelenumi{\arabic{enumi}.}
\tightlist
\item
  bytes;
\item
  Unicode codepoints;
\item
  grapheme-like visible characters;
\item
  words under a declared word-boundary convention;
\item
  subword pieces;
\item
  token IDs;
\item
  learned token embeddings.
\end{enumerate}

A count at one level is not a count at another.

That sounds elementary, but many comparisons silently move between these
levels.

\section{Tokenization as a map}\label{tokenization-as-a-map}

For normalized text space \(\mathcal X\) and vocabulary \(V\), write a
deterministic tokenizer as

\[
\tau:\mathcal X\to V^*.
\]

A stochastic tokenizer instead defines a conditional distribution

\[
q_\tau(s\mid x)
\]

over valid segmentations \(s\) of \(x\).

Detokenization is another map:

\[
\delta:V^*\to\mathcal X'.
\]

Whether

\[
\delta(\tau(x))=x
\]

holds exactly depends on normalization, whitespace handling, byte
conventions, special tokens, and the tokenizer design.

Lossless detokenization is an engineering property, not an automatic
fact about all tokenizers.

\section{Bytes and codepoints are not the same
unit}\label{bytes-and-codepoints-are-not-the-same-unit}

Take the visible string \texttt{é} in NFC form.

It is one Unicode codepoint, U+00E9, but two UTF-8 bytes.

Normalize it into the decomposed form \texttt{e} plus U+0301.

Now it is two codepoints and three UTF-8 bytes.

Thus:

\[
\text{byte length}\neq\text{codepoint length},
\]

and Unicode normalization can change both.

This also warns against treating codepoint count as identical to
user-perceived character count.

\section{Why subwords became
useful}\label{why-subwords-became-useful}

Sennrich, Haddow, and Birch introduced a practical BPE-inspired subword
strategy for neural machine translation to reduce the
rare-word/open-vocabulary problem (Sennrich et al. 2016).

The important conceptual move is between whole-word and character
extremes.

A recurring string can become one reusable unit while rarer forms remain
decomposable.

That gives a finite vocabulary a route to representing strings it did
not store as whole words.

The source establishes this method in its NMT setting.

It does not prove that BPE is universally optimal.

\section{BPE is an algorithm, not merely a
vocabulary}\label{bpe-is-an-algorithm-not-merely-a-vocabulary}

A BPE-style encoder starts from a base symbolization and applies learned
pair merges in a declared order or ranking.

The procedure matters.

A list of piece strings is not always enough to reproduce segmentation.

A reproducible tokenizer therefore needs more than ``vocabulary size =
32k.''

It should bind, at minimum:

\begin{itemize}
\tightlist
\item
  normalization;
\item
  base symbols or pretokenization;
\item
  merge ranks or equivalent encoder state;
\item
  special-token rules;
\item
  vocabulary/ID mapping.
\end{itemize}

This is a recurring Atlas lesson: an interface is more than a tensor
shape or a list of names.

\section{Unigram tokenization treats segmentation
differently}\label{unigram-tokenization-treats-segmentation-differently}

Kudo's unigram method assigns scores to candidate pieces and evaluates
whole segmentations under a unigram-language-model assumption (Kudo
2018).

That makes segmentation ambiguity explicit.

The same vocabulary can support several legal decompositions of the same
string.

One can select the highest-scoring segmentation, or sample alternatives
during training.

Those are different inference rules over the same segmentation space.

\section{Exact ambiguity witness}\label{exact-ambiguity-witness}

Let

\[
V=\{a,b,ab\}
\]

with

\[
p(ab)=\frac12,
\qquad
p(a)=p(b)=\frac14.
\]

For \texttt{abab}, consider:

\[
[ab,ab],
\]

\[
[a,b,ab],
\]

\[
[ab,a,b],
\]

and

\[
[a,b,a,b].
\]

Their unigram product weights are

\[
\frac14,
\quad
\frac1{32},
\quad
\frac1{32},
\quad
\frac1{256}.
\]

After normalization, the probabilities are

\[
\left(
\frac{64}{81},
\frac8{81},
\frac8{81},
\frac1{81}
\right).
\]

The first segmentation is the mode.

It is not the only segmentation admitted by the model.

\section{SentencePiece moves the boundary
again}\label{sentencepiece-moves-the-boundary-again}

SentencePiece was designed to train subword models directly from raw
sentences rather than requiring conventional pretokenized word sequences
(Kudo and Richardson 2018).

This matters because word pretokenization is itself a language-dependent
representation choice.

Removing that requirement does not remove every preprocessing decision.

Normalization, vocabulary construction, segmentation model,
special-token policy, and byte/codepoint handling still remain part of
the interface.

\section{Token count is a useful but narrow
quantity}\label{token-count-is-a-useful-but-narrow-quantity}

Define token length

\[
L_\tau(x)=|\tau(x)|.
\]

It directly controls how many discrete input positions a token-level
Transformer sees.

But it is not entropy.

It is not automatically code length.

It is not a complete compute model.

Attention, embedding lookup, softmax vocabulary cost, architecture,
batching, sequence packing, and hardware all matter downstream.

\section{Token count is not description
length}\label{token-count-is-not-description-length}

Under a declared token probability model \(p\), idealized code length is

\[
\ell_p(s)=-\sum_i\log_2p(v_i).
\]

Take a two-token sequence in which each token has probability \(1/16\).

Its idealized length is 8 bits.

Now take a three-token sequence in which each token has probability
\(1/2\).

Its idealized length is 3 bits.

Therefore:

\[
|s_A|<|s_B|
\not\Rightarrow
\ell_p(s_A)<\ell_p(s_B).
\]

This preserves the Information chapter's firewall: description length is
defined relative to a probability/coding model.

\section{Fertility measures segmentation
expansion}\label{fertility-measures-segmentation-expansion}

For corpus \(C\) with declared reference-word count \(W(C)\), define

\[
F_\tau(C)
=
\frac{\sum_{x\in C}|\tau(x)|}{W(C)}.
\]

A fertility of 1 means one tokenizer unit per reference word on average.

A fertility of 2.5 means two and a half units per reference word on
average.

The number is useful precisely because it is simple.

It is also easy to overinterpret.

Fertility depends on the corpus, normalization, reference-word
convention, and tokenizer.

It is not a fixed property of a language.

\section{Exact fertility witness}\label{exact-fertility-witness}

Take the two-word toy corpus:

\texttt{unbelievable\ cats}.

Tokenizer A returns

\texttt{un\ \textbar{}\ believ\ \textbar{}\ able\ \textbar{}\ cat\ \textbar{}\ s}.

Five tokens over two reference words gives

\[
F_A=\frac52.
\]

Tokenizer B returns

\texttt{unbelievable\ \textbar{}\ cats},

so

\[
F_B=1.
\]

Tokenizer B produces the shorter model sequence.

That alone does not tell us which tokenizer gives better morphology,
likelihood, transfer, robustness, compute efficiency, or downstream
quality.

\section{Morphology gives a different
diagnostic}\label{morphology-gives-a-different-diagnostic}

Suppose a declared linguistic analysis places internal boundaries in a
word.

Let \(B_M(x)\) be those boundaries and \(B_\tau(x)\) the tokenizer's
internal boundaries.

One can define boundary precision and recall:

\[
P_M
=
\frac{|B_M\cap B_\tau|}{|B_\tau|},
\]

\[
R_M
=
\frac{|B_M\cap B_\tau|}{|B_M|}.
\]

For the toy analysis

\texttt{un\ \textbar{}\ believ\ \textbar{}\ able},

an identical tokenization scores perfectly.

A segmentation \texttt{unbel\ \textbar{}\ ievable} need not hit either
declared morphology boundary.

This is a legitimate diagnostic.

It is not a theorem that morphology-aligned tokenization is always
superior for neural language modeling.

\section{Morphology and subword frequency optimize different
objects}\label{morphology-and-subword-frequency-optimize-different-objects}

BPE primarily responds to frequent adjacent symbol patterns under its
training procedure.

Unigram tokenization chooses a probabilistic piece inventory and
segmentation model.

A linguistic morphology may instead aim to identify stems, affixes,
roots, or other language-specific units.

These objectives can agree.

They can also disagree.

A frequent string is not necessarily a morpheme.

A morpheme is not necessarily frequent enough to earn its own vocabulary
item.

\section{Byte-level models remove one tokenizer
bottleneck}\label{byte-level-models-remove-one-tokenizer-bottleneck}

ByT5 studies models that operate directly on UTF-8 bytes and reports
robustness and modeling trade-offs relative to token-level counterparts
(Xue et al. 2022).

A byte alphabet can cover arbitrary UTF-8 text without requiring a
learned word/subword vocabulary.

That removes one kind of out-of-vocabulary boundary.

It does not make representation free.

Byte sequences are often longer.

Longer sequences alter compute and memory costs.

The empirical trade-off belongs to a declared model and workload.

\section{Byte-level does not mean text has no
structure}\label{byte-level-does-not-mean-text-has-no-structure}

A byte-level model still receives an ordered symbol sequence.

UTF-8 itself is an encoding convention.

Unicode normalization still matters.

The model must learn useful regularities across multibyte sequences and
longer spans.

``Token-free'' therefore means free of a learned word/subword tokenizer
in this context, not free of discrete representation choices altogether.

\section{Multilingual tokenization creates allocation
problems}\label{multilingual-tokenization-creates-allocation-problems}

A shared vocabulary has finite capacity.

Across many languages, that capacity must cover different scripts,
frequency distributions, morphology, and corpus sizes.

Rust et al.~conduct controlled comparisons showing that tokenizer
adequacy can materially affect monolingual downstream performance inside
multilingual models (Rust et al. 2021).

Their result motivates measurement.

It does not justify reducing multilingual performance to one tokenizer
statistic.

Training-data scale, architecture, pretraining objective, vocabulary
allocation, and downstream task all remain potential factors.

\section{A multilingual tokenizer report should be
vector-valued}\label{a-multilingual-tokenizer-report-should-be-vector-valued}

For a language/corpus \(C_\ell\), useful quantities can include:

\begin{itemize}
\tightlist
\item
  fertility;
\item
  unknown-token rate, if the tokenizer has unknowns;
\item
  byte-per-token or codepoint-per-token expansion;
\item
  vocabulary coverage;
\item
  morphology-boundary alignment;
\item
  sequence-length tails;
\item
  downstream performance under controlled model/data comparisons.
\end{itemize}

These quantities do not share units.

A single scalar ``tokenizer quality'' score requires an explicit
aggregation rule.

\section{Token IDs are coordinates, not
meaning}\label{token-ids-are-coordinates-not-meaning}

After segmentation, each piece receives an ID.

That ID indexes a learned embedding vector.

The Representation chapter's boundary applies directly:

\[
\text{coordinate}\neq\text{semantics}.
\]

Changing the ID assignment while consistently permuting the embedding
table leaves the representational interface equivalent.

Changing the segmentation itself changes sequence length and
decomposition, which is a different intervention.

\section{Tokenization changes the learning
problem}\label{tokenization-changes-the-learning-problem}

Two tokenizers can detokenize to the same text while presenting the
model with different:

\begin{itemize}
\tightlist
\item
  sequence lengths;
\item
  boundary locations;
\item
  frequency distributions;
\item
  embedding lookup events;
\item
  context allocation;
\item
  prediction targets;
\item
  softmax classes.
\end{itemize}

Thus tokenizer replacement is not a neutral input formatting change.

It changes the model's representation boundary and often its training
objective at the discrete-symbol level.

\section{Comparing tokenizers requires controlled
scope}\label{comparing-tokenizers-requires-controlled-scope}

A credible tokenizer comparison should state:

\begin{itemize}
\tightlist
\item
  normalization;
\item
  tokenizer algorithm and version/state;
\item
  vocabulary size;
\item
  training corpus;
\item
  language mixture;
\item
  special tokens;
\item
  byte/codepoint convention;
\item
  sequence truncation policy;
\item
  model architecture;
\item
  parameter-count accounting;
\item
  training token/byte/character budget;
\item
  evaluation corpus and downstream metric.
\end{itemize}

Without these fields, ``tokenizer A is more efficient'' can conceal
several different interventions.

\section{What this chapter does not
claim}\label{what-this-chapter-does-not-claim-2}

TOKEN does not claim:

\begin{itemize}
\tightlist
\item
  BPE is universally superior to unigram tokenization;
\item
  unigram sampling is always beneficial;
\item
  byte models are universally more robust or efficient;
\item
  low fertility is always good;
\item
  high morphology alignment is always good;
\item
  multilingual vocabulary sharing is always harmful or beneficial;
\item
  token count alone determines compute;
\item
  tokenizer choice alone determines multilingual performance.
\end{itemize}

The chapter supplies typed objects and diagnostics, not a universal
ranking.

\section{Handoff to Tokenization as Compression and
Interface}\label{handoff-to-tokenization-as-compression-and-interface}

ATLAS-CH-TOKENCOMP-001 may now assume:

\begin{itemize}
\tightlist
\item
  tokenization as an explicit representation boundary;
\item
  byte/codepoint/subword separation;
\item
  fertility and relative fertility;
\item
  probability-model code length distinct from token count;
\item
  BPE and unigram algorithms as different tokenizer families;
\item
  morphology-boundary diagnostics;
\item
  multilingual measurement as a vector rather than a single score;
\item
  the exact normalization, segmentation-ambiguity, code-length, and
  fertility witnesses.
\end{itemize}

TOKENCOMP must add the harder systems question:

\begin{quote}
how should vocabulary design trade description length, model compute,
embedding/softmax cost, interoperability, and cross-model tokenizer
interfaces?
\end{quote}

\section{Durable lesson}\label{durable-lesson-2}

Tokenization chooses the discrete units on which the rest of the model
operates.

That makes it neither mere preprocessing nor hidden semantics.

It is a representation boundary with measurable costs, ambiguities, and
language-dependent effects.

\section{References used in this
chapter}\label{references-used-in-this-chapter-50}

\begin{itemize}
\tightlist
\item
  (Sennrich et al. 2016)
\item
  (Kudo 2018)
\item
  (Kudo and Richardson 2018)
\item
  (Xue et al. 2022)
\item
  (Rust et al. 2021)
\end{itemize}

Exact source authority, prerequisite identities, and claim boundaries
are locked in:

sources/source-locks/ATLAS-CH-TOKEN-001.yaml

\chapter{Tokenization as Compression and
Interface}\label{tokenization-as-compression-and-interface}

\textbf{Epistemic status:} audited Tokenization prerequisite +
Atlas-owned exact compression/compute/interoperability constructions.\\
\textbf{Specification:}
manuscript/specifications/ATLAS-CH-TOKENCOMP-001.md\\
\textbf{Derivation packet:}
mathematics/derivations/ATLAS-CH-TOKENCOMP-001-DERIVATIONS.md\\
\textbf{Computational witness:}
mathematics/computational-witnesses/ATLAS-CW-TOKENCOMP-001.md\\
\textbf{Source lock:} sources/source-locks/ATLAS-CH-TOKENCOMP-001.yaml

TOKEN-001 treated tokenization as a representation boundary.

This chapter asks what happens when that boundary is also treated as a
compression and systems interface.

The immediate temptation is to rank tokenizers by one number:

\begin{quote}
fewer tokens is better.
\end{quote}

That rule is too coarse.

Vocabulary size changes how many distinct IDs must be represented.

Segmentation changes sequence length.

A model may pay costs that grow differently with sequence length and
vocabulary size.

Two systems may need to exchange text while using incompatible token
IDs.

The governing rule is:

\[
\boxed{
\text{tokenizer design is a multi-objective representation/interface choice, not a token-count scalar.}
}
\]

\section{The inherited boundary}\label{the-inherited-boundary}

TOKEN-001 already established:

\[
\tau:\mathcal X\to V^*
\]

for a deterministic tokenizer, together with a separate detokenizer:

\[
\delta:V^*\to\mathcal X'.
\]

It also established that:

\begin{itemize}
\tightlist
\item
  byte length is not codepoint length;
\item
  token count is not description length;
\item
  fertility is corpus-relative;
\item
  vocabulary is not the whole segmentation algorithm;
\item
  token IDs and embeddings are representations, not semantics.
\end{itemize}

TOKENCOMP keeps all of those firewalls.

\section{Compression requires a
code}\label{compression-requires-a-code}

A token sequence is not yet a bit string.

To ask how many bits it occupies, we need a code.

One simple toy choice is a fixed-width code over token IDs.

Let:

\[
m=|V|\ge2.
\]

Then a fixed-width ID needs:

\[
w=\lceil\log_2m\rceil
\]

bits.

For a token sequence of length \(n\):

\[
D_{\mathrm{fw}}=nw.
\]

This is a code definition.

It is not an intrinsic property of the text.

\section{Vocabulary size enters the
accounting}\label{vocabulary-size-enters-the-accounting}

Suppose one tokenizer halves the number of tokens but increases
vocabulary size enough to require more bits per ID.

The two effects compete.

So:

\[
\boxed{
\text{shorter token sequence}
\not\Rightarrow
\text{shorter fixed-width ID stream}.
}
\]

This is the first compression/interface obstruction.

\section{Fixed-width length is not entropy
coding}\label{fixed-width-length-is-not-entropy-coding}

TOKEN-001 inherited model-based idealized code length:

\[
\ell_p(s)
=
-\sum_i\log_2p(v_i).
\]

That quantity depends on a declared probability model.

Fixed-width length instead depends only on sequence length and
vocabulary cardinality under the declared code.

They can disagree.

A complete comparison must state which description-length convention is
being used.

\section{A tokenizer is part of the
codebook}\label{a-tokenizer-is-part-of-the-codebook}

The token-ID stream can be decoded only if the recipient already knows
the tokenizer specification.

That can include:

\begin{itemize}
\tightlist
\item
  normalization;
\item
  base alphabet;
\item
  vocabulary;
\item
  token-to-ID mapping;
\item
  merge ranks or unigram scores;
\item
  special-token rules;
\item
  decoder behavior;
\item
  version identity.
\end{itemize}

If that state is not shared, its description or transmission cost
belongs somewhere in the protocol.

Therefore:

\[
\boxed{
\text{ID-stream length}
\neq
\text{complete compressor size}.
}
\]

\section{Compression ratio needs a
baseline}\label{compression-ratio-needs-a-baseline}

A ratio is meaningful only relative to something.

For byte string \(x\) with byte length \(B(x)>0\), define one toy ratio:

\[
R_{\mathrm{fw}}
=
\frac{D_{\mathrm{fw}}}{8B(x)}.
\]

This compares a token-ID stream to an 8-bit-per-byte baseline.

It does not count tokenizer state.

It does not claim optimal compression.

\section{Token count is not
compute}\label{token-count-is-not-compute}

A model may perform several kinds of work.

Some costs can grow strongly with sequence length.

Others can grow with vocabulary size.

Still others depend on hidden dimension, sparsity, batching, memory,
hardware, or implementation.

So the chapter refuses the shortcut:

\[
\text{token count}
=
\text{compute}.
\]

Instead it declares toy cost functions explicitly.

\section{Two exact compute proxies}\label{two-exact-compute-proxies}

For token length \(n\), define:

\[
C_{\mathrm{pair}}(n)=n^2.
\]

This counts all ordered token-position pairs in a toy pair-interaction
model.

For vocabulary size \(m\), define:

\[
C_{\mathrm{vocab}}(n,m)=nm.
\]

This counts one toy dense token-by-vocabulary scoring grid.

These are not universal model FLOP equations.

They are exact counters inside the declared finite model.

\section{The finite witness}\label{the-finite-witness}

Take:

\[
x=\texttt{abababab}.
\]

It is eight ASCII bytes.

Thus the byte baseline contains:

\[
64
\]

bits.

We compare two lossless tokenizers.

\section{Small-vocabulary
tokenizer}\label{small-vocabulary-tokenizer}

Let:

\[
V_S=\{a,b\}.
\]

The tokenizer emits:

\[
[a,b,a,b,a,b,a,b].
\]

So:

\[
n_S=8,
\qquad
m_S=2.
\]

One bit identifies each of two token IDs.

Hence:

\[
D_S=8\cdot1=8\text{ bits}.
\]

The toy compute vector is:

\[
K_S=(8^2,8\cdot2)=(64,16).
\]

\section{Piece tokenizer}\label{piece-tokenizer}

Let:

\[
V_P=\{a,b,ab,x,y,z,w,q\}.
\]

The witness text is segmented as:

\[
[ab,ab,ab,ab].
\]

Thus:

\[
n_P=4,
\qquad
m_P=8.
\]

Three bits are required for one of eight fixed-width IDs.

Therefore:

\[
D_P=4\cdot3=12\text{ bits}.
\]

The toy compute vector is:

\[
K_P=(4^2,4\cdot8)=(16,32).
\]

\section{The orderings conflict}\label{the-orderings-conflict}

The piece tokenizer uses fewer tokens:

\[
4<8.
\]

Its pair proxy is smaller:

\[
16<64.
\]

But its fixed-width ID stream is larger:

\[
12>8.
\]

And its dense-vocabulary proxy is larger:

\[
32>16.
\]

Thus the same design change improves one declared cost while worsening
others.

\section{There is no scalar winner without an
objective}\label{there-is-no-scalar-winner-without-an-objective}

Write:

\[
J_\tau=
(D_{\mathrm{fw}},C_{\mathrm{pair}},C_{\mathrm{vocab}}).
\]

Then:

\[
J_S=(8,64,16),
\]

\[
J_P=(12,16,32).
\]

Neither is componentwise smaller.

So neither Pareto-dominates the other in this toy objective space.

To choose a winner we must add:

\begin{itemize}
\tightlist
\item
  weights;
\item
  constraints;
\item
  workload frequencies;
\item
  a model architecture;
\item
  or another explicit objective.
\end{itemize}

The tokenizer alone does not supply those choices.

\section{Shorter can still matter}\label{shorter-can-still-matter}

The counterexample does not say token length is irrelevant.

A shorter sequence can reduce any declared cost that increases with
sequence length while other variables remain fixed.

The point is narrower:

\begin{quote}
sequence length is one coordinate, not the whole resource model.
\end{quote}

\section{Vocabulary can buy reusable
pieces}\label{vocabulary-can-buy-reusable-pieces}

A larger vocabulary can encode recurring substrings as single tokens.

That can reduce sequence length.

But the vocabulary also creates state:

\begin{itemize}
\tightlist
\item
  more token identities;
\item
  a wider fixed-width ID alphabet in the toy code;
\item
  potentially more tokenizer specification;
\item
  a different representation boundary.
\end{itemize}

The useful question is not ``large or small vocabulary?''

It is:

\begin{quote}
which objective vector is being optimized for which corpus and system?
\end{quote}

\section{Corpus distribution
matters}\label{corpus-distribution-matters}

Our witness contains only repeated ``ab''.

A different corpus can reverse which pieces are useful.

For corpus \(C\), one might compare:

\[
\mathbb E_{x\sim C}[L_\tau(x)]
\]

or:

\[
\mathbb E_{x\sim C}[D_{\mathrm{fw}}(x;\tau)].
\]

Those expectations depend on the corpus distribution.

No finite example creates a language-independent ranking.

\section{Tokenizer
interoperability}\label{tokenizer-interoperability}

Now consider two systems that use different tokenizers.

A direct exchange of token IDs is unsafe unless the two sides share
exactly the relevant tokenizer state.

Integer ID 17 in one vocabulary need not correspond to ID 17 in another.

Even identical token strings can have different IDs.

The interface must name what crosses the boundary.

\section{Canonical transport
domain}\label{canonical-transport-domain}

Let:

\[
\mathcal X_c
\]

be a declared domain of canonical byte strings.

Tokenizer A has:

\[
\tau_A:\mathcal X_c\to V_A^*,
\]

with decoder:

\[
\delta_A:V_A^*\to\mathcal X_c.
\]

Tokenizer B similarly has:

\[
\tau_B,\delta_B.
\]

Assume both round trip exactly on the declared domain:

\[
\delta_A(\tau_A(x))=x,
\]

\[
\delta_B(\tau_B(x))=x.
\]

\section{Lossless translation
theorem}\label{lossless-translation-theorem}

Define:

\[
T_{A\to B}
=
\tau_B\circ\delta_A.
\]

Then:

\[
T_{A\to B}(\tau_A(x))
=
\tau_B(x).
\]

Applying B's decoder:

\[
\delta_B(T_{A\to B}(\tau_A(x)))
=
x.
\]

So the canonical bytes survive translation exactly.

This is an interface theorem.

It is not a model-behavior theorem.

\section{The witness translates
exactly}\label{the-witness-translates-exactly}

For the finite example:

Tokenizer S emits:

\[
[a,b,a,b,a,b,a,b].
\]

Its decoder concatenates these pieces back to:

\[
\texttt{abababab}.
\]

Tokenizer P then emits:

\[
[ab,ab,ab,ab].
\]

The reverse process recovers the S sequence.

The token identities and sequence lengths change.

The canonical byte string does not.

\section{Lossless does not mean
representation-equivalent}\label{lossless-does-not-mean-representation-equivalent}

Suppose system A looks up embeddings using table \(E_A\) and B uses
\(E_B\).

Their internal sequences can be:

\[
E_A(\tau_A(x))
\]

and:

\[
E_B(\tau_B(x)).
\]

Canonical byte equality supplies no equality between these vector
sequences.

Therefore:

\[
\boxed{
\text{lossless transport}
\not\Rightarrow
\text{equal internal representation}.
}
\]

\section{Nor does lossless mean semantic
equivalence}\label{nor-does-lossless-mean-semantic-equivalence}

Canonical byte equality is a syntactic identity under the declared
interface.

It does not prove:

\begin{itemize}
\tightlist
\item
  equal human interpretation under every context;
\item
  equal model probability;
\item
  equal model activation;
\item
  equal downstream answer;
\item
  equal learned semantics.
\end{itemize}

The representation layer has been translated.

Nothing stronger follows automatically.

\section{Direct ID interoperability is
stronger}\label{direct-id-interoperability-is-stronger}

A direct token-ID mapping would avoid decode and retokenize.

But that requires more structure.

A simple permutation works only when the token alphabets and
segmentation semantics align appropriately.

In the witness, one P token ``ab'' corresponds to two S tokens.

So no one-to-one token permutation can implement the translation.

Canonical translation is more general.

\section{Special tokens require protocol
semantics}\label{special-tokens-require-protocol-semantics}

Real token streams may contain control tokens not meant to decode as
ordinary text.

Examples can include boundary, role, padding, or control symbols.

Such objects fall outside the simple canonical-byte theorem unless their
transport semantics are explicitly defined.

A tokenizer interface therefore needs both:

\begin{itemize}
\tightlist
\item
  text/byte semantics;
\item
  control-token semantics.
\end{itemize}

The chapter proves only the first finite case.

\section{Version identity matters}\label{version-identity-matters}

A tokenizer is versioned protocol state.

Changing any of these can change the meaning of an ID stream:

\begin{itemize}
\tightlist
\item
  normalization;
\item
  vocabulary;
\item
  merge order;
\item
  token scores;
\item
  ID assignment;
\item
  special tokens;
\item
  decoder.
\end{itemize}

So ``uses BPE'' is not a sufficient interoperability identifier.

Neither is a vocabulary size.

\section{Three meanings of lingua
franca}\label{three-meanings-of-lingua-franca}

The phrase \textbf{tokenizer lingua franca} can refer to different
architectures.

\subsection{Canonical transport}\label{canonical-transport}

Systems exchange normalized bytes/text and retokenize locally.

\subsection{Shared tokenizer}\label{shared-tokenizer}

Systems adopt one tokenization specification and exchange its IDs.

\subsection{Translation layer}\label{translation-layer}

Systems keep local tokenizers and use explicit adapters.

These are distinct designs.

\section{Canonical bytes as a bounded lingua
franca}\label{canonical-bytes-as-a-bounded-lingua-franca}

A canonical byte domain has one useful property:

different lossless tokenizers can translate through it without requiring
matching vocabularies.

That makes it a candidate \textbf{transport interface}.

But this chapter does not claim:

\begin{itemize}
\tightlist
\item
  that every model should use bytes internally;
\item
  that byte transport is always efficient;
\item
  that a particular normalization should be universal;
\item
  that all control tokens can be represented this way;
\item
  that industry has adopted one standard.
\end{itemize}

``Candidate lingua franca'' is therefore an architectural idea, not an
adoption fact.

\section{Shared tokenizers trade interoperability for
constraints}\label{shared-tokenizers-trade-interoperability-for-constraints}

A shared tokenizer can remove translation ambiguity at the token-ID
layer.

It also couples systems to:

\begin{itemize}
\tightlist
\item
  one vocabulary allocation;
\item
  one normalization;
\item
  one version lifecycle;
\item
  one special-token namespace.
\end{itemize}

That can be desirable in one architecture and restrictive in another.

No universal choice follows.

\section{Translation layers preserve local
freedom}\label{translation-layers-preserve-local-freedom}

A translation layer lets each subsystem retain its own tokenizer.

That preserves local representation choices.

But it adds:

\begin{itemize}
\tightlist
\item
  encode/decode work;
\item
  protocol state;
\item
  version management;
\item
  possible normalization hazards;
\item
  control-token mapping obligations.
\end{itemize}

Again, interoperability has costs.

\section{Compression and interoperability are different
objectives}\label{compression-and-interoperability-are-different-objectives}

Tokenizer S compresses the witness's fixed-width ID stream better than
P.

That tells us nothing by itself about which is easier to standardize
across systems.

Conversely, a widely shared tokenizer could improve ID-level
interoperability without minimizing description length on a particular
corpus.

Thus:

\[
\boxed{
\text{compression quality}
\neq
\text{interoperability quality}.
}
\]

\section{Compute and interoperability are different
objectives}\label{compute-and-interoperability-are-different-objectives}

Tokenizer P reduces the witness's pair proxy.

That does not make its token IDs meaningful to another system.

A compute optimization inside one model can worsen a protocol boundary
if it increases specialization or version coupling.

System objectives must remain typed.

\section{A practical comparison
vector}\label{a-practical-comparison-vector}

A tokenizer/interface comparison might track:

\[
M_\tau(C)=
(
L,
D,
F,
C_{\mathrm{pair}},
C_{\mathrm{vocab}},
I_{\mathrm{roundtrip}},
I_{\mathrm{control}},
Q
),
\]

where the coordinates could denote:

\begin{itemize}
\tightlist
\item
  token length;
\item
  declared description length;
\item
  fertility;
\item
  compute proxies or measured costs;
\item
  round-trip interoperability;
\item
  control-token interoperability;
\item
  downstream quality.
\end{itemize}

These coordinates do not share one unit.

A scalar score requires an explicit objective.

\section{Multilingual scope remains
explicit}\label{multilingual-scope-remains-explicit}

TOKEN-001 already required language/corpus-specific measurement.

TOKENCOMP preserves that rule.

A shared vocabulary may allocate capacity differently across scripts and
languages.

A canonical byte interface may be lossless while token lengths vary
substantially across local tokenizers.

Neither fact creates a universal multilingual optimum.

\section{Modeled compute versus measured
compute}\label{modeled-compute-versus-measured-compute}

The exact witness uses:

\[
n^2
\]

and:

\[
n|V|
\]

as toy counters.

A real system can have:

\begin{itemize}
\tightlist
\item
  different algorithms;
\item
  sparsity;
\item
  caching;
\item
  batching;
\item
  hardware kernels;
\item
  memory effects;
\item
  parallelism;
\item
  compiler transformations.
\end{itemize}

Therefore the toy counters must not be relabeled as measured runtime.

A runtime claim requires a runtime experiment.

\section{What would an empirical comparison
require?}\label{what-would-an-empirical-comparison-require}

At minimum:

\begin{itemize}
\tightlist
\item
  exact tokenizer versions;
\item
  normalization;
\item
  corpus/language distribution;
\item
  vocabulary sizes;
\item
  sequence-length distribution;
\item
  code-length convention;
\item
  model architecture;
\item
  model parameterization affected by vocabulary size;
\item
  training or serving workload;
\item
  hardware/software stack for runtime;
\item
  downstream quality metrics.
\end{itemize}

Without these, ``tokenizer A is more efficient'' is underspecified.

\section{Durable non-implications}\label{durable-non-implications-2}

The chapter preserves:

\[
\text{fewer tokens}
\not\Rightarrow
\text{fewer encoded bits},
\]

\[
\text{fewer tokens}
\not\Rightarrow
\text{lower cost under every compute model},
\]

\[
\text{lossless translation}
\not\Rightarrow
\text{equal tokenization},
\]

\[
\text{equal decoded bytes}
\not\Rightarrow
\text{equal embeddings or behavior},
\]

\[
\text{shared interface}
\not\Rightarrow
\text{universal standard}.
\]

\section{Tokenization as interface
design}\label{tokenization-as-interface-design}

The central systems lesson is not that one tokenizer should dominate.

It is that a tokenizer boundary carries several contracts at once:

\begin{itemize}
\tightlist
\item
  representation;
\item
  compression;
\item
  compute shape;
\item
  serialization;
\item
  versioning;
\item
  interoperability.
\end{itemize}

Those contracts can favor different designs.

A mature tokenizer choice therefore begins by naming the objective
vector rather than by minimizing token count in isolation.

\section{References used in this
chapter}\label{references-used-in-this-chapter-51}

No new external academic authority is added in TOKENCOMP-001.

Tokenizer mechanisms, probability-model code length, multilingual
boundaries, and representation-interface authority are inherited through
the audited TOKEN-001 source lock. TOKENCOMP's new
compression/compute/interoperability results are exact Atlas-owned
constructions.

Exact source identity and claim boundaries are recorded in:

sources/source-locks/ATLAS-CH-TOKENCOMP-001.yaml

\chapter{Data Quality, Mixtures, and
Contamination}\label{data-quality-mixtures-and-contamination}

\textbf{Epistemic status:} established empirical dataset results +
audited Evidence prerequisite + Atlas synthesis + exact finite corpus
witness.\\
\textbf{Specification:} manuscript/specifications/ATLAS-CH-DATA-001.md\\
\textbf{Formal packet:}
mathematics/derivations/ATLAS-CH-DATA-001-DERIVATIONS.md\\
\textbf{Computational witness:}
mathematics/computational-witnesses/DATA-001.md\\
\textbf{Source lock:} sources/source-locks/ATLAS-CH-DATA-001.yaml

A dataset is not a bag of text.

It is a governed sampling object with lineage.

Two collections can contain the same visible records and still train
different models because they differ in:

\begin{itemize}
\tightlist
\item
  duplication;
\item
  filters;
\item
  provenance;
\item
  partition boundaries;
\item
  sampling weights;
\item
  synthetic-data lineage;
\item
  evaluation overlap.
\end{itemize}

The governing rule of this chapter is:

\begin{quote}
data quality is not one scalar property, and contamination is not one
predicate.
\end{quote}

Every empirical claim must retain the exact dataset construction and
evaluation threat model that support it.

\section{Start with the data
object}\label{start-with-the-data-object}

Use:

\texttt{D=(R,S,P,Phi,Delta,mu,E)}.

Here:

\begin{itemize}
\tightlist
\item
  \texttt{R} is the logical record set;
\item
  \texttt{S} contains source and domain labels;
\item
  \texttt{P} records provenance and lineage;
\item
  \texttt{Phi} is the transformation/filtering pipeline;
\item
  \texttt{Delta} contains duplicate and overlap relations;
\item
  \texttt{mu} is the training sampling measure;
\item
  \texttt{E} is the evaluation boundary.
\end{itemize}

This object is richer than a directory of files.

\section{Why lineage matters}\label{why-lineage-matters}

Suppose two records contain identical text.

One may be copied from the other.

Both may descend from the same upstream source.

One may be human-authored and the other generated.

One may have entered through a benchmark-derived instructional set.

The bytes alone do not recover this history.

Provenance is therefore part of the data object.

This follows directly from the Evidence chapter's rule that support and
source identity remain attached to claims.

\section{Exact duplicates}\label{exact-duplicates}

Let \texttt{canon(r)} be a declared canonicalization.

Define:

\texttt{delta\_exact(a,b)\ =\ 1\{canon(a)=canon(b)\}}.

Even ``exact'' requires a method.

Do we normalize whitespace?

Unicode?

HTML?

Case?

Boilerplate?

A deduplication claim without canonicalization semantics is incomplete.

\section{Near duplicates}\label{near-duplicates}

Exact equality misses modified copies.

Let \texttt{T(r)} be a token-set representation.

A simple near-duplicate score is token Jaccard:

\texttt{J(a,b)\ =\ \textbar{}T(a)\ intersect\ T(b)\textbar{}\ /\ \textbar{}T(a)\ union\ T(b)\textbar{}}.

At threshold \texttt{tau}:

\texttt{delta\_near\_tau(a,b)\ =\ 1\{J(a,b)\textgreater{}=tau\}}.

The result depends on both representation and threshold.

Lee et al.~demonstrate why exact and near-duplicate structure both
matter in language-model corpora.

\section{Semantic redundancy is different
again}\label{semantic-redundancy-is-different-again}

Two records can have low lexical overlap while conveying nearly the same
information.

Conversely, boilerplate can create high lexical overlap without carrying
the same substantive content.

So we need at least three distinct relations:

\begin{itemize}
\tightlist
\item
  exact duplicate;
\item
  near duplicate under a declared representation;
\item
  semantic redundancy under a declared semantic method.
\end{itemize}

No one relation silently substitutes for the others.

\section{Deduplication changes the training
measure}\label{deduplication-changes-the-training-measure}

Suppose a record appears ten times.

Uniform record sampling gives the underlying content roughly ten times
the probability mass of a singleton.

Remove nine copies and the empirical training distribution changes.

Therefore:

\begin{quote}
deduplication is not merely storage compression.
\end{quote}

It changes the effective sampling measure unless the sampler explicitly
compensates.

This is one reason deduplication can alter model behavior.

\section{Deduplication itself has
policy}\label{deduplication-itself-has-policy}

A near-duplicate relation can induce clusters.

Then someone must decide:

\begin{itemize}
\tightlist
\item
  which representative survives;
\item
  whether source metadata are merged;
\item
  whether domain labels are preserved;
\item
  whether all cluster members remain addressable;
\item
  whether a canonical source is preferred.
\end{itemize}

The clustering threshold and representative policy can change corpus
composition.

\section{Cross-domain duplicates are especially
revealing}\label{cross-domain-duplicates-are-especially-revealing}

Suppose one article appears on the open web and in a curated book
corpus.

If deduplication keeps only the web copy, the text survives but the
book-domain count falls.

If mixture weights are later computed from retained counts, the
deduplication policy has affected the mixture.

This is not a bug in arithmetic.

It is a consequence of how lineage, clustering, and sampling interact.

\section{``Clean'' is not a mathematical
primitive}\label{clean-is-not-a-mathematical-primitive}

Dataset pipelines often use words such as:

\begin{itemize}
\tightlist
\item
  clean;
\item
  filtered;
\item
  high quality;
\item
  safe;
\item
  curated.
\end{itemize}

These words can hide several separate criteria.

Dodge et al.'s analysis of C4 is important because it shows that
web-corpus filtering choices change both content and population
composition, and that large corpora can contain unexpected material and
benchmark examples.

The Atlas therefore treats ``clean'' as a pipeline label, not a theorem.

\section{Quality is
multidimensional}\label{quality-is-multidimensional}

For application \texttt{A}, write a quality feature vector:

\texttt{Q\_A(r)\ =\ (q\_prov,q\_valid,q\_relevance,q\_fresh,q\_label,...)}.

Possible dimensions include:

\begin{itemize}
\tightlist
\item
  provenance completeness;
\item
  parse/format validity;
\item
  domain relevance;
\item
  freshness;
\item
  label reliability;
\item
  duplication burden;
\item
  policy admissibility.
\end{itemize}

A scalar score may be constructed for a particular application.

The chapter assumes no universal weighting.

\section{Filters can trade one quality for
another}\label{filters-can-trade-one-quality-for-another}

A strict filter can remove malformed text.

It can also remove rare dialects, minority language, domain-specific
notation, or valuable edge cases.

A permissive filter can preserve breadth.

It can also preserve noise.

The correct filter depends on the objective and the evidence.

A change in retained distribution should be measured, not inferred from
the word ``quality.''

\section{Raw corpus proportions}\label{raw-corpus-proportions}

First declare the counting unit. In this simple discussion use retained
records. Suppose domain \texttt{d} has \texttt{n\_d} retained records.

Raw corpus proportion is:

\texttt{p\_d\ =\ n\_d\ /\ sum\_j\ n\_j}.

This describes the retained corpus under the counting unit being used.

It does not yet say how often domain \texttt{d} will appear during
training.

\section{Training mixture weights}\label{training-mixture-weights}

A training sampler can declare separate weights:

\texttt{w\_d\textgreater{}=0}

with:

\texttt{sum\_d\ w\_d=1}.

A two-stage sampler may:

\begin{enumerate}
\def\labelenumi{\arabic{enumi}.}
\tightlist
\item
  sample domain \texttt{d} with probability \texttt{w\_d};
\item
  sample a record inside that domain.
\end{enumerate}

Then record probability depends on both the domain weight and the number
of retained records in that domain.

DoReMi makes mixture proportions an explicit optimizable training
variable in its reported experiments.

\section{Corpus composition and training composition
differ}\label{corpus-composition-and-training-composition-differ}

Three distributions should be kept separate:

\begin{itemize}
\tightlist
\item
  raw source/corpus proportions \texttt{p};
\item
  post-filter/dedup proportions \texttt{p\textquotesingle{}};
\item
  training sampling weights \texttt{w}.
\end{itemize}

In general:

\texttt{p\ !=\ p\textquotesingle{}\ !=\ w}.

The difference is often intentional.

The important thing is to document it.

\section{Mixture weights are an
intervention}\label{mixture-weights-are-an-intervention}

Changing \texttt{w} changes the examples a model expects to see.

That makes mixture weighting part of the training intervention.

A downstream performance claim should therefore retain:

\begin{itemize}
\tightlist
\item
  domain definitions;
\item
  weights;
\item
  sampler;
\item
  number of draws/tokens;
\item
  any dynamic weight schedule.
\end{itemize}

A dataset name alone is not enough.

\section{Synthetic provenance}\label{synthetic-provenance}

Add a lineage field:

\texttt{origin(r)\ in\ \{human,synthetic,mixed,unknown\}}.

This says something about where the record came from.

It does not determine whether the record is correct, useful, novel, or
harmful.

Synthetic data should not be collapsed into a scalar quality judgment.

\section{Recursive generated-data failure
mode}\label{recursive-generated-data-failure-mode}

Shumailov et al.~study recursive training on generated data and report
model-collapse behavior in the investigated settings.

The bounded lesson is:

\begin{quote}
recursive generated-data pipelines can create distributional failure
modes that provenance-aware dataset accounting should expose.
\end{quote}

The chapter does not infer:

\texttt{synthetic\ =\textgreater{}\ bad}.

\section{Synthetic data need lineage
depth}\label{synthetic-data-need-lineage-depth}

For generated records, useful provenance may include:

\begin{itemize}
\tightlist
\item
  generator model/version;
\item
  prompt or conditioning source;
\item
  decoding policy;
\item
  upstream human or synthetic inputs;
\item
  generation date;
\item
  filtering;
\item
  whether the generator was trained on earlier synthetic generations.
\end{itemize}

A single Boolean ``synthetic'' can be too weak to diagnose recursive
lineage.

\section{Train/evaluation
contamination}\label{trainevaluation-contamination}

Let \texttt{T} be training data and \texttt{V} evaluation data.

For overlap predicate \texttt{delta}, define:

\texttt{C\_delta(e;T)=1}

when some training record overlaps evaluation item \texttt{e} under
\texttt{delta}.

Unweighted item contamination rate is:

\texttt{CR\_delta(T,V)\ =\ (1/\textbar{}V\textbar{})\ sum\_(e\ in\ V)\ C\_delta(e;T)}.

If evaluation items or tasks carry unequal weights, those weights define
a different measure and must be stated. This is not one universal
number.

It is relative to the predicate and partitions.

\section{Exact overlap is one threat
model}\label{exact-overlap-is-one-threat-model}

An exact-match threat model asks whether the evaluation item itself,
under a declared canonicalization, occurs in training.

This is precise.

It is also narrow.

Paraphrased, templated, or answer-bearing variants can evade it.

\section{Near overlap is a different threat
model}\label{near-overlap-is-a-different-threat-model}

A near-duplicate threat model may detect modified copies.

But it introduces:

\begin{itemize}
\tightlist
\item
  representation choice;
\item
  threshold choice;
\item
  false positives;
\item
  false negatives.
\end{itemize}

So a near-contamination rate should always travel with the detector
definition.

\section{Benchmark examples in web
corpora}\label{benchmark-examples-in-web-corpora}

Dodge et al.~found examples from evaluation datasets inside C4.

Lee et al.~also measured train-test overlap and showed that
deduplication can reduce such overlap in the studied datasets.

These findings justify treating benchmark overlap as an empirical threat
worth measuring.

They do not make every lexical overlap equally damaging.

\section{Exposure is not
memorization}\label{exposure-is-not-memorization}

Suppose an evaluation item overlaps training.

This establishes an opportunity for exposure.

It does not automatically establish that the trained model memorized the
item.

Memorization is a model-behavior claim.

It needs its own evidence.

\section{Memorization is not causal score
inflation}\label{memorization-is-not-causal-score-inflation}

Even demonstrated memorization does not by itself prove that a benchmark
score was inflated by that memorization.

A causal claim would ask something like:

\begin{quote}
what would the model's score have been if the overlapping material had
not been available?
\end{quote}

That counterfactual generally requires a stronger experimental design or
other causal evidence.

\section{Benchmark pathology is broader than exact
leakage}\label{benchmark-pathology-is-broader-than-exact-leakage}

Potential pathologies include:

\begin{itemize}
\tightlist
\item
  exact item reuse;
\item
  near-duplicate item reuse;
\item
  answer leakage;
\item
  template leakage;
\item
  benchmark-derived instruction examples;
\item
  public benchmark discussions repeated in web data;
\item
  tuning on evaluation feedback.
\end{itemize}

Each has a different detection problem.

``Contaminated'' should therefore name the threat model.

\section{Benchmark version matters}\label{benchmark-version-matters}

A benchmark name can hide version changes.

Items may be corrected, expanded, translated, reordered, or reformatted.

A contamination claim should therefore bind to:

\begin{itemize}
\tightlist
\item
  exact benchmark version;
\item
  exact item set;
\item
  canonicalization;
\item
  detector.
\end{itemize}

Otherwise two analyses can report different numbers while appearing to
discuss the same object.

\section{Exact finite witness}\label{exact-finite-witness}

Use five training records:

\begin{itemize}
\tightlist
\item
  t1: web, \{alpha,beta,gamma\};
\item
  t2: web, exact duplicate of t1;
\item
  t3: books, \{alpha,beta,gamma,delta\};
\item
  t4: code, \{lambda,mu,nu\};
\item
  t5: synthetic, \{theta,iota,kappa\}.
\end{itemize}

Use three evaluation records:

\begin{itemize}
\tightlist
\item
  e1: \{alpha,beta,gamma\};
\item
  e2: \{alpha,beta,gamma,epsilon\};
\item
  e3: \{omega,psi,chi\}.
\end{itemize}

The records never change across the following calculations.

\section{Exact-match contamination in the
witness}\label{exact-match-contamination-in-the-witness}

Under exact token-sequence equality:

\begin{itemize}
\tightlist
\item
  e1 overlaps t1 and t2;
\item
  e2 does not exactly match any training record;
\item
  e3 does not match any training record.
\end{itemize}

Therefore:

\texttt{CR\_exact=1/3}.

This is an exact statement under one overlap predicate.

\section{Near-overlap contamination in the
witness}\label{near-overlap-contamination-in-the-witness}

Use token-set Jaccard with threshold:

\texttt{tau=3/4}.

For e2 and t1:

\texttt{J(e2,t1)\ =\ 3/4}.

Therefore e2 now counts as overlapping.

e1 remains overlapping.

e3 remains non-overlapping.

Hence:

\texttt{CR\_near=2/3}.

The data did not change.

The detection rule did.

\section{Why this does not make either rate
false}\label{why-this-does-not-make-either-rate-false}

The rates answer different questions.

\texttt{1/3} means:

\begin{quote}
one third of evaluation items have an exact match under this
canonicalization.
\end{quote}

\texttt{2/3} means:

\begin{quote}
two thirds have a training neighbor at or above this Jaccard threshold.
\end{quote}

Neither statement should be reported without its predicate.

\section{Exact deduplication changes the
corpus}\label{exact-deduplication-changes-the-corpus}

Training contains five records.

t1 and t2 are exact duplicates.

After exact deduplication there are four representatives.

The raw domain histogram changes from:

\texttt{(2,1,1,1)}

for web, books, code, synthetic

to:

\texttt{(1,1,1,1)}.

The corresponding proportions change from:

\texttt{(2/5,1/5,1/5,1/5)}

to:

\texttt{(1/4,1/4,1/4,1/4)}.

\section{Near-duplicate clustering changes it
again}\label{near-duplicate-clustering-changes-it-again}

At threshold \texttt{3/4}:

\texttt{J(t1,t3)=3/4}.

So a graph-clustering policy over pairwise near-duplicate edges places
t1, t2, and t3 in one connected cluster.

Together with t4 and t5, there are three clusters.

If the cluster representative is chosen from the web copy, the
books-domain representative disappears.

This is not a theorem that such a representative policy is good.

It is a demonstration that deduplication policy can alter domain
balance.

\section{Pairwise near-duplication need not be
transitive}\label{pairwise-near-duplication-need-not-be-transitive}

Suppose A is near B and B is near C.

It need not follow that A is near C.

Therefore a dataset pipeline must distinguish:

\begin{itemize}
\tightlist
\item
  pairwise duplicate edges;
\item
  connected-component clustering;
\item
  complete-link clustering;
\item
  representative-based clustering;
\item
  other grouping policies.
\end{itemize}

The cluster definition is part of the data method.

\section{Mixture weights in the
witness}\label{mixture-weights-in-the-witness}

The raw proportions are:

\texttt{p\_raw=(2/5,1/5,1/5,1/5)}.

The exact-deduplicated proportions are:

\texttt{p\_dedup=(1/4,1/4,1/4,1/4)}.

Now declare training weights:

\texttt{w=(1/2,1/4,1/8,1/8)}.

These are different from both corpus distributions.

For eight expected domain draws:

\texttt{8w=(4,2,1,1)}.

The sampling plan is a training decision, not a passive reflection of
file counts.

\section{Why mixture records should be
first-class}\label{why-mixture-records-should-be-first-class}

If a training run reports only:

\begin{quote}
trained on corpus X
\end{quote}

we may still not know:

\begin{itemize}
\tightlist
\item
  domain weights;
\item
  temperature sampling;
\item
  oversampling;
\item
  epoch boundaries;
\item
  token caps;
\item
  curriculum changes;
\item
  dynamic reweighting.
\end{itemize}

A reproducible run should treat mixture configuration as a first-class
artifact.

\section{Data quality and model quality are not
identical}\label{data-quality-and-model-quality-are-not-identical}

A more carefully curated corpus can still produce a worse model under a
different architecture, optimizer, budget, or objective.

Conversely, a model may perform well despite obvious corpus defects.

Therefore:

\begin{quote}
dataset quality is not defined by one downstream score alone.
\end{quote}

The data object and the training system should be evaluated separately
and jointly.

\section{Deduplication precision and
recall}\label{deduplication-precision-and-recall}

A deduplication detector can make two kinds of error.

False merge:

\begin{itemize}
\tightlist
\item
  distinct useful records are treated as duplicates.
\end{itemize}

Missed duplicate:

\begin{itemize}
\tightlist
\item
  redundant records survive.
\end{itemize}

Therefore deduplication itself has precision/recall-like tradeoffs
relative to a declared duplicate gold standard or evaluation protocol.

One threshold cannot be interpreted without this tradeoff.

\section{Contamination detection also has precision and
recall}\label{contamination-detection-also-has-precision-and-recall}

A contamination detector can:

\begin{itemize}
\tightlist
\item
  miss paraphrased benchmark content;
\item
  flag legitimate domain overlap;
\item
  confuse generic templates with copied items.
\end{itemize}

So the detector's own error profile matters.

A high measured contamination rate can reflect a broad detector.

A low rate can reflect a weak detector.

The measurement process itself needs validation.

\section{Source-level versus span-level
overlap}\label{source-level-versus-span-level-overlap}

A source document can overlap a benchmark without containing the exact
benchmark answer.

A small span can exactly reproduce an answer while the source documents
are otherwise unrelated.

Useful overlap units include:

\begin{itemize}
\tightlist
\item
  source;
\item
  document;
\item
  paragraph;
\item
  span;
\item
  item;
\item
  template;
\item
  semantic neighborhood.
\end{itemize}

The unit should match the evaluation threat model.

\section{Time can be part of
contamination}\label{time-can-be-part-of-contamination}

If an evaluation set was published before the training crawl, exposure
may be possible.

A publication date after the final training cutoff does not by itself
prove non-exposure: the same item or its source may have existed
earlier. Ruling out exposure requires evidence that the relevant
item/content was unavailable to the training pipeline before the cutoff.

Therefore useful lineage includes:

\begin{itemize}
\tightlist
\item
  source timestamps;
\item
  crawl dates;
\item
  benchmark release date;
\item
  preprocessing date;
\item
  training cutoff.
\end{itemize}

Temporal reasoning can eliminate some contamination hypotheses before
expensive similarity analysis.

\section{Public benchmarks create a moving
target}\label{public-benchmarks-create-a-moving-target}

Widely discussed benchmarks can appear in:

\begin{itemize}
\tightlist
\item
  tutorials;
\item
  blog posts;
\item
  solution repositories;
\item
  papers;
\item
  forum discussions;
\item
  synthetic instructional data.
\end{itemize}

Over time, the benchmark becomes part of the public data ecosystem.

That does not make evaluation impossible.

It means the validity claim must become more explicit.

\section{Synthetic data can carry contamination
too}\label{synthetic-data-can-carry-contamination-too}

A generator trained on benchmark material can emit benchmark-derived
examples.

A synthetic record therefore has at least two relevant lineage
questions:

\begin{enumerate}
\def\labelenumi{\arabic{enumi}.}
\tightlist
\item
  who or what generated it?
\item
  what data influenced the generator?
\end{enumerate}

Synthetic generation does not erase upstream provenance concerns.

It can make them harder to trace.

\section{Recursive lineage is a
graph}\label{recursive-lineage-is-a-graph}

For synthetic and transformed data, lineage is naturally a directed
graph.

Nodes can include:

\begin{itemize}
\tightlist
\item
  source documents;
\item
  filtering stages;
\item
  model versions;
\item
  prompts;
\item
  generated records;
\item
  downstream datasets.
\end{itemize}

Edges encode derivation.

A flat source label can lose this structure.

The chapter does not require every corpus to maintain a perfect lineage
graph.

It establishes why one may be needed for strong provenance claims.

\section{Data mixtures can evolve during
training}\label{data-mixtures-can-evolve-during-training}

A fixed vector \texttt{w} is the simplest case.

A curriculum or adaptive sampler can use:

\texttt{w\_t}.

Then the effective training distribution changes with training time.

A later Curriculum chapter will study ordering and adaptive selection.

DATA-001 supplies the prerequisite distinction between corpus contents
and sampling policy.

\section{``More data'' is
underspecified}\label{more-data-is-underspecified}

Adding records can change:

\begin{itemize}
\tightlist
\item
  corpus size;
\item
  domain balance;
\item
  duplication;
\item
  freshness;
\item
  synthetic fraction;
\item
  overlap with evaluation;
\item
  number of unique sources;
\item
  effective sampling weights.
\end{itemize}

So a scaling statement such as:

\begin{quote}
performance improved with more data
\end{quote}

should identify what actually changed.

\section{A practical dataset
ledger}\label{a-practical-dataset-ledger}

Before accepting a dataset-dependent empirical claim, record:

{\def\LTcaptype{none} % do not increment counter
\begin{longtable}[]{@{}
  >{\raggedright\arraybackslash}p{(\linewidth - 2\tabcolsep) * \real{0.5000}}
  >{\raggedright\arraybackslash}p{(\linewidth - 2\tabcolsep) * \real{0.5000}}@{}}
\toprule\noalign{}
\begin{minipage}[b]{\linewidth}\raggedright
Field
\end{minipage} & \begin{minipage}[b]{\linewidth}\raggedright
Question
\end{minipage} \\
\midrule\noalign{}
\endhead
\bottomrule\noalign{}
\endlastfoot
records & What is the logical record unit? \\
sources & Which domains and source IDs contribute? \\
provenance & What lineage is retained? \\
processing & Which filters and transformations ran? \\
duplicate rule & Exact, near, semantic? Which threshold? \\
clustering & How are duplicate edges grouped? \\
representative policy & Which record survives a cluster? \\
raw proportions & What is the corpus histogram? \\
post-filter proportions & How did processing alter it? \\
sampling weights & What distribution does training actually draw
from? \\
synthetic lineage & Which records/models are generated? \\
train/eval boundary & Which exact partitions? \\
overlap detector & Which unit, representation, threshold? \\
benchmark version & Which exact item set? \\
timestamps & Which release/crawl/training dates matter? \\
claim boundary & Exposure, memorization, or causal score effect? \\
\end{longtable}
}

This ledger makes ``the data'' inspectable.

\section{Failure modes}\label{failure-modes-13}

\subsection{Exact-equals-near}\label{exact-equals-near}

Exact duplicates are treated as the whole redundancy problem.

\subsection{Near-equals-semantic}\label{near-equals-semantic}

A lexical threshold is treated as semantic equivalence.

\subsection{Dedup-equals-storage}\label{dedup-equals-storage}

Deduplication is described as compression while its sampling effect is
ignored.

\subsection{Corpus-proportion-equals-mixture}\label{corpus-proportion-equals-mixture}

Raw record counts are treated as training probabilities.

\subsection{Synthetic-equals-bad}\label{synthetic-equals-bad}

Generated provenance is turned into a universal quality judgment.

\subsection{Overlap-equals-memorization}\label{overlap-equals-memorization}

Detected training/evaluation overlap is described as model memorization.

\subsection{Memorization-equals-score-inflation}\label{memorization-equals-score-inflation}

Behavioral recall is treated as proof of a causal benchmark effect.

\subsection{Clean-equals-quality}\label{clean-equals-quality}

A pipeline label substitutes for measured quality dimensions.

\section{What the exact witness
establishes}\label{what-the-exact-witness-establishes-6}

The companion witness proves:

\begin{itemize}
\tightlist
\item
  exact contamination rate \texttt{1/3};
\item
  near-overlap contamination rate \texttt{2/3} under token-Jaccard
  threshold \texttt{3/4};
\item
  exact deduplication changes five records to four representatives;
\item
  near-duplicate graph clustering yields three clusters;
\item
  raw domain proportions are \texttt{(2/5,1/5,1/5,1/5)};
\item
  exact-dedup proportions are \texttt{(1/4,1/4,1/4,1/4)};
\item
  declared training weights are \texttt{(1/2,1/4,1/8,1/8)};
\item
  eight expected domain draws under those weights are
  \texttt{(4,2,1,1)}.
\end{itemize}

These facts require no learned model.

They expose the data semantics directly.

\section{Downstream handoff}\label{downstream-handoff-7}

\textbf{Curriculum Learning --- ATLAS-CH-CURRICULUM-001} may now assume:

\begin{itemize}
\tightlist
\item
  corpus lineage;
\item
  exact/near duplicate distinctions;
\item
  overlap-predicate semantics;
\item
  raw versus effective data distributions;
\item
  explicit mixture weights;
\item
  synthetic lineage;
\item
  scoped contamination evidence.
\end{itemize}

The downstream chapter must independently define:

\begin{itemize}
\tightlist
\item
  ordering;
\item
  difficulty;
\item
  competence;
\item
  pacing;
\item
  automatic curriculum policies.
\end{itemize}

DATA-001 defines what is being sampled.

CURRICULUM-001 will define how sampling changes over learning time.

\section{References used in this
chapter}\label{references-used-in-this-chapter-52}

\begin{itemize}
\tightlist
\item
  Lee et al., \emph{Deduplicating Training Data Makes Language Models
  Better}.
\item
  Dodge et al., \emph{Documenting Large Webtext Corpora: A Case Study on
  the Colossal Clean Crawled Corpus}.
\item
  Xie et al., \emph{DoReMi: Optimizing Data Mixtures Speeds Up Language
  Model Pretraining}.
\item
  Shumailov et al., \emph{AI models collapse when trained on recursively
  generated data}.
\end{itemize}

Exact source identities and authority boundaries are recorded in:

sources/source-locks/ATLAS-CH-DATA-001.yaml

\chapter{Curriculum Learning}\label{curriculum-learning}

\textbf{Epistemic status:} established curriculum-learning literature +
audited Data/Optimization prerequisites + Atlas state-aware synthesis.

A curriculum is often summarized as ``show easy examples first, then
harder ones.''

That description captures one historical intuition, but it is too narrow
for a mathematical atlas.

The deeper object is a controller over training experience.

It decides, explicitly or implicitly, what the learner sees next.

Sometimes the controller is fixed before training begins. Sometimes it
reacts to measurements of the learner. Sometimes it changes only the
probability of sampling different regions of the same corpus. Sometimes
it removes examples from the eligible set. Sometimes it uses a nominal
difficulty score. Sometimes it uses a signal called learning progress.

These mechanisms are not equivalent.

The central lesson of this chapter is:

\begin{quote}
\textbf{Curriculum learning is a control problem over the effective
training distribution. A useful curriculum must be defined by its
available experiences, observations, state, actions, and update rule.
Static difficulty order is one special case, not the definition.}
\end{quote}

The source tradition begins with curriculum learning as an organized
presentation of examples (Bengio et al. 2009), extends to automatic
syllabus selection driven by learning-progress signals (Graves et al.
2017), and includes competence-based scheduling that couples example
difficulty to a model-training schedule (Platanios et al. 2019).

The audited Data chapter supplies the data-distribution and evidence
semantics.

The audited First-Order Optimization chapter supplies the
optimizer-state and training-recipe boundary.

\section{The question is not merely ``easy or
hard?''}\label{the-question-is-not-merely-easy-or-hard}

Suppose a dataset contains two kinds of examples.

One kind has a lower declared difficulty score than the other.

A fixed easy-to-hard syllabus can sort them and move through that
ordering.

But what if the learner has already exhausted what the easy region can
teach?

What if the learner is not yet ready for the nominally hard region?

What if two learners at the same training step occupy different states
because of initialization, stochastic gradients, routing, or earlier
examples?

Then ``difficulty'' alone is not enough to specify the useful next
experience.

A curriculum must answer a stronger question:

\begin{quote}
Given the learner information that the controller is allowed to observe,
what training experience distribution should be used next?
\end{quote}

That is a policy question.

\section{Three states, not one}\label{three-states-not-one}

At training step \(t\), keep at least three categories separate.

\subsection{Model state}\label{model-state}

Let

\[
\theta_t
\]

denote model parameters or another declared model state.

\subsection{Optimizer state}\label{optimizer-state}

Let

\[
u_t
\]

denote optimizer state: momentum, adaptive moments, schedules, or other
quantities required by the optimizer.

The audited Optimization chapter already established that optimizer
motion is not determined by the raw gradient alone.

\subsection{Curriculum state}\label{curriculum-state}

Let

\[
c_t
\]

denote controller state.

It can contain:

\begin{itemize}
\tightlist
\item
  syllabus position;
\item
  current competence threshold;
\item
  sampling weights;
\item
  region statistics;
\item
  progress estimates;
\item
  counters;
\item
  eligibility masks;
\item
  controller-bandit state.
\end{itemize}

These coordinates are not model parameters merely because they evolve
during training.

The complete controller need not observe all of \(\theta_t\) or \(u_t\).

Instead declare an observation

\[
z_t
=
\Omega(\theta_t,u_t,c_t,H_t),
\]

where \(H_t\) is whatever training/evaluation history is available to
the controller.

The observation function matters.

A controller that sees only aggregate loss has a different information
contract from one that sees per-domain validation, gradients, probes, or
held-out competence tests.

\section{Curriculum as a policy}\label{curriculum-as-a-policy}

Let \(\mathcal X\) be the available experience space.

A curriculum action \(a_t\) can mean many things:

\begin{itemize}
\tightlist
\item
  choose one example;
\item
  choose a bucket;
\item
  choose a domain;
\item
  choose a task;
\item
  set sampling weights;
\item
  set an eligibility threshold;
\item
  select a replay source;
\item
  choose a stochastic syllabus arm.
\end{itemize}

Write the curriculum policy as

\[
\pi_t(a\mid z_t,c_t).
\]

Then a sampling kernel produces a training experience

\[
x_t\sim Q(\cdot\mid a_t).
\]

The learner updates through

\[
(\theta_{t+1},u_{t+1})
=
\mathcal T_{\rm train}(\theta_t,u_t,x_t).
\]

The controller may update separately. Let \(f_t\) denote the declared
post-update feedback visible to the controller, such as a loss or
evaluation score. Then

\[
c_{t+1}
=
\mathcal T_{\rm curr}(c_t,z_t,a_t,x_t,f_t).
\]

This decomposition matters because a curriculum is not the optimizer.

It changes which training evidence reaches the optimizer.

\section{Open loop and closed loop}\label{open-loop-and-closed-loop}

The simplest curriculum is open loop.

A deterministic syllabus is

\[
a_t=\sigma(t).
\]

It can be sophisticated.

It can increase difficulty gradually, cycle domains, interleave tasks,
or follow a handcrafted sequence.

It is still open loop if the next action does not depend on current
learner observations.

A stochastic open-loop curriculum can use

\[
a_t\sim\pi_t(a).
\]

Randomness does not make a curriculum adaptive.

A closed-loop curriculum uses feedback:

\[
a_t\sim\pi_t(a\mid z_t,c_t).
\]

Now the same training step \(t\) can produce different curriculum
actions for two learners with different observed states.

This distinction is more fundamental than ``easy first.''

\section{Difficulty is a score, not an
essence}\label{difficulty-is-a-score-not-an-essence}

Let

\[
d:\mathcal X\to\mathbb R
\]

be a declared difficulty score.

Examples of possible scoring rules include:

\begin{itemize}
\tightlist
\item
  sequence length;
\item
  rarity;
\item
  number of compositional steps;
\item
  reference-model loss;
\item
  teacher-model confidence;
\item
  structural depth;
\item
  human difficulty labels;
\item
  corruption level.
\end{itemize}

Each score embeds assumptions.

A long example can be easy for one model and hard for another.

A sample with high loss can be informative, mislabeled,
out-of-distribution, adversarial, or merely underfit.

Therefore:

\[
\boxed{
\text{declared difficulty score}
\neq
\text{intrinsic universal difficulty}
}
\]

Bengio et al.~introduced curriculum learning around meaningful
organization of examples and explored a continuation-method
interpretation (Bengio et al. 2009).

The Atlas keeps that historical role while refusing to turn ``easy''
into an observer-independent property.

\section{Competence schedules}\label{competence-schedules}

A competence-based curriculum can couple the difficulty score to a
threshold.

Let

\[
q_t
\]

be a schedule and define the eligible set

\[
\mathcal X_t
=
\{x\in\mathcal X:d(x)\le q_t\}.
\]

As \(q_t\) grows, more examples become eligible.

Platanios et al.~use estimated example difficulty and model competence
in a curriculum framework for neural machine translation (Platanios et
al. 2019).

The Atlas uses that work as a representative mechanism.

It does not infer:

\[
q_t
=
\text{true competence}.
\]

The controller variable is a scheduling construct.

Any stronger claim about actual competence needs its own measurement
contract.

\section{Automated curricula and learning
progress}\label{automated-curricula-and-learning-progress}

Graves et al.~show one way to automate curriculum construction: use
signals derived from learning progress as reward for a nonstationary
bandit that selects a stochastic syllabus (Graves et al. 2017).

This introduces an important architectural shift.

The curriculum is no longer merely an ordered list.

It becomes a controller that receives a reward-like signal and
reallocates training attention.

A generic finite-difference progress signal for region \(r\) must also
state which metric direction counts as improvement. Let
\(\eta_r\in\{+1,-1\}\). Define

\[
LP_t(r)
=
\eta_r\left[m_t(r)-m_{t-w}(r)\right],
\]

where \(m_t(r)\) is a declared metric, \(w\) is a declared lag,
\(\eta_r=+1\) for higher-is-better metrics, and \(\eta_r=-1\) for
lower-is-better metrics.

This signal has a measurement contract.

It depends on:

\begin{itemize}
\tightlist
\item
  the metric;
\item
  the window;
\item
  evaluation noise;
\item
  sampling frequency;
\item
  smoothing;
\item
  region definition;
\item
  whether higher or lower is better.
\end{itemize}

A different window can reverse which region appears to have the highest
current progress.

\section{Learning progress is not a mechanism
readout}\label{learning-progress-is-not-a-mechanism-readout}

Suppose validation accuracy rises quickly on one region.

That is evidence about the declared behavioral metric.

It does not by itself prove:

\begin{itemize}
\tightlist
\item
  which internal feature was acquired;
\item
  whether the capability will persist;
\item
  whether it is accessible under distribution shift;
\item
  whether the model used the intended mechanism;
\item
  whether later training will erase or mask it.
\end{itemize}

For curriculum design, the practical consequence is simple.

A progress metric can be useful as a controller signal without being a
complete scientific explanation.

That distinction lets us use feedback without overclaiming what feedback
means.

\section{Selection is not the same as
reweighting}\label{selection-is-not-the-same-as-reweighting}

The audited Data chapter separates the retained corpus from the
effective training measure.

Curriculum learning operates on that interface.

Suppose the corpus is \(\mathcal D\) with base sampling measure \(\mu\).

A reweighting curriculum can define

\[
\mu_t'(x)
=
\frac{w_t(x)\mu(x)}
{\sum_{x'}w_t(x')\mu(x')}.
\]

If every weight stays positive, the support is unchanged even though
probability mass moves.

A selector can instead choose

\[
S_t\subseteq\mathcal D
\]

and sample only from \(S_t\).

Now support changes.

This distinction matters for:

\begin{itemize}
\tightlist
\item
  exposure;
\item
  contamination accounting;
\item
  rare domains;
\item
  replay;
\item
  fairness of comparisons;
\item
  target-distribution interpretation.
\end{itemize}

``Curriculum changed the data'' is too vague.

We should say how.

\section{A finite state-aware
separation}\label{a-finite-state-aware-separation}

The smallest useful witness needs no neural network.

Let the learner state be

\[
s=(e,h)\in\{0,1,2\}^2.
\]

Interpret the coordinates only as toy learning levels.

For this exact witness, the curriculum controller observes the entire
toy state:

\[
z=s.
\]

This full-observability assumption is deliberately stronger than the
partial measurements available in most real training systems. The
actions \(E\) and \(H\) denote repeatable experience types, so the same
type may be selected on successive updates.

There are two actions:

\[
E=\text{nominally easy},
\qquad
H=\text{nominally hard}.
\]

Static difficulty is

\[
d(E)=1,
\qquad
d(H)=2.
\]

The easy update is

\[
T_E(e,h)
=
(\min(2,e+1),h).
\]

The hard update is

\[
T_H(e,h)
=
\begin{cases}
(e,\min(2,h+1)),&e\ge1,\\
(e,h),&e=0.
\end{cases}
\]

The hard coordinate has one toy prerequisite: \(e\ge1\).

Define

\[
M(e,h)=e+h
\]

and one-step progress

\[
R(s,a)
=
M(T_a(s))-M(s).
\]

Now compare two learner states.

\subsection{Learner A}\label{learner-a}

At

\[
s_A=(0,0),
\]

easy gives

\[
R(s_A,E)=1.
\]

Hard gives

\[
R(s_A,H)=0.
\]

Easy is uniquely progress-maximizing.

\subsection{Learner B}\label{learner-b}

At

\[
s_B=(2,0),
\]

the easy coordinate is saturated.

Therefore

\[
R(s_B,E)=0.
\]

Hard is available and gives

\[
R(s_B,H)=1.
\]

Hard is uniquely progress-maximizing.

The nominal ranking

\[
d(E)<d(H)
\]

is identical for both learners.

The useful next action is not.

Therefore:

\[
\boxed{
\text{static difficulty order}
\not\Rightarrow
\text{state-uniform next-action optimality}
}
\]

in this finite system.

That is the chapter's exact theorem-grade witness.

\section{What the witness does not
say}\label{what-the-witness-does-not-say}

The witness does not say that real training contains an ``easy
coordinate'' and a ``hard coordinate.''

It does not say that hard examples require a simple prerequisite.

It does not say that immediate progress is the correct objective.

It does not say that a neural model at high performance on easy data is
permanently saturated.

It establishes one logical point only:

\begin{quote}
If training utility depends on learner state, a state-independent
ranking can fail to choose the same useful next experience for all
learner states.
\end{quote}

That is enough to justify treating state-aware curriculum control as a
distinct object.

\section{Equal-budget two-step
replay}\label{equal-budget-two-step-replay}

Start from

\[
s_B=(2,0).
\]

Give both curricula exactly two learner transitions.

A fixed easy-then-hard syllabus does:

\[
(2,0)
\xrightarrow{E}
(2,0)
\xrightarrow{H}
(2,1).
\]

Toy progress is

\[
1.
\]

A state-aware immediate-progress policy does:

\[
(2,0)
\xrightarrow{H}
(2,1)
\xrightarrow{H}
(2,2).
\]

Toy progress is

\[
2.
\]

The comparison is intentionally bounded.

Both execute two learner updates.

The witness does not account for controller computation, evaluation
overhead, measurement latency, or partial-state estimation. Its repeated
\(H,H\) action is legal only because the witness explicitly defines
\(H\) as a repeatable experience type.

\section{Why immediate progress is not
enough}\label{why-immediate-progress-is-not-enough}

A tempting rule is:

\begin{quote}
Always train on the region with the largest current learning progress.
\end{quote}

That rule can be useful.

It is not automatically optimal.

An action with small immediate gain may unlock a later region.

A region with zero observed progress may already be mastered and need
occasional rehearsal.

A region with high progress may be noisy and temporarily flattering the
metric.

A region may need exploration before its learning potential can be
estimated.

These long-horizon questions belong partly to the later Learning
Progress as a Search Operator chapter.

CURRICULUM-001 defines the controller interface that chapter will
consume.

It does not pre-solve search.

\section{Curriculum versus training
dose}\label{curriculum-versus-training-dose}

Suppose two experiments end at the same wall-clock time, but one sees
more tokens.

Or they see the same number of examples, but one repeats easy examples
more often.

Or they see the same corpus, but one uses a different learning-rate
schedule.

Then ``curriculum effect'' is not yet isolated.

A credible comparison should declare or control, as appropriate:

\begin{itemize}
\tightlist
\item
  optimizer-update count;
\item
  examples or tokens processed;
\item
  repeated exposure;
\item
  batch size;
\item
  optimizer;
\item
  learning-rate schedule;
\item
  model initialization;
\item
  seeds;
\item
  controller evaluation cost;
\item
  data support;
\item
  sampling weights;
\item
  stopping rule;
\item
  tuning procedure.
\end{itemize}

The required controls depend on the claim.

The principle is inherited from the Atlas evidence discipline:

\begin{quote}
attribute only what the experiment isolates.
\end{quote}

\section{Curriculum and mixture
optimization}\label{curriculum-and-mixture-optimization}

Curriculum learning can look like time-dependent mixture optimization.

The overlap is real.

If a curriculum chooses domain weights \(w_t\), it changes the effective
sampling distribution.

But the concepts are not identical.

Mixture optimization can seek a fixed or slowly varying distribution
without an easy-to-hard semantics.

Curriculum learning can change support, ordering, eligibility, or task
structure.

The useful abstraction is not a label.

It is the explicit policy over experiences.

\section{Stochastic syllabi}\label{stochastic-syllabi}

A curriculum can be stochastic.

For example,

\[
a_t\sim\pi_t(a\mid z_t,c_t).
\]

Stochasticity can support:

\begin{itemize}
\tightlist
\item
  exploration;
\item
  robustness to noisy progress estimates;
\item
  coverage;
\item
  softened eligibility;
\item
  bandit-style allocation.
\end{itemize}

But stochasticity itself proves nothing about exploration quality.

A random controller can be worse than a deterministic one.

The exploration objective and feedback model must be declared.

\section{When loss is a bad difficulty
proxy}\label{when-loss-is-a-bad-difficulty-proxy}

Current loss is attractive because it is already available.

But high loss can mean many things.

An example may be:

\begin{itemize}
\tightlist
\item
  genuinely complex;
\item
  mislabeled;
\item
  out-of-distribution;
\item
  rare;
\item
  adversarial;
\item
  corrupted;
\item
  already memorized but evaluated under noise;
\item
  incompatible with the current objective.
\end{itemize}

Likewise low loss can mean:

\begin{itemize}
\tightlist
\item
  mastered;
\item
  trivial;
\item
  duplicated;
\item
  leaked;
\item
  overrepresented.
\end{itemize}

Therefore a loss-based curriculum must say what interpretation it
assumes.

The Data chapter's lineage and duplication semantics are directly
relevant here.

\section{The fixed lesson plan and the teacher reading the
room}\label{the-fixed-lesson-plan-and-the-teacher-reading-the-room}

The chapter's pedagogical analogy is a teacher with two possible modes.

The fixed lesson plan is open loop.

It says:

\begin{quote}
Lesson 1, then lesson 2, then lesson 3.
\end{quote}

The adaptive teacher reads declared evidence about the class and may
alter the next exercise.

That is closed loop.

The analogy is useful because it exposes the information path.

It fails if we pretend the teacher sees ``understanding'' directly.

A machine-learning controller sees measurements.

The quality of those measurements is part of the system.

\section{Failure modes}\label{failure-modes-14}

\subsection{Easy-to-hard as dogma}\label{easy-to-hard-as-dogma}

A curriculum can be useful without monotone difficulty.

A learner may need rehearsal, interleaving, or revisiting prerequisite
regions.

\subsection{Competence by naming}\label{competence-by-naming}

A variable called competence is still a variable.

The name does not prove that it measures actual competence.

\subsection{Progress as mechanism}\label{progress-as-mechanism}

A rising metric is not a direct scan of internal learning.

\subsection{Greedy progress as global
policy}\label{greedy-progress-as-global-policy}

One-step gain need not maximize long-horizon capability.

\subsection{Hidden dose changes}\label{hidden-dose-changes}

If one curriculum trains longer or sees more useful examples, ordering
may not be the only causal difference.

\subsection{Controller overfitting}\label{controller-overfitting}

A closed-loop curriculum can optimize its own progress metric rather
than the true downstream objective.

\subsection{Support drift}\label{support-drift}

Aggressive selection can stop exposing the learner to important regions.

\subsection{Noise chasing}\label{noise-chasing}

A controller can repeatedly select regions whose progress estimates are
noisy rather than genuinely useful.

\subsection{Endpoint-only
interpretation}\label{endpoint-only-interpretation}

An endpoint improvement does not explain when, how, or why the
capability emerged.

\subsection{Ignoring controller
cost}\label{ignoring-controller-cost}

A curriculum can reduce learner updates while increasing measurement,
scheduling, or data-pipeline cost.

\section{What later chapters may
assume}\label{what-later-chapters-may-assume}

Learning Progress as a Search Operator may assume:

\begin{itemize}
\tightlist
\item
  a curriculum action space exists;
\item
  open-loop and closed-loop controllers are distinct;
\item
  progress signals have explicit measurement contracts;
\item
  data selection changes the effective training measure;
\item
  learner state and curriculum state are separate;
\item
  static difficulty order need not be state-uniformly optimal;
\item
  curriculum feedback is evidence for control, not automatically
  evidence for mechanism.
\end{itemize}

It must independently add:

\begin{itemize}
\tightlist
\item
  experience-space search;
\item
  exploration;
\item
  long-horizon objectives;
\item
  credit assignment;
\item
  uncertainty over progress estimates;
\item
  search stopping rules.
\end{itemize}

Minimal Curricula and Reasoning Bases remains further downstream.

It may not inherit a claim that the smallest curriculum is the one with
the highest short-term progress.

\section{Epistemic status}\label{epistemic-status-1}

Established external literature:

\begin{itemize}
\tightlist
\item
  curriculum learning as organized training-example presentation and
  continuation-style motivation (Bengio et al. 2009);
\item
  automated stochastic syllabus selection using learning-progress
  signals in the reported settings (Graves et al. 2017);
\item
  difficulty/competence-based scheduling for NMT in the reported setting
  (Platanios et al. 2019).
\end{itemize}

Audited Atlas prerequisites:

\begin{itemize}
\tightlist
\item
  data lineage, sampling measure, mixture semantics, contamination
  boundaries;
\item
  optimizer and schedule mechanics.
\end{itemize}

Atlas-owned synthesis:

\begin{itemize}
\tightlist
\item
  curriculum as an explicit control object over the effective training
  distribution;
\item
  the separation among model, optimizer, curriculum state, and
  curriculum observation;
\item
  the exact finite state-dependent next-action witness;
\item
  the claim firewall between progress signals and mechanism claims.
\end{itemize}

Computational witness:

\begin{itemize}
\tightlist
\item
  exact nine-state transition/progress table;
\item
  unique action reversal between \((0,0)\) and \((2,0)\);
\item
  equal-two-transition comparison from \((2,0)\).
\end{itemize}

Not established here:

\begin{itemize}
\tightlist
\item
  universal curriculum gains;
\item
  universal easy-to-hard optimality;
\item
  direct observability of competence;
\item
  long-horizon optimality of greedy progress;
\item
  mechanism identification from behavioral metrics;
\item
  compute savings in deployed training systems.
\end{itemize}

\section{References used in this
chapter}\label{references-used-in-this-chapter-53}

\begin{itemize}
\tightlist
\item
  (Bengio et al. 2009) --- foundational curriculum-learning formulation
  and continuation-method interpretation.
\item
  (Graves et al. 2017) --- automated curriculum/syllabus selection from
  learning-progress signals.
\item
  (Platanios et al. 2019) --- competence-based curriculum scheduling for
  neural machine translation.
\end{itemize}

Exact provenance and authority boundaries are locked in:

sources/source-locks/ATLAS-CH-CURRICULUM-001.yaml

\chapter{Learning Progress as a Search
Operator}\label{learning-progress-as-a-search-operator}

\textbf{Epistemic status:} audited Curriculum prerequisite + primary
learning-progress/search sources + Atlas synthesis + exact finite
witnesses\\
\textbf{Specification:}
manuscript/specifications/ATLAS-CH-PROGRESSSEARCH-001.md\\
\textbf{Derivation packet:}
mathematics/derivations/ATLAS-CH-PROGRESSSEARCH-001-DERIVATIONS.md\\
\textbf{Computational witness:}
mathematics/computational-witnesses/ATLAS-CW-PROGRESSSEARCH-001.md

A curriculum controller asks what the learner should see next.

A search controller asks a stronger question:

\begin{quote}
Which part of experience space should be investigated next, including
regions we have not yet characterized?
\end{quote}

That extra phrase changes the problem.

Learning progress can rank known regions. Search must also decide when
to leave them.

\section{The inherited progress
signal}\label{the-inherited-progress-signal}

The audited Curriculum chapter already defines an oriented progress
signal

\[
LP_t(r)=\eta_r[m_t(r)-m_{t-w}(r)],
\]

with an explicit metric, lag, orientation, noise boundary, and
observation scope.

PROGRESSSEARCH keeps all of those restrictions.

A large positive value means only that the declared metric moved in the
declared favorable direction over the chosen window.

It does not prove a mechanism was acquired.

\section{From curriculum control to
search}\label{from-curriculum-control-to-search}

Curriculum control can operate over a known menu of examples, buckets,
domains, or distributions.

Search becomes explicit when the controller must also decide which
regions deserve measurement, refinement, or generation.

Let the experience space be \(\mathcal E\). It may be finite,
partitioned, or continuously parameterized.

Define the search state

\[
\mathfrak S_t=(\mathcal E,z_t,h_t,\widehat{LP}_t,u_t,g_t,\Phi_t,H,B_t).
\]

The new pieces are not cosmetic:

\begin{itemize}
\tightlist
\item
  \(h_t\): what the search process has already tried;
\item
  \(u_t\): coverage or uncertainty;
\item
  \(g_t\): declared generalization-state evidence;
\item
  \(\Phi_t\): the search objective or ordering rule;
\item
  \(H\): the horizon used for value/credit;
\item
  \(B_t\): remaining budget.
\end{itemize}

The search operator is

\[
r_t\sim\mathcal S_t(\cdot\mid\mathfrak S_t).
\]

\section{Why progress is
attractive}\label{why-progress-is-attractive}

Progress is different from raw performance.

A region with high performance but no remaining improvement may be
saturated.

A region with low performance and no improvement may be presently
unlearnable, badly measured, or simply unexplored.

A region with moderate performance and rapid improvement may be a
productive frontier.

This intuition appears in progress-sensitive intrinsic-motivation
systems and competence-progress goal exploration (Oudeyer et al. 2007)
(Baranes and Oudeyer 2013).

Portelas et al.~make the search interpretation explicit in a
continuously parameterized environment space, using absolute learning
progress inside a surrogate teacher/bandit problem (Portelas et al.
2020).

The Atlas adopts the structural idea without importing a universal
optimality claim.

\section{Search needs exploration}\label{search-needs-exploration}

If the controller only chooses the largest progress estimate among
regions it has already measured, an unseen region can remain unseen
forever.

That is exploitation, not complete search.

Search therefore needs a declared rule for unobserved or uncertain
regions.

Possible rules include:

\begin{itemize}
\tightlist
\item
  forced coverage;
\item
  random exploration;
\item
  optimism/uncertainty bonuses;
\item
  posterior sampling;
\item
  novelty or diversity constraints;
\item
  hierarchical refinement of promising regions.
\end{itemize}

No one rule is privileged by definition.

\section{Exact exploration witness}\label{exact-exploration-witness}

Use two regions A and B.

A has been measured and gives progress 1.

B is unobserved but, when sampled in the toy system, gives progress 3.

An exploit-only policy restricted to observed regions chooses A twice:

\[
1+1=2.
\]

A coverage-first policy samples B once, observes 3, then selects B
again:

\[
3+3=6.
\]

Therefore:

\[
\boxed{
\text{greedy over observed regions}
\not\Rightarrow
\text{discovery of a better unobserved region}
}.
\]

The result is finite and exact.

It does not say that forced coverage is always best.

\section{Generalization-state evidence is not one
scalar}\label{generalization-state-evidence-is-not-one-scalar}

The GSD boundary distinguishes evidence about:

\begin{itemize}
\tightlist
\item
  acquisition;
\item
  persistence;
\item
  accessibility;
\item
  behavioral expression.
\end{itemize}

Write

\[
g_t(r)=\bigl(g_t^{acq},g_t^{pers},g_t^{access},g_t^{expr}\bigr)(r).
\]

These coordinates can disagree.

A capability can appear acquired but fail persistence checks.

It can persist but become inaccessible to one readout.

It can be accessible but not behaviorally expressed under the tested
distribution.

Search may use these observations, but it must not silently collapse
them into a mechanism label.

\section{Heterogeneous evidence needs an ordering
rule}\label{heterogeneous-evidence-needs-an-ordering-rule}

Progress, uncertainty, persistence evidence, cost, and novelty do not
necessarily share units.

A controller may use a scalarization, lexicographic rule, constrained
optimization, or multi-objective policy.

But that rule is part of the algorithm.

Writing

\[
\Phi_t(\widehat{LP},u,g,c,B,H)
\]

does not itself justify adding the arguments numerically.

\section{Immediate progress can be the wrong
horizon}\label{immediate-progress-can-be-the-wrong-horizon}

Learning progress is often measured locally in time.

Search value can be delayed.

Consider a two-step toy system.

Action G gives immediate progress 2 but moves to a state with no
second-step gain.

Action I gives immediate progress 0 but unlocks a second-step action X
worth 5.

Immediate greed chooses G and obtains total return 2.

The path I,X obtains 5.

Thus:

\[
\boxed{
\text{max immediate progress}
\not\Rightarrow
\text{max finite-horizon return}
}.
\]

\section{Credit assignment is separate from
search}\label{credit-assignment-is-separate-from-search}

Suppose an experience chosen now is followed by several updates before a
later validation gain appears.

Which earlier experience gets credit?

A declared H-step return can be written

\[
G_t^{(H)}=\sum_{k=0}^{H-1}\gamma^k r_{t+k}.
\]

But that is only a temporal aggregation rule.

It does not prove causal attribution.

If several data regions, optimizer updates, or interventions occur in
between, the attribution method must be stated separately.

\section{Absolute learning progress needs
care}\label{absolute-learning-progress-needs-care}

Some systems use

\[
|LP_t(r)|
\]

rather than signed progress.

This can be useful when any strong change signals a region worth
revisiting.

But magnitude is not direction.

A large absolute change may reflect improvement, forgetting,
degradation, noise, or nonstationarity.

Therefore:

\[
\boxed{|LP|\neq\text{beneficial learning}.}
\]

\section{Search changes the training
intervention}\label{search-changes-the-training-intervention}

A search controller can change:

\begin{itemize}
\tightlist
\item
  which regions are sampled;
\item
  how often they are sampled;
\item
  which regions are never sampled;
\item
  how much controller computation is spent;
\item
  the order and dependence structure of training experience.
\end{itemize}

That means search can change the effective training distribution.

A fair experiment must still inherit the Curriculum comparison
obligations for training dose, optimizer schedule, exposure, stopping
rule, and controller overhead.

\section{Search can overfit its
metric}\label{search-can-overfit-its-metric}

A controller may repeatedly select regions that look productive only
because the progress estimate is noisy or easy to game.

A search policy can therefore overfit its own feedback channel.

Useful checks include:

\begin{itemize}
\tightlist
\item
  held-out evaluation;
\item
  alternate windows;
\item
  persistence checks;
\item
  delayed reevaluation;
\item
  cross-region validation;
\item
  robustness to progress-estimator noise.
\end{itemize}

These checks reduce ambiguity but do not turn the signal into a
mechanism oracle.

\section{Search does not mean permanent
novelty}\label{search-does-not-mean-permanent-novelty}

Exploration is not synonymous with choosing something new at every step.

Once a productive region is found, exploitation can be rational.

Likewise, revisiting an old region can be rational when forgetting or
transfer changes its value.

The correct balance depends on the declared objective and horizon.

\section{Search over generated
experience}\label{search-over-generated-experience}

The experience space need not be a fixed dataset.

A region can parameterize:

\begin{itemize}
\tightlist
\item
  a simulator;
\item
  a synthetic-data generator;
\item
  a task template;
\item
  a goal distribution;
\item
  a transformation family;
\item
  a problem generator.
\end{itemize}

In that case, the search operator chooses conditions under which new
experience is produced.

This makes the support itself adaptive.

The provenance of generated experience must therefore remain explicit.

\section{What search has not yet
solved}\label{what-search-has-not-yet-solved}

Even a good progress-search operator does not answer:

\begin{itemize}
\tightlist
\item
  what the smallest sufficient curriculum is;
\item
  which learned mechanisms are transferable;
\item
  which reasoning operations form a reconstructive basis;
\item
  whether the same basis works across model scales;
\item
  whether early mechanisms remain latent after later training.
\end{itemize}

Those are downstream questions.

\section{Handoff to Minimal Curricula and Reasoning
Bases}\label{handoff-to-minimal-curricula-and-reasoning-bases}

ATLAS-CH-MINCURR-001 may now assume:

\begin{itemize}
\tightlist
\item
  explicit experience-space search;
\item
  progress estimates with their observation contract;
\item
  an exploration/coverage state;
\item
  delayed-credit and finite-horizon semantics;
\item
  generalization-state evidence as a non-oracular vector;
\item
  exact examples where exploit-only and immediate-greedy choice fail.
\end{itemize}

MINCURR must add a criterion of minimality and reconstructability.

Finding a high-progress region is not the same as finding a minimal
transferable basis.

\section{Durable lesson}\label{durable-lesson-3}

Learning progress can do more than order a syllabus.

It can become feedback for where to look next.

But once progress becomes a search signal, the missing pieces become
unavoidable:

\[
\text{search}
=
\text{progress signal}
+
\text{exploration}
+
\text{credit}
+
\text{horizon}
+
\text{budget}
+
\text{declared evidence semantics}.
\]

The point is not to worship learning progress.

The point is to turn it into a governed search signal whose limits are
visible.

\section{References used in this
chapter}\label{references-used-in-this-chapter-54}

\begin{itemize}
\tightlist
\item
  (Oudeyer et al. 2007)
\item
  (Baranes and Oudeyer 2013)
\item
  (Portelas et al. 2020)
\end{itemize}

Exact source authority, prerequisite identities, and claim boundaries
are locked in:

sources/source-locks/ATLAS-CH-PROGRESSSEARCH-001.yaml

\chapter{Minimal Curricula and Reasoning
Bases}\label{minimal-curricula-and-reasoning-bases}

\textbf{Epistemic status:} audited Progress Search and Residual
substrates + classical teaching/machine-teaching authority + Atlas-owned
exact finite minimality witness.\\
\textbf{Specification:}
manuscript/specifications/ATLAS-CH-MINCURR-001.md\\
\textbf{Derivation packet:}
mathematics/derivations/ATLAS-CH-MINCURR-001-DERIVATIONS.md\\
\textbf{Computational witness:}
mathematics/computational-witnesses/ATLAS-CW-MINCURR-001.md\\
\textbf{Source lock:} sources/source-locks/ATLAS-CH-MINCURR-001.yaml

A curriculum can be useful without being minimal.

An experience can produce rapid learning progress without belonging to
the smallest basis needed for a declared capability.

A later capability can be expressed without proving that an early
mechanism survived.

Those are three different problems.

This chapter separates them.

The central question is:

\begin{quote}
what is the smallest admissible early basis from which the capability we
care about can be reconstructed?
\end{quote}

The answer is always relative to a declaration.

\section{Progress search is not yet
minimality}\label{progress-search-is-not-yet-minimality}

Learning Progress as a Search Operator supplied a controller for where
to look next.

It can use:

\begin{itemize}
\tightlist
\item
  learning-progress estimates;
\item
  exploration;
\item
  coverage state;
\item
  delayed credit;
\item
  finite horizons;
\item
  generalization-state evidence.
\end{itemize}

That controller can find useful experience.

But it does not answer:

\begin{quote}
which experiences are strictly necessary and jointly sufficient for the
target capability?
\end{quote}

That is the new question.

\section{The Residual supplies the right
discipline}\label{the-residual-supplies-the-right-discipline}

The Residual already established that minimality is not meaningful in
the abstract.

One must declare:

\begin{itemize}
\tightlist
\item
  the capability;
\item
  admissible transformations;
\item
  admissible descriptor class;
\item
  reconstruction class.
\end{itemize}

MINCURR applies the same discipline to early experience and reasoning
bases.

A ``minimal curriculum'' means minimal relative to a target and a
reconstruction contract.

\section{The smallest exact
witness}\label{the-smallest-exact-witness}

Let:

\[
Q=\{q_1,q_2\}.
\]

Consider four possible capability hypotheses:

\[
H=\{h_{00},h_{01},h_{10},h_{11}\}.
\]

Each hypothesis is completely described by two outputs.

For example:

\[
h_{10}(q_1)=1,
\qquad
h_{10}(q_2)=0.
\]

Choose target:

\[
h^\star=h_{11}.
\]

The target capability therefore answers both probes with one.

\section{A declared reconstruction
rule}\label{a-declared-reconstruction-rule}

Use target-consistent labeled experiences:

\[
e_1=(q_1,1),
\]

\[
e_2=(q_2,1).
\]

For a finite set \(T\) of labeled experiences, first restrict it to the
target probe family:

\[
T_Q=\{(q,y)\in T:q\in Q\}.
\]

Define the \textbf{target-relative} version space by

\[
V_H(T)=\{h\in H:\forall(q,y)\in T_Q,\ h(q)=y\}.
\]

An auxiliary experience with probe outside \(Q\) imposes no constraint
on this reconstructor. This is a declared target-identification
convention, not a claim that the experience is useless for learning
other capabilities.

The reconstructor succeeds only when:

\[
|V_H(T)|=1.
\]

This is a very simple learner.

That simplicity is useful because the minimality proof can be exact.

\section{Two examples are
sufficient}\label{two-examples-are-sufficient}

Take:

\[
T^\star=\{e_1,e_2\}.
\]

The first example says:

\[
q_1\mapsto1.
\]

That removes:

\[
h_{00},h_{01}.
\]

The second says:

\[
q_2\mapsto1.
\]

Together the two examples leave only:

\[
\boxed{
V_H(T^\star)=\{h_{11}\}.
}
\]

The target is reconstructed exactly.

\section{No strict smaller set
works}\label{no-strict-smaller-set-works}

With no examples:

\[
|V_H(\varnothing)|=4.
\]

With only \(e_1\):

\[
V_H(\{e_1\})
=
\{h_{10},h_{11}\}.
\]

With only \(e_2\):

\[
V_H(\{e_2\})
=
\{h_{01},h_{11}\}.
\]

Neither singleton identifies the target.

Therefore:

\[
\boxed{
|T^\star|=2
}
\]

is the exact minimum for this declared problem.

\section{This is teaching-set
minimality}\label{this-is-teaching-set-minimality}

Goldman and Kearns formalized teaching complexity through the size of
sets that uniquely identify a target concept within a concept class
(Goldman and Kearns 1995).

The Atlas witness is a tiny exact instance of that idea.

The point is not that every curriculum problem is a teaching-dimension
problem.

The point is that ``minimal examples'' becomes mathematically meaningful
only after the target class and learner semantics are fixed.

\section{Machine teaching makes the optimization direction
explicit}\label{machine-teaching-makes-the-optimization-direction-explicit}

Ordinary learning asks what the learner can infer from given data.

Machine teaching reverses the direction:

\begin{quote}
given the learner and target, what training set should a teacher
provide?
\end{quote}

Zhu frames machine teaching as finding an optimal training set for a
specified learning algorithm and target (Zhu 2015).

MINCURR uses that inverse-design perspective without claiming that real
neural learning is identical to the finite witness.

\section{Minimality is conditional}\label{minimality-is-conditional}

Suppose we already knew the target belonged to:

\[
\{h_{10},h_{11}\}.
\]

Then one example:

\[
(q_2,1)
\]

would suffice.

The minimum has changed because the side information changed.

So:

\[
\boxed{
\text{minimality is relative to the declared hypothesis/capability problem}.
}
\]

\section{Capability identity can be weaker than model
identity}\label{capability-identity-can-be-weaker-than-model-identity}

The finite witness uniquely identifies one hypothesis.

Real systems often do not require that.

Suppose two internal systems behave identically on the entire declared
probe family.

Then capability reconstruction can succeed even if internal states
differ.

This is the Residual connection.

A curriculum can be minimal for:

\[
\text{capability equivalence}
\]

without being sufficient to reconstruct the whole model.

\section{Now add a high-progress
distractor}\label{now-add-a-high-progress-distractor}

Introduce an auxiliary experience:

\[
e_3=(z,1),
\]

where:

\[
z\notin Q.
\]

This auxiliary region is outside the declared target capability family.

Freeze progress scores:

\[
p(e_1)=1,
\]

\[
p(e_2)=1,
\]

\[
p(e_3)=5.
\]

So \(e_3\) is the most attractive region under a pure current-progress
score.

\section{But it tells us nothing about the target
class}\label{but-it-tells-us-nothing-about-the-target-class}

Because \(z\) is outside the target probe family:

\[
V_H(\{e_3\})=H.
\]

Therefore:

\[
|V_H(\{e_3\})|=4.
\]

The auxiliary experience leaves target ambiguity unchanged.

By contrast, either \(e_1\) or \(e_2\) cuts the target version space
from four hypotheses to two.

\section{Search value and basis value can
disagree}\label{search-value-and-basis-value-can-disagree}

We therefore have:

\[
p(e_3)>p(e_1),p(e_2),
\]

but:

\[
e_3\notin T^\star.
\]

So:

\[
\boxed{
\text{highest current learning progress}
\not\Rightarrow
\text{membership in a minimal reconstructive basis}.
}
\]

This does not make progress useless.

It says the two objectives are different.

\section{Search asks where to learn
next}\label{search-asks-where-to-learn-next}

A progress-search controller can ask:

\begin{quote}
where is the learner currently changing fastest?
\end{quote}

A minimal-basis analysis asks:

\begin{quote}
what information or operation is necessary to reconstruct the declared
target capability?
\end{quote}

These questions can agree.

They do not have to.

\section{Minimal set is not optimal
sequence}\label{minimal-set-is-not-optimal-sequence}

The witness has one two-example set:

\[
\{e_1,e_2\}.
\]

Under the declared order-insensitive reconstructor, both orders:

\[
(e_1,e_2)
\]

and:

\[
(e_2,e_1)
\]

work.

Thus:

\[
\boxed{
\text{minimal set}
\not\Rightarrow
\text{unique optimal sequence}.
}
\]

Sequence optimization requires learner dynamics.

\section{Why reasoning bases are harder than teaching
sets}\label{why-reasoning-bases-are-harder-than-teaching-sets}

A teaching set is observable.

A reasoning basis is partly mechanistic.

We might want to say that broad capability was built from a small early
collection of:

\begin{itemize}
\tightlist
\item
  operations;
\item
  circuits;
\item
  heuristics;
\item
  primitives;
\item
  subroutines.
\end{itemize}

But that claim requires more than showing that a small dataset was
sufficient.

It requires evidence about what was acquired and what survived.

\section{Four distinct questions}\label{four-distinct-questions}

For a candidate early mechanism \(m\), ask separately:

\begin{enumerate}
\def\labelenumi{\arabic{enumi}.}
\tightlist
\item
  was it acquired?
\item
  did it persist?
\item
  is it still accessible?
\item
  is it behaviorally expressed?
\end{enumerate}

Those are different propositions.

\section{Acquisition}\label{acquisition}

Acquisition means we have evidence that mechanism \(m\) existed or was
used during an early learning interval.

Possible evidence can include:

\begin{itemize}
\tightlist
\item
  a local mechanistic probe;
\item
  a causal intervention;
\item
  a decodable representation;
\item
  a training-local behavior.
\end{itemize}

This is time-local evidence.

It does not tell us what happens later.

\section{Persistence}\label{persistence}

Persistence means the same declared mechanism remains later.

That immediately raises an identity problem:

\begin{quote}
what counts as ``the same'' mechanism?
\end{quote}

Possible identity criteria include:

\begin{itemize}
\tightlist
\item
  matched functional signature;
\item
  aligned mechanistic subspace;
\item
  matched causal intervention effect;
\item
  reconstructive equivalence under a declared map.
\end{itemize}

The identity criterion must be declared.

\section{Acquisition does not prove
persistence}\label{acquisition-does-not-prove-persistence}

A mechanism can be:

\begin{itemize}
\tightlist
\item
  overwritten;
\item
  transformed;
\item
  replaced;
\item
  absorbed into another computation.
\end{itemize}

Therefore:

\[
\boxed{
\text{acquisition}
\not\Rightarrow
\text{persistence}.
}
\]

Longitudinal evidence is required.

\section{Accessibility}\label{accessibility}

A mechanism can persist without being reachable through the current
interface.

It may be inaccessible to:

\begin{itemize}
\tightlist
\item
  the current prompt;
\item
  routing;
\item
  the active policy;
\item
  the current readout;
\item
  the current controller.
\end{itemize}

Therefore:

\[
\boxed{
\text{persistence}
\not\Rightarrow
\text{accessibility}.
}
\]

Accessibility is interface-relative.

\section{Expression}\label{expression}

A mechanism can be accessible but not normally expressed.

Another route may dominate.

The context may suppress it.

A controller may avoid it.

So:

\[
\boxed{
\text{accessibility}
\not\Rightarrow
\text{expression in every context}.
}
\]

\section{Later expression does not prove early
persistence}\label{later-expression-does-not-prove-early-persistence}

Suppose a later model succeeds on the same behavior.

That could happen because:

\begin{itemize}
\tightlist
\item
  the original mechanism persisted;
\item
  it was relearned;
\item
  a different mechanism replaced it;
\item
  a compensating route emerged.
\end{itemize}

Thus:

\[
\boxed{
\text{later behavioral success}
\not\Rightarrow
\text{persistence of a particular early mechanism}.
}
\]

\section{This matters for developmental
narratives}\label{this-matters-for-developmental-narratives}

It is tempting to tell a story:

\begin{quote}
the model learned primitive \(A\) early, primitive \(B\) later, and
broad reasoning emerged by composing them.
\end{quote}

That story may be correct.

But behavior alone does not establish it.

A mechanistic developmental claim needs evidence at each link.

\section{A minimal reasoning basis must specify
composition}\label{a-minimal-reasoning-basis-must-specify-composition}

Suppose candidate operations are:

\[
r_1,\dots,r_k.
\]

A claim that subset:

\[
R^\star
\]

is sufficient requires an admissible composition class:

\[
\mathfrak M.
\]

Otherwise ``reconstructable from the basis'' is underspecified.

This is exactly the lesson inherited from the Residual.

\section{Strict-smaller failure is
essential}\label{strict-smaller-failure-is-essential}

A set can be sufficient without being minimal.

To establish relative minimality, one must attack strict smaller
candidates.

The finite witness does this exhaustively.

In a larger system, exact enumeration may be impossible.

Then the burden shifts to:

\begin{itemize}
\tightlist
\item
  lower bounds;
\item
  adversarial ablations;
\item
  impossibility arguments;
\item
  certified search.
\end{itemize}

\section{High progress can be
auxiliary}\label{high-progress-can-be-auxiliary}

The \(e_3\) control illustrates a practical danger.

An auxiliary task can be:

\begin{itemize}
\tightlist
\item
  rapidly improving;
\item
  informative for representation learning;
\item
  useful for exploration;
\end{itemize}

yet not belong to the minimal basis for the declared target capability.

The correct conclusion is not that \(e_3\) is worthless.

It is that usefulness is objective-relative.

\section{Low progress can still be
necessary}\label{low-progress-can-still-be-necessary}

The converse can also happen.

An experience region can show little current progress because the
learner already handles it well.

Yet it may still encode a constraint that is necessary to identify or
preserve the target capability.

Therefore progress magnitude alone is not a necessity test.

\section{Minimal basis can change with capability
scope}\label{minimal-basis-can-change-with-capability-scope}

A basis sufficient for:

\[
\mathcal Q=\{q_1,q_2\}
\]

may fail when the target probe family expands.

If we add:

\[
q_3,
\]

new examples or operations may become necessary.

Therefore:

\[
\boxed{
\text{minimal for one probe family}
\not\Rightarrow
\text{minimal for a broader capability}.
}
\]

\section{Minimal does not mean
unique}\label{minimal-does-not-mean-unique}

Many systems can admit several distinct minimal bases of equal size.

A curriculum study should distinguish:

\begin{itemize}
\tightlist
\item
  minimum size;
\item
  one minimizer;
\item
  all minimizers;
\item
  equivalence classes of minimizers.
\end{itemize}

The Atlas witness happens to be unique only because its declared target
example language is tiny.

\section{Operational versus semantic
reconstruction}\label{operational-versus-semantic-reconstruction}

A basis may semantically contain enough information to reconstruct a
capability.

That does not mean an allowed learner can efficiently use it.

So MINCURR inherits the RESIDUAL distinction:

\[
\text{semantic sufficiency}
\neq
\text{operational accessibility}.
\]

A claimed basis should say what reconstruction class is allowed.

\section{A general machine-teaching
formulation}\label{a-general-machine-teaching-formulation}

Let:

\begin{itemize}
\tightlist
\item
  \(L\): declared learner;
\item
  \(\theta^\star\): target;
\item
  \(\mathcal E\): admissible experiences;
\item
  \(c(T)\): curriculum cost.
\end{itemize}

Then a teaching problem can be written:

\[
\min_{T\subseteq\mathcal E} c(T)
\]

subject to:

\[
L(T)=\theta^\star.
\]

For capability-relative reconstruction, replace exact parameter identity
by:

\[
B(L(T),q)=B^\star(q)
\quad
\forall q\in\mathcal Q.
\]

The objective must say which notion is intended.

\section{Early mechanisms require a longitudinal
protocol}\label{early-mechanisms-require-a-longitudinal-protocol}

To claim an early reasoning basis persists, a study should freeze:

\begin{itemize}
\tightlist
\item
  early acquisition criterion;
\item
  mechanism identity criterion;
\item
  later persistence probe;
\item
  accessibility interface;
\item
  behavioral probe family;
\item
  causal test where mechanism claims are made.
\end{itemize}

Then it should evaluate each separately.

\section{A useful empirical disposition
grammar}\label{a-useful-empirical-disposition-grammar}

A future MINCURR experiment can distinguish:

\subsection{MINIMAL\_RECONSTRUCTIVE\_BASIS}\label{minimal_reconstructive_basis}

A candidate basis is sufficient and every admissible strict smaller
candidate is ruled out.

\subsection{SUFFICIENT\_NOT\_MINIMAL}\label{sufficient_not_minimal}

The candidate reconstructs the target but minimality has not been shown
or a smaller basis exists.

\subsection{SEARCH\_RELEVANT\_NOT\_BASIS}\label{search_relevant_not_basis}

An experience region is useful/high-progress but not required for the
declared reconstruction target.

\subsection{PERSISTENT\_INACCESSIBLE}\label{persistent_inaccessible}

A mechanism satisfies the declared persistence criterion but fails the
current accessibility criterion.

\subsection{BEHAVIOR\_ONLY}\label{behavior_only}

Later behavior is present without evidence identifying a persistent
early mechanism.

This grammar prevents several common promotions.

\section{The exact witness in one
line}\label{the-exact-witness-in-one-line}

For the finite target:

\[
h^\star=(1,1),
\]

we have:

\[
V_H(\{e_1,e_2\})=\{h^\star\},
\]

but:

\[
|V_H(T)|>1
\]

for every strict smaller:

\[
T\subsetneq\{e_1,e_2\}.
\]

Meanwhile:

\[
p(e_3)=5
\]

is maximal while:

\[
V_H(\{e_3\})=H.
\]

That is the chapter's core separation.

\section{Non-implications}\label{non-implications-4}

MINCURR rejects:

\[
\text{highest progress}
\not\Rightarrow
\text{minimal-basis membership},
\]

\[
\text{minimal for one target}
\not\Rightarrow
\text{universal minimality},
\]

\[
\text{minimal set}
\not\Rightarrow
\text{optimal sequence},
\]

\[
\text{acquisition}
\not\Rightarrow
\text{persistence},
\]

\[
\text{persistence}
\not\Rightarrow
\text{accessibility},
\]

\[
\text{later behavior}
\not\Rightarrow
\text{early-mechanism persistence},
\]

and:

\[
\text{behavioral equivalence}
\not\Rightarrow
\text{mechanistic identity}.
\]

\section{Atlas connections}\label{atlas-connections-21}

\textbf{Learning Progress as a Search Operator.}\\
Progress can guide experience search without solving minimality.

\textbf{The Residual.}\\
Task-relative sufficiency and declared reconstruction classes provide
the minimality language.

\textbf{Mechanistic diagnostics.}\\
Mechanism persistence needs longitudinal identity and causal evidence.

\textbf{Generalization-state dynamics.}\\
Later behavioral state is evidence about behavior, not automatic
evidence about which early mechanism survived.

\textbf{Compression and intelligence probes.}\\
A minimal reasoning basis is related to compact reconstructive structure
but is not identical to compression ratio.

\section{Closing view}\label{closing-view-30}

The smallest useful curriculum is not simply the region where learning
is happening fastest.

And the smallest early basis is not automatically the mechanism that
later behavior uses.

Minimality is a reconstruction claim.

Persistence is a longitudinal claim.

Accessibility is an interface claim.

Behavior is an expression claim.

The exact finite witness makes the first distinction formal:

\[
\boxed{
T^\star=\{(q_1,1),(q_2,1)\}
}
\]

is minimal for the declared target, while the highest-progress auxiliary
experience contributes nothing to target identification.

The broader research programme begins when the same discipline is
applied to learned operations:

\begin{quote}
identify the smallest reconstructive basis, then separately test whether
it was acquired, whether it persisted, whether it remains accessible,
and whether later capability actually uses it.
\end{quote}

\section{References used in this
chapter}\label{references-used-in-this-chapter-55}

\begin{itemize}
\tightlist
\item
  (Goldman and Kearns 1995)
\item
  (Zhu 2015)
\end{itemize}

Exact source identities and claim boundaries are recorded in:

sources/source-locks/ATLAS-CH-MINCURR-001.yaml

\chapter{Reinforcement Learning and
Control}\label{reinforcement-learning-and-control}

\textbf{Epistemic status:} established reinforcement-learning/control
foundations plus Atlas synthesis and exact finite derivation.\\
\textbf{Specification:}
manuscript/specifications/ATLAS-CH-RLBASE-001.md\\
\textbf{Formal packet:}
mathematics/derivations/ATLAS-CH-RLBASE-001-DERIVATIONS.md\\
\textbf{Computational witness:}
mathematics/computational-witnesses/ATLAS-CW-RLBASE-001.md\\
\textbf{Source lock:} sources/source-locks/ATLAS-CH-RLBASE-001.yaml

A dynamical system tells us how state changes.

A control problem adds a question:

which change should we choose?

Reinforcement learning enters when the consequences of those choices are
uncertain, delayed, only partly modeled, or learned from interaction.

The central object is therefore not ``reward'' by itself.

It is a coupled system containing state, action, dynamics, reward,
policy, and future consequences.

The discipline of reinforcement learning begins by keeping those objects
separate.

\section{A controlled stochastic dynamical
system}\label{a-controlled-stochastic-dynamical-system}

For this chapter, write a discounted Markov decision process as

M = (S, A, P, r, gamma, rho\_0).

Here S is the state space.

A(s) is the set of admissible actions at state s.

P(s'\textbar s,a) is the transition kernel.

r(s,a,s') is the one-step reward associated with the transition.

gamma in {[}0,1) is the discount factor.

rho\_0 is the initial-state distribution.

A policy is another object:

pi(a\textbar s).

The policy says how actions are selected.

The transition kernel says what the environment does after an action.

These should not be fused.

A poor policy can act in a perfectly specified environment.

A good policy can fail in a misspecified environment.

And a learned policy may change while the environment dynamics remain
fixed.

\section{The Markov boundary}\label{the-markov-boundary}

The word Markov does real work.

The MDP state is assumed to contain enough information that, conditional
on the current state and action, the next-state distribution does not
depend on the earlier history.

Informally:

the present state is sufficient for predicting the controlled future.

This is a modeling claim.

It can fail.

If two histories are mapped to the same state representation but require
different predictions or decisions, the chosen state is not Markov for
the purpose at hand.

That failure cannot be repaired merely by applying a more sophisticated
Bellman update.

Sometimes the right repair is a richer state.

Sometimes it is memory.

Sometimes the problem is genuinely partially observable.

\section{Reward is not value}\label{reward-is-not-value}

Reward is local.

Value is prospective.

Assume throughout the finite-MDP discussion that one-step rewards are
finite; because S and A are finite, the reward table is then bounded.

Let

\begin{Shaded}
\begin{Highlighting}[]
\NormalTok{G\_t}
\NormalTok{=}
\NormalTok{sum\_\{k=0\}\^{}\textbackslash{}infty gamma\^{}k R\_\{t+k+1\}}
\end{Highlighting}
\end{Shaded}

be the discounted return. With gamma\textless1 and bounded one-step
reward, this series is absolutely bounded.

For a fixed policy pi, define the state value

\begin{Shaded}
\begin{Highlighting}[]
\NormalTok{V\^{}pi(s)}
\NormalTok{=}
\NormalTok{E\_pi[G\_t | S\_t=s].}
\end{Highlighting}
\end{Shaded}

Define the action value

\begin{Shaded}
\begin{Highlighting}[]
\NormalTok{Q\^{}pi(s,a)}
\NormalTok{=}
\NormalTok{E\_pi[G\_t | S\_t=s,A\_t=a].}
\end{Highlighting}
\end{Shaded}

These quantities summarize expected future reward under declared
conditions.

They are not the same thing as uncertainty.

A value estimate of 10 does not tell us how uncertain that estimate is.

Nor does reward automatically mean ``good'' in a broader human sense.

Reward is part of the formal problem specification.

If the reward function encodes the wrong objective, an optimizer can
solve the mathematical problem and still produce an outcome the designer
did not want.

\section{Bellman's recursion}\label{bellmans-recursion}

Sequential decision problems become tractable when the future can be
folded back into the present.

Condition on the first transition.

Under policy pi,

\begin{Shaded}
\begin{Highlighting}[]
\NormalTok{V\^{}pi(s)}
\NormalTok{=}
\NormalTok{sum\_a pi(a|s)}
\NormalTok{sum\_\{s\textquotesingle{}\} P(s\textquotesingle{}|s,a)}
\NormalTok{[}
\NormalTok{r(s,a,s\textquotesingle{}) + gamma V\^{}pi(s\textquotesingle{})}
\NormalTok{].}
\end{Highlighting}
\end{Shaded}

This is the Bellman expectation equation.

It says that value equals expected immediate reward plus discounted
expected continuation value.

The equation is exact for the declared MDP and policy.

The approximation begins later, when we estimate the quantities with
samples, restricted representations, learned models, or finite
computation.

That distinction matters.

An approximate neural critic does not make the Bellman identity
approximate.

It makes our representation or solution of the identity approximate.

\section{A two-state control
problem}\label{a-two-state-control-problem}

Consider the exact witness used by this chapter.

There are two states.

At s1, the only action gives reward 2 and returns to s1.

At s0, action a gives reward 0 and moves to s1.

Action b gives reward 2 and returns to s0.

Let

gamma = 1/2.

Start with a policy that chooses a at s0.

Then

\begin{Shaded}
\begin{Highlighting}[]
\NormalTok{V\^{}pi(s1)}
\NormalTok{=}
\NormalTok{2 + (1/2)V\^{}pi(s1),}
\end{Highlighting}
\end{Shaded}

so

V\^{}pi(s1)=4.

At s0,

\begin{Shaded}
\begin{Highlighting}[]
\NormalTok{V\^{}pi(s0)}
\NormalTok{=}
\NormalTok{0 + (1/2)4}
\NormalTok{=}
\NormalTok{2.}
\end{Highlighting}
\end{Shaded}

Now ask what would happen if, only at s0, we took action b once and then
followed the old policy.

Its action value is

\begin{Shaded}
\begin{Highlighting}[]
\NormalTok{Q\^{}pi(s0,b)}
\NormalTok{=}
\NormalTok{2 + (1/2)V\^{}pi(s0)}
\NormalTok{=}
\NormalTok{3.}
\end{Highlighting}
\end{Shaded}

Three is greater than two.

So the old policy contains a locally improvable decision.

Switch s0 to action b.

The environment has not changed.

The reward function has not changed.

The transition law has not changed.

Only the policy has changed.

Under the new policy,

\begin{Shaded}
\begin{Highlighting}[]
\NormalTok{V(s0)}
\NormalTok{=}
\NormalTok{2 + (1/2)V(s0)}
\NormalTok{=}
\NormalTok{4.}
\end{Highlighting}
\end{Shaded}

The witness therefore gives the cleanest possible policy-improvement
story:

same world, different controller, larger value.

\section{Policy evaluation and policy
improvement}\label{policy-evaluation-and-policy-improvement}

This pattern generalizes.

Policy evaluation asks:

What is V\^{}pi for the current policy?

Policy improvement asks:

Can we choose actions greedily with respect to Q\^{}pi and obtain a
policy no worse than pi?

In the exact finite discounted setting, the policy-improvement theorem
answers yes.

Alternating exact evaluation and greedy improvement gives policy
iteration.

Applying the optimal Bellman operator directly gives value iteration.

The important point is structural.

Evaluation and improvement are different operations.

One estimates the consequences of a controller.

The other changes the controller.

Blurring them makes it harder to identify where an error entered.

\section{Optimality is relative}\label{optimality-is-relative}

Define the Bellman optimality operator

\begin{Shaded}
\begin{Highlighting}[]
\NormalTok{(TV)(s)}
\NormalTok{=}
\NormalTok{max\_a}
\NormalTok{sum\_\{s\textquotesingle{}\} P(s\textquotesingle{}|s,a)}
\NormalTok{[}
\NormalTok{r(s,a,s\textquotesingle{}) + gamma V(s\textquotesingle{})}
\NormalTok{].}
\end{Highlighting}
\end{Shaded}

The optimal value V* satisfies

V* = T V*.

An optimal policy chooses maximizing actions.

But the star has a scope.

Optimal with respect to which state representation?

Which admissible actions?

Which transition model?

Which reward?

Which discount?

Which horizon?

Which observation model?

An optimal policy for the wrong MDP is still optimal for the wrong MDP.

This is why later chapters will add uncertainty, optionality, correction
capacity, and information value rather than treating a single scalar
Bellman optimum as the entire decision problem.

\section{When the model is known}\label{when-the-model-is-known}

If P and r are known, we can reason with the model directly.

Dynamic programming applies Bellman operators using expected
transitions.

Policy evaluation repeatedly applies the policy Bellman operator.

Value iteration repeatedly applies the optimal Bellman operator.

Policy iteration alternates evaluation and improvement.

This is planning in a known controlled stochastic system.

No interaction sample is needed to define the recursion.

The computational burden may still be large.

But the uncertainty about transition law is conceptually separate.

\section{Learning from sampled
transitions}\label{learning-from-sampled-transitions}

Reinforcement learning often begins when the model is not used directly.

Instead we observe transitions:

(S\_t,A\_t,R\_\{t+1\},S\_\{t+1\}).

A temporal-difference update for a state-value estimate has the form

V(S\_t) \textless- V(S\_t) + alpha {[} R\_\{t+1\} + gamma V(S\_\{t+1\})
- V(S\_t){]}.

The bracketed term compares the old estimate with a one-step
bootstrapped target.

This is where dynamics, information, and stochastic approximation meet.

The Bellman equation supplies the fixed-point structure.

The sampled transition supplies local evidence.

The learning rule decides how quickly to move.

These are three distinct objects.

\section{Q-learning and off-policy
control}\label{q-learning-and-off-policy-control}

Q-learning replaces the next-state policy average with a maximum:

Q(S\_t,A\_t) \textless- Q(S\_t,A\_t) + alpha\_t {[} R\_\{t+1\} + gamma
max\_a Q(S\_\{t+1\},a) - Q(S\_t,A\_t){]}.

This lets the behavior that generated data differ from the greedy policy
defined by the current Q values.

That separation is powerful.

It also creates a distribution question.

Which state-action pairs are actually sampled?

The classical Watkins-Dayan convergence result retains strong
conditions.

It does not say that arbitrary nonlinear approximation, replay,
bootstrapping, nonstationarity, and off-policy data will converge merely
because the update still resembles Q-learning.

The equation can survive while the theorem does not.

\section{Policy gradients}\label{policy-gradients}

Bellman-style control is not the only route.

Suppose a stochastic policy pi\_theta is differentiable in parameters
theta.

Define an expected-return objective J(theta).

Policy-gradient methods estimate

grad\_theta J(theta)

and update the policy parameters directly.

For a sampled trajectory, REINFORCE uses the score-function term
grad\_theta log pi\_theta(A\_t\textbar S\_t) multiplied by an
appropriate sampled return. Under the stated sampling assumptions, its
expectation recovers the policy-objective gradient.

Williams's REINFORCE family is a foundational example.

This viewpoint is useful because the object being optimized is the
policy distribution itself.

Modern actor-critic methods combine policy-gradient and value-estimation
roles.

Again, the distinction matters.

The actor changes the policy.

The critic estimates a value-related quantity.

Neither label guarantees that the learned approximation is correct.

\section{Model-based reinforcement
learning}\label{model-based-reinforcement-learning}

A model can itself be learned.

Then the system may use real experience to improve an estimated
transition/reward model and use that model to generate planning updates.

Dyna is a representative architecture for this integration.

The same experience can support at least three activities:

direct value or policy learning;

model learning;

planning through the learned model.

This can improve data efficiency because one real transition can
influence many simulated updates.

It also creates model bias.

Planning is only as faithful as the learned model in the states and
actions where planning relies on it.

A planner can be internally exact and externally wrong because the model
is wrong.

\section{Partial observability}\label{partial-observability}

The MDP assumes the controller has access to the state used by the
transition model.

Many real systems do not.

In a partially observable Markov decision process, hidden state S\_t
generates observations rather than being directly exposed.

The controller acts on history or on an information state derived from
history.

A standard choice is the belief state

\begin{Shaded}
\begin{Highlighting}[]
\NormalTok{b\_t(s)}
\NormalTok{=}
\NormalTok{P(S\_t=s | history\_t).}
\end{Highlighting}
\end{Shaded}

The belief is a probability distribution over hidden states.

Under the standard POMDP formulation, with the transition/observation
model fixed and the posterior updated from the complete available
history, the exact belief state can recover a Markov control problem in
information space.

That is an elegant transformation.

It is not free.

Exact beliefs may be expensive.

The observation model may be wrong.

Approximate memory may discard information.

And two histories that induce the same approximate representation may
still require different decisions.

\section{What value leaves out}\label{what-value-leaves-out}

Value is useful precisely because it compresses a future distribution
into an expectation.

Compression loses distinctions.

Two actions can have the same expected return and very different
downside tails.

Two policies can have similar value and different information gain.

One policy can preserve many future corrective actions while another
reaches the same expected return through an irreversible commitment.

A value function does not automatically record these differences.

That is not a defect in the definition.

It is a boundary.

The later decision chapters exist because the Atlas does not want one
scalar expectation to silently absorb every notion of good decision
making.

\section{Model-free is not
assumption-free}\label{model-free-is-not-assumption-free}

The phrase model-free can sound stronger than it is.

A model-free algorithm may avoid explicitly representing P.

It still relies on assumptions about the data-generating process, reward
signal, state representation, stationarity, coverage, function class,
optimizer, or learning-rate schedule.

The environment model can be absent from the learned representation
while remaining present in the mathematics of why an estimator should
work.

This is another recurring Atlas lesson:

not representing an object explicitly does not make the object
irrelevant.

\section{Control and uncertainty are different
layers}\label{control-and-uncertainty-are-different-layers}

A value estimate can be wrong.

A transition model can be wrong.

A reward model can be wrong.

An observation model can be wrong.

Future value estimates can be wrong.

These uncertainties interact.

RLBASE-001 does not solve that joint propagation problem.

It establishes the controlled stochastic substrate on which the later
uncertainty chapter can operate.

This is why ATLAS-CH-JOINTUNC-001 depends on both this chapter and the
Information substrate.

The goal there will not be to add several independent error bars
mechanically.

It will be to reason about uncertainty that propagates through one
coupled decision system.

\section{What exploration adds}\label{what-exploration-adds}

The Bellman optimality equation tells us what would be optimal if the
relevant model or values were known.

Learning has another problem.

Which actions should be tried in order to learn what is not yet known?

An action can have two kinds of consequence:

it can produce reward;

it can produce information.

Those are not the same object.

A low-reward action now may improve future decisions by revealing
structure.

A high-reward familiar action may teach us almost nothing.

ATLAS-CH-EXPLORE-001 begins from that gap.

It may assume everything established here:

states, actions, transitions, rewards, policies, returns, values,
Bellman operators, and model-based/model-free distinctions.

Its job is to add the value of information and the
exploration-exploitation problem.

\section{The boundary to remember}\label{the-boundary-to-remember-2}

Reinforcement learning is often summarized as learning to maximize
reward.

That slogan is too compressed for the Atlas.

A more useful statement is:

reinforcement learning studies how a controller can improve sequential
decisions from interaction when consequences are distributed through
time and the relevant dynamics or values are not simply handed to it.

The objects remain distinct.

Dynamics describe what follows an action.

Reward scores a local transition according to the model.

Return accumulates those scores through time.

Value takes an expectation under a policy.

Bellman recursion connects present and future.

Learning estimates or improves the relevant objects from data.

Planning uses a model to reason ahead.

Partial observability moves control into information state.

And optimality remains relative to the problem we chose to write down.

Once those boundaries are explicit, later chapters can ask harder
questions without overloading the word value.

\section{References used in this
chapter}\label{references-used-in-this-chapter-56}

The exact source identities and authority scopes are pinned in:

sources/source-locks/ATLAS-CH-RLBASE-001.yaml

The external basis includes Bellman's Markovian decision-process work,
Sutton and Barto's standard reinforcement-learning treatment, Watkins
and Dayan on Q-learning, Williams on REINFORCE, Sutton on Dyna, and
Kaelbling, Littman, and Cassandra on POMDPs.

\chapter{Exploration and Information
Value}\label{exploration-and-information-value}

\textbf{Epistemic status:} established bandit/exploration foundations
plus Atlas synthesis and one exact finite Bayesian decision witness.\\
\textbf{Specification:}
manuscript/specifications/ATLAS-CH-EXPLORE-001.md\\
\textbf{Formal packet:}
mathematics/derivations/ATLAS-CH-EXPLORE-001-DERIVATIONS.md\\
\textbf{Computational witness:}
mathematics/computational-witnesses/ATLAS-CW-EXPLORE-001.md\\
\textbf{Source lock:} sources/source-locks/ATLAS-CH-EXPLORE-001.yaml

A rational action can look worse now and still be better overall.

That is the core difficulty of exploration.

If every consequence of every action were already known, a controller
could simply choose the action with the highest value under its declared
objective.

Exploration appears when actions do two things at once:

\begin{enumerate}
\def\labelenumi{\arabic{enumi}.}
\tightlist
\item
  they change the world and produce reward;
\item
  they change what the decision-maker knows.
\end{enumerate}

The second effect can alter future choices.

That makes information part of the control problem.

A random action is not automatically exploratory.

An uncertain action is not automatically informative.

An informative action is not automatically useful.

And useful information is not automatically worth its acquisition cost.

This chapter develops the mathematical distinctions required to say
exactly when exploration has decision value.

\section{From control to learning while
acting}\label{from-control-to-learning-while-acting}

The Reinforcement Learning and Control chapter introduced a controlled
stochastic system

\texttt{M=(S,A,P,r,gamma,rho\_0)}

and separated environment dynamics, policy, reward, return, value,
model, uncertainty, and observability.

That substrate answers a question of the form:

\begin{quote}
given a declared control problem, how should future consequences enter
present decisions?
\end{quote}

Exploration adds another question:

\begin{quote}
what should we do when some decision-relevant part of the control
problem is not yet known, and our actions determine what evidence we
will receive?
\end{quote}

The difference is subtle.

In ordinary planning with a known model, an action matters because of
its reward and state transition.

Under epistemic uncertainty, an action can also matter because of the
observation it produces.

An exploratory action is therefore a \textbf{dual-purpose action}:

\begin{itemize}
\tightlist
\item
  it acts;
\item
  it probes.
\end{itemize}

This dual role is what creates the exploration--exploitation tension.

\section{A minimal Bayesian exploration
object}\label{a-minimal-bayesian-exploration-object}

Let \texttt{theta} denote an unknown latent parameter.

It might encode:

\begin{itemize}
\tightlist
\item
  an arm mean;
\item
  a transition probability;
\item
  a reward parameter;
\item
  a hidden environment mode;
\item
  the identity of a better controller;
\item
  another bounded source of model uncertainty.
\end{itemize}

At time \texttt{t}, let

\texttt{b\_t(theta)}

be the current belief over \texttt{theta}.

For action \texttt{a}, let

\texttt{p(y\textbar{}theta,a)}

be the observation law.

After observing \texttt{y}, the belief updates by Bayes' rule:

\texttt{b\_\{t+1\}(theta)\ proportional\ to\ b\_t(theta)p(y\textbar{}theta,a)}.

The important feature is not the formula itself.

It is the direction of dependence:

\texttt{action\ -\textgreater{}\ observation\ -\textgreater{}\ belief\ -\textgreater{}\ future\ action}.

This is the exploration channel.

If an action changes reward but not the future information state, its
value is purely instrumental.

If it changes the information state, the continuation value may differ
even when its immediate reward is lower.

\section{Bayes value makes exploration
explicit}\label{bayes-value-makes-exploration-explicit}

For a finite horizon, write the value of belief state \texttt{b} at time
\texttt{t} as

\texttt{V\_t(b)}.

Let \texttt{C\_t(b)} be the set of actions currently admissible under
any declared constraints.

Then the Bayes recursion is

\texttt{V\_t(b)\ =\ max\_\{a\ in\ C\_t(b)\}\ E{[}\ r(theta,a,Y)\ +\ V\_\{t+1\}(B(b,a,Y))\ {]}}.

The first term is immediate reward.

The second term is continuation value after the observation has changed
the belief.

This is where exploration lives mathematically.

An action can lose on the first term and win on the second.

A purely myopic controller ignores that possibility.

A finite-horizon exploratory controller does not.

\section{The smallest useful
example}\label{the-smallest-useful-example}

Consider a two-step decision.

There are two actions.

A known action \texttt{S} always pays

\texttt{1/2}.

An unknown action \texttt{U} pays

\texttt{theta},

where

\texttt{theta\ in\ \{0,1\}}.

The prior is

\texttt{P(theta=1)=2/5}, \texttt{P(theta=0)=3/5}.

The immediate expected reward of \texttt{U} is

\texttt{2/5}.

So at the first step,

\texttt{2/5\ \textless{}\ 1/2}.

If only the current reward mattered, \texttt{S} would be the rational
choice.

But pulling \texttt{U} reveals \texttt{theta} exactly.

That changes the second decision.

If the first action is \texttt{S}, no information is gained.

The second action is again \texttt{S}.

Total expected reward:

\texttt{1/2+1/2=1}.

If the first action is \texttt{U}, the first expected reward is
\texttt{2/5}.

Then:

\begin{itemize}
\tightlist
\item
  with probability \texttt{2/5}, we learn \texttt{theta=1} and choose
  \texttt{U} next for reward 1;
\item
  with probability \texttt{3/5}, we learn \texttt{theta=0} and choose
  \texttt{S} next for reward \texttt{1/2}.
\end{itemize}

So the expected second reward is

\texttt{(2/5)(1)+(3/5)(1/2)=7/10}.

The total becomes

\texttt{2/5+7/10=11/10}.

Exploration wins:

\texttt{11/10\ \textgreater{}\ 1}.

The lower immediate-reward action is the better two-step action.

That is the entire exploration problem in miniature.

\section{Exploration cost and future information
value}\label{exploration-cost-and-future-information-value}

The example can be decomposed exactly.

The immediate cost of exploring rather than choosing the known action is

\texttt{1/2-2/5=1/10}.

Without new information, the best second-step value is

\texttt{1/2}.

With the observation from \texttt{U}, expected second-step value becomes

\texttt{7/10}.

So the future value of the information is

\texttt{7/10-1/2=1/5}.

The net advantage is therefore

\texttt{1/5-1/10=1/10}.

This decomposition is more useful than the slogan ``exploration can pay
off.''

It exposes the accounting:

\texttt{net\ exploration\ value\ =\ future\ decision\ value\ of\ information\ -\ immediate\ exploration\ cost}

in this particular two-step construction.

A different horizon, observation model, reward model, or constraint set
can reverse the result.

\section{Information gain is not value of
information}\label{information-gain-is-not-value-of-information}

Pulling \texttt{U} reveals \texttt{theta} exactly.

Before the pull, the entropy of \texttt{theta} is

\texttt{h\_2(2/5)}

bits under base-2 logarithms.

After the pull, the entropy is zero.

So the information gain is positive.

Pulling \texttt{S} tells us nothing about \texttt{theta}, so its
information gain is zero.

It is tempting to conclude:

more information is always more valuable.

That is false.

Information gain measures uncertainty reduction.

Decision value measures how the information changes attainable utility
or reward.

The two quantities even have different units.

In the witness:

\begin{itemize}
\tightlist
\item
  information gain is measured in bits;
\item
  value of information is \texttt{1/5} reward units.
\end{itemize}

Howard's decision-theoretic treatment of information value emphasizes
exactly this dependence on consequences: uncertainty reduction matters
economically only through the decisions it changes (Howard 1966).

An observation can be highly informative about something irrelevant.

Another observation can carry only a small amount of information but
resolve the one ambiguity that changes the optimal action.

So the correct question is not merely:

\begin{quote}
how much uncertainty did we remove?
\end{quote}

It is:

\begin{quote}
which future decision did that information change, and by how much?
\end{quote}

\section{Epistemic uncertainty versus
randomness}\label{epistemic-uncertainty-versus-randomness}

Suppose an observed reward can be written schematically as

\texttt{Y=mu\_theta(a)+epsilon}.

There are two conceptually different uncertainties.

\subsection{Epistemic uncertainty}\label{epistemic-uncertainty}

We do not know \texttt{theta}.

Additional data may change our belief about it.

\subsection{Outcome stochasticity}\label{outcome-stochasticity}

Even if \texttt{theta} were known, \texttt{epsilon} could remain random.

More data may help estimate the distribution of \texttt{epsilon}, but it
does not make the underlying process deterministic.

This distinction matters because exploration targets epistemic
uncertainty.

An action can have high variance and almost no learning value.

An action can have low variance and be extremely informative.

``Choose the uncertain action'' is therefore incomplete advice until we
say what kind of uncertainty the action resolves.

\section{Bandits: the cleanest exploration
laboratory}\label{bandits-the-cleanest-exploration-laboratory}

The stochastic multi-armed bandit strips the problem down.

There are \texttt{K} arms with unknown reward laws

\texttt{nu\_1,...,nu\_K}.

At time \texttt{t}, choose one arm

\texttt{A\_t}

and observe only

\texttt{Y\_t\ \textasciitilde{}\ nu\_\{A\_t\}}.

The unchosen rewards are hidden.

There is no nontrivial controlled state transition in the basic model.

That makes the bandit a clean laboratory for the
exploration--exploitation problem.

If we repeatedly pull only the arm that currently looks best, we may
never discover that another arm is better.

If we explore forever, we continue paying opportunity cost after
uncertainty has become small.

The design problem is to trade those pressures over time.

Auer, Cesa-Bianchi, and Fischer gave finite-time analysis for
upper-confidence-bound style allocation in stochastic bandits (Auer,
Cesa-Bianchi, and Fischer 2002).

The chapter uses that work as a representative optimism mechanism, not
as a claim that one UCB formula is universally optimal.

\section{Optimism: act as if plausible good outcomes
matter}\label{optimism-act-as-if-plausible-good-outcomes-matter}

A confidence-based exploration rule often has the schematic form

\texttt{a\_t\ in\ argmax\_i\ {[}\ mu\_hat\_i(t)+beta\_i(t)\ {]}}.

The estimate

\texttt{mu\_hat\_i(t)}

summarizes observed performance.

The bonus

\texttt{beta\_i(t)}

represents uncertainty or confidence width.

An uncertain arm can therefore be chosen because its plausible upper
value remains attractive.

This is \textbf{optimism under uncertainty}.

The logic is:

\begin{quote}
if an action could still plausibly be much better than its current
estimate, sample it enough to resolve that possibility.
\end{quote}

The uncertainty bonus is not free reward.

It is a device for action selection under incomplete knowledge.

And a confidence interval is not the same mathematical object as a
Bayesian posterior.

\section{Posterior sampling: act according to posterior
possibility}\label{posterior-sampling-act-according-to-posterior-possibility}

Thompson's 1933 paper proposed allocating observations according to the
probability that one unknown probability exceeds another (Thompson
1933).

In modern Bayesian language, posterior sampling is often described as:

\begin{enumerate}
\def\labelenumi{\arabic{enumi}.}
\tightlist
\item
  sample a candidate environment parameter from the posterior;
\item
  choose the action optimal under that sampled parameter;
\item
  observe the result;
\item
  update the posterior.
\end{enumerate}

The resulting policy explores because different posterior samples make
different actions look optimal.

This differs from UCB.

UCB deliberately inflates uncertain actions through an upper confidence
quantity.

Posterior sampling randomizes through the posterior itself.

Both can balance exploration and exploitation.

They do so by different mathematical mechanisms.

\section{Information-directed
exploration}\label{information-directed-exploration}

Russo and Van Roy proposed information-directed sampling for online
optimization under partial feedback (Russo and Roy 2014).

Its distinctive idea is to compare expected short-term opportunity cost
with information gained about the optimal action.

The Atlas does not need the full theory here.

The conceptual contribution is enough:

\begin{quote}
not all information is equally decision-relevant.
\end{quote}

Suppose one action reveals a great deal about nuisance variables but
almost nothing about which action is best.

Another reveals only one bit, but that bit identifies the best action.

Raw uncertainty reduction can favor the first.

Decision-directed information can favor the second.

This is why exploration mechanisms should not be collapsed into one
slogan such as ``seek uncertainty.''

\section{Information is useful only through a future
decision}\label{information-is-useful-only-through-a-future-decision}

Define the expected one-step payoff

\texttt{bar\ r(theta,a\textquotesingle{})=E\_\{Y\textquotesingle{}\textasciitilde{}p(.\textbar{}theta,a\textquotesingle{})\}{[}r(theta,a\textquotesingle{},Y\textquotesingle{}){]}}.

The no-new-information future value is

\texttt{J\_0(b)\ =\ max\_\{a\textquotesingle{}\}\ E\_b{[}bar\ r(theta,a\textquotesingle{}){]}}.

Now imagine taking action \texttt{a}, observing \texttt{Y}, updating the
posterior, and then making one more decision.

The expected informed future value is

\texttt{J\_1(b,a)\ =\ E\_Y{[}\ max\_\{a\textquotesingle{}\}E{[}bar\ r(theta,a\textquotesingle{})\textbar{}Y,a{]}\ {]}}.

The one-step value of information is

\texttt{VoI\_b(a)=J\_1(b,a)-J\_0(b)}.

Under the idealized model, this quantity is nonnegative.

Why?

Because after receiving information, we could always ignore it.

The observation cannot reduce the feasible future value unless acquiring
it also changes costs, action availability, timing, or other problem
structure.

Those qualifications matter.

Real exploration can be costly.

It can consume time.

It can alter the state.

It can use finite resources.

It can violate constraints.

So the clean nonnegative value-of-information theorem belongs to the
idealized information channel, not automatically to every physical
probing action.

\section{Exploration in an MDP is harder than exploration in a
bandit}\label{exploration-in-an-mdp-is-harder-than-exploration-in-a-bandit}

In a bandit, an exploratory pull mainly changes what we know about arms.

In an MDP, an exploratory action can also change where we are.

That creates path dependence.

An action may reveal useful information but move the system into a
region from which recovery is difficult.

Another action may preserve future choices but reveal less.

A third may temporarily reduce reward while making an important
transition model identifiable.

This coupling between learning and state evolution is one reason
exploration in reinforcement learning is more difficult than bandit
allocation.

The action is simultaneously:

\begin{itemize}
\tightlist
\item
  an intervention;
\item
  an observation request;
\item
  a state transition;
\item
  a resource expenditure.
\end{itemize}

The RLBASE distinction between environment, policy, reward, value,
model, and observability therefore remains essential.

\section{Exploration is not
randomization}\label{exploration-is-not-randomization}

Random action can generate diverse observations.

That does not make randomization a complete exploration strategy.

Useful exploration asks:

\begin{itemize}
\tightlist
\item
  what uncertainty matters?
\item
  what observation would reduce it?
\item
  what action produces that observation?
\item
  what does that action cost?
\item
  how will the observation change a later decision?
\end{itemize}

Randomness may be part of the answer.

For posterior sampling, randomness is structural.

For epsilon-greedy exploration, random action selection is explicit.

For UCB, the policy can be deterministic given the history.

For information-directed methods, the relevant object is
decision-directed information.

So exploration should be defined by its \textbf{epistemic role}, not by
whether the policy happens to randomize.

\section{Safe exploration starts by naming the
constraint}\label{safe-exploration-starts-by-naming-the-constraint}

The phrase \textbf{safe exploration} sounds stronger than it is.

There is no useful guarantee until ``safe'' has mathematical content.

Different applications may mean:

\begin{itemize}
\tightlist
\item
  remain inside an admissible state set;
\item
  keep the probability of a constraint violation below a threshold;
\item
  limit cumulative cost;
\item
  preserve a path back to a designated region;
\item
  maintain performance above a conservative floor.
\end{itemize}

These are different constraints.

They produce different feasible policy sets.

Moldovan and Abbeel give one explicit formulation for safe exploration
in unknown MDPs using an ergodicity-based safety formulation and
restrictions to guaranteed-safe policy subsets (Moldovan and Abbeel
2012).

The Atlas takes a narrower general lesson from that source:

\begin{quote}
a safety requirement changes which exploratory actions are admissible.
\end{quote}

It is not simply an extra positive bonus attached to uncertainty.

\section{Safety can make information
expensive}\label{safety-can-make-information-expensive}

Suppose the most informative action is outside the admissible set.

Then the constrained agent may have to learn more slowly.

That is not necessarily a defect.

It is the mathematical consequence of placing value on something other
than information and reward alone.

A safety statement should therefore expose at least:

\begin{itemize}
\tightlist
\item
  the constrained quantity;
\item
  the threshold or admissible set;
\item
  the time horizon;
\item
  whether the guarantee is deterministic, probabilistic, or in
  expectation;
\item
  the probability/confidence level when applicable;
\item
  the model assumptions under which the guarantee is derived;
\item
  the resulting feasible action or policy set.
\end{itemize}

Without those details, ``safe'' is rhetorical rather than technical.

\section{Information acquisition can destroy
optionality}\label{information-acquisition-can-destroy-optionality}

An exploratory action can improve knowledge while reducing future
choice.

For example, an irreversible probe may reveal the system's type while
consuming a nonrenewable resource.

A state transition may identify dynamics while entering a region from
which some future actions are unavailable.

This means exploration has at least three potentially competing effects:

\begin{enumerate}
\def\labelenumi{\arabic{enumi}.}
\tightlist
\item
  immediate reward or cost;
\item
  information gain;
\item
  change in future feasible actions.
\end{enumerate}

The later Optionality chapter will study that third object directly.

This chapter does not import its theory.

It records the boundary:

information value is not the whole value of an exploratory action when
the action changes the future action set.

\section{Exploration and regret are related but
distinct}\label{exploration-and-regret-are-related-but-distinct}

Bandit literature often evaluates exploration policies through regret.

That is useful because exploration pays opportunity cost relative to a
comparator.

But regret is not identical to exploration.

A policy can explore for many reasons:

\begin{itemize}
\tightlist
\item
  identify the best arm;
\item
  estimate the model;
\item
  reduce posterior uncertainty;
\item
  protect against misspecification;
\item
  satisfy a constraint while learning;
\item
  gather information needed by a later controller.
\end{itemize}

The downstream Regret chapter will develop Bayesian and minimax regret
and clarify what those quantities control.

This chapter supplies the exploration objects that make that analysis
meaningful.

Dependency direction matters.

Regret theory does not retroactively define exploration here.

\section{Finite horizon versus asymptotic
reasoning}\label{finite-horizon-versus-asymptotic-reasoning}

The exact witness is finite horizon.

The reason to explore is visible after exactly one future decision.

Bandit theory often studies much longer horizons, including asymptotic
behavior.

The distinction matters.

An action worth exploring over 10,000 rounds may not be worth exploring
with one round left.

Conversely, information that helps a single decisive future choice can
have high finite-horizon value even if asymptotic average reward is
irrelevant.

Exploration policy is therefore horizon-dependent.

Any claim that an action is ``worth exploring'' should make the time
scale explicit.

\section{Failure modes}\label{failure-modes-15}

\subsection{Myopic exploitation}\label{myopic-exploitation}

Choose only the action with highest current expected reward.

Failure: never collect evidence that would change the ranking.

\subsection{Undirected exploration}\label{undirected-exploration}

Seek novelty or randomness without identifying the decision-relevant
uncertainty.

Failure: collect information that does not change useful choices.

\subsection{Uncertainty conflation}\label{uncertainty-conflation}

Treat high reward variance as high epistemic information.

Failure: repeatedly sample irreducible noise.

\subsection{Information-value
conflation}\label{information-value-conflation}

Maximize entropy reduction as if every bit had equal utility.

Failure: learn irrelevant facts while ignoring low-entropy
decision-critical uncertainty.

\subsection{Algorithm conflation}\label{algorithm-conflation}

Treat UCB, Thompson sampling, and information-directed sampling as
interchangeable implementations of one formula.

Failure: erase distinct confidence, posterior, and information-ratio
semantics.

\subsection{Safety laundering}\label{safety-laundering}

Call exploration ``safe'' without naming the constraint, horizon,
guarantee type, or assumptions.

Failure: a label substitutes for a theorem.

\subsection{Horizon erasure}\label{horizon-erasure}

Evaluate exploration without saying how many future decisions can
benefit from the information.

Failure: information cost and future value are compared on incompatible
time scales.

\section{A practical exploration
ledger}\label{a-practical-exploration-ledger}

A compact exploration analysis can be written as seven questions.

{\def\LTcaptype{none} % do not increment counter
\begin{longtable}[]{@{}
  >{\raggedright\arraybackslash}p{(\linewidth - 2\tabcolsep) * \real{0.5000}}
  >{\raggedright\arraybackslash}p{(\linewidth - 2\tabcolsep) * \real{0.5000}}@{}}
\toprule\noalign{}
\begin{minipage}[b]{\linewidth}\raggedright
Field
\end{minipage} & \begin{minipage}[b]{\linewidth}\raggedright
Question
\end{minipage} \\
\midrule\noalign{}
\endhead
\bottomrule\noalign{}
\endlastfoot
unknown & What decision-relevant quantity is not known? \\
belief/confidence & How is uncertainty represented? \\
probe & Which action changes the evidence state? \\
information & What does its observation reveal? \\
immediate cost & What reward/resource opportunity is sacrificed now? \\
future decision & Which later choice can improve because of the
information? \\
constraints & Which probes are admissible and under what guarantee? \\
\end{longtable}
}

This ledger prevents ``exploration'' from becoming a synonym for trying
things.

It forces an explicit link between uncertainty and a later choice.

\section{What the exact witness
establishes}\label{what-the-exact-witness-establishes-7}

The two-step witness is intentionally small.

It establishes exactly four numerical facts:

\begin{itemize}
\tightlist
\item
  known-first total value: \texttt{1};
\item
  unknown-first total value: \texttt{11/10};
\item
  immediate exploration cost: \texttt{1/10};
\item
  future decision value of information: \texttt{1/5}.
\end{itemize}

And it establishes one structural fact:

the lower immediate expected reward action is optimal over the two-step
horizon because the observation changes the second decision.

It does not establish that exploration always pays.

It does not establish that information gain alone is a sufficient
exploration objective.

It does not establish that any specific named algorithm will behave
identically in larger problems.

The witness is a microscope, not a universal theorem.

\section{Atlas handoff}\label{atlas-handoff-1}

\textbf{Regret --- ATLAS-CH-REGRET-001.}\\
The next chapter may assume the stochastic bandit object, the
exploration--exploitation distinction, exploration cost, and the
representative differences among optimism, posterior sampling, and
information-directed exploration. It must independently define and
analyze Bayesian and minimax regret.

The larger decision sequence is now visible:

\texttt{control\ -\textgreater{}\ uncertainty\ -\textgreater{}\ exploration\ -\textgreater{}\ regret\ -\textgreater{}\ optionality}.

At this point the Atlas can state the central lesson precisely:

\begin{quote}
information has decision value only when it changes what can rationally
be done later, and that value must be weighed against immediate cost,
horizon, state change, and constraints.
\end{quote}

\section{References used in this
chapter}\label{references-used-in-this-chapter-57}

\begin{itemize}
\tightlist
\item
  Auer, Cesa-Bianchi, and Fischer, \emph{Finite-time Analysis of the
  Multiarmed Bandit Problem} (Auer, Cesa-Bianchi, and Fischer 2002).
\item
  Thompson, \emph{On the Likelihood that One Unknown Probability Exceeds
  Another in View of the Evidence of Two Samples} (Thompson 1933).
\item
  Russo and Van Roy, \emph{Learning to Optimize via Information-Directed
  Sampling} (Russo and Roy 2014).
\item
  Howard, \emph{Information Value Theory} (Howard 1966).
\item
  Moldovan and Abbeel, \emph{Safe Exploration in Markov Decision
  Processes} (Moldovan and Abbeel 2012).
\end{itemize}

Exact source identities, the audited RLBASE prerequisite, and source
authority boundaries are recorded in:

sources/source-locks/ATLAS-CH-EXPLORE-001.yaml

\chapter{Regret}\label{regret}

\textbf{Epistemic status:} established bandit-regret theory + audited
exploration prerequisite + Atlas synthesis + exact finite witnesses.\\
\textbf{Specification:}
manuscript/specifications/ATLAS-CH-REGRET-001.md\\
\textbf{Formal packet:}
mathematics/derivations/ATLAS-CH-REGRET-001-DERIVATIONS.md\\
\textbf{Computational witness:}
mathematics/computational-witnesses/ATLAS-CW-REGRET-001.md\\
\textbf{Source lock:} sources/source-locks/ATLAS-CH-REGRET-001.yaml

Regret measures the price of not acting like a comparator.

That sentence is useful only if the comparator is named.

A regret number does not belong to an algorithm by itself.

It belongs to a complete evaluation contract:

\begin{itemize}
\tightlist
\item
  environment;
\item
  horizon;
\item
  policy;
\item
  reward or loss convention;
\item
  comparator;
\item
  expectation or prior convention.
\end{itemize}

Change the comparator and the regret can change while the learner does
exactly the same thing.

Change the prior and a Bayes-optimal policy can stop being desirable.

Change cumulative interaction cost into final recommendation quality and
the preferred algorithm can change again.

The governing rule of this chapter is:

\begin{quote}
regret is a relational quantity, not an intrinsic property of a
trajectory.
\end{quote}

That rule is the difference between using regret theory and merely
quoting a rate.

\section{From exploration to opportunity
cost}\label{from-exploration-to-opportunity-cost}

The Exploration chapter established the central learning-while-acting
tension.

An exploratory action may sacrifice immediate reward to gain information
that improves later decisions.

Regret asks a complementary question:

\begin{quote}
how much reward was lost relative to a declared comparator over the
interaction?
\end{quote}

This makes regret especially natural for bandits.

If the learner samples an arm that is not truly best, it pays an
opportunity cost.

If that sample reveals useful information, the cost may be worthwhile.

Regret does not say whether the information was worthwhile by itself.

It records the reward shortfall relative to the benchmark.

\section{The stochastic-bandit
object}\label{the-stochastic-bandit-object}

Let environment \texttt{theta} define arm means

\texttt{mu\_theta(a)}

for actions

\texttt{a\ in\ A}.

Let

\texttt{mu\_theta\^{}*=max\_a\ mu\_theta(a)}.

Define the gap

\texttt{Delta\_theta(a)=mu\_theta\^{}*-mu\_theta(a)}.

Let \texttt{Pi} be the declared admissible policy class. A policy
\texttt{pi\ in\ Pi} chooses actions

\texttt{A\_1,...,A\_T}

and observes rewards

\texttt{Y\_1,...,Y\_T}.

Under the stationary stochastic-bandit convention,

\texttt{E{[}Y\_t\ \textbar{}\ H\_\{t-1\},A\_t{]}\ =\ mu\_theta(A\_t)}.

This small structure is enough to expose several regret notions that are
often collapsed.

\section{Pathwise reward shortfall is
random}\label{pathwise-reward-shortfall-is-random}

Define the pathwise mean-benchmark regret

\texttt{R\_T\^{}path(theta)\ =\ T\ mu\_theta\^{}*\ -\ sum\_t\ Y\_t}.

The comparator term uses the expected reward of the best fixed arm.

The learner term uses realized rewards.

So \texttt{R\_T\^{}path} is random.

It can even be negative.

A learner pulling a suboptimal arm can get lucky enough that its
realized reward exceeds

\texttt{T\ mu\_theta\^{}*}.

That does not make the chosen arm optimal.

It means reward noise beat the mean benchmark on that path.

\section{Expected or pseudo-regret}\label{expected-or-pseudo-regret}

A standard stochastic-bandit quantity is

\texttt{bar\ R\_T(pi,theta)\ =\ E\_theta\^{}pi{[}\ sum\_t\ Delta\_theta(A\_t)\ {]}}.

This quantity depends on which arms the learner chooses, not on
favorable or unfavorable reward noise after conditioning on those
choices.

Bubeck and Cesa-Bianchi develop stochastic and adversarial bandit regret
in this comparator-relative framework (Bubeck and Cesa-Bianchi 2012).

Lattimore and Szepesvári give a comprehensive modern treatment across
stochastic, adversarial, Bayesian, minimax, and pure-exploration
settings (Lattimore and Szepesvari 2020).

Terminology varies across subliteratures.

The Atlas therefore writes the formula before attaching the name.

\section{Why expected pathwise regret equals pseudo-regret
here}\label{why-expected-pathwise-regret-equals-pseudo-regret-here}

Under the declared stochastic-bandit model,

\texttt{E{[}Y\_t{]}\ =\ E{[}mu\_theta(A\_t){]}}.

Therefore

\texttt{E{[}R\_T\^{}path(theta){]}\ =\ T\ mu\_theta\^{}*\ -\ sum\_t\ E{[}Y\_t{]}}

\texttt{=\ E{[}\ sum\_t\ (mu\_theta\^{}*-mu\_theta(A\_t))\ {]}}

\texttt{=\ bar\ R\_T(pi,theta)}.

The equality is useful.

The objects remain different.

One is a random pathwise shortfall.

The other is its expected gap accumulation under the declared model.

\section{Pull counts expose the cost of
exploration}\label{pull-counts-expose-the-cost-of-exploration}

Let

\texttt{N\_a(T)}

be the number of times arm \texttt{a} is selected by time \texttt{T}.

Then

\texttt{bar\ R\_T(pi,theta)\ =\ sum\_a\ Delta\_theta(a)\ E{[}N\_a(T){]}}.

This formula turns regret into an accounting identity.

Each suboptimal arm contributes:

\texttt{gap\ x\ expected\ number\ of\ pulls}.

A learning policy must sample uncertain arms enough to discover which
are suboptimal.

Regret theory asks how efficiently that information cost can be paid.

\section{Exact witness: same policy, different pathwise
regret}\label{exact-witness-same-policy-different-pathwise-regret}

Take two Bernoulli arms:

\texttt{mu(A)=1/4}, \texttt{mu(B)=3/4}.

Suppose the learner always pulls A for two rounds.

The pseudo-regret is exactly

\texttt{2(3/4-1/4)=1}.

But the pathwise mean-benchmark regret is

\texttt{3/2-(Y\_1+Y\_2)}.

If both rewards are zero, regret is

\texttt{3/2}.

If exactly one is one, regret is

\texttt{1/2}.

If both rewards are one, regret is

\texttt{-1/2}.

The same policy in the same environment can therefore produce different
pathwise regret values.

The exact expectation is still

\texttt{1}.

\section{Negative pathwise regret is not
paradoxical}\label{negative-pathwise-regret-is-not-paradoxical}

The best arm is best in expectation.

It is not guaranteed to dominate every realized sample path.

This is the same reason a fair or unfavorable gamble can win on one
trial.

So a negative pathwise value does not refute the comparator.

It exposes the difference between:

\begin{itemize}
\tightlist
\item
  stochastic performance on one path;
\item
  expected performance under the environment.
\end{itemize}

A report should say which one it is using.

\section{Bayesian regret adds a
prior}\label{bayesian-regret-adds-a-prior}

Let \texttt{nu} be a prior over environments.

Define Bayesian regret:

\texttt{BR\_T(pi,nu)\ =\ E\_\{theta\textasciitilde{}nu\}{[}\ bar\ R\_T(pi,theta)\ {]}}.

Now the evaluation weights environments by prior probability.

A policy can be excellent under the prior while weak on rare
environments.

That may be rational if the prior is part of the intended decision
problem.

It is not a worst-case statement.

The prior belongs in the claim.

\section{Worst-case regret removes prior
weighting}\label{worst-case-regret-removes-prior-weighting}

For environment class \texttt{Theta}, define

\texttt{W\_T(pi,Theta)\ =\ sup\_\{theta\ in\ Theta\}\ bar\ R\_T(pi,theta)}.

This asks:

\begin{quote}
how bad can the policy be over the declared class?
\end{quote}

Rare environments do not receive less weight.

They receive the same worst-case attention as common ones.

The environment class is therefore just as important as the prior was in
the Bayesian case.

A worst-case guarantee without the class is incomplete.

\section{Minimax regret chooses the best worst-case
policy}\label{minimax-regret-chooses-the-best-worst-case-policy}

Define

\texttt{R\_T\^{}*(Theta,Pi)\ =\ inf\_\{pi\ in\ Pi\}\ sup\_\{theta\ in\ Theta\}\ bar\ R\_T(pi,theta)}.

The policy is chosen to minimize its worst environment-specific expected
regret.

This is a different optimization problem from minimizing

\texttt{BR\_T(pi,nu)}.

The distinction is not philosophical.

It changes optimal policies in the smallest possible examples.

\section{Exact Bayes/minimax
witness}\label{exact-bayesminimax-witness}

Consider one decision and two possible environments.

Under environment \texttt{+}:

\begin{itemize}
\tightlist
\item
  A pays 1;
\item
  B pays 0.
\end{itemize}

Under environment \texttt{-}:

\begin{itemize}
\tightlist
\item
  A pays 0;
\item
  B pays 1.
\end{itemize}

Let the policy choose A with probability \texttt{q}.

Then

\texttt{R(q,+)=1-q}

and

\texttt{R(q,-)=q}.

Worst-case regret is

\texttt{max(1-q,q)}.

The minimax choice is

\texttt{q=1/2}

with regret

\texttt{1/2}.

Now put prior mass

\texttt{9/10}

on \texttt{+}.

Bayesian regret is

\texttt{9/10(1-q)+1/10\ q}.

It is minimized at

\texttt{q=1}.

The Bayes-optimal regret is

\texttt{1/10}.

But that policy has worst-case regret

\texttt{1}.

The criterion changed.

So did the optimal policy.

\section{Bayes performance is bounded by worst-case
performance}\label{bayes-performance-is-bounded-by-worst-case-performance}

For a fixed policy and a prior supported on \texttt{Theta},

\texttt{BR\_T(pi,nu)\ \textless{}=\ W\_T(pi,Theta)}.

An average over a set cannot exceed the set's supremum.

Taking the best policy over the same declared class \texttt{Pi} on each
side gives

\texttt{inf\_\{pi\ in\ Pi\}\ BR\_T(pi,nu)\ \textless{}=\ R\_T\^{}*(Theta,Pi)}.

This inequality does not say that the optimizing policies coincide.

The exact witness shows that they need not.

\section{Minimax is not prior-free
truth}\label{minimax-is-not-prior-free-truth}

A minimax policy protects against the worst environment in its declared
class.

That does not make the class itself uniquely correct.

A class can be:

\begin{itemize}
\tightlist
\item
  too broad;
\item
  too narrow;
\item
  misspecified;
\item
  computationally convenient rather than realistic.
\end{itemize}

Minimax analysis removes a prior.

It does not remove modeling assumptions.

\section{Comparator class is part of
regret}\label{comparator-class-is-part-of-regret}

Regret can be defined against many comparators.

Examples include:

\begin{itemize}
\tightlist
\item
  best fixed action;
\item
  best stationary policy;
\item
  best policy in a hypothesis class;
\item
  best action sequence with at most \texttt{S} switches;
\item
  best dynamic action sequence;
\item
  clairvoyant oracle subject to constraints.
\end{itemize}

A stronger comparator generally makes the benchmark harder to match.

Two regret numbers against different comparator classes are not the same
quantity.

\section{Exact comparator witness}\label{exact-comparator-witness}

Consider reward table:

{\def\LTcaptype{none} % do not increment counter
\begin{longtable}[]{@{}rrr@{}}
\toprule\noalign{}
round & A & B \\
\midrule\noalign{}
\endhead
\bottomrule\noalign{}
\endlastfoot
1 & 1 & 0 \\
2 & 0 & 1 \\
\end{longtable}
}

The learner chooses A twice and earns 1.

Against the best fixed arm:

\begin{itemize}
\tightlist
\item
  A earns 1;
\item
  B earns 1.
\end{itemize}

Regret is

\texttt{0}.

Against an unrestricted per-round comparator:

\begin{itemize}
\tightlist
\item
  choose A on round 1;
\item
  choose B on round 2.
\end{itemize}

Comparator reward is 2.

Regret is

\texttt{1}.

Same learner.

Same reward table.

Different comparator.

Different regret.

\section{The comparator can be too strong to be
useful}\label{the-comparator-can-be-too-strong-to-be-useful}

Suppose a learner must act causally, without future observations.

Now compare it with a clairvoyant sequence that sees all future rewards.

The resulting regret may measure a genuine gap.

But the comparator has access to information the learner never could.

That may be appropriate for one theorem and misleading for another
engineering claim.

A regret analysis should state what information the comparator is
allowed to use.

\section{Cumulative regret measures interaction
cost}\label{cumulative-regret-measures-interaction-cost}

Cumulative regret asks how much opportunity cost was paid throughout
learning.

That is natural when rewards during learning matter.

If every exploratory action costs money, time, safety margin, or service
quality, the cost accumulates.

Bandit regret theory often focuses on this online objective.

A learner should learn while performing well.

\section{Simple regret measures final recommendation
quality}\label{simple-regret-measures-final-recommendation-quality}

Pure-exploration problems often care about a different object.

After \texttt{T} samples, the learner recommends an arm

\texttt{hat\ a\_T}.

Define simple regret:

\texttt{SR\_T\ =\ mu\^{}*\ -\ mu(hat\ a\_T)}.

This ignores how much reward was lost while gathering information.

It asks only:

\begin{quote}
how good is the final recommendation?
\end{quote}

Lattimore and Szepesvári treat pure exploration as a distinct objective
family from cumulative-regret minimization (Lattimore and Szepesvari
2020).

\section{Exact cumulative/simple
witness}\label{exact-cumulativesimple-witness}

Suppose rewards are deterministic:

\begin{itemize}
\tightlist
\item
  A gives 1;
\item
  B gives 0.
\end{itemize}

A learner:

\begin{enumerate}
\def\labelenumi{\arabic{enumi}.}
\tightlist
\item
  pulls B once;
\item
  pulls A once;
\item
  recommends A.
\end{enumerate}

Cumulative regret is

\texttt{1}.

The B pull paid one unit of opportunity cost.

Final simple regret is

\texttt{0}.

The recommendation is perfect.

So cumulative regret and final recommendation error are different
objectives even in a trivial problem.

\section{Exploration can raise cumulative regret and lower simple
regret}\label{exploration-can-raise-cumulative-regret-and-lower-simple-regret}

This is not a contradiction.

Sampling a doubtful arm can be costly during interaction.

The information can improve the final choice.

That is exactly why pure exploration and online reward maximization
require different algorithms and analyses.

The objective must be fixed before the experiment.

Otherwise a method can be declared good after the fact by selecting
whichever regret notion flatters it.

\section{Sublinear regret is an average
guarantee}\label{sublinear-regret-is-an-average-guarantee}

A common goal is

\texttt{bar\ R\_T=o(T)}.

This means

\texttt{bar\ R\_T/T\ -\textgreater{}\ 0}.

Average regret per round vanishes.

Cumulative regret need not vanish.

It need not even stay bounded.

For example,

\texttt{bar\ R\_T=log\ T}

diverges as \texttt{T} grows while

\texttt{log\ T/T\ -\textgreater{}\ 0}.

So ``sublinear regret'' does not mean ``no regret.''

It means the cumulative cost grows more slowly than the horizon.

\section{A good asymptotic rate can hide finite-horizon
cost}\label{a-good-asymptotic-rate-can-hide-finite-horizon-cost}

Asymptotics answer what happens as the horizon becomes large.

Real systems operate at finite horizons.

An algorithm with excellent limiting behavior may have poor constants or
long transients.

Another may be better for the actual deployment horizon and worse
asymptotically.

A regret statement should therefore say whether it is:

\begin{itemize}
\tightlist
\item
  finite time;
\item
  asymptotic order;
\item
  asymptotic constant;
\item
  lower bound;
\item
  upper bound.
\end{itemize}

These are different claim classes.

\section{Lai--Robbins: exploration has an asymptotic information
price}\label{lairobbins-exploration-has-an-asymptotic-information-price}

Lai and Robbins established classical asymptotic efficiency results for
adaptive allocation under regular parametric stochastic-bandit
assumptions (Lai and Robbins 1985).

The broad lesson is important:

a uniformly good learner cannot simply stop sampling every suboptimal
arm immediately.

It must collect enough evidence to distinguish the true environment from
alternatives in which that arm would be optimal.

That information requirement produces logarithmic-scale sampling costs
in the classical setting.

The exact constants and regularity assumptions belong to that model
class.

They are not universal bandit laws.

\section{Finite-time analysis answers another
question}\label{finite-time-analysis-answers-another-question}

Auer, Cesa-Bianchi, and Fischer gave finite-time stochastic-bandit
analysis for upper-confidence methods (Auer, Cesa-Bianchi, and Fischer
2002).

The important Atlas distinction is between:

\begin{itemize}
\tightlist
\item
  asymptotic efficiency;
\item
  explicit finite-\texttt{T} control.
\end{itemize}

A theorem can be strong in one sense and weak in the other.

Do not silently replace one with the other.

\section{Regret does not explain the policy
mechanism}\label{regret-does-not-explain-the-policy-mechanism}

Two policies can have the same regret for different reasons.

One may explore deliberately.

Another may make estimation mistakes.

A third may be constrained from selecting the best action.

A fourth may be randomized for robustness or fairness.

Regret measures comparative reward performance.

It is not a mechanistic explanation of why actions occurred.

The Exploration chapter supplies those mechanism distinctions.

\section{Regret and information are different
units}\label{regret-and-information-are-different-units}

Russo and Van Roy's information-directed sampling explicitly couples
expected single-period regret with information acquisition (Russo and
Roy 2014).

That is useful precisely because the quantities differ.

Regret is measured in reward or loss units.

Information gain is measured in information units.

An information-ratio construction relates them.

It does not identify them.

\section{Low regret does not imply
safety}\label{low-regret-does-not-imply-safety}

Imagine a thousand-round policy with excellent cumulative reward that
takes one action capable of catastrophic failure.

If the reward objective does not encode the catastrophe at the right
scale, low regret can coexist with unacceptable safety.

A safety guarantee requires:

\begin{itemize}
\tightlist
\item
  a constrained set;
\item
  chance constraint;
\item
  risk measure;
\item
  reachability condition;
\item
  conservative baseline;
\item
  or another explicit safety object.
\end{itemize}

Regret alone does not create it.

\section{Low regret does not imply
calibration}\label{low-regret-does-not-imply-calibration}

A decision-maker can choose high-reward actions while maintaining badly
calibrated probabilities.

If only the selected actions matter to reward, prediction quality away
from those decisions may not affect regret.

Calibration therefore requires its own evaluation.

The same is true for uncertainty quality and causal identification.

\section{Low regret does not imply
fairness}\label{low-regret-does-not-imply-fairness}

A reward comparator may average across users or groups.

Low aggregate regret can coexist with severe distributional disparities.

Unless the reward or constraints encode the relevant fairness property,
regret does not measure it.

This is not an objection to regret theory.

It is a reminder that objectives define what is optimized.

\section{Low regret does not imply
robustness}\label{low-regret-does-not-imply-robustness}

A policy can have low regret on the declared environment class and fail
badly outside it.

Worst-case regret is only worst case over

\texttt{Theta}.

If distribution shift moves the world outside \texttt{Theta}, the
theorem no longer applies automatically.

Robustness requires an uncertainty set, perturbation model, or other
shift semantics.

\section{Low regret does not imply
recoverability}\label{low-regret-does-not-imply-recoverability}

A policy can obtain nearly comparator-level reward while entering a
state from which future correction is difficult.

If the comparator does not price that loss of future action freedom,
regret will not record it.

This is the boundary needed by the next chapter.

Optionality is not a synonym for low regret.

\section{Expected regret can ignore tail
structure}\label{expected-regret-can-ignore-tail-structure}

A policy with rare enormous losses can have acceptable expected regret
if those events are sufficiently rare or insufficiently penalized.

Applications sensitive to tails may need high-probability or
risk-sensitive guarantees.

Expected regret should not be promoted into a tail-risk guarantee.

\section{Bayesian regret can hide rare-environment
failure}\label{bayesian-regret-can-hide-rare-environment-failure}

The exact Bayes/minimax witness already demonstrates this.

The prior assigns probability \texttt{1/10} to the environment where
action A is wrong.

Bayes optimization chooses A deterministically.

That is optimal for the declared prior.

Its worst-case regret is still maximal.

Bayesian performance and robustness answer different questions.

\section{Minimax regret can be conservative under a strong
prior}\label{minimax-regret-can-be-conservative-under-a-strong-prior}

The reverse tradeoff also exists.

If the prior is reliable and highly concentrated, a minimax policy may
spend performance protecting against environments believed to be very
unlikely.

Whether that is desirable depends on the application.

The mathematics cannot choose the governance stance.

It can expose the consequence of the stance.

\section{A practical regret ledger}\label{a-practical-regret-ledger}

Before reading or reporting a regret result, record:

{\def\LTcaptype{none} % do not increment counter
\begin{longtable}[]{@{}
  >{\raggedright\arraybackslash}p{(\linewidth - 2\tabcolsep) * \real{0.5000}}
  >{\raggedright\arraybackslash}p{(\linewidth - 2\tabcolsep) * \real{0.5000}}@{}}
\toprule\noalign{}
\begin{minipage}[b]{\linewidth}\raggedright
Field
\end{minipage} & \begin{minipage}[b]{\linewidth}\raggedright
Question
\end{minipage} \\
\midrule\noalign{}
\endhead
\bottomrule\noalign{}
\endlastfoot
horizon & Over how many decisions? \\
environment & Which fixed environment or environment class? \\
policy & Which policy class may the learner use? \\
reward/loss & What is accumulated? \\
comparator & Best fixed action, stationary policy, switching policy,
oracle, or something else? \\
randomness & Pathwise, expected over rewards, or expected over policy
randomization? \\
prior & Is there a Bayesian prior, and what is it? \\
criterion & Bayesian, worst-case, minimax, high-probability, or another
notion? \\
terminal objective & Is simple/final recommendation regret also
relevant? \\
constraints & What important properties are outside the regret
objective? \\
\end{longtable}
}

This ledger prevents rate notation from outrunning semantics.

\section{What the exact witnesses
establish}\label{what-the-exact-witnesses-establish}

The companion witness establishes four finite facts.

First, in a two-round Bernoulli problem:

\begin{itemize}
\tightlist
\item
  pseudo-regret is \texttt{1};
\item
  pathwise regret can be \texttt{3/2}, \texttt{1/2}, or \texttt{-1/2};
\item
  expected pathwise regret is exactly \texttt{1}.
\end{itemize}

Second, in a one-step two-environment problem:

\begin{itemize}
\tightlist
\item
  minimax policy chooses A with probability \texttt{1/2};
\item
  minimax regret is \texttt{1/2};
\item
  under prior \texttt{9/10,1/10}, Bayes-optimal policy chooses A with
  probability \texttt{1};
\item
  Bayes regret is \texttt{1/10};
\item
  its worst-case regret is \texttt{1}.
\end{itemize}

Third, one realized trajectory has regret \texttt{0} against the best
fixed action and \texttt{1} against a dynamic per-round comparator.

Fourth, an exploration sequence has cumulative regret \texttt{1} and
simple regret \texttt{0}.

No asymptotic theorem is inferred from these finite calculations.

\section{Downstream handoff}\label{downstream-handoff-8}

\textbf{Optionality and Correction Capacity ---
ATLAS-CH-OPTIONALITY-001} may now assume:

\begin{itemize}
\tightlist
\item
  environment-specific expected/pseudo-regret;
\item
  pathwise versus expected-regret distinction;
\item
  Bayesian regret;
\item
  worst-case and minimax regret;
\item
  comparator-class dependence;
\item
  cumulative versus simple regret;
\item
  sublinear-regret semantics;
\item
  the negative result that low reward regret does not automatically
  preserve recoverability or future action freedom.
\end{itemize}

The downstream chapter must independently define optionality, viable
future actions, correction capacity, recoverability, and uncertainty
over future opportunities.

It may use regret as one evaluation coordinate.

It may not redefine regret to contain those properties by fiat.

\section{References used in this
chapter}\label{references-used-in-this-chapter-58}

\begin{itemize}
\tightlist
\item
  Lai and Robbins, \emph{Asymptotically Efficient Adaptive Allocation
  Rules} (Lai and Robbins 1985).
\item
  Bubeck and Cesa-Bianchi, \emph{Regret Analysis of Stochastic and
  Nonstochastic Multi-Armed Bandit Problems} (Bubeck and Cesa-Bianchi
  2012).
\item
  Lattimore and Szepesvári, \emph{Bandit Algorithms} (Lattimore and
  Szepesvari 2020).
\item
  Auer, Cesa-Bianchi, and Fischer, \emph{Finite-time Analysis of the
  Multiarmed Bandit Problem} (Auer, Cesa-Bianchi, and Fischer 2002).
\item
  Russo and Van Roy, \emph{Learning to Optimize via Information-Directed
  Sampling} (Russo and Roy 2014).
\end{itemize}

Exact source identities and claim boundaries are recorded in:

sources/source-locks/ATLAS-CH-REGRET-001.yaml

\chapter{Optionality and Correction
Capacity}\label{optionality-and-correction-capacity}

\textbf{Epistemic status:} audited Regret prerequisite + established
irreversibility/viability/empowerment precedents + Atlas synthesis +
exact finite witness.\\
\textbf{Specification:}
manuscript/specifications/ATLAS-CH-OPTIONALITY-001.md\\
\textbf{Formal packet:}
mathematics/derivations/ATLAS-CH-OPTIONALITY-001-DERIVATIONS.md\\
\textbf{Computational witness:}
mathematics/computational-witnesses/ATLAS-CW-OPTIONALITY-001.md\\
\textbf{Source lock:} sources/source-locks/ATLAS-CH-OPTIONALITY-001.yaml

A system can know exactly what went wrong and still be unable to correct
it.

That is the problem this chapter isolates.

The previous Regret chapter made one boundary explicit: low reward
regret does not automatically imply recoverability or preserved
optionality.

Here we make those missing objects visible.

The central distinction is:

\begin{quote}
information tells an agent what correction is needed; optionality
determines whether a correction remains available.
\end{quote}

Those two quantities can move together.

They need not.

\section{The door that information cannot
reopen}\label{the-door-that-information-cannot-reopen}

Imagine two stage-0 decisions.

One decision preserves both exits until new evidence arrives, but
keeping both exits available has a cost.

The other commits to the left exit immediately and costs nothing.

Later, the world reveals whether left or right was needed.

If both policies had different expected values, ordinary value
comparison might settle the choice.

The more interesting case is when the preservation cost is tuned so that
the two policies have exactly the same prior expected return.

Even then, they do not leave the same future decision problem.

One leaves two corrections.

The other leaves one.

That residue is optionality.

The door metaphor is deliberately limited.

Real control systems rarely have literal doors, and ``reopening'' may be
possible at finite cost rather than impossible.

The mathematics therefore needs to represent both hard loss of
feasibility and softer increases in correction cost.

\section{Optionality needs a
contract}\label{optionality-needs-a-contract}

``Keep options open'' is incomplete advice.

Options relative to what?

A usable optionality claim must name at least:

\begin{itemize}
\tightlist
\item
  the state abstraction;
\item
  the future horizon;
\item
  admissible controls;
\item
  viability or safety constraints;
\item
  the uncertainty that may later resolve;
\item
  the outcomes that count as distinct;
\item
  the post-evidence targets;
\item
  the cost allowed for correction.
\end{itemize}

We package these as

\texttt{D\_opt=(Theta,b,S,A,T,K,h,G,U,c\_corr)}.

Here:

\begin{itemize}
\tightlist
\item
  \texttt{Theta} is the environment class;
\item
  \texttt{b} is a belief over environments when one is used;
\item
  \texttt{S} is the state space;
\item
  \texttt{A} is the action correspondence;
\item
  \texttt{T} is the transition model;
\item
  \texttt{K} is the declared viability constraint;
\item
  \texttt{h} is the remaining horizon;
\item
  \texttt{G} identifies functionally distinct consequences;
\item
  \texttt{U} describes post-evidence targets or utilities;
\item
  \texttt{c\_corr} measures correction cost.
\end{itemize}

Without these declarations, ``optionality'' tends to collapse into
metaphor.

\section{Viable continuations come before valuable
continuations}\label{viable-continuations-come-before-valuable-continuations}

Let

\texttt{Pi\_h(s)}

be the admissible continuation policies from state \texttt{s} for the
remaining horizon.

Define the viable set

\texttt{V\_h(s)}

as those continuations that satisfy the declared constraints under the
declared uncertainty semantics.

This is a feasibility object.

It does not say that every viable continuation is good.

That distinction matters because optionality is not simply a larger
value function.

Viability theory provides a mature mathematical language for controlled
evolution under state constraints (Aubin 1991).

The Atlas uses that precedent conservatively.

The finite definitions here are ours, but they inherit one discipline
from viability theory:

\begin{quote}
a continuation should not be counted as available when the governing
constraints make it inadmissible.
\end{quote}

\section{Why counting buttons
fails}\label{why-counting-buttons-fails}

Suppose a console has one button.

Now duplicate the button label ten times without changing anything the
system can actually do.

The raw action count increased by a factor of ten.

The real choice set did not.

To remove this trivial representation dependence, attach a consequence
signature

\texttt{G\_s(pi)}

to each viable continuation.

Declare:

\texttt{pi\ \textasciitilde{}\_s\ pi\textquotesingle{}}

when they have the same consequence signature.

Then define the functional option set:

\texttt{O\_h(s)=V\_h(s)/\textasciitilde{}\_s}.

A finite count

\texttt{N\_h(s)=\textbar{}O\_h(s)\textbar{}}

now counts distinct consequence classes rather than labels.

This is still not absolute.

Change the horizon or the consequence signature and the quotient can
change.

That is not a defect.

It is a reminder that optionality is always relative to what
distinctions matter downstream.

\section{Cardinality is not value}\label{cardinality-is-not-value}

Even functional option count can mislead.

One state might offer:

\begin{itemize}
\tightlist
\item
  two useful recoveries.
\end{itemize}

Another might offer:

\begin{itemize}
\tightlist
\item
  the same two useful recoveries;
\item
  twelve irrelevant actions;
\item
  one catastrophic action.
\end{itemize}

A larger set is not automatically better.

If a weight function \texttt{w} has independent meaning, one may define:

\texttt{W\_h(s)\ =\ sum\_\{o\ in\ O\_h(s)\}\ w(o)}.

But the weights then carry substantive assumptions.

The Atlas therefore treats option count as a descriptive coordinate, not
a universal utility.

\section{Irreversibility enters when uncertainty meets
time}\label{irreversibility-enters-when-uncertainty-meets-time}

Arrow and Fisher's classic analysis made a structural point that remains
useful far beyond its environmental application: when actions are
irreversible and new information may arrive later, preserving the
ability to respond can have decision value (Arrow and Fisher 1974).

The Atlas does not import their model wholesale.

We take only the structural lesson:

\begin{itemize}
\tightlist
\item
  uncertainty may resolve after action;
\item
  an early commitment may remove later responses;
\item
  the value of waiting or preserving flexibility depends on the cost of
  doing so.
\end{itemize}

This last clause is essential.

Irreversibility does not imply ``never commit.''

It makes the timing and cost of commitment part of the decision.

\section{Recoverability is relative, not
mystical}\label{recoverability-is-relative-not-mystical}

Let \texttt{B} be a declared recovery set.

Define:

\texttt{Rec\_h(s;B)=1}

when some admissible viable continuation reaches \texttt{B} within
horizon \texttt{h}.

An action is irreversible relative to \texttt{B,h} when the resulting
state has

\texttt{Rec\_h=0}.

This definition exposes four facts that ordinary language often hides.

First, irreversibility depends on the state abstraction.

Second, it depends on the recovery target.

Third, it depends on the available controls and constraints.

Fourth, it depends on the time budget.

A cracked component may be irreversible at the material level and still
be operationally recoverable by replacement.

A software deployment may be reversible in principle and effectively
irreversible under a five-second deadline.

\section{A bad outcome is not the same as irreversible
action}\label{a-bad-outcome-is-not-the-same-as-irreversible-action}

Suppose a stochastic transition causes a temporary performance drop, but
the original operating region remains reachable in one step.

That is a bad outcome.

It is not irreversible.

Conversely, a commitment may have excellent immediate reward while
permanently deleting one future response.

That can be irreversible without being immediately bad.

The distinction matters because tail risk and optionality are different
evaluation axes.

\section{Correction cost}\label{correction-cost}

After evidence arrives, the relevant question is no longer merely ``what
can I do?''

It is:

\begin{quote}
can I still reach an acceptable response for the world I now believe I
am in?
\end{quote}

For environment \texttt{theta}, let \texttt{B\_theta} be an acceptable
target set.

After evidence resolves the environment to \texttt{theta}, let
\texttt{V\_h(s;theta)} be the viable continuation set under that
resolved environment. Define:

\texttt{k\_h(s,theta)}

as the least declared correction cost among continuations in
\texttt{V\_h(s;theta)} that reach \texttt{B\_theta}.

If no such continuation exists:

\texttt{k\_h(s,theta)=+infinity}.

This one convention cleanly separates two failure modes.

Soft degradation:

\begin{itemize}
\tightlist
\item
  the correction still exists;
\item
  it simply costs more.
\end{itemize}

Hard feasibility loss:

\begin{itemize}
\tightlist
\item
  the target is no longer reachable under the declared contract.
\end{itemize}

The latter is the sharp form of lost correction capacity.

\section{Conditional correction and ex-ante correction
capacity}\label{conditional-correction-and-ex-ante-correction-capacity}

After evidence resolves the environment to \texttt{theta}, define the
conditional correction indicator:

\texttt{C\_\{h,epsilon\}(s,theta)\ =\ 1\{k\_h(s,theta)\textless{}=epsilon\}}.

This is the post-evidence statement: in this environment, can the
declared correction still be executed within tolerance?

Before the evidence arrives, define the ex-ante belief-weighted summary:

\texttt{CC\_\{h,epsilon\}(s;b)\ =\ sum\_theta\ b(theta)\ C\_\{h,epsilon\}(s,theta)}.

Interpretation:

\texttt{CC}

is the pre-evidence probability mass of environments for which a
satisfactory correction will remain reachable within the declared
tolerance.

The quantity lies in \texttt{{[}0,1{]}}.

It is nondecreasing in \texttt{epsilon}.

It exposes hard infeasibility because an infinite correction cost never
enters at finite tolerance.

This is not the only possible optionality functional.

It is useful because its semantics are explicit and testable.

\section{The exact equal-value
witness}\label{the-exact-equal-value-witness}

Now instantiate the opening story exactly.

Let:

\texttt{Theta=\{L,R\}}

with prior:

\texttt{P(L)=P(R)=1/2}.

At stage 0 the admissible action set is \texttt{\{P,C\_L,C\_R\}}.

Preserve \texttt{P}:

\begin{itemize}
\tightlist
\item
  pay immediate reward \texttt{-1/2};
\item
  move to state \texttt{s\_P};
\item
  retain terminal actions \texttt{\{L,R\}}.
\end{itemize}

Commit-left \texttt{C\_L}:

\begin{itemize}
\tightlist
\item
  pay immediate reward \texttt{0};
\item
  move to state \texttt{s\_L};
\item
  retain only terminal action \texttt{\{L\}}.
\end{itemize}

Commit-right \texttt{C\_R} is the symmetric admissible commitment:

\begin{itemize}
\tightlist
\item
  pay immediate reward \texttt{0};
\item
  move to state \texttt{s\_R};
\item
  retain only terminal action \texttt{\{R\}}.
\end{itemize}

The learner-policy comparison remains \texttt{P} versus \texttt{C\_L}.
Declaring \texttt{C\_R} makes the environment-informed comparator below
a member of the declared stage-0 action set.

At stage 1, observe \texttt{theta} perfectly.

Terminal reward is:

\texttt{u\_theta(a)=1}

when the terminal action matches the environment, and zero otherwise.

The preserve policy waits for the observation and matches it.

The commit-left policy can only choose left.

\section{Prior value does not separate the
policies}\label{prior-value-does-not-separate-the-policies}

For preservation:

\texttt{G(pi\_P,L)=1/2};

\texttt{G(pi\_P,R)=1/2}.

Therefore:

\texttt{E{[}G(pi\_P){]}=1/2}.

For commitment:

\texttt{G(pi\_L,L)=1};

\texttt{G(pi\_L,R)=0}.

Therefore:

\texttt{E{[}G(pi\_L){]}=1/2}.

The exact prior action-value gap is zero.

A scalar expected-value comparison is silent.

\section{Bayesian regret can also
tie}\label{bayesian-regret-can-also-tie}

Use a clairvoyant environment-informed comparator that knows
\texttt{theta} before stage 0 and chooses the declared matching
commitment \texttt{C\_theta}.

Its value is:

\texttt{V\_L\^{}*=V\_R\^{}*=1}.

Preservation regret is:

\begin{itemize}
\tightlist
\item
  \texttt{1/2} in L;
\item
  \texttt{1/2} in R.
\end{itemize}

So Bayesian regret is:

\texttt{1/2}.

Commit-left regret is:

\begin{itemize}
\tightlist
\item
  \texttt{0} in L;
\item
  \texttt{1} in R.
\end{itemize}

Its Bayesian regret is also:

\texttt{1/2}.

Thus the two policies tie both in prior expected return and Bayesian
regret.

This does not mean every regret criterion ties.

Worst-case regret is:

\begin{itemize}
\tightlist
\item
  preserve: \texttt{1/2};
\item
  commit-left: \texttt{1}.
\end{itemize}

Optionality is a separate coordinate, not an assertion that all other
coordinates are blind.

\section{The future feasible sets do separate
them}\label{the-future-feasible-sets-do-separate-them}

At \texttt{s\_P}:

\texttt{O\_1(s\_P)=\{L,R\}}.

At \texttt{s\_L}:

\texttt{O\_1(s\_L)=\{L\}}.

So the functional option counts are:

\begin{itemize}
\tightlist
\item
  preserve: \texttt{2};
\item
  commit-left: \texttt{1}.
\end{itemize}

This difference is not created by extra labels.

The two preserved options have distinct terminal consequences.

\section{Correction capacity separates them more
sharply}\label{correction-capacity-separates-them-more-sharply}

Let the post-evidence target in environment \texttt{theta} be:

choose terminal action \texttt{theta}.

Give a successful matching correction cost zero.

If the matching action is unavailable, correction cost is infinity.

Then in the preserved state:

\texttt{k(s\_P,L)=0};

\texttt{k(s\_P,R)=0}.

Therefore the post-evidence indicators satisfy:

\texttt{C\_\{1,0\}(s\_P,L)=C\_\{1,0\}(s\_P,R)=1},

and the ex-ante capacity is:

\texttt{CC\_\{1,0\}(s\_P;b)=1}.

In the committed-left state:

\texttt{k(s\_L,L)=0};

\texttt{k(s\_L,R)=+infinity}.

Therefore:

\texttt{C\_\{1,0\}(s\_L,L)=1};

\texttt{C\_\{1,0\}(s\_L,R)=0};

and the ex-ante capacity is:

\texttt{CC\_\{1,0\}(s\_L;b)=1/2}.

The prior value tie hid a real asymmetry in the ability to act on future
evidence.

\section{Information is identical}\label{information-is-identical}

The environment is revealed perfectly in both branches.

Thus:

\texttt{H(Theta)=1\ bit};

\texttt{H(Theta\textbar{}Y)=0};

\texttt{I(Theta;Y)=1\ bit}.

The information gain is exactly the same.

Yet correction capacity differs by a factor of two.

This is the cleanest lesson in the chapter:

\begin{quote}
learning the answer and retaining a feasible response to the answer are
different resources.
\end{quote}

\section{Empowerment is related, but not
identical}\label{empowerment-is-related-but-not-identical}

Klyubin, Polani, and Nehaniv formalize empowerment as an
information-theoretic capacity of the action-to-future-sensor channel
and explicitly interpret it as keeping options open (Klyubin et al.
2008).

That is an important neighboring idea.

It is not the same object as the Atlas correction capacity.

Empowerment asks how much potential influence an agent can exert and
later perceive through its sensorimotor channel.

Correction capacity here asks whether declared target-specific responses
remain feasible after evidence arrives.

A state can have high influence capacity over many irrelevant dimensions
and still lack the one correction required by a particular future
target.

Conversely, a binary switch can have modest channel capacity yet
preserve exactly the two corrections that matter.

The two concepts can be compared.

They should not be silently identified.

\section{Policy optionality is not environmental
branching}\label{policy-optionality-is-not-environmental-branching}

Suppose the environment can branch into a hundred exogenous states while
the controller has only one admissible response.

The world has many possible futures.

The policy has little optionality.

Now suppose the environment has only two states, but the controller
retains two distinct corrective actions.

Environmental uncertainty is smaller.

Policy optionality is larger.

Keeping these axes separate prevents ``many possible futures'' from
being mistaken for ``many controllable futures.''

\section{Delay is not automatically
recoverability}\label{delay-is-not-automatically-recoverability}

Waiting can preserve options.

It can also consume them.

A system that postpones maintenance may ``delay commitment'' while a
physical degradation process steadily removes future repair paths.

A model that simply waits may cross a deadline after which rollback is
impossible.

Therefore optionality is not procrastination.

The correct question is whether the viable correction set is preserved,
not whether a decision was deferred.

\section{More optionality can be too
expensive}\label{more-optionality-can-be-too-expensive}

Return to the exact witness.

Replace the preservation cost \texttt{1/2} with a parameter \texttt{c}.

Then:

\texttt{Q(P)=1-c}.

Commit-left retains:

\texttt{Q(C\_L)=1/2}.

At:

\texttt{c=3/4},

preservation has expected return:

\texttt{1/4}.

Commit-left has:

\texttt{1/2}.

Preservation still has:

\begin{itemize}
\tightlist
\item
  two functional options;
\item
  correction capacity \texttt{1}.
\end{itemize}

Commit-left still has:

\begin{itemize}
\tightlist
\item
  one functional option;
\item
  correction capacity \texttt{1/2}.
\end{itemize}

Yet commitment now wins on expected return.

This is not a paradox.

Maintaining flexibility has a price.

Optionality becomes a rational control objective only after that price
and the value of future correction are part of the contract.

\section{More labels are not more
options}\label{more-labels-are-not-more-options}

A second counterexample is even simpler.

Suppose one state offers action \texttt{x}.

Another offers \texttt{x\_1} and \texttt{x\_2}.

If both labels have identical transitions, costs, and future
consequences, the second state has twice the raw action count but
exactly the same functional option set.

This is why the quotient by consequence signature is load-bearing.

Without it, optionality can be manufactured by renaming buttons.

\section{Regret and optionality answer different
questions}\label{regret-and-optionality-answer-different-questions}

Regret asks:

\begin{quote}
how much reward was lost relative to this comparator?
\end{quote}

Optionality asks:

\begin{quote}
after this action, which materially distinct viable continuations
remain, and which future corrections can still be executed?
\end{quote}

The exact witness shows that Bayesian regret can tie while correction
capacity differs.

But the objects can also align.

In the witness, worst-case regret prefers preservation.

In other problems, a comparator can explicitly reward recoverability,
making regret sensitive to option destruction.

That does not make optionality redundant.

It means the comparator has been changed to include it.

The Regret chapter's discipline remains intact:

regret is always relative to the declared comparator and reward
contract.

\section{Optionality can be an objective, a constraint, or a
diagnostic}\label{optionality-can-be-an-objective-a-constraint-or-a-diagnostic}

There is no requirement to optimize \texttt{CC} directly.

Three common uses are structurally distinct.

Objective:

maximize a scalarization such as expected reward plus a declared
optionality term.

Constraint:

maximize reward subject to:

\texttt{CC\_\{h,epsilon\}\textgreater{}=alpha}.

Diagnostic:

optimize another objective, but report correction capacity to expose
hidden commitment.

The third use is often valuable because it reveals opportunity
destruction without pretending that all opportunity should be preserved.

\section{Hard and soft correction failure should not be
merged}\label{hard-and-soft-correction-failure-should-not-be-merged}

Suppose two states both retain the desired correction.

From state A, correction costs 1.

From state B, correction costs 100.

Both have nonzero correction capacity at tolerance 100.

Only A has nonzero correction capacity at tolerance 1.

Now compare state C, where correction is impossible.

Its correction cost is infinity.

This is why the tolerance index is useful.

It distinguishes:

\begin{itemize}
\tightlist
\item
  cheap correction;
\item
  expensive correction;
\item
  no correction.
\end{itemize}

A single Boolean ``reversible'' flag often hides too much.

\section{Optionality does not imply
safety}\label{optionality-does-not-imply-safety}

A state can expose many future actions, including dangerous ones.

A high correction-capacity score relative to one target family says
nothing automatically about:

\begin{itemize}
\tightlist
\item
  safety constraints omitted from \texttt{K};
\item
  tail risk;
\item
  fairness;
\item
  robustness to model error;
\item
  calibration;
\item
  institutional authorization.
\end{itemize}

These require separate contracts or connecting theorems.

The chapter therefore resists the phrase:

\begin{quote}
more optionality is safer.
\end{quote}

Sometimes it is.

Sometimes it is not.

\section{Optionality does not imply
knowledge}\label{optionality-does-not-imply-knowledge}

A system can preserve many corrections while having no idea which one is
needed.

That is the mirror image of the main witness.

The witness gives perfect knowledge but limited correction in one
branch.

The opposite case gives broad feasible control but weak information.

Adaptive capability needs both when the task demands both.

\section{A useful evaluation
vector}\label{a-useful-evaluation-vector}

For a policy \texttt{pi}, one may report:

\texttt{E(pi)\ =\ (\ expected\ return,\ Bayesian\ regret,\ worst-case\ regret,\ functional\ option\ count,\ correction\ capacity,\ correction-cost\ profile\ )}.

This is a vector, not a theorem that one coordinate should dominate.

Scalarization or Pareto preference must be declared.

This is often the cleanest way to preserve the semantics of optionality
while still permitting ordinary decision analysis.

\section{What the exact witness
proves}\label{what-the-exact-witness-proves}

The finite witness proves exactly:

\begin{itemize}
\tightlist
\item
  prior expected return can tie;
\item
  Bayesian regret can tie;
\item
  information gain can tie;
\item
  functional option count can differ;
\item
  correction capacity can differ;
\item
  preservation cost can make the higher-optionality policy lower value.
\end{itemize}

It does not prove that correction capacity is unique.

It does not establish an optimal weighting between reward and
flexibility.

It does not promote optionality into a universal safety principle.

Its role is to expose a missing state variable in some adaptive
decisions.

\section{Downstream handoff}\label{downstream-handoff-9}

\textbf{Governed Adaptation --- ATLAS-CH-GOVADAPT-001} may, after this
chapter's audit, assume:

\begin{itemize}
\tightlist
\item
  viable continuation sets \texttt{V\_h(s)};
\item
  functional option quotients \texttt{O\_h(s)};
\item
  horizon-relative recoverability;
\item
  extended-real correction cost \texttt{k\_h(s,theta)};
\item
  conditional correction indicator \texttt{C\_\{h,epsilon\}(s,theta)};
\item
  ex-ante tolerance-indexed correction capacity
  \texttt{CC\_\{h,epsilon\}(s;b)};
\item
  the distinction between information acquisition and ability to exploit
  information;
\item
  the exact equal-value/equal-Bayes-regret separation witness;
\item
  the counterexample showing that preserving options can be too costly.
\end{itemize}

The downstream chapter must keep the target family, horizon,
constraints, consequence equivalence, and uncertainty semantics
explicit.

\section{References used in this
chapter}\label{references-used-in-this-chapter-59}

\begin{itemize}
\tightlist
\item
  Arrow and Fisher, \emph{Environmental Preservation, Uncertainty, and
  Irreversibility} (Arrow and Fisher 1974).
\item
  Aubin, \emph{Viability Theory} (Aubin 1991).
\item
  Klyubin, Polani, and Nehaniv, \emph{Keep Your Options Open: An
  Information-Based Driving Principle for Sensorimotor Systems} (Klyubin
  et al. 2008).
\end{itemize}

Exact source identities and claim boundaries are recorded in:

sources/source-locks/ATLAS-CH-OPTIONALITY-001.yaml

\chapter{Joint Uncertainty
Propagation}\label{joint-uncertainty-propagation}

\textbf{Epistemic status:} audited RL/control and information-theory
substrate + Atlas-owned exact covariance propagation and finite
witness.\\
\textbf{Specification:}
manuscript/specifications/ATLAS-CH-JOINTUNC-001.md\\
\textbf{Derivation packet:}
mathematics/derivations/ATLAS-CH-JOINTUNC-001-DERIVATIONS.md\\
\textbf{Computational witness:}
mathematics/computational-witnesses/ATLAS-CW-JOINTUNC-001.md\\
\textbf{Source lock:} sources/source-locks/ATLAS-CH-JOINTUNC-001.yaml

A sequential decision can be uncertain in several places at once.

The reward can be uncertain.

The transition can be uncertain.

The observation can be uncertain.

The continuation value can be uncertain.

If those uncertainties move together, adding four separate error bars as
if they were independent can be wrong in either direction.

The governing principle is:

\[
\boxed{
\text{joint uncertainty requires the joint law, not only the marginals.}
}
\]

This chapter formalizes that principle in the smallest exact setting
that exposes the issue.

\section{The controlled-system
handoff}\label{the-controlled-system-handoff}

Reinforcement Learning and Control separated:

\begin{itemize}
\tightlist
\item
  state;
\item
  action;
\item
  transition kernel;
\item
  reward;
\item
  policy;
\item
  return;
\item
  value;
\item
  learned model;
\item
  observation structure.
\end{itemize}

It also stated a boundary:

transition, reward, observation, and future-value uncertainty can
interact.

RLBASE did not solve that interaction.

JOINTUNC starts there.

\section{Information-theoretic
discipline}\label{information-theoretic-discipline}

Probability, Information, and Statistical Structure supplied another
discipline:

uncertainty quantities are defined only relative to a declared
probability model.

Dependence is also model-relative.

A covariance or mutual information does not become causal direction
merely because it is nonzero.

JOINTUNC inherits that firewall.

\section{Four local coordinates}\label{four-local-coordinates}

Write a declared local error vector:

\[
\varepsilon
=
\begin{pmatrix}
\varepsilon_R\\
\varepsilon_T\\
\varepsilon_O\\
\varepsilon_V
\end{pmatrix}.
\]

Interpret the coordinates as local scalarized uncertainty in:

\begin{itemize}
\tightlist
\item
  reward;
\item
  transition effect;
\item
  observation or belief state;
\item
  future value.
\end{itemize}

These are not universal latent variables.

They are coordinates chosen for a declared local decision calculation.

\section{Marginals are not enough}\label{marginals-are-not-enough}

Suppose we know:

\[
\operatorname{Var}(\varepsilon_R),
\quad
\operatorname{Var}(\varepsilon_T),
\quad
\operatorname{Var}(\varepsilon_O),
\quad
\operatorname{Var}(\varepsilon_V).
\]

That is not yet a full joint uncertainty description.

We also need the cross-covariances collected in:

\[
\Sigma
=
\operatorname{Cov}(\varepsilon).
\]

The diagonal records marginal variances.

The off-diagonal entries record pairwise covariance.

\section{A declared local decision-error
functional}\label{a-declared-local-decision-error-functional}

Use the exact affine witness:

\[
\boxed{
\delta
=
\varepsilon_R
+
\varepsilon_T
+
\varepsilon_O
+
\frac12\varepsilon_V.
}
\]

The sensitivity vector is:

\[
a=
\begin{pmatrix}
1\\1\\1\\1/2
\end{pmatrix}.
\]

The factor \(1/2\) plays a discount-like role for the future-value
coordinate.

The first three coefficients are normalized to one.

This is a toy local functional, not a universal Bellman uncertainty
equation.

\section{Exact affine propagation}\label{exact-affine-propagation}

For any finite-second-moment error vector:

\[
\boxed{
\operatorname{Var}(a^\top\varepsilon)
=
a^\top\Sigma a.
}
\]

Expanded:

\[
\operatorname{Var}(a^\top\varepsilon)
=
\sum_i a_i^2\operatorname{Var}(\varepsilon_i)
+
2\sum_{i<j}
a_i a_j
\operatorname{Cov}(\varepsilon_i,\varepsilon_j).
\]

The covariance terms are not optional decoration.

They are part of the exact variance.

\section{The naive diagonal
calculation}\label{the-naive-diagonal-calculation}

Give every coordinate unit marginal variance:

\[
\operatorname{Var}(\varepsilon_i)=1.
\]

Then the diagonal-only sum is:

\[
1+1+1+\frac14
=
\boxed{\frac{13}{4}}.
\]

If the weighted covariance terms vanish, this is exact.

If they do not, it is not.

\section{Positive common-shock
dependence}\label{positive-common-shock-dependence}

Let \(U,W\) be independent centered Rademacher variables.

Set:

\[
\varepsilon_R=U,
\]

\[
\varepsilon_T=U,
\]

\[
\varepsilon_O=U,
\]

\[
\varepsilon_V=W.
\]

The first three uncertainties share one common shock.

Their covariance matrix is:

\[
\Sigma_+
=
\begin{pmatrix}
1&1&1&0\\
1&1&1&0\\
1&1&1&0\\
0&0&0&1
\end{pmatrix}.
\]

\section{Positive dependence amplifies
uncertainty}\label{positive-dependence-amplifies-uncertainty}

Now:

\[
\delta_+
=
3U+\frac12W.
\]

Since \(U\) and \(W\) are independent:

\[
\operatorname{Var}(\delta_+)
=
9+\frac14
=
\boxed{\frac{37}{4}}.
\]

The diagonal-only value was:

\[
\frac{13}{4}.
\]

Difference:

\[
\boxed{6}.
\]

Ignoring positive covariance has substantially understated the
propagated variance.

\section{Cancellation dependence}\label{cancellation-dependence}

Now keep the same marginal variances but reverse the observation
coordinate:

\[
\varepsilon_R=U,
\]

\[
\varepsilon_T=U,
\]

\[
\varepsilon_O=-U,
\]

\[
\varepsilon_V=W.
\]

Then:

\[
\Sigma_-
=
\begin{pmatrix}
1&1&-1&0\\
1&1&-1&0\\
-1&-1&1&0\\
0&0&0&1
\end{pmatrix}.
\]

\section{Dependence can also
cancel}\label{dependence-can-also-cancel}

The propagated error is now:

\[
\delta_-
=
U+\frac12W.
\]

Therefore:

\[
\operatorname{Var}(\delta_-)
=
1+\frac14
=
\boxed{\frac54}.
\]

The diagonal-only value was still:

\[
\frac{13}{4}.
\]

Difference:

\[
\boxed{-2}.
\]

Ignoring covariance has now overstated the propagated variance.

\section{Independence control}\label{independence-control}

Let:

\[
U_R,U_T,U_O,U_V
\]

be mutually independent centered Rademacher variables.

Set each uncertainty coordinate equal to its own independent variable.

Then:

\[
\Sigma_0=I_4.
\]

Therefore:

\[
\operatorname{Var}(\delta_0)
=
1+1+1+\frac14
=
\boxed{\frac{13}{4}}.
\]

The diagonal calculation is exact.

\section{One set of marginals, three different
results}\label{one-set-of-marginals-three-different-results}

All three cases have the same marginal variance vector:

\[
(1,1,1,1).
\]

Their propagated variances are:

\[
\boxed{
\frac{37}{4},
\qquad
\frac54,
\qquad
\frac{13}{4}.
}
\]

Therefore:

\[
\boxed{
\text{same marginal variances}
\not\Rightarrow
\text{same propagated uncertainty}.
}
\]

\section{The difference is dependence
structure}\label{the-difference-is-dependence-structure}

The three experiments did not change the marginal error sizes.

They changed only how the error coordinates move together.

That is the central point.

If the uncertainty object is joint, its propagation should be joint.

\section{Independence is stronger than
required}\label{independence-is-stronger-than-required}

The covariance correction vanishes when the relevant weighted
cross-covariances vanish.

Full independence guarantees that.

But full independence is not necessary.

Thus:

\[
\boxed{
\text{uncorrelated}
\not\Rightarrow
\text{independent}.
}
\]

The variance calculation needs second moments.

A claim of independence needs more.

\section{Propagation depends on the decision
functional}\label{propagation-depends-on-the-decision-functional}

Suppose we keep the same covariance matrix but change the decision
quantity.

Use sensitivity vector \(b\) instead of \(a\).

Then:

\[
\operatorname{Var}(b^\top\varepsilon)
=
b^\top\Sigma b.
\]

Different decision functionals can therefore assign different importance
to the same covariance structure.

There is no universal propagated-uncertainty scalar independent of the
quantity being propagated.

\section{Reward uncertainty remains a distinct
object}\label{reward-uncertainty-remains-a-distinct-object}

Reward uncertainty can mean several things:

\begin{itemize}
\tightlist
\item
  random reward realization;
\item
  measurement noise;
\item
  learned reward-model error;
\item
  reward misspecification.
\end{itemize}

Those are not identical.

The finite witness deliberately compresses them into one local scalar
only after the modeling choice is declared.

\section{Transition uncertainty remains distinct
too}\label{transition-uncertainty-remains-distinct-too}

Transition uncertainty can include:

\begin{itemize}
\tightlist
\item
  inherent next-state randomness;
\item
  learned transition-model uncertainty;
\item
  local approximation error;
\item
  uncertainty in a transition-sensitive summary.
\end{itemize}

The chapter does not silently equate those notions.

\section{Observation uncertainty is not transition
uncertainty}\label{observation-uncertainty-is-not-transition-uncertainty}

In a partially observed system, latent-state dynamics and observation
generation are separate stochastic mechanisms.

An observation can be noisy even when the latent transition is known.

A transition can be uncertain even when observation is perfect.

They can also be dependent.

This is why the observation coordinate remains explicit.

\section{Future-value uncertainty is downstream but not
redundant}\label{future-value-uncertainty-is-downstream-but-not-redundant}

Future-value uncertainty can reflect:

\begin{itemize}
\tightlist
\item
  uncertain next state;
\item
  value-function estimation error;
\item
  function approximation;
\item
  finite data;
\item
  model error.
\end{itemize}

It is not automatically absorbed by the transition coordinate.

A declared calculation must say which uncertainty is represented where.

\section{Expected value is not
uncertainty}\label{expected-value-is-not-uncertainty}

RLBASE established value as expected future return under declared
conditions.

That expectation is not an uncertainty estimate.

Two actions can have equal expected value and different propagated
uncertainty.

Two actions can also have equal propagated variance and different
expected value.

Therefore:

\[
\boxed{
\text{expected value}
\neq
\text{uncertainty}.
}
\]

\section{A risk rule must be
declared}\label{a-risk-rule-must-be-declared}

A decision system may intentionally combine expectation and uncertainty.

For example, it may penalize variance.

That is a new decision rule.

It should not be smuggled into the meaning of value.

The utility or risk criterion must be declared explicitly.

\section{Variance is useful but
incomplete}\label{variance-is-useful-but-incomplete}

Variance is a second-moment summary.

Two distributions can share mean and variance while differing in:

\begin{itemize}
\tightlist
\item
  tail probability;
\item
  skewness;
\item
  support;
\item
  multimodality;
\item
  catastrophic-event probability.
\end{itemize}

Therefore:

\[
\boxed{
\text{same variance}
\not\Rightarrow
\text{same risk distribution}.
}
\]

\section{Nonlinear decision maps}\label{nonlinear-decision-maps}

Real decision functionals are often nonlinear.

Let:

\[
F(x)
\]

be differentiable at a nominal state \(x\).

For perturbation \(\varepsilon\):

\[
F(x+\varepsilon)-F(x)
=
\nabla F(x)^\top\varepsilon
+
R.
\]

The remainder \(R\) matters.

\section{The exact remainder
identity}\label{the-exact-remainder-identity}

The variance is exactly:

\[
\operatorname{Var}
\left(
F(x+\varepsilon)-F(x)
\right)
=
\nabla F(x)^\top
\Sigma
\nabla F(x)
+
\operatorname{Var}(R)
+
2\operatorname{Cov}
\left(
\nabla F(x)^\top\varepsilon,
R
\right).
\]

Only the first term is the familiar linearized propagation term.

\section{First-order propagation is
local}\label{first-order-propagation-is-local}

If the remainder contributions are controlled, then:

\[
\operatorname{Var}
\left(
F(x+\varepsilon)-F(x)
\right)
\approx
\nabla F(x)^\top
\Sigma
\nabla F(x).
\]

This is a local approximation.

It is not a license to propagate one covariance matrix globally through
an arbitrary nonlinear agent.

\section{When the local approximation can
fail}\label{when-the-local-approximation-can-fail}

The first-order picture can fail when:

\begin{itemize}
\tightlist
\item
  perturbations are not small;
\item
  curvature is large;
\item
  thresholds or discontinuities matter;
\item
  policies switch;
\item
  state visitation changes;
\item
  tail events dominate;
\item
  the covariance itself changes across the perturbed region.
\end{itemize}

A global theorem needs additional arguments.

\section{Sequential propagation is harder than one
step}\label{sequential-propagation-is-harder-than-one-step}

For:

\[
x_{t+1}=F_t(x_t,\eta_t),
\]

later uncertainty depends on:

\begin{itemize}
\tightlist
\item
  local sensitivities;
\item
  cross-time dependence;
\item
  state-dependent noise;
\item
  policy feedback;
\item
  changing distributions;
\item
  changing approximators.
\end{itemize}

The exact affine witness is a building block.

It does not solve the full recursion.

\section{Why naive independent summation is
dangerous}\label{why-naive-independent-summation-is-dangerous}

The positive common-shock witness gave:

\[
\frac{37}{4}
\]

while naive summation gave:

\[
\frac{13}{4}.
\]

The cancellation witness gave:

\[
\frac54
\]

while naive summation again gave:

\[
\frac{13}{4}.
\]

So the naive rule is not merely slightly noisy.

It can fail in either direction.

\section{Covariance is not
causality}\label{covariance-is-not-causality}

Suppose reward and transition errors have positive covariance.

That says they co-vary under the declared joint law.

It does not tell us whether:

\begin{itemize}
\tightlist
\item
  reward error causes transition error;
\item
  transition error causes reward error;
\item
  a third variable drives both;
\item
  the dependence is induced by conditioning or selection.
\end{itemize}

Therefore:

\[
\boxed{
\text{statistical dependence}
\not\Rightarrow
\text{causal direction}.
}
\]

\section{Conditional structure
matters}\label{conditional-structure-matters}

The relevant covariance can change after conditioning on:

\begin{itemize}
\tightlist
\item
  current state;
\item
  action;
\item
  observation;
\item
  belief state;
\item
  model version;
\item
  policy.
\end{itemize}

A joint uncertainty analysis should say what is conditioned on.

One unconditional covariance matrix may be inappropriate for every local
decision context.

\section{Functional-relative uncertainty is an interface
concept}\label{functional-relative-uncertainty-is-an-interface-concept}

Suppose a downstream system asks:

\begin{quote}
how uncertain is this action-value comparison?
\end{quote}

The upstream uncertainty representation should expose enough joint
structure to answer that question.

A list of independent scalar error bars may not be a sufficient
interface.

This is a systems consequence of the mathematics.

\section{What JOINTUNC does not
claim}\label{what-jointunc-does-not-claim}

It does not claim that:

\[
\text{variance}
=
\text{complete risk}.
\]

It does not claim that:

\[
\text{covariance}
=
\text{cause}.
\]

It does not claim that:

\[
\text{local linearization}
=
\text{global sequential theorem}.
\]

It does not claim that the four witness coordinates are canonical.

It does not claim real RL errors have the toy joint laws.

\section{Non-implications}\label{non-implications-5}

The chapter rejects:

\[
\text{same marginal variances}
\not\Rightarrow
\text{same propagated uncertainty},
\]

\[
\text{zero covariance}
\not\Rightarrow
\text{independence},
\]

\[
\text{statistical dependence}
\not\Rightarrow
\text{causal direction},
\]

\[
\text{local first-order variance}
\not\Rightarrow
\text{global sequential uncertainty},
\]

\[
\text{same variance}
\not\Rightarrow
\text{same risk distribution},
\]

and:

\[
\text{high uncertainty}
\not\Rightarrow
\text{low expected value}.
\]

\section{Practical protocol}\label{practical-protocol}

For a real decision system:

\begin{enumerate}
\def\labelenumi{\arabic{enumi}.}
\tightlist
\item
  define the decision functional;
\item
  identify distinct uncertainty sources;
\item
  define the conditioning context;
\item
  estimate or bound the joint covariance/dependence structure;
\item
  propagate through the declared sensitivity map;
\item
  test whether a local linear approximation is adequate;
\item
  expose remainder or nonlinear risk when it is not;
\item
  report expectation separately from uncertainty;
\item
  retain dependence/causality separation;
\item
  preserve the full distribution when second moments are insufficient.
\end{enumerate}

\section{Atlas connections}\label{atlas-connections-22}

\textbf{Reinforcement Learning and Control.}\\
JOINTUNC extends the controlled stochastic substrate without changing
Bellman semantics.

\textbf{Probability and Information.}\\
Joint distributions and statistical dependence are the mathematical
foundation of the uncertainty object.

\textbf{Uncertainty and Calibration.}\\
Marginal predictive uncertainty can feed JOINTUNC, but calibration of
individual predictors does not by itself specify cross-source
dependence.

\textbf{Distribution Shift.}\\
A covariance structure learned in one regime may change under shift.

\textbf{Agents and Systems.}\\
Joint uncertainty becomes an interface requirement when one subsystem
passes uncertainty to another.

\section{Closing view}\label{closing-view-31}

Uncertainty does not propagate source by source unless the joint
structure permits it.

The exact witness makes that unavoidable.

With the same four marginal variances:

\[
(1,1,1,1),
\]

we obtained:

\[
\frac{37}{4},
\qquad
\frac54,
\qquad
\frac{13}{4}.
\]

Nothing changed except dependence.

The correct conclusion is:

\[
\boxed{
\text{propagate the joint uncertainty of the declared decision object, not an assumed independent sum of marginal errors.}
}
\]

And when the decision map is nonlinear, keep the local approximation
boundary visible.

\section{References used in this
chapter}\label{references-used-in-this-chapter-60}

No new external academic authority is added.

The controlled-decision substrate is inherited through audited
ATLAS-CH-RLBASE-001.

The probability/information substrate is inherited through audited
ATLAS-CH-INFO-001.

Exact prerequisite identities and claim boundaries are recorded in:

sources/source-locks/ATLAS-CH-JOINTUNC-001.yaml

\chapter{Uncertainty and Calibration}\label{uncertainty-and-calibration}

\textbf{Epistemic status:} established probability/statistics +
source-scoped uncertainty methods + Atlas synthesis.

Specification: manuscript/specifications/ATLAS-CH-UNCERTAINTY-001.md
Derivation packet:
mathematics/derivations/ATLAS-CH-UNCERTAINTY-001-DERIVATIONS.md
Computational witness:
mathematics/computational-witnesses/ATLAS-CW-UNCERTAINTY-001.md Source
lock: sources/source-locks/ATLAS-CH-UNCERTAINTY-001.yaml

A model can be wrong in more than one way.

It can be:

\begin{itemize}
\tightlist
\item
  uncertain because the outcome itself is noisy;
\item
  uncertain because several models or parameter settings remain
  plausible;
\item
  confident but systematically miscalibrated;
\item
  calibrated but uninformative;
\item
  uncertain in a Bayesian approximation;
\item
  covered by a conformal set without assigning calibrated probabilities;
\item
  accurate on accepted examples only because it abstains on the rest.
\end{itemize}

These are not synonyms.

This chapter develops a vocabulary for keeping them separate.

The governing principle is:

\begin{quote}
A useful uncertainty statement must name the random object, the
conditioning information, the estimator or diagnostic, and the
guarantee.
\end{quote}

\section{Probability is not yet
calibration}\label{probability-is-not-yet-calibration}

Suppose a classifier emits a score

\[
S\in[0,1]
\]

for a binary outcome

\[
Y\in\{0,1\}.
\]

A probability-like number can be read as a claim.

If the model says

\[
S=0.8,
\]

the calibration question is not whether the model is correct on this one
example.

It is whether examples receiving score 0.8 are positive about 80 percent
of the time under the relevant distribution.

Population calibration can be written as

\[
E[Y\mid S]=S
\]

almost surely.

At an atomic score value,

\[
P(Y=1\mid S=s)=s.
\]

This is a property of the joint law of predictions and outcomes.

It is not a property of the numerical range {[}0,1{]} alone.

\section{Confidence is not
calibration}\label{confidence-is-not-calibration}

A model can emit scores close to one and still be badly calibrated.

Likewise, a model can be calibrated while rarely producing extreme
probabilities.

Guo et al.~study confidence calibration in modern neural networks and
evaluate post-processing methods including temperature scaling (Guo et
al. 2017).

The source-scoped lesson is not that temperature scaling universally
solves uncertainty.

It is that probability calibration is measurable, can fail in trained
neural networks, and can be altered by a post-hoc map without changing
the predicted class ordering.

\section{Calibration is not
accuracy}\label{calibration-is-not-accuracy}

Accuracy asks how often a decision rule is correct.

Calibration asks whether stated probabilities match conditional
frequencies.

A classifier can improve one without necessarily improving the other.

Temperature scaling is a useful example: changing logits by a positive
scalar temperature does not change the argmax class, yet it can change
probability calibration.

This is why uncertainty work should keep:

\begin{itemize}
\tightlist
\item
  decisions;
\item
  scores;
\item
  probability claims;
\item
  losses
\end{itemize}

as separate objects.

\section{Calibration is not
sharpness}\label{calibration-is-not-sharpness}

Consider a population with

\[
P(X=a)=P(X=b)=1/2
\]

and

\[
P(Y=1\mid X=a)=0.9,
\qquad
P(Y=1\mid X=b)=0.1.
\]

The marginal positive rate is 0.5.

A constant predictor

\[
S_0=0.5
\]

is calibrated.

So is the predictor

\[
S_1(a)=0.9,
\qquad
S_1(b)=0.1.
\]

But they do not contain the same information.

For the constant predictor,

\[
E[(Y-S_0)^2]=0.25.
\]

For the sharper conditional predictor,

\[
E[(Y-S_1)^2]=0.09.
\]

Both are calibrated.

Only one distinguishes the two covariate regimes.

Therefore:

\[
\boxed{
\text{calibration}
\not\Rightarrow
\text{sharpness or informativeness}.
}
\]

\section{Proper scoring and calibration are related but
distinct}\label{proper-scoring-and-calibration-are-related-but-distinct}

Log loss and Brier score reward probabilistic predictions using the
realized outcome.

They are not merely calibration errors.

A predictor can have good calibration and poor resolution.

A predictor can have better proper-score risk while both predictors are
calibrated.

The previous exact witness demonstrates this with Brier risk.

The Atlas therefore uses proper scores, calibration diagnostics, and
discrimination metrics as distinct evaluation axes.

\section{Empirical calibration is an
estimate}\label{empirical-calibration-is-an-estimate}

Population calibration is a property of a distribution.

A reliability diagram is a finite-sample diagnostic.

Expected calibration error is a finite-sample, bin-dependent summary.

A common form is

\[
ECE
=
\sum_j
\frac{|B_j|}{n}
\left|
acc(B_j)-conf(B_j)
\right|.
\]

Change the bins and the numerical ECE can change.

Therefore a small empirical ECE is not itself a theorem that

\[
E[Y\mid S]=S
\]

in the deployment population.

\section{One word, several uncertainty
objects}\label{one-word-several-uncertainty-objects}

Suppose a model defines a predictive law

\[
p(y\mid x,D).
\]

Possible uncertainty summaries include:

\begin{itemize}
\tightlist
\item
  predictive variance;
\item
  predictive entropy;
\item
  posterior variance over parameters;
\item
  disagreement among models;
\item
  mutual information between predictions and latent model variables;
\item
  width of a prediction interval;
\item
  conformal set size;
\item
  probability of abstention;
\item
  selective risk.
\end{itemize}

These quantities answer different questions.

A chapter that simply plots an ``uncertainty score'' without defining
the underlying object has not yet specified the claim.

\section{Predictive entropy}\label{predictive-entropy}

For a discrete predictive law,

\[
H(Y\mid x,D)
=
-\sum_y
p(y\mid x,D)
\log p(y\mid x,D).
\]

Entropy summarizes uncertainty in the predictive distribution.

It does not say where that uncertainty came from.

This matters because the same predictive distribution can arise from
very different latent structures.

\section{Aleatoric and epistemic
uncertainty}\label{aleatoric-and-epistemic-uncertainty}

Kendall and Gal use a two-way distinction in Bayesian deep-learning
models for computer vision (Kendall and Gal 2017):

\begin{itemize}
\tightlist
\item
  aleatoric uncertainty: uncertainty associated with noise inherent in
  observations or the conditional data model;
\item
  epistemic uncertainty: uncertainty associated with the model and its
  parameters.
\end{itemize}

This distinction is useful.

It is not an exhaustive ontology of every uncertainty that can arise in
an adaptive system.

Numerical approximation, distribution shift, misspecification,
measurement failure, hidden interventions, and unknown unknowns can
require other descriptions.

The Atlas therefore treats aleatoric/epistemic language as
model-relative.

\section{Entropy decomposition}\label{entropy-decomposition}

Let Theta denote a latent model variable.

Then

\[
I(Y;\Theta)
=
H(Y)-H(Y\mid\Theta).
\]

Equivalently,

\[
H(Y)
=
E_\Theta[H(Y\mid\Theta)]
+
I(Y;\Theta).
\]

In a declared Bayesian predictive model, this can motivate a
decomposition into:

\begin{itemize}
\tightlist
\item
  expected conditional uncertainty;
\item
  model-dependent information gain or epistemic component.
\end{itemize}

The identity is exact.

The interpretation depends on what Theta means.

\section{Same entropy, opposite
decomposition}\label{same-entropy-opposite-decomposition}

Consider two models.

\subsection{Model A}\label{model-a}

Theta is fixed.

\[
Y\sim Bernoulli(1/2).
\]

Then

\[
H(Y)=1\text{ bit},
\]

\[
E[H(Y\mid\Theta)]=1,
\]

\[
I(Y;\Theta)=0.
\]

\subsection{Model B}\label{model-b}

\[
\Theta\sim Bernoulli(1/2),
\]

and

\[
Y=\Theta.
\]

Marginally,

\[
Y\sim Bernoulli(1/2).
\]

So again,

\[
H(Y)=1\text{ bit}.
\]

But now conditioning on Theta makes Y deterministic:

\[
E[H(Y\mid\Theta)]=0,
\]

\[
I(Y;\Theta)=1\text{ bit}.
\]

The predictive distribution is identical.

The decomposition is opposite.

Therefore:

\[
\boxed{
\text{predictive entropy alone}
\not\Rightarrow
\text{epistemic uncertainty}.
}
\]

\section{Reducibility is relative to the information
set}\label{reducibility-is-relative-to-the-information-set}

The phrase ``epistemic uncertainty can be reduced with more data''
requires a model and a data-acquisition regime.

More observations of the same kind may reduce parameter uncertainty.

They may fail to resolve:

\begin{itemize}
\tightlist
\item
  structural misspecification;
\item
  unobserved confounding;
\item
  systematic measurement bias;
\item
  an unidentifiable parameterization;
\item
  a changed environment.
\end{itemize}

So reducibility is not an intrinsic label attached to a scalar variance.

It is relative to the model, observations, and interventions that are
possible.

\section{Bayesian uncertainty}\label{bayesian-uncertainty}

A Bayesian model places a distribution over uncertain quantities and
updates it using data.

For parameters Theta,

\[
p(\Theta\mid D)
\propto
p(D\mid\Theta)p(\Theta).
\]

The posterior predictive is

\[
p(y\mid x,D)
=
\int
p(y\mid x,\Theta)
p(\Theta\mid D)
d\Theta.
\]

This is a precise probabilistic object.

A practical neural-network approximation may only approximate it.

The distinction between target and approximation should remain visible.

\section{Dropout as an approximate Bayesian
construction}\label{dropout-as-an-approximate-bayesian-construction}

Gal and Ghahramani develop a theoretical framework relating dropout
training in neural networks to approximate Bayesian inference in deep
Gaussian-process models (Gal and Ghahramani 2016).

Within that framework, repeated stochastic forward passes can be used to
estimate predictive quantities.

The Atlas preserves two boundaries.

First:

\[
\text{MC dropout}
\neq
\text{exact posterior sampling}.
\]

Second:

\[
\text{dropout used in a network}
\not\Rightarrow
\text{all Bayesian guarantees}.
\]

The interpretation comes from the cited approximation framework, not
from randomness alone.

\section{Deep ensembles}\label{deep-ensembles}

Lakshminarayanan, Pritzel, and Blundell study independently trained
neural networks as a simple and scalable method for predictive
uncertainty estimation (Lakshminarayanan et al. 2017).

Given predictive distributions

\[
p_m(y\mid x),
\qquad
m=1,\dots,M,
\]

an ensemble mixture is

\[
\bar p(y\mid x)
=
\frac1M
\sum_{m=1}^M
p_m(y\mid x).
\]

This is an exact mixture identity.

It does not imply that the members are samples from

\[
p(\Theta\mid D).
\]

Ensemble disagreement can be useful empirically without becoming
posterior variance by definition.

\section{Ensemble spread and epistemic
uncertainty}\label{ensemble-spread-and-epistemic-uncertainty}

It is tempting to call all disagreement epistemic uncertainty.

That is too fast.

Ensemble spread can depend on:

\begin{itemize}
\tightlist
\item
  initialization;
\item
  optimization path;
\item
  training-data resampling;
\item
  architecture diversity;
\item
  explicit priors;
\item
  regularization;
\item
  stochastic training;
\item
  mode coverage.
\end{itemize}

The generated ensemble determines what its spread means.

The Atlas therefore says:

\begin{quote}
ensemble disagreement is an empirical uncertainty diagnostic whose
epistemic interpretation must be justified by the ensemble construction.
\end{quote}

\section{Prediction sets are another
object}\label{prediction-sets-are-another-object}

Sometimes the output is not a probability.

It is a set

\[
C(x)
\]

constructed to contain the future response with a target frequency.

Conformal prediction provides a framework for this kind of set-valued
prediction (Vovk et al. 2005; Lei et al. 2018).

Its signature guarantee is coverage under stated
exchangeability/randomness assumptions.

That is different from probability calibration.

\section{Split conformal rank
logic}\label{split-conformal-rank-logic}

Suppose we have

\[
n
\]

calibration nonconformity scores and one future score.

For target miscoverage

\[
\alpha,
\]

the finite-sample rank correction uses

\[
k
=
\left\lceil
(n+1)(1-\alpha)
\right\rceil.
\]

Under exchangeability, the future score has a symmetric rank among the

\[
n+1
\]

scores.

This rank symmetry is the core of the finite-sample marginal-coverage
logic.

\section{Exact conformal rank
witness}\label{exact-conformal-rank-witness}

Take

\[
n=4,
\qquad
\alpha=0.2.
\]

Then

\[
k
=
\lceil5(0.8)\rceil
=
4.
\]

Assume no ties.

The future score's rank is uniform over

\[
1,2,3,4,5.
\]

The future point is covered unless its score is the unique largest
score.

Hence

\[
P(Y_{new}\in C(X_{new}))
=
\frac45
=
0.8.
\]

This is an exact finite rank argument.

It is marginal.

It is not arbitrary conditional coverage.

\section{Coverage is not
calibration}\label{coverage-is-not-calibration}

Calibration concerns statements such as

\[
P(Y=1\mid S=s)=s.
\]

Conformal coverage concerns

\[
P(Y\in C(X))
\ge
1-\alpha
\]

under the construction's assumptions.

A set can have the target marginal coverage without supplying calibrated
class probabilities.

A calibrated probability model does not automatically produce a
conformal finite-sample coverage theorem.

These are distinct guarantees.

\section{Distribution-free does not mean
assumption-free}\label{distribution-free-does-not-mean-assumption-free}

The phrase ``distribution-free'' can be misleading when detached from
the sampling assumptions.

Conformal validity does not require a parametric family for the outcome
distribution.

It does require the symmetry/exchangeability conditions used by the
theorem.

A deployment process with:

\begin{itemize}
\tightlist
\item
  temporal drift;
\item
  feedback;
\item
  covariate shift;
\item
  intervention;
\item
  adaptive selection
\end{itemize}

can violate the assumptions that justified the original coverage
statement.

This boundary passes directly to the downstream Shift chapter.

\section{Conditional coverage is
stronger}\label{conditional-coverage-is-stronger}

Marginal coverage averages over the input distribution.

A stronger statement would require coverage within every relevant
subgroup or condition.

Generic finite-sample conformal validity does not automatically give
arbitrary exact conditional coverage.

The Atlas will not rewrite a marginal theorem as a subgroup theorem.

\section{Abstention is another control
surface}\label{abstention-is-another-control-surface}

A model need not predict on every example.

Let

\[
A(X)\in\{0,1\}
\]

indicate acceptance.

Coverage is

\[
c
=
P(A(X)=1).
\]

For loss L, selective risk is

\[
R_{sel}
=
E[L\mid A(X)=1].
\]

Selective classification studies this risk-coverage tradeoff; Geifman
and El-Yaniv develop a deep-network construction in this setting
(Geifman and El-Yaniv 2017).

\section{Exact risk-coverage
witness}\label{exact-risk-coverage-witness}

Suppose five equally weighted examples have error indicators

\[
(0,0,0,1,1).
\]

If we accept all five,

\[
c=1,
\]

and

\[
R_{sel}
=
\frac25.
\]

If we accept only the first three,

\[
c=\frac35,
\]

and

\[
R_{sel}=0.
\]

The accepted predictions are perfect.

But forty percent of the population has been rejected.

The improvement is real and conditional on lower coverage.

\section{Abstention is not
calibration}\label{abstention-is-not-calibration}

A selective classifier can reject difficult cases using:

\begin{itemize}
\tightlist
\item
  confidence;
\item
  margin;
\item
  entropy;
\item
  an auxiliary selector;
\item
  a calibrated risk predictor;
\item
  another uncertainty score.
\end{itemize}

Whatever the mechanism, low selective risk does not prove the underlying
probabilities are calibrated.

Conversely, a calibrated predictor need not know which examples to
reject to achieve a desired risk-coverage curve.

The objects remain separate.

\section{Uncertainty and decisions}\label{uncertainty-and-decisions}

Uncertainty matters because decisions have asymmetric costs.

A probability estimate may feed:

\begin{itemize}
\tightlist
\item
  abstention;
\item
  escalation to a human;
\item
  additional sensing;
\item
  active learning;
\item
  a larger conformal set;
\item
  a safer controller;
\item
  deferred execution.
\end{itemize}

But the correct action depends on a loss or utility model.

High entropy does not mathematically imply ``abstain.''

That policy requires a decision rule.

\section{Calibration can be local or global in different
senses}\label{calibration-can-be-local-or-global-in-different-senses}

A single scalar calibration metric can hide structure.

A model may appear calibrated overall while failing on:

\begin{itemize}
\tightlist
\item
  a demographic subgroup;
\item
  a rare class;
\item
  a region of feature space;
\item
  long-context examples;
\item
  high-loss examples;
\item
  a shifted domain.
\end{itemize}

The fix is not to call calibration useless.

It is to state the conditioning structure being tested.

\section{The model can be calibrated for the wrong
world}\label{the-model-can-be-calibrated-for-the-wrong-world}

Suppose calibration is established under distribution P.

Deployment occurs under distribution Q.

Even if the scoring function is unchanged,

\[
P(Y=1\mid S=s)
\]

and

\[
Q(Y=1\mid S=s)
\]

need not agree.

Therefore:

\[
\boxed{
P\text{-calibration}
\not\Rightarrow
Q\text{-calibration}.
}
\]

The same caution applies to:

\begin{itemize}
\tightlist
\item
  ensemble diagnostics;
\item
  conformal coverage;
\item
  selective-risk estimates;
\item
  abstention thresholds.
\end{itemize}

\section{The uncertainty interface}\label{the-uncertainty-interface}

A downstream component should be able to ask:

\begin{enumerate}
\def\labelenumi{\arabic{enumi}.}
\tightlist
\item
  What is random?
\item
  What is conditioned on?
\item
  Is the output a probability, variance, entropy, set, interval, or
  decision?
\item
  Is the quantity predictive or parameter/model-level?
\item
  Is the method Bayesian, approximately Bayesian, ensemble-based,
  conformal, or purely empirical?
\item
  What distributional assumptions support the guarantee?
\item
  Is the guarantee marginal or conditional?
\item
  Is abstention changing the evaluated population?
\item
  Was calibration measured or proved?
\item
  Does deployment match the law under which the statement was
  established?
\end{enumerate}

This is the chapter's practical contract.

\section{Handoff to Distribution Shift and
Robustness}\label{handoff-to-distribution-shift-and-robustness}

The direct consumer is

ATLAS-CH-SHIFT-001.

SHIFT may inherit:

\begin{itemize}
\tightlist
\item
  population calibration definitions;
\item
  the calibration-versus-sharpness witness;
\item
  the predictive entropy decomposition;
\item
  the same-entropy/opposite-decomposition counterexample;
\item
  source-scoped MC-dropout and ensemble semantics;
\item
  conformal rank coverage and its exchangeability boundary;
\item
  selective coverage and selective risk.
\end{itemize}

SHIFT must independently establish what survives when the data law
changes.

In particular it must not assume:

\[
\text{i.i.d. calibration}
\Rightarrow
\text{shifted calibration},
\]

or

\[
\text{exchangeable conformal coverage}
\Rightarrow
\text{coverage under arbitrary drift}.
\]

\section{Failure modes}\label{failure-modes-16}

The most common category errors are:

\begin{itemize}
\tightlist
\item
  calling confidence calibration ``uncertainty estimation'' without
  saying what is uncertain;
\item
  calling predictive entropy epistemic uncertainty;
\item
  treating ensemble members as posterior samples by definition;
\item
  treating stochastic dropout passes as exact Bayesian posterior
  samples;
\item
  treating a conformal set as a Bayesian credible region;
\item
  treating marginal coverage as arbitrary conditional coverage;
\item
  treating low ECE as population calibration;
\item
  treating low selective risk as unchanged-population accuracy;
\item
  treating abstention as calibration;
\item
  transporting an in-distribution guarantee into a shifted environment
  without revalidation.
\end{itemize}

\section{Closing view}\label{closing-view-32}

Uncertainty is not one number.

The Atlas uses a family of objects:

\[
\boxed{
\begin{array}{l}
\text{probability calibration},\\
\text{predictive entropy},\\
\text{model/posterior uncertainty},\\
\text{aleatoric uncertainty},\\
\text{ensemble disagreement},\\
\text{conformal coverage},\\
\text{selective risk and coverage}.
\end{array}
}
\]

Each object comes with its own assumptions and its own question.

The durable rule is:

\begin{quote}
Before using an uncertainty estimate, name what it estimates and name
the guarantee it actually carries.
\end{quote}

\section{References used in this
chapter}\label{references-used-in-this-chapter-61}

\begin{itemize}
\tightlist
\item
  (Guo et al. 2017)
\item
  (Gal and Ghahramani 2016)
\item
  (Kendall and Gal 2017)
\item
  (Lakshminarayanan et al. 2017)
\item
  (Vovk et al. 2005)
\item
  (Lei et al. 2018)
\item
  (Geifman and El-Yaniv 2017)
\end{itemize}

Exact provenance and claim boundaries are locked in:

sources/source-locks/ATLAS-CH-UNCERTAINTY-001.yaml

\chapter{Distribution Shift and
Robustness}\label{distribution-shift-and-robustness}

\textbf{Epistemic status:} audited Uncertainty substrate + primary
covariate-shift and adversarial-robust-optimization sources +
Atlas-owned exact finite witnesses.\\
\textbf{Specification:}
manuscript/specifications/ATLAS-CH-SHIFT-001.md\\
\textbf{Derivation packet:}
mathematics/derivations/ATLAS-CH-SHIFT-001-DERIVATIONS.md\\
\textbf{Computational witness:}
mathematics/computational-witnesses/ATLAS-CW-SHIFT-001.md\\
\textbf{Source lock:} sources/source-locks/ATLAS-CH-SHIFT-001.yaml

A model can remain numerically unchanged while the world in which it is
evaluated changes.

That forces a discipline:

\[
\boxed{\text{name the source law and deployment law before claiming robustness under shift.}}
\]

\section{Source law is not deployment
law}\label{source-law-is-not-deployment-law}

Let P be the law under which a guarantee, diagnostic, or empirical
estimate was established. Let Q be the deployment law. For predictor f
and loss L,

\[
R_P(f)=\mathbb E_P[L(f(X),Y)],
\]

\[
R_Q(f)=\mathbb E_Q[L(f(X),Y)].
\]

These are different mathematical objects unless P=Q or a bridge is
proved.

The Uncertainty chapter established that calibration, conformal
coverage, selective risk, predictive entropy, and other uncertainty
objects are assumption-scoped. SHIFT adds the changed-law layer.

\section{Distribution shift is not one
phenomenon}\label{distribution-shift-is-not-one-phenomenon}

\subsection{Covariate shift}\label{covariate-shift}

\[
P_X\neq Q_X,\qquad P(Y\mid X)=Q(Y\mid X).
\]

Sugiyama, Krauledat, and Müller study this setting and
importance-weighted correction/model-selection methods (Sugiyama et al.
2007).

\subsection{Label/prior shift}\label{labelprior-shift}

The class prior changes while an appropriate class-conditional mechanism
is held fixed.

\subsection{Concept or conditional
shift}\label{concept-or-conditional-shift}

\[
P(Y\mid X)\neq Q(Y\mid X).
\]

Changing only X-weights cannot generally repair this.

\subsection{Structural shift}\label{structural-shift}

A mechanism, topology, measurement process, routing rule, tokenizer,
feature generator, or other declared system component changes.
Structural shift is not identified merely because output statistics
changed.

\section{Covariate shift permits a change of
measure}\label{covariate-shift-permits-a-change-of-measure}

If

\[
P(Y\mid X)=Q(Y\mid X)
\]

and

\[
Q_X\ll P_X,
\]

define

\[
w(x)=\frac{dQ_X}{dP_X}(x).
\]

Then

\[
R_Q(f)=\mathbb E_P[w(X)L(f(X),Y)].
\]

This is not extrapolation. It requires the relevant target support to
lie within source support.

\section{Target-only support is a real
boundary}\label{target-only-support-is-a-real-boundary}

If some region has P\_X(x)=0 and Q\_X(x)\textgreater0, the ordinary
importance ratio is unavailable there. No finite reweighting of observed
source examples creates evidence about a target-only region. Extra
modeling assumptions may permit extrapolation, but they must be
declared.

\section{Exact witness: calibration can fail under pure covariate
shift}\label{exact-witness-calibration-can-fail-under-pure-covariate-shift}

Let X have two states a and b. Outcomes are deterministic:

\[
Y(a)=1,\qquad Y(b)=0.
\]

Under the source law,

\[
P_X(a)=P_X(b)=1/2.
\]

Under deployment,

\[
Q_X(a)=3/4,\qquad Q_X(b)=1/4.
\]

Nothing about Y\textbar X changed.

Use the same predictor before and after deployment:

\[
s(a)=s(b)=1/2.
\]

Under P, the event frequency among all points receiving score 1/2 is
1/2, so the sole score level is perfectly calibrated.

Under Q, the event frequency is 3/4 while the predictor still emits 1/2.
The deployment calibration gap is

\[
\boxed{1/4}.
\]

The model did not change. The conditional outcome mechanism did not
change. Only the covariate mixture changed.

\section{Calibration failure need not mean every risk
changed}\label{calibration-failure-need-not-mean-every-risk-changed}

For the same predictor,

\[
(1/2-Y)^2=1/4
\]

for either label. Thus Brier risk is

\[
\boxed{1/4}
\]

under both P and Q.

So calibration deterioration and generic predictive-risk deterioration
are not the same statement.

\section{Exact importance weights}\label{exact-importance-weights}

For that shift,

\[
w(a)=\frac{3/4}{1/2}=\frac32,
\]

\[
w(b)=\frac{1/4}{1/2}=\frac12.
\]

And

\[
\frac12\frac32+\frac12\frac12=1.
\]

The reweighted source event rate is

\[
\frac12\frac32=\frac34,
\]

which exactly recovers the deployment event rate.

\section{A detected marginal shift need not be a performance
failure}\label{a-detected-marginal-shift-need-not-be-a-performance-failure}

Let X in \{u,v\}, with Y(u)=0 and Y(v)=1. Use exact predictor f(u)=0,
f(v)=1.

Let

\[
P_X=(1/2,1/2)
\]

and

\[
Q_X=(9/10,1/10).
\]

The marginal laws differ substantially:

\[
\operatorname{TV}(P_X,Q_X)=2/5.
\]

Yet the predictor is right on every input under both laws. Therefore

\[
\boxed{R_P^{0/1}=R_Q^{0/1}=0}.
\]

A shift detector can be correct that the input marginal changed while a
task-failure alarm would be wrong.

\section{Concept shift is
different}\label{concept-shift-is-different}

Keep the same X-marginal, but reverse the conditional labels under
deployment. Then

\[
P_X=Q_X
\]

while

\[
P(Y\mid X)\neq Q(Y\mid X).
\]

The source-perfect predictor now has target error one. The covariate
density ratio is identically one, so X-reweighting does nothing.

This is why distribution shift is too coarse a diagnosis.

\section{Calibration, coverage, and selective risk each have their
own transfer
problem}\label{calibration-coverage-and-selective-risk-each-have-their-own-transfer-problem}

A guarantee under P can depend on exchangeability, score distribution,
subgroup mixture, acceptance threshold, calibration level sets, or loss
distribution. Changing P to Q can disturb these in different ways.

Therefore:

\[
P\text{-calibration}\not\Rightarrow Q\text{-calibration},
\]

\[
P\text{-coverage}\not\Rightarrow Q\text{-coverage},
\]

and

\[
P\text{-selective-risk control}\not\Rightarrow Q\text{-selective-risk control}.
\]

Each transfer needs its own assumptions.

\section{Adversarial robustness asks a different
question}\label{adversarial-robustness-asks-a-different-question}

Average-case target risk is

\[
R_Q(f)=\mathbb E_Q[L(f(X),Y)].
\]

Adversarial risk instead has the form

\[
R_{\rm adv}(f)=\mathbb E_P\left[\sup_{x'\in\Delta(X)}L(f(x'),Y)\right].
\]

Madry et al.~study adversarial robustness through robust optimization
against a specified perturbation/adversary class (Madry et al. 2018).

The supremum is the important difference. A deployment distribution and
a worst-case perturbation set are not interchangeable objects.

\section{Exact adversarial
separation}\label{exact-adversarial-separation}

Take X in \{-1,+1\}, with labels matching sign, and a perfect predictor.
Clean risk is zero.

Declare

\[
\Delta(-1)=\Delta(+1)=\{-1,+1\}.
\]

For either original example, the adversary can choose the opposite input
while the original label is retained. Thus worst-case loss is one at
every source point:

\[
\boxed{R_{\rm adv}=1}.
\]

So R\_P=0 can coexist with R\_adv=1.

\section{Robustness is perturbation-set
relative}\label{robustness-is-perturbation-set-relative}

If instead Delta(x)=\{x\}, adversarial risk equals clean risk. The model
did not change; only the allowed perturbation set changed. Therefore
``robust'' is incomplete without the threat/perturbation model.

\section{Robust optimization}\label{robust-optimization}

A robust-training objective is schematically

\[
\min_\theta\mathbb E_P\left[\sup_{\delta\in\Delta(X)}L(f_\theta(X+\delta),Y)\right].
\]

The objective is meaningful only with a declared data law, perturbation
set, loss, and optimization approximation. Robustness to one
perturbation family is not automatically robustness to another.

\section{Structural sensitivity is not an adversarial
synonym}\label{structural-sensitivity-is-not-an-adversarial-synonym}

Suppose a system component changes: tokenizer, routing topology,
communication graph, measurement process, sensor transform,
feature-generation mechanism, or intervention policy.

A structural-sensitivity analysis must name that component and compare
behavior across the change. A change in predictive statistics can signal
changed behavior but cannot by itself identify which mechanism changed.

\section{Average-case shift is not worst-case
shift}\label{average-case-shift-is-not-worst-case-shift}

A law Q assigns probability mass. An adversarial set Delta(x) specifies
possible worst-case alternatives around a point.

One integrates:

\[
\mathbb E_Q[\cdot].
\]

The other maximizes:

\[
\sup_{\delta\in\Delta(x)}[\cdot].
\]

Neither operation silently replaces the other.

\section{Shift magnitude is
metric-relative}\label{shift-magnitude-is-metric-relative}

Even ``large shift'' requires a discrepancy: total variation, KL
divergence, Wasserstein distance, feature discrepancy, calibration
change, risk change, or structural edit magnitude. These can rank
changes differently.

The benign marginal-shift witness has nonzero total variation but zero
task-risk change.

\section{Support overlap determines what reweighting can
say}\label{support-overlap-determines-what-reweighting-can-say}

Importance weighting is strongest when target mass is represented in
source support. Poor overlap creates large weights and unstable
estimation. Zero overlap creates a literal support failure for ordinary
density-ratio correction.

The support condition is therefore part of the claim.

\section{The unchanged model can still face a changed
theorem}\label{the-unchanged-model-can-still-face-a-changed-theorem}

When deployment law changes, model parameters, architecture, scoring
rule, threshold, and code may all remain fixed. Yet the guarantee can
change because its probability law changed.

Governance must bind guarantees to laws, not only model artifacts.

\section{Metric firewall}\label{metric-firewall}

SHIFT keeps separate calibration, conformal/marginal coverage, selective
risk, source predictive risk, deployment predictive risk, adversarial
risk, and structural sensitivity. Every experiment must state which
object it measured.

\section{What SHIFT does not
license}\label{what-shift-does-not-license}

It does not license:

\begin{itemize}
\tightlist
\item
  shift detected implies model failed;
\item
  source calibration implies deployment calibration;
\item
  good target-average risk implies adversarial robustness;
\item
  changed outputs imply an identified structural cause.
\end{itemize}

\section{Operational checklist}\label{operational-checklist}

For every deployment-shift claim: 1. name P; 2. name Q; 3. name the
changed marginal, conditional, or mechanism; 4. name the metric; 5.
check support; 6. state average-case versus worst-case semantics; 7.
separate detection from performance; 8. separate predictive failure from
causal diagnosis; 9. replay an exact or empirical witness; 10. state
what remains unproved.

\section{Durable rule}\label{durable-rule}

\[
\boxed{\text{robustness is always relative to a law, metric, and perturbation/mechanism class.}}
\]

A detected change is evidence of change. It is not by itself evidence of
task failure, calibration failure, adversarial vulnerability, or
structural cause.

\section{References used in this
chapter}\label{references-used-in-this-chapter-62}

\begin{itemize}
\tightlist
\item
  (Sugiyama et al. 2007)
\item
  (Madry et al. 2018)
\end{itemize}

\chapter{Interlude: From Readability to Functional
Evidence}\label{interlude-from-readability-to-functional-evidence}

\textbf{Epistemic status:} audited prerequisites + exact finite witness
+ source-scoped synthesis

A representation may carry a quantity that is easy to read without the
declared output map depending on that quantity. The chapter therefore
separates readout success from functional dependence.

\section{Exact finite separation}\label{exact-finite-separation}

For x in \{-1,+1\}, let h(x)=(x,x) and F(h)=h\_1. Both coordinates
recover x exactly. Removing the second coordinate leaves F unchanged,
while removing the first changes the output. Starting from h(-x),
replacing coordinate 1 with its clean value restores x; replacing
coordinate 2 does not.

This proves a finite separation between perfect readability and use by
the declared output map.

\section{Diagnostic record}\label{diagnostic-record}

A strong component-level result should state the target behavior B,
selected component set S, tested transformation T, metric M, and
reference rule R:

D=(B,S,T,M,R).

The same observed effect can support different interpretations if any of
these fields changes.

\section{Redundancy boundary}\label{redundancy-boundary}

A null one-component result does not prove that no relevant
representation exists. Redundant coordinates can each support the same
output while neither is individually necessary under a single removal
test.

\section{Downstream handoff}\label{downstream-handoff-10}

ATLAS-CH-SPECTRALDIAG-001 may inherit the readability/function
distinction, the diagnostic record, the redundancy warning, and the
exact witness. Any proposed spectral signature still requires its own
functional evidence.

\section{References used in this
chapter}\label{references-used-in-this-chapter-63}

Primary references and exact prerequisite identities are recorded in:

sources/source-locks/ATLAS-CH-MECHDIAG-001.yaml

\chapter{Spectral and Operator
Diagnostics}\label{spectral-and-operator-diagnostics}

\textbf{Epistemic status:} audited non-normal/operator substrate +
audited mechanistic-diagnostics substrate + one primary Koopman source +
Atlas-owned exact finite controls.\\
\textbf{Specification:}
manuscript/specifications/ATLAS-CH-SPECTRALDIAG-001.md\\
\textbf{Derivation packet:}
mathematics/derivations/ATLAS-CH-SPECTRALDIAG-001-DERIVATIONS.md\\
\textbf{Computational witness:}
mathematics/computational-witnesses/ATLAS-CW-SPECTRALDIAG-001.md\\
\textbf{Source lock:}
sources/source-locks/ATLAS-CH-SPECTRALDIAG-001.yaml

A spectrum is never just ``the spectrum.''

It is the spectrum of a particular object:

\begin{itemize}
\tightlist
\item
  a weight matrix;
\item
  a transition matrix;
\item
  a Jacobian at a reference state;
\item
  a Hessian at a parameter point;
\item
  a Koopman operator on a declared observable space;
\item
  a finite empirical approximation to one of these.
\end{itemize}

The governing rule is:

\[
\boxed{\text{name the operator, reference point, and metric before interpreting the diagnostic.}}
\]

\section{Different spectral objects answer different
questions}\label{different-spectral-objects-answer-different-questions}

For a matrix A, eigenvalues summarize invariant linear modes in one
sense.

Singular values summarize Euclidean amplification geometry.

Pseudospectra and resolvent norms expose sensitivity and non-normal
behavior that eigenvalues alone may hide.

A Jacobian spectrum is local to a state.

A Hessian spectrum is local to a parameter point and objective.

A Koopman spectrum belongs to an operator acting on observables.

These are not interchangeable.

\section{The non-normal boundary}\label{the-non-normal-boundary}

The audited NONNORMAL chapter already establishes the crucial warning:

\[
\rho(A)<1
\]

need not prevent finite-horizon amplification when A is non-normal.

Therefore an eigenvalue diagnostic cannot silently substitute for
transient-response or pseudospectral diagnostics.

\section{Exact equal-eigenvalue
control}\label{exact-equal-eigenvalue-control}

Take

\[
D=\begin{pmatrix}1/2&0\\0&1/2\end{pmatrix},
\qquad
N=\begin{pmatrix}1/2&2\\0&1/2\end{pmatrix}.
\]

Both have eigenvalue multiset

\[
\{1/2,1/2\}.
\]

For

\[
e_2=(0,1)^\top,
\]

\[
\|De_2\|_2^2=1/4,
\]

while

\[
\|Ne_2\|_2^2=17/4.
\]

So equal eigenvalues coexist with sharply different finite response.

The diagnostic lesson is simple:

\[
\boxed{\text{eigenvalue equality}\not\Rightarrow\text{response equality}}.
\]

\section{Singular values are different from
eigenvalues}\label{singular-values-are-different-from-eigenvalues}

The singular values of A are the square roots of the eigenvalues of

\[
A^*A.
\]

They describe Euclidean input-output amplification.

For non-normal systems they can be much more directly related to
one-step gain than eigenvalues are.

But singular values still discard orientation information.

\section{Equal singular values can hide interface
differences}\label{equal-singular-values-can-hide-interface-differences}

Let

\[
A=\operatorname{diag}(2,1/2),
\qquad
B=\operatorname{diag}(1/2,2).
\]

Both have the same eigenvalue multiset and the same singular-value
multiset:

\[
\{2,1/2\}.
\]

Yet for a fixed interface vector

\[
e_1=(1,0)^\top,
\]

\[
\|Ae_1\|_2=2,
\]

while

\[
\|Be_1\|_2=1/2.
\]

The difference is not in the global singular-value list.

It is in relative position.

\section{Relative-position
diagnostics}\label{relative-position-diagnostics}

A system often interacts with fixed subspaces, readouts, token
directions, residual channels, control inputs, or protected interfaces.

A useful diagnostic may therefore involve quantities such as

\[
\|Av\|,
\]

principal angles between singular subspaces and a declared interface
subspace, or projected gains such as

\[
\|P_Y A P_X\|.
\]

These are not purely spectral invariants of A.

They are operator-plus-interface diagnostics.

\section{Zero spectral drift need not mean zero behavioral
drift}\label{zero-spectral-drift-need-not-mean-zero-behavioral-drift}

Move from A to B in the previous example.

The sorted eigenvalue multiset is unchanged.

The sorted singular-value multiset is unchanged.

So a drift score using only either multiset is zero.

But the response to e1 changes from 2 to 1/2.

Thus

\[
\boxed{\text{zero spectral-summary drift}\not\Rightarrow\text{zero interface drift}}.
\]

\section{Pseudospectra remain a separate diagnostic
family}\label{pseudospectra-remain-a-separate-diagnostic-family}

Under the inherited 2-norm convention,

\[
\Lambda_\varepsilon(A)
=\{z:\sigma_{\min}(zI-A)\le\varepsilon\}.
\]

This captures resolvent growth and spectral sensitivity that the bare
eigenvalue set may miss.

But a pseudospectrum is still descriptive evidence about a declared
operator.

It does not by itself identify which component of a learned system is
functionally necessary.

\section{Jacobian spectra are
local}\label{jacobian-spectra-are-local}

For nonlinear map

\[
z^+=\Phi(z),
\]

a Jacobian diagnostic at z0 is based on

\[
J_\Phi(z_0).
\]

The reference state is part of the object.

Removing z0 from the report removes part of the diagnostic definition.

\section{Same local Jacobian, different nonlinear
maps}\label{same-local-jacobian-different-nonlinear-maps}

Let

\[
F(x)=x/2,
\]

and

\[
G(x)=x/2+x^2.
\]

At zero,

\[
F'(0)=G'(0)=1/2.
\]

But

\[
F(1/2)=1/4,
\]

while

\[
G(1/2)=1/2.
\]

Therefore a local Jacobian spectrum cannot be promoted into a global
nonlinear theorem.

\section{Transition-local
signatures}\label{transition-local-signatures}

If a system moves through states

\[
z_0,z_1,z_2,\ldots,
\]

then one may study

\[
J_t=D\Phi(z_t).
\]

A spectral sequence

\[
\sigma(J_t)
\]

is a transition-local diagnostic.

Its interpretation requires the state/time index, estimation window, and
operator-estimation method.

\section{Hessian spectra are also local and
incomplete}\label{hessian-spectra-are-also-local-and-incomplete}

For objective L(theta), the Hessian

\[
H_L(\theta_0)
\]

captures second-order curvature at a declared parameter point.

It does not by itself determine the gradient.

\section{Exact Hessian control}\label{exact-hessian-control}

Let

\[
f(x,y)=x^2+y^2,
\]

\[
g(x,y)=x^2+y^2+x.
\]

Both have Hessian

\[
2I
\]

and therefore the same Hessian spectrum

\[
\{2,2\}.
\]

But at the origin

\[
\nabla f=(0,0),
\]

while

\[
\nabla g=(1,0).
\]

The same Hessian spectrum can occur at a stationary point and a
non-stationary point.

\section{Koopman changes the object being
diagonalized}\label{koopman-changes-the-object-being-diagonalized}

For nonlinear state map T and observable h, the Koopman operator acts as

\[
(Uh)(x)=h(T(x)).
\]

It is linear as an operator on observables even when T is nonlinear.

Mezić develops this spectral viewpoint for dynamical systems and model
reduction (Mezić 2005).

The important semantic point is:

\[
\boxed{\text{Koopman spectrum is an observable-operator spectrum.}}
\]

It is not the Jacobian spectrum of T by definition.

\section{Exact two-state Koopman
witness}\label{exact-two-state-koopman-witness}

Let

\[
T(0)=1,
\qquad
T(1)=0.
\]

On indicator observables delta0 and delta1,

\[
U\delta_0=\delta_1,
\qquad
U\delta_1=\delta_0.
\]

So

\[
U=\begin{pmatrix}0&1\\1&0\end{pmatrix}.
\]

Hence

\[
U^2=I,
\]

with trace zero and determinant -1, giving eigenvalues

\[
\boxed{\{1,-1\}}.
\]

\section{Finite Koopman approximations need their own error
language}\label{finite-koopman-approximations-need-their-own-error-language}

In practice one may estimate a finite matrix from data and call it a
Koopman approximation.

That introduces choices:

\begin{itemize}
\tightlist
\item
  observable dictionary;
\item
  sampling distribution;
\item
  time window;
\item
  truncation/rank;
\item
  regression method;
\item
  regularization;
\item
  finite precision.
\end{itemize}

The spectrum of that finite matrix is exactly the spectrum of the fitted
matrix.

Its relation to an underlying infinite-dimensional Koopman operator
requires additional approximation theory.

\section{Spectral drift is
definition-relative}\label{spectral-drift-is-definition-relative}

A drift diagnostic can compare

\begin{itemize}
\tightlist
\item
  eigenvalue sets;
\item
  singular-value lists;
\item
  pseudospectral contours;
\item
  resolvent norms;
\item
  subspace angles;
\item
  projected gains;
\item
  local Jacobian spectra;
\item
  Koopman estimates.
\end{itemize}

These are different drift objects.

A single scalar named ``spectral drift'' is incomplete unless the
underlying object and discrepancy are declared.

\section{Descriptive quality is not functional
necessity}\label{descriptive-quality-is-not-functional-necessity}

The audited MECHDIAG packet supplies a second firewall.

A probe or diagnostic can track a property without showing that the
measured feature is necessary for the system's output.

The same applies to spectral signatures.

A signature can correlate strongly with performance and still be
epiphenomenal.

\section{Predictive utility is not
mechanism}\label{predictive-utility-is-not-mechanism}

Suppose a spectral statistic predicts failure one step before failure
occurs.

That is useful predictive evidence.

It does not yet show that manipulating the associated eigenspace,
singular direction, or pseudospectral region will change the failure.

Prediction and mechanism are different claims.

\section{Mechanistic significance requires
intervention}\label{mechanistic-significance-requires-intervention}

To promote a spectral signature toward mechanism, pair it with a
functional test such as:

\begin{itemize}
\tightlist
\item
  ablation of the implicated subspace;
\item
  controlled perturbation of the operator block;
\item
  substitution of a matched spectral control;
\item
  recovery/rescue after restoring the implicated component;
\item
  another explicitly justified intervention.
\end{itemize}

The intervention must target the object that the spectral interpretation
claims matters.

\section{A matched-spectrum control is especially
valuable}\label{a-matched-spectrum-control-is-especially-valuable}

The A/B control above has the same eigenvalue and singular-value
multisets but different fixed-interface response.

This is useful because it holds the global spectral summaries fixed
while changing a functional interface relation.

Matched controls can therefore falsify over-strong spectral
explanations.

\section{Finite precision and
conditioning}\label{finite-precision-and-conditioning}

Empirical spectra can be sensitive to:

\begin{itemize}
\tightlist
\item
  roundoff;
\item
  non-normality;
\item
  ill-conditioning;
\item
  nearly repeated eigenvalues;
\item
  low-rank truncation;
\item
  finite-sample noise.
\end{itemize}

A plotted spectrum should therefore carry enough numerical context to
distinguish exact structure from estimated structure.

\section{No universal scalar
diagnostic}\label{no-universal-scalar-diagnostic}

Different systems can fail for different reasons.

A scalar that is useful for one operator family can become blind under
another family or another interface.

The exact controls already show this for eigenvalue and singular-value
summaries.

Thus the chapter rejects:

\[
\text{one spectral scalar}\equiv\text{universal health metric}.
\]

\section{Reporting protocol}\label{reporting-protocol}

For every spectral/operator diagnostic report:

\begin{enumerate}
\def\labelenumi{\arabic{enumi}.}
\tightlist
\item
  operator/object identity;
\item
  reference state, parameter point, or time window;
\item
  spectrum/diagnostic type;
\item
  norm and numerical convention;
\item
  estimation/truncation method;
\item
  uncertainty/tolerance;
\item
  interface or relative-position object, if used;
\item
  descriptive versus predictive versus functional claim class;
\item
  intervention evidence, if mechanistic significance is claimed;
\item
  known non-implications.
\end{enumerate}

\section{Durable non-implications}\label{durable-non-implications-3}

\[
\text{same eigenvalues}\not\Rightarrow\text{same transient behavior},
\]

\[
\text{same singular values}\not\Rightarrow\text{same interface behavior},
\]

\[
\text{same local Jacobian spectrum}\not\Rightarrow\text{same nonlinear map},
\]

\[
\text{same Hessian spectrum}\not\Rightarrow\text{same stationarity status},
\]

\[
\text{spectral correlation}\not\Rightarrow\text{functional necessity},
\]

and

\[
\text{functional association}\not\Rightarrow\text{causal mechanism without a suitable intervention}.
\]

\section{Closing view}\label{closing-view-33}

Spectral diagnostics are powerful precisely because they expose
structure that scalar performance metrics can miss.

They become misleading when their object is unnamed or their epistemic
role is inflated.

The durable rule is:

\[
\boxed{\text{a spectral signature is evidence about a declared operator; mechanistic meaning requires a separate functional bridge.}}
\]

\section{References used in this
chapter}\label{references-used-in-this-chapter-64}

\begin{itemize}
\tightlist
\item
  (Horn and Johnson 2012)
\item
  (Trefethen and Embree 2005)
\item
  (Mezić 2005)
\end{itemize}

Mechanistic-diagnostics sources and exact prerequisite identities are
recorded in sources/source-locks/ATLAS-CH-SPECTRALDIAG-001.yaml.

\chapter{Compression and Description
Length}\label{compression-and-description-length}

\textbf{Epistemic status:} mathematical and engineering exposition built
from audited Information and Representation prerequisites, classical
compression sources, and Atlas-owned exact witnesses.\\
\textbf{Derivation packet:}
`mathematics/derivations/ATLAS-CH-COMPRESS-001-DERIVATIONS.md`\\
\textbf{Computational witness:}
`mathematics/computational-witnesses/ATLAS-CW-COMPRESS-001.md`\\
\textbf{Source lock:} `sources/source-locks/ATLAS-CH-COMPRESS-001.yaml`

\section{Compression begins with a
code}\label{compression-begins-with-a-code}

``Smaller'' can mean several different things.

A model can use:

\begin{itemize}
\tightlist
\item
  fewer scalar parameters;
\item
  fewer nonzero parameters;
\item
  fewer distinct parameter values;
\item
  fewer bits per value;
\item
  a lower-rank factorization;
\item
  a shorter entropy-coded file;
\item
  a smaller student network;
\item
  a shorter formal description under a declared code.
\end{itemize}

These are not identical quantities.

The chapter therefore begins with a rule:

\[
\boxed{
\text{compression claim}
\Rightarrow
\text{declare the object, code, and tolerated distortion}.
}
\]

Without those declarations, a compression ratio can be numerically
correct and scientifically ambiguous.

\section{Description length is not parameter
count}\label{description-length-is-not-parameter-count}

Minimum-description-length thinking treats model selection as a coding
problem (Rissanen 1978).

For a declared code \(C\), model \(M\), and data \(D\), a two-part
objective has the form

\[
L_C(M,D)
=
L_C(M)
+
L_C(D\mid M).
\]

The first term describes the model.

The second describes the data given that model.

A shorter model can be worse if it leaves a much longer residual/data
description.

A larger model can be justified if it dramatically shortens the
unexplained data.

This is richer than ``count the parameters.''

\section{A finite MDL witness}\label{a-finite-mdl-witness}

Suppose two hypotheses have declared model-code lengths

\[
L(H_A)=3,
\qquad
L(H_B)=7.
\]

Suppose they encode the observed data equally well:

\[
L(D\mid H_A)
=
L(D\mid H_B)
=
5.
\]

Then

\[
L(H_A,D)=8,
\]

while

\[
L(H_B,D)=12.
\]

Under this declared code, \(H_A\) wins.

Nothing in the arithmetic says that 3 bits and 7 bits are universally
correct descriptions.

The point is that the code belongs inside the claim.

\section{Lossless and lossy are different
questions}\label{lossless-and-lossy-are-different-questions}

Lossless compression preserves the coded object exactly.

Lossy compression preserves it only up to a declared distortion.

For learned systems, several distortions are possible:

\begin{itemize}
\tightlist
\item
  parameter error;
\item
  operator error;
\item
  output error;
\item
  predictive loss;
\item
  calibration change;
\item
  retrieval degradation;
\item
  latency or energy change.
\end{itemize}

Two compression methods can have the same storage ratio and very
different task distortion.

So the meaningful object is not just

\[
\frac{\text{original size}}{\text{compressed size}}.
\]

It is the size--distortion pair.

\section{Exact low-rank
compression}\label{exact-low-rank-compression}

Consider the \(4\times4\) all-ones matrix

\[
W
=
\mathbf 1_4\mathbf 1_4^\top.
\]

Its rank is one.

Let

\[
u=v=(1,1,1,1)^\top.
\]

Then

\[
W=uv^\top.
\]

The dense matrix and the factors implement exactly the same linear map:

\[
Wx=u(v^\top x)
\]

for every \(x\in\mathbb R^4\).

This is exact structural compression, not approximation.

\section{Exact code-length
comparison}\label{exact-code-length-comparison}

Use a deliberately simple codec:

\begin{itemize}
\tightlist
\item
  the \(4\times4\) shape is known;
\item
  entries are binary;
\item
  one format bit distinguishes dense and rank-one representations.
\end{itemize}

Dense form:

\begin{itemize}
\tightlist
\item
  one format bit;
\item
  sixteen matrix-entry bits.
\end{itemize}

Total:

\[
17\text{ bits}.
\]

Rank-one form:

\begin{itemize}
\tightlist
\item
  one format bit;
\item
  four bits for \(u\);
\item
  four bits for \(v\).
\end{itemize}

Total:

\[
9\text{ bits}.
\]

The same exact function has two different description lengths under the
declared code.

That is the chapter's simplest exact compression witness.

\section{Low rank can also be
lossy}\label{low-rank-can-also-be-lossy}

Now take

\[
A=\operatorname{diag}(4,3,1).
\]

Its singular values are

\[
4,\;3,\;1.
\]

By the classical low-rank approximation result (Eckart and Young 1936),
the best rank-one Frobenius approximation keeps only the largest
singular value:

\[
A_1
=
\operatorname{diag}(4,0,0).
\]

The squared error is

\[
\|A-A_1\|_F^2
=
3^2+1^2
=
10.
\]

The best rank-two approximation is

\[
A_2
=
\operatorname{diag}(4,3,0),
\]

with squared error

\[
1.
\]

Low-rank representation is therefore not synonymous with exact
compression.

It is exact only when discarded singular values are zero.

\section{Pruning removes
parameters}\label{pruning-removes-parameters}

Pruning changes the support of the parameterization.

A simple magnitude-pruning rule removes weights whose absolute values
fall below a threshold.

The appeal is obvious:

\begin{itemize}
\tightlist
\item
  fewer nonzero connections;
\item
  potentially less storage;
\item
  potentially less compute.
\end{itemize}

But magnitude is not the same thing as functional importance.

\section{A small weight can still
matter}\label{a-small-weight-can-still-matter}

Let

\[
w=
\begin{pmatrix}
1\\
1/8
\end{pmatrix}
\]

and

\[
x=
\begin{pmatrix}
0\\
8
\end{pmatrix}.
\]

Then

\[
w^\top x=1.
\]

Apply a pruning threshold

\[
\tau=1/4.
\]

The second weight disappears:

\[
w_{\rm pruned}
=
\begin{pmatrix}
1\\
0
\end{pmatrix}.
\]

Now

\[
w_{\rm pruned}^\top x=0.
\]

The small parameter was decisive on this input.

So:

\[
\boxed{
\text{small magnitude}
\not\Rightarrow
\text{functional irrelevance}.
}
\]

Neural pruning methods can work well in practice, including as part of
the Deep Compression pipeline (Han et al. 2016a).

Their success is empirical and distribution-dependent.

\section{Quantization changes
precision}\label{quantization-changes-precision}

Quantization answers a different question.

The connections can remain present while each stored value uses fewer
bits or is mapped to a smaller set of representable levels.

Take

\[
w=
(1/4,1/2,1/2,1/4).
\]

Use the fixed grid

\[
Q=
\{0,1/4,1/2,3/4\}.
\]

Every value of \(w\) lies exactly on the grid.

If raw values use 8 bits each, the payload is

\[
4\times8=32\text{ bits}.
\]

If fixed-grid indices use 2 bits each, the payload is

\[
4\times2=8\text{ bits}.
\]

Parameter distortion is zero in this witness.

The parameter count did not change.

The number of bits per parameter did.

\section{Codebook overhead matters}\label{codebook-overhead-matters}

The previous witness used a fixed grid known by the codec.

If the quantization levels are learned, the levels themselves must be
stored.

That overhead belongs in the accounting.

This is one reason ``4-bit model'' and ``4 bits per raw weight'' are not
necessarily the same storage claim.

Metadata, scales, zero points, codebooks, indices, alignment, and
container formats can matter.

\section{Weight sharing is different
again}\label{weight-sharing-is-different-again}

The same vector contains only two distinct values:

\[
1/4
\quad\text{and}\quad
1/2.
\]

Suppose a codec stores:

\begin{itemize}
\tightlist
\item
  two 8-bit centroids;
\item
  four one-bit centroid indices.
\end{itemize}

The payload is

\[
16+4=20\text{ bits}.
\]

That is shorter than four independent 8-bit values:

\[
32\text{ bits}.
\]

All four connections remain.

Pruning removed a connection.

Weight sharing did not.

\section{Deep Compression combines
mechanisms}\label{deep-compression-combines-mechanisms}

Han, Mao, and Dally combine pruning, trained quantization/weight
sharing, and Huffman coding in one pipeline (Han et al. 2016a).

Those stages attack different redundancies:

\begin{itemize}
\tightlist
\item
  pruning removes selected connections;
\item
  quantization reduces the set of distinct values;
\item
  weight sharing reuses centroids;
\item
  entropy coding exploits frequency structure in the resulting symbols.
\end{itemize}

Calling the whole pipeline ``compression'' is correct.

Treating all its stages as the same mathematical operation is not.

\section{Entropy coding acts on symbol
statistics}\label{entropy-coding-acts-on-symbol-statistics}

After a representation has been chosen, symbols may occur with unequal
frequencies.

A variable-length code can assign shorter codewords to more common
symbols.

This connects compression back to the Information chapter.

Entropy coding does not decide which neural connections are important.

It compresses a symbol stream under a declared statistical model.

\section{Low rank and sparsity are different
structures}\label{low-rank-and-sparsity-are-different-structures}

A matrix can be sparse and full rank.

A matrix can be dense and rank one.

These are different kinds of compressible structure.

Sparsity says many coordinates are zero.

Low rank says many rows or columns are dependent through a
lower-dimensional factorization.

The hardware consequences, approximation errors, and implementation
costs can differ.

\section{Structured
parameterizations}\label{structured-parameterizations}

Low rank and weight sharing are examples of replacing an unconstrained
parameter array with a smaller structured family.

For

\[
W\in\mathbb R^{m\times n},
\]

a dense matrix has

\[
mn
\]

scalar slots.

A rank-\(r\) factorization

\[
W=UV^\top
\]

with

\[
U\in\mathbb R^{m\times r},
\qquad
V\in\mathbb R^{n\times r}
\]

stores

\[
r(m+n)
\]

scalar slots before metadata.

When

\[
r(m+n)<mn,
\]

the factorized form uses fewer scalar slots.

That arithmetic alone does not prove lower wall-clock latency.

Kernel efficiency and memory traffic still matter.

\section{Distillation compresses behavior, not
bits}\label{distillation-compresses-behavior-not-bits}

Knowledge distillation takes a different route (Hinton et al. 2014).

A teacher or ensemble produces predictive targets.

A smaller student is trained to match those targets.

At temperature \(T\), one can form softened teacher probabilities

\[
p_T(i)
=
\frac{\exp(z_{T,i}/T)}
{\sum_j \exp(z_{T,j}/T)}
\]

and analogous student probabilities.

The student is optimized to reproduce teacher behavior under a declared
loss.

Nothing requires the student parameters to be a bit-level encoding of
the teacher parameters.

This is behavioral compression or transfer.

\section{Exact equality on a finite dataset is not global
equality}\label{exact-equality-on-a-finite-dataset-is-not-global-equality}

Suppose a student matches a teacher exactly on every example in a finite
training set.

That proves equality on that set.

It does not prove

\[
f_{\rm student}(x)
=
f_{\rm teacher}(x)
\]

for every possible input.

The relevant domain must be declared.

This is the Representation chapter's equivalence discipline applied to
compression.

\section{Compression ratio is incomplete without retained
behavior}\label{compression-ratio-is-incomplete-without-retained-behavior}

A 100x smaller artifact can be excellent or useless.

The ratio alone does not say.

A serious report pairs size with something like:

\begin{itemize}
\tightlist
\item
  output error;
\item
  predictive loss;
\item
  accuracy;
\item
  calibration;
\item
  robustness;
\item
  retrieval quality;
\item
  task-specific utility.
\end{itemize}

Compression is a tradeoff surface, not a one-number race.

\section{Parameter distortion can disagree with task
distortion}\label{parameter-distortion-can-disagree-with-task-distortion}

A tiny parameter perturbation can have a large functional effect in one
system.

A large parameter change can preserve function through
reparameterization in another.

Therefore

\[
\|\theta-\hat\theta\|
\]

and task degradation are different measurements.

The right distortion depends on the object being preserved.

\section{Compression and representation
equivalence}\label{compression-and-representation-equivalence}

Two encodings can represent the same function differently.

The rank-one witness makes this exact:

\[
W
\]

and

\[
uv^\top
\]

are different stored representations of the same linear operator.

This is why compression is naturally connected to representation theory.

We are looking for shorter encodings inside an equivalence class of
acceptable behavior.

The equivalence relation must be stated.

\section{MDL is not ``smallest network
wins''}\label{mdl-is-not-smallest-network-wins}

MDL balances model description and residual/data description.

A model that is too small can fit poorly and require a longer data
description.

A larger model can be justified if it compresses the data sufficiently.

The principle is not:

\begin{quote}
always pick fewer parameters.
\end{quote}

It is:

\begin{quote}
under a declared coding framework, minimize the relevant total
description.
\end{quote}

\section{Compression can reveal
redundancy}\label{compression-can-reveal-redundancy}

If a model can be compressed with little change in declared behavior,
then some part of its original description was redundant relative to
that behavior and codec.

That is a useful scientific observation.

But redundancy can come from many sources:

\begin{itemize}
\tightlist
\item
  duplicate parameters;
\item
  low-rank structure;
\item
  unused features;
\item
  overprecise numerical representation;
\item
  repeated values;
\item
  symmetries;
\item
  overparameterization;
\item
  training artifacts.
\end{itemize}

Compression alone does not tell us which mechanism caused the
redundancy.

\section{Compressibility is not
interpretability}\label{compressibility-is-not-interpretability}

A highly compressible model can still be opaque.

A concise program can implement a difficult-to-understand computation.

A sparse network can still have distributed causal structure.

A low-rank map can still be semantically uninterpretable.

Therefore:

\[
\boxed{
\text{short description}
\not\Rightarrow
\text{human-understandable mechanism}.
}
\]

That stronger inference belongs downstream, where compression will be
considered as a possible probe rather than a proof.

\section{Compressibility is not
intelligence}\label{compressibility-is-not-intelligence}

The same caution applies to intelligence.

Many simple objects compress extremely well.

Some useful computations require large descriptions under a given code.

A system's compressibility may provide evidence about reusable
structure.

It is not a scalar definition of intelligence.

\section{Failure modes}\label{failure-modes-17}

\subsection{Parameter count equals description
length}\label{parameter-count-equals-description-length}

False.

Precision, codebook structure, sparsity, entropy coding, and metadata
all matter.

\subsection{Pruning equals
quantization}\label{pruning-equals-quantization}

False.

One changes support; the other changes value representation.

\subsection{Quantization equals weight
sharing}\label{quantization-equals-weight-sharing}

Not necessarily.

A fixed quantization grid and a learned shared codebook have different
storage accounting.

\subsection{Low rank equals
sparsity}\label{low-rank-equals-sparsity}

False.

These are different structural constraints.

\subsection{Distillation equals source
coding}\label{distillation-equals-source-coding}

False.

Distillation transfers predictive behavior under a training objective.

\subsection{High compression ratio means good
compression}\label{high-compression-ratio-means-good-compression}

Incomplete.

Retained behavior or distortion must also be reported.

\subsection{Same outputs on one dataset means same
function}\label{same-outputs-on-one-dataset-means-same-function}

False outside the declared dataset.

\subsection{Compressible means
interpretable}\label{compressible-means-interpretable}

False without additional evidence.

\subsection{Compressible means
intelligent}\label{compressible-means-intelligent}

False.

That claim requires an independent theory.

\section{Downstream handoff}\label{downstream-handoff-11}

Compression as Discovery and Intelligence Probe may inherit:

\begin{itemize}
\tightlist
\item
  declared-code discipline;
\item
  exact-versus-lossy distinction;
\item
  mechanism separation among pruning, quantization, sharing, low rank,
  entropy coding, and distillation;
\item
  the requirement to report compression jointly with retained behavior;
\item
  the warning that compressibility is evidence of shorter representation
  under a declared equivalence, not proof of mechanism or intelligence.
\end{itemize}

The downstream chapter must independently justify any stronger claim
about:

\begin{itemize}
\tightlist
\item
  reusable computation;
\item
  mechanistic discovery;
\item
  abstraction;
\item
  intelligence;
\item
  reconstruction from a compressed residual.
\end{itemize}

\section{References used in this
chapter}\label{references-used-in-this-chapter-65}

\begin{itemize}
\tightlist
\item
  (Rissanen 1978) --- shortest-description model selection.
\item
  (Eckart and Young 1936) --- optimal lower-rank matrix approximation.
\item
  (Han et al. 2016a) --- pruning, trained quantization/weight sharing,
  and Huffman coding.
\item
  (Hinton et al. 2014) --- teacher-student knowledge distillation.
\end{itemize}

Exact provenance and claim boundaries are locked in:

`sources/source-locks/ATLAS-CH-COMPRESS-001.yaml`

\chapter{Compression as Discovery and Intelligence
Probe}\label{compression-as-discovery-and-intelligence-probe}

\textbf{Epistemic status:} audited Compression/Residual substrate +
primary MDL/probing/compression-progress sources + Atlas-owned
diagnostic synthesis.\\
\textbf{Specification:}
manuscript/specifications/ATLAS-CH-COMPINTEL-001.md\\
\textbf{Derivation packet:}
mathematics/derivations/ATLAS-CH-COMPINTEL-001-DERIVATIONS.md\\
\textbf{Computational witness:}
mathematics/computational-witnesses/ATLAS-CW-COMPINTEL-001.md\\
\textbf{Source lock:} sources/source-locks/ATLAS-CH-COMPINTEL-001.yaml

Compression is attractive because it turns a vague word such as
structure into a number.

If a model can describe data in fewer bits, perhaps it has found a
pattern.

If a representation allows a target to be transmitted with fewer bits,
perhaps it contains something reusable.

If a system suddenly compresses its observations better, perhaps it has
discovered a regularity.

Those are reasonable hypotheses.

They are not identities.

The central rule of this chapter is therefore:

\[
\boxed{
\text{compression is evidence only relative to a declared code, target, split, and control set.}
}
\]

The stronger the claim, the stronger the controls must be.

\section{Compression is tempting for a good
reason}\label{compression-is-tempting-for-a-good-reason}

Prediction and compression are closely related.

If a model assigns high probability to what actually occurs, an entropy
code can often describe the observations with fewer bits.

If a representation makes a target easy to predict from few examples, a
prequential or MDL-style probe can transmit those targets economically.

This connection appears in several traditions.

Solomonoff's formal induction framework assigns greater prior weight to
sequences generated by shorter programs (Solomonoff 1964).

Minimum-description-length reasoning turns statistical learning into a
coding problem.

Blier and Ollivier explicitly measured deep-learning models through
description length and showed that coding methodology materially changes
the resulting bound (Blier and Ollivier 2018).

Voita and Titov used minimum description length as a probe of learned
representations, replacing a single final-accuracy number with a coding
quantity that also reflects how hard it is to learn the target from the
representation (Voita and Titov 2020).

Schmidhuber proposed compression progress as a signal for curiosity and
discovery by a resource-bounded observer (Schmidhuber 2009).

The shared intuition is productive:

\begin{quote}
regularity that supports prediction can support compression.
\end{quote}

The dangerous step is to reverse or overextend it.

\section{The Atlas boundary}\label{the-atlas-boundary}

The Atlas does not define intelligence as compression.

It does not define understanding as short description.

It does not define mechanism as decodability.

It does not define scientific discovery as any decrease in codelength.

The chapter instead asks a narrower question:

\begin{quote}
when is compression a useful empirical probe for reusable predictive
structure?
\end{quote}

This turns a philosophical slogan into an experimental contract.

\section{Begin with the code}\label{begin-with-the-code}

A codelength is never just a property of the data.

It is a property of the data under a code.

Write

\[
L_C(D)
\]

for the length assigned to data object \(D\) by code \(C\).

A different valid code can give a different answer:

\[
L_{C_1}(D)
\ne
L_{C_2}(D).
\]

That is not a defect.

It is the reason the code belongs inside the notation.

The audited Compression chapter already made the same point for model
and data description.

COMPINTEL extends it to probing and discovery claims.

\section{A short code can be
trivial}\label{a-short-code-can-be-trivial}

Consider eight zeros:

\[
Z=(0,0,0,0,0,0,0,0).
\]

Use a prefix code with two branches:

\begin{itemize}
\tightlist
\item
  one bit means eight zeros;
\item
  otherwise transmit a literal eight-bit string behind a branch bit.
\end{itemize}

Then

\[
L_C(Z)=1.
\]

A literal baseline requires eight bits.

The gain is seven bits.

This is strong compression.

It is not strong intelligence.

No learner was required.

No abstraction was discovered.

No generalization was tested.

No mechanism was explained.

The witness gives an exact counterexample:

\[
\boxed{
\text{compressibility}
\not\Rightarrow
\text{intelligence}.
}
\]

\section{Why training compression is too
weak}\label{why-training-compression-is-too-weak}

A sufficiently flexible system can memorize training data.

That can lower training loss.

It can also lower a training-set code.

Neither fact shows that the system found a rule that applies elsewhere.

A lookup table can describe previously seen labels perfectly.

Its useful content may vanish on unseen examples.

Therefore the first diagnostic upgrade is simple:

\begin{quote}
measure codelength on untouched data.
\end{quote}

\section{Held-out codelength}\label{held-out-codelength}

Let:

\begin{itemize}
\tightlist
\item
  \(R\) be the representation being tested;
\item
  \(T\) be the training partition;
\item
  \(H\) be an untouched held-out partition;
\item
  \(Y_H\) be the held-out targets;
\item
  \(C\) be the code family;
\item
  \(\mathcal P\) be the admissible probe or decoder class.
\end{itemize}

Write

\[
L_C(Y_H\mid R,T)
\]

for the number of bits needed to transmit held-out targets after
learning only from the training partition and the declared
representation.

Now introduce a baseline \(B\).

Define

\[
\boxed{
\Delta_C(R)
=
L_C(Y_H\mid B,T)
-
L_C(Y_H\mid R,T).
}
\]

If

\[
\Delta_C(R)>0,
\]

then \(R\) supports shorter held-out transmission than the baseline
under the declared measurement system.

This is an empirical structure claim.

It is still conditional on the measurement system.

\section{The period-two witness}\label{the-period-two-witness}

Take the training prefix

\[
T=(0,1,0,1).
\]

Suppose the held-out suffix is

\[
H_s=(0,1,0,1).
\]

Use a conditional prefix code:

\begin{itemize}
\tightlist
\item
  one bit means continue the period-two motif inferred from training;
\item
  otherwise use one branch bit plus four literal held-out bits.
\end{itemize}

The structured suffix costs

\[
1
\]

bit.

A four-bit literal baseline costs

\[
4
\]

bits.

Therefore

\[
\boxed{
\Delta_C(H_s)=3.
}
\]

This is held-out compression.

The learned pattern made a correct prediction beyond the training
prefix.

\section{A matched failure control}\label{a-matched-failure-control}

Keep the training prefix exactly the same.

Change only the held-out suffix:

\[
H_c=(0,0,1,1).
\]

The period-two rule fails.

The codec falls back to literal transmission:

\[
1+4=5
\]

bits.

Against the same four-bit baseline,

\[
\boxed{
\Delta_C(H_c)=-1.
}
\]

The comparison is useful because training exposure is identical.

One future continues the learned regularity.

The other does not.

The codelength changes on held-out data.

This is the smallest exact example of what COMPINTEL wants from a
predictive-structure probe.

\section{The result is still
codec-relative}\label{the-result-is-still-codec-relative}

The witness does not mean that one bit is the universal complexity of
the structured suffix.

Another code could assign a different length.

The scientifically relevant claim is narrower:

\begin{quote}
under a declared code whose semantics are held fixed, the learned
regularity reduces held-out description length relative to a declared
baseline.
\end{quote}

A robust result should survive reasonable changes in code or come with
an argument explaining why the comparison is invariant.

\section{MDL probing}\label{mdl-probing}

Voita and Titov make this logic operational for learned representations
(Voita and Titov 2020).

A standard probe often asks:

\begin{quote}
what final accuracy can a classifier reach from this representation?
\end{quote}

An MDL probe asks something closer to:

\begin{quote}
how economically can the target labels be communicated when this
representation is available?
\end{quote}

The difference matters.

Two representations can end at similar accuracy while requiring
different amounts of data or model complexity before the target becomes
easy to recover.

The codelength integrates that learning trajectory.

This makes it a useful diagnostic surface.

It does not make it a mechanism detector.

\section{Same final accuracy can hide different learning
difficulty}\label{same-final-accuracy-can-hide-different-learning-difficulty}

Imagine two probes that both reach perfect final held-out accuracy.

Probe A learns the target after a few examples.

Probe B requires almost the entire training set.

A prequential code charges predictions made along the way.

Probe A can therefore achieve a shorter cumulative description even
though final accuracy is identical.

Hence

\[
\boxed{
\text{same final accuracy}
\not\Rightarrow
\text{same codelength}.
}
\]

This is one reason codelength can reveal differences hidden by accuracy.

\section{Same codelength can hide different
behavior}\label{same-codelength-can-hide-different-behavior}

The converse problem also exists.

Codelength is a scalar summary.

Two models can have equal total code length while:

\begin{itemize}
\tightlist
\item
  failing on different examples;
\item
  having different calibration;
\item
  behaving differently under distribution shift;
\item
  using different internal mechanisms.
\end{itemize}

Therefore

\[
\boxed{
\text{same codelength}
\not\Rightarrow
\text{same behavior}.
}
\]

Compression and behavior belong on separate axes.

\section{Random-label controls}\label{random-label-controls}

A powerful probe can learn surprising things.

Sometimes what it learns is the probing task itself rather than
structure already present in the representation.

Random-label controls help separate these possibilities.

Replace the real target \(Y\) with matched random labels.

Then rerun the same pipeline.

If the representation gives nearly the same codelength advantage on
random labels as on the real target, the intended structure claim is
weak.

A useful diagnostic should show a gap:

\[
\Delta_C^{\mathrm{real}}
>
\Delta_C^{\mathrm{random}}
\]

by a meaningful margin.

Voita and Titov explicitly motivate MDL probing partly through such
random-task distinctions (Voita and Titov 2020).

\section{Shuffle controls}\label{shuffle-controls}

A second control breaks alignment.

Keep the representation vectors.

Keep their marginal distribution.

Shuffle which example receives which representation.

If the compression gain survives, then the signal may come from generic
probe capacity or marginal statistics rather than example-specific
structure.

A representation-target claim should weaken sharply when their
correspondence is destroyed.

\section{Random-representation
controls}\label{random-representation-controls}

Where meaningful, compare against:

\begin{itemize}
\tightlist
\item
  an untrained network;
\item
  a randomly initialized representation;
\item
  a random projection with matched dimension.
\end{itemize}

This helps answer:

\begin{quote}
did training create an economical representation of the target, or would
a generic feature map work just as well?
\end{quote}

The answer is empirical.

The control belongs in the protocol.

\section{Probe capacity is part of the
instrument}\label{probe-capacity-is-part-of-the-instrument}

A probe is a measuring device.

Its capacity matters.

If one representation is paired with a much stronger decoder than
another, a codelength difference can be caused by the instrument.

Therefore a serious COMPINTEL experiment either:

\begin{itemize}
\tightlist
\item
  holds probe capacity fixed;
\item
  sweeps capacity;
\item
  charges model complexity through the code;
\item
  or reports all three.
\end{itemize}

Regularization is part of the same issue.

The measurement apparatus must be visible.

\section{Blier and Ollivier: coding method
matters}\label{blier-and-ollivier-coding-method-matters}

Blier and Ollivier studied the description length of deep-learning
models and found that different coding approaches gave very different
compression bounds (Blier and Ollivier 2018).

The durable lesson for COMPINTEL is not a specific numeric result.

It is methodological:

\begin{quote}
a claim about compressibility can change when the coding method changes.
\end{quote}

That is why codec sensitivity is a mandatory control rather than an
editorial footnote.

\section{Solomonoff: compression and prediction in an idealized
theory}\label{solomonoff-compression-and-prediction-in-an-idealized-theory}

Solomonoff's induction theory gives shorter generating programs greater
prior weight (Solomonoff 1964).

This provides a deep formal connection between economical description
and prediction.

But several boundaries remain.

The ideal theory uses universal computation.

Its complexity quantities are not practical finite compressors.

A neural-network codelength is not automatically Solomonoff complexity.

A short empirical code is not automatically the shortest possible
description.

The source supports the conceptual bridge.

It does not license a universal intelligence score.

\section{Compression progress}\label{compression-progress}

Suppose a learner changes over time.

At stage \(t\), define

\[
\Gamma_t
=
L_{C,t-1}(Y_H\mid R,T)
-
L_{C,t}(Y_H\mid R,T).
\]

Positive \(\Gamma_t\) means that the held-out target became cheaper to
describe.

This can be interesting.

It suggests that the learner has acquired predictive structure.

Schmidhuber proposes a broad version of this idea as a principle for
curiosity and discovery (Schmidhuber 2009).

The Atlas keeps the proposal explicitly hypothesis-level.

\section{When compression progress is
fake}\label{when-compression-progress-is-fake}

A positive \(\Gamma_t\) can arise without discovery.

Examples:

\begin{itemize}
\tightlist
\item
  the held-out set changed;
\item
  labels leaked into training;
\item
  the code stopped charging side information;
\item
  the probe class grew;
\item
  the compressor changed;
\item
  the model was allowed more compute;
\item
  the target behavior was degraded;
\item
  bookkeeping changed.
\end{itemize}

Therefore a discovery interpretation requires fixed or charged
semantics.

The safest longitudinal comparison freezes:

\begin{itemize}
\tightlist
\item
  held-out data;
\item
  target definition;
\item
  code accounting;
\item
  probe class;
\item
  evaluation protocol.
\end{itemize}

\section{Discovery needs controls}\label{discovery-needs-controls}

A compression-progress event is stronger when it satisfies all of the
following:

\begin{enumerate}
\def\labelenumi{\arabic{enumi}.}
\tightlist
\item
  the gain occurs on untouched data;
\item
  random-label controls do not show the same gain;
\item
  shuffled representations destroy the gain;
\item
  matched probe-capacity controls cannot explain it;
\item
  the result survives reasonable codec variants;
\item
  behavior is retained;
\item
  the newly compressed structure helps in another context.
\end{enumerate}

Only then does the word discovery become plausible.

Even then, the statement remains empirical and scoped.

\section{Reuse is stronger than one-task
prediction}\label{reuse-is-stronger-than-one-task-prediction}

Suppose a descriptor is transmitted once and then used for several
tasks.

Let

\[
Y^{(1)},\ldots,Y^{(m)}
\]

be task targets.

If the same descriptor supports low incremental codelength for many
tasks,

\[
\sum_j
L_C(Y_H^{(j)}\mid R,T_j)
\]

can be smaller than the cost of separate task-specific representations.

This is evidence of reuse.

It begins to approach the Atlas notion of portable structure.

But it still does not prove that the descriptor is least.

\section{Compression and the
Residual}\label{compression-and-the-residual}

The Residual chapter asks for a least admissible descriptor sufficient
for a declared capability under a declared factorization preorder.

COMPINTEL asks a different question:

\begin{quote}
which descriptors make relevant behavior economical to recover?
\end{quote}

The two questions are related.

They are not identical.

A shorter code does not prove factorization leastness.

If

\[
L_C(R_1)<L_C(R_2),
\]

we cannot conclude

\[
R_1\preceq R_2.
\]

Codelength defines one ordering.

The Residual preorder defines another.

This distinction prevents compression from swallowing the entire
representation theory.

\section{Invertible recoding}\label{invertible-recoding}

Suppose

\[
R_2=g(R_1)
\]

for invertible \(g\).

Semantically, the two representations contain the same information.

If the admissible decoder class can efficiently apply \(g^{-1}\), a
well-designed probe may treat them equivalently.

But a restricted probe class may not.

For example, a linear probe can fail after a nonlinear invertible
recoding even though no information was lost.

Therefore a probe score is a property of

\[
(R,\mathcal P,C),
\]

not of \(R\) alone.

This is one of the most important interfaces between COMPINTEL and the
Residual.

\section{Recoding stability}\label{recoding-stability}

If a claim is intended to be representation-independent, test it under
admissible recodings.

A robust result should either:

\begin{itemize}
\tightlist
\item
  preserve its codelength ranking;
\item
  or have an explicit explanation for why the decoder class makes the
  recoding operationally harder.
\end{itemize}

A signal that disappears under a harmless coordinate change is still
evidence.

But it is coordinate-sensitive evidence.

It should not be promoted to an intrinsic property.

\section{Probe accessibility is not
mechanism}\label{probe-accessibility-is-not-mechanism}

Suppose a probe decodes a variable \(Y\) from representation \(R\).

Then \(Y\) is accessible to that probe.

That does not show that the original network uses \(Y\).

It may encode the property redundantly.

The property may be a by-product.

The network may ignore it downstream.

Therefore

\[
\boxed{
\text{probe accessibility}
\not\Rightarrow
\text{causal use}.
}
\]

This is where compression diagnostics must hand off to mechanistic
intervention.

\section{What mechanism evidence would
add}\label{what-mechanism-evidence-would-add}

Mechanistic evidence can include:

\begin{itemize}
\tightlist
\item
  ablating the compressed component;
\item
  replacing it with a counterfactual value;
\item
  patching it from another example;
\item
  reconstructing the behavior from the proposed descriptor;
\item
  testing whether destroying the descriptor selectively destroys the
  capability.
\end{itemize}

Now the claim changes.

Instead of saying:

\begin{quote}
the property is economically recoverable,
\end{quote}

we can begin to say:

\begin{quote}
the property participates causally in the computation.
\end{quote}

Compression helps locate candidates.

Intervention tests them.

\section{Codelength and task
distortion}\label{codelength-and-task-distortion}

Compression can also destroy what we care about.

Suppose a lossy compressed model becomes tiny but loses predictive
accuracy.

The codelength improved.

The capability did not.

Therefore the relevant object is at least two-dimensional:

\[
\boxed{
(\text{codelength},\text{retained behavior}).
}
\]

Depending on the setting, retained behavior can mean:

\begin{itemize}
\tightlist
\item
  accuracy;
\item
  calibration;
\item
  reward;
\item
  reconstruction quality;
\item
  safety property;
\item
  transfer performance;
\item
  retrieval fidelity.
\end{itemize}

The distortion axis cannot be omitted.

\section{An evidence ladder}\label{an-evidence-ladder}

The chapter uses five increasingly demanding levels.

\subsection{Level 0: training
compression}\label{level-0-training-compression}

Evidence of fit.

Compatible with memorization.

\subsection{Level 1: held-out
compression}\label{level-1-held-out-compression}

Evidence of predictive structure for the declared target.

\subsection{Level 2: control-separated held-out
compression}\label{level-2-control-separated-held-out-compression}

Evidence that the signal exceeds random-label, shuffle,
random-representation, and probe-capacity baselines.

\subsection{Level 3: recoding-stable
compression}\label{level-3-recoding-stable-compression}

Evidence that the signal is not merely a coordinate accident under the
declared transformation class.

\subsection{Level 4: cross-context
reuse}\label{level-4-cross-context-reuse}

Evidence that one descriptor or learned rule supports several tasks with
low incremental description cost.

\subsection{Level 5: mechanistic
intervention}\label{level-5-mechanistic-intervention}

Evidence that the compressed structure is causally involved in the
capability.

No level is a definition of general intelligence.

\section{Intelligence is broader}\label{intelligence-is-broader}

An intelligent system may need:

\begin{itemize}
\tightlist
\item
  prediction;
\item
  adaptation;
\item
  planning;
\item
  abstraction;
\item
  transfer;
\item
  action;
\item
  uncertainty handling;
\item
  search;
\item
  resource allocation;
\item
  coordination.
\end{itemize}

Compression can interact with many of these.

It need not characterize all of them.

A highly compressible constant sequence is not intelligent.

A capable agent may face genuinely stochastic data that cannot be
strongly compressed.

Therefore the safest Atlas statement is:

\begin{quote}
compression can be a probe of learned regularity and reuse, not a scalar
definition of intelligence.
\end{quote}

\section{A practical COMPINTEL
protocol}\label{a-practical-compintel-protocol}

For a model checkpoint, declare:

\begin{enumerate}
\def\labelenumi{\arabic{enumi}.}
\tightlist
\item
  representation \(R\);
\item
  target \(Y\);
\item
  train/held-out split;
\item
  code \(C\);
\item
  probe class \(\mathcal P\);
\item
  baseline;
\item
  side-information accounting.
\end{enumerate}

Then measure

\[
L_C(Y_H\mid R,T).
\]

Compute

\[
\Delta_C(R).
\]

Then repeat with:

\begin{itemize}
\tightlist
\item
  random labels;
\item
  shuffled representation-target alignment;
\item
  random or untrained representation;
\item
  probe-capacity sweep;
\item
  at least one codec/probe variant.
\end{itemize}

Finally report retained task behavior.

If the claim is representation-independent, add admissible recoding
tests.

If the claim is reuse, add a second task.

If the claim is mechanism, add intervention.

\section{Longitudinal discovery
protocol}\label{longitudinal-discovery-protocol}

For a learner evolving through checkpoints \(t\):

\begin{enumerate}
\def\labelenumi{\arabic{enumi}.}
\tightlist
\item
  freeze a held-out evaluation stream;
\item
  freeze target semantics;
\item
  freeze code accounting;
\item
  compute \(L_{C,t}\) at each checkpoint;
\item
  compute \(\Gamma_t\);
\item
  inspect whether progress coincides with learned representations or
  behaviors;
\item
  rerun random/shuffle controls;
\item
  test transfer to future data or another task.
\end{enumerate}

A sudden positive \(\Gamma_t\) is then a candidate discovery event.

The word candidate matters.

\section{Research use: finding reusable
computation}\label{research-use-finding-reusable-computation}

A practical use of COMPINTEL is comparative.

Suppose several layers, modules, heads, memories, or latent states all
support the same downstream task.

Measure the conditional codelength of the task given each candidate
representation.

Then ask:

\begin{itemize}
\tightlist
\item
  which candidate gives the shortest held-out code?
\item
  which result survives recoding?
\item
  which candidate supports several tasks?
\item
  which one survives interventions?
\end{itemize}

This does not identify mechanism automatically.

It narrows the search.

That is already valuable.

\section{Research use: compression before
interpretation}\label{research-use-compression-before-interpretation}

Interpretability often begins with a human-chosen feature.

COMPINTEL suggests another route.

First locate representations that economically support many predictive
targets.

Then inspect what common structure they encode.

Compression becomes a screening instrument.

Interpretation comes afterward.

This ordering can reduce the temptation to overinterpret visually
appealing but non-predictive features.

\section{Research use: scientific discovery
systems}\label{research-use-scientific-discovery-systems}

A scientific agent can treat unexplained data as a coding problem.

A candidate model that predicts new observations can reduce held-out
description length.

That is a useful objective.

But a scientific claim also needs:

\begin{itemize}
\tightlist
\item
  falsifiability;
\item
  intervention or experimental tests;
\item
  model comparison;
\item
  uncertainty;
\item
  robustness;
\item
  external replication.
\end{itemize}

Compression progress can therefore be one component of discovery
machinery.

It is not the whole scientific method.

\section{Source boundaries}\label{source-boundaries}

Solomonoff supports an idealized connection between short generating
descriptions and induction (Solomonoff 1964).

Blier and Ollivier support empirical deep-model codelength analysis and
show that coding methodology matters (Blier and Ollivier 2018).

Voita and Titov support MDL as a representation-probing method and
demonstrate the value of random baselines (Voita and Titov 2020).

Schmidhuber supports the proposed compression-progress principle for
curiosity and discovery (Schmidhuber 2009).

None of these sources establishes:

\[
\text{compression}
=
\text{intelligence}.
\]

The Atlas does not add that equality.

\section{What the computational witness
establishes}\label{what-the-computational-witness-establishes-2}

The finite replay establishes exactly:

\begin{enumerate}
\def\labelenumi{\arabic{enumi}.}
\tightlist
\item
  a fixed training prefix;
\item
  a prefix-free conditional period-two code;
\item
  one-bit held-out coding for the structured continuation;
\item
  three bits of gain relative to a four-bit literal baseline;
\item
  five-bit fallback for a matched noncontinuing suffix;
\item
  negative one-bit gain for that control;
\item
  one-bit coding of an all-zero eight-bit sequence under a special run
  code;
\item
  seven bits of compression gain for that trivial sequence.
\end{enumerate}

It does not establish optimality of either codec.

It does not establish intelligence.

\section{Non-implications}\label{non-implications-6}

The chapter rejects:

\[
\text{training compression}
\not\Rightarrow
\text{generalization},
\]

\[
\text{held-out compression}
\not\Rightarrow
\text{mechanistic use},
\]

\[
\text{short description}
\not\Rightarrow
\text{unique explanation},
\]

\[
\text{compressibility}
\not\Rightarrow
\text{intelligence},
\]

\[
\text{compression progress}
\not\Rightarrow
\text{scientific discovery},
\]

\[
\text{probe accessibility}
\not\Rightarrow
\text{causal necessity},
\]

\[
\text{same task accuracy}
\not\Rightarrow
\text{same codelength},
\]

and

\[
\text{same codelength}
\not\Rightarrow
\text{same behavior}.
\]

\section{Atlas connections}\label{atlas-connections-23}

\textbf{Mechanistic diagnosis.}\\
Compression identifies economically recoverable structure; intervention
tests whether the structure matters causally.

\textbf{The Residual.}\\
Codelength can compare candidate portable descriptors but does not prove
leastness in the factorization preorder.

\textbf{Token and context compression.}\\
A successful compressor may expose reusable predictive summaries,
provided retained behavior is measured.

\textbf{Adaptive curricula and search.}\\
Compression progress can be used as an exploration signal when held-out
and control semantics are explicit.

\textbf{Research strategy.}\\
A sudden codelength reduction can flag a hypothesis worth testing, not a
result worth certifying.

\section{Closing view}\label{closing-view-34}

Compression is powerful because it forces explanation to pay rent.

A pattern that predicts new data can shorten a code.

A representation that exposes useful structure can reduce the data
needed to learn a target.

A reusable abstraction can amortize its description across tasks.

These are real advantages.

They are not magic.

A code can be badly chosen.

A probe can be too strong.

A system can memorize.

A trivial sequence can compress beautifully.

A causal mechanism can remain hidden behind an accessible correlate.

The disciplined conclusion is:

\[
\boxed{
\text{compression is a probe of predictive structure when the code, split, controls, and retained behavior are explicit.}
}
\]

Discovery and intelligence require more.

\section{References used in this
chapter}\label{references-used-in-this-chapter-66}

\begin{itemize}
\tightlist
\item
  (Solomonoff 1964)
\item
  (Blier and Ollivier 2018)
\item
  (Voita and Titov 2020)
\item
  (Schmidhuber 2009)
\end{itemize}

Exact source identities and claim boundaries are recorded in
sources/source-locks/ATLAS-CH-COMPINTEL-001.yaml.

\chapter{From Models to Agents}\label{from-models-to-agents}

\textbf{Epistemic status:} established representative mechanisms plus
Atlas synthesis and derivation.\\
\textbf{Specification:}
manuscript/specifications/ATLAS-CH-AGENTS-001.md\\
\textbf{Formal packet:}
mathematics/derivations/ATLAS-CH-AGENTS-001-DERIVATIONS.md\\
\textbf{Computational witness:}
mathematics/computational-witnesses/ATLAS-CW-AGENTS-001.md\\
\textbf{Source lock:} sources/source-locks/ATLAS-CH-AGENTS-001.yaml

A model can answer a question without being able to change what question
comes next.

That is the threshold this chapter cares about.

A language model presented with a fixed context computes a continuation.
The continuation may contain a plan, a command, a proof sketch, a search
query, or a vivid description of an action. But text that describes an
action is not yet an action. Nothing outside the model changes merely
because the model wrote a verb.

An agent appears when selected outputs are allowed to participate in a
loop.

An observation enters. A model proposes. A controller chooses among
admissible operations. An interface executes one of them. The
environment or tool returns a consequence. Some part of that consequence
is retained. The next decision is made in light of the changed state.

The important transition is therefore not from ``less intelligent'' to
``more intelligent.''

It is from an open-loop mapping to a closed interaction architecture.

That distinction is modest enough to be useful.

It does not require us to decide whether a system has intentions,
beliefs, or autonomy in a human sense. It asks a narrower systems
question:

Can consequences of selected operations change the state from which
later operations are chosen?

If yes, then we have acquired dynamics that cannot be understood by
inspecting one model call in isolation.

\section{The model is not the loop}\label{the-model-is-not-the-loop}

The Adaptive-System Thesis established an explanatory boundary for the
Atlas.

A learned model may be one component of a larger computational system
that also contains memory, tools, interfaces, evidence, and governance.

Agents make that boundary concrete.

Consider a model

M(c) = y,

where c is supplied context and y is an output.

This object can be extraordinarily capable. It can summarize, prove,
translate, classify, generate code, or propose a sequence of actions.

Yet the equation says nothing about whether any proposed action is
executed.

It says nothing about who may execute it.

It says nothing about what the execution changes.

It says nothing about whether the result is observed.

It says nothing about whether that result changes the next model call.

Those missing arrows are where agent structure lives.

This chapter therefore uses ``model'' and ``agent'' as structural roles,
not prestige labels.

A model can be a component of an agent.

The same model can also be used without an agent loop.

And a system can become more agentic in its interaction structure
without changing model weights at all.

\section{A minimal agent object}\label{a-minimal-agent-object}

For the purposes of the Atlas, write an agent as

A = (M, C, X, O, Γ, U, B, σ).

The components are:

M --- a model or proposal operator;

C --- a controller that selects the next admissible operation;

X --- persistent internal state;

O --- an observation constructor;

Γ --- an execution interface connecting actions to an external state;

U --- an update operator for internal state;

B --- the current budget-and-authority contract;

σ --- a stopping rule.

This is not proposed as the one true definition of an agent.

It is a map legend for the mechanisms this and later chapters need to
discuss.

At step t, let e\_t denote external state, x\_t internal state, and g
the current goal.

The loop is:

o\_t = O(e\_t),

p\_t = M(x\_t, o\_t, g),

a\_t = C(x\_t, o\_t, p\_t; B\_t),

(e\_\{t+1\}, r\_t) = Γ(e\_t, a\_t),

x\_\{t+1\} = U(x\_t, o\_t, a\_t, r\_t).

The budget and authority state are then updated, and σ decides whether
another transition is permitted or required.

Every symbol matters.

The model proposal p\_t need not be the executed action a\_t.

The observation o\_t need not expose the full environment state e\_t.

The returned result r\_t need not be trusted.

The update U may discard information, summarize it, corrupt it, or
preserve it exactly.

The authority contract B\_t may prohibit an operation even if the model
strongly proposes it.

And the stop rule may terminate a competent loop or fail to terminate an
incompetent one.

An agent is therefore not a model plus a longer prompt.

It is a model situated inside a transition system.

\section{Why interaction changes the
problem}\label{why-interaction-changes-the-problem}

The difference becomes sharp when useful information does not exist in
the original context.

Suppose an environment contains one hidden bit s in \{0,1\}.

Before interaction, both possible worlds emit the same observation:
UNKNOWN.

A deterministic one-shot model must therefore produce the same answer in
both worlds.

It cannot guarantee correctness on both because the correct answers
differ.

Now add one authorized tool: QUERY.

QUERY returns the hidden bit.

Nothing about the downstream answer rule has changed. The same model can
simply report an observed bit if one exists.

But the system can now:

\begin{enumerate}
\def\labelenumi{\arabic{enumi}.}
\tightlist
\item
  call QUERY;
\item
  receive the consequence;
\item
  store the returned bit;
\item
  invoke the answer rule on the changed state.
\end{enumerate}

The exact witness in ATLAS-CW-AGENTS-001 checks both possible hidden
states.

Without the query transition, the fixed answer rule succeeds on one of
two states.

With one authorized query and budget one, it succeeds on two of two.

Remove the authority or remove the budget, and performance returns to
one of two.

The point is not that querying is sophisticated.

The point is that capability can depend on the transition structure
around a fixed model.

A system can know more at step t+1 because it did something at step t.

That sentence is the seed from which tool use, search, planning, memory,
and coordination grow.

\section{Tools are transition
channels}\label{tools-are-transition-channels}

A tool is often described by its name: calculator, browser, compiler,
database, theorem prover, shell, simulator, calendar.

For systems analysis, the name is less important than the interface.

For a tool u with argument z, write

Γ\_u(e, z) = (e', r).

The call may leave the external state unchanged and merely return
information.

Or it may change the world.

A search query and a bank transfer are both ``tool calls'' in loose
language. Their transition semantics and authority consequences are
radically different.

This is why the Atlas separates three events:

the model proposes a call;

the controller determines whether the call is admissible;

the execution interface performs the call and returns a result.

Collapsing these steps creates a dangerous fiction in which generated
text appears to carry its own authority.

It does not.

Authority comes from the execution boundary.

The ReAct architecture provides a clear representative mechanism for
interleaving reasoning traces, task actions, and returned observations.

Toolformer provides another representative mechanism: a model can learn
when to call particular APIs, what arguments to supply, and how to
incorporate returned information into later prediction.

These are important constructions.

They should not be inflated into the claim that giving a model more
tools necessarily makes the resulting system better.

A larger action set creates opportunities and failure modes at the same
time.

\section{Planning is computation before
commitment}\label{planning-is-computation-before-commitment}

Not every proposed action should be executed immediately.

Planning inserts additional computation between the current state and
external commitment.

Instead of following one candidate continuation, the system may
construct several.

It may estimate consequences.

It may score alternatives.

It may prune some branches.

It may backtrack.

A generic search layer can be written as

H\_\{k+1\} = Select(Score(Expand(H\_k))),

where H\_k is a set of candidate partial histories at search depth k.

Nothing in this notation requires the candidates to be natural language.

They could be symbolic states, trajectories, programs, proof states,
schedules, or mixed representations.

For language models, Tree of Thoughts is a useful representative example
because it makes the branching structure explicit. Intermediate text
states are generated, evaluated, explored, and sometimes abandoned.

The important abstraction is not the word ``thought.''

It is the separation between proposing one continuation and allocating
computation over several possible continuations before acting.

Planning can improve a decision while still being wrong.

The generator can miss the useful branch.

The evaluator can rank candidates badly.

The search can prune the correct path.

The lookahead model can mispredict the environment.

More search is therefore not identical to more truth.

Planning changes how inference is allocated.

It does not repeal the need for evidence about the quality of the
generator, evaluator, model of consequences, and stopping rule.

\section{Reflection is state update, not
self-transcendence}\label{reflection-is-state-update-not-self-transcendence}

Agent systems are often said to ``reflect.''

The word invites more mystery than necessary.

Suppose an attempt produces feedback f\_t.

A reflective mechanism computes an update

x\_\{t+1\} = U\_reflect(x\_t, f\_t).

The next attempt is conditioned on x\_\{t+1\} rather than x\_t.

That is enough structure to study.

Reflexion is a representative mechanism of this type. Feedback is
converted into verbal material retained in episodic memory and used in
subsequent trials.

The empirical question is whether the update improves later decisions.

The system may reflect accurately.

It may also construct a persuasive explanation of the wrong failure,
preserve it, and become more consistently wrong.

Fluency does not settle that question.

A reflection should therefore be treated as a candidate state update
whose downstream effect can be measured, ablated, substituted, or
contradicted.

This returns us to one of the Atlas's recurring themes:

description is weaker than intervention.

If reflection is claimed to matter, compare the system with and without
the reflective state, or with controlled alternatives, while preserving
the rest of the loop as closely as possible.

\section{Decomposition requires
interfaces}\label{decomposition-requires-interfaces}

Large goals are often decomposed into smaller tasks.

That sounds simple until two subtasks disagree about what passes between
them.

Let a task decomposition be a directed graph

G\_T = (V, E).

Each node v has at least:

\begin{itemize}
\tightlist
\item
  a local goal g\_v;
\item
  admissible inputs I\_v;
\item
  a required return contract Q\_v;
\item
  a local resource budget R\_v;
\item
  a stop condition σ\_v.
\end{itemize}

An edge u -\textgreater{} v means that v may consume some declared
output or state produced by u.

A list is not yet a robust decomposition.

A decomposition becomes operational when the interfaces between pieces
are explicit enough that a downstream node knows what it may assume.

This is the same compositional pressure seen elsewhere in the Atlas.

Tensor shapes do not guarantee semantic compatibility.

Subtask completion labels do not guarantee that a returned object means
what the next task needs.

The more independent the subtasks become, the more important their
boundaries become.

\section{Delegation is decomposition plus
authority}\label{delegation-is-decomposition-plus-authority}

Delegation adds another ingredient.

A child process is not merely asked to compute something.

It is given some capacity to act.

Represent a delegation packet as

D = (g', A', R', Q', σ').

Here g' is the bounded child goal, A' the authority or tool set granted
to the child, R' the resource budget, Q' the required return object, and
σ' the child stop rule.

Suppose the parent currently has authority set A and resource budget R.

The default Atlas rule is:

A' subseteq A

and

R' \textless= R,

with the resource inequality interpreted componentwise when R is a
resource vector, unless a distinct external authority explicitly grants
an enlargement.

This is not a complete security theorem.

It is a clean invariant.

Spawning a child does not create permission out of nothing.

If every delegation edge respects authority inclusion, then authority
cannot expand merely by descending the delegation tree. Set inclusion is
transitive.

If every parent allocates child budgets from its own finite budget, the
total resource assigned to descendants can likewise be bounded by the
root allocation, provided the accounting boundary is actually enforced.

This is what ``bounded delegation'' means in this chapter.

Not timid delegation.

Not necessarily shallow delegation.

Delegation whose authority, resources, return contract, and stopping
condition remain explicit.

\section{The return is not the
truth}\label{the-return-is-not-the-truth}

A delegated child returns something.

Perhaps it is an answer, a proof, a file, a tool result, a plan, a test
report, a witness, or a failure record.

The parent now faces an epistemic problem as well as a control problem.

What is this object allowed to establish?

AGENTS-001 does not solve that question fully. Later chapters on
coordination and evidence exchange will.

But one rule belongs here:

child completion is not automatic parent belief.

The return Q' is an input to the parent's next state transition.

It may be checked, replayed, challenged, combined with other evidence,
rejected, or escalated.

This matters even when the child and parent use the same underlying
model.

A second invocation is not independent verification merely because it
happened later.

\section{Authority and capability are different
coordinates}\label{authority-and-capability-are-different-coordinates}

An agent may be capable of proposing an operation it is not authorized
to perform.

It may be authorized to perform an operation it cannot use competently.

These are different axes.

Let Cap(x) denote the set of operations the system can successfully
formulate or execute from state x under ideal access.

Let Auth(x) denote the operations the current contract permits.

The executable set is constrained by their intersection, together with
the actual interfaces available.

A safe system design cannot infer authority from apparent competence.

Nor should it infer competence from possession of credentials.

This distinction becomes increasingly important as agent systems acquire
access to software environments, communication channels, financial
operations, scientific instruments, or other consequential interfaces.

The model says what it proposes.

The authority boundary says what may happen.

\section{Budgets turn ``keep trying'' into
mathematics}\label{budgets-turn-keep-trying-into-mathematics}

An unbounded instruction such as ``continue until solved'' hides several
decisions.

How many model calls?

How many tool calls?

How much wall-clock time?

How many child tasks?

How deep may recursion go?

How much money, energy, memory, or external rate limit may be consumed?

Budgets make these questions explicit.

Let R\_t be a vector rather than a scalar if resources are
heterogeneous.

For example,

R\_t = (tokens, tool\_calls, wall\_time, money, child\_slots).

Each transition incurs a nonnegative cost vector c\_t.

The update is

R\_\{t+1\} = R\_t - c\_t,

subject to component-wise nonnegativity.

A stop rule can then require termination before a forbidden negative
resource state is reached.

This does not guarantee useful behavior.

It converts one vague control surface into an inspectable one.

\section{Termination is part of the
agent}\label{termination-is-part-of-the-agent}

A loop without a meaningful stop condition is not merely inefficient.

It is incompletely specified.

Stopping may depend on goal satisfaction, proof of impossibility,
resource exhaustion, repeated state, lack of progress, external
cancellation, confidence or risk threshold, authority boundary, or human
review.

A good stop rule can be difficult to design because the system may not
know whether one more step would succeed.

But omitting the problem does not remove it.

It merely replaces an explicit stopping policy with whatever accidental
limit the runtime, user, or infrastructure imposes.

For recursive delegation, termination becomes structural.

A finite, fixed acyclic task graph terminates when each activated node
terminates and execution does not create new nodes outside that graph.

Cyclic or dynamically expanding delegation requires an additional
argument: decreasing resource, decreasing measure, explicit depth bound,
or some other well-founded condition.

``Try again'' is not a termination proof.

\section{Where agent loops fail}\label{where-agent-loops-fail}

The formal decomposition lets us locate failures more precisely.

\subsection{Observation failure}\label{observation-failure}

The useful distinction in the environment never reaches the agent.

Two states that require different actions may look identical through O.

\subsection{Proposal failure}\label{proposal-failure}

The necessary information is present, but M proposes a poor
continuation.

\subsection{Controller failure}\label{controller-failure}

A sound proposal exists, but C selects a worse operation or admits an
operation that should be blocked.

\subsection{Interface failure}\label{interface-failure}

The tool or environment does not behave according to the semantics the
agent assumes.

A returned result may be stale, malformed, adversarial, or simply wrong.

\subsection{State-update failure}\label{state-update-failure}

Useful evidence arrives and is then summarized away, overwritten,
misindexed, or attached to the wrong object.

\subsection{Planning failure}\label{planning-failure}

Search explores the wrong branches, scores them badly, or spends its
budget before reaching a useful continuation.

\subsection{Reflection failure}\label{reflection-failure}

The system explains its failure incorrectly and makes the mistaken
explanation persistent.

\subsection{Delegation failure}\label{delegation-failure}

The child goal is underspecified, authority is too broad, the return
contract is ambiguous, or recursion is allowed to expand without a
decreasing resource.

\subsection{Termination failure}\label{termination-failure}

The system stops before the necessary evidence arrives or continues
after further actions have become wasteful or unsafe.

Calling all of these ``model errors'' would erase the architecture we
need in order to repair them.

\section{Autonomy is not one
number}\label{autonomy-is-not-one-number}

People often ask whether a system is autonomous as though autonomy were
a scalar property.

For engineering purposes, several coordinates matter separately.

How long can the loop continue without external intervention?

Which actions may it execute?

How much state persists?

Which environments can it observe?

Can it create child tasks?

Can it alter its own plan?

Can it spend resources?

Can it change the authority of descendants?

What events force review or termination?

Two systems using the same model may therefore have radically different
autonomy surfaces.

One may be allowed to search a read-only corpus for three steps.

Another may execute code, modify files, communicate externally, spawn
children, and continue for hours.

Calling both ``agents'' does not make those differences disappear.

The useful question is not merely whether an agent exists.

It is what transition system has actually been authorized.

\section{What the representative mechanisms
establish}\label{what-the-representative-mechanisms-establish}

The sources used in this chapter demonstrate several concrete
mechanisms.

ReAct demonstrates that reasoning traces and environment-facing actions
can be interleaved, with observations feeding later steps.

Toolformer demonstrates a learned mechanism for deciding when and how to
call external APIs and incorporate their returns.

Tree of Thoughts demonstrates explicit branching, evaluation, lookahead,
and backtracking over intermediate language-model states.

Reflexion demonstrates a feedback pathway in which verbalized
reflections are retained and influence later episodes.

These results justify discussing such mechanisms as concrete engineering
objects.

They do not prove that every capable agent must use them, every task
benefits from them, self-evaluation is reliable, tool use is safe,
search always improves performance, reflection always corrects errors,
or more autonomy is better.

The Atlas needs mechanisms without mythology.

\section{From one loop to many}\label{from-one-loop-to-many}

A single agent already contains several interacting states: model
context, persistent memory, external environment, tool state, resource
state, and authority state.

The next chapter introduces a different difficulty.

What happens when several such loops share information or act on common
resources?

Then we acquire communication, shared state, concurrency, race
conditions, rendezvous, transactions, provenance across actors, partial
failure, disagreement, and coordination costs.

Those are not decorations on the single-agent loop.

They form a new systems layer.

ATLAS-CH-COORD-001 may therefore assume the agent object

A = (M, C, X, O, Γ, U, B, σ),

together with typed actions, persistent state, bounded delegation,
budget state, authority state, and stop rules.

It must add the mathematics and systems semantics of coordination rather
than pretending that a collection of agents is just one larger prompt.

\section{The boundary to remember}\label{the-boundary-to-remember-3}

The chapter began with a small distinction.

A model can produce text that describes an action.

An agent loop can allow a selected action to change the state from which
later computation proceeds.

Everything else in this chapter elaborates that closure.

Tools enlarge the transition relation.

Planning spends computation before commitment.

Reflection changes persistent state.

Decomposition creates explicit subproblems.

Delegation grants bounded capacity to act on them.

Budgets limit continuation.

Authority limits which transitions may occur.

Stopping rules close the process.

None of these mechanisms guarantees intelligence.

None guarantees truth.

None guarantees safety.

But once they are explicit, they become objects that can be analyzed
rather than properties vaguely attributed to ``the model.''

That is the move from models to agents that the Atlas needs.

\section{References used in this
chapter}\label{references-used-in-this-chapter-67}

External mechanism sources are pinned in:

sources/source-locks/ATLAS-CH-AGENTS-001.yaml

They include:

\begin{itemize}
\tightlist
\item
  ReAct: Synergizing Reasoning and Acting in Language Models, ICLR 2023,
  arXiv:2210.03629;
\item
  Toolformer: Language Models Can Teach Themselves to Use Tools, NeurIPS
  2023, DOI 10.52202/075280-2997;
\item
  Tree of Thoughts: Deliberate Problem Solving with Large Language
  Models, NeurIPS 2023, arXiv:2305.10601;
\item
  Reflexion: Language Agents with Verbal Reinforcement Learning, NeurIPS
  2023, DOI 10.52202/075280-0377.
\end{itemize}

The Atlas formal loop, budget-and-authority contract, and delegation
invariants are chapter-owned synthesis and derivation.

\chapter{Coordination Architectures}\label{coordination-architectures}

\textbf{Epistemic status:} established coordination mechanisms plus
Atlas synthesis and exact finite witness.\\
\textbf{Specification:}
manuscript/specifications/ATLAS-CH-COORD-001.md\\
\textbf{Formal packet:}
mathematics/derivations/ATLAS-CH-COORD-001-DERIVATIONS.md\\
\textbf{Computational witness:}
mathematics/computational-witnesses/ATLAS-CW-COORD-001.md\\
\textbf{Source lock:} sources/source-locks/ATLAS-CH-COORD-001.yaml

One agent can fail because it chose badly.

Two agents can fail even when each chooses correctly.

That is the new phenomenon.

Once several bounded interaction loops communicate or act on shared
resources, correctness is no longer only a property of each local
decision. It depends on what each actor can see, when events become
visible, which operations may interleave, whether messages can repeat,
what state is shared, and which effects are atomic.

Coordination is the mathematics and systems discipline of those
relationships.

It begins where the previous chapter ended.

From Models to Agents gave us a single interaction loop with state,
tools, authority, budgets, and a stop rule. We now place several such
loops in one world.

The result is not simply a larger prompt.

It is a concurrent system.

\section{From one loop to many}\label{from-one-loop-to-many-1}

Let actor i carry an audited single-agent object A\_i.

Each actor can observe, propose, act, update state, spend resources, and
stop.

Now suppose actor P produces an output needed by actor Q.

Several questions appear immediately.

How does Q learn that the output exists?

Can Q see a partial result?

What if P changes it while Q is reading?

What if the message is delivered twice?

What if Q acts before an earlier causal dependency arrives?

What if both actors modify the same state?

What if one fails halfway through a multi-step effect?

What if both wait for one another?

None of these questions is answered by making either model more capable.

They belong to the coordination layer.

For the Atlas, write a coordination system as

C = (I, \{A\_i\}, \{X\_i\}, S, K, ≺, T, F).

Here:

I indexes the participating actors;

A\_i is actor i's single-agent loop;

X\_i is actor-local state;

S is shared coordination state or a shared substrate;

K is the set of communication and coordination operations;

≺ is the causal or application-relevant event-order relation;

T is the declared transaction and atomicity semantics;

F is the failure, retry, timeout, redelivery, and deduplication policy.

This tuple is not the definition of every distributed system.

It is the minimum map legend needed for the architectures in this
chapter.

\section{Communication is not
agreement}\label{communication-is-not-agreement}

Suppose P sends Q a message:

``candidate proof complete.''

Several facts may follow, depending on the system.

The message may have been emitted.

It may have entered a queue.

It may have been delivered.

Q may have read it.

Q may have parsed it.

Q may have accepted its meaning.

Q may have changed state because of it.

The proof may even be correct.

These are different events.

A communication channel transports an object or signal.

It does not confer epistemic authority on the content.

Nor does successful delivery imply that sender and receiver now agree.

This distinction matters for multi-agent systems because a transcript
can look coordinated while the actors are operating with incompatible
state.

The message is part of the coordination story.

Its truth belongs to an evidence story that comes later.

\section{Local state and shared
state}\label{local-state-and-shared-state}

Every actor has some local state X\_i.

A coordination architecture may also expose shared state S.

Shared state is attractive because it creates a common surface.

Write once, and many actors may observe.

But ``shared'' does not mean ``simultaneously identical in every
observer.''

A system may cache.

It may replicate.

It may materialize views at different times.

It may delay propagation.

It may expose snapshots.

It may permit concurrent writers.

The phrase ``shared memory'' therefore hides a second question:

what consistency semantics govern observation of that memory?

Two agents can both have access to the same named store while seeing
different versions of it.

Coordination requires the visibility rule, not merely the storage
location.

\section{Six coordination families}\label{six-coordination-families}

The architectures below solve different coordination problems.

They should not be arranged on a single ladder from primitive to
advanced.

\subsection{Direct message passing}\label{direct-message-passing}

The simplest picture is an addressed channel.

P sends m to Q.

Q receives m.

This architecture makes actor identity and directed communication
explicit.

It can keep local state local.

But it immediately raises delivery questions:

Is the channel reliable?

Can messages be reordered?

Can duplicates occur?

Does a send block?

Can the receiver refuse?

Does one sender's order survive mixing with messages from others?

The Actor tradition is one important lineage for thinking in terms of
independent computational actors and messages.

The useful abstraction is not that every modern agent system is an Actor
system.

It is that independently evolving components can coordinate through
explicit messages rather than through one shared call stack.

\subsection{Rendezvous}\label{rendezvous}

Sometimes the communication event should happen only when both sides
participate.

A rendezvous couples progress.

A sender reaches a communication point.

A receiver reaches a compatible point.

The communication transition occurs when the two synchronize.

Communicating Sequential Processes made this kind of communication a
first-class structuring idea for concurrent programs.

Rendezvous can simplify some ordering questions because synchronization
is explicit.

It can also create blocking and deadlock surfaces.

Coordination is always a trade: making one relation explicit moves
difficulty somewhere else.

\subsection{Blackboard coordination}\label{blackboard-coordination}

A blackboard replaces pairwise addressing with shared visible problem
state.

One component posts a partial hypothesis.

Another notices it and adds a refinement.

A third may detect a conflict.

The blackboard is not merely storage.

It is a coordination surface around which heterogeneous problem-solving
processes organize their activity.

H. Penny Nii's blackboard account is useful because it separates the
blackboard model of problem solving from the many concrete systems built
around it.

For agent systems, the same lesson remains valuable.

A shared workspace can decouple contributors from direct knowledge of
one another.

But the blackboard does not decide which contribution is correct.

And once several writers can modify the same state, scheduling and
consistency become unavoidable.

\subsection{Tuple spaces}\label{tuple-spaces}

Linda sharpens shared coordination into a small algebra.

The shared object is a tuple space: a logically shared multiset of
tuples.

A process may perform operations such as:

out(t): add tuple t;

rd(p): read a tuple matching pattern p without removing it;

in(p): take a matching tuple and remove it.

The important word is matching.

A consumer can coordinate by describing the shape of the object it needs
rather than naming the producer that must send it.

This changes the coupling structure.

Producer and consumer can be separated in space, time, and identity.

The tuple space becomes the rendezvous surface.

The destructive in operation adds another feature: possession can become
exclusive if the take operation is atomic.

A work token can be placed once and consumed by one worker.

That is coordination through state semantics rather than addressed
messaging.

Again, nothing about a tuple's presence makes its content true.

The space coordinates availability, not epistemic status.

\subsection{Event-driven
coordination}\label{event-driven-coordination}

Another architecture treats history as the primary object.

Actors append events.

Other actors subscribe, replay, materialize, or derive later state.

This design has a powerful consequence:

the path by which state was obtained can remain visible.

But event-driven systems acquire their own obligations.

What is the event identity?

Can delivery repeat?

Which order is authoritative?

Can consumers replay from a checkpoint?

Are materialized views deterministic?

What happens when old events are reinterpreted by new code?

An event log can make provenance easier to preserve while making state
reconstruction an explicit part of the architecture.

The Atlas uses the public AETHER project as one bounded GCL example of
this direction.

At its pinned public revision, AETHER describes an append-only causal
journal, deterministic resolution over fixed journal prefixes, recursive
rule execution, and derivation provenance.

That is exact project evidence.

It is not evidence that every coordination system should look like
AETHER.

\subsection{Transactions}\label{transactions}

Transactions answer a different question.

Which set of effects should appear as one unit?

Suppose an operation reads a value, computes a new value, and writes it
back.

If another actor can interleave its own read and write halfway through,
the two locally correct procedures may compose incorrectly.

A transaction boundary can prevent some harmful interleavings by
grouping operations under declared atomicity and isolation semantics.

Jim Gray's transaction work is a foundational reference point for this
view.

But ``transaction'' is not a magic word.

A transaction must have a boundary.

Database writes may be inside it while network calls, emails, actuator
effects, or external APIs remain outside.

Atomicity is only as broad as the set of effects the transaction system
actually controls.

\section{Causality before total
order}\label{causality-before-total-order}

When several actors run concurrently, it is tempting to ask:

What happened first?

Sometimes the question has no application-level answer.

Leslie Lamport's happens-before relation gives a disciplined way to say
what the system actually knows.

Write a ≺ b for the strict happens-before relation generated by three
rules:

\begin{enumerate}
\def\labelenumi{\arabic{enumi}.}
\tightlist
\item
  if a and b occur in one actor and local program order places a before
  b, then a ≺ b;
\item
  if a is the send of a delivered message and b is the matching receive,
  then a ≺ b;
\item
  if a ≺ b and b ≺ c, then a ≺ c.
\end{enumerate}

This gives the causal partial order used in the chapter.

But two independent events may be incomparable.

Neither happened ``before'' the other in the causal sense relevant to
the computation.

A runtime can still assign them timestamps and impose a total order.

That can be useful.

It is also extra structure.

A total order can serialize events that the application did not need to
order.

The distinction matters because unnecessary ordering can create
coordination cost.

If two events are genuinely independent, forcing global agreement about
their order may buy no semantic value.

\section{A small race with a complete
answer}\label{a-small-race-with-a-complete-answer}

The easiest concurrency bug is already enough to justify the chapter.

Let shared state begin at

x = 0.

Actors P and Q both intend to increment x once.

Each performs:

read x;

write the value it read plus one.

Individually, each program is correct.

If P runs completely and then Q runs completely, the final value is 2.

If Q runs completely and then P runs completely, the final value is also
2.

But concurrency introduces four more legal schedules.

The computational witness enumerates all six interleavings that preserve
each actor's local read-before-write order.

Only two finish at 2.

Four finish at 1.

Both actors ``succeeded'' locally.

The shared effect is wrong.

This is the lost-update problem in its smallest useful form.

The lesson is not merely that concurrency is dangerous.

It is more precise:

local correctness plus shared state does not imply compositional
correctness under interleaving.

\section{Atomicity changes the schedule
space}\label{atomicity-changes-the-schedule-space}

Now replace each two-step increment with one atomic operation:

INC: x \textless- x + 1.

There are still two possible actor orders.

P then Q.

Q then P.

Both produce 2.

The arithmetic did not change.

The allowed interleavings changed.

That is what atomicity buys in this example.

It removes intermediate states from the schedule space visible to
competing actors.

This is a recurring systems theme.

Many coordination mechanisms do not improve the intelligence of any
participant.

They change which global executions are possible.

\section{Delivery is not effect}\label{delivery-is-not-effect}

Messages introduce a related distinction.

Suppose a worker receives:

charge account by 10.

If the transport provides at-least-once delivery, the same logical
message may arrive twice.

If the worker blindly applies the command twice, the effect is 20.

The transport may be behaving exactly as designed.

The application is still wrong.

One repair is idempotence.

For state transition E and message m, idempotence means:

E(E(s,m),m) = E(s,m).

Applying the same logical operation twice has the same effect as
applying it once.

Another repair is deduplication.

Give the logical operation a stable identity.

Record that identity durably when its effect commits.

Ignore a later delivery carrying an already-committed identity.

This sounds like exactly-once processing.

It is only as strong as its assumptions.

If the deduplication record commits but an external side effect does
not, or the external side effect occurs but the record does not,
ambiguity returns.

Exactly-once effect is therefore not obtained merely by writing
``exactly once'' on a queue.

It is an end-to-end claim about identity, atomicity, persistence, and
side effects.

\section{Retry is a semantic
operation}\label{retry-is-a-semantic-operation}

Retries are often treated as infrastructure.

They are part of the program.

Suppose an actor sends an operation and sees no reply.

Three states are compatible with that observation:

the operation never arrived;

the operation arrived and failed;

the operation succeeded but the reply was lost.

Retrying is correct in the first case.

It may be correct in the second.

It can duplicate an effect in the third.

A retry policy F must therefore be read together with the operation
semantics.

For a pure read, retry may be harmless.

For a non-idempotent side effect, retry may require an operation
identity, transaction boundary, compensating action, or human review.

The absence of a reply is not the same thing as evidence that nothing
happened.

\section{Shared state is not
authority}\label{shared-state-is-not-authority}

Blackboards, tuple spaces, event logs, databases, and shared memories
can all become central surfaces.

That centrality creates a temptation.

If a fact is in the shared store, treat it as true.

The previous evidence chapter already taught us why that fails.

A coordination substrate answers questions such as:

Who can see this object?

Who can consume it?

Who can overwrite it?

In what order did related events occur?

Has this operation already been applied?

Those are important facts.

They do not establish the semantic truth of the payload.

A bad hypothesis can be perfectly replicated.

A false tuple can be atomically committed.

An incorrect event can have impeccable provenance.

Coordination and evidence meet later.

They must first remain distinct.

\section{Deadlock and livelock}\label{deadlock-and-livelock}

Synchronization can fail even when no data is corrupted.

Consider two actors.

P holds resource a and waits for b.

Q holds b and waits for a.

No actor is locally broken.

The system stops making progress.

That is deadlock.

Livelock is different.

Actors continue changing state but repeatedly interfere so that useful
progress does not occur.

A pair of overly polite processes can continually yield to one another.

An agentic analogue is a pair of workers that repeatedly reassign a
disputed task without resolving it.

Progress therefore needs its own semantics.

``Nothing crashed'' is weaker than ``the system advanced.''

\section{Partial failure changes
reasoning}\label{partial-failure-changes-reasoning}

In a local sequential program, a function call usually returns, raises,
or the process fails.

Distributed coordination adds ambiguous partial outcomes.

One actor can remain healthy while another fails.

A network can partition.

A shared service can be unreachable from one participant and reachable
from another.

A timeout can mean failure, delay, overload, or a lost response after
success.

This makes absence difficult to interpret.

Silence is not a single state.

The coordination architecture must decide what evidence about failure is
available and what actions are justified under uncertainty.

This is one reason distributed protocols become more elaborate than
their local analogues.

They are not merely moving bytes.

They are managing incomplete knowledge about a changing joint state.

\section{Coordination cost is real}\label{coordination-cost-is-real}

Coordination can improve correctness and enlarge what a group can
accomplish.

It also costs something.

Messages consume bandwidth and serialization work.

Synchronization introduces waiting.

Transactions constrain concurrency.

Retries consume repeated work.

Provenance and deduplication require metadata.

Global ordering can create contention.

A useful bookkeeping object is the coordination-cost vector

c\_coord = (c\_comm, c\_sync, c\_retry, c\_meta),

whose components track communication, synchronization, retry/recovery,
and coordination-metadata costs.

The components need not share one physical unit and therefore need not
be collapsed into a scalar.

The vector is a reminder that coordination is not free structure.

This leads to a design question that will recur throughout the Atlas:

What is the weakest coordination semantics that still preserve the claim
or invariant we care about?

If two computations are independent, do not synchronize them merely
because synchronization is available.

If a task requires exclusive consumption, a destructive tuple-space take
may be more natural than global locking.

If a shared decision requires a transaction, pretending eventual
convergence is sufficient may be dangerous.

Architecture begins by identifying the invariant.

\section{The AETHER example}\label{the-aether-example}

AETHER is useful here because it makes one particular coordination
choice explicit.

At the pinned public revision, its center of gravity is an authoritative
semantic kernel built around an append-only causal journal,
deterministic state resolution, recursive rule evaluation, and
provenance/explanation.

Its interface record separates journal, resolver, rules, runtime, and
explanation roles.

It also states an invariant that fixed journal prefix plus compiled
program yields deterministic results.

This is exactly the sort of claim that belongs in a coordination chapter
because it exposes the chosen shared-state semantics.

AETHER does not replace the general theory.

It is one design point.

Compared with direct message passing, it makes shared history more
central.

Compared with a bare blackboard, it makes derivation and replay more
explicit.

Compared with a simple queue, it treats the semantic state reconstructed
from history as a first-class object.

Those are architectural choices with costs and benefits.

The Atlas records them as public project evidence, not as a universal
optimum.

\section{What coordination does not
solve}\label{what-coordination-does-not-solve}

A coordination system can ensure that one worker takes a task token.

It cannot ensure the worker solves the task correctly.

It can serialize two writes.

It cannot ensure the serialized value is scientifically valid.

It can deliver every message.

It cannot ensure the actors interpret the messages the same way.

It can replay every event.

It cannot ensure the original event was truthful.

It can give every actor the same state.

It cannot ensure the common state encodes the right world model.

This boundary is central.

Coordination can preserve relationships among computational events.

It cannot replace the semantics of the work those events carry.

\section{What the next chapters may
assume}\label{what-the-next-chapters-may-assume-1}

Two direct architectural consumers depend on this chapter.

ATLAS-CH-EVIDEX-001 --- Evidence Exchange and Zero-Context Work may
assume stable actor/event identity, local/shared-state separation,
channel semantics, causal order, retry and deduplication policy, and
transaction boundaries.

Its task is to add evidence semantics:

what a returned object establishes;

how provenance follows a handoff;

how an independent actor receives enough context to reproduce or
challenge a claim.

ATLAS-CH-POLITY-001 --- The Computational Polity may assume the full
coordination object

C = (I, \{A\_i\}, \{X\_i\}, S, K, ≺, T, F),

together with the architecture families and failure surfaces developed
here.

Its task is larger:

compose models, agents, memory, tools, humans, validators, and
governance into a coherent computational polity.

Neither chapter needs to rediscover what a race condition or duplicate
effect is.

That is the point of dependency.

\section{The boundary to remember}\label{the-boundary-to-remember-4}

The simplest distributed-system mistake is to believe that because every
component is correct, the system is correct.

The lost-update witness shows otherwise.

Correct local steps can compose into a wrong global state when the
coordination semantics are underspecified.

The reverse lesson is equally important.

Coordination machinery should not be worshipped for its own sake.

Messages, blackboards, tuple spaces, event logs, rendezvous, clocks, and
transactions are ways of constraining or exposing possible interactions.

They are useful when those constraints match an invariant the larger
system actually needs.

The governing question is therefore not:

Which coordination architecture is best?

It is:

Which relationships among actors, events, state, and effects must remain
true, even when execution is concurrent and failure is partial?

Once that question is explicit, coordination stops being plumbing.

It becomes part of the mathematics of adaptive systems.

\section{References used in this
chapter}\label{references-used-in-this-chapter-68}

The source lock pins the mechanism references and Atlas project objects
used by this draft:

sources/source-locks/ATLAS-CH-COORD-001.yaml

External mechanism sources include the 1973 Actor formalism, Hoare's
Communicating Sequential Processes, Carriero and Gelernter's Linda in
Context, Nii's blackboard-system account, Lamport's event-ordering
paper, and Gray's transaction paper.

The AETHER material is exact GCL public project evidence at commit
74b2e322a4453f1665a076bc0682adcd0c5cfb44 and is used only as a bounded
coordination case study.

\chapter{Evidence Exchange and Zero-Context
Work}\label{evidence-exchange-and-zero-context-work}

\textbf{Epistemic status:} provenance standards plus Atlas synthesis,
bounded GCL project evidence, and exact finite witness.\\
\textbf{Specification:}
manuscript/specifications/ATLAS-CH-EVIDEX-001.md\\
\textbf{Formal packet:}
mathematics/derivations/ATLAS-CH-EVIDEX-001-DERIVATIONS.md\\
\textbf{Computational witness:}
mathematics/computational-witnesses/ATLAS-CW-EVIDEX-001.md\\
\textbf{Source lock:} sources/source-locks/ATLAS-CH-EVIDEX-001.yaml

A result that is useful to its author can still be useless to everyone
else.

The missing ingredient is often not intelligence.

It is transfer.

A proof sketch may depend on a source version that was never named.

An experiment may depend on a hidden environment choice.

A contributor may say ``this works'' without stating the exact claim,
the domain on which it was tested, or the assumptions it used.

A second actor can receive the same words and reconstruct a different
task.

At that point the problem is no longer only epistemic and no longer only
coordinative.

It is an evidence-exchange problem.

The Atlas needs a way to move candidate evidence between actors while
preserving enough identity, scope, provenance, and replay information
that a later reader can determine what was actually done.

That is the subject of this chapter.

\section{Evidence must survive
handoff}\label{evidence-must-survive-handoff}

The Evidence chapter gave us a claim-support packet

K = (q, tau, S, Omega, N, D).

The Coordination chapter gave us actors, channels, event identities,
causal order, retries, and transaction boundaries.

Neither object by itself tells us how to package a bounded scientific or
mathematical task for another actor.

That requires an interface.

The interface has two directions.

A dispatch says:

this is the exact obligation;

these are the facts you may import;

these are the sources and tools you may use;

these are the failure conditions;

this is the return grammar;

this is where the result must go.

A return says:

this is the dispatch I answered;

this is my strongest exact conclusion;

this is my support;

this is my provenance;

this is how to test or falsify it;

this is what remains unresolved.

Evidence exchange begins when both directions are explicit.

\section{The dispatch packet}\label{the-dispatch-packet}

Write a dispatch as

D = (delta, Q, B, Sigma, C, Lambda, Gamma, rho).

Here delta is a stable dispatch identity.

Q is the exact bounded obligation.

B is the bootstrap: definitions, protected facts, and immutable
imported-source identities.

Sigma is the allowed source and tool policy.

C is the context and independence class.

Lambda is the set of resource bounds, rejection conditions, and
legitimate stop conditions.

Gamma is the exact return grammar.

rho is the durable return route.

This object is intentionally stricter than an ordinary instruction.

An ordinary instruction can rely on shared context.

A dispatch intended for independent work should not.

\section{What zero-context actually
means}\label{what-zero-context-actually-means}

``Zero context'' can be misunderstood.

It does not mean the worker knows no mathematics.

It does not mean the worker has never seen the topic.

It does not mean the worker is statistically independent from every
other actor.

It means the task itself should not depend on hidden conversational
state.

Fix a declared class of eligible actors A.

Let H be hidden session history.

Zero-context sufficiency constrains the \textbf{authorized task
contract}, not the worker's psychology.

Write

Contract\_A(D,H)

for the obligation, premises, permissions, source policy, stop
conditions, and return requirements that the protocol authorizes when D
is accompanied by hidden history H.

Then require

\begin{Shaded}
\begin{Highlighting}[]
\NormalTok{Contract\_A(D,H1)}
\NormalTok{=}
\NormalTok{Contract\_A(D,H2)}
\end{Highlighting}
\end{Shaded}

for admissible hidden histories H1 and H2.

The exact obligation must be the same.

The imported facts must be the same.

The permitted sources must be the same.

The return grammar must be the same.

The stop conditions must be the same.

A worker can bring background knowledge and can still misunderstand an
instruction.

What it cannot be authorized to obtain from hidden history is the
missing half of the assignment.

\section{Hidden context and explicit
prerequisites}\label{hidden-context-and-explicit-prerequisites}

There is nothing wrong with prerequisites.

Science depends on them.

The problem is invisible prerequisites.

Compare two instructions.

First:

``Use the convergence theorem we discussed yesterday.''

Second:

``Use Theorem T from source S at immutable version h, under hypotheses
H1 through H4.''

The second instruction can still be wrong.

It can name the wrong theorem.

Its hypotheses may not match.

But the dependency is inspectable.

The first instruction delegates part of the task definition to shared
memory.

That works until the worker changes, the session ends, the conversation
is unavailable, or two participants remember different versions.

Zero-context work therefore does not eliminate context.

It compiles relevant context into an explicit packet.

\section{A dispatch is a boundary
object}\label{a-dispatch-is-a-boundary-object}

A good dispatch has to satisfy two constituencies.

The sender needs confidence that the worker is solving the intended
problem.

The worker needs enough information to know what is allowed and what a
valid return looks like.

This is why a dispatch should include excluded assumptions as well as
imported facts.

Suppose a source proves a theorem for finite groups.

The target problem concerns compact groups.

If the dispatch lists only the finite-group theorem and says
``generalize it,'' an actor may quietly assume the missing extension.

A better dispatch says:

finite-group theorem available;

compact-group extension unavailable;

do not import it.

The exclusion is not negative prose.

It defines the search space.

\section{Source policy is part of the
mathematics}\label{source-policy-is-part-of-the-mathematics}

Two actors can receive the same mathematical question and still be
solving different assignments if one is allowed to search external
literature and the other is restricted to a protected packet.

This is why Sigma belongs in the dispatch.

Possible source policies include:

protected packet only;

primary sources allowed;

public sources allowed;

repository-local sources only;

no external source use;

tools allowed but internet disallowed.

These are not administrative details.

They change what evidence can be introduced.

A zero-context dispatch must therefore tell the worker not only what the
problem is, but what evidence universe it is allowed to draw from.

\section{Context class is not a certificate of
independence}\label{context-class-is-not-a-certificate-of-independence}

Suppose a dispatch says

C = independent\_blind.

That can mean the worker is not shown selected prior returns.

This may reduce anchoring or copying.

It may improve adversarial value.

But the label is weaker than the word independent sometimes suggests.

Two workers can use the same foundation model.

They can share training data.

They can share organizational assumptions.

They can unknowingly reproduce the same error.

They can even be the same model in separate sessions.

So an independence label should be read as an information-flow
constraint.

What information was withheld?

What information remained available?

What provenance does the worker actually have?

That is more useful than pretending one adjective proves methodological
independence.

\section{The return object}\label{the-return-object}

Now write a return as

R = (delta, chi, K, Pi, V, Delta).

The first field, delta, links the return to the dispatch.

chi is the declared disposition.

K is the claim-support packet from the Evidence chapter.

Pi is provenance.

V contains verification and falsification hooks.

Delta is the unresolved residual or next bounded obligation.

This structure matters because a return is not merely an answer.

It is a candidate evidence object.

A phrase such as

PROVED

or

COUNTEREXAMPLE

inside chi does not make itself true.

The disposition tells the reviewer what the contributor claims to have
returned.

Review remains a separate operation.

\section{Provenance asks how an object came to
exist}\label{provenance-asks-how-an-object-came-to-exist}

The W3C PROV family gives us a useful conceptual separation.

There are entities.

There are activities.

There are agents.

Entities can be generated by activities.

Activities can use entities.

Agents can bear responsibility for activities or generated objects.

Entities can derive from other entities.

This is exactly the kind of separation evidence exchange needs.

A return can therefore carry provenance like:

which input objects were used;

which actor or execution identity produced the return;

which procedure transformed the inputs;

which output entity was generated;

which derivation links connect the result to its sources.

The Atlas does not need to reproduce the entire PROV specification to
benefit from the distinction.

The key lesson is that payload and lineage are different objects.

\section{Provenance is not truth}\label{provenance-is-not-truth}

A false result can have excellent provenance.

Imagine an arithmetic script that reads the exact intended file, runs
under the exact pinned interpreter, records the exact command, and
deterministically prints the wrong number because the code contains a
bug.

Every provenance field can be correct.

The scientific claim can still fail.

This is an important boundary.

Provenance lets us ask:

where did this result come from?

It does not answer:

is this result true?

The second question requires verification, proof, replication,
adjudication, or some other evidence operation.

Good provenance makes those operations easier.

It does not replace them.

\section{Identity has layers}\label{identity-has-layers}

Evidence exchange often uses identifiers.

Not all identifiers answer the same question.

A filename answers:

what human-facing path or label was used?

A commit-and-path pair answers:

which repository version and file path?

A content hash answers:

which bytes?

A DOI may answer:

which scholarly object or version?

A bibliographic record answers:

which published source is intended?

These identities can reinforce one another.

They can also diverge.

Two files can have the same name and different bytes.

Two archives can contain identical bytes under different bibliographic
descriptions.

A hash can identify bytes perfectly while saying nothing about what
those bytes are supposed to mean.

So source locking should be designed around the failure mode.

If version drift matters, identify the version.

If bibliographic confusion matters, identify the work.

If byte identity matters, fingerprint the bytes.

\section{A tiny ambiguity with an exact
answer}\label{a-tiny-ambiguity-with-an-exact-answer}

The chapter's witness uses two source versions.

Both are called

inputs.txt.

Version one contains:

a=2

b=3.

Version two contains:

a=2

b=4.

The declared procedure is to add a and b.

The returned evidence bytes say:

result=5.

Now consider a handoff that records only

source\_path=inputs.txt.

There are two candidate sources.

Replay version one and we obtain 5.

Replay version two and we obtain 6.

The returned result bytes alone do not tell us which source was used.

The provenance reconstruction is ambiguous.

Now add the exact SHA-256 of version one.

Within the finite candidate set, only one source matches.

Replay reproduces 5.

Nothing about the arithmetic became more intelligent.

The handoff became better identified.

\section{Hashes solve one problem, not every
problem}\label{hashes-solve-one-problem-not-every-problem}

It is tempting to turn that witness into a slogan:

hash everything.

That would be too strong.

A hash is excellent for byte identity.

It does not tell us:

whether the source is authoritative;

whether the source corresponds to the bibliographic work we intended;

whether a PDF scan omitted a page;

whether an executable is safe;

whether a dataset was collected correctly;

whether a theorem's hypotheses match our use.

A strong evidence object can therefore need both content identity and
semantic identity.

The source lock says what the object is.

The fingerprint says which bytes were inspected.

The authority field says what role the source is permitted to play.

These are separate claims.

\section{FAIRness and
exchangeability}\label{fairness-and-exchangeability}

The FAIR principles were written for scientific data management and
stewardship, but several of their ideas apply naturally to evidence
exchange.

Objects benefit from persistent identifiers.

Metadata should be rich enough to describe the object.

Metadata should explicitly identify the data or object it describes.

Objects should carry qualified references to related objects.

Reusable objects should carry detailed provenance.

The principles also explicitly extend beyond conventional data to
algorithms, tools, and workflows.

This is useful for an Atlas concerned with machine-assisted research.

But FAIRness is not a truth metric.

A beautifully described, permanently identified, interoperable object
can still be wrong.

FAIR improves the conditions for discovery and reuse.

Epistemic status remains separate.

\section{Return grammars force useful omissions into
view}\label{return-grammars-force-useful-omissions-into-view}

Free-form prose is flexible.

That is both its strength and its weakness.

A contributor can return three pages of sophisticated reasoning and
forget to say the exact theorem being claimed.

A return grammar reduces that ambiguity.

A useful grammar asks for:

the strongest exact statement;

the derivation or support;

assumptions beyond the bootstrap;

verification or falsification hooks;

the claim boundary;

the next residual.

These fields create friction.

That friction is useful.

It forces the contributor to distinguish what was shown from what merely
seems plausible.

It also makes machine intake easier because important evidence fields
have stable locations.

\section{The return route matters}\label{the-return-route-matters}

Suppose a worker produces a perfect result in a private transient chat.

A later reviewer never sees it.

Did the programme receive evidence?

Operationally, perhaps someone read it.

Durably, perhaps not.

When downstream review depends on a return, the transport should create
a persistent identity.

That could be:

a versioned file;

an issue comment;

an append-only journal event;

a database record;

an object-store artifact;

another durable evidence surface.

The exact technology is secondary.

The requirement is that a later actor can identify the same return
object.

Durability is therefore part of the evidence chain when later inspection
matters.

\section{Receipt is not acceptance}\label{receipt-is-not-acceptance}

Evidence workflows become confused when arrival is treated as
validation.

The Atlas distinguishes four states.

Receipt:

the return reached the declared surface.

Acceptance:

an intake or reviewer determined that it satisfies the structural or
substantive contract required at that stage.

Promotion:

a stronger project object incorporates or relies on some claim from the
return.

Certification:

a separate process confers whatever formal or institutional status
certification means.

These are not synonyms.

A structurally valid counterexample can be accepted even though it
defeats the hoped-for theorem.

A returned proof can be received and rejected.

An accepted proof can still require formal replay before promotion.

A promoted theorem can still lack machine-checked certification.

The state changes must remain explicit.

\section{Structural validity is weaker than mathematical
validity}\label{structural-validity-is-weaker-than-mathematical-validity}

A return can satisfy every field in the grammar.

The dispatch ID can match.

The provenance can be complete.

The claim boundary can be explicit.

The code can run.

The theorem can still be false.

Structural intake answers:

is this object well formed enough to review?

Mathematical or scientific adjudication answers:

does the evidence support the claim?

Those are different layers.

This distinction prevents a parser from accidentally becoming a theorem
prover.

\section{Replayability is graded}\label{replayability-is-graded}

A result is not simply replayable or non-replayable.

We can pin different amounts of the execution.

Level one might identify the source and algorithm.

Level two might add exact code and parameters.

Level three might add package versions, random seeds, and expected
outputs.

Level four might add environment/container identity and
hardware-sensitive details.

Different claims require different levels.

A symbolic proof may care little about GPU model.

A floating-point kernel benchmark may care a great deal.

So a return should pin what is material to the claim, not maximize
metadata without purpose.

The later Replayable Evidence Objects chapter develops this in more
detail.

\section{Zero-context work is compiled
context}\label{zero-context-work-is-compiled-context}

There is a useful way to think about the entire exercise.

A zero-context packet is not context-free.

It is context that has been compiled.

The sender begins with a large working context:

papers;

repository state;

conversation history;

prior failed attempts;

governance boundaries;

current theorem frontier.

Only part of that context is necessary for one bounded task.

The dispatch compiler extracts:

definitions;

protected facts;

forbidden imports;

source identities;

success criteria;

return grammar;

stop conditions.

The worker receives the minimal sufficient task interface rather than
the sender's entire history.

This is an information-design problem.

Too little context produces ambiguity.

Too much context can leak prior answers, waste attention, or defeat
blindness.

\section{Why excluded assumptions are
first-class}\label{why-excluded-assumptions-are-first-class}

Imported facts tell the worker what is available.

Excluded assumptions tell the worker what must still be earned.

This is especially important in theorem work.

Suppose a source states convergence but leaves the topology implicit.

If a downstream theorem needs convergence in a stronger topology, the
missing strength is the problem.

A dispatch that merely says

prove convergence

may invite the worker to restate the source.

A dispatch that says

source-level convergence is imported;

strong-topology convergence is not imported;

show the smallest additional estimate needed

defines a much sharper task.

Zero-context work is strongest when the boundary between granted and
ungranted knowledge is explicit.

\section{Falsification conditions make returns cheaper to
judge}\label{falsification-conditions-make-returns-cheaper-to-judge}

A work packet should not only say what success looks like.

It should say what would make a contribution unusable.

Examples:

uses a forbidden source;

requires an unstated assumption;

changes the theorem domain;

depends on an unpinned artifact;

exceeds the search budget;

returns only numerical evidence where proof is required.

These conditions improve efficiency.

They allow both worker and reviewer to stop early when the task contract
has been violated.

This is the evidence-exchange analogue of an API rejecting an invalid
input rather than producing an ambiguous output.

\section{GCL as a bounded implementation
example}\label{gcl-as-a-bounded-implementation-example}

The public MATHSOLVE repository contains a concrete implementation of
this style.

Its zero-context dispatch template requires an exact bounded obligation,
definitions, imported facts, excluded assumptions, permitted
contribution types, falsification conditions, resource bounds, a return
contract, and an authority boundary.

Its RESULT/1 grammar requires fields for the strongest exact statement,
derivation, assumptions, verification or falsification hooks, claim
boundary, and next residual.

A public Yang-Mills work packet marks itself

ZERO\_CONTEXT

and

independent\_blind,

restricts the worker to a protected packet, states explicit success
criteria, and demands one structured result.

These are useful because they are not hypothetical.

They show how the abstract exchange object can be instantiated in a real
research workflow.

They remain project evidence.

The Atlas does not claim that every research programme should use the
same field names or protocol.

\section{A return must point backward before synthesis can point
forward}\label{a-return-must-point-backward-before-synthesis-can-point-forward}

The dispatch ID delta is not administrative decoration.

It is the primary backward edge.

A return should tell us which obligation it answers.

Without that link, a result can be mathematically interesting and
operationally orphaned.

The same is true of source identities.

Evidence should point backward to what generated it before a synthesis
step points forward to new claims.

This creates a directed provenance structure:

sources and protected facts;

dispatch;

actor activity;

return;

review;

promotion.

Each edge can be inspected.

That is far safer than a narrative chain in which conclusions gradually
detach from their inputs.

\section{Synthesis is a new reasoning
step}\label{synthesis-is-a-new-reasoning-step}

Suppose two independent returns arrive.

One proves a statement for domain Omega\_1.

Another proves a related statement for Omega\_2.

A controller cannot simply paste the prose together and declare a
theorem on a larger domain.

For the special pointwise form

for every x in Omega\_i, q(x),

two valid returns do justify q on Omega\_1 union Omega\_2.

That simple rule does \textbf{not} generalize automatically.

A global property, uniqueness statement, convergence mode, coupled
hypothesis, or domain-dependent definition may require a new composition
theorem.

Synthesis therefore has its own burden.

Are the proposition schemas identical?

Are the definitions aligned?

Do the hypotheses match?

Are the source versions compatible?

Does one argument depend on an assumption excluded by the other?

What logical rule licenses the combined domain?

The return objects make these questions visible.

They do not answer them automatically.

Synthesis is itself an evidence-producing activity.

\section{Stronger prose requires stronger
support}\label{stronger-prose-requires-stronger-support}

Evidence exchange should be monotone in a specific sense.

Passing through more hands must not strengthen a claim unless some new
evidence or argument justifies the strengthening.

A contributor returns:

``I verified the property for n\textless=100.''

A coordinator summarizes:

``the property appears robust.''

A final report writes:

``the property holds.''

The last sentence is stronger.

No new evidence appeared.

This is a promotion failure.

The Evidence chapter already gave us N, the set of stronger conclusions
not licensed.

Evidence exchange preserves that field precisely to prevent confidence
from increasing merely because an object travelled through a workflow.

\section{Conflict is an admissible
output}\label{conflict-is-an-admissible-output}

Independent work is valuable partly because returns may disagree.

If one contributor proves q and another returns a counterexample to q
under the same stated assumptions, the correct evidence state is not

two votes for the average.

It is conflict.

The next operation may be:

compare source locks;

compare hidden assumptions;

replay the counterexample;

check the proof step;

send an adversarial packet;

narrow the theorem.

A good exchange protocol preserves contradictions until they are
resolved.

It does not smooth them away for the sake of a clean dashboard.

\section{Duplicate delivery and evidence
identity}\label{duplicate-delivery-and-evidence-identity}

Coordination systems can retry.

That means the same logical return may appear more than once.

Two comments can carry the same result.

The same payload can be reposted by different actors.

A worker can retry after an uncertain transport failure.

So evidence intake needs stable logical identity.

The dispatch identity helps.

A return identity can help further.

Payload equality alone is not enough because two independent actors can
legitimately return identical text.

Conversely, small formatting changes do not necessarily create a new
logical result.

This is the evidence analogue of the Coordination chapter's distinction
between delivery and effect.

\section{Evidence exchange and
memory}\label{evidence-exchange-and-memory}

Evidence objects often become memory.

A returned proof can enter an archive.

A source lock can enter persistent memory.

A failed attempt can become episodic memory for later search.

A promoted theorem can become semantic project memory.

This creates a new boundary.

Storing a return does not promote it.

Retrieving it later does not revalidate it.

Memory preserves access.

Evidence status must travel with the object.

That is why provenance and epistemic class are not optional metadata in
long-lived research systems.

\section{Evidence exchange and
agents}\label{evidence-exchange-and-agents}

Agentic systems make these issues more urgent because delegation can
happen quickly.

A parent agent can spawn several workers.

Each can search, compute, or prove.

Returns can arrive concurrently.

A summarizer can merge them.

Without explicit exchange semantics, the system can lose:

which worker used which source;

which result was independent;

which assumptions were shared;

which claim was only numerical;

which return was superseded;

which failure was never resolved.

The danger is not only hallucination.

It is provenance collapse.

A coordinated system can become less epistemically legible as it becomes
more capable.

The exchange layer exists to resist that collapse.

\section{Minimal handoff, not maximal
paperwork}\label{minimal-handoff-not-maximal-paperwork}

The goal is not to wrap every thought in bureaucracy.

A tiny task can have a tiny dispatch.

What matters is sufficiency.

If the task is

compute SHA-256 of these exact bytes

then the packet can be short.

If the task is

close the missing theorem-grade convergence bridge in a source with
known omitted proofs

then the packet must be richer.

The metadata burden should scale with the ambiguity and consequence of
the claim.

This is another form of adaptive computation:

spend evidence-handling effort where the epistemic boundary requires it.

\section{What Replay may assume}\label{what-replay-may-assume}

ATLAS-CH-REPLAY-001 may assume:

stable dispatch and return identities;

immutable source/version identities;

the distinction between provenance and truth;

verification/falsification hooks;

graded replayability;

claim-boundary preservation.

Its job is to go further:

pin environments;

capture executable artifacts;

record exact outputs;

build semantic bridges from bytes to claims;

make replay itself a governed evidence object.

It need not rediscover why a mutable pathname is insufficient.

\section{What the Research State Machine may
assume}\label{what-the-research-state-machine-may-assume}

ATLAS-CH-RESEARCHSM-001 may later assume:

bounded work packages;

durable return routes;

context and independence declarations;

receipt versus acceptance versus promotion;

explicit residuals.

Its task is to arrange those objects into programme-level state
transitions.

This chapter supplies the objects that move.

The state-machine chapter supplies the lifecycle.

\section{The boundary to remember}\label{the-boundary-to-remember-5}

Evidence exchange is successful when a result can leave one actor
without losing the conditions under which it means what it means.

That requires more than copying text.

The obligation must have an identity.

The imported facts must be explicit.

The source universe must be declared.

The return must point back to the dispatch.

The claim must carry its scope and non-claims.

The provenance must identify how the object came to be.

The replay hooks must say how another actor can challenge it.

And no workflow step may promote the claim merely because the object
moved successfully.

A good evidence packet therefore behaves like a scientific interface.

It says what goes in.

It says what may be assumed.

It says what comes out.

It says what the output does not establish.

And it leaves enough of a trail that another actor can disagree with it
for the right reasons.

That is the purpose of zero-context work.

Not to remove context.

To make the necessary context explicit, bounded, portable, and
inspectable.

\section{References used in this
chapter}\label{references-used-in-this-chapter-69}

Exact source identities and authority scopes are pinned in:

sources/source-locks/ATLAS-CH-EVIDEX-001.yaml

The external basis includes W3C PROV-DM and the FAIR Guiding Principles.

The bounded implementation examples are exact public MATHSOLVE artifacts
at commit 1273b75457f42da62a8af4c69493d44d293d4567.

\chapter{The Computational Polity}\label{the-computational-polity}

\textbf{Epistemic status:} audited prerequisites + primary scholarly
sources + Atlas synthesis + exact finite witness\\
\textbf{Specification:}
manuscript/specifications/ATLAS-CH-POLITY-001.md\\
\textbf{Derivation packet:}
mathematics/derivations/ATLAS-CH-POLITY-001-DERIVATIONS.md\\
\textbf{Computational witness:}
mathematics/computational-witnesses/ATLAS-CW-POLITY-001.md

A capable model is not automatically a capable system.

A real adaptive system may include several models, a shared memory,
tools, human participants, validators, transaction boundaries, and rules
about who may change what. Once those elements matter to the result,
treating the model alone as the full locus of intelligence hides the
mechanism that actually produced the outcome.

This chapter calls the larger object a \textbf{computational polity}.

\section{From agent to polity}\label{from-agent-to-polity}

Wooldridge and Jennings distinguish agent theory, agent architectures,
and agent languages, already making clear that an intelligent agent is
an architectural object rather than merely a function from prompt to
answer (Wooldridge and Jennings 1995).

The Atlas goes one level outward.

A polity is not just a collection of agents. It is a collection plus the
rules by which local states, shared records, tools, validators, human
roles, and authority compose.

\section{The polity object}\label{the-polity-object}

Write

\[
\Pi=(A,M,U,H,V,C,\Gamma,Q).
\]

Here:

\begin{itemize}
\tightlist
\item
  \(A\) is the set of model or agent roles;
\item
  \(M\) is shared external memory;
\item
  \(U\) is the set of tools or external services;
\item
  \(H\) is the set of human roles;
\item
  \(V\) is the set of validator or reviewer roles;
\item
  \(C\) is the inherited coordination object;
\item
  \(\Gamma\) is the governance relation defining admissible transitions
  and authority;
\item
  \(Q\) is the task or claim contract.
\end{itemize}

These coordinates need not share one data type. The polity is a systems
object, not one giant hidden vector.

\section{Local state and shared
state}\label{local-state-and-shared-state-1}

The Coordination chapter already separates local actor state from shared
coordination state. The External Memory chapter adds version,
provenance, access, synchronization, and deletion semantics for
persistent records.

POLITY therefore inherits a hard boundary:

\[
\text{shared address}\not\Rightarrow\text{shared belief}.
\]

Two participants may refer to the same named store while seeing
different versions, different permissions, different caches, or
different retrieval results.

Common memory is not identical internal state.

\section{Capability is not
authority}\label{capability-is-not-authority}

A role may be able to produce an action without being allowed to commit
it.

Define

\[
\operatorname{Can}(r,a)
\]

for capability and

\[
\operatorname{May}(r,a,\Gamma)
\]

for authorization under governance \(\Gamma\).

In general,

\[
\operatorname{Can}(r,a)\not\Rightarrow\operatorname{May}(r,a,\Gamma).
\]

The reverse implication also fails as a correctness claim. Permission to
act does not prove that the action is wise, true, or safe.

This distinction becomes load-bearing whenever a model proposes a change
that a different role must validate or authorize.

\section{Candidate is not
conclusion}\label{candidate-is-not-conclusion}

A representative polity path is

\[
q
\rightarrow
\text{route}
\rightarrow
\text{candidate}
\rightarrow
\text{evidence record}
\rightarrow
\text{validation}
\rightarrow
\text{authorized commit}.
\]

Every arrow can fail independently.

A good candidate may be routed to the wrong validator. A correct
validator result may be applied to stale evidence. An authorized commit
may fail at an external side effect. A message may be delivered without
being accepted. A shared record may exist without being current.

The polity view makes those failure surfaces visible.

\section{Why humans remain inside the system
boundary}\label{why-humans-remain-inside-the-system-boundary}

Human review is often drawn outside technical diagrams, as though a
person enters only after the computational story is complete.

For a polity, that is a modeling error whenever the human decision
changes system state.

Human roles have inputs, information limits, authority, latency,
provenance, and failure modes. Klein and colleagues emphasize that
automation participating in joint human-agent activity must function as
a team participant rather than as an isolated device (Klein et al.
2004).

The Atlas therefore models human roles explicitly without assuming that
human participation makes the system infallible.

\section{System-level cognition is an old
idea}\label{system-level-cognition-is-an-old-idea}

Hutchins analyzes navigation as a cognitive and computational activity
distributed across people, artifacts, procedures, and representations
rather than exhausted by one person's internal cognition (Hutchins
1995).

The Computational Polity uses that precedent carefully.

It does not claim that every distributed system is intelligent. It
claims only that, when the mechanism of successful computation spans
several components, the correct explanatory boundary may also span those
components.

\section{Exact two-specialist
witness}\label{exact-two-specialist-witness}

Let

\[
Q=\{\alpha,\beta\}.
\]

Define two specialists:

\[
A_\alpha(\alpha)=1,\quad A_\alpha(\beta)=0,
\]

and

\[
A_\beta(\alpha)=0,\quad A_\beta(\beta)=1.
\]

Under a uniform task distribution, each specialist alone has accuracy

\[
\frac12.
\]

Now define a router

\[
R(\alpha)=A_\alpha,\qquad R(\beta)=A_\beta.
\]

The composed polity returns

\[
\Pi(q)=R(q)(q)=1
\]

for both tasks, so

\[
\operatorname{Acc}(\Pi)=1.
\]

The system-level capability is exact and easy to inspect.

\section{The same components can still
fail}\label{the-same-components-can-still-fail}

Now replace the router by one that always chooses \(A_\alpha\).

The component set is unchanged.

The system accuracy returns to

\[
\frac12.
\]

Therefore:

\[
\boxed{\text{more components}\not\Rightarrow\text{greater capability}.}
\]

The composition rule matters.

\section{Shared memory as a polity
institution}\label{shared-memory-as-a-polity-institution}

An accepted result may be stored as

\[
m=(q,y,s,v,\sigma),
\]

where \(q\) is task identity, \(y\) the result, \(s\) source identity,
\(v\) version, and \(\sigma\) status.

This record can be read by later roles without requiring every model to
internalize the result in weights.

But the External Memory boundary remains:

\begin{itemize}
\tightlist
\item
  stored does not mean true;
\item
  present does not mean current;
\item
  shared does not mean universally visible;
\item
  retrieved does not mean correctly used.
\end{itemize}

\section{Validators are roles, not
oracles}\label{validators-are-roles-not-oracles}

A validator can reduce a declared class of error.

It cannot be promoted into an infallible truth function merely because
it has the title `validator'.

A polity therefore records:

\begin{itemize}
\tightlist
\item
  what the validator receives;
\item
  which criterion it applies;
\item
  which authority it has;
\item
  what evidence it preserves;
\item
  what happens when it abstains, fails, or disagrees.
\end{itemize}

The same discipline applies to human review.

\section{Governance is transition
structure}\label{governance-is-transition-structure}

Governance answers questions such as:

\begin{itemize}
\tightlist
\item
  who may create a candidate;
\item
  who may validate it;
\item
  who may commit it;
\item
  which transitions require distinct actors;
\item
  which states are protected;
\item
  which evidence is required before promotion.
\end{itemize}

These are transition constraints.

Governance does not manufacture evidence that is missing from the
underlying claim.

Thus:

\[
\boxed{\text{authorized}\not\Rightarrow\text{correct}.}
\]

and

\[
\boxed{\text{correct-looking}\not\Rightarrow\text{authorized}.}
\]

\section{Tools and external
effects}\label{tools-and-external-effects}

A polity may call search systems, databases, compilers, laboratories,
actuators, or other services.

The inherited coordination boundary matters here: an internal
transaction cannot silently make an external side effect atomic.

A successful tool response and a durably committed external state change
are different events unless the interface explicitly couples them.

\section{Cost is heterogeneous}\label{cost-is-heterogeneous}

A polity consumes more than model FLOPs.

A useful bookkeeping vector is

\[
c_\Pi=(c_{coord},c_{mem},c_{tool},c_{human},c_{val}).
\]

These entries may have different units. They should not be added into
one scalar unless the application declares a conversion or objective.

This prevents `more agents' from being treated as free.

\section{Roles are not identities}\label{roles-are-not-identities}

One physical or software participant may occupy several roles over time.

Conversely, one role may be filled by different participants.

The polity should therefore bind authority and provenance to role and
identity explicitly rather than assuming that labels such as `agent',
`reviewer', or `human' carry permanent semantics.

\section{System-level competence needs a declared
contract}\label{system-level-competence-needs-a-declared-contract}

A polity-level claim should identify at least:

\begin{itemize}
\tightlist
\item
  task set or distribution;
\item
  component roles;
\item
  routing/composition rule;
\item
  memory semantics;
\item
  validation rule;
\item
  governance/authority rule;
\item
  success metric.
\end{itemize}

Without these, `the system can do X' is too underspecified to audit.

\section{What this chapter does not
claim}\label{what-this-chapter-does-not-claim-3}

The polity view does not imply:

\begin{itemize}
\tightlist
\item
  that distributed systems are always better than monolithic models;
\item
  that human review is always superior to automation;
\item
  that more validators always improve reliability;
\item
  that shared memory is globally consistent;
\item
  that governance guarantees correctness;
\item
  that specialization always outweighs coordination cost;
\item
  that system-level behavior is mysterious emergence.
\end{itemize}

The aim is the opposite: make system-level capability reconstructable
from explicit composition.

\section{Handoff to Beyond the Monolithic
Model}\label{handoff-to-beyond-the-monolithic-model}

ATLAS-CH-SYNTHESIS-001 may now assume:

\begin{itemize}
\tightlist
\item
  the polity object \(\Pi\);
\item
  private/shared-state separation;
\item
  the capability/authority distinction;
\item
  candidate, validation, and commit as separate stages;
\item
  the exact specialization witness;
\item
  the requirement that system capability be demonstrated under declared
  composition.
\end{itemize}

SYNTHESIS must still connect this systems view to the rest of the Atlas:
geometry, dynamics, memory, optimization, evidence, adaptation, and
governance.

\section{Durable lesson}\label{durable-lesson-4}

A model can be an important component without being the whole
intelligent system.

When memory, tools, humans, validators, and governance materially
determine the result, the explanatory object should include them.

The computational polity is that larger object.

\section{References used in this
chapter}\label{references-used-in-this-chapter-70}

\begin{itemize}
\tightlist
\item
  (Wooldridge and Jennings 1995)
\item
  (Hutchins 1995)
\item
  (Klein et al. 2004)
\end{itemize}

Exact source authority, prerequisite identities, and claim boundaries
are locked in:

sources/source-locks/ATLAS-CH-POLITY-001.yaml

\chapter{The Machine Under the
Mathematics}\label{the-machine-under-the-mathematics}

\textbf{Epistemic status:} audited Linear Algebra prerequisite +
classical Roofline performance model + current CUDA
programming/best-practices documentation + Atlas synthesis + exact
finite traffic witness.\\
\textbf{Specification:}
manuscript/specifications/ATLAS-CH-HARDWARE-001.md\\
\textbf{Derivation packet:}
mathematics/derivations/ATLAS-CH-HARDWARE-001-DERIVATIONS.md\\
\textbf{Computational witness:}
mathematics/computational-witnesses/ATLAS-CW-HARDWARE-001.md\\
\textbf{Source lock:} sources/source-locks/ATLAS-CH-HARDWARE-001.yaml

\section{Mathematics does not run in zero
time}\label{mathematics-does-not-run-in-zero-time}

A mathematical chapter can write

\[
C=AB.
\]

The hardware must still:

\begin{itemize}
\tightlist
\item
  obtain the inputs;
\item
  schedule the work;
\item
  move values through a memory hierarchy;
\item
  perform arithmetic in some format;
\item
  synchronize where required;
\item
  place the result somewhere.
\end{itemize}

Those operations consume physical resources.

The central boundary of this chapter is therefore:

\[
\boxed{
\text{mathematical work}
\neq
\text{physical execution cost}.
}
\]

Two programs can implement the same map and take different time.

Two programs can execute the same nominal number of floating-point
operations and move different numbers of bytes.

Two kernels can move the same bytes and expose different amounts of
parallel work.

Hardware is not an implementation footnote.

It is part of the effective computational system.

\section{The linear map survives; the implementation
changes}\label{the-linear-map-survives-the-implementation-changes}

The Linear Algebra chapter distinguished an abstract linear map from one
coordinate representation.

Hardware adds another distinction.

Even after the matrix representation has been fixed, there can be many
physical implementations of the same matrix operation.

For example, a matrix product can differ in:

\begin{itemize}
\tightlist
\item
  blocking;
\item
  tile order;
\item
  data layout;
\item
  precision;
\item
  reuse;
\item
  kernel fusion;
\item
  memory placement;
\item
  scheduling.
\end{itemize}

The output may represent the same mathematical operator.

The execution path need not be the same.

This is why algorithmic complexity, FLOP count, and runtime should not
be collapsed.

\section{FLOPs are an accounting
convention}\label{flops-are-an-accounting-convention}

A floating-point operation count is useful only after its convention is
declared.

This chapter's finite witness counts:

\begin{itemize}
\tightlist
\item
  one scalar multiplication as one FLOP;
\item
  one scalar addition as one FLOP.
\end{itemize}

Under that convention, a \(2\times2\) matrix product requires:

\[
8
\]

multiplications and

\[
4
\]

additions, for a total of:

\[
12\ \text{FLOPs}.
\]

Other performance conventions may count fused multiply-add differently.

The point is not that one convention is universally correct.

The point is that comparisons must use one declared convention
consistently.

\section{Throughput is not latency}\label{throughput-is-not-latency}

A machine can sustain a high rate of completed operations while one
individual dependency chain still takes significant time.

These are different quantities.

Latency asks:

\begin{quote}
How long does one operation or dependency chain take?
\end{quote}

Throughput asks:

\begin{quote}
How much work can the machine complete per unit time when enough
parallel work is available?
\end{quote}

Modern accelerators are often designed to tolerate latency by
maintaining many independent operations in flight.

That strategy works only when the workload exposes enough usable
parallelism.

\section{Data movement can dominate
arithmetic}\label{data-movement-can-dominate-arithmetic}

Suppose a kernel performs very little arithmetic on each loaded value.

The arithmetic units may spend much of their time waiting for data.

Suppose instead the same loaded values are reused many times before
eviction.

Then much more arithmetic can be performed per byte transferred.

The relevant quantity is \textbf{arithmetic intensity}.

For a declared memory boundary \(M\), let:

\[
W
\]

be the declared FLOP count and

\[
Q_M
\]

the bytes moved across that boundary.

Then:

\[
\boxed{
I_M=\frac{W}{Q_M}.
}
\]

The subscript matters.

The same computation can have different intensities relative to:

\begin{itemize}
\tightlist
\item
  a host-device link;
\item
  device DRAM;
\item
  an on-chip cache;
\item
  programmer-managed shared memory.
\end{itemize}

Arithmetic intensity is therefore not a property of source code alone.

It is a property of work together with a declared data-movement boundary
and implementation.

\section{The Roofline idea}\label{the-roofline-idea}

Williams, Waterman, and Patterson introduced the Roofline model as a
visual upper-bound model connecting computational throughput, memory
bandwidth, and arithmetic intensity. (Williams et al. 2009)

In its simplest one-level form:

\[
P_{\mathrm{attainable}}
\le
\min(P_{\max},B_M I_M).
\]

Here:

\begin{itemize}
\tightlist
\item
  \(P_{\max}\) is a declared peak compute rate;
\item
  \(B_M\) is a declared bandwidth across memory boundary \(M\);
\item
  \(I_M\) is arithmetic intensity at that boundary.
\end{itemize}

The first term gives a compute roof.

The second gives a bandwidth roof.

The crossover occurs at:

\[
I^\star_M
=
\frac{P_{\max}}{B_M}.
\]

Below this point, the bandwidth-side roof is lower.

Above it, the compute roof is lower.

This does not mean an implementation will reach the roof.

It means the implementation should not be expected to exceed the
declared bound.

\section{A roof is not a
prediction}\label{a-roof-is-not-a-prediction}

Real performance can fall below a Roofline bound because of:

\begin{itemize}
\tightlist
\item
  dependencies;
\item
  poor memory access;
\item
  insufficient parallel work;
\item
  instruction mix;
\item
  synchronization;
\item
  divergence;
\item
  launch overhead;
\item
  cache effects;
\item
  resource constraints;
\item
  implementation quality.
\end{itemize}

So the safe statement is:

\[
\boxed{
\text{Roofline}
=
\text{upper-bound diagnostic},
}
\]

not:

\[
\text{Roofline}
=
\text{runtime oracle}.
\]

The model helps ask whether improving arithmetic throughput is even
relevant before doing so.

\section{The memory hierarchy}\label{the-memory-hierarchy}

The current CUDA programming model exposes several physically and
logically distinct storage levels, including thread-local registers,
shared memory associated with blocks, caches, and global/device memory.
(NVIDIA Corporation 2026b)

The CUDA Best Practices Guide emphasizes that these memory spaces differ
in scope, latency, bandwidth, and intended use. (NVIDIA Corporation
2026a)

The exact hierarchy is architecture-specific.

The general systems lesson is not:

\begin{quote}
every accelerator has this exact hierarchy.
\end{quote}

It is:

\begin{quote}
locality has levels, and moving data between levels has different costs.
\end{quote}

A kernel that repeatedly fetches the same value from a distant level can
behave very differently from one that reuses a nearby copy.

\section{Registers and local state}\label{registers-and-local-state}

Registers are fast, thread-local storage in CUDA's programming model.
(NVIDIA Corporation 2026b)

They are finite.

A kernel that requires more per-thread state can reduce the number of
simultaneously resident threads or spill values into slower storage.

So a transformation that appears to ``save memory traffic'' can still
create another bottleneck if it requires too much local state.

Hardware optimization is usually a trade among constrained resources.

\section{Shared memory and reuse}\label{shared-memory-and-reuse}

CUDA shared memory is an on-chip, programmer-managed storage region
shared among threads in a block. (NVIDIA Corporation 2026b)

The Best Practices Guide uses shared-memory examples to show why staging
data near the compute units can:

\begin{itemize}
\tightlist
\item
  enable coalesced access;
\item
  reduce redundant global-memory loads;
\item
  improve reuse.
\end{itemize}

It also documents bank conflicts as one way shared-memory access can
lose effective bandwidth. (NVIDIA Corporation 2026a)

The general lesson is that ``on chip'' does not mean ``free.''

Access pattern still matters.

\section{Global memory and
bandwidth}\label{global-memory-and-bandwidth}

Device/global memory is large relative to on-chip storage and visible
across the device in the CUDA model. (NVIDIA Corporation 2026b)

Its capacity is useful.

Its distance from execution units makes repeated traffic expensive
relative to reusing data already nearby.

This motivates three recurring questions:

\begin{enumerate}
\def\labelenumi{\arabic{enumi}.}
\tightlist
\item
  How many bytes must cross the boundary?
\item
  How often are the same values reused?
\item
  Can the implementation rearrange work to increase locality?
\end{enumerate}

These questions can matter more than shaving a small number of
arithmetic instructions.

\section{Warps and SIMT}\label{warps-and-simt}

CUDA groups threads into warps of 32 threads in its current programming
model and describes execution using SIMT semantics. (NVIDIA Corporation
2026b)

Threads execute the same kernel but may follow different control-flow
paths.

Divergence can reduce useful lane utilization.

Memory access patterns across a warp also affect how efficiently
global-memory transactions are formed.

These are CUDA-specific execution facts.

They should not be silently promoted into universal laws of all
accelerators.

\section{Parallelism is a resource, not a synonym for
speed}\label{parallelism-is-a-resource-not-a-synonym-for-speed}

A workload can contain many operations and still expose limited
parallelism because of dependencies.

Another workload can expose enormous parallelism but become
bandwidth-limited.

A third can expose both parallelism and reuse but be constrained by
register or shared-memory requirements.

So the Atlas separates:

\begin{itemize}
\tightlist
\item
  work count;
\item
  available parallelism;
\item
  resource residency;
\item
  achieved throughput.
\end{itemize}

These are related.

They are not identical.

\section{Occupancy}\label{occupancy}

In GPU performance discussions, occupancy usually refers to resident
active warps or threads relative to a hardware limit.

It can help hide latency.

It is not itself the objective.

Higher occupancy can coexist with lower performance if, for example:

\begin{itemize}
\tightlist
\item
  each thread does less useful work;
\item
  memory access worsens;
\item
  instruction-level reuse is lost;
\item
  more efficient kernels need more registers or shared memory.
\end{itemize}

The safe rule is:

\[
\boxed{
\text{occupancy}
\neq
\text{performance}.
}
\]

Occupancy is one diagnostic among several.

\section{Matrix and tensor
hardware}\label{matrix-and-tensor-hardware}

Modern NVIDIA GPUs expose warp-level matrix multiply-accumulate
operations that can use Tensor Cores for supported matrix problems of
the form:

\[
D=AB+C.
\]

The programming guide documents tile shapes, supported data types,
fragment layouts, synchronization requirements, and
architecture-dependent restrictions. (NVIDIA Corporation 2026b)

This creates a strong distinction between:

\[
\text{abstract matrix multiplication}
\]

and

\[
\text{matrix multiplication mapped efficiently to specialized hardware}.
\]

A matrix expression does not automatically achieve the device's peak
matrix-unit throughput.

Shapes, layouts, types, alignment, data movement, and surrounding
operations matter.

\section{Tiles connect algebra to
hardware}\label{tiles-connect-algebra-to-hardware}

Large matrix products are decomposed into smaller tiles.

Tiling can improve reuse because a loaded block participates in many
multiply-accumulate operations before replacement.

This is the physical counterpart of a simple algebraic fact:

one matrix entry can contribute to several output entries.

The optimization opportunity comes from arranging the implementation so
the hardware retains that value long enough to reuse it.

The mathematics supplies reuse potential.

The implementation determines whether the machine realizes it.

\section{Precision is several
questions}\label{precision-is-several-questions}

``FP16,'' ``BF16,'' ``FP32,'' and other format names describe
representation and arithmetic choices.

They do not by themselves answer:

\begin{itemize}
\tightlist
\item
  what error accumulates;
\item
  where rounding occurs;
\item
  what the accumulator format is;
\item
  whether overflow or underflow occurs;
\item
  whether the final task tolerates the resulting perturbation.
\end{itemize}

The CUDA Best Practices Guide explicitly separates numerical
accuracy/precision concerns and notes that results can differ across
floating-point precisions because representation and rounding differ.
(NVIDIA Corporation 2026a)

So the Atlas separates:

\[
\text{smaller format}
\]

from

\[
\text{acceptable numerical error}.
\]

The first is a storage/arithmetic choice.

The second is a property that must be checked relative to the
computation and task.

\section{Mixed precision}\label{mixed-precision}

Specialized matrix hardware often supports mixed input and accumulator
formats.

This can be valuable because:

\begin{itemize}
\tightlist
\item
  narrower inputs reduce storage and traffic;
\item
  wider accumulation can reduce some forms of accumulation error.
\end{itemize}

But ``mixed precision'' is not one universal numerical scheme.

The actual behavior depends on the supported operation and formats.

The correct question is always:

\begin{quote}
Which values are represented in which format at which stage?
\end{quote}

\section{Kernels are physical
programs}\label{kernels-are-physical-programs}

A mathematical operator can be implemented by one kernel or several.

Kernel boundaries matter because intermediate values may need to be
materialized in memory.

A fused implementation can sometimes keep intermediates closer to the
arithmetic units.

That can reduce:

\begin{itemize}
\tightlist
\item
  launches;
\item
  synchronization;
\item
  intermediate stores;
\item
  intermediate reloads.
\end{itemize}

But a fused kernel is useful only if it preserves the required semantics
and numerical tolerance.

Performance transformation does not exempt the implementation from
correctness.

\section{Fusion can change the traffic
graph}\label{fusion-can-change-the-traffic-graph}

Consider:

\[
y=f(x),
\qquad
z=g(y).
\]

An unfused implementation may write \(y\) to global memory and later
read it back.

A fused implementation may retain \(y\) in registers or another nearer
storage level.

The mathematical composition remains:

\[
z=(g\circ f)(x).
\]

The physical data-movement graph changes.

This is one reason equal high-level operation graphs can produce
different runtime.

\section{Exact finite witness}\label{exact-finite-witness-1}

Let:

\[
A=
\begin{pmatrix}
1&2\\
3&4
\end{pmatrix},
\qquad
B=
\begin{pmatrix}
5&6\\
7&8
\end{pmatrix}.
\]

Then:

\[
AB=
\begin{pmatrix}
19&22\\
43&50
\end{pmatrix}.
\]

Under the chapter convention:

\[
W=12\ \text{FLOPs}.
\]

\subsection{No cross-output reuse}\label{no-cross-output-reuse}

Charge:

\begin{itemize}
\tightlist
\item
  16 input scalar loads;
\item
  4 output scalar stores.
\end{itemize}

With 4-byte FP32 scalars:

\[
Q_A=80\ \text{bytes}.
\]

Therefore:

\[
I_A=\frac{12}{80}=0.15\ \text{FLOP/byte}.
\]

\subsection{Perfect full-input reuse}\label{perfect-full-input-reuse}

Charge:

\begin{itemize}
\tightlist
\item
  8 total input scalar loads;
\item
  4 output scalar stores.
\end{itemize}

Then:

\[
Q_B=48\ \text{bytes},
\]

and:

\[
I_B=\frac{12}{48}=0.25\ \text{FLOP/byte}.
\]

Same exact matrix product.

Same exact declared FLOP count.

Different charged data movement.

\section{The bound moves even when the mathematics does
not}\label{the-bound-moves-even-when-the-mathematics-does-not}

For a hypothetical one-level machine:

\[
P_{\max}=10\ \mathrm{TFLOP/s},
\qquad
B=1\ \mathrm{TB/s}.
\]

The Roofline bounds become:

\[
P_A\le0.15\ \mathrm{TFLOP/s},
\]

and:

\[
P_B\le0.25\ \mathrm{TFLOP/s}.
\]

The example does not predict a measured GPU speedup.

It proves a narrower point:

\[
\boxed{
\text{same map + same FLOPs}
\not\Rightarrow
\text{same bandwidth-side performance bound}.
}
\]

\section{The witness is intentionally
idealized}\label{the-witness-is-intentionally-idealized}

The finite witness does not charge:

\begin{itemize}
\tightlist
\item
  cache-line overfetch;
\item
  write allocate;
\item
  instruction traffic;
\item
  coherence;
\item
  host-device movement;
\item
  launch overhead;
\item
  synchronization.
\end{itemize}

The perfect-reuse model also assumes the inputs remain available for the
full tiny product.

These are not omissions to hide.

They are part of the declared witness.

A useful performance model states what it counts.

\section{Benchmarking can create false
conclusions}\label{benchmarking-can-create-false-conclusions}

Hardware benchmarks are easy to misread.

A measured number depends on choices such as:

\begin{itemize}
\tightlist
\item
  device model;
\item
  driver/runtime/library versions;
\item
  data type;
\item
  shape;
\item
  warmup;
\item
  synchronization;
\item
  repetitions;
\item
  timing region;
\item
  transfer inclusion;
\item
  compiler options;
\item
  clock/power state;
\item
  kernel selection.
\end{itemize}

A number without this context can be precise and still be
uninterpretable.

\section{Peak versus sustained
performance}\label{peak-versus-sustained-performance}

A vendor peak rate is a hardware specification under declared
operation/type assumptions.

A measured kernel rate is an empirical observation under a workload and
software stack.

An application end-to-end rate includes more than a kernel.

These three quantities should not be mixed.

The safest comparison names which one is being reported.

\section{FLOP/s can reward doing more
work}\label{flops-can-reward-doing-more-work}

A kernel can report more FLOP/s simply because it performs extra
arithmetic that the hardware handles efficiently.

That does not necessarily make the computation finish sooner.

Time-to-solution, energy, memory use, accuracy, and throughput answer
different questions.

A metric is useful only when it matches the declared objective.

\section{Hardware-aware
mathematics}\label{hardware-aware-mathematics}

Hardware considerations can feed back into algorithm design.

One method may use more arithmetic but less memory traffic.

Another may trade recomputation for storage.

A third may select a structured factorization because it maps naturally
to matrix units.

The correct scientific question is not:

\begin{quote}
Is hardware corrupting the mathematics?
\end{quote}

It is:

\begin{quote}
Which mathematically acceptable representation exposes the physical
resources the machine can use efficiently?
\end{quote}

This is a design problem across abstraction levels.

\section{What belongs downstream}\label{what-belongs-downstream}

This chapter is deliberately single-device and kernel-centric.

It does not develop:

\begin{itemize}
\tightlist
\item
  data parallelism;
\item
  tensor parallelism;
\item
  pipeline parallelism;
\item
  expert parallelism;
\item
  collective communication;
\item
  interconnect topology;
\item
  distributed KV caches;
\item
  serving/batching systems;
\item
  multi-device throughput/latency tradeoffs.
\end{itemize}

Those belong in:

\textbf{Scaling, Parallelism, and Serving --- ATLAS-CH-SYSTEMS-001.}

Hardware-001 supplies the substrate that Systems will need:

\begin{itemize}
\tightlist
\item
  bandwidth;
\item
  locality;
\item
  arithmetic intensity;
\item
  precision;
\item
  resource constraints;
\item
  benchmark discipline.
\end{itemize}

\section{Durable boundaries}\label{durable-boundaries-1}

The chapter leaves the reader with nine distinctions:

\begin{enumerate}
\def\labelenumi{\arabic{enumi}.}
\tightlist
\item
  FLOPs are not time.
\item
  Throughput is not latency.
\item
  Equal FLOPs do not imply equal bytes.
\item
  Arithmetic intensity needs a declared memory boundary.
\item
  Peak throughput is not attained throughput.
\item
  Occupancy is not performance.
\item
  Matrix algebra is not guaranteed matrix-unit efficiency.
\item
  Smaller precision is not guaranteed acceptable error.
\item
  Kernel optimization is not distributed-system optimization.
\end{enumerate}

These boundaries are more useful than memorizing one accelerator's
specification sheet.

\section{References used in this
chapter}\label{references-used-in-this-chapter-71}

\begin{itemize}
\tightlist
\item
  (Williams et al. 2009)
\item
  (NVIDIA Corporation 2026b)
\item
  (NVIDIA Corporation 2026a)
\end{itemize}

Exact source identities and authority scopes are pinned in:

sources/source-locks/ATLAS-CH-HARDWARE-001.yaml

\chapter{Scaling, Parallelism, and
Serving}\label{scaling-parallelism-and-serving}

\textbf{Epistemic status:} audited Hardware and Mixture-of-Experts
prerequisites + primary distributed-training/serving systems sources +
Atlas synthesis + exact finite placement/communication witness.\\
\textbf{Specification:}
manuscript/specifications/ATLAS-CH-SYSTEMS-001.md\\
\textbf{Derivation packet:}
mathematics/derivations/ATLAS-CH-SYSTEMS-001-DERIVATIONS.md\\
\textbf{Computational witness:}
mathematics/computational-witnesses/ATLAS-CW-SYSTEMS-001.md\\
\textbf{Source lock:} sources/source-locks/ATLAS-CH-SYSTEMS-001.yaml

A model that fits and runs efficiently on one accelerator is not yet a
distributed system.

Once computation crosses devices, several new objects appear:

\begin{itemize}
\tightlist
\item
  replicas and shards;
\item
  stage boundaries;
\item
  collectives;
\item
  physical links;
\item
  synchronization;
\item
  queues;
\item
  request schedulers;
\item
  distributed memory state;
\item
  stragglers;
\item
  service-level metrics.
\end{itemize}

The governing rule is:

\[
\boxed{
\text{distributed execution cost is a property of computation, communication, placement, scheduling, and measurement together.}
}
\]

This chapter develops that boundary without turning one training stack
or serving system into a universal law.

\section{From hardware to systems}\label{from-hardware-to-systems}

HARDWARE-001 established that:

\begin{itemize}
\tightlist
\item
  FLOPs are not time;
\item
  bytes moved matter;
\item
  arithmetic intensity is boundary-relative;
\item
  peak throughput is not attained throughput;
\item
  latency and throughput are distinct;
\item
  benchmark scope must be declared.
\end{itemize}

Those rules continue to hold across devices.

Distributed execution adds another layer: values may have to move
between devices, and the path itself can become load-bearing.

The same mathematical model can therefore differ in:

\begin{itemize}
\tightlist
\item
  partition;
\item
  placement;
\item
  collective algorithm;
\item
  topology;
\item
  schedule;
\item
  memory replication;
\item
  request batching.
\end{itemize}

Equal model arithmetic does not erase these differences.

\section{A typed systems object}\label{a-typed-systems-object}

Write:

\[
\mathcal S=(R,\Pi,\Gamma,\mathcal C,\mathcal M,\mathcal Q,\mathcal P).
\]

Here:

\begin{itemize}
\tightlist
\item
  \(R\): ranks/devices;
\item
  \(\Pi\): typed parallelism decomposition;
\item
  \(\Gamma\): physical communication topology;
\item
  \(\mathcal C\): collective/point-to-point semantics;
\item
  \(\mathcal M\): memory state and placement;
\item
  \(\mathcal Q\): scheduling/request state;
\item
  \(\mathcal P\): measurement protocol.
\end{itemize}

A claim about one coordinate does not automatically determine the
others.

For example:

\[
\text{tensor-parallel degree}=4
\]

does not tell us:

\begin{itemize}
\tightlist
\item
  which tensor dimension is split;
\item
  which collective follows;
\item
  whether links are intra-node or inter-node;
\item
  how much overlap occurs;
\item
  what latency is measured.
\end{itemize}

\section{Parallelism is not one
thing}\label{parallelism-is-not-one-thing}

Large-model systems commonly combine several different decompositions.

The Atlas keeps four separate:

\begin{enumerate}
\def\labelenumi{\arabic{enumi}.}
\tightlist
\item
  data parallelism;
\item
  tensor parallelism;
\item
  pipeline parallelism;
\item
  expert parallelism.
\end{enumerate}

Narayanan et al.~describe composing data, tensor, and pipeline
parallelism in Megatron-LM and report measured large-cluster trade-offs
in that system.

The Atlas uses that paper as primary evidence that these modes can be
composed and have materially different scaling behavior.

It does not treat the reported configuration as universal.

\section{Data parallelism}\label{data-parallelism}

In synchronous data parallelism, parameter replicas process different
data shards.

Let rank \(r\) compute local gradient:

\[
g_r.
\]

One declared synchronization rule is:

\[
\bar g=\frac1p\sum_{r=0}^{p-1}g_r.
\]

Replicas then apply the same update from \(\bar g\).

The communication is not optional if exact replica synchronization under
this rule is required.

But data parallelism does not by itself specify:

\begin{itemize}
\tightlist
\item
  ring versus tree all-reduce;
\item
  gradient bucketing;
\item
  optimizer-state sharding;
\item
  communication overlap;
\item
  physical placement.
\end{itemize}

Those are additional systems choices.

\section{Tensor parallelism}\label{tensor-parallelism}

Tensor parallelism divides one operator across ranks.

For example, a matrix operation can be partitioned by rows, columns,
heads, or another declared dimension.

Each rank computes a partial result.

The next layer may then require:

\begin{itemize}
\tightlist
\item
  all-reduce;
\item
  all-gather;
\item
  reduce-scatter;
\item
  point-to-point exchange;
\item
  no communication until a later boundary.
\end{itemize}

The exact pattern depends on the partition and surrounding operators.

So:

\[
\boxed{
\text{tensor parallelism}
\neq
\text{one fixed communication pattern}.
}
\]

Megatron-LM is a primary example of practical tensor-parallel
Transformer training.

\section{Pipeline parallelism}\label{pipeline-parallelism}

Pipeline parallelism divides an ordered network into stage groups.

If stages are:

\[
S_1,S_2,\ldots,S_k,
\]

a microbatch flows through them in order.

Different microbatches can occupy different stages concurrently.

GPipe is a primary example of microbatch pipeline parallelism.

The key systems quantities include:

\begin{itemize}
\tightlist
\item
  stage partition;
\item
  microbatch count;
\item
  fill time;
\item
  drain time;
\item
  idle bubbles;
\item
  activation transfers;
\item
  backward transfers during training;
\item
  stage imbalance.
\end{itemize}

More stages can reduce per-stage model memory.

They can also create more communication boundaries and more sensitivity
to imbalance.

\section{A pipeline bubble is
schedule-relative}\label{a-pipeline-bubble-is-schedule-relative}

Suppose one stage finishes while the next required dependency has not
arrived.

That stage is idle.

This idle period belongs to the schedule and dependency graph.

It is not a new mathematical layer in the model.

Pipeline utilization therefore depends on how microbatches are ordered
through stages.

A stage count alone cannot determine it.

\section{Expert parallelism}\label{expert-parallelism-1}

MOE-001 established a chain:

\[
\text{router probabilities}
\to
\text{preferred routes}
\to
\text{capacity/overflow}
\to
\text{accepted dispatch}
\to
\text{expert execution}.
\]

Expert parallelism adds physical placement.

Let:

\[
\operatorname{dev}(e)
\]

be the device hosting expert \(e\).

If token \(i\) originates elsewhere, accepted dispatch induces
communication.

SYSTEMS adds one further object:

\[
\operatorname{path}(\operatorname{src}(i),\operatorname{dev}(e);\Gamma).
\]

Hence:

\[
\boxed{
\text{preference}
\neq
\text{dispatch}
\neq
\text{placement}
\neq
\text{network path}.
}
\]

This distinction prevents a routing diagram from pretending to be a
systems benchmark.

\section{Collective semantics}\label{collective-semantics}

A collective says what distributed values must become.

It does not by itself say how the network achieves that state.

For rank-local values \(x_r\):

\subsection{All-reduce}\label{all-reduce}

Every rank receives the reduction:

\[
y_r=\sum_j x_j.
\]

\subsection{Reduce-scatter}\label{reduce-scatter}

The values are reduced, and declared partitions of the result are
distributed across ranks.

\subsection{All-gather}\label{all-gather}

Declared partitions are gathered so every rank receives the whole
collection.

\subsection{All-to-all}\label{all-to-all}

Each source rank sends destination-specific payloads to every
destination rank.

These are semantic contracts.

Ring, tree, hierarchical, topology-aware, and implementation-specific
algorithms are ways to realize them.

\section{Logical topology and physical
topology}\label{logical-topology-and-physical-topology}

A ring all-reduce has a logical neighbor order.

A cluster has a physical topology.

They are not the same object.

A logical neighbor pair might map to:

\begin{itemize}
\tightlist
\item
  two devices on one board;
\item
  two devices in one node;
\item
  two nodes in one rack;
\item
  two nodes across a slower shared fabric.
\end{itemize}

Placement can therefore change which physical links carry the same
logical collective.

\section{Exact four-rank witness}\label{exact-four-rank-witness}

Use four ranks:

\[
A_0,A_1,B_0,B_1.
\]

Place:

\begin{itemize}
\tightlist
\item
  \(A_0,A_1\) on node A;
\item
  \(B_0,B_1\) on node B.
\end{itemize}

Each rank begins with an 8-element FP32 vector.

Thus:

\[
S=8\cdot4=32\ \text{bytes}.
\]

The collective is sum all-reduce implemented as ring reduce-scatter plus
ring all-gather.

\section{Per-rank ring accounting}\label{per-rank-ring-accounting}

With:

\[
p=4,
\]

each vector is split into four 8-byte chunks.

Reduce-scatter requires three ring steps.

All-gather requires three more.

Each rank sends one chunk per step.

Therefore each rank sends:

\[
6\cdot8=48\ \text{bytes}
\]

and receives 48 bytes.

During reduce-scatter, each rank reduces three 2-scalar chunks:

\[
3\cdot2=6
\]

scalar additions.

The two placements below have exactly the same values.

\section{Grouped placement}\label{grouped-placement}

Use logical ring:

\[
A_0\to A_1\to B_0\to B_1\to A_0.
\]

Only two directed ring edges cross between nodes:

\[
A_1\to B_0,
\qquad
B_1\to A_0.
\]

Each ring edge carries 48 bytes over the complete collective.

Therefore:

\[
Q_{\mathrm{cut}}^G
=
2\cdot48
=
96\ \text{bytes}.
\]

\section{Interleaved placement}\label{interleaved-placement}

Now use:

\[
A_0\to B_0\to A_1\to B_1\to A_0.
\]

Every ring edge crosses the node boundary.

Therefore:

\[
Q_{\mathrm{cut}}^I
=
4\cdot48
=
192\ \text{bytes}.
\]

Nothing changed in the reduction arithmetic.

Nothing changed in per-rank total communication volume.

Only logical-to-physical placement changed.

\section{Equal arithmetic does not imply equal physical
communication}\label{equal-arithmetic-does-not-imply-equal-physical-communication}

Both cases have:

\begin{itemize}
\tightlist
\item
  6 additions per rank;
\item
  24 aggregate additions;
\item
  48 bytes sent per rank;
\item
  48 bytes received per rank.
\end{itemize}

Yet:

\[
Q_{\mathrm{cut}}^I
=
2Q_{\mathrm{cut}}^G.
\]

Therefore:

\[
\boxed{
\text{equal arithmetic and equal per-rank bytes}
\not\Rightarrow
\text{equal cross-cut traffic}.
}
\]

This is the chapter's exact finite control.

\section{A lower bound is not a latency
measurement}\label{a-lower-bound-is-not-a-latency-measurement}

Suppose, only for the toy model, that all cross-node traffic shares
aggregate cut capacity:

\[
B_{\mathrm{cut}}
=
16\ \text{bytes/time-unit}.
\]

Then any execution must satisfy:

\[
T\ge
\frac{Q_{\mathrm{cut}}}{B_{\mathrm{cut}}}.
\]

So:

\[
T_G\ge6,
\qquad
T_I\ge12.
\]

This does not say a real machine takes 6 or 12 units.

Real time can include:

\begin{itemize}
\tightlist
\item
  protocol overhead;
\item
  launch overhead;
\item
  synchronization;
\item
  bidirectional behavior;
\item
  overlap;
\item
  routing;
\item
  congestion;
\item
  software queues;
\item
  device-kernel time.
\end{itemize}

The witness proves an accounting distinction, not a benchmark.

\section{Communication can hide behind
computation}\label{communication-can-hide-behind-computation}

Distributed systems often overlap communication with local arithmetic.

If a communication step lies off the critical path, some of its time can
be hidden.

But hidden is not free.

The bytes still occupy links, buffers, and collective state.

A competing transfer can expose the cost again.

So a serious claim should distinguish:

\begin{itemize}
\tightlist
\item
  communication volume;
\item
  communication duration;
\item
  exposed communication time;
\item
  total step time.
\end{itemize}

\section{Stragglers}\label{stragglers}

A synchronized step often waits for the slowest required participant.

That participant may be slow because of:

\begin{itemize}
\tightlist
\item
  more work;
\item
  slower hardware;
\item
  slower link path;
\item
  cache/memory effects;
\item
  queueing;
\item
  transient interference;
\item
  expert imbalance.
\end{itemize}

Average rank time can therefore fail to predict synchronized step time.

The MoE straggler boundary is a special case of this broader rule.

\section{Training scaling}\label{training-scaling}

For a fixed declared workload and measurement protocol, define
throughput at \(p\) devices:

\[
\mathrm{Th}(p).
\]

A strong-scaling-style efficiency statistic is:

\[
\eta_p=
\frac{\mathrm{Th}(p)}
{p\,\mathrm{Th}(1)}.
\]

The statistic is meaningful only if the comparison states what was held
fixed.

Changing:

\begin{itemize}
\tightlist
\item
  global batch size;
\item
  sequence length;
\item
  precision;
\item
  activation checkpointing;
\item
  optimizer semantics;
\item
  model partition;
\end{itemize}

can change the task being compared.

\section{More devices can reduce
efficiency}\label{more-devices-can-reduce-efficiency}

Additional devices can provide:

\begin{itemize}
\tightlist
\item
  more memory;
\item
  more compute;
\item
  more concurrency.
\end{itemize}

They can also introduce:

\begin{itemize}
\tightlist
\item
  more synchronization;
\item
  more communication;
\item
  more idle imbalance;
\item
  more failure surface.
\end{itemize}

So:

\[
\boxed{
p\uparrow
\not\Rightarrow
\eta_p\uparrow.
}
\]

Nor does low \(\eta_p\) alone identify the cause.

\section{Training and serving optimize different
state}\label{training-and-serving-optimize-different-state}

Training repeatedly performs forward work, backward work, gradient
synchronization, and parameter updates.

Serving receives requests from outside the system.

Autoregressive generation then creates a new state variable: each active
request has a current generated prefix and must be revisited for further
tokens.

The scheduler therefore becomes part of the mechanism.

\section{Request latency}\label{request-latency}

For request \(i\):

\[
L_i=
t_i^{\mathrm{complete}}
-
t_i^{\mathrm{arrival}}.
\]

Different serving questions may care about:

\begin{itemize}
\tightlist
\item
  time to first token;
\item
  time per output token;
\item
  end-to-end completion latency;
\item
  p50 latency;
\item
  p95/p99 tail latency.
\end{itemize}

One number cannot stand for all of them.

\section{Throughput}\label{throughput}

Serving throughput can be measured as:

\begin{itemize}
\tightlist
\item
  requests per second;
\item
  output tokens per second;
\item
  total processed tokens per second.
\end{itemize}

Those are different denominators and can rank systems differently.

A throughput claim must say which one is meant.

\section{Goodput}\label{goodput}

Suppose requests have latency objective \(\tau\).

One possible request goodput is:

\[
G_\tau
=
\frac{|\{i:L_i\le\tau\}|}{T}.
\]

A system can increase raw throughput while violating more latency
objectives.

Thus:

\[
\boxed{
\text{throughput}
\neq
\text{SLO goodput}.
}
\]

\section{Average latency and tail
latency}\label{average-latency-and-tail-latency}

Mean latency is:

\[
\bar L=
\frac1N\sum_i L_i.
\]

A p99 statistic is an order statistic.

The mean does not determine p99.

A system can improve the average while making rare long waits worse.

Tail behavior therefore requires its own evidence.

\section{Iteration-level
scheduling}\label{iteration-level-scheduling}

Autoregressive requests can have different remaining lengths.

A rigid batch can force completed or short requests to wait behind
longer work, depending on system semantics.

Orca introduced iteration-level scheduling, with scheduling decisions at
model-iteration granularity, and selective batching in its reported
system.

The Atlas imports the mechanism and the measured-example status.

It does not import Orca's reported speedup as a universal constant.

\section{Batching is a scheduling
policy}\label{batching-is-a-scheduling-policy}

Batching can increase arithmetic efficiency and throughput by presenting
more parallel work.

It can also increase waiting time before service.

The relevant trade depends on:

\begin{itemize}
\tightlist
\item
  arrival pattern;
\item
  batch formation;
\item
  request lengths;
\item
  device occupancy;
\item
  memory capacity;
\item
  service objectives.
\end{itemize}

Therefore:

\[
\boxed{
\text{larger batch}
\not\Rightarrow
\text{better serving under every metric}.
}
\]

\section{KV-cache state}\label{kv-cache-state}

During autoregressive Transformer serving, previously computed key/value
state can be retained so later token steps do not recompute the entire
prefix attention state.

For request \(i\), use generic accounting:

\[
M_i=\ell_i b_i,
\]

where:

\begin{itemize}
\tightlist
\item
  \(\ell_i\) is retained sequence length;
\item
  \(b_i\) is implementation-specific cache bytes per retained position.
\end{itemize}

Then:

\[
M_{\mathrm{KV}}=\sum_i M_i.
\]

The value \(b_i\) depends on architecture, precision, sharding, and
layout.

\section{Cache allocation can limit
concurrency}\label{cache-allocation-can-limit-concurrency}

Even if model weights fit, active-request cache state can constrain how
many requests remain resident.

Fragmentation and duplication can further reduce usable capacity.

Kwon et al.~identify KV-cache memory-management inefficiency as a
serving bottleneck and introduce PagedAttention/vLLM as one design for
mitigating it.

Their measured gains are evidence about the reported implementation and
workloads.

They are not a hardware-independent theorem.

\section{Prefill and decode are not the same
workload}\label{prefill-and-decode-are-not-the-same-workload}

Processing an input prompt can expose substantial parallel work over
many token positions.

Autoregressive decode advances one or a small number of positions per
active request per iteration.

The arithmetic shape, memory behavior, and batching opportunity can
differ.

A serving benchmark that reports only one aggregate throughput number
can hide this difference.

\section{Utilization}\label{utilization}

A device can be highly utilized while doing work that is not useful
under the service objective.

Examples include:

\begin{itemize}
\tightlist
\item
  padded work;
\item
  work for requests that miss an SLO;
\item
  excess recomputation;
\item
  synchronization overhead;
\item
  speculative work later discarded.
\end{itemize}

Utilization is a resource-activity metric.

It is not automatically useful-work efficiency.

\section{Communication and memory
interact}\label{communication-and-memory-interact}

Tensor or expert parallelism can reduce per-device parameter memory.

That can permit a larger model or larger request batch.

But it can also increase communication.

Conversely, replication can reduce communication on one path while
consuming more memory.

Systems design is therefore multi-objective.

\section{Placement is an optimization
variable}\label{placement-is-an-optimization-variable}

Placement decides where shards, stages, experts, replicas, and cache
state live.

The exact ring witness showed that placement alone can change cross-cut
bytes.

In larger systems, placement can interact with:

\begin{itemize}
\tightlist
\item
  heterogeneous links;
\item
  shared switches;
\item
  NUMA effects;
\item
  rack boundaries;
\item
  expert traffic;
\item
  request locality.
\end{itemize}

The Atlas does not prescribe one placement algorithm.

It requires the placement to be part of the claim.

\section{Measured versus modeled
performance}\label{measured-versus-modeled-performance}

A model may predict:

\begin{itemize}
\tightlist
\item
  bytes;
\item
  message count;
\item
  lower bounds;
\item
  asymptotic scaling;
\item
  critical-path structure.
\end{itemize}

A benchmark measures a concrete realization.

A mismatch is not automatically a failure of arithmetic.

It may reveal omitted effects.

The correct response is to name those effects rather than quietly
relabel a model as a measurement.

\section{Training benchmark record}\label{training-benchmark-record}

A serious training record should identify:

\begin{itemize}
\tightlist
\item
  model/workload;
\item
  batch semantics;
\item
  precision;
\item
  optimizer/update;
\item
  hardware count/type;
\item
  topology;
\item
  parallelism degrees;
\item
  sharding;
\item
  warmup;
\item
  measurement interval;
\item
  throughput unit;
\item
  scaling baseline.
\end{itemize}

Without these, ``X\% scaling efficiency'' is underspecified.

\section{Serving benchmark record}\label{serving-benchmark-record}

A serious serving record should identify:

\begin{itemize}
\tightlist
\item
  model and precision;
\item
  device count/type and topology;
\item
  request/input-length distribution;
\item
  output-length distribution;
\item
  concurrency or arrival process;
\item
  scheduler/batching;
\item
  KV-cache policy;
\item
  latency percentiles;
\item
  throughput/goodput definition;
\item
  SLO if used;
\item
  measurement duration.
\end{itemize}

Without these, ``requests per second'' cannot be interpreted safely.

\section{What the sources
establish}\label{what-the-sources-establish}

Narayanan et al.~provide primary evidence for composed
data/tensor/pipeline training and their measured trade-offs in
Megatron-LM.

Huang et al.~provide primary evidence for microbatch pipeline execution
in GPipe.

Yu et al.~provide primary evidence for iteration-level
scheduling/selective batching in Orca.

Kwon et al.~provide primary evidence for KV-cache memory-management
pressure and PagedAttention/vLLM.

These sources demonstrate representative mechanisms.

They do not close the design space.

\section{The systems layer is not one
scalar}\label{the-systems-layer-is-not-one-scalar}

A distributed system can be described by a resource/performance vector
such as:

\[
K=
(
W,
Q_{\mathrm{local}},
Q_{\mathrm{net}},
M,
L_{50},
L_{99},
\mathrm{Th},
G,
U
).
\]

Here the entries might denote arithmetic work, local bytes, network
bytes, memory, median/tail latency, throughput, goodput, and
utilization.

The exact vector depends on the task.

Collapsing it to one score requires a declared objective.

\section{Failure boundaries}\label{failure-boundaries-1}

The durable non-implications are:

\[
\text{more devices}
\not\Rightarrow
\text{linear speedup},
\]

\[
\text{equal FLOPs}
\not\Rightarrow
\text{equal communication},
\]

\[
\text{equal per-rank bytes}
\not\Rightarrow
\text{equal cross-cut traffic},
\]

\[
\text{high throughput}
\not\Rightarrow
\text{low latency},
\]

\[
\text{low average latency}
\not\Rightarrow
\text{low tail latency},
\]

\[
\text{high utilization}
\not\Rightarrow
\text{high useful-work efficiency},
\]

\[
\text{training efficiency}
\not\Rightarrow
\text{serving efficiency}.
\]

These are not pessimistic slogans.

They are type checks for systems claims.

\section{Systems reasoning}\label{systems-reasoning}

The chapter's systems method is:

\begin{enumerate}
\def\labelenumi{\arabic{enumi}.}
\tightlist
\item
  name the computation;
\item
  name the partition;
\item
  name the communication semantics;
\item
  name the physical placement/topology;
\item
  account for bytes and memory;
\item
  name the schedule;
\item
  define latency/throughput/tail metrics;
\item
  distinguish models from measurements;
\item
  preserve the workload and benchmark boundary;
\item
  compare alternatives only under declared objectives.
\end{enumerate}

That method scales better than treating ``distributed'' as a performance
adjective.

\section{References used in this
chapter}\label{references-used-in-this-chapter-72}

\begin{itemize}
\tightlist
\item
  Deepak Narayanan et al.~(2021), \emph{Efficient Large-Scale Language
  Model Training on GPU Clusters Using Megatron-LM}, arXiv:2104.04473.
\item
  Yanping Huang et al.~(2019), \emph{GPipe: Efficient Training of Giant
  Neural Networks using Pipeline Parallelism}, arXiv:1811.06965.
\item
  Gyeong-In Yu et al.~(2022), \emph{Orca: A Distributed Serving System
  for Transformer-Based Generative Models}, OSDI 22, pp.~521--538.
\item
  Woosuk Kwon et al.~(2023), \emph{Efficient Memory Management for Large
  Language Model Serving with PagedAttention}, arXiv:2309.06180.
\end{itemize}

Exact source identities, inherited prerequisite authority, and claim
boundaries are recorded in:

sources/source-locks/ATLAS-CH-SYSTEMS-001.yaml

\chapter{Experiments as Arguments}\label{experiments-as-arguments}

\textbf{Epistemic status:} Atlas synthesis of established
experimental-design, causal-inference, statistical-reporting,
reproducibility, and machine-learning experimental-method sources, with
one exact finite computational witness.\\
\textbf{Specification:}
manuscript/specifications/ATLAS-CH-EXPERIMENT-001.md\\
\textbf{Derivation packet:}
mathematics/derivations/ATLAS-CH-EXPERIMENT-001-DERIVATIONS.md\\
\textbf{Computational witness:}
mathematics/computational-witnesses/ATLAS-CW-EXPERIMENT-001.md\\
\textbf{Source lock:} sources/source-locks/ATLAS-CH-EXPERIMENT-001.yaml

A machine-learning experiment can be expensive, carefully logged,
repeated across seeds, and still fail to answer the question its authors
think it answers.

That is not because experiments are useless.

It is because \textbf{a run is not yet an argument}.

A run produces observations.

An argument explains why those observations bear on a particular claim.

The distinction is easy to lose when the object on the screen is a
single number: accuracy, loss, reward, latency, FLOPs, energy,
perplexity, exact-match rate.

Numbers look decisive.

But before a number can carry evidentiary force, we must ask what was
changed, what stayed fixed, what population or task family the
comparison refers to, what randomness was sampled, what nuisance
variables moved at the same time, what selection process chose the
reported result, and what claim the result is actually allowed to
support.

This chapter develops experiments as \textbf{structured claim-support
objects}.

The governing idea is:

\begin{quote}
an experiment is a designed contrast whose evidentiary force comes from
the relation among the claim, the estimand, the intervention, the
comparison, the variation model, and the scope.
\end{quote}

The chapter does not propose one universal statistical recipe.

Different scientific questions need different designs.

What must remain stable is the logical discipline connecting design to
claim.

\section{From metric to argument}\label{from-metric-to-argument}

Suppose two systems are evaluated and one obtains score 82 while the
other obtains score 79.

What has been established?

At minimum, under the exact evaluation protocol that produced those
numbers, one recorded outcome is larger than the other.

That statement may be useful.

But many stronger statements do not follow automatically.

It does not yet establish:

\begin{itemize}
\tightlist
\item
  that the architectural difference caused the improvement;
\item
  that the improvement persists across datasets;
\item
  that the improvement persists across seeds;
\item
  that the effect is practically important;
\item
  that a named mechanism explains the difference;
\item
  that the result generalizes outside the benchmark;
\item
  that the experiment would survive a fresh implementation;
\item
  that the comparison was selected before seeing the results.
\end{itemize}

The Evidence chapter already supplied the Atlas rule that support must
retain its scope.

EXPERIMENT-001 specializes that rule for empirical and computational
comparisons.

We will write an experiment-argument object as

\texttt{E\ =\ (q,\ theta,\ U,\ A,\ Z,\ Y,\ g,\ V,\ rho,\ Omega)}.

The notation is deliberately compact.

It forces ten questions into view.

\begin{itemize}
\tightlist
\item
  \texttt{q}: What bounded claim is under examination?
\item
  \texttt{theta}: What quantity is the experiment meant to estimate?
\item
  \texttt{U}: What units, tasks, datasets, or population does the design
  concern?
\item
  \texttt{A}: What intervention or assignment rule creates the
  comparison?
\item
  \texttt{Z}: What nuisance variables, blocks, or controlled conditions
  matter?
\item
  \texttt{Y}: What outcome is measured?
\item
  \texttt{g}: What estimator or comparison rule converts observations
  into a contrast?
\item
  \texttt{V}: What sources of variation or uncertainty are represented?
\item
  \texttt{rho}: What stopping, tuning, selection, and reporting rule
  generated the reported result?
\item
  \texttt{Omega}: What scope is the conclusion allowed to inhabit?
\end{itemize}

This tuple is an Atlas synthesis.

It is not an externally standardized definition of experiment.

Its purpose is diagnostic: if a load-bearing field is missing, the
argument may be weaker than the prose suggests.

\section{An intervention is not an
observation}\label{an-intervention-is-not-an-observation}

The NIST engineering-statistics handbook begins from deliberate
variation of factors and observation of responses, and treats nuisance
factors, blocking, randomization, and replication as central design
concerns (National Institute of Standards and Technology, n.d.).

That intervention logic matters because observational association and
intervention answer different questions.

Suppose high-capacity models obtain higher scores than low-capacity
models in a collection of published results.

That is an association.

Perhaps capacity helps.

Perhaps the high-capacity models were also trained on more data, tuned
more aggressively, evaluated later in the field's history, or
implemented by teams with more compute.

The observation does not identify which difference produced the outcome.

A designed intervention attempts to break that ambiguity.

Change the factor of interest while holding relevant alternatives fixed,
balancing them, randomizing them, or modeling them explicitly.

The resulting contrast can then bear more directly on the intended
claim.

Randomization is one important mechanism for doing this, but it is not
magic.

It addresses assignment-related confounding under the design.

It does not automatically repair:

\begin{itemize}
\tightlist
\item
  bad measurement;
\item
  coding defects;
\item
  selective reporting;
\item
  inappropriate outcomes;
\item
  protocol drift;
\item
  hidden preprocessing;
\item
  nonrepresentative sampling;
\item
  generalization beyond the study scope.
\end{itemize}

Experimental control is therefore specific.

One must always ask: \textbf{what source of ambiguity did this design
actually remove?}

\section{Hypothesis versus
estimand}\label{hypothesis-versus-estimand}

A hypothesis is often verbal:

\begin{quote}
the intervention improves performance.
\end{quote}

The phrase ``improves performance'' is not yet a mathematical object.

Improve what?

Over which population?

By what aggregation?

Under which resource budget?

Relative to what alternative?

An \textbf{estimand} makes the target explicit.

In a simple finite causal setting, imagine units indexed by
\texttt{i=1,...,n}.

Each unit has two potential outcomes:

\begin{itemize}
\tightlist
\item
  \texttt{Y\_i(1)}: the outcome under treatment;
\item
  \texttt{Y\_i(0)}: the outcome under control.
\end{itemize}

The finite-population average treatment effect is

\texttt{Delta\ =\ (1/n)\ sum\_i\ {[}Y\_i(1)-Y\_i(0){]}}.

Rubin's treatment-effect framework made the benefits of randomization
and carefully controlled nonrandomized comparisons explicit in this
language of treatment effects and assignment (Rubin 1974).

The difficult point is immediate.

For a particular unit, we observe only the outcome corresponding to the
treatment actually assigned.

The other outcome is counterfactual.

So an experiment does not merely ``measure the effect.''

It constructs a design under which observed contrasts can support claims
about an otherwise unobserved comparison.

This is why an estimand is not the same thing as a metric.

Accuracy can be a metric.

``The average change in accuracy caused by replacing component X with Y
under a fixed training and evaluation protocol'' is an estimand.

The difference is not stylistic.

It determines what the experiment is trying to learn.

\section{Baselines are comparators; ablations are
interventions}\label{baselines-are-comparators-ablations-are-interventions}

Machine-learning papers often use the words \textbf{baseline} and
\textbf{ablation} as if they were interchangeable.

They are not.

A baseline is any comparator.

An ablation is a structured intervention on a component of a reference
system.

If a model \texttt{S} contains component \texttt{c}, an ablation might
compare \texttt{S} against \texttt{S\_\{-c\}}, where the component is
removed or replaced.

The intended mechanism claim is something like:

\begin{quote}
the presence of component c contributes to the observed behavior.
\end{quote}

But the validity of that interpretation depends on what else changes.

Removing a component might also change:

\begin{itemize}
\tightlist
\item
  parameter count;
\item
  compute;
\item
  activation scale;
\item
  optimization stability;
\item
  training time;
\item
  receptive field;
\item
  memory footprint;
\item
  initialization;
\item
  interface shape.
\end{itemize}

If those effects are not controlled, the ablation is not a pure
intervention on one semantic mechanism.

It is a compound intervention.

This does not make the experiment worthless.

It changes the claim.

A compound ablation can show that one system configuration differs from
another.

It may not isolate why.

The useful rule is:

\begin{quote}
an ablation supports mechanism evidence only to the degree that the
intervention isolates the candidate mechanism from relevant alternative
changes.
\end{quote}

\section{The confounding trap}\label{the-confounding-trap}

A comparison can be numerically exact and logically wrong for the
intended question.

Consider two nuisance strata: easy and hard.

Every easy unit has potential outcomes

\begin{itemize}
\tightlist
\item
  control: 10;
\item
  treatment: 11.
\end{itemize}

Every hard unit has

\begin{itemize}
\tightlist
\item
  control: 0;
\item
  treatment: 1.
\end{itemize}

Treatment helps every unit by exactly 1.

So the exact finite-population average treatment effect is +1.

Now assign units badly.

The treatment group receives:

\begin{itemize}
\tightlist
\item
  1 easy unit;
\item
  9 hard units.
\end{itemize}

The control group receives:

\begin{itemize}
\tightlist
\item
  9 easy units;
\item
  1 hard unit.
\end{itemize}

The treatment mean is

\texttt{(11\ +\ 9*1)/10\ =\ 2}.

The control mean is

\texttt{(9*10\ +\ 0)/10\ =\ 9}.

The naive aggregate contrast is therefore

\texttt{2\ -\ 9\ =\ -7}.

The aggregate says treatment is worse.

But within the easy stratum,

\texttt{11\ -\ 10\ =\ +1}.

Within the hard stratum,

\texttt{1\ -\ 0\ =\ +1}.

If we compare treatment and control within strata and then standardize
with equal stratum weights, the contrast is +1.

The exact computational witness for this chapter reproduces those
values.

Nothing probabilistic is hiding in the example.

The sign reversal is caused entirely by imbalance in the nuisance
variable.

This is a minimal model of confounding.

NIST's blocking guidance describes the corresponding design idea:
nuisance factors can be held constant within blocks so that the factor
of interest can be compared without mixing the block effect into the
contrast (National Institute of Standards and Technology, n.d.).

The lesson is not ``always stratify.''

The lesson is:

\begin{quote}
aggregate differences answer aggregate questions. If treatment
assignment changes the mixture of nuisance conditions, the aggregate
contrast may not estimate the treatment effect you intended.
\end{quote}

\section{Controls do not all play the same
role}\label{controls-do-not-all-play-the-same-role}

The word \textbf{control} can refer to several distinct devices.

A control group may represent the no-treatment alternative.

A baseline model may represent the incumbent method.

A placebo-like condition may preserve everything except the hypothesized
active mechanism.

A negative control may be chosen because no effect is expected.

A positive control may verify that the measurement system can detect a
known effect.

A blocked design may control nuisance variation by comparing like with
like.

These are not synonyms.

The useful question is not ``does the experiment have a control?''

It is:

\begin{quote}
which alternative explanation is this control designed to eliminate?
\end{quote}

A good control has a job.

If its job is not stated, the reader cannot tell what ambiguity it
resolves.

\section{Counterfactuals are the missing half of the
comparison}\label{counterfactuals-are-the-missing-half-of-the-comparison}

Whenever we claim that an intervention \textbf{caused} an outcome, we
are implicitly comparing what happened with what would have happened
otherwise.

That ``otherwise'' is counterfactual.

In randomized experiments, the design uses other assigned units to
estimate the missing alternative.

In matched or blocked observational studies, the design attempts to
build a credible comparison class.

In simulations, the counterfactual may be directly computable because
the same synthetic state can be rerun under two interventions.

In machine-learning ablations, the counterfactual may be ``the same
training/evaluation process without this component.''

But the phrase \textbf{the same} is doing work.

If the counterfactual run changes multiple factors, the claimed
mechanism becomes ambiguous.

A counterfactual therefore requires an invariance contract:

which parts of the world are intended to remain fixed while the
intervention changes?

The stronger the causal interpretation, the more carefully that contract
must be defended.

\section{Seeds: useful, narrow, often
overinterpreted}\label{seeds-useful-narrow-often-overinterpreted}

Random seeds are important in modern machine learning because many
pipelines contain stochastic elements:

\begin{itemize}
\tightlist
\item
  parameter initialization;
\item
  minibatch ordering;
\item
  dropout masks;
\item
  augmentation;
\item
  environment trajectories;
\item
  randomized search;
\item
  sampling.
\end{itemize}

Henderson et al.~documented how nondeterminism, algorithmic variance,
baseline choice, and reporting practice complicate interpretation in
deep reinforcement learning (Henderson et al. 2018).

The general lesson is broader than RL, but the scope must remain
explicit.

Suppose an outcome can be written schematically as

\texttt{Y\ =\ F(a,\ d,\ s,\ h,\ e)}

with

\begin{itemize}
\tightlist
\item
  algorithm \texttt{a};
\item
  data \texttt{d};
\item
  seed \texttt{s};
\item
  hyperparameter/search choice \texttt{h};
\item
  environment and implementation \texttt{e}.
\end{itemize}

If we vary only \texttt{s}, we learn about run-to-run variation
conditional on \texttt{a,d,h,e}.

That can be valuable.

But it does not automatically estimate:

\begin{itemize}
\tightlist
\item
  uncertainty across datasets;
\item
  shift in the target population;
\item
  variability across independent implementations;
\item
  sensitivity to preprocessing;
\item
  uncertainty induced by hyperparameter search;
\item
  uncertainty about the scientific mechanism.
\end{itemize}

``Five seeds'' is therefore not a universal unit of experimental
reliability.

It is a statement about one coordinate of variation.

The experiment should say what that coordinate represents.

\section{Selection rules are part of the
experiment}\label{selection-rules-are-part-of-the-experiment}

Suppose we train 100 configurations and report the best one.

The reported score is not the observation from a prespecified
configuration.

It is the maximum of a search process.

If each observed score is

\texttt{X\_j\ =\ mu\_j\ +\ epsilon\_j},

and the selected run is

\texttt{j*\ =\ argmax\_j\ X\_j},

then the reported statistic is

\texttt{max\_j\ X\_j}.

That selection changes its interpretation.

The same principle applies when we choose:

\begin{itemize}
\tightlist
\item
  the best seed;
\item
  the best checkpoint;
\item
  the best prompt;
\item
  the best benchmark subset;
\item
  the best hyperparameter sweep;
\item
  the best architecture after many failed variants.
\end{itemize}

Selection is not forbidden.

Hidden selection is the problem.

The Atlas therefore places stopping, tuning, selection, and reporting
inside the experiment object as \texttt{rho}.

If \texttt{rho} changes after observing outcomes, the experimental
argument changes too.

\section{Effect size is not statistical
significance}\label{effect-size-is-not-statistical-significance}

A statistically significant result can be tiny.

A practically important effect can fail a significance threshold in a
small or noisy study.

The American Statistical Association's statement on p-values explicitly
separates p-values from effect size and practical importance, and warns
against making scientific conclusions depend only on threshold crossing
(Wasserstein and Lazar 2016).

For the Atlas, this yields a simple discipline.

When an experiment supports a quantitative claim, report the quantity of
interest.

Do not replace it with the fact that some test crossed a threshold.

Depending on the design, useful objects may include:

\begin{itemize}
\tightlist
\item
  mean difference;
\item
  relative change;
\item
  risk difference;
\item
  odds ratio;
\item
  standardized difference;
\item
  confidence or credible interval;
\item
  predictive interval;
\item
  distributional comparison;
\item
  task-level effect distribution.
\end{itemize}

The correct object depends on the estimand.

The principle does not.

\textbf{Magnitude and uncertainty should remain visible.}

A p-value may be part of the argument.

It cannot stand in for the argument.

\section{Benchmark score versus mechanism
evidence}\label{benchmark-score-versus-mechanism-evidence}

Suppose a new architecture beats a baseline on five benchmarks.

That is evidence about comparative performance under those benchmark
protocols.

It is not automatically evidence for the proposed explanation.

A model may improve because of:

\begin{itemize}
\tightlist
\item
  the claimed architectural mechanism;
\item
  increased compute;
\item
  parameter count;
\item
  better optimization;
\item
  hidden regularization;
\item
  preprocessing;
\item
  data leakage;
\item
  implementation quality;
\item
  changed training duration;
\item
  favorable hyperparameter search.
\end{itemize}

Mechanism evidence requires a design that separates the proposed
mechanism from plausible alternatives.

This is why a chapter about experiments must treat ablations, controls,
and counterfactuals as structural objects rather than publication
rituals.

A benchmark answers:

\begin{quote}
what happened under this benchmark protocol?
\end{quote}

A mechanism experiment asks:

\begin{quote}
what change caused or mediated the difference?
\end{quote}

Those are related questions.

They are not the same question.

\section{Reproducibility is valuable and not
sufficient}\label{reproducibility-is-valuable-and-not-sufficient}

The National Academies report adopts a useful terminology (National
Academies of Sciences, Engineering, and Medicine 2019).

In that report:

\begin{itemize}
\tightlist
\item
  \textbf{reproducibility} means obtaining consistent computational
  results using the same input data, code or methods, and analysis
  conditions;
\item
  \textbf{replicability} means obtaining consistent results across
  studies aimed at the same scientific question using newly obtained
  data.
\end{itemize}

The distinction matters because reproducibility is often mistaken for
validity.

A perfectly reproducible pipeline can reproduce a flawed analysis
exactly.

It can reproduce:

\begin{itemize}
\tightlist
\item
  a confounded comparison;
\item
  a data leak;
\item
  a coding bug;
\item
  a biased sample;
\item
  a misleading metric;
\item
  a selectively chosen result.
\end{itemize}

Reproducibility answers an identity-and-reexecution question:

\begin{quote}
can the computation be reconstructed and obtain consistent results?
\end{quote}

Validity asks whether the design and reasoning support the claim.

Replication asks whether a fresh study aimed at the same scientific
question yields consistent findings, subject to its own uncertainty and
design.

Generalization asks whether the result survives beyond the original
context.

These are different evidentiary axes.

No one axis should silently substitute for the others.

\section{Falsifiability belongs in the
design}\label{falsifiability-belongs-in-the-design}

A claim becomes experimentally useful when the design makes clear what
result would count against it.

That does not require reducing every scientific claim to one
null-hypothesis test.

It does require a rejection, revision, or failure condition.

For example:

\begin{itemize}
\tightlist
\item
  if a claimed component matters, an isolated ablation should change the
  target outcome in a declared direction or range;
\item
  if a scaling law is claimed, specified deviations should count against
  the proposed law;
\item
  if a mechanism is claimed to improve stability, instability metrics
  should be defined before observing the candidate;
\item
  if two methods are claimed equivalent, a tolerance and comparison
  domain should be stated.
\end{itemize}

A design that cannot lose is not a strong experimental argument.

A useful experiment risks changing our mind.

\section{Repetition, replication, and
robustness}\label{repetition-replication-and-robustness}

Running the same code twice is not the same as repeating a randomized
study.

Running more seeds is not the same as collecting new data.

Testing more benchmarks is not automatically the same as replication.

Using another implementation is not automatically the same as
generalization.

The words matter because they describe different sources of variation.

A mature experimental report should say which axes were varied.

For example:

\begin{itemize}
\tightlist
\item
  same code, same data, same seed: deterministic replay check;
\item
  same code and data, different seeds: stochastic execution variation;
\item
  same method, new dataset sampled from same frame: replication-like
  evidence;
\item
  new implementation, same intended method: implementation robustness;
\item
  shifted population or domain: external generalization;
\item
  altered nuisance conditions: robustness or stress testing.
\end{itemize}

The more precisely these are named, the less likely one form of evidence
is to be mistaken for another.

\section{The experiment table}\label{the-experiment-table}

A practical way to force the argument into view is to write the
experiment as a table before running it.

{\def\LTcaptype{none} % do not increment counter
\begin{longtable}[]{@{}
  >{\raggedright\arraybackslash}p{(\linewidth - 2\tabcolsep) * \real{0.5000}}
  >{\raggedright\arraybackslash}p{(\linewidth - 2\tabcolsep) * \real{0.5000}}@{}}
\toprule\noalign{}
\begin{minipage}[b]{\linewidth}\raggedright
Field
\end{minipage} & \begin{minipage}[b]{\linewidth}\raggedright
Question
\end{minipage} \\
\midrule\noalign{}
\endhead
\bottomrule\noalign{}
\endlastfoot
claim \texttt{q} & What proposition may change status? \\
estimand \texttt{theta} & What quantity is the experiment trying to
learn? \\
units \texttt{U} & What population, tasks, examples, environments, or
runs are sampled? \\
intervention \texttt{A} & What is deliberately changed or assigned? \\
nuisance \texttt{Z} & What else may affect the outcome? \\
outcome \texttt{Y} & What is measured and why is it relevant? \\
comparison \texttt{g} & How are observations converted into the target
contrast? \\
variation \texttt{V} & Which uncertainty sources are sampled or
estimated? \\
selection \texttt{rho} & How are stopping, tuning, and reporting
determined? \\
scope \texttt{Omega} & Where may the conclusion be applied? \\
\end{longtable}
}

This table is not a bureaucracy.

It is a compression of the experimental argument.

If the table cannot be completed, that is useful information before
compute is spent.

\section{Failure modes}\label{failure-modes-18}

Experiments fail in characteristic ways.

\subsection{Confounded comparison}\label{confounded-comparison}

The treatment groups differ in nuisance factors that also affect the
outcome.

\subsection{Baseline asymmetry}\label{baseline-asymmetry}

One method receives more tuning, compute, data, or engineering effort.

\subsection{Hidden multiplicity}\label{hidden-multiplicity}

Many hypotheses or configurations are tried, but the report presents the
selected result as if it were prespecified.

\subsection{Metric substitution}\label{metric-substitution}

The measured metric drifts away from the actual scientific or
engineering objective.

\subsection{Seed laundering}\label{seed-laundering}

Seed variation is presented as if it covered data, population, or
implementation uncertainty.

\subsection{Ablation overclaim}\label{ablation-overclaim}

A compound system change is interpreted as evidence for one isolated
mechanism.

\subsection{Reproducibility overclaim}\label{reproducibility-overclaim}

Exact rerun is presented as evidence that the scientific conclusion is
valid.

\subsection{Benchmark overclaim}\label{benchmark-overclaim}

Performance on a fixed benchmark is promoted into a universal capability
or mechanism claim.

\subsection{Scope drift}\label{scope-drift}

A result established on one domain is written as if it held across a
broader population.

These failures have one family resemblance:

\textbf{the claim becomes stronger than the design.}

\section{Experiments in adaptive-intelligence
research}\label{experiments-in-adaptive-intelligence-research}

Adaptive systems make experimental design unusually difficult because
the object under study may itself change its behavior in response to:

\begin{itemize}
\tightlist
\item
  data order;
\item
  feedback;
\item
  memory;
\item
  routing;
\item
  learned state;
\item
  tool outputs;
\item
  environment dynamics;
\item
  resource constraints.
\end{itemize}

A static A/B comparison may therefore hide path dependence.

For these systems, the unit of intervention may need to be:

\begin{itemize}
\tightlist
\item
  an episode;
\item
  a trajectory;
\item
  a task sequence;
\item
  a training curriculum;
\item
  a memory state;
\item
  a coordination graph;
\item
  a bounded agent policy.
\end{itemize}

The central discipline remains unchanged.

Name the unit.

Name the intervention.

Name the comparison.

Name the sources of variation.

Name the stopping rule.

Name the scope.

Then state only the claim that this design supports.

\section{What this chapter
establishes}\label{what-this-chapter-establishes-2}

This chapter does not certify any empirical architecture.

It establishes a reading and design grammar.

A downstream reader may now ask of any experiment:

\begin{enumerate}
\def\labelenumi{\arabic{enumi}.}
\tightlist
\item
  What is the claim?
\item
  What is the estimand?
\item
  What creates the contrast?
\item
  What nuisance variables remain?
\item
  What kind of baseline or ablation is used?
\item
  Which randomness is sampled?
\item
  What effect magnitude is observed?
\item
  What uncertainty is represented?
\item
  How was the result selected?
\item
  What is reproducible?
\item
  What is replicated?
\item
  What mechanism, if any, is actually identified?
\item
  Where does the claim stop?
\end{enumerate}

The exact confounding witness demonstrates why these questions matter.

A clean number can point in the wrong direction when the design mixes
treatment with nuisance structure.

The cure is not statistical ornament.

It is experimental clarity.

\section{Atlas handoff}\label{atlas-handoff-2}

\textbf{Replayable Evidence Objects --- ATLAS-CH-REPLAY-001.}\\
Replay may assume that an experiment has a declared claim, estimand,
intervention/comparison structure, variation model, selection rule, and
scope. Replay can then bind the exact data, code, environment,
parameters, seeds, and outputs needed to reconstruct that experiment.

The dependency direction matters.

Replaying an experiment does not create its causal or scientific
validity.

It preserves the experiment that was actually performed.

That preservation is powerful only when the experiment's argument was
explicit in the first place.

\section{References used in this
chapter}\label{references-used-in-this-chapter-73}

\begin{itemize}
\tightlist
\item
  NIST/SEMATECH e-Handbook of Statistical Methods, experimental-design
  sections (National Institute of Standards and Technology, n.d.).
\item
  Rubin, \emph{Estimating causal effects of treatments in randomized and
  nonrandomized studies} (Rubin 1974).
\item
  Wasserstein and Lazar, \emph{The ASA Statement on p-Values: Context,
  Process, and Purpose} (Wasserstein and Lazar 2016).
\item
  National Academies of Sciences, Engineering, and Medicine,
  \emph{Reproducibility and Replicability in Science} (National
  Academies of Sciences, Engineering, and Medicine 2019).
\item
  Henderson et al., \emph{Deep Reinforcement Learning That Matters}
  (Henderson et al. 2018).
\end{itemize}

Exact source identities and authority boundaries are recorded in:

sources/source-locks/ATLAS-CH-EXPERIMENT-001.yaml

\chapter{Replayable Evidence Objects}\label{replayable-evidence-objects}

\textbf{Epistemic status:} methodological synthesis grounded in
reproducibility and provenance literature, exact GCL lifecycle records,
and an executable Atlas replay package.\\
\textbf{Primary figure:} `ATLAS-FIG-REPLAY-001`\\
\textbf{Documentary packet:}
`mathematics/derivations/ATLAS-CH-REPLAY-001-DERIVATIONS.md`\\
\textbf{Computational witness:}
`mathematics/computational-witnesses/ATLAS-CW-REPLAY-001.md`

\section{A result should be
reconstructable}\label{a-result-should-be-reconstructable}

A scientific result is often presented as a sentence:

\begin{quote}
We found \(X\).
\end{quote}

But the sentence is the end of a path.

Before the claim came:

\begin{itemize}
\tightlist
\item
  a source;
\item
  a definition of the object;
\item
  a method;
\item
  code or derivation;
\item
  an execution environment;
\item
  an observation;
\item
  an interpretation;
\item
  some form of review.
\end{itemize}

If those links are lost, the claim may remain readable while the reason
for believing it becomes opaque.

This chapter treats documentary provenance as part of the scientific
object.

Its central distinction is:

\[
\boxed{
\text{replayability is a property of the evidence path}.
}
\]

Truth, verification, certification, and authority are separate
properties.

They may depend on replay.

They are not created by replay.

\section{The chain of custody}\label{the-chain-of-custody}

The useful allegory is a \textbf{chain of custody}.

An exhibit is useful only if its identity and transformation history
remain traceable.

The correspondence is:

\begin{itemize}
\tightlist
\item
  evidence label ↔ stable artifact identity;
\item
  custody record ↔ provenance chain;
\item
  laboratory procedure ↔ replay method;
\item
  analyst conclusion ↔ interpretation;
\item
  formal disposition ↔ review or certification state.
\end{itemize}

The limit is essential.

Scientific truth is not determined by legal procedure.

Institutional acceptance is not mathematical truth.

A correct theorem can exist before any institution recognizes it.

A perfectly documented computation can still be wrong.

The allegory teaches traceability only.

\section{Reproducibility, replication, and Atlas
replay}\label{reproducibility-replication-and-atlas-replay}

Terminology matters because similar words carry different expectations.

The National Academies distinguishes \textbf{reproducibility} from
\textbf{replicability}. Reproducibility concerns obtaining consistent
computational results using the same data, computational steps, methods,
code, and analysis conditions. Replicability concerns obtaining
consistent results in a new study with new data (National Academies of
Sciences, Engineering, and Medicine 2019).

The Atlas uses \textbf{replay} more narrowly:

\begin{quote}
an identity-bound reconstruction of a specified evidence path.
\end{quote}

Replay may be internal or external.

Replay may be deterministic or stochastic.

Replay can be exact at the byte level or approximate at a result level,
provided the contract states which.

The term is local to the Atlas/GCL context.

It is not proposed as a replacement for broader scientific terminology.

\section{A provisional evidence
object}\label{a-provisional-evidence-object}

The Atlas uses

\[
\boxed{
E=(C,S,M,A,O,I,R)
}
\]

as an explanatory model.

The fields are:

\[
C
=
\text{claim identity and exact wording},
\]

\[
S
=
\text{source identities},
\]

\[
M
=
\text{method, derivation, code, or proof object},
\]

\[
A
=
\text{execution environment and artifact identities},
\]

\[
O
=
\text{observations or outputs},
\]

\[
I
=
\text{interpretation and claim boundary},
\]

\[
R
=
\text{review, replay, adjudication, or certification records}.
\]

The tuple is not a universal schema.

It is a checklist against a recurrent failure:

\begin{quote}
``We reproduced it'' without saying what ``it'' was.
\end{quote}

\section{Why claim identity comes
first}\label{why-claim-identity-comes-first}

Suppose two files contain the same output:

\[
86.
\]

One claim might be:

\begin{quote}
this candidate has score 86 under evaluator \(E\).
\end{quote}

Another might be:

\begin{quote}
86 is globally optimal.
\end{quote}

The bytes are identical.

The claims are not.

The first may be supported by replaying an evaluator on an exact
candidate.

The second requires a different argument entirely.

Therefore the evidence object begins with

\[
C,
\]

the exact claim.

A replay object without claim identity is a pile of artifacts waiting
for someone to overinterpret it.

\section{Source identity}\label{source-identity}

A source should be identified strongly enough that another actor can
recover the same object.

A mutable URL is often insufficient.

A branch name can move.

A ``latest'' tag can change.

A webpage can be edited.

A dataset can be refreshed.

Stronger identities include:

\begin{itemize}
\tightlist
\item
  immutable commit SHA;
\item
  Git blob identity;
\item
  content digest;
\item
  DOI plus exact version where applicable;
\item
  archived source snapshot;
\item
  source manifest binding multiple objects.
\end{itemize}

Content addressing does not prove correctness.

It answers a narrower question:

\begin{quote}
Are these the bytes we named?
\end{quote}

That narrow answer is foundational.

\section{Method identity}\label{method-identity}

The method field \(M\) may contain very different kinds of support:

\begin{itemize}
\tightlist
\item
  a hand derivation;
\item
  executable code;
\item
  a proof assistant term;
\item
  an experimental protocol;
\item
  a finite search;
\item
  a numerical solver;
\item
  a statistical analysis;
\item
  a symbolic computation.
\end{itemize}

The important point is not that every method becomes code.

It is that the method remains inspectable enough for its support role.

Sandve et al.~give practical rules for reproducible computational
research, including recording how results were produced and preserving
exact intermediate and external information where needed (Sandve et al.
2013).

The Atlas extends that documentary instinct across computational and
mathematical work.

\section{Environment identity}\label{environment-identity}

A program is not executed in a vacuum.

Results may depend on:

\begin{itemize}
\tightlist
\item
  language version;
\item
  library versions;
\item
  compiler;
\item
  architecture;
\item
  precision;
\item
  accelerator runtime;
\item
  operating system;
\item
  deterministic or nondeterministic kernels;
\item
  external services;
\item
  data state.
\end{itemize}

The necessary level of pinning depends on the claim.

A pure exact-integer computation may need little.

A floating-point distributed training run may need much more.

The goal is not maximal metadata for its own sake.

It is sufficient identity for the evidence path being claimed.

\section{The tiny replay package}\label{the-tiny-replay-package}

This chapter contains an actual executable witness.

The program

`mathematics/computational-witnesses/replay\_examples/deterministic\_fraction\_sum.py`

computes

\[
\frac13+\frac16+\frac12
\]

using exact rational arithmetic.

The expected output is

\[
\boxed{\texttt{1/1}}.
\]

The program SHA-256 is

`e6e841bfe975f283ab948c4f7f9bbbf5ba488d267fda36edc1b270d5dcd062c1`.

The expected-output SHA-256 is

`3117b181de4d46b7ff8adb4c78adec272c019160d4adf27a901e90ec114e1845`.

The identity manifest records:

\begin{itemize}
\tightlist
\item
  code path;
\item
  code hash;
\item
  expected-output path;
\item
  output hash;
\item
  invocation;
\item
  environment class;
\item
  arithmetic model;
\item
  claim boundary.
\end{itemize}

Atlas CI executes the program and compares the output byte-for-byte.

The mathematics is deliberately trivial.

The evidence mechanism is the object under study.

\section{What successful replay
establishes}\label{what-successful-replay-establishes}

For the tiny witness, successful replay establishes:

\begin{enumerate}
\def\labelenumi{\arabic{enumi}.}
\tightlist
\item
  the committed program has the declared hash;
\item
  the expected output has the declared hash;
\item
  the program executes successfully in the CI environment;
\item
  its standard output matches the locked expected bytes.
\end{enumerate}

That is a real result.

It does \textbf{not} establish:

\begin{itemize}
\tightlist
\item
  that every Atlas claim is reproducible;
\item
  that the code implements some unstated theorem;
\item
  that the program is independently reviewed;
\item
  that the result is formally verified;
\item
  that the result is certified.
\end{itemize}

A replay result should say exactly what passed.

\section{A perfectly replayable wrong
program}\label{a-perfectly-replayable-wrong-program}

Imagine the program:

\[
\texttt{print("2+2=5")}.
\]

Its source can be hashed.

Its environment can be pinned.

Its output can replay perfectly forever.

It remains mathematically wrong.

This is the most important counterexample in the chapter.

Replay can establish faithful reconstruction of a computation.

It cannot turn a false computation into a true claim.

\section{A correct argument that is not replayable
enough}\label{a-correct-argument-that-is-not-replayable-enough}

Now take the opposite case.

A mathematician writes a correct argument on paper, but the crucial
source is cited only as:

\begin{quote}
``the recent lemma from that preprint.''
\end{quote}

The reasoning may be mathematically correct.

The documentary path may still be too weak for another actor to
reconstruct efficiently.

Non-replayability does not imply falsehood.

It implies a provenance deficit.

This distinction prevents the Atlas from confusing documentation quality
with truth.

\section{Mutable dependencies}\label{mutable-dependencies}

Suppose a replay instruction says:

\begin{quote}
download `model/latest` and run the script.
\end{quote}

Today, `model/latest` resolves to bytes \(A\).

Tomorrow, it resolves to bytes \(B\).

The name did not change.

The object did.

A replay instruction tied only to the mutable name has lost identity.

Pinning an immutable revision repairs the identity problem.

It still does not prove that the pinned artifact is the \textbf{right}
artifact for the claim.

Identity and semantics remain separate.

\section{Seed is not environment}\label{seed-is-not-environment}

For a stochastic experiment, a random seed is valuable.

It is not always sufficient.

A seeded computation can still differ because of:

\begin{itemize}
\tightlist
\item
  PRNG implementation;
\item
  library changes;
\item
  accelerator nondeterminism;
\item
  thread scheduling;
\item
  reduction order;
\item
  mixed precision;
\item
  compiler transformations;
\item
  external data drift.
\end{itemize}

Therefore

\[
\boxed{
\text{seed}+\text{code}
\neq
\text{complete replay identity}
}
\]

in general.

The environment record should match the sensitivity of the evidence.

\section{Byte identity is not semantic
identity}\label{byte-identity-is-not-semantic-identity}

Consider a program that checks a property for every integer

\[
1\le n\le100.
\]

Suppose:

\begin{itemize}
\tightlist
\item
  the code hash matches;
\item
  the environment matches;
\item
  the output hash matches;
\item
  a second actor replays it successfully.
\end{itemize}

Now suppose the manuscript says:

\begin{quote}
the property holds for every positive integer.
\end{quote}

The replay is exact.

The claim is unsupported.

The failed object is the semantic bridge between observation \(O\) and
interpretation \(I\).

That is why an evidence object must bind both.

\section{Formal proof of the wrong
statement}\label{formal-proof-of-the-wrong-statement}

Formal verification is a powerful support route.

Suppose a proof assistant verifies a theorem

\[
T'.
\]

If the intended human claim was

\[
T,
\]

and the bridge

\[
T\leftrightarrow T'
\]

has not been established, then the formal proof does not automatically
settle \(T\).

The kernel can be perfectly correct about the formal object while the
human interpretation is wrong.

Formal verification therefore strengthens one part of the evidence path.

It does not abolish semantic review.

\section{Finite verification generalized too
far}\label{finite-verification-generalized-too-far}

A program may exhaustively check all cases in a finite set.

That can be an exact finite result.

If prose silently extends it to an infinite domain, the result has
changed class.

For example,

\[
P(n)\text{ verified for }1\le n\le10^6
\]

does not entail

\[
\forall n\in\mathbb N,\;P(n).
\]

Exact finite verification is strong evidence for the finite statement.

It is not a continuum or universal proof unless an additional theorem
supplies the bridge.

\section{Internal replay is not independent
reproduction}\label{internal-replay-is-not-independent-reproduction}

Suppose the same author:

\begin{itemize}
\tightlist
\item
  writes the code;
\item
  runs the code;
\item
  reruns the code;
\item
  confirms the same output.
\end{itemize}

That is useful regression evidence.

It is not independent reproduction.

Independence changes epistemic value only when the relevant independence
is real and documented.

Questions include:

\begin{itemize}
\tightlist
\item
  Did the second actor reuse the same implementation?
\item
  Did they use the same hidden assumptions?
\item
  Did they reconstruct the mathematics independently?
\item
  Did they share the same data-generation error?
\item
  Was the evaluation oracle independently derived?
\end{itemize}

``Independent'' is itself a claim that needs provenance.

\section{Provenance as a graph}\label{provenance-as-a-graph}

W3C PROV provides a general model of entities, activities, agents, and
derivation relations ({Moreau et al.} 2013).

The Atlas specializes the idea into a scientific support path:

\[
\text{source}
\to
\text{method}
\to
\text{execution}
\to
\text{observation}
\to
\text{interpretation}
\to
\text{claim}.
\]

Review and authority transitions are separate.

\begin{figure}
\centering
\pandocbounded{\includegraphics[keepaspectratio,alt={An exact declared graph with a solid evidence path from Source through Method, Execution, Observation, Interpretation, and Claim, followed by dashed separate transitions through Review or replay, Adjudication, and Certification.}]{figures/masters/ATLAS-FIG-REPLAY-001.png}}
\caption{An exact declared graph with a solid evidence path from Source
through Method, Execution, Observation, Interpretation, and Claim,
followed by dashed separate transitions through Review or replay,
Adjudication, and Certification.}
\end{figure}

Every node and directed edge in the figure is declared in the figure
manifest.

The layout is presentation.

The separation of edge classes is content.

\section{Evidence edges versus authority
edges}\label{evidence-edges-versus-authority-edges}

This distinction is worth making explicit.

An evidence edge says something like:

\begin{quote}
this observation was produced by this execution.
\end{quote}

An interpretation edge says:

\begin{quote}
this observation supports this bounded claim under this reasoning.
\end{quote}

A review edge says:

\begin{quote}
this actor checked this aspect of the evidence path.
\end{quote}

A certification edge says:

\begin{quote}
under this institution's rules, this exact revision reached this
authority state.
\end{quote}

These are different relations.

Collapsing them into one arrow marked ``verified'' destroys information.

\section{Current GCL technical-writing
doctrine}\label{current-gcl-technical-writing-doctrine}

The current GCL technical-writing reference states a programme
principle:

\begin{quote}
technical communication is part of the research instrument.
\end{quote}

Its writing invariants include:

\begin{itemize}
\tightlist
\item
  object identity;
\item
  claim identity;
\item
  bounded evidence;
\item
  visible limitations;
\item
  provenance;
\item
  review specificity;
\item
  authority separation;
\item
  public inheritance;
\item
  analogy discipline;
\item
  fail-closed promotion.
\end{itemize}

The exact source is pinned in this chapter's source lock to

`grandchallenge/MATH-PROGRAMME@9c09521f0f7b1b7bbb2830b42f097f227f209d7d`.

The Atlas uses that doctrine as a concrete institutional example.

It does not claim that every scientific organization must use GCL's
exact workflow.

\section{Forge, Solve, Cert as role
separation}\label{forge-solve-cert-as-role-separation}

GCL uses a Forge → Solve → Cert separation.

At a high level:

\begin{itemize}
\tightlist
\item
  \textbf{Forge} fixes source identity and problem/evidence context;
\item
  \textbf{Solve} performs bounded mathematical or computational work and
  adjudication;
\item
  \textbf{Cert} performs separate certification work under its own
  authority.
\end{itemize}

The value is not the names.

The value is that evidence production, problem solving, and
certification are not silently treated as one operation.

This reduces one common failure mode:

\begin{quote}
the actor that produced a result implicitly promotes it to whatever
authority level is convenient.
\end{quote}

Role separation makes that promotion explicit.

\section{The current OPENMATH
lifecycle}\label{the-current-openmath-lifecycle}

The source-locked current OPENMATH controller records:

\[
\texttt{READY}
\to
\texttt{LAUNCHED}
\to
\texttt{RETURNED}
\to
\texttt{CAPTURED}
\to
\texttt{REPLAYED}
\to
\texttt{ADJUDICATED}
\to
\texttt{ADVANCED}.
\]

Its durable substrate is GitHub task and evidence state.

Its return grammar is

\[
\texttt{GCL-CONTRIBUTION-RESULT/1}.
\]

Its authority boundary is recorded as

\[
\texttt{MATHFORGE\_TO\_MATHSOLVE\_TO\_MATHCERT}.
\]

Most importantly for this chapter, the controller explicitly prohibits
itself from:

\begin{itemize}
\tightlist
\item
  claiming mathematical correctness not established by Solve
  adjudication;
\item
  performing MATHCERT certification.
\end{itemize}

That is documentary evidence of authority separation, not merely a
slogan.

\section{Replay evidence can exist while certification remains
false}\label{replay-evidence-can-exist-while-certification-remains-false}

A current source-locked MATHCERT OPENMATH intake contains:

\begin{itemize}
\tightlist
\item
  source authority;
\item
  source lock;
\item
  statement identity;
\item
  evaluator contract;
\item
  proof receipt;
\item
  search receipt;
\item
  replay receipt;
\item
  candidate byte identity;
\item
  independent-executor requirements.
\end{itemize}

It also records:

\[
\boxed{
\texttt{certification\_effect:false}.
}
\]

This is a particularly useful example because it blocks an easy
conceptual collapse.

The presence of replay evidence does not itself create certification.

A certification transition is another event with another authority
basis.

\section{Fail closed}\label{fail-closed}

Suppose a replay requires:

\begin{itemize}
\tightlist
\item
  exact source identity;
\item
  exact candidate identity;
\item
  exact method identity;
\item
  exact expected output.
\end{itemize}

If any of those identities disagree, the correct response is not to
``approximately continue'' while preserving the same status.

The support path has changed.

GCL's current writing and OPENMATH doctrine uses a fail-closed pattern:
identity or provenance mismatches block the applicable promotion.

That does not mean every scientific workflow must stop on every
mismatch.

It means a workflow should not silently preserve an authority state
whose prerequisites no longer hold.

\section{Evidence object versus
archive}\label{evidence-object-versus-archive}

An archive stores artifacts.

An evidence object stores relationships among artifacts and claims.

The difference is important.

A directory containing:

\begin{itemize}
\tightlist
\item
  `result.csv`;
\item
  `final.py`;
\item
  `notes.md`;
\item
  `plot.png`
\end{itemize}

may be an archive.

It becomes stronger evidence when we can answer:

\begin{itemize}
\tightlist
\item
  Which code produced the CSV?
\item
  Which source revision did the code use?
\item
  Which rows support which claim?
\item
  Which figure came from which data?
\item
  What interpretation was applied?
\item
  What review checked that interpretation?
\item
  What version was actually certified?
\end{itemize}

Evidence is structured provenance plus bounded meaning.

\section{The evidence object as a reconstruction
test}\label{the-evidence-object-as-a-reconstruction-test}

The chapter's acceptance test is simple:

\begin{quote}
Can another competent actor reconstruct why we believe this, and can
they see exactly what that reconstruction does not establish?
\end{quote}

This is stronger than asking:

\begin{quote}
Are all the files present?
\end{quote}

It asks for an intelligible support path.

For a computational result, another actor should be able to reconstruct
the computation.

For a derivation, another actor should be able to reconstruct the
mathematical steps.

For a literature-dependent claim, another actor should be able to
recover the exact source.

For an institutional status, another actor should be able to identify
the exact authority record.

\section{Five distinctions to keep
visible}\label{five-distinctions-to-keep-visible-1}

\subsection{Replay versus truth}\label{replay-versus-truth}

Replay means the evidence path can be reconstructed under its stated
conditions.

It does not make a false method true.

\subsection{Replay versus
replication}\label{replay-versus-replication}

Replay may reuse the same source, data, and method.

Replication asks a different question and may use new data or an
independently rebuilt method (National Academies of Sciences,
Engineering, and Medicine 2019).

\subsection{Replay versus formal
verification}\label{replay-versus-formal-verification}

Formal verification checks a formal statement under a formal trust base.

Replay reconstructs an evidence path.

Either may support the other.

They are not synonyms.

\subsection{Replay versus
certification}\label{replay-versus-certification}

Certification is a governed disposition under some authority model.

Replay can be an input to certification.

It is not certification by implication.

\subsection{Certification versus
truth}\label{certification-versus-truth}

Institutional certification is a status granted under defined rules.

It does not redefine mathematical truth.

The Atlas must keep epistemic and institutional categories separate even
when they interact.

\section{What a durable evidence object should make
inspectable}\label{what-a-durable-evidence-object-should-make-inspectable}

A mature evidence object should make it possible to inspect:

\begin{itemize}
\tightlist
\item
  exact claim;
\item
  source identities;
\item
  source revisions;
\item
  code or derivation identity;
\item
  environment;
\item
  input identities;
\item
  output identities;
\item
  figure/data lineage;
\item
  assumptions;
\item
  limitations;
\item
  negative evidence;
\item
  replay status;
\item
  review scope;
\item
  independence claims;
\item
  adjudication state;
\item
  certification state;
\item
  supersession history.
\end{itemize}

Not every result needs every field.

But a missing field should be a conscious absence, not an accidental
one.

\section{Atlas connections}\label{atlas-connections-24}

\textbf{Experiments as Arguments.}\\
A controlled experiment becomes stronger when its exact support path can
be replayed.

\textbf{Formal Methods.}\\
Formal proofs can enter the evidence graph as machine-checkable support
objects without collapsing the semantic bridge.

\textbf{Research as a State Machine.}\\
Replay, adjudication, and certification become explicit state
transitions.

\textbf{Governed Adaptation.}\\
A system that changes itself needs provenance strong enough to know
which evidence applies to which revision.

\textbf{Memory.}\\
Institutional memory becomes more useful when it stores not only
conclusions but also support paths.

The recurring Atlas shift is:

\[
\boxed{
\text{research result}
\longrightarrow
\text{replayable evidence object}.
}
\]

\section{Closing view}\label{closing-view-35}

A claim is easier to trust when the path behind it remains visible.

But visibility is not truth.

A replayable wrong program is still wrong.

A correct theorem with poor provenance can still be correct.

A formally verified narrow statement can still be misinterpreted.

A successful replay can still await independent review.

A result can be well supported and still uncertified.

The chain-of-custody allegory is only an allegory.

The evidence object is the object.

And the final test remains:

\begin{quote}
Can another competent actor reconstruct why we believe this, and can
they see exactly what that reconstruction does not establish?
\end{quote}

\section{References used in this
chapter}\label{references-used-in-this-chapter-74}

\begin{itemize}
\tightlist
\item
  (National Academies of Sciences, Engineering, and Medicine 2019)
\item
  (Sandve et al. 2013)
\item
  ({Moreau et al.} 2013)
\end{itemize}

See `sources/source-locks/ATLAS-CH-REPLAY-001.yaml` for exact source
identities, current GCL lifecycle records, and claim scope.

\chapter{Formal Methods and Machine-Checkable
Claims}\label{formal-methods-and-machine-checkable-claims}

\textbf{Epistemic status:} established formal-methods foundations +
audited replay/provenance prerequisite + Atlas synthesis + exact
state-machine witness.\\
\textbf{Specification:}
manuscript/specifications/ATLAS-CH-FORMAL-001.md\\
\textbf{Formal packet:}
mathematics/derivations/ATLAS-CH-FORMAL-001-DERIVATIONS.md\\
\textbf{Computational witness:}
mathematics/computational-witnesses/ATLAS-CW-FORMAL-001.md\\
\textbf{Source lock:} sources/source-locks/ATLAS-CH-FORMAL-001.yaml

Formal methods are powerful because they force us to say exactly what is
being claimed.

They are also easy to overstate because a machine can check the wrong
formal statement perfectly.

The central problem is therefore not only:

\begin{quote}
can this theorem be checked?
\end{quote}

It is also:

\begin{quote}
what object was formalized, what semantics does it use, what does the
proof actually range over, what trusted machinery checks it, and how is
the formal result connected back to the intended system?
\end{quote}

The governing distinction of this chapter is:

\begin{quote}
machine-checked correctness is always correctness of a declared formal
object under a declared trust base.
\end{quote}

The semantic bridges around that object remain part of the scientific
argument.

\section{The formal support graph}\label{the-formal-support-graph}

Use the Atlas object set

\texttt{F=(R,S,M,P,K,I,W)}.

Here:

\begin{itemize}
\tightlist
\item
  \texttt{R}: intended requirement;
\item
  \texttt{S}: formal specification/property;
\item
  \texttt{M}: formal model and semantics;
\item
  \texttt{P}: proof object, derivation, certificate, or model-checking
  result;
\item
  \texttt{K}: trusted checker/kernel and surrounding trust base;
\item
  \texttt{I}: implementation;
\item
  \texttt{W}: deployed world or environment.
\end{itemize}

These objects should not be arranged as one untyped implication chain.

Instead, use typed obligations:

\texttt{Formalizes(S,R)}

asks whether the formal statement captures the intended requirement.

\texttt{Interprets(M,S)}

binds the statement to its formal semantics/model.

\texttt{Checks(K,P,S,M)}

records that checker \texttt{K} accepts support object \texttt{P} for
the declared formal claim.

\texttt{Conforms(I,M)}

asks whether the implementation refines or conforms to the formal model
strongly enough for the property being transferred.

\texttt{AssumptionsHold(W,M)}

asks whether the deployed world satisfies the environmental assumptions
used by the model.

This graph prevents a common collapse.

A proof checker may establish \texttt{Checks}.

It does not automatically establish \texttt{Formalizes},
\texttt{Conforms}, or \texttt{AssumptionsHold}.

\section{Formal specification is not human
intent}\label{formal-specification-is-not-human-intent}

Suppose the intended requirement is:

\begin{quote}
the controller must never issue an unsafe command.
\end{quote}

A formal specification might encode:

\texttt{forall\ reachable\ s,\ command(s)\ in\ SafeSet}.

That may be precise.

But perhaps the formalization omitted:

\begin{itemize}
\tightlist
\item
  a sensor failure mode;
\item
  an actuator delay;
\item
  an environmental transition;
\item
  a unit conversion;
\item
  an emergency override.
\end{itemize}

The theorem can be completely correct relative to its formal world and
still fail to establish the real intended requirement.

This is not a defect peculiar to theorem provers.

It is the ordinary scientific problem of model adequacy expressed
sharply.

Formalization moves ambiguity into explicit definitions.

It does not make interpretation disappear.

\section{Hoare logic: contracts as proof
objects}\label{hoare-logic-contracts-as-proof-objects}

Hoare's axiomatic account of programming introduced reasoning in the
form

\texttt{\{P\}\ C\ \{Q\}}

for command \texttt{C}, precondition \texttt{P}, and postcondition
\texttt{Q} (Hoare 1969).

Under partial-correctness semantics, the triple means:

if \texttt{P} holds before execution and \texttt{C} terminates, then
\texttt{Q} holds afterward.

For

\texttt{x\ :=\ x+2},

a useful triple is

\texttt{\{Even(x)\}\ x:=x+2\ \{Even(x)\}}.

If

\texttt{x=2k},

then after the assignment

\texttt{x\textquotesingle{}=2k+2=2(k+1)}.

The postcondition follows.

This small example contains the essential formal-methods pattern:

\begin{itemize}
\tightlist
\item
  state a contract;
\item
  define semantics;
\item
  prove the transition preserves the contract.
\end{itemize}

\section{Invariants turn local preservation into global
claims}\label{invariants-turn-local-preservation-into-global-claims}

Consider a transition system with:

\begin{itemize}
\tightlist
\item
  initial predicate \texttt{Init(x)};
\item
  next-state relation \texttt{Next(x,x\textquotesingle{})};
\item
  invariant candidate \texttt{Inv(x)}.
\end{itemize}

To prove that every reachable state satisfies \texttt{Inv}, establish:

\begin{enumerate}
\def\labelenumi{\arabic{enumi}.}
\tightlist
\item
  initialization: \texttt{Init(x)\ =\textgreater{}\ Inv(x)};
\item
  preservation:
  \texttt{Inv(x)\ and\ Next(x,x\textquotesingle{})\ =\textgreater{}\ Inv(x\textquotesingle{})}.
\end{enumerate}

Then induction over the reachability derivation gives the global
invariant.

This is different from observing many successful executions.

The proof quantifies over every reachable state admitted by the formal
transition system.

\section{Exact witness: an infinite-state
counter}\label{exact-witness-an-infinite-state-counter}

Let

\texttt{x\_0=0}

and

\texttt{x\_\{n+1\}=x\_n+2}.

The claimed invariant is

\texttt{Even(x\_n)}.

The base case is immediate:

\texttt{0=2*0}.

For the inductive step, if

\texttt{x\_n=2k},

then

\texttt{x\_\{n+1\}=2k+2=2(k+1)}.

Therefore

\texttt{x\_n=2n}

for every natural number \texttt{n}.

Every reachable specified state is even.

This is a universal statement over an infinite family of
executions/states.

\section{Three tests are not that
theorem}\label{three-tests-are-not-that-theorem}

Now test only:

\texttt{0\ -\textgreater{}\ 2}, \texttt{2\ -\textgreater{}\ 4},
\texttt{4\ -\textgreater{}\ 6}.

All pass.

These tests establish exactly those executions.

They do not, by themselves, establish the universal invariant.

The difference is logical, not rhetorical.

A finite list of successful examples has the shape:

\texttt{P(x\_1),...,P(x\_m)}.

The invariant theorem has the shape:

\texttt{forall\ x\ in\ Reachable,\ P(x)}.

The quantifier changed.

An additional theorem is needed to bridge them.

\section{A passing implementation can still be
wrong}\label{a-passing-implementation-can-still-be-wrong}

Define an implementation:

\texttt{ImplStep(x)=x+2} when \texttt{x\textless{}6};

\texttt{ImplStep(x)=x+1} when \texttt{x\textgreater{}=6}.

Starting from zero:

\texttt{0\ -\textgreater{}\ 2\ -\textgreater{}\ 4\ -\textgreater{}\ 6}.

The three tests pass.

The next transition is:

\texttt{6\ -\textgreater{}\ 7}.

The invariant fails.

At exactly the same state, the specification says:

\texttt{6\ -\textgreater{}\ 8}.

So three facts coexist:

\begin{enumerate}
\def\labelenumi{\arabic{enumi}.}
\tightlist
\item
  the formal invariant proof is correct for the specification;
\item
  the tested implementation prefix passes;
\item
  the implementation does not conform to the specification.
\end{enumerate}

This is the chapter's central exact counterexample.

\section{Proof of specification is not proof of
implementation}\label{proof-of-specification-is-not-proof-of-implementation}

Suppose we prove:

\texttt{Even(x)\ -\textgreater{}\ Even(SpecStep(x))}.

That theorem says nothing directly about \texttt{ImplStep}.

To transfer the property, we need a bridge such as:

\texttt{ImplStep(x)=SpecStep(x)}

or a refinement relation strong enough to preserve the property.

Possible routes include:

\begin{itemize}
\tightlist
\item
  verified compilation;
\item
  refinement proofs;
\item
  equivalence proofs;
\item
  code extraction;
\item
  executable contracts;
\item
  verified interpreters;
\item
  constrained conformance testing.
\end{itemize}

Each route adds assumptions.

Formal methods do not remove bridges.

They make them inspectable.

\section{Tests and proofs are
complementary}\label{tests-and-proofs-are-complementary}

Testing is not a weak version of proof.

It is a different support route.

Tests are excellent for:

\begin{itemize}
\tightlist
\item
  finding concrete counterexamples;
\item
  exercising real implementations;
\item
  regression checking;
\item
  exposing integration failures;
\item
  validating environment assumptions.
\end{itemize}

Proof is excellent for:

\begin{itemize}
\tightlist
\item
  universal claims over a formal domain;
\item
  invariant preservation;
\item
  exact algebraic properties;
\item
  eliminating untested cases inside the formal model.
\end{itemize}

One failing test can refute a universal claim immediately.

Many successful tests usually do not establish the universal claim.

The asymmetry matters.

\section{Exhaustive finite checking is stronger than
sampling}\label{exhaustive-finite-checking-is-stronger-than-sampling}

Suppose a state space is genuinely finite.

If every reachable state is enumerated and every transition is checked,
then a finite computation can establish a universal property over that
finite model.

That is not the same as sampling.

The support route depends on:

\begin{itemize}
\tightlist
\item
  the exact finite model;
\item
  transition semantics;
\item
  completeness of exploration;
\item
  property checked;
\item
  correctness of the checker.
\end{itemize}

The word \textbf{finite} does not weaken the result.

It defines the result's domain.

\section{Model checking}\label{model-checking}

State-machine formal methods often write a system as:

\begin{itemize}
\tightlist
\item
  initial-state predicate;
\item
  transition relation;
\item
  safety properties;
\item
  liveness/fairness assumptions where needed.
\end{itemize}

TLA+ is one mature framework for this style, and Lamport's
\emph{Specifying Systems} develops specification, safety/liveness
reasoning, composition, and TLC model checking (Lamport 2002).

A model checker can explore all reachable states of a finite model.

That can be an exact result about the model.

It is still not automatically a proof about arbitrary deployed code.

The implementation/model bridge remains separate.

\section{The model can be the wrong
model}\label{the-model-can-be-the-wrong-model}

Imagine a distributed protocol model with perfect channels.

The safety theorem may hold.

The deployed system may lose, duplicate, or reorder messages in ways
omitted from the formal model.

The theorem is not thereby false.

Its assumptions are simply not the deployment assumptions.

This distinction is essential:

\begin{quote}
a correct theorem with an inadequate model is still a correct theorem
about the inadequate model.
\end{quote}

Formal verification therefore strengthens the need for explicit model
identity.

\section{Trusted kernels}\label{trusted-kernels}

Proof assistants introduce another architectural distinction.

Robin Milner's LCF work established the influential pattern of a small
trusted proof core surrounded by more complex proof-producing machinery
(Milner 1979).

The intuition is simple.

Tactics and automation may be sophisticated and buggy.

If they can only produce theorem objects through a small trusted
interface, then bugs in automation should fail to manufacture accepted
false theorems unless they compromise the trusted core or exploit
assumptions already present in the logic.

This reduces the trusted computing base.

It does not eliminate trust.

\section{Lean as a proof-assistant
example}\label{lean-as-a-proof-assistant-example}

Lean is an interactive theorem prover based on dependent type theory
with a small trusted kernel (Moura et al. 2015).

Its architecture separates convenient proof construction from final
checking.

A user may rely on:

\begin{itemize}
\tightlist
\item
  notation;
\item
  elaboration;
\item
  type-class inference;
\item
  tactics;
\item
  automation;
\item
  libraries.
\end{itemize}

The final elaborated object is checked by the kernel against the formal
environment.

This is the key trust pattern:

\begin{quote}
automation may be large while the final logical checker remains small.
\end{quote}

\section{Small trusted core does not mean no
trust}\label{small-trusted-core-does-not-mean-no-trust}

Even a small kernel has assumptions.

A formal result can depend on:

\begin{itemize}
\tightlist
\item
  the type theory or logic;
\item
  explicit axioms;
\item
  the kernel implementation;
\item
  imported definitions;
\item
  library theorems;
\item
  parser/elaborator behavior if not independently reconstructed;
\item
  compiler/runtime/hardware when executing extracted code;
\item
  external certificates or oracles.
\end{itemize}

So ``Lean checked it'' is meaningful but incomplete.

A mature statement says:

\begin{itemize}
\tightlist
\item
  which theorem;
\item
  under which imports/axioms;
\item
  in which toolchain;
\item
  checked by which kernel;
\item
  with what semantic bridge to the intended claim.
\end{itemize}

\section{Proof terms versus theorem
meaning}\label{proof-terms-versus-theorem-meaning}

In a propositions-as-types setting, a proof term inhabits the
proposition representing the theorem.

Kernel acceptance therefore establishes a precise syntactic-semantic
fact inside the formal environment.

But theorem meaning still depends on the chosen definitions.

If a theorem formalizes:

\texttt{Safe\ :=\ x\textless{}100}

while the intended policy meant:

\texttt{Safe\ :=\ x\textless{}=10},

the kernel may correctly prove the wrong policy.

Machine checking is strongest when formal statement identity is
preserved alongside the proof.

This directly inherits the claim-identity discipline from REPLAY-001.

\section{Replayable proof is stronger than an unpinned proof
artifact}\label{replayable-proof-is-stronger-than-an-unpinned-proof-artifact}

REPLAY-001 established that source, method, environment, output, and
interpretation should be identity-bound.

Formal proof benefits from exactly the same discipline.

A proof artifact can depend on:

\begin{itemize}
\tightlist
\item
  theorem prover version;
\item
  library revision;
\item
  imported modules;
\item
  generated code;
\item
  solver version;
\item
  external certificate format.
\end{itemize}

If those dependencies float, another actor may not be able to
reconstruct the check.

Thus:

\texttt{formal\ proof}

and

\texttt{replayable\ formal\ proof}

are distinct documentary states.

\section{Replay is not proof}\label{replay-is-not-proof}

The converse is equally important.

A program can reproduce the same wrong output forever.

A checker can replay a finite search exactly.

A build can remain byte-for-byte stable.

None of those facts automatically imply a theorem.

REPLAY-001 already established this boundary.

FORMAL-001 adds the complementary rule:

\begin{quote}
proof and replay are orthogonal support dimensions.
\end{quote}

A result can be:

\begin{itemize}
\tightlist
\item
  replayable but unproved;
\item
  proved but poorly replayable;
\item
  both;
\item
  neither.
\end{itemize}

\section{Executable checkers and proof
objects}\label{executable-checkers-and-proof-objects}

Not every machine-checkable result uses a proof assistant.

Some systems use compact certificates.

A producer may perform expensive search.

A smaller checker verifies the returned certificate.

Examples can include:

\begin{itemize}
\tightlist
\item
  SAT/SMT proof traces;
\item
  primality certificates;
\item
  optimization certificates;
\item
  typed intermediate representations;
\item
  proof-carrying code.
\end{itemize}

The epistemic advantage is architectural:

the producer can be less trusted if the checker fully validates the
support object.

But the checker still validates only its declared relation.

\section{Static types are formal claims with bounded
scope}\label{static-types-are-formal-claims-with-bounded-scope}

A type checker proves something.

What it proves depends on the type system.

A well-typed program may guarantee:

\begin{itemize}
\tightlist
\item
  certain operations receive values of the right form;
\item
  some memory or effect discipline;
\item
  absence of certain representation errors.
\end{itemize}

It does not automatically prove:

\begin{itemize}
\tightlist
\item
  business correctness;
\item
  security policy;
\item
  termination;
\item
  liveness;
\item
  numerical accuracy.
\end{itemize}

A formal mechanism should always be paired with its actual claim class.

\section{Runtime contracts are path-local
evidence}\label{runtime-contracts-are-path-local-evidence}

A runtime assertion such as

\texttt{assert\ even(x)}

can detect a violation on an executed path.

If the assertion never fires across many tests, that is useful evidence.

But unless all possible executions are covered or an additional theorem
gives completeness, the assertion history remains execution-local.

This is another example of the quantifier discipline:

observed paths are not automatically all paths.

\section{Property-based testing is still
testing}\label{property-based-testing-is-still-testing}

Property-based testing can generate large numbers of structured inputs.

It often finds counterexamples that humans would not write manually.

This is valuable.

But if generation samples an infinite or enormous domain, passing
campaigns remain sampling evidence.

If the generator is exhaustive over a finite domain, the claim class
changes.

The difference should be stated, not implied.

\section{SMT and SAT trust paths}\label{smt-and-sat-trust-paths}

Automated solvers are common inside formal methods.

There are several trust architectures:

\begin{itemize}
\tightlist
\item
  trust the solver directly;
\item
  ask the solver to emit a proof trace;
\item
  independently check the trace;
\item
  reconstruct the result in a proof assistant;
\item
  let an untrusted tactic invoke the solver and produce a kernel-checked
  term.
\end{itemize}

The phrase

\begin{quote}
the solver proved it
\end{quote}

does not tell us which architecture was used.

The trust path belongs in the evidence record.

\section{Formal proof can expose missing
requirements}\label{formal-proof-can-expose-missing-requirements}

One practical benefit of formalization is that it forces vague language
into explicit predicates and transitions.

That process often reveals that the requirement itself was incomplete.

Questions emerge:

\begin{itemize}
\tightlist
\item
  What exactly counts as failure?
\item
  Which state variables matter?
\item
  Is liveness required or only safety?
\item
  Is termination required?
\item
  Which environment actions are possible?
\item
  What fairness assumption is being made?
\end{itemize}

Formal methods therefore have value even before a theorem is completed.

They refine the question.

\section{Invariants are contracts over
evolution}\label{invariants-are-contracts-over-evolution}

In adaptive systems, invariants are especially useful because the system
changes over time.

Examples include:

\begin{itemize}
\tightlist
\item
  a resource bound is never exceeded;
\item
  authority never escalates without a permitted transition;
\item
  provenance is never discarded during promotion;
\item
  a representation remains normalized;
\item
  a safety region remains invariant;
\item
  an evidence state cannot jump directly from unreviewed to certified.
\end{itemize}

This prepares the next chapter.

Research processes themselves can be modeled as governed transition
systems.

\section{Safety and liveness are
different}\label{safety-and-liveness-are-different}

A safety property says, roughly:

\begin{quote}
something bad never happens.
\end{quote}

An invariant is a common safety-property form.

A liveness property says:

\begin{quote}
something good eventually happens.
\end{quote}

A system can be perfectly safe by doing nothing forever.

Therefore a governance or research state machine often needs both:

\begin{itemize}
\tightlist
\item
  forbidden transitions must never occur;
\item
  legitimate work must still be able to progress.
\end{itemize}

Formal verification should not hide this distinction under one word such
as \textbf{correct}.

\section{Formal certification is overloaded
language}\label{formal-certification-is-overloaded-language}

The word \textbf{certified} can mean several things:

\begin{itemize}
\tightlist
\item
  a theorem was machine checked;
\item
  a proof certificate was checked;
\item
  an implementation property was verified;
\item
  a compliance body issued a certificate;
\item
  an institution promoted an artifact into an authoritative state.
\end{itemize}

These are different events.

The Atlas reserves the stronger authority language for the actual
governed transition.

A proof assistant accepting a theorem is not automatically institutional
certification.

\section{Machine-checked does not mean
assumption-free}\label{machine-checked-does-not-mean-assumption-free}

A theorem can use axioms.

A theorem can rely on imported propositions.

A model can assume fairness.

A verified controller can assume bounded disturbance.

A proof can be exact while its premises are conditional.

The appropriate statement is therefore:

\begin{quote}
theorem T is machine checked under assumptions A in formal environment
E.
\end{quote}

Conditional correctness is still correctness.

Its assumptions must remain visible.

\section{The trust-base ledger}\label{the-trust-base-ledger}

For a machine-checkable claim, record:

{\def\LTcaptype{none} % do not increment counter
\begin{longtable}[]{@{}
  >{\raggedright\arraybackslash}p{(\linewidth - 2\tabcolsep) * \real{0.5000}}
  >{\raggedright\arraybackslash}p{(\linewidth - 2\tabcolsep) * \real{0.5000}}@{}}
\toprule\noalign{}
\begin{minipage}[b]{\linewidth}\raggedright
Field
\end{minipage} & \begin{minipage}[b]{\linewidth}\raggedright
Question
\end{minipage} \\
\midrule\noalign{}
\endhead
\bottomrule\noalign{}
\endlastfoot
requirement & What human/system claim is intended? \\
formal statement & What exact proposition/property was encoded? \\
semantics & What model gives the symbols meaning? \\
proof/check & What derivation, trace, model-check run, or certificate
supports it? \\
checker & What validates the support object? \\
axioms/imports & What is assumed rather than derived? \\
toolchain & What exact version/environment is needed to replay
checking? \\
implementation bridge & Why does the deployed artifact satisfy the
formal model? \\
environment bridge & Why do real conditions satisfy formal
assumptions? \\
authority & What review/certification state was actually reached? \\
\end{longtable}
}

This ledger is the formal-methods analogue of REPLAY-001's evidence
object.

\section{Failure modes}\label{failure-modes-19}

\subsection{Proof-of-the-wrong-specification}\label{proof-of-the-wrong-specification}

A theorem is valid but the formal predicate omits the real requirement.

\subsection{Model-equals-system}\label{model-equals-system}

A finite/abstract model is treated as identical to deployment.

\subsection{Tests-equal-proof}\label{tests-equal-proof}

A passing finite suite is promoted into a universal claim.

\subsection{Automation-equals-trust}\label{automation-equals-trust}

A tactic or solver is trusted merely because it is sophisticated.

\subsection{Kernel-equals-no-assumptions}\label{kernel-equals-no-assumptions}

A small trusted kernel is treated as eliminating axioms and semantic
assumptions.

\subsection{Type-safe-equals-correct}\label{type-safe-equals-correct}

A bounded static property is promoted into total system correctness.

\subsection{Machine-checked-equals-certified}\label{machine-checked-equals-certified}

A proof check is treated as an institutional authority transition.

\subsection{Formal-equals-replayable}\label{formal-equals-replayable}

A theorem artifact is accepted without pinning toolchain, imports, or
source identity.

\section{Why the exact witness
matters}\label{why-the-exact-witness-matters}

The counter witness is deliberately simple because the distinction is
conceptual.

The invariant proof covers every specified reachable state.

The finite tests cover three implementation transitions.

The implementation then diverges.

Nothing subtle is happening.

The point is that all three evidence statements can be true
simultaneously.

That makes the boundary hard to evade:

\begin{quote}
proof, testing, and conformance are different relations.
\end{quote}

\section{Downstream handoff}\label{downstream-handoff-12}

\textbf{Research as a State Machine --- ATLAS-CH-RESEARCHSM-001} may now
assume:

\begin{itemize}
\tightlist
\item
  transition-system semantics;
\item
  inductive invariants;
\item
  initial/preservation proof obligations;
\item
  safety versus liveness distinction;
\item
  machine-checking versus semantic adequacy;
\item
  trusted-kernel/trust-base discipline;
\item
  replay and formal proof as complementary support routes;
\item
  authority transitions as separate from theorem checking.
\end{itemize}

The downstream chapter must independently formalize:

\begin{itemize}
\tightlist
\item
  Forge -\textgreater{} Solve -\textgreater{} Cert;
\item
  bounded work packages;
\item
  independent actors;
\item
  idempotence;
\item
  promotion gates;
\item
  failure/rollback states;
\item
  authority transitions.
\end{itemize}

\section{References used in this
chapter}\label{references-used-in-this-chapter-75}

\begin{itemize}
\tightlist
\item
  Hoare, \emph{An Axiomatic Basis for Computer Programming} (Hoare
  1969).
\item
  Milner, \emph{LCF: A Way of Doing Proofs with a Machine} (Milner
  1979).
\item
  de Moura et al., \emph{The Lean Theorem Prover (System Description)}
  (Moura et al. 2015).
\item
  Lamport, \emph{Specifying Systems: The TLA+ Language and Tools for
  Hardware and Software Engineers} (Lamport 2002).
\end{itemize}

Exact source identities and authority boundaries are recorded in:

sources/source-locks/ATLAS-CH-FORMAL-001.yaml

\chapter{Research as a State Machine}\label{research-as-a-state-machine}

\textbf{Epistemic status:} audited Replay/Formal prerequisites + public
GCL workflow case study + Atlas synthesis + exact finite witness.\\
\textbf{Specification:}
manuscript/specifications/ATLAS-CH-RESEARCHSM-001.md\\
\textbf{Formal packet:}
mathematics/derivations/ATLAS-CH-RESEARCHSM-001-DERIVATIONS.md\\
\textbf{Computational witness:}
mathematics/computational-witnesses/ATLAS-CW-RESEARCHSM-001.md\\
\textbf{Source lock:} sources/source-locks/ATLAS-CH-RESEARCHSM-001.yaml

A sentence such as ``the research task is done'' hides too much.

A result can arrive without being preserved. A preserved artifact can
remain unchecked. A successful replay can reconstruct a computation
without deciding the surrounding claim. A bounded review can support a
result without settling every later use of it.

This chapter replaces those ambiguities with explicit objects and
transitions.

\section{From evidence objects to research
state}\label{from-evidence-objects-to-research-state}

The Replayable Evidence chapter established that source identity,
execution, observation, interpretation, review, and certification are
different relations.

The Formal Methods chapter established that a transition system is
useful only when its state, transition relation, invariants, and model
boundary are explicit.

Research workflows need both ideas.

Evidence objects tell us what moves through the workflow.

A state machine tells us which changes are permitted, what each change
means, and which facts remain unchanged when a transition fails.

The object of this chapter is therefore not a task list.

It is a typed lifecycle for evidence-bearing research objects.

\section{A bounded work package}\label{a-bounded-work-package}

Let

\texttt{W=(id,Q,B,S,D,R,A,Z)}.

The coordinates record:

\begin{itemize}
\tightlist
\item
  stable work-package identity;
\item
  bounded question or obligation;
\item
  explicit bootstrap facts;
\item
  allowed source universe;
\item
  required deliverables;
\item
  durable return route;
\item
  authority granted to the worker;
\item
  stop conditions and explicit non-authorities.
\end{itemize}

The point of this tuple is not administrative completeness.

It is semantic closure.

A competent contributor should be able to decide, from the package
itself, what problem is being attempted, which premises are authorized,
what a valid return must contain, and what the return is not allowed to
decide.

\section{The research state is a
product}\label{the-research-state-is-a-product}

Represent the live research state as

\texttt{x=(q,E,U,J,P,C,L)}.

Here:

\begin{itemize}
\tightlist
\item
  \texttt{q} records execution or transport phase;
\item
  \texttt{E} contains preserved immutable evidence identities;
\item
  \texttt{U} records replay or checking state;
\item
  \texttt{J} records adjudication;
\item
  \texttt{P} records programme disposition;
\item
  \texttt{C} records certification state;
\item
  \texttt{L} is append-only history.
\end{itemize}

A representative transport path is:

\texttt{READY\ -\textgreater{}\ LAUNCHED\ -\textgreater{}\ RETURNED\ -\textgreater{}\ CAPTURED\ -\textgreater{}\ REPLAYED\ -\textgreater{}\ ADJUDICATED}.

That path is useful.

It is not the whole machine.

Programme disposition and certification are separate coordinates because
neither is merely another name for successful transport.

\section{State change requires a declared
condition}\label{state-change-requires-a-declared-condition}

For each event type, define a condition \texttt{g\_e(x)} and a partial
state update \texttt{delta\_e(x)}.

The update is applied only when the declared condition holds.

A return can arrive without satisfying its schema. A schema-valid return
can fail replay. A replayed result can fail adjudication. A positively
adjudicated result can still await a separate programme decision.

The state machine keeps these cases separate.

\section{Forge, Solve, and Cert are
roles}\label{forge-solve-and-cert-are-roles}

The public GCL workflow uses Forge, Solve, and Cert as a separation of
work.

At the level needed here:

\begin{itemize}
\tightlist
\item
  Forge fixes problem and source identity;
\item
  Solve performs bounded mathematical or computational work and
  adjudication;
\item
  Cert performs a separate certification function.
\end{itemize}

The names are project-specific. The structural lesson is that evidence
production, adjudication, and later certification need not be one
operation.

\section{Receipt is not
incorporation}\label{receipt-is-not-incorporation}

Suppose a durable return reaches its declared route.

That establishes receipt.

It does not automatically establish structural validity, replay success,
scientific correctness, programme incorporation, successor selection, or
certification.

The earlier evidence chapters stated this boundary conceptually.

The state-machine view makes it operational: each distinction occupies a
different coordinate or transition.

\section{Idempotence belongs to canonical
effect}\label{idempotence-belongs-to-canonical-effect}

A retryable research system needs two views of state.

The first is canonical state: which evidence objects and dispositions
currently count.

The second is history: what attempts, retries, failures, and recoveries
occurred.

Let \texttt{Canon(x)} omit retry-log multiplicity.

For a stable event key \texttt{k}, canonical idempotence means

\texttt{Canon(F\_k(F\_k(x)))=Canon(F\_k(x))}.

The history may still record both attempts.

This distinction matters.

If the entire history had to be identical after retry, the system would
have to hide the fact that a retry occurred.

If canonical state were allowed to multiply with every retry, the same
evidence could acquire multiple effects merely because transport
repeated.

Neither behavior is desirable.

\section{Duplicate is an identity
relation}\label{duplicate-is-an-identity-relation}

A return is not a duplicate merely because it belongs to the same work
package.

Use the exact return identity

\texttt{r=(d,h)}

for dispatch identity \texttt{d} and immutable result identity
\texttt{h}.

Two returns are exact duplicates only when both coordinates agree.

If the dispatch is the same but the result identity differs, the
programme has two evidence objects.

It may later decide that one supersedes, refutes, strengthens, or merely
differs from the other.

But that semantic decision must not be smuggled into transport by
treating every second return as ``the same thing.''

\section{Exact finite retry
witness}\label{exact-finite-retry-witness}

Take dispatch

\texttt{d=17}

and result identity

\texttt{r\_A=(17,A)}.

The first capture creates the singleton evidence set

\texttt{E\_1=\{r\_A\}}.

Receive the identical result again.

Set union gives

\texttt{E\_2=E\_1\ union\ \{r\_A\}=E\_1}.

Therefore canonical evidence count remains

\texttt{1}.

Now introduce

\texttt{r\_B=(17,B)}

with

\texttt{B!=A}.

Then

\texttt{E\_3=\{r\_A,r\_B\}}

has cardinality

\texttt{2}.

The witness therefore separates retry from genuinely distinct evidence
using only stable identity and finite set arithmetic.

\section{Advancement needs its own
guard}\label{advancement-needs-its-own-guard}

Receipt and replay are evidence-state transitions.

Programme advancement is a different act.

For one finite model, define five Boolean conditions:

\begin{itemize}
\tightlist
\item
  \texttt{I(e)}: exact identity and provenance are valid;
\item
  \texttt{S(e)}: the structural return contract is valid;
\item
  \texttt{R(e)}: the required replay or check has completed;
\item
  \texttt{A(e)}: adjudication supports the bounded claim;
\item
  \texttt{D(e)}: the programme has issued an explicit advancement
  disposition.
\end{itemize}

Then define:

\texttt{G\_adv(e)=I(e)\ S(e)\ R(e)\ A(e)\ D(e)}.

This is a deliberately strict finite witness.

It is not presented as the only possible research-governance formula.

\section{A missing programme disposition keeps the gate
closed}\label{a-missing-programme-disposition-keeps-the-gate-closed}

Evaluate:

\texttt{b=(1,1,1,1,0)}.

Then:

\texttt{G\_adv(b)=0}.

The evidence may be well identified.

Its structure may be valid.

Replay may have succeeded.

Adjudication may even support the bounded claim.

The successor is still not authorized under this witness because
programme disposition is absent.

This is exactly the distinction encoded by the public GCL Frontier
Advancement Gate.

\section{A complete witness vector opens the finite
gate}\label{a-complete-witness-vector-opens-the-finite-gate}

Now evaluate:

\texttt{b\textquotesingle{}=(1,1,1,1,1)}.

Then:

\texttt{G\_adv(b\textquotesingle{})=1}.

The arithmetic is trivial.

The semantic point is not.

A transition becomes legal because every declared condition is
satisfied, not because the underlying prose sounds persuasive.

\section{A missing required check also keeps the gate
closed}\label{a-missing-required-check-also-keeps-the-gate-closed}

Take:

\texttt{b\_R=(1,1,0,1,1)}.

Then:

\texttt{G\_adv(b\_R)=0}.

A strong result statement cannot substitute for a required check that
the governing process explicitly declared mandatory.

This is one value of machine-readable workflow state:

\begin{quote}
missing conditions remain visible as missing conditions.
\end{quote}

\section{Advancement can be canonically idempotent
too}\label{advancement-can-be-canonically-idempotent-too}

Let:

\texttt{K\_adv}

be the canonical set of evidence identities that have received the
relevant advancement effect.

For identity \texttt{m}:

\texttt{K\_adv\ union\ \{m\}\ union\ \{m\}=K\_adv\ union\ \{m\}}.

So repeated exact application of the same advancement effect need not
create a second canonical promotion.

The history may still record both attempts.

This matters whenever delivery, CI, merge automation, or network retries
can repeat an otherwise valid event.

\section{Idempotence is scoped, not
magical}\label{idempotence-is-scoped-not-magical}

The word \textbf{idempotent} is easy to overuse.

This chapter means something narrow:

\begin{quote}
a declared canonical mutation keyed by stable identity does not multiply
its canonical effect when the exact same event is replayed.
\end{quote}

It does not mean:

\begin{itemize}
\tightlist
\item
  two independent reviews must reach the same judgment;
\item
  rerunning a stochastic experiment must return identical measurements;
\item
  policy changes have no effect;
\item
  a new result identity should be discarded.
\end{itemize}

Idempotence is a property of a specified operation.

\section{Different evidence under one dispatch is not a
retry}\label{different-evidence-under-one-dispatch-is-not-a-retry}

Return:

\texttt{r\_A=(17,A)}

and later:

\texttt{r\_B=(17,B)}.

Because:

\texttt{A\ !=\ B},

these are different immutable result identities.

The second object may:

\begin{itemize}
\tightlist
\item
  conflict with the first;
\item
  strengthen it;
\item
  supersede it under an explicit policy;
\item
  address a different subclaim.
\end{itemize}

The transport layer should preserve that distinction instead of silently
replacing one object by dispatch key.

\section{Conflict is a scientific relation, not an intake
shortcut}\label{conflict-is-a-scientific-relation-not-an-intake-shortcut}

Suppose two preserved returns disagree.

The state machine should first preserve both identities.

Only then should adjudication decide the semantic relation.

This ordering matters:

\texttt{preserve\ -\textgreater{}\ compare\ -\textgreater{}\ adjudicate}

is different from:

\texttt{receive\ second\ -\textgreater{}\ overwrite\ first}.

The first route retains the evidence needed to understand disagreement.

\section{Actor roles are not actor
independence}\label{actor-roles-are-not-actor-independence}

Research workflows often name roles:

\begin{itemize}
\tightlist
\item
  contributor;
\item
  reviewer;
\item
  adjudicator;
\item
  certifier.
\end{itemize}

Those labels are useful.

They do not prove independence.

Two differently named actors can share:

\begin{itemize}
\tightlist
\item
  information;
\item
  incentives;
\item
  training data;
\item
  organizational control;
\item
  hidden communication.
\end{itemize}

The chapter therefore treats independence as a governed relation, not a
typography choice.

\section{A narrow separation-of-duty
predicate}\label{a-narrow-separation-of-duty-predicate}

Suppose one policy requires only that the producer and reviewer
identities differ.

Define:

\texttt{Sep(a\_prod,a\_rev)=1\{a\_prod\ !=\ a\_rev\}}.

Then:

\texttt{Sep(A,A)=0}

and:

\texttt{Sep(A,B)=1}.

This proves identity separation only.

It does not prove epistemic independence.

\section{Same-actor review can fail a declared
gate}\label{same-actor-review-can-fail-a-declared-gate}

If the governing policy requires:

\texttt{Sep=1},

then a review by the same actor who produced the artifact does not
satisfy that particular condition.

This is a policy statement, not a universal moral claim about
self-review.

Self-review can still be valuable.

It simply does not substitute for a distinct-actor condition when such a
condition has been explicitly required.

\section{Certification is another state
change}\label{certification-is-another-state-change}

Certification should not be hidden inside words such as accepted,
merged, green, or adjudicated.

Let certification state be coordinate:

\texttt{C}.

A certification transition has its own exact target, support
requirements, and authority predicate.

For example:

\texttt{G\_cert(o)=T(o)\ S\_cert(o)\ H(o)},

where:

\begin{itemize}
\tightlist
\item
  \texttt{T(o)} binds the exact target revision;
\item
  \texttt{S\_cert(o)} binds the required support state;
\item
  \texttt{H(o)} binds the governing authority conditions.
\end{itemize}

A positive programme disposition does not automatically set
certification state.

\section{Exact target identity
matters}\label{exact-target-identity-matters}

A review of revision:

\texttt{r\_1}

does not automatically certify later revision:

\texttt{r\_2}.

If \texttt{r\_2} changes the load-bearing artifact, the target identity
has changed.

The state machine should bind evidence, replay, adjudication, and
certification to the exact object they support.

This prevents stale evidence from leaking across revisions.

\section{Forge, Solve, and Cert can be understood as authority
compartments}\label{forge-solve-and-cert-can-be-understood-as-authority-compartments}

The public GCL lifecycle controller records:

\texttt{MATHFORGE\_TO\_MATHSOLVE\_TO\_MATHCERT}

as an authority boundary.

At the level needed here:

\begin{itemize}
\tightlist
\item
  Forge fixes source/problem identity;
\item
  Solve performs bounded work, evidence preservation, replay, and
  adjudication;
\item
  Cert performs a distinct certification function.
\end{itemize}

The chapter does not claim this three-part split is universally optimal.

It uses the public workflow as one concrete example of authority
compartmentalization.

\section{Frontier advancement is not contributor
authority}\label{frontier-advancement-is-not-contributor-authority}

The public Frontier Advancement Gate makes one boundary explicit:

a contributor may return a next residual.

That residual is evidence.

It is not programme scheduling authority.

The programme separately chooses a disposition such as advance,
independently verify, formalize, adversarially attack, source-block,
refute, or close without successor.

This distinction prevents a useful contributor from silently becoming
the scheduler of protected research state.

\section{Independent-intelligence intake is an authority
boundary}\label{independent-intelligence-intake-is-an-authority-boundary}

The public Controlled Epistemic Interface describes a narrow crossing:

\texttt{dispatch\ -\textgreater{}\ bounded\ contribution\ -\textgreater{}\ immutable\ intake\ -\textgreater{}\ separate\ adjudication\ -\textgreater{}\ protected\ route}.

Its compact boundary is:

\texttt{receive\ !=\ believe\ !=\ admit\ !=\ certify}.

That is a state-machine statement.

Every verb denotes a different possible transition.

\section{Zero-context is a work-context
property}\label{zero-context-is-a-work-context-property}

A bounded dispatch can be designed so that a contributor need not absorb
the repository's local framing before attempting the task.

This can reduce one kind of institutional-context coupling.

It does not prove statistical independence, cryptographic independence,
absence of shared pretraining, or absence of common mathematical
culture.

The source itself makes this boundary explicit.

\section{Fail closed means preserve protected
state}\label{fail-closed-means-preserve-protected-state}

Suppose a transition guard fails.

The default protected behavior in this chapter is:

\begin{quote}
do not perform the protected canonical mutation.
\end{quote}

That does not mean nothing happens.

The system may still append a failure receipt, capture diagnostics, open
a repair issue, preserve a candidate artifact, or retry an allowed
check.

Fail-closed promotion and active recovery are compatible.

\section{Safety is not liveness}\label{safety-is-not-liveness}

This distinction is essential.

A system can preserve every safety invariant and still make no progress.

Consider a state in which evidence identities remain correct, duplicate
retries do not multiply effects, invalid transitions are blocked,
certification remains protected, and the advancement disposition stays
false forever.

No declared safety rule is violated.

Yet the claim never advances.

Safety answers:

\begin{quote}
what bad transitions are forbidden?
\end{quote}

Liveness answers:

\begin{quote}
what good transition eventually occurs?
\end{quote}

They need separate arguments.

\section{A finite safety-without-liveness
witness}\label{a-finite-safety-without-liveness-witness}

Let the machine be in an adjudicated state with:

\texttt{D(e)=0}.

Suppose the only future event available is an exact retry that leaves:

\texttt{D(e)=0}.

Every retry is canonically idempotent.

Every protected invariant remains true.

But:

\texttt{eventually(advanced(e))}

is false on this execution.

That is enough to show:

\begin{quote}
fail-closed safety alone does not guarantee progress.
\end{quote}

\section{Recovery needs typed
classes}\label{recovery-needs-typed-classes}

After a failure or block, ``try again'' is too vague.

At least five different operations should be distinguished.

\subsection{Retry}\label{retry}

The same event/evidence identity is attempted again.

\subsection{New evidence}\label{new-evidence}

A different immutable result identity arrives.

\subsection{Supersession}\label{supersession}

A governed relation explicitly marks one artifact as replacing another
for a declared purpose.

\subsection{Implementation repair}\label{implementation-repair}

Tooling or workflow code changes while protected evidence identity
remains intact.

\subsection{Policy change}\label{policy-change}

The governing guard or authority rule itself changes through an
authorized route.

These operations have different provenance and authority meanings.

\section{Recoverable tooling failure is not scientific
refutation}\label{recoverable-tooling-failure-is-not-scientific-refutation}

A CI outage, rate limit, parser defect, or transient connector failure
does not by itself falsify the research claim.

Likewise, a mathematical refutation is not a tooling failure.

The state machine should type these failures differently.

This prevents infrastructure noise from silently becoming epistemic
judgment.

\section{Protected repair requires
replay}\label{protected-repair-requires-replay}

If implementation code changes after a failed run, old execution
evidence should not be relabeled as evidence for the repaired
implementation.

The repaired artifact needs fresh replay at its new exact identity when
the governing process requires it.

This is the operational form of:

\begin{quote}
repair, then re-evaluate the repaired object.
\end{quote}

\section{Public lifecycle case
study}\label{public-lifecycle-case-study}

The pinned GCL lifecycle controller records the concrete sequence:

\texttt{READY\ -\textgreater{}\ LAUNCHED\ -\textgreater{}\ RETURNED\ -\textgreater{}\ CAPTURED\ -\textgreater{}\ REPLAYED\ -\textgreater{}\ ADJUDICATED\ -\textgreater{}\ ADVANCED}.

It also records durable GitHub evidence, protected merges, allowed and
prohibited actions, fail-closed behavior on identity/provenance
mismatch, and event-driven advancement plus a recovery poll.

This is evidence that one real workflow can be encoded as a state
machine.

It is not proof that every research programme should use the same
states.

\section{The case study also exposes authority
limits}\label{the-case-study-also-exposes-authority-limits}

The same controller explicitly prohibits actions such as claiming
mathematical correctness not established by Solve adjudication, MATHCERT
certification, topology redesign, and direct protected-main pushes.

That is important.

A controller's permissions are part of its semantics.

Automation should be described not only by what it can do, but also by
what it is forbidden to do.

\section{State machines make authority
inspectable}\label{state-machines-make-authority-inspectable}

An informal workflow often hides authority in convention.

A typed transition machine can make authority explicit through event
type, actor, target, guard, and permitted mutation.

This does not make governance automatically correct.

It makes the governance claim inspectable.

\section{The machine does not prove the
science}\label{the-machine-does-not-prove-the-science}

Suppose every lifecycle transition is valid.

That establishes workflow conformance.

It does not establish mathematical truth, novelty, completeness of
review, absence of omitted premises, correctness of every human
judgment, or adequacy of the governance policy.

Formal workflow correctness and scientific correctness remain different
claims.

\section{Machine state should not collapse scientific
uncertainty}\label{machine-state-should-not-collapse-scientific-uncertainty}

A programme may legitimately store unresolved status, conflicting
evidence, source-blocked status, conditional support, falsification
under one assumption, or an open question under another.

A state machine should preserve those distinctions.

Governance is not improved by forcing every question into a binary
accepted/rejected flag.

\section{Bounded work packages protect handoff
quality}\label{bounded-work-packages-protect-handoff-quality}

A work package is strongest when a new contributor can determine the
exact target, premises, allowed sources, required output, return route,
authority boundary, and stop condition.

That is why the bounded package belongs inside the state-machine
chapter.

It is the unit on which safe delegation operates.

\section{Independent actors need bounded interfaces
too}\label{independent-actors-need-bounded-interfaces-too}

Actor separation is useful only if the transferred work is well defined.

Otherwise two nominally independent contributors may solve different
problems while appearing to review one another.

Stable identities and bounded work packages give the separation
predicate something concrete to refer to.

\section{Idempotence enables automation under
retries}\label{idempotence-enables-automation-under-retries}

Distributed automation is rarely exactly-once at the transport layer.

Events can be delivered twice, retried after timeout, replayed after a
crash, or awakened by multiple triggers.

Canonical idempotence allows those operational duplicates without
multiplying protected effects.

That is one reason stable identity is not administrative decoration.

\section{History should remain append-only where
possible}\label{history-should-remain-append-only-where-possible}

An append-only history preserves failed attempts, retries, repairs,
policy changes, and supersession events.

Canonical state can still expose the current protected interpretation.

The two views answer different questions.

Canonical state asks:

\begin{quote}
what counts now?
\end{quote}

History asks:

\begin{quote}
how did we get here?
\end{quote}

\section{A practical research-state
ledger}\label{a-practical-research-state-ledger}

Before accepting a protected state change, record:

{\def\LTcaptype{none} % do not increment counter
\begin{longtable}[]{@{}
  >{\raggedright\arraybackslash}p{(\linewidth - 2\tabcolsep) * \real{0.5000}}
  >{\raggedright\arraybackslash}p{(\linewidth - 2\tabcolsep) * \real{0.5000}}@{}}
\toprule\noalign{}
\begin{minipage}[b]{\linewidth}\raggedright
Field
\end{minipage} & \begin{minipage}[b]{\linewidth}\raggedright
Question
\end{minipage} \\
\midrule\noalign{}
\endhead
\bottomrule\noalign{}
\endlastfoot
work identity & Which exact bounded package is this? \\
target & Which exact artifact/revision is affected? \\
event type & Receipt, replay, adjudication, promotion, certification, or
other? \\
actor & Who/what is attempting the transition? \\
evidence identity & Which immutable result is involved? \\
guard & Which conditions must be true? \\
authority & Who is permitted to perform this mutation? \\
canonical effect & What protected state changes? \\
history effect & What attempt/receipt is appended? \\
retry key & How are exact repeats recognized? \\
conflict rule & What happens to a different result under the same
dispatch? \\
failure class & Tooling, evidence, policy, authority, or science? \\
recovery & Retry, repair, new evidence, supersession, or policy
change? \\
certification & Is this a separate transition? \\
frontier effect & Who may authorize a successor? \\
\end{longtable}
}

This ledger turns vague workflow verbs into inspectable state
transitions.

\section{Failure modes}\label{failure-modes-20}

\subsection{Done-as-one-state}\label{done-as-one-state}

Receipt, replay, adjudication, promotion, and certification are
collapsed into one label.

\subsection{Retry-multiplies-authority}\label{retry-multiplies-authority}

The same immutable result gains multiple effects because transport
repeated.

\subsection{Dispatch-key-overwrite}\label{dispatch-key-overwrite}

A different result under the same work package silently replaces prior
evidence.

\subsection{Replay-equals-truth}\label{replay-equals-truth}

A successful execution is treated as proof of the surrounding scientific
claim.

\subsection{Reviewer-label-equals-independence}\label{reviewer-label-equals-independence}

Different role names are treated as evidence of independent judgment.

\subsection{Contributor-equals-scheduler}\label{contributor-equals-scheduler}

A suggested next residual is treated as authorization to advance the
research frontier.

\subsection{Safe-equals-live}\label{safe-equals-live}

Fail-closed preservation is treated as proof the programme will
eventually make progress.

\subsection{Repair-reuses-stale-evidence}\label{repair-reuses-stale-evidence}

A changed implementation is promoted using evidence bound to an older
revision.

\section{What the finite witness
establishes}\label{what-the-finite-witness-establishes}

The companion witness proves only these exact finite facts:

\begin{itemize}
\tightlist
\item
  first capture of \texttt{(17,A)} creates one canonical evidence
  identity;
\item
  exact retry leaves the evidence count at one;
\item
  distinct result \texttt{(17,B)} raises the preserved count to two;
\item
  advancement vector \texttt{(1,1,1,1,0)} evaluates to zero;
\item
  advancement vector \texttt{(1,1,1,1,1)} evaluates to one;
\item
  missing-check vector \texttt{(1,1,0,1,1)} evaluates to zero;
\item
  repeated exact advancement has one canonical set-insertion effect;
\item
  same-actor separation predicate \texttt{Sep(A,A)} evaluates to zero;
\item
  a fail-closed execution can preserve all declared safety invariants
  while never reaching advancement.
\end{itemize}

It does not prove that the five-factor gate is universal or that
identity-separated actors are epistemically independent.

\section{Downstream handoff}\label{downstream-handoff-13}

\textbf{Governed Adaptation --- ATLAS-CH-GOVADAPT-001} may now assume
bounded work packages, product research state, canonical/history
separation, exact-identity idempotent retry, duplicate/conflict
distinction, guarded programme advancement, separate certification
state, explicit actor-separation predicates, fail-closed safety versus
liveness, typed recovery operations, and exact-target evidence binding.

It must independently answer the harder question:

\begin{quote}
when a system is allowed to modify its own behavior, policy, tools,
memory, or authority, which state-machine invariants and
correction-capacity constraints must survive that adaptation?
\end{quote}

\section{References used in this
chapter}\label{references-used-in-this-chapter-76}

\begin{itemize}
\tightlist
\item
  Atlas prerequisite: \emph{Replayable Evidence Objects}
  (\texttt{ATLAS-CH-REPLAY-001}) and its audited source lock.
\item
  Atlas prerequisite: \emph{Formal Methods and Machine-Checkable Claims}
  (\texttt{ATLAS-CH-FORMAL-001}) and \texttt{AUDIT-025}.
\item
  GCL public lifecycle controller:
  \texttt{grandchallenge/MATH-PROGRAMME@fdd7a3fe3df7b2d699753347080c1cbc2127e02d},
  \texttt{governance/openmath\_unattended\_lifecycle\_controller.json}.
\item
  GCL public Frontier Advancement Gate at the same protected commit.
\item
  GCL public Controlled Epistemic Interface at the same protected
  commit.
\end{itemize}

Exact source identities, blob hashes, and authority boundaries are
recorded in:

sources/source-locks/ATLAS-CH-RESEARCHSM-001.yaml

\chapter{Governed Adaptation}\label{governed-adaptation-1}

\textbf{Epistemic status:} audited Optionality and
Research-State-Machine prerequisites + bounded
corrigibility/interruptibility literature + Atlas synthesis.

Adaptive systems change.

Parameters move. Policies update. Data mixtures shift. Components are
replaced. New evidence arrives. Research programmes advance from one
protected state to another.

The difficult question is not whether change is possible. It is:

\begin{quote}
\textbf{When a system is allowed to change, what must remain true so
that later correction, provenance, and bounded authority are not
silently destroyed by the change itself?}
\end{quote}

This chapter calls that problem \textbf{governed adaptation}.

The central distinction is:

\[
\boxed{
\text{can change}
\neq
\text{authorized to change}
\neq
\text{safe to promote}
}
\]

The audited Optionality chapter supplies a language for correction
capacity and recoverability.

The audited Research as a State Machine chapter supplies typed states,
guarded transitions, exact identity, idempotence, authority separation,
and fail-closed promotion.

The present chapter composes those ideas into an adaptation contract.

\section{Adaptation as a governed
transition}\label{adaptation-as-a-governed-transition}

Let the current protected state be

\[
g=(r,x,E,A,P,C,L).
\]

Here:

\begin{itemize}
\tightlist
\item
  \(r\) is the exact protected revision identity;
\item
  \(x\) is the operative state;
\item
  \(E\) is canonical evidence;
\item
  \(A\) is the authority relation;
\item
  \(P\) is promotion/programme state;
\item
  \(C\) is certification state;
\item
  \(L\) is durable history.
\end{itemize}

A candidate revision is

\[
q=(id,r_{parent},\Delta,\hat r).
\]

The candidate declares which protected parent it assumes, what bounded
change it proposes, and which exact result identity it produces.

If the parent changed, the proposal may already be stale.

\section{Authority is typed}\label{authority-is-typed}

A common governance error is to collapse every form of permission into
one word.

The Atlas keeps at least five relations separate:

\begin{itemize}
\tightlist
\item
  proposal authority;
\item
  execution authority;
\item
  promotion authority;
\item
  recovery authority;
\item
  certification authority.
\end{itemize}

These can belong to the same actor under one policy. They are still
different relations.

A contributor who may propose a change does not thereby own canonical
promotion.

An executor who can stage a revision does not thereby certify it.

A recovery operator does not thereby acquire policy-writing authority.

\section{Execution is not
promotion}\label{execution-is-not-promotion}

A candidate can be executed in a staging state:

\[
g\xrightarrow{execute(q)}g_q.
\]

This creates a candidate state.

Promotion is separate:

\[
g_q\xrightarrow{promote}g^+.
\]

This separation creates a place to ask:

\begin{itemize}
\tightlist
\item
  Did the exact candidate pass its required checks?
\item
  Are those checks bound to this revision?
\item
  Were the relevant actors authorized?
\item
  Is a recovery path still available?
\item
  Did a repair invalidate earlier evidence?
\item
  Is certification required or merely promotion?
\end{itemize}

The Research-State-Machine prerequisite already established that
successful execution is not sufficient for promotion.

Governed adaptation applies that rule to changing systems.

\section{Evidence belongs to exact
targets}\label{evidence-belongs-to-exact-targets}

Let an evidence object be

\[
e=(id,target,claim,payload,provenance).
\]

Define:

\[
Binds(e,r)
\iff
target(e)=r.
\]

Suppose revision \(r_1\) passes validation.

A repair then produces:

\[
r_2\neq r_1.
\]

The old validation remains historical evidence.

It is not exact-target evidence for \(r_2\).

\[
Binds(e_{r_1},r_2)=0.
\]

The practical rule is simple:

\begin{quote}
\textbf{After a material revision changes exact identity, replay the
load-bearing evidence on the new exact target.}
\end{quote}

This protects against validating one object and promoting another.

\section{Canonical state and
history}\label{canonical-state-and-history}

Adaptation can overwrite the evidence needed to understand the change.

A governed system therefore separates:

\begin{itemize}
\tightlist
\item
  canonical current state;
\item
  append-only transition history;
\item
  evidence objects;
\item
  interpretation;
\item
  review;
\item
  promotion;
\item
  certification.
\end{itemize}

A new revision should not erase the identity of its predecessor merely
because the predecessor is no longer current.

History answers: how did we get here?

Canonical state answers: what is protected now?

Those are different questions.

\section{Correction capacity enters the
gate}\label{correction-capacity-enters-the-gate}

Optionality introduced conditional correction feasibility:

\[
C_{h,\epsilon}(x,\theta)
\]

and ex-ante correction capacity:

\[
CC_{h,\epsilon}(x;b)
=
\mathbb E_{\theta\sim b}
[
C_{h,\epsilon}(x,\theta)
].
\]

Governed adaptation can use this as one explicit admissibility
condition.

For a declared floor \(\kappa\), require:

\[
CC_{h,\epsilon}(x_q;b)\ge\kappa.
\]

This does not say that preserving more optionality is always better.

It says something narrower:

\begin{quote}
If the governing process declares preservation of specified correction
paths to be a requirement, then candidate evaluation must measure those
paths rather than assume they survived.
\end{quote}

\section{A finite admissibility
gate}\label{a-finite-admissibility-gate}

One possible gate is:

\[
Adm(g,q)
=
ParentOK
\cdot
EvidenceOK
\cdot
AuthorityOK
\cdot
Sep
\cdot
Inv
\cdot
\mathbf 1\{CC_{h,\epsilon}\ge\kappa\}.
\]

The factors mean:

\begin{itemize}
\tightlist
\item
  exact parent identity matches;
\item
  required evidence binds to the candidate;
\item
  required authority predicates hold;
\item
  any declared separation relation holds;
\item
  the declared invariant holds;
\item
  the declared correction-capacity floor holds.
\end{itemize}

The product is Boolean.

If any required factor is zero, the candidate is not admissible under
this particular gate.

The Atlas does not claim this factorization is universal.

Its value is that each reason for refusal is typed.

\section{Equal immediate value, unequal correction
capacity}\label{equal-immediate-value-unequal-correction-capacity}

Consider a deliberately small witness.

Let:

\[
\Theta=\{N,D\},
\]

where:

\begin{itemize}
\tightlist
\item
  \(N\): no later defect is discovered;
\item
  \(D\): a later defect is discovered.
\end{itemize}

Use symmetric prior:

\[
b(N)=b(D)=1/2.
\]

Let \(x_0\) be the common protected reference state. Candidate \(A\)
leads to \(x_A\), and candidate \(B\) leads to \(x_B\).

Define

\[
\Delta U(q)=U(x_q)-U(x_0).
\]

For the two candidates,

\[
\Delta U(A)=\Delta U(B)=1.
\]

Both pass the same immediate checks.

Only one preserves a governed recovery path.

\subsection{Candidate A}\label{candidate-a}

If \(N\) occurs, no correction is required.

If \(D\) occurs, an authorized recovery path returns to the accepted
target.

Therefore:

\[
C_{h,0}(x_A,N)=1,
\qquad
C_{h,0}(x_A,D)=1.
\]

Hence:

\[
CC_{h,0}(x_A;b)=1.
\]

\subsection{Candidate B}\label{candidate-b}

If \(N\) occurs, no correction is required.

If \(D\) occurs, the declared correction target is unreachable within
the authorized horizon.

Therefore:

\[
C_{h,0}(x_B,N)=1,
\qquad
C_{h,0}(x_B,D)=0.
\]

Hence:

\[
CC_{h,0}(x_B;b)=1/2.
\]

Immediate value ties.

Correction capacity does not.

\section{Why a utility-only gate misses the
difference}\label{why-a-utility-only-gate-misses-the-difference}

Define:

\[
G_U(q)=\mathbf 1\{\Delta U(q)\ge1\}.
\]

Then:

\[
G_U(A)=G_U(B)=1.
\]

Now declare:

\[
\kappa=1
\]

and define:

\[
G_C(q)
=
G_U(q)
\mathbf 1\{CC_{h,0}(x_q;b)\ge1\}.
\]

Then:

\[
G_C(A)=1,
\qquad
G_C(B)=0.
\]

The witness establishes only this finite separation.

It does not tell us that every real system should use \(\kappa=1\).

It tells us that an objective containing only immediate value can omit a
property the governing process may care about.

\section{Raw action count is not
enough}\label{raw-action-count-is-not-enough}

Suppose both candidate states expose two commands:

\begin{itemize}
\tightlist
\item
  continue;
\item
  rollback.
\end{itemize}

A superficial option count says both states have two actions.

But suppose rollback in one state:

\begin{itemize}
\tightlist
\item
  lacks authority;
\item
  points to an incompatible artifact;
\item
  cannot restore external state;
\item
  exceeds the permitted horizon;
\item
  or fails the declared target semantics.
\end{itemize}

Then the labels are present while the functional correction path is
absent.

This is why Optionality treats functional consequences and correction
feasibility separately from raw action count.

\section{A backup is not recovery}\label{a-backup-is-not-recovery}

``We have a backup'' is an artifact statement.

``We can recover'' is a transition claim.

Recovery can fail because:

\begin{itemize}
\tightlist
\item
  the backup is corrupt;
\item
  dependencies changed;
\item
  external data cannot be restored;
\item
  credentials expired;
\item
  migrations were one-way;
\item
  the recovery actor lacks authority;
\item
  restored bytes do not recreate the declared consequence class.
\end{itemize}

For target class \([r_\star]\), define governed recovery by:

\[
RecAuth_h(g,[r_\star])=1
\]

only when an authorized path of length at most \(h\) reaches the
declared target class.

\section{Corrigibility and intervention
precedents}\label{corrigibility-and-intervention-precedents}

Corrigibility was introduced as a problem of cooperation with corrective
intervention despite incentives that can make capable agents resist
shutdown or preference modification (Soares et al. 2015).

Orseau and Armstrong study safe interruptibility in reinforcement
learning, asking how intervention can occur without training the agent
to seek or avoid interruption (Orseau and Armstrong 2016).

Hadfield-Menell and colleagues analyze an off-switch game in which an
agent's incentive to preserve human intervention depends on uncertainty
about the underlying objective and on treating human action as
informative (Hadfield-Menell et al. 2017).

These are important precedents.

They are not interchangeable.

The Atlas keeps separate:

\begin{itemize}
\tightlist
\item
  corrigibility;
\item
  interruptibility;
\item
  human intervention-channel preservation;
\item
  rollback;
\item
  correction capacity;
\item
  governance authority.
\end{itemize}

The literature motivates the intervention problem. It does not imply the
Atlas gate.

\section{Preserving correction can have
cost}\label{preserving-correction-can-have-cost}

Optionality already established that preserving more options is not
universally optimal.

The same warning applies here.

A recovery path can impose:

\begin{itemize}
\tightlist
\item
  storage cost;
\item
  compatibility burden;
\item
  slower migrations;
\item
  operational complexity;
\item
  delayed optimization.
\end{itemize}

Governed adaptation can therefore treat correction capacity as:

\begin{itemize}
\tightlist
\item
  an objective;
\item
  a constraint;
\item
  or a diagnostic.
\end{itemize}

It is not an unconditional command to preserve every possible past
state.

\section{Bounded authority}\label{bounded-authority}

Authority to alter operative state \(x\) does not automatically include
authority to alter the authority relation \(A\).

If a candidate changes the governance policy itself, that is a different
transition class.

The system needs whatever authority the governing process assigns to
policy change.

Thus:

\[
\boxed{
\text{authority over system state}
\neq
\text{authority over the rule that grants authority}
}
\]

Without this distinction, bounded delegation becomes unrestricted
delegation by implication.

\section{Safety is not liveness}\label{safety-is-not-liveness-1}

A strict gate can refuse unsupported candidates.

That is a safety property.

But a gate can also refuse everything forever.

Consider:

\[
g_0\to g_0\to g_0\to\cdots.
\]

The invariant remains true.

No candidate advances.

Therefore:

\[
\text{fail-closed safety}
\not\Rightarrow
\text{eventual progress}.
\]

A programme claiming eventual adaptation needs a separate liveness
argument.

\section{Governance conformance is not substantive
truth}\label{governance-conformance-is-not-substantive-truth}

Suppose a candidate:

\begin{itemize}
\tightlist
\item
  has exact provenance;
\item
  passes all required checks;
\item
  has proper authority;
\item
  preserves the declared recovery path;
\item
  receives promotion.
\end{itemize}

What has been shown?

That the declared governance contract was satisfied.

What has not automatically been shown?

\begin{itemize}
\tightlist
\item
  that a scientific hypothesis is true;
\item
  that the system is globally safer;
\item
  that every objective improved;
\item
  that the evidence was complete;
\item
  that no undeclared failure mode exists.
\end{itemize}

Process integrity matters because it lets us state exactly what was
established.

\section{Repairs invalidate exact-target
evidence}\label{repairs-invalidate-exact-target-evidence}

A common operational pattern is:

\begin{enumerate}
\def\labelenumi{\arabic{enumi}.}
\tightlist
\item
  candidate \(r_1\) is tested;
\item
  test fails;
\item
  candidate is repaired to \(r_2\);
\item
  someone remembers that ``validation already ran.''
\end{enumerate}

That is insufficient.

The correct sequence is:

\begin{enumerate}
\def\labelenumi{\arabic{enumi}.}
\tightlist
\item
  bind evidence to \(r_1\);
\item
  repair to \(r_2\);
\item
  treat \(r_1\)'s exact-target validation as stale for promotion of
  \(r_2\);
\item
  replay on \(r_2\);
\item
  promote only if the new exact target passes.
\end{enumerate}

This closes the identity gap between evidence and mutation.

\section{Recovery must name a
target}\label{recovery-must-name-a-target}

``Return to normal'' is not a formal recovery target.

A governed recovery should name:

\begin{itemize}
\tightlist
\item
  target revision or consequence class;
\item
  maximum horizon;
\item
  authority;
\item
  acceptable cost or tolerance;
\item
  evidence required after restoration.
\end{itemize}

A rollback can restore bytes while failing to restore behavior.

It can restore behavior while failing to restore external state.

It can restore a previous revision while invalidating later evidence.

Recovery semantics belong in the transition system.

\section{State-relative authority}\label{state-relative-authority}

Authority can be state-relative.

Permission granted for one bounded candidate should not silently cover a
materially different object.

If the candidate's parent, scope, target identity, or governing policy
changes materially, the applicable authority contract should be
re-evaluated.

This is the governance analogue of exact-target evidence.

\section{Why this matters for adaptive
intelligence}\label{why-this-matters-for-adaptive-intelligence}

An adaptive system that cannot preserve correction channels can become
more capable while becoming harder to steer.

A research programme that cannot preserve provenance can become more
productive while becoming less auditable.

A deployment pipeline that cannot preserve exact-target evidence can
become faster while losing the ability to say what was actually
validated.

Governed adaptation is the attempt to preserve the infrastructure of
correction while permitting change.

It is not anti-adaptation.

It is a theory of admissible change under declared constraints.

\section{Failure modes}\label{failure-modes-21}

\subsection{Capability as authority}\label{capability-as-authority}

A system can perform an action and therefore treats itself as permitted
to perform it.

This confuses mechanism with governance.

\subsection{Promotion by execution}\label{promotion-by-execution}

A candidate runs successfully and is treated as canonical.

This skips the promotion gate.

\subsection{Stale green evidence}\label{stale-green-evidence}

A repaired candidate inherits validation from its predecessor.

This violates exact-target binding.

\subsection{Cosmetic rollback}\label{cosmetic-rollback}

A rollback command exists but cannot restore the declared target.

This confuses interface with recoverability.

\subsection{Option counting}\label{option-counting}

A system counts command labels rather than functional correction paths.

\subsection{Corrigibility by naming}\label{corrigibility-by-naming}

An intervention channel exists and the system is therefore called
corrigible.

The stronger property has not been established.

\subsection{Safety without progress}\label{safety-without-progress}

A fail-closed system never violates the gate but also never advances.

\subsection{Governance as truth}\label{governance-as-truth}

A compliant process outcome is treated as proof that a scientific or
empirical claim is correct.

\subsection{Unbounded delegation}\label{unbounded-delegation}

Authority over operative state is interpreted as authority over the
governance contract.

\subsection{Recovery without
revalidation}\label{recovery-without-revalidation}

A prior revision is restored but its current environment and
dependencies are not rechecked.

\section{Downstream handoff}\label{downstream-handoff-14}

Frontier Questions of Adaptive Intelligence may assume:

\begin{itemize}
\tightlist
\item
  governed adaptation is a typed transition problem;
\item
  authority is separated by role;
\item
  candidate evidence binds to exact target identity;
\item
  authorized recovery paths are distinct from stored backups;
\item
  correction capacity can be an explicit gate dimension;
\item
  immediate value can tie while correction capacity differs;
\item
  fail-closed safety does not imply liveness;
\item
  intervention precedents do not collapse into a universal theorem.
\end{itemize}

FRONTIER must independently determine which unresolved questions deserve
programme priority.

\section{Epistemic status}\label{epistemic-status-2}

Established prerequisite layer:

\begin{itemize}
\tightlist
\item
  viable continuation and correction-capacity semantics;
\item
  typed governed states and guarded transitions;
\item
  exact identity, idempotence, recovery classes, and authority
  separation.
\end{itemize}

External precedent layer:

\begin{itemize}
\tightlist
\item
  corrective intervention/corrigibility problem (Soares et al. 2015);
\item
  safe interruptibility in a declared RL setting (Orseau and Armstrong
  2016);
\item
  off-switch incentive analysis under objective uncertainty
  (Hadfield-Menell et al. 2017).
\end{itemize}

Atlas-owned synthesis:

\begin{itemize}
\tightlist
\item
  governed adaptive state;
\item
  authority separation;
\item
  exact-target adaptation gate;
\item
  authorized recovery-path semantics;
\item
  correction-capacity gate composition;
\item
  equal-immediate-value finite witness.
\end{itemize}

Not established here:

\begin{itemize}
\tightlist
\item
  universal safety;
\item
  universal adaptation optimality;
\item
  complete corrigibility;
\item
  universal rollback feasibility;
\item
  substantive truth from governance conformance.
\end{itemize}

\section{References used in this
chapter}\label{references-used-in-this-chapter-77}

\begin{itemize}
\tightlist
\item
  (Soares et al. 2015) --- corrigibility and corrective intervention.
\item
  (Orseau and Armstrong 2016) --- safe interruptibility in reinforcement
  learning.
\item
  (Hadfield-Menell et al. 2017) --- incentives around preserving a human
  off-switch.
\end{itemize}

Exact provenance and authority boundaries are locked in:

sources/source-locks/ATLAS-CH-GOVADAPT-001.yaml

\chapter{Frontier Questions of Adaptive
Intelligence}\label{frontier-questions-of-adaptive-intelligence}

\textbf{Epistemic status:} audited Atlas substrate + project-local
research programme + architecture-stage questions + explicit open proof
and experiment obligations.

\textbf{Programme specification:} `governance/tranches/CH171-SPEC.md`\\
\textbf{Documentary evidence matrix:}
`governance/tranches/CH171-EVIDENCE-MATRIX.md`\\
\textbf{Source lock:} `sources/source-locks/ATLAS-CH-FRONTIER-001.yaml`

\section{A frontier is not a
result}\label{a-frontier-is-not-a-result}

The Atlas has now accumulated enough audited substrate to ask sharper
research questions.

That does not mean those questions are solved.

This chapter separates audited substrate, bounded evidence,
architecture-stage programmes, open proof obligations, open experiment
obligations, and conjectural cross-programme connections.

The governing rule is:

\[
\boxed{\text{research programme}\neq\text{established result}.}
\]

The value of a frontier chapter is that each open direction has a named
next obligation.

\section{The three governing
prerequisites}\label{the-three-governing-prerequisites}

Governed Adaptation contributes exact-target evidence, typed authority,
execution-versus-promotion separation, recovery discipline, and
correction-capacity boundaries.

The Residual contributes a reconstructive question: what is the least
admissible structure that survives declared representation changes while
remaining sufficient for a declared capability?

Boundary Contracts contribute the interface question: what semantic,
geometric, sensitivity, and numerical information must cross a component
boundary for composition to remain meaningful?

These chapters constrain how the programme is stated. They do not solve
the programme.

\section{Geometry-derived
optimization}\label{geometry-derived-optimization}

The second-order chapter is draft-v0.1 at the protected ledger snapshot
and provides curvature, Hessian, Fisher, natural-gradient, trust-region,
quasi-Newton, and proximal distinctions.

The live architecture-stage programmes are Matrix-Aware Optimization,
Optimization on Manifolds, and Variational and Divergence-Derived
Optimization.

The research question is not whether geometric language can describe an
optimizer. It is whether a declared matrix, manifold, conditioning, or
variational structure predicts an update rule whose behavior changes for
a reason that survives comparison with matched controls.

A smallest useful result would bind a problem family, a structural
assumption, an update rule, a matched baseline, a predicted consequence,
and evidence that the structure caused the consequence.

A successful run is evidence. It is not yet a mechanism.

\section{Operator-valued position}\label{operator-valued-position}

The Geometry of Position chapter already establishes exact rotary
relative-phase identities and their limits.

The Relative-Position Operator programme begins beyond that substrate.

Its canonical contract asks whether positional structure is better
exposed through low-dimensional operators, frequency modes, DC
components, and head-specific specialization.

The smallest useful advance is not another restatement of RoPE algebra.

It is a result showing that a compact operator description supports a
reproducible prediction, intervention, compression, or transfer property
that ordinary positional-vector language does not already explain.

Until then:

\[
\boxed{\text{operator notation}\neq\text{operator discovery}.}
\]

\section{Adaptive depth as error
control}\label{adaptive-depth-as-error-control-2}

Depth as Computational Time and Networks as Numerical Schemes are
already audited.

That makes a stronger question possible:

\begin{quote}
Can adaptive compute behave like adaptive time stepping, where more
computation is allocated because a declared error or quality signal says
it is needed?
\end{quote}

Variable depth alone is not error control.

A valid test needs an observable error or quality surrogate, a
controller responding to that signal, a stopping or depth action, and a
fixed-depth control at matched compute or matched quality.

If the error signal does not explain the depth decision, the weaker
phrase ``adaptive compute'' is more accurate.

\section{Persistent and shared
memory}\label{persistent-and-shared-memory}

The External-Memory Thesis is already audited.

The frontier question is no longer simply whether retrieval helps.

It is:

\begin{quote}
Which information belongs in weights, which belongs in active context,
and which belongs in persistent shared memory?
\end{quote}

A shared memory layer introduces requirements for identity, indexing,
provenance, updates, consistency, transactions, and access across models
or agents.

The thesis does not imply that all useful knowledge should leave
parameters.

It turns knowledge placement into an allocation problem under declared
costs and governance.

\section{Learning progress as a search
operator}\label{learning-progress-as-a-search-operator-1}

The Curriculum Learning chapter already distinguishes open-loop
schedules from state-aware selection and shows that progress depends on
learner state.

The next programme asks whether learning progress can become a search
signal over experiences.

The key risk is myopia.

An experience can produce high immediate gain while harming later
transfer, causing forgetting, consuming disproportionate compute, or
steering the learner toward a narrow local mechanism.

So the frontier question is not ``can progress be maximized?''

It is:

\begin{quote}
Under what conditions does progress feedback guide experience search
toward the declared long-horizon objective?
\end{quote}

That requires horizon-aware controls and separation of dose, difficulty,
and compute confounds.

\section{Minimal curricula and reasoning
bases}\label{minimal-curricula-and-reasoning-bases-1}

The Minimal Curricula and Reasoning Bases chapter remains
architecture-stage.

Its question is deliberately stronger than ordinary curriculum
optimization:

\begin{quote}
What is the smallest early set of mechanisms or reasoning operations
from which broad later capability can be reconstructed?
\end{quote}

``Smallest'' requires a declared mechanism vocabulary, capability
family, reconstruction process, resource budget, and equivalence
relation.

A short syllabus is not automatically a minimal curriculum.

A small dataset is not automatically a minimal reasoning basis.

The programme must also distinguish acquisition, persistence,
accessibility, and behavioral expression.

\section{The Residual remains open at frontier
scale}\label{the-residual-remains-open-at-frontier-scale}

The Residual chapter has an audited formalization.

It does not claim a universal frontier-scale Residual.

The object is relative to a declared state space, transformation family,
capability probes, behavior profile, descriptor class, and admissible
recovery maps.

Open questions remain: existence, useful uniqueness, finite
dimensionality, efficient learnability, transfer across architecture
families, and operational reconstruction at acceptable cost.

The next serious advance is a nontrivial declared family in which
existence can be proved or falsified.

\section{Benign nonconvexity}\label{benign-nonconvexity-1}

The Atlas source inventory names benign nonconvexity, landscape
topology, saddle structure, conditioning, and optimizer dynamics as open
research objects.

At present benign nonconvexity is a programme theme, not a theorem.

A meaningful claim needs a declared landscape class, a declared
optimizer or dynamics, and a declared target property.

That property might concern reachability, absence of bad attractors,
quotient simplification, correction capacity, or another explicit
mechanism.

A converged training run is not sufficient.

The conjectural connection is that optimizer geometry and landscape
structure may have to be analyzed together.

\section{Boundary contracts}\label{boundary-contracts-1}

Boundary Contracts are already audited substrate.

The next frontier is to determine which exposed boundary variables are
actually sufficient for composition.

A useful contract might expose semantics, invariants, local
sensitivities, error budgets, and low-order separator state.

The audited chapter also shows that a low-dimensional separator can
forget state that matters downstream.

So the open question is:

\begin{quote}
Under which assumptions do local interface obligations compose into a
useful system-level guarantee?
\end{quote}

Local compatibility alone is not a global guarantee.

\section{Source boundaries}\label{source-boundaries-1}

Research vocabulary and source evidence are different objects.

When project terminology is stronger than the exact source-locked
artifact, the Atlas keeps the stronger wording at programme status and
binds factual implementation claims to the inspected source.

This rule applies across the programme:

\[
\boxed{\text{programme vocabulary}\neq\text{source evidence}.}
\]

\section{Cross-programme questions}\label{cross-programme-questions}

Several connections are worth testing but are not inherited results.

The evidence matrix records them explicitly:

\begin{itemize}
\tightlist
\item
  Residual structure may be relevant to reconstructive curriculum
  questions.
\item
  Boundary contracts may carry compact state needed by downstream
  composition.
\item
  Learning-progress feedback may help select experiences that expose
  reusable computation.
\item
  Adaptive depth may become an error-control mechanism when a valid
  error signal exists.
\item
  Shared memory may hold reconstructive state that need not be
  duplicated in every model.
\end{itemize}

Each remains conjectural until a declared proof or experiment supports
the stronger relation.

\section{How a programme item
advances}\label{how-a-programme-item-advances}

A programme item advances when new evidence supports a more specific
claim.

Governed Adaptation supplies the discipline:

\begin{itemize}
\tightlist
\item
  bind evidence to the exact target;
\item
  state the claim class;
\item
  preserve provenance;
\item
  use the relevant comparison or proof obligation;
\item
  replay evidence after identity-changing repair;
\item
  separate execution from promotion.
\end{itemize}

This prevents one bounded result from silently becoming a stronger
general claim.

\section{Negative results are
progress}\label{negative-results-are-progress}

A research programme should become smaller when it learns.

A counterexample can remove a conjecture. A failed replication can
shrink a regime. A lower bound can eliminate an approach. A source audit
can weaken an unsupported claim. A proof can reduce a broad question to
one missing estimate.

Progress is not only accumulation. It is reduction of uncertainty in the
programme graph.

\section{Preferred research
pattern}\label{preferred-research-pattern}

The Atlas prefers:

broad question -\textgreater{} bounded mechanism -\textgreater{}
smallest named proof or experiment -\textgreater{} exact evidence
-\textgreater{} revised programme.

A frontier chapter should expose handles, not merely aspirations.

\section{What this chapter does not
claim}\label{what-this-chapter-does-not-claim-4}

This chapter does not claim that the architecture-stage programmes
recorded in the evidence matrix are already solved.

It does not promote bounded examples into universal results.

It does not convert programme vocabulary into public implementation
evidence.

The exact open boundaries remain those recorded in the documentary
matrix and source lock.

\section{Why collect the programme
now?}\label{why-collect-the-programme-now}

The substrate has matured enough to make the questions less vague.

Optimization can now be decomposed into curvature, matrix structure,
manifolds, and variational derivation.

Position can now be separated from exact rotary algebra and tested as an
operator hypothesis.

Adaptive compute can now be asked as an error-control question.

Minimal knowledge can now be stated in reconstructive terms.

The frontier becomes sharper as the substrate becomes less ambiguous.

\section{Downstream handoff}\label{downstream-handoff-15}

ATLAS-CH-SYNTHESIS-001 may inherit:

\begin{itemize}
\tightlist
\item
  the programme taxonomy;
\item
  the audited/open distinction;
\item
  the named proof and experiment obligations;
\item
  governed advancement rules.
\end{itemize}

It may not treat unresolved programme items as completed pillars merely
because they appear here.

The final synthesis must still distinguish what is known, engineered,
measured, hypothesized, and open.

\section{References used in this
chapter}\label{references-used-in-this-chapter-78}

This chapter introduces no new broad external literature survey.

Its authority is project-local and source-locked:

\begin{itemize}
\tightlist
\item
  audited Governed Adaptation;
\item
  audited Residual formalization;
\item
  audited Boundary Contracts;
\item
  the Atlas Map;
\item
  the Atlas source inventory;
\item
  the Chapter-Family Rollout;
\item
  the exact Chapter Ledger status snapshot.
\end{itemize}

Exact provenance and claim boundaries are locked in:

`sources/source-locks/ATLAS-CH-FRONTIER-001.yaml`

\chapter{Beyond the Monolithic Model}\label{beyond-the-monolithic-model}

\textbf{Epistemic status:} audited Frontier and Computational Polity
prerequisites + Atlas-owned exact systems-composition witness.\\
\textbf{Specification:}
manuscript/specifications/ATLAS-CH-SYNTHESIS-001.md\\
\textbf{Derivation packet:}
mathematics/derivations/ATLAS-CH-SYNTHESIS-001-DERIVATIONS.md\\
\textbf{Computational witness:}
mathematics/computational-witnesses/ATLAS-CW-SYNTHESIS-001.md\\
\textbf{Source lock:} sources/source-locks/ATLAS-CH-SYNTHESIS-001.yaml

The Atlas does not end by declaring one universal model architecture.

It ends by making the system boundary explicit.

A deployed capability can depend on representation, dynamics, memory,
tools, coordination, evidence, adaptation, and governance at once. When
those interfaces are part of the mechanism, they belong inside the
explanatory object.

The governing rule is:

\[
\boxed{
\text{system capability is a property of declared composition, not component count alone.}
}
\]

\section{A typed synthesis object}\label{a-typed-synthesis-object}

Write

\[
\Sigma=(X,D,M,T,C,E,A,G,Q,K).
\]

Here:

\begin{itemize}
\tightlist
\item
  \(X\) is model/representation and private state;
\item
  \(D\) is dynamics and update law;
\item
  \(M\) is persistent/shared memory;
\item
  \(T\) is tools and external services;
\item
  \(C\) is coordination, routing, and composition;
\item
  \(E\) is evidence and validation;
\item
  \(A\) is adaptation;
\item
  \(G\) is governance and authority;
\item
  \(Q\) is the task/success contract;
\item
  \(K\) is the heterogeneous cost vector.
\end{itemize}

These are typed objects. They need not share a vector space, timescale,
optimizer, or authority model.

\section{Frontier status survives
synthesis}\label{frontier-status-survives-synthesis}

FRONTIER supplies the exact programme grammar:

\begin{itemize}
\tightlist
\item
  AUDIT\_BOUND\_SUBSTRATE;
\item
  DRAFT\_SUBSTRATE;
\item
  BOUNDED\_EVIDENCE;
\item
  ARCHITECTURE\_STAGE;
\item
  OPEN\_PROOF\_OBLIGATION;
\item
  OPEN\_EXPERIMENT\_OBLIGATION;
\item
  CONJECTURAL\_CONNECTION.
\end{itemize}

SYNTHESIS preserves these categories. Narrative coherence does not
promote an open item into established substrate.

\section{Polity becomes a systems
interface}\label{polity-becomes-a-systems-interface}

POLITY already separates private state, shared state, candidate
production, validation, authority, commit, and cost.

SYNTHESIS extends the explanatory boundary without erasing those
separations.

Capability is not authority.

Evidence is not governance.

Stored state is not truth.

Component count is not composition quality.

\section{Exact finite composition}\label{exact-finite-composition}

Let

\[
Q=\{\alpha,\beta\}.
\]

Define specialists

\[
A_\alpha(\alpha)=1,\qquad A_\alpha(\beta)=0,
\]

\[
A_\beta(\alpha)=0,\qquad A_\beta(\beta)=1.
\]

Each specialist alone succeeds on one of the two tasks.

Use router

\[
R(\alpha)=A_\alpha,\qquad R(\beta)=A_\beta.
\]

The routed system returns the correct answer \(1\) on both tasks.

Thus task accuracy is

\[
\boxed{1}.
\]

\section{Evidence and provenance}\label{evidence-and-provenance}

Each candidate is

\[
c=(q,y,s),
\]

where \(s\) is the producing specialist.

The validator accepts only when:

\begin{enumerate}
\def\labelenumi{\arabic{enumi}.}
\tightlist
\item
  \(y=1\);
\item
  \(s\) is the declared routed specialist for \(q\).
\end{enumerate}

This finite validator checks both answer and provenance. It is not
promoted into a universal oracle.

\section{Governance and commit}\label{governance-and-commit}

Governance authorizes commit only after validator acceptance.

The transition is therefore

\[
\text{candidate}\rightarrow
\text{validation}\rightarrow
\text{authorization}\rightarrow
\text{commit}.
\]

These are different events.

Correct-looking output cannot silently bypass the authority boundary
without changing the system.

\section{Shared memory}\label{shared-memory}

Committed records are

\[
m=(q,y,s,\mathrm{validated}).
\]

They persist in shared memory keyed by task.

Later recall therefore has access to both accepted answer and
provenance.

For the full composition:

\begin{itemize}
\tightlist
\item
  task accuracy \(=2/2\);
\item
  validated commit coverage \(=2/2\);
\item
  persistent recall coverage \(=2/2\);
\item
  unauthorized commits \(=0\).
\end{itemize}

\section{Router ablation}\label{router-ablation}

Keep the same specialists but use

\[
R'(\alpha)=A_\alpha,\qquad R'(\beta)=A_\alpha.
\]

Now beta receives answer \(0\).

The validator rejects the beta candidate.

Task accuracy and validated commit coverage each fall to

\[
\boxed{\frac12}.
\]

The component set is unchanged. The composition rule changed.

\section{Validator ablation}\label{validator-ablation}

Restore correct routing but remove the validator while keeping
governance unchanged.

The system can still produce correct transient answers on both tasks.

But governance requires positive validation evidence.

Without it, authorized commit coverage becomes

\[
\boxed{0}.
\]

Thus

\[
\boxed{
\text{answer capability}\not\Rightarrow\text{authorized durable state}.
}
\]

\section{Memory ablation}\label{memory-ablation}

Keep routing, validation, and governance.

Remove persistent shared memory.

Immediate answers can remain correct.

Later persistent recall coverage becomes

\[
\boxed{0}.
\]

Thus

\[
\boxed{
\text{immediate success}\not\Rightarrow\text{persistent shared recall}.
}
\]

\section{Governance ablation}\label{governance-ablation}

Consider

\[
c_{\mathrm{bad}}=(\beta,1,A_\alpha).
\]

The answer value looks correct but the source is wrong.

The validator rejects it.

With governance enforced, the rejected candidate cannot commit.

If the commit gate is bypassed, it can be written.

Therefore

\[
\boxed{
\text{correct-looking content}\not\Rightarrow\text{authorized state transition}.
}
\]

\section{Capability and authority}\label{capability-and-authority}

A component may be capable of producing an output it is not allowed to
commit.

A role may be authorized to commit a validated result without being
capable of producing that result itself.

So

\[
\operatorname{Can}\neq\operatorname{May}.
\]

This distinction becomes more important as the system boundary expands.

\section{Evidence and governance}\label{evidence-and-governance}

Evidence supports a claim.

Governance constrains transitions.

Governance cannot manufacture missing evidence.

Strong evidence does not automatically grant authority.

Hence

\[
\boxed{
\text{evidence}\neq\text{authorization}.
}
\]

\section{Memory and truth}\label{memory-and-truth}

Shared memory can preserve accepted records.

It cannot make a record true merely by storing it.

The inherited boundary remains:

\begin{itemize}
\tightlist
\item
  stored does not mean true;
\item
  present does not mean current;
\item
  shared does not mean universally visible;
\item
  retrieved does not mean correctly used.
\end{itemize}

\section{Adaptation and permission}\label{adaptation-and-permission}

A system may be able to change weights, memory, routing, prompts, tools,
or policies.

That ability belongs to \(A\).

Whether the change is permitted belongs to \(G\).

Thus adaptation and authority must remain separate.

\section{Tools inside the system
boundary}\label{tools-inside-the-system-boundary}

If a capability depends on search, compilers, proof checkers, databases,
laboratories, or external services, those interfaces belong in the
system-level claim.

A tool response is not automatically the same event as a durable
external state change.

The transaction boundary must remain explicit.

\section{Heterogeneous cost}\label{heterogeneous-cost}

Retain the polity cost vector

\[
K=(c_{\mathrm{coord}},c_{\mathrm{mem}},c_{\mathrm{tool}},c_{\mathrm{human}},c_{\mathrm{val}}).
\]

A composite system can improve one capability while increasing several
costs.

No scalar comparison is justified until the application declares a
scalarization or objective.

\section{Specialization does not imply net
benefit}\label{specialization-does-not-imply-net-benefit}

The finite specialist witness has a clear composition gain.

Real specialization can add routing latency, coordination overhead,
memory contention, validation cost, and new failure modes.

Therefore

\[
\boxed{
\text{specialization benefit}\not\Rightarrow\text{net system benefit under every cost model}.
}
\]

\section{The monolithic model remains
legitimate}\label{the-monolithic-model-remains-legitimate}

Nothing here proves that one model cannot internalize memory-like state,
routing, tool use, validation, or adaptation.

The synthesis claim is not that monoliths are obsolete.

It is that the explanatory boundary should follow the actual mechanism.

\section{A monolith can be internally
composite}\label{a-monolith-can-be-internally-composite}

One externally packaged model may contain conditional computation,
recurrent state, modular subspaces, iterative computation, routing, or
memory.

So ``monolithic'' is not a primitive mathematical category.

The relevant question is which internal or external interfaces are
necessary for the declared capability.

\section{A distributed system can still be poorly
composed}\label{a-distributed-system-can-still-be-poorly-composed}

Many components can wrap one weak central capability.

A collection of strong components can fail under bad routing or invalid
interfaces.

Therefore

\[
\boxed{
\text{distributed}\not\Rightarrow\text{capable}.
}
\]

Composition must be demonstrated.

\section{Evidence levels remain
visible}\label{evidence-levels-remain-visible}

The Atlas contains exact derivations, finite computational witnesses,
audited engineering contracts, bounded empirical evidence, hypotheses,
open proof obligations, and open experiment obligations.

SYNTHESIS does not convert these into one evidence level.

\section{Open research remains
open}\label{open-research-remains-open}

FRONTIER names research programmes and explicit obligations.

Some downstream chapters have supplied stronger local results.

Any unresolved obligation remains unresolved until separately discharged
and audited.

The final chapter is not a certification shortcut.

\section{Durable Atlas disciplines}\label{durable-atlas-disciplines}

Across the Atlas, several practices recur:

\begin{enumerate}
\def\labelenumi{\arabic{enumi}.}
\tightlist
\item
  type the mathematical object;
\item
  declare the state and update law;
\item
  separate local and global claims;
\item
  expose boundary contracts;
\item
  separate residual, error, and task metrics;
\item
  distinguish memory from truth;
\item
  distinguish capability from authority;
\item
  bind evidence to exact identity;
\item
  replay before promotion;
\item
  preserve open obligations.
\end{enumerate}

These are disciplines, not one prescribed architecture.

\section{Architecture remains a design
choice}\label{architecture-remains-a-design-choice}

Different systems may place memory, routing, tools, validation,
adaptation, human review, and authority at different boundaries.

The right placement depends on tasks, failure models, costs, and
governance constraints.

The Atlas does not prescribe one universal placement.

\section{Systems claims need a
contract}\label{systems-claims-need-a-contract}

A serious claim should identify:

\begin{itemize}
\tightlist
\item
  task set/distribution;
\item
  component roles;
\item
  private/shared state;
\item
  routing/composition;
\item
  memory semantics;
\item
  evidence rule;
\item
  adaptation rule;
\item
  governance rule;
\item
  success metric;
\item
  cost model;
\item
  failure model.
\end{itemize}

Without these, ``the system can do X'' is underspecified.

\section{Ablation is part of
synthesis}\label{ablation-is-part-of-synthesis}

If a system-level capability is claimed to depend on composition,
perturb the interface.

The finite witness shows distinct failures:

\begin{itemize}
\tightlist
\item
  routing ablation harms task capability;
\item
  validator ablation harms governed commit;
\item
  memory ablation harms persistence;
\item
  governance ablation harms authorization integrity.
\end{itemize}

Different interfaces support different properties.

\section{Research and engineering
consequence}\label{research-and-engineering-consequence}

A single benchmark score can hide mechanism.

A typed systems view can expose where information lives, where
uncertainty enters, which boundary matters, which role holds authority,
and which evidence supports promotion.

That makes claims more auditable even when the underlying system remains
complex.

\section{No universal architecture
theorem}\label{no-universal-architecture-theorem}

The Atlas does not prove

\[
\text{distributed}>\text{monolithic},
\]

or

\[
\text{more modules}>\text{fewer modules}.
\]

Such comparisons require declared tasks, metrics, costs, and failure
models.

\section{The durable synthesis}\label{the-durable-synthesis}

The strongest justified statement is

\[
\boxed{
\text{adaptive intelligence can be analyzed as a typed system of representations, dynamics, memory, tools, coordination, evidence, adaptation, and governance.}
}
\]

When those interfaces are part of the mechanism, they belong inside the
explanatory boundary.

The final lesson is therefore not ``beyond models.''

It is beyond pretending that every consequential system property must
live inside one undifferentiated model boundary.

\section{References used in this
chapter}\label{references-used-in-this-chapter-79}

No new external academic authority is added.

Programme-status authority is inherited through audited
ATLAS-CH-FRONTIER-001.

System-composition authority is inherited through audited
ATLAS-CH-POLITY-001.

Exact prerequisite identities and claim boundaries are recorded in:

sources/source-locks/ATLAS-CH-SYNTHESIS-001.yaml

\protect\phantomsection\label{refs}
\begin{CSLReferences}{1}{1}
\bibitem[\citeproctext]{ref-AbsilMahonySepulchre2008}
Absil, P.-A., R. Mahony, and R. Sepulchre. 2008. \emph{Optimization
Algorithms on Matrix Manifolds}. Princeton University Press.
\url{https://doi.org/10.1515/9781400830244}.

\bibitem[\citeproctext]{ref-AchilleSoatto2018}
Achille, Alessandro, and Stefano Soatto. 2018. {``Emergence of
Invariance and Disentanglement in Deep Representations.''} \emph{Journal
of Machine Learning Research} 19 (50): 1--34.
\url{https://www.jmlr.org/papers/v19/17-646.html}.

\bibitem[\citeproctext]{ref-Amari1998NaturalGradient}
Amari, Shun-ichi. 1998. {``Natural Gradient Works Efficiently in
Learning.''} \emph{Neural Computation} 10 (2): 251--76.
\url{https://doi.org/10.1162/089976698300017746}.

\bibitem[\citeproctext]{ref-Arnoldi1951}
Arnoldi, Walter Edwin. 1951. {``The Principle of Minimized Iterations in
the Solution of the Matrix Eigenvalue Problem.''} \emph{Quarterly of
Applied Mathematics} 9 (1): 17--29.
\url{https://www.ams.org/qam/1951-09-01/}.

\bibitem[\citeproctext]{ref-ArrowFisher1974}
Arrow, Kenneth J., and Anthony C. Fisher. 1974. {``Environmental
Preservation, Uncertainty, and Irreversibility.''} \emph{The Quarterly
Journal of Economics} 88 (2): 312--19.
\url{https://doi.org/10.2307/1883074}.

\bibitem[\citeproctext]{ref-Aubin1991Viability}
Aubin, Jean-Pierre. 1991. \emph{Viability Theory}. Birkhäuser Boston.

\bibitem[\citeproctext]{ref-AuerEtAl2002Bandit}
Auer, Peter, Nicolò Cesa-Bianchi, and Paul Fischer. 2002. {``Finite-Time
Analysis of the Multiarmed Bandit Problem.''} \emph{Machine Learning}
47: 235--56. \url{https://doi.org/10.1023/A:1013689704352}.

\bibitem[\citeproctext]{ref-AuerEtAl2002Nonstochastic}
Auer, Peter, Nicolò Cesa-Bianchi, Yoav Freund, and Robert E. Schapire.
2002. {``The Nonstochastic Multiarmed Bandit Problem.''} \emph{SIAM
Journal on Computing} 32 (1): 48--77.
\url{https://doi.org/10.1137/S0097539701398375}.

\bibitem[\citeproctext]{ref-BaKirosHinton2016}
Ba, Jimmy Lei, Jamie Ryan Kiros, and Geoffrey E. Hinton. 2016. {``Layer
Normalization.''} \emph{arXiv Preprint}.
\url{https://arxiv.org/abs/1607.06450}.

\bibitem[\citeproctext]{ref-BabenkoLempitsky2012}
Babenko, Artem, and Victor Lempitsky. 2012. {``The Inverted
Multi-Index.''} \emph{2012 IEEE Conference on Computer Vision and
Pattern Recognition}, 3069--76.
\url{https://doi.org/10.1109/CVPR.2012.6248038}.

\bibitem[\citeproctext]{ref-BaranesOudeyer2013GoalExploration}
Baranes, Adrien, and Pierre-Yves Oudeyer. 2013. {``Active Learning of
Inverse Models with Intrinsically Motivated Goal Exploration in
Robots.''} \emph{Robotics and Autonomous Systems} 61 (1): 49--73.
\url{https://doi.org/10.1016/j.robot.2012.05.008}.

\bibitem[\citeproctext]{ref-BaydinEtAl2018}
Baydin, Atilim Gunes, Barak A. Pearlmutter, Alexey Andreyevich Radul,
and Jeffrey Mark Siskind. 2018. {``Automatic Differentiation in Machine
Learning: A Survey.''} \emph{Journal of Machine Learning Research} 18
(153): 1--43. \url{https://www.jmlr.org/papers/v18/17-468.html}.

\bibitem[\citeproctext]{ref-BengioCourvilleVincent2013}
Bengio, Yoshua, Aaron Courville, and Pascal Vincent. 2013.
{``Representation Learning: A Review and New Perspectives.''} \emph{IEEE
Transactions on Pattern Analysis and Machine Intelligence} 35 (8):
1798--828. \url{https://doi.org/10.1109/TPAMI.2013.50}.

\bibitem[\citeproctext]{ref-BengioEtAl2009Curriculum}
Bengio, Yoshua, Jérôme Louradour, Ronan Collobert, and Jason Weston.
2009. {``Curriculum Learning.''} \emph{Proceedings of the 26th Annual
International Conference on Machine Learning}, 41--48.
\url{https://doi.org/10.1145/1553374.1553380}.

\bibitem[\citeproctext]{ref-BernsteinNewhouse2024OldOptimizer}
Bernstein, Jeremy, and Laker Newhouse. 2024. \emph{Old Optimizer, New
Norm: An Anthology}. \url{https://arxiv.org/abs/2409.20325}.

\bibitem[\citeproctext]{ref-BlierOllivier2018DescriptionLength}
Blier, Léonard, and Yann Ollivier. 2018. {``The Description Length of
Deep Learning Models.''} \emph{Advances in Neural Information Processing
Systems 31}.
\url{https://proceedings.neurips.cc/paper/2018/hash/3b712de48137572f3849aabd5666a4e3-Abstract.html}.

\bibitem[\citeproctext]{ref-BorgeaudEtAl2022}
Borgeaud, Sebastian, Arthur Mensch, Jordan Hoffmann, et al. 2022.
{``Improving Language Models by Retrieving from Trillions of Tokens.''}
\emph{International Conference on Machine Learning}, Proceedings of
machine learning research, vol. 162: 2206--37.
\url{https://proceedings.mlr.press/v162/borgeaud22a.html}.

\bibitem[\citeproctext]{ref-BubeckCesaBianchi2012}
Bubeck, Sebastien, and Nicolo Cesa-Bianchi. 2012. {``Regret Analysis of
Stochastic and Nonstochastic Multi-Armed Bandit Problems.''}
\emph{Foundations and Trends in Machine Learning} 5: 1--122.
\url{https://doi.org/10.1561/2200000024}.

\bibitem[\citeproctext]{ref-ChenRubanovaBettencourtDuvenaud2018}
Chen, Ricky T. Q., Yulia Rubanova, Jesse Bettencourt, and David
Duvenaud. 2018. {``Neural Ordinary Differential Equations.''}
\emph{Advances in Neural Information Processing Systems 31}.

\bibitem[\citeproctext]{ref-ChenEtAl2023PositionInterpolation}
Chen, Shouyuan, Sherman Wong, Liangjian Chen, and Yuandong Tian. 2023.
{``Extending Context Window of Large Language Models via Positional
Interpolation.''} \emph{arXiv Preprint arXiv:2306.15595}, ahead of
print. \url{https://doi.org/10.48550/arXiv.2306.15595}.

\bibitem[\citeproctext]{ref-ChoromanskiEtAl2021Performer}
Choromanski, Krzysztof, Valerii Likhosherstov, David Dohan, et al. 2021.
{``Rethinking Attention with Performers.''} \emph{International
Conference on Learning Representations}.
\url{https://openreview.net/forum?id=Ua6zuk0WRH}.

\bibitem[\citeproctext]{ref-Codd1970}
Codd, E. F. 1970. {``A Relational Model of Data for Large Shared Data
Banks.''} \emph{Communications of the ACM} 13 (6): 377--87.
\url{https://doi.org/10.1145/362384.362685}.

\bibitem[\citeproctext]{ref-CohenWelling2016}
Cohen, Taco, and Max Welling. 2016. {``Group Equivariant Convolutional
Networks.''} \emph{Proceedings of the 33rd International Conference on
Machine Learning}, Proceedings of machine learning research, vol. 48:
2990--99. \url{https://proceedings.mlr.press/v48/cohenc16.html}.

\bibitem[\citeproctext]{ref-CoverThomas2006}
Cover, Thomas M., and Joy A. Thomas. 2006. \emph{Elements of Information
Theory}. 2nd ed. John Wiley \& Sons.
\url{https://doi.org/10.1002/047174882X}.

\bibitem[\citeproctext]{ref-Curry2013Sheaves}
Curry, Justin Michael. 2013. {``Sheaves, Cosheaves and Applications.''}
PhD thesis, University of Pennsylvania.
\url{https://arxiv.org/abs/1303.3255}.

\bibitem[\citeproctext]{ref-DeCandiaEtAl2007}
DeCandia, Giuseppe, Deniz Hastorun, Madan Jampani, et al. 2007.
{``Dynamo: Amazon's Highly Available Key-Value Store.''}
\emph{Proceedings of the 21st ACM Symposium on Operating Systems
Principles}, 205--20. \url{https://doi.org/10.1145/1294261.1294281}.

\bibitem[\citeproctext]{ref-DevlinChangLeeToutanova2019}
Devlin, Jacob, Ming-Wei Chang, Kenton Lee, and Kristina Toutanova. 2019.
{``BERT: Pre-Training of Deep Bidirectional Transformers for Language
Understanding.''} \emph{Proceedings of NAACL-HLT 2019}, 4171--86.
\url{https://doi.org/10.18653/v1/N19-1423}.

\bibitem[\citeproctext]{ref-DormandPrince1980}
Dormand, J. R., and P. J. Prince. 1980. {``A Family of Embedded
Runge--Kutta Formulae.''} \emph{Journal of Computational and Applied
Mathematics} 6 (1): 19--26.
\url{https://doi.org/10.1016/0771-050X(80)90013-3}.

\bibitem[\citeproctext]{ref-DuchiHazanSinger2011}
Duchi, John, Elad Hazan, and Yoram Singer. 2011. {``Adaptive Subgradient
Methods for Online Learning and Stochastic Optimization.''}
\emph{Journal of Machine Learning Research} 12: 2121--59.
\url{https://www.jmlr.org/papers/v12/duchi11a.html}.

\bibitem[\citeproctext]{ref-EckartYoung1936}
Eckart, Carl, and Gale Young. 1936. {``The Approximation of One Matrix
by Another of Lower Rank.''} \emph{Psychometrika} 1 (3): 211--18.
\url{https://doi.org/10.1007/BF02288367}.

\bibitem[\citeproctext]{ref-EdelmanAriasSmith1998}
Edelman, Alan, Tomás A. Arias, and Steven T. Smith. 1998. {``The
Geometry of Algorithms with Orthogonality Constraints.''} \emph{SIAM
Journal on Matrix Analysis and Applications} 20 (2): 303--53.
\url{https://doi.org/10.1137/S0895479895290954}.

\bibitem[\citeproctext]{ref-ElhageEtAl2022Superposition}
Elhage, Nelson, Tristan Hume, Catherine Olsson, et al. 2022. {``Toy
Models of Superposition.''} \emph{Transformer Circuits Thread}.
\url{https://transformer-circuits.pub/2022/toy_model/index.html}.

\bibitem[\citeproctext]{ref-FigurnovEtAl2017}
Figurnov, Michael, Maxwell D. Collins, Yukun Zhu, et al. 2017.
{``Spatially Adaptive Computation Time for Residual Networks.''}
\emph{Proceedings of the IEEE Conference on Computer Vision and Pattern
Recognition}.

\bibitem[\citeproctext]{ref-GalGhahramani2016DropoutBayes}
Gal, Yarin, and Zoubin Ghahramani. 2016. {``Dropout as a Bayesian
Approximation: Representing Model Uncertainty in Deep Learning.''}
\emph{Proceedings of the 33rd International Conference on Machine
Learning}, Proceedings of machine learning research, vol. 48: 1050--59.
\url{https://proceedings.mlr.press/v48/gal16.html}.

\bibitem[\citeproctext]{ref-GaleElsenHooker2019}
Gale, Trevor, Erich Elsen, and Sara Hooker. 2019. \emph{The State of
Sparsity in Deep Neural Networks}.
\url{https://arxiv.org/abs/1902.09574}.

\bibitem[\citeproctext]{ref-GeifmanElYaniv2017Selective}
Geifman, Yonatan, and Ran El-Yaniv. 2017. {``Selective Classification
for Deep Neural Networks.''} \emph{Advances in Neural Information
Processing Systems 30}, 4878--87.
\url{https://papers.nips.cc/paper_files/paper/2017/hash/4a8423d5e91fda00bb7e46540e2b0cf1-Abstract.html}.

\bibitem[\citeproctext]{ref-GlorotBengio2010}
Glorot, Xavier, and Yoshua Bengio. 2010. {``Understanding the Difficulty
of Training Deep Feedforward Neural Networks.''} \emph{Proceedings of
the Thirteenth International Conference on Artificial Intelligence and
Statistics}, Proceedings of machine learning research, vol. 9: 249--56.
\url{https://proceedings.mlr.press/v9/glorot10a.html}.

\bibitem[\citeproctext]{ref-GoldmanKearns1995Teaching}
Goldman, Sally A., and Michael J. Kearns. 1995. {``On the Complexity of
Teaching.''} \emph{Journal of Computer and System Sciences} 50 (1):
20--31. \url{https://doi.org/10.1006/jcss.1995.1003}.

\bibitem[\citeproctext]{ref-GolubVanLoan2013}
Golub, Gene H., and Charles F. Van Loan. 2013. \emph{Matrix
Computations}. 4th ed. Johns Hopkins University Press.

\bibitem[\citeproctext]{ref-Graves2016ACT}
Graves, Alex. 2016. \emph{Adaptive Computation Time for Recurrent Neural
Networks}. \url{https://doi.org/10.48550/arXiv.1603.08983}.

\bibitem[\citeproctext]{ref-GravesEtAl2017Curriculum}
Graves, Alex, Marc G. Bellemare, Jacob Menick, Rémi Munos, and Koray
Kavukcuoglu. 2017. {``Automated Curriculum Learning for Neural
Networks.''} \emph{Proceedings of the 34th International Conference on
Machine Learning}, Proceedings of machine learning research, vol. 70:
1311--20. \url{https://proceedings.mlr.press/v70/graves17a.html}.

\bibitem[\citeproctext]{ref-GriewankWalther2008}
Griewank, Andreas, and Andrea Walther. 2008. \emph{Evaluating
Derivatives: Principles and Techniques of Algorithmic Differentiation}.
2nd ed. Society for Industrial; Applied Mathematics.
\url{https://doi.org/10.1137/1.9780898717761}.

\bibitem[\citeproctext]{ref-GrigsbyLindseyRolnick2023}
Grigsby, J. Elisenda, Kathryn Lindsey, and David Rolnick. 2023.
{``Hidden Symmetries of ReLU Networks.''} \emph{Proceedings of the 40th
International Conference on Machine Learning}, Proceedings of machine
learning research, vol. 202: 11734--60.
\url{https://proceedings.mlr.press/v202/grigsby23a.html}.

\bibitem[\citeproctext]{ref-GuoPleissSunWeinberger2017Calibration}
Guo, Chuan, Geoff Pleiss, Yu Sun, and Kilian Q. Weinberger. 2017. {``On
Calibration of Modern Neural Networks.''} \emph{Proceedings of the 34th
International Conference on Machine Learning}, Proceedings of machine
learning research, vol. 70: 1321--30.
\url{https://proceedings.mlr.press/v70/guo17a.html}.

\bibitem[\citeproctext]{ref-GuptaKorenSinger2018Shampoo}
Gupta, Vineet, Tomer Koren, and Yoram Singer. 2018. {``Shampoo:
Preconditioned Stochastic Tensor Optimization.''} \emph{Proceedings of
the 35th International Conference on Machine Learning}, Proceedings of
machine learning research, vol. 80: 1842--50.
\url{https://proceedings.mlr.press/v80/gupta18a.html}.

\bibitem[\citeproctext]{ref-HaberRuthotto2018}
Haber, Eldad, and Lars Ruthotto. 2018. {``Stable Architectures for Deep
Neural Networks.''} \emph{Inverse Problems} 34 (1): 014004.
\url{https://doi.org/10.1088/1361-6420/aa9a90}.

\bibitem[\citeproctext]{ref-HadfieldMenellEtAl2017OffSwitch}
Hadfield-Menell, Dylan, Anca Dragan, Pieter Abbeel, and Stuart Russell.
2017. {``The Off-Switch Game.''} \emph{Proceedings of the Twenty-Sixth
International Joint Conference on Artificial Intelligence}, 220--27.
\url{https://doi.org/10.24963/ijcai.2017/32}.

\bibitem[\citeproctext]{ref-HairerLubichWanner2006}
Hairer, Ernst, Christian Lubich, and Gerhard Wanner. 2006.
\emph{Geometric Numerical Integration: Structure-Preserving Algorithms
for Ordinary Differential Equations}. 2nd ed. Vol. 31. Springer Series
in Computational Mathematics. Springer.
\url{https://doi.org/10.1007/3-540-30666-8}.

\bibitem[\citeproctext]{ref-HairerWanner1996}
Hairer, Ernst, and Gerhard Wanner. 1996. \emph{Solving Ordinary
Differential Equations II: Stiff and Differential-Algebraic Problems}.
2nd ed. Springer Series in Computational Mathematics. Springer.
\url{https://doi.org/10.1007/978-3-642-05221-7}.

\bibitem[\citeproctext]{ref-HanMaoDally2016DeepCompression}
Han, Song, Huizi Mao, and William J. Dally. 2016a. {``Deep Compression:
Compressing Deep Neural Networks with Pruning, Trained Quantization and
Huffman Coding.''} \emph{4th International Conference on Learning
Representations (ICLR 2016)}. \url{https://arxiv.org/abs/1510.00149}.

\bibitem[\citeproctext]{ref-HanMaoDally2016}
Han, Song, Huizi Mao, and William J. Dally. 2016b. {``Deep Compression:
Compressing Deep Neural Networks with Pruning, Trained Quantization and
Huffman Coding.''} \emph{International Conference on Learning
Representations}. \url{https://arxiv.org/abs/1510.00149}.

\bibitem[\citeproctext]{ref-HeEtAl2015}
He, Kaiming, Xiangyu Zhang, Shaoqing Ren, and Jian Sun. 2015. {``Delving
Deep into Rectifiers: Surpassing Human-Level Performance on ImageNet
Classification.''} \emph{Proceedings of the IEEE International
Conference on Computer Vision}, 1026--34.
\url{https://doi.org/10.1109/ICCV.2015.123}.

\bibitem[\citeproctext]{ref-HeZhangRenSun2016}
He, Kaiming, Xiangyu Zhang, Shaoqing Ren, and Jian Sun. 2016. {``Deep
Residual Learning for Image Recognition.''} \emph{2016 IEEE Conference
on Computer Vision and Pattern Recognition}, 770--78.
\url{https://doi.org/10.1109/CVPR.2016.90}.

\bibitem[\citeproctext]{ref-HendersonEtAl2018}
Henderson, Peter, Riashat Islam, Philip Bachman, Joelle Pineau, Doina
Precup, and David Meger. 2018. {``Deep Reinforcement Learning That
Matters.''} \emph{Proceedings of the AAAI Conference on Artificial
Intelligence} 32. \url{https://doi.org/10.1609/aaai.v32i1.11694}.

\bibitem[\citeproctext]{ref-HerbsterWarmuth1998Tracking}
Herbster, Mark, and Manfred K. Warmuth. 1998. {``Tracking the Best
Expert.''} \emph{Machine Learning} 32: 151--78.
\url{https://doi.org/10.1023/A:1007424614876}.

\bibitem[\citeproctext]{ref-Higham1986Polar}
Higham, Nicholas J. 1986. {``Computing the Polar Decomposition---with
Applications.''} \emph{SIAM Journal on Scientific and Statistical
Computing} 7 (4): 1160--74. \url{https://doi.org/10.1137/0907079}.

\bibitem[\citeproctext]{ref-Higham2002}
Higham, Nicholas J. 2002. \emph{Accuracy and Stability of Numerical
Algorithms}. 2nd ed. SIAM.
\url{https://doi.org/10.1137/1.9780898718027}.

\bibitem[\citeproctext]{ref-HintonVinyalsDean2015Distillation}
Hinton, Geoffrey, Oriol Vinyals, and Jeff Dean. 2014. {``Distilling the
Knowledge in a Neural Network.''} \emph{NIPS 2014 Deep Learning
Workshop}. \url{https://arxiv.org/abs/1503.02531}.

\bibitem[\citeproctext]{ref-Hoare1969}
Hoare, C. A. R. 1969. {``An Axiomatic Basis for Computer Programming.''}
\emph{Communications of the ACM} 12 (10): 576--80.
\url{https://doi.org/10.1145/363235.363259}.

\bibitem[\citeproctext]{ref-HochreiterSchmidhuber1997}
Hochreiter, Sepp, and Jürgen Schmidhuber. 1997. {``Long Short-Term
Memory.''} \emph{Neural Computation} 9 (8): 1735--80.
\url{https://doi.org/10.1162/neco.1997.9.8.1735}.

\bibitem[\citeproctext]{ref-HornJohnson2012}
Horn, Roger A., and Charles R. Johnson. 2012. \emph{Matrix Analysis}.
2nd ed. Cambridge University Press.
\url{https://doi.org/10.1017/CBO9781139020411}.

\bibitem[\citeproctext]{ref-Howard1966Information}
Howard, Ronald A. 1966. {``Information Value Theory.''} \emph{IEEE
Transactions on Systems Science and Cybernetics} 2 (1): 22--26.
\url{https://doi.org/10.1109/TSSC.1966.300074}.

\bibitem[\citeproctext]{ref-Hutchins1995CognitionWild}
Hutchins, Edwin. 1995. \emph{Cognition in the Wild}. MIT Press.

\bibitem[\citeproctext]{ref-Iserles2008}
Iserles, Arieh. 2008. \emph{A First Course in the Numerical Analysis of
Differential Equations}. 2nd ed. Vol. 44. Cambridge Texts in Applied
Mathematics. Cambridge University Press.
\url{https://doi.org/10.1017/CBO9780511995569}.

\bibitem[\citeproctext]{ref-JordanEtAl2024Muon}
Jordan, Keller, Yuchen Jin, Vlado Boza, et al. 2024. \emph{Muon: An
Optimizer for Hidden Layers in Neural Networks}.
\url{https://kellerjordan.github.io/posts/muon/}.

\bibitem[\citeproctext]{ref-KatharopoulosEtAl2020}
Katharopoulos, Angelos, Apoorv Vyas, Nikolaos Pappas, and François
Fleuret. 2020. {``Transformers Are RNNs: Fast Autoregressive
Transformers with Linear Attention.''} \emph{Proceedings of the 37th
International Conference on Machine Learning}, Proceedings of machine
learning research, vol. 119: 5156--65.
\url{https://proceedings.mlr.press/v119/katharopoulos20a.html}.

\bibitem[\citeproctext]{ref-KendallGal2017Uncertainty}
Kendall, Alex, and Yarin Gal. 2017. {``What Uncertainties Do We Need in
Bayesian Deep Learning for Computer Vision?''} \emph{Advances in Neural
Information Processing Systems 30}.
\url{https://papers.nips.cc/paper_files/paper/2017/hash/2650d6089a6d640c5e85b2b88265dc2b-Abstract.html}.

\bibitem[\citeproctext]{ref-Khalil2002}
Khalil, Hassan K. 2002. \emph{Nonlinear Systems}. 3rd ed. Prentice Hall.

\bibitem[\citeproctext]{ref-KhandelwalEtAl2020}
Khandelwal, Urvashi, Omer Levy, Dan Jurafsky, Luke Zettlemoyer, and Mike
Lewis. 2020. {``Generalization Through Memorization: Nearest Neighbor
Language Models.''} \emph{International Conference on Learning
Representations}. \url{https://openreview.net/forum?id=HklbjCEKDB}.

\bibitem[\citeproctext]{ref-KingmaBa2015}
Kingma, Diederik P., and Jimmy Ba. 2015. {``Adam: A Method for
Stochastic Optimization.''} \emph{3rd International Conference on
Learning Representations}. \url{https://arxiv.org/abs/1412.6980}.

\bibitem[\citeproctext]{ref-KirkpatrickEtAl2017EWC}
Kirkpatrick, James, Razvan Pascanu, Neil Rabinowitz, et al. 2017.
{``Overcoming Catastrophic Forgetting in Neural Networks.''}
\emph{Proceedings of the National Academy of Sciences} 114 (13):
3521--26. \url{https://doi.org/10.1073/pnas.1611835114}.

\bibitem[\citeproctext]{ref-KleinEtAl2004TeamPlayer}
Klein, Gary, David D. Woods, Jeffrey M. Bradshaw, Robert R. Hoffman, and
Paul J. Feltovich. 2004. {``Ten Challenges for Making Automation a Team
Player in Joint Human-Agent Activity.''} \emph{IEEE Intelligent Systems}
19 (6): 91--95. \url{https://doi.org/10.1109/MIS.2004.74}.

\bibitem[\citeproctext]{ref-KlyubinPolaniNehaniv2008}
Klyubin, Alexander S., Daniel Polani, and Chrystopher L. Nehaniv. 2008.
{``Keep Your Options Open: An Information-Based Driving Principle for
Sensorimotor Systems.''} \emph{PLoS ONE} 3 (12): e4018.
\url{https://doi.org/10.1371/journal.pone.0004018}.

\bibitem[\citeproctext]{ref-KornblithShlensLe2019}
Kornblith, Simon, Jonathon Shlens, and Quoc V. Le. 2019. {``Do Better
ImageNet Models Transfer Better?''} \emph{Proceedings of the IEEE/CVF
Conference on Computer Vision and Pattern Recognition}, 2661--71.
\url{https://openaccess.thecvf.com/content_CVPR_2019/html/Kornblith_Do_Better_ImageNet_Models_Transfer_Better_CVPR_2019_paper.html}.

\bibitem[\citeproctext]{ref-Kudo2018SubwordRegularization}
Kudo, Taku. 2018. {``Subword Regularization: Improving Neural Network
Translation Models with Multiple Subword Candidates.''}
\emph{Proceedings of the 56th Annual Meeting of the Association for
Computational Linguistics (Volume 1: Long Papers)}, 66--75.
\url{https://doi.org/10.18653/v1/P18-1007}.

\bibitem[\citeproctext]{ref-KudoRichardson2018SentencePiece}
Kudo, Taku, and John Richardson. 2018. {``SentencePiece: A Simple and
Language Independent Subword Tokenizer and Detokenizer for Neural Text
Processing.''} \emph{Proceedings of the 2018 Conference on Empirical
Methods in Natural Language Processing: System Demonstrations}, 66--71.
\url{https://doi.org/10.18653/v1/D18-2012}.

\bibitem[\citeproctext]{ref-Kuznetsov2004}
Kuznetsov, Yuri A. 2004. \emph{Elements of Applied Bifurcation Theory}.
3rd ed. Vol. 112. Applied Mathematical Sciences. Springer.
\url{https://doi.org/10.1007/978-1-4757-3978-7}.

\bibitem[\citeproctext]{ref-LaiRobbins1985}
Lai, T. L., and Herbert Robbins. 1985. {``Asymptotically Efficient
Adaptive Allocation Rules.''} \emph{Advances in Applied Mathematics} 6
(1): 4--22. \url{https://doi.org/10.1016/0196-8858(85)90002-8}.

\bibitem[\citeproctext]{ref-LakshminarayananPritzelBlundell2017DeepEnsembles}
Lakshminarayanan, Balaji, Alexander Pritzel, and Charles Blundell. 2017.
{``Simple and Scalable Predictive Uncertainty Estimation Using Deep
Ensembles.''} \emph{Advances in Neural Information Processing Systems
30}.
\url{https://papers.nips.cc/paper/2017/hash/9ef2ed4b7fd2c810847ffa5fa85bce38-Abstract.html}.

\bibitem[\citeproctext]{ref-Lamport2002SpecifyingSystems}
Lamport, Leslie. 2002. \emph{Specifying Systems: The TLA+ Language and
Tools for Hardware and Software Engineers}. Addison-Wesley.

\bibitem[\citeproctext]{ref-Lanczos1950}
Lanczos, Cornelius. 1950. {``An Iteration Method for the Solution of the
Eigenvalue Problem of Linear Differential and Integral Operators.''}
\emph{Journal of Research of the National Bureau of Standards} 45:
255--82. \url{https://doi.org/10.6028/jres.045.026}.

\bibitem[\citeproctext]{ref-LattimoreSzepesvari2020}
Lattimore, Tor, and Csaba Szepesvari. 2020. \emph{Bandit Algorithms}.
Cambridge University Press. \url{https://doi.org/10.1017/9781108571401}.

\bibitem[\citeproctext]{ref-LeCunBottouBengioHaffner1998}
LeCun, Yann, Léon Bottou, Yoshua Bengio, and Patrick Haffner. 1998.
{``Gradient-Based Learning Applied to Document Recognition.''}
\emph{Proceedings of the IEEE} 86 (11): 2278--324.
\url{https://doi.org/10.1109/5.726791}.

\bibitem[\citeproctext]{ref-Lee2018Riemannian}
Lee, John M. 2018. \emph{Introduction to Riemannian Manifolds}. 2nd ed.
Vol. 176. Graduate Texts in Mathematics. Springer.
\url{https://doi.org/10.1007/978-3-319-91755-9}.

\bibitem[\citeproctext]{ref-LehmannCasella1998}
Lehmann, E. L., and George Casella. 1998. \emph{Theory of Point
Estimation}. 2nd ed. Springer. \url{https://doi.org/10.1007/b98854}.

\bibitem[\citeproctext]{ref-LeiEtAl2018ConformalRegression}
Lei, Jing, Max G'Sell, Alessandro Rinaldo, Ryan J. Tibshirani, and Larry
Wasserman. 2018. {``Distribution-Free Predictive Inference for
Regression.''} \emph{Journal of the American Statistical Association}
113 (523): 1094--111.
\url{https://doi.org/10.1080/01621459.2017.1307116}.

\bibitem[\citeproctext]{ref-LessardRechtPackard2016}
Lessard, Laurent, Benjamin Recht, and Andrew Packard. 2016. {``Analysis
and Design of Optimization Algorithms via Integral Quadratic
Constraints.''} \emph{SIAM Journal on Optimization} 26 (1): 57--95.
\url{https://doi.org/10.1137/15M1009597}.

\bibitem[\citeproctext]{ref-LewisEtAl2020RAG}
Lewis, Patrick, Ethan Perez, Aleksandra Piktus, et al. 2020.
{``Retrieval-Augmented Generation for Knowledge-Intensive NLP Tasks.''}
\emph{Advances in Neural Information Processing Systems 33}.

\bibitem[\citeproctext]{ref-LocatelloEtAl2019}
Locatello, Francesco, Stefan Bauer, Mario Lucic, et al. 2019.
{``Challenging Common Assumptions in the Unsupervised Learning of
Disentangled Representations.''} \emph{Proceedings of the 36th
International Conference on Machine Learning}, Proceedings of machine
learning research, vol. 97: 4114--24.
\url{https://proceedings.mlr.press/v97/locatello19a.html}.

\bibitem[\citeproctext]{ref-LopezPazRanzato2017GEM}
Lopez-Paz, David, and Marc'Aurelio Ranzato. 2017. {``Gradient Episodic
Memory for Continual Learning.''} \emph{Advances in Neural Information
Processing Systems 30}. \url{https://arxiv.org/abs/1706.08840}.

\bibitem[\citeproctext]{ref-LoshchilovGinsburg2026TrainingNGPT}
Loshchilov, Ilya, and Boris Ginsburg. 2026. {``Training nGPT.''}
\emph{arXiv Preprint}. \url{https://arxiv.org/abs/2608.01284}.

\bibitem[\citeproctext]{ref-LoshchilovEtAl2024nGPT}
Loshchilov, Ilya, Cheng-Ping Hsieh, Simeng Sun, and Boris Ginsburg.
2024. {``nGPT: Normalized Transformer with Representation Learning on
the Hypersphere.''} \emph{arXiv Preprint}.
\url{https://arxiv.org/abs/2410.01131}.

\bibitem[\citeproctext]{ref-LoshchilovHutter2019AdamW}
Loshchilov, Ilya, and Frank Hutter. 2019. {``Decoupled Weight Decay
Regularization.''} \emph{7th International Conference on Learning
Representations}. \url{https://arxiv.org/abs/1711.05101}.

\bibitem[\citeproctext]{ref-LuZhongLiDong2018}
Lu, Yiping, Aoxiao Zhong, Quanzheng Li, and Bin Dong. 2018. {``Beyond
Finite Layer Neural Networks: Bridging Deep Architectures and Numerical
Differential Equations.''} \emph{Proceedings of the 35th International
Conference on Machine Learning}, Proceedings of machine learning
research, vol. 80: 3276--85.
\url{https://proceedings.mlr.press/v80/lu18d.html}.

\bibitem[\citeproctext]{ref-MadryEtAl2018Adversarial}
Madry, Aleksander, Aleksandar Makelov, Ludwig Schmidt, Dimitris Tsipras,
and Adrian Vladu. 2018. {``Towards Deep Learning Models Resistant to
Adversarial Attacks.''} \emph{International Conference on Learning
Representations}. \url{https://openreview.net/forum?id=rJzIBfZAb}.

\bibitem[\citeproctext]{ref-MalkovYashunin2020}
Malkov, Yu A., and D. A. Yashunin. 2020. {``Efficient and Robust
Approximate Nearest Neighbor Search Using Hierarchical Navigable Small
World Graphs.''} \emph{IEEE Transactions on Pattern Analysis and Machine
Intelligence} 42 (4): 824--36.
\url{https://doi.org/10.1109/TPAMI.2018.2889473}.

\bibitem[\citeproctext]{ref-MallyaLazebnik2018PackNet}
Mallya, Arun, and Svetlana Lazebnik. 2018. {``PackNet: Adding Multiple
Tasks to a Single Network by Iterative Pruning.''} \emph{2018 IEEE/CVF
Conference on Computer Vision and Pattern Recognition}, 7765--73.
\url{https://doi.org/10.1109/CVPR.2018.00810}.

\bibitem[\citeproctext]{ref-McCloskeyCohen1989}
McCloskey, Michael, and Neal J. Cohen. 1989. {``Catastrophic
Interference in Connectionist Networks: The Sequential Learning
Problem.''} In \emph{Psychology of Learning and Motivation}, vol. 24.
Academic Press. \url{https://doi.org/10.1016/S0079-7421(08)60536-8}.

\bibitem[\citeproctext]{ref-McLachlanQuispel2002}
McLachlan, Robert I., and G. Reinout W. Quispel. 2002. {``Splitting
Methods.''} \emph{Acta Numerica} 11: 341--434.
\url{https://doi.org/10.1017/S0962492902000053}.

\bibitem[\citeproctext]{ref-MengEtAl2022}
Meng, Kevin, David Bau, Alex Andonian, and Yonatan Belinkov. 2022.
{``Locating and Editing Factual Associations in {GPT}.''} \emph{Advances
in Neural Information Processing Systems} 35: 17359--72.
\url{https://proceedings.neurips.cc/paper_files/paper/2022/hash/6f1d43d5a82a37e897f454479b637600-Abstract.html}.

\bibitem[\citeproctext]{ref-Mezic2005Spectral}
Mezić, Igor. 2005. {``Spectral Properties of Dynamical Systems, Model
Reduction and Decompositions.''} \emph{Nonlinear Dynamics} 41: 309--25.
\url{https://doi.org/10.1007/s11071-005-2824-x}.

\bibitem[\citeproctext]{ref-Milner1979LCF}
Milner, Robin. 1979. {``LCF: A Way of Doing Proofs with a Machine.''}
\emph{Mathematical Foundations of Computer Science 1979}, Lecture notes
in computer science, vol. 74: 146--59.
\url{https://doi.org/10.1007/3-540-09526-8_11}.

\bibitem[\citeproctext]{ref-MoldovanAbbeel2012}
Moldovan, Teodor Mihai, and Pieter Abbeel. 2012. {``Safe Exploration in
Markov Decision Processes.''} \emph{Proceedings of the 29th
International Conference on Machine Learning}.
\url{https://arxiv.org/abs/1205.4810}.

\bibitem[\citeproctext]{ref-W3CPROVDM2013}
{Moreau, Luc, Paolo Missier, et al.} 2013. \emph{PROV-DM: The PROV Data
Model}. W3C Recommendation. World Wide Web Consortium.
\url{https://www.w3.org/TR/2013/REC-prov-dm-20130430/}.

\bibitem[\citeproctext]{ref-DeMouraEtAl2015Lean}
Moura, Leonardo de, Soonho Kong, Jeremy Avigad, Floris van Doorn, and
Jakob von Raumer. 2015. {``The Lean Theorem Prover (System
Description).''} \emph{Automated Deduction - CADE-25}, Lecture notes in
computer science, vol. 9195: 378--88.
\url{https://doi.org/10.1007/978-3-319-21401-6_26}.

\bibitem[\citeproctext]{ref-NASEM2019Reproducibility}
National Academies of Sciences, Engineering, and Medicine. 2019.
\emph{Reproducibility and Replicability in Science}. The National
Academies Press. \url{https://doi.org/10.17226/25303}.

\bibitem[\citeproctext]{ref-NISTDOEHandbook}
National Institute of Standards and Technology. n.d. \emph{NIST/SEMATECH
e-Handbook of Statistical Methods: Process Improvement and Design of
Experiments}. NIST Information Technology Laboratory.
\url{https://www.itl.nist.gov/div898/handbook/pri/pri.htm}.

\bibitem[\citeproctext]{ref-Nesterov2004}
Nesterov, Yurii. 2004. \emph{Introductory Lectures on Convex
Optimization: A Basic Course}. Kluwer Academic Publishers.

\bibitem[\citeproctext]{ref-NocedalWright2006}
Nocedal, Jorge, and Stephen J. Wright. 2006. \emph{Numerical
Optimization}. 2nd ed. Springer.
\url{https://doi.org/10.1007/978-0-387-40065-5}.

\bibitem[\citeproctext]{ref-NVIDIACUDABestPractices2026}
NVIDIA Corporation. 2026a. \emph{CUDA c++ Best Practices Guide}.
\url{https://docs.nvidia.com/cuda/cuda-c-best-practices-guide/}.

\bibitem[\citeproctext]{ref-NVIDIACUDAProgrammingGuide2026}
NVIDIA Corporation. 2026b. \emph{CUDA Programming Guide}.
\url{https://docs.nvidia.com/cuda/cuda-programming-guide/}.

\bibitem[\citeproctext]{ref-OlshausenField1996}
Olshausen, Bruno A., and David J. Field. 1996. {``Emergence of
Simple-Cell Receptive Field Properties by Learning a Sparse Code for
Natural Images.''} \emph{Nature} 381 (6583): 607--9.
\url{https://doi.org/10.1038/381607a0}.

\bibitem[\citeproctext]{ref-OrseauArmstrong2016}
Orseau, Laurent, and Stuart Armstrong. 2016. {``Safely Interruptible
Agents.''} \emph{Proceedings of the Thirty-Second Conference on
Uncertainty in Artificial Intelligence}, 557--66.
\url{https://www.auai.org/uai2016/proceedings/papers/68.pdf}.

\bibitem[\citeproctext]{ref-OudeyerKaplanHafner2007IntrinsicMotivation}
Oudeyer, Pierre-Yves, Frederic Kaplan, and Verena V. Hafner. 2007.
{``Intrinsic Motivation Systems for Autonomous Mental Development.''}
\emph{IEEE Transactions on Evolutionary Computation} 11 (2): 265--86.
\url{https://doi.org/10.1109/TEVC.2006.890271}.

\bibitem[\citeproctext]{ref-PanYang2010}
Pan, Sinno Jialin, and Qiang Yang. 2010. {``A Survey on Transfer
Learning.''} \emph{IEEE Transactions on Knowledge and Data Engineering}
22 (10): 1345--59. \url{https://doi.org/10.1109/TKDE.2009.191}.

\bibitem[\citeproctext]{ref-ParikhBoyd2014Proximal}
Parikh, Neal, and Stephen Boyd. 2014. {``Proximal Algorithms.''}
\emph{Foundations and Trends in Optimization} 1 (3): 127--239.
\url{https://doi.org/10.1561/2400000003}.

\bibitem[\citeproctext]{ref-PascanuMikolovBengio2013}
Pascanu, Razvan, Tomas Mikolov, and Yoshua Bengio. 2013. {``On the
Difficulty of Training Recurrent Neural Networks.''} \emph{Proceedings
of the 30th International Conference on Machine Learning}, Proceedings
of machine learning research, vol. 28: 1310--18.
\url{https://proceedings.mlr.press/v28/pascanu13.html}.

\bibitem[\citeproctext]{ref-Pearlmutter1994Hessian}
Pearlmutter, Barak A. 1994. {``Fast Exact Multiplication by the
Hessian.''} \emph{Neural Computation} 6 (1): 147--60.
\url{https://doi.org/10.1162/neco.1994.6.1.147}.

\bibitem[\citeproctext]{ref-PengEtAl2023YaRN}
Peng, Bowen, Jeffrey Quesnelle, Honglu Fan, and Enrico Shippole. 2023.
{``YaRN: Efficient Context Window Extension of Large Language Models.''}
\emph{arXiv Preprint arXiv:2309.00071}, ahead of print.
\url{https://doi.org/10.48550/arXiv.2309.00071}.

\bibitem[\citeproctext]{ref-PetersNiculaeMartins2019Entmax}
Peters, Ben, Vlad Niculae, and André F. T. Martins. 2019. {``Sparse
Sequence-to-Sequence Models.''} \emph{Proceedings of the 57th Annual
Meeting of the Association for Computational Linguistics}, 1504--19.
\url{https://doi.org/10.18653/v1/P19-1146}.

\bibitem[\citeproctext]{ref-PlataniosEtAl2019Curriculum}
Platanios, Emmanouil Antonios, Otilia Stretcu, Graham Neubig, Barnabas
Poczos, and Tom M. Mitchell. 2019. {``Competence-Based Curriculum
Learning for Neural Machine Translation.''} \emph{Proceedings of the
2019 Conference of the North American Chapter of the Association for
Computational Linguistics: Human Language Technologies, Volume 1},
1162--72. \url{https://doi.org/10.18653/v1/N19-1119}.

\bibitem[\citeproctext]{ref-Polyak1964}
Polyak, B. T. 1964. {``Some Methods of Speeding up the Convergence of
Iteration Methods.''} \emph{USSR Computational Mathematics and
Mathematical Physics} 4 (5): 1--17.
\url{https://doi.org/10.1016/0041-5553(64)90137-5}.

\bibitem[\citeproctext]{ref-PortelasEtAl2020TeacherAlgorithms}
Portelas, Remy, Cedric Colas, Katja Hofmann, and Pierre-Yves Oudeyer.
2020. {``Teacher Algorithms for Curriculum Learning of Deep RL in
Continuously Parameterized Environments.''} \emph{Proceedings of the
Conference on Robot Learning}, Proceedings of machine learning research,
vol. 100: 835--53.
\url{https://proceedings.mlr.press/v100/portelas20a.html}.

\bibitem[\citeproctext]{ref-RadfordEtAl2019}
Radford, Alec, Jeffrey Wu, Rewon Child, David Luan, Dario Amodei, and
Ilya Sutskever. 2019. \emph{Language Models Are Unsupervised Multitask
Learners}. OpenAI.

\bibitem[\citeproctext]{ref-RaffelEtAl2020T5}
Raffel, Colin, Noam Shazeer, Adam Roberts, et al. 2020. {``Exploring the
Limits of Transfer Learning with a Unified Text-to-Text Transformer.''}
\emph{Journal of Machine Learning Research} 21 (140): 1--67.
\url{https://www.jmlr.org/papers/v21/20-074.html}.

\bibitem[\citeproctext]{ref-RaoEtAl2021DynamicViT}
Rao, Yongming, Wenliang Zhao, Benlin Liu, Jiwen Lu, Jie Zhou, and
Cho-Jui Hsieh. 2021. {``DynamicViT: Efficient Vision Transformers with
Dynamic Token Sparsification.''} \emph{Advances in Neural Information
Processing Systems} 34.

\bibitem[\citeproctext]{ref-RebuffiEtAl2017iCaRL}
Rebuffi, Sylvestre-Alvise, Alexander Kolesnikov, Georg Sperl, and
Christoph H. Lampert. 2017. {``iCaRL: Incremental Classifier and
Representation Learning.''} \emph{2017 IEEE Conference on Computer
Vision and Pattern Recognition}, 2001--10.
\url{https://doi.org/10.1109/CVPR.2017.587}.

\bibitem[\citeproctext]{ref-ReddySchmidHenningson1993}
Reddy, Satish C., Peter J. Schmid, and Dan S. Henningson. 1993.
{``Pseudospectra of the Orr--Sommerfeld Operator.''} \emph{SIAM Journal
on Applied Mathematics} 53 (1): 15--47.
\url{https://doi.org/10.1137/0153002}.

\bibitem[\citeproctext]{ref-Rissanen1978MDL}
Rissanen, Jorma. 1978. {``Modeling by Shortest Data Description.''}
\emph{Automatica} 14 (5): 465--71.
\url{https://doi.org/10.1016/0005-1098(78)90005-5}.

\bibitem[\citeproctext]{ref-RobbinsMonro1951}
Robbins, Herbert, and Sutton Monro. 1951. {``A Stochastic Approximation
Method.''} \emph{The Annals of Mathematical Statistics} 22 (3):
400--407. \url{https://doi.org/10.1214/aoms/1177729586}.

\bibitem[\citeproctext]{ref-RobertsonZaragoza2009}
Robertson, Stephen, and Hugo Zaragoza. 2009. {``The Probabilistic
Relevance Framework: BM25 and Beyond.''} \emph{Foundations and Trends in
Information Retrieval} 4 (1--2): 1--174.
\url{https://doi.org/10.1561/1500000019}.

\bibitem[\citeproctext]{ref-Rubin1974Causal}
Rubin, Donald B. 1974. {``Estimating Causal Effects of Treatments in
Randomized and Nonrandomized Studies.''} \emph{Journal of Educational
Psychology} 66 (5): 688--701. \url{https://doi.org/10.1037/h0037350}.

\bibitem[\citeproctext]{ref-RumelhartHintonWilliams1986}
Rumelhart, David E., Geoffrey E. Hinton, and Ronald J. Williams. 1986.
{``Learning Representations by Back-Propagating Errors.''} \emph{Nature}
323: 533--36. \url{https://doi.org/10.1038/323533a0}.

\bibitem[\citeproctext]{ref-RussoVanRoy2014}
Russo, Daniel, and Benjamin Van Roy. 2014. {``Learning to Optimize via
Information-Directed Sampling.''} \emph{Advances in Neural Information
Processing Systems 27}. \url{https://arxiv.org/abs/1403.5556}.

\bibitem[\citeproctext]{ref-RustEtAl2021Tokenizer}
Rust, Phillip, Jonas Pfeiffer, Ivan Vulić, Sebastian Ruder, and Iryna
Gurevych. 2021. {``How Good Is Your Tokenizer? On the Monolingual
Performance of Multilingual Language Models.''} \emph{Proceedings of the
59th Annual Meeting of the Association for Computational Linguistics and
the 11th International Joint Conference on Natural Language Processing
(Volume 1: Long Papers)}, 3118--35.
\url{https://doi.org/10.18653/v1/2021.acl-long.243}.

\bibitem[\citeproctext]{ref-Saad2003}
Saad, Yousef. 2003. \emph{Iterative Methods for Sparse Linear Systems}.
2nd ed. Society for Industrial; Applied Mathematics.
\url{https://doi.org/10.1137/1.9780898718003}.

\bibitem[\citeproctext]{ref-SakoeChiba1978DTW}
Sakoe, Hiroaki, and Seibi Chiba. 1978. {``Dynamic Programming Algorithm
Optimization for Spoken Word Recognition.''} \emph{IEEE Transactions on
Acoustics, Speech, and Signal Processing} 26 (1): 43--49.
\url{https://doi.org/10.1109/TASSP.1978.1163055}.

\bibitem[\citeproctext]{ref-SaltonWongYang1975}
Salton, Gerard, Anita Wong, and C. S. Yang. 1975. {``A Vector Space
Model for Automatic Indexing.''} \emph{Communications of the ACM} 18
(11): 613--20. \url{https://doi.org/10.1145/361219.361220}.

\bibitem[\citeproctext]{ref-SandveEtAl2013}
Sandve, Geir Kjetil, Anton Nekrutenko, James Taylor, and Eivind Hovig.
2013. {``Ten Simple Rules for Reproducible Computational Research.''}
\emph{PLOS Computational Biology} 9 (10): e1003285.
\url{https://doi.org/10.1371/journal.pcbi.1003285}.

\bibitem[\citeproctext]{ref-Schmidhuber2009CompressionProgress}
Schmidhuber, Jürgen. 2009. {``Driven by Compression Progress: A Simple
Principle Explains Essential Aspects of Subjective Beauty, Novelty,
Surprise, Interestingness, Attention, Curiosity, Creativity, Art,
Science, Music, Jokes.''} \emph{Computational Creativity: An
Interdisciplinary Approach}, Dagstuhl seminar proceedings, vol. 09291:
1--35. \url{https://doi.org/10.4230/DagSemProc.09291.14}.

\bibitem[\citeproctext]{ref-SennrichHaddowBirch2016BPE}
Sennrich, Rico, Barry Haddow, and Alexandra Birch. 2016. {``Neural
Machine Translation of Rare Words with Subword Units.''}
\emph{Proceedings of the 54th Annual Meeting of the Association for
Computational Linguistics (Volume 1: Long Papers)}, 1715--25.
\url{https://doi.org/10.18653/v1/P16-1162}.

\bibitem[\citeproctext]{ref-ShawUszkoreitVaswani2018Relative}
Shaw, Peter, Jakob Uszkoreit, and Ashish Vaswani. 2018.
{``Self-Attention with Relative Position Representations.''}
\emph{Proceedings of the 2018 Conference of the North American Chapter
of the Association for Computational Linguistics: Human Language
Technologies, Volume 2 (Short Papers)}, 464--68.
\url{https://doi.org/10.18653/v1/N18-2074}.

\bibitem[\citeproctext]{ref-Shoemake1985}
Shoemake, Ken. 1985. {``Animating Rotation with Quaternion Curves.''}
\emph{ACM SIGGRAPH Computer Graphics} 19 (3): 245--54.
\url{https://doi.org/10.1145/325165.325242}.

\bibitem[\citeproctext]{ref-SoaresEtAl2015Corrigibility}
Soares, Nate, Benja Fallenstein, Eliezer Yudkowsky, and Stuart
Armstrong. 2015. {``Corrigibility.''} \emph{AAAI Workshop on Artificial
Intelligence and Ethics}.
\url{https://intelligence.org/files/Corrigibility.pdf}.

\bibitem[\citeproctext]{ref-Solomonoff1964InductiveInferenceI}
Solomonoff, Ray J. 1964. {``A Formal Theory of Inductive Inference. Part
i.''} \emph{Information and Control} 7 (1): 1--22.
\url{https://doi.org/10.1016/S0019-9958(64)90223-2}.

\bibitem[\citeproctext]{ref-SrivastavaGreffSchmidhuber2015}
Srivastava, Rupesh Kumar, Klaus Greff, and Jürgen Schmidhuber. 2015.
{``Highway Networks.''} \emph{arXiv Preprint}.
\url{https://arxiv.org/abs/1505.00387}.

\bibitem[\citeproctext]{ref-SuEtAl2021RoFormer}
Su, Jianlin, Yu Lu, Shengfeng Pan, Ahmed Murtadha, Bo Wen, and Yunfeng
Liu. 2021. {``RoFormer: Enhanced Transformer with Rotary Position
Embedding.''} \emph{arXiv Preprint arXiv:2104.09864}, ahead of print.
\url{https://doi.org/10.48550/arXiv.2104.09864}.

\bibitem[\citeproctext]{ref-SugiyamaKrauledatMuller2007Covariate}
Sugiyama, Masashi, Matthias Krauledat, and Klaus-Robert Müller. 2007.
{``Covariate Shift Adaptation by Importance Weighted Cross
Validation.''} \emph{Journal of Machine Learning Research} 8: 985--1005.
\url{https://www.jmlr.org/papers/v8/sugiyama07a.html}.

\bibitem[\citeproctext]{ref-SutskeverEtAl2013}
Sutskever, Ilya, James Martens, George Dahl, and Geoffrey Hinton. 2013.
{``On the Importance of Initialization and Momentum in Deep Learning.''}
\emph{Proceedings of the 30th International Conference on Machine
Learning}, Proceedings of machine learning research, vol. 28: 1139--47.
\url{https://proceedings.mlr.press/v28/sutskever13.html}.

\bibitem[\citeproctext]{ref-SutskeverVinyalsLe2014}
Sutskever, Ilya, Oriol Vinyals, and Quoc V. Le. 2014. {``Sequence to
Sequence Learning with Neural Networks.''} \emph{Advances in Neural
Information Processing Systems 27}, 3104--12.

\bibitem[\citeproctext]{ref-Thompson1933}
Thompson, William R. 1933. {``On the Likelihood That One Unknown
Probability Exceeds Another in View of the Evidence of Two Samples.''}
\emph{Biometrika} 25 (3--4): 285--94.
\url{https://doi.org/10.1093/biomet/25.3-4.285}.

\bibitem[\citeproctext]{ref-TielemanHinton2012}
Tieleman, Tijmen, and Geoffrey Hinton. 2012. \emph{Lecture
6.5---RMSProp: Divide the Gradient by a Running Average of Its Recent
Magnitude}. Coursera: Neural Networks for Machine Learning.

\bibitem[\citeproctext]{ref-TishbyPereiraBialek2000}
Tishby, Naftali, Fernando C. Pereira, and William Bialek. 2000. {``The
Information Bottleneck Method.''} \emph{arXiv Preprint}.
\url{https://arxiv.org/abs/physics/0004057}.

\bibitem[\citeproctext]{ref-TrefethenBau1997}
Trefethen, Lloyd N., and III David Bau. 1997. \emph{Numerical Linear
Algebra}. Society for Industrial; Applied Mathematics.
\url{https://doi.org/10.1137/1.9780898719574}.

\bibitem[\citeproctext]{ref-TrefethenEmbree2005}
Trefethen, Lloyd N., and Mark Embree. 2005. \emph{Spectra and
Pseudospectra: The Behavior of Nonnormal Matrices and Operators}.
Princeton University Press. \url{https://doi.org/10.2307/j.ctvzxx9kj}.

\bibitem[\citeproctext]{ref-TrefethenEtAl1993}
Trefethen, Lloyd N., Anne E. Trefethen, Satish C. Reddy, and Tobin A.
Driscoll. 1993. {``Hydrodynamic Stability Without Eigenvalues.''}
\emph{Science} 261 (5121): 578--84.
\url{https://doi.org/10.1126/science.261.5121.578}.

\bibitem[\citeproctext]{ref-TsaiEtAl2019}
Tsai, Yao-Hung Hubert, Shaojie Bai, Makoto Yamada, Louis-Philippe
Morency, and Ruslan Salakhutdinov. 2019. {``Transformer Dissection: An
Unified Understanding for Transformer's Attention via the Lens of
Kernel.''} \emph{Proceedings of EMNLP-IJCNLP}, 4344--53.
\url{https://doi.org/10.18653/v1/D19-1443}.

\bibitem[\citeproctext]{ref-VaswaniEtAl2017}
Vaswani, Ashish, Noam Shazeer, Niki Parmar, et al. 2017. {``Attention Is
All You Need.''} \emph{Advances in Neural Information Processing Systems
30}. \url{https://papers.nips.cc/paper/7181-attention-is-all-you-need}.

\bibitem[\citeproctext]{ref-VanDeVenTuytelaarsTolias2022}
Ven, Gido M. van de, Tinne Tuytelaars, and Andreas S. Tolias. 2022.
{``Three Types of Incremental Learning.''} \emph{Nature Machine
Intelligence} 4: 1185--97.
\url{https://doi.org/10.1038/s42256-022-00568-3}.

\bibitem[\citeproctext]{ref-Vershynin2018}
Vershynin, Roman. 2018. \emph{High-Dimensional Probability: An
Introduction with Applications in Data Science}. Vol. 47. Cambridge
Series in Statistical and Probabilistic Mathematics. Cambridge
University Press. \url{https://doi.org/10.1017/9781108231596}.

\bibitem[\citeproctext]{ref-VoitaTitov2020MDLProbing}
Voita, Elena, and Ivan Titov. 2020. {``Information-Theoretic Probing
with Minimum Description Length.''} \emph{Proceedings of the 2020
Conference on Empirical Methods in Natural Language Processing},
183--96. \url{https://doi.org/10.18653/v1/2020.emnlp-main.14}.

\bibitem[\citeproctext]{ref-VovkGammermanShafer2005Conformal}
Vovk, Vladimir, Alexander Gammerman, and Glenn Shafer. 2005.
\emph{Algorithmic Learning in a Random World}. Springer.
\url{https://doi.org/10.1007/b106715}.

\bibitem[\citeproctext]{ref-WainwrightJordan2008}
Wainwright, Martin J., and Michael I. Jordan. 2008. {``Graphical Models,
Exponential Families, and Variational Inference.''} \emph{Foundations
and Trends in Machine Learning} 1 (1--2): 1--305.
\url{https://doi.org/10.1561/2200000001}.

\bibitem[\citeproctext]{ref-WangEtAl2020Linformer}
Wang, Sinong, Belinda Z. Li, Madian Khabsa, Han Fang, and Hao Ma. 2020.
\emph{Linformer: Self-Attention with Linear Complexity}.
\url{https://arxiv.org/abs/2006.04768}.

\bibitem[\citeproctext]{ref-WangIsola2020Hypersphere}
Wang, Tongzhou, and Phillip Isola. 2020. {``Understanding Contrastive
Representation Learning Through Alignment and Uniformity on the
Hypersphere.''} \emph{Proceedings of the 37th International Conference
on Machine Learning}, Proceedings of machine learning research, vol.
119: 9929--39. \url{https://proceedings.mlr.press/v119/wang20k.html}.

\bibitem[\citeproctext]{ref-WangEtAl2018SkipNet}
Wang, Xin, Fisher Yu, Zi-Yi Dou, Trevor Darrell, and Joseph E. Gonzalez.
2018. {``SkipNet: Learning Dynamic Routing in Convolutional Networks.''}
\emph{Proceedings of the European Conference on Computer Vision},
409--24. \url{https://arxiv.org/abs/1711.09485}.

\bibitem[\citeproctext]{ref-WassersteinLazar2016}
Wasserstein, Ronald L., and Nicole A. Lazar. 2016. {``The ASA Statement
on p-Values: Context, Process, and Purpose.''} \emph{The American
Statistician} 70 (2): 129--33.
\url{https://doi.org/10.1080/00031305.2016.1154108}.

\bibitem[\citeproctext]{ref-WilliamsWatermanPatterson2009}
Williams, Samuel, Andrew Waterman, and David Patterson. 2009.
{``Roofline: An Insightful Visual Performance Model for Multicore
Architectures.''} \emph{Communications of the ACM} 52 (4): 65--76.
\url{https://doi.org/10.1145/1498765.1498785}.

\bibitem[\citeproctext]{ref-WooldridgeJennings1995Agents}
Wooldridge, Michael, and Nicholas R. Jennings. 1995. {``Intelligent
Agents: Theory and Practice.''} \emph{The Knowledge Engineering Review}
10 (2): 115--52. \url{https://doi.org/10.1017/S0269888900008122}.

\bibitem[\citeproctext]{ref-XiongEtAl2020}
Xiong, Ruibin, Yunchang Yang, Di He, et al. 2020. {``On Layer
Normalization in the Transformer Architecture.''} \emph{Proceedings of
the 37th International Conference on Machine Learning}, Proceedings of
machine learning research, vol. 119: 10524--33.
\url{https://proceedings.mlr.press/v119/xiong20b.html}.

\bibitem[\citeproctext]{ref-XiongEtAl2021Nystromformer}
Xiong, Yunyang, Zhanpeng Zeng, Rudrasis Chakraborty, et al. 2021.
{``Nyströmformer: A Nyström-Based Algorithm for Approximating
Self-Attention.''} \emph{Proceedings of the AAAI Conference on
Artificial Intelligence} 35: 14138--48.
\url{https://doi.org/10.1609/aaai.v35i16.17664}.

\bibitem[\citeproctext]{ref-XueEtAl2022ByT5}
Xue, Linting, Aditya Barua, Noah Constant, et al. 2022. {``ByT5: Towards
a Token-Free Future with Pre-Trained Byte-to-Byte Models.''}
\emph{Transactions of the Association for Computational Linguistics} 10:
291--306. \url{https://doi.org/10.1162/tacl_a_00461}.

\bibitem[\citeproctext]{ref-YosinskiEtAl2014}
Yosinski, Jason, Jeff Clune, Yoshua Bengio, and Hod Lipson. 2014. {``How
Transferable Are Features in Deep Neural Networks?''} \emph{Advances in
Neural Information Processing Systems} 27: 3320--28.
\url{https://proceedings.neurips.cc/paper_files/paper/2014/hash/532a2f85b6977104bc93f8580abbb330-Abstract.html}.

\bibitem[\citeproctext]{ref-ZhouDeLaTorre2012GTW}
Zhou, Feng, and Fernando De la Torre. 2012. {``Generalized Time Warping
for Multi-Modal Alignment of Human Motion.''} \emph{2012 IEEE Conference
on Computer Vision and Pattern Recognition}, 1282--89.
\url{https://doi.org/10.1109/CVPR.2012.6247812}.

\bibitem[\citeproctext]{ref-ZhouEtAl2022ExpertChoice}
Zhou, Yanqi, Tao Lei, Hanxiao Liu, et al. 2022. {``Mixture-of-Experts
with Expert Choice Routing.''} \emph{Advances in Neural Information
Processing Systems} 35: 7103--14.
\url{https://doi.org/10.52202/068431-0515}.

\bibitem[\citeproctext]{ref-Zhu2015MachineTeaching}
Zhu, Xiaojin. 2015. {``Machine Teaching: An Inverse Problem to Machine
Learning and an Approach Toward Optimal Education.''} \emph{Proceedings
of the Twenty-Ninth AAAI Conference on Artificial Intelligence},
4083--87. \url{https://doi.org/10.1609/aaai.v29i1.9761}.

\bibitem[\citeproctext]{ref-Ziyin2024Symmetry}
Ziyin, Liu. 2024. {``Symmetry Induces Structure and Constraint of
Learning.''} \emph{Proceedings of the 41st International Conference on
Machine Learning}, Proceedings of machine learning research, vol. 235:
62847--66. \url{https://proceedings.mlr.press/v235/ziyin24a.html}.

\bibitem[\citeproctext]{ref-ZophEtAl2022STMoE}
Zoph, Barret, Irwan Bello, Sameer Kumar, et al. 2022. \emph{ST-MoE:
Designing Stable and Transferable Sparse Expert Models}.
\url{https://arxiv.org/abs/2202.08906}.

\end{CSLReferences}

\backmatter
\end{document}
